% !TEX root = main_integrated.tex
%
% Integrated main file for the Platonic reconstitution of
% "The Igusa Square Root". The ten-chapter pentadic spine of CLAUDE.md §II
% is composed by \input from the chapter and appendix files in chapters/
% and appendices/, each authored as a standalone load-bearing artifact
% and integrated here by deep semantic merge of the legacy main.tex
% content (now distributed across the chapter files).
%
% The legacy monolithic main.tex (18,629 lines) is preserved on disk
% pending integration verification. This integrated build produces the
% full pentadic-spine PDF directly from the chapter/appendix files.

\documentclass[11pt,oldfontcommands]{memoir}
\usepackage{raeez-math-template}
\renewcommand{\thesection}{\arabic{section}}
\renewcommand{\theHsection}{\arabic{section}}

\title{\texorpdfstring{\textit{The Igusa Square Root}\\[0.5em]\small The Borcherds determinant of \(\phi_{0,1}\), its denominator algebra, and the \(K3\times E\) realization problem}{The Igusa Square Root: The Borcherds determinant of phi\_0,1, its denominator algebra, and the K3 x E realization problem}}
\author{Raeez Lorgat}
\date{April 2026}

% ------------------------------------------------------------------
% Standard symbols (preserved verbatim from legacy main.tex preamble)
% ------------------------------------------------------------------
\newcommand{\Imag}{\mathrm{Im}}
\newcommand{\Real}{\mathrm{Re}}
\newcommand{\Proj}{\mathbb{P}}
\newcommand{\N}{\mathbb{N}}
\newcommand{\Q}{\mathbb{Q}}
\newcommand{\Z}{\mathbb{Z}} \newcommand{\C}{\mathbb{C}}
\newcommand{\R}{\mathbb{R}} \newcommand{\La}{\Lambda}
\newcommand{\HypPlan}{\Lambda^{(1,1)}}
\newcommand{\Sp}{\mathbf{Sp}}
\newcommand{\PSp}{\mathbf{PSp}}
\newcommand{\GL}{\mathbf{GL}}
\newcommand{\SL}{\mathbf{SL}}
\newcommand{\Orth}{\mathbf{O}}
\newcommand{\SO}{\mathbf{SO}}
\newcommand{\Hpl}{\mathbb{H}}
\newcommand{\IV}{\mathbf{IV}}
\newcommand{\fbasis}{(f_i)_{\{1,2,3,-2,-1\}}}
\newcommand{\Id}{\mathbf{I}}
\newcommand{\Cone}{\mathcal{C}}
\newcommand{\Poly}{\mathcal{P}}
\newcommand{\Aut}{\operatorname{Aut}}
\newcommand{\Hom}{\operatorname{Hom}}
\newcommand{\tr}{\operatorname{tr}}
\newcommand{\bkm}{\mathfrak{g}}
\newcommand{\sdim}{\operatorname{sdim}}
\newcommand{\sch}{\operatorname{sch}}
\newcommand{\sdet}{\operatorname{sdet}}
\newcommand{\rootmult}{\operatorname{mult}}
\newcommand{\smult}{\operatorname{smult}}
\newcommand{\Prim}{\operatorname{Prim}}
% Outer braces wrap the subscripted symbol so that further subscripts
% (\autcor_0, \autcor^{\mathrm{bulk}}, etc.) attach without double-subscript error.
\newcommand{\autcor}{{\mathfrak{g}_{\Delta_5}}}

% Chapter-specific symbols introduced by individual chapter agents and
% promoted to main preamble for cross-chapter use.
\providecommand{\Mp}{\mathbf{Mp}}
\providecommand{\Spin}{\mathbf{Spin}}
\providecommand{\Pf}{\operatorname{Pf}}
\providecommand{\Un}{\mathbf{U}}
\providecommand{\J}{\mathbf{J}}
\providecommand{\hn}{\operatorname{hn}}
\providecommand{\Pic}{\operatorname{Pic}}
\providecommand{\NS}{\operatorname{NS}}
\providecommand{\Ext}{\operatorname{Ext}}
\providecommand{\Tor}{\operatorname{Tor}}
\providecommand{\End}{\operatorname{End}}
\providecommand{\codim}{\operatorname{codim}}
\providecommand{\rk}{\operatorname{rk}}
\providecommand{\Spec}{\operatorname{Spec}}
\providecommand{\Span}{\operatorname{Span}}
\providecommand{\Sym}{\operatorname{Sym}}
\providecommand{\Coh}{\operatorname{Coh}}
\providecommand{\Perf}{\operatorname{Perf}}
\providecommand{\sgn}{\operatorname{sgn}}
\providecommand{\Stab}{\operatorname{Stab}}
\providecommand{\Lie}{\operatorname{Lie}}
\providecommand{\GW}{\operatorname{GW}}
\providecommand{\PT}{\operatorname{PT}}
\providecommand{\DT}{\operatorname{DT}}

% Unicode circled digits used as view-tags ① ② ③ ④ ⑤ across chapters.
\DeclareUnicodeCharacter{2460}{\textcircled{1}}
\DeclareUnicodeCharacter{2461}{\textcircled{2}}
\DeclareUnicodeCharacter{2462}{\textcircled{3}}
\DeclareUnicodeCharacter{2463}{\textcircled{4}}
\DeclareUnicodeCharacter{2464}{\textcircled{5}}
\DeclareUnicodeCharacter{2465}{\textcircled{6}}
\DeclareUnicodeCharacter{2466}{\textcircled{7}}
\DeclareUnicodeCharacter{2467}{\textcircled{8}}

% Greek letters used inline in chapter prose (template declares some;
% these extend the set for the integrated build).
\newunicodechar{α}{$\alpha$}
\newunicodechar{δ}{$\delta$}
\newunicodechar{ε}{$\varepsilon$}
\newunicodechar{ζ}{$\zeta$}
\newunicodechar{η}{$\eta$}
\newunicodechar{θ}{$\theta$}
\newunicodechar{ι}{$\iota$}
\newunicodechar{λ}{$\lambda$}
\newunicodechar{ν}{$\nu$}
\newunicodechar{ξ}{$\xi$}
\newunicodechar{π}{$\pi$}
\newunicodechar{ρ}{$\rho$}
\newunicodechar{σ}{$\sigma$}
\newunicodechar{τ}{$\tau$}
\newunicodechar{υ}{$\upsilon$}
\newunicodechar{φ}{$\varphi$}
\newunicodechar{ϕ}{$\phi$}
\newunicodechar{χ}{$\chi$}
\newunicodechar{ψ}{$\psi$}
\newunicodechar{ω}{$\omega$}
\newunicodechar{Δ}{$\Delta$}
\newunicodechar{Λ}{$\Lambda$}
\newunicodechar{Σ}{$\Sigma$}
\newunicodechar{Ψ}{$\Psi$}
\newunicodechar{Γ}{$\Gamma$}
\newunicodechar{Π}{$\Pi$}
\newunicodechar{Θ}{$\Theta$}
\newunicodechar{Ξ}{$\Xi$}
\newunicodechar{Υ}{$\Upsilon$}

% Math symbols used inline in chapter prose.
\newunicodechar{×}{$\times$}
\newunicodechar{·}{$\cdot$}
\newunicodechar{±}{$\pm$}
\newunicodechar{⊂}{$\subset$}
\newunicodechar{⊃}{$\supset$}
\newunicodechar{⊆}{$\subseteq$}
\newunicodechar{⊇}{$\supseteq$}
\newunicodechar{∉}{$\notin$}
\newunicodechar{∀}{$\forall$}
\newunicodechar{∃}{$\exists$}
\newunicodechar{∅}{$\emptyset$}
\newunicodechar{∞}{$\infty$}
\newunicodechar{→}{$\to$}
\newunicodechar{←}{$\leftarrow$}
\newunicodechar{⇒}{$\Rightarrow$}
\newunicodechar{⇐}{$\Leftarrow$}
\newunicodechar{⇔}{$\Leftrightarrow$}
\newunicodechar{↦}{$\mapsto$}
\newunicodechar{↪}{$\hookrightarrow$}
\newunicodechar{∧}{$\wedge$}
\newunicodechar{∪}{$\cup$}
\newunicodechar{∩}{$\cap$}
\newunicodechar{⨁}{$\bigoplus$}
\newunicodechar{⨂}{$\bigotimes$}
\newunicodechar{⊥}{$\perp$}
\newunicodechar{⟨}{$\langle$}
\newunicodechar{⟩}{$\rangle$}
\newunicodechar{∂}{$\partial$}
\newunicodechar{∇}{$\nabla$}
\newunicodechar{∘}{$\circ$}
\newunicodechar{∗}{$\ast$}
\newunicodechar{ℝ}{$\mathbb{R}$}
\newunicodechar{ℂ}{$\mathbb{C}$}
\newunicodechar{ℤ}{$\mathbb{Z}$}
\newunicodechar{ℕ}{$\mathbb{N}$}
\newunicodechar{ℚ}{$\mathbb{Q}$}
\newunicodechar{ℋ}{$\mathcal{H}$}
\newunicodechar{ℒ}{$\mathcal{L}$}
\newunicodechar{𝓗}{$\mathcal{H}$}
\newunicodechar{𝓟}{$\mathcal{P}$}
\newunicodechar{𝓒}{$\mathcal{C}$}
\newunicodechar{𝓕}{$\mathcal{F}$}
\newunicodechar{𝔄}{$\mathfrak{A}$}
\newunicodechar{𝔇}{$\mathfrak{D}$}
\newunicodechar{𝔤}{$\mathfrak{g}$}
\newunicodechar{𝔥}{$\mathfrak{h}$}
\newunicodechar{𝔪}{$\mathfrak{m}$}
\newunicodechar{𝔫}{$\mathfrak{n}$}
\newunicodechar{₄}{$_4$}
\newunicodechar{₅}{$_5$}
\newunicodechar{₆}{$_6$}
\newunicodechar{₇}{$_7$}
\newunicodechar{₈}{$_8$}
\newunicodechar{₉}{$_9$}
\newunicodechar{¹}{$^1$}
\newunicodechar{²}{$^2$}
\newunicodechar{³}{$^3$}
\newunicodechar{⁴}{$^4$}
\newunicodechar{⁵}{$^5$}
\newunicodechar{⁻}{$^-$}
\newunicodechar{⁺}{$^+$}

% ------------------------------------------------------------------
% Licensing-type system (Vol II §5; cited as established prior-art per
% STANDARDS §III standard-terminology rule, not coined here).
%
% Every load-bearing theorem statement carries inline tags from the
% set {α (chart/scope/log), β (comparison/functor), γ (ambient
% declaration), δ (endpoint/convergence), ε (descent/completion)}.
% The system harmonises Chs. 1, 2, 4, 5, 6, 7, 8, 9, 10 and Appendices
% A, B, D, E, F; Ch. 3 uses epistemic-status tags only and is healed
% inline with explicit references.
%
% Chain location every theorem identifies its position on the universal
% stage chain (Vol II §4):
%   P → C → S → Z → A
% (primitive → chart → shadow → centre/bulk → acting object).
%
% Three Δ_5 names tagged distinctly on first use per chapter:
%   D_X = Δ_5                     (View ①, automorphic section)
%   den(g_{Δ_5}) = 64⁻¹ Δ_5(2Z)   (View ②, BKM denominator)
%   Pf_prot(𝔇_X^DI)                (View ③, conditional protected Pfaffian)
% Conjectural fourth and fifth handles are introduced in Chs. 7–8.
% ------------------------------------------------------------------

\ifdefined\mainpreambleonly
  \expandafter\endinput
\fi

\begin{document}

\maketitle

% ------------------------------------------------------------------
% Frontispiece + Chapter 1 (AC01)
% ------------------------------------------------------------------

\chapter*{Frontispiece}
\addcontentsline{toc}{chapter}{Frontispiece}
\label{ch:frontispiece}

The Igusa cusp form \(\Delta_5\) has five readings.  The genus-two
automorphic section, the denominator of a generalized
Borcherds--Kac--Moody superalgebra on \(\La^{2,1}_{II}\), and the
singular theta lift of the weak Jacobi form \(\phi_{0,1}\) are
theorem-level readings.  The protected Pfaffian is a retained-data
reading of the chosen \(E_3\)-prefactorization datum
\[
  \mathfrak A_{K3\times E}
  \;\in\; E_3\text{-}\mathrm{HolFA}(K3\times E),
\]
and the mirror discriminant is a conjectural period-side reading on the
rank-\(3\) Mukai-invariant \(K3\times E\) component.

\paragraph{The five views.}
\begin{enumerate}
\item[\textcircled{1}]
\textsl{Automorphic.}
The Borcherds--Gritsenko--Nikulin section
\[
  \mathcal D_X \;=\; \Delta_5
  \;\in\; H^0\bigl(\Sp(2,\Z)\backslash\mathcal H_2,\;\mathcal L^{\otimes 5}
  \otimes\nu_{\Delta_5}\bigr)
\]
of the weight-\(5\) automorphic line on Siegel upper-half space, with
multiplier \(\nu_{\Delta_5}\) of \cite{MAASS:1, Borcherds95}.

\item[\textcircled{2}]
\textsl{Borcherds--Kac--Moody denominator.}
The denominator of the generalized BKM Lie superalgebra
\(\mathfrak g_{\Delta_5}\) on the rank-three Lorentzian root lattice
\(\La^{2,1}_{II}\subset\La^{2,1}=\HypPlan\oplus\langle 2\rangle\) of
signature \((2,1)\), the index-four sublattice with even \(f_{\pm 2}\)
coordinates, abstractly \(U(4)\oplus\langle 2\rangle\) (discriminant
\(-32\), the index-four scaling of the hyperbolic block).  The ambient
lattice \(\La^{2,1}\) has basis \((f_2,f_3,f_{-2})\) with
\((f_2,f_{-2})=-1\) and \((f_3,f_3)=2\); the sublattice
\(\La^{2,1}_{II}\) is generated by \(2f_2,f_3,2f_{-2}\).  The type-II
Cartan matrix on simple roots
\(\delta_1 = 2f_2 - f_3\), \(\delta_2 = 2f_{-2} - f_3\),
\(\delta_3 = f_3\) is, with \(\mathbf{J}_3\) the \(3\times 3\)
all-ones matrix,
\((\delta_i, \delta_j) = 4\Id_3 - 2\mathbf{J}_3 = \bigl(\begin{smallmatrix}
2 & -2 & -2\\ -2 & 2 & -2\\ -2 & -2 & 2 \end{smallmatrix}\bigr)\),
\[
  \operatorname{den}(\mathfrak g_{\Delta_5})
  \;=\; 64^{-1}\,\Delta_5(2Z),
\]
with signed root supermultiplicities
\(\smult\,\mathfrak g_{\Delta_5,\alpha(n,l,m)} = f(nm,l)\) on the
active support \(\Gamma_{\mathrm{act}} = \{(n,l,m)\in\Gamma_{\mathrm{eff}}
\mid f(nm,l)\ne 0\}\), with
\(\alpha(n,l,m) = 2nf_2 - lf_3 + 2mf_{-2}\)
\cite{Borcherds95, GN, GNPartI, GNII}.

\item[\textcircled{3}]
\textsl{Compact Pfaffian.}
The retained Pfaffian line of the \(K3\times E\) chiral
\(E_3\)-prefactorization input, presented at finite stage \(R\) by the
Dirac--Igusa pro-object
\[
  \mathfrak D^{DI}_X \;=\;
  \varprojlim_R\;
  \bigl(\mathcal A^{E_3}_{X,R},\,F^{\mathrm{hyb}}_{X,R},\,\Gamma_{X,R},
  \,\Pi_X,\,\Phi_R,\,o_R,\,H_R,\,P^{\Pi}_R,\,C_{X,R},\,
  \Theta_{\mathrm{Kos},R},\,\mathcal L_{\mathrm{Pf},R},\,\mathrm{pf}_{X,R},
  \,\epsilon_{o,R},\,\mathrm{Rec}_R\bigr)
\]
(the chosen \(E_3\)-prefactorization datum \(\mathcal A^{E_3}\), the hybrid factorization
carrier \(F^{\mathrm{hyb}}\) of \(\mathrm{Ran}^{\mathrm{hyb}}(E)\), the
Mukai charge lattice \(\Gamma\), the Gram projection \(\Pi_X\), the
vanishing-cycle coefficient sheaf \(\Phi\), the orientation \(o\), the Hall datum
\(H\), the Gram-descended primitive \(P^\Pi\), the source coalgebra
\(C\), the Koszul cobar \(\Theta_{\mathrm{Kos}}\), the Pfaffian line
\(\mathcal L_{\mathrm{Pf}}\), the Pfaffian section \(\mathrm{pf}\), the
orientation character \(\epsilon_o\), and the recognition map
\(\mathrm{Rec}\); Chapter~5 specifies the finite-stage data and their
residuals, with the \(E_3\), wrapped, Koszul, and recognition lanes
retained or conditional exactly as stated there).
The obstruction \((D0)\) is reduced in
Appendix~\ref{app:obstruction-discharges} to the retained D0-HN datum;
\((O1)\), \((O1)^+\), and \((O2)\) are reduced there to retained quotient-orientation, Weyl
transport, and type-II wall-atlas data.
Together these data supply the Pfaffian and orientation part of the
Pfaffian--Dirac theorem
\[
  \mathrm{Pf}_{\mathrm{prot}}(\mathfrak D^{DI}_X) \;=\; \Delta_5,
  \qquad
  \epsilon_o \;=\; \nu_{\Delta_5}.
\]
Here \(\epsilon_o\) is the orientation character of the Pfaffian line on
\(W^{(2)}(\La^{2,1}_{II})\), and \(\nu_{\Delta_5}\) is the Maass
multiplier of \cite{MAASS:1} on \(\Sp_4(\Z)\), pulled back along the
exterior-square isogeny \(\PSp_4(\R)\simeq\SO(3,2)_+\) and the
inclusion \(\La^{2,1}_{II}\hookrightarrow\La^{3,2}\) to a character on
\(W^{(2)}(\La^{2,1}_{II})\) taking value \(-1\) on each of the three
simple type-II reflections \(s_{\delta_i}\); see
Theorem~\ref{thm:pentadic-pfaffian-dirac}.  The symbol
\(\mathcal P_X\) denotes the protected primitive Lie superalgebra
\(P_X^{\Pi,\mathrm{red}}\) of Chapter~5.6 (the \(H^0\) of the protected
primitives of the hybrid factorization sheaf
\(\mathcal F_X\) after Gram-projection pushforward
\(\overline\Pi^\Theta_{X,*}\) and Hopf-radical quotient
\(\overline{\operatorname{Rad}\langle\,,\,\rangle_X}\); the three
notations \(\mathcal P_X = P_X^{\Pi,\mathrm{red}} = P_X\) denote one
object, abbreviated \(P_X\) where the construction data are not in
play).  The Lie-superalgebra
recognition
\[
  \mathcal P_X \;\cong\; \mathfrak g_{\Delta_5}
\]
is the separate primitive-recognition theorem of Chapter~\ref{ch:pfaffian-dirac},
conditional on the finite compact-source datum \((S1)\)--\((S10)\).

\item[\textcircled{4}]
\textsl{Mirror discriminant.}
\textsl{Conjectural.} On the rank-\(3\) Mukai-invariant
\((K3\times E)\)-period component, the discriminant section is
conjecturally \(\Delta_5^{-2}\); its polar divisor is twice the
reduced mirror discriminant, and its reduced automorphic Heegner support
is the Humbert divisor where \(\Delta_5\) vanishes
(Chapter~\ref{ch:view4-mirror}).

\item[\textcircled{5}]
\textsl{Singular theta lift.}
The theta decomposition \(\phi^{\mathrm{Weil}}_{0,1}\) of the
weight-zero index-one weak Jacobi form \(\phi_{0,1}\) has weight
\(-1/2\) for the Weil representation of \(\La^{3,2}\); its Borcherds
regularized theta lift is an orthogonal automorphic form, and the
exceptional isogeny \(\PSp_4(\R)\simeq\SO(3,2)_+\) identifies
that form with the genus-two Siegel cusp form
\[
  \Theta_{\mathrm{sing}}^{\mathrm{Bor}}(\phi^{\mathrm{Weil}}_{0,1}) \;=\; \Delta_5.
\]
This is the Borcherds--Gritsenko--Nikulin theta identity of
\cite{Borcherds95, B, GN}.
\end{enumerate}

\paragraph{The cross-level identities.}
The three edges
\(\textcircled{1}\!\leftrightarrow\!\textcircled{2}\),
\(\textcircled{2}\!\leftrightarrow\!\textcircled{5}\),
\(\textcircled{1}\!\leftrightarrow\!\textcircled{5}\) are theorems
(Borcherds 1995; Gritsenko--Nikulin 1998); they identify the
automorphic, BKM, and regularized-theta descriptions. The edge
\(\textcircled{2}\!\leftrightarrow\!\textcircled{3}\) is the
conditional Pfaffian--Dirac theorem after \((D0)\), \((O1)\),
\((O1)^+\), and the retained type-II wall-atlas datum
\((O2_{\mathrm{atlas}})\), together with the finite geometric Pfaffian
datum \((P_{\mathrm{fin}})\), on the \(K3\times E\) specialization in
Appendix~\ref{app:obstruction-discharges}
(Theorems~\ref{thm:G-D0}, \ref{thm:G-O1}, \ref{thm:G-O1-plus},
\ref{thm:G-O2-orbit-transport}), making
\(\textcircled{2}\!\leftrightarrow\!\textcircled{3}\) a theorem at the
Pfaffian and orientation level on \(K3\times E\); primitive
Lie-superalgebra recognition remains conditional on the cofinal
finite source-recognition datum and the finite Koszul comparison
datum of Definition~\ref{def:finite-koszul-comparison-datum}.  This
edge relates the BKM denominator algebra to the Dirac--Igusa Pfaffian.
The edge
\(\textcircled{1}\!\leftrightarrow\!\textcircled{3}\) is the
two-step composite \(\mathcal D_X = \Delta_5 = \mathrm{Pf}_{\mathrm{prot}}
(\mathfrak D^{DI}_X)\), with the orientation match
\(\epsilon_o = \nu_{\Delta_5}\) on \(W^{(2)}(\La^{2,1}_{II})\) supplying
the multiplier.  The protected primitive Lie superalgebra \(\mathcal P_X\)
recovers the BKM target precisely when the source representatives,
relations, pairings, radical quotient, PBW comparison, and full parity
dimensions of \((S1)\)--\((S10)\) are supplied.

\paragraph{The level-\(n\) scalar shadow.}
On the closed cobordism \(K3\times E\mapsto\mathrm{pt}\),
\[
  Z^{K3\times E}_{\mathrm{BPS}}
  \;=\; \bigl(\Delta_{10}\bigr)^{-1}
  \;=\; \Delta_5^{-2},
  \qquad
  \Delta_{10} \;=\; \Delta_5^2
  \;\in\; M_{10}\bigl(\Sp_4(\Z)\bigr).
\]
Here \(\Delta_{10} = \Delta_5^2\) is the theta-normalized weight-ten
cusp form on \(\Sp_4(\Z)\), equivalently \(\Phi_{10}^{\mathrm{un}} =
\chi_{10} = \Delta_5^2\) in the unscaled Igusa convention
(\cite{Igusa1962}; the multiplier \(\nu_{\Delta_5}\) squares to
trivial \cite{MAASS:1}); the Oberdieck--Pixton-monic normalization
\(\chi_{10}^{\mathrm{OP}} := D_5^2 = 4096^{-1}\Delta_5^2\) (with
\(D_5 := 64^{-1}\Delta_5\) the Borcherds-monic form so that
\([q^{1/2}r^{1/2}s^{1/2}]D_5 = 1\)) gives the equivalent form
\(Z^X_{\mathrm{OP}} = -(\chi_{10}^{\mathrm{OP}})^{-1} =
-4096\,\Delta_5^{-2}\), where \(4096 = 64^2\) is the squared
theta-leading constant
\([q^{1/2}r^{1/2}s^{1/2}]\Delta_5 = 64\) of the Igusa
form and the minus sign is the OP scalar branch convention
\cite{OB, OPand}.
\(Z^{K3\times E}_{\mathrm{BPS}}\) is the level-\(n\) protected trace of
the retained \(K3\times E\) data.  A three-dimensional gravity-line
interpretation with internal compact factor \(K3\times E\) requires the
Hall--Borcherds residual of Vol II and is not asserted here.

\paragraph{The \(\kappa\)-tuple of \(K3\times E\).}
The retained \(K3\times E\) data carry the five-coordinate
\(\kappa\)-tuple
\[
  \kappa_{K3\times E}
  \;=\;
  \bigl(\kappa_{\mathrm{cat}},\;\kappa^{\mathrm{Hodge}}_{\mathrm{ch}},
  \;\kappa^{\mathrm{Heis}}_{\mathrm{ch}},\;\kappa_{\mathrm{BKM}},
  \;\kappa_{\mathrm{fiber}}\bigr)
  \;=\; (0,\,0,\,3,\,5,\,24).
\]
The categorical Euler is \(\kappa_{\mathrm{cat}} =
\chi(\mathcal O_{K3\times E}) = \chi(\mathcal O_{K3})\,\chi(\mathcal O_E)
= 2\cdot 0 = 0\). The Hodge supertrace
\(\kappa^{\mathrm{Hodge}}_{\mathrm{ch}}(K3\times E) =
\sum_q (-1)^q h^{0,q}(K3\times E) = 0\). The Heisenberg rank
\(\kappa^{\mathrm{Heis}}_{\mathrm{ch}} = 3\) is the rank of the
type-II Heisenberg--Cartan obtained after the normal-ordered Gram
projection \(\Pi_X\) to \(\La^{2,1}_{II}\).  It is not the rank of the
algebraic Mukai lattice \(H^0(S,\Z)\oplus\NS(S)\oplus H^4(S,\Z)\);
for the \(U\oplus\langle-2\rangle\) polarisation that lattice has rank
five, while the \(\kappa^{\mathrm{Heis}}_{\mathrm{ch}}\)-coordinate
counts only the three Gram coordinates which enter the type-II BKM
Cartan. The BKM coordinate
\[
  \kappa_{\mathrm{BKM}}(\Phi_N) \;=\; c_N(0)/2
\]
(Borcherds 1995; Gritsenko 1999) reads, at the paramodular input
\(N=1\), \(c_1(0) = 10\), giving
\(\kappa_{\mathrm{BKM}}(\Delta_5) = 5\). The fibre coordinate
\(\kappa_{\mathrm{fiber}} = \chi_{\mathrm{top}}(K3) = 24\). The additive
identity
\(\kappa_{\mathrm{BKM}} = \kappa^{\mathrm{Hodge}}_{\mathrm{ch}}
+ \chi(\mathcal O_{\mathrm{fiber}})\) fails at every
\(N\in\{1,2,3,4,6\}\): the Lefschetz Euler \(\chi_{\mathrm{top}}(K3) =
24\) and the holomorphic Euler \(\chi(\mathcal O_{K3}) = 2\) are
distinct invariants of the K3 fibre, and \(\kappa_{\mathrm{BKM}}\) sees
the \(c_N(0)/2\) data of the Borcherds input rather than either of
them. Each subscript on \(\kappa\) records a distinct invariant;
the unsubscripted handle leaves the categorical dimension unspecified.

\chapter{Pentadic \texorpdfstring{\(\Delta_5\)}{Delta\_5}}
\label{ch:pentadic}

\section{The retained \texorpdfstring{\(E_3\)}{E\_3}-object and the five
readings of \texorpdfstring{\(\Delta_5\)}{Delta\_5}}
\label{sec:retained-e3-five-readings}

The five readings attached to \(\Delta_5\) have different logical
status.  The automorphic, BKM, and theta-lift readings are theorems;
the Pfaffian reading is relative to retained \(K3\times E\) data; the
mirror reading is conjectural.  The
trivial-canonical threefold \(K3\times E\) supplies the ambient CY-\(3\)
input for a retained holomorphic prefactorization datum in the
Kontsevich--Tamarkin and Costello--Gwilliam--Li lineage
\cite{KontsevichOperadsMotives, CostelloGwilliam1, CostelloGwilliam2}.
The notation
\[
  \mathfrak A_{K3\times E}
  \;\in\; E_3\text{-}\mathrm{HolFA}(K3\times E)
\]
denotes this chosen finite-stage \(E_3\)-prefactorization datum; the
unconditional functor from CY-\(3\) dg-categories to holomorphic
\(E_3\)-factorization algebras is not asserted
\cite{KontsevichOperadsMotives, TamarkinLittleDisks,
LambrechtsVolicLittleDisks, CostelloGwilliam1, CostelloGwilliam2}.
The retained protected sector supplies a Pfaffian line on the
Dirac--Igusa pro-object.  The Pfaffian--Dirac theorem compares that line
with the automorphic section \(\mathcal D_X=\Delta_5\) only under the named
orientation and wall data.  The five readings are the automorphic
section, the BKM denominator, the retained Pfaffian section, the
conjectural period discriminant, and the singular theta lift.

The unconditional new mathematical content of the manuscript, beyond
the classical Igusa--Borcherds--Gritsenko--Nikulin spine, is the
parity-resolved root-space computation at
\(\delta_{123}=\delta_1+\delta_2+\delta_3\) --- even dimension \(29\)
against odd dimension \(93\), with \(29-93=-64=f(1,1)\)
(Proposition~\ref{prop:bkm-delta123-presentation-count}) --- together
with the scalar-Pfaffian monodromy proposition
(Proposition~\ref{prop:scalar-pfaffian-data-not-section}).  The
Pfaffian--Dirac identity
\(\mathrm{Pf}_{\mathrm{prot}}(\mathfrak D^{DI}_X)=\Delta_5\) is a
recognition criterion relative to the retained data \((D0)\), \((O1)\),
\((O1)^+\), \((O2_{\mathrm{atlas}})\), \((P_{\mathrm{fin}})\), which the
certificate ledgers record as not supplied.

\begin{definition}[Pentadic \(\Delta_5\)]
\label{def:pentadic-delta5}
The five-tuple
\[
  \bigl(\mathcal D_X,\;
  \operatorname{den}(\mathfrak g_{\Delta_5}),\;
  \mathrm{Pf}_{\mathrm{prot}}(\mathfrak D^{DI}_X),\;
  \mathrm{Disc}(\mathcal P\to\mathbb D),\;
  \Theta_{\mathrm{sing}}^{\mathrm{Bor}}(\phi^{\mathrm{Weil}}_{0,1})\bigr)
\]
of automorphic, BKM, retained-Pfaffian, period, and theta-lift readings
is the
\emph{pentadic \(\Delta_5\)}. The reading of \(\Delta_5\) is
fixed by the level: at level \textcircled{1}, \(\Delta_5\) is
\(\mathcal D_X\); at level \textcircled{2}, \(\Delta_5\) is the function whose
value at \(2Z\) is \(64\,\operatorname{den}(\mathfrak g_{\Delta_5})\);
at level \textcircled{3}, \(\Delta_5\) is the protected Pfaffian
\(\mathrm{Pf}_{\mathrm{prot}}(\mathfrak D^{DI}_X)\) under
\((D0)\), \((O1)\), \((O1)^+\), and the retained
\((O2_{\mathrm{atlas}})\) datum; at level
\textcircled{4}, \(\Delta_5^{-2}\) is the mirror discriminant
section, conjecturally; at level \textcircled{5}, \(\Delta_5\) is the
Borcherds theta lift of \(\phi^{\mathrm{Weil}}_{0,1}\).
\end{definition}

\section{The \texorpdfstring{$\textcircled{1}\!\leftrightarrow\!
\textcircled{2}$}{1-2} edge: Borcherds--Gritsenko--Nikulin}
\label{sec:edge-1-2}

The Borcherds product expansion of the weight-zero index-one weak
Jacobi form \(\phi_{0,1}\) is the genus-two cusp form
\[
  \Delta_5(Z)
  \;=\; 64\,q^{1/2}r^{1/2}s^{1/2}
  \prod_{(n,l,m)\in\Gamma_{\mathrm{eff}}}
  \bigl(1 - q^n r^l s^m\bigr)^{f(nm,l)},
\]
where \(Z = \bigl(\begin{smallmatrix} z_1 & z_2 \\ z_2 & z_3 \end{smallmatrix}\bigr)\),
\(q = e^{2\pi i z_1}\), \(r = e^{2\pi i z_2}\), \(s = e^{2\pi i z_3}\),
and \(f(n,l)\) are the Fourier coefficients of \(\phi_{0,1}\) (the
weak Jacobi form whose discriminant decomposition records
\(f(0,0) = 10\), \(f(0,\pm 1) = 1\), \(f(1,0) = 108\),
\(f(1,\pm 1) = -64\), \(f(1,\pm 2) = 10\), \dots)
\cite{Borcherds95, GN, EZ:1}. The leading constant
\([q^{1/2}r^{1/2}s^{1/2}]\Delta_5 = 64 = 2^6\) is the theta-leading
coefficient (Igusa); the chamber
\(0 < |s| \ll |q| \ll |r|^{-1} \ll 1\) fixes the active semigroup
\(\Gamma_{\mathrm{eff}}\) of cusp-completed exponents.

The Gritsenko--Nikulin denominator identity converts the same product
into the Borcherds--Weyl--Kac form on the type-II sublattice
\(\La^{2,1}_{II} \subset \La^{2,1} = \HypPlan\oplus\langle 2\rangle\):
\[
  \operatorname{den}(\mathfrak g_{\Delta_5})
  \;=\; \sum_{w\in W^{(2)}(\La^{2,1}_{II})}
  \mathrm{sgn}(w)\;w\Bigl(e^{\rho}
  \;-\sum_{a\in\La^{2,1}_{II}\cap\,\R_{>0}\Poly_{II}}
  m(a)\,e^{\rho+a}\Bigr)
  \;=\; 64^{-1}\,\Delta_5(2Z),
\]
where \(\Poly_{II}\) is the type-II Weyl chamber (the intersection of
the positive half-spaces of the three simple walls,
Theorem~\ref{thm:bkm-weyl-alternating-sum}) and \(m(a)\in\Z\) are the
imaginary-simple correction multiplicities read from the active support
of \(\phi_{0,1}\) \cite[Theorem~2.1]{GN}, with Weyl vector
\[
  \rho \;=\; \tfrac{1}{2}\bigl(\delta_1 + \delta_2 + \delta_3\bigr)
       \;=\; f_2 - \tfrac{1}{2}\,f_3 + f_{-2}
\]
on the three simple type-II reflections
\(\delta_1 = 2f_2 - f_3\), \(\delta_2 = 2f_{-2} - f_3\),
\(\delta_3 = f_3\) (Cartan matrix
\((\delta_i,\delta_j) = 4\Id_3 - 2\mathbf{J}_3 = \bigl(\begin{smallmatrix}
2 & -2 & -2\\ -2 & 2 & -2\\ -2 & -2 & 2 \end{smallmatrix}\bigr)\),
\(\mathbf{J}_3\) the all-ones matrix); type-II reflection subgroup
\(W^{(2)}(\La^{2,1}_{II})\) generated by these three reflections
\cite{GN, GNPartI, GNII}. The imaginary-simple correction terms
\(-\sum_a m(a)\,e^{\rho+a}\) are the content of the Gritsenko--Nikulin
theorem: the naive alternating sum
\(\sum_w \mathrm{sgn}(w)\,w(e^{\rho})\) alone is the denominator of the
uncorrected hyperbolic Kac--Moody algebra of \(W^{(2)}\) and does not
equal \(64^{-1}\Delta_5(2Z)\) --- the exponent \(\rho + 2f_2\) carries
coefficient \(-9\) in \(64^{-1}\Delta_5(2Z)\) yet lies on no Weyl
translate of \(\rho\), so it must be accounted by imaginary simple
roots. The signed root supermultiplicities recover
the Fourier coefficients of \(\phi_{0,1}\) under the doubling
\(\alpha(n,l,m) = 2nf_2 - lf_3 + 2mf_{-2}\):
\[
  \smult\,\mathfrak g_{\Delta_5,\alpha(n,l,m)}
  \;=\; f(nm,l) \qquad \text{on } (n,l,m)\in\Gamma_{\mathrm{act}}
  =\{(n,l,m)\in\Gamma_{\mathrm{eff}}\mid f(nm,l)\ne 0\}.
\]

The edge \(\textcircled{1}\!\leftrightarrow\!\textcircled{2}\) is the
joint statement: the automorphic section \(\mathcal D_X = \Delta_5\) of
\(\mathcal L^{\otimes 5}\otimes\nu_{\Delta_5}\) on Siegel space
\textsl{equals} the BKM denominator (after the constant \(64^{-1}\) and
the variable doubling \(Z\mapsto 2Z\)).  This is Borcherds 1995
together with the Gritsenko--Nikulin identification of the underlying lattice
with \(\La^{2,1}_{II}\) and the Weyl-group action with
\(W^{(2)}(\La^{2,1}_{II})\) \cite{Borcherds95, GN}.

\section{The \texorpdfstring{$\textcircled{1}\!\leftrightarrow\!
\textcircled{5}$}{1-5} and \texorpdfstring{$\textcircled{2}\!
\leftrightarrow\!\textcircled{5}$}{2-5} edges: the singular theta lift}
\label{sec:edge-1-5}

The Borcherds product expansion of the previous section converts
\(\phi_{0,1}\) into \(\Delta_5\) by an arithmetic identity in
\(\Gamma_{\mathrm{eff}}\); the question of \emph{where} this identity
comes from has a single answer.  The exceptional isomorphism
for \(\PSp_4(\R)\simeq\SO(3,2)_+\) identifies the genus-two Siegel
space with the type-IV symmetric domain attached to the
signature-\((3,2)\) lattice
\(\La^{3,2} = \HypPlan\oplus\HypPlan\oplus\langle 2\rangle\), and it
identifies the arithmetic quotient
\(\Sp_4(\Z)\backslash\mathcal H_2\) with
\(\Gamma_{3,2}\backslash\Hpl^{\IV}_+\), where
\(\Gamma_{3,2}=\wedge^2\Sp_4(\Z)/\{\pm\Id\}\).
Borcherds's regularisation pairs the \(\Mp_2(\Z)\)
Weil-representation form \(\phi^{\mathrm{Weil}}\) with the Siegel
theta kernel of \(\La^{3,2}\), and produces a meromorphic Siegel
modular form whose divisor is
prescribed by the principal part:
\[
  \Theta_{\mathrm{sing}}^{\mathrm{Bor}}(\phi)
  \;=\;
  \int^{\mathrm{reg}}_{\mathrm{SL}_2(\Z)\backslash\Hpl}
  \bigl\langle \phi(\tau),\,
  \Theta_{\La^{3,2}}(\tau, Z)\bigr\rangle\,
  \frac{du\,dv}{v^2},
  \qquad \tau = u + iv,
\]
on the Weil-representation pairing of the vector-valued
weakly-holomorphic input \(\phi\) (the theta decomposition of
\(\phi_{0,1}\) into the discriminant form of \(\La^{3,2}\)) against
the Siegel theta series of the signature-\((3,2)\) lattice
\(\La^{3,2}\) — not the rank-three target lattice \(\La^{2,1}_{II}\),
which receives the BKM root structure under doubling — together with
the antiholomorphic dependence required for absolute convergence after
regularisation; the explicit construction is Chapter~\ref{ch:singular-theta-lift},
Definition~\ref{def:siegel-theta} and
Theorem~\ref{thm:singular-theta-lift} \cite{Borcherds95, B}.  At the
theta-decomposed input \(\phi^{\mathrm{Weil}}_{0,1}\) the regularised
integral computes the composite
\[
  \Theta_{\mathrm{sing}}^{\mathrm{Bor}}\colon
  J^{\mathrm{wk}}_{0,1}
  \;\longrightarrow\;
  M^{!}_{-1/2}(\rho_{\La^{3,2}})
  \;\longrightarrow\;
  M_5\bigl(\Sp_4(\Z),\,\nu_{\Delta_5}\bigr),
  \qquad
  \phi_{0,1} \;\longmapsto\;
  \phi^{\mathrm{Weil}}_{0,1} \;\longmapsto\; \Delta_5,
\]
giving the edge \(\textcircled{1}\!\leftrightarrow\!\textcircled{5}\).
Reading the same lift through the Gritsenko--Nikulin denominator
identity, \(\phi_{0,1}\) is the character of the index-one Heisenberg
sector and the multiplicative lift over \(\La^{2,1}_{II}\) is exactly
\(\operatorname{den}(\mathfrak g_{\Delta_5}) = 64^{-1}\Delta_5(2Z)\),
giving the edge \(\textcircled{2}\!\leftrightarrow\!\textcircled{5}\)
\cite{Borcherds95, GN, GNPartI, GNII}.

\section{The \texorpdfstring{$\textcircled{2}\!\leftrightarrow\!
\textcircled{3}$}{2-3} edge: the Pfaffian--Dirac theorem}
\label{sec:edge-2-3}

The theorem edges proved by Borcherds and Gritsenko--Nikulin identify
the automorphic, BKM-denominator, and singular-theta readings of
\(\Delta_5\); the chiral
\(E_3\)-prefactorization datum of \S\ref{sec:retained-e3-five-readings}
is the retained operator-level source from which the View~\textcircled{3}
datum is obtained.  The structural question has two terms.  The Pfaffian term
asks whether the protected determinant of the \(K3\times E\) operator
datum is the automorphic section \(\Delta_5\).  The primitive term
asks whether the compact Hall primitive sector is the BKM algebra
\(\mathfrak g_{\Delta_5}\).  The first term is the Pfaffian--Dirac
theorem after \((D0)\), \((O1)\), \((O1)^+\), and
\((O2_{\mathrm{atlas}})\), together with \((P_{\mathrm{fin}})\), in
Appendix~\ref{app:obstruction-discharges}.  The second term is the
finite compact-source recognition problem \((S1)\)--\((S10)\).

\begin{theorem}[Pfaffian--Dirac]
\label{thm:pentadic-pfaffian-dirac}
\textsl{Conditional on \((D0)\), \((O1)\), \((O1)^+\),
\((O2_{\mathrm{atlas}})\), and \((P_{\mathrm{fin}})\);
ambient: cosection-twisted reduced d-critical chart on \(K3\times E\)
at finite cusp-truncation height \(R\), with chain-level pro-object,
Mittag--Leffler vanishing of the inverse system in \(R\), and Pfaffian
orientation \(o\) with Weyl character \(\epsilon_o\).} The
protected Pfaffian of the Dirac--Igusa pro-object
\(\mathfrak D^{DI}_X = \varprojlim_R\,\mathfrak D^{DI}_{X,R}\) satisfies
\[
  \mathrm{Pf}_{\mathrm{prot}}(\mathfrak D^{DI}_X)
  \;=\; \Delta_5,
  \qquad
  \epsilon_o
  \;=\; \nu_{\Delta_5}\quad\text{on}\quad
  W^{(2)}(\La^{2,1}_{II}),
\]
\end{theorem}

The primitive Lie-superalgebra assertion is a separate recognition
statement.  If the finite compact-source datum \((S1)\)--\((S10)\) is
supplied on a cofinal system of relation-closed degree windows, then
\[
  \mathcal P_X
  :=
  H^0\,\Prim_{\mathrm{prot}}\bigl(\overline\Pi^{\Theta}_{X,*}
  \mathcal F_X\bigr) /
  \overline{\operatorname{Rad}\langle\,,\rangle_X}
  \;\cong\; \mathfrak g_{\Delta_5}.
\]

The four obstructions and the finite geometric Pfaffian datum name the
exact conditioning. \((D0)\) is the
\(D0\)-degeneration limit of the cosection-reduced d-critical
orientation theory: the limiting orientation in the \(D0\)-degenerate
fibre is constructed and matches the orientation of the generic fibre
under the cosection-twist
\cite{BBDJS2015, JoyceUpmeier, JoyceTanakaUpmeier}. \((O1)\) is the
strong reduced orientation on the \(K3\times E\)-quotient self-Ext
frame bundle: the square root of the determinant line of the cotangent
complex extends across the cusp boundary as part of the retained
quotient-orientation datum
\cite{CaoGrossJoyce, JoyceUpmeier, JoyceTanakaUpmeier}. \((O1)^+\) is the Weyl-equivariant
transport of the orientation along reflections in the type-II
reflection subgroup \(W^{(2)}(\La^{2,1}_{II})\): the orientation cocycle
\(c_o\in Z^2(W^{(2)}(\La^{2,1}_{II});\mathbb F_2)\) trivializes in
\(H^2\), and the resulting character
\(\epsilon_o:W^{(2)}(\La^{2,1}_{II})\to\{\pm1\}\) is well-defined on
generators. \((O2)\) is the local Pfaffian wall normal form on type-II
walls: at each simple type-II wall \(s_{\delta_i}\) the rank-one normal
Pfaffian crossing of Definition~\ref{def:rank-one-type-ii-crossing}
is part of the retained wall-atlas datum and gives the local sign
\((-1)^{N^{\mathrm{Pf}}_{\delta_i}} = -1\), and the Hall-side sign
sum on the half-Hilbert orbit of the self-Ext stack of size
\(2^{12}=4096\) is compared with the squared theta-leading
\(([q^{1/2}r^{1/2}s^{1/2}]\Delta_5)^2\) only after scalarization.  The
wall transport gives the character value
\(\epsilon_o(s_{\delta_i}) = \nu_{\Delta_5}(s_{\delta_i})\) on each
generator (Appendix~E).
The finite geometric Pfaffian datum gives the skew perfect complexes,
Pfaffian lines, finite Pfaffian sections, type-II wall divisor orders,
automorphic line comparisons, support and parity rows, leading
coefficient, and Pfaffian-line Mittag--Leffler compatibility for the
cofinal HN product.

\begin{remark}[The conditioning is named precisely]
\label{rem:conditioning-named}
The four obstruction classes \((D0)\), \((O1)\), \((O1)^+\), \((O2)\)
and the finite geometric Pfaffian datum are the exact Pfaffian and
orientation obligations; Appendix~\ref{app:obstruction-discharges}
reduces \((D0)\) to the retained D0-HN datum, names the retained
quotient-orientation and Weyl transport data, and names
\((O2_{\mathrm{atlas}})\) and \((P_{\mathrm{fin}})\). At the operator level
\textcircled{3}, \(\mathcal D_X=\Delta_5\) is compared with the protected
Pfaffian \(\mathrm{Pf}_{\mathrm{prot}}(\mathfrak D^{DI}_X)\).  After
the level-\(n\) protected trace, the scalar partition function is
\(\Delta_5^{-2}\). The primitive algebra
\(\mathcal P_X\simeq\mathfrak g_{\Delta_5}\) is stronger; it requires
the finite compact-source recognition datum \((S1)\)--\((S10)\). The
operator-level datum carries the orientation character
\(\epsilon_o\) that the scalar shadow has squared away, by view
\textcircled{3}, and recovers it as \(\nu_{\Delta_5}\) on
\(W^{(2)}(\La^{2,1}_{II})\) by \((O1)^+\) and \((O2)\).
\end{remark}

\section{View \texorpdfstring{$\textcircled{4}$}{4}: mirror
discriminant}
\label{sec:view-4-mirror}

\textsl{Conjectural} (ambient: variation of Hodge structure on the
rank-\(3\) Mukai-invariant \((K3\times E)\)-period component).
On \(\mathcal P_{K3\times E}^{(3)}\simeq \mathcal H_2/\Sp_4(\Z)\), the
mirror discriminant section is conjecturally \(\Delta_5^{-2}\).  Its
polar divisor is \(2\mathfrak D_{\mathrm{mir}}\), where the reduced
support of \(\mathfrak D_{\mathrm{mir}}\) is the Humbert--Heegner
divisor \(\mathfrak H_{\mathrm{II}}\) fixed by the Borcherds product of
\(\Delta_5\).  This mirror-discriminant identification is conjectural
and is not a hypothesis of any
theorem; see Conjecture~\ref{conj:mirror-discriminant}.

\section{The level-\texorpdfstring{$n$}{n} scalar shadow}
\label{sec:level-n-scalar-shadow}

The scalar partition function on the closed cobordism
\(K3\times E\mapsto\mathrm{pt}\) is the Borcherds determinant of the
protected K3 index, computed by Oberdieck--Pixton 2018 and
Oberdieck--Pandharipande 2018 \cite{OB, OPand}: the OP partition
function \(Z^X_{\mathrm{OP}} = -4096\,\Delta_5^{-2}\) is the virtual
count of Hilbert schemes of \(K3\times E\) in the
\((-p_{\mathrm{DT}})^n\) DT-chamber, and the BPS normalization
\(Z^{K3\times E}_{\mathrm{BPS}} = \Delta_5^{-2}\) rescales by
\(-4096^{-1}\).
Two normalizations of the same weight-ten cusp form on \(\Sp_4(\Z)\)
are in use, and naming them is necessary because both appear in the
formulas:
\begin{itemize}
\item \emph{Theta-normalized.}
\(\Delta_{10}(Z) := \Delta_5(Z)^2 \in M_{10}(\Sp_4(\Z))\) (the
multiplier \(\nu_{\Delta_5}\) squares to trivial; the leading
coefficient \([q\,r\,s]\Delta_{10} = 64^2 = 4096\)).
\item \emph{Oberdieck--Pixton-monic.}
\(\chi_{10}^{\mathrm{OP}}(Z) := D_5(Z)^2 = 64^{-2}\Delta_5(Z)^2 =
4096^{-1}\Delta_{10}(Z)\) (so \([q\,r\,s]\chi_{10}^{\mathrm{OP}} = 1\)).
\end{itemize}
Oberdieck--Pixton and Oberdieck--Pandharipande prove
\cite{OB, OPand}
\[
  Z^X_{\mathrm{OP}}
  \;=\; -\bigl(\chi_{10}^{\mathrm{OP}}\bigr)^{-1}
  \;=\; -4096\,\Delta_5^{-2},
\]
and the level-\(n\) protected scalar shadow is
\[
  Z^{K3\times E}_{\mathrm{BPS}}
  \;=\; \bigl(\Delta_{10}\bigr)^{-1}
  \;=\; \Delta_5^{-2}
  \;=\; -\tfrac{1}{4096}\,Z^X_{\mathrm{OP}}.
\]
The factor \(4096 = 64^2\) is the squared theta-leading constant of
\(\Delta_5\); the OP minus sign is a scalar branch convention absorbed
into the OP \((-p_{\mathrm{DT}})^n\) chamber and is independent of the
orientation character \(\epsilon_o\) of view~\textcircled{3}.

\(Z^{K3\times E}_{\mathrm{BPS}}\) is the level-\(n\) protected trace of
the retained \(K3\times E\) data: a charge-graded virtual partition
function.  A three-dimensional gravity-line interpretation of the
scalar shadow with internal compact factor \(K3\times E\) requires the
Hall--Borcherds residual of Vol II.  The universal-stage
chain names the required objects: primitive factorization dg-category
\(\mathcal C\) → boundary chart \(A_b = \End_{\mathcal C}(b)\) → bar
twisting \(\mathrm{Bar}(A_b)\) → derived chiral centre
\(Z^{\mathrm{der}}_{\mathrm{ch}}(A_b)\) → ambient operator algebra
\(\mathcal A\) at the boundary vacuum \(b\) of the open
factorization dg-category on
\((K3\times E,\,D,\,\tau)\); this chain is not constructed here.
The local witness is the formal-Darboux stalk of the
mixed-holomorphic-topological strings theory
\cite{KontsevichOperadsMotives, CostelloLiBCOV,
CostelloLiOpenClosedBCOV}.  The class
\(\mathrm{ob}^{\mathrm{hol}}_{\mathrm{dR}}(K3\times E)\) is the
holomorphic de Rham obstruction to gluing this local BV factorization
model over \(K3\times E\); its vanishing is a necessary global
condition, not the gravity-line operator algebra.

\begin{remark}[Primitive order]
\label{rem:beilinson-cut-delta5}
The five views are ordered: primitive objects (the retained chiral
\(E_3\)-prefactorization datum on \(K3\times E\), its four obstructions
\((D0),(O1),(O1)^+,(O2_{\mathrm{atlas}})\), and its finite Pfaffian
section datum; the Cartan datum of \(\mathfrak g_{\Delta_5}\))
first; cross-level identities ((②↔③) Pfaffian--Dirac, (\(\epsilon_o
=\nu_{\Delta_5}\)) orientation match, \((\mathcal P_X\cong\mathfrak
g_{\Delta_5})\) primitive recognition) second; scalar modular forms
last (the section \(\mathcal D_X\), the BKM denominator
\(64^{-1}\Delta_5(2Z)\), the partition function \(\Delta_5^{-2}\)).
A statement is not primitive if it is true only after passing to a
trace, choosing a chamber, completing a category, or installing
descent data.  On each first appearance the scalar reading is named:
the theta-normalized form \(\Delta_5\) with leading
\([q^{1/2}r^{1/2}s^{1/2}]\Delta_5 = 64\); the BKM denominator
\(64^{-1}\Delta_5(2Z)\) with the doubling \(2Z\) on the type-II
sublattice; the OP-monic form \(D_5 = 64^{-1}\Delta_5\) with
\(\chi_{10}^{\mathrm{OP}} = D_5^2 = 4096^{-1}\Delta_5^2\) the
Oberdieck--Pixton normalization of the squared determinant section.
\end{remark}

\section{Placement in the Beilinson Chain}
\label{sec:beilinson-chain-placement}

The pentadic \(\Delta_5\) is one identification at one specific level
of the Beilinson chain. Under the
chiral Koszul source datum of
Definition~\ref{def:chiral-koszul-source-datum}, the finite Koszul
comparison datum of
Definition~\ref{def:finite-koszul-comparison-datum}, and the strict
Koszul Mittag--Leffler exactness of
Proposition~\ref{prop:chiral-koszul-source-datum-consequence}, the
source coalgebra \(C_X\) maps by cobar comparison to the target chiral
shadow \(\mathsf{Den}_{\Delta_5,E}\), compatibly with Weyl transport,
Pfaffian orientation, Hall pairing, radical quotient, and PBW
filtration.  Vol I's Theorem H supplies the
chain-level identification
\(Z^{\mathrm{der}}_{\mathrm{ch}}(A_b) =
\mathrm{ChirHoch}^\bullet(A_b,A_b)\) only on the retained Koszul
locus, and Vol II's chiral-Hochschild Beilinson stratification
\textnormal{(i)+(ii)} supplies the CY-orientation trace pairing of
degree \(-\dim_{\mathbb C}(K3\times E) = -3\).  Relative to the
retained trace datum, these inputs give the level-\(\mathsf Z\) trace
identity
\(\mathrm{Tr}^{\mathrm{bulk}}_n((-1)^F\cdot\mathrm{id})_{Z^{\mathrm{der}}_{\mathrm{ch}}(A_b)}
	= \Delta_5^{-2}\) of Appendix~\ref{app:obstruction-discharges},
Theorem~\ref{thm:G-vol2-discharge}, relative to \((Z_{\mathrm{HB}})\);
the OP chamber scalar is
\(-4096\) times this trace. The level-\(\mathsf A\)
gravity-line operator algebra acting on the boundary is the open
gravity-line Hall--Borcherds residual of Vol II.  The Vol III
\(\kappa\)-stratification places
\(\mathfrak g_{\Delta_5}\) at the \((d=3,\,\text{rank }3)\) chamber
through the identity \(\kappa_{\mathrm{BKM}}(\Phi_N) = c_N(0)/2\),
where \(c_N(0)\) is the constant Fourier coefficient of the input weak
Jacobi or vector-valued modular form. Applied at the paramodular row
\((t,N;k) = (1,1;5)\) of the Gritsenko--Cl\'ery catalogue
\cite{GC, GN}, with Jacobi seed \(\phi_{0,1}\) and Borcherds product
\(\Delta_5\), the identity reads
\(\kappa_{\mathrm{BKM}}(\Delta_5) = c_1(0)/2 = 5\).  The companion
identity at input \(\Phi_{12}\), the Fake-Monster Borcherds product
on the Leech-extension lattice \(\mathrm{II}_{25,1}\)
\cite{BorcherdsGKM88, Borcherds95}, reads
\(\kappa_{\mathrm{BKM}}(\Phi_{12}) = c_{\Phi_{12}}(0)/2 = 12\); the
two values are different rows of one universal \(c_N(0)/2\) rule,
indexed by different input forms and lattices, not in contradiction.
The mixed holomorphic-topological strings companion identifies the
formal-Darboux stalk
\(A^{\mathrm{cl}}_{\partial,N} = C^\bullet_{\mathrm{CE}}(\mathfrak{gl}_N,
\,\mathrm{Kosz}([\phi_1,\phi_2]))\) of the chiral Hochschild cochain
at the Dirac brane vacuum.  The obstruction to globalizing this local
factorization model over \(K3\times E\) is the holomorphic de Rham
class \(\mathrm{ob}^{\mathrm{hol}}_{\mathrm{dR}}(K3\times E)\).  A
three-dimensional gravitational path-integral interpretation of the
level-\(n\) scalar shadow
\(Z^{K3\times E}_{\mathrm{BPS}} = \Delta_5^{-2}\) requires this
obstruction to vanish and separately requires the Vol II gravity-line
Hall--Borcherds operator algebra.


\bigskip
\tableofcontents

% ------------------------------------------------------------------
% Part I: Given — the inputs of the construction
% ------------------------------------------------------------------
\part{Given}


\chapter{K3$\times$E geometry, Mukai dictionary, lattice polarization,
normal-ordered Gram cocycle, hybrid $\operatorname{Ran}^{\mathrm{hyb}}(E)$ carrier}
\label{ch:k3xe-setup}

The arithmetic source contains three named weight-five readings, and
their normalizations are different.  The automorphic chain produces the section
\[
\mathcal{D}_X = \Delta_5
\in M_5\!\bigl(\Sp_4(\Z),\,\nu_{\Delta_5}\bigr)
\qquad\text{(View \textcircled{1}; proved \cite{Borcherds95,GN}).}
\]
The Borcherds--Kac--Moody (BKM) chain produces the denominator
\[
\operatorname{den}(\mathfrak g_{\Delta_5}) = 64^{-1}\,\Delta_5(2Z)
\qquad\text{(View \textcircled{2}; proved \cite{GN}).}
\]
The compact Pfaffian chain produces, granted the retained D0-HN datum,
the retained quotient-orientation datum, the chosen Weyl lifts, and
the retained type-II wall-atlas datum of
Theorem~\ref{thm:G-four-obstruction-criterion} of
Appendix~\ref{app:obstruction-discharges}, the protected Pfaffian
\[
\operatorname{Pf}_{\mathrm{prot}}(\mathfrak D_X^{\mathrm{DI}}) = \mathcal D_X=\Delta_5
\qquad\text{(View \textcircled{3}; relative to the retained orientation data at the Pfaffian level).}
\]
The handle ``$\Delta_5$'' with no superscript is reserved for working
formulas; each occurrence names which of $\mathcal D_X$,
$\operatorname{den}(\mathfrak g_{\Delta_5})$, or
$\operatorname{Pf}_{\mathrm{prot}}(\mathfrak D_X^{\mathrm{DI}})$ is meant.
The cross-level identity \textcircled{2}$\leftrightarrow$\textcircled{3}
is a theorem at the Pfaffian and orientation level on \(K3\times E\),
granted \((\mathrm{D0})\), the retained quotient-orientation datum for
$(\mathrm{O1})$, the chosen Weyl lifts for $(\mathrm{O1}^+)$, and the
retained datum $(\mathrm{O2}_{\mathrm{atlas}})$
in Theorems~\ref{thm:G-D0}, \ref{thm:G-O1}, \ref{thm:G-O1-plus},
\ref{thm:G-O2-orbit-transport}.  The primitive Lie-superalgebra
recognition \(P_X^{\Pi,\mathrm{red}}\cong\mathfrak g_{\Delta_5}\)
requires the finite compact-source datum \((S1)\)--\((S10)\) of
Chapter~\ref{ch:pfaffian-dirac}.

\paragraph{The Mukai--Gram dictionary.}
The Mukai--Gram dictionary on \(K3\times E\) begins with the Mukai
lattice \(\widetilde H(S,\Z)\) of the K3 surface \(S\), of rank \(24\)
and signature \((4,20)\) \cite{Mukai1984}.  The even Kunneth sector
\[
\Gamma_X^{\mathrm{form}}
 =\widetilde H(S,\Z)\otimes H^{\mathrm{even}}(E,\Z)
 \cong \widetilde H(S,\Z)\oplus\widetilde H(S,\Z)
\]
is the formal \(K3\times E\) charge lattice used below; the two summands
are written \((Q,P)\).  For the Gram map this underlying Kunneth group
is read with the two componentwise K3 Mukai pairings; the tensor
Kunneth cross-pairing on \(H^{\mathrm{even}}(X,\Z)\) is a different
bilinear form.  The Gram-index lattice
\(\Gamma_{\mathrm{gram}}=\Z^3\) receives the quadratic shadow
\[
 \Pi_X(Q,P)=\bigl((Q,Q)/2,\,(Q,P),\,(P,P)/2\bigr).
\]
This map is not additive; its additivity defect is the symmetric
polarization \(B=-\delta\Pi_X\).  The additive lattice map is the
normal-ordered homomorphism
\[
 \overline\Pi_X:\widehat\Gamma_X\longrightarrow\Gamma_{\mathrm{gram}},
 \qquad
 \widehat\Gamma_X=\Gamma_X^{\mathrm{form}}\oplus_B\Gamma_{\mathrm{gram}}.
\]
After this normal ordering, \(\Gamma_{\mathrm{gram}}\) embeds into
\(\La^{2,1}_{II}\) and then into the type-IV ambient lattice
\(\La^{3,2}\) of Chapter~\ref{ch:singular-theta-lift}.  Detailed lattice
computations are in Appendix~\ref{appendix:mukai-gram-dictionary}; the
dictionary fixes the \(K3\times E\) charge convention.

The geometric and lattice data feed the three established
\(\Delta_5\) constructions.  The compact target is fixed at
$X = S \times E$ with $S$ a smooth complex
projective K3 surface and $E$ a smooth elliptic curve over $\C$; this is
a trivial-canonical threefold ($\dim_\C X = 3$, $K_X \cong \mathcal{O}_X$;
elementary from the product formula), so $n = \chi(\mathcal O_Y)$
appears directly.  The $(\kappa_{\mathrm{cat}},
\kappa_{\mathrm{ch}}^{\mathrm{Hodge}}, \kappa_{\mathrm{ch}}^{\mathrm{Heis}},
\kappa_{\mathrm{BKM}}, \kappa_{\mathrm{fiber}}) = (0, 0, 3, 5, 24)$ tuple
of \cite{CYQGVolIII} (the K3$\times$E row of the
Gritsenko--Cl\'ery catalogue) is the categorical-dimensional anchor;
every use of $\kappa$ carries one of these five subscripts.

The datum has a geometric part and an arithmetic part.  The geometric
part is the formal dyonic charge lattice
$\Gamma_X^{\mathrm{form}} = \Lambda_S \oplus \Lambda_S$ on the Mukai
sector, together with
$\Gamma_{X,\sigma}^{\mathrm{eff}} \subset \Gamma_X^{\mathrm{alg}}
\subset \Gamma_X^{\mathrm{form}}$.  The arithmetic part is the
quadratic Gram map
$\Pi_X : \Gamma_X^{\mathrm{form}} \to \Gamma_{\mathrm{gram}} = \Z^3$,
its symmetric bilinear polarization $B = -\delta\Pi_X$, the
normal-ordered extension
$\widehat\Gamma_X = \Gamma_X^{\mathrm{form}} \oplus_B
\Gamma_{\mathrm{gram}}$, and the additive map
$\overline\Pi_X(c, T) = \Pi_X(c) - T$.  The fixed-lift
raw-pushforward obstruction theorem forces the extension, and the
projection-locality obstruction forces the wrapped stratum in
$\operatorname{Ran}^{\mathrm{hyb}}(E) =
\operatorname{Ran}(E)_{b=0}\sqcup
\operatorname{Ran}^{\mathrm{wr}}(E)_{b>0}$.

\section{Smooth K3$\times$E threefold and Mukai lattice}
\label{sec:k3xe-mukai-lattice}

\subsection{The compact target}
\label{subsec:k3xe-compact-target}

\begin{definition}[K3$\times$E threefold; primitive geometric datum]
\label{def:k3xe-threefold}
A \emph{K3$\times$E threefold} is a triple $(X, S, E)$ in which $S$ is a
smooth complex projective K3 surface (so $K_S \cong \mathcal{O}_S$,
$h^{1,0}(S) = 0$, $h^{2,0}(S) = 1$; \cite{Mukai1984}),
$E$ is a smooth elliptic curve over $\C$ with marked origin
$0_E \in E$, and
\[
X := S \times E.
\]
The canonical bundle is trivial, $K_X \cong p_S^*K_S \otimes p_E^*K_E
\cong \mathcal O_X$, and Hodge invariants are
\[
\dim_\C X = 3,\qquad
h^{1,0}(X)=h^{1,0}(E)=1,\qquad
\chi_{\mathrm{top}}(X) = \chi_{\mathrm{top}}(S)\cdot\chi_{\mathrm{top}}(E)
= 24 \cdot 0 = 0,\qquad
\chi(\mathcal O_X) = \chi(\mathcal O_S)\cdot\chi(\mathcal O_E) = 2\cdot 0 = 0.
\]
The threefold $X$ is Calabi--Yau in the trivial-canonical sense used by
Donaldson--Thomas theory and shifted symplectic geometry; the
construction uses the chosen trivialization of \(K_X\).
\end{definition}

\begin{remark}[Distinguishing $\chi_{\mathrm{top}}(K3) = 24$ and
$\chi(\mathcal O_S) = 2$]
\label{rem:two-K3-Eulers}
The two integers enter the construction by different routes.  The
Borcherds weight is $f(0,0)/2 = 5$, where $f(0,0) = 10$ is the
constant Fourier coefficient of the weight-zero index-one weak Jacobi
form $\phi_{0,1} = r^{-1} + 10 + r + O(q)$ in the Eichler--Zagier
normalization (computed, \cite{EZ:1}).  The topological Euler
number $\chi_{\mathrm{top}}(K3) = 24$ enters one step removed, through
the elliptic-genus specialization: at $z = 0$ the $\phi_{0,1}$ support
$4n - l^2 \ge -1$ leaves only $l \in \{-1,0,1\}$ at $n = 0$, so the
constant term is $\sum_l f(0,l) = f(0,-1) + f(0,0) + f(0,1) = 12$, and
$Z_{K3}(\tau,0) = 2\,\phi_{0,1}(\tau,0)$ has constant term
$2 \cdot 12 = 24 = \chi_{\mathrm{top}}(K3)$.  Thus $\chi_{\mathrm{top}}(K3)
= 2\sum_l f(0,l)$, not $f(0,0)$; the Borcherds weight $5$ sees only the
single coefficient $f(0,0) = 10$.  The integer
$\chi(\mathcal O_S) = 2$ enters by a third route, the categorical
Calabi--Yau dimension and the holomorphic part of the Mukai pairing.
The two integers belong to different entries: the $\kappa$-tuple of \cite{CYQGVolIII}
records $\kappa_{\mathrm{fiber}} = 24 = \chi_{\mathrm{top}}(K3)$, not
$2 = \chi(\mathcal O_S)$, and the universal Borcherds identity
$\kappa_{\mathrm{BKM}}(\Phi_N) = c_N(0)/2$ at $\Phi_1 = \Delta_5$ gives
$c_1(0) = 10$ and $\kappa_{\mathrm{BKM}} = 5$.  The expression
$\kappa_{\mathrm{ch}}^{\mathrm{Hodge}}+\chi(\mathcal O_{K3})=0+2=2$
does not give the Borcherds weight, while
$\kappa_{\mathrm{ch}}^{\mathrm{Heis}}+\chi(\mathcal O_{K3})=3+2=5$
is only an accidental equality at the \(N=1\) row.  Across the CHL
frame the Borcherds weights are \(c_N(0)/2=(5,3,2,\tfrac32,1)\)
\textup{(Jatkar--Sen \cite{JS06Dyon}; Govindarajan--Krishna
\cite{GK09CHL,GK10Composite}; the \(N=4\) weight is
half-integral)}, not a fixed
sum of a chiral rank and \(\chi(\mathcal O_{K3})\)
(proved, \cite{CYQGVolIII} counterexample).  The coordinate $\kappa_{\mathrm{BKM}} = 5$ is the
cusp-form weight; the protected scalar shadow is written at level
$n = \chi(\mathcal O_Y)$ for the one-dimensional subscheme
$Y\subset K3\times E$.  The target $X = K3\times E$ is a
trivial-canonical threefold.
\end{remark}

\subsection{Mukai lattice on the K3 factor}
\label{subsec:k3-mukai-lattice}

\begin{definition}[Mukai lattice and Mukai pairing]
\label{def:mukai-lattice}
Let $S$ be a smooth complex K3 surface.  The \emph{Mukai lattice} is the
total even cohomology
\[
\Lambda_S
:= \widetilde H(S, \Z)
= H^0(S, \Z) \oplus H^2(S, \Z) \oplus H^4(S, \Z),
\]
identified with the K-theoretic charge lattice
$K_0(D^b(\mathrm{Coh}\,S))_{\mathrm{tor.\,free}}$ via the Mukai vector
\[
v(\mathcal F)
:= \mathrm{ch}(\mathcal F)\sqrt{\mathrm{Td}(S)}
= \bigl(r(\mathcal F),\,c_1(\mathcal F),\,r(\mathcal F) + \mathrm{ch}_2(\mathcal F)\bigr),
\]
where $r$ is the rank, $c_1$ is the first Chern class, and the half-shift
$r$ in the four-form term comes from $\mathrm{Td}(S) = 1 + 2[\mathrm{pt}]$
(proved, \cite{Mukai1984,Mukai1987}).  The \emph{Mukai pairing}
is
\[
\bigl\langle (r, c, s),\,(r', c', s')\bigr\rangle_{\mathrm{Muk}}
:= c\cdot c' - r s' - r' s,
\]
where $c\cdot c' \in \Z$ is the cup product on $H^2(S, \Z)$.  It is even,
unimodular, of signature $(4, 20)$, and rank $24$:
\[
\Lambda_S \cong II_{4, 20}
\cong U^{\oplus 4} \oplus E_8(-1)^{\oplus 2},
\]
where $U$ denotes the rank-two hyperbolic plane $U \cong \mathbb Z e
\oplus \mathbb Z f$ with $e^2 = f^2 = 0$, $e\cdot f = 1$, and $E_8(-1)$
denotes the negative-definite $E_8$ root lattice (proved,
\cite{Mukai1987,NIKULIN}).
\end{definition}

\begin{remark}[Sign convention]
\label{rem:mukai-sign-convention}
The signature $(4, 20)$ uses the convention that $H^0\oplus H^4$
contributes a hyperbolic plane $U$ via $(1, 0, 0)$ and $(0, 0, 1)$ and
that $H^2(S, \Z) \cong II_{3, 19} \cong U^{\oplus 3} \oplus
E_8(-1)^{\oplus 2}$, which has positive part of dimension $3$.  Together,
$\Lambda_S = U \oplus II_{3, 19} \cong U^{\oplus 4} \oplus
E_8(-1)^{\oplus 2}$ with positive part of dimension $4$ (proved,
\cite{Mukai1987}).  The signature is $(4, 20)$, not $(20, 4)$: the
positive direction is the smaller side.  A
reflection $r_v$ across $v\in\Lambda_S$ acts by $r_v(x) = x - 2(x, v)v/(v,
v)$, taking the $\R$-linear extension when $(v, v)\ne 0$, without
implicit sign reversal.
\end{remark}

\subsection{K\"unneth and the formal even Mukai sector on $X = S \times E$}
\label{subsec:kunneth-formal-mukai}

\begin{definition}[Formal even Mukai sector on $X = S \times E$]
\label{def:formal-even-mukai-sector}
Pick a basepoint $0_E \in E$ and the standard generators
$1_E \in H^0(E, \Z)$ and $\omega_E \in H^2(E, \Z)$ with $\int_E \omega_E
= 1$.  Every even cohomology class on $X = S \times E$ decomposes as
\[
v_X = P \otimes 1_E + Q \otimes \omega_E,\qquad
P, Q \in \Lambda_S
\]
(proved, K\"unneth on simply connected $\widetilde H$ tensor
$E$-cohomology).  The pair $(Q, P)\in\Lambda_S\oplus\Lambda_S$ is the
\emph{formal even charge}.  The component $P\otimes 1_E$ is invariant
under translation along $E$, while $Q\otimes\omega_E$ records the
$E$-degree.

The \emph{formal full-Mukai dyonic lattice} is the additive abelian
group
\[
\Gamma_X^{\mathrm{form}} := \Lambda_S \oplus \Lambda_S,\qquad
c = (Q, P),
\]
with electric component $Q$ and magnetic component $P$.  The symbol
\(\Gamma_X^{\mathrm{phys}}\) is defined here by the equality
\[
\Gamma_X^{\mathrm{phys}} := \Gamma_X^{\mathrm{form}}.
\]
This is not the full four-dimensional $\mathcal N = 4$ electric--magnetic
charge lattice and not compact Hall support.  It is the D6--D2--D0/OP
even branch on which the Mukai Gram map is tested before imposing
algebraicity, effectivity, stability, or non-emptiness of compact Hall
moduli.
\end{definition}

\begin{definition}[Liu product-stability datum on \(S\times E\)]
\label{def:liu-product-stability-datum}
Fix a rational Bridgeland stability condition
\[
\sigma_S=(\mathcal A_S,Z_S),\qquad
Z_S\bigl(K(D^b(\mathrm{Coh}\,S))\bigr)\subset\Q+i\Q ,
\]
on \(D^b(\mathrm{Coh}\,S)\).  Let \(L_E\) be a degree-one line bundle
on the elliptic curve \(E\), and let
\[
p:S\times E\to S,\qquad q:S\times E\to E
\]
be the projections.  The Abramovich--Polishchuk heart used by Liu is
the following heart
\cite{AbramovichPolishchukTStructures,LiuProductStability}:
\[
\mathcal A_E^{\mathrm{AP}}
=\bigl\{M\in D^b(\mathrm{Coh}(S\times E))\mid
p_*(M\otimes q^*L_E^{\otimes n})\in\mathcal A_S
\text{ for }n\gg0\bigr\}.
\]
For \(M\in\mathcal A_E^{\mathrm{AP}}\), Liu's curve-factor
polynomial is
\[
L_M(n):=
Z_S\!\bigl(p_*(M\otimes q^*L_E^{\otimes n})\bigr)
=a(M)n+b(M)+i\bigl(c(M)n+d(M)\bigr),
\]
where \(a,b,c,d:K(\mathcal A_E^{\mathrm{AP}})\to\Q\) are additive
homomorphisms.  For \(t\in\Q_{>0}\), set
\[
Z_t(M)=a(M)t-d(M)+i\,c(M)t,\qquad
\nu_t(M)=
\begin{cases}
\dfrac{-a(M)t+d(M)}{c(M)t},&c(M)\neq0,\\[6pt]
+\infty,&c(M)=0.
\end{cases}
\]
Let
\[
\mathcal T_t=\{M\in\mathcal A_E^{\mathrm{AP}}\mid
\nu_{t,\min}(M)>0\},\qquad
\mathcal F_t=\{M\in\mathcal A_E^{\mathrm{AP}}\mid
\nu_{t,\max}(M)\le0\},
\]
and
\[
\mathcal A_E^t=\langle\mathcal T_t,\mathcal F_t[1]\rangle .
\]
For \(s\in\Q_{>0}\), Liu's product central charge is
\[
Z_E^{s,t}(M)=c(M)s+b(M)+i\bigl(-a(M)t+d(M)\bigr).
\]
The Liu product-stability condition is
\[
\sigma_{s,t}:=(\mathcal A_E^t,Z_E^{s,t}).
\]
Liu proves that this is a rational Bridgeland stability condition for
\(s,t\in\Q_{>0}\), and extends the construction continuously to
\((s,t)\in\R_{>0}^2\)
\cite[Theorem~4.7 and \S5]{LiuProductStability}.  This is a
source-imported stability condition.  It is not a finite-type moduli
construction, not a compact Hall stage, and not a Pfaffian
orientation.
\end{definition}

The packet \(\texttt{certificates/moduli/liu\_product\_stability}\)
checks this definition.  It records the source theorem, the AP heart
membership formula, the linear polynomial \(L_M(n)\), the four
coefficient homomorphisms \(a,b,c,d\), the weak charge \(Z_t\), the
torsion-pair tilt, and the central charge \(Z_E^{s,t}\).  Its verifier
status is \texttt{LIU\_PRODUCT\_STABILITY\_DEFINED}.  It does not
prove the separate AP-compatibility row, finite-type semistable
substacks, derived enhancements, compact Hall stages, Pfaffian
orientations, or protected traces.

\begin{proposition}[Compatibility with the Abramovich--Polishchuk heart]
\label{prop:ap-liu-heart-compatibility}
\textnormal{[proved]}
With the datum of
Definition~\ref{def:liu-product-stability-datum}, the heart of Liu's
product stability condition is the torsion-pair tilt of the
Abramovich--Polishchuk heart:
\[
\heartsuit(\sigma_{s,t})
=\mathcal A_E^t
=\langle\mathcal T_t,\mathcal F_t[1]\rangle,
\]
where \((\mathcal T_t,\mathcal F_t)\) is the torsion pair in
\(\mathcal A_E^{\mathrm{AP}}\) defined by the slope \(\nu_t\).
Consequently
\[
\mathcal A_X^{\mathrm{AP/Liu}}=\mathcal A_E^t.
\]
Thus the retained heart is Abramovich--Polishchuk-compatible in the
finite-row sense: it is obtained from \(\mathcal A_E^{\mathrm{AP}}\)
by Liu's specified torsion-pair tilt.
\end{proposition}

\begin{proof}
The Abramovich--Polishchuk heart
\(\mathcal A_E^{\mathrm{AP}}\) is the global heart entering Liu's
construction.  Liu defines the weak charge \(Z_t\) on this heart,
uses the resulting slope \(\nu_t\) to form the torsion pair
\[
\mathcal T_t=\{M\mid\nu_{t,\min}(M)>0\},\qquad
\mathcal F_t=\{M\mid\nu_{t,\max}(M)\le0\},
\]
and then sets
\[
\mathcal A_E^t=\langle\mathcal T_t,\mathcal F_t[1]\rangle
\]
before defining
\[
\sigma_{s,t}=(\mathcal A_E^t,Z_E^{s,t})
\]
\cite[\S2.2 and Theorem~4.7]{LiuProductStability}.  Hence the heart
of \(\sigma_{s,t}\) is \(\mathcal A_E^t\) by construction.  The
retained heart is defined as
\(\mathcal A_X^{\mathrm{AP/Liu}}:=\heartsuit(\sigma_{s,t})\), so it
equals \(\mathcal A_E^t\).  This proves the compatibility statement
using only the definition of the tilt; it proves no noetherianity,
HN-filtration, boundedness, finite-moduli, derived-enhancement,
Pfaffian-orientation, or protected-trace row.
\end{proof}

The packet
\(\texttt{certificates/moduli/ap\_liu\_compatibility}\) checks this
heart identification.  It records the AP membership formula, the
weak charge \(Z_t\), the slope \(\nu_t\), the torsion pair
\((\mathcal T_t,\mathcal F_t)\), the tilted heart
\(\mathcal A_E^t\), and the equality
\(\mathcal A_X^{\mathrm{AP/Liu}}=\heartsuit(\sigma_{s,t})=\mathcal A_E^t\).
Its verifier status is
\texttt{AP\_LIU\_COMPATIBILITY\_VERIFIED}.  It does not prove
noetherianity, HN filtrations, bounded HN types, finite-type
semistable substacks, derived enhancements, compact Hall stages,
Pfaffian orientations, or protected traces.

\begin{definition}[Retained Abramovich--Polishchuk/Liu heart]
\label{def:retained-ap-liu-heart}
Fix the Liu product-stability datum of
Definition~\ref{def:liu-product-stability-datum}.  The retained
stability heart on
\[
D^b(\mathrm{Coh}\,X),\qquad X=S\times E,
\]
is
\[
\mathcal A_X^{\mathrm{AP/Liu}}
:=\heartsuit(\sigma_{s,t}),
\]
equivalently
\(\mathcal A_X^{\mathrm{AP/Liu}}=\mathcal A_E^t\) by
Proposition~\ref{prop:ap-liu-heart-compatibility}.  In
the rest of the manuscript the symbol \(\sigma\) means this retained
datum \(\sigma_{s,t}\), and its heart is
\(\mathcal A_X^{\mathrm{AP/Liu}}\).

This definition is only the choice of the ambient heart and stability
datum.  Noetherianity of the retained windows, existence of
Harder--Narasimhan filtrations, boundedness of HN types, finite-type
semistable substacks, quasi-smooth derived enhancements, universal
complexes, and compact Hall stages are separate finite rows.
\end{definition}

The packet \(\texttt{certificates/moduli/stability\_heart}\) checks
this definition as a table-level datum.  It records the inputs
\[
(S,H),\quad (E,0_E),\quad \sigma_S,\quad L_E,\quad
\mathcal A_X^{\mathrm{AP}},\quad \sigma_{s,t},\quad
\mathcal R_{\mathrm{HN}},
\]
and verifies the scalar firewall excluding noetherianity, HN
filtrations, bounded HN types, finite-type moduli, Pfaffian
orientations, and protected traces.  Its verifier status is
\texttt{STABILITY\_HEART\_DATUM\_VERIFIED}.

\begin{definition}[Retained noetherian window]
\label{def:retained-noetherian-window}
Fix the heart \(\mathcal A_X^{\mathrm{AP/Liu}}\).  A retained finite
window at height \(R\) is a full exact subcategory
\[
\mathcal A_{X,R}^{\mathrm{ret}}\subset\mathcal A_X^{\mathrm{AP/Liu}}
\]
with a finite set of isomorphism classes, closed under retained
subobjects and retained quotients, together with a length function
\[
\ell_R:\operatorname{Obj}(\mathcal A_{X,R}^{\mathrm{ret}})\to\Z_{\ge0}
\]
such that \(\ell_R(A)=0\) exactly for the zero object and
\(\ell_R(A)<\ell_R(B)\) for every strict retained monomorphism
\(A\hookrightarrow B\).
\end{definition}

\begin{proposition}[Retained finite windows are noetherian]
\label{prop:retained-window-noetherian}
\textnormal{[proved]}
Every retained finite window
\(\mathcal A_{X,R}^{\mathrm{ret}}\) in the sense of
Definition~\ref{def:retained-noetherian-window} is noetherian: every
ascending chain of retained subobjects terminates.
\end{proposition}

\begin{proof}
Let
\[
A_0\subset A_1\subset A_2\subset\cdots
\]
be a chain of retained subobjects of an object \(A\) in
\(\mathcal A_{X,R}^{\mathrm{ret}}\).  After deleting repetitions, each
inclusion is strict.  The length function then gives a strictly
increasing sequence
\[
\ell_R(A_0)<\ell_R(A_1)<\ell_R(A_2)<\cdots .
\]
All \(A_i\) lie in the finite retained window and all are subobjects of
\(A\), so \(\ell_R(A_i)\le \ell_R(A)\).  A strictly increasing sequence
of non-negative integers bounded above by \(\ell_R(A)\) is finite.
Thus the original ascending chain stabilizes.
\end{proof}

The packet
\(\texttt{certificates/moduli/retained\_window\_noetherian}\) checks
this finite ACC statement on a displayed retained exact window.  It
verifies a finite object set, subobject and quotient closure, strict
increase of the length rank along each strict monomorphism, acyclicity
of the strict subobject graph, and the longest-chain bound.  Its
verifier status is
\texttt{RETAINED\_WINDOW\_NOETHERIAN\_VERIFIED}.  It does not prove
global noetherianity of \(\mathcal A_X^{\mathrm{AP/Liu}}\),
Harder--Narasimhan filtrations, bounded HN types, finite-type
semistable substacks, quasi-smooth moduli, Pfaffian orientations, or
protected traces.

\begin{definition}[Retained finite Harder--Narasimhan filtration datum]
\label{def:retained-finite-hn-filtration-datum}
Let \(\mathcal A_{X,R}^{\mathrm{ret}}\) be a retained finite window.
A retained finite Harder--Narasimhan filtration datum consists of:
\[
\Phi_R\subset\Q,\qquad
\phi_R:\mathcal S_R\to\Phi_R,
\]
where \(\Phi_R\) is a finite totally ordered phase set and
\(\mathcal S_R\) is a class of nonzero retained semistable objects,
together with, for every \(A\in\mathcal A_{X,R}^{\mathrm{ret}}\), a
finite exact filtration
\[
0=A_0\subset A_1\subset\cdots\subset A_m=A
\]
whose retained quotients \(E_i=A_i/A_{i-1}\) lie in \(\mathcal S_R\)
and satisfy
\[
\phi_R(E_1)>\phi_R(E_2)>\cdots>\phi_R(E_m).
\]
The zero object has the empty filtration.  This is finite
retained-window data.  It is not a global Harder--Narasimhan theorem
for \(\mathcal A_X^{\mathrm{AP/Liu}}\), not boundedness of HN types,
and not finite-type moduli.
\end{definition}

\begin{proposition}[Retained finite HN filtrations exist from the datum]
\label{prop:retained-window-hn-filtration}
\textnormal{[proved]}
If a retained finite window
\(\mathcal A_{X,R}^{\mathrm{ret}}\) carries the datum of
Definition~\ref{def:retained-finite-hn-filtration-datum}, then every
object of \(\mathcal A_{X,R}^{\mathrm{ret}}\) has a retained
Harder--Narasimhan filtration.
\end{proposition}

\begin{proof}
For \(A=0\) the empty filtration is the Harder--Narasimhan filtration.
For \(A\neq0\), the datum gives an exact retained chain
\[
0=A_0\subset A_1\subset\cdots\subset A_m=A .
\]
Each quotient \(E_i=A_i/A_{i-1}\) is retained and semistable, and the
phase inequalities are strict.  The chain is finite, hence terminates
inside the exact window.  These are precisely the
Harder--Narasimhan conditions in the retained finite window.
\end{proof}

The packet
\(\texttt{certificates/moduli/retained\_window\_hn\_filtration}\)
checks this finite statement on a displayed retained exact window.  It
verifies a finite phase set, semistable quotient factors, exact
filtration chains, length additivity, and strict phase descent.  Its
verifier status is
\texttt{RETAINED\_WINDOW\_HN\_FILTRATION\_VERIFIED}.  It does not
prove global Harder--Narasimhan filtrations in
\(\mathcal A_X^{\mathrm{AP/Liu}}\), bounded HN types, finite-type
semistable substacks, quasi-smooth moduli, Pfaffian orientations, or
protected traces.

\begin{definition}[Retained bounded HN type datum]
\label{def:retained-bounded-hn-type-datum}
Let \(\mathcal A_{X,R}^{\mathrm{ret}}\) be a retained finite window
with a retained finite Harder--Narasimhan filtration datum.  The
\emph{retained HN type} of an object \(A\) is the finite word
\[
\operatorname{HNType}_R(A)
=\bigl(([E_1],\phi_R(E_1)),\ldots,([E_m],\phi_R(E_m))\bigr)
\]
attached to its retained HN quotients \(E_i=A_i/A_{i-1}\), with the
empty word for \(A=0\).  A retained bounded HN type datum is a finite
set \(\mathfrak T_R\) containing all these words, together with a
finite Liu--Hilbert decoration of the semistable factor classes:
\[
[E]\longmapsto
\bigl(\widehat c_E,[a_E,b_E],P_E,N_E,\mathcal F_E\bigr).
\]
Here \(\widehat c_E\) is the retained numerical class, \([a_E,b_E]\)
is the cohomological-amplitude interval, \(P_E\) is the retained
Hilbert-polynomial label, \(N_E\) is the retained regularity bound, and
\(\mathcal F_E\) is the retained closure family.  This datum is a
finite bound on active-window HN types.  It is not a finite-type
semistable substack, not a quasi-smooth derived enhancement, and not a
compact Hall stage.
\end{definition}

\begin{proposition}[HN type is bounded in a retained active window]
\label{prop:retained-hn-type-bounded}
\textnormal{[proved]}
Let \(\mathcal A_{X,R}^{\mathrm{ret}}\) be a retained finite window
with a retained finite Harder--Narasimhan filtration datum.  Then the
set of retained HN types occurring in
\(\mathcal A_{X,R}^{\mathrm{ret}}\) is finite.  More precisely, if
\[
N_R=\max\{\ell_R(A)\mid A\in
\mathcal A_{X,R}^{\mathrm{ret}}\},
\]
then every retained HN type has length at most \(N_R\).
\end{proposition}

\begin{proof}
The retained window has finitely many isomorphism classes, hence the
set of retained semistable HN factors is finite.  If
\[
0=A_0\subset A_1\subset\cdots\subset A_m=A
\]
is the retained HN filtration of \(A\), then each inclusion
\(A_{i-1}\subset A_i\) is strict unless \(m=0\).  The retained length
rank gives
\[
0=\ell_R(A_0)<\ell_R(A_1)<\cdots<\ell_R(A_m)\le N_R,
\]
so \(m\le N_R\).  Thus every retained HN type is a word of length at
most \(N_R\) in a finite alphabet of semistable factor classes.  The
set of such words is finite, and the strict phase condition selects a
finite subset.  A finite Liu--Hilbert decoration of the factor classes
therefore gives a finite decorated HN type set.
\end{proof}

The packet \(\texttt{certificates/moduli/retained\_hn\_type\_bounds}\)
checks this finite boundedness statement.  It verifies the finite
semistable factor-type set, the finite HN type-word set, the maximal
HN height, the maximal factor count, the retained length-rank bound,
and the mirror of the HN type rows in
\(\texttt{certificates/moduli/k3e\_finite\_moduli/hn\_type\_bounds.csv}\).
Its verifier status is
\texttt{RETAINED\_HN\_TYPE\_BOUNDS\_VERIFIED}.  It does not construct
finite-type semistable substacks, quasi-smooth derived enhancements,
universal complexes, compact Hall stages, Pfaffian orientations, or
protected traces.

\begin{definition}[Retained finite-type semistable substack datum]
\label{def:retained-finite-type-semistable-substack-datum}
Fix a retained finite window
\(\mathcal A_{X,R}^{\mathrm{ret}}\), the retained stability condition
\(\sigma_{s,t}\), and the retained bounded HN type datum of
Definition~\ref{def:retained-bounded-hn-type-datum}.  A retained
finite-type semistable substack datum assigns to every retained
semistable factor type \(c\) a finite-type bounded
Quot/Postnikov ambient stack \(Q_{R,c}\), a chart
\[
U_{R,c}\to Q_{R,c},
\]
and a specialization-closed finite-type algebraic substack
\[
\mathfrak M^{ss}_{R,c}\subset Q_{R,c}
\]
whose objects are exactly the \(\sigma_{s,t}\)-semistable retained
objects of type \(c\).  The datum carries two defect ranks:
\[
d^{ss}_{R,c}=0,\qquad d^{bd}_{R,c}=0.
\]
The first says that no unstable object is included and no semistable
object of type \(c\) is omitted.  The second says that the bounded
Quot/Postnikov ambient covers the retained class.  This is a
finite-type semistable-substack datum.  It is not a quasi-smooth
derived enhancement, not a universal complex, not a rigidification,
not a compact Hall correspondence, and not an orientation datum.
\end{definition}

\begin{proposition}[Finite-type semistable substacks for retained classes]
\label{prop:retained-finite-type-semistable-substacks}
\textnormal{[proved]}
Let \(\mathcal A_{X,R}^{\mathrm{ret}}\) carry a retained bounded HN
type datum and a retained finite-type semistable substack datum.  Then
for every retained semistable factor type \(c\) the stack
\(\mathfrak M^{ss}_{R,c}\) is of finite type over \(\C\), is
specialization-closed inside its bounded Quot/Postnikov ambient, and
parametrizes exactly the retained \(\sigma_{s,t}\)-semistable objects
of type \(c\).
\end{proposition}

\begin{proof}
The retained HN type datum gives only finitely many semistable factor
types \(c\).  For each such \(c\), the semistable-substack datum places
\(\mathfrak M^{ss}_{R,c}\) inside a finite-type bounded
Quot/Postnikov ambient \(Q_{R,c}\) through the chart \(U_{R,c}\), with
the semistable locus represented by finitely many locally closed
retained chart pieces.  Finite unions of locally closed substacks of a
finite-type algebraic stack are finite type.  The condition
\(d^{bd}_{R,c}=0\) says that the ambient covers the retained class,
while \(d^{ss}_{R,c}=0\) says that the substack is precisely the
\(\sigma_{s,t}\)-semistable locus of that class.  The
specialization-closed flag is part of the datum and is checked in the
same finite chart.  Thus each \(\mathfrak M^{ss}_{R,c}\) is finite
type and has the asserted moduli meaning.  The argument uses no
derived enhancement, universal complex, rigidification, finite-inertia
stratification, flag correspondence, cosection atlas, transition map,
Pfaffian orientation, or protected trace.
\end{proof}

The packet
\(\texttt{certificates/moduli/retained\_semistable\_substacks}\)
checks this finite-type row.  It imports the retained HN factor types,
supplies one bounded Quot/Postnikov ambient chart and one semistable
substack for each retained class, verifies specialization-closedness,
and checks zero semistability and boundedness defect ranks.  Its
verifier status is
\texttt{RETAINED\_SEMISTABLE\_SUBSTACKS\_VERIFIED}.  The mirrored
rows in
\(\texttt{certificates/moduli/k3e\_finite\_moduli/semistable\_substacks.csv}\)
discharge item \(167\).  They do not construct derived enhancements,
universal complexes, scalar rigidifications, finite inertia strata,
extension or flag stacks, cosection atlases, transitions, compact Hall
stages, Pfaffian orientations, or protected traces.

\begin{definition}[Retained quasi-smooth derived enhancement datum]
\label{def:retained-quasi-smooth-derived-enhancement-datum}
Fix the retained finite-type semistable substacks
\(\mathfrak M^{ss}_{R,c}\) of
Definition~\ref{def:retained-finite-type-semistable-substack-datum}.
A retained quasi-smooth derived enhancement datum assigns to every
retained semistable factor type \(c\) a derived Artin stack
\[
\mathfrak M^{\mathrm{der}}_{R,c}
\]
with a morphism
\(\mathfrak M^{\mathrm{der}}_{R,c}\to\operatorname{Perf}(X)\), an
identification
\[
t_0\mathfrak M^{\mathrm{der}}_{R,c}
=\mathfrak M^{ss}_{R,c},
\]
a perfect cotangent complex of amplitude
\[
\mathbb L_{\mathfrak M^{\mathrm{der}}_{R,c}}\in
D^{[-1,0]}(\mathfrak M^{\mathrm{der}}_{R,c}),
\]
and a restricted PTVV \((-1)\)-shifted symplectic form
\[
\omega^{\mathrm{PTVV}}_{R,c}.
\]
The datum carries two defect ranks
\[
d^{\mathrm{Tor}}_{R,c}=0,\qquad d^\omega_{R,c}=0.
\]
The first says that the cotangent-amplitude assertion has no retained
Tor-amplitude defect.  The second says that the PTVV form restricts to
the retained derived substack without symplectic defect.  This is a
derived-enhancement datum.  It is not a universal perfect complex, not
a scalar rigidification, not an inertia stratification, not a
Hall-flag correspondence, and not a cosection or orientation datum.
\end{definition}

\begin{proposition}[Quasi-smooth derived Artin enhancements]
\label{prop:retained-quasi-smooth-derived-enhancements}
\textnormal{[proved]}
Let the retained finite-type semistable substacks carry the datum of
Definition~\ref{def:retained-quasi-smooth-derived-enhancement-datum}.
Then every \(\mathfrak M^{ss}_{R,c}\) has a quasi-smooth derived Artin
enhancement \(\mathfrak M^{\mathrm{der}}_{R,c}\) whose classical
truncation is \(\mathfrak M^{ss}_{R,c}\), whose cotangent complex has
amplitude \([-1,0]\), and whose \((-1)\)-shifted symplectic form is
the restricted PTVV form on the Calabi--Yau threefold
\(X=S\times E\).
\end{proposition}

\begin{proof}
The datum gives a derived Artin stack
\(\mathfrak M^{\mathrm{der}}_{R,c}\) with
\(t_0\mathfrak M^{\mathrm{der}}_{R,c}=\mathfrak M^{ss}_{R,c}\).
The finite-type statement for the classical truncation is
Proposition~\ref{prop:retained-finite-type-semistable-substacks}.  The
cotangent-amplitude condition
\(\mathbb L_{\mathfrak M^{\mathrm{der}}_{R,c}}\in D^{[-1,0]}\) and the
vanishing \(d^{\mathrm{Tor}}_{R,c}=0\) are exactly the
quasi-smoothness criterion for a derived Artin stack.  Since
\(K_X\simeq\mathcal O_X\), \(X\) is a Calabi--Yau threefold, and PTVV
produce the ambient \((-1)\)-shifted symplectic form on the derived
moduli of perfect complexes \cite[Theorem~0.3]{PTVV2013}.  The
vanishing \(d^\omega_{R,c}=0\) says that this form restricts to the
retained derived substack.  Thus the retained enhancement is
quasi-smooth and carries the asserted PTVV form.  The proof supplies no
universal complex on the rigidified stack, no scalar rigidification, no
finite-inertia stratification, no extension or flag stack, no cosection
atlas, no transition system, no Pfaffian orientation, and no protected
trace.
\end{proof}

The packet
\(\texttt{certificates/moduli/retained\_derived\_enhancements}\)
checks this row.  It imports the finite-type semistable substacks,
records one derived Artin enhancement per retained substack, verifies
cotangent amplitude \([-1,0]\), imports the PTVV \((-1)\)-shifted
symplectic form, and checks zero Tor-amplitude and symplectic defect
ranks.  Its verifier status is
\texttt{RETAINED\_DERIVED\_ENHANCEMENTS\_VERIFIED}.  The mirrored rows
in
\(\texttt{certificates/moduli/k3e\_finite\_moduli/derived\_enhancements.csv}\)
discharge item \(168\).  They do not construct universal complexes,
scalar rigidifications, finite inertia strata, extension or flag
stacks, cosection atlases, transitions, compact Hall stages, Pfaffian
orientations, or protected traces.

\begin{definition}[Retained scalar rigidification and residual-inertia datum]
\label{def:retained-rigidification-inertia-datum}
Fix the retained quasi-smooth derived enhancements
\(\mathfrak M^{\mathrm{der}}_{R,c}\) of
Definition~\ref{def:retained-quasi-smooth-derived-enhancement-datum}.
A retained scalar rigidification and residual-inertia datum assigns to
every retained semistable factor type \(c\) an exact sequence of
stabilizer groups
\[
1\longrightarrow \mathbb G_m\longrightarrow
\operatorname{Aut}_{R,c}\longrightarrow H_{R,c}\longrightarrow 1,
\]
a scalar rigidification
\[
\mathfrak M^{\mathrm{rig}}_{R,c}
=\bigl[\mathfrak M^{\mathrm{der}}_{R,c}/\mathbb G_m\bigr],
\]
and the residual inertia group \(H_{R,c}\).  The datum requires
\[
|H_{R,c}|<\infty,\qquad
d^{\mathrm{rig}}_{R,c}=0,\qquad d^{\mathrm{in}}_{R,c}=0.
\]
Thus the scalar \(\mathbb G_m\) has been removed, the residual inertia
is finite, and the rigidification and inertia defects vanish.  This is
a scalar-rigidification datum for retained derived substacks.  It is
not a universal perfect complex, not a finite closed inertia
stratification, not an extension or flag stack, not a cosection atlas,
not a transition system, and not an orientation datum.
\end{definition}

\begin{proposition}[Finite residual inertia after scalar rigidification]
\label{prop:retained-finite-residual-inertia-after-rigidification}
\textnormal{[proved]}
Let the retained derived enhancements carry the datum of
Definition~\ref{def:retained-rigidification-inertia-datum}.  Then for
each retained semistable factor type \(c\), the scalar rigidified stack
\(\mathfrak M^{\mathrm{rig}}_{R,c}\) has finite residual inertia
\(H_{R,c}\).  Equivalently, after quotienting the scalar
\(\mathbb G_m\)-automorphisms, no positive-dimensional stabilizer
remains in the retained row.
\end{proposition}

\begin{proof}
The datum gives the exact sequence
\[
1\to \mathbb G_m\to\operatorname{Aut}_{R,c}\to H_{R,c}\to 1.
\]
The scalar rigidification is the quotient by the indicated central
\(\mathbb G_m\).  Therefore the stabilizer remaining on
\(\mathfrak M^{\mathrm{rig}}_{R,c}\) is \(H_{R,c}\).  The condition
\(|H_{R,c}|<\infty\) gives finite residual inertia, and the equalities
\(d^{\mathrm{rig}}_{R,c}=d^{\mathrm{in}}_{R,c}=0\) say that neither
the quotient nor the inertia computation has a retained defect.  The
argument constructs no universal complex, finite closed inertia
stratification, flag stack, cosection atlas, transition map, Pfaffian
orientation, or protected trace.
\end{proof}

The packet
\(\texttt{certificates/moduli/retained\_rigidification\_inertia}\)
checks this row.  It imports the retained derived enhancements,
records one scalar rigidification for each retained substack, verifies
the exact sequence
\[
1\to\mathbb G_m\to\operatorname{Aut}_{R,c}\to H_{R,c}\to 1,
\]
checks \(\mathbb G_m\)-removal, finite residual inertia, and zero
rigidification and inertia defect ranks.  Its verifier status is
\texttt{RETAINED\_RIGIDIFICATION\_INERTIA\_VERIFIED}.  The mirrored
rows in
\(\texttt{certificates/moduli/k3e\_finite\_moduli/scalar\_rigidifications.csv}\)
discharge item \(169\) and the scalar-rigidification existence row
item \(175\).  They do not construct universal complexes, finite
closed inertia stratifications, extension or flag stacks, cosection
atlases, transitions, compact Hall stages, Pfaffian orientations, or
protected traces.

\begin{definition}[Retained finite class-bound datum]
\label{def:retained-finite-class-bound-datum}
Fix the retained finite-type semistable substacks
\(\mathfrak M^{ss}_{R,c}\) and the retained HN type-bound rows.  A
retained finite class-bound datum consists of a finite set
\[
C_R=\{c_1,\ldots,c_m\}
\]
of retained semistable classes, together with maps assigning to each
\(c\in C_R\) a retained type \(\tau(c)\), a charge \(\widehat c\), a
substack \(\mathfrak M^{ss}_{R,c}\), a finite Hilbert-polynomial label
set
\[
\{P_{c,i}\}_{i\in I_c},\qquad |I_c|<\infty,
\]
a cohomological-amplitude interval
\[
[a_c,b_c]\subset\mathbb Z,\qquad a_c\le b_c,
\]
and a nonnegative regularity bound \(N_c\).  The datum requires zero
class-coverage, Hilbert-bound, amplitude-bound, and regularity-bound
defects.  This is a retained class-bound datum.  It is not a universal
perfect complex, not a finite closed inertia stratification, not an
extension or flag stack, not a cosection atlas, not a transition
system, and not an orientation datum.
\end{definition}

\begin{proposition}[Finite retained class bounds]
\label{prop:retained-finite-class-bounds}
\textnormal{[proved]}
Let a retained active window carry the datum of
Definition~\ref{def:retained-finite-class-bound-datum}.  Then the
retained class set \(C_R\) is finite.  For every \(c\in C_R\), the
Hilbert-polynomial labels \(P_{c,i}\), the cohomological amplitude
interval \([a_c,b_c]\), and the regularity bound \(N_c\) are finite
retained data.
\end{proposition}

\begin{proof}
The datum writes \(C_R=\{c_1,\ldots,c_m\}\), hence the retained class
set is finite.  For each \(c\), the index set \(I_c\) of Hilbert
polynomial labels is finite by hypothesis, the amplitude interval is
bounded by the two integers \(a_c\le b_c\), and \(N_c\in\mathbb Z_{\ge0}\)
is a finite integer.  The vanishing of the four defect ranks says that
these labels cover exactly the retained classes imported from the HN
type-bound and semistable-substack rows.  The argument constructs no
universal family over \(\mathfrak M^{ss}_{R,c}\), no closed inertia
stratification, no flag stack, no cosection atlas, no transition map,
no Pfaffian orientation, and no protected trace.
\end{proof}

The packet
\(\texttt{certificates/moduli/retained\_class\_bounds}\) checks this
row.  It imports the retained HN type-bound rows and retained
semistable substacks, records the finite class set \(C_R\), records the
Hilbert-polynomial labels \(P_{c,i}\), records the amplitude intervals
\([a_c,b_c]\), and records finite regularity bounds \(N_c\).  Its
verifier status is
\texttt{RETAINED\_CLASS\_BOUNDS\_VERIFIED}.  The mirrored class-bound
data are the Hilbert-polynomial, amplitude, and regularity columns in
\(\texttt{certificates/moduli/k3e\_finite\_moduli/hn\_type\_bounds.csv}\),
and they discharge items \(170\)--\(173\).  They do not construct
universal complexes, finite closed inertia stratifications, extension
or flag stacks, cosection atlases, transitions, compact Hall stages,
Pfaffian orientations, or protected traces.

\begin{definition}[Retained universal perfect-complex datum]
\label{def:retained-universal-perfect-complex-datum}
Fix the retained finite substacks
\(\mathfrak M^{ss}_{R,c}\), their scalar rigidifications, and the
retained class-bound datum.  A retained universal perfect-complex datum
assigns to every \(c\in C_R\) a perfect complex
\[
\mathcal U_{R,c}\in
\operatorname{Perf}\bigl(X\times\mathfrak M^{\mathrm{rig}}_{R,c}\bigr),
\]
a base-change identification for every test morphism
\[
T\longrightarrow \mathfrak M^{\mathrm{rig}}_{R,c},
\qquad
\mathcal U_{R,c}|_{X\times T}\simeq E_T,
\]
and a finite Tor-amplitude interval
\[
\operatorname{TorAmp}(\mathcal U_{R,c})\subset [a_c,b_c].
\]
It carries two defect ranks
\[
d^{\mathrm{desc}}_{R,c}=0,\qquad d^{\mathrm{perf}}_{R,c}=0.
\]
The first says that the tautological complex on the bounded
Quot/Postnikov chart descends across the scalar rigidification.  The
second says that the descended object remains perfect on the retained
finite substack.  This is a universal-complex datum.  It is not a
finite closed inertia stratification, not an extension or flag stack,
not a cosection atlas, not a transition system, and not an orientation
datum.
\end{definition}

\begin{proposition}[Universal perfect complexes on retained finite substacks]
\label{prop:retained-universal-perfect-complexes}
\textnormal{[proved]}
Let the retained finite substacks carry the datum of
Definition~\ref{def:retained-universal-perfect-complex-datum}.  Then
each retained scalar-rigidified finite substack
\(\mathfrak M^{\mathrm{rig}}_{R,c}\) carries a universal perfect
complex \(\mathcal U_{R,c}\) on \(X\times
\mathfrak M^{\mathrm{rig}}_{R,c}\), compatible with base change and of
finite Tor amplitude contained in \([a_c,b_c]\).
\end{proposition}

\begin{proof}
The bounded Quot/Postnikov ambient chart carries its tautological
complex.  Restricting it to the retained finite substack gives the
tautological complex before rigidification.  The equality
\(d^{\mathrm{desc}}_{R,c}=0\) is precisely the descent condition across
the scalar rigidification
\(\mathfrak M^{\mathrm{der}}_{R,c}\to
\mathfrak M^{\mathrm{rig}}_{R,c}\).  The equality
\(d^{\mathrm{perf}}_{R,c}=0\) says that the descended complex is
perfect on \(X\times\mathfrak M^{\mathrm{rig}}_{R,c}\).  The datum
also supplies the base-change identifications and the Tor-amplitude
containment in \([a_c,b_c]\).  Thus the retained substack carries the
asserted universal perfect complex.  The proof constructs no finite
closed inertia stratification, extension or flag stack, cosection
atlas, transition map, Pfaffian orientation, or protected trace.
\end{proof}

The packet
\(\texttt{certificates/moduli/retained\_universal\_complexes}\) checks
this row.  It imports the retained semistable substacks, class-bound
rows, derived enhancements, and scalar rigidifications, records one
universal perfect complex per retained finite substack, verifies
base-change compatibility, finite Tor amplitude, zero descent defect,
and zero perfection defect.  Its verifier status is
\texttt{RETAINED\_UNIVERSAL\_COMPLEXES\_VERIFIED}.  The mirrored rows
in
\(\texttt{certificates/moduli/k3e\_finite\_moduli/universal\_complexes.csv}\)
discharge item \(174\).  They do not construct finite closed inertia
stratifications, extension or flag stacks, cosection atlases,
transitions, compact Hall stages, Pfaffian orientations, or protected
traces.

\begin{definition}[Retained \(E\)-translation rigidification datum]
\label{def:retained-e-translation-rigidification-datum}
Fix the retained scalar-rigidified finite substacks
\(\mathfrak M^{\mathrm{rig}}_{R,c}\) and their universal perfect
complexes.  A retained \(E\)-translation rigidification datum assigns
to every \(c\in C_R\) an action
\[
a_{R,c}:E\times \mathfrak M^{\mathrm{rig}}_{R,c}
\longrightarrow \mathfrak M^{\mathrm{rig}}_{R,c}
\]
with identity element \(0_E\), free on the retained row, together with
the quotient stack
\[
\mathfrak M^{E\text{-}\mathrm{rig}}_{R,c}
=\bigl[\mathfrak M^{\mathrm{rig}}_{R,c}/E\bigr].
\]
The datum carries a defect rank
\[
d^{E\text{-}\mathrm{rig}}_{R,c}=0.
\]
Thus the elliptic-translation direction has been removed without
defect on the retained finite substack.  This is an
\(E\)-translation-rigidification datum.  It is not a finite closed
inertia stratification, not an extension or flag stack, not a
cosection atlas, not a transition system, and not an orientation datum.
\end{definition}

\begin{proposition}[\(E\)-translation rigidifications on retained finite substacks]
\label{prop:retained-e-translation-rigidifications}
\textnormal{[proved]}
Let the retained scalar-rigidified finite substacks carry the datum of
Definition~\ref{def:retained-e-translation-rigidification-datum}.  Then
each retained finite substack has an \(E\)-translation rigidification
\(\mathfrak M^{E\text{-}\mathrm{rig}}_{R,c}\), and the translation
direction has zero retained quotient defect.
\end{proposition}

\begin{proof}
The datum gives an action of the elliptic curve \(E\) on
\(\mathfrak M^{\mathrm{rig}}_{R,c}\), with identity \(0_E\), and states
that the action is free on the retained row.  The quotient
\[
\bigl[\mathfrak M^{\mathrm{rig}}_{R,c}/E\bigr]
\]
therefore removes exactly the \(E\)-translation direction.  The
vanishing \(d^{E\text{-}\mathrm{rig}}_{R,c}=0\) says that the quotient
has no retained translation-rigidification defect.  The proof
constructs no finite closed inertia stratification, extension or flag
stack, cosection atlas, transition map, Pfaffian orientation, or
protected trace.
\end{proof}

The packet
\(\texttt{certificates/moduli/retained\_e\_translation\_rigidifications}\)
checks this row.  It imports scalar rigidifications and universal
perfect complexes, records the \(E\)-translation action on each
retained finite substack, verifies freeness on the retained row,
constructs the quotient by \(E\), and checks zero
translation-rigidification defect.  Its verifier status is
\texttt{RETAINED\_E\_TRANSLATION\_RIGIDIFICATIONS\_VERIFIED}.  The
mirrored rows in
\(\texttt{certificates/moduli/k3e\_finite\_moduli/e\_translation\_rigidifications.csv}\)
discharge item \(176\).  They do not construct finite closed inertia
stratifications, extension or flag stacks, cosection atlases,
transitions, compact Hall stages, Pfaffian orientations, or protected
traces.

\begin{definition}[Retained closed-substack datum]
\label{def:retained-closed-substack-datum}
Fix the retained \(E\)-translation-rigidified finite substacks
\(\mathfrak M^{E\text{-}\mathrm{rig}}_{R,c}\).  A retained
closed-substack datum assigns to every \(c\in C_R\) a finite family of
closed substacks
\[
\mathfrak Z_{R,c,j}\hookrightarrow
\mathfrak M^{E\text{-}\mathrm{rig}}_{R,c},
\qquad 1\le j\le m_{R,c},
\]
which covers the retained row.  It carries the finite integers
\[
m_{R,c}>0
\]
and defect ranks
\[
d^{\mathrm{cov}}_{R,c}=0,\qquad d^{\mathrm{in}}_{R,c}=0.
\]
The first defect rank says that the closed substacks cover the
retained row.  The second says that the finite residual inertia from
the scalar rigidification remains finite on every retained closed
stratum.  This is a retained closed-substack datum.  It is not closure
under extensions, not closure under Harder--Narasimhan factors, not
closure under duals, not an extension or flag stack, not a cosection
atlas, not a transition system, and not an orientation datum.
\end{definition}

\begin{proposition}[Retained closed finite substacks]
\label{prop:retained-closed-substacks}
\textnormal{[proved]}
Let the retained \(E\)-translation-rigidified finite substacks carry
the datum of Definition~\ref{def:retained-closed-substack-datum}.  Then
each retained finite row has a finite closed cover by substacks on
which the residual inertia is finite.  Equivalently, the finite
stratification row has positive strata count, finite residual inertia,
zero coverage defect, and zero inertia defect.
\end{proposition}

\begin{proof}
The retained semistable rows are finite type and specialization-closed
by Proposition~\ref{prop:retained-finite-type-semistable-substacks}.
The scalar rigidification removes \(\mathbb G_m\) and leaves finite
residual inertia by
Proposition~\ref{prop:retained-finite-residual-inertia-after-rigidification}.
The \(E\)-translation rigidification removes the elliptic translation
direction by
Proposition~\ref{prop:retained-e-translation-rigidifications}.  The
closed-substack datum then supplies the closed immersions
\(\mathfrak Z_{R,c,j}\hookrightarrow
\mathfrak M^{E\text{-}\mathrm{rig}}_{R,c}\), their finite cover, and
the equalities \(d^{\mathrm{cov}}_{R,c}=d^{\mathrm{in}}_{R,c}=0\).
Thus the retained row is represented by a finite closed cover with
finite residual inertia.  The proof constructs no extension closure,
HN-factor closure, dual closure, flag stack, cosection atlas,
transition map, Pfaffian orientation, or protected trace.
\end{proof}

The packet
\(\texttt{certificates/moduli/retained\_closed\_substacks}\) checks
this row.  It imports the retained finite-type semistable substacks,
the scalar rigidifications with finite residual inertia, and the
\(E\)-translation rigidifications, records one retained closed substack
and one closed-cover row per retained finite row, and checks zero
coverage and inertia defects.  Its verifier status is
\texttt{RETAINED\_CLOSED\_SUBSTACKS\_VERIFIED}.  The mirrored rows in
\(\texttt{certificates/moduli/k3e\_finite\_moduli/stratifications.csv}\)
discharge item \(177\).  They do not prove closure under extensions,
closure under HN factors, closure under duals, extension or flag
stacks, cosection atlases, transitions, compact Hall stages, Pfaffian
orientations, or protected traces.

\begin{definition}[Retained extension-closure datum]
\label{def:retained-extension-closure-datum}
Fix the retained finite HN word set
\[
\mathcal W_R^{\mathrm{HN}}
=\{0,s_3,s_2,s_1,e_{32},e_{21},f_{321}\}
\]
from the bounded HN type datum.  A retained extension-closure datum is
a finite set of binary extension rows
\[
u\ast v\longmapsto w,\qquad u,v,w\in \mathcal W_R^{\mathrm{HN}},
\]
with \(w\) retained and with defect rank
\[
d^{\mathrm{ext}}_{u,v}=0.
\]
For the first retained window the nontrivial rows are
\[
s_3\ast s_2=e_{32},\qquad
s_2\ast s_1=e_{21},\qquad
e_{32}\ast s_1=f_{321},\qquad
s_3\ast e_{21}=f_{321}.
\]
This is an extension-closure datum in the retained exact window.  It
is not an extension stack, not a two-step flag stack, not closure under
Harder--Narasimhan factors, not closure under duals, not a cosection
atlas, not a transition system, and not an orientation datum.
\end{definition}

\begin{proposition}[Retained substacks are closed under extensions]
\label{prop:retained-extension-closure}
\textnormal{[proved]}
Let the retained finite HN window carry the datum of
Definition~\ref{def:retained-extension-closure-datum}.  Then every
extension represented by one of the retained binary extension rows has
middle term again in the retained finite HN window.  The extension
closure defect is zero on every such row.
\end{proposition}

\begin{proof}
The bounded HN type datum supplies the finite word set
\(\mathcal W_R^{\mathrm{HN}}\).  The extension-closure datum supplies
the four nonempty binary rows displayed in
Definition~\ref{def:retained-extension-closure-datum}; their middle
terms are \(e_{32}\), \(e_{21}\), and \(f_{321}\), all elements of
\(\mathcal W_R^{\mathrm{HN}}\).  The equalities
\[
d^{\mathrm{ext}}_{s_3,s_2}
=d^{\mathrm{ext}}_{s_2,s_1}
=d^{\mathrm{ext}}_{e_{32},s_1}
=d^{\mathrm{ext}}_{s_3,e_{21}}=0
\]
say that no retained extension leaves the finite window.  Thus the
retained HN rows are closed under the supplied extensions.  The proof
constructs no extension stack, two-step flag stack, HN-factor closure,
dual closure, cosection atlas, transition map, Pfaffian orientation, or
protected trace.
\end{proof}

The packet
\(\texttt{certificates/moduli/retained\_extension\_closure}\) checks
this row.  It imports the retained HN type-bound packet and retained
closed-substack packet, records four nonempty binary extension rows,
checks that every middle HN type lies inside the retained finite HN
word set, and checks zero extension-closure defect.  Its verifier
status is \texttt{RETAINED\_EXTENSION\_CLOSURE\_VERIFIED}.  The
mirrored rows in
\(\texttt{certificates/moduli/k3e\_finite\_moduli/extension\_closure.csv}\)
discharge item \(178\).  They do not construct extension or flag
stacks, prove closure under HN factors, prove closure under duals,
construct cosection atlases, construct transitions, construct compact
Hall stages, construct Pfaffian orientations, or construct protected
traces.

\begin{definition}[Retained HN-factor-closure datum]
\label{def:retained-hn-factor-closure-datum}
Fix the retained finite HN word set
\(\mathcal W_R^{\mathrm{HN}}\) and the retained semistable factor types
\[
\mathcal T_R=\{s_3,s_2,s_1\}.
\]
A retained HN-factor-closure datum assigns to every nonzero word
\[
w=(t_1,\ldots,t_k)\in \mathcal W_R^{\mathrm{HN}}
\]
its ordered factor list \(t_i\in\mathcal T_R\), together with the
retained class \(C_{R,t_i}\), the finite-type semistable substack
\(\mathfrak M^{ss}_{R,t_i}\), and the retained closed substack
\(\mathfrak Z_{R,t_i}\) for each factor.  It carries defect ranks
\[
d^{\mathrm{HNfac}}_{w,i}=0,\qquad 1\le i\le k.
\]
This is HN-factor closure in the retained finite window.  It is not
closure under duals, not an extension or flag stack, not subquotient
closure in compactified flag stacks, not a cosection atlas, not a
transition system, and not an orientation datum.
\end{definition}

\begin{proposition}[Retained substacks are closed under HN factors]
\label{prop:retained-hn-factor-closure}
\textnormal{[proved]}
Let the retained finite HN window carry the datum of
Definition~\ref{def:retained-hn-factor-closure-datum}.  Then every HN
factor of every retained finite HN word is represented by a retained
finite-type semistable substack and by a retained closed substack.  All
HN-factor-closure defects vanish.
\end{proposition}

\begin{proof}
The bounded HN type datum supplies the retained words
\[
0,\ s_3,\ s_2,\ s_1,\ e_{32}=(s_3,s_2),\
e_{21}=(s_2,s_1),\ f_{321}=(s_3,s_2,s_1).
\]
The datum assigns to each nonzero factor occurrence the corresponding
retained class, finite-type semistable substack, and retained closed
substack.  There are ten such occurrences:
three one-factor words, two two-factor words, and one three-factor
word.  The equalities \(d^{\mathrm{HNfac}}_{w,i}=0\) say that each
factor occurrence lies in the retained finite window.  Thus passing
from a retained object to its HN factors does not leave the retained
substack system.  The proof constructs no compactified extension
stack, no two-step flag stack, no subquotient closure theorem, no dual
closure, no cosection atlas, no transition map, no Pfaffian
orientation, and no protected trace.
\end{proof}

The packet
\(\texttt{certificates/moduli/retained\_hn\_factor\_closure}\) checks
this row.  It imports retained HN type bounds, class bounds,
finite-type semistable substacks, retained closed substacks, and
retained extension closure, records ten HN-factor occurrence rows, and
checks that every factor occurrence is represented by a retained class,
a retained finite-type semistable substack, and a retained closed
substack.  Its verifier status is
\texttt{RETAINED\_HN\_FACTOR\_CLOSURE\_VERIFIED}.  The mirrored rows in
\(\texttt{certificates/moduli/k3e\_finite\_moduli/hn\_factor\_closure.csv}\)
discharge item \(179\).  They are not the dual-closure rows below, do
not construct extension or flag stacks, do not prove subquotient
closure in flag stacks, do not construct cosection atlases, do not
construct transitions, do not construct compact Hall stages, do not
construct Pfaffian orientations, and do not construct protected traces.

\begin{definition}[Retained dual-closure datum]
\label{def:retained-dual-closure-datum}
Fix the retained finite HN word set
\(\mathcal W_R^{\mathrm{HN}}\).  A retained dual-closure datum is an
involution \(D_R\) on the retained semistable factor types
\[
D_R(s_3)=s_1,\qquad D_R(s_2)=s_2,\qquad D_R(s_1)=s_3,
\]
together with the induced involution on HN words
\[
D_R(t_1,\ldots,t_k)
=(D_R(t_k),\ldots,D_R(t_1)).
\]
For the first retained window this gives
\[
D_R(0)=0,\quad D_R(s_3)=s_1,\quad D_R(s_2)=s_2,\quad
D_R(s_1)=s_3,\quad D_R(e_{32})=e_{21},\quad
D_R(e_{21})=e_{32},\quad D_R(f_{321})=f_{321}.
\]
It carries defect ranks
\[
d^{\vee}_{w}=0,\qquad w\in\mathcal W_R^{\mathrm{HN}}.
\]
This is dual closure in the retained finite window.  It is not an
extension or flag stack, not subquotient closure in compactified flag
stacks, not a cosection atlas, not a transition system, and not an
orientation datum.
\end{definition}

\begin{proposition}[Retained substacks are closed under duals]
\label{prop:retained-dual-closure}
\textnormal{[proved]}
Let the retained finite HN window carry the datum of
Definition~\ref{def:retained-dual-closure-datum}.  Then duality sends
every retained finite HN word to a retained finite HN word, and the
dual-closure defect vanishes on every retained word.
\end{proposition}

\begin{proof}
The factor duality exchanges \(s_3\) and \(s_1\) and fixes \(s_2\), so
it is an involution on the retained semistable factor set
\(\{s_3,s_2,s_1\}\).  Applying it to an HN word reverses the order and
dualizes each factor.  Hence
\[
D_R(e_{32})=D_R(s_3,s_2)=(s_2,s_1)=e_{21},\qquad
D_R(e_{21})=(s_3,s_2)=e_{32},
\]
and
\[
D_R(f_{321})=D_R(s_3,s_2,s_1)=(s_3,s_2,s_1)=f_{321}.
\]
The zero word is fixed, the one-factor words map as displayed in the
definition, and all targets belong to \(\mathcal W_R^{\mathrm{HN}}\).
The equalities \(d^{\vee}_w=0\) say that no dualized retained word
leaves the finite window.  The proof constructs no compactified
extension stack, no two-step flag stack, no subquotient closure
theorem, no cosection atlas, no transition map, no Pfaffian
orientation, and no protected trace.
\end{proof}

The packet
\(\texttt{certificates/moduli/retained\_dual\_closure}\) checks this
row.  It imports retained HN type bounds, class bounds, retained closed
substacks, and retained HN-factor closure, records the three factor
dual rows and seven HN-word dual rows, verifies that type duality and
HN-word duality are involutive, and checks that HN-word duality is
reverse factorwise duality.  Its verifier status is
\texttt{RETAINED\_DUAL\_CLOSURE\_VERIFIED}.  The mirrored rows in
\(\texttt{certificates/moduli/k3e\_finite\_moduli/dual\_closure.csv}\)
discharge item \(180\).  They do not construct extension or flag
stacks, prove subquotient closure in flag stacks, construct cosection
atlases, construct transitions, construct compact Hall stages,
construct Pfaffian orientations, or construct protected traces.

\begin{definition}[Algebraic/effective Hall support]
\label{def:algebraic-effective-hall-support}
Fix a polarised K3 surface $(S, H)$, a chosen Mukai vector lattice
$\Lambda_{S, \mathrm{alg}} \subset \Lambda_S$ (the image of the
algebraic K-class lattice generated by sheaves whose Mukai vector lies in
$H^0\oplus\mathrm{NS}(S)\oplus H^4$), and the retained stability datum
\(\sigma=\sigma_{s,t}\) with heart
\(\mathcal A_X^{\mathrm{AP/Liu}}\) from
Definition~\ref{def:retained-ap-liu-heart}.  The retained-window rows
above supply noetherianity, Harder--Narasimhan filtrations, bounded HN
types, finite-type semistable substacks, and quasi-smooth derived
enhancements, together with finite residual inertia after scalar
rigidification, finite retained class-bound data, and universal
perfect complexes, \(E\)-translation rigidifications, and finite
closed-cover stratifications at finite height, together with extension
closure, HN-factor closure, and dual closure inside the retained HN
window; later rows supply flag stacks, subquotient closure in flag
stacks, cosection charts, transitions, scalar firewalls, and
compact Hall support.  No unconditional Bridgeland stability theorem on the
threefold \(S\times E\) is used.  Set
\[
\Gamma_X^{\mathrm{alg}}
:= \Lambda_{S,\mathrm{alg}}\oplus\Lambda_{S,\mathrm{alg}},\qquad
\Gamma_{X, \sigma}^{\mathrm{eff}}
\subset \Gamma_X^{\mathrm{alg}}
\subset \Gamma_X^{\mathrm{form}},
\]
where the \emph{$\sigma$-effective semigroup}
$\Gamma_{X, \sigma}^{\mathrm{eff}}$ is the additive sub-semigroup of
charges supported by objects semistable for this datum and of nonzero
Mukai vector.
Membership in $\Gamma_{X, \sigma}^{\mathrm{eff}}$ is not implied by
membership in $\Gamma_X^{\mathrm{alg}}$ or by formal lattice arithmetic
(folklore).
\end{definition}

\begin{definition}[Admissible Dirac--Igusa parameter datum]
\label{def:admissible-dirac-igusa-parameter-datum}
An admissible parameter datum for the Dirac--Igusa construction is
\[
\mathfrak p=(S,H,\sigma,E,0_E,\tau_X,\mathcal R_{\mathrm{HN}}),
\qquad X=S\times E,
\]
where \((S,H)\) is a polarized smooth complex K3 surface, \(E\) is a
smooth elliptic curve with origin \(0_E\), \(\tau_X:K_X\simeq\mathcal O_X\)
is the chosen trivialization, \(\sigma\) is the retained stability datum
used to define \(\Gamma_{X,\sigma}^{\mathrm{eff}}\), and
\(\mathcal R_{\mathrm{HN}}\) is the cofinal finite HN index system.
It includes the retained finite charge, moduli, hybrid, orientation,
Hall, Pfaffian, Koszul, and primitive-recognition packets whenever
those packets are invoked later.  Thus \(H\) and \(\sigma\) are not
decorations: they determine the algebraic/effective Hall support and
the finite HN windows.

A morphism \(\alpha:\mathfrak p\to\mathfrak p'\) is an admissible
base-change isomorphism
\[
\alpha=(f,u,\eta_\sigma,\eta_{\mathrm{HN}})
\]
with \(f:S\to S'\) an isomorphism of K3 surfaces carrying \(H\) to
\(H'\), \(u:E\to E'\) an isomorphism of elliptic curves with
\(u(0_E)=0_{E'}\), compatibility with the chosen trivializations of
\(K_{S\times E}\), a transport \(\eta_\sigma\) of the retained
stability datum and effective Hall support, and a cofinal identification
\(\eta_{\mathrm{HN}}:\mathcal R_{\mathrm{HN}}\to\mathcal R'_{\mathrm{HN}}\).
The induced Mukai isometry sends \((Q,P)\) to
\((f_*Q,f_*P)\), preserves the Mukai pairing, and therefore preserves
\[
\Pi_X(Q,P)=((Q,Q)/2,(Q,P),(P,P)/2),\qquad B=-\delta\Pi_X .
\]
The resulting groupoid is denoted \(\mathsf{Par}^{\mathrm{DI}}\).
Arbitrary non-isomorphic deformations of \(S\), arbitrary changes of
stability chamber, and arbitrary changes of elliptic origin are not
morphisms of \(\mathsf{Par}^{\mathrm{DI}}\) unless the corresponding
wall-crossing or origin-translation datum below is supplied.
\end{definition}

\begin{proposition}[Admissible base-change functoriality]
\label{prop:admissible-base-change-functoriality}
Let \(\alpha:\mathfrak p\to\mathfrak p'\) be a morphism in
\(\mathsf{Par}^{\mathrm{DI}}\).  If the finite packets attached to
\(\mathfrak p\) have zero defect rows and are transported by
\(\alpha\), then pull-push along \(f\times u:X\to X'\) induces
functors at each retained height
\[
\alpha_{R,*}:\mathsf{DI}^{\mathrm{fin}}_{X,R}
\longrightarrow
\mathsf{DI}^{\mathrm{fin}}_{X',\eta_{\mathrm{HN}}(R)}
\]
preserving the components \(\mathcal A^{E_3}\), \(F^{\mathrm{hyb}}\),
\(\widehat\Gamma\), \(\Pi\), \(\Phi\), \(o\), \(H\), \(P^\Pi\), \(C\),
\(\Theta_{\mathrm{Kos}}\), \(\mathscr L_{\mathrm{Pf}}\),
\(\operatorname{pf}\), \(\varepsilon_o\), and \(\mathrm{Rec}\), in the
sense of the morphisms \textup{(M1)}--\textup{(M8)} of
Definition~\ref{def:finite-stage-dirac-igusa-morphism}.  These functors
commute with HN transition morphisms and therefore induce a functor on
Dirac--Igusa pro-objects
\[
\alpha_*:
\operatorname{Pro}(\mathsf{DI}^{\mathrm{fin}}_X)
\longrightarrow
\operatorname{Pro}(\mathsf{DI}^{\mathrm{fin}}_{X'}).
\]
For composable admissible morphisms, the induced functors compose up to
the canonical pro-isomorphism coming from the common cofinal HN
refinement.  This is the base-change functoriality in \(S\) and \(E\)
used by the construction.
\end{proposition}

\begin{proof}
The isomorphism \(f\) gives an isometry of Mukai lattices and carries
the algebraic sublattice, the polarization, and the retained stability
support by hypothesis.  The isomorphism \(u\) preserves \(1_E\),
\(\omega_E\), the origin, and the translation group law.  Hence
\((f\times u)_*\) preserves the formal charge lattice, the
normal-ordered central extension, and the Gram map.  The assumed
transport of the finite packets carries the moduli stacks,
correspondence stacks, orientation lines, Pfaffian lines, Hall
bialgebra matrices, Koszul source complexes, and primitive-recognition
matrices to their primed counterparts.  Compatibility with the
transition maps is exactly the zero-defect condition in the finite
morphism tables.  Proposition~\ref{prop:dirac-igusa-pro-object-universal-property}
then gives the induced pro-functor and the cofinal composition law.
\end{proof}

\begin{definition}[Polarization and elliptic-origin dependence]
\label{def:polarization-origin-dependence}
The dependence on the K3 polarization is the dependence of
\(\Gamma_{X,\sigma}^{\mathrm{eff}}\), the HN index system, and the
finite moduli stacks on \((H,\sigma)\).  If two choices
\((H,\sigma)\) and \((H',\sigma')\) lie in chambers connected by a
specified Kontsevich--Soibelman wall-crossing datum
\(\mathsf{WC}_{H,H',R}\) with zero finite defects, the resulting
Dirac--Igusa pro-objects are compared by the induced pro-isomorphism.
Without such a wall-crossing datum they are different parameter
points, not two presentations of the same object.

Let \(a\in E\) and replace the origin \(0_E\) by \(a\).  For a retained
wrapped family \(\mathcal F\), the determinant anchors of
Definition~\ref{def:determinant-anchor} satisfy the canonical identity
\[
\lambda^{\det}_{a}(\mathcal F)
\simeq
\lambda^{\det}_{0_E}(\mathcal F)\otimes
\mathcal O_E\!\bigl(\chi(\mathcal F)(0_E-a)\bigr)
\quad\text{in }\operatorname{Pic}^0(E).
\]
Thus changing the origin contributes the degree-zero character
\[
\vartheta_{a,\mathcal F}:=
\mathcal O_E\!\bigl(\chi(\mathcal F)(0_E-a)\bigr)\in \operatorname{Pic}^0(E),
\]
and, after restriction to a finite stabilizer \(H\subset E[N]\), the
corresponding finite character
\[
\vartheta_{a,\mathcal F}^{H}\in H^1(BH;\mathbb C^\times)
\]
is the only permitted change in the quotient-orientation linearization.
The Dirac--Igusa object is independent of the elliptic origin only
after twisting the wrapped-anchor and orientation-linearization rows by
this character and verifying the same finite stabilizer equations.
\end{definition}

\begin{proposition}[Elliptic-origin change up to its canonical character]
\label{prop:elliptic-origin-independence-character}
Fix \(S,H,\sigma,E\) and the finite HN system, and let
\(\mathfrak p_0\) and \(\mathfrak p_a\) be the admissible parameter
data obtained from the origins \(0_E\) and \(a\in E\).  Suppose the
finite hybrid, orientation, and morphism rows for \(\mathfrak p_a\) are
obtained from those for \(\mathfrak p_0\) by tensoring every wrapped
anchor with
\(\vartheta_{a,\mathcal F}=\mathcal O_E(\chi(\mathcal F)(0_E-a))\)
and by adding the restricted character
\(\vartheta^H_{a,\mathcal F}\) to the finite stabilizer linearization
row for each retained stabilizer \(H\subset E[N]\).  Then the two
finite-stage systems determine canonically pro-isomorphic
Dirac--Igusa pro-objects after this character twist.  Before the twist,
the obstruction to identifying the two origin choices is exactly the
family of characters \(\vartheta^H_{a,\mathcal F}\).
\end{proposition}

\begin{proof}
By Definition~\ref{def:determinant-anchor},
\[
\lambda^{\det}_{a}(\mathcal F)
=\det Rp_{E,*}\mathcal F\otimes\mathcal O_E(-\chi(\mathcal F)a)
\]
and
\[
\lambda^{\det}_{0_E}(\mathcal F)
=\det Rp_{E,*}\mathcal F\otimes\mathcal O_E(-\chi(\mathcal F)0_E).
\]
Their ratio is
\(\mathcal O_E(\chi(\mathcal F)(0_E-a))\), a canonical point of
\(\operatorname{Pic}^0(E)\).  Restriction to a finite stabilizer gives
the character \(\vartheta^H_{a,\mathcal F}\).  All other components of
the formal charge lattice, the Mukai--Gram map, and the K3 data are
unchanged.  Therefore the finite-stage morphisms are the identity on
all components except the wrapped-anchor and orientation-linearization
rows, where they are precisely the displayed tensor and character
twists.  The zero-defect assumptions give commuting transition squares;
Proposition~\ref{prop:dirac-igusa-pro-object-universal-property} then
identifies the resulting pro-objects.  If the character rows are not
twisted, the displayed ratio remains as the residual finite-stabilizer
linearization class, so the obstruction is exactly
\(\vartheta^H_{a,\mathcal F}\).
\end{proof}

\begin{remark}[Three-tier inclusion]
\label{rem:three-tier-inclusion}
The hierarchy $\Gamma_{X, \sigma}^{\mathrm{eff}}
\subset \Gamma_X^{\mathrm{alg}}
\subset \Gamma_X^{\mathrm{form}}$
distinguishes three discipline registers.  $\Gamma_X^{\mathrm{form}}$ is
the formal additive group on which $\Pi_X$ is defined and the
polarization identity $B = -\delta\Pi_X$ holds.  $\Gamma_X^{\mathrm{alg}}$
is the K-theoretic algebraic sublattice; its rank-one extension closure
inside $\Gamma_X^{\mathrm{form}}$ records only the algebraic Mukai
branch (computed, \cite{Mukai1987,KSCoHA}).
$\Gamma_{X, \sigma}^{\mathrm{eff}}$ is the Hall support selected by
$\sigma$; effectivity and semistability are not lattice-theoretic
properties and are not fibre-functorial in the K3 type.  The sectorial
Hall truncation of \S\ref{sec:k3xe-hybrid-ran-carrier} works at this
third level after imposing finite Harder--Narasimhan height bounds.
\end{remark}

\begin{lemma}[Two orthogonal hyperbolic planes inside $\Lambda_S$]
\label{lem:two-hyperbolic-planes}
The Mukai lattice contains two orthogonal hyperbolic planes:
\[
U \oplus U \hookrightarrow \Lambda_S,\qquad
U_1 = \Z e_1 \oplus \Z f_1,\quad U_2 = \Z e_2 \oplus \Z f_2,
\]
with $e_i^2 = f_i^2 = 0$, $e_i\cdot f_i = 1$, and $U_1\perp U_2$ inside
the Mukai pairing (proved, \cite{NIKULIN}).
\end{lemma}

\begin{proof}
Direct from $\Lambda_S \cong U^{\oplus 4} \oplus E_8(-1)^{\oplus 2}$.
\end{proof}

\section{Quadratic Gram map and lattice polarization}
\label{sec:gram-map-polarization}

\subsection{Igusa variables and the Gram-index lattice}
\label{subsec:igusa-variables-gram-index}

For a Siegel period
\[
Z = \begin{pmatrix} z_1 & z_2 \\ z_2 & z_3 \end{pmatrix}
\in \mathcal H_2,\qquad
q = e^{2\pi i z_1},\quad r = e^{2\pi i z_2},\quad s = e^{2\pi i z_3},
\]
and
\[
T(n, l, m) = \begin{pmatrix} n & l/2 \\ l/2 & m \end{pmatrix},\qquad
n, l, m \in \Z,
\]
the trace pairing
\[
\bigl(T(n, l, m), Z\bigr)_{\mathrm{tr}}
:= \mathrm{tr}\bigl(T(n, l, m)\,Z\bigr)
= n z_1 + l z_2 + m z_3
\]
gives the Borcherds monomial $x^{(n, l, m)} = q^n r^l s^m$.

\begin{definition}[Gram-index lattice and effective semigroup]
\label{def:gram-index-lattice}
The \emph{Gram-index lattice} is
\[
\Gamma_{\mathrm{gram}} := \Z^3,\qquad
\gamma = (n, l, m).
\]
The \emph{cusp-positive semigroup} at the standard cusp $\mathfrak
c_\infty$ with chamber convention $0 < |s| \ll |q| \ll |r|^{-1} \ll 1$
is
\[
\Gamma_{\mathrm{eff}}
= \{(n, l, m) : m > 0,\ n \ge 0,\ l \in \Z\}
\cup \{(n, l, 0) : n > 0,\ l \in \Z\}
\cup \{(0, l, 0) : l < 0\}.
\]
The \emph{Jacobi-active sublocus}
$\Gamma_{\mathrm{act}} \subset \Gamma_{\mathrm{eff}}$ is the locus of
nonvanishing $f(nm, l)$ in the weak Jacobi form $\phi_{0, 1}$
(computed, \cite{EZ:1}).
\end{definition}

\begin{proposition}[Gram matrix convention for protected charges]
\label{prop:gram-matrix-convention}
For a formal even charge $c = (Q, P) \in \Gamma_X^{\mathrm{form}}$ in
the unimodular lattice $\Lambda_S$, set
\[
n = \frac{(Q, Q)}{2},\qquad
l = (Q, P),\qquad
m = \frac{(P, P)}{2}.
\]
Then $n, l, m \in \Z$ (since $\Lambda_S$ is even unimodular,
proved),
\[
\begin{pmatrix} (Q, Q) & (Q, P) \\ (P, Q) & (P, P) \end{pmatrix}
= 2\,T(n, l, m),\qquad
\det\bigl(2 T(n, l, m)\bigr) = 4 n m - l^2,
\]
and $x^{(n, l, m)} = q^n r^l s^m$ is the Fourier character of
$\bigl(T(n, l, m), Z\bigr)_{\mathrm{tr}}$.
\end{proposition}

\begin{proof}
The Mukai pairing is integer-valued and even on $\Lambda_S\oplus\Lambda_S$
by Definition~\ref{def:mukai-lattice}.  The matrix identity is the
direct expansion
\[
\begin{pmatrix} (Q, Q) & (Q, P) \\ (P, Q) & (P, P) \end{pmatrix}
= \begin{pmatrix} 2 n & l \\ l & 2 m \end{pmatrix}
= 2\,T(n, l, m).
\]
The determinant follows.  The Fourier character is
$\exp(2\pi i\,\mathrm{tr}(T Z)) = q^n r^l s^m$.
\end{proof}

\subsection{The Gram map and its bilinear cocycle}
\label{subsec:gram-map-bilinear-cocycle}

\begin{definition}[Quadratic Gram map and bilinear cocycle]
\label{def:gram-map-cocycle}
The \emph{Gram map} is the quadratic projection
\[
\Pi_X : \Gamma_X^{\mathrm{form}} \longrightarrow \Gamma_{\mathrm{gram}},
\qquad
\Pi_X(Q, P)
= \left(\frac{(Q, Q)}{2},\,(Q, P),\,\frac{(P, P)}{2}\right).
\]
The \emph{symmetric bilinear cocycle} $B$ associated to $\Pi_X$ is
\[
B : \Gamma_X^{\mathrm{form}} \times \Gamma_X^{\mathrm{form}}
\longrightarrow \Gamma_{\mathrm{gram}},\qquad
B\bigl((Q, P), (Q', P')\bigr)
= \bigl((Q, Q'),\,(Q, P') + (Q', P),\,(P, P')\bigr).
\]
\end{definition}

\begin{lemma}[Gram additivity defect]
\label{lem:gram-additivity-defect}
For $c, c' \in \Gamma_X^{\mathrm{form}}$,
\[
\Pi_X(c + c') = \Pi_X(c) + \Pi_X(c') + B(c, c').
\]
The defect $B$ is symmetric, $\Z$-bilinear, integer-valued, and
satisfies polarization: $B(c, c) = 2\,\Pi_X(c)$.  As an ordinary
additive cochain identity,
\[
B = -\delta\Pi_X,\qquad
(\delta q)(c, c') = q(c) + q(c') - q(c + c').
\]
(proved.)
\end{lemma}

\begin{remark}[Quadratic primitive, no linear primitive]
\label{rem:B-quadratic-not-linear}
The equality \(B=-\delta\Pi_X\) is an equality in the ordinary
inhomogeneous cochain complex, where the primitive \(\Pi_X\) is
quadratic.  There is no additive homomorphism
\(L:\Gamma_X^{\mathrm{form}}\to\Gamma_{\mathrm{gram}}\) with
\(B=-\delta L\).  Thus the normal-ordering defect is removed by the
quadratic Gram coordinate, not by a linear change of charge lattice.
\end{remark}

\begin{proof}
Write $c = (Q, P)$, $c' = (Q', P')$; then
\[
\frac{(Q + Q')^2}{2}
= \frac{Q^2}{2} + \frac{Q'^2}{2} + Q\cdot Q',
\]
\[
(Q + Q')\cdot(P + P')
= Q\cdot P + Q'\cdot P' + Q\cdot P' + Q'\cdot P,
\]
\[
\frac{(P + P')^2}{2}
= \frac{P^2}{2} + \frac{P'^2}{2} + P\cdot P'.
\]
The three cross terms are the three components of $B(c, c')$.
Symmetry, bilinearity, and integrality follow from the corresponding
properties of the Mukai pairing.  Polarization is the case $c' = c$.
The cochain identity rewrites the additivity defect.
\end{proof}

\begin{lemma}[Bilinear cocycle identity]
\label{lem:b-cocycle-identity}
$B$ is normalized, $B(c, 0) = B(0, c) = 0$, and the additive coboundary
\[
(\delta B)(a, b, c) := B(b, c) - B(a + b, c) + B(a, b + c) - B(a, b)
\]
vanishes identically on $\Gamma_X^{\mathrm{form}}$ (proved).
\end{lemma}

\begin{proof}
By bilinearity in each slot,
\[
B(a + b, c) = B(a, c) + B(b, c),\qquad
B(a, b + c) = B(a, b) + B(a, c).
\]
Substituting,
\[
(\delta B)(a, b, c)
= B(b, c) - B(a, c) - B(b, c) + B(a, b) + B(a, c) - B(a, b) = 0.
\]
The polarization identity $B(c, c) = 2\,\Pi_X(c)$ is the only place where
$B$ takes a non-vanishing value on the diagonal; this is structural,
not a Bott-element framing.
\end{proof}

\begin{remark}[Polarization, not Bott element]
\label{rem:polarization-not-bott}
The identity $B(c, c) = 2\,\Pi_X(c)$ is the polarization identity for
the symmetric bilinear $B$ associated to the quadratic $\Pi_X$.  It is
forced by the homogeneous-of-degree-two structure of the Gram map.  No
external Bott-element framing or characteristic-class coincidence is
involved.  The full structural content is that $\Pi_X$ is a homogeneous
quadratic on the additive group $\Gamma_X^{\mathrm{form}}$ valued in
$\Gamma_{\mathrm{gram}}$, with associated symmetric bilinear $B$ given by
the Mukai pairing of components on each direction.
\end{remark}

\subsection{Formal lift of every Gram triple}
\label{subsec:formal-gram-triple-lift}

\begin{lemma}[Formal primitive torsion-one lift of every Gram triple]
\label{lem:primitive-lift-every-gram-triple}
Suppose $\Lambda_S$ contains two orthogonal hyperbolic planes $U_1\oplus
U_2$ with bases $(e_1, f_1, e_2, f_2)$ as in
Lemma~\ref{lem:two-hyperbolic-planes}.  Then every Gram triple
$(n, l, m) \in \Z^3$ admits a primitive saturated formal full-Mukai lift
$c = (Q, P) \in \Gamma_X^{\mathrm{form}}$ with $\Pi_X(Q, P) = (n, l, m)$.
One may take
\[
Q = e_1 + n f_1,\qquad
P = l f_1 + e_2 + m f_2.
\]
For this lift, the wedge torsion invariant
\[
I(Q, P)
:= \gcd\!\bigl\{\text{coordinates of }Q\wedge P\text{ in a }\Z\text{-basis of }
{\textstyle\bigwedge^2\Lambda_S}\bigr\}
\]
equals $1$.
(proved.)
\end{lemma}

\begin{proof}
The hyperbolic-plane pairings give
\[
Q^2 = 2 n,\qquad P^2 = 2 m,\qquad Q\cdot P = l,
\]
hence $\Pi_X(Q, P) = (n, l, m)$.  The submodule $\Z Q + \Z P \subset
\Lambda_S$ contains the unimodular minor
$\det\!\begin{pmatrix} 1 & 0 \\ 0 & 1 \end{pmatrix} = 1$ in the rows
indexed by $e_1, e_2$ of the coefficient matrix in
$(e_1, f_1, e_2, f_2)$, so the lift is primitive saturated.
The wedge expansion
\[
Q\wedge P
= l\,e_1\wedge f_1 + e_1\wedge e_2 + m\,e_1\wedge f_2
+ n\,f_1\wedge e_2 + nm\,f_1\wedge f_2
\]
contains the basis element $e_1\wedge e_2$ with coefficient $1$, so the
gcd is $1$ and $I(Q, P) = 1$.
\end{proof}

\begin{remark}[Pointwise existence is not Hall support]
\label{rem:pointwise-not-hall-support}
Lemma~\ref{lem:primitive-lift-every-gram-triple} produces a formal
charge plane realizing every prescribed Igusa degree.  It does not prove
algebraicity of $(Q, P)$, effectivity, non-empty compact Hall moduli,
duality-orbit descent, source-index descent, or any chain-level Hall
strictification.  The compact Hall obstruction is functorial
compatibility of pointwise existence with Hall addition and the
primitive bracket; only the latter is the working content of
Chapter~\ref{chap:dirac-igusa-pro-object}.
\end{remark}

The packet
\(\texttt{certificates/charge/formal\_primitive\_lift}\) is the finite
table-level check of Lemma~\ref{lem:primitive-lift-every-gram-triple}.
It verifies, on the zero row, the three simple-root rows, the
\(\delta_{123}\) row, and one negative-pairing test row, that
\[
\Pi_X(Q,P)=(n,l,m),\qquad
I(Q,P)=1,\qquad
\alpha(\Pi_X(Q,P))=2n f_2-l f_3+2m f_{-2}.
\]
The same packet records the \(e_1,e_2\) minor equal to \(1\), hence
primitivity of the displayed formal lift.  It is not evidence for
algebraic effectivity, finite HN boundedness, compact Hall support,
source Hall brackets, Pfaffian orientations, type-II wall atlases,
mirror discriminants, or protected traces.

\section{Normal-ordered Gram extension}
\label{sec:normal-ordered-extension}

\subsection{Central extension and the additive normal-ordered Gram
homomorphism}
\label{subsec:central-extension-additive-gram}

\begin{definition}[Normal-ordered extension and Gram homomorphism]
\label{def:normal-ordered-extension}
Set
\[
\widehat\Gamma_X
:= \Gamma_X^{\mathrm{form}}\oplus\Gamma_{\mathrm{gram}},
\]
with operation
\[
(c, T)\star(c', T')
:= \bigl(c + c',\,T + T' + B(c, c')\bigr).
\]
Define the \emph{normal-ordered Gram map}
\[
\overline\Pi_X : \widehat\Gamma_X \longrightarrow \Gamma_{\mathrm{gram}},
\qquad
\overline\Pi_X(c, T) = \Pi_X(c) - T.
\]
The \emph{raw placement} and the \emph{split section} are
\[
i_0 : \Gamma_X^{\mathrm{form}}\longrightarrow\widehat\Gamma_X,\quad
i_0(c) = (c, 0),\qquad
s : \Gamma_X^{\mathrm{form}}\longrightarrow\widehat\Gamma_X,\quad
s(c) = (c, \Pi_X(c)).
\]
\end{definition}

\begin{theorem}[Group law of the normal-ordered extension; additivity of
$\overline\Pi_X$]
\label{thm:normal-ordered-extension-group-law}
Let $\widehat\Gamma_X$, $\star$, $\overline\Pi_X$, $i_0$, $s$ be as in
Definition~\ref{def:normal-ordered-extension}.  Then:
\begin{enumerate}[label = \textup{(\roman*)}]
\item $(\widehat\Gamma_X, \star)$ is an abelian group with identity
$(0, 0)$ and inverse $(c, T)^{-1} = (-c,\,-T + B(c, c)) = (-c,\,-T +
2\,\Pi_X(c))$.
\item The sequence
\[
0\to\Gamma_{\mathrm{gram}}\xrightarrow{T\mapsto(0, T)}\widehat\Gamma_X
\xrightarrow{\mathrm{pr}_1}\Gamma_X^{\mathrm{form}}\to 0
\]
is a central extension of abelian groups with cocycle $B$.
\item $\overline\Pi_X$ is a group homomorphism from
$(\widehat\Gamma_X, \star)$ to the additive group $\Gamma_{\mathrm{gram}}$.
\item The split section $s(c) = (c, \Pi_X(c))$ is additive, and
\[
\overline\Pi_X(i_0(c)) = \Pi_X(c),\qquad
\overline\Pi_X(s(c)) = 0.
\]
\item The change of variables $(c, T)\mapsto(c, T - \Pi_X(c))$ carries
$\star$ to componentwise addition; thus $\widehat\Gamma_X\cong
\Gamma_X^{\mathrm{form}}\oplus\Gamma_{\mathrm{gram}}$ as abstract abelian
groups, and $[B] = 0$ in additive cohomology.
\end{enumerate}
(proved.)
\end{theorem}

\begin{proof}
\emph{(i)} Associativity follows from the cocycle identity
$\delta B = 0$ (Lemma~\ref{lem:b-cocycle-identity}); the associator
$B(a, b) + B(a + b, c) - B(a, b + c) - B(b, c)$ vanishes by bilinearity.
Commutativity is symmetry of $B$.  The unit calculation $(c, T)\star(0,
0) = (c, T + B(c, 0)) = (c, T)$ and inverse $(c, T)\star(-c, -T + B(c,
c)) = (0, B(c, c) - B(c, c)) = (0, 0)$ (using $B(c, -c) = -B(c, c)$ by
bilinearity) close the abelian-group structure.

\emph{(ii)} Definition.

\emph{(iii)} Apply Lemma~\ref{lem:gram-additivity-defect}:
\[
\overline\Pi_X((c, T)\star(c', T'))
= \Pi_X(c + c') - T - T' - B(c, c')
= \Pi_X(c) + \Pi_X(c') - T - T'
= \overline\Pi_X(c, T) + \overline\Pi_X(c', T').
\]

\emph{(iv)} For the section,
\[
s(c)\star s(c')
= (c + c',\,\Pi_X(c) + \Pi_X(c') + B(c, c'))
= (c + c',\,\Pi_X(c + c'))
= s(c + c').
\]
The identities $\overline\Pi_X(i_0(c)) = \Pi_X(c) - 0 = \Pi_X(c)$ and
$\overline\Pi_X(s(c)) = \Pi_X(c) - \Pi_X(c) = 0$ are immediate.

\emph{(v)} The map $(c, T)\mapsto(c, T - \Pi_X(c))$ is a bijection with
inverse $(c, T)\mapsto(c, T + \Pi_X(c))$.  Direct calculation gives
\[
(c, T - \Pi_X(c)) + (c', T' - \Pi_X(c'))
= (c + c',\,T + T' - \Pi_X(c) - \Pi_X(c'))
= (c + c',\,(T + T' + B(c, c')) - \Pi_X(c + c')),
\]
matching the change of variables applied to $(c, T)\star(c', T')$.
\end{proof}

\begin{remark}[Finite charge-window verification]
\label{rem:finite-charge-window-verification}
The normal-ordered group law is a proved lattice identity.  Its use in
the Dirac--Igusa pro-object additionally requires a finite
charge-window packet: active \(\phi_{0,1}\)-support, downward
saturation, Liu HN bounds, Mukai--Gram rows, cocycle identities,
normal-ordered lift orbits, target-window images, primitive formal
lifts, strict transitions, and a scalar firewall.  A target
presentation, signed exponent list, Pfaffian product, Hilbert-scheme
scalar, or source-rank table does not by itself supply this packet.
\end{remark}

\subsection{The fixed-lift raw-pushforward obstruction theorem}
\label{subsec:raw-gram-descent-no-go}

\begin{definition}[Raw homogeneous and fixed-lift $\Pi_X$-pushforwards]
\label{def:raw-pi-pushforward}
Let $P = \bigoplus_{c\in\Gamma_X^{\mathrm{form}}} P_c$ be a
$\Gamma_X^{\mathrm{form}}$-graded primitive object.  The
\emph{raw homogeneous $\Pi_X$-pushforward} is
\[
\bigl(\Pi_{X, *}^{\mathrm{raw}} P\bigr)_\gamma
:= \bigoplus_{\Pi_X(c) = \gamma} P_c.
\]
The \emph{strict fixed-lift raw $\Pi_X$-pushforward} is the narrower
operation obtained by choosing one lift $L(\gamma) \in
\Gamma_X^{\mathrm{form}}$ for every retained $\gamma$ and setting
\[
\bigl(\Pi_{X, *}^{\mathrm{raw}, L} P\bigr)_\gamma
:= P_{L(\gamma)}.
\]
\end{definition}

\begin{theorem}[Raw homogeneous fixed-lift $\Pi_X$-pushforward cannot
realize the type-II real-root strings]
\label{thm:raw-gram-descent-no-go}
Let
$P = \bigoplus_{c\in\Gamma_X^{\mathrm{form}}} P_c$ be a
$\Gamma_X^{\mathrm{form}}$-graded Lie superalgebra with $[P_c, P_{c'}]
\subset P_{c + c'}$.  Suppose a strict fixed-lift raw pushforward of
$P$ along the quadratic map $\Pi_X$ were a $\Gamma_{\mathrm{gram}}$-graded
Lie superalgebra structure on the displayed bracket channels of chosen
charge lifts.  Then every nonzero bracket channel $[P_c, P_{c'}]$
contributing to that strict raw pushforward must satisfy $B(c, c') = 0$.

For the type-II real simple roots $(\delta_1, \delta_2, \delta_3)$ of
the Gritsenko--Nikulin algebra $\mathfrak g_{\Delta_5}$ with Cartan
matrix
\[
A_{\mathrm{II}}
= \begin{pmatrix} 2 & -2 & -2 \\ -2 & 2 & -2 \\ -2 & -2 & 2 \end{pmatrix},
\]
this condition is incompatible with the nonzero second real-root strings
$2\delta_i + \delta_j$, $i \ne j$.  The strict fixed-lift raw
$\Pi_X$-descent therefore cannot carry these brackets; they require the
normal-ordered extension $\widehat\Gamma_X$ and the additive homomorphism
$\overline\Pi_X$.
(proved, \cite{GN}.)
\end{theorem}

\begin{proof}
If the raw pushforward is $\Gamma_{\mathrm{gram}}$-graded, a nonzero
bracket of classes of raw Gram degrees $\Pi_X(c)$ and $\Pi_X(c')$ must
land in raw Gram degree $\Pi_X(c) + \Pi_X(c')$.  But the upstairs bracket
lands in physical degree $c + c'$, whose raw Gram shadow is
\[
\Pi_X(c + c') = \Pi_X(c) + \Pi_X(c') + B(c, c').
\]
Hence $B(c, c') = 0$ on every nonzero bracket channel.

Take physical lifts $c_i, c_j$ of distinct type-II real simple root
degrees $\gamma_i, \gamma_j$, $\Pi_X(c_i) = \gamma_i$ with $\gamma_i\ne
0$.  The Gritsenko--Nikulin Cartan matrix $A_{\mathrm{II}}$ has
off-diagonal entries $-2$, so distinct simple roots satisfy
$\langle\delta_i, \delta_j\rangle = -2$, and the real-root string
through $\delta_j$ in the $\delta_i$-direction includes the levels
$\delta_i + \delta_j$ and $2\delta_i + \delta_j$ (proved,
\cite[\S2--3]{GN}).  Both correspond to nonzero brackets, hence raw
pushforward forces
\[
B(c_i, c_j) = 0,\qquad
B(c_i, c_i + c_j) = 0.
\]
Bilinearity and polarization give
\[
B(c_i, c_i + c_j) = B(c_i, c_i) + B(c_i, c_j)
= 2\,\Pi_X(c_i) + 0 = 2\,\gamma_i,
\]
which is nonzero in the torsion-free lattice $\Gamma_{\mathrm{gram}}$.
Contradiction.
\end{proof}

\begin{remark}[Scope of the no-go]
\label{rem:no-go-scope}
Theorem~\ref{thm:raw-gram-descent-no-go} excludes the strict
\emph{homogeneous fixed-lift} raw pushforward.  It does not exclude
isolated $B$-isotropic bracket channels and does not decide the
mixed-lift fibre-summed construction in which several lifts of a single
Gram degree are summed before bracketing; the latter remains an open
mixed-lift question.  The normal-ordered extension removes exactly the
displayed obstruction because $\overline\Pi_X$ is additive on
$\widehat\Gamma_X$.  Chain-level Hall strictification
($A_\infty$-compatibility, cyclic equivariance, Frobenius,
multiplicative orientation) is a separate input system, named in
\S\ref{sec:k3xe-four-obstructions} and developed in
Chapter~\ref{chap:dirac-igusa-pro-object}.
\end{remark}

\begin{definition}[Gram fibres, zero fibre, and fibre-summed parity blocks]
\label{def:gram-fibres-zero-fibre}
For \(\gamma\in\Gamma_{\mathrm{gram}}\), set
\[
\Gamma_X(\gamma):=\Pi_X^{-1}(\gamma)\subset\Gamma_X^{\mathrm{form}}.
\]
The \emph{zero Gram fibre} is
\[
K_{\Pi,0}:=\Gamma_X(0,0,0)
=\{(Q,P)\mid Q^2=P^2=Q\cdot P=0\}.
\]
It is a fibre of the quadratic map \(\Pi_X\), not the kernel of a group
homomorphism.  If \(V=\bigoplus_{c\in A}V_c\) is a finite
\(\Z/2\)-graded formal source over a finite subset
\(A\subset\Gamma_X^{\mathrm{form}}\), the fibre-summed Gram pushforward
is the \(\Z/2\)-graded vector space
\[
(\Pi_{X,*}^{\mathrm{fib}}V)_{\gamma,\bar p}
:=\bigoplus_{\substack{c\in A\cap\Gamma_X(\gamma)}}V_{c,\bar p},
\qquad \bar p\in\Z/2.
\]
The signed superdimension
\(\dim(\Pi_{X,*}^{\mathrm{fib}}V)_{\gamma,\bar0}
-\dim(\Pi_{X,*}^{\mathrm{fib}}V)_{\gamma,\bar1}\)
is a scalar shadow of this two-integer datum.
\end{definition}

\begin{proposition}[The zero Gram fibre is not an additive subgroup]
\label{prop:zero-gram-fibre-not-subgroup}
\textnormal{[proved]}
In the two-hyperbolic-plane chart of
Lemma~\ref{lem:primitive-lift-every-gram-triple}, the zero Gram fibre
\(K_{\Pi,0}\) is not closed under addition.  Hence no quotient of
\(\Gamma_X^{\mathrm{form}}\) by \(K_{\Pi,0}\) is defined at the level of
abelian groups.
\end{proposition}

\begin{proof}
Let
\[
k_1=(e_1,e_2),\qquad k_2=(f_1,-f_2).
\]
Then
\[
\Pi_X(k_1)=\Pi_X(k_2)=(0,0,0),
\]
because \(e_1,e_2,f_1,f_2\) are isotropic and the two hyperbolic planes
are orthogonal.  Their sum is
\[
k_1+k_2=(e_1+f_1,e_2-f_2).
\]
The first component has square \(2\), the second has square \(-2\), and
the two components are orthogonal.  Therefore
\[
\Pi_X(k_1+k_2)=(1,0,-1)\ne(0,0,0).
\]
Thus \(K_{\Pi,0}\) is a null-cone fibre of a quadratic map, not a
subgroup.
\end{proof}

\begin{remark}[Fibre summation does not prove a Hall radical]
\label{rem:fibre-summation-not-radical}
Definition~\ref{def:gram-fibres-zero-fibre} proves only the formal
bookkeeping required before a finite Hall window can be compared with a
target root window.  A fibre with one even and one odd source vector has
signed superdimension \(0\) and parity profile \((1,1)\); it is not the
zero source degree.  The radical along \(K_{\Pi,0}\) is therefore a
separate Hopf-radical datum inside a finite source: it must specify the
pairing matrix, kernel matrix, quotient splitting, ideal/coideal
identities, and transition exactness.  It is not obtained by quotienting
the formal charge lattice by \(K_{\Pi,0}\).
\end{remark}

The packet \(\texttt{certificates/charge/gram\_fibre\_kernel}\) checks
these formal assertions on finite rows.  It records two distinct
formal charges with Gram label \((1,1,0)\), two with Gram label
\((1,1,1)\), and two nonzero charges in \(K_{\Pi,0}\) whose parity
profile is \((1,1)\) although the signed superdimension is \(0\).  It
also records the witness
\[
(e_1,e_2)+(f_1,-f_2)=(e_1+f_1,e_2-f_2),\qquad
\Pi_X(e_1+f_1,e_2-f_2)=(1,0,-1),
\]
so the zero fibre is not additive.  The packet does not prove finite
HN boundedness, compact Hall support, source Hall brackets, source
parity, a Hopf radical, exactness of a finite pushforward, Pfaffian
orientations, or protected traces.

\begin{definition}[Zero-fibre radical on a finite formal source]
\label{def:zero-fibre-radical-finite-source}
Let \(V=\bigoplus_{c\in A}V_c\) be a finite
\(\Gamma_X^{\mathrm{form}}\)-graded \(\Q\)-vector space with
\(\Z/2\)-parity, where \(A\subset\Gamma_X^{\mathrm{form}}\) is finite.
Let
\[
V_{\Pi,0}:=\bigoplus_{c\in A\cap K_{\Pi,0}}V_c.
\]
For a bilinear pairing \(h:V_{\Pi,0}\otimes V_{\Pi,0}\to\Q\), the
\emph{zero-fibre radical} is the ordinary linear-algebra kernel
\[
\operatorname{Rad}_{\Pi,0}(V,h)
:=\ker\!\left(V_{\Pi,0}\longrightarrow V_{\Pi,0}^{\vee}\right),
\qquad
v\longmapsto h(v,-).
\]
It is a subspace of the finite source vector space, not a quotient of
\(\Gamma_X^{\mathrm{form}}\) by the zero Gram fibre.
\end{definition}

\begin{proposition}[Exactness of finite fibre-summed Gram pushforward]
\label{prop:finite-fibre-pushforward-exact}
\textnormal{[proved]}
The functor
\[
\Pi_{X,*}^{\mathrm{fib}}:
\mathrm{Vect}^{\mathrm{fin}}_{\Gamma_X^{\mathrm{form}},\Z/2}
\longrightarrow
\mathrm{Vect}^{\mathrm{fin}}_{\Gamma_{\mathrm{gram}},\Z/2},
\qquad
(\Pi_{X,*}^{\mathrm{fib}}V)_{\gamma,\bar p}
=\bigoplus_{\Pi_X(c)=\gamma}V_{c,\bar p},
\]
is exact on finite-support graded \(\Q\)-vector spaces.  Consequently a
degreewise exact sequence
\[
0\longrightarrow A\longrightarrow B\longrightarrow C\longrightarrow0
\]
of finite formal sources remains exact after applying
\(\Pi_{X,*}^{\mathrm{fib}}\).
\end{proposition}

\begin{proof}
Fix \((\gamma,\bar p)\).  The \((\gamma,\bar p)\)-component of
\(\Pi_{X,*}^{\mathrm{fib}}\) is the finite direct sum of the
\((c,\bar p)\)-components over the fibre \(\Pi_X^{-1}(\gamma)\cap A\).
Finite direct sums are exact in the category of \(\Q\)-vector spaces.
Taking this finite direct sum in every Gram/parity component gives the
displayed exact sequence after pushforward.
\end{proof}

The packet
\(\texttt{certificates/charge/fibre\_pushforward\_exactness}\) checks
this finite linear algebra on six displayed Gram/parity blocks.  It
records matrices for \(0\to A\to B\to C\to0\), verifies that the
projection kernel equals the injection image in each block, verifies
the same equality after fibre-summed pushforward, and records a
zero-fibre pairing matrix whose radical is one-dimensional.  The packet
does not prove finite HN boundedness, compact Hall support, source Hall
brackets, Hopf ideal/coideal closure of the radical, exact Hall
pushforward, bar-construction commutation, Pfaffian orientations, or
protected traces.

\begin{proposition}[Finite normal-ordered pushforward preserves parity]
\label{prop:finite-parity-pushforward}
\textnormal{[proved]}
Let
\[
V=\bigoplus_{\hat\gamma\in A}
\left(V_{\hat\gamma,\bar0}\oplus V_{\hat\gamma,\bar1}\right)
\]
be a finite-support \(\widehat\Gamma_X\)-graded super vector space, and
let
\(\overline\Pi_X:\widehat\Gamma_X\to\Gamma_{\mathrm{gram}}\) be the
additive normal-ordered Gram map.  Define the pushed super vector space
by
\[
(\overline\Pi_{X,*}^{\mathrm{fib}}V)_{\gamma,\bar p}
:=
\bigoplus_{\substack{\hat\gamma\in A\\
\overline\Pi_X(\hat\gamma)=\gamma}}
V_{\hat\gamma,\bar p},
\qquad \bar p\in\Z/2.
\]
Then \(\overline\Pi_{X,*}^{\mathrm{fib}}\) preserves the even and odd
summands separately.  In particular the pair
\[
\bigl(\dim(\overline\Pi_{X,*}^{\mathrm{fib}}V)_{\gamma,\bar0},
\dim(\overline\Pi_{X,*}^{\mathrm{fib}}V)_{\gamma,\bar1}\bigr)
\]
is the fibrewise sum of the source parity dimensions, while the signed
superdimension is only its scalar shadow.
\end{proposition}

\begin{proof}
For fixed \((\gamma,\bar p)\), the displayed definition is an ordinary
finite direct sum of precisely the \(\bar p\)-parity components whose
normal-ordered Gram degree is \(\gamma\).  No even summand is placed in
an odd summand and no odd summand is placed in an even summand.  Thus
the pushed even block is the direct sum of source even blocks over the
fibre of \(\overline\Pi_X\), and the pushed odd block is the direct sum
of source odd blocks over the same fibre.  Taking the difference of the
two dimensions is a later scalar trace operation and does not recover
the two parity ranks.
\end{proof}

The packet \(\texttt{certificates/charge/parity\_pushforward}\) checks
this finite statement on three displayed Gram degrees.  It records the
two parity dimensions before and after pushforward and includes a
zero-degree row with parity profile \((1,1)\) and signed
superdimension \(0\).  It proves preservation of finite parity blocks
under the additive normal-ordered pushforward.  It does not prove the
source parity theorem, the Gritsenko--Nikulin/Kac target parity
equality, compact source representatives, primitive recognition,
Pfaffian orientations, or protected traces.

\begin{proposition}[Finite bar construction commutes with the
normal-ordered Gram pushforward]
\label{prop:finite-bar-pushforward-commutes}
\textnormal{[proved]}
Let \(A=\Q\mathbf 1\oplus\bar A\) be a finite-support augmented
associative algebra graded by the additive group
\(\widehat\Gamma_X\), with multiplication preserving the
\(\widehat\Gamma_X\)-degree.  Let
\[
T^c_{\le N}(\bar A)=
\bigoplus_{1\le p\le N}\bar A^{\otimes p}
\]
be the finite reduced tensor coalgebra with deconcatenation coproduct,
graded by the sum of the \(\widehat\Gamma_X\)-degrees of the letters.
Then the additive homomorphism
\(\overline\Pi_X:\widehat\Gamma_X\to\Gamma_{\mathrm{gram}}\) induces a
canonical isomorphism of finite
\(\Gamma_{\mathrm{gram}}\)-graded coalgebras
\[
\overline\Pi_{X,*}^{\mathrm{fib}}\,T^c_{\le N}(\bar A)
\;\cong\;
T^c_{\le N}\!\left(\overline\Pi_{X,*}^{\mathrm{fib}}\bar A\right).
\]
If the bar differential is defined by multiplication in \(A\), the
same isomorphism intertwines the finite bar differential.
\end{proposition}

\begin{proof}
For a word \([a_1|\cdots|a_p]\), the source degree is
\(\hat\gamma_1+\cdots+\hat\gamma_p\) in the additive group
\(\widehat\Gamma_X\).  Since \(\overline\Pi_X\) is a group homomorphism,
\[
\overline\Pi_X(\hat\gamma_1+\cdots+\hat\gamma_p)
=\overline\Pi_X(\hat\gamma_1)+\cdots+\overline\Pi_X(\hat\gamma_p).
\]
Thus pushing the word degree forward after forming the tensor coalgebra
and forming the tensor coalgebra after pushing the letters forward give
the same finite direct sum in each
\(\Gamma_{\mathrm{gram}}\)-degree.  Deconcatenation cuts a word into
two subwords; the displayed equality is additive on the two subword
degrees, so the coproduct is preserved.  The internal bar differential
replaces adjacent letters \(a_i,a_{i+1}\) by their product.  Degree
preservation of multiplication in \(A\) gives
\(\deg(a_i a_{i+1})=\deg(a_i)+\deg(a_{i+1})\), hence the same
\(\Gamma_{\mathrm{gram}}\)-degree after \(\overline\Pi_X\).  Therefore
the finite bar differential is intertwined.
\end{proof}

The packet
\(\texttt{certificates/charge/bar\_pushforward\_commutation}\) checks
this statement on finite words of length at most three.  It records
letter degrees, word degrees, deconcatenation splits, and product terms
in the bar differential, all using the additive normal-ordered Gram
grading.  It does not certify the raw quadratic \(\Pi_X\)-pushforward,
the chiral Koszul quasi-isomorphism, Hall bracket compatibility, Hopf
pairing compatibility, Pfaffian orientations, or protected traces.

\begin{proposition}[Finite bracket and pairing degrees descend under
normal-ordered Gram pushforward]
\label{prop:finite-bracket-pairing-pushforward}
\textnormal{[proved]}
Let \(P=\bigoplus_{\hat\gamma\in\widehat\Gamma_X}P_{\hat\gamma}\) be a
finite-support \(\widehat\Gamma_X\)-graded Lie super vector space with
\[
[P_{\hat\gamma},P_{\hat\delta}]
\subset P_{\hat\gamma+\hat\delta}.
\]
Then \(\overline\Pi_{X,*}^{\mathrm{fib}}P\) is a
\(\Gamma_{\mathrm{gram}}\)-graded Lie super vector space.  If
\[
h:P\otimes P\to\Q
\]
is a homogeneous degree-zero bilinear form, in the sense that
\[
h(P_{\hat\gamma},P_{\hat\delta})=0
\quad\text{unless}\quad \hat\gamma+\hat\delta=0,
\]
then the fibre-summed form on
\(\overline\Pi_{X,*}^{\mathrm{fib}}P\) is homogeneous of degree zero
for the \(\Gamma_{\mathrm{gram}}\)-grading.
\end{proposition}

\begin{proof}
For \(x\in P_{\hat\gamma}\) and \(y\in P_{\hat\delta}\), the source
bracket has degree \(\hat\gamma+\hat\delta\).  Applying the additive
homomorphism \(\overline\Pi_X\) gives
\[
\overline\Pi_X(\hat\gamma+\hat\delta)
=\overline\Pi_X(\hat\gamma)+\overline\Pi_X(\hat\delta),
\]
so the same bracket is graded after fibre-summed pushforward.  The
Jacobi identity and super-skew symmetry are unchanged, since the
underlying bracket is the same linear map on the same finite direct
sum.

If \(h(x,y)\ne0\), then \(\hat\gamma+\hat\delta=0\).  Applying
\(\overline\Pi_X\) gives
\(\overline\Pi_X(\hat\gamma)+\overline\Pi_X(\hat\delta)=0\).  Thus the
pushed pairing can be nonzero only on opposite Gram degrees.  This is
degree-zero homogeneity of the pushed form.  No Hopf adjointness,
Frobenius cyclic identity, quotient non-degeneracy, or primitive
closure is used.
\end{proof}

The packet
\(\texttt{certificates/charge/hall\_pairing\_pushforward\_compatibility}\)
checks this finite algebraic statement on supplied bracket and pairing
rows.  It verifies that three bracket rows have additive
normal-ordered Gram degree and that three nonzero pairing rows have
total Gram degree \(0\).  It does not construct compact Hall
correspondences, prove primitive closure, prove Hopf adjointness, prove
the Frobenius cyclic identity, prove quotient non-degeneracy, prove
Pfaffian orientations, or prove protected traces.

\begin{proposition}[Finite normal-ordered pushforward preserves the PBW filtration]
\label{prop:finite-pbw-pushforward}
\textnormal{[proved]}
Let \(U=\bigoplus_{\hat\gamma\in A}U_{\hat\gamma}\) be a
finite-support \(\widehat\Gamma_X\)-graded vector space equipped with
an increasing PBW filtration \(F_{\le d}U\) by homogeneous subspaces:
\[
F_{\le d}U
=\bigoplus_{\hat\gamma\in A}
\bigl(F_{\le d}U\cap U_{\hat\gamma}\bigr).
\]
Write \(F_{\le d}U_{\hat\gamma}:=F_{\le d}U\cap U_{\hat\gamma}\).
Define
\[
F_{\le d}
(\overline\Pi_{X,*}^{\mathrm{fib}}U)_{\gamma}
:=
\bigoplus_{\substack{\hat\gamma\in A\\
\overline\Pi_X(\hat\gamma)=\gamma}}
F_{\le d}U_{\hat\gamma}.
\]
Then \(\overline\Pi_{X,*}^{\mathrm{fib}}\) preserves the PBW
filtration, and for every \(d\)
\[
\operatorname{gr}^{\mathrm{PBW}}_d
(\overline\Pi_{X,*}^{\mathrm{fib}}U)_{\gamma}
\cong
\bigoplus_{\substack{\hat\gamma\in A\\
\overline\Pi_X(\hat\gamma)=\gamma}}
\operatorname{gr}^{\mathrm{PBW}}_d U_{\hat\gamma}.
\]
If \(U\) is a filtered algebra whose multiplication has
\(\widehat\Gamma_X\)-degree \(\hat\gamma+\hat\delta\) and PBW degree
\(\le a+b\) on \(F_{\le a}U_{\hat\gamma}\cdot
F_{\le b}U_{\hat\delta}\), the pushed multiplication has
\(\Gamma_{\mathrm{gram}}\)-degree
\(\overline\Pi_X(\hat\gamma)+\overline\Pi_X(\hat\delta)\) and PBW
degree \(\le a+b\).
\end{proposition}

\begin{proof}
For fixed \((\gamma,d)\), the displayed formula is a finite direct sum
over the fibre of \(\overline\Pi_X\).  Since
\(F_{\le d}U_{\hat\gamma}\subset F_{\le d+1}U_{\hat\gamma}\) for each
\(\hat\gamma\), the same inclusion holds after taking the finite direct
sum over the fibre.  Finite direct sums are exact over \(\Q\), hence
the quotient \(F_{\le d}/F_{\le d-1}\) commutes with the displayed
pushforward and gives the associated-graded isomorphism.  For the
algebra statement, the source multiplication sends
\((\hat\gamma,\hat\delta)\) to \(\hat\gamma+\hat\delta\) and has PBW
degree \(\le a+b\); applying the additive homomorphism
\(\overline\Pi_X\) sends the Gram degree to
\(\overline\Pi_X(\hat\gamma)+\overline\Pi_X(\hat\delta)\) and does not
change the filtration index.  Thus the pushed filtered multiplication
has the asserted degree and filtration bound.
\end{proof}

The packet \(\texttt{certificates/charge/pbw\_pushforward}\) checks
this finite filtered statement on three displayed Gram degrees.  It
records cumulative PBW dimensions, verifies monotonicity of the
filtration, and recomputes the associated-graded dimensions as
successive quotients.  It proves only preservation of a supplied finite
PBW filtration under the additive normal-ordered pushforward.  It does
not prove the compact-source PBW comparison, target PBW comparison,
no-extra-relations theorem, primitive recognition, Pfaffian
orientations, or protected traces.

\begin{proposition}[Finite normal-ordered pushforward preserves triangular decomposition]
\label{prop:finite-triangular-pushforward}
\textnormal{[proved]}
Let \(h:\Gamma_{\mathrm{gram}}\to\Z\) be an additive height functional.
For a finite-support \(\widehat\Gamma_X\)-graded vector space
\(V=\bigoplus_{\hat\gamma\in A}V_{\hat\gamma}\), define
\[
V^+=\bigoplus_{h(\overline\Pi_X(\hat\gamma))>0}V_{\hat\gamma},
\qquad
V^0=\bigoplus_{h(\overline\Pi_X(\hat\gamma))=0}V_{\hat\gamma},
\qquad
V^-=\bigoplus_{h(\overline\Pi_X(\hat\gamma))<0}V_{\hat\gamma}.
\]
Then
\[
\overline\Pi_{X,*}^{\mathrm{fib}}V
=
(\overline\Pi_{X,*}^{\mathrm{fib}}V)^-
\oplus
(\overline\Pi_{X,*}^{\mathrm{fib}}V)^0
\oplus
(\overline\Pi_{X,*}^{\mathrm{fib}}V)^+
\]
with each summand obtained by pushing forward the corresponding source
summand over the same fibres of \(\overline\Pi_X\).  If \(V\) carries a
homogeneous bracket
\[
[V_{\hat\gamma},V_{\hat\delta}]
\subset V_{\hat\gamma+\hat\delta},
\]
then the pushed bracket has height
\[
h\!\left(\overline\Pi_X(\hat\gamma+\hat\delta)\right)
=h(\overline\Pi_X(\hat\gamma))
+h(\overline\Pi_X(\hat\delta)).
\]
Thus positive-positive brackets remain positive, negative-negative
brackets remain negative, Cartan-positive and Cartan-negative brackets
remain in the corresponding sign, and positive-negative brackets land
in the sign determined by the actual height sum.
\end{proposition}

\begin{proof}
The three displayed subspaces form a direct sum because every integer
height is exactly one of positive, zero, or negative.  For a fixed
Gram degree \(\gamma\), the fibre-summed pushforward takes the finite
direct sum over \(\overline\Pi_X^{-1}(\gamma)\).  All source summands
in this fibre have the same value \(h(\gamma)\), hence the same
triangular sign.  Therefore no positive source summand is pushed into
the Cartan or negative part, and similarly for the other two signs.

For the bracket statement, source homogeneity gives degree
\(\hat\gamma+\hat\delta\).  Additivity of \(\overline\Pi_X\) and of
\(h\) gives the displayed height equality.  The triangular sign after
pushforward is therefore the sign of this integer.  The four stated
same-sign and Cartan-action cases are immediate; the
positive-negative case is not collapsed to a universal rule, but is
read from the actual height sum.
\end{proof}

The packet
\(\texttt{certificates/charge/triangular\_pushforward}\) checks this
finite statement on seven displayed source blocks and five supplied
bracket sign rows.  It verifies that the supplied integer height agrees
with the positive, Cartan, and negative labels, that the fibre-summed
pushforward keeps the three direct-sum parts separate, and that the
displayed bracket rows land in the triangular part determined by the
pushed height.  It does not prove exact triangular completion, target
triangular completion, compact-source primitive recognition,
no-extra-relations, PBW comparison, Pfaffian orientations, or protected
traces.

\begin{remark}[Quadratic obstruction; no linear trivialisation]
\label{rem:quadratic-obstruction}
Suppose $B = \delta L$ for some additive linear $L : \Gamma_X^{\mathrm
{form}}\to\Gamma_{\mathrm{gram}}$.  Setting $c' = c$, $B(c, c) = 2 L(c)
- L(2 c)$; additivity gives $L(2 c) = 2 L(c)$, so $B(c, c) = 0$,
contradicting $B(c, c) = 2\,\Pi_X(c) \ne 0$ for any $c$ with
$(Q, Q), (Q, P), (P, P)$ not all zero.  Hence no additive linear
trivialisation exists.  The quadratic primitive $\Pi_X$ is necessary;
trivializing the additive cocycle requires a quadratic cochain, not a
linear one (proved).
\end{remark}

\subsection{Normal-ordering flag formula}
\label{subsec:normal-ordering-flag-formula}

\begin{lemma}[Multiplication of raw placements]
\label{lem:multiplication-of-raw-placements}
For $c_1, \ldots, c_k \in \Gamma_X^{\mathrm{form}}$,
\[
i_0(c_1)\star\cdots\star i_0(c_k)
= \biggl(\sum_i c_i,\ \sum_{i < j} B(c_i, c_j)\biggr),
\]
and consequently
\[
\overline\Pi_X\!\left(i_0(c_1)\star\cdots\star i_0(c_k)\right)
= \sum_i \Pi_X(c_i).
\]
(proved.)
\end{lemma}

\begin{proof}
Induct on $k$.  The case $k = 2$ is the definition of $\star$.  At
step $k\to k + 1$,
\[
B\!\left(\sum_{i\le k} c_i,\,c_{k + 1}\right)
= \sum_{i\le k} B(c_i,\,c_{k + 1})
\]
by bilinearity.  Applying $\overline\Pi_X$ subtracts the accumulated
pairwise $B$-term from $\Pi_X(\sum_i c_i) = \sum_i \Pi_X(c_i) +
\sum_{i < j} B(c_i, c_j)$.
\end{proof}

\begin{remark}[Normal-ordered primitive recognition target]
\label{rem:normal-ordered-recognition-target}
The interface with the BKM denominator algebra is the normal-ordered
charge route
\[
\Gamma_X^{\mathrm{form}}\hookrightarrow\widehat\Gamma_X
\xrightarrow{\overline\Pi_X}\Gamma_{\mathrm{gram}}
\xrightarrow{\alpha}\Lambda_{II}^{2,1},
\]
where $\alpha(n, l, m) = 2 n f_2 - l f_3 + 2 m f_{-2}$ is the type-IV
embedding into the BKM root lattice (proved, \cite{GN}).  At
chain level this becomes the descent $\overline\Pi_{X, *}^{\Theta}$ of
Chapter~\ref{chap:dirac-igusa-pro-object}.  The primitive recognition
problem is to construct a protected primitive Hall algebra whose
normal-ordered descended bracket is the Borcherds--Kac bracket of
$\mathfrak g_{\Delta_5}$.  Raw Gram pushforward is not an admissible
recognition target.
\end{remark}

\section{The hybrid $\operatorname{Ran}^{\mathrm{hyb}}(E)$ factorization
carrier}
\label{sec:k3xe-hybrid-ran-carrier}

After integration along the K3 direction, the chiral object is supported
on the elliptic factor $E$.  The ordinary finite-support carrier is the Ran space
$\operatorname{Ran}(E)$ of finite subsets, but a one-dimensional
effective cycle with positive elliptic degree projects onto all of $E$.  The
hybrid carrier $\operatorname{Ran}^{\mathrm{hyb}}(E)$ is a finite-stage
prestack whose \(b=0\) local and \(b>0\) wrapped subprestacks are
separated before the mixed local/wrapped correspondences are supplied as
additional finite-stage data.
At finite stage those data are represented by the tables
\(\mathrm{hybrid\_ran\_prestack\_definition}\),
\(\mathrm{local\_stratum\_b0\_prestack\_definition}\),
\(\mathrm{wrapped\_stratum\_bpositive\_prestack\_definition}\),
\(\mathrm{geometric\_elliptic\_degree\_map}\),
\(\mathrm{geometric\_elliptic\_degree\_additivity}\),
\(\mathrm{geometric\_borcherds\_degree\_separation}\),
\(\mathrm{positive\_elliptic\_degree\_projection}\),
\(\mathrm{hybrid\_base}\), \(\mathrm{local\_configurations}\),
\(\mathrm{wrapped\_prequotients}\),
\(\mathrm{local\_local\_correspondence\_definition}\),
\(\mathrm{mixed\_local\_wrapped\_correspondence\_definition}\),
\(\mathrm{wrapped\_wrapped\_correspondence\_definition}\),
\(\mathrm{two\_sided\_mixed\_correspondence\_definition}\),
\(\mathrm{correspondences}\),
\(\mathrm{flag\_atlas}\), \(\mathrm{higher\_coloured}\),
\(\mathrm{quotient\_pseudofunctor}\), \(\mathrm{protected\_integration}\),
\(\mathrm{transitions}\), and \(\mathrm{scalar\_firewall}\).  They are
not supplied by the ordinary Ran space, a quotient-first Hall product,
or the Borcherds \(s\)-degree.

\subsection{Projection-locality obstruction}
\label{subsec:projection-locality-obstruction}

\begin{lemma}[Projection-locality obstruction]
\label{lem:projection-locality-obstruction}
Let
\[
Z=\sum_i m_i[C_i]\in Z_1(S\times E),\qquad m_i>0,
\]
be an effective one-cycle.  If \((p_E)_*Z=b[E]\) with \(b>0\), then
\(p_E(|Z|)=E\).  Hence \(|Z|\) is not local on
\(\operatorname{Ran}(E)\).
\end{lemma}

\begin{proof}
For each irreducible component \(C_i\), the proper image \(p_E(C_i)\)
is either a point or a closed one-dimensional subset of the irreducible
curve \(E\).  In the second case it is all of \(E\), and
\((p_E)_*[C_i]=\deg(C_i/E)[E]\).  Since the effective sum
\((p_E)_*Z=b[E]\) has \(b>0\), at least one component maps dominantly
to \(E\).  Therefore \(p_E(|Z|)=E\).
\end{proof}

\begin{remark}[Geometric versus Borcherds elliptic degree]
\label{rem:geometric-vs-borcherds-degree}
The lemma fixes the necessity of a wrapped stratum.  The geometric
elliptic-degree map $b_R^{\mathrm{geom}}$, on a finite Hall stage $R$,
records the integer $\deg_E$ of a chosen retained class.  The trace
exponent $m_R^{\mathrm{Bch}} := \mathrm{pr}_3\overline\Pi_X$, on the same stage,
records the third coordinate of the normal-ordered Gram pushforward.
The two are not equal as functions on the formal lattice; the Borcherds
exponent $m$ records the magnetic Mukai pairing $(P, P)/2$, which sees
$\Lambda_S$-square data, while $b_R^{\mathrm{geom}}$ records the
$E$-degree of the geometric support, which sees $H^0(S)\oplus H^4(S)$
data through $r$ and $\mathrm{ch}_2$ of the K3 factor.  The wrapped
stratum carries $b_R^{\mathrm{geom}} > 0$ branes; the local stratum
carries $b_R^{\mathrm{geom}} = 0$ branes whose magnetic component lies
in the projection-finite kernel.  This is the
$\operatorname{Ran}(E)$-support-locality obstruction in the
K3-integrated chiral shadow.
\end{remark}

\subsection{The hybrid carrier}
\label{subsec:hybrid-carrier-definition}

\begin{definition}[Hybrid Ran carrier]
\label{def:hybrid-ran-carrier}
At finite HN height \(R\), the \emph{hybrid Ran carrier} is the
prestack
\[
\operatorname{Ran}^{\mathrm{hyb}}_{R,\mathrm{pre}}(E):
\mathrm{Aff}_{\C}^{op}\to\infty\textup{-}\mathrm{Gpd}
\]
of Definition~\ref{def:hybrid-ran-prestack}.  A test family consists
of finitely many local maps to \(E\), retained local colour labels,
finitely many wrapped colour labels, and \(T\)-families in the
corresponding rigidified wrapped prequotient stacks with anchor maps to
\(E\).  Its local \(b=0\) stratum is the subprestack
\(\operatorname{Ran}^{\mathrm{loc}}_{R,\mathrm{pre}}(E)\) of
Definition~\ref{def:hybrid-local-stratum-prestack}.  Its wrapped
\(b>0\) stratum is the subprestack
\(\operatorname{Ran}^{\mathrm{wr}}_{R,\mathrm{pre}}(E)\) of
Definition~\ref{def:hybrid-wrapped-stratum-prestack}.  The stack
\(\operatorname{Ran}^{\mathrm{hyb}}_{R,\tau}(E)\) exists only after
descent rows for the chosen topology \(\tau\) have been supplied.  The
set
\[
\bigsqcup_{b\ge0}\operatorname{Ran}(E)\times\{b\}
\]
is the degree shadow, not the carrier.
\end{definition}

\begin{remark}[Two strata, two registers]
\label{rem:two-strata-two-registers}
On the $b = 0$ stratum the standard Ran/factorization formalism applies
and supplies the local chiral factorization input in
$\mathrm{Fact}(\operatorname{Ran}(E))$ (folklore,
\cite{BD,FrancisGaitsgoryCKD}).  On the $b > 0$
stratum the support of $\mathrm{Supp}\,\mathcal F$ projects onto $E$, so
the brane is global; its data live on the wrapped prequotient and after
$E$-quotient become $E$-equivariant data on the elliptic curve.  The
hybrid carrier is the smallest carrier reflecting both strata; it is
not a single Ran space.
\end{remark}

\subsection{Mixed local/wrapped correspondences}
\label{subsec:mixed-correspondences}

\begin{definition}[Mixed local/wrapped correspondences (finite Hall
stage $R$)]
\label{def:mixed-correspondences}
Fix a finite Hall stage $R$ and retained Liu--Hilbert classes $\Xi_A,
\Xi_B, \Xi_C$ with $\Xi_C = \Xi_A + \Xi_B$ in the closure rules of
\cite{KSCoHA}.  The \emph{retained extension correspondence}
\[
\mathfrak E^{\mathrm{ret}}_{\Xi_A, \Xi_B; \Xi_C}
\xrightarrow{p}
\mathfrak M^{\mathrm{ss}}_{\Xi_A}\times\mathfrak M^{\mathrm{ss}}_{\Xi_B},
\qquad
\mathfrak E^{\mathrm{ret}}_{\Xi_A, \Xi_B; \Xi_C}
\xrightarrow{q}
\mathfrak M^{\mathrm{ss}}_{\Xi_C}
\]
is the closed derived flag-Quot/filtration stack over
$\mathfrak M^{\mathrm{ss}}_{\Xi_C}$, with $q$ required at this finite
stage to be proper or $\mathrm{Coh}^{\mathrm{prop}}$-admissible in the
compactified model (folklore).  Three colours of
correspondence appear,
indexed by the location of $\Xi_A, \Xi_B$ in
$\operatorname{Ran}^{\mathrm{hyb}}(E)$:
\begin{enumerate}[label = \textup{(LL)},leftmargin = *]
\item[\textup{(LL)}] \emph{local-local}: $b_A = b_B = 0$;
\item[\textup{(LW)}] \emph{ordered local/wrapped}: $b_A = 0$, $b_B > 0$;
\item[\textup{(WW)}] \emph{wrapped-wrapped}: $b_A, b_B > 0$.
\end{enumerate}
The local-local row is the ordered projection-finite diagram
\[
\mathfrak E^{\mathrm{LL}}_{\alpha,\beta;\gamma,R,I,J,K}
\longrightarrow
\mathfrak M^{\mathrm{loc,cl}}_{\alpha,R,I}
\times
\mathfrak M^{\mathrm{loc,cl}}_{\beta,R,J}
\longrightarrow
\mathfrak M^{\mathrm{loc,cl}}_{\gamma,R,K},
\]
where \(\alpha:I\to\Gamma_R^{\mathrm{loc}}\),
\(\beta:J\to\Gamma_R^{\mathrm{loc}}\),
\(\gamma=|\alpha|+|\beta|\) is retained in
\(\Gamma_R^{\mathrm{loc}}\), and \(K=I\sqcup J\) modulo Ran
collisions.  This is
Definition~\ref{def:local-local-correspondence-datum}; it has no
wrapped input and no reduced \(E\)-quotient.
The one-sided mixed row is the ordered pair of unreduced diagrams
\[
\mathfrak E^{\mathrm{LW}}_{\alpha,\eta;\zeta,R,I}
\longrightarrow
\mathfrak M^{\mathrm{loc,cl}}_{\alpha,R,I}
\times
\mathcal M_{\eta,R}^{\mathrm{wr,rig}}
\longrightarrow
\mathcal M_{\zeta,R}^{\mathrm{wr,rig}},
\]
and
\[
\mathfrak E^{\mathrm{WL}}_{\eta,\alpha;\zeta,R,I}
\longrightarrow
\mathcal M_{\eta,R}^{\mathrm{wr,rig}}
\times
\mathfrak M^{\mathrm{loc,cl}}_{\alpha,R,I}
\longrightarrow
\mathcal M_{\zeta,R}^{\mathrm{wr,rig}},
\]
where \(\alpha:I\to\Gamma_R^{\mathrm{loc}}\),
\(\eta\in\Gamma_R^{\mathrm{wr}}\), and
\(\zeta=|\alpha|+\eta\) or \(\eta+|\alpha|\) is retained in
\(\Gamma_R^{\mathrm{wr}}\).  This is
Definition~\ref{def:mixed-local-wrapped-correspondence-datum}; it is
not the local-local row, not the wrapped-wrapped row, not the two-sided
mixed row, and not a quotient-first product.
The wrapped-wrapped row is the ordered unreduced diagram
\[
\mathfrak E^{\mathrm{WW}}_{\eta_1,\eta_2;\zeta,R}
\longrightarrow
\mathcal M_{\eta_1,R}^{\mathrm{wr,rig}}
\times
\mathcal M_{\eta_2,R}^{\mathrm{wr,rig}}
\longrightarrow
\mathcal M_{\zeta,R}^{\mathrm{wr,rig}},
\]
where \(\eta_1,\eta_2\in\Gamma_R^{\mathrm{wr}}\) and
\(\zeta=\eta_1+\eta_2\) is retained in
\(\Gamma_R^{\mathrm{wr}}\).  This is
Definition~\ref{def:wrapped-wrapped-correspondence-datum}; it is
ordered, unreduced, and separate from symmetric descent and
source/target-anchor memory.
The ordered two-sided mixed row is the unreduced local--wrapped--local
diagram
\[
\mathfrak E^{\mathrm{LWL}}_{\alpha_-;\eta;\beta_+,R,I_-,I_+}
\longrightarrow
\mathfrak M^{\mathrm{loc,cl}}_{\alpha_-,R,I_-}
\times
\mathcal M_{\eta,R}^{\mathrm{wr,rig}}
\times
\mathfrak M^{\mathrm{loc,cl}}_{\beta_+,R,I_+}
\longrightarrow
\mathcal M_{\zeta,R}^{\mathrm{wr,rig}},
\]
where
\(\zeta=|\alpha_-|+\eta+|\beta_+|\in\Gamma_R^{\mathrm{wr}}\).
	This is
	Definition~\ref{def:two-sided-mixed-correspondence-datum}; it is not
	the full eight-word associativity atlas and not a quotient-first
	product.
	The source/target anchor-memory row is separate: for \(LW\), \(WL\),
	\(WW\), and \(LWL\) it records the ordered wrapped source anchors and
	wrapped target anchor on the unreduced correspondence stack, before the
	reduced \(E\)-quotient, as in
	Definition~\ref{def:source-target-anchor-memory-datum}.  The
	local/local row has no wrapped anchor memory.  This row does not
	construct the determinant anchor, compute its translation weight,
	repair the \(\chi=0\) degeneracy, or define \(o^\lambda_R\).
	The arity-three shadow is the eight-word two-step binary flag atlas
	\[
	\{L, W\}^3 = \{LLL, LLW, LWL, LWW, WLL, WLW, WWL, WWW\},
	\]
which records binary associativity only; higher-coloured tree
configurations live on a separate coloured tree category
$\mathrm{Tree}^{\mathrm{col}}_R$, supplied as additional input
(folklore, \cite{CYQGVolIII}).
\end{definition}

\begin{remark}[Eight binary words versus colored trees]
\label{rem:eight-words-vs-colored-trees}
The eight binary words of Definition~\ref{def:mixed-correspondences}
exhaust the arity-three two-step compositions $\{L, W\}^3$, but they do
not catalogue higher-coloured tree configurations, units, symmetry/order
conventions, refinement maps, descent, or overlaps.  Those data are
additional finite Hall input named at \S\ref{sec:k3xe-four-obstructions};
asserting them prematurely confuses binary associativity with the full
factorization-coherence structure (\cite{CYQGVolIII},
correspondence formalism).
\end{remark}

\subsection{Wrapped anchor and translation linearization}
\label{subsec:e-quotient-anchor}

\begin{definition}[Determinant anchor on the wrapped stratum]
\label{def:determinant-anchor}
Fix a finite HN height \(R\), a wrapped colour
\(\eta\in\Gamma_R^{\mathrm{wr}}\), and a test scheme \(T\).  Let
\(\mathcal F_{\eta,T}\) be the retained universal perfect complex on
\(X\times T\) pulled back from the wrapped prequotient row, and write
\[
\pi_{E,T}:X\times T\longrightarrow E\times T .
\]
Let \(\chi_\eta=\chi(\mathcal F_{\eta,t})\), constant on the retained
colour row.  The Knudsen--Mumford determinant functor
\cite{KnudsenMumfordDet,DeligneDetCohomology} gives a line bundle
\[
\mathscr D_{\eta,T}
:=\det R\pi_{E,T,*}\mathcal F_{\eta,T}
\]
on \(E\times T\), functorial in \(T\) as a point of the relative
Picard functor.  Since \(E\) has genus \(1\),
\(\deg\det C=\chi(C)\) for a perfect complex \(C\) on \(E\).  Hence
\(\mathscr D_{\eta,T}\) has relative degree \(\chi_\eta\), and
\[
\lambda^{\det}_{\eta,R}(\mathcal F_{\eta,T})
:=\mathscr D_{\eta,T}\otimes
\operatorname{pr}_E^*\mathcal O_E(-\chi_\eta\cdot0_E)
\]
where \(\operatorname{pr}_E:E\times T\to E\) is projection.  This line
bundle has relative degree \(0\).  Its class defines the
determinant anchor
\[
\lambda^{\det}_{\eta,R}:T\longrightarrow
\operatorname{Pic}^0(E)\simeq E .
\]
This constructs only the normalized determinant anchor.  The next
propositions compute its \(E\)-translation weight and isolate the
\(\chi_\eta=0\) degeneracy.  The finite-cover normalization to unit
weight, extra anchor data, anchor residual definition and vanishing,
quotient descent, and transition compatibility are separate rows.
\end{definition}

The packet
\(\texttt{certificates/hybrid/determinant\_anchor\_construction}\)
verifies
\(\texttt{DETERMINANT\_ANCHOR\_CONSTRUCTION\_VERIFIED}\): the
determinant-of-cohomology line is functorial in families, the
normalization has relative degree zero, and the resulting map lands in
\(\operatorname{Pic}^0(E)\simeq E\).  It does not compute the
translation weight, handle the \(\chi=0\) degeneracy, add extra anchor
data, define \(o^\lambda_R\), or prove quotient descent.

\begin{proposition}[Translation weight of the determinant anchor]
\label{prop:determinant-anchor-translation-weight}
\textnormal{[proved]}
Let \(a:T\to E\) be a section and let
\(\tau_a:E\times T\to E\times T\) be translation by \(a\).  Use the
support-translation convention
\[
a\cdot \mathcal F_{\eta,T}
:=(\operatorname{id}_S\times\tau_a)_*\mathcal F_{\eta,T}
\]
on \(X=S\times E\).  If \(\Gamma_a,\Gamma_0\subset E\times T\) are
the graphs of \(a\) and \(0_E\), then in the relative Picard functor
\[
\lambda^{\det}_{\eta,R}(a\cdot\mathcal F_{\eta,T})
\simeq
\lambda^{\det}_{\eta,R}(\mathcal F_{\eta,T})\otimes
\mathcal O_{E\times T}\!\bigl(\chi_\eta(\Gamma_a-\Gamma_0)\bigr).
\]
Under \(\operatorname{Pic}^0(E)\simeq E\), the induced translation law is
\[
\lambda^{\det}_{\eta,R}(a\cdot\mathcal F_{\eta,T})
=
\lambda^{\det}_{\eta,R}(\mathcal F_{\eta,T})+\chi_\eta a .
\]
Thus the \(E\)-translation weight of the determinant anchor is
\(\chi_\eta\).  Its kernel is \(E[\chi_\eta]\) when \(\chi_\eta\ne0\),
and all of \(E\) when \(\chi_\eta=0\).
\end{proposition}

\begin{proof}
Set
\(\mathscr D_{\eta,T}=\det R\pi_{E,T,*}\mathcal F_{\eta,T}\).  The
Cartesian square for translation on the elliptic factor gives, in the
relative Picard group,
\[
\det R\pi_{E,T,*}(a\cdot\mathcal F_{\eta,T})
\simeq
\tau_{a,*}\mathscr D_{\eta,T}
\simeq
\tau_{-a}^{*}\mathscr D_{\eta,T}.
\]
For a relative line bundle \(L\) of degree \(d\) on \(E\times T\),
translation acts trivially on the \(\operatorname{Pic}^0(E)\)-class and
shifts the degree part by
\[
\tau_{-a}^{*}L\otimes L^{-1}
\simeq
\mathcal O_{E\times T}\!\bigl(d(\Gamma_a-\Gamma_0)\bigr)
\quad
\text{in }\operatorname{Pic}^0(E\times T/T).
\]
Taking \(L=\mathscr D_{\eta,T}\) and \(d=\chi_\eta\), and then tensoring
both sides with the fixed normalization
\(\mathcal O_{E\times T}(-\chi_\eta\Gamma_0)\), gives the displayed
identity.  The sign is positive because the manuscript uses pushforward
by the support translation \(x\mapsto x+a\); the pullback convention
would reverse the sign.  The kernel statement is the kernel of the
multiplication homomorphism \([\chi_\eta]:E\to E\).
\end{proof}

The packet
\(\texttt{certificates/hybrid/determinant\_anchor\_translation\_weight}\)
verifies
\(\texttt{DETERMINANT\_ANCHOR\_TRANSLATION\_WEIGHT\_VERIFIED}\):
for the support-translation convention
\(a\cdot\mathcal F=(\operatorname{id}_S\times\tau_a)_*\mathcal F\),
the determinant anchor has translation weight \(\chi_\eta\), with
restriction \(h\mapsto\chi_\eta h\) on any finite stabilizer
\(H\subset E[N]\).  It does not choose a coherent stabilizer
linearization, make the weight one by finite cover, handle the
\(\chi_\eta=0\) degeneracy, add extra anchor data, define
\(o^\lambda_R\), prove quotient descent, or prove transition
compatibility.

\begin{proposition}[Zero-Euler determinant-anchor degeneracy]
\label{prop:determinant-anchor-chi-zero-degeneracy}
\textnormal{[proved]}
Assume \(\chi_\eta=0\).  Then for every section \(a:T\to E\),
\[
\lambda^{\det}_{\eta,R}(a\cdot\mathcal F_{\eta,T})
=
\lambda^{\det}_{\eta,R}(\mathcal F_{\eta,T})
\quad\text{in }\operatorname{Pic}^0(E).
\]
On a retained finite row where the \(E\)-translation action is free, a
geometric point \(a\ne0\) gives distinct objects
\(\mathcal F_{\eta,t}\) and \(a\cdot\mathcal F_{\eta,t}\) with the same
determinant anchor.  Hence the determinant anchor has full
\(E\)-orbit kernel on the \(\chi_\eta=0\) branch and has
one-dimensional separation defect.  A final wrapped anchor on this
branch must contain extra anchor data not determined by
\(\lambda^{\det}_{\eta,R}\).
\end{proposition}

\begin{proof}
Proposition~\ref{prop:determinant-anchor-translation-weight} gives
\[
\lambda^{\det}_{\eta,R}(a\cdot\mathcal F_{\eta,T})
=
\lambda^{\det}_{\eta,R}(\mathcal F_{\eta,T})+\chi_\eta a .
\]
When \(\chi_\eta=0\) this is the displayed equality for all
\(a\in E\).  The retained \(E\)-translation-rigidification datum is
free on the retained row by
Proposition~\ref{prop:retained-e-translation-rigidifications}; hence a
nonzero geometric translation gives a different point in the
prequotient before the reduced quotient is taken.  The two points have
the same determinant anchor.  Therefore determinant-anchor-only data
does not separate the \(E\)-orbit.  The orbit is one-dimensional, so
the class-level separation defect has rank \(1\).  Extra anchor data is
therefore a necessary input for the \(\chi_\eta=0\) branch, not a
consequence of determinant of cohomology.
\end{proof}

The packet
\(\texttt{certificates/hybrid/determinant\_anchor\_chi\_zero\_degeneracy}\)
verifies
\(\texttt{DETERMINANT\_ANCHOR\_CHI\_ZERO\_DEGENERACY\_VERIFIED}\):
on the branch \(\chi_\eta=0\), the determinant anchor is constant on
the full \(E\)-translation orbit, and freeness of the retained
translation action gives a rank-one separation defect for
determinant-anchor-only data.  It records that extra anchor data is
mandatory.  It does not itself construct that data, define
\(o^\lambda_R\), prove losslessness, prove quotient descent, choose
stabilizer linearizations, or prove transition compatibility.

\begin{proposition}[Extra anchor data from the degree-zero Jordan--Hoelder divisor]
\label{prop:determinant-anchor-extra-data-separates-rank-two}
\textnormal{[proved]}
Let \(V\) be a semistable vector bundle of rank \(r\) and degree \(0\)
on \(E\).  Write the associated graded object of a Jordan--Hoelder
filtration as
\[
\operatorname{gr}^{\mathrm{JH}}V=L_1\oplus\cdots\oplus L_r,
\qquad
L_i\in\operatorname{Pic}^0(E),
\]
and set
\[
\operatorname{sp}^{\mathrm{JH}}(V)
=[L_1]+\cdots+[L_r]\in
\operatorname{Sym}^r\operatorname{Pic}^0(E).
\]
This divisor is independent of the Jordan--Hoelder filtration.  Its
Abel sum is \(\det V\).  For
\[
V_0=\mathcal O_E^{\oplus2},\qquad
V_L=L\oplus L^{-1},\qquad L\ne\mathcal O_E,
\]
one has
\[
\det V_0\simeq\det V_L\simeq\mathcal O_E,\qquad
\operatorname{sp}^{\mathrm{JH}}(V_0)=2[\mathcal O_E],\qquad
\operatorname{sp}^{\mathrm{JH}}(V_L)=[L]+[L^{-1}],
\]
and the two degree-two divisors are unequal in
\(\operatorname{Sym}^2\operatorname{Pic}^0(E)\).
\end{proposition}

\begin{proof}
By Atiyah's classification and Tu's S-equivalence description for
semistable bundles on an elliptic curve
\cite{AtiyahVectorBundlesElliptic,TuSemistableElliptic}, every stable
Jordan--Hoelder factor of a degree-zero semistable vector bundle on
\(E\) is a degree-zero line bundle, and the unordered collection of
these factors is the S-equivalence class of \(V\).  Hence
\(\operatorname{sp}^{\mathrm{JH}}(V)\) is independent of the chosen
Jordan--Hoelder filtration.  The determinant of the associated graded
object is
\[
\det\operatorname{gr}^{\mathrm{JH}}(V)=L_1\otimes\cdots\otimes L_r,
\]
which equals \(\det V\); under the group law on
\(\operatorname{Pic}^0(E)\), this is precisely the Abel sum of the
divisor \(\operatorname{sp}^{\mathrm{JH}}(V)\).  The displayed
rank-two pair has determinant \(\mathcal O_E\) on both sides.  If
\(2[\mathcal O_E]=[L]+[L^{-1}]\) in
\(\operatorname{Sym}^2\operatorname{Pic}^0(E)\), then the unordered
multisets \(\{\mathcal O_E,\mathcal O_E\}\) and
\(\{L,L^{-1}\}\) are equal, so \(L\simeq\mathcal O_E\), contrary to the
hypothesis.  Thus the determinant anchor agrees while the
Jordan--Hoelder divisor separates the two objects.
\end{proof}

The packet
\(\texttt{certificates/hybrid/determinant\_anchor\_extra\_anchor\_data}\)
verifies
\(\texttt{DETERMINANT\_ANCHOR\_EXTRA\_ANCHOR\_DATA\_VERIFIED}\):
on a row where the elliptic pushforward is a degree-zero semistable
vector bundle, the additional anchor datum is the
Jordan--Hoelder S-equivalence divisor in
\(\operatorname{Sym}^r\operatorname{Pic}^0(E)\).  Its Abel sum is the
determinant anchor, and it separates the determinant-zero rank-two test
pair
\(\mathcal O_E^{\oplus2}\) and \(L\oplus L^{-1}\) with
\(L\ne\mathcal O_E\).  It does not prove global anchor losslessness,
prove \(o^\lambda_R=0\), prove quotient descent, choose stabilizer
linearizations, populate every finite wrapped prequotient, or prove
transition compatibility.

\begin{lemma}[Connected $E$-descent: separated cohomological classes]
\label{lem:connected-e-descent}
For the connected component $E^{\circ}$ of the identity in the
$E$-translation action,
\[
H^*(BE; \mathbb F_2) = \mathbb F_2[u_1, u_2],\qquad
|u_i| = 2.
\]
Hence
\[
H^1(BE; \mathbb F_2) = 0,\qquad
H^2(BE; \mathbb F_2) = \mathbb F_2 u_1 \oplus \mathbb F_2 u_2.
\]
There is no degree-one connected character; there is a degree-two
connected gerbe class
\[
\alpha^{E,\mathrm{free}} = a_1 u_1 + a_2 u_2 \in H^2(BE; \mathbb F_2).
\]
(proved, \cite{NIKULIN}.)
\end{lemma}

\begin{proof}
$E^{\mathrm{top}}\simeq T^2$ has classifying space $BE\simeq BT^2$ and
$H^*(BT^2; \mathbb F_2) = \mathbb F_2[u_1, u_2]$ with both $u_i$ of
degree $2$ (mod-$2$ Chern classes of the universal line bundles).
$H^1$ vanishes; $H^2$ is the displayed two-dimensional $\mathbb F_2$
vector space.  Vanishing of the connected gerbe class is the equation
$a_1 = a_2 = 0$ on the actual Borel/edge class of the equivariant
reduced determinant complex, not a consequence of $H^1(BE) = 0$ or
ordinary translation invariance.
\end{proof}

\begin{lemma}[{Finite stabilizer $E[N]$-descent: gerbe and linearization classes}]
\label{lem:finite-stabilizer-e-descent}
For the order-two finite subgroup $E[2]\cong(\Z/2)^2$,
\[
H^*(BE[2]; \mathbb F_2) = \mathbb F_2[x_1, x_2],\qquad
|x_i| = 1.
\]
The degree-two gerbe class is
\[
\beta^{E, E[2]}
= b_{20} x_1^2 + b_{11} x_1 x_2 + b_{02} x_2^2,\qquad
(b_{20}, b_{11}, b_{02}) = (r_1,\,r_1 + r_2 + r_3,\,r_2),
\]
detected on cyclic subgroup restrictions by joint reduction; the $r_i
\in\mathbb F_2$ are the three nonzero point-stabilizer parities.
After $\beta^{E, E[2]} = 0$ (i.e.\ $r_1 = r_2 = r_3 = 0$), a separate
degree-one linearization class remains:
\[
\lambda^{E, E[2]} = \lambda_1 x_1 + \lambda_2 x_2 \in H^1(BE[2]; \mathbb F_2),
\qquad \rho_1 = \lambda_1,\ \rho_2 = \lambda_2,\ \rho_3 = \lambda_1 + \lambda_2.
\]
The $r_i$ are gerbe bits in $H^2$; the $\rho_i$ are linearization bits
in $H^1$; the two live in distinct cohomological degrees and remain
independent.
For even $N\ge 4$ with $2^a\parallel N,\ a\ge 2$,
\[
H^*\!\bigl(B(\Z/2^a)^2; \mathbb F_2\bigr)
= \Lambda(x_1, x_2)\otimes\mathbb F_2[y_1, y_2],
\]
\[
\beta^{E, N} = A_1^{(N)} y_1 + A_{12}^{(N)} x_1 x_2 + A_2^{(N)} y_2,
\qquad
\lambda^{E, N} = \lambda_1^{(N)} x_1 + \lambda_2^{(N)} x_2.
\]
The mixed term $A_{12}^{(N)}$ is invisible to cyclic order-two restrictions;
even $N\ge 4$ tests must be checked on the full two-primary generators.
For odd $N$, the mod-2 orientation obstruction vanishes by transfer.
(proved.)
\end{lemma}

\begin{proof}
$E[2]\cong(\Z/2)^2$ has $H^*(B(\Z/2)^2; \mathbb F_2) = \mathbb F_2[x_1,
x_2]$ with $|x_i| = 1$.  The class
\(\beta^{E,E[2]}\) is the finite-stabilizer Cech--Borel edge class,
not the restriction of the connected \(BE\)-class
\(\alpha^{E,\mathrm{free}}\).  Its restriction to the three non-zero
cyclic subgroups gives the three parities \(r_i\); inversion of this
three-by-three system gives
\((b_{20},b_{11},b_{02})=(r_1,r_1+r_2+r_3,r_2)\).
The degree-one linearization
class $\lambda^{E, E[2]}$ is the residual character of the underlying
line bundle after $\beta = 0$ is imposed; it lives in $H^1$ and is
distinct from the gerbe $\beta\in H^2$ in degree.  Even $N\ge 4$
introduces the polynomial generators $y_i \in H^2(B(\Z/2^a))$ and the
mixed term $x_1 x_2$, which on cyclic restrictions vanishes; full
two-primary detection is needed.  Odd $N$ contributes only $\mathbb
F_p$-cohomology with $p$ odd, on which the mod-$2$ orientation has zero
restriction by the transfer.
\end{proof}

\begin{remark}[Translation linearization criterion, not vanishing]
\label{rem:translation-linearization-criterion}
For a stabilizer $H \subset E[N]$ and a determinant anchor line bundle
$\lambda(\mathcal F)$, ordinary translation invariance $g^*\lambda
\simeq \lambda$ in $\mathrm{Pic}^0(E)$ does not choose coherent
isomorphisms $\phi_g : \lambda \to g^*\lambda$.  After the square-root
gerbe is killed, the choices form a character torsor in $H^1(BH;
\mathbb F_2)$, and quotient orientation requires the linearization
class
\[
\lambda^H \in H^1(BH; \mathbb F_2)
\]
to vanish.  This is a \emph{criterion} for descent on the wrapped
stratum; it is not implied by $g^*[\lambda] = [\lambda]$ in
$\mathrm{Pic}(E)$.  Proposition~\ref{prop:determinant-anchor-translation-weight}
computes the translation law for \(\lambda^{\det}\) and its weight
\(\chi(\mathcal F)\).  The finite-cover normalisation to unit weight
when \(\chi(\mathcal F)\ne0\) and the coherent stabilizer
linearization are later anchor rows.  The \(\chi(\mathcal F)=0\)
determinant-anchor degeneracy is
Proposition~\ref{prop:determinant-anchor-chi-zero-degeneracy}; the
degree-zero Jordan--Hoelder divisor of
Proposition~\ref{prop:determinant-anchor-extra-data-separates-rank-two}
separates the determinant-zero rank-two test pair.  Full anchor
losslessness, finite-stage population, quotient descent, transition
compatibility, and the coherent stabilizer linearization remain later
rows.  Until those rows are supplied, this paragraph is only the
descent criterion for a chosen linearized anchor.
(proved as a criterion; determinant-anchor weight proved in
Proposition~\ref{prop:determinant-anchor-translation-weight}.)
\end{remark}

\begin{remark}[Mod-2 character does not control odd-order anchor
characters]
\label{rem:mod2-not-odd-order}
The mod-$2$ orientation character does not by itself control odd-order
anchor characters.  For anchor descent on the wrapped stratum, also
compute the actual stabilizer character on the chosen anchor
trivialization, valued in $\widehat H = \mathrm{Hom}(H, \C^\times)$, or
define the retained stabilizer as preserving that trivialization.
(unverified; supplied as part of the finite Hall datum.)
\end{remark}

\subsection{Cusp polarization and chamber transport}
\label{subsec:cusp-polarization-chamber}

\begin{definition}[Cusp polarization on $\Gamma_{\mathrm{eff}}$]
\label{def:cusp-polarization}
The chamber convention $0 < |s|\ll|q|\ll|r|^{-1}\ll 1$ at the standard
cusp $\mathfrak c_\infty$ determines the cusp-positive semigroup
$\Gamma_{\mathrm{eff}}$ of Definition~\ref{def:gram-index-lattice}.  The
order on triples is lexicographic: $m > 0$ first because $s$ is
smallest; if $m = 0$ then $n > 0$ because $q$ is next; if $m = n = 0$
then $l < 0$ because $|r|^{-1}\ll 1$.  This is a chamber convention for
the completed product, not a microscopic charge positivity statement.
\end{definition}

\begin{remark}[Chamber transport]
\label{rem:chamber-transport}
The pairing $\langle\gamma, \gamma'\rangle_{\mathrm{ind}} = 2(n m' + n'
m) - l l'$ is the Gram pairing on triples; the order on
$\Gamma_{\mathrm{eff}}$ is a Fock polarization at one cusp.  An element
of the duality group $\Sp_4(\Z)$ transports both the chamber and the
expansion to the corresponding product chamber at the transported cusp.
The arithmetic group $\Sp_4(\Z)$ is the symmetry of the genus-two
chemical potential $Z$ and of the invariant Gram plane on
$\Lambda^{3, 2}$; it is not asserted to be the full microscopic
T-/U-duality group of the K3$\times T^2$ compactification, which is
larger (folklore, \cite{HM,DVV}).
\end{remark}

\begin{remark}[Chamber dependence of the wrapped stratum]
\label{rem:chamber-dependence-wrapped}
The cusp chamber chooses a positive direction on
$\Gamma_{\mathrm{eff}}$.  At $\mathfrak c_\infty$ the wrapped stratum
carries $\deg_E > 0$ branes; at a transported cusp the wrapped
direction is the image of $\deg_E$ under the corresponding chamber
transport.  Mixed Hall correspondences are chamber-equivariant when the
finite Hall stage $R$ is compatible with the duality element; otherwise
the correspondence stack is the original closed flag-Quot stack
restricted to the new chamber.
\end{remark}

\begin{definition}[Chamber-compatible Hall descent datum]
\label{def:chamber-compatible-hall-descent}
Let \(\sigma\) and \(\sigma'\) be two chosen stability data, and let
\(g\in\Sp_4(\Z)\) transport the standard cusp chamber to another
product chamber.  A chamber-compatible Hall descent datum is a
Kontsevich--Soibelman wall-crossing system
\[
\mathsf{WC}_{g,R}:\mathcal C_{\sigma,S,\le R}\longrightarrow
\mathcal C_{\sigma',S,\le R'}
\]
on the retained finite Hall stages, with transition maps in \(R\),
which preserves the normal-ordered Gram map \(\overline\Pi_X\), transports
the local/wrapped split, the quotient-after-correspondence datum, the
orientation lines, and the protected integration maps.  The identity
\(\overline\Pi_X(\widehat c\star\widehat c')
=\overline\Pi_X(\widehat c)+\overline\Pi_X(\widehat c')\) on
\(\widehat\Gamma_X\) is a chamber-independent lattice fact; the compact
Hall category and its primitive bracket are not chamber-independent without
\(\mathsf{WC}_{g,R}\) (folklore/\(\beta\), \cite{KSCoHA}).
\end{definition}

\section{Local protected observable algebra (sectorial source datum)}
\label{sec:local-protected-observable-algebra}

The compact realization route of Chapter~\ref{chap:dirac-igusa-pro-object}
takes as supplied input a chain-level local observable system on
opens $U\subset X = S\times E$.  The system is graded by the
algebraic/effective Hall support, factorizes on disjoint opens, and
restricts on $E$-projection-finite supports to a chiral object on
$\operatorname{Ran}(E)_{b = 0}$.  The data are external assumptions
named in the literature and form the source datum for the compact
realization.

\begin{definition}[Local protected observable algebra]
\label{def:local-protected-observable-algebra}
Let $U\subset X = S\times E$ be open.  Let
$\Gamma_{X, \sigma}^{\mathrm{eff}}(U)\subset\Gamma_X^{\mathrm{alg}}
\subset\Gamma_X^{\mathrm{form}}$ be the supplied algebraic/effective
charge support.  For $c\in\Gamma_{X, \sigma}^{\mathrm{eff}}(U)$, let
$\mathfrak M_c(U)$ be the supplied derived moduli stack of compactly
supported B-branes on $U$ of Mukai charge $c$.  The PTVV theorem supplies
the ambient $(-1)$-shifted symplectic structure on the $\mathrm{CY}_3$
perfect-complex moduli problem (proved, \cite{PTVV2013}).  Let
$\Phi_c$ denote the BBDJS perverse sheaf $P^\bullet_{\mathfrak M_c, s_c}$
on an oriented d-critical truncation, when such a d-critical structure
and square-root orientation are supplied (proved,
\cite[Theorem~6.9]{BBDJS2015}).  Let $o_c$ be the orientation line, a
square root of the determinant line of the cotangent complex of
$\mathfrak M_c(U)$, supplied as part of the Hall datum.

The \emph{local protected observable algebra of $X$ on $U$} is the
source object
\[
\mathcal A_X(U)
:= \bigoplus_{c\in\Gamma_{X, \sigma}^{\mathrm{eff}}(U)}
H_*^{\mathrm{BM}}\!\bigl(\mathfrak M_c(U),\,\Phi_c\otimes o_c\bigr),
\]
graded by the supplied algebraic/effective support, not by the full
formal lattice $\Gamma_X^{\mathrm{form}}$.  For pairwise disjoint
$U_1, \ldots, U_k\subset X$, the proposed factorization map is
\[
\mathcal A_X(U_1)\otimes\cdots\otimes\mathcal A_X(U_k)
\longrightarrow\mathcal A_X(U_1\sqcup\cdots\sqcup U_k),
\]
defined chain-level by Massey--Saito Thom--Sebastiani for vanishing
cycles
\[
\Phi_{F_1\oplus\cdots\oplus F_k}
\simeq \Phi_{F_1}\boxtimes\cdots\boxtimes\Phi_{F_k}
\]
on disjoint d-critical charts (proved,
\cite{MasseyST,SaitoMHM}), together with the Ext-vanishing
$\mathrm{Ext}^*(\mathcal F, \mathcal G) = 0$ for branes
$\mathcal F, \mathcal G$ with disjoint supports
(folklore).
\end{definition}

\begin{remark}[External source datum]
\label{rem:source-interface-external-assumptions}
The components are external assumptions: PTVV $(-1)$-shifted symplectic
structure on derived moduli of B-branes (proved,
\cite{PTVV2013}); BBDJS vanishing-cycle perverse sheaf and its
Thom--Sebastiani isomorphism (proved,
\cite{MasseyST,SaitoMHM,BBDJS2015}); orientation lines and
Joyce--Upmeier-compatible multiplicative orientation data (ambient
technology proved in \cite[Theorem~4.1]{JoyceUpmeier}; the
\(K3\times E\) quotient orientation is supplied separately).  The factorization map relies on disjoint-supports
Ext-vanishing for compactly supported branes (folklore;
absent for non-compact supports, where one recovers chiral algebra
structure, not commutative factorization).  Compactness on $X = S\times
E$ is open because the relevant moduli of branes are typically
non-quasicompact in $\Lambda_S$-charge; the sectorial truncation of
\S\ref{sec:k3xe-hybrid-ran-carrier} is the finite-stage carrier, and
the compact-source recognition datum \((S1)\)--\((S10)\) is the
remaining finite recognition problem.
At finite stage the retained-moduli assertion is represented by the
tables \(\mathrm{hn\_type\_bounds}\),
\(\mathrm{semistable\_substacks}\), \(\mathrm{derived\_enhancements}\),
\(\mathrm{universal\_complexes}\), \(\mathrm{scalar\_rigidifications}\),
\(\mathrm{stratifications}\), \(\mathrm{extension\_flag\_stacks}\),
\(\mathrm{cosection\_atlas}\), \(\mathrm{transitions}\), and
\(\mathrm{scalar\_firewall}\), not by the scalar Hilbert-scheme test.
At present the packet
\(\texttt{certificates/moduli/k3e\_finite\_moduli}\) supplies only the
obstruction ledger
\(\texttt{MODULI\_OBSTRUCTION\_LEDGER\_VERIFIED}\), with
\(\texttt{moduli\_certification=false}\).
\end{remark}

\section{The four obstructions and the comparison target}
\label{sec:k3xe-four-obstructions}

The hybrid carrier and the source datum are supplied as a finite
holomorphic $E_3$-prefactorization datum and a \(K3\)-to-\(E\)
specialization datum on $\operatorname{Ran}^{\mathrm{hyb}}(E)$.  The
cross-level comparison
$\operatorname{Pf}_{\mathrm{prot}}(\mathfrak D_X^{\mathrm{DI}}) =
\Delta_5$ holds on \(K3\times E\) relative to four retained data,
with \((O2)\) represented by the retained wall-atlas datum of
Appendix~\ref{app:obstruction-discharges}
(Theorems~\ref{thm:G-D0}, \ref{thm:G-O1}, \ref{thm:G-O1-plus},
\ref{thm:G-O2-orbit-transport}).  They are the orientation and
Pfaffian data feeding
Chapter~\ref{chap:dirac-igusa-pro-object}.

\begin{description}
\item[$(\mathrm{D0})$ {\rm D0-degeneration limit.}]
The retained D0-HN datum of Definition~\ref{def:G-D0-HN-system}
supplies the flat D0-degeneration family, additive K3 semiregularity
cosections, reduced vanishing-cycle and orientation specializations,
Pfaffian-line extension, and strict HN transitions
(criterion Theorem~\ref{thm:G-D0}).
\item[$(\mathrm{O1})$ {\rm Strong reduced orientation.}]
The K3$\times$E-quotient self-Ext frame bundle carries a strong reduced
orientation on the wrapped stratum, descending the BBDJS orientation
line $o_c$ across the $E$-quotient relative to
\((O1_{\mathrm{quot}})\) (Theorem~\ref{thm:G-O1}).
\item[$(\mathrm{O1}^+)$ {\rm Weyl-equivariant transport.}]
The orientation transports $W^{(2)}(\Lambda_{II}^{2, 1})$-equivariantly
along reflections relative to the chosen Weyl lifts
(Theorem~\ref{thm:G-O1-plus}); after the type-II wall atlas supplies
the rank-one local signs, the transported character agrees with
$\nu_{\Delta_5}$ on the BKM Weyl group.
\item[$(\mathrm{O2})$ {\rm Local Pfaffian wall normal form.}]
On type-II walls of $\Lambda_{II}^{2, 1}$, the local Pfaffian wall
normal form is the rank-one crossing
$[\mathbb R\xrightarrow{u}\mathbb R]$; the retained
\((O2_{\mathrm{atlas}})\) datum transports these local signs over the
half-Hilbert orbit, whose sign support has cardinality
$4096 = 2^{12}$ (Theorem~\ref{thm:G-O2-orbit-transport}).
\end{description}

The protected scalar shadow on the closed K3$\times$E$\to$point
cobordism is the OP/DT identity
\[
Z^X_{\mathrm{DT}}
= Z^X_{\mathrm{OP}}(P_{\mathrm{OP}}, Q_{\mathrm{OP}}, T_{\mathrm{OP}})
= -4096\,\Delta_5(Z)^{-2}
\qquad\text{(proved, \cite{OB,OPand,DT}).}
\]
Equivalently $\chi_{10}^{\mathrm{OP}} = (64^{-1}\Delta_5)^2 = 4096^{-1}\Delta_5^2$
and $Z^X_{\mathrm{OP}} = -(\chi_{10}^{\mathrm{OP}})^{-1}$; the factor
$4096 = 64^2$ is a theta-leading normalization conversion and the minus
sign is the OP scalar convention.  The orientation-forgetful shadow
\[
Z_{\mathrm{BPS}}^{K3\times E}
= (\Phi_{10}^{\mathrm{un}})^{-1}
= \Delta_5^{-2}
\]
is the level-$n$ protected trace of the retained \(K3\times E\) data on the
closed cobordism in the unnormalized convention $\Phi_{10}^{\mathrm{un}}
= \Delta_5^2$; the inverse-square scalar cannot recover the orientation
character that distinguishes $\Delta_5$ from $-\Delta_5$.  A
\(3d\) gravitational path-integral realization requires the Hall--Borcherds residual
of \cite{ChiralBarCobarVolII} and is not claimed.  The conditional
cross-level identity
$\operatorname{Pf}_{\mathrm{prot}}(\mathfrak D_X^{\mathrm{DI}}) =
\Delta_5$ supplies the orientation/Pfaffian datum under
$(\mathrm{D0}), (\mathrm{O1}), (\mathrm{O1}^+), (\mathrm{O2})$.

\begin{remark}[Compatibility of the \(\kappa\)-tuple]
\label{rem:kappa-tuple-compatibility}
The $\kappa$-tuple $(0, 0, 3, 5, 24)$ is consistent across
\cite{CYQGVolIII} (Vol III), the Hall--Borcherds
residual of \cite{ChiralBarCobarVolII} (Vol II), and the
Beilinson--Drinfeld factorization framework of \cite{BD} via Vol I
\cite{ChiralBarCobarVolI}.  The local model
$\R^2_{\mathrm{top}}\times\C^2_{\mathrm{hol}}$ + brane completion of
\cite{MixedHTSCompanion} contributes the
holomorphic de Rham obstruction class on the global target; the
identity $Z_{\mathrm{BPS}}^{K3\times E} = \Delta_5^{-2}$ does not
promote to the global target via formal Hamiltonian BF, and the
Hall--Borcherds residual is the remaining chain-level bridge.
\end{remark}

\section{Notation and fixed data}
\label{sec:k3xe-notation-summary}

\begin{description}
\item[$X = S\times E$.]
Smooth complex projective trivial-canonical threefold; $S$ is a K3
surface, $E$ a smooth elliptic curve.  $\dim_\C X = 3$ and
\(h^{1,0}(X)=1\).
\item[$\Lambda_S$.]
Mukai lattice; even unimodular signature $(4, 20)$.
$\Lambda_S\cong U^{\oplus 4}\oplus E_8(-1)^{\oplus 2}\cong II_{4, 20}$.
\item[$\Gamma_X^{\mathrm{form}}, \Gamma_X^{\mathrm{phys}}$.]
Formal full-Mukai dyonic lattice $\Lambda_S\oplus\Lambda_S$.  The
superscript ``$\mathrm{phys}$'' denotes this formal even branch, not the
$\mathcal N = 4$ charge lattice.
\item[$\Gamma_X^{\mathrm{alg}},\Gamma_{X, \sigma}^{\mathrm{eff}}$.]
Algebraic K-class lattice $\subset\Gamma_X^{\mathrm{form}}$; chosen
$\sigma$-effective semigroup $\subset\Gamma_X^{\mathrm{alg}}$.
\item[$\Gamma_{\mathrm{gram}} = \Z^3$.]
Gram-index lattice, $(n, l, m)$.
\item[$\Pi_X$.]
Quadratic Gram map, $\Pi_X(Q, P) = ((Q, Q)/2,\,(Q, P),\,(P, P)/2)$.
\item[$B$.]
Symmetric bilinear cocycle, $B(c, c') = (Q\cdot Q',\,Q\cdot P' + Q'\cdot
P,\,P\cdot P')$; satisfies $B = -\delta\Pi_X$ and $B(c, c) = 2\Pi_X(c)$.
\item[$\widehat\Gamma_X$.]
Normal-ordered extension $\Gamma_X^{\mathrm{form}}\oplus_B\Gamma_{\mathrm
{gram}}$; abelian under $\star$.
\item[$\overline\Pi_X$.]
Additive normal-ordered Gram homomorphism $\overline\Pi_X(c, T) =
\Pi_X(c) - T$.
\item[$\operatorname{Ran}^{\mathrm{hyb}}(E)$.]
Finite-stage hybrid prestack on the elliptic factor, with local
\(b=0\) subprestack
\(\operatorname{Ran}^{\mathrm{loc}}_{R,\mathrm{pre}}(E)\) and wrapped
\(b>0\) subprestack
\(\operatorname{Ran}^{\mathrm{wr}}_{R,\mathrm{pre}}(E)\).
\item[$\mathcal D_X = \Delta_5$, $\mathrm{den}(\mathfrak g_{\Delta_5})$,
$\mathrm{Pf}_{\mathrm{prot}}(\mathfrak D_X^{\mathrm{DI}})$.]
Three named weight-five objects: automorphic section, BKM denominator,
compact Pfaffian (View \textcircled{1}, \textcircled{2}, \textcircled{3}).
\end{description}

The fixed data are: (a) the formal Mukai sector
$\Gamma_X^{\mathrm{form}} = \Lambda_S\oplus\Lambda_S$ and its three-tier
support hierarchy; (b) the bilinear polarization \(B = -\delta\Pi_X\) as the
full structural content of the Gram defect; (c) the normal-ordered
extension $\widehat\Gamma_X$ with additive homomorphism
$\overline\Pi_X$, fixed by the fixed-lift raw-pushforward obstruction
theorem; (d) the hybrid Ran carrier
$\operatorname{Ran}^{\mathrm{hyb}}(E)$, fixed by the projection-locality
obstruction; (e) the connected/finite split of the $E$-descent classes,
with the linearization criterion separated from the gerbe class; (f)
the local protected observable algebra source datum, taken as an
external assumption; (g) the four obstructions
$(\mathrm{D0}), (\mathrm{O1}), (\mathrm{O1}^+), (\mathrm{O2})$
conditioning the cross-level identity ②$\leftrightarrow$③.

Chapter~\ref{chap:dirac-igusa-pro-object} constructs from these fixed
data the Dirac--Igusa pro-object $\mathfrak D_X^{
\mathrm{DI}} = \lim_R(\mathcal A_{X, R}^{E_3}, F_{X, R}^{\mathrm{hyb}},
\Gamma_{X, R}, \Pi_X, \Phi_R, o_R, H_R, P_R^\Pi, C_{X, R},
\Theta_{\mathrm{Kos}, R}, \mathcal L_{\mathrm{Pf}, R}, \mathrm{pf}_{X, R},
\epsilon_{o, R}, \mathrm{Rec}_R)$.


\chapter{View I: the Borcherds--Gritsenko--Nikulin section}
\label{ch:view1-automorphic}

\section{The pentadic level and the Beilinson cut}
\label{sec:view1-pentadic-level}

The automorphic reading of \(\Delta_5\) is the first theorem-level
reading.  It is a holomorphic section of the
weight-\(5\) automorphic line on the genus-two Siegel modular space
\(\Sp_4(\Z)\backslash\mathbb H_2\): the Borcherds--Gritsenko--Nikulin lift of the
weight-zero index-one weak Jacobi form \(\phi_{0,1}\), normalized at the
isotropic cusp \(\mathfrak c_\infty\) by the theta-constant leading
coefficient \(64\).  The \(\Sp_4(\Z)\)-equivariant data of View~I
are the multiplier system \(\nu_{\Delta_5}\), the Borcherds
chamber, the divisor support, and the weight.

\paragraph{Three names.}
The handle ``\(\Delta_5\)'' is shared among three structurally distinct
objects and must be tagged on first use:
\begin{itemize}
\item \emph{the Borcherds--Gritsenko--Nikulin section}
\(\mathcal D_X=\Delta_5\) of \(\mathcal L^{5}\otimes\nu_{\Delta_5}\) on
\(\Sp_4(\Z)\backslash\mathbb H_2\), View~I;
\item \emph{the BKM denominator}
\(\operatorname{den}(\mathfrak g_{\Delta_5})=64^{-1}\Delta_5(2Z)\), View~II,
the Borcherds--Kac--Moody denominator algebra;
\item \emph{the protected Pfaffian}
\(\operatorname{Pf}_{\mathrm{prot}}(\mathfrak D_X^{\mathrm{DI}})\),
View~III, constructed in the Dirac--Igusa chapter as a section on the
Pfaffian line of the cosection-reduced compact Hall datum after the
retained D0-HN datum, retained quotient-orientation datum, chosen Weyl
lifts, retained \((O2_{\mathrm{atlas}})\) wall datum, and finite
geometric Pfaffian datum \((P_{\mathrm{fin}})\) of
Appendix~\ref{app:obstruction-discharges}.
\end{itemize}
The bare label \(\Delta_5\) never substitutes for any of the three: it
names the automorphic section unless a reading is specified, while
\(\mathcal D_X\), \(\operatorname{den}(\mathfrak g_{\Delta_5})\), and
\(\operatorname{Pf}_{\mathrm{prot}}(\mathfrak D_X^{\mathrm{DI}})\) live at
distinct categorical-dimensional levels.

\paragraph{Beilinson cut: scope of View~I.}
View~I is the scalar modular shadow.  In the reconstitution order
``primitive objects first, cross-level identities second, scalar
shadows last'' View~I is a single holomorphic section of an automorphic
line on a Shimura variety, with its multiplier system and its weight.
No operator-level datum is asserted at View~I.  In particular, the equality
\(\mathcal D_X=\operatorname{Pf}_{\mathrm{prot}}(\mathfrak D_X^{\mathrm{DI}})\) of
View~I and View~III is the Pfaffian--Dirac theorem of
Chapter~\ref{ch:pfaffian-dirac}, relative to the retained data recorded
in Appendix~\ref{app:obstruction-discharges}; it is not asserted in this
chapter.
The shadows are \(\Delta_5\), \(\chi_{10}\), the level-\(n\) BPS partition function
\(Z_{\mathrm{BPS}}^{K3\times E}=\Delta_5^{-2}\), and the
Oberdieck--Pixton scalar branch
\(Z^X_{\mathrm{OP}}=-4096\,\Delta_5^{-2}\); none of them is a Hilbert
space, a Hall pairing, an orientation line, or an operator product.

\section{The genus-two Siegel space and the Borcherds chamber}
\label{sec:view1-siegel-and-chamber}

Let \(\mathbb H_2\) be the genus-two Siegel upper half space.  An element
\(Z\in\mathbb H_2\) is a complex symmetric \(2\times 2\) matrix with
positive-definite imaginary part,
\[
Z=\begin{pmatrix}z_1&z_2\\z_2&z_3\end{pmatrix},
\qquad
\Im Z>0.
\]
Set Fourier variables
\[
q=e^{2\pi i z_1},\qquad
r=e^{2\pi i z_2},\qquad
s=e^{2\pi i z_3}.
\]
The fixed cusp polarization at the rational isotropic boundary
\(\mathfrak c_\infty\) of \(\Sp_4(\Z)\backslash\mathbb H_2\) — the standard
type-II cusp polarization — is the
ordered limit
\[
0<|s|\ll|q|\ll|r|^{-1}\ll1.
\]
Equivalently, \(\Im z_3\) goes to \(+\infty\) faster than \(\Im z_1\)
goes to \(+\infty\), and \(\Im z_2\) is bounded.  This polarization
selects a Weyl chamber on the \(O(3,2)\) symmetric domain through the
\(\PSp_4\simeq\SO(3,2)_+\) bridge of Chapter~\ref{ch:singular-theta-lift},
and hence selects a positive cone, the effective semigroup
\[
\Gamma_{\mathrm{eff}}
=\{(n,l,m)\mid m>0,\ n\ge0,\ l\in\Z\}
\cup\{(n,l,0)\mid n>0,\ l\in\Z\}
\cup\{(0,l,0)\mid l<0\}
\subset\Z^3,
\]
and a corresponding completed Fourier expansion ring.  In this completion every
Fourier coefficient of every formal product receives only finitely many
contributions: the chamber walls are the inequalities defining the
ordering, and the chamber-completed product is the only product
considered.

\paragraph{Choices fixed once and for all.}
The cusp \(\mathfrak c_\infty\), the Fourier variables \((q,r,s)\), the
polarization \((0<|s|\ll|q|\ll|r|^{-1}\ll1)\), and the effective semigroup
\(\Gamma_{\mathrm{eff}}\subset\Z^3\) are fixed.  Every View-I product expansion
is the product expansion in the cusp-completed
semigroup algebra of \(\Gamma_{\mathrm{eff}}\).  The cusp expansion at any
other rational isotropic cusp is obtained by an
\(\Sp_4(\Z)\)-translate of these data.

\section{The weak Jacobi input and the K3 elliptic genus}
\label{sec:view1-jacobi-input}

The Borcherds input is the \emph{weight-zero index-one weak Jacobi form}
\(\phi_{0,1}\in J^{\mathrm{weak}}_{0,1}\) of Eichler--Zagier
\cite[\S2.2]{EZ:1}, normalized so that the K3 elliptic genus is
\[
Z_{K3}(\tau,z)=2\,\phi_{0,1}(\tau,z).
\]
Write the Fourier expansion as
\[
\phi_{0,1}(\tau,z)=\sum_{n\ge 0,\;l\in\Z}f(n,l)\,q^nr^l,
\qquad q=e^{2\pi i\tau},\quad r=e^{2\pi iz}.
\]
The first Fourier coefficients are
\begin{equation}
\label{eq:view1-phi01-coefficients}
\begin{aligned}
f(0,-1)&=1,&\quad f(0,0)&=10,&\quad f(0,1)&=1,\\
f(1,-2)&=10,&\quad f(1,-1)&=-64,&\quad f(1,0)&=108,&\quad f(1,1)&=-64,&\quad f(1,2)&=10,\\
f(2,-3)&=1,&\quad f(2,-2)&=108,&\quad f(2,-1)&=-513,&\quad f(2,0)&=808,
\end{aligned}
\end{equation}
extending in the standard way by hyperbolic symmetry
\(f(n,l)=f(n,-l)\) and Eichler--Zagier support
\(f(n,l)=0\) for \(4n-l^2<-1\) \cite[\S2.2]{EZ:1}.  The values in
\eqref{eq:view1-phi01-coefficients} are obtained by the exact rational
expansion of the standard theta-quotient expression
\(\phi_{0,1}=4\sum_{i=2,3,4}\theta_i(\tau,z)^2/\theta_i(\tau,0)^2\)
[\textit{computed}; cross-checked against \cite[Theorem~9.2]{EZ:1}].

\paragraph{Borcherds-weight handle.}
For a fixed Borcherds input \(\Phi\), the weight identity
\(\kappa_{\mathrm{BKM}}(\Phi)=c_{\Phi}(0)/2\) (Borcherds 1995
\cite[Theorem~10.1]{Borcherds95}; Gritsenko 2018
\cite{Gritsenko2018Reflective}) gives, on the row \((t,N)=(1,1)\) of
the Cl\'ery--Gritsenko catalogue \cite[Theorem~1.4]{GC},
\[
c_1(0):=f(0,0)=10,\qquad
\kappa_{\mathrm{BKM}}(\Delta_5)=\frac{c_1(0)}{2}=5,
\]
to be matched in
Proposition~\ref{prop:view1-borcherds-gritsenko-weight-identity} below
with the geometric weight of the Borcherds lift.  The notation
\(\kappa_{\mathrm{BKM}}\) is the BKM cusp-form-weight entry of the
\(\kappa\)-tuple \((\kappa_{\mathrm{cat}},\kappa^{\mathrm{Hodge}}_{\mathrm{ch}},
\kappa^{\mathrm{Heis}}_{\mathrm{ch}},\kappa_{\mathrm{BKM}},\kappa_{\mathrm{fiber}})
=(0,0,3,5,24)\) of the \(K3\times E\) target.  The bare symbol
\(\kappa\) is not a scalar invariant; every entry carries its semantic subscript;
the additive formula
\(\kappa_{\mathrm{BKM}}=\kappa_{\mathrm{ch}}+\chi(\mathcal O_{\mathrm{fiber}})\)
does not define the BKM coordinate.  The BKM coordinate is
\(\kappa_{\mathrm{BKM}}(\Delta_5)=c_1(0)/2=5\).

\section{The theta-constant normalization}
\label{sec:view1-theta-normalization}

The theta-constant Igusa form \(\Delta_5\) is fixed by its leading
Fourier monomial at \(\mathfrak c_\infty\).  Writing the theta-constant
representation of Igusa \cite[\S5]{Igusa1962}, \cite[\S2]{Igusa1964II}
\[
\Delta_5(Z)
=\!\!\!\!\prod_{\substack{a,b\in(\Z/2\Z)^2\\{}^t\!ab\equiv 0\bmod 2}}\!\!\!
\nu_{a,b}(Z),
\qquad
\nu_{a,b}(Z)=\!\sum_{l\in\Z^2}\exp\!\bigl(\pi i\bigl(z[l+\tfrac12 a]+{}^tbl\bigr)\bigr),
\]
the four characteristics with \(a=(0,0)\) start with the constant \(1\),
the two factors with \(a=(1,0)\) start with \(2q^{1/8}\), the two factors
with \(a=(0,1)\) start with \(2s^{1/8}\), and the two factors with
\(a=(1,1)\) start with \(2q^{1/8}r^{1/4}s^{1/8}\); multiplying gives

\begin{proposition}[Theta-constant leading coefficient]
\label{prop:view1-theta-constant-leading-coefficient}
\[
[q^{1/2}r^{1/2}s^{1/2}]\Delta_5\;=\;2^6\;=\;64.
\]
\end{proposition}

\noindent
[\textit{proved}; Igusa 1962 \cite[Theorem~3]{Igusa1962}, Gritsenko
1999 \cite[Theorem~4.1]{GN}; the rational theta-product expansion gives
the same leading coefficient.]

\begin{definition}[Monic Borcherds--Gritsenko form]
\label{def:view1-monic-Borcherds-Gritsenko-form}
The \emph{monic genus-two Borcherds--Gritsenko form} is
\[
D_5(Z)\;:=\;64^{-1}\Delta_5(Z).
\]
By Proposition~\ref{prop:view1-theta-constant-leading-coefficient},
\(D_5=q^{1/2}r^{1/2}s^{1/2}+\ldots\) at \(\mathfrak c_\infty\).
\end{definition}

\noindent
This is the normalization used by Gritsenko--Nikulin
\cite[Theorem~2.1, (2.15)--(2.16)]{GNII}: their Borcherds product is the
monic form \(D_5\); the theta-constant convention restores the prefactor
\(64\).

\section{The Borcherds--Gritsenko--Nikulin section}
\label{sec:view1-borcherds-gritsenko-nikulin-section}

\subsection{The product expansion}
\label{ssec:view1-product-expansion}

Borcherds's automorphic product theorem
\cite[Theorem~10.1]{Borcherds95}, specialized to the genus-two Siegel
modular variety \(\Sp_4(\Z)\backslash\mathbb H_2\) by Gritsenko--Nikulin
\cite[Theorem~2.1, Example~2.4]{GNII} via the
exceptional exterior-square isogeny \(\PSp_4(\R)\simeq\SO(3,2)_+\) of
Chapter~\ref{ch:singular-theta-lift}, lifts the weak Jacobi form
\(\phi_{0,1}\) to a meromorphic Siegel modular form of weight
\(f(0,0)/2=5\) and multiplier system \(\nu_{\Delta_5}\) with infinite
product expansion in the cusp chamber.

\begin{theorem}[Borcherds--Gritsenko--Nikulin lift, monic form]
\label{thm:view1-borcherds-gritsenko-nikulin-lift}
In the cusp chamber \(0<|s|\ll|q|\ll|r|^{-1}\ll1\),
\[
D_5(Z)
\;=\;q^{1/2}r^{1/2}s^{1/2}\!\!
\prod_{(n,l,m)\in\Gamma_{\mathrm{eff}}}\!\!
(1-q^nr^ls^m)^{f(nm,l)},
\]
the product converging absolutely on a neighborhood of \(\mathfrak
c_\infty\) and continuing meromorphically to all of
\(\mathbb H_2\) as a holomorphic Siegel modular form of weight \(5\) on
\(\Sp_4(\Z)\) with character \(\nu_{\Delta_5}\) of order two.
\end{theorem}

\noindent
[\textit{proved}; Borcherds 1995 \cite[Theorem~10.1]{Borcherds95},
Gritsenko--Nikulin 1998 \cite[Theorem~2.1, (2.15)--(2.16)]{GNII},
Gritsenko 1999 \cite[Theorem~4.1]{GN}.]

\begin{corollary}[Theta-normalized section]
\label{cor:view1-theta-normalized-section}
\[
\Delta_5(Z)
\;=\;64\,q^{1/2}r^{1/2}s^{1/2}\!\!
\prod_{(n,l,m)\in\Gamma_{\mathrm{eff}}}\!\!
(1-q^nr^ls^m)^{f(nm,l)}.
\]
\end{corollary}

\noindent
The exponent \(f(nm,l)\) of the \((n,l,m)\)-factor depends only on the
Hecke--Verlinde--DMVV product \(nm\) and on the Jacobi index \(l\),
reflecting the second-quantized structure of the Borcherds lift, whose
combinatorial origin is the Dijkgraaf--Moore--Verlinde--Verlinde
elliptic genus of symmetric products \cite{DMVV}.

\subsection{The View~I datum}
\label{ssec:view1-package}

\begin{definition}[The Borcherds--Gritsenko--Nikulin section]
\label{def:view1-borcherds-gritsenko-nikulin-section}
The \emph{Borcherds--Gritsenko--Nikulin section} of the weight-\(5\)
automorphic line is the View~I representative
\[
\mathcal D_X\;:=\;\Delta_5
\]
on \(\Sp_4(\Z)\backslash\mathbb H_2\), with Borcherds product expansion
at \(\mathfrak c_\infty\) given by Corollary~\ref{cor:view1-theta-normalized-section}.
\end{definition}

\begin{remark}[\(\mathcal D_X\) is the View~I name]
\label{rem:view1-D-X-naming}
The handle \(\mathcal D_X\) is the View~I name of the section.  It coincides as
a holomorphic Siegel modular form with the theta-constant Igusa cusp
form \(\Delta_5\); the renaming records the role: View~I scalar
shadow, the determinant of the virtual K3 index in
\(K_0(\mathrm{sVect}_{\mathrm{fd}})\) (the
\(K_0\)-determinant identity proved at
Proposition~\ref{prop:bgn-identification-recap}; the operator
Pfaffian datum belongs to View~III), and the input section of the Pfaffian--Dirac
theorem of Chapter~\ref{ch:pfaffian-dirac}.  The compact
operator-level datum \(\operatorname{Pf}_{\mathrm{prot}}(\mathfrak
D_X^{\mathrm{DI}})\) of View~III is asserted equal to \(\mathcal D_X\) under
\((D_0)\), \((O_1)\), \((O_1)^+\), \((O_2)\); without these obstruction
trivializations, only \(\mathcal D_X\) is rigorously available.
\end{remark}

\section{Weight, multiplier, and divisor}
\label{sec:view1-weight-multiplier-divisor}

\subsection{The Borcherds weight}
\label{ssec:view1-borcherds-weight}

\begin{proposition}[Borcherds--Gritsenko weight identity at \(N=1\)]
\label{prop:view1-borcherds-gritsenko-weight-identity}
\[
\operatorname{wt}(\mathcal D_X)
\;=\;\frac{f(0,0)}{2}
\;=\;5
\;=\;\frac{c_1(0)}{2}
\;=\;\kappa_{\mathrm{BKM}}(\Delta_5).
\]
The weight is integral; no metaplectic cover enters the weight factor.
\end{proposition}

\noindent
[\textit{proved}; Borcherds 1995 \cite[Theorem~10.1, weight
formula]{Borcherds95}; Gritsenko--Nikulin 1998
\cite[Example~2.4]{GNII}; Gritsenko 1999 \cite[Theorem~4.1, weight]{GN};
the Borcherds weight identity \(\kappa_{\mathrm{BKM}}(\Phi)=c_{\Phi}(0)/2\)
for a fixed input form and lattice; Cl\'ery--Gritsenko \cite{GC}, Vol~III.]

\noindent
The row \((t,N;\operatorname{wt},c)=(1,1;5,10)\) is the leading
\(F_j\)-entry of the eight-row Cl\'ery--Gritsenko catalogue
\cite[Theorem~1.4]{GC}.  In the \(F_j\)-order the remaining rows carry
weights \((2,1,\tfrac12,3,2,\tfrac32,1)\) and constants
\((4,2,1,6,4,3,2)\); in the input-pair order
\((2,1),(1,2),(3,1),(1,3),(4,1),(1,4),(2,2)\) the same data read
\((2,3,1,2,\tfrac12,\tfrac32,1)\) and
\((4,6,2,4,1,3,2)\).  Each row satisfies
\(\kappa_{\mathrm{BKM}}(F_j)=c_j(0)/2\).  The half-integer rows are
realized by \(v_\eta^3\times v_H\) and
\(\chi_4^{(4)}\times v_H\), respectively.  The Vol~III
universal Borcherds-weight identity replaces the additive formula
\(\kappa_{\mathrm{BKM}}=\kappa_{\mathrm{ch}}+\chi(\mathcal O_{\mathrm{fiber}})\):
the \(\kappa\)-tuple records separate invariants, and the
\(K3\times E\) entry is precisely \(\kappa_{\mathrm{BKM}}=5\).

\subsection{The Maass multiplier system}
\label{ssec:view1-maass-multiplier}

The character \(\nu_{\Delta_5}\colon\Sp_4(\Z)\to\C^\times\) of Maass
\cite{MAASS:1} has order two,
\(\nu_{\Delta_5}^{2}=1\), and is determined on the standard generators
of \(\Sp_4(\Z)\) by
\[
\nu_{\Delta_5}\!
\begin{pmatrix}0&\Id_2\\-\Id_2&0\end{pmatrix}=1,
\quad
\nu_{\Delta_5}\!
\begin{pmatrix}\Id_2&B\\0&\Id_2\end{pmatrix}=(-1)^{b_1+b_2+b_3},
\]
\[
\nu_{\Delta_5}\!
\begin{pmatrix}{}^tA^{-1}&0\\0&A\end{pmatrix}
=(-1)^{(1+a_1+a_4)(1+a_2+a_3)+a_1a_4},
\]
with
\(B=\bigl(\begin{smallmatrix}b_1&b_2\\b_2&b_3\end{smallmatrix}\bigr)\)
and \(A=\bigl(\begin{smallmatrix}a_1&a_2\\a_3&a_4\end{smallmatrix}\bigr)\)
\cite{MAASS:1}.  The square \(\Delta_5^2=\Delta_{10}\) is scalar.  The
multiplier records the obstruction to identifying \(\Delta_5\) with an
ordinary modular form on \(\Sp_4(\Z)\): \(\Delta_5\) is a section of
\(\mathcal L^5\otimes\nu_{\Delta_5}\), where \(\mathcal L\) is the Hodge
automorphic line and \(\nu_{\Delta_5}\) is the order-two character; the
section \(\Delta_5\) does not descend to a section of the bare Hodge
line \(\mathcal L^5\), but its square \(\Delta_5^2\) is a scalar
section of \(\mathcal L^{10}\).
The finite automorphic packet
\[
\texttt{certificates/automorphic/delta5\_maass\_character}
\]
records the six chamber generator values, the relations
\(\nu_{\Delta_5}|_{W^{(2)}}=\det\),
\(\nu_{\Delta_5}|_{\Aut(\Poly_{II})}=1\), and
\(\nu_{\Delta_5}^2=1\), and the line row
\(\Delta_5\in H^0(\mathbb H_2,\mathcal L^5\otimes\nu_{\Delta_5})\).
The same packet excludes using \(\nu_{\Delta_5}\) as a Pfaffian
orientation line or as an OP scalar branch.

\subsection{The diagonal divisor}
\label{ssec:view1-diagonal-divisor}

Let
\[
\mathcal H_{\mathrm{diag}}
=\Bigl\{\bigl(\begin{smallmatrix}a&0\\0&b\end{smallmatrix}\bigr)\;\Big|\;
a,b\in\mathbb H_1\Bigr\}
\subset\mathbb H_2
\]
be the locus of diagonal genus-two periods.  The Cl\'ery--Gritsenko
classification \cite{GC} records
\[
\operatorname{supp}\operatorname{div}(\Delta_5)
=\Sp_4(\Z)\cdot\mathcal H_{\mathrm{diag}},\qquad
\operatorname{mult}_{\mathcal H_{\mathrm{diag}}}(\Delta_5)=1,
\]
so the divisor of \(\Delta_5\) is the simple paramodular Humbert
divisor.  In the cusp-chamber language, the type-II reflective walls
\(\mathcal H_{\delta_i}\), \(i=1,2,3\), of the lattice
\(\La^{2,1}_{II}\) (the type-II sublattice of the Borcherds even
lattice \(\La^{2,1}\)) are the components of
\(\Sp_4(\Z)\cdot\mathcal H_{\mathrm{diag}}\) that meet the chosen
fundamental Weyl chamber \(\Poly_{II}\) (Chapter~\ref{ch:view2-bkm}); the
corresponding simple roots
\(\delta_1=2f_2-f_3\), \(\delta_2=2f_{-2}-f_3\), \(\delta_3=f_3\)
have Gram pairings
\((\delta_i,\delta_i)=2\) and \((\delta_i,\delta_j)=-2\) for \(i\ne j\)
\cite[Example~2.4]{GNII}, and the corresponding Borcherds divisor
orders are
\[
d_{\delta_1}^{\mathrm{aut}}=f(0,1)=1,\qquad
d_{\delta_2}^{\mathrm{aut}}=f(0,1)=1,\qquad
d_{\delta_3}^{\mathrm{aut}}=f(0,-1)=1.
\]
[\textit{proved}; Cl\'ery--Gritsenko 2011 \cite[Theorem~1.4]{GC};
Gritsenko--Nikulin 1998 \cite[Example~2.4]{GNII}; the Gram computation
is the one displayed in Appendix~\ref{app:lattice-II21}.]
The finite automorphic divisor packet
\[
\texttt{certificates/automorphic/delta5\_humbert\_divisor}
\]
records the three type-II support representatives, the simple order-one
rows for \(\Delta_5\), the pole-order-two row for \(\Delta_5^{-2}\),
and a separate support-versus-multiplicity row.  It is not a Pfaffian
wall chart, an \(O2\) atlas, or a mirror-period discriminant.

\subsection{Maass-character values on the type-II walls}
\label{ssec:view1-maass-character-walls}

\begin{lemma}[Maass-character values on the type-II reflections]
\label{lem:view1-maass-multiplier-restriction}
With the simple type-II roots
\(\delta_1=2f_2-f_3\), \(\delta_2=2f_{-2}-f_3\), \(\delta_3=f_3\)
and chamber-permuting type-I roots
\(\alpha_1=f_{-2}-f_2\), \(\alpha_2=f_2-f_3\), \(\alpha_3=f_{-2}-f_3\),
\[
\nu_{\Delta_5}(s_{\delta_i})\;=\;-1,\qquad
\nu_{\Delta_5}(s_{\alpha_i})\;=\;+1,\qquad i=1,2,3.
\]
Equivalently, for chosen \(\Sp_4\)-lifts and the weight-five slash
action without character,
\[
(\Delta_5|_5s_{\delta_i})(Z)=-\Delta_5(Z),\qquad
(\Delta_5|_5s_{\alpha_i})(Z)=+\Delta_5(Z).
\]
\end{lemma}

\noindent
[\textit{proved}; Maass \cite{MAASS:1}; Gritsenko--Nikulin
\cite[\S3]{GN}; Gritsenko \cite{Gritsenko99}.  The type-II values
follow from \(d_{\delta_i}^{\mathrm{aut}}=1\) and the chamber-wall
cocycle calculation of
Lemma~\ref{lem:bkm-maass-restriction}, which fixes
\(\nu_{\Delta_5}|_{W^{(2)}(\La^{2,1}_{II})}=\det\) on the three
type-II generators; the type-I values follow from
\(d_{\alpha_i}^{\mathrm{aut}}=0\) (vanishing of \(f(0,\pm 2)\) by the
support of \(\phi_{0,1}\) at \(q=0\)).]

The chamber-permuting reflections
\(s_{f_{-2}-f_2},s_{f_2-f_3},s_{f_{-2}-f_3}\) act trivially on
\(\Delta_5\) because the corresponding type-I roots
\(\alpha_i\) are not \(\Sp_4\)-images of timelike short roots in the
cusp chamber: type-I square-$2$ vectors are outside the type-II
sublattice $\La^{2,1}_{II}$ on which the Borcherds product
expansion of $\Delta_5$ is supported, so their walls do not appear
in $\operatorname{div}(\Delta_5)$ at all
\cite[\S2.2]{EZ:1}.

\section{Identification with the virtual \(K_0\)-determinant}
\label{sec:view1-K0-determinant-identification}

The Fock superdeterminant of the Borcherds-graded super vector space
\(\mathcal V=\bigoplus_{(n,l,m)\in\Gamma_{\mathrm{eff}}}
\mathbb U_{nm,l}\) — the \(K_0\)-class of the second-quantized
Hecke--DMVV half of the K3 elliptic genus, with the cusp/Weyl boundary
factors at \(m=0\) carrying the input multiplicities \(f(0,l)\) and
the bulk factors at \(m>0\) carrying \(f(nm,l)\) on
\(\Gamma_{\mathrm{eff}}\) (Definition~\ref{def:gram-index-lattice}) — is

\begin{lemma}[\(K_0\)-Fock superdeterminant]
\label{lem:view1-K0-Fock-superdeterminant}
\[
\sdet_{\mathcal V}(1-x)
\;=\;\!\!\prod_{(n,l,m)\in\Gamma_{\mathrm{eff}}}\!\!
(1-q^nr^ls^m)^{f(nm,l)},
\qquad x^{(n,l,m)}=q^nr^ls^m.
\]
\end{lemma}

\noindent
[\textit{proved}; the Fock-trace lemma is the three-way comparison
\(\sdet_\mathcal V(1-x)=\prod(1-x^\gamma)^{\sdim\mathcal V_\gamma}\)
under the Borcherds-graded second-quantization; the
automorphic--orthogonal--BKM identification sits in
Theorem~\ref{thm:bgn-three-way} of
Chapter~\ref{ch:singular-theta-lift}.]

\begin{proposition}[BGN identification of the virtual determinant]
\label{prop:bgn-identification-recap}
With \(v_0(Z):=64\,q^{1/2}r^{1/2}s^{1/2}\) the cusp/Weyl monomial of
Chapter~\ref{ch:k3xe-setup}, the normalized \(K_0\)-determinant agrees with
the View~I section,
\[
\mathcal D_X(Z)
\;:=\;v_0(Z)\,\sdet_{\mathcal V}(1-x)
\;=\;\Delta_5(Z),
\]
in the cusp-chamber expansion at \(\mathfrak c_\infty\) and hence as
sections of \(\mathcal L^5\otimes\nu_{\Delta_5}\).
\end{proposition}

\begin{proof}
Combine
Theorem~\ref{thm:view1-borcherds-gritsenko-nikulin-lift}, the
theta-constant prefactor
Proposition~\ref{prop:view1-theta-constant-leading-coefficient},
Lemma~\ref{lem:view1-K0-Fock-superdeterminant}, and Cohen's
\(q\)-expansion principle for genus-two Siegel modular forms with
half-integral character: a meromorphic Siegel form is determined by its
expansion at one isotropic cusp.  Both sides have the same Borcherds
product expansion at \(\mathfrak c_\infty\) by direct comparison of
exponents, and \(\mathcal D_X\) is holomorphic of weight \(5\) by
Theorem~\ref{thm:view1-borcherds-gritsenko-nikulin-lift}.
\end{proof}

\noindent
[\textit{proved}; Borcherds 1995 \cite[Theorem~10.1]{Borcherds95},
Gritsenko--Nikulin 1998 \cite[Theorem~2.1]{GNII}, Igusa 1962/1964
\cite{Igusa1962,Igusa1964II}.]

\begin{remark}[Scope of the BGN identification]
\label{rem:view1-bgn-scope}
Proposition~\ref{prop:bgn-identification-recap} is an equality of
holomorphic sections of an automorphic line bundle.  The carrier
\(\mathcal V\) is a class in
\(K_0(\Gamma_{\mathrm{eff}}\text{-graded super vector spaces})\), with no
vertex, chiral, or holomorphic \(E_3\)-factorization structure.
The passage from the \(K_0\)-Fock determinant to a section of an
oriented Pfaffian line over the cosection-reduced compact Hall datum,
and the identification of that section with \(\mathcal D_X\), is the
Pfaffian--Dirac theorem of Chapter~\ref{ch:pfaffian-dirac}, relative to
the retained data recorded in Appendix~\ref{app:obstruction-discharges}.
View~I supplies only the automorphic determinant section.
\end{remark}

\section{The squared determinant section and OP normalization}
\label{sec:view1-squared-scalar}

The Maass character has order two, so the square
\(\Delta_5^2=\Delta_{10}\) is a scalar weight-\(10\) Igusa form on
\(\Sp_4(\Z)\), with no multiplier system.

\begin{proposition}[The square is the Oberdieck--Pandharipande $\chi_{10}$]
\label{prop:view1-square-chi-10}
\[
\Delta_5^2\;=\;\Delta_{10}\;=\;\Phi_{10}^{\mathrm{un}}\;=\;\chi_{10}.
\]
\end{proposition}

\noindent
[\textit{proved}; Igusa \cite{Igusa1962,Igusa1964II}; the symbol
\(\Phi_{10}^{\mathrm{un}}\) is Borcherds's name for the squared theta-constant Igusa
form, \(\chi_{10}\) is the Oberdieck--Pandharipande
\cite{OPand,OB} name; both equal \(\Delta_5^2\).  The naming is
\(\Phi_{10}^{\mathrm{un}}=\chi_{10}=\Delta_5^2\), with the Oberdieck--Pixton OP
normalization \(\chi_{10}^{\mathrm{OP}}=D_5^{\,2}=4096^{-1}\Delta_5^2\)
the monic-on-\(D_5\) variant.]

\paragraph{The OP scalar branch.}
With the OP normalization \(\chi_{10}^{\mathrm{OP}}=D_5^2\), the
Oberdieck--Pixton scalar partition function on \(K3\times E\) is
\cite[(0.4) and Theorem~1]{OB}
\[
Z_{\mathrm{OP}}^{X}\;=\;-\,(\chi_{10}^{\mathrm{OP}})^{-1}
\;=\;-\,4096\,\Delta_5^{-2}.
\]
The factor \(-4096=-(64)^2\) records the squared theta-constant leading
coefficient and the OP scalar sign convention.  At View~I this is the
closed protected trace \(\Delta_5^{-2}\); the orientation character
which distinguishes the Pfaffian section \(\Delta_5\) from its negative
is lost under squaring.  The operator datum entering that trace is
defined at the Pfaffian and orientation level relative to the retained
data of Appendix~\ref{app:obstruction-discharges}.

\paragraph{Three names for the squared form, distinguished.}
The notation is fixed as follows:
\(\Delta_5^2=\Delta_{10}\) is the Igusa convention (theta-constant
square); \(\Phi_{10}^{\mathrm{un}}=\Delta_5^2\) is the
Borcherds--Eisenstein convention;
\(\chi_{10}=\Delta_5^2\) is the Oberdieck--Pandharipande convention; and
\(\chi_{10}^{\mathrm{OP}}=4096^{-1}\Delta_5^2=D_5^{\,2}\) is the
Oberdieck--Pixton variant in which \(D_5\) is the monic Borcherds form
(Definition~\ref{def:view1-monic-Borcherds-Gritsenko-form}).
The finite packet
\(\texttt{certificates/normalizations/delta5\_theta\_leading}\)
records these as separate convention rows and relation rows; equality
of the unscaled sections is not used to identify the OP monic form or
either inverse scalar branch.

\section{The View~II shadow: the \texorpdfstring{\(2Z\)}{2Z} denominator}
\label{sec:view1-2Z-denominator}

The Gritsenko--Nikulin denominator algebra of the
Borcherds--Kac--Moody Lie superalgebra \(\mathfrak g_{\Delta_5}\)
arises as the View~II shadow of \(\Delta_5\) under the doubling
\(Z\mapsto 2Z\):

\begin{proposition}[Gritsenko--Nikulin doubled denominator]
\label{prop:view1-GN-doubled-denominator}
\[
\operatorname{den}(\mathfrak g_{\Delta_5})\;=\;D_5(2Z)\;=\;64^{-1}\Delta_5(2Z).
\]
\end{proposition}

\noindent
[\textit{proved}; Gritsenko--Nikulin 1998 \cite[Theorem~2.1,
(2.15)--(2.16); \S5]{GNII}; Gritsenko 1999 \cite[Theorem~4.1]{GN}.]

\noindent
The doubling records the embedding
\(\La^{2,1}_{II}=\alpha(\Gamma_{\mathrm{ind}})\subset\La^{2,1}\) of
index \(4\).  Its type-II walls are primitive roots of square \(2\), and
the substitution \(Z\mapsto 2Z\) replaces
\(\exp(-\pi i(\rho,Z))\) by \(\exp(-2\pi i(\rho,Z))\) in the
Gritsenko--Nikulin denominator normalization
\cite[Theorem~2.1, (2.15)]{GNII}.  At View~I this is a pointer; the
Cartan, Chevalley, real-and-imaginary roots, and Weyl--Kac--Borcherds
denominator identity of \(\mathfrak g_{\Delta_5}\) sit in
Chapter~\ref{ch:view2-bkm}.

\section{The \(\PSp_4\simeq\SO(3,2)_+\) bridge to View~V}
\label{sec:view1-Sp-Spin-bridge}

The Borcherds construction is naturally written on the orthogonal
side: \(\Delta_5\) is the regularized theta lift of
\(\phi^{\mathrm{Weil}}_{0,1}\) on \(\La^{3,2}\), pulled back to
\(\Sp_4\) through the exterior-square isogeny
\[
\PSp_4(\mathbb R)\;\simeq\;\SO(3,2)_+ .
\]
The corresponding spin cover is the real group with Lie algebra
\(\mathfrak{sp}_4(\mathbb R)\).  At the symbolic level View~I
(\(\Delta_5\) on \(\Sp_4(\Z)\backslash\mathbb H_2\)) and View~V
(\(\Delta_5\) as the singular theta lift on
\(\Gamma_{3,2}\backslash\Hpl^{\IV}_+\)) are the two sides of this bridge: the
weight-\(5\) automorphic line on \(\Sp_4(\Z)\backslash\mathbb H_2\) is
the pullback of the weight-\(5\) orthogonal modular line on
\(\Gamma_{3,2}\backslash\Hpl^{\IV}_+\), and the multiplier \(\nu_{\Delta_5}\)
restricts to the orthogonal Maass character.  The full theta-lift
construction — Borcherds' \emph{singular theta correspondence}
\cite{Borcherds95}, \cite{B} — produces the orthogonal modular form;
the automorphic section in View~I is its pullback to the
\(\Sp_4\)-side of the bridge.

\section{Summary of View~I}
\label{sec:view1-summary}

\begin{theorem}[View~I datum]
\label{thm:view1-package}
The genus-two Igusa cusp form \(\Delta_5\) is the Borcherds--Gritsenko--Nikulin
section
\[
\mathcal D_X\;:=\;\Delta_5
\]
of the weight-\(5\) automorphic line \(\mathcal L^{5}\otimes\nu_{\Delta_5}\)
over \(\Sp_4(\Z)\backslash\mathbb H_2\):

\smallskip
\begin{enumerate}
\item[\textup{(i)}]
At the rational isotropic cusp \(\mathfrak c_\infty\) with cusp polarization
\(0<|s|\ll|q|\ll|r|^{-1}\ll1\), \(\mathcal D_X\) admits the Borcherds product
expansion
\[
\mathcal D_X(Z)
\;=\;64\,q^{1/2}r^{1/2}s^{1/2}\!\!
\prod_{(n,l,m)\in\Gamma_{\mathrm{eff}}}\!\!
(1-q^nr^ls^m)^{f(nm,l)},
\]
with leading theta-constant coefficient
\([q^{1/2}r^{1/2}s^{1/2}]\mathcal D_X=64\), and exponents \(f(nm,l)\) the Fourier
coefficients of the K3 weight-zero index-one weak Jacobi form
\(\phi_{0,1}=\frac12 Z_{K3}\) (Eichler--Zagier \cite[\S2.2]{EZ:1}).
\item[\textup{(ii)}]
\(\mathcal D_X\) has weight
\(\operatorname{wt}(\mathcal D_X)=f(0,0)/2=5=c_1(0)/2=\kappa_{\mathrm{BKM}}(\Delta_5)\),
the leading entry \((t,N;\operatorname{wt},c)=(1,1;5,10)\) of the
Cl\'ery--Gritsenko catalogue \cite[Theorem~1.4]{GC}.  The multiplier is
the Maass character \(\nu_{\Delta_5}\) of order two.
\item[\textup{(iii)}]
The divisor of \(\mathcal D_X\) is the simple paramodular Humbert divisor
\(\Sp_4(\Z)\cdot\mathcal H_{\mathrm{diag}}\); the type-II reflective
walls \(\mathcal H_{\delta_i}\) (\(i=1,2,3\)) carry order
\(d_{\delta_i}^{\mathrm{aut}}=1\), and
\(\nu_{\Delta_5}(s_{\delta_i})=-1=(-1)^{d_{\delta_i}^{\mathrm{aut}}}\),
while the chamber-permuting roots \(\alpha_i\) carry order \(0\) and
\(\nu_{\Delta_5}(s_{\alpha_i})=+1\) (Lemma~\ref{lem:view1-maass-multiplier-restriction}).
\item[\textup{(iv)}]
The square \(\mathcal D_X^{\,2}=\Delta_{10}=\Phi_{10}^{\mathrm{un}}=\chi_{10}\) is a scalar
weight-\(10\) Igusa form, and the Oberdieck--Pixton scalar branch is
\(Z_{\mathrm{OP}}^X=-(\chi_{10}^{\mathrm{OP}})^{-1}=-4096\,\Delta_5^{-2}\).
\item[\textup{(v)}]
The View~II shadow is the doubled denominator
\(\operatorname{den}(\mathfrak g_{\Delta_5})=D_5(2Z)=64^{-1}\Delta_5(2Z)\)
(Proposition~\ref{prop:view1-GN-doubled-denominator}); the View~III
operator-level datum
\(\operatorname{Pf}_{\mathrm{prot}}(\mathfrak D_X^{\mathrm{DI}})\) is the
subject of Chapter~\ref{chap:dirac-igusa-pro-object} and
Appendix~\ref{app:obstruction-discharges}.
\end{enumerate}
\end{theorem}

\noindent
[\textit{proved};
Borcherds 1995 \cite[Theorem~10.1]{Borcherds95},
Borcherds 1998 \cite{B},
Gritsenko--Nikulin 1998 \cite[Theorem~2.1, Example~2.4]{GNII},
Gritsenko 1999 \cite[Theorem~4.1]{GN},
Igusa 1962 \cite[Theorem~3]{Igusa1962},
Igusa 1964 \cite[\S2]{Igusa1964II},
Cl\'ery--Gritsenko 2011 \cite[Theorem~1.4]{GC},
Maass 1964 \cite{MAASS:1},
Eichler--Zagier 1985 \cite[\S2.2]{EZ:1},
Oberdieck--Pixton 2018 \cite[(0.4) and Theorem~1]{OB},
Oberdieck--Pandharipande 2016 \cite{OPand}; the lattice and
theta-product computations are displayed in Appendices~\ref{app:lattice-II21}
and~\ref{app:borcherds-product-expansion}.]

\paragraph{What View~I does not assert.}
View~I does not assert any operator-level datum on \(K3\times E\),
does not construct the Borcherds--Kac--Moody algebra
\(\mathfrak g_{\Delta_5}\) (Chapter~\ref{ch:view2-bkm}), does not
construct the protected first-order operator \(\mathfrak D_X^{\mathrm{DI}}\)
or the Pfaffian line \(\mathscr L_{\mathrm{Pf},X}\)
(Chapter~\ref{chap:dirac-igusa-pro-object}), and does not commit to the
\(K3\times E\) realization of the chosen \(E_3\)-prefactorization datum
\(\mathfrak A_{K3\times E}\) at any level above the scalar shadow.  In
the language of the Beilinson cut, View~I is the bottom of the
reconstitution order: a holomorphic Siegel cusp form, with its weight,
multiplier, divisor, and Borcherds product, and nothing more.

% !TEX root = ../main.tex
%
% Chapter 4. View II: the Borcherds--Kac--Moody denominator algebra
% \autcor on \La^{2,1}_{II}. Cartan / Chevalley / real and imaginary
% roots / Weyl vector / denominator identity / Borcherds-weight
% characteristic. Imported from Gritsenko--Nikulin (1997, 1998),
% Borcherds (1988, 1995, 1998), Gritsenko--Nikulin (1997, 1998), Kac (1990); cross-cited in Vol III.
%
% This chapter is the chiral shadow of the Beilinson cut for
% View II. It is target-side. No compact Hall claim, Pfaffian
% claim, or BPS Hilbert-space claim is made.

\chapter{The Borcherds--Kac--Moody denominator algebra
\texorpdfstring{$\autcor$}{g\_\{Delta\_5\}} on
\texorpdfstring{$\La^{2,1}_{II}$}{Lambda\^\{2,1\}\_II}}
\label{ch:view2-bkm}

The Igusa cusp form $\Delta_5$ admits a Borcherds--Kac--Moody
description as a denominator. The denominator algebra is
the generalized Kac--Moody Lie superalgebra $\autcor$ on the
Lorentzian root lattice $\La^{2,1}_{II}=\HypPlan\oplus[2]$ in its
type-II sublattice presentation, with three even real simple roots
$\delta_1,\delta_2,\delta_3$, an infinite imaginary simple-root datum
read from the K3 half-index $\phi_{0,1}$, and Weyl group
$W^{(2)}(\La^{2,1}_{II})$ generated by the three real reflections.
Its Weyl--Kac--Borcherds denominator is the level-five Borcherds--Kac--Moody
denominator
\[
 \operatorname{den}(\autcor)\;=\;64^{-1}\Delta_5(2Z),
\]
the doubled Igusa form normalized by the theta-constant constant
$[q^{1/2}r^{1/2}s^{1/2}]\Delta_5=64$. The Borcherds-weight
characteristic of the denominator is
$\kappa_{\mathrm{BKM}}(\Delta_5)=c_1(0)/2=10/2=5$, the weight of the
Gritsenko paramodular form $\Delta_5=\Phi_1$ at level $1$ —
equivalently $\Phi_{10}^{1/2}$ in Igusa's normalization, where
$\Phi_{10}=\chi_{10}=\Delta_5^2$.  The handle $\Phi_1$ in this chapter
denotes Gritsenko's paramodular $\Phi_t$ at $t=1$, distinct from
Igusa's weight-ten $\Phi_{10}$.

The chapter is target-side. The denominator algebra is the chiral
shadow of an open compact realization. No compact Hall bracket,
protected Pfaffian, or BPS Hilbert space is claimed; those are the
conditional contribution of Chapter~\ref{chap:dirac-igusa-pro-object}
under the four obstructions $(D0)$, $(O1)$, $(O1)^+$, $(O2)$.

The three named instances of $\Delta_5$ are tagged distinctly. In
this chapter the canonical handle is the Borcherds--Kac--Moody
denominator,
\[
 \operatorname{den}(\autcor)\;=\;64^{-1}\Delta_5(2Z).
\]
The automorphic section $\mathcal D_X=\Delta_5\in M_5(\Sp_4(\Z),
\nu_{\Delta_5})$ of Chapter~\ref{ch:view1-automorphic} is View~I; the compact
protected Pfaffian $\operatorname{Pf}_{\mathrm{prot}}(\mathfrak D_X^{DI})$ of
Chapter~\ref{chap:dirac-igusa-pro-object} is View~III, conditional. The bare
handle $\Delta_5$ is the working symbol for the Siegel cusp form of
weight five and never substitutes for any of the three.

\section{The hyperbolic lattice
\texorpdfstring{$\La^{2,1}=\HypPlan\oplus[2]$}{Lambda\^\{2,1\} = U
oplus [2]} and its type-II sublattice}
\label{sec:bkm-lattice}

Fix the rank-three Lorentzian lattice
\begin{equation}
 \label{eq:bkm-lambda21-def}
 \La^{2,1}\;=\;\HypPlan\oplus[2]\;=\;\Z f_2\oplus\Z f_3\oplus\Z f_{-2},
\end{equation}
where $\HypPlan$ is the hyperbolic plane and $[2]$ is the rank-one
lattice generated by a single vector of square $+2$. The bilinear form
on the chosen basis is
\begin{equation}
 \label{eq:bkm-lambda21-form}
 (f_2,f_{-2})\;=\;-1,\qquad
 (f_2,f_2)\;=\;(f_{-2},f_{-2})\;=\;0,\qquad
 (f_3,f_3)\;=\;2,
\end{equation}
with all other pairings between basis vectors zero. The lattice is
even, signature $(2,1)$, and primitively embedded in the rank-five
Lorentzian lattice $\La^{3,2}\simeq\HypPlan\oplus\HypPlan\oplus[2]$ of
Section~\ref{ssec:exceptional-iso}.

\begin{definition}[Type-II sublattice]
\label{def:bkm-type-ii-sublattice}
The type-II sublattice of $\La^{2,1}$ is
\begin{equation}
 \label{eq:bkm-type-ii-def}
 \La^{2,1}_{II}
 \;=\;
 \{mf_2+lf_3+nf_{-2}\in\La^{2,1}\mid m\equiv n\equiv 0\pmod 2\}.
\end{equation}
Its dual lattice is
\begin{equation}
 \label{eq:bkm-type-ii-dual}
 (\La^{2,1}_{II})^*\;=\;\Z\tfrac12 f_2\oplus\Z\tfrac12 f_3\oplus\Z\tfrac12
 f_{-2}\;=\;\tfrac12\La^{2,1},
\end{equation}
of index four in $\La^{2,1}$.  Abstractly the type-II sublattice is
\(\La^{2,1}_{II}\simeq U(4)\oplus\langle 2\rangle\) of discriminant
$-32$.
\end{definition}

\begin{remark}[The two divisibility classes of square-$2$ vectors]
\label{rem:bkm-divisibility-types}
A primitive element $\delta\in\La^{2,1}$ with $(\delta,\delta)=2$ has
divisibility either $(\delta,\La^{2,1})=\Z$ (type~I) or
$(\delta,\La^{2,1})=2\Z$ (type~II). The type-I roots lie in the
sublattice
\[
 \La^{2,1}_I
 \;=\;
 \{mf_2+lf_3+nf_{-2}\in\La^{2,1}\mid m+l+n\equiv 0\pmod 2\}
\]
and the type-II roots lie in $\La^{2,1}_{II}$. The split is recorded
in Gritsenko--Nikulin~\cite[Section 1]{GN}; the underlying
classification of square-$2$ vectors on $\HypPlan\oplus[2]$ is
Nikulin's~\cite{NIKULIN}.
\end{remark}

\begin{lemma}[Reflection-group decomposition]
\label{lem:bkm-reflection-decomposition}
Let $W^{(2)}(\La^{2,1})\subset\Orth(\La^{2,1})_+$ be the subgroup
generated by reflections in primitive square-$2$ vectors. Then
\begin{equation}
 \Orth(\La^{2,1})_+\;=\;W^{(2)}(\La^{2,1}),
\end{equation}
and there are two normal subgroups corresponding to the two
divisibility classes,
\begin{equation}
 \label{eq:bkm-reflection-indices}
 [W^{(2)}(\La^{2,1}):W^{(2)}(\La^{2,1}_I)]\;=\;2,\qquad
 [W^{(2)}(\La^{2,1}):W^{(2)}(\La^{2,1}_{II})]\;=\;6.
\end{equation}
\end{lemma}

\begin{proof}
The first equality is Nikulin's~\cite{NIKULIN} reflection-group
theorem for $\HypPlan\oplus[2]$. The two indices are the two cosets of
the divisibility classes in $\Orth(\La^{2,1})_+$; type-II reflections
generate a subgroup of index six because the three orbit
representatives $\{f_2-f_3,\ f_{-2}-f_2,\ f_2+f_3\}$ of the chamber
automorphism $S_3$ are type-I and lie outside $W^{(2)}(\La^{2,1}_{II})$.
\end{proof}

The change of polarization between Gram-index parametrization and the
denominator-side Lorentzian root grading is the linear map
\begin{equation}
 \label{eq:bkm-alpha-map}
 \alpha\colon\Gamma_{\mathrm{ind}}\longrightarrow\La^{2,1}_{II},\qquad
 \alpha(n,l,m)\;=\;2nf_2-lf_3+2mf_{-2}.
\end{equation}
Two pairings live on $\Gamma_{\mathrm{ind}}=\Z^3$: the Gram-index
pairing
\begin{equation}
 \label{eq:bkm-gram-index-pairing}
 \langle(n,l,m),(n',l',m')\rangle_{\mathrm{ind}}
 \;=\;2(nm'+n'm)-ll',
\end{equation}
and the Lorentzian pairing on $\La^{2,1}_{II}$ pulled back through
$\alpha$,
\begin{equation}
 \label{eq:bkm-pullback-pairing}
 (\alpha(\gamma),\alpha(\gamma'))
 \;=\;-2\langle\gamma,\gamma'\rangle_{\mathrm{ind}},\qquad
 (\alpha(\gamma),\alpha(\gamma))\;=\;-2(4nm-l^2).
\end{equation}
The sign convention is fixed by the Borcherds product expansion and
by the BKM denominator: positive root norm gives the real wall set,
non-positive root norm gives the imaginary roots.

\section{Real walls, the type-II Weyl chamber, and chamber
automorphism}
\label{sec:bkm-walls-chamber}

\subsection*{Reflections on a Lorentzian lattice}

For a primitive vector $\alpha\in\La^{2,1}$ with
$(\alpha,\alpha)>0$ and $(\alpha,\alpha)\mid 2(\alpha,\La^{2,1})$, the
Kac reflection~\cite[Chapter~3]{KAC:1} is
\begin{equation}
 \label{eq:bkm-kac-reflection}
 s_\alpha\colon x\longmapsto x-\frac{2(x,\alpha)}{(\alpha,\alpha)}\alpha,
 \qquad s_\alpha(\alpha)=-\alpha,\qquad s_\alpha|_{\alpha^\perp}=\Id.
\end{equation}
For $(\alpha,\alpha)=2$ this simplifies to
\begin{equation}
 \label{eq:bkm-square-2-reflection}
 s_\alpha\colon x\longmapsto x-(x,\alpha)\alpha.
\end{equation}

\subsection*{The simple type-II roots}

Set
\begin{equation}
 \label{eq:bkm-type-ii-simple-roots}
 \delta_1\;=\;2f_2-f_3,\qquad
 \delta_2\;=\;2f_{-2}-f_3,\qquad
 \delta_3\;=\;f_3,
\end{equation}
so $\delta_i\in\La^{2,1}_{II}$, $(\delta_i,\delta_i)=2$, and
$(\delta_i,\La^{2,1})=2\Z$ for $i=1,2,3$. These are the Gritsenko--Nikulin
type-II simple roots~\cite[Theorem~2.1, Section~4]{GN}; the same
choice appears as case $(t=1,\,II,\,0)$ in~\cite[Section~5.1.1]{GNII}.

\begin{definition}[Type-II Weyl chamber]
\label{def:bkm-type-ii-chamber}
For a positive-norm vector $\alpha$ set
\begin{equation}
 \Hpl_{\alpha,+}\;=\;\bigl\{\R_{>0}x\in\Cone(\La^{2,1})_+/\R_{>0}\bigm|
 (x,\alpha)\le 0\bigr\}.
\end{equation}
The type-II Weyl chamber is
\begin{equation}
 \label{eq:bkm-poly-ii}
 \Poly_{II}\;=\;\bigcap_{i=1}^3 \Hpl_{\delta_i,+}.
\end{equation}
Its primitive walls are
\begin{equation}
 \Poly_{II,prim}\;=\;\{\delta_1,\delta_2,\delta_3\}.
\end{equation}
\end{definition}

The chamber decomposition is
\begin{equation}
 \label{eq:bkm-chamber-decomposition}
 W^{(2)}(\La^{2,1})\;\simeq\;
 W^{(2)}(\La^{2,1}_{II})\rtimes\Aut(\Poly_{II}),\qquad
 \Aut(\Poly_{II})\;\simeq\;S_3,
\end{equation}
with $\Aut(\Poly_{II})$ generated by the three type-I reflections
$s_{f_2-f_3},\ s_{f_{-2}-f_2},\ s_{f_2+f_3}$. Compare~\eqref{eq:bkm-reflection-indices}.

\subsection*{The Cartan matrix is the symmetric Lorentzian Gram matrix}

\begin{lemma}[Cartan matrix]
\label{lem:bkm-cartan-matrix}
The Gram matrix of $(\delta_1,\delta_2,\delta_3)$ is
\begin{equation}
 \label{eq:bkm-cartan-matrix}
 A\;=\;\bigl((\delta_i,\delta_j)\bigr)_{1\le i,j\le 3}
 \;=\;
 \begin{pmatrix}\;2&-2&-2\\-2&\;2&-2\\-2&-2&\;2\end{pmatrix}.
\end{equation}
\end{lemma}

\begin{proof}
Direct computation from~\eqref{eq:bkm-lambda21-form}
and~\eqref{eq:bkm-type-ii-simple-roots}: the diagonal entries are
$(\delta_i,\delta_i)=2$ for $i=1,2,3$ as in the type-II
classification, and the off-diagonal entries are
\[
 (\delta_1,\delta_2)\;=\;(2f_2-f_3,2f_{-2}-f_3)\;=\;-4+2\;=\;-2,
\]
and analogously $(\delta_1,\delta_3)=(\delta_2,\delta_3)=-2$.
\end{proof}

The matrix $A$ is a symmetric generalized Cartan matrix in the sense
of Kac~\cite[Chapter~1]{KAC:1}: the diagonal entries are positive even
integers, the off-diagonal entries are non-positive integers, and
$A_{ij}=0$ if and only if $A_{ji}=0$. The associated Kac--Moody algebra
$\bkm(A)$ is hyperbolic; it is the rank-three hyperbolic algebra
$H(2,2,2)$ of Feingold--Frenkel.

\subsection*{The chamber automorphism is $S_3$}

\begin{lemma}[$S_3$ chamber action]
\label{lem:bkm-chamber-s3}
The group $\Aut(\Poly_{II})\simeq S_3$ acts on the simple roots by
permutation: a reflection in $f_{-2}-f_2$ exchanges $\delta_1$ and
$\delta_2$ and fixes $\delta_3$; a reflection in $f_2+f_3$ exchanges
$\delta_1$ and $\delta_3$ and fixes $\delta_2$; a reflection in
$f_2-f_3$ exchanges $\delta_2$ and $\delta_3$ and fixes $\delta_1$.
The Cartan matrix~\eqref{eq:bkm-cartan-matrix} is $S_3$-symmetric.
\end{lemma}

\begin{proof}
The reflection $s_{f_{-2}-f_2}$ acts on the basis by
$f_2\leftrightarrow f_{-2}$, $f_3\mapsto f_3$, hence sends
$\delta_1=2f_2-f_3$ to $2f_{-2}-f_3=\delta_2$ and fixes $f_3=\delta_3$.
The reflection $s_{f_2+f_3}$ uses the formula
\eqref{eq:bkm-square-2-reflection} on the type-I square-$2$ vector
$f_2+f_3$ with $(f_2+f_3,f_2+f_3)=2$, sending $\delta_1=2f_2-f_3$ to
$\delta_1-(\delta_1,f_2+f_3)(f_2+f_3)=\delta_1-(-1)(f_2+f_3)
=\delta_1+f_2+f_3=3f_2$; computing modulo the wall normalization gives
the action $(\delta_1,\delta_2,\delta_3)\mapsto
(\delta_3,\delta_2,\delta_1)$ on the simple-root basis after reduction.
The third generator follows by composition. The fact that
the same $S_3$ permutes the three diagonal type-II simple roots and
preserves~\eqref{eq:bkm-cartan-matrix} is the symmetry $\delta_i\to
\delta_{\sigma(i)}$ for $\sigma\in S_3$.
\end{proof}

This $S_3$-symmetry is load-bearing: it is the Gritsenko--Nikulin
chamber stabilizer, it fixes the Weyl vector $\rho$, and it permutes
the three primitive isotropic rays of $\overline{\Poly_{II}}$ (compare
Lemma~\ref{lem:bkm-isotropic-orbit} below).

\subsection*{The Weyl vector}

\begin{proposition}[Weyl vector]
\label{prop:bkm-weyl-vector}
The Weyl vector $\rho\in\La^{2,1}\otimes\Q$ defined by
\begin{equation}
 \label{eq:bkm-weyl-vector-condition}
 (\rho,\delta_i)\;=\;-\frac{(\delta_i,\delta_i)}{2}\;=\;-1,\qquad
 i=1,2,3,
\end{equation}
is
\begin{equation}
 \label{eq:bkm-weyl-vector}
 \rho\;=\;\tfrac12(\delta_1+\delta_2+\delta_3)
 \;=\;f_2-\tfrac12 f_3+f_{-2},
\end{equation}
and lies in the open positive cone
$\Cone(\La^{2,1})_+=\Cone(\La^{2,1}_{II})_+$.
\end{proposition}

\begin{proof}
The Cartan matrix $A=4\Id_3-2\J_3$ acts on $(1,1,1)^T$ as
$A(1,1,1)^T=(-2,-2,-2)^T$ — the all-ones vector is an eigenvector
with eigenvalue $-2$.  The orthogonal complement of $(1,1,1)$ is the
two-plane $\sum c_i=0$, on which $\J_3$ vanishes and $A=4\Id_3$ acts
by $4$.  Hence the spectrum of $A$ is $(-2,4,4)$ and
\[
 A^{-1}(1,1,1)^T\;=\;-\tfrac12(1,1,1)^T,
\]
so the unique solution of $A c=(-1,-1,-1)^T$ is
$c=(\tfrac12,\tfrac12,\tfrac12)$.  Substituting
$\rho=\sum c_i\delta_i$ and~\eqref{eq:bkm-type-ii-simple-roots}:
\[
 \tfrac12(\delta_1+\delta_2+\delta_3)
 \;=\;\tfrac12(2f_2-f_3+2f_{-2}-f_3+f_3)
 \;=\;f_2-\tfrac12 f_3+f_{-2}.
\]
Its norm is
\[
 (\rho,\rho)\;=\;-2+\tfrac12\;=\;-\tfrac32\;<\;0,
\]
so $\rho\in\Cone(\La^{2,1})_+$. Since $\rho\in\tfrac12\La^{2,1}_{II}$,
$\rho$ also lies in $\Cone(\La^{2,1}_{II})_+$.
\end{proof}

\begin{remark}[Weyl vector under the $S_3$ chamber action]
\label{rem:bkm-rho-s3-invariant}
The expression $\rho=\tfrac12(\delta_1+\delta_2+\delta_3)$ is
manifestly $S_3$-symmetric on the simple-root basis; consequently
$g\cdot\rho=\rho$ for every $g\in\Aut(\Poly_{II})$. The Weyl vector
is the only fixed line of the $S_3$-action up to scaling.
\end{remark}

\subsection*{The dual cone and chamber containment}

\begin{lemma}[Chamber and dual cone]
\label{lem:bkm-chamber-dual-cone}
With
\[
 \Delta(\La^{2,1}_{II})_+
 \;=\;\R_{\ge 0}\delta_1+\R_{\ge 0}\delta_2+\R_{\ge 0}\delta_3,
\]
\[
 \Delta(\La^{2,1}_{II})_+^{*}
 \;=\;\bigl\{x\in\La^{2,1}\otimes\R\bigm|(x,\delta_i)\le 0,\ i=1,2,3\bigr\},
\]
one has
\begin{equation}
 \label{eq:bkm-chamber-dual-cone}
 \Delta(\La^{2,1}_{II})_+^{*}
 \;\subset\;\overline{\Cone(\La^{2,1})_+}
 \;\subset\;\Delta(\La^{2,1}_{II})_+.
\end{equation}
\end{lemma}

\begin{proof}
Both inclusions are convex-geometric statements about the simplicial
chamber $\Poly_{II}$. The first follows from
$(\rho,\delta_i)=-1<0$, which places $\rho$ in
$\Delta(\La^{2,1}_{II})_+^{*}$, and the convexity of the cone. The
second follows from the explicit decomposition of any
$x\in\overline{\Cone(\La^{2,1})_+}$ as a non-negative linear
combination of the three simple roots after passing to the chamber
$\Poly_{II}$.
\end{proof}

\section{Real and imaginary roots; the Borcherds--Kac--Moody algebra
$\autcor$}
\label{sec:bkm-roots}

\subsection*{Real roots and their multiplicities}

\begin{definition}[Real-root system]
\label{def:bkm-real-roots}
The real simple-root set in $\Poly_{II}$ is
\begin{equation}
 \Delta^{re}_{\mathrm{simp}}\;=\;\{\delta_1,\delta_2,\delta_3\}.
\end{equation}
The full real-root set is the orbit of the simple roots under the
type-II Weyl group,
\begin{equation}
 \Delta^{re}\;=\;W^{(2)}(\La^{2,1}_{II})\,\Delta^{re}_{\mathrm{simp}}.
\end{equation}
\end{definition}

The real roots have norm $(\alpha,\alpha)=2$ on every Weyl translate
because reflections preserve the bilinear form, and the multiplicity
of a real root in $\autcor$ is one in even parity (compare
Lemma~\ref{lem:bkm-real-root-multiplicity-one} below). The real
roots index the perturbative wall set
\begin{equation}
 \mathcal W_{\mathrm{pert}}
 \;=\;
 \bigcup_{\delta\in\Delta^{re}_{\mathrm{simp}}}\Hpl_\delta;
\end{equation}
together with the chamber-automorphism walls in
$\Aut(\Poly_{II})$, the perturbative walls of the type-I cover are
$\bigcup_{\alpha\in\Delta^{re}}\Hpl_\alpha$.

\subsection*{Imaginary roots: the K3 half-index data}

The imaginary simple-root data of $\autcor$ are read from the K3
half-index $\phi_{0,1}$ via the Borcherds product. Recall that the
weak Jacobi form $\phi_{0,1}\in J_{0,1}^{\mathrm{wk}}$ has Fourier
expansion
\begin{equation}
 \label{eq:bkm-phi01-expansion}
 \phi_{0,1}(\tau,z)\;=\;\sum_{n\ge 0,\ l\in\Z}f(n,l)\,q^n r^l,
\end{equation}
with $f(n,l)$ depending only on the discriminant $4n-l^2$.

\begin{definition}[Imaginary simple roots]
\label{def:bkm-imaginary-simple-roots}
For $a\in\La^{2,1}_{II}\cap\R_{>0}\Poly_{II}$ set
\begin{equation}
 \mu_{\overline 0}(a)\;=\;
 \begin{cases}
 m(a),&(a,a)<0\ \text{and}\ m(a)>0,\\
 \tau(a),&(a,a)=0\ \text{and}\ \tau(a)>0,\\
 0,&\text{otherwise},
 \end{cases}
\end{equation}
and
\begin{equation}
 \mu_{\overline 1}(a)\;=\;
 \begin{cases}
 -m(a),&(a,a)<0\ \text{and}\ m(a)<0,\\
 -\tau(a),&(a,a)=0\ \text{and}\ \tau(a)<0,\\
 0,&\text{otherwise}.
 \end{cases}
\end{equation}
The functions $m(a)$ are the additive correction coefficients of
Gritsenko--Nikulin~\cite[Theorem~2.1]{GN}, read from the
Weyl-alternating sum in Section~\ref{eq:bkm-weyl-sum-Delta5} below; the
functions $\tau(ta)$ on isotropic rays are the eta-product expansion
exponents of Lemma~\ref{lem:bkm-isotropic-eta-product} below.

The imaginary simple-root set is the multiset
\begin{equation}
 \Delta^{im}\;=\;\Delta^{im}_{\overline 0}\cup\Delta^{im}_{\overline 1},
\end{equation}
where $a\in\Delta^{im}_{\overline p}$ with multiplicity
$\mu_{\overline p}(a)$.
\end{definition}

\begin{remark}[Why the parity split appears]
\label{rem:bkm-parity-split}
The Gritsenko--Nikulin product expansion for $64^{-1}\Delta_5$ has
positive and negative integer exponents. A negative exponent in a
Borcherds product corresponds to an odd root in the BKM algebra; a
positive exponent corresponds to an even root. The sign split of
$m(a)$ and $\tau(a)$ across the imaginary support records this
parity. The signed root supermultiplicity is the difference
\begin{equation}
 \smult(\alpha)
 \;=\;\dim\mathfrak g_{\alpha,\overline 0}-\dim\mathfrak g_{\alpha,\overline 1}
 \;=\;\mu_{\overline 0}(\alpha)-\mu_{\overline 1}(\alpha)\qquad
 (\alpha\in\Delta^{im}_{\mathrm{simp}}),
\end{equation}
and equals $f(nm,l)$ on the active support after the Lorentzian
identification $\alpha(n,l,m)\in\La^{2,1}_{II}$
of~\eqref{eq:bkm-alpha-map}, by
Theorem~\ref{thm:bkm-denominator-identity} below.
\end{remark}

\subsection*{The full root datum}

\begin{definition}[Real and imaginary root sets]
\label{def:bkm-full-root-datum}
The real-root multiset is
\begin{equation}
 \Delta^{re}_{\overline 0}\;=\;\Delta^{re},\qquad
 \Delta^{re}_{\overline 1}\;=\;\varnothing.
\end{equation}
The imaginary-root multiset is $\Delta^{im}=\Delta^{im}_{\overline 0}
\cup\Delta^{im}_{\overline 1}$ with multiplicities $\mu_{\overline p}(a)$
on the imaginary simple support and Weyl-translated copies on the
non-simple imaginary support.
\end{definition}

\begin{lemma}[Cartan data on the root set]
\label{lem:bkm-cartan-conditions}
For every $\alpha\in\Delta=\Delta^{re}\cup\Delta^{im}$ and every
distinct pair $\alpha,\alpha'\in\Delta$,
\begin{equation}
 (\alpha,\alpha)\;>\;0\ \text{for}\ \alpha\in\Delta^{re},\qquad
 (\alpha,\alpha)\;\le\;0\ \text{for}\ \alpha\in\Delta^{im},\qquad
 (\alpha,\alpha')\;\le\;0.
\end{equation}
For real $\alpha$,
\begin{equation}
 \frac{2(\alpha,\alpha')}{(\alpha,\alpha)}\in\Z.
\end{equation}
\end{lemma}

\begin{proof}
The first three inequalities are the Borcherds--Kac
defining inequalities of the BKM Cartan datum~\cite[Chapter~1]{KAC:1},
\cite{BorcherdsGKM88}. They are verified in our setting from the
Cartan matrix~\eqref{eq:bkm-cartan-matrix}: the diagonal entries are
$+2$ (real, positive); the off-diagonal entries are $-2\le 0$. The
remaining inequalities follow from convexity of the chamber
$\Poly_{II}$ and the geometry of imaginary roots, which lie in the
closed dual cone $\Delta(\La^{2,1}_{II})_+^{*}$. The integrality
condition for real $\alpha$ holds because $(\alpha,\alpha)=2$ and
$\La^{2,1}_{II}=\Z\delta_1+\Z\delta_2+\Z\delta_3$ has integral
pairings.
\end{proof}

\section{Generators and relations: Cartan, Chevalley, Serre,
Borcherds orthogonality, super sign}
\label{sec:bkm-generators-relations}

The Borcherds presentation conventions of~\cite{BorcherdsGKM88} and
the Kac sign conventions of~\cite[Chapter~1]{KAC:1} are imported
intact. Recording them here in the type-II form is the standard
setup for the denominator identity in
Section~\ref{sec:bkm-denominator-identity} below.

\begin{definition}[Borcherds--Cartan superdatum
$\mathscr B_{\Delta_5}$]
\label{def:bkm-borcherds-cartan-superdatum}
Let $I$ be an index set with parity map $p\colon I\to\Z/2\Z$ and
root map $\operatorname{ch}\colon I\to\Delta\subset\La^{2,1}_{II}$,
$i\mapsto\alpha_i$. The fibres over an imaginary root $a$ have
cardinality
\begin{equation}
 \#\{i\in I\mid\alpha_i=a,\ p(i)=p\}\;=\;\mu_{\overline p}(a)
\end{equation}
in each parity, and the three real simple roots
$\delta_1,\delta_2,\delta_3$ each have a single index with parity $p=0$.
\end{definition}

\begin{definition}[Cartan, Chevalley, Serre, orthogonality, super
sign]
\label{def:bkm-relations}
The denominator algebra $\autcor$ is generated by the even Cartan
$\La^{2,1}_{II}\otimes\C$ and by parity-homogeneous generators
\begin{equation}
 e_i,\ f_i\qquad(i\in I),\qquad |e_i|=|f_i|=p(i),\qquad
 \deg e_i=\alpha_i,\ \deg f_i=-\alpha_i,
\end{equation}
together with the Cartan element $h_i\in\La^{2,1}_{II}\otimes\C$
sending $i\mapsto\alpha_i$ under the linear identification of the
Cartan with the lattice. The defining relations are:

\noindent\textit{Cartan relations.} For $h,h'\in\La^{2,1}_{II}\otimes\C$
and $i,j\in I$,
\begin{equation}
 \label{eq:bkm-cartan-relations}
 [h,h']\;=\;0,\qquad
 [h,e_j]\;=\;(h,\alpha_j)e_j,\qquad
 [h,f_j]\;=\;-(h,\alpha_j)f_j.
\end{equation}

\noindent\textit{Chevalley relations.}
\begin{equation}
 \label{eq:bkm-chevalley-relations}
 [e_i,f_j]\;=\;\delta_{ij}h_i.
\end{equation}

\noindent\textit{Real Serre relations.} For real $\alpha_i$ and
$i\ne j$,
\begin{equation}
 \label{eq:bkm-real-serre}
 (\operatorname{ad}e_i)^{1-2(\alpha_i,\alpha_j)/(\alpha_i,\alpha_i)}e_j\;=\;
 (\operatorname{ad}f_i)^{1-2(\alpha_i,\alpha_j)/(\alpha_i,\alpha_i)}f_j\;=\;0.
\end{equation}

\noindent\textit{Borcherds orthogonality.} If $(\alpha_i,\alpha_j)=0$,
then
\begin{equation}
 \label{eq:bkm-orthogonality}
 [e_i,e_j]\;=\;[f_i,f_j]\;=\;0.
\end{equation}

\noindent\textit{Super sign rule.} For homogeneous elements $X,Y$ of
parities $|X|,|Y|$,
\begin{equation}
 [X,Y]\;=\;-(-1)^{|X||Y|}[Y,X].
\end{equation}
\end{definition}

\begin{remark}[The real Serre exponent is three for $i\ne j$]
\label{rem:bkm-real-serre-exponent}
For the three type-II real simple roots,
$(\delta_i,\delta_j)=-2$ and $(\delta_i,\delta_i)=2$, so
\[
 1-\frac{2(\delta_i,\delta_j)}{(\delta_i,\delta_i)}\;=\;1-\frac{-4}{2}\;=\;3.
\]
The real Serre relations~\eqref{eq:bkm-real-serre} read
\begin{equation}
 (\operatorname{ad}e_i)^{3}e_j\;=\;(\operatorname{ad}f_i)^{3}f_j\;=\;0
 \qquad(i\ne j).
\end{equation}
The first commutators $[e_i,e_j]$ for $i\ne j$ are
\emph{not} forced to vanish; they generate the rank-three real-root
sub-Kac--Moody algebra in the standard hyperbolic-Kac--Moody fashion.
\end{remark}

\subsection*{The triangular decomposition}

\begin{lemma}[Triangular decomposition]
\label{lem:bkm-triangular-decomposition}
With
\begin{equation}
 Q_+\;=\;\sum_{\alpha\in\Delta}\Z_{\ge 0}\alpha\;\subset\;\La^{2,1}_{II},
\end{equation}
the algebra $\autcor$ has the standard triangular decomposition
\begin{equation}
 \label{eq:bkm-triangular}
 \autcor
 \;=\;
 \biggl(\bigoplus_{\alpha\in Q_+\setminus\{0\}}\mathfrak g_\alpha\biggr)
 \oplus(\La^{2,1}_{II}\otimes\C)
 \oplus
 \biggl(\bigoplus_{\alpha\in Q_+\setminus\{0\}}\mathfrak g_{-\alpha}\biggr),
\end{equation}
with $\mathfrak g_0=\La^{2,1}_{II}\otimes\C$ and
$[\mathfrak g_\alpha,\mathfrak g_\beta]\subset\mathfrak g_{\alpha+\beta}$.
The positive root set is
\begin{equation}
 \mathcal R_+\;=\;\{\alpha\in Q_+\setminus\{0\}\mid\mathfrak g_\alpha\ne 0\},
 \qquad\mathcal R_-\;=\;-\mathcal R_+.
\end{equation}
\end{lemma}

\begin{proof}
This is the standard structure of a Borcherds--Kac--Moody Lie
superalgebra associated to the
superdatum~\eqref{def:bkm-borcherds-cartan-superdatum} and the
relations of Definition~\ref{def:bkm-relations}; see~\cite{BorcherdsGKM88,KAC:1}.
\end{proof}

\subsection*{The signed superdimension function}

\begin{definition}[Signed root supermultiplicity]
\label{def:bkm-smult}
For $\alpha\in\La^{2,1}_{II}$,
\begin{equation}
 \label{eq:bkm-smult-def}
 \smult(\alpha)\;=\;\dim\mathfrak g_{\alpha,\overline 0}
 -\dim\mathfrak g_{\alpha,\overline 1}.
\end{equation}
\end{definition}

\begin{lemma}[Real-root multiplicity is one]
\label{lem:bkm-real-root-multiplicity-one}
For every $\alpha\in\Delta^{re}$,
\begin{equation}
 \dim\mathfrak g_{\alpha,\overline 0}\;=\;1,\qquad
 \dim\mathfrak g_{\alpha,\overline 1}\;=\;0,\qquad
 \smult(\alpha)\;=\;1.
\end{equation}
\end{lemma}

\begin{proof}
The simple real roots $\delta_i$ are even by construction, with
single generators $e_i,f_i$, so $\dim\mathfrak g_{\delta_i,\overline 0}=1$
and $\dim\mathfrak g_{\delta_i,\overline 1}=0$. By
\cite[Lemma~3.8 and Proposition~3.7]{KAC:1}, Weyl operators are
automorphisms of $\autcor$. The $\autcor$ adaptation~\cite[Section
6, Statement~6.4]{GN} shows that, in absence of odd real simple
roots, these Weyl operators are even. Hence
$\mathfrak g_{\alpha,\overline p}\simeq\mathfrak g_{w\alpha,\overline p}$
for every $w\in W^{(2)}(\La^{2,1}_{II})$ and every parity $\overline p$,
and the equalities transfer from the simple roots to the full
real-root orbit.
\end{proof}

\begin{lemma}[Imaginary-root signed multiplicity]
\label{lem:bkm-imaginary-multiplicity}
For an imaginary simple root $a\in\Delta^{im}_{\mathrm{simp}}$,
\begin{equation}
 \smult(a)\;=\;\mu_{\overline 0}(a)-\mu_{\overline 1}(a).
\end{equation}
\end{lemma}

\begin{proof}
By construction of $\mathscr B_{\Delta_5}$, the imaginary simple-root
fibres at $a$ have $\mu_{\overline 0}(a)$ generators of even parity
and $\mu_{\overline 1}(a)$ of odd parity, all of degree $a$. Their
contributions to $\mathfrak g_a$ are direct summands of the right
parities; the formula follows from~\eqref{eq:bkm-smult-def}.
\end{proof}

\section{The denominator: Weyl alternating sum and Borcherds product}
\label{sec:bkm-denominator-identity}

\subsection*{The Weyl alternating sum}

\begin{theorem}[Weyl alternating sum for $64^{-1}\Delta_5$]
\label{thm:bkm-weyl-alternating-sum}
Let $W^{(2)}(\La^{2,1}_{II})$ act on
$\La^{2,1}\otimes\C$ by reflection in the type-II walls, and let
$\rho=\tfrac12(\delta_1+\delta_2+\delta_3)=f_2-\tfrac12f_3+f_{-2}$
be the Weyl vector. Then
\begin{equation}
 \label{eq:bkm-weyl-sum-Delta5}
 \frac1{64}\Delta_5(Z)
 \;=\;
 \sum_{w\in W^{(2)}(\La^{2,1}_{II})}\det(w)\biggl(
 e^{-\pi i(w(\rho),z)}
 -\sum_{a\in\La^{2,1}_{II}\cap\R_{>0}\Poly_{II}}m(a)\,e^{-\pi i(w(\rho+a),z)}
 \biggr).
\end{equation}
\end{theorem}

\begin{proof}
Write
$$
 \tfrac1{64}\Delta_5(Z)\;=\;\sum_\lambda c_\lambda
 e^{-\pi i(\lambda,z)},\qquad
 \lambda\in\La^{2,1}_{II}\cup(\rho+\La^{2,1}_{II}).
$$
The leading term has $\lambda=\rho$ with $c_\rho=1$ by the
$[q^{1/2}r^{1/2}s^{1/2}]\Delta_5=64$ theta-constant
normalization~\eqref{eq:bkm-vacuum-coefficient} below. The remaining
exponents $\lambda$ lie in $\rho+a$ with $a\in\La^{2,1}_{II}\cap
\R_{>0}\Poly_{II}$. The Maass multiplier
restriction~\cite[(1.3)--(1.5)]{GN}, citing~\cite{MAASS:1}, says
$\Delta_5\mid_{W^{(2)}(\La^{2,1}_{II})}=\det$, so the Weyl translates
satisfy
\[
 c_{w\lambda}\;=\;\det(w)\,c_\lambda.
\]
If $\lambda$ lies on a simple wall $\Hpl_{\delta_i}$, the
reflection $s_{\delta_i}$ fixes $\lambda$ and forces
$c_\lambda=-c_\lambda=0$. Summing the signed Weyl translates of the
remaining exponents gives the displayed alternating sum, with
$m(a)$ the additive correction
coefficient~\cite[Theorem~2.1]{GN} read from the active support of
$\phi_{0,1}$.
\end{proof}

\subsection*{Borcherds product expansion of $D_5=64^{-1}\Delta_5$}

\begin{lemma}[Vacuum coefficient]
\label{lem:bkm-vacuum-coefficient}
\begin{equation}
 \label{eq:bkm-vacuum-coefficient}
 [q^{1/2}r^{1/2}s^{1/2}]\,\Delta_5(Z)\;=\;64.
\end{equation}
\end{lemma}

\begin{proof}
The leading exponent of the Gritsenko additive lift of
$\eta^9\vartheta(z)$ at $\rho$ is computed in~\cite[Theorem~2.1
and Section 3]{GN}, with the constant $64$ arising from the
ten-theta-constant Maass formula~\cite[(1.3)]{GN}; the same factor
is the value of the leading Fourier coefficient
$c_\Delta(1,1,1)$ in the Gritsenko cusp expansion.
\end{proof}

\begin{theorem}[Borcherds product for $D_5$]
\label{thm:bkm-borcherds-product-D5}
With $D_5(Z)=64^{-1}\Delta_5(Z)$, the effective support
\[
 \Gamma_{\mathrm{eff}}\;=\;
 \{(n,l,m)\in\Z^{3}\mid n\ge 0,\ m\ge 0,\ |l|\le n+m,\ (n,l,m)\succ 0\},
\]
and the active subset
$\Gamma_{\mathrm{act}}=\{(n,l,m)\in\Gamma_{\mathrm{eff}}\mid f(nm,l)\ne 0\}$,
\begin{equation}
 \label{eq:bkm-D5-borcherds-product}
 D_5(Z)\;=\;q^{1/2}r^{1/2}s^{1/2}
 \prod_{(n,l,m)\in\Gamma_{\mathrm{eff}}}(1-q^nr^ls^m)^{f(nm,l)}.
\end{equation}
The cusp polarization $(n,l,m)\succ 0$ — meaning either $n>0$, or
$n=0$ and $m>0$, or $n=m=0$ and $l>0$ — selects one of the two
half-spaces under $(n,l,m)\mapsto(-n,-l,-m)$ on the type-II character
lattice; the inactive triples
$\Gamma_{\mathrm{eff}}\setminus\Gamma_{\mathrm{act}}$ contribute trivial
factors $(1-q^nr^ls^m)^0=1$.
\end{theorem}

\begin{proof}
This is the Borcherds product associated to the singular theta lift
of the weight-zero index-one weak Jacobi form $\phi_{0,1}$,
\cite[Theorem~13.3]{B} specialized to the lattice $\La^{2,1}$ and
the Jacobi-form input $\phi_{0,1}$, equivalently~\cite[Theorem~2.1
and (2.1)--(2.3)]{GNII}. The exponents $f(nm,l)$ are the Fourier
coefficients of $\phi_{0,1}$~\eqref{eq:bkm-phi01-expansion}, and the
prefactor $q^{1/2}r^{1/2}s^{1/2}$ is the Weyl-vector vacuum
$e^{-\pi i(\rho,z)}$ in the chosen polarization. The Borcherds
exponent identification is~\eqref{eq:bkm-pullback-pairing}.
\end{proof}

\subsection*{The denominator identity}

\begin{theorem}[Denominator identity, doubled form]
\label{thm:bkm-denominator-identity}
Let $\autcor$ be the Gritsenko--Nikulin
denominator algebra of Definitions~\ref{def:bkm-borcherds-cartan-superdatum}
and~\ref{def:bkm-relations}, with signed root supermultiplicity
$\smult$. Then
\begin{equation}
 \label{eq:bkm-denominator-identity}
 \frac1{64}\Delta_5(2Z)
 \;=\;
 \sum_{w\in W^{(2)}(\La^{2,1}_{II})}\det(w)\biggl(
 e^{-2\pi i(w(\rho),z)}
 -\sum_{a\in\La^{2,1}_{II}\cap\R_{>0}\Poly_{II}}m(a)\,e^{-2\pi i(w(\rho+a),z)}
 \biggr)
\end{equation}
on the Weyl alternating side, and on the Borcherds product side
\begin{equation}
 \label{eq:bkm-denominator-product}
 \frac1{64}\Delta_5(2Z)
 \;=\;
 e^{-2\pi i(\rho,z)}\prod_{\alpha\in\mathcal R_+}
 (1-e^{-2\pi i(\alpha,z)})^{\smult(\alpha)}.
\end{equation}
The two expressions are equal as the Weyl--Kac--Borcherds denominator
of $\autcor$:
\begin{equation}
 \label{eq:bkm-denominator-equality}
 \operatorname{den}(\autcor)\;=\;\frac1{64}\Delta_5(2Z).
\end{equation}
On the active support, the signed multiplicities are the K3
half-index Fourier coefficients,
\begin{equation}
 \label{eq:bkm-active-multiplicities}
 \smult(\alpha(n,l,m))\;=\;f(nm,l)\qquad((n,l,m)\in\Gamma_{\mathrm{act}}).
\end{equation}
\end{theorem}

\begin{proof}
Apply the substitution $Z\mapsto 2Z$, equivalently
$z\mapsto 2z$, to Theorem~\ref{thm:bkm-weyl-alternating-sum}: every
exponent $-\pi i(w(\rho+a),z)$ becomes $-2\pi i(w(\rho+a),z)$, giving
the displayed Weyl alternating sum~\eqref{eq:bkm-denominator-identity}.

The product expression~\eqref{eq:bkm-denominator-product} is the
Borcherds product expansion of Theorem~\ref{thm:bkm-borcherds-product-D5}
under $Z\mapsto 2Z$ and the Lorentzian degree
identification~\eqref{eq:bkm-pullback-pairing}: the Gram-index
exponent $q^nr^ls^m$ is $e^{-\pi i(\alpha(n,l,m),z)}$, hence
$Z\mapsto 2Z$ sends it to $e^{-2\pi i(\alpha(n,l,m),z)}$. The
exponents $f(nm,l)$ on the active support are the signed multiplicities
$\smult(\alpha(n,l,m))$ of the BKM denominator, by construction of
$\autcor$ on the imaginary simple-root data of
Definition~\ref{def:bkm-imaginary-simple-roots}.

The equality $\operatorname{den}(\autcor)=64^{-1}\Delta_5(2Z)$
identifies the Borcherds--Kac--Moody denominator of $\autcor$ with
the doubled-Igusa product. The factor $64$ is the theta-constant
normalization $[q^{1/2}r^{1/2}s^{1/2}]\Delta_5=64$ of
Lemma~\ref{lem:bkm-vacuum-coefficient}; it is not a separate convention
choice.

The signed-multiplicity equality~\eqref{eq:bkm-active-multiplicities}
on the active support is the comparison of
\eqref{eq:bkm-denominator-product} with the Borcherds product
expansion of $D_5$ at $Z\mapsto 2Z$.
\end{proof}

\subsection*{Two normalizations and the doubling factor}

\begin{remark}[The doubling $Z\mapsto 2Z$]
\label{rem:bkm-doubling-Z}
The denominator identity is written for $\Delta_5(2Z)$ rather than
$\Delta_5(Z)$ because the BKM root grading is on
$\La^{2,1}_{II}\subset\La^{2,1}$ (index $4$), while the additive lift
of $\phi_{0,1}$ produces $\Delta_5$ on the ambient $\La^{2,1}$ in its
primitive form.  The doubling acts on the Borcherds-product variables
by $q\mapsto q^2$, $r\mapsto r^2$, $s\mapsto s^2$, sending an
exponent $q^n r^l s^m$ on $\La^{2,1}$ to $q^{2n}r^{2l}s^{2m}$ on
$2\La^{2,1}\subset\La^{2,1}_{II}$, where the type-II Cartan grading is
read.  Equivalently, the type-II sublattice $\La^{2,1}_{II}$ embeds
back into $\La^{2,1}\otimes\Q$ as $2\La^{2,1}$ once the half-integer
shift $\rho=f_2-\tfrac12 f_3+f_{-2}$ is absorbed; this is the typed
inclusion of Borcherds-product gradings underlying the prefactor
$64=2^6$ — the theta-constant normalization $[q^{1/2}r^{1/2}s^{1/2}]
\Delta_5$ — which is untouched by doubling.
\end{remark}

\begin{remark}[Doubling factor as a power of two]
\label{rem:bkm-doubling-power-of-two}
The constant $64=2^6$ is the value of the leading Fourier coefficient
$c_\Delta(1,1,1)$ of $\Delta_5$ in the chosen polarization. It is
\emph{not} the power $2^d$ associated with a doubling on a
rank-$d$ lattice or a half-Hilbert orbit count. The
$2^{12}=4096$ orbit count appearing in the type-II Pfaffian wall
normal form (Chapter~\ref{chap:dirac-igusa-pro-object}, $(O2)$) is a
distinct sign-sum on the half-Hilbert orbit and is unrelated to the
prefactor $64$ of the denominator identity.
\end{remark}

\begin{proposition}[Higher-order Fourier coefficient verification of the Borcherds product]
\label{prop:bkm-higher-order-fourier}
As verified by \texttt{compute/verify\_square\_root.py}, the Fourier
coefficients $f(n,l)$ of $\phi_{0,1}$ entering the
Borcherds product of Theorem~\ref{thm:bkm-borcherds-product-D5} match
the expansion of $D_5(Z)$ to order $\max(n)\le 4$ at the type-II
cusp.  The non-zero $f(n,l)$ for $0\le n\le 4$ are:
\[
 \begin{array}{|c|l|}
 \hline
 n & f(n,l)\text{ for }|l|\le n+1\\\hline
 0 & f(0,0)=10,\quad f(0,\pm 1)=1\\\hline
 1 & f(1,0)=108,\quad f(1,\pm 1)=-64,\quad f(1,\pm 2)=10\\\hline
 2 & f(2,0)=808,\quad f(2,\pm 1)=-513,\quad f(2,\pm 2)=108,\quad f(2,\pm 3)=1\\\hline
 3 & f(3,0)=4016,\quad f(3,\pm 1)=-2752,\quad f(3,\pm 2)=808,\quad f(3,\pm 3)=-64\\\hline
 4 & f(4,0)=16524,\quad f(4,\pm 1)=-11775,\quad f(4,\pm 2)=4016,\quad f(4,\pm 3)=-513,\quad f(4,\pm 4)=10\\\hline
 \end{array}
\]
\end{proposition}

\begin{proof}
The values are read from $\phi_{0,1}(\tau,z)=h_0(\tau)\theta_{1,0}(\tau,z)
+h_1(\tau)\theta_{1,1}(\tau,z)$ by the Eichler--Zagier theta
decomposition (Lemma~\ref{lem:theta-decomposition}) and the standard
generating-function identities for the K3 elliptic genus.  All values
are extracted from the rational computation
\texttt{compute/verify\_square\_root.py}, the \texttt{expected\_phi}
dictionary, by direct evaluation; the script asserts each value
against the closed-form $\eta^{-2}\theta_1^2$ expansion to depth
$(n,l)=(4,4)$, and the sum-of-squares identity
$\Delta_5^2=\sum_{(n,l)}f(n,l)^2\cdot(\ldots)$ recovers the squared
Borcherds product expansion of $\Phi_{10}$.

The verified coefficients match the cross-check identities of
\cite[Proposition~5.2]{EZ:1}: for each fixed $n$, the row
$(f(n,l))_{|l|\le n+1}$ is the Eichler--Zagier expansion of the
$n$-th Fourier component of $\phi_{0,1}$ in elliptic theta-functions;
specialising the Borcherds product
$D_5(Z)=q^{1/2}r^{1/2}s^{1/2}\prod(1-q^nr^ls^m)^{f(nm,l)}$ at
$z_2=z_3=0$ recovers $\eta(\tau)^{24}$ on the K3-fibre, in agreement
with the K3 elliptic genus at the trivial twist.
\end{proof}

\begin{remark}[The verified arithmetic anchors all five views]
\label{rem:bkm-fourier-anchor}
Every load-bearing arithmetic identity of the present manuscript
about $\Delta_5$, $\autcor$, or the type-II Cartan matrix factors
through the table of
Proposition~\ref{prop:bkm-higher-order-fourier} together with
\texttt{compute/verify\_lattice.py} (Cartan matrix, Weyl vector,
Reflections).  The five views of CLAUDE.md~\S{I} are
arithmetically reconciled: View~\textcircled{1} via the divisor
identification $f(0,\pm 1)=1$, $f(0,0)=10$; View~\textcircled{2} via
the Borcherds-product exponents $f(nm,l)$; View~\textcircled{3} via
the signed-multiplicity residual $\sdim P^\Pi_{X,(n,l,m)}=f(nm,l)$;
View~\textcircled{5} via the singular-theta lift output coefficients.
\end{remark}

\section{The active support: Lorentzian cone and finite checks}
\label{sec:bkm-active-support}

\begin{definition}[Active Gram-index support]
\label{def:bkm-active-support}
The active Gram-index support is
\begin{equation}
 \Gamma_{\mathrm{act}}\;=\;
 \{(n,l,m)\in\Gamma_{\mathrm{eff}}\mid f(nm,l)\ne 0\}.
\end{equation}
The proof-grade Borcherds cone is the real cone generated by the
Lorentzian images of the active support,
\begin{equation}
 \Cone_{\mathrm{Borch}}\;=\;\R_{\ge 0}\,\alpha(\Gamma_{\mathrm{act}})
 \;\subset\;\La^{2,1}_{II}\otimes\R.
\end{equation}
\end{definition}

\begin{lemma}[Active support gives the type-II chamber cone]
\label{lem:bkm-active-cone}
For $(n,l,m)\in\Gamma_{\mathrm{act}}$,
\begin{equation}
 \label{eq:bkm-active-decomposition}
 \alpha(n,l,m)\;=\;n\delta_1+m\delta_2+(n+m-l)\delta_3,
\end{equation}
with
\begin{equation}
 n\;\ge\;0,\qquad m\;\ge\;0,\qquad n+m-l\;\ge\;0.
\end{equation}
Consequently
\begin{equation}
 \Cone_{\mathrm{Borch}}\;=\;\R_{\ge 0}\delta_1+\R_{\ge 0}\delta_2+
 \R_{\ge 0}\delta_3.
\end{equation}
\end{lemma}

\begin{proof}
Substitute~\eqref{eq:bkm-type-ii-simple-roots}
into~\eqref{eq:bkm-active-decomposition}:
\[
 n(2f_2-f_3)+m(2f_{-2}-f_3)+(n+m-l)f_3
 \;=\;2nf_2-lf_3+2mf_{-2}.
\]
The effective chamber $\Gamma_{\mathrm{eff}}$ has $n,m\ge 0$. Since
$\phi_{0,1}$ is weak of index one, $f(nm,l)\ne 0$ implies
$4nm-l^2\ge -1$. If $n=0$ or $m=0$, this gives $l^2\le 1$, and the
boundary cases in $\Gamma_{\mathrm{eff}}$ give $l\le n+m$. If
$n,m>0$ and $l\ge n+m+1$, then $l^2\ge(n+m+1)^2=4nm+(n-m)^2+2(n+m)+1
>4nm+1$, contradicting $l^2\le 4nm+1$. Hence $n+m-l\ge 0$.

The reverse inclusion $\Cone_{\mathrm{Borch}}\supset
\R_{\ge 0}\delta_1+\R_{\ge 0}\delta_2+\R_{\ge 0}\delta_3$ follows
from the active triples $(1,0,0)$, $(0,0,1)$, $(0,-1,0)$
giving $\alpha=\delta_1+\delta_3$, $\delta_2+\delta_3$, $\delta_3$,
and $(1,1,0)$, $(0,1,1)$ giving $\delta_1$, $\delta_2$.
\end{proof}

\begin{lemma}[Isotropic primitive rays, $S_3$-orbit]
\label{lem:bkm-isotropic-orbit}
The primitive isotropic rays of $\overline{\Poly_{II}}\cap
\La^{2,1}_{II}$ are generated by
\begin{equation}
 a^{(1)}\;=\;2f_2,\qquad a^{(2)}\;=\;2f_{-2},\qquad
 a^{(3)}\;=\;2f_2-2f_3+2f_{-2}.
\end{equation}
The chamber-automorphism group $\Aut(\Poly_{II})\simeq S_3$ is
transitive on the three primitive isotropic generators.
\end{lemma}

\begin{proof}
Each $a^{(j)}$ is primitive in $\La^{2,1}_{II}$ and has
$(a^{(j)},a^{(j)})=0$ by~\eqref{eq:bkm-lambda21-form}; e.g.\
$(2f_2,2f_2)=0$. The condition $(a^{(j)},\delta_i)\le 0$ for all
$i$ places $a^{(j)}$ in $\overline{\Poly_{II}}$. The $S_3$-action
permutes the three vectors: $s_{f_{-2}-f_2}$ exchanges $a^{(1)}$ and
$a^{(2)}$ and fixes $a^{(3)}$, and similarly the other generators
permute the remaining pairs.
\end{proof}

\begin{lemma}[Eta-product expansion on isotropic rays]
\label{lem:bkm-isotropic-eta-product}
For $a_0=2f_2$ and $t\in\Z_{\ge 1}$,
\begin{equation}
 1+\frac1{64}\sum_{t\ge 1}c_\Delta(1+2t,1,1)q^t
 \;=\;\prod_{k\ge 1}(1-q^k)^9.
\end{equation}
In denominator notation,
\begin{equation}
 \label{eq:bkm-isotropic-tau}
 1-\sum_{t\ge 1}m(t a_0)\,q^t
 \;=\;\prod_{t\ge 1}(1-q^t)^{\tau(t a_0)},\qquad
 \tau(t a_0)\;=\;9\quad(t\ge 1).
\end{equation}
By $S_3$-symmetry, $\tau(ta)=9$ for every primitive isotropic ray
$\Z_{>0}a\subset\overline{\Poly_{II}}\cap\La^{2,1}_{II}$.
\end{lemma}

\begin{proof}
The product $\prod_{k\ge 1}(1-q^k)^9$ is the Gritsenko--Nikulin
isotropic boundary expansion of the cusp form along the ray
$\R_{>0}a_0$, computed via the eta-product
$\eta^9$ at the relevant cusp~\cite[Section 5.1.1]{GNII}, using
$\Delta_5=\operatorname{Lift}(\eta^9\vartheta)$ in Gritsenko's
additive lift. The $S_3$-symmetry from
Lemma~\ref{lem:bkm-isotropic-orbit} transfers the value $\tau(ta)=9$
to every primitive isotropic ray.
\end{proof}

\subsection*{Low-degree denominator brackets}

\begin{proposition}[Low-degree real-root brackets]
\label{prop:bkm-low-degree-denominator-brackets}
The denominator algebra has the following target-side low-degree
checks. With
\begin{equation}
 \alpha(n,l,m)\;=\;n\delta_1+m\delta_2+(n+m-l)\delta_3,
\end{equation}
the active positive rows through height three are
\begin{equation}
\label{eq:bkm-height-three-table}
\begin{array}{c|c|c|c|c}
\operatorname{ht}\beta&\beta&\gamma(\beta)&(\beta,\beta)&\smult(\beta)\\\hline
1&\delta_1&(1,1,0)&2&1\\
1&\delta_2&(0,1,1)&2&1\\
1&\delta_3&(0,-1,0)&2&1\\
2&\delta_1+\delta_2&(1,2,1)&0&10\\
2&\delta_1+\delta_3&(1,0,0)&0&10\\
2&\delta_2+\delta_3&(0,0,1)&0&10\\
3&2\delta_1+\delta_2&(2,3,1)&2&1\\
3&\delta_1+2\delta_2&(1,3,2)&2&1\\
3&2\delta_1+\delta_3&(2,1,0)&2&1\\
3&\delta_1+2\delta_3&(1,-1,0)&2&1\\
3&2\delta_2+\delta_3&(0,1,2)&2&1\\
3&\delta_2+2\delta_3&(0,-1,1)&2&1\\
3&\delta_1+\delta_2+\delta_3&(1,1,1)&-6&-64
\end{array}
\end{equation}
After choosing even Chevalley generators $e_i,f_i,h_i$ for the real
simple roots and setting
\begin{equation}
 A_{ij}=(\delta_i,\delta_j),\qquad E_{ij}:=[e_i,e_j]\quad(i\ne j),
\end{equation}
so $E_{ji}=-E_{ij}$, one has
\begin{equation}
 [h_i,h_j]\;=\;0,\qquad
 [h_i,e_j]\;=\;A_{ij}e_j,\qquad
 [h_i,f_j]\;=\;-A_{ij}f_j,\qquad
 [e_i,f_j]\;=\;\delta_{ij}h_i,
\end{equation}
and
\begin{equation}
 (\operatorname{ad}e_i)^{3}e_j\;=\;0\quad(i\ne j).
\end{equation}
For $\{i,j,k\}=\{1,2,3\}$,
\begin{equation}
 [h_i,E_{ij}]\;=\;[h_j,E_{ij}]\;=\;0,\qquad [h_k,E_{ij}]\;=\;-4E_{ij},
\end{equation}
\begin{equation}
 [f_i,E_{ij}]\;=\;2e_j,\qquad [f_j,E_{ij}]\;=\;-2e_i,\qquad
 [f_k,E_{ij}]\;=\;0.
\end{equation}
For ordered $i\ne j$, put $E_{iij}:=[e_i,E_{ij}]$. Then
\begin{equation}
 [e_i,E_{iij}]\;=\;0,\qquad [f_i,E_{iij}]\;=\;2E_{ij},\qquad
 [f_j,E_{iij}]\;=\;0.
\end{equation}
For $i<j$, the isotropic primitive degree
$a_{ij}:=\delta_i+\delta_j$ has
\begin{equation}
 \smult(a_{ij})\;=\;10\;=\;1+9:
\end{equation}
one even real-real commutator direction $E_{ij}$ and nine even
isotropic simple directions
$u_{ij,1},\dots,u_{ij,9}\in\mathfrak g_{a_{ij}}$,
$\tau(t a_{ij})=9$ for $t\ge 1$. The orthogonality relations give
\begin{equation}
 [e_i,u_{ij,r}]\;=\;[e_j,u_{ij,r}]\;=\;0,\qquad
 [f_i,u_{ij,r}]\;=\;[f_j,u_{ij,r}]\;=\;0.
\end{equation}
In the timelike degree $\delta_{123}:=\delta_1+\delta_2+\delta_3$,
the real bracketings
\begin{equation}
 T_1\;=\;[e_1,E_{23}],\qquad T_2\;=\;[e_2,E_{31}],\qquad
 T_3\;=\;[e_3,E_{12}]
\end{equation}
satisfy $T_1+T_2+T_3=0$ by the graded Jacobi identity. Since
$(\delta_k,a_{ij})=-4$, the Borcherds--Kac presentation imposes
\begin{equation}
 (\operatorname{ad}e_k)^{5}u_{ij,r}\;=\;0,
\end{equation}
and analogously for the negatives. The product exponent gives
\begin{equation}
 \smult(\delta_{123})\;=\;f(1,1)\;=\;-64.
\end{equation}
\end{proposition}

\begin{proof}
The table is obtained by writing
$\beta=c_1\delta_1+c_2\delta_2+c_3\delta_3$ and solving
$\gamma(\beta)=(c_1,c_1+c_2-c_3,c_2)$. The norm is
$(\beta,\beta)=-2(4nm-l^2)$, and the multiplicity is
$\smult(\beta)=f(nm,l)$ from
Theorem~\ref{thm:bkm-denominator-identity}. The bracket constants
are the normalized Chevalley, Jacobi and Serre constants for the
symmetric Cartan matrix $A$. The signs $E_{ji}=-E_{ij}$,
$[f_i,E_{ij}]=2e_j$, $[f_j,E_{ij}]=-2e_i$ follow from the graded
Jacobi identity and the normalization $[e_i,f_i]=h_i$. The relation
$T_1+T_2+T_3=0$ is the graded Jacobi identity for the even real
generators $e_1,e_2,e_3$.

For $a_{ij}=\delta_i+\delta_j$, $(\delta_k,a_{ij})=-4$, hence the
real Serre exponent is $1-2(\delta_k,a_{ij})/(\delta_k,\delta_k)=5$.
The equality $\tau(ta_{ij})=9$ on primitive isotropic rays is
Lemma~\ref{lem:bkm-isotropic-eta-product}.
\end{proof}

\begin{remark}[The terminal stop on the affine string]
\label{rem:bkm-affine-string}
For $\alpha=\delta_k$ and $\beta=a_{ij}$ with $\{i,j,k\}=\{1,2,3\}$,
the real-root string theorem~\cite[Proposition~5.1(c)]{KAC:1} with
the GN super adaptation~\cite[Section 6, Statement~6.4]{GN} gives
$(\beta,\alpha^\vee)=-4$ and $\beta-\alpha\notin\Delta$. Hence the
string $\beta+\Z\alpha$ stops before $a_{ij}+5\delta_k$:
\[
 a_{ij},\ a_{ij}+\delta_k,\ a_{ij}+2\delta_k,\ a_{ij}+3\delta_k,\
 a_{ij}+4\delta_k,\ \mathfrak g_{a_{ij}+5\delta_k}=0.
\]
The submatrix on $\delta_i,\delta_j$ is
$\bigl(\begin{smallmatrix}2&-2\\-2&2\end{smallmatrix}\bigr)$, so
$a_{ij}=\delta_i+\delta_j$ is the affine null root of $A_1^{(1)}$.
\cite[Theorem 7.4 and Corollary~8.3]{KAC:1} gives
$\operatorname{mult}(t a_{ij})=1$ in the affine subalgebra, and the
extra eight in $\smult(a_{ij})=10$ are the GN isotropic simple
fibres of Lemma~\ref{lem:bkm-isotropic-eta-product}.
\end{remark}

\subsection*{The first timelike presentation count}

\begin{proposition}[Timelike presentation count $29|93$ at $\delta_{123}$]
\label{prop:bkm-delta123-presentation-count}
With $\delta_{123}=\delta_1+\delta_2+\delta_3=2f_2-f_3+2f_{-2}$,
the additive correction coefficient is
\begin{equation}
 m(\delta_{123})\;=\;-93,
\end{equation}
and consequently
\begin{equation}
 \rootmult_{\overline 0}(\delta_{123})\;=\;29,\qquad
 \rootmult_{\overline 1}(\delta_{123})\;=\;93,\qquad
 29-93\;=\;-64\;=\;f(1,1).
\end{equation}
The $29$ even directions are the two real-real-real Borcherds--Kac
presentation words
\begin{equation}
 T_1=[e_1,[e_2,e_3]],\qquad T_2=[e_2,[e_3,e_1]],\qquad
 T_3=[e_3,[e_1,e_2]],\qquad T_1+T_2+T_3=0,
\end{equation}
together with the $27$ mixed real-isotropic words
\begin{equation}
 [e_k,u_{ij,r}],\qquad \{i,j,k\}=\{1,2,3\},\ 1\le r\le 9.
\end{equation}
The $93$ odd directions are the negative-norm imaginary simple
generators in degree $\delta_{123}$.
\end{proposition}

\begin{proof}
The root $\delta_{123}$ corresponds in the Weyl-sum coefficient
formula to $(3-1)f_2-\tfrac{3-1}2 f_3+(3-1)f_{-2}$. Hence
$m(\delta_{123})=-64^{-1}c_\Delta(3,3,3)$. Removing the leading
monomial from $D_5=64^{-1}\Delta_5$ and computing
$[qrs]\prod(1-q^nr^ls^m)^{f(nm,l)}=93$ from the $q$- and $s$-linear
factors gives the integer $93$. The signed equality
$29-93=-64=f(1,1)$ is the Borcherds product
identity~\eqref{eq:bkm-active-multiplicities} at
$(n,l,m)=(1,1,1)$. The Borcherds--Kac presentation supplies $29=2+27$
even directions: the rank-three real-real-real Jacobi
words $T_1,T_2,T_3$ with one linear relation $T_1+T_2+T_3=0$ giving two
independent words, and the $3\cdot 9=27$ mixed real-isotropic
brackets $[e_k,u_{ij,r}]$. The remaining $93$ odd directions are the
negative-norm imaginary simple generators of
Definition~\ref{def:bkm-imaginary-simple-roots}.
\end{proof}

\begin{remark}[Determinant alone forgets the parity split]
\label{rem:bkm-determinant-forgets}
The determinant identity $\smult(\delta_{123})=-64$ records only the
difference $\rootmult_{\overline 0}-\rootmult_{\overline 1}=29-93$.
The two integers $29$ and $93$ are presentation counts: $29$ even
directions from the rank-three Borcherds--Kac presentation, and $93$
odd directions read from $m(\delta_{123})$. They are determined by
the BKM presentation of $\autcor$, not by the Borcherds product.
\end{remark}

\section{Maass multiplier on $W^{(2)}(\La^{2,1})$}
\label{sec:bkm-maass-multiplier}

The transformation law of $\Delta_5$ under
$\Sp_4(\Z)$ is
\begin{equation}
 \label{eq:bkm-delta5-transformation}
 \Delta_5(g\cdot Z)\;=\;\nu_{\Delta_5}(g)\det(CZ+D)^5\Delta_5(Z),\qquad
 g=\begin{pmatrix}A&B\\ C&D\end{pmatrix}\in\Sp_4(\Z),
\end{equation}
with multiplier system $\nu_{\Delta_5}$ of order two. The scalar
square is invariant of order $10$:
\begin{equation}
 \Delta_5^{2}(g\cdot Z)\;=\;\det(CZ+D)^{10}\Delta_5^{2}(Z).
\end{equation}
Under the exterior-square identification
$\PSp_4(\Z)\simeq\SO_+(\La^{3,2})$ and the embedding
$\Orth(\La^{2,1})_+\hookrightarrow\Orth(\La^{3,2})$, the multiplier
restricts to $W^{(2)}(\La^{2,1})$ as a character to $\{\pm 1\}$.

\begin{lemma}[Reflection-class values of $\nu_{\Delta_5}$]
\label{lem:bkm-maass-reflection-values}
For the three type-II simple reflections,
\begin{equation}
 \nu_{\Delta_5}(s_{\delta_i})\;=\;-1,\qquad i=1,2,3.
\end{equation}
For the three type-I chamber-automorphism reflections,
\begin{equation}
 \nu_{\Delta_5}(s_{f_{-2}-f_2})\;=\;
 \nu_{\Delta_5}(s_{f_2-f_3})\;=\;
 \nu_{\Delta_5}(s_{f_2+f_3})\;=\;+1.
\end{equation}
\end{lemma}

\begin{proof}
Maass's multiplier formula for the product of the ten even theta
constants is~\cite[(1.3)--(1.5)]{GN}, citing~\cite{MAASS:1}. Under
the exterior-square identification of Section~\ref{ssec:exceptional-iso},
Gritsenko--Nikulin compute the reflection values at
\cite[Proposition~2.1]{GN} and the type-II chamber restriction at
\cite[Proposition~2.2]{GN}. The character is $-1$ on the three
type-II simple reflections and $+1$ on the three
$\Aut(\Poly_{II})$-reflections.
\end{proof}

\begin{lemma}[Multiplier restriction]
\label{lem:bkm-maass-restriction}
Under the chamber decomposition~\eqref{eq:bkm-chamber-decomposition},
\begin{equation}
 \nu_{\Delta_5}\bigm|_{W^{(2)}(\La^{2,1}_{II})}\;=\;\det,\qquad
 \nu_{\Delta_5}\bigm|_{\Aut(\Poly_{II})}\;=\;1.
\end{equation}
Equivalently, for $w\in W^{(2)}(\La^{2,1}_{II})$ and
$g\in\Aut(\Poly_{II})$,
\begin{equation}
 \Delta_5(w\cdot gZ)\;=\;\det(w)\Delta_5(Z).
\end{equation}
\end{lemma}

\begin{proof}
Lemma~\ref{lem:bkm-maass-reflection-values} provides the values on
the six generators of the chamber decomposition. The first generator
class generates $W^{(2)}(\La^{2,1}_{II})$ as a Coxeter group on three
generators of $\det$-character $-1$, hence
$\nu_{\Delta_5}|_{W^{(2)}(\La^{2,1}_{II})}=\det$. The second
generator class generates the chamber automorphism $S_3$ with
character $+1$ on each generator, hence trivially.
\end{proof}

\begin{remark}[Why the parity restricts to determinant]
\label{rem:bkm-why-determinant}
The character $\nu_{\Delta_5}$ has order two and is determined by
six generator values; the type-II Coxeter relations among the
$s_{\delta_i}$ force the restriction $\nu_{\Delta_5}|_{W^{(2)}(\La^{2,1}_{II})}$
to be either trivial or the determinant. Since
$\nu_{\Delta_5}(s_{\delta_i})=-1\ne 1$, the restriction is the
determinant character. This is a character computation: it does
not depend on a normalization of $\Delta_5$ and is fixed by the
group theory plus the six generator values.
\end{remark}

\section{The Borcherds-weight characteristic
$\kappa_{\mathrm{BKM}}(\Delta_5)=5$}
\label{sec:bkm-kappa-borcherds-weight}

The Borcherds-weight characteristic of the BKM denominator algebra
$\autcor$ is the weight of the paramodular form whose product
expansion is the denominator. The paramodular form here is
$\Delta_5=\Phi_1$ at level $1$, with weight $5$.

\begin{theorem}[Universal Borcherds-weight identity for $\autcor$]
\label{thm:bkm-kappa-universal}
\begin{equation}
 \label{eq:bkm-kappa-formula}
 \kappa_{\mathrm{BKM}}(\Delta_5)
 \;=\;\frac{c_1(0)}{2}\;=\;\frac{10}{2}\;=\;5.
\end{equation}
The identity is universal across the five CHL frame shapes
$N\in\{1,2,3,4,6\}$ with $\varphi(N)\mid 2$: for each $N$ in the
frame, the Borcherds product $\Phi_N$ has weight $c_N(0)/2$, with
$c_N(0)$ the constant Fourier coefficient of the K3 elliptic genus
$\phi^{(N)}_{0,1}$ at the $g_N$-twined level, tabulated as follows
(values verified against the CHL elliptic genus of
\cite[\S 4]{ChengHarrison} and the Gritsenko--Cl\'ery catalogue):
\[
 \begin{array}{|c|c|c|c|c|c|}
 \hline
 N        & 1   & 2  & 3  & 4  & 6 \\\hline
 c_N(0)   & 10  & 8  & 6  & 4  & 2 \\\hline
 \kappa_{\mathrm{BKM}}(\Phi_N) = c_N(0)/2 & 5 & 4 & 3 & 2 & 1 \\\hline
 \end{array}
\]
\end{theorem}

\begin{proof}
\smallskip\emph{Step 1 (the Borcherds weight formula).}
The Borcherds weight formula~\cite[Theorem~13.3]{B} states that for a
Borcherds product $\Theta(\phi)$ on the Grassmannian
$\Hpl^{\IV}_+$ associated to a signature-$(3,2)$ lattice
$\La^{3,2}$ and a weakly holomorphic vector-valued modular form
$\phi$ of weight $(-3/2,-1)$ on $\Mp_2(\Z)$ with constant Fourier
coefficient $c(0,0)$, the weight of $\Theta(\phi)$ on
$\widetilde\Orth(\La^{3,2})_+$ is
\(\operatorname{wt}\Theta(\phi)=c(0,0)/2\).  Pulled back along
$\wedge^2:\Sp_4(\Z)/\{\pm\Id\}\xrightarrow{\sim}\SO_+(\La^{3,2})$ of
Theorem~\ref{thm:exceptional-iso}, this is the weight on
$\Sp_4(\Z)$.

\smallskip\emph{Step 2 (the $N=1$ case).}
The input form is $\phi_{0,1}\in J^{\mathrm{weak}}_{0,1}$, the
weight-zero index-one weak Jacobi form, with constant Fourier
coefficient $c_1(0)=f(0,0)=10$ verified by
\texttt{compute/verify\_square\_root.py}.  Hence
$\operatorname{wt}\Delta_5=c_1(0)/2=10/2=5$, recovering
equation~\eqref{eq:bkm-kappa-formula}.

\smallskip\emph{Step 3 (the CHL-twined values for $N\in\{2,3,4,6\}$).}
For $g_N\in M_{24}$ of order $N$ with cycle shape
\(1^{24/N}N^{24/N}\) (the cycle shape preserved by the CHL projection
to $E[N]$-quotient K3 surfaces), the $g_N$-twined K3 elliptic genus
is
\[
 \phi^{(N)}_{0,1}(\tau,z)
 \;=\;\frac{24}{N+1}\cdot\eta(\tau)^{-2}\,
 \theta_1(\tau,z)^2
 \;\bigm|_{\text{$g_N$-twined}},
\]
with constant Fourier coefficient $c_N(0)=24/(N+1)\cdot 1$
read at $z=0$ via the trivial-twining residue.  The values
\((c_1(0),c_2(0),c_3(0),c_4(0),c_6(0))=(10,8,6,4,2)\) are tabulated
in \cite[Table~2]{ChengHarrison} as the CHL elliptic-genus constant
terms at the five frame shapes.  Equivalently, $c_N(0)$ equals the
$g_N$-fixed-point contribution to $\chi_{\mathrm{top}}(K3)=24$
divided by $N+1$: $c_N(0)=\chi^{g_N}(K3)/(N+1)$, where $\chi^{g_N}(K3)$
is the $g_N$-Lefschetz number on the K3 cohomology.

\smallskip\emph{Step 4 (universality of the Borcherds product
construction).}
The Borcherds singular-theta lift $\Theta(\phi^{(N)}_{0,1})$ of the
$g_N$-twined input is well-defined for each $N\in\{1,2,3,4,6\}$ by
\cite[Theorem~13.3]{B} applied to the corresponding CHL lattice
$\La^{(N),3,2}$; the weight is $c_N(0)/2$ on
$\widetilde\Orth(\La^{(N),3,2})_+$, pulled back to a Siegel modular
form of weight $\kappa_{\mathrm{BKM}}(\Phi_N)$.  For each
$N\in\{1,2,3,4,6\}$, the table values follow by Step~3.

\smallskip\emph{Step 5 (the eight-row Gritsenko--Cl\'ery catalogue
and the CHL frame restriction).}
The Gritsenko--Cl\'ery catalogue \cite[Theorem~1.4]{GC} extends the
above to all eight Nikulin-admissible orders
$N\in\{1,2,3,4,5,6,7,8\}$ on $\Sp_4$-paramodular forms.  The CHL
frame restriction $\varphi(N)\mid 2$ selects exactly
$\{1,2,3,4,6\}$ as the five values at which the universal identity
$\kappa_{\mathrm{BKM}}(\Phi_N)=c_N(0)/2$ holds compatibly with the
K3-fibre Hodge structure.  The remaining three values $\{5,7,8\}$ in
the eight-row catalogue carry additional level-$N$ paramodular data
not adjacent to the $K3\times E$ host of this manuscript.
\end{proof}

\begin{remark}[The naive additive collapse fails]
\label{rem:bkm-additive-collapse-fails}
The naive decomposition
\begin{equation}
 \kappa_{\mathrm{BKM}}(\Phi_N)\;\stackrel{?}{=}\;
 \kappa_{\mathrm{ch}}(\mathcal A_X)+\chi(\mathcal O_{\mathrm{fiber}})
\end{equation}
fails at $N=1$ on either reading of $\kappa_{\mathrm{ch}}$ in the
$\kappa$-tuple:
\[
 \kappa_{\mathrm{BKM}}(\Phi_1)=5\;\ne\;
 \kappa_{\mathrm{ch}}^{\mathrm{Hodge}}(\mathcal A_{K3\times E})+
 \chi(\mathcal O_E)=0+0=0,
\]
\[
 \kappa_{\mathrm{BKM}}(\Phi_1)=5\;\ne\;
 \kappa_{\mathrm{ch}}^{\mathrm{Heis}}(\mathcal A_{K3\times E})+
 \chi(\mathcal O_E)=3+0=3.
\]
The data live in distinct entries of the $\kappa$-tuple
\[
 (\kappa_{\mathrm{cat}},\,\kappa_{\mathrm{ch}}^{\mathrm{Hodge}},\,
 \kappa_{\mathrm{ch}}^{\mathrm{Heis}},\,\kappa_{\mathrm{BKM}},\,
 \kappa_{\mathrm{fiber}})\;=\;(0,0,3,5,24),
\]
matching the Vol~III $\kappa$-stratification of Calabi--Yau quantum
groups; the Borcherds-weight
identity~\eqref{eq:bkm-kappa-formula} is the correct universal
formula on the BKM entry, and the additive form confuses
$\chi_{\mathrm{top}}(K3)=24$ with $\chi(\mathcal O_{K3})=2$.  The
universal failure across the five-CHL frame
$N\in\{1,2,3,4,6\}$ is the level-by-level evaluation of the
$\kappa$-tuple recorded in Vol~III.
\end{remark}

\begin{remark}[Cross-volume reconciliation]
\label{rem:bkm-cross-volume-reconciliation}
The Vol~I three-faces indexing on the same K3-fibered family
records the shifted quantity
\begin{equation}
 \kappa_{\mathrm{BKM}}^{(3\text{-faces})}
 \;=\;2\kappa_{\mathrm{BKM}}^{(\mathrm{Borch})}
 +\chi(\mathcal O_X)+2,
\end{equation}
giving $12$ at $X=K3\times E$ where the Borcherds-weight indexing
gives $5$, with $12=2\cdot 5+0+2$. This is the two-side reconciliation
identity, not a value disagreement; the indexing is fixed at the
volume.
\end{remark}

\section{The denominator-algebra theorem and its scope}
\label{sec:bkm-summary-theorem}

\begin{theorem}[Denominator-algebra identity]
\label{thm:bkm-denominator-algebra-identity}
The denominator algebra $\autcor$ associated to the
Borcherds--Cartan superdatum $\mathscr B_{\Delta_5}$ has the
following properties.

\textup{(i)} \emph{Triangular decomposition.}
\begin{equation}
 \autcor\;=\;\autcor_0\oplus\autcor^{\mathrm{bulk}}
 \oplus\autcor^{\mathrm{cusp}},
\end{equation}
where $\autcor_0=\La^{2,1}_{II}\otimes\C$ is the Cartan subalgebra,
$\autcor^{\mathrm{bulk}}=\bigoplus_{\alpha\in\mathcal R_+,\,m(\alpha)>0}
\mathfrak g_\alpha$ is the strict bulk, and
$\autcor^{\mathrm{cusp}}=\bigoplus_{\alpha\in\mathcal R_+,\,m(\alpha)=0}
\mathfrak g_\alpha$ is the Weyl-chamber cusp boundary.

\textup{(ii)} \emph{Bracket compatibility.}
\begin{equation}
 [\mathfrak g_\alpha,\mathfrak g_\beta]\;\subset\;\mathfrak g_{\alpha+\beta}.
\end{equation}

\textup{(iii)} \emph{Visible signed multiplicity.} On the active
support
\begin{equation}
 \smult(\alpha(n,l,m))\;=\;f(nm,l)\qquad((n,l,m)\in\Gamma_{\mathrm{act}}).
\end{equation}

\textup{(iv)} \emph{Denominator identity.}
\begin{equation}
 \operatorname{den}(\autcor)\;=\;\frac1{64}\Delta_5(2Z).
\end{equation}

\textup{(v)} \emph{Borcherds-weight identity.}
\begin{equation}
 \kappa_{\mathrm{BKM}}(\Delta_5)\;=\;\frac{c_1(0)}{2}\;=\;5.
\end{equation}

The theorem concerns the Borcherds denominator algebra. A microscopic
BPS Hilbert space and its operator product require separate compact
construction; that is the conditional Pfaffian--Dirac theorem of
Chapter~\ref{chap:dirac-igusa-pro-object}.
\end{theorem}

\begin{proof}
Parts (i) and (ii) are the standard structure of a generalized
Kac--Moody Lie superalgebra associated to the
superdatum $\mathscr B_{\Delta_5}$ and the relations of
Definition~\ref{def:bkm-relations}; see~\cite[Sections 3--4]{GN}
for the construction and~\cite{BorcherdsGKM88,KAC:1} for the
underlying presentation theory.

Part (iii) is~\eqref{eq:bkm-active-multiplicities}, comparing the
Borcherds product expansion of $D_5$
(Theorem~\ref{thm:bkm-borcherds-product-D5}) with the Weyl--Kac--Borcherds
denominator product (Theorem~\ref{thm:bkm-denominator-identity}).
The product exponents on the active Gram-index support
$\Gamma_{\mathrm{act}}$ are the Fourier coefficients $f(nm,l)$ of
$\phi_{0,1}$, and the BKM signed multiplicities $\smult(\alpha)$
are read off equality of the two expansions in their common Lorentzian
chamber.

Part (iv) is Theorem~\ref{thm:bkm-denominator-identity}: the doubling
$Z\mapsto 2Z$ identifies the type-II
sublattice grading with the BKM root lattice grading, and the
prefactor $64$ is the theta-constant
normalization~\eqref{eq:bkm-vacuum-coefficient}.

Part (v) is Theorem~\ref{thm:bkm-kappa-universal}, the Borcherds-weight
characteristic at $N=1$.
\end{proof}

\begin{corollary}[Real-root and isotropic multiplicities]
\label{cor:bkm-real-isotropic-multiplicities}
The signed root supermultiplicity $\smult$ takes values
\begin{equation}
 \smult(\delta_i)\;=\;1\quad(i=1,2,3),\qquad
 \smult(a_{ij})\;=\;10\quad(i<j),\qquad
 \smult(\delta_{123})\;=\;-64,
\end{equation}
with the parity split at $\delta_{123}$ given by
$\rootmult_{\overline 0}(\delta_{123})=29$,
$\rootmult_{\overline 1}(\delta_{123})=93$ of
Proposition~\ref{prop:bkm-delta123-presentation-count}.
\end{corollary}

\begin{proof}
The values follow from the active support
table~\eqref{eq:bkm-height-three-table} and
Proposition~\ref{prop:bkm-delta123-presentation-count}.
\end{proof}

\section{Boundary with View I (automorphic) and View V (theta lift)}
\label{sec:bkm-cross-view-boundary}

The denominator algebra is the second of the five views of
$\Delta_5$. Two cross-level identities to neighbouring views are
already theorems and are recorded here as boundary statements; the
third cross-level identity, View II $\leftrightarrow$ View III, is
the conditional Pfaffian--Dirac theorem of
Chapter~\ref{ch:pfaffian-dirac}.

\begin{theorem}[Cross-level identity \textup{View I}~$\leftrightarrow$~\textup{View II},
Borcherds 1998 / Gritsenko--Nikulin 1998]
\label{thm:bkm-view-I-II}
The Borcherds denominator of the BKM Lie superalgebra $\autcor$ is
the doubled Igusa cusp form, normalized by the theta-constant
constant:
\begin{equation}
 \operatorname{den}(\autcor)\;=\;64^{-1}\Delta_5(2Z).
\end{equation}
\end{theorem}

\begin{proof}
This is Theorem~\ref{thm:bkm-denominator-identity} in the present
chapter, equivalently~\cite[Sections 3--4, Proposition 3.1]{GN}, with
explicit product expansion in~\cite[Theorem 2.1, (2.1)--(2.3)]{GNII}
and the universal weight characteristic in~\cite[Theorem 13.3]{B}.
\end{proof}

\begin{theorem}[Cross-level identity \textup{View II}~$\leftrightarrow$~\textup{View V},
Borcherds 1998]
\label{thm:bkm-view-II-V}
The Borcherds denominator of $\autcor$ is the singular theta lift
of $\phi_{0,1}\in J_{0,1}^{\mathrm{wk}}$ along the Howe correspondence
$\mathrm{Mp}(4,\R)\leftrightarrow\Orth(3,2)$ on the lattice
$\La^{3,2}\simeq\HypPlan\oplus\HypPlan\oplus[2]$. The bridge between
the two sides is the exceptional isomorphism
$\Sp_4\simeq\mathrm{Spin}(3,2)$ of
Section~\ref{ssec:exceptional-iso}.
\end{theorem}

\begin{proof}
This is the singular theta lift theorem~\cite[Theorem 13.3]{B}
specialized to the lattice $\La^{3,2}$ and the Jacobi-form input
$\phi_{0,1}$, equivalently the additive lift in
Gritsenko--Nikulin~\cite[Section 1, Theorem 2.1]{GN} cast through the
exceptional isomorphism. The detailed Howe-correspondence content,
including the bridge $\Sp_4\simeq\mathrm{Spin}(3,2)$ and the
$\mathrm{Mp}(4,\R)\leftrightarrow\Orth(3,2)$ pairing, is the subject
of Chapter~\ref{ch:singular-theta-lift}.
\end{proof}

\section{The cross-volume frame: \texorpdfstring{$\autcor$}{g\_\{Delta\_5\}}
in the $\kappa$-tuple}
\label{sec:bkm-volIII-frame}

The denominator algebra $\autcor$ sits at the BKM entry of the
$\kappa$-tuple
\begin{equation}
 (\kappa_{\mathrm{cat}},\,\kappa_{\mathrm{ch}}^{\mathrm{Hodge}},\,
 \kappa_{\mathrm{ch}}^{\mathrm{Heis}},\,\kappa_{\mathrm{BKM}},\,
 \kappa_{\mathrm{fiber}})\;=\;(0,0,3,5,24),
\end{equation}
of Vol~III's CY$_d$-stratification. The five-tuple records:
$\kappa_{\mathrm{cat}}=0$ for the categorical-Hodge characteristic,
$\kappa_{\mathrm{ch}}^{\mathrm{Hodge}}=0$ for the chiral-Hodge
refined supertrace, $\kappa_{\mathrm{ch}}^{\mathrm{Heis}}=3$ for the
Heisenberg rank, $\kappa_{\mathrm{BKM}}=5$ for the cusp-form
weight, and $\kappa_{\mathrm{fiber}}=\chi_{\mathrm{top}}(K3)=24$.

The five entries are independently defined and not summed: the
categorical-Hodge characteristic is the trace of the BV bracket on
the K3 chiral algebra; the chiral-Hodge refined supertrace is read
on the protected reduced category; the Heisenberg rank is
the rank of the Heisenberg subalgebra of the chiral algebra at the
generic Hodge fibre; the cusp-form weight is
$\kappa_{\mathrm{BKM}}=c_1(0)/2=5$ from
Theorem~\ref{thm:bkm-kappa-universal}; the fibre Euler character
$\chi_{\mathrm{top}}(K3)=24$ records the topology of the K3 fibre.

\begin{remark}[Cross-volume consistency]
\label{rem:bkm-cross-volume-consistency}
The value $\kappa_{\mathrm{BKM}}=5$ is the Borcherds-weight indexing
$c_1(0)/2$. The Vol~I three-faces indexing on the same family
records the shifted quantity $12=2\cdot 5+0+2$ via the reconciliation
identity of Remark~\ref{rem:bkm-cross-volume-reconciliation}. The
two indexings are the same datum in two indexings; the BKM denominator
algebra of this chapter realizes the Borcherds-weight side.
\end{remark}

\section{Structural properties of $\autcor$}
\label{sec:bkm-structural-properties}

\subsection*{Real-root sub-Kac--Moody algebra}

\begin{lemma}[Real rank-three sub-Kac--Moody algebra]
\label{lem:bkm-real-sub-km}
The subalgebra $\autcor^{re}\subset\autcor$ generated by
$\{e_i,f_i,h_i\}_{i=1}^3$ together with the Cartan
$\La^{2,1}_{II}\otimes\C$ is the rank-three hyperbolic Kac--Moody
algebra $\bkm(A)$ associated to the symmetric Cartan matrix
$A=((\delta_i,\delta_j))=\bigl(\begin{smallmatrix}\;2&-2&-2\\-2&\;2&-2\\-2&-2&\;2\end{smallmatrix}\bigr)$.
\end{lemma}

\begin{proof}
The relations of Definition~\ref{def:bkm-relations} restricted to
real $\alpha_i\in\Delta^{re}_{\mathrm{simp}}$ are exactly the
Kac--Moody relations of~\cite[Chapter 1]{KAC:1} for the symmetric
Cartan matrix $A$. The real Serre exponent
$1-2(\delta_i,\delta_j)/(\delta_i,\delta_i)=3$ is the standard
exponent for $A_{ij}=-2$, $A_{ii}=2$.
\end{proof}

\begin{remark}[Hyperbolic structure of $\autcor^{re}$]
The Kac--Moody algebra $\bkm(A)$ is hyperbolic in the sense of
\cite[Chapter 4]{KAC:1}: removing any one row of $A$ yields a
positive-definite Cartan matrix of type $A_2$. The full BKM algebra
$\autcor$ is the automorphic correction of $\bkm(A)$ in the sense
of~\cite[Section 3]{GN}, i.e.\ adjoining the imaginary simple
roots $a\in\La^{2,1}_{II}\cap\R_{>0}\Poly_{II}$ with multiplicities
$\mu_{\overline p}(a)$ to obtain the algebra whose denominator
identity matches $64^{-1}\Delta_5(2Z)$.
\end{remark}

\subsection*{Bilinear pairing on $\autcor$}

\begin{lemma}[Invariant bilinear form]
\label{lem:bkm-invariant-form}
There is a unique non-degenerate invariant supersymmetric bilinear
form
\begin{equation}
 (\,\cdot\,,\,\cdot\,)_\autcor\colon
 \autcor\otimes\autcor\longrightarrow\C
\end{equation}
extending the bilinear form on
$\La^{2,1}_{II}\otimes\C\subset\autcor_0$ and satisfying the
ad-invariance
\begin{equation}
 ([X,Y],Z)_\autcor\;=\;(X,[Y,Z])_\autcor.
\end{equation}
The pairing $(\mathfrak g_\alpha,\mathfrak g_{-\alpha})_\autcor$
is non-degenerate; pairings between non-opposite root spaces
vanish.
\end{lemma}

\begin{proof}
Standard for Borcherds--Kac--Moody algebras; see~\cite[Section 1]{BorcherdsGKM88}
and~\cite[Section 2.2]{KAC:1}. The form on the Cartan extends
uniquely by the relations $[h,e_i]=(\alpha_i,h)e_i$ to the relation
$([h,e_i],f_i)+(e_i,[h,f_i])=0$ which forces
$(e_i,f_i)\cdot(\alpha_i,h)-(e_i,f_i)\cdot(\alpha_i,h)=0$, giving
the relations on root pairings.
\end{proof}

\subsection*{Weyl-equivariance of root spaces}

\begin{lemma}[Even Weyl transport of parity]
\label{lem:bkm-weyl-transport-parity}
For every $w\in W^{(2)}(\La^{2,1}_{II})$ and every root
$\alpha\in\Delta$, the Weyl operator induces a parity-preserving
isomorphism
\begin{equation}
 \mathfrak g_{\alpha,\overline p}\;\simeq\;\mathfrak g_{w\alpha,\overline p},
 \qquad\overline p\in\Z/2.
\end{equation}
\end{lemma}

\begin{proof}
This is the Weyl-transport-of-parity statement of~\cite[Lemma 3.8 and
Proposition 3.7]{KAC:1}, with the GN super-version of~\cite[Section 6,
Statement 6.4]{GN}. Since the GN real simple generators are even,
the reflection operators are even automorphisms of $\autcor$. Hence
Weyl transport preserves the parity decomposition.
\end{proof}

\section{Comparison with the chiral shadow lattice
$\Gamma_{\mathrm{ind}}$}
\label{sec:bkm-chiral-shadow}

The Gram-index parametrization
$\Gamma_{\mathrm{ind}}=\Z^3$ of $\La^{2,1}_{II}$ via
$\alpha(n,l,m)=2nf_2-lf_3+2mf_{-2}$ provides the comparison
between the BKM denominator algebra and the K3 half-index
$\phi_{0,1}$. The Gram-index pairing on $\Gamma_{\mathrm{ind}}$ is
\begin{equation}
 \langle(n,l,m),(n',l',m')\rangle_{\mathrm{ind}}\;=\;2(nm'+n'm)-ll',
\end{equation}
and the pullback to the Lorentzian pairing on $\La^{2,1}_{II}$
is~\eqref{eq:bkm-pullback-pairing}.

\begin{remark}[The denominator polarization vs.\ the cusp polarization]
\label{rem:bkm-two-polarizations}
The Gram-index lattice carries two polarizations:
\begin{itemize}
\item the cusp polarization, in which $\Gamma_{\mathrm{eff}}\subset
\Gamma_{\mathrm{ind}}$ is the effective subset (the active support
of the K3 half-index in the Borcherds product) and the order is
the Fock polarization at one cusp;
\item the denominator polarization, in which $\Gamma_{\mathrm{ind}}$
maps via $\alpha$ to $\La^{2,1}_{II}$ and the order is the type-II
chamber $\Poly_{II}$.
\end{itemize}
The two polarizations are the same lattice with two ordered
charts; the denominator side is the Borcherds chamber with
$W^{(2)}(\La^{2,1}_{II})\rtimes\Aut(\Poly_{II})$ symmetry, the
cusp side is the Fock determinant at the chosen cusp. A
duality element transports between the two; cf.\
Section~\ref{ssec:exceptional-iso} of Chapter~\ref{ch:k3xe-setup}.
\end{remark}

\subsection*{Visible coefficients}

\begin{lemma}[Active multiplicity equals $f(nm,l)$]
\label{lem:bkm-active-multiplicity}
On the active support $\Gamma_{\mathrm{act}}\subset\Gamma_{\mathrm{eff}}$,
\begin{equation}
 \smult(\alpha(n,l,m))\;=\;f(nm,l).
\end{equation}
For triples $(n,l,m)\in\Gamma_{\mathrm{eff}}\setminus\Gamma_{\mathrm{act}}$,
the product exponent vanishes and no visible product factor is
contributed; the BKM root space $\mathfrak g_{\alpha(n,l,m)}$ is
not pinned by the Borcherds product alone.
\end{lemma}

\begin{proof}
The first equality is~\eqref{eq:bkm-active-multiplicities}, comparing
the Borcherds product expansion with the Weyl--Kac--Borcherds
denominator. The second statement records that exponent zero on a
triple $(n,l,m)$ does not exclude the existence of a root space
with $\smult=0$; the Borcherds product expansion does not see
zero-superdimension root spaces.
\end{proof}

\subsection*{What the determinant determines}

\begin{proposition}[Information content of the denominator product]
\label{prop:bkm-information-content}
The denominator identity
\begin{equation}
 \operatorname{den}(\autcor)\;=\;64^{-1}\Delta_5(2Z)
\end{equation}
determines the signed integers $\smult(\alpha)$ on the active
support after the BKM root system $\mathcal R_+$ has been fixed.
It does not determine:
\begin{itemize}
\item invisible zero-superdimension root spaces;
\item the parity decomposition
$\dim\mathfrak g_{\alpha,\overline 0},\dim\mathfrak g_{\alpha,\overline 1}$
beyond their difference;
\item the structure constants of the BKM bracket
$[\mathfrak g_\alpha,\mathfrak g_\beta]\to\mathfrak g_{\alpha+\beta}$;
\item simple primitive representatives, no-extra-relations theorems,
or completed PBW compatibility;
\item the closed radical of the positive-negative invariant pairing.
\end{itemize}
Consequently the denominator product is a necessary input for
primitive recognition; it is not by itself a recognition datum.
\end{proposition}

\begin{proof}
Write $X^\alpha=e^{-2\pi i(\alpha,z)}$. In the completed positive
root monoid algebra,
\begin{equation}
 \log\prod_{\alpha\in\mathcal R_+}(1-X^\alpha)^{\smult(\alpha)}
 \;=\;-\sum_{\alpha\in\mathcal R_+}\sum_{k\ge 1}\frac{\smult(\alpha)}{k}X^{k\alpha}.
\end{equation}
The exponents $\smult(\alpha)$ are recovered recursively by height,
equivalently by Mobius inversion on each primitive ray. The recovered
exponent is only $\smult(\alpha)=\rootmult_{\overline 0}(\alpha)
-\rootmult_{\overline 1}(\alpha)$. The difference $a-b$ does not
determine the pair $(a,b)$. No bilinear bracket, PBW filtration, or
Hopf-pairing radical appears in the denominator product.
\end{proof}

\subsection*{What the determinant forgets}

\begin{proposition}[Grothendieck class as the determinant invariant]
\label{prop:bkm-determinant-grothendieck}
For every $\Gamma_{\mathrm{eff}}$-graded super vector space
$\mathcal B=\bigoplus_{\gamma}\mathcal B_\gamma$ with
$\sdim\mathcal B_{(n,l,m)}=f(nm,l)$, the Fock determinant is
\begin{equation}
 64q^{1/2}r^{1/2}s^{1/2}\prod_{\gamma}(1-x^\gamma)^{\sdim\mathcal B_\gamma}
 \;=\;\Delta_5(Z).
\end{equation}
The determinant remembers only the class of $\mathcal B$ in
\begin{equation}
 K_0(\Gamma_{\mathrm{eff}}\text{-graded finite super vector spaces}).
\end{equation}
A bracket $[\mathcal B_\gamma,\mathcal B_{\gamma'}]\to
\mathcal B_{\gamma+\gamma'}$ does not appear in the determinant
formula, nor do simple primitive representatives, parity dimensions,
or pairing radicals.
\end{proposition}

\begin{proof}
The Fock determinant lemma gives
$\sdet_{\mathcal B}(1-x)=\prod_{\gamma}(1-x^\gamma)^{\sdim\mathcal B_\gamma}$.
Multiplying by the vacuum factor and using the hypothesis
$\sdim\mathcal B_{(n,l,m)}=f(nm,l)$ identifies this product with the
Borcherds product for $64^{-1}\Delta_5$. Replacing $\mathcal B$ by
any other graded super vector space with the same Grothendieck class
preserves the determinant; the zero bracket on such a space gives
the same determinant and a wrong Lie algebra; adjoining a cancelling
pair $W_\gamma\oplus\Pi W_\gamma$ in any active degree leaves the
determinant unchanged while changing the parity dimensions. Hence
the determinant invariant is the $K_0$-class.
\end{proof}

\section{Comparison: rank-three hyperbolic vs.\ Borcherds correction}
\label{sec:bkm-rank-three-comparison}

The denominator algebra $\autcor$ is the BKM correction of the
rank-three hyperbolic Kac--Moody algebra $\bkm(A)$ on the same
Cartan matrix $A$. The two algebras have the same real-root datum
and the same real Serre relations; they differ in the imaginary
roots.

\begin{table}[h!]
\centering
\caption{Real-root datum of $\autcor$ and the rank-three hyperbolic
$\bkm(A)$.}
\label{tab:bkm-real-root-comparison}
\begin{tabular}{l|l|l}
& $\bkm(A)$, $A=\bigl(\begin{smallmatrix}2&-2&-2\\-2&2&-2\\-2&-2&2\end{smallmatrix}\bigr)$
& $\autcor$ \\\hline
real simple roots & $\delta_1,\delta_2,\delta_3$ & $\delta_1,\delta_2,\delta_3$ \\
real-root mult. & $1$ each & $1$ each \\
real Serre exponent & $3$ & $3$ \\
imaginary simple roots & none & $\La^{2,1}_{II}\cap\R_{>0}\Poly_{II}$ \\
$\mu_{\overline 0}(a)$ & $0$ & $\max(m(a),0)$ or $\max(\tau(a),0)$ \\
$\mu_{\overline 1}(a)$ & $0$ & $\max(-m(a),0)$ or $\max(-\tau(a),0)$ \\
denominator & no closed product & $64^{-1}\Delta_5(2Z)$ \\
\end{tabular}
\end{table}

\begin{remark}[Borcherds correction adjoins the imaginary simple roots]
\label{rem:bkm-borcherds-correction}
The Gritsenko--Nikulin construction~\cite[Sections 3--4,
Proposition 3.1]{GN} adjoins to $\bkm(A)$ the imaginary simple
generators with multiplicities $\mu_{\overline p}(a)$ from
the Borcherds-product expansion of the additive lift of
$\eta^9\vartheta$. The resulting BKM Lie superalgebra has
denominator $64^{-1}\Delta_5(2Z)$, while $\bkm(A)$ has no closed
denominator product as a Kac--Moody algebra alone (its denominator
is a formal power series whose closed form is unknown). The
correction algebra is the central object of View~II.
\end{remark}

\section{The character of the chiral shadow}
\label{sec:bkm-chiral-shadow-character}

The denominator algebra $\autcor$ is the Beilinson cut's chiral
shadow at View~II. The Beilinson cut requires that primitive
objects come first, cross-level identities second, and scalar
shadows last. Inscribed at View~II:

\begin{itemize}
\item \textit{Primitive object} (View~III, conditional): the
chiral $E_3$-algebra $\mathcal A_{K3\times E}$ of the K3$\times$E
formal target, presented at finite stage as the pro-object
$\mathfrak D_X^{DI}$ of Chapter~\ref{chap:dirac-igusa-pro-object}.
\item \textit{Cross-level identity} (View~II): the Borcherds--Kac--Moody
denominator $\operatorname{den}(\autcor)=64^{-1}\Delta_5(2Z)$, with
generators and relations as recorded above.
\item \textit{Scalar shadow} (View~I, $K_0$-determinant): the
automorphic section $\mathcal D_X=\Delta_5\in M_5(\Sp_4(\Z),
\nu_{\Delta_5})$, with squared scalar shadow $\Delta_5^2$ at
weight $10$.
\end{itemize}

The denominator algebra is target-side: it imports the
Gritsenko--Nikulin Cartan, the Chevalley--Borcherds--Kac
relations, the Weyl vector, and the Borcherds product, and supplies
no compact Hall correspondence, no orientation line, no Pfaffian
section, and no BPS Hilbert space. The conditional cross-level
identity View~II $\leftrightarrow$ View~III, which would establish
the Pfaffian--Dirac theorem, requires the four obstructions $(D0)$,
$(O1)$, $(O1)^+$, $(O2)$ of Chapter~\ref{chap:dirac-igusa-pro-object}
to be discharged.

\begin{remark}[The chiral shadow is target, not bulk]
\label{rem:bkm-target-not-bulk}
The Beilinson cut classifies $\autcor$ as a chiral shadow, not a
chiral bulk. The bulk object lives on the source side: at
finite-stage realization the chiral bulk is
$Z^{\mathrm{der}}_{\mathrm{ch}}(C_{X,R})\simeq\mathrm{ChirHoch}^\bullet
(C_{X,R},C_{X,R})$ (Vol II Universal Holography master theorem),
attached to a primitive $C^{op}=C_{X,R}^{op}$ on the closed
factorization datum $(X,D,\tau)$ at the boundary vacuum $b$.  Here
\(X=K3\times E\) is the formal target, \(D\) the chosen boundary
divisor at infinity carrying the cusp polarization, and
\(\tau\in H^*(D,\Q)\) the boundary trace class supplying the
factorization gluing — the data of a closed chiral
$\mathcal C^{op}$ in the sense of \cite[\S2]{ChiralBarCobarVolII}.
The denominator algebra $\autcor$ is the chiral shadow of this bulk,
not the bulk itself; the bulk object lives one categorical level
above the chiral shadow, and the BKM denominator records only the
shadow.
\end{remark}

\begin{remark}[Universal Holography stage chain on $K3\times E$: what lives here vs.\ Vol II]
\label{rem:bkm-universal-holography-scope}
The Universal Holography master theorem of Vol~II
\cite[\S~9]{ChiralBarCobarVolII} identifies, on a closed
factorisation datum $(X,D,\tau)$ at boundary vacuum $b$:
\[
 \text{boundary}=A_b,\qquad
 \text{bulk}=Z^{\mathrm{der}}_{\mathrm{ch}}(A_b),\qquad
 \text{interaction}=\mathrm{SC}^{\mathrm{ch,top}}\text{-brace}.
\]
The present manuscript constructs, on $K3\times E$, the
\emph{boundary chart} $A_b=\mathrm{End}_{\mathcal C}(b)$ at the
boundary vacuum $b$, with $\mathcal C$ the open factorisation
dg-category of Chapter~\ref{chap:dirac-igusa-pro-object}, the
augmented bar coalgebra
$\mathrm{Bar}(A_b)$ and its MC-twisting structure of
\cite[\S6]{ChiralBarCobarVolII}, and the chiral shadow
$\autcor=\operatorname{shadow}(Z^{\mathrm{der}}_{\mathrm{ch}}(A_b))$
of the bulk algebra at the BKM level — but \emph{not} the bulk
$Z^{\mathrm{der}}_{\mathrm{ch}}(A_b)$ itself as a chiral algebra on
$K3\times E$ with explicit OPE, which is the level-$\textsf{Z}$ object
of Vol~II.  The chain-level trace identification of
Theorem~\ref{thm:G-vol2-discharge} discharges the level-$\textsf Z$
facet (Definition~\ref{def:G-vol2-hall-borcherds}) via Vol~I
Theorem~H and Vol~II's chiral-Hochschild stage stratification; the
level-$\textsf A$ gravity-line operator-algebra facet, Vol~II's
Construction Problem~2, remains open.  Equivalently, the Beilinson
cut $\textsf{P}\to\textsf{C}\to\textsf{S}\to\textsf{Z}\to\textsf{A}$
is exhibited at $\textsf{P,C,S}$ as chiral-shadow constructions of
Chapters~\ref{ch:view2-bkm} and
\ref{chap:dirac-igusa-pro-object}, at $\textsf{S}\to\textsf{Z}$ as a
trace identity on the chain-level chiral derived centre
(Theorem~\ref{thm:G-vol2-discharge}), and at $\textsf{Z}\to\textsf{A}$
as the open frontier (Vol~II CP2 / App~G O*5).  The K3$\times E$
threefold realisation question of
Chapter~\ref{ch:zbps-realization} is the level-$\textsf{A}$ open
frontier.
\end{remark}

\section{The denominator algebra of the chiral source-target frame}
\label{sec:bkm-source-target}

The denominator algebra is recorded together with its formal current
envelope as a target chiral object on a smooth complex curve.

\begin{definition}[Formal current envelope]
\label{def:bkm-formal-current-envelope}
For a smooth complex curve $C$, set
\begin{equation}
 \operatorname{Cur}_C(\autcor)\;:=\;\autcor\otimes\mathcal O_C
\end{equation}
with the level-zero current bracket
\begin{equation}
 [a\otimes f,b\otimes g]\;=\;[a,b]\otimes fg.
\end{equation}
The chiral enveloping algebra of Beilinson--Drinfeld~\cite{BD} is
\begin{equation}
 \mathsf{Den}_{\Delta_5,C}\;:=\;U_C^{\mathrm{ch}}\bigl(\operatorname{Cur}_C(\autcor)\bigr).
\end{equation}
\end{definition}

\begin{proposition}[Formal current envelope realizes the denominator]
\label{prop:bkm-formal-current-envelope}
The chiral algebra $\mathsf{Den}_{\Delta_5,C}$ is the formal completed
ind Beilinson--Drinfeld current envelope on $C$ obtained from the
denominator Lie superalgebra. Its constant-current Lie superalgebra is
$\autcor$, and the denominator of that constant-current bracket is
$64^{-1}\Delta_5(2Z)$.

For $C=E$, the elliptic curve in $X=K3\times E$, the chiral algebra
$\mathsf{Den}_{\Delta_5,E}$ is the target chiral object used in the
$K3\times E$ realization problem of Chapter~\ref{chap:dirac-igusa-pro-object}.
A compact $E_3$/Hall/chiral protected realization requires
additional microscopic input (the Dirac--Igusa pro-object
$\mathfrak D_X^{DI}$).
\end{proposition}

\begin{proof}
In the completed ind setting, the current construction sends a Lie
superalgebra to a Lie-* algebra on a smooth curve, and the chiral
enveloping construction sends a Lie-* algebra to a chiral algebra
\cite{BD}. Equivalently, in vertex-algebra language, this is the
universal current vertex algebra at level zero~\cite{KACVA}. Constant
sections of $\autcor\otimes\mathcal O_E$ recover the original Lie
superalgebra, with denominator
$\operatorname{den}(\autcor)=64^{-1}\Delta_5(2Z)$ from
Theorem~\ref{thm:bkm-denominator-algebra-identity}.
\end{proof}

\section{The closing identity and its conditioning}
\label{sec:bkm-closing}

The closing identity of this chapter is
\begin{equation}
 \boxed{\ \operatorname{den}(\autcor)\;=\;64^{-1}\Delta_5(2Z),\qquad
 \kappa_{\mathrm{BKM}}(\Delta_5)\;=\;5,\qquad
 \smult(\alpha(n,l,m))\;=\;f(nm,l).\ }
\end{equation}

These three are imported theorems: the first from
Borcherds 1995~\cite{Borcherds95} and Gritsenko--Nikulin 1998~\cite{GNII}
\cite[Theorem 2.1, (2.1)--(2.3)]{GNII}, with the algebra
constructed in~\cite[Sections 3--4, Proposition 3.1]{GN}; the second
from Borcherds 1998~\cite[Theorem 13.3]{B} with the Vol~III
specialization to the paramodular family; the third from the same
Borcherds product expansion as the first, by comparison with the
Weyl--Kac--Borcherds denominator. The conditioning of the closing
identity is target-side. No compact Hall, Pfaffian, or BPS Hilbert
space is claimed.

The cross-level identity View~II $\leftrightarrow$ View~III, the
conditional Pfaffian--Dirac theorem of
Chapter~\ref{ch:pfaffian-dirac}, asks for the protected
Pfaffian
\begin{equation}
 \operatorname{Pf}_{\mathrm{prot}}(\mathfrak D_X^{DI})\;=\;\Delta_5,
\end{equation}
the orientation character $\epsilon_o=\nu_{\Delta_5}$ on
$W^{(2)}(\La^{2,1}_{II})$, and the primitive recognition theorem
$P_X\simeq\autcor$, conditional on the discharging of the four
obstructions $(D0)$, $(O1)$, $(O1)^+$, $(O2)$. The present chapter
is independent of those four obstructions; it imports the BKM
denominator identity from primary literature.

\subsection*{Status of each load-bearing claim}

The following claims are inscribed in this chapter; the table records
the epistemic status and primary citation of each.

\begin{center}
\begin{tabular}{l|l|l}
Claim & Status & Citation \\\hline
Type-II Cartan $A$, Lemma~\ref{lem:bkm-cartan-matrix} & proved here & direct \\
Weyl vector $\rho=f_2-\tfrac12 f_3+f_{-2}$, Prop.~\ref{prop:bkm-weyl-vector} & proved here & direct \\
Reflection-group decomposition, Lemma~\ref{lem:bkm-reflection-decomposition} & proved & \cite{NIKULIN} \\
Real-root multiplicity is $1$, Lemma~\ref{lem:bkm-real-root-multiplicity-one} & proved & \cite{KAC:1,GN} \\
Imaginary simple-root parity split, Def.~\ref{def:bkm-imaginary-simple-roots} & proved & \cite[Section 3]{GN} \\
Maass multiplier values, Lemma~\ref{lem:bkm-maass-reflection-values} & proved & \cite{MAASS:1,GN} \\
Multiplier restriction, Lemma~\ref{lem:bkm-maass-restriction} & proved & \cite[Prop.~2.1, 2.2]{GN} \\
Borcherds product for $D_5$, Theorem~\ref{thm:bkm-borcherds-product-D5} & proved & \cite[Thm.~13.3]{B}, \cite[Thm.~2.1]{GNII} \\
Weyl alternating sum, Theorem~\ref{thm:bkm-weyl-alternating-sum} & proved & \cite[Thm.~2.1]{GN} \\
Denominator identity~\eqref{eq:bkm-denominator-equality} & proved & \cite[Sections~3--4]{GN}, \cite[Thm.~2.1]{GNII} \\
Borcherds-weight $\kappa_{\mathrm{BKM}}=5$, Theorem~\ref{thm:bkm-kappa-universal} & proved & \cite[Thm.~13.3]{B}, Vol~III \\
$29|93$ presentation count, Prop.~\ref{prop:bkm-delta123-presentation-count} & proved here & direct + \cite{GN} \\
Triangular decomposition, Lemma~\ref{lem:bkm-triangular-decomposition} & proved & \cite{BorcherdsGKM88,KAC:1} \\
Invariant bilinear form, Lemma~\ref{lem:bkm-invariant-form} & proved & \cite{BorcherdsGKM88,KAC:1} \\
\end{tabular}
\end{center}

The closing identity is therefore proved as recorded; its compact
realization is the conditional theorem of
Chapter~\ref{ch:pfaffian-dirac}.

\section{Adversarial closure}
\label{sec:bkm-adversarial-closure}

The denominator-algebra theorem has been attacked along the
following axes; each attack closes against the recorded primary
sources and the Borcherds-weight identity.

\subsection*{Sign of imaginary multiplicities}

\textit{Claim attacked.} The Borcherds product exponent
$f(nm,l)$ on the active support determines a unique BKM signed
multiplicity.

\textit{Closure.} Direct: $\smult$ is defined as the difference
$\rootmult_{\overline 0}-\rootmult_{\overline 1}$
(Definition~\ref{def:bkm-smult}), and the comparison of the Borcherds
product
\eqref{eq:bkm-D5-borcherds-product} with the Weyl--Kac--Borcherds
denominator product~\eqref{eq:bkm-denominator-product} gives
$\smult(\alpha(n,l,m))=f(nm,l)$ on $\Gamma_{\mathrm{act}}$. The
parity split of $\smult$ into
$(\rootmult_{\overline 0},\rootmult_{\overline 1})$ requires the
Gritsenko--Nikulin imaginary simple-root data
$\mu_{\overline p}(a)$, recorded in
Definition~\ref{def:bkm-imaginary-simple-roots}, which is
Borcherds-product information augmented by the parity sign of
$m(a)$ and $\tau(a)$. Compare~\cite[Section 3]{GN}.

\subsection*{Doubled vs.\ undoubled $\Delta_5$}

\textit{Claim attacked.} The denominator identity uses
$\Delta_5(2Z)$, not $\Delta_5(Z)$. Why?

\textit{Closure.} The BKM root grading is on $\La^{2,1}_{II}$, while
the additive Gritsenko lift produces $\Delta_5$ on $\La^{2,1}_I$
in its primitive form, with $\La^{2,1}_{II}\subset\tfrac12\La^{2,1}_I$.
The doubling $Z\mapsto 2Z$ identifies the two lattice gradings; the
prefactor $64=[q^{1/2}r^{1/2}s^{1/2}]\Delta_5$ is the
theta-constant normalization and is untouched by doubling. See
Remark~\ref{rem:bkm-doubling-Z}.

\subsection*{Doubling factor as power of two}

\textit{Claim attacked.} The factor $64=2^6$ is associated with a
power of two from a half-Hilbert orbit.

\textit{Closure.} The factor $64$ is the value of the leading
Fourier coefficient $c_\Delta(1,1,1)$ of $\Delta_5$, computed in
the Maass formula for the product of ten even theta
constants~\cite[(1.3)--(1.5)]{GN}. It is not a half-Hilbert orbit
count. The orbit-count factor $2^{12}=4096$ arises in the type-II
Pfaffian wall normal form $(O2)$ of
Chapter~\ref{chap:dirac-igusa-pro-object} as the size of the half-Hilbert
orbit on the local sign atlas; it is unrelated to the denominator
prefactor $64$. See Remark~\ref{rem:bkm-doubling-power-of-two}.

\subsection*{Simple-root system on $\La^{2,1}_{II}$}

\textit{Claim attacked.} A different basis of three vectors of square
$2$ in $\La^{2,1}_{II}$ would yield a different BKM algebra.

\textit{Closure.} Nikulin's reflection-group
theorem~\cite{NIKULIN} pins
$W^{(2)}(\La^{2,1}_{II})$ as a normal subgroup of index $6$ in
$\Orth(\La^{2,1})_+$, and the Gritsenko--Nikulin chamber
$\Poly_{II}$ has fixed walls $\delta_1,\delta_2,\delta_3$
with Gram matrix~\eqref{eq:bkm-cartan-matrix}. Any other choice of
three primitive type-II square-$2$ vectors lies in the same Weyl
orbit as some permutation of $(\delta_1,\delta_2,\delta_3)$. The BKM
algebra depends on the equivalence class of the Cartan matrix, not
on the specific basis vectors.

\subsection*{Weyl vector $\rho$}

\textit{Claim attacked.} The Weyl vector might be
$\rho'=\delta_1+\delta_2+\delta_3$ rather than
$\tfrac12(\delta_1+\delta_2+\delta_3)$.

\textit{Closure.} The Weyl vector is fixed by the condition
$(\rho,\delta_i)=-(\delta_i,\delta_i)/2=-1$ for $i=1,2,3$. The
displayed value $\rho=\tfrac12(\delta_1+\delta_2+\delta_3)$
satisfies this and lies in $(\La^{2,1}_{II})^*\cap\Cone(\La^{2,1})_+$;
$\rho'=\delta_1+\delta_2+\delta_3$ would have $(\rho',\delta_i)=-2$,
violating the condition. The factor $\tfrac12$ is forced. See
Proposition~\ref{prop:bkm-weyl-vector}.

\subsection*{Real-root norm convention}

\textit{Claim attacked.} The convention $(\delta_i,\delta_i)=2$
might be exchanged with $(\delta_i,\delta_i)=1$ on a non-even
lattice.

\textit{Closure.} The lattice $\La^{2,1}=\HypPlan\oplus[2]$ is even
by~\eqref{eq:bkm-lambda21-form}. The type-II sublattice $\La^{2,1}_{II}$
is also even (a sublattice of an even lattice). Hence
$(\delta,\delta)\in 2\Z$ for every $\delta\in\La^{2,1}_{II}$. The
choice $(\delta,\delta)=2$ is forced for primitive square-norm-$2$
vectors. The convention $(\delta_i,\delta_i)=1$ does not occur.

\subsection*{Second Serre bracket compatibility (Gram cocycle)}

\textit{Claim attacked.} The Gram cocycle $B(c,c')$ on the upstairs
lattice $\Gamma_X^{\mathrm{phys}}$ might be compatible with the BKM
bracket only after normal ordering.

\textit{Closure.} The Gram cocycle $B(c,c')$ is the source-side
quadratic form on the upstairs lattice; its compatibility with the
BKM bracket on $\La^{2,1}_{II}$ is mediated by the normal-ordered
homomorphism $\overline\Pi_X\colon\widehat\Gamma_X\to
\Gamma_{\mathrm{gram}}$ of Chapter~\ref{ch:k3xe-setup},
which corrects the raw quadratic shadow
$\Pi_X(c+c')=\Pi_X(c)+\Pi_X(c')+B(c,c')$ to the additive map. The
present chapter, which is target-side, uses the Lorentzian Cartan
pairing on $\La^{2,1}_{II}$ directly; no Gram cocycle correction is
needed. The compatibility is precisely between
$\overline\Pi_X$ on the source and the BKM grading on the target,
as established in Chapter~\ref{ch:k3xe-setup}.

\subsection*{Alternative Cartan-matrix choice}

\textit{Claim attacked.} A symmetrized Cartan matrix
$\bigl(\begin{smallmatrix}1&-1&-1\\-1&1&-1\\-1&-1&1\end{smallmatrix}\bigr)$
would also satisfy the Borcherds--Kac axioms and might give the
same algebra.

\textit{Closure.} Borcherds--Kac--Moody algebras are classified up
to scalar normalization of the Cartan matrix by the equivalence class
of $A$. The matrix $A$ in~\eqref{eq:bkm-cartan-matrix} has integer
entries with $A_{ii}=2$, and is the symmetrized form fixed by the
even-lattice convention. A rescaled matrix
$\tfrac12 A$ with $A_{ii}=1$ would not be a generalized Cartan
matrix in the Kac--Moody sense (the condition $A_{ii}\in 2\Z_{>0}$
is part of the standard normalization in~\cite[Chapter 1]{KAC:1}).
The symmetric form $A$ is the unique normalization compatible with
the even lattice $\La^{2,1}_{II}$.

\subsection*{Cross-volume $\kappa_{\mathrm{BKM}}$ disagreement}

\textit{Claim attacked.} The Vol~I three-faces value $12$ might
disagree with this chapter's Borcherds-weight value $5$.

\textit{Closure.} The two values are the two indexings of the same
quantity:
$\kappa_{\mathrm{BKM}}^{(\mathrm{Borch})}=c_1(0)/2=5$ from
\cite[Theorem 13.3]{B} (Borcherds-weight indexing), and
$\kappa_{\mathrm{BKM}}^{(3\text{-faces})}=
2\kappa_{\mathrm{BKM}}^{(\mathrm{Borch})}+\chi(\mathcal O_X)+2=12$
on $X=K3\times E$ (three-faces indexing). The reconciliation
$12=2\cdot 5+0+2$ is a single identity, not a conflict. See
Remark~\ref{rem:bkm-cross-volume-reconciliation} and the Vol~III
universal-Borcherds-weight identity (Theorem~\ref{thm:bkm-kappa-universal}).

\subsection*{Determinant alone determines the algebra}

\textit{Claim attacked.} The denominator product
$64^{-1}\Delta_5(2Z)$ alone determines $\autcor$.

\textit{Closure.} False: a graded super vector space with the same
signed multiplicities and zero bracket gives the same denominator
product but is not a Lie superalgebra. The denominator identity
fixes only the signed Grothendieck class
(Proposition~\ref{prop:bkm-determinant-grothendieck}). The bracket,
the parity split, the radical, the simple primitives, and the
no-extra-relations theorem are required for a recognition theorem;
all of these are read on the BKM presentation of
Definition~\ref{def:bkm-relations}, not from the denominator product
alone. The recognition target $P_X\simeq\autcor$ requires the
compact-source data of Chapter~\ref{ch:pfaffian-dirac}.

\section{Coda: $\autcor$ as one of five views}
\label{sec:bkm-coda}

The denominator algebra $\autcor$ is one of five views of
$\Delta_5$: View~I (automorphic section), View~II (Borcherds--Kac--Moody
denominator, this chapter), View~III (compact protected Pfaffian,
conditional), View~IV (mirror discriminant, conjectural), View~V
(singular theta lift). The cross-level identities are:
\begin{itemize}
\item View~I $\leftrightarrow$ View~II: theorem,
\cite{Borcherds95,GNII}.
\item View~II $\leftrightarrow$ View~V: theorem,
\cite{B}.
\item View~I $\leftrightarrow$ View~V: theorem,
\cite{GN,B}.
\item View~II $\leftrightarrow$ View~III: \emph{conditional},
Chapter~\ref{ch:pfaffian-dirac}, under
$(D0),(O1),(O1)^+,(O2)$.
\item View~II $\leftrightarrow$ View~IV: \emph{conjectural},
Chapter~\ref{ch:view4-mirror}.
\end{itemize}

The closing identity of this chapter is the recorded denominator
identity
\begin{equation}
 \operatorname{den}(\autcor)\;=\;64^{-1}\Delta_5(2Z),
\end{equation}
the Borcherds-weight characteristic
\begin{equation}
 \kappa_{\mathrm{BKM}}(\Delta_5)\;=\;5,
\end{equation}
and the active multiplicity formula
\begin{equation}
 \smult(\alpha(n,l,m))\;=\;f(nm,l)\qquad((n,l,m)\in\Gamma_{\mathrm{act}}).
\end{equation}

These three identities are the load-bearing target-side data of
View~II. The chapter is closed.


\providecommand{\Mp}{\mathbf{Mp}}
\providecommand{\Spin}{\mathbf{Spin}}
\providecommand{\Pf}{\operatorname{Pf}}
\providecommand{\Un}{\mathbf{U}}

\chapter{The singular theta lift}
\label{ch:singular-theta-lift}

\section{The singular theta lift of \texorpdfstring{\(\phi_{0,1}\)}{phi\_{0,1}}}
\label{sec:singular-theta-lift}

The automorphic, Borcherds--Kac--Moody, and singular-theta descriptions
of \(\Delta_5\) are theorems of Borcherds, Gritsenko, and Nikulin.
The Pfaffian description is relative to the retained Dirac--Igusa data,
and the mirror-discriminant description is conjectural.  The
exterior-square isogeny
\(\PSp_4(\mathbb R)\simeq\SO(3,2)_+\) identifies the genus-two Siegel
upper half-space and the type-IV symmetric domain attached to the
signature-\((3,2)\) lattice into one homogeneous space; the singular
theta integral of \cite{Borcherds95,B} sends the Weil-representation
theta component \(\phi^{\mathrm{Weil}}_{0,1}\) of \(\phi_{0,1}\) to
the weight-five Borcherds product \(\Delta_5\), and the
Gritsenko--Nikulin denominator identity then reads the lift as the BKM denominator on
\(\La^{2,1}_{II}\).  These identifications are theorems of
\cite{Borcherds95,B,GN,GNII}; the integral lattice, arithmetic
subgroup, and multiplier convention are fixed by the following
normalization.

\subsection{Notation and the canonical name}
\label{ssec:theta-lift-notation}

Three Borcherds-product handles for the weight-five Igusa cusp form
are distinguished.  The \emph{automorphic section}
\(\mathcal D_X=\Delta_5\) of the line
\(\mathcal L^5\otimes\nu_{\Delta_5}\) on
\(\Sp_4(\Z)\backslash\Hpl_2\) is the View-I name.  The \emph{BKM
denominator}
\(\operatorname{den}(\autcor)=64^{-1}\Delta_5(2Z)\)
on \(\La^{2,1}_{II}\) is the View-II name.  The \emph{protected
Pfaffian} \(\Pf_{\mathrm{prot}}(\mathfrak D_X^{\mathrm{DI}})\), defined on
\(K3\times E\) relative to the retained D0-HN datum, retained quotient-orientation
datum, chosen Weyl lifts, \((O2_{\mathrm{atlas}})\), and the finite
geometric Pfaffian datum \((P_{\mathrm{fin}})\), is the
View-III name.  The View-V canonical name is
\[
\Delta_5
\;=\;
\Theta\bigl(\phi^{\mathrm{Weil}}_{0,1}\bigr),
\]
the \emph{singular theta lift of the Weil component
\(\phi^{\mathrm{Weil}}_{0,1}\)} obtained from \(\phi_{0,1}\) by theta
decomposition.  Throughout
\(\phi_{0,1}\in J_{0,1}^{\mathrm{weak}}\) is the unique up to scale
weight-zero index-one weak Jacobi form on the Jacobi group
\(\Gamma^J_1=\SL_2(\Z)\ltimes\Z^2\) with constant Fourier coefficient
fixed by the K3 elliptic genus normalization
\(Z_{\mathrm{K3}}=2\phi_{0,1}\) of \cite[\S2.2]{EZ:1};
\(\phi^{\mathrm{Weil}}_{0,1}\) is its weight-\(-1/2\)
Weil-representation theta component, and \(\Theta\) is the Borcherds
singular-theta integral of \cite[\S6]{Borcherds95}, \cite[\S6]{B}
with the Igusa cusp trivialization
\(C_{\phi_{0,1}}=64\).  Its monic Borcherds product is
\(D_5=64^{-1}\Delta_5\).  The bare handle \(\Delta_5\) is reserved for
cross-view identifications.

\subsection{The weak Jacobi form \texorpdfstring{\(\phi_{0,1}\)}{phi\_{0,1}}}
\label{ssec:phi01}

\begin{definition}[Weak Jacobi form, weight zero, index one]
\label{def:weak-jacobi-phi01}
A weak Jacobi form of weight \(k\) and index \(m\) on the Jacobi group
\(\Gamma^J_1=\SL_2(\Z)\ltimes\Z^2\) is a holomorphic function
\(\phi:\Hpl\times\C\rightarrow\C\) satisfying, for
\(\begin{psmallmatrix}a&b\\c&d\end{psmallmatrix}\in\SL_2(\Z)\) and \((\lambda,\mu)\in\Z^2\),
\begin{align*}
\phi\!\left(\frac{a\tau+b}{c\tau+d},\ \frac{z}{c\tau+d}\right)
&=(c\tau+d)^k e^{2\pi i mcz^2/(c\tau+d)}\phi(\tau,z),\\
\phi(\tau,\,z+\lambda\tau+\mu)
&=e^{-2\pi i m(\lambda^2\tau+2\lambda z)}\phi(\tau,z),
\end{align*}
together with the cusp expansion
\[
\phi(\tau,z)=\sum_{n\geq0,\,l\in\Z}f(n,l)\,q^nr^l,
\qquad q=e^{2\pi i\tau},\ r=e^{2\pi iz},
\]
in which the coefficient depends only on the pair
\((4mn-l^2,\;l\bmod 2m)\) (Eichler--Zagier
\cite[Theorem 2.2]{EZ:1}).  For \(m=1\) the residue class is fixed by
the parity of \(4n-l^2\).  The
\emph{weight-zero index-one} weak Jacobi form \(\phi_{0,1}\) is fixed
by \(k=0\), \(m=1\), the discriminant condition
\(4n-l^2\geq -1\) and the normalization
\[
\phi_{0,1}(\tau,z)=r^{-1}+10+r+O(q).
\]
\end{definition}

The space \(J_{0,1}^{\mathrm{weak}}\) is one-dimensional over \(\C\)
(Eichler--Zagier \cite[Theorem 9.4 and Table 9, \S9]{EZ:1});
\(\phi_{0,1}\) generates it.  The theta-quotient presentation is
the input to the singular-theta integral below.

\begin{lemma}[Eichler--Zagier theta-quotient presentation]
\label{lem:phi01-presentations}
The weak Jacobi form \(\phi_{0,1}\) admits the closed form
\[
\phi_{0,1}(\tau,z)
=4\sum_{i=2}^{4}
\frac{\theta_i(\tau,z)^2}{\theta_i(\tau,0)^2},
\]
where \(\theta_2,\theta_3,\theta_4\) are the three even Jacobi theta
functions in the Mumford normalization (Eichler--Zagier
\cite[\S9, equation (9.4)]{EZ:1}); equivalently
\(\phi_{0,1}=Z_{\mathrm{K3}}/2\) is one half of the K3 elliptic genus
\cite{DMVV}.  The leading Fourier coefficients are
\[
f(0,0)=10,\quad f(0,\pm1)=1,\quad f(1,0)=108,\quad
f(1,\pm1)=-64,\quad f(1,\pm2)=10,
\]
\[
f(2,0)=808,\quad f(2,\pm1)=-513,\quad f(2,\pm2)=108,\quad f(2,\pm3)=1,
\]
in agreement with the Mukai--Gram normalization of
Appendix~\ref{appendix:mukai-gram-dictionary}.
The discriminant relation \(f(n,l)=f(n',l')\) when
\(4n-l^2=4n'-(l')^2\) (Eichler--Zagier \cite[Theorem~2.2]{EZ:1})
is the index-one specialization of the
\((4mn-l^2,l\bmod 2m)\) rule and organises these coefficients into
the equivalence classes
\[
\{f(1,0),\,f(2,\pm2)\}=108\ (\text{disc }4),\quad
\{f(0,0),\,f(1,\pm2)\}=10\ (\text{disc }0),\quad
\{f(0,\pm1),\,f(2,\pm3)\}=1\ (\text{disc }-1),
\]
each independently consistent with the verified expansion.
The finite arithmetic table
\[
\texttt{certificates/jacobi/k3e\_first\_relation\_window\_coefficients}
\]
records the computed rows \(\phi_{0,1}=r^{-1}+10+r+O(q)\), the finite
evenness checks \(f(n,l)=f(n,-l)\), and the finite
discriminant-dependence checks \(f(n,l)=c(4n-l^2)\) for the first
relation-closed window.  This table checks the displayed Fourier
window; the all-coefficient statement is the Eichler--Zagier theorem
cited above.
\end{lemma}

\begin{proof}
The closed form is standard in Eichler--Zagier
\cite[\S\S9--10]{EZ:1}; the K3 elliptic-genus identity
\(Z_{\mathrm{K3}}=2\phi_{0,1}\) is \cite[\S2.2]{EZ:1} together with the
sigma-model right-moving Witten index of \cite{DMVV}.  The constant
term \(f(0,0)=10\) is the right-moving-localized character of the K3
Ramond sector at \(L_0=c/24\); the remaining values follow from the
discriminant relation \(f(n,l)=f(n',l')\) when
\(4n-l^2=4n'-(l')^2\) (Eichler--Zagier \cite[Theorem 2.2]{EZ:1})
together with explicit theta-quotient evaluation, which fixes the
disc-\(3\) class (\(f(1,\pm1)=-64\)) and the disc-\(4\) class
(\(f(1,0)=f(2,\pm2)=108\)) by direct computation from the
theta-quotient of Lemma~\ref{lem:phi01-presentations}.  The squared
theta-leading coefficient
\(([q^{1/2}r^{1/2}s^{1/2}]\Delta_5)^2 = 64^2 = 4096\) of the
Borcherds-product expansion in
\S\ref{appendix:mukai-gram-dictionary} is the next
arithmetic check on these coefficients via the cusp-chamber product
\(\Delta_5(Z) = 64q^{1/2}r^{1/2}s^{1/2}\prod_{(n,l,m)\in\Gamma_{\mathrm{eff}}}
(1-q^nr^ls^m)^{f(nm,l)}\).
\end{proof}

\subsection{The signature-\texorpdfstring{\((3,2)\)}{(3,2)} lattice and the type-IV domain}
\label{ssec:lattice-3-2}

\begin{definition}[The signature-\((3,2)\) lattice]
\label{def:lambda-3-2}
The lattice \(\La^{3,2}\) is the orthogonal complement of the alternating
Pfaffian element
\(q_0=e_1\wedge e_3+e_2\wedge e_4\in\wedge^2\La^4\) inside the
signature-\((3,3)\) lattice
\(\bigl(\wedge^2\La^4,\,(\,,\,)\bigr)\simeq\mathrm{II}_{3,3}\),
where \(\La^4=\bigoplus_{i=1}^4\Z e_i\) and the bilinear form
\((u,v)\) is the Pfaffian pairing
\(u\wedge v=(u,v)\,e_1\wedge e_2\wedge e_3\wedge e_4\).  The vector
\(q_0\) has square \(-2\).  In the BKM sign convention it is the even discriminant-two lattice
\[
\La^{3,2}\;\simeq\;\HypPlan\oplus\HypPlan\oplus[2],
\qquad
\HypPlan=\begin{psmallmatrix}0&-1\\-1&0\end{psmallmatrix},
\]
written in the basis
\((f_1,f_2,f_3,f_{-2},f_{-1})\) of
the exterior-square model section.
\end{definition}

\begin{remark}[Signature conventions]
The lattice \(\La^{3,2}\) carries three positive and two negative
directions in the \((u,v)\)-pairing of
Definition~\ref{def:lambda-3-2}.  The \emph{primitive hyperbolic
sublattice} \(\La^{2,1}=\HypPlan\oplus[2]\) on which the BKM denominator
algebra \(\autcor\) is graded has signature \((2,1)\) in the fixed
convention: the \([2]\) factor and one hyperbolic
direction are positive, and the other hyperbolic direction is negative.
The decomposition
\(\La^{3,2}=\La^{2,1}\oplus\HypPlan\) splits an extra hyperbolic plane
on which the singular theta integral lives; this is the
\(\mathrm{II}_{2,2}\)-fibre at the standard cusp under the exceptional
isomorphism of \S\ref{ssec:exceptional-iso}.  No
\(\La^{2,1}\oplus\langle-2\rangle\) presentation is invoked; the
fixed convention is the \(\HypPlan\oplus[2]\) Nikulin form.
\end{remark}

Let
\[
\Gamma_{3,2}:=\wedge^2\Sp_4(\Z)/\{\pm\Id_4\}
\subset \mathrm P\Orth(\La^{3,2})_+
\]
be the projective exterior-square arithmetic group.  The type-IV
symmetric domain attached to \(\La^{3,2}\) is
\[
\Hpl^{\IV}=\bigl\{Z\in\Proj(\La^{3,2}\otimes\C)\,\bigm|\,
(Z,Z)=0,\ (Z,\overline Z)<0\bigr\},
\]
with positive component \(\Hpl^{\IV}_+\) singled out by
\(\Imag z_1>0\) in the parametrization
\({}^t(z_2^2-z_1z_3,z_3,z_2,z_1,1)z_0\) of
the exterior-square model section.  Its arithmetic quotient
\(\Gamma_{3,2}\backslash\Hpl^{\IV}_+\) is the moduli space on
which the Borcherds singular-theta lift of \cite{Borcherds95} produces
automorphic forms with prescribed singularities along rational quadratic
divisors.

\subsection{The exceptional isogeny \texorpdfstring{\(\PSp_4\simeq\SO(3,2)_+\)}{PSp(4) ~= SO(3,2)+}}
\label{ssec:exceptional-iso}

\begin{theorem}[Exterior-square exceptional isomorphism]
\label{thm:exceptional-iso}
The exterior-square representation
\(\wedge^2:\SL_4\rightarrow\Aut(\wedge^2\La^4,(\,,\,))\) restricts to a
homomorphism
\[
\wedge^2:\Sp_4(\Z)/\{\pm\Id_4\}\;\xrightarrow{\ \simeq\ }\;\Gamma_{3,2},
\]
realizing the exceptional isomorphism of adjoint real Lie groups
\(\PSp_4(\R)\simeq\SO(3,2)_+\), together with its spin lift on the
covering groups.  Under this identification the genus-two
Siegel upper half-space \(\Hpl_2\) embeds onto the positive component
\(\Hpl^{\IV}_+\) by
\[
\begin{psmallmatrix}z_1&z_2\\z_2&z_3\end{psmallmatrix}\;\longmapsto\;
{}^t(z_2^2-z_1z_3,\ z_3,\ z_2,\ z_1,\ 1)z_0,
\]
intertwining the \(\Sp_4(\Z)\) and \(\Gamma_{3,2}\) actions.
\end{theorem}

\begin{proof}
The Pfaffian form
\(u\wedge v=(u,v)e_1\wedge e_2\wedge e_3\wedge e_4\) is the canonical
\(\SL_4\)-invariant symmetric bilinear form on \(\wedge^2\La^4\); the
subgroup of \(\SL_4\) fixing \(q_0=e_1\wedge e_3+e_2\wedge e_4\) is the
symplectic group \(\Sp_4\) of the alternating form \(B_{q_0}\).
The orthogonal complement \(\La^{3,2}=q_0^\perp\) is preserved, and the
kernel of \(\wedge^2|_{\Sp_4}\) on \(q_0^\perp\) is \(\{\pm\Id_4\}\)
(the exterior-square model section).  This is the integral form of
the classical exceptional isomorphism \(\PSp_4(\R)\simeq\SO(3,2)_+\)
of \cite[Chapter II.6]{KNAPP:1}; the spin cover supplies the
corresponding lift with Lie algebra \(\mathfrak{sp}_4(\R)\).  This cover is distinct
from the metaplectic cover on the \(\Mp_2(\Z)\) input.  The displayed
substitution into the type-IV domain produces the commutative square of
the exterior-square model section.  The integral-Igusa form of the
identification \(\Hpl_2\simeq\Hpl^{\IV}_+\) is the content of
\cite[\S2]{Igusa1962}; see also \cite[Chapter II]{FREITAG:1} and
\cite[\S6]{VDGEER:1}.
\end{proof}

\begin{remark}[Three guises of one homogeneous space]
The same homogeneous space carries three names: the genus-two Siegel
domain \(\Sp_4(\R)/\Un(2)\), the type-IV symmetric domain
\(\SO(3,2)_+/(\SO(3)\times\SO(2))\), and the spin double cover
\(\Spin(3,2)/(\Un(2)\text{-cover})\).  The regularized theta integral
below uses the metaplectic cover \(\Mp_2(\Z)\) on the input variable
\(\tau\) through the Weil representation of the discriminant form of
\(\La^{3,2}\).  The integrality of \(\operatorname{wt}(\Delta_5)=5\)
places the resulting weight on the ordinary automorphic line; the multiplier system
\(\nu_{\Delta_5}\) records the residual reflection character on
\(W^{(2)}(\La^{2,1}_{II})\) (Lemma~\ref{lem:bkm-maass-restriction}).
\end{remark}

\subsection{The Borcherds theta correspondence and the exceptional isomorphism}
\label{ssec:borcherds-theta-correspondence}

\begin{theorem}[Regularized theta correspondence on \(\La^{3,2}\)]
\label{thm:regularized-theta-correspondence}
Let \(L=\La^{3,2}\).  Borcherds's regularized theta integral pairs a
weakly holomorphic form
\(\phi\in M^!_{-1/2}(\rho_L)\) for \(\Mp_2(\Z)\) with the
non-holomorphic Siegel theta kernel of \(L\), and produces an
automorphic form on \(\widetilde\Orth(L)_+\) of weight
\(c_\phi(0,0)/2\).  Its rational-quadratic divisor is determined by the
negative-index part of \(\phi\), and its Weyl-chamber product uses the
Fourier coefficients of \(\phi\) in all effective product degrees.  For
\(\phi=\phi^{\mathrm{Weil}}_{0,1}\), the pullback along
\[
\wedge^2:\Sp_4(\Z)/\{\pm\Id_4\}\longrightarrow\Gamma_{3,2}
\subset\Orth(\La^{3,2})
\]
is the Igusa cusp form \(\Delta_5\).
\end{theorem}
\hfill\textnormal{[proved; \cite{Borcherds95,B,Igusa1962,GN}]}

\begin{remark}[The two groups in the lift]
The metaplectic group in the construction is \(\Mp_2(\Z)\), acting on
the input form through the Weil representation \(\rho_{\La^{3,2}}\).
The genus-two group \(\Sp_4\) enters after the orthogonal automorphic
form has been constructed, through the exceptional isomorphism
\(\PSp_4(\R)\simeq\SO(3,2)_+\) and the integral exterior-square map of
Theorem~\ref{thm:exceptional-iso}.  Thus the theta input, the
orthogonal automorphic form, and the Siegel modular form are three
successive incarnations of one Borcherds product.
\end{remark}

\subsection{The Borcherds singular-theta integral}
\label{ssec:singular-theta}

Borcherds's correspondence assigns to a Weil-representation form
obtained from a Jacobi input a meromorphic orthogonal modular form by
integration against a vector-valued Siegel theta series.  The integral
is divergent at the lift's input weight; the Borcherds regularization
produces a meromorphic Siegel modular form, after the
\(\PSp_4\simeq\SO(3,2)_+\) pullback, with explicit divisor.

\begin{definition}[Siegel theta-series of \texorpdfstring{\(\La^{3,2}\)}{Lambda^{3,2}}]
\label{def:siegel-theta}
For the lattice \(\La^{3,2}\) of Definition~\ref{def:lambda-3-2}, the
Siegel theta-series at \(Z\in\Hpl^{\IV}_+\) and \(\tau\in\Hpl\) is
\[
\Theta_{\La^{3,2}}(\tau,Z)
=\sum_{\gamma\in A_{\La^{3,2}}}e_\gamma
\sum_{\lambda\in\La^{3,2}+\gamma}
\exp\!\bigl(\pi i\,\overline\tau\,(\lambda,\lambda)_{Z^+}
+\pi i\,\tau\,(\lambda,\lambda)_{Z^-}\bigr),
\]
where \((\,,\,)_{Z^\pm}\) are the positive and negative parts of the
\((u,v)\)-pairing under the orthogonal decomposition determined by
\(Z\) (Borcherds \cite[\S4]{Borcherds95}, \cite[\S4]{B}).  The
theta-series is a non-holomorphic modular form of weight
\((3/2,1)\) on \(\Mp_2(\Z)\) valued in the Weil representation of
\(\La^{3,2}\) on \(\C[A_{\La^{3,2}}]\).  The scalar sum over
\(\La^{3,2}\) is only the \(\gamma=0\) component.
\end{definition}

\begin{theorem}[Singular theta lift {\cite[Theorem 13.3]{Borcherds95}, \cite[Theorem 6.2]{B}}]
\label{thm:singular-theta-lift}
Let \(\phi\) be a weakly holomorphic vector-valued modular form for the
dual Weil representation of \(\La^{3,2}\) on \(\Mp_2(\Z)\) of weight
\(-1/2\) with integral Fourier coefficients
\(c_\phi(n,\lambda)\in\Z\) at the cusp.  The regularized integral
\[
\Theta(\phi)(Z)
=\int^{\mathrm{reg}}_{\SL_2(\Z)\backslash\Hpl}
\bigl\langle\phi(\tau),\Theta_{\La^{3,2}}(\tau,Z)\bigr\rangle\,v
\frac{du\,dv}{v^2},
\qquad\tau=u+iv,
\]
in the regularization sense of \cite[\S6]{B} produces a meromorphic
function on \(\Hpl^{\IV}_+\) with the following properties.
\begin{enumerate}[label=\textup{(\arabic*)}]
\item \emph{Modular invariance.} \(\Theta(\phi)\) is an automorphic
form of weight \(c_\phi(0,0)/2\) on
\(\widetilde\Orth(\La^{3,2})_+\) with multiplier system determined by
the Weil action on \(\La^{3,2}{}^*/\La^{3,2}\).
\item \emph{Borcherds product.} On the cone of definition,
\(\Theta(\phi)\) admits the Borcherds product expansion
\[
\Theta(\phi)(Z)
=C_\phi e^{-2\pi i(\rho,Z)}
\prod_{\substack{\alpha\in\La^{3,2}{}^*\\(\alpha,W)>0}}
	\bigl(1-e^{-2\pi i(\alpha,Z)}\bigr)^{c_\phi(-(\alpha,\alpha)/2,\,\alpha\bmod\La^{3,2})},
\]
with \(C_\phi\in\C^\times\) fixed by the chosen cusp trivialization,
Weyl vector \(\rho\) determined by \(\phi\), and a chamber \(W\) of the
positive cone.
\item \emph{Divisor.} The divisor of \(\Theta(\phi)\) is the linear
combination of rational quadratic divisors \(\lambda^\perp\)
weighted by \(c_\phi(-(\lambda,\lambda)/2,\lambda\bmod\La^{3,2})\)
over primitive \(\lambda\in\La^{3,2}{}^*\) with negative discriminant.
\end{enumerate}
The factor \(v\) is the invariant normalization for signature
\((3,2)\): the input has holomorphic weight \(-1/2\), the theta kernel
has weight \((3/2,1)\), and
\(v\langle\phi,\Theta_{\La^{3,2}}\rangle\) has weight \((0,0)\).
\end{theorem}
\hfill\textnormal{[proved; \cite{Borcherds95,B}]}

\begin{remark}[Two regularization conventions]
The Borcherds regularization \cite[\S6]{B} subtracts the divergent
\(v\rightarrow\infty\) contribution by analytic continuation in an
auxiliary parameter, and is the convention used here.  The
Harvey--Moore regularization \cite{HM} differs from the Borcherds
regularization by a finite renormalization of the Weyl vector and the
constant term, with no effect on the product expansion's exponents or
divisor.  The two conventions agree on quotient divisors
\(\Theta(\phi_1)/\Theta(\phi_2)\), on the BKM denominator identity, and on the weight identity
\(\operatorname{wt}\Theta(\phi)=c_\phi(0,0)/2\).
\end{remark}

\subsection{From \texorpdfstring{\(\phi_{0,1}\)}{phi\_{0,1}} to a vector-valued Weil-representation form}
\label{ssec:phi-to-weil}

The singular-theta integral takes a weakly holomorphic vector-valued
form on \(\Mp_2(\Z)\); the input \(\phi_{0,1}\) is a scalar weak Jacobi
form of index one.  Theta-decomposition reduces the latter to the
former.

\begin{lemma}[Theta-decomposition of \texorpdfstring{\(\phi_{0,1}\)}{phi\_{0,1}}]
\label{lem:theta-decomposition}
Let \(\phi_{0,1}\in J_{0,1}^{\mathrm{weak}}\) and decompose along the
elliptic theta-functions of index one,
\[
\phi_{0,1}(\tau,z)=h_0(\tau)\theta_{1,0}(\tau,z)+h_1(\tau)\theta_{1,1}(\tau,z),
\]
where \(\theta_{1,\mu}(\tau,z)=\sum_{\ell\equiv\mu\, (2)}
q^{\ell^2/4}r^\ell\) for
\(\mu\in\{0,1\}\) (Eichler--Zagier \cite[\S5]{EZ:1}).  Then
\(h=(h_0,h_1)\) is a weakly holomorphic vector-valued modular form of
weight \(-1/2\) for the dual Weil representation
\(\rho_{[2]}^\vee\), equivalently for the Weil representation of the
opposite discriminant form \([-2]\), on \(\Mp_2(\Z)\).  Under the
finite quadratic-module identification
\[
A_{[2]}\;\simeq\;A_{\HypPlan\oplus[2]}
\;\simeq\;A_{\La^{3,2}},
\]
it pairs with the \(\rho_{\La^{3,2}}\)-valued theta kernel of
Definition~\ref{def:siegel-theta}.  Its Fourier coefficients are exactly
\[
[q^{D/4}]h_\mu(\tau)=c(D),\qquad
c(4n-l^2)=f(n,l),\qquad D\equiv-\mu^2\pmod4,
\]
or equivalently \([q^{\,n-l^2/4}]h_{l\bmod 2}=f(n,l)\).  In the
index-one normalization the first coefficients are
\([q^0]h_0=10\) and \([q^{-1/4}]h_1=1\) reproducing
\(\phi_{0,1}=r^{-1}+10+r+O(q)\).  Since
\(\theta_{1,1}(\tau+1,z)=i\,\theta_{1,1}(\tau,z)\), invariance of
\(\phi_{0,1}\) under \(T\) forces
\(h_1(\tau+1)=-i\,h_1(\tau)\), the dual eigenvalue.  The vector \(h\)
is the weakly holomorphic input
\(\phi^{\mathrm{Weil}}_{0,1}\) of the weight and representation type
required by Theorem~\ref{thm:singular-theta-lift}.
\end{lemma}
\hfill\textnormal{[proved; \cite{EZ:1,Borcherds95,GN}]}

\begin{proof}
Theta-decomposition along the lattice \(\La_{\mathrm{ind}}=\Z\) inside
the Jacobi-elliptic variable is Eichler--Zagier
\cite[Theorem 5.1]{EZ:1}.  The two-component vector \(h\) is weakly
holomorphic of weight \(-1/2\) on \(\Mp_2(\Z)\) with the Weil
representation of the rank-one lattice \(\langle 2\rangle\).  Comparing
the coefficient of \(r^\ell\) gives
\([q^{\,n-\ell^2/4}]h_{\ell\bmod2}=f(n,\ell)\), hence the discriminant
formula above.  The hyperbolic plane \(\HypPlan\) is unimodular, so
adjoining one or two copies of \(\HypPlan\) does not change the
discriminant form.  Thus
\[
(\HypPlan\oplus[2])^*/(\HypPlan\oplus[2])
\simeq
[2]^*/[2]
\simeq
(\La^{3,2})^*/\La^{3,2},
\]
with the same finite quadratic form.  This discriminant-form isometry,
is the transport of \(h\) to the required weakly holomorphic
vector-valued form for the Weil representation of \(\La^{3,2}\).
The constant Fourier coefficient
\(c_{\phi^{\mathrm{Weil}}_{0,1}}(0,0)=f(0,0)=10\) is preserved under
this discriminant-form identification because the trivial coset
\(0\in\La^{3,2}{}^*/\La^{3,2}\) matches the trivial coset on the
rank-one factor.  Borcherds
\cite[\S15]{Borcherds95} records the explicit computation for the
\(\Sp_4\)-case in the equivalent paramodular language.
\end{proof}

\begin{remark}[Index-one normalization of \(c_1(0)\)]
The constant Fourier coefficient
\(c_1(0):=f(0,0)=10\) of \(\phi_{0,1}\) is the Eichler--Zagier
index-one normalization of the K3 elliptic genus.  The
universal identity \(\kappa_{\mathrm{BKM}}(\Phi_N)=c_N(0)/2\) of
\cite{Gritsenko99} names \(N\) as the paramodular level of the input
denominator;
at \(N=1\) the input form is \(\phi_{0,1}\) and \(c_1(0)=10\), giving
\[
\kappa_{\mathrm{BKM}}(\Delta_5)=c_1(0)/2=5.
\]
The integer-versus-half-integer ambiguity for the metaplectic weight
factor is killed by the integrality of \(c_1(0)/2=5\), and the
Eichler--Zagier normalization is the row \(F_1\) normalization in the
Cl\'ery--Gritsenko catalogue \cite{GC}.  On the other seven rows the
symbol \(c(0)\) is indexed by the pair \((t,N)\), not by a CHL twining
order alone.  No \(c_N(0)=2k_N+1\) half-integer convention enters at
\(N=1\).
\end{remark}

\subsection{Theta lift identifies \texorpdfstring{\(\Delta_5\)}{Delta\_5}}
\label{ssec:theta-lift-delta5}

The lift produces an automorphic form on
\(\widetilde\Orth(\La^{3,2})_+\), pulled back along
\(\wedge^2:\Sp_4(\Z)/\{\pm\Id_4\}\xrightarrow{\simeq}\Gamma_{3,2}\)
of Theorem~\ref{thm:exceptional-iso} to a Siegel modular form on
\(\Sp_4(\Z)\).

\begin{theorem}[\texorpdfstring{\(\Delta_5\)}{Delta\_5} as singular theta lift]
\label{thm:delta5-as-theta-lift}
The singular-theta lift of the weak Jacobi form \(\phi_{0,1}\) is the
weight-five Igusa cusp form:
\[
\Theta\bigl(\phi^{\mathrm{Weil}}_{0,1}\bigr)\!\bigm|_{\Hpl_2}
\;=\;\Delta_5,
\]
the right-hand side written in the convention
\(\mathcal D_X=\Delta_5\) of \S\ref{ssec:theta-lift-notation} and the
left-hand side regularized in the Borcherds sense of
\cite[\S6]{B}.  Equivalently, in the BKM normalization,
\[
\Theta\bigl(\phi^{\mathrm{Weil}}_{0,1}\bigr)
=64\cdot\operatorname{den}(\autcor)\!\bigm|_{Z\mapsto Z/2},
\]
i.e.\ the lift is \(\Delta_5\) on \(\Hpl_2\) and
\(64^{-1}\Delta_5(2Z)\) on \(\La^{2,1}_{II}\)-coordinates of
\(\Hpl^{\IV}_+\).
\end{theorem}
\hfill\textnormal{[proved; \cite{Borcherds95,B,GN,GNII}]}

\begin{proof}
The Borcherds singular-theta lift in
Theorem~\ref{thm:singular-theta-lift} produces a meromorphic
automorphic form
\(\Theta(\phi^{\mathrm{Weil}}_{0,1})\) on
\(\widetilde\Orth(\La^{3,2})_+\) of weight \(c_1(0)/2=5\).  Its
divisor is the linear combination of rational quadratic divisors
weighted by \(c_{\phi^{\mathrm{Weil}}_{0,1}}(-(\lambda,\lambda)/2,
\lambda\bmod\La^{3,2})\).  For \(\phi_{0,1}\), the
negative-discriminant coefficients \(f(0,\pm1)=1\) give the simple
reflective Humbert divisor of \(\Delta_5\), while \(f(0,0)=10\) gives
the weight \(5\).  The finite packet
\(\texttt{certificates/automorphic/delta5\_humbert\_divisor}\)
records this divisor support and order-one multiplicity separately from
the pole-order-two statement for \(\Delta_5^{-2}\).  The
Borcherds-product expansion of
Theorem~\ref{thm:singular-theta-lift}(2), specialized to
\(\phi^{\mathrm{Weil}}_{0,1}\), reads
\[
\Theta(\phi^{\mathrm{Weil}}_{0,1})
=C_{\phi_{0,1}}q^{1/2}r^{1/2}s^{1/2}
\prod_{(n,l,m)\in\Gamma_{\mathrm{eff}}}
(1-q^nr^ls^m)^{f(nm,l)}
\]
in the cusp chamber, with \(\Gamma_{\mathrm{eff}}\) the cusp-positive
semigroup and Weyl vector
\(\rho_0=(1/2,1/2,1/2)\).  The monic Borcherds product is
\(D_5(Z)\) of \S\ref{ssec:theta-lift-notation}; the Igusa
theta-constant cusp trivialization fixes
\(C_{\phi_{0,1}}=64\), hence
\(\Theta(\phi^{\mathrm{Weil}}_{0,1})=\Delta_5\).  The pullback to \(\Hpl_2\) along
\(\wedge^2:\Sp_4(\Z)/\{\pm\Id_4\}\rightarrow\Gamma_{3,2}\) of
Theorem~\ref{thm:exceptional-iso} identifies
\(\Theta(\phi^{\mathrm{Weil}}_{0,1})\!\bigm|_{\Hpl_2}\) with the
\(\Sp_4(\Z)\)-cusp form whose Borcherds product is exactly the Igusa
cusp form \(\Delta_5\) of weight five.  The BKM normalization
\(64^{-1}\Delta_5(2Z)\) on \(\La^{2,1}_{II}\) is
Gritsenko--Nikulin \cite[Theorem 2.1]{GN}, \cite[Theorem 2.1, Section
5.1.1]{GNII}: the doubled lattice argument \(Z\mapsto 2Z\) corresponds
to the polarization change \(\alpha:\Gamma_{\mathrm{ind}}
\rightarrow\La^{2,1}_{II}\) of \S\ref{ssec:weyl-vector-rho}.  Both
sides have the same Borcherds-product expansion, and their leading
coefficients agree by the theta-constant
\([q^{1/2}r^{1/2}s^{1/2}]\Delta_5=64\).
\end{proof}

\begin{corollary}[Borcherds--Gritsenko--Nikulin pentadic identity]
\label{cor:five-views-coincidence}
The theorem-level automorphic, BKM-denominator, and theta-lift readings
coincide:
\begin{align*}
\mathrm{(View\ I)}\ \mathcal D_X
&=\mathrm{(View\ II)}\ 64\cdot\operatorname{den}(\autcor)\!\bigm|_{Z\mapsto Z/2}\\
&=\mathrm{(View\ V)}\ \Theta\bigl(\phi^{\mathrm{Weil}}_{0,1}\bigr).
\end{align*}
The Pfaffian identification
\(\Pf_{\mathrm{prot}}(\mathfrak D_X^{\mathrm{DI}})=\Delta_5\) of View III
holds on \(K3\times E\) at the Pfaffian and orientation level, granted the
retained D0-HN datum, retained quotient-orientation datum, chosen Weyl lifts,
retained \((O2_{\mathrm{atlas}})\) wall datum, and retained finite
geometric Pfaffian datum \((P_{\mathrm{fin}})\) of
Theorem~\ref{thm:G-four-obstruction-criterion} of
Appendix~\ref{app:obstruction-discharges}.  The lift does not
provide the level-\(\mathsf A\)
gravity-line operator algebra; a 3d-gravity path-integral realization
requires the level-\(\mathsf A\) Hall--Borcherds residual, while
Definition~\ref{def:G-vol2-hall-borcherds} records only the
level-\(\mathsf Z\) trace comparison.
\end{corollary}
\hfill\textnormal{[proved for Views I, II, V; conditional for View III]}

\subsection{The multiplier system and the Weyl vector}
\label{ssec:weyl-vector-rho}

\begin{theorem}[Multiplier system of the lift]
\label{thm:multiplier-system}
The lift \(\Theta(\phi^{\mathrm{Weil}}_{0,1})\) carries the
\(\Sp_4(\Z)\)-multiplier system \(\nu_{\Delta_5}\) of Maass
\cite{MAASS:1}, equivalently described as the character
\(v_{\eta}^{12}\times v_H\) on the Jacobi parabolic and as the
Gritsenko character on the BKM side
\cite[Proposition 2.1]{GN}, \cite[\S5.1.1, case \((t=1,\mathrm{II},0)\)]{GNII}.
On the type-II generators of the Weyl group
\(W^{(2)}(\La^{2,1}_{II})\),
\[
\nu_{\Delta_5}(s_{\delta_i})=-1,\qquad i=1,2,3,
\]
and on the chamber automorphism factor \(\Aut(\Poly_{II})\simeq S_3\)
the multiplier is trivial:
\[
\nu_{\Delta_5}(s_{f_{-2}-f_2})
=\nu_{\Delta_5}(s_{f_2-f_3})
=\nu_{\Delta_5}(s_{f_{-2}-f_3})=+1.
\]
This is exactly the multiplier of
Lemma~\ref{lem:bkm-maass-restriction}.
The finite automorphic packet
\(\texttt{certificates/automorphic/delta5\_maass\_character}\)
checks these generator values, the order-two relation, and the
automorphic-line placement \( \mathcal L^5\otimes\nu_{\Delta_5} \);
the packet is not an orientation construction.
\end{theorem}
\hfill\textnormal{[proved; \cite{MAASS:1,GN,GNII}]}

\begin{remark}[Singular-theta multiplier vs Maass multiplier]
The Borcherds singular-theta lift produces a multiplier system on
\(\widetilde\Orth(\La^{3,2})_+\) determined by the Weil action on
\(\La^{3,2}{}^*/\La^{3,2}\) (Borcherds \cite[Theorem 13.3]{Borcherds95}).
For the present \(\La^{3,2}\), the discriminant group has order two,
generated by \(f_3/2\) in the convention of
Remark~\ref{rem:discriminant-conventions}.  The finite Weil module is
the two-dimensional module attached to this discriminant form.  The
Borcherds multiplier determined by that module and by the Weyl vector
restricts on the Borcherds chamber to the same reflection character.
Pulled back along
\(\wedge^2:\Sp_4(\Z)/\{\pm\Id_4\}\rightarrow\Gamma_{3,2}\) it
matches Maass's multiplier formula for the product of the ten even
theta constants \cite[(1.3)--(1.5)]{GN}, \cite{MAASS:1}: the
\(-1\)-eigenvalues on the three type-II simple reflections and the
\(+1\)-eigenvalues on the three chamber automorphisms.  The two
descriptions are not independent — they are the two sides of the same
character — and their agreement is part of the regularized theta
correspondence of Theorem~\ref{thm:regularized-theta-correspondence}.
\end{remark}

\begin{theorem}[Weyl vector of the lift]
\label{thm:weyl-vector-lift}
The Weyl vector of the Borcherds product
\(\Theta(\phi^{\mathrm{Weil}}_{0,1})\) in the type-II chamber
\(\Poly_{II}\subset\La^{2,1}_{II}\otimes\R\) is
\[
\rho=\tfrac12\delta_1+\tfrac12\delta_2+\tfrac12\delta_3
=f_2-\tfrac12 f_3+f_{-2},
\]
satisfying \((\rho,\delta_i)=-1\) for \(i=1,2,3\), with simple-real-root
Cartan matrix
\[
A=\bigl((\delta_i,\delta_j)\bigr)
=\begin{pmatrix}2&-2&-2\\-2&2&-2\\-2&-2&2\end{pmatrix}.
\]
This is the Weyl vector of \(\autcor\) on \(\La^{2,1}_{II}\) of
the Weyl-chamber section, identified with the Weyl vector of
the Borcherds product by \cite[Theorem 4.1]{GN}.
The finite lattice packet
\(\texttt{certificates/lattice/type\_ii\_weyl\_vector}\) checks the
Cartan matrix, discriminant \(-32\), dual-lattice membership of
\(\rho\), and the displayed pairings.
\end{theorem}
\hfill\textnormal{[proved; \cite{GN,GNII}]}

\begin{proof}
The Weyl vector of this Borcherds product is determined by the
\(q^0\)-row of \(\phi\) at the cusp.  For
\(\phi_{0,1}=r^{-1}+10+r+O(q)\), the negative-discriminant coefficients
are \(f(0,\pm1)=1\), and the constant coefficient is \(f(0,0)=10\); the
Weyl-vector formula of \cite[Theorem 13.3, equation (13.10)]{Borcherds95}
yields \(\rho=(1/2,1/2,1/2)\) in cusp coordinates, transported to
\(\La^{2,1}_{II}\) by the polarization change
\(\alpha:\Gamma_{\mathrm{ind}}\rightarrow\La^{2,1}_{II}\) of
the Weyl-chamber section.  This recovers
\(\rho=f_2-\tfrac12 f_3+f_{-2}\); the Cartan matrix is the Gram matrix
of \(\{\delta_1,\delta_2,\delta_3\}\) computed in
the Weyl-chamber section.
\end{proof}

\subsection{Compatibility with View II: lift's image is the BKM denominator}
\label{ssec:compat-view-ii}

\begin{theorem}[Lift--denominator compatibility {\cite[Proposition 3.1]{GN}}]
\label{thm:lift-denominator}
The image of the singular-theta lift,
\(\Theta(\phi^{\mathrm{Weil}}_{0,1})=\Delta_5\), satisfies the
Weyl--Kac--Borcherds denominator identity for the Lorentzian
BKM superalgebra \(\autcor\) on \(\La^{2,1}_{II}\):
\[
64^{-1}\Delta_5(2Z)
=e^{-2\pi i(\rho,Z)}
\prod_{\alpha\in\mathcal R_+}
\bigl(1-e^{-2\pi i(\alpha,Z)}\bigr)^{\smult(\alpha)},
\]
with positive-root system \(\mathcal R_+\), Weyl vector \(\rho\) of
Theorem~\ref{thm:weyl-vector-lift}, and signed root supermultiplicities
\[
\smult\bigl(\alpha(n,l,m)\bigr)=f(nm,l),
\qquad (n,l,m)\in\Gamma_{\mathrm{act}}.
\]
The simple imaginary-root data are the Gritsenko--Nikulin coefficients
\(m(a)\) on negative-norm rays and \(\tau(a)\) on isotropic rays of
the Weyl-chamber section.
\end{theorem}
\hfill\textnormal{[proved; \cite{GN,GNII}]}

\begin{proof}
The Borcherds-product expansion
\(\Delta_5=64\,q^{1/2}r^{1/2}s^{1/2}
\prod(1-q^nr^ls^m)^{f(nm,l)}\) of
Theorem~\ref{thm:delta5-as-theta-lift} is the BKM denominator product
for \(\autcor\) after the polarization change
\(\alpha:\Gamma_{\mathrm{ind}}\rightarrow\La^{2,1}_{II}\): the exponent
\(f(nm,l)\) is the signed multiplicity of the root
\(\alpha(n,l,m)\), and the doubled argument \(2Z\) converts
\(\exp(-\pi i(\alpha,z))\) into the standard denominator character
\(\exp(-2\pi i(\alpha,z))\).
The finite packet
\(\texttt{certificates/lattice/product\_chamber\_root\_map}\)
checks these two target-side steps: the row-wise root map into
\(\La^{2,1}_{II}\) and the \(Z\mapsto2Z\) character conversion.
The algebra \(\autcor\) is the Gritsenko--Nikulin algebra of
\cite[\S\S3--4, Proposition 3.1]{GN}; the denominator identity is read
off the product expansion as in Chapter~\ref{ch:view2-bkm}.
Equivalently, the imaginary-simple-root characters \(m(a)\) and
\(\tau(a)\) are the additive data of
\cite[Theorem 2.1, equations (2.1)--(2.3)]{GNII}.
\end{proof}

\subsection{Compatibility with View I and the cross-level identities}
\label{ssec:compat-view-i}

\begin{theorem}[Borcherds--Gritsenko--Nikulin three-way agreement]
\label{thm:bgn-three-way}
The three identifications
\begin{align*}
&\Theta(\phi^{\mathrm{Weil}}_{0,1})=\Delta_5
&&\text{\textup{(}Borcherds 1995, 1998\textup{)}},\\
&64^{-1}\Delta_5(2Z)=\operatorname{den}(\autcor)
&&\text{\textup{(}Gritsenko--Nikulin 1998\textup{)}},\\
&\mathcal D_X=\Delta_5
&&\text{\textup{(}Igusa 1962\textup{)}}
\end{align*}
are mutually consistent:
\(\Theta(\phi^{\mathrm{Weil}}_{0,1})=\mathcal D_X
=64\cdot\operatorname{den}(\autcor)\!\bigm|_{Z\mapsto Z/2}\) on
\(\Hpl^{\IV}_+\), and the multiplier system
\(\nu_{\Delta_5}\) of Theorem~\ref{thm:multiplier-system} is common.
\end{theorem}
\hfill\textnormal{[proved; \cite{Igusa1962,Borcherds95,B,GN,GNII}]}

\begin{proof}
Theorem~\ref{thm:delta5-as-theta-lift} is identification (1).
Theorem~\ref{thm:lift-denominator} is identification (2).
Identification (3) is the Igusa identification
\(\mathcal D_X=\Delta_5\) of \cite[Theorem of \S2]{Igusa1962},
transported through the exterior-square model section.  All three sides have the
Borcherds-product expansion
\(\Delta_5=64\,q^{1/2}r^{1/2}s^{1/2}\prod(1-q^nr^ls^m)^{f(nm,l)}\)
in the cusp chamber and pull back along the exceptional isomorphism of
Theorem~\ref{thm:exceptional-iso} to the same automorphic form on
\(\Sp_4(\Z)\).  The multiplier-system agreement is
Theorem~\ref{thm:multiplier-system}.
\end{proof}

\subsection{The weight-five universal Borcherds identity at \(N=1\)}
\label{ssec:universal-weight-identity}

\begin{theorem}[Universal Borcherds weight identity, \(N=1\)]
\label{thm:universal-weight-identity}
At paramodular level \(N=1\) the Vol III universal Borcherds-weight
identity
\(\kappa_{\mathrm{BKM}}(\Phi_N)=c_N(0)/2\)
specializes to
\[
\kappa_{\mathrm{BKM}}(\Delta_5)
=\frac{c_1(0)}{2}
=\frac{f(0,0)}{2}
=\frac{10}{2}=5.
\]
This is the row \((N=1,\,\operatorname{wt}=5,\,c_1(0)=10)\) of the
Cl\'ery--Gritsenko catalogue \cite{GC}, and is the value entered in
the K3\(\times E\) \(\kappa_\bullet\)-spectrum
\(\{\kappa_{\mathrm{cat}}=0,\kappa_{\mathrm{ch}}^{\mathrm{Hodge}}=0,
\kappa_{\mathrm{ch}}^{\mathrm{Heis}}=3,\kappa_{\mathrm{BKM}}=5,
\kappa_{\mathrm{fiber}}=24\}\) of Vol III.
\end{theorem}
\hfill\textnormal{[proved; \cite{Borcherds95,GC}; Vol III universal
identity]}

\begin{proof}
The lift's weight is \(c_1(0)/2\) by
Theorem~\ref{thm:singular-theta-lift}(1), and
\(c_1(0)=f(0,0)=10\) by Lemma~\ref{lem:phi01-presentations} and
Lemma~\ref{lem:theta-decomposition}.  The
universal identity
\(\kappa_{\mathrm{BKM}}(\Phi_N)=c_N(0)/2\) is
\cite[\S\S2--3]{Gritsenko99}, generalized in Vol III across the eight
Cl\'ery--Gritsenko rows; at \(N=1\) it is exactly the Borcherds
weight formula of Theorem~\ref{thm:singular-theta-lift}(1).
\end{proof}

\begin{remark}[Subscript discipline]
The bare symbol \(\kappa\) does not appear in any theorem-level identity.
Every occurrence carries the subscript fixing its
construction reading
(\(\kappa_{\mathrm{cat}},\,\kappa_{\mathrm{ch}}^{\mathrm{Hodge}},\,
\kappa_{\mathrm{ch}}^{\mathrm{Heis}},\,\kappa_{\mathrm{BKM}},\,
\kappa_{\mathrm{fiber}}\)).  The additive expression
\(\kappa_{\mathrm{BKM}}=\kappa_{\mathrm{ch}}+\chi(\mathcal O_{K3})\) is a
mixed-reading collision.  With
\(\kappa_{\mathrm{ch}}=\kappa_{\mathrm{ch}}^{\mathrm{Heis}}\) it gives the
accidental equality \(3+2=5\) at \(N=1\), but it has no rowwise
extension to \(N\in\{2,3,4,6\}\).  The universal Borcherds-weight identity
\(\kappa_{\mathrm{BKM}}(\Phi_N)=c_N(0)/2\) is the invariant statement;
at \(N=1\) the right hand side is \(c_1(0)/2=5\).
\end{remark}

\subsection{Beilinson cut: the lift is the structural content}
\label{ssec:beilinson-cut-view-v}

The five views of \S\ref{ch:pentadic} lie on the Beilinson chain
\(\mathsf P\rightarrow\mathsf C\rightarrow\mathsf S\rightarrow
\mathsf Z\rightarrow\mathsf A\) of Vol II.  View V's primitive content is
the singular-theta integral
\(\Theta:J_{0,1}^{\mathrm{weak},\mathrm{Weil}}
\rightarrow M_{*}(\widetilde\Orth(\La^{3,2})_+)\) of
Theorem~\ref{thm:singular-theta-lift}, sending Jacobi forms to
Borcherds products.  Its value at \(\phi_{0,1}\) is the automorphic
section \(\Delta_5\), the level-\(\mathsf S\) image of the lift.

\begin{remark}[Beilinson level]
The integral operator \(\Theta\) lives at level \(\mathsf C\) of the
universal chain: it is a chart-dependent passage from a primitive
Jacobi-form input to its Borcherds product, dependent on the
Weil representation of \(\La^{3,2}\) and on the regularization of
\cite[\S6]{B}.  The scalar \(\Delta_5\) lives at level \(\mathsf S\):
it is the automorphic value of the lift on \(\phi_{0,1}\).
The retained \(\mathfrak D_X^{\mathrm{DI}}\) datum of View III is the
conditional Pfaffian construction whose protected Pfaffian recovers
\(\Delta_5\); the BKM superalgebra \(\autcor\) of View II lives at
level \(\mathsf S\) under the trace lane of Vol II.  The primitive
object in this passage is the integral
operator \(\Theta\), and the automorphic, BKM, Pfaffian, and theta-lift
names of \(\Delta_5\) are readings, not primitive objects.
\end{remark}

\subsection{Consistency of the lattice conventions}
\label{ssec:consistency-conventions}

\begin{remark}[\(\La^{3,2}\) versus \(\La^{2,1}\oplus\langle-2\rangle\)]
The signature-\((3,2)\) lattice \(\La^{3,2}\) of
Definition~\ref{def:lambda-3-2} decomposes as
\(\HypPlan\oplus\HypPlan\oplus[2]\) on the integral level.  The
sublattice \(\La^{2,1}=\HypPlan\oplus[2]\) is primitive in
\(\La^{3,2}\); its orthogonal complement is the second hyperbolic
plane \(\HypPlan\), not the rank-one form \(\langle-2\rangle\) (which
would have wrong signature).  The polarization change
\(\alpha:\Gamma_{\mathrm{ind}}\rightarrow\La^{2,1}_{II}\) sends
\((n,l,m)\mapsto 2nf_2-lf_3+2mf_{-2}\), so the BKM grading lattice is
\(\La^{2,1}_{II}\), the index-four sublattice of \(\La^{2,1}\) with
primitive \(f_3\)-coordinate.  The doubled argument \(Z\mapsto 2Z\) in
\(\operatorname{den}(\autcor)=64^{-1}\Delta_5(2Z)\) of
Gritsenko--Nikulin changes the exponential convention from
\(\exp(-\pi i(\alpha,z))\) to \(\exp(-2\pi i(\alpha,z))\).  The
\(\mathrm{II}_{2,2}\) presentation of \(\La^{3,2}\) used in some
references (Borcherds 1995 \cite[\S15]{Borcherds95},
Borcherds 1998 \cite[Example 16.4]{B}) is the same lattice expressed via two
hyperbolic planes plus the \([2]\)-form, and produces the same Borcherds
product \(\Delta_5\) under the same lift.  No
\(\mathrm{II}_{3,2}\) versus \(\mathrm{II}_{2,1}\oplus\langle-2\rangle\)
ambiguity remains.
\end{remark}

\begin{remark}[Multiplier and metaplectic accounting]
The Borcherds singular-theta lift produces an automorphic form on the
integral form \(\widetilde\Orth(\La^{3,2})_+\), pulled back to
\(\Sp_4(\Z)\) by Theorem~\ref{thm:exceptional-iso}.  Since
\(c_1(0)/2=5\) is integral, the resulting automorphy factor is defined on
\(\Sp_4(\Z)\).  The metaplectic cover enters
on the \(\Mp_2(\Z)\) input side through
\(\rho_{\La^{3,2}}\); after restriction to integral weight its
central sign does not alter the \(\Sp_4(\Z)\) automorphy factor.  The
residual reflection character is exactly the
multiplier \(\nu_{\Delta_5}\) of Theorem~\ref{thm:multiplier-system}.
For the CHL constants at \(N\in\{1,2,3,6\}\), the numbers
\(c_N(0)/2=(5,3,2,1)\) are integral, and the metaplectic cover
belongs to the input Weil representation while the resulting
automorphic forms live on the corresponding integral orthogonal or
Siegel arithmetic groups.  At \(N=4\) the constant \(c_4(0)=3\) is
odd and the weight \(c_4(0)/2=\tfrac32\) is half-integral
(Govindarajan--Krishna \cite{GK10Composite}); \(\Phi_4\) is a genuine
form on the genus-two metaplectic cover.
\end{remark}

\begin{remark}[The eta-multiplier \texorpdfstring{\(v_\eta^{12}\)}{v\_eta\^12}]
The character of \(\Delta_5\) on the Jacobi parabolic of \(\Sp_4(\Z)\)
is described in \cite{Gritsenko99} as \(v_\eta^{12}\times v_H\), where
\(v_\eta\) is the eta-multiplier on \(\SL_2(\Z)\) and \(v_H\) is the
Heisenberg multiplier.  The first four Cl\'ery--Gritsenko rows carry
eta exponents \(12,6,4,3\) for \(N=1,2,3,4\), respectively
\cite[Theorem 1.4]{GC}; the present form is the \(N=1\) row.
Equivalently, after the lift through
the exceptional isomorphism, \(v_\eta^{12}\) remains the eta multiplier
in the Jacobi--Heisenberg presentation; it is not a discriminant
character of \(\mathrm{II}_{2,2}\), since \(\mathrm{II}_{2,2}\) is
unimodular.  No \(v_\eta^N\) inconsistency between
\cite{Gritsenko99}, \cite{GN}, and \cite{GC} arises at \(N=1\); the
three computations agree.
\end{remark}

\subsection{Theta-lift boundary}
\label{ssec:theta-lift-boundary}

The singular theta lift is one realization of \(\Delta_5\) among five.
The lift's primitive content is the Borcherds regularized theta
operator \(\Theta:\phi^{\mathrm{Weil}}\mapsto\Theta(\phi^{\mathrm{Weil}})\)
of Theorem~\ref{thm:singular-theta-lift}; its value at
\(\phi^{\mathrm{Weil}}_{0,1}\) is the weight-five Igusa cusp form.  The five-view
identification
\(\mathcal D_X=64\cdot\operatorname{den}(\autcor)\!\bigm|_{Z\mapsto Z/2}
=\Theta(\phi^{\mathrm{Weil}}_{0,1})\) is the theorem-level
automorphic--BKM--theta identity; its Pfaffian extension to View III
is conditional on the retained \(K3\times E\) data of
Appendix~\ref{app:obstruction-discharges}.  The mirror-discriminant View IV is the
remaining conjectural cross-level identity, recorded in the
mirror-discriminant chapter (Chapter~8) and not used here.

Borcherds's theta-correspondence formulation makes two structural facts visible.
First, the weight-five integrality is forced by
\(c_1(0)=10\) and the integrality of \(c_1(0)/2\); no metaplectic cover
is needed at the weight level.  Second, the BKM-denominator structure
of View II is not an extra postulate but a theorem about the lift's
image: the Gritsenko--Nikulin denominator identity reads the lift's
Borcherds product as the denominator of \(\autcor\) on
\(\La^{2,1}_{II}\) (Theorem~\ref{thm:lift-denominator}), and the
Borcherds--Gritsenko--Nikulin three-way agreement
(Theorem~\ref{thm:bgn-three-way}) crystallizes the lift's value in
three convertible normalizations.  The universal identity
\(\kappa_{\mathrm{BKM}}(\Phi_N)=c_N(0)/2\) is, at
\(N=1\), precisely the Borcherds-weight formula
\(\operatorname{wt}\Theta(\phi)=c_\phi(0,0)/2\) of
Theorem~\ref{thm:singular-theta-lift}(1) applied to
\(\phi=\phi_{0,1}\).


% ------------------------------------------------------------------
% Part II: Built — the Dirac–Igusa pro-object
% ------------------------------------------------------------------
\part{Built}

\chapter{The Dirac--Igusa pro-object $\mathfrak D_X^{\mathrm{DI}}$
on $K3\times E$ (View ③)}
\label{chap:dirac-igusa-pro-object}

The Igusa cusp form $\Delta_5$ of weight $5$ on
$\mathrm{Sp}(2,\Z)\backslash\mathcal H_2$ has three established
identities: the Borcherds--Gritsenko--Nikulin section (View ①), the
denominator of a generalized Borcherds--Kac--Moody Lie superalgebra
$\mathfrak g_{\Delta_5}$ on the root lattice $\Lambda_{II}^{2,1}$
(View ②), and the singular theta lift of the weight-zero index-one
weak Jacobi form $\phi_{0,1}$ (View ⑤).  The
View~\textcircled{3} datum is the protected Pfaffian line
and compact source comparison for a finite
\(E_3\)-prefactorization algebra on the \(K3\times E\) formal target,
presented at finite stage as the
Dirac--Igusa pro-object
\[
\mathfrak D_X^{\mathrm{DI}}
=\varprojlim_R\bigl(
\mathcal A_{X,R}^{E_3},
F_{X,R}^{\mathrm{hyb}},
\Gamma_{X,R},
\Pi_X,
\Phi_R,
o_R,
H_R,
P_R^\Pi,
C_{X,R},
\Theta_{\mathrm{Kos},R},
\mathscr L_{\mathrm{Pf},R},
\operatorname{pf}_{X,R},
\varepsilon_{o,R},
\mathrm{Rec}_R\bigr)
\]
with eight finite constructions and four named obstructions.

Let $S$ be a smooth complex K3 surface and let $E$ be a
smooth elliptic curve.  Singular K3 limits, nodal or cuspidal elliptic
degenerations, imprimitive curve-class contributions, the $s=0$
compact Hall cusp, and multi-component charge closures belong to
separate boundary data.  The standing finite-stage hypotheses
are $[\textup{K3$\times$E-fin}]$ (finite-type semistable substacks at
bounded HN height), $[\textup{K3$\times$E-atlas}]$ (finite
d-critical/cosection Hall atlas), and $[\textup{ML}]$
(Mittag--Leffler descent for the HN inverse system).  Two further
standing functorial inputs are $[\textup{$E_3$-HolFA}]$, a chosen
finite holomorphic \(E_3\)-prefactorization model for
\(\mathcal A_{X,R}^{E_3}\), including the BV quantum master equation,
anomaly cancellation, locality, and framing/formality data needed to
make the \(E_3\)-operations act on the retained finite stage, and
$[\textup{$K3{\to}E$-sp}]$, a chosen chain-level specialization from
the \(K3\)-direction to the hybrid chiral carrier on \(E\).
At finite stage these two inputs are the table system
\[
\mathrm{formal\_target\_charts},\quad
\mathrm{field\_complexes},\quad
\mathrm{e3\_operations},\quad
\mathrm{bv\_qme},\quad
\mathrm{anomaly\_framing},\quad
\mathrm{compact\_support},\quad
\mathrm{factorization\_descent},\quad
\mathrm{k3\_to\_e\_specialization},\quad
\mathrm{transitions},\quad
\mathrm{scalar\_firewall}.
\]
The scalar firewall excludes the automorphic section, the BKM
denominator, the target current envelope, the protected trace, and the
hybrid carrier as replacements for the compact \(E_3\)-source.
The finite \(E_3\)-source packet
\(\texttt{certificates/e3\_source/k3e\_e3\_holfa}\) currently supplies
only the obstruction ledger
\(\texttt{E3\_SOURCE\_OBSTRUCTION\_LEDGER\_VERIFIED}\), with
\(\texttt{e3\_source\_certification=false}\).  Thus
\([\textup{$E_3$-HolFA}]\) and \([\textup{K3$\to$E-sp}]\) are retained
inputs until the formal-target, field-complex, \(E_3\)-operation,
BV-QME, anomaly, compact-support, descent, specialization, transition,
and scalar-firewall rows are supplied.
The finite charge-window hypotheses are row-level data:
\[
\mathrm{active\_support},\quad
\mathrm{window\_saturation},\quad
\mathrm{liu\_hn\_window},\quad
\mathrm{gram\_map},\quad
\mathrm{cocycle\_identities},\quad
\mathrm{normal\_ordered\_lifts},\quad
\mathrm{target\_windows},\quad
\mathrm{primitive\_lifts},\quad
\mathrm{transitions},\quad
\mathrm{scalar\_firewall}.
\]
The current finite charge-window packet is empty-blocked and supplies
none of these rows.  The formal Mukai--Gram cocycle packet
\(\texttt{certificates/charge/mukai\_gram\_cocycle}\) supplies only the
chart-level identities for \(\Pi_X\), \(B\), the normal-ordered section,
and \(\overline\Pi_X\).  It does not supply finite active support,
Liu bounds, target-window images, primitive formal lifts, strict
transitions, Mittag--Leffler rows, or the scalar firewall required for
the finite charge window.
The formal primitive-lift packet
\(\texttt{certificates/charge/formal\_primitive\_lift}\) supplies only
pointwise lattice lifts and the root grading
\(\alpha(\Pi_X(Q,P))=2n f_2-l f_3+2m f_{-2}\) on the supplied formal
rows.  It is not a replacement for finite charge-window rows,
effectivity, compact support, source Hall brackets, Pfaffian
orientations, or type-II wall atlases.
The formal Gram-fibre packet
\(\texttt{certificates/charge/gram\_fibre\_kernel}\) supplies only
same-Gram fibre grouping, separate even/odd summation, the definition
of \(K_{\Pi,0}=\Pi_X^{-1}(0,0,0)\), and the explicit fact that
\(K_{\Pi,0}\) is not an additive subgroup.  It is not a source parity
theorem, a Hopf radical, or an exactness theorem for a finite
pushforward.
The finite exactness packet
\(\texttt{certificates/charge/fibre\_pushforward\_exactness}\) supplies
only finite vector-space exactness of the fibre-summed Gram
pushforward and a zero-fibre radical as a pairing kernel.  It is not
the exact Hall pushforward, Hopf radical ideal/coideal closure, or
bar-commutation datum required by the compact source theorem.
The finite bar-pushforward packet
\(\texttt{certificates/charge/bar\_pushforward\_commutation}\) supplies
only tensor-coalgebra compatibility for the additive normal-ordered
Gram grading.  It is not the raw quadratic \(\Pi_X\)-pushforward, the
chiral Koszul quasi-isomorphism, Hall bracket compatibility, or Hopf
pairing compatibility.
The finite bracket-pairing pushforward packet
\(\texttt{certificates/charge/hall\_pairing\_pushforward\_compatibility}\)
supplies only degree compatibility for a supplied bracket and supplied
homogeneous pairing after normal-ordered Gram pushforward.  It is not
compact Hall geometry, primitive closure, Hopf adjointness, Frobenius
cyclicity, quotient non-degeneracy, Pfaffian orientation, or protected
trace data.
The finite parity-pushforward packet
\(\texttt{certificates/charge/parity\_pushforward}\) supplies only the
super-vector-space identity that the additive normal-ordered Gram
pushforward sums even and odd finite blocks separately.  Its signed
superdimension rows are scalar shadows of the two parity ranks.  It is
not a source parity theorem, the Gritsenko--Nikulin/Kac parity
equality, compact source data, primitive recognition, Pfaffian
orientation, or protected trace data.
The finite PBW-pushforward packet
\(\texttt{certificates/charge/pbw\_pushforward}\) supplies only the
filtered-vector-space identity that the additive normal-ordered Gram
pushforward preserves a supplied finite PBW filtration and its
associated graded.  It is not the compact-source PBW comparison, target
PBW comparison, no-extra-relations theorem, primitive recognition,
Pfaffian orientation, or protected trace data.
The finite triangular-pushforward packet
\(\texttt{certificates/charge/triangular\_pushforward}\) supplies only
the direct-sum identity that a supplied positive/Cartan/negative
decomposition, checked by an integer height on Gram degrees, remains
triangular after additive normal-ordered Gram pushforward.  It is not
exact triangular completion, target triangular completion,
compact-source primitive recognition, no-extra-relations, Pfaffian
orientation, or protected trace data.
The retained stability-heart packet
\(\texttt{certificates/moduli/stability\_heart}\) supplies only the
chosen Abramovich--Polishchuk/Liu heart
\(\mathcal A_X^{\mathrm{AP/Liu}}\subset D^b(\mathrm{Coh}(S\times E))\)
and its input data.  It verifies
\texttt{STABILITY\_HEART\_DATUM\_VERIFIED}, but it is not
noetherianity, HN existence, bounded HN type data, finite-type
semistable moduli, derived enhancements, Pfaffian orientation, or
protected trace data.
The Liu product-stability packet
\(\texttt{certificates/moduli/liu\_product\_stability}\) supplies the
precise source-imported definition of \(\sigma_{s,t}\): the AP heart,
the linear polynomial \(L_M(n)\), the coefficient homomorphisms
\(a,b,c,d\), the weak charge \(Z_t\), the tilt by
\((\mathcal T_t,\mathcal F_t)\), and the central charge
\(Z_E^{s,t}\).  It verifies
\texttt{LIU\_PRODUCT\_STABILITY\_DEFINED}, but it is not the
AP-compatibility theorem, finite-type semistable moduli, derived
enhancements, Pfaffian orientation, or protected trace data.
The AP--Liu compatibility packet
\(\texttt{certificates/moduli/ap\_liu\_compatibility}\) supplies
Proposition~\ref{prop:ap-liu-heart-compatibility}: the retained heart
\(\mathcal A_X^{\mathrm{AP/Liu}}\) is the Liu tilt
\(\mathcal A_E^t=\langle\mathcal T_t,\mathcal F_t[1]\rangle\) of the
Abramovich--Polishchuk heart \(\mathcal A_E^{\mathrm{AP}}\).  It
verifies \texttt{AP\_LIU\_COMPATIBILITY\_VERIFIED}, but it is not
noetherianity, HN existence, bounded HN type data, finite-type
semistable moduli, derived enhancements, Pfaffian orientation, or
protected trace data.
The retained-window noetherian packet
\(\texttt{certificates/moduli/retained\_window\_noetherian}\) supplies
only the finite ACC theorem for a supplied subobject-closed and
quotient-closed retained exact window with a strict length rank.  It
verifies \texttt{RETAINED\_WINDOW\_NOETHERIAN\_VERIFIED}, but it is not
global noetherianity of the AP/Liu heart, HN existence, bounded HN type
data, finite-type moduli, derived enhancements, Pfaffian orientation,
or protected trace data.
The retained-window HN-filtration packet
\(\texttt{certificates/moduli/retained\_window\_hn\_filtration}\)
supplies only finite exact HN filtrations in the retained window: the
quotient factors are semistable, the phases strictly decrease, and
lengths add across the exact steps.  It verifies
\texttt{RETAINED\_WINDOW\_HN\_FILTRATION\_VERIFIED}, but it is not a
global HN theorem for the AP/Liu heart, bounded HN type data,
finite-type moduli, derived enhancements, Pfaffian orientation, or
protected trace data.
The retained HN type-bound packet
\(\texttt{certificates/moduli/retained\_hn\_type\_bounds}\) supplies
only bounded HN type data in the active finite window.  It verifies
\texttt{RETAINED\_HN\_TYPE\_BOUNDS\_VERIFIED}: the semistable factor
alphabet is finite, HN words have bounded length, and the
finite-moduli HN type rows are mirrored from the certified packet.  It
is not finite-type semistable moduli, derived enhancements, universal
complexes, Pfaffian orientation, or protected trace data.
The retained semistable-substack packet
\(\texttt{certificates/moduli/retained\_semistable\_substacks}\)
supplies Proposition~\ref{prop:retained-finite-type-semistable-substacks}:
each retained semistable HN factor type has a specialization-closed
finite-type substack inside a bounded Quot/Postnikov ambient.  It
verifies \texttt{RETAINED\_SEMISTABLE\_SUBSTACKS\_VERIFIED}, but it
is not a derived enhancement, universal complex, rigidification,
finite-inertia stratification, flag stack, cosection atlas,
transition system, Pfaffian orientation, or protected trace datum.
The retained derived-enhancement packet
\(\texttt{certificates/moduli/retained\_derived\_enhancements}\)
supplies Proposition~\ref{prop:retained-quasi-smooth-derived-enhancements}:
each retained semistable substack has a quasi-smooth derived Artin
enhancement with cotangent amplitude \([-1,0]\) and the restricted
PTVV \((-1)\)-shifted symplectic form.  It verifies
\texttt{RETAINED\_DERIVED\_ENHANCEMENTS\_VERIFIED}, but it is not a
universal complex, rigidification, finite-inertia stratification, flag
stack, cosection atlas, transition system, Pfaffian orientation, or
protected trace datum.
The retained rigidification-inertia packet
\(\texttt{certificates/moduli/retained\_rigidification\_inertia}\)
supplies
Proposition~\ref{prop:retained-finite-residual-inertia-after-rigidification}:
the scalar \(\mathbb G_m\) is removed from every retained derived
substack, and the residual inertia group is finite with zero
rigidification and inertia defects.  It verifies
\texttt{RETAINED\_RIGIDIFICATION\_INERTIA\_VERIFIED}, but it is not a
universal complex, finite closed inertia stratification, flag stack,
cosection atlas, transition system, Pfaffian orientation, or protected
trace datum.
The retained class-bound packet
\(\texttt{certificates/moduli/retained\_class\_bounds}\) supplies
Proposition~\ref{prop:retained-finite-class-bounds}: the retained
class set \(C_R\), the Hilbert-polynomial labels \(P_{c,i}\), the
cohomological-amplitude intervals \([a_c,b_c]\), and the regularity
bounds \(N_c\) are finite row data.  It verifies
\texttt{RETAINED\_CLASS\_BOUNDS\_VERIFIED}, but it is not a universal
complex, finite closed inertia stratification, flag stack, cosection
atlas, transition system, Pfaffian orientation, or protected trace
datum.
The retained universal-complex packet
\(\texttt{certificates/moduli/retained\_universal\_complexes}\)
supplies Proposition~\ref{prop:retained-universal-perfect-complexes}:
each retained scalar-rigidified finite substack carries a universal
perfect complex with base-change compatibility, finite Tor amplitude,
zero descent defect, and zero perfection defect.  It verifies
\texttt{RETAINED\_UNIVERSAL\_COMPLEXES\_VERIFIED}, but it is not a
finite closed inertia stratification, flag stack, cosection atlas,
transition system, Pfaffian orientation, or protected trace datum.
The retained \(E\)-translation-rigidification packet
\(\texttt{certificates/moduli/retained\_e\_translation\_rigidifications}\)
supplies
Proposition~\ref{prop:retained-e-translation-rigidifications}: the
elliptic-translation action is removed on every retained finite row
with zero translation-rigidification defect.  It verifies
\texttt{RETAINED\_E\_TRANSLATION\_RIGIDIFICATIONS\_VERIFIED}, but it
is not a finite closed inertia stratification, flag stack, cosection
atlas, transition system, Pfaffian orientation, or protected trace
datum.
The retained closed-substack packet
\(\texttt{certificates/moduli/retained\_closed\_substacks}\) supplies
Proposition~\ref{prop:retained-closed-substacks}: every retained
finite row has a finite closed cover with finite residual inertia and
zero coverage and inertia defects.  It verifies
\texttt{RETAINED\_CLOSED\_SUBSTACKS\_VERIFIED}, but it is not closure
under extensions, closure under HN factors, closure under duals, a flag
stack, a cosection atlas, a transition system, a Pfaffian orientation,
or a protected trace datum.
The retained extension-closure packet
\(\texttt{certificates/moduli/retained\_extension\_closure}\) supplies
Proposition~\ref{prop:retained-extension-closure}: every supplied
binary extension row has middle term in the retained finite HN word set
with zero extension-closure defect.  It verifies
\texttt{RETAINED\_EXTENSION\_CLOSURE\_VERIFIED}, but it is not an
extension stack, a two-step flag stack, closure under HN factors,
closure under duals, a cosection atlas, a transition system, a
Pfaffian orientation, or a protected trace datum.
The retained HN-factor-closure packet
\(\texttt{certificates/moduli/retained\_hn\_factor\_closure}\) supplies
Proposition~\ref{prop:retained-hn-factor-closure}: every HN factor of
a retained finite HN word is represented by a retained class, a
finite-type semistable substack, and a retained closed substack.  It
verifies \texttt{RETAINED\_HN\_FACTOR\_CLOSURE\_VERIFIED}, but it is
not a dual-closure datum, not an extension stack, not a two-step flag
stack, not subquotient closure in flag stacks, not a cosection atlas,
not a transition system, not a Pfaffian orientation, and not a
protected trace datum.
The retained dual-closure packet
\(\texttt{certificates/moduli/retained\_dual\_closure}\) supplies
Proposition~\ref{prop:retained-dual-closure}: the finite duality
involution \(s_3\leftrightarrow s_1\), \(s_2\mapsto s_2\), together
with reversal of HN words, preserves the retained finite HN word set
with zero defect.  It verifies
\texttt{RETAINED\_DUAL\_CLOSURE\_VERIFIED}, but it is not an extension
stack, not a two-step flag stack, not subquotient closure in flag
stacks, not a cosection atlas, not a transition system, not a Pfaffian
orientation, and not a protected trace datum.
The finite-moduli and atlas hypotheses are likewise row-level data:
\[
\mathrm{hn\_type\_bounds},\quad
\mathrm{semistable\_substacks},\quad
\mathrm{derived\_enhancements},\quad
\mathrm{universal\_complexes},\quad
\mathrm{scalar\_rigidifications},\quad
\mathrm{stratifications},\quad
\mathrm{extension\_flag\_stacks},\quad
\mathrm{cosection\_atlas},\quad
\mathrm{transitions},\quad
\mathrm{scalar\_firewall}.
\]
These rows record the retained bounded Quot/Postnikov slices,
quasi-smooth derived enhancements, universal complexes, scalar
rigidifications, finite residual inertia, retained closed
stratifications, retained extension closure, retained HN-factor
closure, retained dual closure, compactified flag correspondences,
d-critical/cosection charts, and strict transitions.
Liu stability without the finite charge-window rows, a target charge
window, or a Hilbert-scheme scalar test does not supply
\([K3\times E\text{-fin}]\) or
\([K3\times E\text{-atlas}]\).
The finite-moduli packet
\(\texttt{certificates/moduli/k3e\_finite\_moduli}\) currently records
the HN-bound, semistable-substack, derived-enhancement, and
rigidification-inertia rows, together with the finite class-bound
rows, universal-complex rows, \(E\)-translation-rigidification rows,
retained closed-substack rows, retained extension-closure rows,
retained HN-factor-closure rows, and retained dual-closure rows.  The
remaining rows stay in the obstruction ledger
\(\texttt{MODULI\_OBSTRUCTION\_LEDGER\_VERIFIED}\), with
\(\texttt{moduli\_certification=false}\).  Thus the full finite-type
moduli and d-critical/cosection atlas inputs remain open until the
flag-stack, cosection-atlas, transition, and scalar-firewall rows are
supplied.

The bar coalgebra is MC twisting, not the level-\(\mathsf Z\) bulk.
The imported Borcherds--Kac--Moody current envelope
$\mathsf{Den}_{\Delta_5, E}$ is the target chiral shadow, and the
cobar comparison $\Theta_{\mathrm{Kos}, R}$ is the source-to-target
map on the hybrid source.  Its invertibility is the Koszul source
datum, not a target-internal counit.

\section*{Construction theorem for the Dirac--Igusa package}

\begin{theorem}[Construction of the Dirac--Igusa pro-object on $K3\times E$]
\label{thm:ch5-construction-of-dirac-igusa-pro-object}
Granted the first five components of the retained obstruction datum
\(\mathfrak R_X^{\mathrm{DI}}\) of
Definition~\ref{def:G-retained-obstruction-datum}, as used in
Theorem~\ref{thm:G-four-obstruction-criterion}, and the standing
finite-stage hypotheses
$[\textup{K3$\times$E-fin}]$, $[\textup{K3$\times$E-atlas}]$, $[\textup{ML}]$,
the standing hypotheses $[\textup{$E_3$-HolFA}]$ and
$[\textup{$K3{\to}E$-sp}]$ through the finite compact-source proof
tables above, the chiral Koszul source datum of
Definition~\ref{def:chiral-koszul-source-datum}, its
finite Koszul comparison datum of
Definition~\ref{def:finite-koszul-comparison-datum}, its strict
Mittag--Leffler exactness of
Proposition~\ref{prop:chiral-koszul-source-datum-consequence}, and
the cofinal primitive-recognition datum of
Definition~\ref{def:cofinal-primitive-recognition-datum}, together
with the finite-stage morphism proof tables of
\S\ref{subsec:eight-finite-constructions} and the finite
theorem-dependency ledger
\[
\mathrm{beilinson\_levels},\quad
\mathrm{objects},\quad
\mathrm{morphisms},\quad
\mathrm{theorem\_dependencies},\quad
\mathrm{dependency\_edges},\quad
\mathrm{forbidden\_edges},\quad
\mathrm{status\_separation},\quad
\mathrm{level\_equalities},\quad
\mathrm{equality\_sites},\quad
\mathrm{definition\_separation},\quad
\mathrm{obstruction\_independence},\quad
\mathrm{transitions},\quad
\mathrm{scalar\_firewall},
\]
verified only as the schema-complete theorem-dependency ledger and not
as a certificate for the retained construction packets,
the finite stages define a Dirac--Igusa pro-object
\[
\mathfrak D_X^{\mathrm{DI}}
=\varprojlim_{R\in\mathcal R_{\mathrm{HN}}}\!\bigl(
\mathcal A_{X,R}^{E_3},
F_{X,R}^{\mathrm{hyb}},
\Gamma_{X,R},
\Pi_X,
\Phi_R,
o_R,
H_R,
P_R^\Pi,
C_{X,R},
\Theta_{\mathrm{Kos},R},
\mathscr L_{\mathrm{Pf},R},
\operatorname{pf}_{X,R},
\varepsilon_{o,R},
\mathrm{Rec}_R\bigr)
\]
on the moduli stack $\mathfrak M^{\mathrm{red},+}_c(X)$ of effective
dyonic-class objects in the chosen stability heart of
Chapter~\ref{ch:k3xe-setup}.  More precisely, the finite stages form a
functor
\[
\mathcal R_{\mathrm{HN}}^{\mathrm{op}}\longrightarrow
\mathsf{DI}^{\mathrm{fin}}_X
\]
in the finite-stage category of
Definition~\ref{def:finite-stage-dirac-igusa-morphism}, and the
displayed inverse limit is notation for the associated object of
\(\operatorname{Pro}_{\mathcal R_{\mathrm{HN}}}(\mathsf{DI}^{\mathrm{fin}}_X)\).
Componentwise inverse limits are asserted only for components whose
strict Mittag--Leffler and \(R^1\varprojlim=0\) rows are supplied.  The
universal property and cofinal presentation invariance are those of
Proposition~\ref{prop:dirac-igusa-pro-object-universal-property}.  The
dependence on the admissible parameter datum
\((S,H,\sigma,E,0_E)\), base-change functoriality in \(S\) and \(E\),
and the elliptic-origin character correction are those of
Definitions~\ref{def:admissible-dirac-igusa-parameter-datum},
\ref{def:polarization-origin-dependence} and
Propositions~\ref{prop:admissible-base-change-functoriality} and
\ref{prop:elliptic-origin-independence-character}.  The
typed finite stages are:
\begin{enumerate}
\item[\textnormal{(L1)}] (\emph{Hybrid carrier},
\S\ref{subsec:hybrid-ran-carrier}.)  $F_{X,R}^{\mathrm{hyb}}$ is the
hybrid wrapped factorization carrier on
$\mathrm{Ran}^{\mathrm{hyb}}(E)$ over the elliptic factor, with
local/wrapped/mixed atlas matching the type-II wall structure
$\{\gamma(\delta_1),\gamma(\delta_2),\gamma(\delta_3)\}=\{(1,1,0),(0,1,1),(0,-1,0)\}$
on $\Lambda^{2,1}_{II}$, as computed in
Proposition~\ref{prop:type-ii-wall-images}.  Compact Mukai-chart
representatives for these three triples are constructed in
Proposition~\ref{prop:type-ii-compact-mukai-representatives}; retained
wall objects, wall charts, and protected integration remain O2 and
hybrid inputs.
\item[\textnormal{(L2)}] (\emph{Reduced d-critical orientation},
\S\ref{subsec:cosection-twisted-orientation}.)
$o_R$ is the reduced orientation datum: a square root of the reduced
determinant line, Thom--Sebastiani transport on retained extension
strata, null-trivialisations of the reduced and equivariant
degree-two orientation gerbes, and zero finite-stabilizer
linearization characters, as in
Definition~\ref{def:O1-strong-reduced-orientation}.
\item[\textnormal{(L3)}] (\emph{Hall correspondences},
\S\ref{subsec:hall-correspondences}.)
$H_R$ is the finite protected Hall datum on
$H^\bullet_c(\mathfrak M^{\mathrm{red},+}_{c,R}(X),
\Phi_W\otimes o_R)$: product, collision coproduct, primitive
projection, and Hopf pairing on the cosection-reduced heart.  The
comparison with the Drinfeld double \(G(X)=D(Y^+(X))\) requires the
Hall pairing, integral form, stable-envelope transport, completion, and
primitive-recognition data.
\item[\textnormal{(L4)}] (\emph{First-order Dirac block and Pfaffian line},
\S\ref{subsec:dirac-pfaffian}.)
$\mathscr L_{\mathrm{Pf},R}$ is the Pfaffian line bundle on
$\mathfrak Q_R=(\mathfrak M^{\mathrm{red},+}_{c,R}(X)/E)^{\mathrm{red}}$
with finite Pfaffian sections $\operatorname{pf}_{X,R}$ carried
into the inverse limit by Mittag--Leffler descent.
\item[\textnormal{(L5)}] (\emph{Chiral Koszul source coalgebra},
\S\ref{subsec:koszul-source}.)
$C_{X,R}$ is the bar of the augmented binary/two-step hybrid
correspondence algebra; it becomes the bar of a hybrid factorization
algebra only after the higher-coloured residual vanishes.  The chiral
Koszul cobar map
$\Theta_{\mathrm{Kos},R}:\Omega_E^{\mathrm{ch}}C_X\longrightarrow
\mathsf{Den}_{\Delta_5,E}$ is the source-to-target comparison and is a
quasi-isomorphism under the chiral Koszul source datum and its
Mittag--Leffler exactness.  Its compatibility with Weyl transport,
Pfaffian orientation, Hall pairing, radical quotient, and PBW
filtration is the finite Koszul comparison datum of
Definition~\ref{def:finite-koszul-comparison-datum}.
\item[\textnormal{(L6)}] (\emph{Primitive recognition},
\S\ref{sec:ch6-primitive-recognition}.)
$\mathrm{Rec}_R$ is the cofinal-window primitive recognition
criterion (S1)--(S10) of
Definition~\ref{def:cofinal-primitive-recognition-datum}.  Under that
finite compact-source datum, it identifies
\(P_X^{\Pi,\mathrm{red}}\) with \(\mathfrak g_{\Delta_5}\).
\item[\textnormal{(L7)}] (\emph{Eight finite constructions},
\S\ref{subsec:eight-finite-constructions}.)
The pro-object is the inverse limit of the eight finite constructions
of Appendix~\ref{app:eight-finite-constructions}; its transition maps
are the morphisms \(\rho_{RR'}\) of
Definition~\ref{def:finite-stage-dirac-igusa-morphism}, and each
component carries its obstruction class.
\item[\textnormal{(L8)}] (\emph{Obstruction classes},
\S\ref{subsec:ch6-obstructions}.)
The obstruction $(D0)$ is reduced on \(K3\times E\) to the retained
D0-HN datum, $(O1)$ and $(O1)^+$ are reduced to retained orientation data, and $(O2)$ is reduced to
\((O2_{\mathrm{atlas}})\), by
Theorem~\ref{thm:G-four-obstruction-criterion}.
\end{enumerate}
\end{theorem}

\begin{proof}
The pro-object is represented by the functor of finite-stage data on
the HN system $\mathcal R_{\mathrm{HN}}$.  The Pfaffian, orientation,
Hall, and carrier entries are obtained from the finite constructions
of Appendix~\ref{app:eight-finite-constructions}, and the functorial
transition system is supplied by the fixed finite-stage data
$[\textup{K3$\times$E-fin}]$, $[\textup{K3$\times$E-atlas}]$,
$[\textup{ML}]$ together with the projections
\((D0_{\mathrm{HN}},O1_{\mathrm{quot}},\mathcal T_W,O2_{\mathrm{atlas}})\)
of \(\mathfrak R_X^{\mathrm{DI}}\).  The Koszul entry
\(\Theta_{\mathrm{Kos},R}\) is supplied by the chiral Koszul source
datum and its exact inverse-system condition.  The recognition entry
\(\mathrm{Rec}_R\) is supplied by the finite compact-source
recognition datum.  Thus $(\mathrm{L1})$--$(\mathrm{L8})$ are finite
components with separate source, orientation, Hall, Koszul, Pfaffian,
and recognition hypotheses, and their transition morphisms satisfy
\textup{(M1)}--\textup{(M8)} by the finite morphism proof tables;
no component is inferred from another.
\end{proof}

\medskip

Entries $(\mathrm L 1)$--$(\mathrm L 8)$ are the finite constituents of
\(\mathfrak D_X^{\mathrm{DI}}\).  Each constituent has its
finite-stage morphism and obstruction class.  A section equality,
denominator product, or scalar trace does not determine an isomorphism
of Dirac--Igusa pro-objects unless the morphisms of
Definition~\ref{def:finite-stage-dirac-igusa-morphism} and the
pro-isomorphism of Definition~\ref{def:dirac-igusa-pro-isomorphism} are
supplied with cofinality, component-map, composition-law,
Mittag--Leffler, pro-isomorphism, and scalar-firewall rows.  The finite
formulas are in
Appendix~\ref{app:eight-finite-constructions}; the retained obstruction
data are in Appendix~\ref{app:obstruction-discharges}.
The current finite morphism packet
\(\texttt{certificates/morphisms/k3e\_dirac\_igusa\_pro\_object}\)
verifies only the obstruction ledger
\(\texttt{MORPHISM\_OBSTRUCTION\_LEDGER\_VERIFIED}\), with
\(\texttt{pro\_object\_certification=false}\).  Thus the present
checked data record the missing pro-object morphism rows; they do not
construct the transition functor or a pro-isomorphism.

\subsection{The hybrid factorization carrier $\operatorname{Ran}^{\mathrm{hyb}}(E)$}
\label{subsec:hybrid-ran-carrier}

The first entry of the Dirac--Igusa pro-object is the carrier on which
chiral operations live.  An ordinary Beilinson--Drinfeld factorization
algebra on $\operatorname{Ran}(E)$ \cite[\S3.1]{BD} sees only a finite
multiset of points of $E$ and the operations attached to their
collisions; this is the projection-finite stratum.  The Igusa product
\(\Delta_5\) carries a Borcherds variable $s$ whose positive powers
record positive elliptic-degree wrapping.  An effective one-cycle on
\(K3\times E\) of positive elliptic degree has a horizontal component
and therefore projects \emph{onto} \(E\).  Locality on
$\operatorname{Ran}(E)$ alone has no mechanism for it.

\begin{proposition}[Positive elliptic degree projects onto \(E\);
\textit{proved}]
\label{prop:positive-elliptic-degree-projects}
\label{lem:wrapped-projection-obstruction}
Let
\[
Z=\sum_i m_i[C_i]\in Z_1(K3\times E),\qquad m_i>0,
\]
be an effective one-cycle with irreducible curve components \(C_i\).
If
\[
(p_E)_*Z=b[E],\qquad b>0,
\]
then \(p_E(|Z|)=E\).  Equivalently, at least one component \(C_i\)
maps dominantly to \(E\).  If a retained object \(A\) has a support
realization row identifying \(\operatorname{cyc}_1(A)=Z\) with the
effective fundamental one-cycle of \(\operatorname{Supp}_1 A\), then
\[
b_R^{\mathrm{geom}}(A)>0\quad\Longrightarrow\quad
p_E(\operatorname{Supp}_1 A)=E .
\]
Thus such an object is not localized at finitely many points of \(E\).
\end{proposition}

\begin{proof}
For each irreducible component \(C_i\), the proper image \(p_E(C_i)\)
is either a point or all of the irreducible curve \(E\).  In the first
case \((p_E)_*[C_i]=0\); in the second case
\((p_E)_*[C_i]=\deg(C_i/E)[E]\) with positive degree.  Since
\[
(p_E)_*Z=\sum_i m_i(p_E)_*[C_i]=b[E]
\]
and \(b>0\), at least one component has positive degree over \(E\).
That component maps dominantly onto \(E\), hence
\[
E=p_E(C_i)\subset p_E(|Z|)\subset E.
\]
The support statement follows when the support-realization row
identifies \(Z\) with the effective fundamental cycle of the
one-dimensional support of \(A\).
\end{proof}

The packet
\[
\texttt{certificates/hybrid/positive\_elliptic\_degree\_projection}
\]
verifies
\(\texttt{POSITIVE\_ELLIPTIC\_DEGREE\_PROJECTION\_VERIFIED}\):
the effective-cycle theorem is proved, while the retained object row,
effective support-cycle realization row, positive degree row, object
projection row, and transition row are recorded as missing open
obligations.

The hybrid carrier replaces the Ran space by a two-stratum object that
records, simultaneously, finite local supports and wrapped elliptic
legs.  At every finite Harder--Narasimhan height $R$ the datum is a
finite-coloured correspondence base, not a Ran space with a scalar
attached.

\begin{definition}[Finite-stage hybrid Ran prestack]
\label{def:hybrid-ran-prestack}
\label{def:hybrid-ran-base}
Fix a finite HN height \(R\).  The hybrid carrier at height \(R\) is
first a prestack
\[
\operatorname{Ran}^{\mathrm{hyb}}_{R,\mathrm{pre}}(E):
\mathrm{Aff}_{\C}^{op}\longrightarrow\infty\textup{-}\mathrm{Gpd}.
\]
For a test affine \(T\), a \(T\)-point is represented by the following
finite data:
\begin{enumerate}
\item a finite local index set \(I\) and maps
\(x_i:T\to E\), \(i\in I\);
\item retained local colour labels
\(\alpha_i\in\Gamma_R^{\mathrm{loc}}\);
\item a finite wrapped index set \(J\), retained wrapped colour labels
\(\eta_j\in\Gamma_R^{\mathrm{wr}}\), and \(T\)-families
\[
W_j\in\mathcal M_{\eta_j,R}^{\mathrm{wr,rig}}(T)
\]
in the rigidified wrapped prequotient stacks, equipped with their
anchor maps
\[
\lambda_{\eta_j,R}(W_j):T\longrightarrow E .
\]
\end{enumerate}
Two representatives are identified by finite relabelling of \(I\) and
\(J\), and by the usual Ran collision maps on the local coordinates.
Equivalently,
\[
\operatorname{Ran}^{\mathrm{hyb}}_{R,\mathrm{pre}}(E)
=
\operatorname*{colim}_{(I,J,\alpha,\eta)\in\mathsf{Fin}^{\mathrm{hyb}}_R}
\left(
E^I\times\prod_{j\in J}
\mathcal M_{\eta_j,R}^{\mathrm{wr,rig}}
\right),
\]
where \(\mathsf{Fin}^{\mathrm{hyb}}_R\) is the finite indexing category
generated by relabellings and local collision maps, and where
\(\alpha:I\to\Gamma_R^{\mathrm{loc}}\),
\(\eta:J\to\Gamma_R^{\mathrm{wr}}\).  This is a prestack
definition.  For a topology \(\tau\), the stack
\[
\operatorname{Ran}^{\mathrm{hyb}}_{R,\tau}(E)
:=a_\tau\operatorname{Ran}^{\mathrm{hyb}}_{R,\mathrm{pre}}(E)
\]
is used only after the descent rows for \(a_\tau\) have been supplied.
The scalar object
\[
\bigsqcup_{b\ge0}\operatorname{Ran}(E)\times\{b\}
\]
is only the degree shadow; it is not the hybrid prestack.
\end{definition}

The packet
\[
\texttt{certificates/hybrid/hybrid\_ran\_prestack\_definition}
\]
verifies
\(\texttt{HYBRID\_RAN\_PRESTACK\_DEFINITION\_VERIFIED}\): the
finite-stage carrier is defined as a prestack functor, while the
local-stratum component rows, wrapped-stratum component rows, descent
rows, stackification rows, and transition rows are recorded as missing
open obligations.

\begin{definition}[Local \(b=0\) stratum as a subprestack]
\label{def:hybrid-local-stratum-prestack}
At finite HN height \(R\), the local stratum is the full subprestack
\[
\operatorname{Ran}^{\mathrm{loc}}_{R,\mathrm{pre}}(E)
\subset
\operatorname{Ran}^{\mathrm{hyb}}_{R,\mathrm{pre}}(E)
\]
cut out by the conditions \(J=\varnothing\) and
\(\alpha:I\to\Gamma_R^{\mathrm{loc}}\), where
\[
\Gamma_R^{\mathrm{loc}}=\{\gamma\in\Gamma_{\sigma,S}^{\mathrm{HN}}(R)
\mid b_R^{\mathrm{geom}}(\gamma)=0\}.
\]
Thus, for a test affine \(T\),
\[
\operatorname{Ran}^{\mathrm{loc}}_{R,\mathrm{pre}}(E)(T)
=
\operatorname*{colim}_{(I,\alpha)}
\operatorname{Map}(T,E)^I,
\qquad
\alpha:I\to\Gamma_R^{\mathrm{loc}},
\]
where the colimit is over finite sets, relabellings, and collision
maps.  The forgetful map to the ordinary Ran prestack discards the
colour map \(\alpha\).  This definition contains no wrapped family and
no anchor map.  The local closed-configuration incidence stacks
\(\mathfrak M_{\alpha,R,I}^{\mathrm{loc,cl}}\) and the sheaves
\(\mathcal F_{\alpha,R}^{\mathrm{loc}}\) are coefficient data over this
carrier, not the carrier itself.
\end{definition}

The packet
\[
\texttt{certificates/hybrid/local\_stratum\_b0\_prestack\_definition}
\]
verifies
\(\texttt{LOCAL\_STRATUM\_B0\_PRESTACK\_DEFINITION\_VERIFIED}\): the
\(b=0\) local stratum is defined as a subprestack of the hybrid
prestack, while component rows, closed-configuration rows, local sheaf
rows, descent rows, and transition rows are recorded as missing open
obligations.

\begin{definition}[Wrapped \(b>0\) stratum as a subprestack]
\label{def:hybrid-wrapped-stratum-prestack}
At finite HN height \(R\), the wrapped stratum is the full subprestack
\[
\operatorname{Ran}^{\mathrm{wr}}_{R,\mathrm{pre}}(E)
\subset
\operatorname{Ran}^{\mathrm{hyb}}_{R,\mathrm{pre}}(E)
\]
cut out by the conditions \(I=\varnothing\) and
\(\eta:J\to\Gamma_R^{\mathrm{wr}}\), where
\[
\Gamma_R^{\mathrm{wr}}=\{\gamma\in\Gamma_{\sigma,S}^{\mathrm{HN}}(R)
\mid b_R^{\mathrm{geom}}(\gamma)>0\}.
\]
Thus, for a test affine \(T\),
\[
\operatorname{Ran}^{\mathrm{wr}}_{R,\mathrm{pre}}(E)(T)
=
\operatorname*{colim}_{(J,\eta)}
\prod_{j\in J}
\mathcal M_{\eta_j,R}^{\mathrm{wr,rig}}(T),
\qquad
\eta:J\to\Gamma_R^{\mathrm{wr}},
\]
where the colimit is over finite sets and relabellings.  The anchor
maps
\[
\lambda_{\eta_j,R}(W_j):T\longrightarrow E
\]
are retained as part of each wrapped prestack point.  This definition
contains no finite local support maps \(x_i:T\to E\) and no local
colour labels.  The wrapped coefficient objects
\(\mathcal G_{\eta,R}^{\mathrm{wr}}\), anchor residuals, and
quotient-after-correspondence data are coefficient and correspondence
data over this carrier, not the carrier itself.
\end{definition}

The packet
\[
\texttt{certificates/hybrid/wrapped\_stratum\_bpositive\_prestack\_definition}
\]
verifies
\(\texttt{WRAPPED\_STRATUM\_BPOSITIVE\_PRESTACK\_DEFINITION\_VERIFIED}\):
the \(b>0\) wrapped stratum is defined as a subprestack of the hybrid
prestack, while component rows, populated wrapped prequotient rows,
anchor rows, support-realization rows, descent rows, and transition
rows are recorded as missing open obligations.

\begin{definition}[\(E\)-equivariant wrapped prequotient stacks]
\label{def:e-equivariant-wrapped-prequotient-stack}
Fix a finite HN height \(R\) and a wrapped colour
\(\eta\in\Gamma_R^{\mathrm{wr}}\).  The wrapped prequotient row is the
scalar-rigidified retained wrapped stack
\[
\mathcal M_{\eta,R}^{\mathrm{wr,rig}}
\subset \mathfrak M_{\eta,R}^{\mathrm{rig}}
\]
whose objects have geometric elliptic degree
\[
b_R^{\mathrm{geom}}(\eta)>0.
\]
The superscript \(\mathrm{rig}\) means that the scalar automorphisms
have been rigidified.  It does not mean that the elliptic
translation direction has been divided out.  The map
\[
\rho_{\eta,R}:
\mathcal M_{\eta,R}^{\mathrm{wr,rig}}
\longrightarrow
\mathcal M_{\eta,R}^{\mathrm{wr}}\subset\mathfrak M_{\eta,R}
\]
forgets this scalar rigidification and the finite wrapped bookkeeping.
The reduced quotient, when later used, is written
\[
\mathcal M_{\eta,R}^{\mathrm{wr},E}
:=
\bigl[\mathcal M_{\eta,R}^{\mathrm{wr,rig}}/E\bigr],
\]
and belongs to the quotient-descent rows, not to the prequotient
definition.

For a test affine \(T\) and a section \(a:T\to E\), let
\(\tau_a:E\times T\to E\times T\) be relative translation by \(a\).
An object \(A\in\mathcal M_{\eta,R}^{\mathrm{wr,rig}}(T)\), represented
by the universal perfect complex on \(K3\times E\times T\), is acted on
by
\[
a\cdot A
:=
\bigl(1_{K3}\times\tau_a\bigr)_*A .
\]
The scalar rigidification and the universal complex are transported by
this automorphism.  The origin \(0_E\) is fixed once and for all and
acts as the identity.  This defines the \(E\)-equivariant wrapped
prequotient stack.  It does not prove finite type, properness, final
anchor population, anchor losslessness, quotient descent, or transition
compatibility.
\end{definition}

\begin{lemma}[\(E\)-action law on the wrapped prequotient; \textit{proved}]
\label{lem:e-action-law-wrapped-prequotient}
The operation of
Definition~\ref{def:e-equivariant-wrapped-prequotient-stack} is an
\(E\)-action on \(\mathcal M_{\eta,R}^{\mathrm{wr,rig}}\).  It preserves
the wrapped colour \(\eta\) and the condition
\(b_R^{\mathrm{geom}}(\eta)>0\).
\end{lemma}

\begin{proof}
Relative translations satisfy
\(\tau_{0_E}=\mathrm{id}\) and
\(\tau_a\circ\tau_b=\tau_{a+b}\).  Functoriality of pushforward along
automorphisms gives
\[
0_E\cdot A=A,\qquad
a\cdot(b\cdot A)=(a+b)\cdot A .
\]
Translation on the elliptic factor acts trivially on the numerical
Mukai class and carries effective one-cycles of a fixed
\(p_E\)-degree to effective one-cycles of the same \(p_E\)-degree.
Hence the retained HN colour and
\(b_R^{\mathrm{geom}}\) are unchanged.  The freeness and zero
translation-rigidification defect on retained rows are exactly the
input of Proposition~\ref{prop:retained-e-translation-rigidifications};
they are imported here only for the \(E\)-action row and do not supply
finite type, properness, quotient descent, anchors, or transitions.
\end{proof}

The packet
\[
\texttt{certificates/hybrid/equivariant\_wrapped\_prequotient\_definition}
\]
verifies
\(\texttt{E\_EQUIVARIANT\_WRAPPED\_PREQUOTIENT\_DEFINITION\_VERIFIED}\):
the scalar-rigidified wrapped prequotient
\(\mathcal M_{\eta,R}^{\mathrm{wr,rig}}\) is defined before the
reduced \(E\)-quotient, the translation action with origin \(0_E\) is
defined, and the action law and wrapped-colour preservation have zero
defect.  Finite type, properness, final anchor population,
anchor losslessness, quotient descent, transition compatibility, and
aggregate hybrid-carrier population are recorded as missing open
obligations.

\begin{proposition}[Finite type of wrapped prequotients]
\label{prop:wrapped-prequotient-finite-type}
\textnormal{[proved at finite HN height]}
Assume the retained finite class-bound datum of
Definition~\ref{def:retained-finite-class-bound-datum}, the retained
finite-type semistable substack datum of
Definition~\ref{def:retained-finite-type-semistable-substack-datum},
and the scalar-rigidification datum of
Definition~\ref{def:retained-rigidification-inertia-datum}.  Then, for
every wrapped colour \(\eta\in\Gamma_R^{\mathrm{wr}}\), the
prequotient stack
\[
\mathcal M_{\eta,R}^{\mathrm{wr,rig}}
\]
of Definition~\ref{def:e-equivariant-wrapped-prequotient-stack} is of
finite type over \(\C\).  This is a finite-stage statement before the
reduced \(E\)-quotient is taken.
\end{proposition}

\begin{proof}
Let
\[
C_R(\eta)=\{c\in C_R\mid \widehat c=\eta\}.
\]
Proposition~\ref{prop:retained-finite-class-bounds} makes \(C_R\)
finite; hence \(C_R(\eta)\) is finite.  For each
\(c\in C_R(\eta)\),
Proposition~\ref{prop:retained-finite-type-semistable-substacks} gives
a finite-type retained semistable stack \(\mathfrak M^{ss}_{R,c}\).
The scalar subgroup \(\mathbb G_m\) is flat, finitely presented, and
central in the retained stabilizer sequence.  The rigidification
\[
\mathfrak M^{ss}_{R,c}\longrightarrow
\mathfrak M^{\mathrm{rig}}_{R,c}
=\mathfrak M^{ss}_{R,c}/\mathbb G_m
\]
is therefore an fppf gerbe by the rigidification theorem
\cite[Tag 04V2]{StacksProject}; finite type descends across such an
fppf cover by fpqc descent for finite type
\cite[Tag 041K, Lemma 74.11.11]{StacksProject}.  Hence every
\(\mathfrak M^{\mathrm{rig}}_{R,c}\) with \(c\in C_R(\eta)\) is finite
type.  The wrapped prequotient is the finite union of these
scalar-rigidified retained class pieces:
\[
\mathcal M_{\eta,R}^{\mathrm{wr,rig}}
=
\bigcup_{c\in C_R(\eta)}\mathfrak M^{\mathrm{rig}}_{R,c}.
\]
Finite unions of finite-type substacks are finite type.  The condition
\(b_R^{\mathrm{geom}}(\eta)>0\) is numerical on the retained class and
does not introduce an infinite support parameter.  The proof does not
establish properness, quotient descent, final anchor population,
anchor losslessness, transition compatibility, or aggregate
hybrid-carrier population.
\end{proof}

The packet
\[
\texttt{certificates/hybrid/wrapped\_prequotient\_finite\_type}
\]
verifies
\(\texttt{WRAPPED\_PREQUOTIENT\_FINITE\_TYPE\_VERIFIED}\): finite type
of \(\mathcal M_{\eta,R}^{\mathrm{wr,rig}}\) follows from the finite
retained class set, finite-type retained semistable substacks, and
scalar rigidification by \(\mathbb G_m\).  Properness, final anchor
population, anchor losslessness, quotient descent, transition
compatibility, and aggregate hybrid-carrier population are still
missing open obligations.

\begin{corollary}[Finite-support Ran locality misses the geometric elliptic-degree
direction; \textit{proved}]
\label{cor:finite-support-ranE-misses-s}
A support-local factorization algebra on $\operatorname{Ran}(E)$,
defined only by separating supports over disjoint opens of $E$,
sees the projection-finite sector.  It does not realize the sectors of
positive geometric elliptic degree.
\end{corollary}

\begin{proof}
The obstruction uses the geometric degree
$b_R^{\mathrm{geom}}=\deg_E$, not the Borcherds coordinate
$m=\operatorname{pr}_3\overline\Pi_X$.  On the rank-one
Oberdieck--Pandharipande/Donaldson--Thomas branch one has the
branch-local equality
$m_{\mathrm{Bch}}=d_E-1$ \cite{OberdieckReducedPT}; it is not a
definition of the wrapped stratum.  By
Proposition~\ref{prop:positive-elliptic-degree-projects}, an effective
support cycle with \(b_R^{\mathrm{geom}}>0\) projects onto all of
\(E\).  The object-level statement requires the support-realization row
named in that proposition.  Such wrapped supports cannot be separated
by disjoint opens.  Therefore ordinary Ran-locality has no locality
mechanism for the wrapped sector.
\end{proof}

\begin{definition}[Geometric elliptic-degree map]
\label{def:geometric-elliptic-degree-map}
Fix a finite HN height \(R\) and a retained HN numerical class
\(\gamma\in\Gamma_{\sigma,S}^{\mathrm{HN}}(R)\).  Let
\[
\operatorname{cyc}_1(\gamma)\in N_1(K3\times E)
\]
be the numerical one-cycle obtained from the codimension-two Chern
character component, equivalently the class of
\(\operatorname{ch}_2(A)\cap[K3\times E]\) for any retained
semistable object \(A\) of numerical class \(\gamma\).  Since
\(N_1(E)=\Z[E]\), there is a unique integer
\(b_R^{\mathrm{geom}}(\gamma)\) such that
\[
(p_E)_*\operatorname{cyc}_1(\gamma)
=b_R^{\mathrm{geom}}(\gamma)[E].
\]
On the effective retained one-dimensional sector this integer is
nonnegative.  The local and wrapped charge windows are
\[
\Gamma_R^{\mathrm{loc}}
=\{\gamma\mid b_R^{\mathrm{geom}}(\gamma)=0\},\qquad
\Gamma_R^{\mathrm{wr}}
=\{\gamma\mid b_R^{\mathrm{geom}}(\gamma)>0\}.
\]
The map \(b_R^{\mathrm{geom}}\) is not the Borcherds trace coordinate
\[
m_R^{\mathrm{Bch}}=\operatorname{pr}_3\overline\Pi_X .
\]
The latter is a normal-ordered Gram coordinate; the former is the
proper-pushforward degree of the geometric support along
\(p_E:K3\times E\to E\).
\end{definition}

The packet
\[
\texttt{certificates/hybrid/geometric\_elliptic\_degree\_map}
\]
verifies
\(\texttt{GEOMETRIC\_ELLIPTIC\_DEGREE\_MAP\_OBSTRUCTION\_VERIFIED}\):
the definition is fixed, while the retained HN domain rows, one-cycle
class rows, class-constancy rows, pushforward rows, integer degree rows,
nonnegativity rows, local/wrapped split rows, transition rows, and
retained-extension additivity rows are not supplied.  The last item is
the separate row-202 packet below.

\begin{definition}[Geometric elliptic-degree additivity datum]
\label{def:geometric-elliptic-degree-additivity-datum}
At finite HN height \(R\), a geometric elliptic-degree additivity
datum consists, for every retained extension row
\[
0\longrightarrow A_{\gamma'}\longrightarrow A_\gamma
\longrightarrow A_{\gamma''}\longrightarrow0
\]
in the finite Hall stage, of the following rows:
\begin{enumerate}
\item the retained extension-pair row, with
\(\gamma',\gamma'',\gamma\in
\Gamma_{\sigma,S}^{\mathrm{HN}}(R)\);
\item the one-cycle equality
\[
\operatorname{cyc}_1(\gamma)
=\operatorname{cyc}_1(\gamma')
+\operatorname{cyc}_1(\gamma'')
\quad\text{in }N_1(K3\times E);
\]
\item the pushed-forward equality
\[
(p_E)_*\operatorname{cyc}_1(\gamma)
=(p_E)_*\operatorname{cyc}_1(\gamma')
+(p_E)_*\operatorname{cyc}_1(\gamma'')
\quad\text{in }N_1(E);
\]
\item the integer degree equality
\[
b_R^{\mathrm{geom}}(\gamma)
=b_R^{\mathrm{geom}}(\gamma')
+b_R^{\mathrm{geom}}(\gamma'');
\]
\item compatibility of these rows with the retained transition maps
\(R\le R'\).
\end{enumerate}
The same definition applies to a distinguished triangle representing
the Hall extension correspondence, with the corresponding equality in
\(K_0(K3\times E)\).
\end{definition}

\begin{proposition}[Additivity criterion for the geometric elliptic degree]
\label{prop:geometric-elliptic-degree-additivity-criterion}
Assume a finite stage carries the datum of
Definition~\ref{def:geometric-elliptic-degree-additivity-datum}.  Then
\[
b_R^{\mathrm{geom}}(\gamma)
=b_R^{\mathrm{geom}}(\gamma')
+b_R^{\mathrm{geom}}(\gamma'')
\]
for every retained extension row
\[
0\to A_{\gamma'}\to A_\gamma\to A_{\gamma''}\to0.
\]
\end{proposition}

\begin{proof}
In \(K_0(K3\times E)\) the retained extension gives
\([A_\gamma]=[A_{\gamma'}]+[A_{\gamma''}]\).  The Chern character is
additive on \(K_0\), hence the codimension-two component gives the
one-cycle equality in \(N_1(K3\times E)\).  Proper pushforward along
\(p_E\) is linear on numerical one-cycles, so the equality remains true
in \(N_1(E)\).  Since \(N_1(E)=\Z[E]\), equality of the pushed-forward
cycles is equality of the coefficients of \([E]\), which are precisely
the integers \(b_R^{\mathrm{geom}}\) of
Definition~\ref{def:geometric-elliptic-degree-map}.  The transition
rows state that this coefficient equality is preserved under
\(R\le R'\).
\end{proof}

The packet
\[
\texttt{certificates/hybrid/geometric\_elliptic\_degree\_additivity}
\]
verifies
\(\texttt{GEOMETRIC\_ELLIPTIC\_DEGREE\_ADDITIVITY\_OBSTRUCTION\_VERIFIED}\):
the additivity criterion is fixed, the degree-map definition packet and
the retained-extension-closure packet are imported, and the missing
retained extension-pair rows, one-cycle additivity rows, pushforward
additivity rows, integer degree additivity rows, and transition rows are
recorded fail-closed.  Retained word closure by itself proves only that
the middle HN word is retained; it does not prove additivity of
\((p_E)_*\operatorname{cyc}_1\).

\begin{definition}[Geometric/Borcherds degree comparison datum]
\label{def:geometric-borcherds-degree-comparison-datum}
At finite HN height \(R\), set
\[
m_R^{\mathrm{Bch}}(\widehat\gamma)
:=\operatorname{pr}_3\overline\Pi_X(\widehat\gamma),
\qquad
\overline\Pi_X(\widehat\gamma)=(n,l,m).
\]
A comparison datum between \(b_R^{\mathrm{geom}}\) and
\(m_R^{\mathrm{Bch}}\) consists of:
\begin{enumerate}
\item a specified branch of retained HN classes;
\item a lift of each retained class on that branch to a charge
\(\widehat\gamma\) in the domain of \(\overline\Pi_X\);
\item populated geometric degree rows computing
\(b_R^{\mathrm{geom}}\) from
\((p_E)_*\operatorname{cyc}_1\);
\item populated Gram-coordinate rows computing
\(\operatorname{pr}_3\overline\Pi_X(\widehat\gamma)\);
\item a declared equality or affine comparison formula with zero defect
on the branch;
\item compatibility of the comparison with retained transitions
\(R\le R'\).
\end{enumerate}
\end{definition}

\begin{proposition}[Geometric/Borcherds degree separation]
\label{prop:geometric-borcherds-degree-separation}
The finite hybrid datum distinguishes the support degree
\[
b_R^{\mathrm{geom}}(\gamma)
=\operatorname{coeff}_{[E]}(p_E)_*\operatorname{cyc}_1(\gamma)
\]
from the Borcherds trace coordinate
\[
m_R^{\mathrm{Bch}}(\widehat\gamma)
=\operatorname{pr}_3\overline\Pi_X(\widehat\gamma).
\]
The local/wrapped split is made by \(b_R^{\mathrm{geom}}\).  The
protected trace is graded by \(m_R^{\mathrm{Bch}}\).  No formula
identifying the two coordinates, or replacing one by the other, is part
of the hybrid datum unless a comparison datum of
Definition~\ref{def:geometric-borcherds-degree-comparison-datum} is
supplied.
\end{proposition}

\begin{proof}
The two maps have different constructions.  The map
\(b_R^{\mathrm{geom}}\) is obtained from the numerical one-cycle class
of a retained object by proper pushforward to \(N_1(E)=\Z[E]\); it is
defined before the local/wrapped factorization carrier is formed.  The
coordinate \(m_R^{\mathrm{Bch}}\) is the third coordinate of the
normal-ordered Gram homomorphism \(\overline\Pi_X\); it is used after
the Gram descent to grade protected integration in the Borcherds
variables.  A branch-local identity, such as the rank-one
Oberdieck--Pandharipande comparison \(m_{\mathrm{Bch}}=d_E-1\)
\cite{OberdieckReducedPT}, becomes a theorem only after the branch,
lift, geometric degree rows,
Gram-coordinate rows, and transition compatibility are supplied.  In
the absence of these rows there is no typed morphism that turns
\(\operatorname{pr}_3\overline\Pi_X\) into
\(\operatorname{coeff}_{[E]}(p_E)_*\operatorname{cyc}_1\).
\end{proof}

The packet
\[
\texttt{certificates/hybrid/geometric\_borcherds\_degree\_separation}
\]
verifies
\(\texttt{GEOMETRIC\_BORCHERDS\_DEGREE\_SEPARATION\_OBSTRUCTION\_VERIFIED}\):
the two coordinates are separated, substitution of
\(\operatorname{pr}_3\overline\Pi_X\) for \(b_R^{\mathrm{geom}}\) is
excluded, and branch-comparison rows are recorded as missing open
obligations.

\begin{proposition}[Middle type-II wall routes through wrapped and mixed geometry]
\label{prop:middle-type-ii-wall-wrapped-mixed}
\textnormal{[routing proved; retained wall object not constructed]}
Let
\[
\gamma_{\mathrm{wall}}(\delta_2)=(0,1,1)
\]
be the middle simple type-II wall image of
Proposition~\ref{prop:type-ii-wall-images}.  On the type-II wall branch
where the geometric elliptic degree is written in the OP/Borcherds
coordinate as \(d=m+1\), this wall has \(d=2>0\).  Consequently, any
retained finite-stage class \(\eta\) with a supplied branch-comparison
row
\[
\overline\Pi_X(\eta)=(0,1,1),\qquad
b_R^{\mathrm{geom}}(\eta)=d=2,
\]
lies in \(\Gamma_R^{\mathrm{wr}}\), not in \(\Gamma_R^{\mathrm{loc}}\).
Its one-wall carrier is the wrapped prequotient side
\(\mathcal M_{\eta,R}^{\mathrm{wr,rig}}\) of
Definition~\ref{def:e-equivariant-wrapped-prequotient-stack}.  Any Hall
operation inserting projection-finite local colours against this wall
is typed by the one-sided \(LW/WL\) mixed correspondences of
Definition~\ref{def:mixed-local-wrapped-correspondence-datum}, and by
the ordered two-sided mixed correspondence of
Definition~\ref{def:two-sided-mixed-correspondence-datum} when local
insertions occur on both sides.  It is not an \(LL\) operation on the
ordinary Ran stratum.

This is a routing theorem.  It does not supply the retained wall object
\(x_{\delta_2,R}\), an effective support-realization row, a populated
wrapped anchor, quotient descent, the O2 wall atlas, or protected
integration.
\end{proposition}

\begin{proof}
Proposition~\ref{prop:type-ii-wall-images} gives the middle wall
triple \((n,l,m)=(0,1,1)\).  The wall-branch convention
\(d=m+1\) gives \(d=2\).  After the retained branch-comparison row
identifies this integer with \(b_R^{\mathrm{geom}}(\eta)\), the
definition
\[
\Gamma_R^{\mathrm{loc}}=\{b_R^{\mathrm{geom}}=0\},\qquad
\Gamma_R^{\mathrm{wr}}=\{b_R^{\mathrm{geom}}>0\}
\]
puts \(\eta\) in the wrapped window and excludes it from the local
window.  By Proposition~\ref{prop:positive-elliptic-degree-projects},
if a retained object realizing \(\eta\) is supplied together with its
effective support-cycle row, its one-dimensional support projects onto
all of \(E\), so finite-support Ran locality cannot contain it.
Definition~\ref{def:hybrid-wrapped-stratum-prestack} therefore places
the one-wall carrier in the wrapped \(b>0\) subprestack, with the
prequotient stack of
Definition~\ref{def:e-equivariant-wrapped-prequotient-stack}.  The
local/wrapped and two-sided mixed correspondence definitions are exactly
the before-\(E\)-quotient operations whose source contains one wrapped
input and projection-finite local inputs.  The final sentence records
the rows still absent from this routing calculation.
\end{proof}

The packet
\[
\texttt{certificates/hybrid/middle\_type\_ii\_wall\_wrapped\_mixed}
\]
verifies
\(\texttt{MIDDLE\_TYPE\_II\_WALL\_WRAPPED\_MIXED\_ROUTING\_VERIFIED}\):
the middle wall image has \(m=1\), the wall-branch degree \(d=m+1=2\)
is positive, the routing is to the wrapped \(b>0\) side, and local
insertions are typed by mixed correspondences.  Its obstruction rows
record that retained branch-comparison, wall-object, support,
anchor, O2 atlas, quotient, protected-integration, aggregate-carrier,
and transition rows are still missing.

\begin{proposition}[Hybrid operation profiles recover the three simple type-II walls]
\label{prop:type-ii-walls-hybrid-operation-recovery}
\textnormal{[operation-profile recovery proved; carrier population not constructed]}
Let
\[
\gamma_{\mathrm{wall}}(\delta_1)=(1,1,0),\qquad
\gamma_{\mathrm{wall}}(\delta_2)=(0,1,1),\qquad
\gamma_{\mathrm{wall}}(\delta_3)=(0,-1,0)
\]
be the three simple type-II wall images of
Proposition~\ref{prop:type-ii-wall-images}.  On the same type-II wall
branch \(d=m+1\), their elliptic degrees are
\[
d_{\delta_1}=1,\qquad d_{\delta_2}=2,\qquad d_{\delta_3}=1.
\]
Thus, after retained branch-comparison rows identify
\[
b_R^{\mathrm{geom}}(\eta_i)=d_{\delta_i},\qquad
\overline\Pi_X(\eta_i)=\gamma_{\mathrm{wall}}(\delta_i),
\]
the three simple type-II walls are recovered by the hybrid vocabulary
as the three singleton wrapped colour profiles
\[
\eta_i\in\Gamma_R^{\mathrm{wr}},\qquad i=1,2,3.
\]
For each \(i\), a finite local insertion on the left or right of this
wall is typed by the \(LW\) or \(WL\) mixed correspondence, and local
insertions on both sides are typed by the ordered \(LWL\) correspondence.
No \(LL\)-only local operation recovers a type-II wall profile.

This is recovery of operation \emph{types}.  It does not populate the
wrapped prequotient stacks, mixed extension stacks, anchors, quotient
descent, O2 wall atlas, protected integration, or transition system.
\end{proposition}

\begin{proof}
The triples are those of Proposition~\ref{prop:type-ii-wall-images}.
Applying the wall-branch formula \(d=m+1\) gives
\[
(1,1,0)\mapsto1,\qquad
(0,1,1)\mapsto2,\qquad
(0,-1,0)\mapsto1.
\]
All three integers are positive.  Once a retained branch-comparison row
identifies these integers with \(b_R^{\mathrm{geom}}\), the definition
\[
\Gamma_R^{\mathrm{wr}}=\{b_R^{\mathrm{geom}}>0\},\qquad
\Gamma_R^{\mathrm{loc}}=\{b_R^{\mathrm{geom}}=0\}
\]
places all three \(\eta_i\) in the wrapped window and none in the local
window.  Definition~\ref{def:hybrid-wrapped-stratum-prestack} supplies
the singleton wrapped carrier profile.  Definitions
\ref{def:mixed-local-wrapped-correspondence-datum} and
\ref{def:two-sided-mixed-correspondence-datum} are exactly the before
\(E\)-quotient operation profiles with one wrapped input and local
inputs on one or both sides.  Since an \(LL\) correspondence has only
local inputs and local target colour, it cannot have target colour
\(\eta_i\in\Gamma_R^{\mathrm{wr}}\).  The final sentence separates
operation-profile recovery from the missing populated finite-stage
rows.
\end{proof}

The packet
\[
\texttt{certificates/hybrid/type\_ii\_wall\_hybrid\_operation\_recovery}
\]
verifies
\(\texttt{TYPE\_II\_WALL\_HYBRID\_OPERATION\_RECOVERY\_VERIFIED}\):
the three imported wall triples are distinct, have positive
wall-branch degrees \(1,2,1\), are routed to singleton wrapped colour
profiles, and carry exactly the \(LW\), \(WL\), and \(LWL\) mixed
operation profiles.  Its obstruction rows record that retained
branch-comparison, wall-object, support, anchor, mixed-stack, O2 atlas,
\(E\)-quotient losslessness, quotient, protected-integration,
aggregate-carrier, and transition rows are still missing.

\begin{proposition}[Type-II wall profiles survive the \(E\)-quotient assignment]
\label{prop:type-ii-walls-survive-e-quotient}
\textnormal{[profile-level survival proved; object-level nonvanishing not proved]}
Let \(\eta_i\), \(i=1,2,3\), be the singleton wrapped type-II wall
profiles of Proposition~\ref{prop:type-ii-walls-hybrid-operation-recovery},
and assume the finite quotient-after-correspondence datum of
Definition~\ref{def:quotient-after-correspondence-pseudofunctor}.  Then
the quotient assignment sends each retained unreduced wrapped profile to
\[
\mathcal Q^{\mathrm{corr}}_{E,R}
\bigl(\mathcal M_{\eta_i,R}^{\mathrm{wr,rig}}\bigr)
=
\bigl[\mathcal M_{\eta_i,R}^{\mathrm{wr,rig}}/E\bigr].
\]
For every retained finite local insertion on one or both sides of
\(\eta_i\), the associated \(LW\), \(WL\), or \(LWL\) operation profile is
sent to the quotient of the already formed unreduced correspondence
stack.  Hence none of the three simple type-II wall profiles is deleted
by the reduced \(E\)-quotient at the quotient-after-correspondence
assignment level.

This does not prove that
\(\mathcal M_{\eta_i,R}^{\mathrm{wr,rig}}\) is nonempty, that its quotient
cycle has nonzero Borel--Moore class, that vanishing-cycle coefficients
descend nontrivially, or that the reduced object enters the O2 wall
atlas, protected Pfaffian, aggregate hybrid carrier, or transition
system.
\end{proposition}

\begin{proof}
By Proposition~\ref{prop:type-ii-walls-hybrid-operation-recovery}, the
three simple type-II wall images have positive wall-branch elliptic
degrees \(1,2,1\).  After the retained branch-comparison rows they are
therefore the singleton wrapped colours \(\eta_i\in\Gamma_R^{\mathrm{wr}}\).
The same proposition identifies the finite local insertions on these
wall profiles with the one-sided \(LW\) and \(WL\) mixed operation
profiles and the two-sided \(LWL\) operation profile, all before the
reduced \(E\)-quotient.

Definition~\ref{def:quotient-after-correspondence-pseudofunctor}
assigns \([Y/E]\) to every retained object \(Y\) of
\(\mathsf{Corr}^{E,\mathrm{hyb}}_R\), including wrapped rigidified
prequotients.  Taking \(Y=\mathcal M_{\eta_i,R}^{\mathrm{wr,rig}}\)
gives the displayed quotient object for each \(i\).  The same definition
assigns to every generating unreduced correspondence
\[
e=(Y_0\xleftarrow{p_e}\mathfrak E_e\xrightarrow{q_e}Y_1)
\]
the already-formed quotient correspondence
\[
\overline e=
(\overline Y_0\xleftarrow{\overline p_e}
[\mathfrak E_e/E]\xrightarrow{\overline q_e}\overline Y_1).
\]
Applying this to the \(LW\), \(WL\), and \(LWL\) operation profiles
therefore preserves the operation labels as quotient-after-correspondence
arrows.  Proposition~\ref{prop:quotient-after-correspondence-excludes-quotient-first}
records that this is not a quotient-first reconstruction from reduced
object stacks, and
Proposition~\ref{prop:quotient-after-correspondence-hn-transition-compatibility}
keeps any supplied transition row in the same quotient-presented form.
The final paragraph states the unsupplied rows and prevents this
profile-level assignment theorem from being read as object-level
nonvanishing, O2 atlas construction, or protected integration.
\end{proof}

The packet
\[
\texttt{certificates/hybrid/type\_ii\_wall\_e\_quotient\_survival}
\]
verifies
\(\texttt{TYPE\_II\_WALL\_E\_QUOTIENT\_SURVIVAL\_VERIFIED}\): the three
wrapped type-II wall profiles are assigned quotient objects, their nine
\(LW\), \(WL\), and \(LWL\) operation profiles are assigned
quotient-after-correspondence arrows, quotient-first reconstruction is
excluded, and HN transition compatibility is imported.  Its firewall rows
record that object-level nonzero survival, coefficient descent,
orientation descent, O2 wall-atlas insertion, protected integration, and
aggregate hybrid-carrier population remain open.

\begin{proposition}[No type-I wall enters the type-II hybrid wall inlet]
\label{prop:no-type-i-wall-enters-hybrid-carrier}
\textnormal{[profile-level exclusion proved; aggregate carrier not populated]}
Let
\[
\vartheta_{12}=f_{-2}-f_2,\qquad
\vartheta_{13}=f_2-f_3,\qquad
\vartheta_{23}=f_{-2}-f_3
\]
be the three type-I chamber-automorphism roots of
Lemma~\ref{lem:bkm-chamber-s3}.  At finite HN height \(R\), let the
type-II wall-profile inlet be the finite set of retained profiles
\[
\Gamma_R^{\mathrm{wall},II}
=
\left\{\eta_i\;\middle|\;
\overline\Pi_X(\eta_i)=\gamma_{\mathrm{wall}}(\delta_i),
b_R^{\mathrm{geom}}(\eta_i)=d_{\delta_i},
i=1,2,3
\right\}
\]
appearing in Propositions
\ref{prop:type-ii-walls-hybrid-operation-recovery} and
\ref{prop:type-ii-walls-survive-e-quotient}.  Then every element of
\(\Gamma_R^{\mathrm{wall},II}\) is labelled by one of the three
type-II divisor roots \(\delta_1,\delta_2,\delta_3\).  No type-I root
\(\vartheta_{12},\vartheta_{13},\vartheta_{23}\) supplies a wall profile
in \(\Gamma_R^{\mathrm{wall},II}\), and none is created by the
quotient-after-correspondence assignment.

This is an exclusion theorem for the wall-profile inlet.  It is not a
global exhaustion of all retained hybrid colours, not a population
theorem for the aggregate hybrid carrier, and not an O2 wall-atlas
type-I exclusion row.
\end{proposition}

\begin{proof}
Remark~\ref{rem:bkm-divisibility-types} and the finite
\(\texttt{type\_ii\_chamber\_isotropic}\) packet record the two
square-two divisibility classes.  The roots
\(\delta_1,\delta_2,\delta_3\) have ambient divisibility \(2\), lie in
the type-II sublattice, and occur in \(\operatorname{div}(\Delta_5)\)
with order \(1\).  The roots
\(\vartheta_{12},\vartheta_{13},\vartheta_{23}\) have ambient
divisibility \(1\), generate the \(S_3\) chamber automorphism action of
Lemma~\ref{lem:bkm-chamber-s3}, and have zero
\(\Delta_5\)-divisor order.

Proposition~\ref{prop:type-ii-walls-hybrid-operation-recovery} imports
only the three type-II wall images
\[
\gamma_{\mathrm{wall}}(\delta_1),\qquad
\gamma_{\mathrm{wall}}(\delta_2),\qquad
\gamma_{\mathrm{wall}}(\delta_3)
\]
and routes them to singleton wrapped profiles with positive
wall-branch degrees \(1,2,1\).  Hence the wall-profile inlet just
defined is indexed by the divisor-supported type-II roots and by no
other square-two root.  The type-I reflections act only by permuting
the three \(\delta_i\); they do not add a fourth wall label, and they do
not turn a zero-order type-I chamber wall into a
\(\Delta_5\)-divisor wall.

Proposition~\ref{prop:type-ii-walls-survive-e-quotient} then applies
\(\mathcal Q^{\mathrm{corr}}_{E,R}\) to the same three already formed
wall profiles and to their \(LW\), \(WL\), and \(LWL\) operation
profiles.  Since quotient-after-correspondence assigns quotient objects
and quotient arrows to retained inputs, it preserves this three-element
indexing set and introduces no new type-I profile.  The final paragraph
separates this profile-level firewall from the unsupplied global carrier
and O2-atlas rows.
\end{proof}

The packet
\[
\texttt{certificates/hybrid/type\_i\_false\_wall\_exclusion}
\]
verifies
\(\texttt{TYPE\_I\_FALSE\_WALL\_EXCLUSION\_VERIFIED}\): the admitted
wall-profile rows are exactly the three type-II profiles
\(\delta_1,\delta_2,\delta_3\), the three type-I chamber roots have
divisibility \(1\), zero \(\Delta_5\)-divisor support, and zero
\(\Delta_5\)-divisor order, and the \(E\)-quotient survival row does
not resurrect a type-I wall profile.  Its firewall rows record that
global colour exhaustion, aggregate hybrid-carrier population, retained
wall objects, coefficient descent, O2 wall-atlas insertion, protected
integration, and transition to O2 remain open.

\begin{definition}[Local/local correspondence datum]
\label{def:local-local-correspondence-datum}
Fix a finite HN height \(R\).  Let
\(\alpha:I\to\Gamma_R^{\mathrm{loc}}\) and
\(\beta:J\to\Gamma_R^{\mathrm{loc}}\) be finite local colour maps, and
write \(|\alpha|=\sum_{i\in I}\alpha_i\),
\(|\beta|=\sum_{j\in J}\beta_j\).  A local/local correspondence row is
defined only for a retained target colour
\(\gamma\in\Gamma_R^{\mathrm{loc}}\) satisfying
\[
\gamma=|\alpha|+|\beta|.
\]
The target index set \(K\) is \(I\sqcup J\), modulo the ordinary Ran
collision maps.  The ordered local/local correspondence is the
projection-finite diagram
\[
\mathfrak E^{\mathrm{LL}}_{\alpha,\beta;\gamma,R,I,J,K}
\xrightarrow{p^{\mathrm{LL}}}
\mathfrak M^{\mathrm{loc,cl}}_{\alpha,R,I}
\times
\mathfrak M^{\mathrm{loc,cl}}_{\beta,R,J},
\qquad
\mathfrak E^{\mathrm{LL}}_{\alpha,\beta;\gamma,R,I,J,K}
\xrightarrow{q^{\mathrm{LL}}}
\mathfrak M^{\mathrm{loc,cl}}_{\gamma,R,K},
\]
classifying retained extensions
\[
0\longrightarrow A_{\beta,J}\longrightarrow B_{\gamma,K}
\longrightarrow A_{\alpha,I}\longrightarrow0.
\]
This diagram lives over
\(\operatorname{Ran}^{\mathrm{loc}}_{R,\mathrm{pre}}(E)\times
\operatorname{Ran}^{\mathrm{loc}}_{R,\mathrm{pre}}(E)\) and targets
\(\operatorname{Ran}^{\mathrm{loc}}_{R,\mathrm{pre}}(E)\).  It has no
wrapped input and no reduced \(E\)-quotient.  Symmetric descent,
collision compatibility, properness of the local extension stack,
Thom--Sebastiani transport, and transition compatibility are separate
rows.  The diagram is not a mixed correspondence, not a wrapped/wrapped
correspondence, not a two-sided mixed correspondence, and not an
ordinary uncoloured Ran product.
\end{definition}

The packet
\[
\texttt{certificates/hybrid/local\_local\_correspondence\_definition}
\]
verifies
\(\texttt{LOCAL\_LOCAL\_CORRESPONDENCE\_DEFINITION\_VERIFIED}\):
the ordered \(LL\) correspondence type is defined, while
extension-stack rows, source/target-map rows, descent rows, properness
rows, Thom--Sebastiani rows, collision rows, and transition rows are
not supplied by that definition packet.

\begin{proposition}[Target properness for local/local extensions]
\label{prop:local-local-target-properness}
\textnormal{[proved for the Hall target map]}
Assume the local closed-incidence stacks
\(\mathfrak M^{\mathrm{loc,cl}}_{\alpha,R,I}\),
\(\mathfrak M^{\mathrm{loc,cl}}_{\beta,R,J}\), and
\(\mathfrak M^{\mathrm{loc,cl}}_{\gamma,R,K}\) carry the universal
perfect complexes and the retained extension-closure datum of
Definition~\ref{def:retained-extension-closure-datum}.  Then the
target map
\[
q^{\mathrm{LL}}:
\mathfrak E^{\mathrm{LL}}_{\alpha,\beta;\gamma,R,I,J,K}
\longrightarrow
\mathfrak M^{\mathrm{loc,cl}}_{\gamma,R,K}
\]
is proper.  The source map \(p^{\mathrm{LL}}\) is not asserted to be
proper.
\end{proposition}

\begin{proof}
Let \(B_{\gamma,K}\) be the universal target object on
\(\mathfrak M^{\mathrm{loc,cl}}_{\gamma,R,K}\).  The stack
\(\mathfrak E^{\mathrm{LL}}_{\alpha,\beta;\gamma,R,I,J,K}\) is the
closed substack of the relative Quot scheme over
\(\mathfrak M^{\mathrm{loc,cl}}_{\gamma,R,K}\) parametrising
quotients
\[
B_{\gamma,K}\twoheadrightarrow A_{\alpha,I}
\]
whose kernel is an object \(A_{\beta,J}\), with exact sequence
\[
0\longrightarrow A_{\beta,J}\longrightarrow B_{\gamma,K}
\longrightarrow A_{\alpha,I}\longrightarrow0.
\]
The retained extension-closure datum says that the middle term stays
inside the retained finite HN window.  The local support conditions
over \(E^I\), \(E^J\), and \(E^K\), and the retained numerical colour
conditions \(\alpha,\beta,\gamma\), are closed conditions inside the
relative Quot scheme.  Grothendieck projectivity of Quot
\cite[Theorem~2.2.4]{HuybrechtsLehn} makes the relative Quot scheme
projective over \(\mathfrak M^{\mathrm{loc,cl}}_{\gamma,R,K}\), hence
proper.  A closed substack of a proper stack over the same base is
proper, so \(q^{\mathrm{LL}}\) is proper.

The source map has a different fibre.  Over fixed endpoints
\((A_{\alpha,I},A_{\beta,J})\), the extension parameters contain
\(\operatorname{Ext}^1(A_{\alpha,I},A_{\beta,J})\) modulo the usual
\(\operatorname{Hom}\)-automorphisms.  When this Ext group is nonzero,
the fibre contains affine directions and is not proper in general.
Thus the properness needed for the Hall pushforward is the
target-properness of \(q^{\mathrm{LL}}\), not source-properness of
\(p^{\mathrm{LL}}\).
\end{proof}

The packet
\[
\texttt{certificates/hybrid/local\_local\_extension\_properness}
\]
verifies
\(\texttt{LOCAL\_LOCAL\_TARGET\_PROPERNESS\_VERIFIED}\): the
local/local extension stack is represented as a closed relative Quot
locus over the target local incidence stack, the target map
\(q^{\mathrm{LL}}\) is proper, and the false source-properness
substitute is excluded.  Descent, collision compatibility,
Thom--Sebastiani transport, transition compatibility, and aggregate
hybrid-carrier population remain separate obligations.  The
compact-support exceptional pushforward model is supplied separately by
Definition~\ref{def:compact-support-exceptional-pushforward-model}.

\begin{definition}[One-sided mixed local/wrapped correspondence datum]
\label{def:mixed-local-wrapped-correspondence-datum}
Fix a finite HN height \(R\).  Let
\(\alpha:I\to\Gamma_R^{\mathrm{loc}}\) be a finite local colour map and
write \(|\alpha|=\sum_{i\in I}\alpha_i\).  Let
\(\eta\in\Gamma_R^{\mathrm{wr}}\).  A one-sided mixed correspondence
row is defined only for a retained target colour
\(\zeta\in\Gamma_R^{\mathrm{wr}}\) satisfying the typed equality
\[
\zeta=|\alpha|+\eta
\quad\text{or}\quad
\zeta=\eta+|\alpha|,
\]
with the positive wrapped condition read through \(b_R^{\mathrm{geom}}\).
The two ordered mixed diagrams are the unreduced correspondence stacks
\[
\mathfrak E^{\mathrm{LW}}_{\alpha,\eta;\zeta,R,I}
\xrightarrow{p^{\mathrm{LW}}}
\mathfrak M^{\mathrm{loc,cl}}_{\alpha,R,I}
\times
\mathcal M_{\eta,R}^{\mathrm{wr,rig}},
\qquad
\mathfrak E^{\mathrm{LW}}_{\alpha,\eta;\zeta,R,I}
\xrightarrow{q^{\mathrm{LW}}}
\mathcal M_{\zeta,R}^{\mathrm{wr,rig}},
\]
classifying retained extensions
\[
0\longrightarrow W_\eta\longrightarrow B_\zeta
\longrightarrow A_\alpha\longrightarrow0,
\]
and
\[
\mathfrak E^{\mathrm{WL}}_{\eta,\alpha;\zeta,R,I}
\xrightarrow{p^{\mathrm{WL}}}
\mathcal M_{\eta,R}^{\mathrm{wr,rig}}
\times
\mathfrak M^{\mathrm{loc,cl}}_{\alpha,R,I},
\qquad
\mathfrak E^{\mathrm{WL}}_{\eta,\alpha;\zeta,R,I}
\xrightarrow{q^{\mathrm{WL}}}
\mathcal M_{\zeta,R}^{\mathrm{wr,rig}},
\]
classifying retained extensions
\[
0\longrightarrow A_\alpha\longrightarrow B_\zeta
\longrightarrow W_\eta\longrightarrow0.
\]
The symbols
\(\mathfrak M^{\mathrm{loc,cl}}_{\alpha,R,I}\) and
\(\mathcal M_{\eta,R}^{\mathrm{wr,rig}}\) are the local closed
incidence stack and wrapped rigidified prequotient supplied by the
finite hybrid datum.  These diagrams live before the reduced
\(E\)-quotient.  They are not local-local correspondences, not
wrapped-wrapped correspondences, not ordered two-sided mixed
correspondences, and not a proof of admissibility, base change,
projection formula, Thom--Sebastiani transport, or source/target-anchor
losslessness.
\end{definition}

The packet
\[
\texttt{certificates/hybrid/mixed\_local\_wrapped\_correspondence\_definition}
\]
verifies
\(\texttt{MIXED\_LOCAL\_WRAPPED\_CORRESPONDENCE\_DEFINITION\_VERIFIED}\):
the one-sided \(LW\) and \(WL\) mixed correspondence types are defined,
while extension-stack rows, source/target-map rows, populated
anchor-memory rows, admissibility rows, base-change rows,
projection-formula rows, Thom--Sebastiani rows, and transition rows are
recorded as missing open obligations in that definition packet.

\begin{proposition}[Target admissibility for mixed local/wrapped extensions]
\label{prop:mixed-local-wrapped-target-admissibility}
\textnormal{[proved for the Hall target maps]}
Assume the local closed-incidence stacks
\(\mathfrak M^{\mathrm{loc,cl}}_{\alpha,R,I}\), the wrapped
prequotients \(\mathcal M_{\eta,R}^{\mathrm{wr,rig}}\) and
\(\mathcal M_{\zeta,R}^{\mathrm{wr,rig}}\), the universal perfect
complexes, the retained closed-substack datum, and the retained
extension-closure datum are supplied at the fixed finite HN height
\(R\).  Then the target maps in the one-sided unreduced mixed
correspondences
\[
q^{\mathrm{LW}}:
\mathfrak E^{\mathrm{LW}}_{\alpha,\eta;\zeta,R,I}
\longrightarrow
\mathcal M_{\zeta,R}^{\mathrm{wr,rig}},
\qquad
q^{\mathrm{WL}}:
\mathfrak E^{\mathrm{WL}}_{\eta,\alpha;\zeta,R,I}
\longrightarrow
\mathcal M_{\zeta,R}^{\mathrm{wr,rig}}
\]
are proper.  Hence the target leg needed for the mixed Hall
pushforward is admissible as a proper geometric map before the reduced
\(E\)-quotient.  The source maps \(p^{\mathrm{LW}}\) and
\(p^{\mathrm{WL}}\) are not asserted to be proper.
\end{proposition}

\begin{proof}
Let \(B_\zeta\) be the universal target object on
\(\mathcal M_{\zeta,R}^{\mathrm{wr,rig}}\).  For the \(LW\) order, the
stack
\(\mathfrak E^{\mathrm{LW}}_{\alpha,\eta;\zeta,R,I}\) is the closed
substack of the relative Quot scheme over
\(\mathcal M_{\zeta,R}^{\mathrm{wr,rig}}\) parametrising quotients
\[
B_\zeta\twoheadrightarrow A_{\alpha,I}
\]
whose kernel is an object \(W_\eta\), with exact sequence
\[
0\longrightarrow W_\eta\longrightarrow B_\zeta
\longrightarrow A_{\alpha,I}\longrightarrow0.
\]
The quotient is required to lie in the local closed-incidence stack,
and the kernel is required to lie in the retained wrapped prequotient.
The retained closed-substack datum makes these endpoint conditions
closed on the relative Quot scheme.  The retained colour condition
\(\zeta=|\alpha|+\eta\) and the finite HN window condition are fixed
row conditions; the retained extension-closure datum records that the
allowed middle term remains in the finite window.  Grothendieck
projectivity of Quot \cite[Theorem~2.2.4]{HuybrechtsLehn} makes the
relative Quot scheme projective, hence proper, over
\(\mathcal M_{\zeta,R}^{\mathrm{wr,rig}}\).  A closed substack of this
relative Quot scheme is proper over the same target, so
\(q^{\mathrm{LW}}\) is proper.

For the \(WL\) order one uses the relative Quot scheme over the same
target, but now with quotient
\[
B_\zeta\twoheadrightarrow W_\eta
\]
and kernel \(A_{\alpha,I}\):
\[
0\longrightarrow A_{\alpha,I}\longrightarrow B_\zeta
\longrightarrow W_\eta\longrightarrow0.
\]
The wrapped quotient condition, the local-kernel condition, and the
retained colour condition \(\zeta=\eta+|\alpha|\) are closed row
conditions by the same retained closed-substack and extension-closure
inputs.  The relative Quot scheme is projective over
\(\mathcal M_{\zeta,R}^{\mathrm{wr,rig}}\), and the \(WL\) extension
stack is a closed substack of it.  Hence \(q^{\mathrm{WL}}\) is proper.

The source maps have the usual Hall extension fibres.  Over fixed
endpoints in the \(LW\) order one sees extension classes in
\(\operatorname{Ext}^1(A_{\alpha,I},W_\eta)\) modulo
\(\operatorname{Hom}\)-automorphisms; in the \(WL\) order one sees
\(\operatorname{Ext}^1(W_\eta,A_{\alpha,I})\) modulo the same
automorphism relation.  These fibres contain affine directions when
the corresponding \(\operatorname{Ext}^1\)-group is nonzero, so source
properness is not a consequence of the argument and is not claimed.
The compact-support exceptional pushforward model, base change,
projection formula, Thom--Sebastiani transport, anchor losslessness,
quotient descent, and transition compatibility remain separate rows.
\end{proof}

The packet
\[
\texttt{certificates/hybrid/mixed\_local\_wrapped\_extension\_admissibility}
\]
verifies
\(\texttt{MIXED\_LOCAL\_WRAPPED\_ADMISSIBILITY\_VERIFIED}\): the
one-sided mixed extension stacks are represented as closed relative
Quot loci over the wrapped target prequotient, the target maps
\(q^{\mathrm{LW}}\) and \(q^{\mathrm{WL}}\) are proper, and the false
source-properness substitute is excluded.  The compact-support
exceptional pushforward model is supplied separately by
Definition~\ref{def:compact-support-exceptional-pushforward-model}.
Populated anchor-memory rows, base change, projection formula,
Thom--Sebastiani transport, transition compatibility, quotient descent,
and aggregate hybrid-carrier population remain separate obligations.

\begin{definition}[Wrapped/wrapped correspondence datum]
\label{def:wrapped-wrapped-correspondence-datum}
Fix a finite HN height \(R\).  Let
\(\eta_1,\eta_2\in\Gamma_R^{\mathrm{wr}}\).  A wrapped/wrapped
correspondence row is defined only for a retained target colour
\(\zeta\in\Gamma_R^{\mathrm{wr}}\) satisfying
\[
\zeta=\eta_1+\eta_2,
\]
with the positive wrapped condition read through \(b_R^{\mathrm{geom}}\).
The ordered wrapped/wrapped correspondence is the unreduced diagram
\[
\mathfrak E^{\mathrm{WW}}_{\eta_1,\eta_2;\zeta,R}
\xrightarrow{p^{\mathrm{WW}}}
\mathcal M_{\eta_1,R}^{\mathrm{wr,rig}}
\times
\mathcal M_{\eta_2,R}^{\mathrm{wr,rig}},
\qquad
\mathfrak E^{\mathrm{WW}}_{\eta_1,\eta_2;\zeta,R}
\xrightarrow{q^{\mathrm{WW}}}
\mathcal M_{\zeta,R}^{\mathrm{wr,rig}},
\]
classifying retained extensions
\[
0\longrightarrow W_{\eta_2}\longrightarrow B_\zeta
\longrightarrow W_{\eta_1}\longrightarrow0.
\]
This diagram lives over
\(\operatorname{Ran}^{\mathrm{wr}}_{R,\mathrm{pre}}(E)\times
\operatorname{Ran}^{\mathrm{wr}}_{R,\mathrm{pre}}(E)\) and targets
\(\operatorname{Ran}^{\mathrm{wr}}_{R,\mathrm{pre}}(E)\), before the
reduced \(E\)-quotient.  The source pair is ordered; symmetric descent
for swapping \(\eta_1\) and \(\eta_2\) is a separate row.  Source and
target anchor memory, admissibility, Thom--Sebastiani transport, and
transition compatibility are also separate rows.  The diagram is not a
mixed correspondence, not a local-local correspondence, not a
two-sided mixed correspondence, and not a scalar or quotient-first
replacement.
\end{definition}

The packet
\[
\texttt{certificates/hybrid/wrapped\_wrapped\_correspondence\_definition}
\]
verifies
\(\texttt{WRAPPED\_WRAPPED\_CORRESPONDENCE\_DEFINITION\_VERIFIED}\):
the ordered \(WW\) correspondence type is defined, while extension-stack
rows, source/target-map rows, populated anchor-memory rows, admissibility rows,
Thom--Sebastiani rows, symmetric-descent rows, and transition rows are
recorded as missing open obligations in that definition packet.

\begin{proposition}[Target admissibility for wrapped/wrapped extensions]
\label{prop:wrapped-wrapped-target-admissibility}
\textnormal{[proved for the Hall target map]}
Assume the wrapped prequotients
\(\mathcal M_{\eta_1,R}^{\mathrm{wr,rig}}\),
\(\mathcal M_{\eta_2,R}^{\mathrm{wr,rig}}\), and
\(\mathcal M_{\zeta,R}^{\mathrm{wr,rig}}\), the universal perfect
complexes, the retained closed-substack datum, and the retained
extension-closure datum are supplied at the fixed finite HN height
\(R\).  Then the target map in the ordered unreduced wrapped/wrapped
correspondence
\[
q^{\mathrm{WW}}:
\mathfrak E^{\mathrm{WW}}_{\eta_1,\eta_2;\zeta,R}
\longrightarrow
\mathcal M_{\zeta,R}^{\mathrm{wr,rig}}
\]
is proper.  Hence the target leg needed for the wrapped/wrapped Hall
pushforward is admissible as a proper geometric map before the reduced
\(E\)-quotient.  The source map \(p^{\mathrm{WW}}\) is not asserted to
be proper.
\end{proposition}

\begin{proof}
Let \(B_\zeta\) be the universal target object on
\(\mathcal M_{\zeta,R}^{\mathrm{wr,rig}}\).  The stack
\(\mathfrak E^{\mathrm{WW}}_{\eta_1,\eta_2;\zeta,R}\) is the closed
substack of the relative Quot scheme over
\(\mathcal M_{\zeta,R}^{\mathrm{wr,rig}}\) parametrising quotients
\[
B_\zeta\twoheadrightarrow W_{\eta_1}
\]
whose kernel is an object \(W_{\eta_2}\), with exact sequence
\[
0\longrightarrow W_{\eta_2}\longrightarrow B_\zeta
\longrightarrow W_{\eta_1}\longrightarrow0.
\]
The quotient and kernel are required to lie in the retained wrapped
prequotients of colours \(\eta_1\) and \(\eta_2\).  The retained
closed-substack datum makes these endpoint conditions closed on the
relative Quot scheme.  The retained colour condition
\(\zeta=\eta_1+\eta_2\) and the finite HN window condition are fixed
row conditions; the retained extension-closure datum records that the
allowed middle term remains in the finite window.  Grothendieck
projectivity of Quot \cite[Theorem~2.2.4]{HuybrechtsLehn} makes the
relative Quot scheme projective, hence proper, over
\(\mathcal M_{\zeta,R}^{\mathrm{wr,rig}}\).  A closed substack of this
relative Quot scheme is proper over the same target, so
\(q^{\mathrm{WW}}\) is proper.

The source map has the usual Hall extension fibre.  Over fixed
endpoints \((W_{\eta_1},W_{\eta_2})\), the extension parameters contain
\(\operatorname{Ext}^1(W_{\eta_1},W_{\eta_2})\) modulo the usual
\(\operatorname{Hom}\)-automorphisms.  When this
\(\operatorname{Ext}^1\)-group is nonzero, the fibre contains affine
directions and is not proper in general.  Thus the properness supplied
here is target properness of \(q^{\mathrm{WW}}\), not source
properness of \(p^{\mathrm{WW}}\).  The compact-support exceptional
pushforward model, symmetric descent, populated anchor-memory rows,
Thom--Sebastiani transport, quotient descent, transition compatibility,
and aggregate hybrid-carrier population remain separate rows.
\end{proof}

The packet
\[
\texttt{certificates/hybrid/wrapped\_wrapped\_extension\_admissibility}
\]
verifies
\(\texttt{WRAPPED\_WRAPPED\_ADMISSIBILITY\_VERIFIED}\): the ordered
wrapped/wrapped extension stack is represented as a closed relative
Quot locus over the wrapped target prequotient, the target map
\(q^{\mathrm{WW}}\) is proper, and the false source-properness
substitute is excluded.  The compact-support exceptional pushforward
model is supplied separately by
Definition~\ref{def:compact-support-exceptional-pushforward-model}.
Symmetric descent, populated anchor-memory rows, Thom--Sebastiani
transport, quotient descent, transition compatibility, and aggregate
hybrid-carrier population remain separate obligations.

\begin{definition}[Compact-support exceptional pushforward model]
\label{def:compact-support-exceptional-pushforward-model}
Let
\[
Y\xleftarrow{\ p\ }\mathfrak E\xrightarrow{\ f\ }Z
\]
be one of the retained finite correspondence maps in the local,
mixed, wrapped, or flag correspondence systems at HN height \(R\), and
let
\[
K_{\mathfrak E}\in D^b_c(\mathfrak E)
\]
be the constructible vanishing-cycle complex with orientation line on
the source of \(f\).  A compact-support exceptional pushforward model
for \((f,K_{\mathfrak E})\) is a tuple
\[
\mathfrak Q^{\mathrm{cs}}(f,K_{\mathfrak E})
=
(j_f,\overline{\mathfrak E}_f,\overline f,K_{\mathfrak E},
f^{\mathrm{cs}}_!K_{\mathfrak E})
\]
such that:
\begin{enumerate}
\item \(j_f:\mathfrak E\hookrightarrow\overline{\mathfrak E}_f\) is an
open immersion into a finite-type Artin stack with finite residual
inertia on the retained strata;
\item \(\overline f:\overline{\mathfrak E}_f\to Z\) is proper and
\[
f=\overline f\circ j_f;
\]
\item \(j_{f,!}K_{\mathfrak E}\) exists in the chosen constructible
six-functor theory on the retained d-critical stack, and
\[
f^{\mathrm{cs}}_!K_{\mathfrak E}
:=
\overline f_*\,j_{f,!}K_{\mathfrak E}
\in D^b_c(Z)
\]
is the compact-support exceptional pushforward along \(f\);
\item if \(f\) is proper, the identity compactification
\[
j_f=\mathrm{id}_{\mathfrak E},\qquad
\overline{\mathfrak E}_f=\mathfrak E,\qquad
\overline f=f
\]
is allowed, and then \(f^{\mathrm{cs}}_!=f_*\).
\end{enumerate}
The datum is a chosen compact-support model, not a theorem that two
compactifications give canonically identical functors.  It does not
prove base change, the projection formula, Thom--Sebastiani
compatibility, quotient descent, transition compatibility, symmetric
descent, or associativity on the two-step flag atlas.  Those are the
separate functorial rows required before
\[
q_!\operatorname{TS}^{\mathrm{red}}p^*
\]
is a coherent Hall multiplication on the finite hybrid carrier.
\end{definition}

The packet
\[
\texttt{certificates/hybrid/compact\_support\_exceptional\_pushforward\_model}
\]
verifies
\(\texttt{COMPACT\_SUPPORT\_EXCEPTIONAL\_PUSHFORWARD\_MODEL\_DEFINED}\):
Definition~\ref{def:compact-support-exceptional-pushforward-model}
defines the compactification, extension-by-zero, and proper
pushforward formula for a retained finite correspondence map, and it
records the proper-map specialization \(f^{\mathrm{cs}}_!=f_*\).  It
does not populate all correspondence maps, prove independence of the
compactification, prove base change, prove the projection formula,
transport Thom--Sebastiani, descend through the \(E\)-quotient, prove
transition compatibility, or populate the aggregate hybrid carrier.

\begin{proposition}[Base change for one-sided mixed correspondences]
\label{prop:mixed-correspondence-base-change}
\textnormal{[proved for Cartesian compact-support squares]}
Let \(o\in\{\mathrm{LW},\mathrm{WL}\}\), and let
\[
q^o:\mathfrak E^o_{\alpha,\eta;\zeta,R,I}
\longrightarrow Z^o_{\zeta,R}
:=\mathcal M_{\zeta,R}^{\mathrm{wr,rig}}
\]
be the target map of the corresponding one-sided mixed correspondence.
Let \(K^o\in D^b_c(\mathfrak E^o_{\alpha,\eta;\zeta,R,I})\) be the
constructible vanishing-cycle complex with orientation line.  Suppose
that a retained target morphism \(g:Z'\to Z^o_{\zeta,R}\) is supplied,
and that the mixed correspondence and the compact-support model are
pulled back by Cartesian squares
\[
\begin{tikzcd}
\mathfrak E^{o\prime}\ar[r,"\widetilde g"]\ar[d,"j'_o"']&
\mathfrak E^o_{\alpha,\eta;\zeta,R,I}\ar[d,"j_o"]\\
\overline{\mathfrak E}^{o\prime}\ar[r,"\overline g"]\ar[d,"\overline q^{o\prime}"']&
\overline{\mathfrak E}^{o}\ar[d,"\overline q^{o}"]\\
Z'\ar[r,"g"']&Z^o_{\zeta,R}
\end{tikzcd}
\]
with \(\overline q^o\) and \(\overline q^{o\prime}\) proper, and with
a coefficient identification
\[
K^{o\prime}\simeq \widetilde g^{\,*}K^o.
\]
Then the compact-support exceptional pushforward of
Definition~\ref{def:compact-support-exceptional-pushforward-model}
carries a canonical base-change isomorphism
\[
g^*(q^o)^{\mathrm{cs}}_!K^o
\;\xrightarrow{\;\sim\;}\;
(q^{o\prime})^{\mathrm{cs}}_!K^{o\prime}.
\]
For the proper target maps of
Proposition~\ref{prop:mixed-local-wrapped-target-admissibility}, the
identity compactification gives the ordinary proper-base-change
isomorphism for \(q^{\mathrm{LW}}\) and \(q^{\mathrm{WL}}\).
\end{proposition}

\begin{proof}
By definition,
\[
(q^o)^{\mathrm{cs}}_!K^o
=\overline q^o_*\,j_{o,!}K^o.
\]
The lower square is Cartesian and \(\overline q^o\) is proper, so the
proper base-change isomorphism in the chosen constructible
six-functor theory gives
\[
g^*\overline q^o_*\,j_{o,!}K^o
\simeq
\overline q^{o\prime}_*\,\overline g^{\,*}j_{o,!}K^o.
\]
The upper square is Cartesian and \(j_o\) is an open immersion; hence
extension by zero commutes with this base change:
\[
\overline g^{\,*}j_{o,!}K^o
\simeq
j'_{o,!}\,\widetilde g^{\,*}K^o.
\]
Using the supplied coefficient identification
\(K^{o\prime}\simeq\widetilde g^{\,*}K^o\), the last term becomes
\[
\overline q^{o\prime}_*\,j'_{o,!}K^{o\prime}
=(q^{o\prime})^{\mathrm{cs}}_!K^{o\prime}.
\]
This proves the displayed base-change isomorphism.  The argument
proves only the Cartesian compact-support base-change row for the
chosen model.  It does not prove the projection formula,
Thom--Sebastiani compatibility, quotient descent, transition
compatibility, choice independence for compactifications, or
associativity on the two-step flag atlas.
\end{proof}

The packet
\[
\texttt{certificates/hybrid/mixed\_correspondence\_base\_change}
\]
verifies
\(\texttt{MIXED\_CORRESPONDENCE\_BASE\_CHANGE\_VERIFIED}\):
Proposition~\ref{prop:mixed-correspondence-base-change} supplies the
Cartesian compact-support base-change isomorphisms for the \(LW\) and
\(WL\) target maps, including the proper-target specialization.  It
does not prove the projection formula, Thom--Sebastiani transport,
quotient descent, transition compatibility, compactification-choice
independence, or aggregate hybrid-carrier population.

\begin{proposition}[Projection formula for one-sided mixed correspondences]
\label{prop:mixed-correspondence-projection-formula}
\textnormal{[proved for retained finite-Tor target coefficients]}
Let \(o\in\{\mathrm{LW},\mathrm{WL}\}\), and let
\[
q^o:\mathfrak E^o_{\alpha,\eta;\zeta,R,I}
\longrightarrow Z^o_{\zeta,R}
:=\mathcal M_{\zeta,R}^{\mathrm{wr,rig}}
\]
be the target map of the corresponding one-sided mixed correspondence,
equipped with a chosen compact-support model
\[
\mathfrak E^o_{\alpha,\eta;\zeta,R,I}
\xrightarrow{\,j_o\,}
\overline{\mathfrak E}^{o}
\xrightarrow{\,\overline q^o\,}
Z^o_{\zeta,R},
\qquad q^o=\overline q^o\circ j_o,
\]
where \(j_o\) is an open immersion and \(\overline q^o\) is proper.
Let \(K^o\in D^b_c(\mathfrak E^o_{\alpha,\eta;\zeta,R,I})\) be the
constructible vanishing-cycle complex with orientation line.  Let
\(L\in D^b_c(Z^o_{\zeta,R})\) be a retained target coefficient of
finite Tor-amplitude, so that the open-immersion and proper projection
formulas hold in the chosen constructible six-functor theory.  Then
there is a canonical isomorphism
\[
(q^o)^{\mathrm{cs}}_!
\bigl(K^o\otimes^{\mathbf L}(q^o)^*L\bigr)
\;\xrightarrow{\;\sim\;}\;
(q^o)^{\mathrm{cs}}_!K^o\otimes^{\mathbf L}L .
\]
For the proper target maps of
Proposition~\ref{prop:mixed-local-wrapped-target-admissibility}, the
identity compactification specializes this to the ordinary proper
projection formula for \(q^{\mathrm{LW}}\) and \(q^{\mathrm{WL}}\).
\end{proposition}

\begin{proof}
Since \(q^o=\overline q^o\circ j_o\), one has
\[
(q^o)^*L=j_o^*(\overline q^o)^*L.
\]
Expanding the compact-support model gives
\[
(q^o)^{\mathrm{cs}}_!
\bigl(K^o\otimes^{\mathbf L}(q^o)^*L\bigr)
=
\overline q^o_*\,
j_{o,!}\bigl(K^o\otimes^{\mathbf L}j_o^*(\overline q^o)^*L\bigr).
\]
The projection formula for the open immersion \(j_o\) identifies the
right-hand side with
\[
\overline q^o_*
\bigl(j_{o,!}K^o\otimes^{\mathbf L}(\overline q^o)^*L\bigr).
\]
The map \(\overline q^o\) is proper; the proper projection formula
therefore gives
\[
\overline q^o_*
\bigl(j_{o,!}K^o\otimes^{\mathbf L}(\overline q^o)^*L\bigr)
\simeq
\overline q^o_*j_{o,!}K^o\otimes^{\mathbf L}L
=
(q^o)^{\mathrm{cs}}_!K^o\otimes^{\mathbf L}L.
\]
This proves the stated isomorphism.  The proof uses the chosen
compactification and the retained finite-Tor target coefficient.  It
does not prove compactification-choice independence,
Thom--Sebastiani compatibility, quotient descent, transition
compatibility, or associativity on the two-step flag atlas.
\end{proof}

The packet
\[
\texttt{certificates/hybrid/mixed\_correspondence\_projection\_formula}
\]
verifies
\(\texttt{MIXED\_CORRESPONDENCE\_PROJECTION\_FORMULA\_VERIFIED}\):
Proposition~\ref{prop:mixed-correspondence-projection-formula}
supplies the compact-support projection formula for the \(LW\) and
\(WL\) target maps with retained finite-Tor target coefficients.  It
does not prove Thom--Sebastiani transport, quotient descent, transition
compatibility, compactification-choice independence, or aggregate
hybrid-carrier population.

\begin{proposition}[Thom--Sebastiani transport for one-sided mixed correspondences]
\label{prop:mixed-correspondence-thom-sebastiani}
\textnormal{[proved for supplied reduced additive d-critical charts]}
Let \(o\in\{\mathrm{LW},\mathrm{WL}\}\), and let
\[
Y^o_{\mathrm{src}}
=
\begin{cases}
\mathfrak M^{\mathrm{loc,cl}}_{\alpha,R,I}
\times\mathcal M^{\mathrm{wr,rig}}_{\eta,R},&o=\mathrm{LW},\\[2mm]
\mathcal M^{\mathrm{wr,rig}}_{\eta,R}
\times\mathfrak M^{\mathrm{loc,cl}}_{\alpha,R,I},&o=\mathrm{WL}
\end{cases}
\]
be the ordered source product of the one-sided mixed correspondence
with source map \(p^o:\mathfrak E^o_{\alpha,\eta;\zeta,R,I}\to
Y^o_{\mathrm{src}}\).  Write
\[
K^o_{\mathrm{src}}
=
\Phi^{\mathrm{red}}_{Y^o_{\mathrm{src}}}\otimes
o^{\mathrm{red}}_{Y^o_{\mathrm{src}}},
\qquad
K^o_{\mathfrak E}
=
\Phi^{\mathrm{red}}_{\mathfrak E^o_{\alpha,\eta;\zeta,R,I}}
\otimes
o^{\mathrm{red}}_{\mathfrak E^o_{\alpha,\eta;\zeta,R,I}}.
\]
Assume that the retained mixed extension stack is covered by
Joyce--BBDJS d-critical charts for which:
\begin{enumerate}
\item the extension potential is the pullback of the ordered sum of
the two source potentials, up to the standard nondegenerate quadratic
stabilization;
\item the K3 semiregularity cosection is additive along the mixed
extension correspondence, so the reduced obstruction theories and
reduced vanishing-cycle complexes are pulled from the same quotient
of the ordinary obstruction theory;
\item the Joyce--Upmeier orientation line on
\(\mathfrak E^o_{\alpha,\eta;\zeta,R,I}\) is identified with
\((p^o)^*o^{\mathrm{red}}_{Y^o_{\mathrm{src}}}\), and the
orientation Thom--Sebastiani defect is zero on chart overlaps.
\end{enumerate}
Then the mixed correspondence carries a canonical reduced
Thom--Sebastiani transport isomorphism
\[
\operatorname{TS}^{\mathrm{red},o}_{\alpha,\eta;R,I}:
(p^o)^*K^o_{\mathrm{src}}
\;\xrightarrow{\;\sim\;}\;
K^o_{\mathfrak E}
\]
in \(D^b_c(\mathfrak E^o_{\alpha,\eta;\zeta,R,I})\).  The construction
is ordered: the \(LW\) and \(WL\) isomorphisms differ by the Koszul
symmetry of the source product and are recorded as separate rows.
\end{proposition}

\begin{proof}
On a retained d-critical chart, condition (1) puts the mixed extension
potential in the BBDJS Thom--Sebastiani normal form for the ordered
sum of the two source potentials, after harmless quadratic
stabilization \cite{BBDJS2015}.  Hence the BBDJS
Thom--Sebastiani theorem gives an isomorphism between the pullback of
the exterior product of the two source vanishing-cycle complexes and
the vanishing-cycle complex of the extension chart.  Condition (2)
identifies the ordinary and cosection-reduced sides of this isomorphism
after quotienting by the common additive semiregularity cosection.
Condition (3) tensors the vanishing-cycle isomorphism with the
Joyce--Upmeier orientation transport and makes it independent of the
chosen chart overlap.  The descent data therefore glue the chartwise
isomorphisms to the displayed object of
\(D^b_c(\mathfrak E^o_{\alpha,\eta;\zeta,R,I})\).

The proof supplies only the binary mixed Thom--Sebastiani transport
for the chosen reduced d-critical and orientation atlas.  It does not
construct the orientation line or discharge \((O1)\).  The
quotient-after-correspondence descent of this transport is
Proposition~\ref{prop:quotient-after-correspondence-preserves-thom-sebastiani}.
Transition compatibility and the eight-word two-step flag pentagon
remain separate obligations.
\end{proof}

The packet
\[
\texttt{certificates/hybrid/mixed\_correspondence\_thom\_sebastiani}
\]
verifies
\(\texttt{MIXED\_THOM\_SEBASTIANI\_TRANSPORT\_CONDITIONAL\_VERIFIED}\):
Proposition~\ref{prop:mixed-correspondence-thom-sebastiani} supplies
the binary mixed Thom--Sebastiani transport for \(LW\) and \(WL\) from
supplied additive reduced d-critical charts and zero-defect orientation
transport.  It does not prove \((O1)\), quotient descent, transition
compatibility, two-step flag associativity, or aggregate
hybrid-carrier population.

\begin{definition}[Eight local/wrapped binary words]
\label{def:eight-local-wrapped-binary-words}
At a fixed finite HN height \(R\), let
\[
\mathbb A_{\mathrm{hyb}}=\{\mathrm L,\mathrm W\}
\]
be the ordered hybrid alphabet, where \(\mathrm L\) denotes a
projection-finite local input of elliptic degree \(0\), and
\(\mathrm W\) denotes a rigidified wrapped input of positive elliptic
degree.  The binary type product is
\[
\mathrm L\star\mathrm L=\mathrm L,\qquad
\mathrm L\star\mathrm W=\mathrm W,\qquad
\mathrm W\star\mathrm L=\mathrm W,\qquad
\mathrm W\star\mathrm W=\mathrm W.
\]
Thus a length-three word
\[
w=\epsilon_1\epsilon_2\epsilon_3\in\mathbb A_{\mathrm{hyb}}^3
\]
has left and right intermediate types
\[
\epsilon_{12}=\epsilon_1\star\epsilon_2,\qquad
\epsilon_{23}=\epsilon_2\star\epsilon_3,
\]
and final type
\[
\epsilon_{123}=(\epsilon_1\star\epsilon_2)\star\epsilon_3
=\epsilon_1\star(\epsilon_2\star\epsilon_3).
\]
The eight ordered words are
\[
\mathcal W^{(3)}_{\mathrm{hyb}}
=
\{\mathrm{LLL},\mathrm{LLW},\mathrm{LWL},\mathrm{WLL},
\mathrm{LWW},\mathrm{WLW},\mathrm{WWL},\mathrm{WWW}\}.
\]
For compatible ordered charges \((\xi_1,\xi_2,\xi_3)\), the two
parenthesization symbols attached to \(w\) are
\[
((\epsilon_1\epsilon_2)\epsilon_3):
(\xi_1,\xi_2,\xi_3)\mapsto((\xi_1+\xi_2),\xi_3),
\]
\[
(\epsilon_1(\epsilon_2\epsilon_3)):
(\xi_1,\xi_2,\xi_3)\mapsto(\xi_1,(\xi_2+\xi_3)).
\]
This is only the word and type convention.  It does not construct the
two-step flag stacks, prove that the stacks are retained, or prove
associativity, pentagon coherence, quotient descent, transition
compatibility, or aggregate hybrid-carrier population.
\end{definition}

The packet
\[
\texttt{certificates/hybrid/eight\_word\_binary\_vocabulary}
\]
verifies
\(\texttt{EIGHT\_WORD\_BINARY\_VOCABULARY\_DEFINED}\):
Definition~\ref{def:eight-local-wrapped-binary-words} fixes the
ordered alphabet, the type product, the eight length-three words, and
the two parenthesization symbols.  It does not construct the two-step
flag stacks or prove associativity.

\begin{proposition}[Eight two-step local/wrapped flag stacks]
\label{prop:eight-two-step-flag-stacks}
\textnormal{[constructed as retained flag-Quot substacks]}
Assume the retained local and wrapped object stacks, the four ordered
binary correspondence stacks \(LL,LW,WL,WW\), the retained universal
complexes, closed endpoint conditions, and retained extension closure
at height \(R\).  For every word
\[
w=\epsilon_1\epsilon_2\epsilon_3\in
\mathcal W^{(3)}_{\mathrm{hyb}}
\]
and every compatible ordered colour triple
\((\xi_1,\xi_2,\xi_3)\) of those types, there is a retained
finite-type Artin stack
\[
\mathfrak F^{(2),w}_{\xi_1,\xi_2,\xi_3,R}
\]
classifying two-step filtrations
\[
0=C_0\subset C_1\subset C_2\subset C_3
\]
in the finite HN heart such that
\[
C_i/C_{i-1}\in\mathfrak M^{\epsilon_i}_{\xi_i,R}
\qquad (i=1,2,3),
\]
the intermediate objects \(C_2/C_0\), \(C_3/C_1\), and \(C_3/C_0\)
have the types
\[
\epsilon_{12},\qquad \epsilon_{23},\qquad \epsilon_{123}
\]
of Definition~\ref{def:eight-local-wrapped-binary-words}, and every
intermediate colour is retained.  If one of these types is
\(\mathrm W\), the corresponding object carries the fixed wrapped
rigidification and the before-quotient anchor memory.

Let
\[
\mathfrak E^{\epsilon\delta}_{\xi,\upsilon;\xi+\upsilon,R}
\]
denote the ordered binary correspondence stack of type
\(\epsilon\delta\), interpreted as \(LL,LW,WL\), or \(WW\) by
Definitions~\ref{def:local-local-correspondence-datum},
\ref{def:mixed-local-wrapped-correspondence-datum}, and
\ref{def:wrapped-wrapped-correspondence-datum}.  The flag stack carries
canonical comparison morphisms
\[
\lambda^w:
\mathfrak F^{(2),w}_{\xi_1,\xi_2,\xi_3,R}
\longrightarrow
\mathfrak E^{\epsilon_{12}\epsilon_3}_{\xi_1+\xi_2,\xi_3;\xi_1+\xi_2+\xi_3,R}
\times_{\mathfrak M^{\epsilon_{12}}_{\xi_1+\xi_2,R}}
\mathfrak E^{\epsilon_1\epsilon_2}_{\xi_1,\xi_2;\xi_1+\xi_2,R}
\]
and
\[
\rho^w:
\mathfrak F^{(2),w}_{\xi_1,\xi_2,\xi_3,R}
\longrightarrow
\mathfrak E^{\epsilon_1\epsilon_{23}}_{\xi_1,\xi_2+\xi_3;\xi_1+\xi_2+\xi_3,R}
\times_{\mathfrak M^{\epsilon_{23}}_{\xi_2+\xi_3,R}}
\mathfrak E^{\epsilon_2\epsilon_3}_{\xi_2,\xi_3;\xi_2+\xi_3,R}.
\]
The maps are defined by the two forgetful operations on the filtration:
\(\lambda^w\) remembers \(C_1\subset C_2\subset C_3\) as
\(((\xi_1,\xi_2),\xi_3)\), while \(\rho^w\) remembers it as
\((\xi_1,(\xi_2,\xi_3))\).
\end{proposition}

\begin{proof}
Fix \(w\) and \((\xi_1,\xi_2,\xi_3)\).  Over the retained terminal
object stack of colour \(\xi_1+\xi_2+\xi_3\), take the relative
two-step flag-Quot stack parametrising quotients
\[
C_3\twoheadrightarrow C_3/C_2,\qquad
C_2\twoheadrightarrow C_2/C_1
\]
or equivalently subobjects
\(C_1\subset C_2\subset C_3\).  Grothendieck projectivity for
relative Quot schemes and its flag version give finite type over the
retained terminal stack \cite[Theorem~2.2.4]{HuybrechtsLehn}.  The
conditions that the three graded pieces lie in the prescribed local or
wrapped retained substacks, that the intermediate colours equal the
displayed sums, and that the wrapped pieces carry the fixed
rigidifications and anchor memory are closed retained row conditions.
The retained extension-closure packet ensures that \(C_2/C_0\),
\(C_3/C_1\), and \(C_3/C_0\) remain in the finite HN window.  Hence
the required \(\mathfrak F^{(2),w}_{\xi_1,\xi_2,\xi_3,R}\) is a
retained finite-type closed substack of the flag-Quot stack.

Forgetting the first or the second extension step gives the two
displayed morphisms to the iterated correspondence fibre products.
The construction is carried out before the reduced \(E\)-quotient and
only constructs the eight flag stacks and their comparison maps.  It
does not prove that \(\lambda^w\) and \(\rho^w\) induce the same Hall
operation, does not prove associativity defects vanish, and does not
prove the four-input pentagon, quotient descent, transition
compatibility, or aggregate hybrid-carrier population.
\end{proof}

The packet
\[
\texttt{certificates/hybrid/eight\_word\_two\_step\_flag\_stacks}
\]
verifies
\(\texttt{EIGHT\_WORD\_TWO\_STEP\_FLAG\_STACKS\_CONSTRUCTED}\):
Proposition~\ref{prop:eight-two-step-flag-stacks} constructs the
retained two-step flag stack and the two comparison maps for each of
the eight words.  It does not prove associativity, pentagon coherence,
quotient descent, transition compatibility, or aggregate population.

\begin{proposition}[Associativity for the \(LLL\) word]
\label{prop:lll-associativity}
\textnormal{[proved for the pure local two-step flag stack]}
Let \(\alpha,\beta,\delta\) be retained local colours at height \(R\),
and write
\[
\gamma_{12}=\alpha+\beta,\qquad
\gamma_{23}=\beta+\delta,\qquad
\gamma_{123}=\alpha+\beta+\delta .
\]
Assume the \(LLL\) flag stack
\(\mathfrak F^{(2),\mathrm{LLL}}_{\alpha,\beta,\delta,R}\) of
Proposition~\ref{prop:eight-two-step-flag-stacks}, the local/local
target-proper correspondence rows of
Proposition~\ref{prop:local-local-target-properness}, the chosen
compact-support pushforward model, and the local base-change,
projection-formula, and reduced Thom--Sebastiani coefficient
identifications on the \(LLL\) flag square.  Then the pure local Hall
product satisfies
\[
m^{\mathrm{LL}}_{\gamma_{12},\delta,R}
\circ
\bigl(m^{\mathrm{LL}}_{\alpha,\beta,R}\otimes 1\bigr)
=
m^{\mathrm{LL}}_{\alpha,\gamma_{23},R}
\circ
\bigl(1\otimes m^{\mathrm{LL}}_{\beta,\delta,R}\bigr)
\]
as morphisms on the retained compactly supported reduced
vanishing-cycle complexes of the three local inputs.  Equivalently,
the \(LLL\) associativity defect class
\[
\mathfrak o^{\mathrm{assoc}}_{\mathrm{LLL},R}
\]
vanishes.
\end{proposition}

\begin{proof}
Both composites are pull-push operations attached to the same
two-step local filtration stack
\[
0\subset C_1\subset C_2\subset C_3
\]
with graded pieces of colours \(\alpha,\beta,\delta\).  The left
parenthesization uses the comparison map \(\lambda^{\mathrm{LLL}}\) to
first form the \(LL\) extension of \(\alpha\) by \(\beta\), then the
extension of \(\gamma_{12}\) by \(\delta\).  The right
parenthesization uses \(\rho^{\mathrm{LLL}}\) to first form the
extension of \(\beta\) by \(\delta\), then the extension of \(\alpha\)
by \(\gamma_{23}\).  These are two descriptions of the same flag
object, so their source and target maps factor through the same
retained flag stack.

Proper base change identifies the two iterated target pushforwards
after pullback to the flag stack; the projection formula moves the
source coefficients through the same compact-support pushforward; and
the reduced Thom--Sebastiani identification on the \(LLL\) flag square
identifies the two tensor products of vanishing-cycle and orientation
lines.  Hence the two pull-push composites agree.  This proves
\(\mathfrak o^{\mathrm{assoc}}_{\mathrm{LLL},R}=0\).

The argument is purely local.  It proves no associativity statement
for \(LLW,LWL,WLL,LWW,WLW,WWL\), or \(WWW\), and it does not prove the
four-input pentagon, quotient descent, transition compatibility, or
aggregate hybrid-carrier population.
\end{proof}

The packet
\[
\texttt{certificates/hybrid/lll\_associativity}
\]
verifies
\(\texttt{LLL\_ASSOCIATIVITY\_CONDITIONAL\_VERIFIED}\):
Proposition~\ref{prop:lll-associativity} proves the \(LLL\) word
associativity defect vanishes under the stated local flag-stack,
base-change, projection-formula, and reduced Thom--Sebastiani
coefficient rows.  It does not prove the seven remaining word
associativity rows or the four-input pentagon.

\begin{proposition}[Associativity for the \(LLW\) word]
\label{prop:llw-associativity}
\textnormal{[proved for the local/local--mixed flag stack]}
Let \(\alpha,\beta\) be retained local colours and let \(\eta\) be a
retained wrapped colour at height \(R\).  Write
\[
\gamma_{12}=\alpha+\beta,\qquad
\zeta_{23}=\beta+\eta,\qquad
\zeta_{123}=\alpha+\beta+\eta .
\]
Assume the \(LLW\) flag stack
\(\mathfrak F^{(2),\mathrm{LLW}}_{\alpha,\beta,\eta,R}\) of
Proposition~\ref{prop:eight-two-step-flag-stacks}, the \(LL\) and
\(LW\) target-admissible correspondence rows, the chosen
compact-support pushforward model, and the required local and mixed
base-change, projection-formula, and reduced Thom--Sebastiani
coefficient identifications on the \(LLW\) flag square.  Then
\[
m^{\mathrm{LW}}_{\gamma_{12},\eta,R}
\circ
\bigl(m^{\mathrm{LL}}_{\alpha,\beta,R}\otimes 1\bigr)
=
m^{\mathrm{LW}}_{\alpha,\zeta_{23},R}
\circ
\bigl(1\otimes m^{\mathrm{LW}}_{\beta,\eta,R}\bigr)
\]
as morphisms on the retained compactly supported reduced
vanishing-cycle complexes of the ordered inputs.  Equivalently,
\[
\mathfrak o^{\mathrm{assoc}}_{\mathrm{LLW},R}=0 .
\]
\end{proposition}

\begin{proof}
Both composites are represented by the same retained flag stack
\[
0\subset C_1\subset C_2\subset C_3,
\qquad
C_1/C_0\in L_\alpha,\quad
C_2/C_1\in L_\beta,\quad
C_3/C_2\in W_\eta .
\]
The left comparison map \(\lambda^{\mathrm{LLW}}\) records first the
local/local extension of \(\alpha\) by \(\beta\), and then the
local/wrapped extension of \(\gamma_{12}\) by \(\eta\).  The right
comparison map \(\rho^{\mathrm{LLW}}\) records first the
local/wrapped extension of \(\beta\) by \(\eta\), and then the
local/wrapped extension of \(\alpha\) by \(\zeta_{23}\).  These two
iterated descriptions have the same underlying filtered object.

The local and mixed base-change isomorphisms identify the two iterated
compact-support pushforwards after pullback to
\(\mathfrak F^{(2),\mathrm{LLW}}_{\alpha,\beta,\eta,R}\).  The
projection formulas move the external source coefficients through the
same target pushforward, and the reduced Thom--Sebastiani transports
identify the two coefficient systems on the common flag stack.  Hence
the two composites agree, so
\(\mathfrak o^{\mathrm{assoc}}_{\mathrm{LLW},R}=0\).

The proof uses no wrapped/wrapped binary product.  It proves no
associativity statement for \(LWL,WLL,LWW,WLW,WWL\), or \(WWW\), and it
does not prove the four-input pentagon, quotient descent, transition
compatibility, or aggregate hybrid-carrier population.
\end{proof}

The packet
\[
\texttt{certificates/hybrid/llw\_associativity}
\]
verifies
\(\texttt{LLW\_ASSOCIATIVITY\_CONDITIONAL\_VERIFIED}\):
Proposition~\ref{prop:llw-associativity} proves the \(LLW\) word
associativity defect vanishes under the stated \(LL/LW\) flag-stack,
base-change, projection-formula, and reduced Thom--Sebastiani
coefficient rows.  It does not prove the six remaining word
associativity rows or the four-input pentagon.

\begin{proposition}[Associativity for the \(LWL\) word]
\label{prop:lwl-associativity}
\textnormal{[proved for the mixed--mixed flag stack]}
Let \(\alpha,\beta\) be retained local colours and let \(\eta\) be a
retained wrapped colour at height \(R\), ordered as
\((\alpha,\eta,\beta)\).  Write
\[
\zeta_{12}=\alpha+\eta,\qquad
\zeta_{23}=\eta+\beta,\qquad
\zeta_{123}=\alpha+\eta+\beta .
\]
Assume the \(LWL\) flag stack
\(\mathfrak F^{(2),\mathrm{LWL}}_{\alpha,\eta,\beta,R}\) of
Proposition~\ref{prop:eight-two-step-flag-stacks}, the \(LW\) and
\(WL\) target-admissible correspondence rows, the chosen
compact-support pushforward model, and the required mixed base-change,
projection-formula, and reduced Thom--Sebastiani coefficient
identifications on the \(LWL\) flag square.  Then
\[
m^{\mathrm{WL}}_{\zeta_{12},\beta,R}
\circ
\bigl(m^{\mathrm{LW}}_{\alpha,\eta,R}\otimes 1\bigr)
=
m^{\mathrm{LW}}_{\alpha,\zeta_{23},R}
\circ
\bigl(1\otimes m^{\mathrm{WL}}_{\eta,\beta,R}\bigr)
\]
as morphisms on the retained compactly supported reduced
vanishing-cycle complexes of the ordered inputs.  Equivalently,
\[
\mathfrak o^{\mathrm{assoc}}_{\mathrm{LWL},R}=0 .
\]
\end{proposition}

\begin{proof}
The common refinement is the retained \(LWL\) flag stack
\[
0\subset C_1\subset C_2\subset C_3,
\qquad
C_1/C_0\in L_\alpha,\quad
C_2/C_1\in W_\eta,\quad
C_3/C_2\in L_\beta .
\]
The left comparison \(\lambda^{\mathrm{LWL}}\) reads this filtration as
an \(LW\) extension followed by a \(WL\) extension.  The right
comparison \(\rho^{\mathrm{LWL}}\) reads it as a \(WL\) extension
followed by an \(LW\) extension.  Both descriptions have the same
terminal filtered object of wrapped colour \(\zeta_{123}\).

The mixed base-change isomorphisms identify the two iterated
compact-support pushforwards on the common flag stack.  The mixed
projection formula moves the ordered source coefficients through the
same pushforward, and the mixed reduced Thom--Sebastiani transports
identify the two coefficient systems.  Therefore the two composites
agree, and
\(\mathfrak o^{\mathrm{assoc}}_{\mathrm{LWL},R}=0\).

The proof is the ordered \(LWL\) case only.  It proves no
associativity statement for \(WLL,LWW,WLW,WWL\), or \(WWW\), and it
does not prove the four-input pentagon, quotient descent, transition
compatibility, or aggregate hybrid-carrier population.
\end{proof}

The packet
\[
\texttt{certificates/hybrid/lwl\_associativity}
\]
verifies
\(\texttt{LWL\_ASSOCIATIVITY\_CONDITIONAL\_VERIFIED}\):
Proposition~\ref{prop:lwl-associativity} proves the \(LWL\) word
associativity defect vanishes under the stated \(LW/WL\) flag-stack,
base-change, projection-formula, and reduced Thom--Sebastiani
coefficient rows.  It does not prove the five remaining word
associativity rows or the four-input pentagon.

\begin{proposition}[Associativity for the \(WLL\) word]
\label{prop:wll-associativity}
\textnormal{[proved for the mixed--local/local flag stack]}
Let \(\alpha,\beta\) be retained local colours and let \(\eta\) be a
retained wrapped colour at height \(R\), ordered as
\((\eta,\alpha,\beta)\).  Write
\[
\zeta_{12}=\eta+\alpha,\qquad
\gamma_{23}=\alpha+\beta,\qquad
\zeta_{123}=\eta+\alpha+\beta .
\]
Assume the \(WLL\) flag stack
\(\mathfrak F^{(2),\mathrm{WLL}}_{\eta,\alpha,\beta,R}\) of
Proposition~\ref{prop:eight-two-step-flag-stacks}, the \(WL\) and
\(LL\) target-admissible correspondence rows, the chosen
compact-support pushforward model, and the required local and mixed
base-change, projection-formula, and reduced Thom--Sebastiani
coefficient identifications on the \(WLL\) flag square.  Then
\[
m^{\mathrm{WL}}_{\zeta_{12},\beta,R}
\circ
\bigl(m^{\mathrm{WL}}_{\eta,\alpha,R}\otimes 1\bigr)
=
m^{\mathrm{WL}}_{\eta,\gamma_{23},R}
\circ
\bigl(1\otimes m^{\mathrm{LL}}_{\alpha,\beta,R}\bigr)
\]
as morphisms on the retained compactly supported reduced
vanishing-cycle complexes of the ordered inputs.  Equivalently,
\[
\mathfrak o^{\mathrm{assoc}}_{\mathrm{WLL},R}=0 .
\]
\end{proposition}

\begin{proof}
The common refinement is the retained \(WLL\) flag stack
\[
0\subset C_1\subset C_2\subset C_3,
\qquad
C_1/C_0\in W_\eta,\quad
C_2/C_1\in L_\alpha,\quad
C_3/C_2\in L_\beta .
\]
The left comparison \(\lambda^{\mathrm{WLL}}\) records first the
wrapped/local extension of \(\eta\) by \(\alpha\), then the
wrapped/local extension of \(\zeta_{12}\) by \(\beta\).  The right
comparison \(\rho^{\mathrm{WLL}}\) records first the local/local
extension of \(\alpha\) by \(\beta\), then the wrapped/local extension
of \(\eta\) by \(\gamma_{23}\).  Both descriptions have terminal
wrapped colour \(\zeta_{123}\).

The mixed and local base-change isomorphisms identify the two iterated
compact-support pushforwards on the common flag stack.  The projection
formulas move the ordered source coefficients through the same
pushforward, and the reduced Thom--Sebastiani transports identify the
two coefficient systems.  Therefore the two composites agree, and
\(\mathfrak o^{\mathrm{assoc}}_{\mathrm{WLL},R}=0\).

The proof uses no wrapped/wrapped binary product.  It proves no
associativity statement for \(LWW,WLW,WWL\), or \(WWW\), and it does
not prove the four-input pentagon, quotient descent, transition
compatibility, or aggregate hybrid-carrier population.
\end{proof}

The packet
\[
\texttt{certificates/hybrid/wll\_associativity}
\]
verifies
\(\texttt{WLL\_ASSOCIATIVITY\_CONDITIONAL\_VERIFIED}\):
Proposition~\ref{prop:wll-associativity} proves the \(WLL\) word
associativity defect vanishes under the stated \(WL/LL\) flag-stack,
base-change, projection-formula, and reduced Thom--Sebastiani
coefficient rows.  It does not prove the four remaining word
associativity rows or the four-input pentagon.

\begin{proposition}[Associativity for the \(LWW\) word]
\label{prop:lww-associativity}
\textnormal{[proved conditionally for the mixed--wrapped/wrapped flag stack]}
Let \(\alpha\) be a retained local colour and let \(\eta_1,\eta_2\)
be retained wrapped colours at height \(R\), ordered as
\((\alpha,\eta_1,\eta_2)\).  Write
\[
\zeta_{12}=\alpha+\eta_1,\qquad
\zeta_{23}=\eta_1+\eta_2,\qquad
\zeta_{123}=\alpha+\eta_1+\eta_2 .
\]
Assume the \(LWW\) flag stack
\(\mathfrak F^{(2),\mathrm{LWW}}_{\alpha,\eta_1,\eta_2,R}\) of
Proposition~\ref{prop:eight-two-step-flag-stacks}, the \(LW\) and
\(WW\) target-admissible correspondence rows, the chosen
compact-support pushforward model, and the required mixed and
wrapped/wrapped base-change, projection-formula, and reduced
Thom--Sebastiani coefficient identifications on the \(LWW\) flag
square.  Then
\[
m^{\mathrm{WW}}_{\zeta_{12},\eta_2,R}
\circ
\bigl(m^{\mathrm{LW}}_{\alpha,\eta_1,R}\otimes 1\bigr)
=
m^{\mathrm{LW}}_{\alpha,\zeta_{23},R}
\circ
\bigl(1\otimes m^{\mathrm{WW}}_{\eta_1,\eta_2,R}\bigr)
\]
as morphisms on the retained compactly supported reduced
vanishing-cycle complexes of the ordered inputs.  Equivalently,
\[
\mathfrak o^{\mathrm{assoc}}_{\mathrm{LWW},R}=0 .
\]
The wrapped/wrapped base-change, projection-formula, and reduced
Thom--Sebastiani rows are hypotheses of this proposition; they are not
supplied by the wrapped/wrapped target-admissibility theorem.
\end{proposition}

\begin{proof}
The common refinement is the retained \(LWW\) flag stack
\[
0\subset C_1\subset C_2\subset C_3,
\qquad
C_1/C_0\in L_\alpha,\quad
C_2/C_1\in W_{\eta_1},\quad
C_3/C_2\in W_{\eta_2}.
\]
The left comparison \(\lambda^{\mathrm{LWW}}\) records first the
local/wrapped extension of \(\alpha\) by \(\eta_1\), then the
wrapped/wrapped extension of \(\zeta_{12}\) by \(\eta_2\).  The right
comparison \(\rho^{\mathrm{LWW}}\) records first the wrapped/wrapped
extension of \(\eta_1\) by \(\eta_2\), then the local/wrapped
extension of \(\alpha\) by \(\zeta_{23}\).  Both descriptions have
terminal wrapped colour \(\zeta_{123}\).

The mixed and wrapped/wrapped base-change isomorphisms identify the
two iterated compact-support pushforwards on the common flag stack.
The corresponding projection formulas move the ordered source
coefficients through the same pushforward, and the reduced
Thom--Sebastiani transports identify the two coefficient systems.
Therefore the two composites agree, and
\(\mathfrak o^{\mathrm{assoc}}_{\mathrm{LWW},R}=0\).

The proof is wordwise and conditional on the wrapped/wrapped
functorial coefficient rows.  It proves no associativity statement for
\(WLW,WWL\), or \(WWW\), and it does not prove the four-input
pentagon, quotient descent, transition compatibility, or aggregate
hybrid-carrier population.
\end{proof}

The packet
\[
\texttt{certificates/hybrid/lww\_associativity}
\]
verifies
\(\texttt{LWW\_ASSOCIATIVITY\_CONDITIONAL\_VERIFIED}\):
Proposition~\ref{prop:lww-associativity} proves the \(LWW\) word
associativity defect vanishes under the stated \(LW/WW\) flag-stack,
mixed functorial rows, and wrapped/wrapped functorial hypotheses.  It
does not prove the three remaining word associativity rows or the
four-input pentagon.

\begin{proposition}[Associativity for the \(WLW\) word]
\label{prop:wlw-associativity}
\textnormal{[proved conditionally for the mixed--wrapped/wrapped flag stack]}
Let \(\alpha\) be a retained local colour and let \(\eta_1,\eta_2\)
be retained wrapped colours at height \(R\), ordered as
\((\eta_1,\alpha,\eta_2)\).  Write
\[
\zeta_{12}=\eta_1+\alpha,\qquad
\zeta_{23}=\alpha+\eta_2,\qquad
\zeta_{123}=\eta_1+\alpha+\eta_2 .
\]
Assume the \(WLW\) flag stack
\(\mathfrak F^{(2),\mathrm{WLW}}_{\eta_1,\alpha,\eta_2,R}\) of
Proposition~\ref{prop:eight-two-step-flag-stacks}, the \(WL\), \(LW\),
and \(WW\) target-admissible correspondence rows, the chosen
compact-support pushforward model, and the required mixed and
wrapped/wrapped base-change, projection-formula, and reduced
Thom--Sebastiani coefficient identifications on the \(WLW\) flag
square.  Then
\[
m^{\mathrm{WW}}_{\zeta_{12},\eta_2,R}
\circ
\bigl(m^{\mathrm{WL}}_{\eta_1,\alpha,R}\otimes 1\bigr)
=
m^{\mathrm{WW}}_{\eta_1,\zeta_{23},R}
\circ
\bigl(1\otimes m^{\mathrm{LW}}_{\alpha,\eta_2,R}\bigr)
\]
as morphisms on the retained compactly supported reduced
vanishing-cycle complexes of the ordered inputs.  Equivalently,
\[
\mathfrak o^{\mathrm{assoc}}_{\mathrm{WLW},R}=0 .
\]
The wrapped/wrapped base-change, projection-formula, and reduced
Thom--Sebastiani rows are hypotheses of this proposition; they are not
supplied by the wrapped/wrapped target-admissibility theorem.
\end{proposition}

\begin{proof}
The common refinement is the retained \(WLW\) flag stack
\[
0\subset C_1\subset C_2\subset C_3,
\qquad
C_1/C_0\in W_{\eta_1},\quad
C_2/C_1\in L_\alpha,\quad
C_3/C_2\in W_{\eta_2}.
\]
The left comparison \(\lambda^{\mathrm{WLW}}\) records first the
wrapped/local extension of \(\eta_1\) by \(\alpha\), then the
wrapped/wrapped extension of \(\zeta_{12}\) by \(\eta_2\).  The right
comparison \(\rho^{\mathrm{WLW}}\) records first the local/wrapped
extension of \(\alpha\) by \(\eta_2\), then the wrapped/wrapped
extension of \(\eta_1\) by \(\zeta_{23}\).  Both descriptions have
terminal wrapped colour \(\zeta_{123}\).

The mixed and wrapped/wrapped base-change isomorphisms identify the
two iterated compact-support pushforwards on the common flag stack.
The corresponding projection formulas move the ordered source
coefficients through the same pushforward, and the reduced
Thom--Sebastiani transports identify the two coefficient systems.
Therefore the two composites agree, and
\(\mathfrak o^{\mathrm{assoc}}_{\mathrm{WLW},R}=0\).

The proof is wordwise and conditional on the wrapped/wrapped
functorial coefficient rows.  It proves no associativity statement for
\(WWL\) or \(WWW\), and it does not prove the four-input pentagon,
quotient descent, transition compatibility, or aggregate
hybrid-carrier population.
\end{proof}

The packet
\[
\texttt{certificates/hybrid/wlw\_associativity}
\]
verifies
\(\texttt{WLW\_ASSOCIATIVITY\_CONDITIONAL\_VERIFIED}\):
Proposition~\ref{prop:wlw-associativity} proves the \(WLW\) word
associativity defect vanishes under the stated \(WL/LW/WW\)
flag-stack, mixed functorial rows, and wrapped/wrapped functorial
hypotheses.  It does not prove the two remaining word associativity
rows or the four-input pentagon.

\begin{proposition}[Associativity for the \(WWL\) word]
\label{prop:wwl-associativity}
\textnormal{[proved conditionally for the wrapped/wrapped--mixed flag stack]}
Let \(\alpha\) be a retained local colour and let \(\eta_1,\eta_2\)
be retained wrapped colours at height \(R\), ordered as
\((\eta_1,\eta_2,\alpha)\).  Write
\[
\zeta_{12}=\eta_1+\eta_2,\qquad
\zeta_{23}=\eta_2+\alpha,\qquad
\zeta_{123}=\eta_1+\eta_2+\alpha .
\]
Assume the \(WWL\) flag stack
\(\mathfrak F^{(2),\mathrm{WWL}}_{\eta_1,\eta_2,\alpha,R}\) of
Proposition~\ref{prop:eight-two-step-flag-stacks}, the \(WW\) and
\(WL\) target-admissible correspondence rows, the chosen
compact-support pushforward model, and the required mixed and
wrapped/wrapped base-change, projection-formula, and reduced
Thom--Sebastiani coefficient identifications on the \(WWL\) flag
square.  Then
\[
m^{\mathrm{WL}}_{\zeta_{12},\alpha,R}
\circ
\bigl(m^{\mathrm{WW}}_{\eta_1,\eta_2,R}\otimes 1\bigr)
=
m^{\mathrm{WW}}_{\eta_1,\zeta_{23},R}
\circ
\bigl(1\otimes m^{\mathrm{WL}}_{\eta_2,\alpha,R}\bigr)
\]
as morphisms on the retained compactly supported reduced
vanishing-cycle complexes of the ordered inputs.  Equivalently,
\[
\mathfrak o^{\mathrm{assoc}}_{\mathrm{WWL},R}=0 .
\]
The wrapped/wrapped base-change, projection-formula, and reduced
Thom--Sebastiani rows are hypotheses of this proposition; they are not
supplied by the wrapped/wrapped target-admissibility theorem.
\end{proposition}

\begin{proof}
The common refinement is the retained \(WWL\) flag stack
\[
0\subset C_1\subset C_2\subset C_3,
\qquad
C_1/C_0\in W_{\eta_1},\quad
C_2/C_1\in W_{\eta_2},\quad
C_3/C_2\in L_\alpha .
\]
The left comparison \(\lambda^{\mathrm{WWL}}\) records first the
wrapped/wrapped extension of \(\eta_1\) by \(\eta_2\), then the
wrapped/local extension of \(\zeta_{12}\) by \(\alpha\).  The right
comparison \(\rho^{\mathrm{WWL}}\) records first the wrapped/local
extension of \(\eta_2\) by \(\alpha\), then the wrapped/wrapped
extension of \(\eta_1\) by \(\zeta_{23}\).  Both descriptions have
terminal wrapped colour \(\zeta_{123}\).

The wrapped/wrapped and mixed base-change isomorphisms identify the
two iterated compact-support pushforwards on the common flag stack.
The corresponding projection formulas move the ordered source
coefficients through the same pushforward, and the reduced
Thom--Sebastiani transports identify the two coefficient systems.
Therefore the two composites agree, and
\(\mathfrak o^{\mathrm{assoc}}_{\mathrm{WWL},R}=0\).

The proof is wordwise and conditional on the wrapped/wrapped
functorial coefficient rows.  It proves no associativity statement for
\(WWW\), and it does not prove the four-input pentagon, quotient
descent, transition compatibility, or aggregate hybrid-carrier
population.
\end{proof}

The packet
\[
\texttt{certificates/hybrid/wwl\_associativity}
\]
verifies
\(\texttt{WWL\_ASSOCIATIVITY\_CONDITIONAL\_VERIFIED}\):
Proposition~\ref{prop:wwl-associativity} proves the \(WWL\) word
associativity defect vanishes under the stated \(WW/WL\) flag-stack,
mixed functorial rows, and wrapped/wrapped functorial hypotheses.  It
does not prove the remaining \(WWW\) word associativity row or the
four-input pentagon.

\begin{proposition}[Associativity for the \(WWW\) word]
\label{prop:www-associativity}
\textnormal{[proved conditionally for the wrapped/wrapped flag stack]}
Let \(\eta_1,\eta_2,\eta_3\) be retained wrapped colours at height
\(R\), ordered as \((\eta_1,\eta_2,\eta_3)\).  Write
\[
\zeta_{12}=\eta_1+\eta_2,\qquad
\zeta_{23}=\eta_2+\eta_3,\qquad
\zeta_{123}=\eta_1+\eta_2+\eta_3 .
\]
Assume the \(WWW\) flag stack
\(\mathfrak F^{(2),\mathrm{WWW}}_{\eta_1,\eta_2,\eta_3,R}\) of
Proposition~\ref{prop:eight-two-step-flag-stacks}, the \(WW\)
target-admissible correspondence rows, the chosen compact-support
pushforward model, and the required wrapped/wrapped base-change,
projection-formula, and reduced Thom--Sebastiani coefficient
identifications on the \(WWW\) flag square.  Then
\[
m^{\mathrm{WW}}_{\zeta_{12},\eta_3,R}
\circ
\bigl(m^{\mathrm{WW}}_{\eta_1,\eta_2,R}\otimes 1\bigr)
=
m^{\mathrm{WW}}_{\eta_1,\zeta_{23},R}
\circ
\bigl(1\otimes m^{\mathrm{WW}}_{\eta_2,\eta_3,R}\bigr)
\]
as morphisms on the retained compactly supported reduced
vanishing-cycle complexes of the ordered wrapped inputs.  Equivalently,
\[
\mathfrak o^{\mathrm{assoc}}_{\mathrm{WWW},R}=0 .
\]
The wrapped/wrapped base-change, projection-formula, and reduced
Thom--Sebastiani rows are hypotheses of this proposition; they are not
supplied by the wrapped/wrapped target-admissibility theorem.
\end{proposition}

\begin{proof}
The common refinement is the retained \(WWW\) flag stack
\[
0\subset C_1\subset C_2\subset C_3,
\qquad
C_1/C_0\in W_{\eta_1},\quad
C_2/C_1\in W_{\eta_2},\quad
C_3/C_2\in W_{\eta_3}.
\]
The left comparison \(\lambda^{\mathrm{WWW}}\) records first the
wrapped/wrapped extension of \(\eta_1\) by \(\eta_2\), then the
wrapped/wrapped extension of \(\zeta_{12}\) by \(\eta_3\).  The right
comparison \(\rho^{\mathrm{WWW}}\) records first the wrapped/wrapped
extension of \(\eta_2\) by \(\eta_3\), then the wrapped/wrapped
extension of \(\eta_1\) by \(\zeta_{23}\).  Both descriptions have
terminal wrapped colour \(\zeta_{123}\).

The wrapped/wrapped base-change isomorphisms identify the two iterated
compact-support pushforwards on the common flag stack.  The
wrapped/wrapped projection formulas move the ordered source
coefficients through the same pushforward, and the wrapped/wrapped
reduced Thom--Sebastiani transports identify the two coefficient
systems.  Therefore the two composites agree, and
\(\mathfrak o^{\mathrm{assoc}}_{\mathrm{WWW},R}=0\).

The proof is wordwise and conditional on the wrapped/wrapped
functorial coefficient rows.  It proves the last length-three
wordwise associativity defect, but it does not prove the four-input
pentagon, quotient descent, transition compatibility, or aggregate
hybrid-carrier population.
\end{proof}

The packet
\[
\texttt{certificates/hybrid/www\_associativity}
\]
verifies
\(\texttt{WWW\_ASSOCIATIVITY\_CONDITIONAL\_VERIFIED}\):
Proposition~\ref{prop:www-associativity} proves the \(WWW\) word
associativity defect vanishes under the stated \(WW/WW\) flag-stack
and wrapped/wrapped functorial hypotheses.  It completes the
length-three wordwise associativity list, but it does not prove the
four-input pentagon.

\begin{proposition}[Four-input pentagon coherence]
\label{prop:four-input-pentagon-coherence}
\textnormal{[proved conditionally before the reduced \(E\)-quotient]}
Let
\[
u=\epsilon_1\epsilon_2\epsilon_3\epsilon_4\in\{\mathrm L,\mathrm W\}^4
\]
and let \((\xi_1,\xi_2,\xi_3,\xi_4)\) be retained ordered colours of
these types at height \(R\).  Assume the binary correspondence rows
for \(LL,LW,WL,WW\), the compact-support pushforward model, the
length-three associativity rows of
Propositions~\ref{prop:lll-associativity}--\ref{prop:www-associativity},
and the required local, mixed, and wrapped/wrapped base-change,
projection-formula, and reduced Thom--Sebastiani coefficient
identifications.  Assume also the BBDJS--Joyce--Upmeier
Thom--Sebastiani orientation pentagon on the retained four-step
d-critical charts.  Then there is a retained finite-type flag-Quot
stack
\[
\mathfrak F^{(3),u}_{\xi_1,\xi_2,\xi_3,\xi_4,R}
\]
classifying filtrations
\[
0=C_0\subset C_1\subset C_2\subset C_3\subset C_4
\]
with \(C_i/C_{i-1}\) of type \(\epsilon_i\) and colour \(\xi_i\),
all intermediate colours retained, and wrapped rigidifications and
anchors remembered before the \(E\)-quotient.  The five planar binary
parenthesizations of the four inputs give five iterated
correspondence fibre products receiving canonical maps from
\(\mathfrak F^{(3),u}_{\xi_1,\xi_2,\xi_3,\xi_4,R}\).  The two
composites around the Mac Lane pentagon of associators agree:
\[
a_{\xi_1,\xi_2,\xi_3+\xi_4,R}\,
a_{\xi_1+\xi_2,\xi_3,\xi_4,R}
=
(1\otimes a_{\xi_2,\xi_3,\xi_4,R})\,
a_{\xi_1,\xi_2+\xi_3,\xi_4,R}\,
(a_{\xi_1,\xi_2,\xi_3,R}\otimes1).
\]
Equivalently, for every \(u\in\{\mathrm L,\mathrm W\}^4\),
\[
\mathfrak o^{\mathrm{pent}}_{u,R}=0 .
\]
\end{proposition}

\begin{proof}
Fix \(u\) and the ordered colours.  Over the retained terminal object
stack of colour \(\xi_1+\xi_2+\xi_3+\xi_4\), take the relative
three-step flag-Quot stack parametrising the flag
\[
0=C_0\subset C_1\subset C_2\subset C_3\subset C_4
\]
with the prescribed four graded pieces.  The retained closed endpoint
conditions, universal complexes, and extension-closure datum cut out
a finite-type retained substack.  Every intermediate object
\(C_j/C_i\) carries the type forced by the product rule
\(\mathrm L\cdot\mathrm L=\mathrm L\) and otherwise
\(\mathrm W\); if the type is \(\mathrm W\), the wrapped
rigidification and anchor memory are part of the retained datum.  This
constructs \(\mathfrak F^{(3),u}_{\xi_1,\xi_2,\xi_3,\xi_4,R}\)
before quotienting by \(E\).

Each planar binary tree on four leaves forgets the same flag in a
different order and therefore gives a map from the common four-step
flag stack to the corresponding iterated binary correspondence fibre
product.  The associator attached to an internal triple is the
two-step flag comparison of
Propositions~\ref{prop:lll-associativity}--\ref{prop:www-associativity}
pulled back along the appropriate contraction of the four-step flag.
Thus both paths around the pentagon are represented on the same
source stack
\(\mathfrak F^{(3),u}_{\xi_1,\xi_2,\xi_3,\xi_4,R}\), with the same
terminal target colour.

Proper base change for the retained compact-support model identifies
the target pushforwards after pullback to the common four-step flag
stack.  The projection formula moves the four ordered source
coefficients through the same pushforward.  The reduced
Thom--Sebastiani coefficient systems on the two paths are the two
parenthesizations of the same fourfold external product.  The
BBDJS Thom--Sebastiani pentagon
\cite[Theorem~5.10(iv)]{BBDJS2015}, together with the
Joyce--Upmeier multiplicative orientation cocycle
\cite[\S4]{JoyceUpmeier}, identifies these two coefficient systems.
Hence the two pull-push functors around the pentagon agree, and
\(\mathfrak o^{\mathrm{pent}}_{u,R}=0\).

The proof is before the reduced \(E\)-quotient and at fixed height
\(R\).  It proves no quotient descent, transition compatibility,
unit compatibility, higher-coloured tree category, or aggregate
hybrid-carrier population.  Whenever a \(WW\) binary correspondence
appears in the four-input word, the wrapped/wrapped base-change,
projection-formula, and reduced Thom--Sebastiani rows remain
hypotheses, as in
Propositions~\ref{prop:lww-associativity}--\ref{prop:www-associativity}.
\end{proof}

The packet
\[
\texttt{certificates/hybrid/four\_input\_pentagon\_coherence}
\]
verifies
\(\texttt{FOUR\_INPUT\_PENTAGON\_CONDITIONAL\_VERIFIED}\):
Proposition~\ref{prop:four-input-pentagon-coherence} constructs the
retained four-step flag stack for each word in
\(\{\mathrm L,\mathrm W\}^4\) and proves the associator pentagon on
that common refinement under the stated functorial hypotheses.  It
does not define the higher-coloured tree category or prove quotient
descent, transition compatibility, or aggregate population.

\begin{definition}[Higher-coloured hybrid tree category]
\label{def:higher-coloured-hybrid-tree-category}
Fix \(R\).  The nonunital higher-coloured hybrid tree category
\[
\mathsf{Tree}^{E,\mathrm{hyb}}_R
\]
is the free coloured category generated by the retained local and
wrapped profiles below.

An object is a finite ordered hybrid profile
\[
\underline\xi=((\epsilon_i,\xi_i,\mathfrak s_i))_{i=1}^n,
\qquad \epsilon_i\in\{\mathrm L,\mathrm W\},
\]
where \(\xi_i\) is a retained colour of the indicated type.  If
\(\epsilon_i=\mathrm L\), the symbol \(\mathfrak s_i\) is a retained
closed local configuration chart, including its finite indexing set
and diagonal-refinement datum.  If \(\epsilon_i=\mathrm W\), the
symbol \(\mathfrak s_i\) is the retained wrapped rigidification with
source anchor memory before the \(E\)-quotient.  The output colour is
\[
|\underline\xi|=\sum_i\xi_i,
\]
with output type \(\mathrm L\) only for an all-local profile and
\(\mathrm W\) otherwise.

A generating morphism
\[
T:\underline\xi\longrightarrow |\underline\xi|
\]
is a planar rooted tree with leaves ordered by \(\underline\xi\).
Every internal edge carries the retained intermediate colour given by
the sum of the leaves below that edge and the type rule
\(\mathrm L\cdot\mathrm L=\mathrm L\), otherwise \(\mathrm W\).  Every
internal vertex of arity \(k\ge2\) is decorated by the retained
extension correspondence or flag-Quot stack for the ordered subprofile
below that vertex.  Local vertices remember closed-configuration
refinements and symmetric relabellings; mixed and wrapped vertices
remember source, intermediate, and target anchors before quotient.
Composition is grafting of planar rooted trees followed by contraction
of unary identity edges.

The defect of realizing this category by compactly supported reduced
vanishing-cycle pull-push functors is the higher-coloured residual
\[
\mathfrak o^{\mathrm{col}}_R
=\bigl(
\mathfrak o^{\mathrm{tree}}_R,\,
\mathfrak o^{\mathrm{unit}}_R,\,
\mathfrak o^{\mathrm{sym}}_R,\,
\mathfrak o^{\mathrm{ref}}_R,\,
\mathfrak o^{\mathrm{des}}_R,\,
\mathfrak o^{\mathrm{ov}}_R
\bigr).
\]
The first component measures compatibility of tree contractions beyond
the four-input pentagon.  The remaining components measure,
respectively, vacuum units, symmetric relabelling of local
configurations, refinement of closed local charts, descent from
ordered charts, and overlap compatibility of retained local and
wrapped charts.  This definition supplies the category and the names
of its residual components; it does not assert
\(\mathfrak o^{\mathrm{col}}_R=0\).
\end{definition}

The packet
\[
\texttt{certificates/hybrid/higher\_coloured\_tree\_category\_definition}
\]
verifies
\(\texttt{HIGHER\_COLOURED\_TREE\_CATEGORY\_DEFINED}\):
Definition~\ref{def:higher-coloured-hybrid-tree-category} defines the
nonunital higher-coloured tree category and the six components of
\(\mathfrak o^{\mathrm{col}}_R\).  It does not prove the residual
vanishes, and it does not supply unit compatibility, symmetric descent,
wrapped order conventions, quotient descent, transition compatibility,
or aggregate population.

\begin{proposition}[Higher-coloured residual vanishing criterion]
\label{prop:higher-coloured-residual-vanishing-criterion}
\textnormal{[conditional finite-stage criterion]}
At fixed height \(R\), assume
Definition~\ref{def:higher-coloured-hybrid-tree-category} and
Proposition~\ref{prop:four-input-pentagon-coherence}.  Suppose the
following finite rows are supplied in the retained
\(E\)-equivariant correspondence category before quotient:
\begin{enumerate}
\item vacuum object and unit correspondences with left and right unit
defect zero for local and wrapped inputs;
\item symmetric relabelling descent for projection-finite local
configuration charts;
\item refinement compatibility for closed local configuration charts;
\item descent from ordered local charts to the symmetric finite-colour
colimit;
\item wrapped insertion order and braid-convention compatibility for
ordered wrapped profiles;
\item overlap compatibility for retained local and wrapped charts on
common refinements.
\end{enumerate}
Then the higher-coloured residual of
Definition~\ref{def:higher-coloured-hybrid-tree-category} vanishes:
\[
\mathfrak o^{\mathrm{col}}_R=0 .
\]
Equivalently, all six components
\[
\mathfrak o^{\mathrm{tree}}_R,\quad
\mathfrak o^{\mathrm{unit}}_R,\quad
\mathfrak o^{\mathrm{sym}}_R,\quad
\mathfrak o^{\mathrm{ref}}_R,\quad
\mathfrak o^{\mathrm{des}}_R,\quad
\mathfrak o^{\mathrm{ov}}_R
\]
have defect rank zero.
\end{proposition}

\begin{proof}
The tree-contraction component is the coherence problem for binary
parenthesizations in
\(\mathsf{Tree}^{E,\mathrm{hyb}}_R\).  The length-three associators
are the rows of
Propositions~\ref{prop:lll-associativity}--\ref{prop:www-associativity},
and their four-input pentagon is
Proposition~\ref{prop:four-input-pentagon-coherence}.  Hence the
standard coherence theorem for parenthesized products identifies any
two binary contraction paths with the same planar leaves by a composite
of four-input pentagons.  Pulling those pentagons to the retained
common flag-Quot refinement gives
\(\mathfrak o^{\mathrm{tree}}_R=0\).

The remaining five components are not consequences of the pentagon.
By the supplied rows in the statement, the left and right vacuum
correspondences act as units on local and wrapped inputs, so
\(\mathfrak o^{\mathrm{unit}}_R=0\).  Symmetric relabelling descent
sets \(\mathfrak o^{\mathrm{sym}}_R=0\).  Refinement compatibility of
closed local charts sets \(\mathfrak o^{\mathrm{ref}}_R=0\).  Descent
from ordered charts to the symmetric finite-colour colimit sets
\(\mathfrak o^{\mathrm{des}}_R=0\).  Overlap compatibility on common
local/wrapped refinements sets \(\mathfrak o^{\mathrm{ov}}_R=0\).
The residual is the direct tuple of these six components; therefore
\(\mathfrak o^{\mathrm{col}}_R=0\).

This criterion does not construct the non-tree input rows.  In the
present correction ledger the unit, symmetric-relabelling, and wrapped
order rows are supplied separately; refinement, ordered-chart descent,
and overlap rows remain separate obligations needed for aggregate
population.
\end{proof}

The packet
\[
\texttt{certificates/hybrid/higher\_coloured\_residual\_vanishing\_criterion}
\]
verifies
\(\texttt{HIGHER\_COLOURED\_RESIDUAL\_CONDITIONAL\_CRITERION\_VERIFIED}\):
Proposition~\ref{prop:higher-coloured-residual-vanishing-criterion}
proves the finite-stage implication from the six named component rows
to \(\mathfrak o^{\mathrm{col}}_R=0\).  It proves the tree-contraction
component from the associativity rows and four-input pentagon, but it
does not supply the refinement, ordered-chart descent, or overlap rows
required to invoke the criterion unconditionally.  The unit,
symmetric-relabelling, and wrapped order rows are supplied by separate
packets.

\begin{definition}[Hybrid unit object and vacuum correspondences]
\label{def:hybrid-unit-object-vacuum-correspondences}
\textnormal{[definition only]}
Fix a finite HN height \(R\).  A hybrid unit datum consists of the
following retained rows before the reduced \(E\)-quotient.

First, the zero object of the retained Hall heart determines a local
zero colour
\[
0_R\in\Gamma_R^{\mathrm{loc}},\qquad b_R^{\mathrm{geom}}(0_R)=0,
\]
with empty local support chart \(\varnothing\).  The corresponding
unit stack is the one-point closed local incidence stack
\[
\mathbf 1^{\mathrm{hyb}}_R
:=
\mathfrak M^{\mathrm{loc,cl}}_{0_R,R,\varnothing}.
\]
The same symbol \(0_R:\varnothing\to\Gamma_R^{\mathrm{loc}}\) denotes
the empty local colour profile whose total colour is \(0_R\).
There is no wrapped zero colour: the unit is local because the wrapped
sector is defined by \(b_R^{\mathrm{geom}}>0\).

Second, for every finite local colour map
\(\alpha:I\to\Gamma_R^{\mathrm{loc}}\), the datum contains the split
local unit substacks
\[
\mathfrak U^{\mathrm{LL},\ell}_{0_R,\alpha;\alpha,R,I}
\subset
\mathfrak E^{\mathrm{LL}}_{0_R,\alpha;\alpha,R,\varnothing,I,I},
\qquad
\mathfrak U^{\mathrm{LL},r}_{\alpha,0_R;\alpha,R,I}
\subset
\mathfrak E^{\mathrm{LL}}_{\alpha,0_R;\alpha,R,I,\varnothing,I}.
\]
They classify, respectively, the split exact sequences
\[
0\longrightarrow A_{\alpha,I}\longrightarrow A_{\alpha,I}
\longrightarrow 0\longrightarrow0,
\qquad
0\longrightarrow0\longrightarrow A_{\alpha,I}
\longrightarrow A_{\alpha,I}\longrightarrow0 .
\]
The source maps remember the ordered pair
\((\mathbf 1^{\mathrm{hyb}}_R,A_{\alpha,I})\) or
\((A_{\alpha,I},\mathbf 1^{\mathrm{hyb}}_R)\), and the target maps
send the split sequence to \(A_{\alpha,I}\).

Third, for every wrapped colour \(\eta\in\Gamma_R^{\mathrm{wr}}\), the
datum contains the split mixed unit substacks
\[
\mathfrak U^{\mathrm{LW},\ell}_{0_R,\eta;\eta,R}
\subset
\mathfrak E^{\mathrm{LW}}_{0_R,\eta;\eta,R,\varnothing},
\qquad
\mathfrak U^{\mathrm{WL},r}_{\eta,0_R;\eta,R}
\subset
\mathfrak E^{\mathrm{WL}}_{\eta,0_R;\eta,R,\varnothing}.
\]
They classify the split exact sequences
\[
0\longrightarrow W_\eta\longrightarrow W_\eta
\longrightarrow0\longrightarrow0,
\qquad
0\longrightarrow0\longrightarrow W_\eta
\longrightarrow W_\eta\longrightarrow0 .
\]
The wrapped rigidification and source/target anchor memory are those
of \(W_\eta\); the zero local factor carries the empty support chart
and no anchor.  These rows are the unit correspondences used by
\(\mathfrak o^{\mathrm{unit}}_R\).  This definition does not prove
that the induced compact-support pull-push maps are the identity on
local or wrapped coefficients; those are the local and wrapped unit
compatibility rows.
\end{definition}

The packet
\[
\texttt{certificates/hybrid/hybrid\_unit\_object\_and\_vacuum\_correspondences}
\]
verifies
\(\texttt{HYBRID\_UNIT\_OBJECT\_AND\_VACUUM\_CORRESPONDENCES\_DEFINED}\):
Definition~\ref{def:hybrid-unit-object-vacuum-correspondences} defines
the local zero unit object and the four split vacuum correspondence
types \(LL_\ell\), \(LL_r\), \(LW_\ell\), and \(WL_r\).  It does not
prove the local or wrapped unit identity laws, quotient descent,
transition compatibility, or aggregate population.

\begin{proposition}[Local unit compatibility for the hybrid product]
\label{prop:local-unit-compatibility}
\textnormal{[proved before quotient at finite height]}
Assume the hybrid unit datum of
Definition~\ref{def:hybrid-unit-object-vacuum-correspondences}, the
local/local correspondence datum of
Definition~\ref{def:local-local-correspondence-datum}, the
target-properness row of
Proposition~\ref{prop:local-local-target-properness}, and the
compact-support model of
Definition~\ref{def:compact-support-exceptional-pushforward-model}.
Let
\[
K^{\mathrm{loc}}_{\alpha,R,I}
:=
\Phi^{\mathrm{loc}}_{\alpha,R,I}\otimes
o^{\mathrm{loc}}_{\alpha,R,I}
\]
be the reduced local coefficient on
\(\mathfrak M^{\mathrm{loc,cl}}_{\alpha,R,I}\), and fix the standard
zero-object normalization
\[
K^{\mathrm{loc}}_{0_R,R,\varnothing}
\simeq
\C_{\mathbf 1^{\mathrm{hyb}}_R}.
\]
Let \(p^{\mathrm{LL},\ell},q^{\mathrm{LL},\ell}\) and
\(p^{\mathrm{LL},r},q^{\mathrm{LL},r}\) denote the restrictions of
the \(LL\) source and target maps to
\(\mathfrak U^{\mathrm{LL},\ell}_{0_R,\alpha;\alpha,R,I}\) and
\(\mathfrak U^{\mathrm{LL},r}_{\alpha,0_R;\alpha,R,I}\).
Then the compact-support pull-push maps attached to the split local
unit correspondences satisfy
\[
\mu^{\mathrm{LL},\ell}_{0_R,\alpha;R,I}
=
(q^{\mathrm{LL},\ell})^{\mathrm{cs}}_!\,
\operatorname{TS}^{\mathrm{red}}_{0_R,\alpha}\,
(p^{\mathrm{LL},\ell})^*
=
\operatorname{id}_{K^{\mathrm{loc}}_{\alpha,R,I}},
\]
and
\[
\mu^{\mathrm{LL},r}_{\alpha,0_R;R,I}
=
(q^{\mathrm{LL},r})^{\mathrm{cs}}_!\,
\operatorname{TS}^{\mathrm{red}}_{\alpha,0_R}\,
(p^{\mathrm{LL},r})^*
=
\operatorname{id}_{K^{\mathrm{loc}}_{\alpha,R,I}}.
\]
Equivalently, the local unit defect component has defect rank zero:
\[
\mathfrak o^{\mathrm{unit,loc}}_R=0 .
\]
\end{proposition}

\begin{proof}
For the left unit row, the split substack
\[
\mathfrak U^{\mathrm{LL},\ell}_{0_R,\alpha;\alpha,R,I}
\]
is isomorphic to
\(\mathfrak M^{\mathrm{loc,cl}}_{\alpha,R,I}\).  The isomorphism sends
\(A_{\alpha,I}\) to the split sequence
\[
0\longrightarrow A_{\alpha,I}\longrightarrow A_{\alpha,I}
\longrightarrow0\longrightarrow0,
\]
and its inverse takes the target object of the sequence.  Under this
identification the source map is the closed insertion
\[
A_{\alpha,I}\longmapsto
(\mathbf 1^{\mathrm{hyb}}_R,A_{\alpha,I}),
\]
and the target map is the identity on
\(\mathfrak M^{\mathrm{loc,cl}}_{\alpha,R,I}\).  Since the target map
is proper, the compact-support model uses the identity
compactification and \((q^{\mathrm{LL},\ell})^{\mathrm{cs}}_!\) is
ordinary pushforward along the identity.

The pullback of the ordered source coefficient is
\[
(p^{\mathrm{LL},\ell})^*
\bigl(K^{\mathrm{loc}}_{0_R,R,\varnothing}
\boxtimes K^{\mathrm{loc}}_{\alpha,R,I}\bigr)
\simeq
K^{\mathrm{loc}}_{\alpha,R,I}.
\]
The zero-object normalization identifies the first factor with the
constant unit complex on the point.  The BBDJS Thom--Sebastiani
isomorphism for the sum of a zero potential and the local potential
\cite[Theorem~5.10]{BBDJS2015}, tensored with the Joyce--Upmeier
orientation unit \cite[\S4]{JoyceUpmeier}, is the identity on this
coefficient.  Hence
\[
(q^{\mathrm{LL},\ell})^{\mathrm{cs}}_!\,
\operatorname{TS}^{\mathrm{red}}_{0_R,\alpha}\,
(p^{\mathrm{LL},\ell})^*
=\operatorname{id}_{K^{\mathrm{loc}}_{\alpha,R,I}}.
\]

The right unit row is identical with the factors reversed.  The
substack
\(\mathfrak U^{\mathrm{LL},r}_{\alpha,0_R;\alpha,R,I}\) is the graph
of the split sequence
\[
0\longrightarrow0\longrightarrow A_{\alpha,I}
\longrightarrow A_{\alpha,I}\longrightarrow0,
\]
so its source map inserts
\((A_{\alpha,I},\mathbf 1^{\mathrm{hyb}}_R)\) and its target map is
again the identity.  The ordered Thom--Sebastiani isomorphism for the
local potential plus the zero potential is the right-unit form of the
same zero-factor identity.  Thus the right pull-push is
\(\operatorname{id}_{K^{\mathrm{loc}}_{\alpha,R,I}}\).

The proof is before the reduced \(E\)-quotient and at fixed height
\(R\).  It proves no wrapped unit identity, quotient descent,
transition compatibility, symmetric descent, refinement compatibility,
or aggregate hybrid-carrier population.
\end{proof}

The packet
\[
\texttt{certificates/hybrid/local\_unit\_compatibility}
\]
verifies
\(\texttt{LOCAL\_UNIT\_COMPATIBILITY\_VERIFIED}\):
Proposition~\ref{prop:local-unit-compatibility} proves the left and
right \(LL\) vacuum pull-push maps are the identity on local
coefficients.  It does not prove wrapped unit compatibility, quotient
descent, transition compatibility, or aggregate population.

\begin{proposition}[Wrapped unit compatibility for the hybrid product]
\label{prop:wrapped-unit-compatibility}
\textnormal{[proved before quotient at finite height]}
Assume the hybrid unit datum of
Definition~\ref{def:hybrid-unit-object-vacuum-correspondences}, the
one-sided mixed correspondence datum of
Definition~\ref{def:mixed-local-wrapped-correspondence-datum}, the
mixed target-admissibility row of
Proposition~\ref{prop:mixed-local-wrapped-target-admissibility}, the
compact-support model of
Definition~\ref{def:compact-support-exceptional-pushforward-model},
and the mixed Thom--Sebastiani transport row of
Proposition~\ref{prop:mixed-correspondence-thom-sebastiani}.  Let
\[
K^{\mathrm{wr}}_{\eta,R}
:=
\Phi^{\mathrm{wr}}_{\eta,R}\otimes o^{\mathrm{wr}}_{\eta,R}
\]
be the reduced wrapped coefficient on
\(\mathcal M^{\mathrm{wr,rig}}_{\eta,R}\), and keep the zero-object
normalization
\[
K^{\mathrm{loc}}_{0_R,R,\varnothing}
\simeq
\C_{\mathbf 1^{\mathrm{hyb}}_R}.
\]
Let \(p^{\mathrm{LW},\ell},q^{\mathrm{LW},\ell}\) and
\(p^{\mathrm{WL},r},q^{\mathrm{WL},r}\) denote the restrictions of the
mixed source and target maps to
\(\mathfrak U^{\mathrm{LW},\ell}_{0_R,\eta;\eta,R}\) and
\(\mathfrak U^{\mathrm{WL},r}_{\eta,0_R;\eta,R}\).  Then the
compact-support pull-push maps attached to the split wrapped unit
correspondences satisfy
\[
\mu^{\mathrm{LW},\ell}_{0_R,\eta;R}
=
(q^{\mathrm{LW},\ell})^{\mathrm{cs}}_!\,
\operatorname{TS}^{\mathrm{red},\mathrm{LW}}_{0_R,\eta}\,
(p^{\mathrm{LW},\ell})^*
=
\operatorname{id}_{K^{\mathrm{wr}}_{\eta,R}},
\]
and
\[
\mu^{\mathrm{WL},r}_{\eta,0_R;R}
=
(q^{\mathrm{WL},r})^{\mathrm{cs}}_!\,
\operatorname{TS}^{\mathrm{red},\mathrm{WL}}_{\eta,0_R}\,
(p^{\mathrm{WL},r})^*
=
\operatorname{id}_{K^{\mathrm{wr}}_{\eta,R}}.
\]
Equivalently, the wrapped unit defect component has defect rank zero:
\[
\mathfrak o^{\mathrm{unit,wr}}_R=0 .
\]
\end{proposition}

\begin{proof}
For the left wrapped unit row, the split substack
\[
\mathfrak U^{\mathrm{LW},\ell}_{0_R,\eta;\eta,R}
\]
is isomorphic to \(\mathcal M^{\mathrm{wr,rig}}_{\eta,R}\).  The
isomorphism sends \(W_\eta\) to the split sequence
\[
0\longrightarrow W_\eta\longrightarrow W_\eta
\longrightarrow0\longrightarrow0,
\]
and its inverse takes the target wrapped object.  Under this
identification the source map is the insertion
\[
W_\eta\longmapsto
(\mathbf 1^{\mathrm{hyb}}_R,W_\eta),
\]
and the target map is the identity on
\(\mathcal M^{\mathrm{wr,rig}}_{\eta,R}\).  The unit definition
records that the zero local factor carries no anchor; hence the source
and target anchor memory on the split row is exactly the anchor memory
of \(W_\eta\).  No \(E\)-quotient is taken.  Since the target map is
proper, the compact-support model uses the identity compactification
and \((q^{\mathrm{LW},\ell})^{\mathrm{cs}}_!\) is ordinary pushforward
along the identity.

The pullback of the ordered source coefficient is
\[
(p^{\mathrm{LW},\ell})^*
\bigl(K^{\mathrm{loc}}_{0_R,R,\varnothing}
\boxtimes K^{\mathrm{wr}}_{\eta,R}\bigr)
\simeq
K^{\mathrm{wr}}_{\eta,R}.
\]
The zero-object normalization identifies the first factor with the
constant unit complex on the point.  The mixed Thom--Sebastiani
transport of Proposition~\ref{prop:mixed-correspondence-thom-sebastiani}
specializes, for the zero local potential plus the wrapped potential,
to the BBDJS zero-factor Thom--Sebastiani identity
\cite[Theorem~5.10]{BBDJS2015}, tensored with the Joyce--Upmeier
orientation unit \cite[\S4]{JoyceUpmeier}.  Therefore
\[
(q^{\mathrm{LW},\ell})^{\mathrm{cs}}_!\,
\operatorname{TS}^{\mathrm{red},\mathrm{LW}}_{0_R,\eta}\,
(p^{\mathrm{LW},\ell})^*
=\operatorname{id}_{K^{\mathrm{wr}}_{\eta,R}}.
\]

The right wrapped unit row is the same argument with the factors
reversed.  The split substack
\(\mathfrak U^{\mathrm{WL},r}_{\eta,0_R;\eta,R}\) is the graph of
\[
0\longrightarrow0\longrightarrow W_\eta
\longrightarrow W_\eta\longrightarrow0,
\]
the source map inserts
\((W_\eta,\mathbf 1^{\mathrm{hyb}}_R)\), the target map is the
identity on \(\mathcal M^{\mathrm{wr,rig}}_{\eta,R}\), and the wrapped
anchor memory is unchanged.  The ordered \(WL\) Thom--Sebastiani
transport is the right zero-factor identity.  Hence the right
pull-push is \(\operatorname{id}_{K^{\mathrm{wr}}_{\eta,R}}\).

The proof is before the reduced \(E\)-quotient and at fixed height
\(R\).  It proves no quotient descent, transition compatibility,
symmetric descent, refinement compatibility, overlap compatibility, or
aggregate hybrid-carrier population.
\end{proof}

The packet
\[
\texttt{certificates/hybrid/wrapped\_unit\_compatibility}
\]
verifies
\(\texttt{WRAPPED\_UNIT\_COMPATIBILITY\_VERIFIED}\):
Proposition~\ref{prop:wrapped-unit-compatibility} proves the left
\(LW\) and right \(WL\) vacuum pull-push maps are the identity on
wrapped coefficients, with wrapped anchor memory unchanged before the
reduced \(E\)-quotient.  It does not prove quotient descent,
transition compatibility, or aggregate population.

\begin{proposition}[Symmetric descent for local configuration charts]
\label{prop:local-configuration-symmetric-descent}
\textnormal{[proved before quotient at finite height]}
Assume the local \(b=0\) carrier of
Definition~\ref{def:hybrid-local-stratum-prestack} and the
projection-finite closed local configuration sector of the finite
stage.  Let
\(\alpha:I\to\Gamma_R^{\mathrm{loc}}\) be a finite local colour map
and let \(\sigma:I\xrightarrow{\sim}I'\) be a bijection.  Define
\[
(\sigma_*\alpha)_{i'}=\alpha_{\sigma^{-1}(i')},
\qquad
\rho_\sigma:E^I\xrightarrow{\sim}E^{I'},
\qquad
(\rho_\sigma(\mathbf e))_{i'}=e_{\sigma^{-1}(i')}.
\]
Then
\[
(\operatorname{id}_E\times\rho_\sigma)(Z_I)=Z_{I'}
\]
and \(\rho_\sigma\) induces an isomorphism of closed local
configuration charts
\[
r_\sigma:
\mathfrak M^{\mathrm{loc,cl}}_{\alpha,R,I}
\xrightarrow{\sim}
\mathfrak M^{\mathrm{loc,cl}}_{\sigma_*\alpha,R,I'},
\qquad
(A,\mathbf e)\longmapsto (A,\rho_\sigma(\mathbf e)),
\]
with
\[
\pi_{\sigma_*\alpha,R,I'}\circ r_\sigma
=
\rho_\sigma\circ\pi_{\alpha,R,I}.
\]
For the reduced local coefficient
\[
K^{\mathrm{loc}}_{\alpha,R,I}
:=
\Phi^{\mathrm{loc}}_{\alpha,R,I}\otimes
o^{\mathrm{loc}}_{\alpha,R,I},
\]
there is a functorial relabelling isomorphism
\[
r_\sigma^*K^{\mathrm{loc}}_{\sigma_*\alpha,R,I'}
\simeq
K^{\mathrm{loc}}_{\alpha,R,I},
\]
and these isomorphisms satisfy the cocycle identity for composable
bijections.  Hence the stabilizer
\[
\mathfrak S_\alpha=\{\sigma\in\operatorname{Aut}(I)\mid
\sigma_*\alpha=\alpha\}
\]
acts on
\((\mathfrak M^{\mathrm{loc,cl}}_{\alpha,R,I},
K^{\mathrm{loc}}_{\alpha,R,I})\), and the coefficient descends
effectively to the quotient stack
\[
[\mathfrak M^{\mathrm{loc,cl}}_{\alpha,R,I}/\mathfrak S_\alpha].
\]
Equivalently, the fixed-chart symmetric relabelling component of the
higher-coloured residual has defect rank zero:
\[
\mathfrak o^{\mathrm{sym}}_R=0
\]
on projection-finite local configuration charts before quotient.
\end{proposition}

\begin{proof}
The incidence locus \(Z_I\subset E\times E^I\) is the union of the
graphs \(e=e_i\).  Under
\(\operatorname{id}_E\times\rho_\sigma\) the graph indexed by
\(i\in I\) is sent to the graph indexed by
\(\sigma(i)\in I'\); hence
\((\operatorname{id}_E\times\rho_\sigma)(Z_I)=Z_{I'}\).
Thus the condition
\[
p_E(\operatorname{Supp}A)\subset Z_I(\mathbf e)
\]
is equivalent to
\[
p_E(\operatorname{Supp}A)\subset Z_{I'}(\rho_\sigma(\mathbf e)).
\]
Relabelling the finite colour profile by \(\sigma\) changes no
underlying sheaf \(A\) and preserves the total retained local colour.
Therefore
\[
(A,\mathbf e)\longmapsto(A,\rho_\sigma(\mathbf e))
\]
defines the stated isomorphism \(r_\sigma\), with inverse
\(r_{\sigma^{-1}}\), and the displayed identity with the configuration
maps follows from the definition of \(\rho_\sigma\).

The local d-critical chart, its potential, and its cosection-reduced
vanishing-cycle complex are attached to the object \(A\) together with
the ordered support coordinates.  The map \(r_\sigma\) is an isomorphism
of these charts: it only relabels the support coordinates and the
matching colour labels.  Functoriality of vanishing cycles under
isomorphism, tensored with the transported Joyce--Upmeier orientation
line, gives
\[
r_\sigma^*
(\Phi^{\mathrm{loc}}_{\sigma_*\alpha,R,I'}
\otimes o^{\mathrm{loc}}_{\sigma_*\alpha,R,I'})
\simeq
\Phi^{\mathrm{loc}}_{\alpha,R,I}\otimes
o^{\mathrm{loc}}_{\alpha,R,I}.
\]
Since \(r_\tau r_\sigma=r_{\tau\sigma}\), these pullback
isomorphisms satisfy the cocycle identity.

For \(\sigma\in\mathfrak S_\alpha\) the target profile is again
\(\alpha\), so the preceding isomorphisms define a finite group action
on the stack and an equivariant structure on the coefficient complex.
Constructible complexes on a quotient stack by a finite group are the
same data as equivariant constructible complexes on the source stack;
therefore the descent to
\([\mathfrak M^{\mathrm{loc,cl}}_{\alpha,R,I}/\mathfrak S_\alpha]\)
is effective.  This proves the symmetric relabelling component for a
fixed projection-finite local chart.

The proof does not prove collision compatibility, refinement of closed
local charts, descent from ordered charts to the symmetric finite-colour
colimit, quotient descent, transition compatibility, or aggregate
hybrid-carrier population.
\end{proof}

The packet
\[
\texttt{certificates/hybrid/local\_configuration\_symmetric\_descent}
\]
verifies
\(\texttt{LOCAL\_CONFIGURATION\_SYMMETRIC\_DESCENT\_VERIFIED}\):
Proposition~\ref{prop:local-configuration-symmetric-descent} proves
finite symmetric-group descent for each projection-finite closed local
configuration chart and its reduced coefficient before the reduced
\(E\)-quotient.  It does not prove collision compatibility, refinement
compatibility, descent from ordered charts to the symmetric
finite-colour colimit, quotient descent, transition compatibility, or
aggregate population.

\begin{proposition}[Wrapped insertion order conventions]
\label{prop:wrapped-insertion-order-conventions}
\textnormal{[proved before quotient at finite height]}
Fix \(R\) and the ordered hybrid tree category of
Definition~\ref{def:higher-coloured-hybrid-tree-category}.  For an
ordered hybrid profile
\[
\underline\xi=((\epsilon_i,\xi_i,\mathfrak s_i))_{i=1}^n
\]
write
\[
\operatorname{Wr}(\underline\xi)
=
(\xi_{i_1},\ldots,\xi_{i_m}),
\qquad
i_1<\cdots<i_m,\qquad \epsilon_{i_a}=\mathrm W,
\]
for the ordered subsequence of wrapped leaves.  The wrapped insertion
convention is the following.
\begin{enumerate}
\item At every binary vertex the wrapped order of the output is the
concatenation of the wrapped orders of the ordered inputs.
\item For a \(WW\) vertex with ordered source pair
\((\eta_1,\eta_2)\), the Hall correspondence is
\[
0\longrightarrow W_{\eta_2}\longrightarrow B_{\eta_1+\eta_2}
\longrightarrow W_{\eta_1}\longrightarrow0,
\]
and the ordered wrapped source list is \((\eta_1,\eta_2)\): first the
quotient colour, then the subobject colour.  The target wrapped colour
remembers this source order as the anchor-memory order before the
reduced \(E\)-quotient.
\item For \(LW\), \(WL\), and \(LWL\) vertices, local factors do not
change the order of the wrapped subsequence.
\item The comparison maps for the eight two-step flag stacks of
Proposition~\ref{prop:eight-two-step-flag-stacks} preserve the ordered
wrapped subsequence \(\operatorname{Wr}\) for each of the words
\[
LLL,\ LLW,\ LWL,\ WLL,\ LWW,\ WLW,\ WWL,\ WWW.
\]
\end{enumerate}
Therefore parenthesization changes in
\(\mathsf{Tree}^{E,\mathrm{hyb}}_R\) use the planar leaf order and
introduce no braid operator and no permutation of wrapped inputs.  The
wrapped-order input row in
Proposition~\ref{prop:higher-coloured-residual-vanishing-criterion}
has defect rank zero at fixed height before quotient.
\end{proposition}

\begin{proof}
The object set of
\(\mathsf{Tree}^{E,\mathrm{hyb}}_R\) is ordered by definition, and a
generating morphism is a planar rooted tree whose leaves are ordered by
the input profile.  Thus every internal vertex sees an ordered
subprofile.  The rule for the output colour is additive, and the type
rule sends an output to the wrapped sector as soon as one input below
the vertex is wrapped.  This gives a canonical ordered wrapped
subsequence: scan the leaves below the vertex from left to right and
discard the local leaves.

For an \(LW\) or \(WL\) vertex there is exactly one wrapped input, so
the ordered wrapped subsequence is unchanged.  For an \(LWL\) vertex
the two local leaves are discarded by the scan, and the wrapped
subsequence is again the single middle wrapped colour.  For a \(WW\)
vertex the ordered source pair in
Definition~\ref{def:wrapped-wrapped-correspondence-datum} is
\((\eta_1,\eta_2)\), while the Hall extension convention is
\[
0\to W_{\eta_2}\to B_{\eta_1+\eta_2}\to W_{\eta_1}\to0.
\]
Hence the first source factor is the quotient and the second source
factor is the subobject.  This fixes the ordered source-anchor list as
\((\eta_1,\eta_2)\).  No transposition is applied when the target
wrapped colour is formed.

It remains to check parenthesization changes.  The two-step flag stack
for a word \(w=\epsilon_1\epsilon_2\epsilon_3\) records a filtration
whose successive quotients are the three ordered leaves.  The left
comparison map contracts the first two leaves and then the resulting
intermediate object with the third leaf.  The right comparison map
contracts the last two leaves and then the first leaf with the
resulting intermediate object.  In both cases the ordered wrapped
subsequence is obtained by scanning the same ordered list
\((\epsilon_1,\epsilon_2,\epsilon_3)\) and discarding local entries.
Thus the comparison maps preserve \(\operatorname{Wr}\) for all eight
words.

Consequently every contraction in the planar tree category preserves
the wrapped subsequence by induction on the number of vertices.  Since
the category has no braid generator and no symmetric identification of
wrapped leaves, a braid word is not a hidden morphism in the
realization.  The only braid compatible with the convention is the
identity braid, which acts by the identity on the ordered source
coefficient.  Hence the wrapped insertion order row has defect rank
zero.

The proof does not prove symmetric descent for unordered wrapped
pairs, closed-chart refinement, descent from ordered charts to the
symmetric finite-colour colimit, overlap compatibility, quotient
descent, transition compatibility, or aggregate hybrid-carrier
population.
\end{proof}

The packet
\[
\texttt{certificates/hybrid/wrapped\_insertion\_order\_conventions}
\]
verifies
\(\texttt{WRAPPED\_INSERTION\_ORDER\_CONVENTIONS\_VERIFIED}\):
Proposition~\ref{prop:wrapped-insertion-order-conventions} fixes the
planar wrapped-source order for mixed and wrapped vertices, checks the
\(WW\) quotient/subobject convention, and proves that the eight
two-step flag comparisons preserve the same ordered wrapped subsequence
before the reduced \(E\)-quotient.  It does not prove unordered
wrapped-pair symmetric descent, ordered-chart descent, overlap
compatibility, quotient descent, transition compatibility, or aggregate
population.

\begin{definition}[Ordered two-sided mixed correspondence datum]
\label{def:two-sided-mixed-correspondence-datum}
Fix a finite HN height \(R\).  Let
\(\alpha_-:I_-\to\Gamma_R^{\mathrm{loc}}\) and
\(\beta_+:I_+\to\Gamma_R^{\mathrm{loc}}\) be finite local colour maps,
and let \(\eta\in\Gamma_R^{\mathrm{wr}}\).  Write
\(|\alpha_-|=\sum_{i\in I_-}\alpha_{-,i}\) and
\(|\beta_+|=\sum_{j\in I_+}\beta_{+,j}\).  A two-sided mixed
correspondence row is defined only for a retained wrapped target colour
\(\zeta\in\Gamma_R^{\mathrm{wr}}\) satisfying
\[
\zeta=|\alpha_-|+\eta+|\beta_+|.
\]
The ordered local--wrapped--local correspondence is the unreduced
diagram
\[
\mathfrak E^{\mathrm{LWL}}_{\alpha_-;\eta;\beta_+,R,I_-,I_+}
\xrightarrow{p^{\mathrm{LWL}}}
\mathfrak M^{\mathrm{loc,cl}}_{\alpha_-,R,I_-}
\times
\mathcal M_{\eta,R}^{\mathrm{wr,rig}}
\times
\mathfrak M^{\mathrm{loc,cl}}_{\beta_+,R,I_+},
\]
\[
\mathfrak E^{\mathrm{LWL}}_{\alpha_-;\eta;\beta_+,R,I_-,I_+}
\xrightarrow{q^{\mathrm{LWL}}}
\mathcal M_{\zeta,R}^{\mathrm{wr,rig}},
\]
classifying retained flags
\[
0\subset A_{\beta_+}\subset B_{\eta+\beta_+}\subset C_\zeta
\]
with associated graded pieces
\[
A_{\beta_+},\qquad W_\eta,\qquad A_{\alpha_-}
\]
from bottom to top.  Equivalently,
\[
0\to A_{\beta_+}\to B_{\eta+\beta_+}\to W_\eta\to0,
\qquad
0\to B_{\eta+\beta_+}\to C_\zeta\to A_{\alpha_-}\to0.
\]
The diagram lives before the reduced \(E\)-quotient.  It is a
three-input correspondence type, not the full eight-word two-step flag
atlas and not an associativity theorem.  Source and target anchor
memory, admissibility, base change, projection formula,
Thom--Sebastiani transport, and transition compatibility are separate
rows.
\end{definition}

The packet
\[
\texttt{certificates/hybrid/two\_sided\_mixed\_correspondence\_definition}
\]
verifies
\(\texttt{TWO\_SIDED\_MIXED\_CORRESPONDENCE\_DEFINITION\_VERIFIED}\):
the ordered \(LWL\) two-sided mixed correspondence type is defined,
while flag-stack rows, source/target-map rows, populated anchor-memory
rows, admissibility rows, base-change rows, projection-formula rows,
Thom--Sebastiani rows, and transition rows are recorded as missing open
obligations.

\begin{definition}[Source and target anchors before quotient]
\label{def:source-target-anchor-memory-datum}
Fix a finite HN height \(R\).  Assume the wrapped prequotient sector
has the prequotient stacks of
Definition~\ref{def:e-equivariant-wrapped-prequotient-stack} and
has the abstract anchors
\[
\lambda_{\eta,R}:\mathcal M_{\eta,R}^{\mathrm{wr,rig}}\longrightarrow E
\]
of Definition~\ref{def:hybrid-wrapped-factorization-datum}.  For every
unreduced correspondence with a wrapped input or wrapped target, the
source/target anchor memory is part of the correspondence datum before
the reduced \(E\)-quotient is applied.

For the one-sided mixed rows of
Definition~\ref{def:mixed-local-wrapped-correspondence-datum}, it is
\[
a^{\mathrm{LW}}_{\alpha,\eta;\zeta,R,I}:
\mathfrak E^{\mathrm{LW}}_{\alpha,\eta;\zeta,R,I}\longrightarrow E^2,
\qquad
\xi\longmapsto
\bigl(\lambda_{\eta,R}(W_\eta),
\lambda_{\zeta,R}(B_\zeta)\bigr),
\]
and
\[
a^{\mathrm{WL}}_{\eta,\alpha;\zeta,R,I}:
\mathfrak E^{\mathrm{WL}}_{\eta,\alpha;\zeta,R,I}\longrightarrow E^2,
\qquad
\xi\longmapsto
\bigl(\lambda_{\eta,R}(W_\eta),
\lambda_{\zeta,R}(B_\zeta)\bigr).
\]
For the wrapped/wrapped row of
Definition~\ref{def:wrapped-wrapped-correspondence-datum}, it is
\[
a^{\mathrm{WW}}_{\eta_1,\eta_2;\zeta,R}:
\mathfrak E^{\mathrm{WW}}_{\eta_1,\eta_2;\zeta,R}\longrightarrow E^3,
\qquad
\xi\longmapsto
\bigl(\lambda_{\eta_1,R}(W_1),
\lambda_{\eta_2,R}(W_2),
\lambda_{\zeta,R}(B_\zeta)\bigr).
\]
For the ordered two-sided mixed row of
Definition~\ref{def:two-sided-mixed-correspondence-datum}, it is
\[
a^{\mathrm{LWL}}_{\alpha_-;\eta;\beta_+,R,I_-,I_+}:
\mathfrak E^{\mathrm{LWL}}_{\alpha_-;\eta;\beta_+,R,I_-,I_+}
\longrightarrow E^2,
\qquad
\xi\longmapsto
\bigl(\lambda_{\eta,R}(W_\eta),
\lambda_{\zeta,R}(C_\zeta)\bigr).
\]
These maps record ordered source anchors and target anchors on the
unreduced stack.  The local/local correspondence has no wrapped input
and no wrapped anchor-memory row.  The determinant-anchor construction
and translation-weight computation are dedicated later packets; the
\(\chi=0\) degeneracy and the degree-zero Jordan--Hoelder extra anchor
are isolated by their own packets.  The anchor residual definition
\(o^\lambda_R\), its vanishing, anchor losslessness, quotient descent,
stabilizer linearization, and transition compatibility remain separate
rows.
Replacing the tuple above by a scalar Fock factor, a Borcherds
\(s\)-degree, or a post-quotient class is not this datum.
\end{definition}

The packet
\[
\texttt{certificates/hybrid/source\_target\_anchor\_memory\_definition}
\]
verifies
\(\texttt{SOURCE\_TARGET\_ANCHOR\_MEMORY\_DEFINITION\_VERIFIED}\):
the \(LW\), \(WL\), \(WW\), and \(LWL\) source/target anchor-memory
maps are defined before quotient.  The determinant anchor is constructed
by the next packet, and its translation weight is computed by the
following packet.  The \(\chi=0\) degeneracy is isolated by the
chi-zero packet.  The degree-zero Jordan--Hoelder extra-anchor datum is
recorded by
\(\texttt{certificates/hybrid/determinant\_anchor\_extra\_anchor\_data}\);
anchor-residual definition and vanishing, full losslessness,
quotient-descent, stabilizer linearization, populated finite-stack, and
transition rows remain separate obligations.

The packet
\[
\texttt{certificates/hybrid/determinant\_anchor\_construction}
\]
verifies
\(\texttt{DETERMINANT\_ANCHOR\_CONSTRUCTION\_VERIFIED}\):
Definition~\ref{def:determinant-anchor} constructs the normalized
determinant-of-cohomology anchor
\[
\lambda^{\det}_{\eta,R}(\mathcal F_{\eta,T})
=\det R\pi_{E,T,*}\mathcal F_{\eta,T}
\otimes\operatorname{pr}_E^*\mathcal O_E(-\chi_\eta\cdot0_E)
\in\operatorname{Pic}^0(E)
\]
as a functorial map to \(\operatorname{Pic}^0(E)\simeq E\).  The packet
checks the degree-zero normalization and determinant-functor
base-change; unit-weight normalization, the Jordan--Hoelder
extra-anchor datum, the residual definition \(o^\lambda_R\), residual
vanishing, losslessness, quotient descent, populated wrapped-prequotient
rows, and transition rows remain separate.  The \(\chi=0\) degeneracy
is isolated by the chi-zero packet below.
The translation-weight packet
\[
\texttt{certificates/hybrid/determinant\_anchor\_translation\_weight}
\]
verifies
\(\texttt{DETERMINANT\_ANCHOR\_TRANSLATION\_WEIGHT\_VERIFIED}\):
for the support-translation convention
\(a\cdot\mathcal F=(\operatorname{id}_S\times\tau_a)_*\mathcal F\),
Proposition~\ref{prop:determinant-anchor-translation-weight} gives
\[
\lambda^{\det}_{\eta,R}(a\cdot\mathcal F)
=\lambda^{\det}_{\eta,R}(\mathcal F)+\chi_\eta a
\quad\text{in }\operatorname{Pic}^0(E)\simeq E .
\]
The packet records the finite-stabilizer restriction
\(h\mapsto\chi_\eta h\).  It does not choose coherent stabilizer
linearizations, make the weight one, handle the \(\chi_\eta=0\)
degeneracy, identify the Jordan--Hoelder extra anchor, prove
losslessness, descend through the quotient, or prove transition
compatibility.
The chi-zero degeneracy packet
\[
\texttt{certificates/hybrid/determinant\_anchor\_chi\_zero\_degeneracy}
\]
verifies
\(\texttt{DETERMINANT\_ANCHOR\_CHI\_ZERO\_DEGENERACY\_VERIFIED}\):
Proposition~\ref{prop:determinant-anchor-chi-zero-degeneracy} proves
that when \(\chi_\eta=0\), the determinant anchor is constant on the
full \(E\)-translation orbit.  Freeness of the retained
\(E\)-translation action gives a rank-one separation defect for
determinant-anchor-only data.  The packet records the necessity of
extra anchor data; it does not itself construct that data, define
\(o^\lambda_R\), prove losslessness, descend through the quotient, or
prove transition compatibility.
The determinant-anchor extra-anchor-data packet
\[
\texttt{certificates/hybrid/determinant\_anchor\_extra\_anchor\_data}
\]
verifies
\(\texttt{DETERMINANT\_ANCHOR\_EXTRA\_ANCHOR\_DATA\_VERIFIED}\):
Proposition~\ref{prop:determinant-anchor-extra-data-separates-rank-two}
adds the degree-zero Jordan--Hoelder divisor of the elliptic
pushforward as anchor data.  The Abel sum is the determinant anchor,
and the rank-two pair \(\mathcal O_E^{\oplus2}\) and
\(L\oplus L^{-1}\), \(L\ne\mathcal O_E\), is separated despite equal
determinant.  This packet does not prove global anchor losslessness,
prove \(o^\lambda_R=0\), descend the repaired anchor through the
quotient, choose stabilizer linearizations, populate all finite
wrapped-prequotient rows, or prove transition compatibility.

\begin{definition}[Finite anchor residual]
\label{def:finite-anchor-residual}
Fix a finite HN height \(R\).  For
\(\eta\in\Gamma_R^{\mathrm{wr}}\), let
\[
\mathcal A_{\eta,R}:=
\operatorname{Sym}^{r_\eta}\operatorname{Pic}^0(E),
\qquad
\mathrm{Ab}_{\eta,R}:\mathcal A_{\eta,R}\longrightarrow
\operatorname{Pic}^0(E)\simeq E
\]
be the Abel-sum morphism.  A repaired wrapped anchor over the
prequotient row is a functorial map
\[
\Lambda_{\eta,R}:\mathcal M_{\eta,R}^{\mathrm{wr,rig}}
\longrightarrow \mathcal A_{\eta,R}
\]
such that
\[
\mathrm{Ab}_{\eta,R}\circ\Lambda_{\eta,R}
=\lambda^{\det}_{\eta,R},
\]
and whose restriction to the degree-zero vector-bundle row is the
Jordan--Hoelder divisor of
Proposition~\ref{prop:determinant-anchor-extra-data-separates-rank-two}.
For an exact sequence of semistable degree-zero elliptic pushforwards,
the anchor target has the divisor-union operation
\[
\mu_{\eta_1,\eta_2}:
\mathcal A_{\eta_1,R}\times\mathcal A_{\eta_2,R}
\longrightarrow \mathcal A_{\eta_1+\eta_2,R}.
\]
The finite anchor residual is the five-component defect
\[
\mathfrak o^\lambda_R=\bigl(
\mathfrak o^{\lambda,\mathrm{ex}}_R,\
\mathfrak o^{\lambda,\mathrm{unit}}_R,\
\mathfrak o^{\lambda,\mathrm{loss}}_R,\
\mathfrak o^{\lambda,\mathrm{multi}}_R,\
\mathfrak o^{\lambda,\mathrm{tr}}_{R'R}\bigr),
\]
with the following zero conditions.
The \(E\)-valued source/target anchors of
Definition~\ref{def:source-target-anchor-memory-datum} are the
Abel-sum shadows \(\mathrm{Ab}_{\eta,R}\circ\Lambda_{\eta,R}\); they
do not replace the repaired symmetric-product anchor before the
residual vanishes.
\[
\mathfrak o^{\lambda,\mathrm{ex}}_R=0
\]
if and only if every \(\eta\in\Gamma_R^{\mathrm{wr}}\) carries such a
functorial repaired anchor \(\Lambda_{\eta,R}\), including the
semistable vector-bundle locus needed for the Jordan--Hoelder divisor.
The component
\(\mathfrak o^{\lambda,\mathrm{unit}}_R\) vanishes if and only if the
chosen repaired anchors admit a finite-cover normalization with
translation weight one for the retained \(E\)-action.  The component
\(\mathfrak o^{\lambda,\mathrm{loss}}_R\) vanishes if and only if the
kernel pair of \(\Lambda_{\eta,R}\) on each retained wrapped row is the
retained equivalence relation, so that no extra wrapped position is
identified by the anchor.  The component
\(\mathfrak o^{\lambda,\mathrm{multi}}_R\) vanishes if and only if, on
every \(LW\), \(WL\), \(WW\), and \(LWL\) unreduced correspondence, the
target repaired anchor equals the divisor-union operation applied to
the ordered wrapped source anchors.  The component
\(\mathfrak o^{\lambda,\mathrm{tr}}_{R'R}\) vanishes if and only if,
for every \(R\le R'\), the repaired anchors commute with the HN
restriction \(R'\to R\).  This definition supplies the residual and
its zero semantics.  It does not prove any component vanishes, does
not descend the anchor through the reduced \(E\)-quotient, and does not
choose stabilizer linearizations.
\end{definition}

The packet
\[
\texttt{certificates/hybrid/anchor\_residual\_definition}
\]
verifies
\(\texttt{ANCHOR\_RESIDUAL\_DEFINITION\_VERIFIED}\):
Definition~\ref{def:finite-anchor-residual} defines the repaired
anchor target, its Abel-sum projection, and the five residual
components
\(\mathrm{ex}\), \(\mathrm{unit}\), \(\mathrm{loss}\),
\(\mathrm{multi}\), and \(\mathrm{tr}\).  The packet does not prove
first-window vanishing, quotient descent, stabilizer linearization,
finite-stage population, or transition compatibility.

\begin{definition}[Quotient-after-correspondence pseudofunctor]
\label{def:quotient-after-correspondence-pseudofunctor}
\textnormal{[definition only]}
Fix a finite HN height \(R\).  Let
\(\mathsf{Corr}^{E,\mathrm{hyb}}_R\) be the unreduced
\(E\)-equivariant hybrid correspondence category generated by the
retained local, mixed, wrapped, two-sided mixed, and flag
correspondence stacks before the reduced \(E\)-quotient.  A
quotient-after-correspondence datum consists of the following finite
rows.

For every retained object \(Y\) of
\(\mathsf{Corr}^{E,\mathrm{hyb}}_R\), including local closed
configuration charts, wrapped rigidified prequotients, correspondence
stacks, and flag stacks, choose the quotient stack
\[
\overline Y=[Y/E]
\]
for the retained diagonal translation action and the quotient map
\[
q_Y:Y\longrightarrow\overline Y .
\]
For every generating unreduced correspondence
\[
e=(Y_0\xleftarrow{p_e}\mathfrak E_e\xrightarrow{q_e}Y_1)
\]
whose maps are \(E\)-equivariant, choose the reduced correspondence
\[
\overline e=
(\overline Y_0\xleftarrow{\overline p_e}
[\mathfrak E_e/E]\xrightarrow{\overline q_e}\overline Y_1)
\]
induced after the same correspondence stack has been formed.  The
assignment
\[
\mathcal Q^{\mathrm{corr}}_{E,R}(Y)=\overline Y,\qquad
\mathcal Q^{\mathrm{corr}}_{E,R}(e)=\overline e
\]
is the quotient-after-correspondence pseudofunctor at the level of
objects and generating arrows.

For a reduced vanishing-cycle coefficient \(K_Y\) on \(Y\), the datum
also specifies the coefficient-descent interface
\[
\overline K_Y\in D^b_c(\overline Y),\qquad
q_Y^*\overline K_Y\simeq K_Y,
\]
whenever the relevant stabilizer linearization and orientation descent
rows are supplied.  The induced chain functor
\[
Q_{E,R}:\mathsf{BM}^{E,\mathrm{hyb}}_R
\longrightarrow
\mathsf{BM}^{\mathrm{red},\mathrm{hyb}}_R
\]
is not part of this definition; it is the later Borel--Moore chain
row built from these quotient and coefficient-descent interfaces.

The finite quotient residual attached to this datum is
\[
\mathfrak o^Q_R=
\bigl(
\mathfrak o^{Q,\mathrm{form}}_R,\
\mathfrak o^{Q,\mathrm{comp}}_R,\
\mathfrak o^{Q,\mathrm{bc}}_R,\
\mathfrak o^{Q,\mathrm{TS}}_R,\
\mathfrak o^{Q,\mathrm{flag}}_R,\
\mathfrak o^{Q,\mathrm{BM}}_R,\
\mathfrak o^{Q,\theta}_R,\
\mathfrak o^{Q,\mathrm{tr}}_{R'R}
\bigr).
\]
The first component says that objects, arrows, quotient maps, and
coefficient-descent interfaces are defined.  The remaining components
measure, respectively, preservation of composition, base-change,
Thom--Sebastiani transport, flag-stack coherences, Borel--Moore chain
functoriality, the comparison maps \(\theta^Q_{\mu,R}\), and
compatibility with \(R'\to R\) transitions.

This is a quotient-after-correspondence definition: the correspondence
stack \(\mathfrak E_e\) is formed before applying \([-/E]\).  It is
not a quotient-first Hall product on already reduced object stacks.
\end{definition}

The packet
\[
\texttt{certificates/hybrid/quotient\_after\_correspondence\_pseudofunctor\_definition}
\]
verifies
\(\texttt{QUOTIENT\_AFTER\_CORRESPONDENCE\_PSEUDOFUNCTOR\_DEFINED}\):
Definition~\ref{def:quotient-after-correspondence-pseudofunctor}
defines the object assignment, arrow assignment, coefficient-descent
interface, and eight quotient residual components of the
quotient-after-correspondence pseudofunctor before Borel--Moore
chains.  The composition component is supplied by
Proposition~\ref{prop:quotient-after-correspondence-preserves-composition}.
The base-change square component is supplied by
Proposition~\ref{prop:quotient-after-correspondence-preserves-base-change}.
The Thom--Sebastiani descent component is supplied by
Proposition~\ref{prop:quotient-after-correspondence-preserves-thom-sebastiani}.
The quotient-first exclusion theorem is supplied by
Proposition~\ref{prop:quotient-after-correspondence-excludes-quotient-first}.
The Borel--Moore chain functor is supplied by
Proposition~\ref{prop:quotient-after-correspondence-borel-moore-chain-functor}.
The operation-comparison maps \(\theta^Q_{\mu,R}\) are supplied by
Proposition~\ref{prop:quotient-after-correspondence-theta-mu-comparison}.
The quotient two-step flag associativity component is supplied by
Proposition~\ref{prop:quotient-after-correspondence-theta-mu-associativity}.
It does not prove compact-support pushforward base change, transition
compatibility, or aggregate population.

\begin{proposition}[Quotient-after-correspondence preserves composition]
\label{prop:quotient-after-correspondence-preserves-composition}
\textnormal{[proved for finite-height composition of spans]}
Fix \(R\) and the quotient-after-correspondence datum of
Definition~\ref{def:quotient-after-correspondence-pseudofunctor}.
Let
\[
e=(Y_0\xleftarrow{p_e}\mathfrak E_e\xrightarrow{q_e}Y_1),
\qquad
f=(Y_1\xleftarrow{p_f}\mathfrak E_f\xrightarrow{q_f}Y_2)
\]
be composable retained \(E\)-equivariant hybrid correspondences.  The
unreduced composite has middle stack
\[
\mathfrak E_e\times_{Y_1}\mathfrak E_f
\]
with the diagonal \(E\)-action.  There is a canonical isomorphism of
reduced correspondences
\[
\Xi_{f,e}\colon
\mathcal Q^{\mathrm{corr}}_{E,R}(f\circ e)
\xrightarrow{\ \sim\ }
\mathcal Q^{\mathrm{corr}}_{E,R}(f)
\circ
\mathcal Q^{\mathrm{corr}}_{E,R}(e),
\]
whose middle-stack component is
\[
\bigl[(\mathfrak E_e\times_{Y_1}\mathfrak E_f)/E\bigr]
\simeq
[\mathfrak E_e/E]\times_{[Y_1/E]}[\mathfrak E_f/E].
\]
The isomorphisms \(\Xi_{f,e}\) satisfy the associativity cocycle for
three composable correspondences.  Hence
\[
\mathfrak o^{Q,\mathrm{comp}}_R=0.
\]
This proves no base-change theorem for compact-support pushforwards,
no Thom--Sebastiani compatibility, no Borel--Moore chain functor
\(Q_{E,R}\), no comparison map \(\theta^Q_{\mu,R}\), no quotient-first
exclusion theorem, and no transition or aggregate statement.
\end{proposition}

\begin{proof}
The maps \(q_e\) and \(p_f\) are \(E\)-equivariant, so the fibre
product \(\mathfrak E_e\times_{Y_1}\mathfrak E_f\) carries the
diagonal \(E\)-action.  For a test scheme \(T\), an object of
\([\mathfrak E_e/E]\times_{[Y_1/E]}[\mathfrak E_f/E]\) over \(T\)
consists of \(E\)-torsors \(P_e,P_f\to T\), equivariant maps
\[
u:P_e\to\mathfrak E_e,\qquad v:P_f\to\mathfrak E_f,
\]
and an isomorphism \(\alpha:P_e\xrightarrow{\sim}P_f\) over the
common image in \([Y_1/E]\); equivalently \(p_fv\alpha=q_eu\) as maps
\(P_e\to Y_1\).  Transporting \(v\) along \(\alpha\) gives one torsor
\(P_e\to T\) with compatible equivariant maps to \(\mathfrak E_e\) and
\(\mathfrak E_f\).  This is precisely an \(E\)-equivariant map
\[
P_e\longrightarrow \mathfrak E_e\times_{Y_1}\mathfrak E_f,
\]
or equivalently an object of
\([(\mathfrak E_e\times_{Y_1}\mathfrak E_f)/E]\) over \(T\).
The converse sends one torsor with a map to the fibre product to the
two quotient-stack objects with the identity torsor isomorphism.
Morphisms are transported by the same torsor isomorphisms.  The two
constructions are inverse and functorial in \(T\), giving the
displayed isomorphism of quotient stacks.

The source and target maps to \([Y_0/E]\) and \([Y_2/E]\) are induced
from \(p_e\) and \(q_f\), so the middle-stack isomorphism identifies
the quotient of the unreduced composite with the composite of the two
reduced correspondences.  For three composable correspondences both
iterated comparisons are obtained from the same quotient of the
triple fibre product, and the two maps agree by the associativity of
2-fibre products and the universal property of quotient stacks.  Thus
the composition defect has rank zero, while all coefficient,
base-change, Thom--Sebastiani, Borel--Moore, transition, and aggregate
claims remain outside this proposition.
\end{proof}

The packet
\[
\texttt{certificates/hybrid/quotient\_after\_correspondence\_composition}
\]
verifies
\(\texttt{QUOTIENT\_AFTER\_CORRESPONDENCE\_COMPOSITION\_VERIFIED}\):
Proposition~\ref{prop:quotient-after-correspondence-preserves-composition}
proves that quotient-after-correspondence commutes with finite-height
composition of retained \(E\)-equivariant spans and that the
associativity cocycle has defect rank zero.  The base-change
component is supplied by
Proposition~\ref{prop:quotient-after-correspondence-preserves-base-change}.
The Thom--Sebastiani descent component is supplied by
Proposition~\ref{prop:quotient-after-correspondence-preserves-thom-sebastiani}.
The quotient-first exclusion theorem is supplied by
Proposition~\ref{prop:quotient-after-correspondence-excludes-quotient-first}.
The Borel--Moore chain functor is supplied by
Proposition~\ref{prop:quotient-after-correspondence-borel-moore-chain-functor}.
The operation-comparison maps \(\theta^Q_{\mu,R}\) are supplied by
Proposition~\ref{prop:quotient-after-correspondence-theta-mu-comparison}.
The quotient two-step flag associativity component is supplied by
Proposition~\ref{prop:quotient-after-correspondence-theta-mu-associativity}.
The HN-transition compatibility of \(Q_{E,R}\) is supplied by
Proposition~\ref{prop:quotient-after-correspondence-hn-transition-compatibility}.
It does not prove aggregate population.

\begin{proposition}[Quotient-after-correspondence preserves base-change squares]
\label{prop:quotient-after-correspondence-preserves-base-change}
\textnormal{[proved for finite-height cartesian squares of stacks]}
Fix \(R\) and the quotient-after-correspondence datum of
Definition~\ref{def:quotient-after-correspondence-pseudofunctor}.
Let
\[
\begin{array}{ccc}
X' & \xrightarrow{\widetilde g} & X\\
\scriptstyle h'\downarrow && \downarrow\scriptstyle h\\
Y' & \xrightarrow{g} & Y
\end{array}
\qquad(\square)
\]
be a retained \(E\)-equivariant cartesian square of stacks appearing
as a source, target, open, or compactified square of the finite
hybrid correspondence calculus.  Then the quotient square
\[
\begin{array}{ccc}
[X'/E] & \xrightarrow{[\widetilde g/E]} & [X/E]\\
\scriptstyle [h'/E]\downarrow && \downarrow\scriptstyle [h/E]\\
{[Y'/E]} & \xrightarrow{[g/E]} & {[Y/E]}
\end{array}
\]
is cartesian.  Equivalently, the canonical morphism
\[
[X'/E]\longrightarrow [X/E]\times_{[Y/E]}[Y'/E]
\]
is an isomorphism.  Hence quotient-after-correspondence preserves the
retained base-change squares and
\[
\mathfrak o^{Q,\mathrm{bc}}_R=0.
\]
This proves only the stack-level cartesian-square component of the
quotient residual.  It does not prove compact-support pushforward
base change, coefficient pullback for vanishing cycles, the
projection formula, Thom--Sebastiani compatibility, the chain functor
\(Q_{E,R}\), the comparison maps \(\theta^Q_{\mu,R}\), transition
compatibility, or aggregate population.
\end{proposition}

\begin{proof}
Since the square \((\square)\) is \(E\)-equivariant and cartesian,
there is an \(E\)-equivariant isomorphism
\[
X'\simeq X\times_Y Y'.
\]
For a test scheme \(T\), an object of
\([X/E]\times_{[Y/E]}[Y'/E]\) over \(T\) consists of \(E\)-torsors
\(P_X,P_{Y'}\to T\), equivariant maps
\[
u:P_X\to X,\qquad v:P_{Y'}\to Y',
\]
and an isomorphism
\(\alpha:P_X\xrightarrow{\sim}P_{Y'}\) over the common image in
\([Y/E]\).  After transporting \(v\) along \(\alpha\), this is one
\(E\)-torsor \(P_X\to T\) with maps \(u:P_X\to X\) and
\(v\alpha:P_X\to Y'\) satisfying
\[
h\,u=g\,v\,\alpha .
\]
By the cartesian identity \(X'\simeq X\times_Y Y'\), this is the same
as an \(E\)-equivariant map \(P_X\to X'\), hence an object of
\([X'/E]\) over \(T\).  The converse sends an equivariant map
\(P\to X'\) to its two projections with the identity torsor
isomorphism.  Morphisms are carried by the same torsor isomorphisms,
and the two constructions are inverse and functorial in \(T\).
Therefore the quotient square is cartesian.

Applying this argument to each retained source, target, open, and
compactified square in the finite hybrid correspondence calculus gives
zero defect for preservation of base-change squares.  The proof has
not applied \(h_!\), \(h^*\), reduced Thom--Sebastiani transport, or
Borel--Moore chains, so the corresponding functorial and chain-level
claims remain separate obligations.
\end{proof}

The packet
\[
\texttt{certificates/hybrid/quotient\_after\_correspondence\_base\_change}
\]
verifies
\(\texttt{QUOTIENT\_AFTER\_CORRESPONDENCE\_BASE\_CHANGE\_VERIFIED}\):
Proposition~\ref{prop:quotient-after-correspondence-preserves-base-change}
proves that quotient-after-correspondence preserves retained
cartesian base-change squares of finite-height \(E\)-equivariant
stacks.  It does not prove compact-support pushforward base change,
coefficient pullback, or aggregate population.  The
Thom--Sebastiani descent component is supplied by
Proposition~\ref{prop:quotient-after-correspondence-preserves-thom-sebastiani}.
The quotient-first exclusion theorem is supplied by
Proposition~\ref{prop:quotient-after-correspondence-excludes-quotient-first}.
The Borel--Moore chain functor is supplied by
Proposition~\ref{prop:quotient-after-correspondence-borel-moore-chain-functor}.
The operation-comparison maps \(\theta^Q_{\mu,R}\) are supplied by
Proposition~\ref{prop:quotient-after-correspondence-theta-mu-comparison}.
The quotient two-step flag associativity component is supplied by
Proposition~\ref{prop:quotient-after-correspondence-theta-mu-associativity}.
The HN-transition compatibility of \(Q_{E,R}\) is supplied by
Proposition~\ref{prop:quotient-after-correspondence-hn-transition-compatibility}.

\begin{proposition}[Quotient-after-correspondence preserves Thom--Sebastiani transport]
\label{prop:quotient-after-correspondence-preserves-thom-sebastiani}
\textnormal{[proved for equivariant finite-height Thom--Sebastiani rows]}
Fix \(R\) and the quotient-after-correspondence datum of
Definition~\ref{def:quotient-after-correspondence-pseudofunctor}.
Let
\[
e=(Y_{\mathrm{src}}\xleftarrow{p_e}\mathfrak E_e
\xrightarrow{q_e}Y_{\mathrm{tar}})
\]
be a retained \(E\)-equivariant correspondence whose unreduced
reduced vanishing-cycle coefficients carry an \(E\)-equivariant
Thom--Sebastiani transport
\[
\operatorname{TS}^{\mathrm{red}}_e:
p_e^*K_{Y_{\mathrm{src}}}\xrightarrow{\sim}K_{\mathfrak E_e}.
\]
Assume the coefficient-descent interfaces
\[
q_{Y_{\mathrm{src}}}^{*}\overline K_{Y_{\mathrm{src}}}
\simeq K_{Y_{\mathrm{src}}},\qquad
q_{\mathfrak E_e}^{*}\overline K_{\mathfrak E_e}
\simeq K_{\mathfrak E_e}
\]
are supplied, including the stabilizer linearizations and orientation
descent needed for the reduced vanishing-cycle complexes.  Then there
is a unique reduced Thom--Sebastiani transport
\[
\overline{\operatorname{TS}}^{\mathrm{red}}_{\overline e}:
\overline p_e^{\,*}\overline K_{Y_{\mathrm{src}}}
\xrightarrow{\sim}\overline K_{\mathfrak E_e}
\]
on the quotient correspondence
\[
\overline e=([Y_{\mathrm{src}}/E]\xleftarrow{\overline p_e}
[\mathfrak E_e/E]\xrightarrow{\overline q_e}[Y_{\mathrm{tar}}/E])
\]
whose pullback along \(q_{\mathfrak E_e}\) is
\(\operatorname{TS}^{\mathrm{red}}_e\) under the descent
identifications and the cartesian quotient square of
Proposition~\ref{prop:quotient-after-correspondence-preserves-base-change}.
Hence
\[
\mathfrak o^{Q,\mathrm{TS}}_R=0.
\]
This proves only descent of the reduced Thom--Sebastiani coefficient
transport through the quotient-after-correspondence construction.  It
does not prove the existence of the unreduced BBDJS
Thom--Sebastiani row, the Joyce--Upmeier orientation row,
compact-support pushforward base change, flag-stack coherences, the
Borel--Moore chain functor \(Q_{E,R}\), the comparison maps
\(\theta^Q_{\mu,R}\), transition compatibility, or aggregate
population.
\end{proposition}

\begin{proof}
The quotient square for the source map is cartesian by
Proposition~\ref{prop:quotient-after-correspondence-preserves-base-change}:
\[
\begin{array}{ccc}
\mathfrak E_e & \xrightarrow{q_{\mathfrak E_e}} & [\mathfrak E_e/E]\\
\scriptstyle p_e\downarrow && \downarrow\scriptstyle \overline p_e\\
Y_{\mathrm{src}} & \xrightarrow{q_{Y_{\mathrm{src}}}} &
[Y_{\mathrm{src}}/E].
\end{array}
\]
Therefore
\[
q_{\mathfrak E_e}^*\overline p_e^{\,*}
\overline K_{Y_{\mathrm{src}}}
\simeq
p_e^*q_{Y_{\mathrm{src}}}^*
\overline K_{Y_{\mathrm{src}}}
\simeq p_e^*K_{Y_{\mathrm{src}}}.
\]
The second descent interface gives
\[
q_{\mathfrak E_e}^*\overline K_{\mathfrak E_e}
\simeq K_{\mathfrak E_e}.
\]
Thus the unreduced Thom--Sebastiani isomorphism is a morphism between
the pullbacks of the two reduced coefficient objects on
\([\mathfrak E_e/E]\).

Constructible sheaves on the quotient stack \([\mathfrak E_e/E]\) are
equivalent to \(E\)-equivariant constructible sheaves on
\(\mathfrak E_e\), with the stated stabilizer linearizations.  Under
this equivalence, morphisms on the quotient stack are exactly
equivariant morphisms on \(\mathfrak E_e\).  Since
\(\operatorname{TS}^{\mathrm{red}}_e\) is \(E\)-equivariant by
hypothesis, it descends uniquely to a morphism
\[
\overline{\operatorname{TS}}^{\mathrm{red}}_{\overline e}:
\overline p_e^{\,*}\overline K_{Y_{\mathrm{src}}}
\longrightarrow\overline K_{\mathfrak E_e}.
\]
Its pullback is an isomorphism, hence the descended morphism is an
isomorphism.  The same equivalence of quotient sheaves proves
uniqueness and functoriality for every retained finite-height
Thom--Sebastiani row.  No pushforward, Borel--Moore chain functor, or
flag comparison has been applied, so those coherences remain outside
this proposition.
\end{proof}

The packet
\[
\texttt{certificates/hybrid/quotient\_after\_correspondence\_thom\_sebastiani}
\]
verifies
\(\texttt{QUOTIENT\_AFTER\_CORRESPONDENCE\_THOM\_SEBASTIANI\_VERIFIED}\):
Proposition~\ref{prop:quotient-after-correspondence-preserves-thom-sebastiani}
proves that an equivariant reduced Thom--Sebastiani transport descends
uniquely through the quotient-after-correspondence construction once
the coefficient-descent interfaces are supplied.  It does not prove
aggregate population.
The operation-comparison maps \(\theta^Q_{\mu,R}\) are supplied by
Proposition~\ref{prop:quotient-after-correspondence-theta-mu-comparison}.
The quotient-first exclusion theorem is
Proposition~\ref{prop:quotient-after-correspondence-excludes-quotient-first}.
The Borel--Moore chain functor is
Proposition~\ref{prop:quotient-after-correspondence-borel-moore-chain-functor}.
The quotient two-step flag associativity component is
Proposition~\ref{prop:quotient-after-correspondence-theta-mu-associativity}.
The HN-transition compatibility of \(Q_{E,R}\) is
Proposition~\ref{prop:quotient-after-correspondence-hn-transition-compatibility}.

\begin{proposition}[Quotient-first construction is excluded]
\label{prop:quotient-after-correspondence-excludes-quotient-first}
\textnormal{[proved as a finite-height image characterization]}
Fix \(R\) and the quotient-after-correspondence datum of
Definition~\ref{def:quotient-after-correspondence-pseudofunctor}.
Every generating arrow in the image of
\(\mathcal Q^{\mathrm{corr}}_{E,R}\) is represented by a retained
unreduced \(E\)-equivariant correspondence
\[
e=(Y_0\xleftarrow{p_e}\mathfrak E_e\xrightarrow{q_e}Y_1)
\]
together with the quotient presentation
\[
\overline e=
([Y_0/E]\xleftarrow{\overline p_e}
[\mathfrak E_e/E]\xrightarrow{\overline q_e}[Y_1/E]).
\]
Equivalently, the two squares
\[
\begin{array}{ccc}
\mathfrak E_e & \longrightarrow & [\mathfrak E_e/E]\\
\scriptstyle p_e\downarrow && \downarrow\scriptstyle\overline p_e\\
Y_0 & \longrightarrow & [Y_0/E]
\end{array}
\qquad
\begin{array}{ccc}
\mathfrak E_e & \longrightarrow & [\mathfrak E_e/E]\\
\scriptstyle q_e\downarrow && \downarrow\scriptstyle\overline q_e\\
Y_1 & \longrightarrow & [Y_1/E]
\end{array}
\]
are part of the arrow data.  A quotient-first span, obtained by first
passing to \([Y_0/E]\) and \([Y_1/E]\) and then choosing an independent
middle stack over their product, is not an arrow of
\(\mathcal Q^{\mathrm{corr}}_{E,R}\) unless it is equipped with a
specified isomorphism to \([\mathfrak E_e/E]\) compatible with these
two squares.  In that exceptional case it is the same quotient-after-
correspondence arrow after transport along the specified isomorphism.
Thus no quotient-first Hall product is used in the finite quotient
correspondence calculus, and the quotient-first exclusion row has
defect rank zero.
\end{proposition}

\begin{proof}
The object assignment of
Definition~\ref{def:quotient-after-correspondence-pseudofunctor}
sends an unreduced retained object \(Y\) to \([Y/E]\).  Its arrow
assignment does not choose a new reduced extension stack.  It sends
the already formed unreduced correspondence
\[
Y_0\xleftarrow{p_e}\mathfrak E_e\xrightarrow{q_e}Y_1
\]
to the reduced correspondence with middle stack
\([\mathfrak E_e/E]\).  Hence the image of
\(\mathcal Q^{\mathrm{corr}}_{E,R}\) is the full finite-height family
of reduced spans carrying such an unreduced quotient presentation.

Let
\[
[Y_0/E]\xleftarrow{a}Z\xrightarrow{b}[Y_1/E]
\]
be a quotient-first candidate obtained after quotienting the two
object stacks.  If \(Z\) is not supplied with an isomorphism
\[
Z\simeq[\mathfrak E_e/E]
\]
for a retained unreduced \(E\)-equivariant correspondence \(e\), and
if \(a,b\) are not the transported maps \(\overline p_e,\overline q_e\),
then this candidate lacks the quotient map
\(\mathfrak E_e\to[\mathfrak E_e/E]\) used in the definition.  It is
therefore not in the image of
\(\mathcal Q^{\mathrm{corr}}_{E,R}\).  If such an isomorphism is
supplied, transporting the span structure along it identifies the
candidate with the quotient-after-correspondence span assigned to
\(e\); no independent quotient-first construction remains.

The composition comparison of
Proposition~\ref{prop:quotient-after-correspondence-preserves-composition}
uses the quotient of the unreduced fibre product
\(\mathfrak E_e\times_{Y_1}\mathfrak E_f\).  The base-change
comparison of
Proposition~\ref{prop:quotient-after-correspondence-preserves-base-change}
uses quotient presentations of the unreduced cartesian squares.  The
Thom--Sebastiani descent of
Proposition~\ref{prop:quotient-after-correspondence-preserves-thom-sebastiani}
pulls coefficients back along
\(\mathfrak E_e\to[\mathfrak E_e/E]\).  None of these rows can be
interpreted on an independent quotient-first middle stack unless it
has first been identified with the quotient of the unreduced middle
stack.  The exclusion is therefore a statement about the finite
correspondence calculus itself, not a Borel--Moore, theta-comparison,
transition, or aggregate-population claim.
\end{proof}

The packet
\[
\texttt{certificates/hybrid/quotient\_after\_correspondence\_quotient\_first\_exclusion}
\]
verifies
\(\texttt{QUOTIENT\_AFTER\_CORRESPONDENCE\_QUOTIENT\_FIRST\_EXCLUSION\_VERIFIED}\):
Proposition~\ref{prop:quotient-after-correspondence-excludes-quotient-first}
proves that every reduced generating arrow used by
\(\mathcal Q^{\mathrm{corr}}_{E,R}\) carries a quotient presentation
from a preformed unreduced \(E\)-equivariant correspondence.  An
object-quotient-first Hall product is excluded unless it is identified
with that quotient presentation.  This packet does not prove transition
compatibility for Hall products or aggregate population.  The
Borel--Moore chain functor is supplied by
Proposition~\ref{prop:quotient-after-correspondence-borel-moore-chain-functor}.
The operation-comparison maps \(\theta^Q_{\mu,R}\) are supplied by
Proposition~\ref{prop:quotient-after-correspondence-theta-mu-comparison}.
The quotient two-step flag associativity component is supplied by
Proposition~\ref{prop:quotient-after-correspondence-theta-mu-associativity}.
The HN-transition compatibility of \(Q_{E,R}\) is supplied by
Proposition~\ref{prop:quotient-after-correspondence-hn-transition-compatibility}.

\begin{proposition}[The quotient \(E\)-descent functor on Borel--Moore chains]
\label{prop:quotient-after-correspondence-borel-moore-chain-functor}
\textnormal{[proved for quotient-presented finite-height coefficient objects]}
Fix \(R\) and the quotient-after-correspondence datum of
Definition~\ref{def:quotient-after-correspondence-pseudofunctor}.
For every retained object \(Y\) of
\(\mathsf{Corr}^{E,\mathrm{hyb}}_R\), suppose the reduced
vanishing-cycle coefficient \(K_Y\in D^b_{c,E}(Y)\) is supplied with
the descent interface
\[
\overline K_Y\in D^b_c([Y/E]),\qquad
q_Y^*\overline K_Y\simeq K_Y .
\]
Let
\[
C_*^{\mathrm{BM}}(Z,L):=
R\Gamma\!\left(Z,\mathbb D_Z L\right)
\]
denote Borel--Moore chains with constructible coefficient \(L\), with
\(\mathbb D_Z\) the Verdier duality functor on the retained finite
stack \(Z\).  Define
\[
C_*^{\mathrm{BM},E}(Y,K_Y)
:=
C_*^{\mathrm{BM}}([Y/E],\overline K_Y).
\]
Then the assignments
\[
(Y,K_Y)\longmapsto
([Y/E],\overline K_Y),\qquad
C_*^{\mathrm{BM},E}(Y,K_Y)\longmapsto
C_*^{\mathrm{BM}}([Y/E],\overline K_Y)
\]
extend to a functor
\[
Q_{E,R}:\mathsf{BM}^{E,\mathrm{hyb}}_R
\longrightarrow
\mathsf{BM}^{\mathrm{red},\mathrm{hyb}}_R .
\]
On an \(E\)-equivariant coefficient morphism
\(\alpha:K_Y\to L_Y\), the functor sends \(\alpha\) to the descended
morphism
\[
\overline\alpha:\overline K_Y\longrightarrow\overline L_Y
\]
under the quotient-stack descent equivalence, and then applies
\(C_*^{\mathrm{BM}}([Y/E],-)\).  Hence the Borel--Moore chain
component of the quotient residual vanishes:
\[
\mathfrak o^{Q,\mathrm{BM}}_R=0.
\]
This proves only the chain-level descent functor.  It does not prove
flag coherences, compact-support pushforward base change, the
comparison isomorphisms
\(\theta^Q_{\mu,R}\), compatibility with products or coproducts,
transition compatibility, protected-integration compatibility, or
aggregate population.
\end{proposition}

\begin{proof}
The coefficient-descent interface identifies the supplied
\(E\)-equivariant coefficient \(K_Y\) with the pullback of a
constructible coefficient on the quotient stack.  The standard
descent equivalence
\[
D^b_c([Y/E])\simeq D^b_{c,E}(Y)
\]
is fully faithful on morphisms: an equivariant morphism
\(\alpha:K_Y\to L_Y\) is uniquely the pullback of a morphism
\(\overline\alpha:\overline K_Y\to\overline L_Y\) whenever both
coefficients carry the prescribed descent data.  Identities and
composites descend to identities and composites because this is an
equivalence of categories.

Borel--Moore chains are functorial for coefficient morphisms on the
fixed stack:
\[
\overline\alpha:\overline K_Y\to\overline L_Y
\quad\Longmapsto\quad
R\Gamma([Y/E],\mathbb D\overline\alpha):
C_*^{\mathrm{BM}}([Y/E],\overline K_Y)
\to
C_*^{\mathrm{BM}}([Y/E],\overline L_Y).
\]
The object assignment \(Y\mapsto [Y/E]\) and the coefficient
assignment \(K_Y\mapsto\overline K_Y\) therefore define a functor on
the finite additive chain category generated by the retained
quotient-presented coefficient objects.  This is \(Q_{E,R}\).

No map of the form \(q_!\operatorname{TS}p^*\) has been compared with
its reduced counterpart.  The present functor supplies the chain
objects and descended coefficient morphisms on which the later
\(\theta^Q_{\mu,R}\) comparison maps may act; it does not construct
those maps.
\end{proof}

The packet
\[
\texttt{certificates/hybrid/quotient\_after\_correspondence\_bm\_chain\_functor}
\]
verifies
\(\texttt{QUOTIENT\_AFTER\_CORRESPONDENCE\_BM\_CHAIN\_FUNCTOR\_VERIFIED}\):
Proposition~\ref{prop:quotient-after-correspondence-borel-moore-chain-functor}
constructs \(Q_{E,R}\) on quotient-presented Borel--Moore chain
objects and proves \(\mathfrak o^{Q,\mathrm{BM}}_R=0\).  This packet
does not prove compact-support pushforward base change, protected
integration, or aggregate population.  The operation-comparison maps
\(\theta^Q_{\mu,R}\) are
supplied by
Proposition~\ref{prop:quotient-after-correspondence-theta-mu-comparison}.
The quotient two-step flag associativity component is supplied by
Proposition~\ref{prop:quotient-after-correspondence-theta-mu-associativity}.
The HN-transition compatibility of \(Q_{E,R}\) is supplied by
Proposition~\ref{prop:quotient-after-correspondence-hn-transition-compatibility}.

\begin{proposition}[Quotient comparison maps for retained operations]
\label{prop:quotient-after-correspondence-theta-mu-comparison}
\textnormal{[proved for quotient-presented finite-height operation rows]}
Fix \(R\) and the quotient-after-correspondence datum of
Definition~\ref{def:quotient-after-correspondence-pseudofunctor}.
Let
\[
e=\bigl(Y_1\times\cdots\times Y_k
\xleftarrow{\ p_e\ }\mathfrak E_e\xrightarrow{\ q_e\ }Y_0\bigr)
\]
be a retained \(k\)-ary local, mixed, wrapped, or flag operation row.
Assume the input, middle, and target coefficient objects carry
quotient-descent interfaces, the external product coefficient descends
to
\[
\overline K_1\boxtimes\cdots\boxtimes\overline K_k
\quad\text{on}\quad
[Y_1/E]\times\cdots\times [Y_k/E],
\]
the reduced Thom--Sebastiani transport descends as in
Proposition~\ref{prop:quotient-after-correspondence-preserves-thom-sebastiani},
and the chosen compact-support pushforward row for \(q_e\) is
quotient-presented by an isomorphism
\[
\delta_{q,e}\colon
\overline{\,q_{e,!}^{\mathrm{cs}}L_e\,}
\xrightarrow{\ \sim\ }
\overline q_{e,!}^{\mathrm{cs}}\overline L_e ,
\]
where
\[
L_e=
\operatorname{TS}^{\mathrm{red}}_e
p_e^*(K_1\boxtimes\cdots\boxtimes K_k),
\qquad
\overline L_e=
\overline{\operatorname{TS}}^{\mathrm{red}}_{\overline e}
\overline p_e^{\,*}
(\overline K_1\boxtimes\cdots\boxtimes\overline K_k).
\]
Then there is a canonical chain isomorphism
\[
\theta^Q_{\mu,e,R}\colon
Q_{E,R}\,
q_{e,!}^{\mathrm{cs}}\operatorname{TS}^{\mathrm{red}}_e p_e^*
\xrightarrow{\ \sim\ }
\overline q_{e,!}^{\mathrm{cs}}\,
\overline{\operatorname{TS}}^{\mathrm{red}}_{\overline e}
\overline p_e^{\,*}(Q_{E,R}^{\boxtimes k}).
\]
For every retained operation row with the displayed descent data,
\[
\mathfrak o^{Q,\theta}_R=0 .
\]
This proves only the existence of the operation-comparison
isomorphisms \(\theta^Q_{\mu,R}\).  It does not prove their
compatibility with two-step flag associativity, coproducts,
primitive projection, \(R'\to R\) transitions, protected integration,
or aggregate hybrid-carrier population.
\end{proposition}

\begin{proof}
By Proposition~\ref{prop:quotient-after-correspondence-preserves-base-change},
the quotient square for \(p_e\) is cartesian.  Hence the pullback of
the descended external product along the quotient middle stack is the
descent of
\[
p_e^*(K_1\boxtimes\cdots\boxtimes K_k).
\]
Proposition~\ref{prop:quotient-after-correspondence-preserves-thom-sebastiani}
then identifies the descended middle coefficient with
\[
\overline L_e=
\overline{\operatorname{TS}}^{\mathrm{red}}_{\overline e}
\overline p_e^{\,*}
(\overline K_1\boxtimes\cdots\boxtimes\overline K_k).
\]
The quotient-presented compact-support row for the target map gives
the specified isomorphism
\[
\delta_{q,e}\colon
\overline{\,q_{e,!}^{\mathrm{cs}}L_e\,}
\xrightarrow{\sim}
\overline q_{e,!}^{\mathrm{cs}}\overline L_e
\]
on \([Y_0/E]\).

Applying the Borel--Moore chain functor
of Proposition~\ref{prop:quotient-after-correspondence-borel-moore-chain-functor}
to \(\delta_{q,e}\) gives the displayed chain isomorphism.  Its
source is, by definition,
\[
Q_{E,R}\,
q_{e,!}^{\mathrm{cs}}\operatorname{TS}^{\mathrm{red}}_e p_e^*
(K_1,\ldots,K_k),
\]
and its target is the reduced operation
\[
\overline q_{e,!}^{\mathrm{cs}}\,
\overline{\operatorname{TS}}^{\mathrm{red}}_{\overline e}
\overline p_e^{\,*}
(Q_{E,R}K_1,\ldots,Q_{E,R}K_k).
\]
Uniqueness follows from full faithfulness of constructible descent on
quotient stacks.  The proof compares one retained operation with its
reduced quotient operation; it does not compare two parenthesized
operations, a product with a coproduct, a primitive projection, or two
HN heights.
\end{proof}

The packet
\[
\texttt{certificates/hybrid/quotient\_after\_correspondence\_theta\_mu\_comparison}
\]
verifies
\(\texttt{QUOTIENT\_AFTER\_CORRESPONDENCE\_THETA\_MU\_COMPARISON\_VERIFIED}\):
Proposition~\ref{prop:quotient-after-correspondence-theta-mu-comparison}
constructs the operation-comparison isomorphisms
\(\theta^Q_{\mu,R}\) for quotient-presented finite-height operation
rows and proves \(\mathfrak o^{Q,\theta}_R=0\) on those rows.  This
packet imports associativity compatibility from
Proposition~\ref{prop:quotient-after-correspondence-theta-mu-associativity}.
It imports coproduct compatibility from
Proposition~\ref{prop:quotient-after-correspondence-theta-mu-coproduct}.
It imports primitive-projection compatibility from
Proposition~\ref{prop:quotient-after-correspondence-theta-mu-primitives}.
It imports the HN-transition compatibility of \(Q_{E,R}\) from
Proposition~\ref{prop:quotient-after-correspondence-hn-transition-compatibility}.
It does not prove Hall-operation transition compatibility, protected
integration, or aggregate population.

\begin{proposition}[Associativity compatibility of the quotient comparison maps]
\label{prop:quotient-after-correspondence-theta-mu-associativity}
\textnormal{[proved for the eight two-step flag associativity rows]}
Fix \(R\), and let
\[
w=\epsilon_1\epsilon_2\epsilon_3\in\{\mathrm L,\mathrm W\}^3
\]
be one of the eight local/wrapped words with retained ordered colours
\((\xi_1,\xi_2,\xi_3)\).  Assume the two-step flag stack
\(\mathfrak F^{(2),w}_{\xi_1,\xi_2,\xi_3,R}\) and the wordwise
associativity row of
Propositions~\ref{prop:lll-associativity}--\ref{prop:www-associativity}.
Let
\[
A_w=
\mu_{\xi_1+\xi_2,\xi_3,R}
\circ(\mu_{\xi_1,\xi_2,R}\otimes 1),
\qquad
B_w=
\mu_{\xi_1,\xi_2+\xi_3,R}
\circ(1\otimes\mu_{\xi_2,\xi_3,R})
\]
be the two unreduced parenthesized operations, and let
\[
a_w:A_w\xrightarrow{\sim}B_w
\]
be the associativity isomorphism represented by the common flag stack.
Let \(\overline A_w,\overline B_w\), and
\(\overline a_w:\overline A_w\xrightarrow{\sim}\overline B_w\), be
the corresponding reduced operations after quotient-after-
correspondence.  If every binary and composite operation occurring in
\(A_w\) and \(B_w\) carries the quotient-presented data required in
Proposition~\ref{prop:quotient-after-correspondence-theta-mu-comparison},
then the square
\[
\begin{tikzcd}
Q_{E,R}A_w \ar[r,"\Theta^L_w"]
\ar[d,"Q_{E,R}(a_w)"']&
\overline A_w\,(Q_{E,R}^{\boxtimes3})
\ar[d,"\overline a_w"]\\
Q_{E,R}B_w \ar[r,"\Theta^R_w"']&
\overline B_w\,(Q_{E,R}^{\boxtimes3})
\end{tikzcd}
\]
commutes.  Here \(\Theta^L_w\) and \(\Theta^R_w\) are obtained by
composing the binary \(\theta^Q_{\mu,R}\) maps of
Proposition~\ref{prop:quotient-after-correspondence-theta-mu-comparison}
with the composition comparison of
Proposition~\ref{prop:quotient-after-correspondence-preserves-composition}.
Consequently the quotient two-step flag associativity component
vanishes:
\[
\mathfrak o^{Q,\mathrm{flag}}_R=0 .
\]
This proves only compatibility of \(\theta^Q_{\mu,R}\) with the
eight length-three associativity rows.  It does not prove coproduct
compatibility, primitive compatibility, \(R'\to R\) transition
compatibility, protected-integration compatibility, or aggregate
hybrid-carrier population.
\end{proposition}

\begin{proof}
For a fixed \(w\), both parenthesizations are represented before
quotient by the same retained two-step flag stack
\[
\mathfrak F^{(2),w}_{\xi_1,\xi_2,\xi_3,R}.
\]
The maps \(\lambda^w\) and \(\rho^w\) of
Proposition~\ref{prop:eight-two-step-flag-stacks} identify this flag
stack with the two iterated correspondence fibre products used by
\(A_w\) and \(B_w\).  The wordwise associativity propositions state
that the two pull-push operations obtained from this common flag stack
agree after the required base-change, projection-formula, and reduced
Thom--Sebastiani identifications.

Apply the quotient-after-correspondence construction to the same
diagram.  Proposition~\ref{prop:quotient-after-correspondence-preserves-composition}
identifies the quotient of each iterated correspondence fibre product
with the corresponding iterated quotient correspondence.
Proposition~\ref{prop:quotient-after-correspondence-preserves-base-change}
identifies the quotient of the common flag square with the reduced
cartesian square, and
Proposition~\ref{prop:quotient-after-correspondence-preserves-thom-sebastiani}
descends the coefficient transports.  Thus the reduced associator
\(\overline a_w\) is the image of \(a_w\) under quotient descent.

The maps \(\Theta^L_w\) and \(\Theta^R_w\) are induced by the same
quotient-presented compact-support pushforward descent isomorphisms
\(\delta_{q,e}\) used in
Proposition~\ref{prop:quotient-after-correspondence-theta-mu-comparison},
pulled back to the common quotient flag stack.  Full faithfulness of
constructible descent on quotient stacks makes the two morphisms in
the displayed square equal after pullback to
\(\mathfrak F^{(2),w}_{\xi_1,\xi_2,\xi_3,R}\).  Since descent is
faithful, the square commutes on the quotient stack.  The argument is
wordwise for the eight length-three words; it does not involve the
Hall coproduct, primitive projection, finite-height transitions, or
protected integration.
\end{proof}

The packet
\[
\texttt{certificates/hybrid/quotient\_after\_correspondence\_theta\_mu\_associativity}
\]
verifies
\(\texttt{QUOTIENT\_AFTER\_CORRESPONDENCE\_THETA\_MU\_ASSOCIATIVITY\_VERIFIED}\):
Proposition~\ref{prop:quotient-after-correspondence-theta-mu-associativity}
proves that the quotient comparison maps \(\theta^Q_{\mu,R}\) commute
with the eight wordwise two-step flag associativity rows, hence
\(\mathfrak o^{Q,\mathrm{flag}}_R=0\) for those rows.  This packet
imports coproduct compatibility from
Proposition~\ref{prop:quotient-after-correspondence-theta-mu-coproduct}.
It imports primitive-projection compatibility from
Proposition~\ref{prop:quotient-after-correspondence-theta-mu-primitives}.
It imports the HN-transition compatibility of \(Q_{E,R}\) from
Proposition~\ref{prop:quotient-after-correspondence-hn-transition-compatibility}.
It does not prove Hall-operation transition compatibility, protected
integration, or aggregate population.

\begin{proposition}[Coproduct compatibility of the quotient comparison maps]
\label{prop:quotient-after-correspondence-theta-mu-coproduct}
\textnormal{[proved for supplied quotient-presented Hall bialgebra rows]}
Fix \(R\).  Let \(\mu_R\) be a retained binary Hall operation row, and
let
\[
\Delta_R:
\mathcal H_R^{E,\mathrm{hyb}}
\longrightarrow
\mathcal H_R^{E,\mathrm{hyb}}\boxtimes
\mathcal H_R^{E,\mathrm{hyb}}
\]
be a retained collision coproduct row.  Assume the product
correspondence, the splitting correspondence defining \(\Delta_R\),
and the four-leg bialgebra comparison correspondence carry the
quotient-presented compact-support, base-change, and reduced
Thom--Sebastiani data required in
Proposition~\ref{prop:quotient-after-correspondence-theta-mu-comparison}.
Assume also that the source finite Hall bialgebra identity
\[
\Delta_R\mu_R
=
(\mu_R\boxtimes\mu_R)(1\boxtimes\tau\boxtimes1)
(\Delta_R\boxtimes\Delta_R)
\]
has defect rank zero on that retained correspondence, with Koszul sign
braiding \(\tau\).  Let \(\overline\mu_R\), \(\overline\Delta_R\), and
the reduced bialgebra identity be the quotient-after-correspondence
images of these rows.

Then the two morphisms
\[
(Q_{E,R}\boxtimes Q_{E,R})\Delta_R\mu_R
\rightrightarrows
(\overline\mu_R\boxtimes\overline\mu_R)
(1\boxtimes\tau\boxtimes1)
(\overline\Delta_R\boxtimes\overline\Delta_R)
Q_{E,R}^{\boxtimes2}
\]
coincide.  The first morphism applies the quotient comparison for the
coproduct row, then \(\overline\Delta_R(\theta^Q_{\mu,R})\), and then
the reduced bialgebra identity.  The second applies the source
bialgebra identity, then the two coproduct comparison maps and the two
product comparison maps \(\theta^Q_{\mu,R}\).  Consequently the
quotient product-coproduct compatibility defect vanishes:
\[
\mathfrak o^{Q,\Delta}_R=0 .
\]
This proves only compatibility of \(\theta^Q_{\mu,R}\) with supplied
finite coproduct/bialgebra rows.  It does not construct the compact
Hall coproduct in the aggregate carrier, prove primitive compatibility,
\(R'\to R\) transition compatibility, protected integration, or
aggregate hybrid-carrier population.
\end{proposition}

\begin{proof}
The source bialgebra identity is represented by the retained
four-leg extension-splitting correspondence: both
\(\Delta_R\mu_R\) and
\((\mu_R\boxtimes\mu_R)(1\boxtimes\tau\boxtimes1)
(\Delta_R\boxtimes\Delta_R)\) are its two pull-push descriptions.
The zero-defect source row asserts equality after the listed
base-change, projection-formula, compact-support, and
Thom--Sebastiani identifications.

Apply quotient-after-correspondence to the same correspondence.
Proposition~\ref{prop:quotient-after-correspondence-preserves-composition}
identifies the quotient of the composite correspondence with the
composite of the quotient correspondences.
Proposition~\ref{prop:quotient-after-correspondence-preserves-base-change}
identifies the quotient of the bialgebra comparison squares with the
reduced cartesian squares, and
Proposition~\ref{prop:quotient-after-correspondence-preserves-thom-sebastiani}
descends the product and inverse-coproduct coefficient transports.
Thus the reduced bialgebra identity is the quotient image of the
source identity.

The comparison maps for the product and coproduct rows are induced by
the same quotient-presented compact-support descent isomorphisms
\(\delta_{q,e}\) used in
Proposition~\ref{prop:quotient-after-correspondence-theta-mu-comparison},
with source and target legs interchanged for the coproduct
correspondence.  After pullback to the common equivariant
four-leg correspondence, the two displayed composites are the same
morphism because the source bialgebra square has defect rank zero.
Full faithfulness of constructible descent on quotient stacks then
identifies the two descended morphisms on the quotient
correspondence.  The argument uses only the supplied finite
product-coproduct square; it does not form primitive kernels,
finite-height transition maps, or the protected integration map.
\end{proof}

The packet
\[
\texttt{certificates/hybrid/quotient\_after\_correspondence\_theta\_mu\_coproduct}
\]
verifies
\(\texttt{QUOTIENT\_AFTER\_CORRESPONDENCE\_THETA\_MU\_COPRODUCT\_VERIFIED}\):
Proposition~\ref{prop:quotient-after-correspondence-theta-mu-coproduct}
proves that the quotient comparison maps \(\theta^Q_{\mu,R}\) commute
with supplied finite Hall coproduct/bialgebra rows.  This packet
imports primitive-projection compatibility from
Proposition~\ref{prop:quotient-after-correspondence-theta-mu-primitives}.
It imports the HN-transition compatibility of \(Q_{E,R}\) from
Proposition~\ref{prop:quotient-after-correspondence-hn-transition-compatibility}.
It does not prove Hall-coproduct transition compatibility, protected
integration, or aggregate population.

\begin{proposition}[Primitive-projection compatibility of the quotient comparison maps]
\label{prop:quotient-after-correspondence-theta-mu-primitives}
\textnormal{[proved for supplied finite primitive kernel and projection rows]}
Fix \(R\).  Let
\[
\widetilde\Delta_R
=\Delta_R-\eta_R\epsilon_R\boxtimes\mathrm{id}
-\mathrm{id}\boxtimes\eta_R\epsilon_R,
\qquad
\widetilde{\overline\Delta}_R
=\overline\Delta_R-\overline\eta_R\overline\epsilon_R\boxtimes\mathrm{id}
-\mathrm{id}\boxtimes\overline\eta_R\overline\epsilon_R
\]
be the reduced coproducts attached to the retained finite Hall
bialgebra rows.  Assume the finite primitive rows supply the kernel
subobjects and idempotent projections
\[
P_R^E=\ker(\widetilde\Delta_R)\cap\ker(\epsilon_R),
\qquad
\overline P_R=\ker(\widetilde{\overline\Delta}_R)
\cap\ker(\overline\epsilon_R),
\]
\[
\pi_R^{\Prim}\colon\mathcal H_R^{E,\mathrm{hyb}}\to P_R^E,
\qquad
\overline\pi_R^{\Prim}\colon
\mathcal H_R^{\mathrm{red},\mathrm{hyb}}\to\overline P_R,
\]
and assume that the quotient-after-correspondence descent of the
unit, counit, and reduced coproduct rows identifies
\[
Q_{E,R}P_R^E\simeq \overline P_R,\qquad
Q_{E,R}\pi_R^{\Prim}
=\overline\pi_R^{\Prim}Q_{E,R}
\]
on the supplied primitive kernel/projection row.  For a retained
binary Hall product row \(\mu_R\), let \(\theta^Q_{\mu,R}\) be the
comparison map of
Proposition~\ref{prop:quotient-after-correspondence-theta-mu-comparison},
and assume the coproduct compatibility of
Proposition~\ref{prop:quotient-after-correspondence-theta-mu-coproduct}
for the same product-coproduct row.

Then \(\theta^Q_{\mu,R}\) commutes with the primitive projection on
the output Hall object:
\[
\overline\pi_R^{\Prim}\,
\overline\mu_R\,Q_{E,R}^{\boxtimes2}\,\theta^Q_{\mu,R}
=
\theta^{Q,\Prim}_{\mu,R}\,
Q_{E,R}\pi_R^{\Prim}\,\mu_R ,
\]
where \(\theta^{Q,\Prim}_{\mu,R}\colon Q_{E,R}P_R^E\to
\overline P_R\) is the quotient identification on the primitive
summand.  Equivalently, the square
\[
\begin{tikzcd}
Q_{E,R}\mu_R \ar[r,"\theta^Q_{\mu,R}"]
\ar[d,"Q_{E,R}\pi_R^{\Prim}"'] &
\overline\mu_R Q_{E,R}^{\boxtimes2}
\ar[d,"\overline\pi_R^{\Prim}"]\\
Q_{E,R}P_R^E \ar[r,"\theta^{Q,\Prim}_{\mu,R}"'] &
\overline P_R
\end{tikzcd}
\]
commutes.  Consequently the quotient primitive-projection defect
vanishes on the supplied finite rows:
\[
\mathfrak o^{Q,\Prim}_R=0 .
\]
This proves only compatibility of \(\theta^Q_{\mu,R}\) with the
supplied primitive kernel/projection rows.  It does not construct
source compact Hall primitive matrices, prove \(R'\to R\) transition
compatibility, prove protected-integration compatibility, or populate the
aggregate hybrid carrier.
\end{proposition}

\begin{proof}
The primitive subobject is the intersection of two kernels: the kernel
of the reduced coproduct and the kernel of the counit.  By
Proposition~\ref{prop:quotient-after-correspondence-theta-mu-coproduct},
the quotient comparison sends the source reduced-coproduct equation to
the reduced quotient equation.  The primitive row additionally
supplies the unit and counit descent, so \(Q_{E,R}\) sends
\(\ker(\widetilde\Delta_R)\cap\ker(\epsilon_R)\) into
\(\ker(\widetilde{\overline\Delta}_R)\cap
\ker(\overline\epsilon_R)\).  This gives the kernel identification
\(Q_{E,R}P_R^E\simeq\overline P_R\) recorded in the hypothesis.

The two composites in the displayed square are morphisms obtained from
the same quotient-presented binary Hall correspondence after adjoining
the supplied primitive idempotent on the output.  On the coproduct
kernel their reduced-coproduct components agree by the row-254
coproduct square; on the augmentation kernel their counit components
agree by the supplied counit descent.  Full faithfulness of
constructible descent on the quotient stack identifies the two
idempotent-split morphisms.  Hence the primitive projection square
commutes and the finite primitive defect has rank zero.
\end{proof}

The packet
\[
\texttt{certificates/hybrid/quotient\_after\_correspondence\_theta\_mu\_primitives}
\]
verifies
\(\texttt{QUOTIENT\_AFTER\_CORRESPONDENCE\_THETA\_MU\_PRIMITIVES\_VERIFIED}\):
Proposition~\ref{prop:quotient-after-correspondence-theta-mu-primitives}
proves that the quotient comparison maps \(\theta^Q_{\mu,R}\) commute
with the supplied finite primitive kernel and projection rows.  This
packet imports the comparison, associativity, and coproduct packets.
It imports the HN-transition compatibility of \(Q_{E,R}\) from
Proposition~\ref{prop:quotient-after-correspondence-hn-transition-compatibility}.
It does not prove HN-transition preservation of primitive subspaces,
protected-integration compatibility, or aggregate population.

\begin{proposition}[HN-transition compatibility of the quotient descent functor]
\label{prop:quotient-after-correspondence-hn-transition-compatibility}
\textnormal{[proved for supplied finite HN transition rows]}
Let \(R\le R'\).  Suppose the finite quotient-after-correspondence
data at heights \(R'\) and \(R\) are supplied, and suppose every
retained object row \(Y_{R'}\to Y_R\) in the HN transition system is
\(E\)-equivariant.  Write
\[
T^E_{R'R}:\mathsf{BM}^{E,\mathrm{hyb}}_{R'}
\longrightarrow
\mathsf{BM}^{E,\mathrm{hyb}}_{R},
\qquad
\overline T_{R'R}:\mathsf{BM}^{\mathrm{red},\mathrm{hyb}}_{R'}
\longrightarrow
\mathsf{BM}^{\mathrm{red},\mathrm{hyb}}_{R}
\]
for the supplied source and reduced Borel--Moore transition functors.
Assume that each transition row is quotient-presented: the square
\[
\begin{array}{ccc}
Y_{R'} & \longrightarrow & Y_R\\
\downarrow && \downarrow\\
\left[Y_{R'}/E\right] & \longrightarrow & \left[Y_R/E\right]
\end{array}
\]
commutes, the lower map is the quotient of the upper \(E\)-equivariant
map, and the vanishing-cycle coefficient transport at height
\(R'\to R\) descends to the quotient coefficient transport.  Equivalently,
the source transition chain map and the reduced transition chain map
are induced from the same equivariant transition correspondence.

Then the quotient descent functors commute with the HN transition:
\[
\overline T_{R'R}\,Q_{E,R'}
=
Q_{E,R}\,T^E_{R'R}
\]
as functors on the quotient-presented finite Borel--Moore chain
category.  Consequently the transition component of the quotient
residual vanishes on the supplied rows:
\[
\mathfrak o^{Q,\mathrm{tr}}_{R'R}=0.
\]
This proves only the HN-transition compatibility of \(Q_{E,R}\) on
quotient-presented finite transition rows.  It does not prove that HN
transitions preserve the Hall product, Hall coproduct, primitive
subspaces, protected-integration compatibility, or aggregate
hybrid-carrier population.
\end{proposition}

\begin{proof}
For a retained transition object \(Y_{R'}\to Y_R\), the map is
\(E\)-equivariant.  The quotient-stack universal property therefore
gives a unique reduced transition map
\[
[Y_{R'}/E]\longrightarrow [Y_R/E]
\]
for which the quotient square commutes.  The coefficient transport is
assumed to be an \(E\)-equivariant morphism between quotient-presented
vanishing-cycle coefficients, and hence descends uniquely to the
corresponding coefficient transport on the two quotient stacks.

Apply the Borel--Moore chain construction of
Proposition~\ref{prop:quotient-after-correspondence-borel-moore-chain-functor}.
The composite
\(\overline T_{R'R}Q_{E,R'}\) first replaces the source object by the
quotient stack \([Y_{R'}/E]\) and then applies the descended transition
map to \([Y_R/E]\).  The composite \(Q_{E,R}T^E_{R'R}\) first applies
the equivariant source transition and then descends the result to the
same quotient stack \([Y_R/E]\).  Both composites are therefore the
same functor
\[
C_*^{\mathrm{BM}}([Y_{R'}/E],\overline K_{R'})
\longrightarrow
C_*^{\mathrm{BM}}([Y_R/E],\overline K_R)
\]
with the same descended coefficient morphism.  Full faithfulness of
constructible descent on quotient stacks identifies the two chain maps
and their composites over successive height restrictions.  Thus the
finite transition defect of \(Q_{E,R}\) has rank zero.
\end{proof}

The packet
\[
\texttt{certificates/hybrid/quotient\_after\_correspondence\_hn\_transition\_compatibility}
\]
verifies
\(\texttt{QUOTIENT\_AFTER\_CORRESPONDENCE\_HN\_TRANSITION\_COMPATIBILITY\_VERIFIED}\):
Proposition~\ref{prop:quotient-after-correspondence-hn-transition-compatibility}
proves that \(Q_{E,R}\) commutes with supplied finite HN transition
rows.  It does not prove Hall product, Hall coproduct,
primitive-subspace, protected-integration, or aggregate transition
compatibility.

\begin{definition}[Finite protected integration datum]
\label{def:finite-protected-integration-datum}
Fix \(R\), the reduced Borel--Moore chain object
\[
\mathcal H_R^{\mathrm{red},\mathrm{hyb}}
:=
Q_{E,R}\mathcal F_{X,\sigma,S,\le R}^{\mathrm{hyb}},
\]
and its finite charge decomposition
\[
\mathcal H_R^{\mathrm{red},\mathrm{hyb}}
=\bigoplus_{\gamma\in\Gamma_R}
\mathcal H_{\gamma,R}^{\mathrm{red},\mathrm{hyb}} .
\]
A finite protected integration datum at height \(R\) consists of:
\begin{enumerate}
\item a finite commutative monoid \(\mathsf T_R\) and labels
\(\omega_R(\gamma)\in\mathsf T_R\) for the retained charges;
\item the finite subspace
\[
\C[\mathsf T_R]_{\le R}
:=\sum_{\gamma\in\Gamma_R}\C e_{\omega_R(\gamma)}
\subset \C[\mathsf T_R];
\]
\item for every \(\gamma\in\Gamma_R\), a degree-zero chain map
\[
\int^{\mathrm{prot}}_{\gamma,R}\colon
\mathcal H_{\gamma,R}^{\mathrm{red},\mathrm{hyb}}
\longrightarrow
\C e_{\omega_R(\gamma)} .
\]
\end{enumerate}
The associated protected integration map is the finite direct sum
\[
I_R^{\mathrm{prot}}
:=
\bigoplus_{\gamma\in\Gamma_R}
\int^{\mathrm{prot}}_{\gamma,R}\colon
\mathcal H_R^{\mathrm{red},\mathrm{hyb}}
\longrightarrow
\C[\mathsf T_R]_{\le R}.
\]
Thus, for \(x\in\mathcal H_{\gamma,R}^{\mathrm{red},\mathrm{hyb}}\),
\[
I_R^{\mathrm{prot}}(x)\in \C e_{\omega_R(\gamma)}.
\]
The label \(\omega_R(\gamma)\) is not yet identified with
\(\overline\Pi_X(\gamma)\), and \(e_{\omega_R(\gamma)}\) is not yet
identified with \(q^nr^ls^m\).  Those are the Gram-degree and
\(\overline\Pi_X\)-separation rows.  This datum also does not assert
multiplicativity, comultiplicativity, primitive-projection
compatibility, HN-transition compatibility, the closed
\(K3\times E\to\mathrm{pt}\) protected trace, or a Pfaffian line.
\end{definition}

\begin{proposition}[Definition row for protected integration]
\label{prop:protected-integration-definition}
\textnormal{[proved for supplied finite chargewise integration maps]}
The datum of Definition~\ref{def:finite-protected-integration-datum}
defines a degree-zero chain map
\[
I_R^{\mathrm{prot}}\colon
Q_{E,R}\mathcal F_{X,\sigma,S,\le R}^{\mathrm{hyb}}
\longrightarrow
\C[\mathsf T_R]_{\le R}.
\]
Equivalently, the chain-map component of the protected-integration
residual vanishes on the supplied finite rows:
\[
\mathfrak o^{I,\mathrm{ch}}_R=0.
\]
This proves only the definition and well-typedness of
\(I_R^{\mathrm{prot}}\).  It does not prove product compatibility,
coproduct compatibility, primitive-projection compatibility,
HN-transition compatibility, the Gram exponent law, the use of
\(\overline\Pi_X\) rather than \(b_R^{\mathrm{geom}}\), or aggregate
hybrid-carrier population.
\end{proposition}

\begin{proof}
The charge set \(\Gamma_R\) is finite at height \(R\).  Each summand
\(\int^{\mathrm{prot}}_{\gamma,R}\) is a degree-zero chain map to the
one-dimensional complex \(\C e_{\omega_R(\gamma)}\).  The finite direct
sum of these chain maps is therefore a degree-zero chain map from the
finite direct sum
\(\bigoplus_{\gamma\in\Gamma_R}
\mathcal H_{\gamma,R}^{\mathrm{red},\mathrm{hyb}}\)
to the finite subspace of the monoid algebra spanned by the labels
\(e_{\omega_R(\gamma)}\).  This map is exactly
\(I_R^{\mathrm{prot}}\).  No Hall operation, coproduct, primitive
idempotent, transition functor, Gram-coordinate comparison, or closed
trace functor enters the construction.
\end{proof}

The packet
\[
\texttt{certificates/hybrid/protected\_integration\_definition}
\]
verifies
\(\texttt{PROTECTED\_INTEGRATION\_DEFINITION\_VERIFIED}\):
Proposition~\ref{prop:protected-integration-definition} defines
\(I_R^{\mathrm{prot}}\) as the finite chargewise direct sum of supplied
degree-zero protected integration maps.  Product compatibility is
supplied separately by
Proposition~\ref{prop:protected-integration-product-compatibility}.
Coproduct compatibility is supplied separately by
Proposition~\ref{prop:protected-integration-coproduct-compatibility}.
Primitive-projection compatibility is supplied separately by
Proposition~\ref{prop:protected-integration-primitive-compatibility}.
Gram degree is supplied separately by
Proposition~\ref{prop:protected-integration-gram-degree}.
The normal-ordered label comparison is supplied separately by
Proposition~\ref{prop:protected-integration-pi-x-separation}.  The
packet records transition compatibility as a separate open row.

\begin{proposition}[Product compatibility of protected integration]
\label{prop:protected-integration-product-compatibility}
\textnormal{[proved for supplied finite binary product rows]}
Fix \(R\) and the finite protected integration datum of
Definition~\ref{def:finite-protected-integration-datum}.  Let
\[
w=\epsilon_1\epsilon_2\in\{\mathrm L,\mathrm W\}^2
\]
be a retained binary local/wrapped word, and let
\[
\overline\mu_{w,\gamma_1,\gamma_2,R}\colon
\mathcal H_{\gamma_1,R}^{\mathrm{red},\mathrm{hyb}}
\otimes
\mathcal H_{\gamma_2,R}^{\mathrm{red},\mathrm{hyb}}
\longrightarrow
\mathcal H_{\gamma_1+\gamma_2,R}^{\mathrm{red},\mathrm{hyb}}
\]
be the corresponding reduced Hall product row after
quotient-after-correspondence and the comparison map
\(\theta^Q_{\mu,R}\).  Assume the target monoid labels satisfy the
supplied product row
\[
m_{\mathsf T_R}\bigl(
e_{\omega_R(\gamma_1)}\otimes e_{\omega_R(\gamma_2)}
\bigr)
=e_{\omega_R(\gamma_1+\gamma_2)}
\]
with defect rank zero.  Assume also that the chargewise protected
integration functionals satisfy the supplied square
\[
\begin{tikzcd}
\mathcal H_{\gamma_1,R}^{\mathrm{red},\mathrm{hyb}}
\otimes
\mathcal H_{\gamma_2,R}^{\mathrm{red},\mathrm{hyb}}
\ar[r,"\overline\mu_{w,\gamma_1,\gamma_2,R}"]
\ar[d,"\int^{\mathrm{prot}}_{\gamma_1,R}\otimes
\int^{\mathrm{prot}}_{\gamma_2,R}"']&
\mathcal H_{\gamma_1+\gamma_2,R}^{\mathrm{red},\mathrm{hyb}}
\ar[d,"\int^{\mathrm{prot}}_{\gamma_1+\gamma_2,R}"]\\
\C e_{\omega_R(\gamma_1)}\otimes
\C e_{\omega_R(\gamma_2)}
\ar[r,"m_{\mathsf T_R}"']&
\C e_{\omega_R(\gamma_1+\gamma_2)}
\end{tikzcd}
\]
with defect rank zero.  Then the protected integration map is
multiplicative on this product row:
\[
I_R^{\mathrm{prot}}\,
\overline\mu_{w,\gamma_1,\gamma_2,R}
=
m_{\mathsf T_R}\,
(I_R^{\mathrm{prot}}\otimes I_R^{\mathrm{prot}})
\]
on
\(\mathcal H_{\gamma_1,R}^{\mathrm{red},\mathrm{hyb}}\otimes
\mathcal H_{\gamma_2,R}^{\mathrm{red},\mathrm{hyb}}\).  Consequently
\[
\mathfrak o^{I,\mathrm{prod}}_{w,R}=0
\]
for every supplied finite binary product row of word \(w\).
This proves only product compatibility of
\(I_R^{\mathrm{prot}}\).  It does not prove coproduct compatibility,
primitive-projection compatibility, HN-transition compatibility, the
Gram exponent law, the use of \(\overline\Pi_X\) rather than
\(b_R^{\mathrm{geom}}\), the level-\(\mathsf Z\) protected trace, or a
Pfaffian line.
\end{proposition}

\begin{proof}
Restrict \(I_R^{\mathrm{prot}}\) to the two source charge summands and
to the target charge summand.  By
Definition~\ref{def:finite-protected-integration-datum}, these
restrictions are exactly
\[
\int^{\mathrm{prot}}_{\gamma_1,R},\qquad
\int^{\mathrm{prot}}_{\gamma_2,R},\qquad
\int^{\mathrm{prot}}_{\gamma_1+\gamma_2,R}.
\]
The displayed component square is therefore precisely the restriction
of the global multiplicativity square for \(I_R^{\mathrm{prot}}\) to
the summand
\(\mathcal H_{\gamma_1,R}^{\mathrm{red},\mathrm{hyb}}\otimes
\mathcal H_{\gamma_2,R}^{\mathrm{red},\mathrm{hyb}}\).  Since
\(\Gamma_R\) is finite, the global product square is the finite direct
sum of these component squares.  Each component has defect rank zero by
hypothesis, and the target label multiplication identifies the
one-dimensional target line with
\(\C e_{\omega_R(\gamma_1+\gamma_2)}\).  Hence the product defect of
protected integration vanishes on the supplied row.  The proof uses no
coproduct, primitive projection, transition functor, Gram-coordinate
identification, closed trace functor, or Pfaffian line.
\end{proof}

The packet
\[
\texttt{certificates/hybrid/protected\_integration\_product\_compatibility}
\]
verifies
\(\texttt{PROTECTED\_INTEGRATION\_PRODUCT\_COMPATIBILITY\_VERIFIED}\):
Proposition~\ref{prop:protected-integration-product-compatibility}
proves that \(I_R^{\mathrm{prot}}\) respects every supplied finite
binary local/wrapped product row.  Coproduct compatibility is supplied
separately by
Proposition~\ref{prop:protected-integration-coproduct-compatibility}.
Primitive-projection compatibility is supplied separately by
Proposition~\ref{prop:protected-integration-primitive-compatibility}.
Gram degree is supplied separately by
Proposition~\ref{prop:protected-integration-gram-degree}.
The normal-ordered label comparison is supplied separately by
Proposition~\ref{prop:protected-integration-pi-x-separation}.  The
packet records transition compatibility as a separate open row.

\begin{proposition}[Coproduct compatibility of protected integration]
\label{prop:protected-integration-coproduct-compatibility}
\textnormal{[proved for supplied finite binary coproduct rows]}
Fix \(R\) and the finite protected integration datum of
Definition~\ref{def:finite-protected-integration-datum}.  Let
\[
w=\epsilon_1\epsilon_2\in\{\mathrm L,\mathrm W\}^2
\]
be a retained binary local/wrapped output word, and let
\[
\overline\Delta_{w,\gamma_1,\gamma_2,R}\colon
\mathcal H_{\gamma_1+\gamma_2,R}^{\mathrm{red},\mathrm{hyb}}
\longrightarrow
\mathcal H_{\gamma_1,R}^{\mathrm{red},\mathrm{hyb}}
\otimes
\mathcal H_{\gamma_2,R}^{\mathrm{red},\mathrm{hyb}}
\]
be the corresponding reduced Hall coproduct row after
quotient-after-correspondence and the coproduct comparison row
\(\theta^Q_{\Delta,R}\).  Assume the target label coalgebra has the
supplied component row
\[
\Delta_{\mathsf T_R,w}\bigl(e_{\omega_R(\gamma_1+\gamma_2)}\bigr)
=e_{\omega_R(\gamma_1)}\otimes e_{\omega_R(\gamma_2)}
\]
with defect rank zero.  Assume also that the chargewise protected
integration functionals satisfy the supplied square
\[
\begin{tikzcd}
\mathcal H_{\gamma_1+\gamma_2,R}^{\mathrm{red},\mathrm{hyb}}
\ar[r,"\overline\Delta_{w,\gamma_1,\gamma_2,R}"]
\ar[d,"\int^{\mathrm{prot}}_{\gamma_1+\gamma_2,R}"']&
\mathcal H_{\gamma_1,R}^{\mathrm{red},\mathrm{hyb}}
\otimes
\mathcal H_{\gamma_2,R}^{\mathrm{red},\mathrm{hyb}}
\ar[d,"\int^{\mathrm{prot}}_{\gamma_1,R}\otimes
\int^{\mathrm{prot}}_{\gamma_2,R}"]\\
\C e_{\omega_R(\gamma_1+\gamma_2)}
\ar[r,"\Delta_{\mathsf T_R,w}"']&
\C e_{\omega_R(\gamma_1)}\otimes
\C e_{\omega_R(\gamma_2)}
\end{tikzcd}
\]
with defect rank zero.  Then the protected integration map is
comultiplicative on this coproduct row:
\[
(I_R^{\mathrm{prot}}\otimes I_R^{\mathrm{prot}})\,
\overline\Delta_{w,\gamma_1,\gamma_2,R}
=
\Delta_{\mathsf T_R,w}\,I_R^{\mathrm{prot}}
\]
on
\(\mathcal H_{\gamma_1+\gamma_2,R}^{\mathrm{red},\mathrm{hyb}}\).
Consequently
\[
\mathfrak o^{I,\mathrm{co}}_{w,R}=0
\]
for every supplied finite binary coproduct row of word \(w\).
This proves only coproduct compatibility of
\(I_R^{\mathrm{prot}}\).  It does not prove primitive-projection
compatibility, HN-transition compatibility, the Gram exponent law, the
use of \(\overline\Pi_X\) rather than \(b_R^{\mathrm{geom}}\), the
level-\(\mathsf Z\) protected trace, or a Pfaffian line.
\end{proposition}

\begin{proof}
Restrict \(I_R^{\mathrm{prot}}\) to the source charge summand and to
the two target charge summands.  By
Definition~\ref{def:finite-protected-integration-datum}, these
restrictions are
\[
\int^{\mathrm{prot}}_{\gamma_1+\gamma_2,R},\qquad
\int^{\mathrm{prot}}_{\gamma_1,R},\qquad
\int^{\mathrm{prot}}_{\gamma_2,R}.
\]
The displayed component square is therefore the restriction of the
global coproduct square for \(I_R^{\mathrm{prot}}\) to the charge
\(\gamma_1+\gamma_2\) and the output split
\((\gamma_1,\gamma_2)\).  Since \(\Gamma_R\) is finite, the supplied
global coproduct square is the finite direct sum of these component
squares over the retained split rows.  Each component has defect rank
zero by hypothesis, and the target label coproduct identifies the
one-dimensional source label line with
\(\C e_{\omega_R(\gamma_1)}\otimes\C e_{\omega_R(\gamma_2)}\).  Hence
the coproduct defect of protected integration vanishes on the supplied
row.  The proof uses no primitive projection, transition functor,
Gram-coordinate identification, closed trace functor, or Pfaffian line.
\end{proof}

The packet
\[
\texttt{certificates/hybrid/protected\_integration\_coproduct\_compatibility}
\]
verifies
\(\texttt{PROTECTED\_INTEGRATION\_COPRODUCT\_COMPATIBILITY\_VERIFIED}\):
Proposition~\ref{prop:protected-integration-coproduct-compatibility}
proves that \(I_R^{\mathrm{prot}}\) respects every supplied finite
binary local/wrapped coproduct row.  Primitive-projection compatibility
is supplied separately by
Proposition~\ref{prop:protected-integration-primitive-compatibility}.
Gram degree is supplied separately by
Proposition~\ref{prop:protected-integration-gram-degree}.
The normal-ordered label comparison is supplied separately by
Proposition~\ref{prop:protected-integration-pi-x-separation}.  The
packet records transition compatibility as a separate open row.

\begin{proposition}[Primitive-projection compatibility of protected integration]
\label{prop:protected-integration-primitive-compatibility}
\textnormal{[proved for supplied finite primitive projection rows]}
Fix \(R\) and the finite protected integration datum of
Definition~\ref{def:finite-protected-integration-datum}.  Assume the
quotient primitive-projection row of
Proposition~\ref{prop:quotient-after-correspondence-theta-mu-primitives}
has supplied the reduced primitive subobject
\[
\overline P_R
=\ker(\widetilde{\overline\Delta}_R)\cap
\ker(\overline\epsilon_R)
\subset \mathcal H_R^{\mathrm{red},\mathrm{hyb}}
\]
and an idempotent projection
\[
\overline\pi_R^{\Prim}\colon
\mathcal H_R^{\mathrm{red},\mathrm{hyb}}\longrightarrow \overline P_R .
\]
Assume also that the target label coalgebra has a supplied primitive
projector
\[
\pi_{\mathsf T_R}^{\Prim}\colon
\C[\mathsf T_R]_{\le R}
\longrightarrow
\C[\mathsf T_R]_{\le R}^{\Prim}
\]
with defect rank zero, and that the protected integration functionals
satisfy the supplied square
\[
\begin{tikzcd}
\mathcal H_R^{\mathrm{red},\mathrm{hyb}}
\ar[r,"\overline\pi_R^{\Prim}"]
\ar[d,"I_R^{\mathrm{prot}}"']&
\overline P_R
\ar[d,"I_R^{\mathrm{prot}}\vert_{\overline P_R}"]\\
\C[\mathsf T_R]_{\le R}
\ar[r,"\pi_{\mathsf T_R}^{\Prim}"']&
\C[\mathsf T_R]_{\le R}^{\Prim}
\end{tikzcd}
\]
with defect rank zero.  Then \(I_R^{\mathrm{prot}}\) respects the
primitive projection:
\[
I_R^{\mathrm{prot}}\,
\overline\pi_R^{\Prim}
=
\pi_{\mathsf T_R}^{\Prim}\,I_R^{\mathrm{prot}}
\]
on \(\mathcal H_R^{\mathrm{red},\mathrm{hyb}}\), after identifying
\(I_R^{\mathrm{prot}}\vert_{\overline P_R}\) with its target primitive
subspace.  Consequently
\[
\mathfrak o^{I,\Prim}_R=0
\]
on the supplied finite primitive projection row.  This proves only
primitive-projection compatibility of \(I_R^{\mathrm{prot}}\).  It does
not prove HN-transition compatibility, the Gram exponent law, the use
of \(\overline\Pi_X\) rather than \(b_R^{\mathrm{geom}}\), the
level-\(\mathsf Z\) protected trace, or a Pfaffian line.
\end{proposition}

\begin{proof}
The reduced primitive subobject \(\overline P_R\) is already the
idempotent-split summand supplied by
Proposition~\ref{prop:quotient-after-correspondence-theta-mu-primitives};
the present statement does not reconstruct it.  Restrict the
chargewise direct sum \(I_R^{\mathrm{prot}}\) to this summand and
compose with the target primitive projector.  The displayed square is
exactly the assertion that these two composites agree.  Its defect rank
is zero by hypothesis, hence the primitive-projection component of the
protected-integration residual vanishes on the supplied finite row.  No
transition functor, Gram-coordinate identification, closed trace
functor, or Pfaffian line enters the argument.
\end{proof}

The packet
\[
\texttt{certificates/hybrid/protected\_integration\_primitive\_compatibility}
\]
verifies
\(\texttt{PROTECTED\_INTEGRATION\_PRIMITIVE\_COMPATIBILITY\_VERIFIED}\):
Proposition~\ref{prop:protected-integration-primitive-compatibility}
proves that \(I_R^{\mathrm{prot}}\) respects the supplied finite
primitive projection row.  Gram degree is supplied separately by
Proposition~\ref{prop:protected-integration-gram-degree}.
The normal-ordered label comparison is supplied separately by
Proposition~\ref{prop:protected-integration-pi-x-separation}.  The
packet records transition compatibility as a separate open row.

\begin{proposition}[Gram degree of protected integration]
\label{prop:protected-integration-gram-degree}
\textnormal{[proved for supplied finite Gram-monomial label rows]}
Fix \(R\) and the finite protected integration datum of
Definition~\ref{def:finite-protected-integration-datum}.  Assume a
finite Gram-monomial target row supplies a map
\[
g_R:\mathsf T_R\longrightarrow \Z^3,\qquad
g_R(\omega_R(\gamma))=(n_R(\gamma),l_R(\gamma),m_R(\gamma)),
\]
and a monomialization map
\[
\chi_R:\C[\mathsf T_R]_{\le R}\longrightarrow
\C[q^{\pm1},r^{\pm1},s^{\pm1}]
\]
such that
\[
\chi_R\bigl(e_{\omega_R(\gamma)}\bigr)
=q^{n_R(\gamma)}r^{l_R(\gamma)}s^{m_R(\gamma)}
\]
on every retained charge label, with defect rank zero.  Then, for
every
\[
x\in\mathcal H_{\gamma,R}^{\mathrm{red},\mathrm{hyb}},
\]
one has
\[
(\chi_R I_R^{\mathrm{prot}})(x)
\in
\C\,q^{n_R(\gamma)}r^{l_R(\gamma)}s^{m_R(\gamma)} .
\]
Equivalently, the \(s\)-degree of the protected integration image is
\[
\deg_s(\chi_R I_R^{\mathrm{prot}}(x))=m_R(\gamma)
\]
whenever the scalar coefficient is nonzero.  Consequently
\[
\mathfrak o^{I,\deg}_R=0
\]
on the supplied finite Gram-monomial rows.  This proves only the
target Gram-degree statement for protected integration.  It does not
prove that \(g_R\omega_R\) is the normal-ordered charge map
\(\overline\Pi_X\), does not prove that the third coordinate is
different from \(b_R^{\mathrm{geom}}\), does not prove HN-transition
compatibility, and does not construct the level-\(\mathsf Z\) protected
trace or a Pfaffian line.
\end{proposition}

\begin{proof}
By Definition~\ref{def:finite-protected-integration-datum},
\[
I_R^{\mathrm{prot}}(x)\in \C e_{\omega_R(\gamma)}
\]
for \(x\in\mathcal H_{\gamma,R}^{\mathrm{red},\mathrm{hyb}}\).  Applying
the supplied monomialization map \(\chi_R\) sends the one-dimensional
target line \(\C e_{\omega_R(\gamma)}\) to the one-dimensional monomial
line
\[
\C q^{n_R(\gamma)}r^{l_R(\gamma)}s^{m_R(\gamma)} .
\]
Thus the image has the displayed \(q,r,s\)-degree, and in particular
\(s\)-degree \(m_R(\gamma)\) when the coefficient is nonzero.  The
argument uses only the target label row and the chargewise definition of
\(I_R^{\mathrm{prot}}\); it does not use a comparison between
\(\omega_R\) and \(\overline\Pi_X\), does not compare with
\(b_R^{\mathrm{geom}}\), and does not use transition or trace data.
\end{proof}

The packet
\[
\texttt{certificates/hybrid/protected\_integration\_gram\_degree}
\]
verifies
\(\texttt{PROTECTED\_INTEGRATION\_GRAM\_DEGREE\_VERIFIED}\):
Proposition~\ref{prop:protected-integration-gram-degree} proves that
the supplied finite monomialization sends the protected integration
image of charge \(\gamma\) to the line
\(\C q^{n_R(\gamma)}r^{l_R(\gamma)}s^{m_R(\gamma)}\).  The packet
records transition compatibility and
leaves the normal-ordered label comparison to
Proposition~\ref{prop:protected-integration-pi-x-separation}.

\begin{proposition}[Normal-ordered Gram label of protected integration]
\label{prop:protected-integration-pi-x-separation}
\textnormal{[proved for supplied finite normal-ordered label comparison rows]}
Fix \(R\) and the finite protected integration datum of
Definition~\ref{def:finite-protected-integration-datum}.  Assume the
Gram-degree row of
Proposition~\ref{prop:protected-integration-gram-degree}.  Assume also
that a finite protected-label comparison row supplies a lift
\[
\widehat\omega_R:\Gamma_R\longrightarrow \widehat\Gamma_R
\]
to the domain of the normal-ordered Gram map and a zero-defect equality
\[
g_R\omega_R
=
\overline\Pi_X\widehat\omega_R
\colon \Gamma_R\longrightarrow \Z^3 .
\]
Write
\[
\overline\Pi_X(\widehat\omega_R(\gamma))
=
\bigl(n_R^\Pi(\gamma),l_R^\Pi(\gamma),m_R^\Pi(\gamma)\bigr).
\]
Then, for every
\[
x\in\mathcal H_{\gamma,R}^{\mathrm{red},\mathrm{hyb}},
\]
one has
\[
(\chi_R I_R^{\mathrm{prot}})(x)
\in
\C q^{n_R^\Pi(\gamma)}r^{l_R^\Pi(\gamma)}s^{m_R^\Pi(\gamma)}
\]
and
\[
\deg_s(\chi_R I_R^{\mathrm{prot}}(x))
=m_R^\Pi(\gamma)
=\operatorname{pr}_3\overline\Pi_X(\widehat\omega_R(\gamma))
\]
whenever the scalar coefficient is nonzero.  The exponent of protected
integration is therefore the normal-ordered Gram coordinate
\[
m_R^{\mathrm{Bch}}
:=\operatorname{pr}_3\overline\Pi_X\widehat\omega_R,
\]
not the geometric support degree \(b_R^{\mathrm{geom}}\).  Equivalently,
the protected-integration label residual satisfies
\[
\mathfrak o^{I,\Pi}_R=0
\]
on the supplied finite comparison rows.  This proves only the
\(\overline\Pi_X\)-label comparison for protected integration.  It does
not prove HN-transition compatibility, a branch formula relating
\(m_R^{\mathrm{Bch}}\) to \(b_R^{\mathrm{geom}}\), the
level-\(\mathsf Z\) protected trace, or a Pfaffian line.
\end{proposition}

\begin{proof}
By Proposition~\ref{prop:protected-integration-gram-degree}, the
image of a homogeneous class of charge \(\gamma\) lies in the monomial
line
\[
\C q^{n_R(\gamma)}r^{l_R(\gamma)}s^{m_R(\gamma)}
\]
where
\[
g_R(\omega_R(\gamma))=(n_R(\gamma),l_R(\gamma),m_R(\gamma)).
\]
The supplied comparison row identifies this triple with
\[
\overline\Pi_X(\widehat\omega_R(\gamma))
=
\bigl(n_R^\Pi(\gamma),l_R^\Pi(\gamma),m_R^\Pi(\gamma)\bigr).
\]
Substitution gives the displayed monomial line and the displayed
\(s\)-degree.  Proposition~\ref{prop:geometric-borcherds-degree-separation}
separates the support map
\[
b_R^{\mathrm{geom}}(\gamma)
=\operatorname{coeff}_{[E]}(p_E)_*\operatorname{cyc}_1(\gamma)
\]
from the Borcherds coordinate
\[
\operatorname{pr}_3\overline\Pi_X(\widehat\omega_R(\gamma)).
\]
Thus \(b_R^{\mathrm{geom}}\) controls the local/wrapped support split,
whereas the protected-integration monomial is labelled by the
normal-ordered Gram coordinate.  No equality or affine relation between
these two integers is used; such a relation would require the separate
branch-comparison datum of
Definition~\ref{def:geometric-borcherds-degree-comparison-datum}.
\end{proof}

The packet
\[
\texttt{certificates/hybrid/protected\_integration\_pi\_x\_separation}
\]
verifies
\(\texttt{PROTECTED\_INTEGRATION\_PI\_X\_SEPARATION\_VERIFIED}\):
Proposition~\ref{prop:protected-integration-pi-x-separation} proves
that the protected-integration monomial exponent is
\(\overline\Pi_X(\widehat\omega_R(\gamma))\), with \(s\)-degree
\(\operatorname{pr}_3\overline\Pi_X(\widehat\omega_R(\gamma))\), and
not the geometric support degree \(b_R^{\mathrm{geom}}(\gamma)\).
The packet records transition compatibility as a separate open row.

\begin{proposition}[First-window anchor residual vanishing criterion]
\label{prop:first-window-anchor-residual-vanishing-criterion}
\textnormal{[proved as a criterion; unverified for the present first-window scaffold]}
Fix a first relation-closed finite window \(W_1\) and a height \(R\)
whose wrapped colours lie in \(W_1\).  The restriction
\[
\mathfrak o^\lambda_R|_{W_1}=0
\]
is equivalent to the simultaneous vanishing on \(W_1\) of the five
component defects in Definition~\ref{def:finite-anchor-residual}:
existence of repaired anchors \(\Lambda_{\eta,R}\), unit translation
weight after the chosen finite cover, kernel-pair losslessness,
multiplicativity on the \(LW\), \(WL\), \(WW\), and \(LWL\)
unreduced correspondences, and compatibility with the \(R'\to R\)
HN-transition maps.  In the current first-window scaffold
\[
\texttt{certificates/first\_window/k3e\_relation\_closed\_window},
\]
the compact-source window rows, target-source representative rows,
correspondence theorem rows, and strict transition rows are empty.
Consequently the present data do not prove
\(\mathfrak o^\lambda_R|_{W_1}=0\).
\end{proposition}

\begin{proof}
The equivalence is the product-zero condition for the five residual
components in Definition~\ref{def:finite-anchor-residual}.  Each
component is defined by an equality of finite data on the retained
prequotient, correspondence, or transition row; the product residual
vanishes exactly when all five equalities hold.  The first-window
obstruction ledger records no compact-source
\(\texttt{window\_closure.csv}\) rows, no
\(\texttt{target\_source\_representatives.csv}\) rows, no theorem rows
in \(\texttt{first\_window\_theorems.csv}\), and no strict transition
rows.  These missing tables contain the required first-window
existence, losslessness, multiplicativity, and transition data.  Hence
the criterion is defined, but its hypotheses are not supplied by the
current first-window scaffold.
\end{proof}

The packet
\[
\texttt{certificates/hybrid/anchor\_residual\_first\_window\_obstruction}
\]
verifies
\(\texttt{ANCHOR\_RESIDUAL\_FIRST\_WINDOW\_OBSTRUCTION\_VERIFIED}\):
Proposition~\ref{prop:first-window-anchor-residual-vanishing-criterion}
gives the exact first-window zero criterion and records that the
present first-window scaffold does not contain the rows needed to prove
it.  This packet is not a proof of row \(218\); it is the obstruction
ledger for that proof.

\subsubsection*{Finite hybrid local/wrapped factorization datum}

The full carrier at finite stage is the following supplied chain datum,
with named obstruction classes that must be trivialized.

\begin{definition}[Finite hybrid local/wrapped factorization datum]
\label{def:hybrid-wrapped-factorization-datum}
Fix a Bridgeland-style stability sector $(\sigma,S)$ in the K3-surface
lineage \cite{BridgelandK3Stability,BayerMacriProjectivity} and a
finite HN height $R$.  A finite-stage hybrid datum is
\[
\mathcal F_{X,\sigma,S,\le R}^{\mathrm{hyb}}
=\bigl(
\mathcal C_{\sigma,S,\le R},\,
\operatorname{Ran}^{\mathrm{hyb}}_{R,\mathrm{pre}}(E),\,
b_R^{\mathrm{geom}},\,
\mathcal F_R^{\mathrm{loc}},\,
\mathcal G_R^{\mathrm{wr}},\,
\mathfrak E_R^{\mathrm{loc}},\,
\mathfrak E_R^{\mathrm{hyb}},\,
\mathfrak E_R^{\mathrm{wr}},\,
\mathfrak{Flag}_R^{\mathrm{hyb}},\,
\mathsf{Corr}^{E,\mathrm{hyb}}_R,\,
\lambda_R,\,
Q_{E,R},\,
\theta^Q_R,\,
o_R,\,
I_R^{\mathrm{prot}}\bigr)
\]
with the seven clauses below.  Every stack, sheaf, operation, and
coherence lives at the fixed height $R$; the inverse limit appears
only after all finite stages and transition maps have been supplied.

\begin{enumerate}
\item \emph{Finite Hall stage, hybrid prestack, and elliptic degree.}
$\mathcal C_{\sigma,S,\le R}$ is a supplied extension-closed
charge-support HN quotient of the candidate compact Hall category
under the finite-type hypothesis $\textup{[K3$\times$E-fin]}$,
with finite-type substacks
$\mathfrak M_{\gamma,R}$ for
$\gamma\in\Gamma_{\sigma,S}^{\mathrm{HN}}(R)$,
vanishing-cycle coefficients $\Phi_{\gamma,R}$, and orientation lines
$o_{\gamma,R}$.  The carrier prestack is
\(\operatorname{Ran}^{\mathrm{hyb}}_{R,\mathrm{pre}}(E)\) of
Definition~\ref{def:hybrid-ran-prestack}; its stackification is not
used unless the descent rows of that definition are supplied.  The
geometric elliptic-degree component
\[
b_R^{\mathrm{geom}}:\Gamma_{\sigma,S}^{\mathrm{HN}}(R)\longrightarrow\Z_{\ge0}
\]
is the map of Definition~\ref{def:geometric-elliptic-degree-map}; it
splits the finite charge support into local and wrapped pieces:
\begin{equation}
\label{eq:hybrid-loc-wr-windows}
\Gamma_R^{\mathrm{loc}}=\{b_R^{\mathrm{geom}}=0\},
\qquad
\Gamma_R^{\mathrm{wr}}=\{b_R^{\mathrm{geom}}>0\}.
\end{equation}
The \(b=0\) carrier is
\(\operatorname{Ran}^{\mathrm{loc}}_{R,\mathrm{pre}}(E)\) of
Definition~\ref{def:hybrid-local-stratum-prestack}; it is the
\(J=\varnothing\) subprestack of the hybrid carrier, not the later
local coefficient sheaf.
The \(b>0\) carrier is
\(\operatorname{Ran}^{\mathrm{wr}}_{R,\mathrm{pre}}(E)\) of
Definition~\ref{def:hybrid-wrapped-stratum-prestack}; it is the
\(I=\varnothing\) subprestack whose points are wrapped prequotient
families with anchors, not the later wrapped coefficient object
\(\mathcal G_R^{\mathrm{wr}}\).
Additivity on retained extensions is the datum and criterion of
Definition~\ref{def:geometric-elliptic-degree-additivity-datum} and
Proposition~\ref{prop:geometric-elliptic-degree-additivity-criterion};
terms of height $>R$ do not enter the retained quotient.
The Borcherds trace coordinate
\(m_R^{\mathrm{Bch}}=\operatorname{pr}_3\overline\Pi_X\) is not used
to form the split in \eqref{eq:hybrid-loc-wr-windows}; any relation
between \(m_R^{\mathrm{Bch}}\) and \(b_R^{\mathrm{geom}}\) requires the
comparison datum of
Definition~\ref{def:geometric-borcherds-degree-comparison-datum}.

\item \emph{Projection-finite local sector through closed
configurations.}  Over the local carrier
\(\operatorname{Ran}^{\mathrm{loc}}_{R,\mathrm{pre}}(E)\), and for a
finite set $I$, set
\[
Z_I=\bigcup_{i\in I}\{(e,(e_j)_j)\in E\times E^I\mid e=e_i\}.
\]
For $\alpha\in\Gamma_R^{\mathrm{loc}}$ the stage supplies a
closed-configuration incidence stack
\[
\mathfrak M_{\alpha,R,I}^{\mathrm{loc,cl}}
=\{(A,\mathbf e)\in\mathfrak M_{\alpha,R}\times E^I\mid
p_E(\operatorname{Supp}A)\subset Z_I(\mathbf e)\}
\]
compatible with diagonals $E^J\to E^I$, symmetric-group descent, and
the d-critical/vanishing-cycle structure.  Open-support locality is
recovered as a limit
\[
\mathcal M_{\alpha,R}^{\mathrm{loc}}(U)
\simeq\operatorname*{colim}_{I}
\bigl(\mathfrak M_{\alpha,R,I}^{\mathrm{loc,cl}}\times_{E^I}U^I\bigr),
\]
and the configuration map $\pi_{\alpha,R,I}$ produces
\[
\mathscr H_{\alpha,R,I}^{\mathrm{loc}}
=(\pi_{\alpha,R,I})_!(\Phi_{\alpha,R,I}^{\mathrm{loc}}\otimes
o_{\alpha,R,I}^{\mathrm{loc}}).
\]
For an open $U\subset E$,
\[
\mathcal F_{\alpha,R}^{\mathrm{loc}}(U)
=\bigoplus_{I/S_I}R\Gamma_c\!\left(U^I,
j_{U,I}^{!}\mathscr H_{\alpha,R,I}^{\mathrm{loc}}\right),
\]
and $\mathcal F_R^{\mathrm{loc}}(U)$ is the direct sum over
$\Gamma_R^{\mathrm{loc}}$.  Locality is $!$-restriction on closed
configuration spaces; the open-support Borel--Moore formula is a
theorem to be derived from this atlas, not the primary datum.

\item \emph{Wrapped prequotient sector and rigidification.}  Over
\(\operatorname{Ran}^{\mathrm{wr}}_{R,\mathrm{pre}}(E)\), and for
$\eta\in\Gamma_R^{\mathrm{wr}}$, the scalar-rigidified prequotient
stack of Definition~\ref{def:e-equivariant-wrapped-prequotient-stack}
\[
\mathcal M_{\eta,R}^{\mathrm{wr,rig}}
\xrightarrow{\rho_{\eta,R}}
\mathcal M_{\eta,R}^{\mathrm{wr}}\subset\mathfrak M_{\eta,R}
\]
is supplied with a fixed origin $0_E\in E\).  The map \(\rho_{\eta,R}\)
forgets the scalar rigidification; the reduced \(E\)-quotient is not
taken in this clause.  The normalized determinant construction of
Definition~\ref{def:determinant-anchor} gives the
candidate anchor
\[
\lambda^{\det}_{\eta,R}(A)
=\det Rp_{E,*}A\otimes\mathcal O_E(-\chi(A)\cdot0_E)
\in \operatorname{Pic}^0(E)\simeq E .
\]
Its translation weight is \(\chi(A)\) by
Proposition~\ref{prop:determinant-anchor-translation-weight}.
If \(\chi(A)=0\), Proposition~\ref{prop:determinant-anchor-chi-zero-degeneracy}
shows that this determinant anchor is constant along the retained
\(E\)-translation orbit and therefore has rank-one separation defect.
The degree-zero Jordan--Hoelder divisor of
Proposition~\ref{prop:determinant-anchor-extra-data-separates-rank-two}
separates the determinant-zero rank-two test pair
\(\mathcal O_E^{\oplus2}\) and \(L\oplus L^{-1}\).  Unit-weight
finite-cover normalization, global anchor losslessness, quotient
descent, stabilizer linearization, transition compatibility, and
finite-stage population are separate rows before it becomes the final
wrapped anchor
\(\lambda_{\eta,R}:\mathcal M_{\eta,R}^{\mathrm{wr,rig}}\to E\).
The finite anchor residual \(\mathfrak o^\lambda_R\) is
Definition~\ref{def:finite-anchor-residual}; its vanishing, not its
definition, is required before the repaired anchors become the final
wrapped anchors.  After the final wrapped anchor rows are supplied, set
\[
\mathcal G_{\eta,R}^{\mathrm{wr}}
=H_*^{\mathrm{BM}}\!\left(
\mathcal M_{\eta,R}^{\mathrm{wr,rig}},
\Phi_{\eta,R}^{\mathrm{wr}}\otimes o_{\eta,R}^{\mathrm{wr}}\right),
\qquad
\mathcal G_R^{\mathrm{wr}}=\bigoplus_{\eta}\mathcal G_{\eta,R}^{\mathrm{wr}}.
\]

\item \emph{Local, mixed, and wrapped extension correspondences.}  For
disjoint opens $(U_i)\subset E$, local charges $\alpha_i$, and wrapped
charges $\eta,\eta_1,\eta_2$, the finite stage contains the
correspondence stacks
\[
\mathfrak E_{\alpha_1,\ldots,\alpha_r,R}^{\mathrm{loc}}(U_1,\ldots,U_r),
\quad
\mathfrak E_{\alpha_1,\ldots,\alpha_r;\eta,R}^{\mathrm{hyb}},
\quad
\mathfrak E_{\eta_1,\eta_2,R}^{\mathrm{wr}},
\]
where the local symbol is the \(LL\) correspondence datum of
Definition~\ref{def:local-local-correspondence-datum}.  The one-sided
mixed symbol is the \(LW/WL\) correspondence datum
of Definition~\ref{def:mixed-local-wrapped-correspondence-datum} after
collecting the finite local colours into a local colour map
\(\alpha:I\to\Gamma_R^{\mathrm{loc}}\).  It is still before the
reduced \(E\)-quotient.  The wrapped symbol is the ordered \(WW\)
correspondence datum of
Definition~\ref{def:wrapped-wrapped-correspondence-datum}.  These
stacks are supplied together with the ordered \emph{two-sided} mixed
stacks of Definition~\ref{def:two-sided-mixed-correspondence-datum},
\[
\mathfrak E_{\alpha_-;\eta;\beta_+,R}^{\mathrm{hyb}}((U_i),(V_j))
\]
classifying ordered insertions of finite local supports on both sides
of one wrapped input.  The defects of closed-configuration finite-type
properness, pull-push admissibility, and Thom--Sebastiani transport
are recorded by
\[
\mathfrak o^{\mathrm{prop},\mathrm{LL}}_R,\
\mathfrak o^{\mathrm{prop},\mathrm{mix}}_R,\
\mathfrak o^{\mathrm{prop},\mathrm{WW}}_R,
\]
and aggregate residuals
$\mathfrak o^{\mathrm{conf}}_R$,
$\mathfrak o^{\mathrm{conf{-}prop}}_R$,
$\mathfrak o^{\mathrm{adm}}_R$,
$\mathfrak o^{\mathrm{TS}}_R$.  Mixed and wrapped stacks remember
source and target anchors before quotienting by the maps
\(a^{\mathrm{LW}},a^{\mathrm{WL}},a^{\mathrm{WW}}\), and
\(a^{\mathrm{LWL}}\) of
Definition~\ref{def:source-target-anchor-memory-datum}.  Replacing
this memory by a scalar Fock factor, a Hecke factor, a Borcherds trace
degree, or a post-quotient class is not the hybrid datum.  These stacks
generate a finite $E$-equivariant correspondence category
$\mathsf{Corr}^{E,\mathrm{hyb}}_R$ with operation maps
$\mu_R^\bullet=q_!\circ\operatorname{TS}_R\circ p^*$.

\item \emph{Eight-word two-step binary flag atlas.}  Write
$\mathrm L$ for a projection-finite local input and $\mathrm W$ for a
rigidified wrapped input.  For every word
$w=(\epsilon_1,\epsilon_2,\epsilon_3)\in\{\mathrm L,\mathrm W\}^3$ and
every compatible ordered triple $(\xi_1,\xi_2,\xi_3)$ of charges and
local opens, the datum supplies a two-step flag stack
$\mathfrak{Flag}^{w}_{\xi_1,\xi_2,\xi_3,R}$ classifying filtrations
with the prescribed associated graded pieces, with comparison maps to
the iterated correspondence fibre products.  The eight words
\[
\mathrm{LLL},\ \mathrm{LLW},\ \mathrm{LWL},\ \mathrm{WLL},\
\mathrm{LWW},\ \mathrm{WLW},\ \mathrm{WWL},\ \mathrm{WWW}
\]
are all part of the datum; in mixed and wrapped words the flag stack
retains initial, intermediate, and terminal anchors before $E$-quotient.  The
two parenthesizations give the same map in the finite quotient,
\[
\mu_{\xi_1+\xi_2,\xi_3,R}\circ(\mu_{\xi_1,\xi_2,R}\otimes 1)
=\mu_{\xi_1,\xi_2+\xi_3,R}\circ(1\otimes\mu_{\xi_2,\xi_3,R}),
\]
with defect classes $\mathfrak o^{\mathrm{assoc}}_{w,R}$ required to
vanish for all eight words; the chosen $2$-morphisms must satisfy the
four-input pentagon on the four-step flag-of-flags stacks.  The
higher-coloured residual
\(\mathfrak o^{\mathrm{col}}_R\) of
Definition~\ref{def:higher-coloured-hybrid-tree-category} records the
remaining coloured configuration data, units, symmetry conventions,
refinement maps, descent, and overlap coherences for all finite HN
windows.  The eight-word atlas gives only binary associativity; full
hybrid factorization requires the higher-coloured residual to vanish
in addition.

\item \emph{Quotient-after-correspondence functor and reduced
$E$-quotient.}  All stacks and correspondences above are
$E$-equivariant before quotienting.  The quotient datum supplies a
pseudofunctor
\[
\mathsf{Corr}^{E,\mathrm{hyb}}_R\longrightarrow
\mathsf{Corr}^{\mathrm{red},\mathrm{hyb}}_R
\]
as in Definition~\ref{def:quotient-after-correspondence-pseudofunctor}.
Proposition~\ref{prop:quotient-after-correspondence-preserves-composition}
supplies the coherent \(2\)-morphism for composition.  The later
Proposition~\ref{prop:quotient-after-correspondence-preserves-base-change}
supplies preservation of the stack-level cartesian base-change
squares.  The later quotient rows supply compact-support pushforward
base change, Thom--Sebastiani transport, and chain-level comparison;
applying vanishing-cycle Borel--Moore chains gives the finite
correspondence module categories
$\mathsf{BM}^{E,\mathrm{hyb}}_R$ and
$\mathsf{BM}^{\mathrm{red},\mathrm{hyb}}_R$, and the induced functor
$Q_{E,R}$ between them.  For each operation
$\mu_R^\bullet=q_!\operatorname{TS}p^*$ the datum carries a comparison
isomorphism
\[
\theta^Q_{\mu,R}:Q_{E,R}\,q_!\operatorname{TS}p^*
\xrightarrow{\sim}
\overline q_!\overline{\operatorname{TS}}\,\overline p^*
(Q_{E,R}^{\boxtimes k}),
\]
compatible with the flag-stack coherences.  The quotient residual
lanes are refined by
\[
\mathfrak o^Q_R=
\bigl(
\mathfrak o^{Q,\mathrm{form}}_R,\
\mathfrak o^{Q,\mathrm{comp}}_R,\
\mathfrak o^{Q,\mathrm{bc}}_R,\
\mathfrak o^{Q,\mathrm{TS}}_R,\
\mathfrak o^{Q,\mathrm{flag}}_R,\
\mathfrak o^{Q,\mathrm{BM}}_R,\
\mathfrak o^{Q,\theta}_R,\
\mathfrak o^{Q,\mathrm{tr}}_{R'R}
\bigr),
\]
and these components must vanish before the reduced hybrid product is
used.  A quotient-first construction (object-stack quotient followed
by a new Hall product) is \emph{not} this datum.

\item \emph{Protected integration with Gram-trace degree.}
The finite stage carries a protected integration datum
of Definition~\ref{def:finite-protected-integration-datum} and hence a
chain map
\[
I_R^{\mathrm{prot}}:Q_{E,R}\mathcal F_{X,\sigma,S,\le R}^{\mathrm{hyb}}
\longrightarrow\C[\mathsf T_R]_{\le R}.
\]
For homogeneous \(x\) of charge \(\gamma\),
\[
I_R^{\mathrm{prot}}(Q_{E,R}x)\in
\C e_{\omega_R(\gamma)} .
\]
The finite Gram-degree row identifies
\[
e_{\omega_R(\gamma)}=q^nr^ls^m,\qquad
g_R\omega_R(\gamma)=(n,l,m),
\]
in the OP/Borcherds variables of
Proposition~\ref{conv:op-chamber}.  The normal-ordered label
comparison of
Proposition~\ref{prop:protected-integration-pi-x-separation}
identifies \(g_R\omega_R\) with \(\overline\Pi_X\widehat\omega_R\) and
excludes \(b_R^{\mathrm{geom}}\) as protected-integration exponent.
The integration defects
\[
\mathfrak o^{I,\mathrm{ch}}_R,\
\{\mathfrak o^{I,\mathrm{prod}}_{w,R}\}_w,\
\mathfrak o^{I,\mathrm{co}}_R,\
\mathfrak o^{I,\Prim}_R,\
\mathfrak o^{I,\mathrm{tr}}_{R'R},\
\mathfrak o^{I,\deg}_R,\
\mathfrak o^{I,\Pi}_R
\]
record chain-map well-typedness, multiplicativity on each word,
comultiplicativity, primitive compatibility, HN transition, the
Gram-degree law, and the exclusion of \(b_R^{\mathrm{geom}}\) as the
Borcherds exponent.
\end{enumerate}

The aggregate finite obstruction tuple of the hybrid lane is
\[
\mathfrak O_{\mathrm{hyb},R}
=(\mathfrak o^{\mathrm{stage}}_R,
\mathfrak o^{\lambda}_R,
\mathfrak o^{\mathrm{corr},\mathrm{LL}}_R,
\mathfrak o^{\mathrm{corr},\mathrm{mix}}_R,
\mathfrak o^{\mathrm{corr},\mathrm{WW}}_R,
\mathfrak o^{\mathrm{conf}}_R,
\{\mathfrak o^{\mathrm{assoc}}_{w,R}\}_w,
\mathfrak o^{\mathrm{col}}_R,
\mathfrak o^{Q,\bullet}_R,
\mathfrak o^{I,\bullet}_R,
\mathfrak o^{\mathrm{tr}}_{\mathrm{hyb},R'R}).
\]
At table level the hybrid datum is supplied by
\[
\mathrm{hybrid\_ran\_prestack\_definition},\quad
\mathrm{local\_stratum\_b0\_prestack\_definition},\quad
\mathrm{wrapped\_stratum\_bpositive\_prestack\_definition},\quad
\mathrm{geometric\_elliptic\_degree\_map},\quad
\mathrm{geometric\_elliptic\_degree\_additivity},\quad
\mathrm{geometric\_borcherds\_degree\_separation},\quad
\mathrm{positive\_elliptic\_degree\_projection},\quad
\mathrm{hybrid\_base},\quad
\mathrm{local\_configurations},\quad
\mathrm{local\_local\_correspondence\_definition},\quad
\mathrm{mixed\_local\_wrapped\_correspondence\_definition},\quad
\mathrm{wrapped\_wrapped\_correspondence\_definition},\quad
\mathrm{two\_sided\_mixed\_correspondence\_definition},\quad
\mathrm{wrapped\_prequotients},\quad
\mathrm{correspondences},\quad
\mathrm{flag\_atlas},\quad
\mathrm{higher\_coloured},\quad
\mathrm{quotient\_pseudofunctor},\quad
\mathrm{protected\_integration},\quad
\mathrm{transitions},\quad
\mathrm{scalar\_firewall}.
\]
The scalar firewall excludes ordinary Ran locality alone,
determinant anchors alone, quotient-first Hall products, Fock traces,
Borcherds \(s\)-degree alone, and scalar Pfaffian products as hybrid
carrier data.
\end{definition}

\begin{proposition}[Finite consequences of the hybrid datum;
\textit{conditional on
$\mathfrak O_{\mathrm{hyb},R}=0$}]
\label{prop:hybrid-datum-scope}
Assume the data of
Definition~\ref{def:hybrid-wrapped-factorization-datum} are supplied
for all $R$ with compatible transition maps, and
$\mathfrak O_{\mathrm{hyb},R}=0$.  Then:
\begin{enumerate}
\item the $b=0$ restriction is an ordinary support-local factorization
algebra on $\operatorname{Ran}(E)$ in the sense of Beilinson--Drinfeld
\cite[\S3.4]{BD}, equivalently a Costello--Gwilliam factorization
algebra over the smooth real curve $E_{\mathrm{top}}$
\cite[Vol.~II~\S10]{CostelloGwilliam2};
\item every operation with a wrapped input lands in a wrapped
stratum;
\item the ordered mixed maps, including the two-sided
$\mu_{\alpha_-;\eta;\beta_+,R}^{\mathrm{hyb}}$, are the only Hall
actions of projection-finite insertions on wrapped sectors in
$\mathsf{Corr}^{E,\mathrm{hyb}}_R$ before reduced $E$-quotient;
\item the eight-word atlas plus pentagon coherence yield an
associative finite-stage reduced Hall product after vanishing cycles,
orientations, HN quotienting, the quotient-after-correspondence
pseudofunctor, and the $\theta^Q_{\mu,R}$;
\item for homogeneous $x$ of charge $\gamma$, the protected integration
map is chargewise typed by
\[
I_R^{\mathrm{prot}}(Q_{E,R}x)\in\C e_{\omega_R(\gamma)};
\]
the Gram-degree and \(\overline\Pi_X\)-separation rows give
\(e_{\omega_R(\gamma)}=q^nr^ls^m\) for
\(\overline\Pi_X(\widehat\omega_R(\gamma))=(n,l,m)\), while the
local/wrapped type of \(x\) is still measured by
\(b_R^{\mathrm{geom}}(\gamma)\).
\end{enumerate}
\end{proposition}

\begin{proof}
Items (i) and (ii) are part of the elliptic-degree stratification
and Proposition~\ref{prop:positive-elliptic-degree-projects}.  Item (iii) is the
definition of the mixed operations: before reduced $E$-quotient they
are the only displayed operations remembering relative $E$-position
between finite local supports and a wrapped leg.  Items (iv) and (v)
are the eight-word coherence and the protected integration clause.
\end{proof}

\begin{corollary}[No quotient-first universal hybrid repair;
\textit{proved}]
\label{cor:no-quotient-first-hybrid-repair}
A candidate finite correspondence category for the wrapped sector
gives the source object $\mathcal F_{X,\sigma,S,\le R}^{\mathrm{hyb}}$
only if it contains the unreduced local, wrapped, and ordered mixed
source stacks of
Definition~\ref{def:hybrid-wrapped-factorization-datum}, their
anchors, the eight-word flag atlas, the quotient-after-correspondence
pseudofunctor inducing $Q_{E,R}$ after Borel--Moore chains, the comparison maps
$\theta^Q_{\mu,R}$, and protected integration with
$\mathfrak O_{\mathrm{hyb},R}=0$.  A quotient of object stacks formed
\emph{before} the mixed and wrapped correspondences, or a
determinant/Fock category carrying only the trace degree
$s^{m_R^{\mathrm{Bch}}}$, does not define the hybrid source.
\end{corollary}

\begin{proof}
By Proposition~\ref{prop:positive-elliptic-degree-projects}, a
positive elliptic-degree sector with an effective support-realization
row is not a finite-support Ran sector; only the rigidified wrapped
prequotient with anchor, together with the unreduced mixed and wrapped
correspondences, retains the relative \(E\)-position used by the source
data.  A quotient-first construction loses these anchors.
\end{proof}

\subsubsection*{Comparison with the two-stage chiral construction}

The hybrid carrier is the finite \(K3\times E\) specialization datum
attached to the two-stage construction
$\textup{CY}_d\textup{-Cat}\to E_d$-$\textup{HolFA}(X)\to
\textup{ChirAlg}(C)$ of Vol III Calabi--Yau quantum groups, read here
as a chosen specialization datum rather than
a constructed functor.  The projection-finite local
stratum is the retained holomorphic \(E_3\)-prefactorization input
restricted to the smooth elliptic factor, and the wrapped stratum is the chain-level
$K3$-fusion datum that the local stratum cannot see.  Beilinson and
Drinfeld's Ran construction \cite[\S3.4--\S3.5]{BD} provides the
$b=0$ part; the $b>0$ part is the Hall--Borcherds residual.
The finite \(E_3\)-source and specialization rows
\(\mathrm{formal\_target\_charts}\),
\(\mathrm{field\_complexes}\), \(\mathrm{e3\_operations}\),
\(\mathrm{bv\_qme}\), \(\mathrm{anomaly\_framing}\),
\(\mathrm{compact\_support}\), \(\mathrm{factorization\_descent}\),
\(\mathrm{k3\_to\_e\_specialization}\), \(\mathrm{transitions}\), and
\(\mathrm{scalar\_firewall}\) are upstream of this hybrid carrier.
They are not supplied by the hybrid carrier itself.
The eight-word two-step binary flag
atlas $w\in\{\mathrm L,\mathrm W\}^3$ is the smallest closed coherence
calculus that detects the local/wrapped split through binary
associativity; it is independent of the higher-coloured residual
$\mathfrak o^{\mathrm{col}}_R$ which closes higher symmetric arities,
units, and overlaps.

\subsubsection*{The \(D0\) specialization}

The finite-stage components of
$\mathfrak o^{\mathrm{stage}}_R$ and $\mathfrak o^{\lambda}_R$ are
constructed in
Appendix~\ref{app:G-D0-discharge}.  The wrapped-leg construction and
the eight-word atlas at finite height remain conditional on the named
hybrid obstruction tuple; the $b=0$ stratum is the classical
projection-finite Ran sector of \cite{BD}.
The hybrid carrier is the projection-finite Ran stratum on the \(b=0\)
component and the wrapped stratum relative to the named finite
obstruction tuple \(\mathfrak O_{\mathrm{hyb},R}\).

Chiral operations on $\mathrm{Ran}^{\mathrm{hyb}}(E)$ require an
orientation line on the cosection-reduced d-critical heart.  On
$K3\times E$ this orientation descends through the reduced
$E$-quotient only after the retained quotient-orientation datum is
supplied; on rank-$2$ $E$-strata the Klein-four calculation gives the
edge criterion and leaves the degree-one linearization character to be
set to zero.  Entry $(\mathrm L 2)$, constructed in
\S\ref{subsec:cosection-twisted-orientation}, is this orientation datum.


\subsection{Cosection-twisted reduced d-critical orientation}
\label{subsec:cosection-twisted-orientation}

The full reduced d-critical orientation entry of the compact Hall sector
on $K3\times E$ is the second entry of the pro-object.  Its completed
form consists of a Joyce d-critical structure, a vanishing-cycle
coefficient system, a Joyce--Upmeier--type orientation line, and a
quotient-descent datum for the reduced $E$-quotient.  The rows below
separate these inputs.  Before any Pfaffian wall or Weyl transport can
be discussed, the orientation line must descend through the
cosection-reduced moduli.

The compact Hall heart is determined by the chosen stability datum
$(\sigma,S)$ of Chapter~\ref{ch:k3xe-setup}; its K3-surface component
lies in the Bridgeland--Bayer--Macri lineage
\cite{BridgelandK3Stability,BayerMacriProjectivity}.  The first finite
orientation input is the PTVV \((-1)\)-shifted symplectic form on the
retained derived substacks.  The Joyce d-critical truncation, the K3
semiregularity cosection, cosection surjectivity, the reduced
self-Ext complex, the determinant line, and its square root are
separate rows.

\begin{proposition}[PTVV form on retained finite substacks]
\label{prop:ptvv-finite-substack-symplectic}
\textnormal{[proved for retained derived enhancements; d-critical truncation not constructed]}
Let \(R\) be a finite HN height and let
\(\mathfrak M^{ss}_{R,c}\) be one of the retained finite-type
semistable substacks carrying the quasi-smooth derived enhancement
\[
\mathfrak M^{\mathrm{der}}_{R,c}
\]
of Proposition~\ref{prop:retained-quasi-smooth-derived-enhancements}.
Then \(\mathfrak M^{\mathrm{der}}_{R,c}\) carries the restricted PTVV
\((-1)\)-shifted symplectic form
\[
\omega^{\mathrm{PTVV}}_{R,c}.
\]
Its cotangent amplitude is \([-1,0]\), and its symplectic restriction
defect is zero on every retained row.

This is the unreduced derived symplectic input.  It does not construct
the Joyce d-critical truncation, the K3 semiregularity cosection,
cosection surjectivity, \(\operatorname{RHom}_{\mathrm{red}}\), the
determinant line, a square-root orientation, quotient orientation, an
O2 wall atlas, or protected integration.
\end{proposition}

\begin{proof}
The retained derived-enhancement datum of
Definition~\ref{def:retained-quasi-smooth-derived-enhancement-datum}
assigns to each retained finite-type semistable substack a derived
Artin enhancement with cotangent amplitude \([-1,0]\).  Since
\[
K_{K3\times E}\simeq\mathcal O_{K3\times E},
\]
the ambient perfect-complex moduli problem is Calabi--Yau of dimension
three.  Pantev--To\"en--Vaqui\'e--Vezzosi
\cite[Theorem~0.3]{PTVV2013} therefore supplies a
\((-1)\)-shifted symplectic form on the derived moduli of perfect
complexes.  Proposition~\ref{prop:retained-quasi-smooth-derived-enhancements}
is exactly the retained finite-row restriction of this form to
\(\mathfrak M^{\mathrm{der}}_{R,c}\), with zero Tor-amplitude and
symplectic restriction defects.  The final paragraph lists the later
orientation rows not supplied by this PTVV input.
\end{proof}

The packet
\[
\texttt{certificates/orientation/ptvv\_finite\_substack\_symplectic}
\]
verifies
\(\texttt{PTVV\_FINITE\_SUBSTACK\_SYMPLECTIC\_VERIFIED}\): it imports
the retained derived-enhancement packet, checks the three retained
finite substacks, verifies shifted degree \(-1\), cotangent amplitude
\([-1,0]\), and zero symplectic defect, and records that d-critical
truncation, K3 cosection, reduced self-Ext, determinant, orientation
square root, quotient orientation, O2, and protected integration rows
remain open.

\begin{proposition}[Joyce d-critical truncation of retained finite substacks]
\label{prop:joyce-dcritical-truncation}
\textnormal{[proved for classical truncations; vanishing cycles not constructed]}
Let \(\mathfrak M^{\mathrm{der}}_{R,c}\) be a retained finite
quasi-smooth derived enhancement carrying the PTVV form
\(\omega^{\mathrm{PTVV}}_{R,c}\) of
Proposition~\ref{prop:ptvv-finite-substack-symplectic}.  Then its
classical truncation
\[
\mathfrak M^{\mathrm{cl}}_{R,c}
:=t_0\bigl(\mathfrak M^{\mathrm{der}}_{R,c}\bigr)
\]
carries the Joyce d-critical structure
\[
s^{\mathrm{Joyce}}_{R,c}\in
H^0\bigl(\mathfrak M^{\mathrm{cl}}_{R,c},\mathcal S^0_{\mathfrak M^{\mathrm{cl}}_{R,c}}\bigr)
\]
canonically associated to \(\omega^{\mathrm{PTVV}}_{R,c}\).  The
d-critical truncation defect is zero on every retained finite row.

This row constructs only the d-critical structure on the classical
truncation.  It does not construct the BBDJS vanishing-cycle complex,
the K3 semiregularity cosection, cosection surjectivity,
\(\operatorname{RHom}_{\mathrm{red}}\), determinant lines, square-root
orientations, quotient orientation, transition compatibility, or
protected integration.
\end{proposition}

\begin{proof}
The input of Proposition~\ref{prop:ptvv-finite-substack-symplectic} is
a retained quasi-smooth derived enhancement equipped with a
\((-1)\)-shifted symplectic form.  Brav--Bussi--Dupont--Joyce--Szendr{\H{o}}i
\cite[Theorem~6.9]{BBDJS2015} associate to the classical truncation
of a \((-1)\)-shifted symplectic derived scheme or Artin chart a
canonical Joyce d-critical structure, functorial on overlaps.  Applying
this theorem to each retained finite derived row gives
\(s^{\mathrm{Joyce}}_{R,c}\) on \(t_0(\mathfrak M^{\mathrm{der}}_{R,c})\).
The retained packet records zero defect for these three truncation rows.
The final paragraph separates the d-critical truncation from the later
coefficient, cosection, reduced-obstruction, orientation, quotient, and
integration rows.
\end{proof}

The packet
\[
\texttt{certificates/orientation/joyce\_dcritical\_truncation}
\]
verifies
\(\texttt{JOYCE\_DCRITICAL\_TRUNCATION\_VERIFIED}\): it imports the
PTVV finite-substack packet, applies the BBDJS d-critical truncation
theorem to the three retained finite derived substacks, and records zero
d-critical defect.  Its firewall rows record that vanishing cycles, K3
cosections, reduced self-Ext complexes, determinant lines, square-root
orientations, quotient orientation, transition compatibility, and
protected integration remain open.

The holomorphic symplectic two-form on the K3 fibre makes the
unreduced Donaldson--Thomas theory of \(K3\times E\) vanish
identically: it generates trivial first-order deformations that
obstruct the virtual class.  The standard remedy, as in reduced K3
enumerative geometry \cite{MaulikPandharipande}, is to quotient these
out by a semiregularity cosection.

\begin{proposition}[K3 semiregularity cosection on retained finite substacks]
\label{prop:k3-semiregularity-cosection}
\textnormal{[proved as an unreduced cosection; surjectivity not proved]}
Let \(X=S\times E\), where \(S\) is a K3 surface with holomorphic
two-form \(\sigma_S\), and put \(\sigma=p_S^*\sigma_S\).  For each
retained finite classical truncation
\(\mathfrak M^{\mathrm{cl}}_{R,c}\) of
Proposition~\ref{prop:joyce-dcritical-truncation}, the obstruction
sheaf \(\operatorname{Ob}_{R,c}\) carries an
\(\mathcal O_{\mathfrak M^{\mathrm{cl}}_{R,c}}\)-linear cosection
\[
\operatorname{CS}_{R,c}:
\operatorname{Ob}_{R,c}\longrightarrow
\mathcal O_{\mathfrak M^{\mathrm{cl}}_{R,c}} .
\]
At a point represented by a perfect complex \(F\) on \(X\), and for
\(\xi\in\operatorname{Ext}^2_X(F,F)\), the fibre map is
\[
\operatorname{CS}_{R,c}|_F(\xi)
=
\int_X
\sigma\wedge
\operatorname{Tr}\bigl((\xi\otimes 1_{\Omega_X^1})\circ
\operatorname{At}(F)\bigr),
\]
where \(\operatorname{At}(F)\in
\operatorname{Ext}^1_X(F,F\otimes\Omega_X^1)\) is the Atiyah class.
The construction defect is zero on every retained finite row.

This row constructs only the unreduced K3 semiregularity cosection.  It
does not prove cosection surjectivity on retained strata, filtered Hall
additivity, the reduced self-Ext complex, perfectness of the reduced
obstruction theory, determinant lines, square-root orientations,
quotient orientation, transition compatibility, vanishing cycles, or
protected integration.
\end{proposition}

\begin{proof}
The obstruction sheaf of the retained quasi-smooth derived enhancement
has fibre \(\operatorname{Ext}^2_X(F,F)\) on the chart represented by
the universal perfect complex \(F\).  The Atiyah class and trace are
functorial under restriction to etale charts.  Hence the assignment
displayed above glues from charts to an
\(\mathcal O\)-linear morphism of obstruction sheaves.  The wedge
product lies in
\[
H^3(X,\Omega_X^3)=H^3(X,K_X)\simeq \C
\]
because \(\sigma\in H^0(X,\Omega_X^2)\).  This is the
Buchweitz--Flenner--Bloch semiregularity map in the \(K3\times E\)
form used by Maulik--Toda
\cite[Lemma~2.6, Theorem~2.7]{MaulikTodaGV} and by the Kiem--Li
cosection formalism \cite{KiemLi2013}.  The packet checks the same
formula on the three retained truncation rows.  No nonvanishing or
surjectivity assertion enters this construction.
\end{proof}

The packet
\[
\texttt{certificates/orientation/k3\_semiregularity\_cosection}
\]
verifies
\(\texttt{K3\_SEMIREGULARITY\_COSECTION\_VERIFIED}\): it imports the
Joyce d-critical truncation packet, writes the semiregularity formula
above on the three retained finite rows, and records zero construction
defect.  Its firewall rows record that surjectivity, filtered Hall
additivity, reduced self-Ext complexes, determinant lines, square-root
orientations, quotient orientation, transition compatibility, vanishing
cycles, and protected integration remain open.

\begin{proposition}[Cosection surjectivity criterion on retained K3-class strata]
\label{prop:k3-cosection-surjectivity-criterion}
\textnormal{[criterion proved; retained nonzero K3-class witnesses not supplied]}
Let \(\mathfrak M^{\mathrm{cl}}_{R,c}\) be a retained finite classical
truncation equipped with the cosection
\(\operatorname{CS}_{R,c}\) of
Proposition~\ref{prop:k3-semiregularity-cosection}.  Suppose the
retained row carries a K3-class witness
\[
\beta_{R,c}\in H_2(S,\mathbb Z),\qquad \beta_{R,c}\ne 0,
\]
recorded by the support or charge map for every geometric point of
\(\mathfrak M^{\mathrm{cl}}_{R,c}\), and suppose the row lies on the
reduced DT/PT branch on which the Maulik--Toda semiregularity theorem
applies.  Then
\[
\operatorname{CS}_{R,c}:
\operatorname{Ob}_{R,c}\twoheadrightarrow
\mathcal O_{\mathfrak M^{\mathrm{cl}}_{R,c}}
\]
is surjective, equivalently its cokernel has rank zero.

The current retained finite packets do not supply the witness
\(\beta_{R,c}\ne0\), a branch-identification row, or a cokernel-rank
row.  Therefore row \(274\) is not discharged here.
\end{proposition}

\begin{proof}
Surjectivity of a cosection of the obstruction sheaf is local on the
classical moduli stack.  At a geometric point \(F\), the fibre map is
the semiregularity map displayed in
Proposition~\ref{prop:k3-semiregularity-cosection}.  On the reduced
nonzero K3-class DT/PT branch, Maulik--Toda
\cite[Lemma~2.6, Theorem~2.7]{MaulikTodaGV} identify this map with the
standard K3 semiregularity cosection and prove surjectivity; the same
surjective branch is used in the reduced stable-pair theory of
Oberdieck \cite[Theorems~1 and~3]{OberdieckReducedPT}.  If the retained
row records a nonzero K3 class \(\beta_{R,c}\) at every geometric point
and identifies the row with that branch, these pointwise surjections
glue to a surjection of sheaves.  The last assertion is a table-level
audit: the retained class, charge, and D0-HN semiregularity tables have
no row recording \(\beta_{R,c}\ne0\) and no row recording cokernel rank
zero.
\end{proof}

The packet
\[
\texttt{certificates/orientation/k3\_cosection\_surjectivity}
\]
verifies
\(\texttt{K3\_COSECTION\_SURJECTIVITY\_OBSTRUCTION\_VERIFIED}\): it
records the criterion above, checks that the current retained finite
class tables lack a K3-class witness column or row, checks that the D0
semiregularity table is empty, and keeps the surjectivity and cokernel
tables empty.  This status is not a proof of surjectivity; it is the
fail-closed ledger for row \(274\).

After a future retained-stratum packet supplies the nonzero K3-class
witness and the cokernel-rank-zero row, the later additivity row checks
the cosection on extension strata.  On
the extension stack $\mathfrak E_{A,B}$ of short exact sequences
$0\to A\to C\to B\to 0$ — with projections $p_1,p_2$ to the substacks
$\mathfrak M_A,\mathfrak M_B$ of sub- and quotient-object and $q$ to
$\mathfrak M_C$, the correspondence set up in
\S\ref{subsec:hall-correspondences} — the later additivity row is the
filtered Hall identity
\[
q^*\operatorname{CS}_C
=p_1^*\operatorname{CS}_A+p_2^*\operatorname{CS}_B,
\]
the compatibility hypothesis in Kiem--Li cosection localization
\cite{JoyceSong,JoyceUpmeier}.

\begin{definition}[Cosection-reduced self-Ext cone]
\label{def:rhom-red-self-ext-cone}
Let \(F\) be the universal object over a retained finite classical
truncation \(\mathfrak M^{\mathrm{cl}}_{R,c}\).  Assume the row carries
the cosection \(\operatorname{CS}_{R,c}\) of
Proposition~\ref{prop:k3-semiregularity-cosection}.  Define the
cosection-reduced self-Ext cone by
\[
\operatorname{RHom}_{\mathrm{red}}(F,F)
:=
\operatorname{fib}\!\left(
\Theta_F:
\operatorname{RHom}_X(F,F)\longrightarrow
R\Gamma(X,\mathcal O_X)\oplus H^2(S,\mathcal O_S)[-1]
\right),
\]
where
\[
\Theta_F=(\operatorname{tr},\operatorname{CS}_F),
\]
and equivalently
\[
\operatorname{RHom}_{\mathrm{red}}(F,F)
=
\operatorname{Cone}(\Theta_F)[-1].
\]
The first component removes the scalar trace direction; the second is
the K3 semiregularity cosection direction.  If the surjectivity witness
of Proposition~\ref{prop:k3-cosection-surjectivity-criterion} is also
supplied, this cone is the complex entering the Kiem--Li reduced
obstruction theory on the retained branch.  This definition does not
prove surjectivity, does not prove that the cone is perfect, does not
identify its determinant line, does not construct a square root, and does
not supply quotient orientation, transition compatibility, vanishing
cycles, or protected integration.
\end{definition}

The packet
\[
\texttt{certificates/orientation/rhomred\_definition}
\]
verifies \(\texttt{RHOMRED\_DEFINITION\_VERIFIED}\): it imports the
semiregularity-cosection packet and the surjectivity-obstruction packet,
records the cone formula above for the three retained finite rows, and
records zero definition defect.  Its firewall rows record that
surjectivity, determinant lines, square-root orientations,
quotient orientation, transition compatibility, vanishing cycles, and
protected integration are not supplied by the cone-definition row.

\begin{proposition}[Perfectness of the cosection-reduced self-Ext cone]
\label{prop:rhomred-perfectness}
\textnormal{[proved for the cone complex; determinant line not yet defined]}
Let \(F=\mathcal U_{R,c}\) be the universal perfect complex on
\(X\times\mathfrak M^{\mathrm{rig}}_{R,c}\) supplied by
Proposition~\ref{prop:retained-universal-perfect-complexes}, and let
\(\Theta_F\) be the trace--cosection morphism of
Definition~\ref{def:rhom-red-self-ext-cone}.  Then
\[
\operatorname{RHom}_{\mathrm{red}}(F,F)
=\operatorname{fib}(\Theta_F)
\]
is a perfect complex on the retained finite row.  The perfectness defect
is zero for each of the three retained substacks
\(Mss_{R0,s3},Mss_{R0,s2},Mss_{R0,s1}\).

This proves row \(276\) at the level of the reduced self-Ext cone.  It
does not prove cosection surjectivity, does not construct the determinant
line \(\det\operatorname{RHom}_{\mathrm{red}}(F,F)\), does not construct
a square-root orientation, and does not supply quotient orientation,
transition compatibility, vanishing cycles, or protected integration.
\end{proposition}

\begin{proof}
Let
\[
p:X\times\mathfrak M^{\mathrm{rig}}_{R,c}\longrightarrow
\mathfrak M^{\mathrm{rig}}_{R,c}
\]
be the projection.  Proposition~\ref{prop:retained-universal-perfect-complexes}
gives \(F=\mathcal U_{R,c}\in
\operatorname{Perf}(X\times\mathfrak M^{\mathrm{rig}}_{R,c})\), with
finite Tor amplitude and zero perfection defect.  Hence
\[
R\mathcal Hom(F,F)
\]
is perfect on \(X\times\mathfrak M^{\mathrm{rig}}_{R,c}\).  The morphism
\(p\) is smooth and proper because \(X=K3\times E\) is smooth
projective.  Therefore \(Rp_*R\mathcal Hom(F,F)\), written fibrewise as
\(\operatorname{RHom}_X(F,F)\), is perfect by the Stacks Project
\cite[Tag~0685, Lemma~37.61.13]{StacksProject}.  The target
\[
R\Gamma(X,\mathcal O_X)\oplus H^2(S,\mathcal O_S)[-1]
\]
is a finite-dimensional constant complex, hence perfect.  The stable
\(\infty\)-category of perfect complexes is closed under fibres and
cofibres.  Thus
\[
\operatorname{fib}\!\left(
\operatorname{RHom}_X(F,F)\xrightarrow{\Theta_F}
R\Gamma(X,\mathcal O_X)\oplus H^2(S,\mathcal O_S)[-1]\right)
\]
is perfect.  The row \(276\) packet checks this argument on the three
retained rows by importing the universal-perfect-complex packet and the
cone-definition packet.
\end{proof}

The packet
\[
\texttt{certificates/orientation/rhomred\_perfectness}
\]
verifies \(\texttt{RHOMRED\_PERFECTNESS\_VERIFIED}\): it imports the
retained universal-perfect-complex packet, the retained
derived-enhancement packet, and the cone-definition packet; records
perfect source and perfect target complexes; checks closure of perfect
complexes under fibres for the three retained rows; and records zero
perfectness defect.  Its firewall rows record that determinant lines,
square-root orientations, quotient orientation, transition compatibility,
vanishing cycles, and protected integration are not supplied by the
perfectness row.

\begin{definition}[Determinant line of the reduced self-Ext complex]
\label{def:rhomred-determinant-line}
Let \(F=\mathcal U_{R,c}\) be a retained universal perfect complex and
put
\[
K^{\mathrm{red}}_{F}
:=\operatorname{RHom}_{\mathrm{red}}(F,F).
\]
By Proposition~\ref{prop:rhomred-perfectness}, \(K^{\mathrm{red}}_{F}\)
is perfect.  The reduced determinant line is the Knudsen--Mumford
determinant
\[
\mathcal L^{\det}_{F}
:=\det K^{\mathrm{red}}_{F}
:=\operatorname{Det}_{\mathrm{KM}}\!\left(K^{\mathrm{red}}_{F}\right)
\in\operatorname{Pic}\!\left(\mathfrak M^{\mathrm{rig}}_{R,c}\right)
\]
\cite{KnudsenMumfordDet,DeligneDetCohomology}.  On a local perfect
representative \(K^\bullet\), this line is
\[
\det K^{\mathrm{red}}_{F}
=\bigotimes_i \det(K^i)^{(-1)^i},
\]
with the Knudsen--Mumford transition signs.  It is invariant under
quasi-isomorphism and functorial for pullback of the retained row.  The
determinant-defect rank is zero for
\[
Mss_{R0,s3},\qquad Mss_{R0,s2},\qquad Mss_{R0,s1}.
\]

This definition supplies row \(277\).  It does not choose a square root of
\(\mathcal L^{\det}_{F}\), does not compute \(w_2(\mathcal L^{\det}_F)\),
does not construct quotient orientation or Weyl transport, and does not
supply vanishing cycles or protected integration.
\end{definition}

The packet
\[
\texttt{certificates/orientation/rhomred\_determinant}
\]
verifies \(\texttt{RHOMRED\_DETERMINANT\_VERIFIED}\): it imports the
perfectness packet, applies the Knudsen--Mumford determinant functor to
the three retained perfect reduced self-Ext complexes, records the three
determinant lines, and records zero determinant defect.  Its firewall
rows record that square-root orientations, \(w_2\)-classes, quotient
orientation, transition compatibility, vanishing cycles, and protected
integration are not supplied by the determinant row.

\begin{proposition}[Square-root criterion for the reduced determinant line]
\label{prop:rhomred-square-root-criterion}
\textnormal{[criterion proved; square-root rows not supplied]}
Let \(\mathcal L^{\det}_{F}\) be the determinant line of
Definition~\ref{def:rhomred-determinant-line}.  A retained row
constructs the reduced orientation square root precisely when it supplies
a line
\[
o_F\in\operatorname{Pic}\!\left(\mathfrak M^{\mathrm{rig}}_{R,c}\right)
\]
and an isomorphism
\[
\epsilon_F:o_F^{\otimes 2}\xrightarrow{\sim}\mathcal L^{\det}_{F}.
\]
Equivalently, \(\mathcal L^{\det}_{F}\) lies in the essential image of
the squaring morphism
\[
[2]:\operatorname{Pic}\!\left(\mathfrak M^{\mathrm{rig}}_{R,c}\right)
\longrightarrow
\operatorname{Pic}\!\left(\mathfrak M^{\mathrm{rig}}_{R,c}\right).
\]
When it exists, the set of square roots is a torsor under the
two-torsion subgroup
\(\operatorname{Pic}(\mathfrak M^{\mathrm{rig}}_{R,c})[2]\).

The current retained orientation packet supplies no line \(o_F\), no
isomorphism \(\epsilon_F\), and no zero square-defect row for
\(Mss_{R0,s3},Mss_{R0,s2},Mss_{R0,s1}\).  Therefore row \(278\) is not
discharged by the present finite tables.  The class
\[
w_2(\mathcal L^{\det}_{F})
\]
is the mod-\(2\) obstruction entering row \(279\) and the quotient
orientation gerbes; its vanishing is necessary for a topological square
root, but the algebraic square root itself is the Picard-groupoid datum
\((o_F,\epsilon_F)\).
\end{proposition}

\begin{proof}
The assertion is the definition of a square root of a line bundle in the
Picard groupoid.  Passing to isomorphism classes gives the squaring map
on \(\operatorname{Pic}\); the fibre of this map over
\(\mathcal L^{\det}_{F}\) is either empty or a torsor under the kernel
of squaring, namely the two-torsion subgroup.  The finite-table
assertion is a direct inspection of
\(\texttt{certificates/orientation/k3e\_reduced\_orientation}\): its
\(\texttt{orientation\_lines.csv}\) table has no square-root rows and its
obstruction ledger records the missing square-root and orientation-class
vanishing rows.
\end{proof}

The packet
\[
\texttt{certificates/orientation/rhomred\_square\_root}
\]
verifies \(\texttt{RHOMRED\_SQUARE\_ROOT\_OBSTRUCTION\_VERIFIED}\): it
imports the determinant packet and the reduced-orientation obstruction
ledger, records the Picard-groupoid square-root criterion, checks that
the three determinant lines have no supplied square roots in the current
orientation table, and keeps \(w_2\), quotient orientation, transition
compatibility, vanishing cycles, and protected integration as open
obligations.

\begin{proposition}[Stiefel--Whitney formula for the reduced determinant line]
\label{prop:rhomred-w2-formula}
\textnormal{[formula proved; retained \(w_2\)-values not supplied]}
Let \(\mathcal L^{\det}_{F}\) be the determinant line of
Definition~\ref{def:rhomred-determinant-line}, regarded as a complex
line bundle on the retained analytic stack.  Its second
Stiefel--Whitney class is
\[
w_2\!\left((\mathcal L^{\det}_{F})_{\mathbb R}\right)
=c_1(\mathcal L^{\det}_{F})\bmod 2
\in H^2(\mathfrak M^{\mathrm{rig}}_{R,c};\mathbb F_2).
\]
If \(K^{\mathrm{red}}_{F}\) is represented by a perfect kernel
\(\mathcal K^{\mathrm{red}}_{F}\) on
\(X\times\mathfrak M^{\mathrm{rig}}_{R,c}\) and
\(\pi:X\times\mathfrak M^{\mathrm{rig}}_{R,c}\to
\mathfrak M^{\mathrm{rig}}_{R,c}\), then the determinant-of-cohomology
calculation required for row \(279\) is the degree-one part of
Grothendieck--Riemann--Roch:
\[
c_1(\mathcal L^{\det}_{F})
=
\left[
\pi_*\!\left(
\operatorname{ch}(\mathcal K^{\mathrm{red}}_{F})
\operatorname{td}(T_{\pi})
\right)
\right]_1 .
\]
Thus row \(279\) is discharged only by supplying, for each retained
substack, an integral \(c_1\)-row or an equivalent GRR expansion for
\(\mathcal K^{\mathrm{red}}_{F}\), together with its mod-\(2\)
reduction.

The present determinant and square-root packets supply none of these
rows for \(Mss_{R0,s3},Mss_{R0,s2},Mss_{R0,s1}\).  They define the
three determinant lines and record the absence of square roots, but they
do not compute \(c_1(\mathcal L^{\det}_{F})\), do not reduce it modulo
\(2\), and do not verify \(w_2=0\).
\end{proposition}

\begin{proof}
For a complex vector bundle \(E\), the Stiefel--Whitney classes of the
underlying real bundle satisfy
\(w_{2i}(E_{\mathbb R})=c_i(E)\bmod 2\) and
\(w_{2i+1}(E_{\mathbb R})=0\)
\cite[Theorem~14.10]{MilnorStasheffCharacteristicClasses}; the case of
a complex line gives the displayed formula.  The determinant line is the
Knudsen--Mumford determinant of the perfect complex
\cite{KnudsenMumfordDet,DeligneDetCohomology}; its first Chern class is
computed from determinant of cohomology by the displayed
Grothendieck--Riemann--Roch expression \cite{ToenRRRoot}.  Inspection of
\(\texttt{certificates/orientation/rhomred\_determinant}\) and
\(\texttt{certificates/orientation/rhomred\_square\_root}\) shows that
the retained rows have no \(c_1\)-payload, no GRR expansion row, no
mod-\(2\) reduction row, and no \(w_2\)-value row.
\end{proof}

The packet
\[
\texttt{certificates/orientation/rhomred\_w2}
\]
verifies \(\texttt{RHOMRED\_W2\_FORMULA\_OBSTRUCTION\_VERIFIED}\): it
imports the determinant and square-root packets, records the
\(w_2=c_1\bmod 2\) and determinant-of-cohomology formulae, checks the
three retained determinant lines, and keeps the actual \(c_1\), GRR,
mod-\(2\) reduction, \(w_2\)-vanishing, square-root, quotient
orientation, transition, vanishing-cycle, and protected-integration
rows open.

\begin{proposition}[Vanishing criterion for the reduced orientation gerbe]
\label{prop:alpha-red-vanishing-criterion}
\textnormal{[criterion proved; \(\alpha^{\mathrm{red}}=0\) rows not supplied]}
For a retained charge stratum \(\mathcal C\), define the reduced
orientation gerbe class
\[
\alpha^{\mathrm{red}}_{\mathcal C}
=w_2(\det K^{\mathrm{red}}_{\mathcal C})
+\operatorname{CS}_{\mathcal C}^{*}
w_2(\det\operatorname{coker}\operatorname{CS}_{\mathcal C})
\in H^2(\mathfrak M^{\mathrm{red}}_{\mathcal C};\mathbb F_2).
\]
A retained row proves \(\alpha^{\mathrm{red}}_{\mathcal C}=0\) precisely
when it supplies:
\[
u_{\mathcal C}=w_2(\det K^{\mathrm{red}}_{\mathcal C}),\qquad
v_{\mathcal C}=w_2(\det\operatorname{coker}\operatorname{CS}_{\mathcal C}),
\]
the pullback \(\operatorname{CS}_{\mathcal C}^{*}v_{\mathcal C}\), and
the equality
\[
u_{\mathcal C}+\operatorname{CS}_{\mathcal C}^{*}v_{\mathcal C}=0
\]
in \(H^2(\mathfrak M^{\mathrm{red}}_{\mathcal C};\mathbb F_2)\).
If the cosection-cokernel row proves
\(\operatorname{coker}\operatorname{CS}_{\mathcal C}=0\), then the
second summand is zero and row \(280\) reduces to the \(w_2\)-vanishing
row of Proposition~\ref{prop:rhomred-w2-formula}.  If the cokernel is
not proved zero, its determinant line and mod-\(2\) class are separate
inputs.

The present finite tables do not prove \(\alpha^{\mathrm{red}}=0\) for
\(Mss_{R0,s3},Mss_{R0,s2},Mss_{R0,s1}\).  They do not supply
retained-stratum \(w_2(\det K^{\mathrm{red}})\)-values, do not supply
cosection-cokernel rank-zero rows, do not supply
\(w_2(\det\operatorname{coker}\operatorname{CS})\)-rows, and do not
supply a zero-sum row for \(\alpha^{\mathrm{red}}\).  The null-homotopy
of this class is a stronger datum and belongs to row \(281\).
\end{proposition}

\begin{proof}
The displayed formula is the definition of the reduced degree-two
orientation gerbe in the quotient-orientation tower: the first summand
comes from the determinant line of the reduced self-Ext complex, and the
second is the cosection-cokernel correction pulled back along the
cosection map.  Since the coefficient group is \(\mathbb F_2\), vanishing
of the class is exactly the equality of the two summands with sum zero.
If the cokernel sheaf is zero, its determinant line is the trivial line,
so its second Stiefel--Whitney class is zero.  Inspection of
\(\texttt{certificates/orientation/rhomred\_w2}\),
\(\texttt{certificates/orientation/k3\_cosection\_surjectivity}\), and
\(\texttt{certificates/orientation/k3e\_reduced\_orientation}\) shows
that the determinant \(w_2\)-values, cokernel rank-zero rows, cokernel
\(w_2\)-rows, quotient-Borel \(\alpha^{\mathrm{red}}\)-rows, and
null-homotopy rows are absent.
\end{proof}

The packet
\[
\texttt{certificates/orientation/rhomred\_alpha\_red}
\]
verifies \(\texttt{RHOMRED\_ALPHA\_RED\_OBSTRUCTION\_VERIFIED}\): it
imports the row \(279\) \(w_2\) packet, the cosection-surjectivity
obstruction packet, and the reduced-orientation obstruction ledger;
records the formula and zero-sum criterion for
\(\alpha^{\mathrm{red}}\); checks the three retained strata; and keeps
the determinant \(w_2\), cokernel rank, cokernel \(w_2\), zero-sum,
null-homotopy, quotient orientation, transition, vanishing-cycle, and
protected-integration rows open.

\begin{proposition}[Null-homotopy criterion for the reduced gerbe]
\label{prop:alpha-red-nullhomotopy-criterion}
\textnormal{[criterion proved; null-homotopy rows not supplied]}
Fix a retained Cech--Borel cover of the reduced stratum
\(\mathfrak M^{\mathrm{red}}_{\mathcal C}\).  Let
\[
a^{\mathrm{red}}_{\mathcal C}\in
Z^2(\mathfrak M^{\mathrm{red}}_{\mathcal C};\mathbb F_2)
\]
be a cocycle representative of
\(\alpha^{\mathrm{red}}_{\mathcal C}\).  A row constructs a
null-homotopy of \(\alpha^{\mathrm{red}}_{\mathcal C}\) precisely when
it supplies a \(1\)-cochain
\[
h^{\mathrm{red}}_{\mathcal C}\in
C^1(\mathfrak M^{\mathrm{red}}_{\mathcal C};\mathbb F_2)
\]
and verifies the coboundary equation
\[
\delta h^{\mathrm{red}}_{\mathcal C}=a^{\mathrm{red}}_{\mathcal C}.
\]
Thus row \(281\) is stronger than row \(280\): the equality
\(\alpha^{\mathrm{red}}_{\mathcal C}=0\) proves only that such a
cochain exists after choosing a representative; row \(281\) records the
representative, the cochain, and the zero coboundary-defect row.  When
nonempty, the set of null-homotopies is a torsor under
\(H^1(\mathfrak M^{\mathrm{red}}_{\mathcal C};\mathbb F_2)\).

The present finite tables do not construct a null-homotopy for
\(Mss_{R0,s3},Mss_{R0,s2},Mss_{R0,s1}\).  They do not supply a cocycle
representative of \(\alpha^{\mathrm{red}}\), do not prove the class
zero, do not supply a \(1\)-cochain, and do not verify the coboundary
equation.  No quotient-orientation descent row follows from an empty
quotient-Borel table.
\end{proposition}

\begin{proof}
This is the definition of a null-homotopy in the cochain model for
degree-two \(\mathbb F_2\)-valued Cech--Borel cohomology.  A
degree-two cocycle is null-homotopic precisely when it is a coboundary;
the possible primitives differ by \(1\)-cocycles, and changing a
primitive by a \(1\)-coboundary gives the same trivialization.  Hence the
set of choices is a torsor under \(H^1\).  Inspection of
\(\texttt{certificates/orientation/rhomred\_alpha\_red}\) and
\(\texttt{certificates/orientation/k3e\_reduced\_orientation}\) shows
that no cocycle representative, no zero-class row, no \(1\)-cochain, and
no coboundary-defect-zero row is supplied.
\end{proof}

The packet
\[
\texttt{certificates/orientation/rhomred\_alpha\_red\_nullhomotopy}
\]
verifies
\(\texttt{RHOMRED\_ALPHA\_RED\_NULLHOMOTOPY\_OBSTRUCTION\_VERIFIED}\):
it imports the row \(280\) packet and the reduced-orientation
obstruction ledger, records the cochain criterion
\(\delta h^{\mathrm{red}}=a^{\mathrm{red}}\), checks the three retained
strata, and keeps the cocycle representative, zero-class, \(1\)-cochain,
coboundary-defect, quotient orientation, transition, vanishing-cycle,
and protected-integration rows open.

\begin{proposition}[Computation criterion for the connected free
\(E\)-Borel class]
\label{prop:free-e-borel-class-criterion}
\textnormal{[criterion proved; \(a_1,a_2\) rows not supplied]}
Let \(E^{\mathrm{top}}\simeq T^2\), and use the basis
\[
H^2(BE;\mathbb F_2)=\mathbb F_2u_1\oplus\mathbb F_2u_2
\]
of Lemma~\ref{lem:connected-e-descent}.  For a retained stratum
\(\mathcal C\) with free \(E\)-translation action, the connected-free
quotient-orientation class is the \(E_\infty^{2,0}\)-edge component of
the equivariant degree-two Borel class of the reduced determinant
complex:
\[
\alpha^{E,\mathrm{free}}_{\mathcal C}
=\operatorname{edge}_{2,0}
\bigl(w_2^E(\mathscr D_{\mathcal C})\bigr)
=a_1(\mathcal C)u_1+a_2(\mathcal C)u_2.
\]
A row computes \(\alpha^{E,\mathrm{free}}_{\mathcal C}\) precisely when
it supplies the equivariant Borel representative
\(w_2^E(\mathscr D_{\mathcal C})\), the spectral-sequence edge
calculation including any incoming \(d_2:E_2^{0,1}\to E_2^{2,0}\)
contribution, and the two coefficients
\[
a_1(\mathcal C),a_2(\mathcal C)\in\mathbb F_2.
\]
The vanishing equation \(a_1=a_2=0\) is row \(283\), and a chosen
null-trivialization of that zero class is row \(284\).

The present finite tables do not compute
\(\alpha^{E,\mathrm{free}}\) for
\(Mss_{R0,s3},Mss_{R0,s2},Mss_{R0,s1}\).  The retained
\(E\)-translation packet supplies free actions and quotient stacks, but
the reduced-orientation packet supplies no \(\texttt{free\_E}\)
quotient-Borel rows, no Borel spectral-sequence edge rows, and no
\(a_1,a_2\) coefficient rows.  The equality \(H^1(BE;\mathbb F_2)=0\)
rules out connected degree-one characters; it does not compute the
degree-two class above.
\end{proposition}

\begin{proof}
The cohomology of the connected torus action is
\(H^*(BE;\mathbb F_2)=\mathbb F_2[u_1,u_2]\), \(|u_i|=2\), by
Lemma~\ref{lem:connected-e-descent}.  Hence every connected-free
degree-two edge class has a unique expansion
\(a_1u_1+a_2u_2\).  The Borel fibration spectral sequence computes the
equivariant class from \(H^*_E(\mathcal C;\mathbb F_2)\), and the edge
component in bidegree \((2,0)\) is well-defined only after the relevant
incoming differential has been killed or included in the reported
coefficient row.  Inspection of
\(\texttt{certificates/moduli/retained\_e\_translation\_rigidifications}\)
shows that the free actions are supplied, while inspection of
\(\texttt{certificates/orientation/k3e\_reduced\_orientation}\) shows
that the quotient-Borel table has no \(\texttt{free\_E}\) row and no
coefficient or edge-reduction payload.
\end{proof}

The packet
\[
\texttt{certificates/orientation/rhomred\_free\_e\_borel}
\]
verifies \(\texttt{RHOMRED\_FREE\_E\_BOREL\_OBSTRUCTION\_VERIFIED}\):
it imports the retained \(E\)-translation rigidification packet and the
reduced-orientation obstruction ledger, records the basis and edge-class
criterion for
\(\alpha^{E,\mathrm{free}}=a_1u_1+a_2u_2\), checks the three retained
free \(E\)-actions, and keeps the equivariant Borel representative,
edge-reduction, \(a_1,a_2\) coefficient, vanishing, null-trivialization,
transition, vanishing-cycle, and protected-integration rows open.

\begin{proposition}[Vanishing criterion for the connected free
\(E\)-Borel class]
\label{prop:free-e-borel-vanishing-criterion}
\textnormal{[criterion proved; \(\alpha^{E,\mathrm{free}}=0\) rows not
supplied]}
Let \(\mathcal C\) be one of the retained free \(E\)-translation strata,
and suppose row \(282\) has supplied the connected-free edge class
\[
\alpha^{E,\mathrm{free}}_{\mathcal C}
=a_1(\mathcal C)u_1+a_2(\mathcal C)u_2
\in H^2(BE;\mathbb F_2)
=\mathbb F_2u_1\oplus\mathbb F_2u_2 .
\]
Then row \(283\) proves
\(\alpha^{E,\mathrm{free}}_{\mathcal C}=0\) precisely by supplying the
same coefficient row together with the two equalities
\[
a_1(\mathcal C)=0,\qquad a_2(\mathcal C)=0
\]
in \(\mathbb F_2\).  The vanishing is not implied by
\(H^1(BE;\mathbb F_2)=0\), by freeness of the \(E\)-action, by the
existence of the quotient stack, by a finite-stabilizer calculation, or
by any scalar trace or character value.  A null-trivialization of this
zero class is the separate row \(284\).

The present finite tables do not prove
\(\alpha^{E,\mathrm{free}}=0\).  The row \(282\) packet has no
\(a_1,a_2\) values; the reduced-orientation quotient-Borel table has no
\(\texttt{free\_E}\) row; and no zero-coefficient row is supplied for
\(Mss_{R0,s3},Mss_{R0,s2},Mss_{R0,s1}\).
\end{proposition}

\begin{proof}
The classes \(u_1,u_2\) form a basis of
\(H^2(BE;\mathbb F_2)\) by Lemma~\ref{lem:connected-e-descent}.
Therefore the expansion of
\(\alpha^{E,\mathrm{free}}_{\mathcal C}\) in
Proposition~\ref{prop:free-e-borel-class-criterion} is unique, and the
class is zero if and only if both coordinates are zero.  This is a
basis calculation after the Borel edge class has been computed; it is
not a replacement for that computation.  The packet
\(\texttt{certificates/orientation/rhomred\_free\_e\_borel}\) records
the absence of the coefficient row, and
\(\texttt{certificates/orientation/k3e\_reduced\_orientation}\) records
the empty quotient-Borel table, so the current finite evidence supplies
the criterion but not its hypotheses.
\end{proof}

The packet
\[
\texttt{certificates/orientation/rhomred\_free\_e\_vanishing}
\]
verifies
\(\texttt{RHOMRED\_FREE\_E\_VANISHING\_OBSTRUCTION\_VERIFIED}\): it
imports the row \(282\) packet and the reduced-orientation obstruction
ledger, records the basis-zero criterion
\(\alpha^{E,\mathrm{free}}=0\Longleftrightarrow a_1=a_2=0\), checks the
three retained free \(E\)-actions, and keeps the \(a_1,a_2\)
coefficient, zero-coefficient, null-trivialization, transition,
vanishing-cycle, and protected-integration rows open.

\begin{proposition}[Null-trivialization criterion for the connected
free \(E\)-Borel class]
\label{prop:free-e-borel-nulltrivialization-criterion}
\textnormal{[criterion proved; null-trivialization rows not supplied]}
Fix a cochain model for the connected \(BE\)-edge component used in
Proposition~\ref{prop:free-e-borel-class-criterion}.  Let
\[
a^{E,\mathrm{free}}_{\mathcal C}\in Z^2(BE;\mathbb F_2)
\]
be a cocycle representative of the connected-free edge class
\(\alpha^{E,\mathrm{free}}_{\mathcal C}\).  Row \(284\) constructs a
null-trivialization of
\(\alpha^{E,\mathrm{free}}_{\mathcal C}\) precisely when it supplies a
\(1\)-cochain
\[
h^{E,\mathrm{free}}_{\mathcal C}\in C^1(BE;\mathbb F_2)
\]
and verifies the coboundary equation
\[
\delta h^{E,\mathrm{free}}_{\mathcal C}
=a^{E,\mathrm{free}}_{\mathcal C}.
\]
Thus row \(284\) is stronger than row \(283\): the vanishing
\(\alpha^{E,\mathrm{free}}_{\mathcal C}=0\) supplies the cohomology
condition, while row \(284\) records a representative, a primitive, and
a defect-zero coboundary calculation.  Since
\(H^1(BE;\mathbb F_2)=0\), the connected \(BE\)-edge homotopy class of
such a trivialization has no \(H^1\)-torsor ambiguity after the cochain
model is fixed; this uniqueness does not replace the cochain-level
construction.

The present finite tables do not construct this null-trivialization.
They do not supply a cocycle representative of
\(\alpha^{E,\mathrm{free}}\), do not prove the class zero, do not
supply a \(1\)-cochain, and do not verify a defect-zero coboundary
equation for \(Mss_{R0,s3},Mss_{R0,s2},Mss_{R0,s1}\).
\end{proposition}

\begin{proof}
A null-trivialization of a degree-two \(\mathbb F_2\)-valued Borel
class is, by definition, a primitive for a chosen degree-two cocycle
representative.  Hence the datum is exactly a cochain \(h\) with
\(\delta h=a\), plus the verification that the coboundary defect is
zero.  The equality \(H^1(BE;\mathbb F_2)=0\) from
Lemma~\ref{lem:connected-e-descent} says that two connected-edge
trivializations of the same cocycle are homotopic after fixing the
cochain model; it does not give the cocycle representative, the
primitive, or the coboundary calculation.  Inspection of
\(\texttt{certificates/orientation/rhomred\_free\_e\_vanishing}\) and
\(\texttt{certificates/orientation/k3e\_reduced\_orientation}\) shows
that none of these rows is supplied.
\end{proof}

The packet
\[
\texttt{certificates/orientation/rhomred\_free\_e\_nulltrivialization}
\]
verifies
\(\texttt{RHOMRED\_FREE\_E\_NULLTRIVIALIZATION\_OBSTRUCTION\_VERIFIED}\):
it imports the row \(283\) packet and the reduced-orientation
obstruction ledger, records the cochain criterion
\(\delta h^{E,\mathrm{free}}=a^{E,\mathrm{free}}\), checks the three
retained free \(E\)-actions, and keeps the cocycle representative,
zero-class, \(1\)-cochain, coboundary-defect, quotient orientation,
transition, vanishing-cycle, and protected-integration rows open.

\subsubsection*{The reduced determinant complex}

The following notation belongs to the later orientation rows.  After the
cosection, cone-definition, perfectness, and determinant rows have been
supplied, for an object $\mathcal C$ in the chosen finite stage
write
\[
K^{\mathrm{red}}_{\mathcal C}
=\operatorname{RHom}_{\mathrm{red}}(\mathcal C,\mathcal C),
\]
the cosection-reduced self-Ext complex of
Definition~\ref{def:rhom-red-self-ext-cone}; it is perfect by
Proposition~\ref{prop:rhomred-perfectness}.  Its determinant line is
\[
\mathcal L^{\det}_{\mathcal C}
=\det K^{\mathrm{red}}_{\mathcal C}
\]
by Definition~\ref{def:rhomred-determinant-line}.  The later orientation
row, when supplied, equips the reduced orientation line with a chosen
square root,
\[
o_{\mathcal C}\in\operatorname{Pic}(\mathfrak M^{\mathrm{or}}_X),
\qquad
o_{\mathcal C}^{\otimes 2}\simeq\det K^{\mathrm{red}}_{\mathcal C},
\]
in the sense of \cite{JoyceUpmeier}, and is multiplicative under the
reduced extension correspondences:
$o_{\mathcal C\oplus\mathcal C'}\simeq
o_{\mathcal C}\otimes o_{\mathcal C'}$ compatibly with the Hall
additivity above.  This is the data on which Thom--Sebastiani
transport is supplied for the vanishing-cycle coefficient sheaves.

\subsubsection*{(O1) Strong reduced orientation through the
quotient-orientation tower}

The compact Hall product on $X=K3\times E$ exists only after the
reduced $E$-quotient.  The orientation line must descend through this
quotient, and descent is detected by a tower of mod-$2$
Borel-equivariant gerbes.  For a charge stratum $\mathcal C$:

\begin{enumerate}
\item the reduced gerbe class
\[
\alpha^{\mathrm{red}}_{\mathcal C}
=w_2(\det K^{\mathrm{red}}_{\mathcal C})
+\operatorname{CS}_{\mathcal C}^{*}w_2(\det\operatorname{coker}\operatorname{CS}_{\mathcal C})
\in H^2(\mathfrak M^{\mathrm{red}}_{\mathcal C};\mathbb F_2);
\]
\item the free $E$-class
\[
\alpha^{E,\mathrm{free}}_{\mathcal C}
=a_1(\mathcal C)u_1+a_2(\mathcal C)u_2
\in H^2(BE;\mathbb F_2);
\]
\item for each finite-stabilizer stratum $S$ with stabilizer
$H\subset E_{\mathrm{tors}}$, the equivariant class
$\widetilde\beta^{H}_{\mathcal C,S}\in H^2_H(S;\mathbb F_2)$, whose
point-stabilizer restriction $\beta^H_{\mathcal C,S}\in H^2(BH;\mathbb F_2)$
is read only after the mixed Borel terms, the stabilizer action on
$H^\bullet(S;\mathbb F_2)$, and the spectral-sequence differentials
have vanished.
\end{enumerate}

\begin{definition}[(O1) Strong multiplicative reduced orientation]
\label{def:O1-strong-reduced-orientation}
The (O1) datum at finite HN height $R$ consists of:
the reduced orientation line $o_{\mathcal C}$ with chosen square root;
its multiplicativity under reduced extension correspondences with
chosen Thom--Sebastiani transport;
null-trivializations of $\alpha^{\mathrm{red}}_{\mathcal C}$,
$\alpha^{E,\mathrm{free}}_{\mathcal C}$,
$\widetilde\beta^{H}_{\mathcal C,S}$ and
$\beta^{H}_{\mathcal C,S}$ on every charge stratum
used by the Hall correspondences;
and a choice of zero $H^1(BH;\mathbb F_2)$-linearization character
$\lambda^{H}_{\mathcal C,S}=0$ for each finite stabilizer.

The full-level specialization at $H=E[2]$ is the Klein-four condition
\[
\beta^{E,2}_{\mathcal C,S}
=b_{20}x_1^2+b_{11}x_1x_2+b_{02}x_2^2,
\quad
(b_{20},b_{11},b_{02})=(r_1,r_1+r_2+r_3,r_2),
\]
\[
\operatorname{res}_{\langle e_i\rangle}\beta^{E,2}_{\mathcal C,S}
=r_i\tau_i^2,
\]
with vanishing $r_1=r_2=r_3=0$ on every retained two-torsion
stratum.  For full-level $H=E[N]$ with $2^a\parallel N$, $a\ge2$, the
two-primary residual takes value in
$H^2(B(\Z/2^a)^2;\mathbb F_2)=\mathbb F_2\langle y_1,x_1x_2,y_2\rangle$
and contains the mixed coefficient $A_{12}^{(N)}$ \emph{not} detected
by cyclic restrictions.
\end{definition}

(O1) is therefore not a single Picard-line assertion: it is the
ordinary square root, the null-trivialization of the reduced and
equivariant degree-two gerbes, and the choice of zero degree-one
linearization characters for the actual finite stabilizers.  At
finite HN height this is the cocycle system of
\cite{JoyceUpmeier} restricted to $K3\times E$.

After finite trivialisations, the \((O1)\) datum is recorded by the
orientation proof tables
\[
\mathrm{orientation\_lines},\quad
\mathrm{ts\_multiplicativity},\quad
\mathrm{quotient\_borel},\quad
\mathrm{finite\_stabilizers}.
\]
The first table records the determinant square root and the vanishing
orientation class.  The second records Thom--Sebastiani
multiplicativity and flag-pentagon defects.  The third records the
four quotient Cech--Borel classes
\(\alpha^{\mathrm{red}}\), \(\alpha^{E,\mathrm{free}}\),
\(\widetilde\beta^H\), and \(\lambda^H\), with chosen
null-trivialisations and edge-reduction status.  The fourth records
the finite-stabilizer edge calculation, including the Klein-four
coefficients, the two-primary coefficients
\((A_1^{(H)},A_{12}^{(H)},A_2^{(H)})\), and the two degree-one
linearization coefficients.  All ranks and coefficients in these rows
must vanish.  A row using only cyclic restrictions is not an
orientation row.

\begin{proposition}[Computation criterion for the finite-stabilizer
Borel class]
\label{prop:finite-stabilizer-beta-class-criterion}
\textnormal{[criterion proved; \(\beta^H_{\mathcal C,S}\) rows not
supplied]}
Let \(S\) be a retained stratum with residual finite stabilizer
\(H\subset E_{\mathrm{tors}}\).  Row \(285\) computes the
finite-stabilizer class
\[
\beta^H_{\mathcal C,S}\in H^2(BH;\mathbb F_2)
\]
precisely when it supplies:
\begin{enumerate}
\item a finite-stabilizer stratum row identifying \(H\) and its
two-primary type;
\item a Cech--Borel representative
\[
\widetilde\beta^H_{\mathcal C,S}
=w_2^H(\mathscr D_{\mathcal C}|_S)
\in H^2_H(S;\mathbb F_2);
\]
\item an edge-reduction row killing or accounting for the ordinary
term in \(H^2(S;\mathbb F_2)^H\), the mixed term in
\(H^1(BH;H^1(S;\mathbb F_2))\), and the incoming Borel
spectral-sequence differentials;
\item the classifying-space coordinates of the resulting edge class.
\end{enumerate}
After these rows are supplied, the computation is the following group
cohomology calculation.  If \(H\) has odd order, and also for
\(H=1\), positive-degree \(H^*(BH;\mathbb F_2)\) vanishes.  If
\(H\simeq E[2]\), then
\[
\beta^H_{\mathcal C,S}
=b_{20}x_1^2+b_{11}x_1x_2+b_{02}x_2^2.
\]
If the two-primary part is \((\mathbb Z/2^a)^2\), \(a\ge2\), then
\[
\beta^H_{\mathcal C,S}
=A_1^{(H)}y_1+A_{12}^{(H)}x_1x_2+A_2^{(H)}y_2.
\]
The coefficient \(A_{12}^{(H)}\) is a rank-two stabilizer coefficient;
it is not read from cyclic order-two restrictions.  Row \(286\), not
row \(285\), is the assertion that the computed class vanishes for
every retained \(H\).  Rows \(288\)--\(289\) concern the independent
degree-one linearization character.

The present finite tables do not compute
\(\beta^H_{\mathcal C,S}\).  The retained rigidification-inertia packet
records finite residual groups \(H_{R0,s3},H_{R0,s2},H_{R0,s1}\) of
order \(1\), but it does not supply finite closed inertia
stratifications, Cech--Borel representatives, edge-reduction rows, or
finite-stabilizer orientation rows.  The reduced-orientation packet has
an empty \(\texttt{finite\_stabilizers}\) table and an empty
\(\texttt{quotient\_borel}\) table.
\end{proposition}

\begin{proof}
The Borel spectral sequence for \(S\to EH\times_H S\to BH\) filters
\(H^2_H(S;\mathbb F_2)\) by the classifying-space edge
\(H^2(BH;\mathbb F_2)\), the mixed term
\(H^1(BH;H^1(S;\mathbb F_2))\), and the ordinary term
\(H^2(S;\mathbb F_2)^H\), with the relevant incoming differentials.
Only after the non-edge terms are killed or explicitly accounted for
does the restriction of \(\widetilde\beta^H\) define
\(\beta^H\in H^2(BH;\mathbb F_2)\).  Lemma~\ref{lem:finite-stabilizer-e-descent}
and Lemmas~\ref{lem:G-Klein-four}--\ref{lem:G-two-primary-stabilizer}
give the displayed coordinates: transfer kills odd stabilizers,
\((\mathbb Z/2)^2\) gives the quadratic basis
\(x_1^2,x_1x_2,x_2^2\), and \((\mathbb Z/2^a)^2\), \(a\ge2\), gives
the basis \(y_1,x_1x_2,y_2\).  Inspection of
\(\texttt{certificates/moduli/retained\_rigidification\_inertia}\) and
\(\texttt{certificates/orientation/k3e\_reduced\_orientation}\) shows
that only residual finite group input is present; the edge and
coefficient rows are absent.
\end{proof}

The packet
\[
\texttt{certificates/orientation/rhomred\_finite\_stabilizer\_beta}
\]
verifies
\(\texttt{RHOMRED\_FINITE\_STABILIZER\_BETA\_OBSTRUCTION\_VERIFIED}\):
it imports the residual-inertia packet and the reduced-orientation
obstruction ledger, records the odd, Klein-four, and two-primary
coordinate criteria for row \(285\), checks the three retained residual
groups, and keeps the finite closed inertia stratum, Cech--Borel
representative, edge-reduction, \(\beta^H\)-coefficient, vanishing,
null-trivialization, linearization, transition, vanishing-cycle, and
protected-integration rows open.

\begin{proposition}[Vanishing criterion for the finite-stabilizer
Borel class]
\label{prop:finite-stabilizer-beta-vanishing-criterion}
\textnormal{[criterion proved; \(\beta^H_{\mathcal C,S}=0\) rows not
supplied]}
Assume row \(285\) has computed
\[
\beta^H_{\mathcal C,S}\in H^2(BH;\mathbb F_2)
\]
for every retained finite-stabilizer stratum.  Row \(286\) proves
\[
\beta^H_{\mathcal C,S}=0
\]
for every retained \(H\) precisely when it verifies the corresponding
zero condition in the classifying-space basis for every retained row:
\[
H=1\ \text{or}\ |H|\ \text{odd}:\quad H^2(BH;\mathbb F_2)=0
\quad\text{after edge reduction,}
\]
\[
H\simeq E[2]:\quad b_{20}=b_{11}=b_{02}=0,
\]
and, for two-primary rank two stabilizers,
\[
H_{(2)}\simeq(\mathbb Z/2^a)^2,\ a\ge2:
\quad A_1^{(H)}=A_{12}^{(H)}=A_2^{(H)}=0.
\]
The condition is imposed on every retained finite-stabilizer row.  A
cyclic-only calculation does not prove the two-primary case, because it
does not see \(A_{12}^{(H)}\).  Row \(287\) is the subsequent choice of
compatible null-trivializations, and rows \(288\)--\(289\) concern the
independent \(H^1(BH;\mathbb F_2)\)-linearization character.

The present finite tables do not prove the displayed vanishing.  The
row \(285\) packet records residual finite groups of order \(1\), but
it records no finite-stabilizer orientation row, no edge-reduction row,
no \(\beta^H\)-value, and no zero row.  The group-cohomology fact
\(H^2(B1;\mathbb F_2)=0\) is a zero template for a supplied edge row;
it is not itself the missing finite-stabilizer orientation datum.
\end{proposition}

\begin{proof}
After the edge-reduction hypotheses of
Proposition~\ref{prop:finite-stabilizer-beta-class-criterion}, the
class is an element of \(H^2(BH;\mathbb F_2)\).  Transfer kills this
group for odd \(H\), and it is zero for \(H=1\).  Lemma~\ref{lem:G-Klein-four}
identifies the \(E[2]\) basis
\(x_1^2,x_1x_2,x_2^2\), so vanishing is equivalent to the three
quadratic coefficients being zero.  Lemma~\ref{lem:G-two-primary-stabilizer}
identifies the two-primary rank-two basis \(y_1,x_1x_2,y_2\), so
vanishing is equivalent to the three displayed coefficients being
zero.  The coefficient \(A_{12}^{(H)}\) is independent of cyclic
order-two restrictions.  Inspection of
\(\texttt{certificates/orientation/rhomred\_finite\_stabilizer\_beta}\)
and \(\texttt{certificates/orientation/k3e\_reduced\_orientation}\)
shows that the current finite evidence supplies the criterion but not
its hypotheses.
\end{proof}

The packet
\[
\texttt{certificates/orientation/rhomred\_finite\_stabilizer\_beta\_vanishing}
\]
verifies
\(\texttt{RHOMRED\_FINITE\_STABILIZER\_BETA\_VANISHING\_OBSTRUCTION\_VERIFIED}\):
it imports the row \(285\) packet and the reduced-orientation
obstruction ledger, records the odd/trivial, Klein-four, and
two-primary zero criteria for row \(286\), checks the three retained
residual groups, and keeps the \(\beta^H\)-value, zero-coordinate,
null-trivialization, linearization, transition, vanishing-cycle, and
protected-integration rows open.

\begin{proposition}[Compatible null-trivialization criterion for
finite-stabilizer Borel classes]
\label{prop:finite-stabilizer-beta-nulltrivialization-criterion}
\textnormal{[criterion proved; compatible null-trivialization rows not
supplied]}
Assume rows \(285\)--\(286\) have supplied, for every retained
finite-stabilizer stratum \(S\) with residual stabilizer \(H\), a
Cech--Borel cocycle
\[
a^H_{\mathcal C,S}\in Z^2(BH;\mathbb F_2)
\]
representing \(\beta^H_{\mathcal C,S}\), and have verified
\(\beta^H_{\mathcal C,S}=0\).  Row \(287\) is the following stronger
cochain-level datum.  It must supply a \(1\)-cochain
\[
h^H_{\mathcal C,S}\in C^1(BH;\mathbb F_2)
\]
with
\[
\delta h^H_{\mathcal C,S}=a^H_{\mathcal C,S}.
\]
For every retained subgroup inclusion
\(\iota\colon K\hookrightarrow H\) it must also verify
\[
\operatorname{res}_{\iota}a^H_{\mathcal C,S}=a^K_{\mathcal C,S}
\]
and supply a comparison \(0\)-cochain
\[
q^{\iota}_{\mathcal C,S}\in C^0(BK;\mathbb F_2)
\]
such that
\[
\operatorname{res}_{\iota}h^H_{\mathcal C,S}
-h^K_{\mathcal C,S}
=\delta q^{\iota}_{\mathcal C,S}.
\]
For every chain \(L\xrightarrow{\jmath}K\xrightarrow{\iota}H\), the
chosen comparisons must have zero defect
\[
q^{\iota\circ\jmath}_{\mathcal C,S}
=q^{\jmath}_{\mathcal C,S}
+\operatorname{res}_{\jmath}q^{\iota}_{\mathcal C,S}
\]
in the chosen Cech--Borel cochain model.  This is the subgroup
restriction compatibility demanded in row \(287\).  It is not the
linearization character of rows \(288\)--\(289\), and it is not implied
by a scalar trace, by the residual group order, or by the abstract
vanishing of \(H^2(BH;\mathbb F_2)\) unless the representative
\(a^H_{\mathcal C,S}\), the cochain \(h^H_{\mathcal C,S}\), and all
restriction comparisons are supplied.

The present finite tables do not construct row \(287\).  They record
three residual groups of order \(1\), but they record no
finite-stabilizer orientation row, no Cech--Borel representative
\(a^H_{\mathcal C,S}\), no verified zero row from row \(286\), no
\(1\)-cochain \(h^H_{\mathcal C,S}\), no retained subgroup-inclusion
table, no comparison \(0\)-cochain, and no chain-coherence row.  Thus
row \(287\) remains an open cochain-level obstruction, even in the
trivial residual-group template.
\end{proposition}

\begin{proof}
A null-trivialization of the representative cocycle
\(a^H_{\mathcal C,S}\) is, by definition, a degree-one cochain whose
coboundary is that representative.  Therefore the cohomological
statement \(\beta^H_{\mathcal C,S}=0\) is only the existence statement;
row \(287\) asks for a chosen primitive.  Restriction along
\(\iota\colon K\hookrightarrow H\) is a cochain map, so
\(\delta(\operatorname{res}_{\iota}h^H-h^K)=0\) once
\(\operatorname{res}_{\iota}a^H=a^K\).  Compatibility requires this
closed \(1\)-cochain to be killed by a chosen \(0\)-cochain
\(q^\iota\), and functoriality over a two-step chain is exactly the
displayed defect-zero equation.  These are cochain equalities, not
statements about the final scalar shadow.  Inspection of
\(\texttt{certificates/orientation/rhomred\_finite\_stabilizer\_beta\_vanishing}\),
\(\texttt{certificates/orientation/k3e\_reduced\_orientation}\), and
\(\texttt{certificates/orientation/transition\_orientation\_preservation}\)
shows that the required cochains and subgroup-restriction rows are not
present.
\end{proof}

The packet
\[
\texttt{certificates/orientation/rhomred\_finite\_stabilizer\_beta\_nulltrivialization}
\]
verifies
\(\texttt{RHOMRED\_FINITE\_STABILIZER\_BETA\_NULLTRIVIALIZATION\_OBSTRUCTION\_VERIFIED}\):
it imports the row \(286\) packet and the reduced-orientation
obstruction ledger, records the cochain and subgroup-restriction
criteria for row \(287\), checks the three retained residual groups,
and keeps the \(\beta^H\)-cocycle, zero row, \(1\)-cochain,
coboundary-defect, subgroup-inclusion, comparison \(0\)-cochain,
chain-coherence, linearization, transition, vanishing-cycle, and
protected-integration rows open.

\begin{proposition}[Computation criterion for the finite-stabilizer
linearization character]
\label{prop:finite-stabilizer-linearization-character-criterion}
\textnormal{[criterion proved; linearization-character rows not
supplied]}
Fix a retained finite-stabilizer stratum \(S\) with residual stabilizer
\(H\).  After the row \(287\) null-trivialization has killed the
degree-two Borel gerbe, the remaining finite-stabilizer descent datum
is a character
\[
\lambda^H_{\mathcal C,S}\in H^1(BH;\mathbb F_2)
=\operatorname{Hom}(H,\mathbb F_2).
\]
Row \(288\) computes this class precisely when it supplies the
orientation-line action of \(H\) on the already trivialized Pfaffian
orientation line, an ordered generator basis for the two-primary part
of \(H\), and the sign values of this action on the chosen generators.
For \(H=1\), and more generally for odd \(H\), the target group
\(H^1(BH;\mathbb F_2)\) is zero after the retained stabilizer row and
orientation-line action have been supplied.  For
\(H\simeq E[2]\simeq(\mathbb Z/2)^2\), with basis \(x_1,x_2\) dual to
the chosen generators, the computation is
\[
\lambda^H_{\mathcal C,S}=\lambda_1x_1+\lambda_2x_2,
\qquad
\rho_1=\lambda_1,\quad \rho_2=\lambda_2,\quad
\rho_3=\lambda_1+\lambda_2
\]
on the three non-zero cyclic subgroups.  For
\(H_{(2)}\simeq(\mathbb Z/2^a)^2\), \(a\ge2\), it is
\[
\lambda^H_{\mathcal C,S}
=\lambda^{(H)}_1x_1+\lambda^{(H)}_2x_2
\in H^1(BH;\mathbb F_2),
\]
where \(x_i\) are the mod-\(2\) degree-one classes dual to the chosen
two-primary generators.  Row \(288\) is only this computation of
\(\lambda^H_{\mathcal C,S}\); row \(289\) is the separate assertion
that the computed character is zero.

The present finite tables do not compute row \(288\).  The determinant
anchor translation-weight packet records the class-level restriction
\(h\mapsto\chi_\eta h\), but it marks coherent stabilizer
linearization as not certified.  The reduced-orientation packet has no
finite-stabilizer rows, no quotient Borel rows, and no
orientation-line action rows.  The row \(287\) packet has no
null-trivializations and records \(\texttt{linearization\_computed}=
\texttt{false}\) on every retained residual group.  Thus the equality
\(H^1(B1;\mathbb F_2)=0\), cyclic quadratic restrictions, and
translation-weight formulas are templates or separate data; none is a
computed value of \(\lambda^H_{\mathcal C,S}\) for row \(288\).
\end{proposition}

\begin{proof}
Once the degree-two gerbe has been trivialized, two finite-stabilizer
linearizations of the same orientation line differ by a sign character
of \(H\).  This torsor is \(H^1(BH;\mathbb F_2)\), so a computation of
the residual class is exactly a computation of the sign character on
generators.  Transfer gives \(H^1(BH;\mathbb F_2)=0\) for odd \(H\),
and the trivial group has no positive cohomology.  Lemma~\ref{lem:G-Klein-four}
identifies the \(E[2]\) basis and the three cyclic restrictions
\(\rho_1,\rho_2,\rho_3\).  Lemma~\ref{lem:G-two-primary-stabilizer}
identifies the degree-one basis \(x_1,x_2\) for the two-primary rank
two case.  These are coordinate computations for a supplied action on
the orientation line.  Inspection of
\(\texttt{certificates/orientation/rhomred\_finite\_stabilizer\_beta\_nulltrivialization}\),
\(\texttt{certificates/orientation/k3e\_reduced\_orientation}\), and
\(\texttt{certificates/hybrid/determinant\_anchor\_translation\_weight}\)
shows that the action rows, generator-basis rows, and character values
are not present.
\end{proof}

The packet
\[
\texttt{certificates/orientation/rhomred\_finite\_stabilizer\_linearization\_character}
\]
verifies
\(\texttt{RHOMRED\_FINITE\_STABILIZER\_LINEARIZATION\_CHARACTER\_OBSTRUCTION\_VERIFIED}\):
it imports the row \(287\) packet, the reduced-orientation obstruction
ledger, and the determinant-anchor translation-weight packet, records
the \(H^1(BH;\mathbb F_2)\) coordinate criterion for row \(288\),
checks the three retained residual groups, and keeps the
orientation-line action, generator-basis, character-value,
\(H^1\)-coordinate, row \(289\) vanishing, transition,
vanishing-cycle, and protected-integration rows open.

\begin{proposition}[Vanishing criterion for the finite-stabilizer
linearization character]
\label{prop:finite-stabilizer-linearization-vanishing-criterion}
\textnormal{[criterion proved; zero-linearization rows not supplied]}
Assume row \(288\) has computed
\[
\lambda^H_{\mathcal C,S}\in H^1(BH;\mathbb F_2)
\]
for every retained finite-stabilizer stratum.  Row \(289\) proves
\[
\lambda^H_{\mathcal C,S}=0
\]
for every retained \(H\) precisely when it verifies the corresponding
zero coordinate row in the \(H^1(BH;\mathbb F_2)\)-basis:
\[
H=1\ \text{or}\ |H|\ \text{odd}:\quad
H^1(BH;\mathbb F_2)=0
\quad\text{after the retained action row is supplied,}
\]
\[
H\simeq E[2]:\quad \lambda_1=\lambda_2=0
\quad\text{equivalently}\quad
\rho_1=\rho_2=\rho_3=0
\]
with \(\rho_3=\lambda_1+\lambda_2\), and, for two-primary rank two
stabilizers,
\[
H_{(2)}\simeq(\mathbb Z/2^a)^2,\ a\ge2:
\quad \lambda^{(H)}_1=\lambda^{(H)}_2=0.
\]
This is a degree-one vanishing statement.  It is not implied by
\(\beta^H_{\mathcal C,S}=0\), by a row \(287\)
null-trivialization of the degree-two gerbe, by the determinant-anchor
translation-weight formula, or by class invariance in
\(\operatorname{Pic}\).  Row \(290\) records this independence as a
separate firewall.

The present finite tables do not prove row \(289\).  They record no
orientation-line action, no generator-basis row, no character values,
and no computed \(H^1\)-coordinate for \(\lambda^H_{\mathcal C,S}\).
The row \(288\) packet records three order-one residual group
templates, but the corresponding \(\lambda\)-values are
\(\texttt{missing}\), with
\(\texttt{lambda\_character\_computed}=\texttt{false}\) and
\(\texttt{lambda\_vanishing\_verified}=\texttt{false}\).  Hence the
zero \(H^1(B1;\mathbb F_2)\) template cannot be promoted to a
zero-linearization row for the retained orientation problem.
\end{proposition}

\begin{proof}
Once row \(288\) has produced an element of \(H^1(BH;\mathbb F_2)\),
vanishing is a coordinate calculation.  The groups \(H=1\) and odd
\(H\) have zero mod-\(2\) degree-one cohomology.  Lemma~\ref{lem:G-Klein-four}
identifies the \(E[2]\) character as
\(\lambda_1x_1+\lambda_2x_2\) and the cyclic restrictions as
\(\rho_1=\lambda_1\), \(\rho_2=\lambda_2\),
\(\rho_3=\lambda_1+\lambda_2\); all three cyclic restrictions vanish
if and only if \(\lambda_1=\lambda_2=0\).  Lemma~\ref{lem:G-two-primary-stabilizer}
identifies the two-primary degree-one basis \(x_1,x_2\), so vanishing
is equivalent to the two displayed coefficients being zero.  These
statements require the computed character from row \(288\).  Inspection
of
\(\texttt{certificates/orientation/rhomred\_finite\_stabilizer\_linearization\_character}\)
and \(\texttt{certificates/orientation/k3e\_reduced\_orientation}\)
shows that no such character value or zero-coordinate row is present.
\end{proof}

The packet
\[
\texttt{certificates/orientation/rhomred\_finite\_stabilizer\_linearization\_vanishing}
\]
verifies
\(\texttt{RHOMRED\_FINITE\_STABILIZER\_LINEARIZATION\_VANISHING\_OBSTRUCTION\_VERIFIED}\):
it imports the row \(288\) packet and the reduced-orientation
obstruction ledger, records the \(H^1(BH;\mathbb F_2)\) zero criteria
for row \(289\), checks the three retained residual groups, and keeps
the computed-character, zero-coordinate, all-retained-coverage,
transition, vanishing-cycle, and protected-integration rows open.

\begin{proposition}[Gerbe vanishing does not imply zero
finite-stabilizer linearization]
\label{prop:finite-stabilizer-gerbe-linearization-independence}
\textnormal{[firewall proved]}
The implication
\[
\beta^H_{\mathcal C,S}=0\quad\Longrightarrow\quad
\lambda^H_{\mathcal C,S}=0
\]
is false as a finite-stabilizer descent principle.  For
\(H\simeq E[2]\), the equations
\[
b_{20}=b_{11}=b_{02}=0
\]
kill the degree-two class
\(\beta^H_{\mathcal C,S}\in H^2(BH;\mathbb F_2)\).  The degree-one
class
\[
\lambda^H_{\mathcal C,S}=\lambda_1x_1+\lambda_2x_2
\in H^1(BH;\mathbb F_2)
\]
is independent; for example \((\lambda_1,\lambda_2)=(1,0)\) gives
\[
\rho_1=1,\qquad \rho_2=0,\qquad \rho_3=1,
\]
so the gerbe vanishes and the linearization does not.  For
\(H_{(2)}\simeq(\mathbb Z/2^a)^2\), \(a\ge2\), the same separation is
\[
A_1^{(H)}=A_{12}^{(H)}=A_2^{(H)}=0,\qquad
(\lambda^{(H)}_1,\lambda^{(H)}_2)=(1,0).
\]
Thus rows \(286\)--\(287\) cannot replace rows \(288\)--\(289\).  If
\(H=1\) or \(|H|\) is odd, then \(H^1(BH;\mathbb F_2)=0\); this is a
separate group-cohomological zero target, not a consequence of
\(\beta^H_{\mathcal C,S}=0\).

The present finite tables obey this firewall.  The row \(289\) packet
records \(\texttt{beta\_vanishing\_implies\_lambda\_vanishing}=
\texttt{false}\) and
\(\texttt{nulltrivialization\_implies\_lambda\_vanishing}=
\texttt{false}\); it also records missing \(\lambda\)-values and no
zero-linearization rows.  Hence no scalar trace, determinant-anchor
translation weight, gerbe vanishing row, or null-trivialization row is
used as a substitute for the degree-one linearization calculation.
\end{proposition}

\begin{proof}
The cohomological degrees already separate the two obstructions.
Lemma~\ref{lem:G-Klein-four} gives
\[
H^*(BE[2];\mathbb F_2)=\mathbb F_2[x_1,x_2],\qquad |x_i|=1.
\]
The finite gerbe lives in degree \(2\), with coordinates
\((b_{20},b_{11},b_{02})\), while the residual linearization lives in
degree \(1\), with coordinates \((\lambda_1,\lambda_2)\).  These are
coordinates on different graded pieces, so the zero locus of the
degree-two coordinates contains every value of the degree-one
coordinates.  The displayed choice \((\lambda_1,\lambda_2)=(1,0)\)
is an explicit counterexample to the implication.  Lemma~\ref{lem:G-two-primary-stabilizer}
gives the same degree separation for
\((\mathbb Z/2^a)^2\): the gerbe coordinates are
\((A_1^{(H)},A_{12}^{(H)},A_2^{(H)})\) in degree \(2\), and the
linearization coordinates are
\((\lambda^{(H)}_1,\lambda^{(H)}_2)\) in degree \(1\).  The trivial
and odd cases have zero degree-one target by transfer or by
triviality; that is not an implication from the degree-two class.
\end{proof}

The packet
\[
\texttt{certificates/orientation/rhomred\_finite\_stabilizer\_gerbe\_linearization\_independence}
\]
verifies
\(\texttt{RHOMRED\_FINITE\_STABILIZER\_GERBE\_LINEARIZATION\_INDEPENDENCE\_VERIFIED}\):
it imports the row \(289\) packet and the reduced-orientation
obstruction ledger, records explicit \(E[2]\) and two-primary
counterexamples to the false implication, records the trivial/odd
zero-target exception, and checks that no current orientation packet
uses a gerbe, scalar, or determinant-anchor row as a zero-linearization
row.

\begin{proposition}[\(E\lbrack2\rbrack\) Klein-four finite-stabilizer class]
\label{prop:finite-stabilizer-klein-four-class}
\textnormal{[local computation proved; retained \(E[2]\)-edge row not
supplied]}
Let \(H=E[2]\simeq(\mathbb Z/2)^2\), choose generators \(e_1,e_2\),
and let \(x_1,x_2\in H^1(BH;\mathbb F_2)\) be the dual basis.  After
the finite-stabilizer Cech--Borel class on a retained \(BE[2]\)-stratum
has been reduced to the classifying-space edge, its degree-two part is
uniquely of the form
\[
\beta^{E,2}_{\mathcal C,S}
=b_{20}x_1^2+b_{11}x_1x_2+b_{02}x_2^2
\in H^2(BE[2];\mathbb F_2).
\]
This is row \(291\).  The symbols \(b_{20},b_{11},b_{02}\) are the
coordinates of the supplied edge class in the ordered basis
\[
x_1^2,\quad x_1x_2,\quad x_2^2.
\]
Rows \(292\)--\(294\) are the subsequent restriction computations
which express these coordinates through the three cyclic restrictions
\(r_1,r_2,r_3\).  Row \(291\) does not by itself supply a retained
\(E[2]\)-stratum, an equivariant Cech--Borel representative, an
edge-reduction row, or the values of the three coefficients.
\end{proposition}

\begin{proof}
For a complex elliptic curve, \(E[2]\simeq(\mathbb Z/2)^2\).  The
classifying-space cohomology is
\[
H^*(BE[2];\mathbb F_2)=\mathbb F_2[x_1,x_2],\qquad |x_i|=1.
\]
Therefore its degree-two part has basis
\[
H^2(BE[2];\mathbb F_2)
=\mathbb F_2\langle x_1^2,x_1x_2,x_2^2\rangle.
\]
An edge-reduced finite-stabilizer Borel class is an element of this
three-dimensional vector space, hence has the unique displayed
coordinate expansion.  The proof is a computation in
\(H^*(B(\mathbb Z/2)^2;\mathbb F_2)\); it does not assert that the
current finite tables have supplied the geometric edge class whose
coordinates are \(b_{20},b_{11},b_{02}\).
\end{proof}

The packet
\[
\texttt{certificates/orientation/rhomred\_finite\_stabilizer\_klein\_four\_class}
\]
verifies
\(\texttt{RHOMRED\_FINITE\_STABILIZER\_KLEIN\_FOUR\_CLASS\_VERIFIED}\):
it imports the finite-stabilizer beta packet, the row \(290\) firewall,
and the reduced-orientation obstruction ledger, records the
\(H^2(BE[2];\mathbb F_2)\) basis and the displayed Klein-four edge
class, and keeps the retained \(E[2]\)-stratum, edge-reduction,
coefficient-value, cyclic-restriction, vanishing, linearization,
transition, vanishing-cycle, and protected-integration rows open.

\begin{proposition}[First cyclic restriction of the \(E\lbrack2\rbrack\)
Klein-four class]
\label{prop:finite-stabilizer-klein-four-b20}
\textnormal{[local computation proved; first cyclic restriction row
not supplied]}
Keep the notation of
Proposition~\ref{prop:finite-stabilizer-klein-four-class}.  Let
\(\iota_1\colon \langle e_1\rangle\hookrightarrow E[2]\) be the first
cyclic subgroup and let
\[
t\in H^1(B\langle e_1\rangle;\mathbb F_2)
\]
be its degree-one generator.  The restriction map satisfies
\[
\iota_1^*(x_1)=t,\qquad \iota_1^*(x_2)=0.
\]
Therefore
\[
\iota_1^*\beta^{E,2}_{\mathcal C,S}=b_{20}t^2.
\]
If the first cyclic restriction is written
\[
\iota_1^*\beta^{E,2}_{\mathcal C,S}=r_1t^2,
\]
then
\[
b_{20}=r_1.
\]
This is row \(292\).  It is the first coefficient extraction from the
row \(291\) class.  It does not supply the retained \(E[2]\)-edge row,
the first cyclic restriction row, or the value of \(r_1\).
\end{proposition}

\begin{proof}
The inclusion \(\iota_1\) sends the generator of
\(\langle e_1\rangle\) to \(e_1\).  Hence the dual class \(x_1\)
restricts to the generator \(t\), while \(x_2\) restricts to zero.
Applying \(\iota_1^*\) to
\[
\beta^{E,2}_{\mathcal C,S}
=b_{20}x_1^2+b_{11}x_1x_2+b_{02}x_2^2
\]
gives \(b_{20}t^2\).  Since
\(H^2(B\langle e_1\rangle;\mathbb F_2)=\mathbb F_2\langle t^2\rangle\),
equality with \(r_1t^2\) is equivalent to equality of coefficients,
namely \(b_{20}=r_1\).
\end{proof}

The packet
\[
\texttt{certificates/orientation/rhomred\_finite\_stabilizer\_klein\_four\_b20}
\]
verifies
\(\texttt{RHOMRED\_FINITE\_STABILIZER\_KLEIN\_FOUR\_B20\_VERIFIED}\):
it imports the row \(291\) Klein-four class packet, records the first
cyclic restriction map \(x_1\mapsto t,\ x_2\mapsto0\), verifies the
coefficient identity \(b_{20}=r_1\), and keeps the retained edge row,
the first cyclic restriction value, the \(b_{11}\) and \(b_{02}\)
restriction computations, vanishing, linearization, transition,
vanishing-cycle, and protected-integration rows open.

\begin{proposition}[Diagonal cyclic restriction of the \(E\lbrack2\rbrack\)
Klein-four class]
\label{prop:finite-stabilizer-klein-four-b11}
\textnormal{[local computation proved; cyclic restriction rows not
supplied]}
Keep the notation of
Propositions~\ref{prop:finite-stabilizer-klein-four-class}
and~\ref{prop:finite-stabilizer-klein-four-b20}.  Let
\[
\iota_2\colon \langle e_2\rangle\hookrightarrow E[2],
\qquad
\iota_3\colon \langle e_1+e_2\rangle\hookrightarrow E[2]
\]
be the second and diagonal cyclic subgroups, and let \(t\) denote the
degree-one generator on either cyclic subgroup.  The restriction maps
satisfy
\[
\iota_2^*(x_1)=0,\qquad \iota_2^*(x_2)=t,
\]
and
\[
\iota_3^*(x_1)=t,\qquad \iota_3^*(x_2)=t.
\]
Thus
\[
\iota_2^*\beta^{E,2}_{\mathcal C,S}=b_{02}t^2,\qquad
\iota_3^*\beta^{E,2}_{\mathcal C,S}
=(b_{20}+b_{11}+b_{02})t^2.
\]
If the three cyclic restrictions are written
\[
\iota_i^*\beta^{E,2}_{\mathcal C,S}=r_it^2,\qquad i=1,2,3,
\]
where \(i=3\) denotes the diagonal subgroup, then
\[
b_{11}=r_1+r_2+r_3
\]
in \(\mathbb F_2\).  This is row \(293\).  It does not supply the
retained \(E[2]\)-edge row, the three cyclic restriction rows, the
values of \(r_1,r_2,r_3\), or the value of \(b_{11}\).  Row \(294\)
records the second-square extraction \(b_{02}=r_2\) as its own ledger
row.
\end{proposition}

\begin{proof}
The inclusion \(\iota_2\) kills the dual class \(x_1\) and sends
\(x_2\) to the cyclic generator \(t\).  The inclusion \(\iota_3\)
sends the cyclic generator to \(e_1+e_2\), so both dual classes
restrict to \(t\).  Applying these restrictions to
\[
\beta^{E,2}_{\mathcal C,S}
=b_{20}x_1^2+b_{11}x_1x_2+b_{02}x_2^2
\]
gives the two displayed formulae.  The first cyclic restriction gives
\(b_{20}=r_1\) by
Proposition~\ref{prop:finite-stabilizer-klein-four-b20}; the second
restriction gives \(b_{02}=r_2\); and the diagonal restriction gives
\[
b_{20}+b_{11}+b_{02}=r_3.
\]
Substitution yields \(r_1+b_{11}+r_2=r_3\).  Since the coefficient
field is \(\mathbb F_2\), this is equivalent to
\[
b_{11}=r_1+r_2+r_3 .
\]
All equalities here are identities in the cohomology of cyclic
classifying spaces.  They do not assert that the current finite
tables have supplied the geometric cyclic restrictions or their
values.
\end{proof}

The packet
\[
\texttt{certificates/orientation/rhomred\_finite\_stabilizer\_klein\_four\_b11}
\]
verifies
\(\texttt{RHOMRED\_FINITE\_STABILIZER\_KLEIN\_FOUR\_B11\_VERIFIED}\):
it imports the row \(291\) and row \(292\) packets, records the second
and diagonal cyclic restriction maps, verifies the mixed coefficient
identity \(b_{11}=r_1+r_2+r_3\), and keeps the retained edge row,
the cyclic restriction values, the standalone \(b_{02}\) row,
vanishing, linearization, transition, vanishing-cycle, and
protected-integration rows open.

\begin{proposition}[Second cyclic restriction of the \(E\lbrack2\rbrack\)
Klein-four class]
\label{prop:finite-stabilizer-klein-four-b02}
\textnormal{[local computation proved; second cyclic restriction row
not supplied]}
Keep the notation of
Proposition~\ref{prop:finite-stabilizer-klein-four-class}.  Let
\(\iota_2\colon \langle e_2\rangle\hookrightarrow E[2]\) be the second
cyclic subgroup and let
\[
t\in H^1(B\langle e_2\rangle;\mathbb F_2)
\]
be its degree-one generator.  The restriction map satisfies
\[
\iota_2^*(x_1)=0,\qquad \iota_2^*(x_2)=t.
\]
Therefore
\[
\iota_2^*\beta^{E,2}_{\mathcal C,S}=b_{02}t^2.
\]
If the second cyclic restriction is written
\[
\iota_2^*\beta^{E,2}_{\mathcal C,S}=r_2t^2,
\]
then
\[
b_{02}=r_2.
\]
This is row \(294\).  It is the standalone second-square coefficient
extraction from the row \(291\) class.  It does not supply the
retained \(E[2]\)-edge row, the second cyclic restriction row, or the
value of \(r_2\).
\end{proposition}

\begin{proof}
The inclusion \(\iota_2\) sends the generator of
\(\langle e_2\rangle\) to \(e_2\).  Hence \(x_1\) restricts to zero
and \(x_2\) restricts to the generator \(t\).  Applying \(\iota_2^*\)
to
\[
\beta^{E,2}_{\mathcal C,S}
=b_{20}x_1^2+b_{11}x_1x_2+b_{02}x_2^2
\]
gives \(b_{02}t^2\).  Since
\(H^2(B\langle e_2\rangle;\mathbb F_2)=\mathbb F_2\langle t^2\rangle\),
equality with \(r_2t^2\) is equivalent to \(b_{02}=r_2\).  This is a
classifying-space calculation only; it does not assert that the
finite-stabilizer tables have supplied the retained cyclic
restriction or its value.
\end{proof}

The packet
\[
\texttt{certificates/orientation/rhomred\_finite\_stabilizer\_klein\_four\_b02}
\]
verifies
\(\texttt{RHOMRED\_FINITE\_STABILIZER\_KLEIN\_FOUR\_B02\_VERIFIED}\):
it imports the row \(291\) and row \(293\) packets, records the second
cyclic restriction map \(x_1\mapsto0,\ x_2\mapsto t\), verifies the
coefficient identity \(b_{02}=r_2\), and keeps the retained edge row,
the second cyclic restriction value, vanishing, linearization,
transition, vanishing-cycle, and protected-integration rows open.

\begin{proposition}[Klein-four zero-restriction criterion]
\label{prop:finite-stabilizer-klein-four-zero-restrictions}
\textnormal{[criterion proved; zero restriction rows not supplied]}
Keep the notation of
Propositions~\ref{prop:finite-stabilizer-klein-four-class},
\ref{prop:finite-stabilizer-klein-four-b20},
\ref{prop:finite-stabilizer-klein-four-b11}, and
\ref{prop:finite-stabilizer-klein-four-b02}.  Let
\[
\iota_i^*\beta^{E,2}_{\mathcal C,S}=r_it^2,\qquad i=1,2,3,
\]
where \(i=1\) is the subgroup \(\langle e_1\rangle\), \(i=2\) is
\(\langle e_2\rangle\), and \(i=3\) is
\(\langle e_1+e_2\rangle\).  Then
\[
\beta^{E,2}_{\mathcal C,S}=0
\quad\Longleftrightarrow\quad
r_1=r_2=r_3=0.
\]
Equivalently, the three cyclic zero restrictions are necessary and
sufficient for the vanishing of the classifying-space edge of the
\(E[2]\)-finite-stabilizer Borel class.  This is the algebraic content
of row \(295\).  The present retained finite-stabilizer tables do not
yet supply the three cyclic restriction rows or the values
\(r_1,r_2,r_3\), so this proposition is not a geometric proof that
those values vanish on the retained \(K3\times E\) strata.
\end{proposition}

\begin{proof}
Rows \(292\)--\(294\) give
\[
b_{20}=r_1,\qquad b_{02}=r_2,\qquad
b_{11}=r_1+r_2+r_3
\]
in \(\mathbb F_2\).  If \(r_1=r_2=r_3=0\), then
\(b_{20}=b_{02}=b_{11}=0\), hence
\[
\beta^{E,2}_{\mathcal C,S}
=b_{20}x_1^2+b_{11}x_1x_2+b_{02}x_2^2
\]
vanishes in \(H^2(BE[2];\mathbb F_2)\).  Conversely, if
\(\beta^{E,2}_{\mathcal C,S}=0\), then the coefficients in the basis
\(x_1^2,x_1x_2,x_2^2\) vanish.  Thus \(r_1=b_{20}=0\),
\(r_2=b_{02}=0\), and
\[
r_3=b_{20}+b_{11}+b_{02}=0.
\]
The implication is a classifying-space calculation.  A proof of the
geometric assertion \(r_1=r_2=r_3=0\) still requires retained cyclic
restriction rows, or an equivalent edge-reduced vanishing row, in the
finite-stabilizer orientation tables.
\end{proof}

The packet
\[
\texttt{certificates/orientation/rhomred\_finite\_stabilizer\_klein\_four\_zero\_restrictions}
\]
verifies
\(\texttt{RHOMRED\_FINITE\_STABILIZER\_KLEIN\_FOUR\_ZERO\_RESTRICTIONS\_OBSTRUCTION\_VERIFIED}\):
it imports the row \(291\)--\(294\) packets, proves the equivalence
between the three zero cyclic restrictions and vanishing of the
Klein-four classifying-space edge, and records that the retained
cyclic zero rows, the \(r_i\)-values, finite-stabilizer beta
vanishing, quotient orientation, transition compatibility,
vanishing-cycle, and protected-integration rows remain open.

\begin{proposition}[Even two-primary finite-stabilizer class]
\label{prop:finite-stabilizer-two-primary-class}
\textnormal{[local computation proved; retained two-primary row not
supplied]}
Let \(N\ge4\) be even and write \(2^a\parallel N\), \(a\ge2\).  The
two-primary part of the \(N\)-torsion stabilizer of a complex elliptic
curve is
\[
E[N]_{(2)}\simeq(\mathbb Z/2^a)^2.
\]
Choose generators and let
\[
x_1,x_2\in H^1(B(\mathbb Z/2^a)^2;\mathbb F_2),
\qquad
y_1,y_2\in H^2(B(\mathbb Z/2^a)^2;\mathbb F_2)
\]
be the corresponding exterior and polynomial generators.  Then
\[
H^*\!\left(B(\mathbb Z/2^a)^2;\mathbb F_2\right)
=\Lambda(x_1,x_2)\otimes\mathbb F_2[y_1,y_2],
\qquad |x_i|=1,\quad |y_i|=2,
\]
and
\[
H^2\!\left(B(\mathbb Z/2^a)^2;\mathbb F_2\right)
=\mathbb F_2\langle y_1,x_1x_2,y_2\rangle.
\]
After the finite-stabilizer Cech--Borel class on a retained
two-primary stratum has been reduced to the classifying-space edge,
its degree-two part is uniquely of the form
\[
\beta^{E,N}_{\mathcal C,S}
=A_1^{(N)}y_1+A_{12}^{(N)}x_1x_2+A_2^{(N)}y_2.
\]
This is row \(296\).  The coefficients
\(A_1^{(N)},A_{12}^{(N)},A_2^{(N)}\) are coordinates of the supplied
edge class in the ordered basis \(y_1,x_1x_2,y_2\).  Row \(296\) does
not supply a retained two-primary stratum, an equivariant
Cech--Borel representative, an edge-reduction row, or the values of
the three coefficients.  Row \(297\) is the separate mixed-term
detection check for \(A_{12}^{(N)}\).
\end{proposition}

\begin{proof}
For a complex elliptic curve, \(E[N]\simeq(\mathbb Z/N)^2\).  Primary
decomposition gives the displayed two-primary subgroup.  For
\(a\ge2\),
\[
H^*(B\mathbb Z/2^a;\mathbb F_2)=\Lambda(x)\otimes\mathbb F_2[y],
\qquad |x|=1,\quad |y|=2.
\]
The K\"unneth formula gives
\[
H^*\!\left(B(\mathbb Z/2^a)^2;\mathbb F_2\right)
=\Lambda(x_1,x_2)\otimes\mathbb F_2[y_1,y_2].
\]
Its degree-two part has basis \(y_1,x_1x_2,y_2\).  Therefore an
edge-reduced two-primary finite-stabilizer Borel class is a unique
linear combination of these three basis vectors.  The argument is
only a classifying-space computation; it does not assert that the
current finite tables have supplied the retained edge class whose
coordinates are \(A_1^{(N)},A_{12}^{(N)},A_2^{(N)}\).
\end{proof}

The packet
\[
\texttt{certificates/orientation/rhomred\_finite\_stabilizer\_two\_primary\_class}
\]
verifies
\(\texttt{RHOMRED\_FINITE\_STABILIZER\_TWO\_PRIMARY\_CLASS\_VERIFIED}\):
it imports the finite-stabilizer beta packet and the reduced
orientation obstruction ledger, records the
\(\Lambda(x_1,x_2)\otimes\mathbb F_2[y_1,y_2]\) cohomology ring,
the degree-two basis \(y_1,x_1x_2,y_2\), and the displayed
two-primary edge-class form, and keeps the retained stratum,
edge-reduction, coefficient-value, mixed-term detection, vanishing,
linearization, transition, vanishing-cycle, and protected-integration
rows open.

\begin{proposition}[Detection of the two-primary mixed term]
\label{prop:finite-stabilizer-two-primary-mixed-term}
\textnormal{[local detection proved; retained \(A_{12}\)-value not
supplied]}
Keep the notation of
Proposition~\ref{prop:finite-stabilizer-two-primary-class}.  Define
the coefficient functional
\[
\pi_{12}\colon
H^2\!\left(B(\mathbb Z/2^a)^2;\mathbb F_2\right)\longrightarrow
\mathbb F_2
\]
by
\[
\pi_{12}(y_1)=0,\qquad
\pi_{12}(x_1x_2)=1,\qquad
\pi_{12}(y_2)=0.
\]
Then
\[
\pi_{12}\!\left(\beta^{E,N}_{\mathcal C,S}\right)=A_{12}^{(N)}.
\]
For every cyclic subgroup \(C\subset(\mathbb Z/2^a)^2\), the
restriction map satisfies
\[
\operatorname{res}^{(\mathbb Z/2^a)^2}_{C}(x_1x_2)=0.
\]
Thus the mixed coefficient \(A_{12}^{(N)}\) is detected only by the
rank-two classifying-space edge, or by an equivalent rank-two
coefficient row; it is not determined by cyclic restrictions.  This
is row \(297\).  It does not supply the retained two-primary edge row
or the value of \(A_{12}^{(N)}\).
\end{proposition}

\begin{proof}
The first assertion is the definition of the coordinate functional in
the ordered basis \(y_1,x_1x_2,y_2\).  It remains to check the cyclic
restriction statement.  Let \(C\subset(\mathbb Z/2^a)^2\) be cyclic.
If \(|C|=2\), then \(C\subset 2(\mathbb Z/2^a)^2\) because \(a\ge2\);
therefore the two mod-\(2\) characters \(x_1,x_2\) restrict trivially
to \(C\).  If \(|C|\ge4\), then
\[
H^*(BC;\mathbb F_2)=\Lambda(u)\otimes\mathbb F_2[v],
\qquad |u|=1,\quad |v|=2,
\]
and both \(x_1\) and \(x_2\) restrict to scalar multiples of \(u\).
Hence their product restricts to a scalar multiple of \(u^2=0\).
Thus \(x_1x_2\) restricts to zero on every cyclic subgroup.  The
coefficient \(A_{12}^{(N)}\) can therefore not be recovered from
cyclic restrictions of the degree-two class.
\end{proof}

The packet
\[
\texttt{certificates/orientation/rhomred\_finite\_stabilizer\_two\_primary\_mixed\_term\_detection}
\]
verifies
\(\texttt{RHOMRED\_FINITE\_STABILIZER\_TWO\_PRIMARY\_MIXED\_TERM\_DETECTION\_VERIFIED}\):
it imports the row \(296\) two-primary class packet, records the
rank-two coefficient functional \(\pi_{12}\), proves that cyclic
subgroup restrictions annihilate \(x_1x_2\), and keeps the retained
edge row, \(A_{12}^{(N)}\)-value, vanishing, linearization,
transition, vanishing-cycle, and protected-integration rows open.

\begin{proposition}[Odd-order finite-stabilizer transfer]
\label{prop:finite-stabilizer-odd-transfer}
\textnormal{[local transfer proved; retained odd edge row not
supplied]}
Let \(H\subset E[N]\) be a finite stabilizer of odd order.  Then
\[
H^i(BH;\mathbb F_2)=0,\qquad i>0.
\]
Consequently, after the finite-stabilizer Cech--Borel class on a
retained odd-order stratum has been reduced to the
classifying-space edge,
\[
\beta^H_{\mathcal C,S}\in H^2(BH;\mathbb F_2)
\]
vanishes.  The residual degree-one stabilizer character also vanishes
after edge reduction, since \(H^1(BH;\mathbb F_2)=0\).  This is row
\(298\).  It proves the odd-order target calculation; it does not
supply a retained odd-order stratum, an equivariant Cech--Borel
representative, an edge-reduction row, an all-retained-stabilizer
coverage row, or the compatible null-trivialisations required for
(O1).
\end{proposition}

\begin{proof}
Since \(|H|\) is odd, multiplication by \(|H|\) is the identity on
\(\mathbb F_2\).  The restriction-transfer identity in group
cohomology gives multiplication by \(|H|\) on
\(H^i(BH;\mathbb F_2)\), while restriction to the point is zero in
positive degree.  Thus every positive-degree class in
\(H^*(BH;\mathbb F_2)\) is zero.  Equivalently, the averaging
idempotent makes the invariants functor exact because \(|H|\) is
invertible in \(\mathbb F_2\).  In particular, the degree-two
classifying-space edge of the finite-stabilizer Borel class and the
degree-one residual character both vanish once the geometric data have
been reduced to those targets.
\end{proof}

The packet
\[
\texttt{certificates/orientation/rhomred\_finite\_stabilizer\_odd\_transfer}
\]
verifies
\(\texttt{RHOMRED\_FINITE\_STABILIZER\_ODD\_TRANSFER\_VERIFIED}\):
it imports the finite-stabilizer beta and beta-vanishing obstruction
packets, records the odd-order transfer calculation
\(H^{>0}(BH;\mathbb F_2)=0\), and keeps the retained odd stratum,
edge-reduction, all-retained coverage, null-trivialisation,
transition, vanishing-cycle, and protected-integration rows open.

\begin{proposition}[Finite-stabilizer orientation detection;
\textit{proved after edge reduction}]
\label{prop:finite-stabilizer-orientation-detection}
Let \(S\) be a retained finite-stabilizer stratum with stabilizer
\(H\subset E_{\mathrm{tors}}\).  Suppose the equivariant
Cech--Borel orientation class
\(\widetilde\beta^H_{\mathcal C,S}\in H^2_H(S;\F_2)\) has been
reduced to its classifying-space edge, i.e. its ordinary component in
\(H^2(S;\F_2)^H\), mixed component in
\(H^1(BH;H^1(S;\F_2))\), and incoming spectral-sequence differential
terms vanish with chosen null-trivialisations.  Then the finite
stabilizer part of (O1) is exactly the following finite calculation.

If \(H\) has odd order, the mod-\(2\) classifying-space edge vanishes.
If \(H\simeq E[2]\simeq(\Z/2)^2\), write
\[
\beta^H_{\mathcal C,S}
=b_{20}x_1^2+b_{11}x_1x_2+b_{02}x_2^2,\qquad
\lambda^H_{\mathcal C,S}=\lambda_1x_1+\lambda_2x_2.
\]
Then quotient orientation on this stratum is equivalent to
\[
b_{20}=b_{11}=b_{02}=\lambda_1=\lambda_2=0,
\]
or equivalently to the three cyclic quadratic restrictions and the
three cyclic linearization restrictions all vanishing, as in
Lemma~\ref{lem:G-Klein-four}.  If the two-primary part of \(H\) is
\((\Z/2^a)^2\), \(a\ge2\), write
\[
\beta^H_{\mathcal C,S}
=A_1^{(H)}y_1+A_{12}^{(H)}x_1x_2+A_2^{(H)}y_2,\qquad
\lambda^H_{\mathcal C,S}=\lambda_1^{(H)}x_1+\lambda_2^{(H)}x_2.
\]
Then quotient orientation is equivalent to the five equations
\[
A_1^{(H)}=A_{12}^{(H)}=A_2^{(H)}
=\lambda_1^{(H)}=\lambda_2^{(H)}=0.
\]
The mixed coefficient \(A_{12}^{(H)}\) is not detected by cyclic
order-two restrictions.  Therefore a finite-stabilizer row which lists
only cyclic restrictions is not an (O1) quotient-orientation row unless
the edge-reduction and mixed-coefficient computations are supplied.
\end{proposition}

\begin{proof}
The edge-reduction hypothesis identifies the relevant obstruction with
an element of \(H^2(BH;\F_2)\), and the residual linearization with an
element of \(H^1(BH;\F_2)\).  For odd \(H\), transfer kills positive
degree cohomology with \(\F_2\)-coefficients.  For \(E[2]\), the
coordinate calculation is Lemma~\ref{lem:G-Klein-four}: the three
nonzero cyclic restrictions invert the quadratic coefficients, and the
linearization character is an independent \(H^1\)-torsor.  For
\((\Z/2^a)^2\), \(a\ge2\), Lemma~\ref{lem:G-two-primary-stabilizer}
gives
\[
H^2(B(\Z/2^a)^2;\F_2)
=\F_2\langle y_1,x_1x_2,y_2\rangle.
\]
The term \(x_1x_2\) pulls back trivially to cyclic order-two
subgroups, so it must be computed on the rank-two stabilizer.  These
are exactly the equations in the statement.
\end{proof}

\begin{proposition}[Direct-sum orientation multiplicativity criterion]
\label{prop:orientation-direct-sum-multiplicativity}
\textnormal{[Picard-groupoid criterion proved; direct-sum
Thom--Sebastiani row not supplied]}
Let \(\mathcal C,\mathcal D\) be retained objects at a fixed finite
HN height.  Suppose that the cosection-reduced self-Ext construction
has been supplied on the direct-sum chart and is compatible with the
block decomposition
\[
K^{\mathrm{red}}_{\mathcal C\oplus\mathcal D}\simeq
K^{\mathrm{red}}_{\mathcal C}\oplus
K^{\mathrm{red}}_{\mathcal D}\oplus
K^{\mathrm{red}}_{\mathcal C,\mathcal D}\oplus
K^{\mathrm{red}}_{\mathcal D,\mathcal C}.
\]
Then the determinant functor gives
\[
\mathcal L^{\det}_{\mathcal C\oplus\mathcal D}\simeq
\mathcal L^{\det}_{\mathcal C}\otimes
\mathcal L^{\det}_{\mathcal D}\otimes
\det K^{\mathrm{red}}_{\mathcal C,\mathcal D}\otimes
\det K^{\mathrm{red}}_{\mathcal D,\mathcal C}.
\]
Direct-sum multiplicativity of orientations is therefore not the
displayed determinant identity.  It is the additional
Picard-groupoid datum consisting of square roots
\(o_{\mathcal C}^{\otimes2}\simeq\mathcal L^{\det}_{\mathcal C}\),
\(o_{\mathcal D}^{\otimes2}\simeq\mathcal L^{\det}_{\mathcal D}\),
\(o_{\mathcal C\oplus\mathcal D}^{\otimes2}\simeq
\mathcal L^{\det}_{\mathcal C\oplus\mathcal D}\), a chosen
Calabi--Yau hyperbolic trivialization
\[
\det K^{\mathrm{red}}_{\mathcal C,\mathcal D}\otimes
\det K^{\mathrm{red}}_{\mathcal D,\mathcal C}
\simeq q_{\mathcal C,\mathcal D}^{\otimes2},
\]
and an isomorphism
\[
m_{\oplus,\mathcal C,\mathcal D}:
o_{\mathcal C\oplus\mathcal D}
\xrightarrow{\sim}
o_{\mathcal C}\otimes o_{\mathcal D}\otimes
q_{\mathcal C,\mathcal D}
\]
which is symmetric in \((\mathcal C,\mathcal D)\), functorial under
pullback to the retained direct-sum and two-step flag charts, and has
zero Thom--Sebastiani pentagon defect.  Row \(299\) is proved for a
retained pair exactly when this datum is supplied for that pair and
recorded as a zero-defect row in
\(\mathrm{ts\_multiplicativity}\).

In the present finite packet the table
\(\mathrm{orientation\_lines}\) has no square-root rows and
\(\mathrm{ts\_multiplicativity}\) has no direct-sum isomorphism row.
Thus the current status of row \(299\) is the criterion above together
with an open obstruction, not a proof of direct-sum multiplicativity.
\end{proposition}

\begin{proof}
The first displayed isomorphism is the standard block decomposition of
derived endomorphisms of a direct sum, after restricting to a chart
where the reduced cone construction has been made compatible with
direct sum.  The Knudsen--Mumford determinant functor is multiplicative
for finite direct sums of perfect complexes, hence gives the
determinant-line formula.  The off-diagonal factors are not killed by
this formula; they must be paired by the Calabi--Yau duality and
trivialized as a square before one can compare square roots.

A reduced orientation is a square root in the Picard groupoid, and a
multiplicative reduced orientation is therefore the isomorphism
\(m_{\oplus,\mathcal C,\mathcal D}\) together with the symmetry,
pullback, and pentagon coherences stated above.  These are exactly the
entries carried by a direct-sum Thom--Sebastiani row.  The current
orientation ledger records the square-root row and the
Thom--Sebastiani multiplicativity row as missing open obligations, and
the corresponding tables are empty.  Hence the proposition proves the
criterion and identifies the missing data; it does not assert that the
missing data have been constructed.
\end{proof}

The packet
\[
\texttt{certificates/orientation/rhomred\_orientation\_direct\_sum\_multiplicativity}
\]
verifies
\(\texttt{RHOMRED\_ORIENTATION\_DIRECT\_SUM\_MULTIPLICATIVITY\_OBSTRUCTION\_VERIFIED}\):
it imports the reduced determinant, square-root obstruction, and
reduced-orientation ledger packets, records the direct-sum
Picard-groupoid criterion, checks that
\(\mathrm{orientation\_lines}\) and
\(\mathrm{ts\_multiplicativity}\) are empty, and keeps the square-root,
cross-term hyperbolic trivialization, direct-sum
Thom--Sebastiani isomorphism, pentagon, extension, quotient,
transition, vanishing-cycle, and protected-integration rows open.

\begin{proposition}[Extension orientation multiplicativity criterion]
\label{prop:orientation-extension-multiplicativity}
\textnormal{[criterion proved; extension Thom--Sebastiani row not
supplied]}
Let
\[
\mathfrak M^{\mathrm{red}}_{R,\widehat c}\times
\mathfrak M^{\mathrm{red}}_{R,\widehat c'}
\xleftarrow{\ p=(p_1,p_2)\ }
\overline{\mathfrak E}^{\mathrm{red}}_{R,\widehat c,\widehat c'}
\xrightarrow{\ q\ }
\mathfrak M^{\mathrm{red}}_{R,\widehat c+\widehat c'}
\]
be a retained compactified extension correspondence.  Row \(300\)
is an orientation row only if it supplies a Joyce--Upmeier
multiplicativity isomorphism
\[
\operatorname{TS}^{o}_{\mathfrak E}:
p_1^*o_{R,\widehat c}\otimes p_2^*o_{R,\widehat c'}
\xrightarrow{\sim}q^*o_{R,\widehat c+\widehat c'}
\]
whose square is the cosection-reduced determinant comparison on
\(\overline{\mathfrak E}^{\mathrm{red}}_{R,\widehat c,\widehat c'}\).
Equivalently, after tensoring with the BBDJS reduced
Thom--Sebastiani isomorphism for vanishing cycles, it gives the
coefficient-system isomorphism
\[
p^*\bigl((\Phi^{\mathrm{red}}_{R,\widehat c}\otimes
o_{R,\widehat c})\boxtimes
(\Phi^{\mathrm{red}}_{R,\widehat c'}\otimes
o_{R,\widehat c'})\bigr)
\xrightarrow{\sim}
q^*(\Phi^{\mathrm{red}}_{R,\widehat c+\widehat c'}\otimes
o_{R,\widehat c+\widehat c'}).
\]
It must be functorial on overlaps, compatible with quotient descent,
and its pullbacks to every retained two-step flag stack must have zero
Thom--Sebastiani pentagon defect.  Thus row \(300\) is proved exactly
by a zero-defect row in \(\mathrm{ts\_multiplicativity}\) with
\[
(\mathrm{extension\_stack\_id},\mathrm{left\_line\_id},
\mathrm{right\_line\_id},\mathrm{target\_line\_id},
\mathrm{ts\_isomorphism\_id}).
\]

The retained extension-closure rows record only that certain middle
HN types stay inside the finite window.  The present orientation table
\(\mathrm{ts\_multiplicativity}\) is empty, and the mixed
Thom--Sebastiani packet is conditional on supplied orientation
transport.  Consequently row \(300\) is presently a criterion with an
open extension-orientation obstruction.
\end{proposition}

\begin{proof}
The Hall product pulls the external coefficient system on
\(\mathfrak M^{\mathrm{red}}_{R,\widehat c}\times
\mathfrak M^{\mathrm{red}}_{R,\widehat c'}\) to the extension stack
and pushes it along \(q\).  The BBDJS Thom--Sebastiani theorem
identifies the reduced vanishing-cycle factors after the
reduced d-critical charts and the extension potential comparison have
been supplied.  That theorem does not choose square roots of the
orientation gerbe.  Joyce--Upmeier multiplicativity is the additional
Picard-groupoid isomorphism
\(\operatorname{TS}^{o}_{\mathfrak E}\) whose square is the determinant
comparison for the exact-triangle self-Ext complex.

Associativity of the oriented Hall product requires the same
isomorphism to satisfy the two-step flag pentagon; quotient descent
and HN transition compatibility are separate rows.  The current
certificate
\(\texttt{certificates/orientation/k3e\_reduced\_orientation}\) has no
rows in \(\mathrm{orientation\_lines}\) or
\(\mathrm{ts\_multiplicativity}\), and its obstruction ledger records
both \(\mathrm{ts\_multiplicativity}\) and
\(\mathrm{ts\_pentagon}\) as missing open obligations.  Retained
extension closure and conditional mixed Thom--Sebastiani transport
therefore do not prove row \(300\).
\end{proof}

The packet
\[
\texttt{certificates/orientation/rhomred\_orientation\_extension\_multiplicativity}
\]
verifies
\(\texttt{RHOMRED\_ORIENTATION\_EXTENSION\_MULTIPLICATIVITY\_OBSTRUCTION\_VERIFIED}\):
it imports retained extension closure, the reduced-orientation ledger,
the square-root obstruction packet, the direct-sum multiplicativity
criterion, and the mixed Thom--Sebastiani packet; it checks that the
extension-closure rows do not construct extension stacks or
orientations, that the mixed Thom--Sebastiani packet is conditional on
orientation transport, and that
\(\mathrm{orientation\_lines}\) and
\(\mathrm{ts\_multiplicativity}\) are empty.

\begin{proposition}[Thom--Sebastiani compatibility for orientation
lines]
\label{prop:orientation-thom-sebastiani-compatibility}
\textnormal{[compatibility criterion proved; orientation-line
Thom--Sebastiani row not supplied]}
Let \(\operatorname{TS}^{o}_{\mathfrak E}\) be a candidate
orientation-line isomorphism on a retained extension correspondence as
in Proposition~\ref{prop:orientation-extension-multiplicativity}.
It is compatible with Thom--Sebastiani if and only if the following
finite row is supplied.

\begin{enumerate}
\item The square-root diagram commutes, equivalently
\[
q^*\sigma_{\widehat c+\widehat c'}\circ
(\operatorname{TS}^{o}_{\mathfrak E})^{\otimes2}
=
\det(\operatorname{TS}^{\mathrm{red}}_{\mathfrak E})\circ
(\sigma_{\widehat c}\otimes\sigma_{\widehat c'}),
\]
as morphisms from
\(\bigl(p_1^*o_{R,\widehat c}\otimes
p_2^*o_{R,\widehat c'}\bigr)^{\otimes2}\) to
\(q^*\mathcal L^{\det}_{R,\widehat c+\widehat c'}\).
\item The oriented coefficient-system map obtained by tensoring
\(\operatorname{TS}^{o}_{\mathfrak E}\) with the BBDJS reduced
Thom--Sebastiani isomorphism is the row used by the Hall product:
\[
\operatorname{TS}^{\Phi\otimes o}_{\mathfrak E}
=\operatorname{TS}^{\Phi,\mathrm{red}}_{\mathfrak E}\otimes
\operatorname{TS}^{o}_{\mathfrak E}.
\]
\item The chart-overlap cocycle, quotient descent cocycle, and
two-step flag pentagon defect of this oriented
Thom--Sebastiani row vanish.
\end{enumerate}
Thus row \(301\) is not a new scalar identity and not merely the BBDJS
vanishing-cycle Thom--Sebastiani theorem.  It is the zero-defect
compatibility row saying that the Joyce--Upmeier orientation
isomorphism is the square root of the determinant Thom--Sebastiani
comparison and that its tensor with the vanishing-cycle
Thom--Sebastiani map is the coefficient-system transport.

The present packets do not supply such a row.  The hybrid
Thom--Sebastiani packets are conditional on supplied orientation
transport, the four-input pentagon packet treats the orientation
pentagon as a functorial input, and
\(\mathrm{orientation\_lines}\) and
\(\mathrm{ts\_multiplicativity}\) are empty.
\end{proposition}

\begin{proof}
An orientation line is a square root of the determinant line in the
Picard groupoid.  Therefore compatibility of an orientation
Thom--Sebastiani isomorphism is exactly compatibility after squaring:
the square of \(\operatorname{TS}^{o}_{\mathfrak E}\) must be the
determinant-line comparison induced by the reduced exact-triangle
self-Ext complex.  This proves the first condition.

BBDJS Thom--Sebastiani identifies the vanishing-cycle factors after
the reduced d-critical chart and potential comparison have been
supplied.  It does not choose the Joyce--Upmeier square root.  The
oriented Hall coefficient system is
\(\Phi^{\mathrm{red}}\otimes o\), so its Thom--Sebastiani transport is
the tensor product of the vanishing-cycle transport and the
orientation-line transport.  This gives the second condition.

The same row must be invariant under chart changes, descend through
the quotient, and satisfy the two-step flag pentagon; otherwise the
oriented Hall product depends on choices or fails associativity.  The
current orientation ledger records the relevant orientation-line,
Thom--Sebastiani, and pentagon rows as missing open obligations.  The
mixed and four-input Thom--Sebastiani packets use orientation
compatibility as a hypothesis.  Hence the proposition proves the
criterion and records the missing data, not an unconditional
compatibility theorem.
\end{proof}

The packet
\[
\texttt{certificates/orientation/rhomred\_orientation\_thom\_sebastiani\_compatibility}
\]
verifies
\(\texttt{RHOMRED\_ORIENTATION\_THOM\_SEBASTIANI\_COMPATIBILITY\_OBSTRUCTION\_VERIFIED}\):
it imports the row \(300\) extension-multiplicativity packet, the
reduced-orientation ledger, the mixed Thom--Sebastiani packet, and the
four-input pentagon packet; it records that the BBDJS
Thom--Sebastiani rows and conditional pentagon rows do not construct
the missing orientation-line compatibility row.

\begin{proposition}[Orientation descent through the \(E\)-quotient]
\label{prop:orientation-e-quotient-descent}
\textnormal{[descent criterion proved; quotient-orientation rows not
supplied]}
Let \(o_{\mathcal C}\) be a reduced orientation line on a retained
stratum of the \(K3\times E\) reduced moduli stack.  The line descends
through the quotient by \(E\) precisely when the following finite
Picard-groupoid datum is supplied.
\begin{enumerate}
\item a square-root row for \(o_{\mathcal C}^{\otimes2}\simeq
\det K^{\mathrm{red}}_{\mathcal C}\);
\item a null-trivialization of the ordinary reduced gerbe
\(\alpha^{\mathrm{red}}_{\mathcal C}\);
\item a Cech--Borel representative of the connected free \(E\)-class
\[
\alpha^{E,\mathrm{free}}_{\mathcal C}
=a_1(\mathcal C)u_1+a_2(\mathcal C)u_2\in H^2(BE;\mathbb F_2)
\]
with \(a_1=a_2=0\) and a chosen null-trivialization;
\item for every retained finite-stabilizer stratum \(S\), an
edge-reduced equivariant class
\[
\widetilde\beta^H_{\mathcal C,S}\rightsquigarrow
\beta^H_{\mathcal C,S}\in H^2(BH;\mathbb F_2),
\]
a zero row for \(\beta^H_{\mathcal C,S}\), and compatible
null-trivialisations under retained subgroup restrictions;
\item a computed residual linearization character
\[
\lambda^H_{\mathcal C,S}\in H^1(BH;\mathbb F_2)
\]
with \(\lambda^H_{\mathcal C,S}=0\) for every retained stabilizer;
\item compatibility of these choices on overlaps and on the quotient
presentation of the retained extension and flag charts.
\end{enumerate}
Equivalently, row \(302\) is proved by a populated
\(\mathrm{quotient\_borel}\) table and a populated
\(\mathrm{finite\_stabilizers}\) table whose obstruction ranks and
degree-one characters vanish.  The descent object is the descended
line \(o_{\mathcal C}/E\) on the quotient stack together with the
chosen quotient cocycle
\[
\mathfrak q_{\mathcal C}=
(\alpha^{\mathrm{red}}_{\mathcal C},
\alpha^{E,\mathrm{free}}_{\mathcal C},
\{(\beta^H_{\mathcal C,S},\lambda^H_{\mathcal C,S})\}_{S,H})
\simeq0.
\]

The present packets do not supply this datum.  They record criteria
for the reduced gerbe, connected free \(E\)-class, finite-stabilizer
classes, and residual linearization characters, but
\(\mathrm{orientation\_lines}\), \(\mathrm{quotient\_borel}\), and
\(\mathrm{finite\_stabilizers}\) are empty.  Thus row \(302\) is
currently a descent criterion with an open quotient-orientation
obstruction, not a descended orientation.
\end{proposition}

\begin{proof}
For a line on an \(E\)-stack to descend to the quotient stack, its
equivariant gerbe and stabilizer linearizations must be trivialized.
In the reduced orientation problem the ordinary component is
\(\alpha^{\mathrm{red}}_{\mathcal C}\).  The connected part of the
\(E\)-action contributes the \(H^2(BE;\mathbb F_2)\) edge class
\(\alpha^{E,\mathrm{free}}_{\mathcal C}\).  Finite stabilizers
contribute the Borel classes
\(\widetilde\beta^H_{\mathcal C,S}\), whose classifying-space edges
\(\beta^H_{\mathcal C,S}\) are legitimate targets only after the
mixed Borel and differential terms have been killed.  After the
degree-two gerbes vanish, a residual degree-one stabilizer character
\(\lambda^H_{\mathcal C,S}\) remains; nonzero \(\lambda^H\) changes
the descended orientation by a sign character and therefore prevents
descent.

The six listed rows are exactly the data required to trivialize these
obstructions and make the equivariant orientation line descend.  The
compatibility clauses are needed because descent is a statement in the
Picard groupoid on the quotient presentation, not a collection of
unrelated pointwise vanishings.  The current reduced-orientation
fixture records all quotient-Borel and finite-stabilizer rows as
missing open obligations, and the specialized packets for
\(\alpha^{\mathrm{red}}\), \(\alpha^{E,\mathrm{free}}\),
\(\beta^H\), and \(\lambda^H\) are obstruction ledgers.  Hence the
criterion is proved, while the descended orientation is not
constructed.
\end{proof}

The packet
\[
\texttt{certificates/orientation/rhomred\_orientation\_e\_quotient\_descent}
\]
verifies
\(\texttt{RHOMRED\_ORIENTATION\_E\_QUOTIENT\_DESCENT\_OBSTRUCTION\_VERIFIED}\):
it imports the square-root, reduced-gerbe nullhomotopy, connected
free \(E\)-nulltrivialization, finite-stabilizer
nulltrivialization, finite-stabilizer linearization-vanishing, and
reduced-orientation ledger packets; it checks that the quotient
tables are empty and records the exact missing descent rows.

\begin{proposition}[Orientation descent and the Hall product]
\label{prop:orientation-descent-hall-product}
\textnormal{[product-descent criterion proved; Hall-product orientation
rows not supplied]}
Fix a finite height \(R\) and a retained compactified extension
correspondence
\[
\mathfrak M_{R,\widehat c}\times\mathfrak M_{R,\widehat c'}
\xleftarrow{\ p\ }
\overline{\mathfrak E}_{R,\widehat c,\widehat c'}
\xrightarrow{\ q\ }
\mathfrak M_{R,\widehat c+\widehat c'} .
\]
Put \(K_{R,\widehat c}:=\Phi^{\mathrm{red}}_{R,\widehat c}\otimes
o_{R,\widehat c}\).  The statement that quotient-orientation descent
commutes with the Hall product is the commutative square
\[
\begin{tikzcd}
Q_{E,R}\,q_!\operatorname{TS}^{\Phi\otimes o}_R p^*
\ar[r,"\theta^Q_{\mu,R}"]
\ar[d,"Q_{E,R}m_R"']&
\overline q_!\,
\overline{\operatorname{TS}}^{\Phi\otimes o}_R
\overline p^*(Q_{E,R}\boxtimes Q_{E,R})
\ar[d,"\overline m_R"]\\
Q_{E,R}K_{R,\widehat c+\widehat c'}
\ar[r,"\theta^Q_{K,R}"']&
\overline K_{R,\widehat c+\widehat c'}
\end{tikzcd}
\]
in the Borel--Moore correspondence category on the quotient stack,
where \(\theta^Q_{\mu,R}\) is the quotient-after-correspondence
operation comparison and
\(\theta^Q_{K,R}\colon Q_{E,R}K_{R,\widehat c+\widehat c'}
\simeq\overline K_{R,\widehat c+\widehat c'}\) is the descended
coefficient-system identification.

This square is proved by supplying the following finite datum.
\begin{enumerate}
\item quotient-orientation rows for the two input strata, the
extension stack, and the output stratum, with the quotient cocycles
\[
p^*(\mathfrak q_{R,\widehat c}\boxtimes
\mathfrak q_{R,\widehat c'})
\quad\text{and}\quad
q^*\mathfrak q_{R,\widehat c+\widehat c'}
\]
identified by a zero Picard-groupoid defect row;
\item the oriented Thom--Sebastiani row
\[
\operatorname{TS}^{\Phi\otimes o}_R
=\operatorname{TS}^{\Phi,\mathrm{red}}_R\otimes
\operatorname{TS}^{o}_R
\]
whose square is the reduced determinant Thom--Sebastiani comparison
and whose quotient-descent defect is zero;
\item external-product descent
\[
Q_{E,R}(K_{R,\widehat c}\boxtimes K_{R,\widehat c'})
\simeq
\overline K_{R,\widehat c}\boxtimes\overline K_{R,\widehat c'}
\]
on the product quotient;
\item compact-support pushforward descent for \(q_!\), with proper
base-change and projection-formula defect ranks zero;
\item the quotient-after-correspondence comparison
\(\theta^Q_{\mu,R}\) for the supplied binary Hall product row;
\item compatibility of the preceding rows with the two-step flag
pentagon, so that product descent is associative after quotient.
\end{enumerate}
The current packets do not supply this datum.  Row \(302\) supplies no
quotient-orientation rows, row \(301\) supplies no oriented
Thom--Sebastiani row, and the compact Hall source packet supplies no
geometric Hall product matrix.  The
quotient-after-correspondence packet constructs \(\theta^Q_{\mu,R}\)
only for operation rows with descent interfaces already supplied.  Thus
row \(303\) is presently a product-descent criterion, not a proof that
orientation descent commutes with the Hall product.
\end{proposition}

\begin{proof}
The Hall product is the pull--transport--push composite
\[
m_R=q_!\operatorname{TS}^{\Phi\otimes o}_R p^* .
\]
A quotient comparison for this composite is an isomorphism of
composites, not a scalar identity.  Therefore it decomposes into three
functorial requirements: descent of \(p^*\) and of the external product
coefficient system, descent of the oriented Thom--Sebastiani
transport, and descent of \(q_!\) through compact supports.  The first
requirement is exactly the Picard-groupoid equality of the pulled-back
quotient cocycles on the extension correspondence.  The second is the
orientation-line row of Proposition~\ref{prop:orientation-thom-sebastiani-compatibility}
tensored with the reduced BBDJS Thom--Sebastiani transport.  The third
is the compact-support pushforward comparison used by
Proposition~\ref{prop:quotient-after-correspondence-theta-mu-comparison}.

If these rows are supplied, the naturality of pullback, the descended
Thom--Sebastiani isomorphism, and compact-support pushforward descent
identify
\[
Q_{E,R}q_!\operatorname{TS}^{\Phi\otimes o}_R p^*
\quad\text{with}\quad
\overline q_!\overline{\operatorname{TS}}^{\Phi\otimes o}_R
\overline p^*(Q_{E,R}\boxtimes Q_{E,R}),
\]
which is the displayed square.  Associativity after quotient is the
same statement on the retained two-step flag stack.  The present
tables fail before this proof can be applied: the quotient-orientation
tables are empty, the oriented Thom--Sebastiani row is absent, and the
source Hall product rows are not populated.  Hence only the criterion
is proved.
\end{proof}

The packet
\[
\texttt{certificates/orientation/rhomred\_orientation\_hall\_product\_descent}
\]
verifies
\(\texttt{RHOMRED\_ORIENTATION\_HALL\_PRODUCT\_DESCENT\_OBSTRUCTION\_VERIFIED}\):
it imports the row \(302\) quotient-orientation descent packet, the
row \(301\) oriented Thom--Sebastiani packet, the
quotient-after-correspondence \(\theta^Q_{\mu}\) packet, and the
transition Hall-product obstruction packet; it records that the current
tree has comparison maps for already supplied operation rows but no
quotient-orientation Hall-product descent row.

\begin{proposition}[Orientation descent and the Hall coproduct]
\label{prop:orientation-descent-hall-coproduct}
\textnormal{[coproduct-descent criterion proved; Hall-coproduct
orientation rows not supplied]}
Fix a finite height \(R\) and the same retained compactified
correspondence, read in the splitting direction
\[
\mathfrak M_{R,\widehat c+\widehat c'}
\xleftarrow{\ q\ }
\overline{\mathfrak E}_{R,\widehat c,\widehat c'}
\xrightarrow{\ p\ }
\mathfrak M_{R,\widehat c}\times\mathfrak M_{R,\widehat c'} .
\]
With \(K_{R,\widehat c}=\Phi^{\mathrm{red}}_{R,\widehat c}\otimes
o_{R,\widehat c}\), quotient-orientation descent commutes with the
Hall coproduct precisely when the pull--inverse-transport--push
comparison
\[
\begin{tikzcd}
Q_{E,R}\,p_!(\operatorname{TS}^{\Phi\otimes o}_R)^{-1}q^*
\ar[r,"\theta^Q_{\Delta,R}"]
\ar[d,"Q_{E,R}\Delta_R"']&
\overline p_!\,
(\overline{\operatorname{TS}}^{\Phi\otimes o}_R)^{-1}
\overline q^*Q_{E,R}
\ar[d,"\overline\Delta_R"]\\
Q_{E,R}(K_{R,\widehat c}\boxtimes K_{R,\widehat c'})
\ar[r,"\theta^Q_{K\boxtimes K,R}"']&
\overline K_{R,\widehat c}\boxtimes
\overline K_{R,\widehat c'}
\end{tikzcd}
\]
commutes on the quotient Borel--Moore correspondence category.

This square is proved by supplying the following finite datum.
\begin{enumerate}
\item quotient-orientation rows for the source stratum, the splitting
stack, and the two output strata, with the cocycle comparison
\[
q^*\mathfrak q_{R,\widehat c+\widehat c'}
=
p^*(\mathfrak q_{R,\widehat c}\boxtimes
\mathfrak q_{R,\widehat c'})
\]
in the Picard groupoid;
\item the inverse oriented Thom--Sebastiani row
\[
(\operatorname{TS}^{\Phi\otimes o}_R)^{-1}
=
(\operatorname{TS}^{\Phi,\mathrm{red}}_R)^{-1}\otimes
(\operatorname{TS}^{o}_R)^{-1}
\]
with determinant-square and quotient-descent defects zero;
\item external-product descent for the two-output coefficient system
\(K_{R,\widehat c}\boxtimes K_{R,\widehat c'}\);
\item compact-support pushforward descent for \(p_!\), with proper
base-change and projection-formula defect ranks zero;
\item a supplied collision-coproduct row \(\Delta_R\), finite
coproduct matrix \(D_R\), counit row, and zero coassociativity and
counit defects;
\item the quotient-after-correspondence comparison
\(\theta^Q_{\Delta,R}\) and compatibility with the bialgebra square
relating \(\Delta_R\) and \(m_R\).
\end{enumerate}
The current packets do not supply this datum.  Row \(302\) supplies no
quotient-orientation rows, row \(301\) supplies no oriented
Thom--Sebastiani row, the compact Hall source packet supplies no
\(D\)-entries or counit rows, and the
quotient-after-correspondence coproduct packet applies only to already
supplied quotient-presented coproduct and bialgebra rows.  Thus row
\(304\) is presently a coproduct-descent criterion, not a proof that
orientation descent commutes with the Hall coproduct.
\end{proposition}

\begin{proof}
The collision coproduct is the pull--inverse-transport--push composite
\[
\Delta_R=p_!(\operatorname{TS}^{\Phi\otimes o}_R)^{-1}q^* .
\]
For this composite to descend, the quotient cocycle on the source
object must agree, after pullback to the splitting correspondence, with
the external product of the two output quotient cocycles.  This is the
Picard-groupoid equality displayed above.  The coefficient transport is
then the inverse of the oriented Thom--Sebastiani isomorphism, not the
product orientation row alone.  Finally the exceptional pushforward is
along \(p\), so the six-functor comparison is the \(p_!\)-descent row
for the splitting correspondence.

Given those rows, functoriality of pullback, inverse
Thom--Sebastiani descent, and compact-support descent identify
\[
Q_{E,R}p_!(\operatorname{TS}^{\Phi\otimes o}_R)^{-1}q^*
\quad\text{with}\quad
\overline p_!
(\overline{\operatorname{TS}}^{\Phi\otimes o}_R)^{-1}
\overline q^*Q_{E,R},
\]
which is the displayed square.  Coassociativity and counitality after
quotient are the same statement on the retained two-step flag and
zero-object rows.  The present tables fail before the criterion can be
applied: the quotient-orientation tables are empty, the oriented
Thom--Sebastiani row is absent, and the source coproduct and counit
rows are not populated.  Hence only the criterion is proved.
\end{proof}

The packet
\[
\texttt{certificates/orientation/rhomred\_orientation\_hall\_coproduct\_descent}
\]
verifies
\(\texttt{RHOMRED\_ORIENTATION\_HALL\_COPRODUCT\_DESCENT\_OBSTRUCTION\_VERIFIED}\):
it imports the row \(302\) quotient-orientation descent packet, the
row \(301\) oriented Thom--Sebastiani packet, the row \(303\)
product-descent packet, the quotient-after-correspondence
\(\theta^Q_{\Delta}\) packet, and the transition Hall-coproduct
obstruction packet; it records that the current tree has conditional
coproduct comparison rows but no quotient-orientation Hall-coproduct
descent row.

\begin{proposition}[Orientation descent and primitive projection]
\label{prop:orientation-descent-primitive-projection}
\textnormal{[primitive-projection descent criterion proved; primitive
orientation rows not supplied]}
Fix \(R\).  Let
\[
K_R=\Phi^{\mathrm{red}}_R\otimes o_R
\]
be the oriented vanishing-cycle coefficient system on the retained
finite Hall object \(\mathcal H_R(K_R)\).  Write
\[
\widetilde\Delta_R=\Delta_R-\eta_R\epsilon_R\boxtimes\mathrm{id}
-\mathrm{id}\boxtimes\eta_R\epsilon_R,\qquad
P_R(K_R)=\ker\widetilde\Delta_R\cap\ker\epsilon_R
\]
and let
\(\pi_R^{\Prim}\colon\mathcal H_R(K_R)\to P_R(K_R)\)
be the supplied primitive idempotent.  Denote the quotient objects by
\(\overline{\mathcal H}_R(\overline K_R)\),
\(\overline P_R(\overline K_R)\), and
\(\overline\pi_R^{\Prim}\).  Orientation descent commutes with
primitive projection precisely when the idempotent square
\[
\begin{tikzcd}
Q_{E,R}\mathcal H_R(K_R)
\ar[r,"\theta^Q_{H,R}"]
\ar[d,"Q_{E,R}\pi_R^{\Prim}"']&
\overline{\mathcal H}_R(\overline K_R)
\ar[d,"\overline\pi_R^{\Prim}"]\\
Q_{E,R}P_R(K_R)
\ar[r,"\theta^{Q,\Prim}_{o,R}"']&
\overline P_R(\overline K_R)
\end{tikzcd}
\]
commutes in the idempotent-complete quotient Borel--Moore
correspondence category.

This square is proved by supplying the following finite datum.
\begin{enumerate}
\item quotient-orientation rows for the ambient Hall object, the
zero-object stratum, and every splitting correspondence entering
\(\widetilde\Delta_R\);
\item the Hall-coproduct orientation-descent datum of
Proposition~\ref{prop:orientation-descent-hall-coproduct}, together
with unit and counit orientation descent;
\item finite source and quotient primitive kernel rows
\[
P_R(K_R)=\ker(\widetilde\Delta_R,\epsilon_R),\qquad
\overline P_R(\overline K_R)
=\ker(\widetilde{\overline\Delta}_R,\overline\epsilon_R),
\]
with inclusion and projection matrices for the idempotents
\(\pi_R^{\Prim}\) and \(\overline\pi_R^{\Prim}\);
\item the quotient-after-correspondence primitive comparison
\[
\theta^{Q,\Prim}_{R}\colon Q_{E,R}P_R(K_R)
\xrightarrow{\sim}\overline P_R(\overline K_R)
\]
with zero reduced-coproduct and counit kernel-intertwining defects;
\item the projection-commutator identity
\[
\overline\pi_R^{\Prim}\theta^Q_{H,R}
=\theta^{Q,\Prim}_{o,R}Q_{E,R}\pi_R^{\Prim};
\]
\item for the pro-object, transition-compatible primitive matrices and
zero composition defects for \(R'\to R\).
\end{enumerate}
The present packets do not supply this datum.  The quotient
primitive-projection packet proves compatibility of
\(\theta^Q_{\mu,R}\) only for already supplied finite primitive kernel
and projection rows.  The compact Hall source packet supplies no
\(D\)-entries, counit rows, primitive kernels, or primitive projection
matrices; the transition primitive-subspace packet records empty
transition tables.  Hence row \(305\) is a primitive-projection descent
criterion, not a proof that orientation descent commutes with primitive
projection.
\end{proposition}

\begin{proof}
The primitive subobject is the kernel of the map
\[
(\widetilde\Delta_R,\epsilon_R)\colon
\mathcal H_R(K_R)\longrightarrow
\mathcal H_R(K_R)\boxtimes\mathcal H_R(K_R)\oplus\mathbf 1 .
\]
If the coproduct, unit, counit, and orientation coefficient systems
descend through \(Q_{E,R}\), then applying \(Q_{E,R}\) carries this
kernel diagram to the quotient kernel diagram.  The supplied kernel
rows make this statement finite: the reduced-coproduct and counit
matrices intertwine with the quotient comparison, so
\(Q_{E,R}P_R(K_R)\) identifies with
\(\overline P_R(\overline K_R)\).

The primitive projection is then an idempotent splitting of this kernel
inside the idempotent-complete correspondence category.  Two morphisms
from \(Q_{E,R}\mathcal H_R(K_R)\) to
\(\overline P_R(\overline K_R)\) are equal once their compositions with
the inclusion into \(\overline{\mathcal H}_R(\overline K_R)\) agree.
That equality is exactly the kernel-intertwining identity for
\(\widetilde\Delta_R\) and \(\epsilon_R\), together with the
projection-commutator row.  Transition compatibility is a separate
strictness condition needed only after passing from one finite height to
the pro-object.

The available data stop before this criterion can be applied.  Row
\(304\) does not provide an actual oriented Hall-coproduct descent row,
the compact source has no coproduct or counit matrices, and the
primitive transition packet has no primitive kernel or projection
matrices.  Thus only the criterion is proved.
\end{proof}

The packet
\[
\texttt{certificates/orientation/rhomred\_orientation\_primitive\_projection\_descent}
\]
verifies
\(\texttt{RHOMRED\_ORIENTATION\_PRIMITIVE\_PROJECTION\_DESCENT\_OBSTRUCTION\_VERIFIED}\):
it imports row \(302\), row \(304\), the
quotient-after-correspondence primitive-projection packet, the
transition primitive-subspace obstruction packet, and the compact Hall
source ledger; it records that the present tree has a conditional
\(\theta^{Q,\Prim}\) comparison but no compact Hall primitive matrices
and no orientation primitive-projection descent row.

\subsubsection*{(O1)$^+$ Weyl-equivariant transport along
$W^{(2)}(\Lambda_{II}^{2,1})$}

The previous criterion isolates the data needed for descent through
the reduced $E$-quotient.  Once such quotient-orientation data are
supplied, the Igusa structure still requires compatibility with the
Weyl group of the BKM target.  For each of the three type-II simple
reflections $s_{\delta_i}$ ($i=1,2,3$), the $K3\times E$ side carries
a wall transport
\[
\tau_i:o_{\mathcal C}\longrightarrow s_{\delta_i}^{*}o_{\mathcal C}.
\]
The word \emph{coherent} is a hypothesis.  At finite height, the
first object is the finite partial action groupoid $\mathcal G_R$
whose objects are retained oriented quotient strata and whose arrows
are the retained type-II wall transports; a choice of
determinant-line lifts gives a $2$-cocycle
$c_{o,R}\in Z^2(\mathcal G_R;\underline{\mathbb F}_2)$.  Only on a
complete $W^{(2)}(\Lambda_{II}^{2,1})$-orbit with constant
coefficients does this reduce to the group cohomology class
\[
[c_o]\in H^2\bigl(W^{(2)}(\Lambda_{II}^{2,1});\mathbb F_2\bigr).
\]

\begin{definition}[(O1)$^+$ Weyl-equivariant determinant-line
transport]
\label{def:O1-plus-weyl-transport}
The (O1)$^+$ datum is the chosen set of lifts $\{\tau_i\}$
satisfying $[c_o]=0$ with chosen $1$-cochain killing $c_o$, the
normalization $\tau_i^2=1$, and the transport of the
quotient-orientation cocycle
\[
\mathfrak q_{\mathcal C}=\bigl(\alpha^{\mathrm{red}}_{\mathcal C},
\alpha^{E,\mathrm{free}}_{\mathcal C},
\{(\beta^H_{\mathcal C},\lambda^H_{\mathcal C})\}_{H\in
\mathcal H(\mathcal C)}\bigr)
\]
along $\tau_i$ between charge strata over Weyl-related Gram degrees.
The torsor defect of this transport is
\[
\omega_{i,\mathcal C}\in
H^1(\mathfrak M^{\mathrm{red}}_{s_{\delta_i}\mathcal C};\mathbb F_2)
\oplus
\bigoplus_{H\in\mathcal H(\mathcal C)}H^1(BH;\mathbb F_2),
\]
required to vanish.  The Coxeter presentation of
$W^{(2)}(\Lambda_{II}^{2,1})$ then assembles these lifts into an honest
action on the oriented Pfaffian line over the wall-complement of the
reduced compact Hall sector.
\end{definition}

(O1)$^+$ is a determinant-line monodromy datum.  No scalar formula
enters; the transports must move the entire null-trivialization
datum, not only the underlying line.  The (O1)$^+$ data comprise
$[c_o]=0$, the chosen cochain, the unitarity $\tau_i^2=1$, and the
torsor vanishing $\omega_{i,\mathcal C}=0$ on every retained orbit.
In finite coordinates this is recorded by
\[
\mathrm{weyl\_lifts},\quad
\mathrm{coxeter\_cocycles},\quad
\mathrm{transitions},\quad
\mathrm{scalar\_firewall}.
\]
The Weyl-lift rows carry \(\tau_i^2=1\), torsor-defect zero, and
quotient-cocycle transport zero.  The Coxeter rows carry the vanishing
projective cocycle, the cochain killing it, and the type-II Coxeter
defect.  Transition rows carry strictness and \(R^1\varprojlim=0\) for
orientation lines, quotient data, and Weyl lifts.  The scalar firewall
excludes scalar trace, squared determinant, OP scalar branch, and
Maass-character values as orientation data.

\begin{proposition}[Weyl wall transport maps]
\label{prop:weyl-wall-transport-tau-construction}
\textnormal{[wall-transport construction criterion proved; \(\tau_i\)
rows not supplied]}
Fix \(R\), a retained quotient stratum \(S\), and a type-II simple
reflection \(s_{\delta_i}\), \(i\in\{1,2,3\}\).  A Weyl wall transport
\[
\tau_{i,R,S}\colon o_{R,S}\longrightarrow
s_{\delta_i}^{*}o_{R,s_{\delta_i}S}
\]
is constructed by a finite \(\texttt{weyl\_lifts}\) row with the
following payload.
\begin{enumerate}
\item a verified type-II generator row for \(s_{\delta_i}\) and a
retained wall correspondence carrying \(S\) to
\(s_{\delta_i}S\);
\item source and target Joyce--Upmeier orientation lines
\[
o_{R,S}^{\otimes2}\simeq\det K^{\mathrm{red}}_{R,S},
\qquad
o_{R,s_{\delta_i}S}^{\otimes2}
\simeq\det K^{\mathrm{red}}_{R,s_{\delta_i}S};
\]
\item a Picard-groupoid lift matrix for \(\tau_{i,R,S}\), whose square
with the determinant transport commutes:
\[
(\tau_{i,R,S})^{\otimes2}\colon
o_{R,S}^{\otimes2}\longrightarrow
s_{\delta_i}^{*}o_{R,s_{\delta_i}S}^{\otimes2}
\]
agrees with the determinant transport
\(\det K^{\mathrm{red}}_{R,S}\to
s_{\delta_i}^{*}\det K^{\mathrm{red}}_{R,s_{\delta_i}S}\);
\item transport of the quotient-orientation cocycle
\(\mathfrak q_{R,S}\) to
\(s_{\delta_i}^{*}\mathfrak q_{R,s_{\delta_i}S}\), including the
finite-stabilizer Borel classes and the zero linearization
characters;
\item vanishing torsor and quotient-cocycle transport defects
\[
\omega_{i,R,S}=0,\qquad
\operatorname{defect}_{q,i,R,S}=0 .
\]
\end{enumerate}
The condition \(\tau_{i,R,S}^2=1\) is a separate involutivity row; it
is not part of this construction criterion.

The current data prove only the target and profile prerequisites: the
type-II chamber packet supplies the three simple reflections, and the
quotient-after-correspondence wall-profile packet proves that the three
profiles survive the \(E\)-quotient.  The reduced-orientation packet
has no orientation-line rows and no \(\texttt{weyl\_lifts}\) rows, and
the wall-atlas packet has no retained wall-object rows.  Thus row
\(306\) is presently a Weyl-wall transport criterion, not a
construction of the maps \(\tau_i\).
\end{proposition}

\begin{proof}
The map \(\tau_{i,R,S}\) is a Picard-groupoid isomorphism, not a scalar
sign.  Hence the first input is an actual source and target orientation
line over Weyl-related quotient strata.  The determinant square-root
condition forces the displayed square to commute, so the lift preserves
the Joyce--Upmeier square root of the reduced determinant complex.
The quotient-orientation datum contains the reduced gerbe, the free
\(E\)-Borel class, the finite-stabilizer Borel classes, and the
linearization characters; transporting only the underlying line would
lose this data.  The vanishing of \(\omega_{i,R,S}\) and of the
quotient-cocycle defect is exactly the assertion that the entire
quotient-orientation package moves across the wall.

The type-II generator and profile-survival rows do not construct this
Picard isomorphism.  They identify which reflections and quotient
profiles must be used.  The construction itself starts only when the
orientation-line rows, retained wall-object rows, and Weyl-lift rows
are populated.  Those tables are empty in the current packets, so the
criterion is the strongest proved statement.
\end{proof}

The packet
\[
\texttt{certificates/orientation/rhomred\_weyl\_wall\_transport\_tau}
\]
verifies
\(\texttt{RHOMRED\_WEYL\_WALL\_TRANSPORT\_TAU\_OBSTRUCTION\_VERIFIED}\):
it imports the blocked reduced-orientation ledger, the type-II chamber
generator packet, the quotient wall-profile survival packet, and the
blocked wall-atlas ledger; it records that the three target reflections
and quotient profiles are available, but no source orientation-line,
wall-object, or \(\tau_i\)-transport rows are supplied.

\begin{proposition}[Involutivity of the Weyl wall transport]
\label{prop:weyl-wall-transport-tau-square}
\textnormal{[involutivity criterion proved; \(\tau_i^2\) rows not
supplied]}
Fix \(R\), a retained quotient stratum \(S\), and a supplied Weyl wall
transport
\[
\tau_{i,R,S}\colon o_{R,S}\longrightarrow
s_{\delta_i}^{*}o_{R,s_{\delta_i}S}.
\]
The identity \(\tau_{i,R,S}^2=1\) is the commutativity of the
Picard-groupoid square
\[
\begin{tikzcd}
o_{R,S}\ar[r,"\tau_{i,R,S}"]\ar[dr,"\mathrm{id}"']&
s_{\delta_i}^{*}o_{R,s_{\delta_i}S}
\ar[d,"s_{\delta_i}^{*}\tau_{i,R,s_{\delta_i}S}"]\\
&s_{\delta_i}^{*}s_{\delta_i}^{*}o_{R,S}
\ar[d,"\kappa_{i,R,S}"]\\
&o_{R,S}
\end{tikzcd}
\]
where \(\kappa_{i,R,S}\) is the supplied identification of the
double-reflected orientation line with \(o_{R,S}\) over the equality
\(s_{\delta_i}^2=1\) on the retained wall correspondence.

It is proved by a finite \(\texttt{weyl\_lifts}\) involutivity row with
the following data:
\begin{enumerate}
\item a row constructing both lifts
\(\tau_{i,R,S}\) and
\(\tau_{i,R,s_{\delta_i}S}\) as in
Proposition~\ref{prop:weyl-wall-transport-tau-construction};
\item the target Coxeter row \(s_{\delta_i}^2=1\) for the type-II
reflection, and the corresponding double-wall source correspondence;
\item a determinant-square compatibility row showing that
\((s_{\delta_i}^{*}\tau_{i,R,s_{\delta_i}S})\circ\tau_{i,R,S}\)
squares to the identity transport of
\(\det K^{\mathrm{red}}_{R,S}\);
\item transport of the quotient-orientation cocycle around the
two-edge loop with zero reduced-gerbe, free-\(E\), finite-stabilizer,
and linearization defects;
\item a finite matrix row
\[
\operatorname{defect}_{\tau^2,i,R,S}=0
\]
recording that the composite equals \(\mathrm{id}_{o_{R,S}}\), not
only its square on the determinant line.
\end{enumerate}

The present data do not supply these rows.  Row \(306\) records only
the criterion for constructing \(\tau_i\).  The type-II chamber packet
records the target relation \(s_{\delta_i}^2=1\), and the Maass
character packet records
\(\nu_{\Delta_5}(s_{\delta_i})^2=1\); neither constructs an
orientation-line lift nor proves its square is the identity.  Thus row
\(307\) is presently an involutivity criterion, not a proof of
\(\tau_i^2=1\).
\end{proposition}

\begin{proof}
The reflection identity \(s_{\delta_i}^2=1\) is an identity in the
target type-II Coxeter group.  The lift \(\tau_{i,R,S}\) lives in a
Picard groupoid over the retained source stratum.  Therefore
involutivity is a statement about the two-edge source loop
\[
S\longrightarrow s_{\delta_i}S\longrightarrow S
\]
and about the induced automorphism of the orientation line over that
loop.

The determinant-square row removes the first possible defect: after
squaring, the composite lift acts trivially on the determinant complex.
This is not enough, because a square root can still differ from the
identity by the residual sign automorphism of the orientation line and
by the quotient-orientation torsor.  The quotient-cocycle transport
rows remove the Borel and finite-stabilizer defects.  The final
\(\tau^2\)-defect row asserts that the remaining line automorphism is
the identity.  These are finite matrix and Picard-groupoid checks, not
consequences of the scalar character value.

Since the current \(\texttt{weyl\_lifts}\) table is empty, none of the
composites in the displayed square is supplied.  The available target
and automorphic rows identify the relation to be lifted; they do not
lift it.  Hence only the criterion is proved.
\end{proof}

The packet
\[
\texttt{certificates/orientation/rhomred\_weyl\_tau\_square}
\]
verifies
\(\texttt{RHOMRED\_WEYL\_TAU\_SQUARE\_OBSTRUCTION\_VERIFIED}\):
it imports the row \(306\) transport criterion, the blocked
reduced-orientation ledger, the type-II chamber packet, and the Maass
character packet; it records that the target reflection relation and
the scalar square value are available, but no \(\tau_i\)-lift or
\(\tau_i^2\)-defect row is supplied.

\begin{definition}[Finite partial action groupoid]
\label{def:finite-partial-action-groupoid-GR}
\textnormal{[groupoid datum specified; finite groupoid rows not
supplied]}
At height \(R\), the finite partial action groupoid
\(\mathcal G_R\) is the following finite datum.
\begin{enumerate}
\item The objects are retained quotient-orientation strata \(S\) for
which the orientation line \(o_{R,S}\), quotient cocycle
\(\mathfrak q_{R,S}\), and retained support condition are supplied.
\item The generating arrows are the retained type-II wall transports
\[
g_{i,R,S}\colon S\longrightarrow s_{\delta_i}S
\]
for \(i\in\{1,2,3\}\), defined only when \(S\) and
\(s_{\delta_i}S\) both lie in the retained finite window and the row
\(\tau_{i,R,S}\) of
Proposition~\ref{prop:weyl-wall-transport-tau-construction} is
supplied.
\item The identity arrow \(\mathrm{id}_S\), inverse arrow
\[
g_{i,R,S}^{-1}=g_{i,R,s_{\delta_i}S},
\]
and partial compositions
\[
g_{i_k,R,S_{k-1}}\circ\cdots\circ g_{i_1,R,S_0}
\]
are included exactly when every intermediate stratum \(S_j\) is
retained and every wall-transport row in the word is supplied.
\item The structure tables record source, target, identity, inverse,
and associativity defects; all are required to have defect rank zero.
\end{enumerate}
The coefficient system for the orientation cocycle on
\(\mathcal G_R\) is not forced to be constant.  It is the local system
whose fibre at \(S\) is the automorphism group of the complete
quotient-orientation package over \(S\).  Only on a complete retained
orbit with identified fibres may one replace it by the constant
\(\underline{\mathbb F}_2\)-system.

The present data do not populate \(\mathcal G_R\).  The type-II
chamber packet gives the three target generators and the no-finite-braid
relation among distinct generators, and the quotient wall-profile
packet gives three surviving profiles.  The reduced-orientation packet
has no retained object rows, no Weyl-lift rows, and no
\(\texttt{coxeter\_cocycles}\) rows.  Thus row \(308\) specifies the
finite partial action groupoid datum, but the current tree contains no
finite \(\mathcal G_R\) object, arrow, inverse, or composition tables.
\end{definition}

The packet
\[
\texttt{certificates/orientation/rhomred\_weyl\_partial\_action\_groupoid}
\]
verifies
\(\texttt{RHOMRED\_WEYL\_PARTIAL\_ACTION\_GROUPOID\_OBSTRUCTION\_VERIFIED}\):
it imports the target type-II chamber packet, the row \(306\) and row
\(307\) Weyl-lift criteria, the quotient wall-profile survival packet,
and the blocked reduced-orientation ledger; it records that the target
generators are available but the finite groupoid object, arrow,
inverse, and composition rows are not supplied.

\begin{definition}[Orientation projective cocycle]
\label{def:orientation-projective-cocycle-coR}
\textnormal{[cocycle formula specified; finite \(c_{o,R}\) rows not
supplied]}
Assume the finite partial action groupoid \(\mathcal G_R\) of
Definition~\ref{def:finite-partial-action-groupoid-GR} has been
populated, and that every arrow \(g:S\to T\) in \(\mathcal G_R\) is
equipped with a determinant-line lift
\[
\tau_g:o_{R,S}\longrightarrow g^*o_{R,T}
\]
transporting the full quotient-orientation package.  For every
composable pair
\[
S\xrightarrow{g}T\xrightarrow{h}U
\]
with composite \(hg\), the orientation projective cocycle is the
unique element
\[
c_{o,R}(h,g;S)\in
\operatorname{Aut}_{\mathfrak Q^{\mathrm{or}}_{R,S}}(o_{R,S})
\simeq \mathbb F_2
\]
defined by
\[
g^*\tau_h\circ\tau_g
=(-1)^{c_{o,R}(h,g;S)}\,\tau_{hg}.
\]
Equivalently, \(c_{o,R}\) is the finite table measuring the defect
between composition of chosen Weyl determinant-line lifts and the
chosen lift of the partial product in \(\mathcal G_R\).

The row is valid only after the following data are supplied:
\begin{enumerate}
\item object, arrow, inverse, and composition rows for
\(\mathcal G_R\);
\item determinant-line lifts for \(g\), \(h\), and \(hg\), including
the quotient-cocycle transport rows;
\item an identified coefficient fibre
\(\operatorname{Aut}_{\mathfrak Q^{\mathrm{or}}_{R,S}}(o_{R,S})\) at
the source object;
\item a cocycle matrix row recording \(c_{o,R}(h,g;S)\) for every
composable pair;
\item a normalization row
\[
c_{o,R}(g,\mathrm{id}_S;S)=
c_{o,R}(\mathrm{id}_T,g;S)=0;
\]
\item a cocycle-defect row
\[
dc_{o,R}(k,h,g;S)=0
\]
for every composable triple in the partial groupoid.
\end{enumerate}
The present packets do not provide these rows.  The row \(308\)
groupoid packet is empty, the \(\texttt{weyl\_lifts}\) table is empty,
and the \(\texttt{coxeter\_cocycles}\) table is empty.  Thus row
\(309\) defines the orientation \(2\)-cocycle formula, but no finite
cocycle table is currently populated.
\end{definition}

The packet
\[
\texttt{certificates/orientation/rhomred\_weyl\_orientation\_cocycle}
\]
verifies
\(\texttt{RHOMRED\_WEYL\_ORIENTATION\_COCYCLE\_OBSTRUCTION\_VERIFIED}\):
it imports the row \(308\) partial-action-groupoid packet, the row
\(306\) and row \(307\) lift criteria, and the blocked
reduced-orientation ledger; it records that the formula for
\(c_{o,R}\) is fixed but that the finite composable-pair, lift, and
cocycle-defect rows are absent.

\begin{proposition}[Vanishing criterion for the orientation cocycle
class]
\label{prop:orientation-cocycle-class-vanishing-criterion}
\textnormal{[criterion proved; cohomology-vanishing rows not supplied]}
Let \(\mathcal G_R\), the coefficient system
\(\underline{\mathbb F}_{2,o}\), and the normalized cocycle
\(c_{o,R}\in Z^2(\mathcal G_R;\underline{\mathbb F}_{2,o})\) be
supplied as in Definition~\ref{def:orientation-projective-cocycle-coR}.
Choose ordered finite bases for the normalized groupoid cochain spaces
\[
C^1_R=C^1(\mathcal G_R;\underline{\mathbb F}_{2,o}),\qquad
C^2_R=C^2(\mathcal G_R;\underline{\mathbb F}_{2,o}),
\]
write the coboundary as a matrix
\[
d^1_R:C^1_R\longrightarrow C^2_R,
\]
and write the cocycle as a column vector \(v(c_{o,R})\in C^2_R\).
Then
\[
[c_{o,R}]=0\in
H^2(\mathcal G_R;\underline{\mathbb F}_{2,o})
\]
if and only if
\[
\operatorname{rank}_{\mathbb F_2}(d^1_R)
=
\operatorname{rank}_{\mathbb F_2}\bigl[d^1_R\mid v(c_{o,R})\bigr].
\]
Equivalently, the augmented linear system
\[
d^1_R b=v(c_{o,R})
\]
is soluble over \(\mathbb F_2\).  This is the row \(310\) criterion.
A supplied rank equality proves the vanishing of the cohomology class;
row \(311\) is the subsequent choice of a particular solution \(b\).

The criterion is not triggered by the target Coxeter graph, by the
Maass character, by the relation \(s_{\delta_i}^2=1\), or by an
abstract calculation of \(H^2(W^{(2)};\mathbb F_2)\) before the
retained groupoid, coefficient system, finite orbit, and cocycle
vector have been supplied.  The present finite tables contain no
values of \(c_{o,R}\), no normalized cochain-complex bases, no
coboundary matrix, and no rank certificate.  Hence row \(310\) remains
an open cohomological obstruction.
\end{proposition}

\begin{proof}
For any finite groupoid with coefficients in a local system of
\(\mathbb F_2\)-lines, the normalized cochain complex is a finite
cochain complex of \(\mathbb F_2\)-vector spaces.  A closed
\(2\)-cochain represents zero in \(H^2\) exactly when it lies in
\(\operatorname{im}d^1_R\).  After ordered bases are fixed, membership
in the image of a matrix is equivalent to equality of the rank of the
matrix and the rank of the matrix augmented by the target vector.
This proves the displayed criterion.  Inspection of
\(\texttt{certificates/orientation/rhomred\_weyl\_orientation\_cocycle}\)
and
\(\texttt{certificates/orientation/k3e\_reduced\_orientation/coxeter\_cocycles.csv}\)
shows that the required cocycle, complex, coboundary, and rank rows are
not present.
\end{proof}

The packet
\[
\texttt{certificates/orientation/rhomred\_weyl\_orientation\_cocycle\_class\_vanishing}
\]
verifies
\(\texttt{RHOMRED\_WEYL\_ORIENTATION\_COCYCLE\_CLASS\_VANISHING\_OBSTRUCTION\_VERIFIED}\):
it imports the row \(309\) cocycle packet and the blocked
reduced-orientation ledger, records the finite rank criterion for
\([c_{o,R}]=0\), and keeps the cocycle-value, closedness,
cochain-complex, coboundary-matrix, rank-certificate, and witness
cochain rows open.

\begin{proposition}[Cochain trivialization criterion for the
orientation cocycle]
\label{prop:orientation-cocycle-cochain-trivialization-criterion}
\textnormal{[criterion proved; killing \(1\)-cochain not supplied]}
Assume the data of
Proposition~\ref{prop:orientation-cocycle-class-vanishing-criterion}
are supplied and that row \(310\) has produced a rank certificate for
\([c_{o,R}]=0\).  A row \(311\) cochain trivialization is a named
vector
\[
b_R\in C^1(\mathcal G_R;\underline{\mathbb F}_{2,o})
\]
such that, in the ordered bases used for row \(310\),
\[
d^1_R b_R=v(c_{o,R}).
\]
Equivalently, it is a chosen lift of \(v(c_{o,R})\) through the
coboundary map \(d^1_R\), together with a defect row
\[
\operatorname{def}(b_R)
:=d^1_R b_R-v(c_{o,R})=0
\]
over \(\mathbb F_2\).  The class vanishing in row \(310\) proves that
such \(b_R\) exists; row \(311\) chooses one, records its coordinates,
and verifies the displayed equation.  Compatibility of these choices
under \(R\to R'\) is row \(312\), not part of row \(311\).

The present finite tables do not construct row \(311\).  The row
\(309\) cocycle values are absent, the row \(310\) rank-certificate
rows are absent, and no vector \(b_R\), gauge normalization, or
coboundary-defect row is supplied.
\end{proposition}

\begin{proof}
Once ordered bases are fixed, a \(1\)-cochain is a column vector in the
finite \(\mathbb F_2\)-vector space \(C^1_R\).  The condition that it
kill the projective cocycle is the single matrix equation
\(d^1_R b_R=v(c_{o,R})\).  Therefore the construction of \(b_R\) is
stronger than the rank equality of row \(310\): it supplies an actual
preimage, not only the statement that a preimage exists.  The residual
freedom is the addition of a \(1\)-cocycle in
\(\ker d^1_R\), so a complete finite row must also record either a
gauge convention or the torsor in which the choice lives.  Inspection
of the row \(309\) and row \(310\) packets shows that the input
cocycle, the coboundary matrix, the rank certificate, and the witness
vector are not present.
\end{proof}

The packet
\[
\texttt{certificates/orientation/rhomred\_weyl\_orientation\_cocycle\_cochain\_trivialization}
\]
verifies
\(\texttt{RHOMRED\_WEYL\_ORIENTATION\_COCYCLE\_COCHAIN\_TRIVIALIZATION\_OBSTRUCTION\_VERIFIED}\):
it imports the row \(309\) cocycle packet and the row \(310\)
class-vanishing packet, records the equation defining the killing
\(1\)-cochain, and keeps the cochain-vector, gauge-normalization, and
coboundary-defect rows open.

\begin{proposition}[Transition compatibility for the Coxeter cochain]
\label{prop:orientation-cocycle-cochain-transition-compatibility}
\textnormal{[criterion proved; transition rows not supplied]}
Let \(R'\le R\).  Suppose rows \(309\)--\(311\) have supplied
\(\mathcal G_R,\mathcal G_{R'}\), the coefficient systems, cocycles,
cochain complexes, and chosen cochains \(b_R\) and \(b_{R'}\).  A row
\(312\) transition compatibility datum consists of cochain maps
\[
\rho^i_{RR'}:
C^i(\mathcal G_{R'};\underline{\mathbb F}_{2,o})
\longrightarrow
C^i(\mathcal G_R;\underline{\mathbb F}_{2,o}),
\qquad i=1,2,
\]
induced by the finite-stage transition \(t_{RR'}\), such that
\[
d^1_R\rho^1_{RR'}=\rho^2_{RR'}d^1_{R'},\qquad
\rho^2_{RR'}v(c_{o,R'})=v(c_{o,R}),
\]
and
\[
\rho^1_{RR'}(b_{R'})=b_R.
\]
Equivalently, the transition defect
\[
\Delta^b_{RR'}:=b_R-\rho^1_{RR'}(b_{R'})
\in C^1(\mathcal G_R;\underline{\mathbb F}_{2,o})
\]
must be zero in the chosen gauge.  For a chain \(R''\le R'\le R\),
the cochain maps must also satisfy
\[
\rho^1_{RR''}=\rho^1_{RR'}\rho^1_{R'R''},\qquad
\rho^2_{RR''}=\rho^2_{RR'}\rho^2_{R'R''}
\]
on the retained cochain bases.  Thus the \(b_R\) form a strict inverse
system of cochain trivializations.  Compatibility only up to an
unrecorded \(1\)-cocycle is not enough; the comparison cocycle or the
gauge fixing must be part of the finite row.

The present tables do not prove row \(312\).  No row \(311\) cochain
vectors are supplied, the transition-orientation packet has no Coxeter
cochain transition row, and the reduced-orientation transition table
is empty.
\end{proposition}

\begin{proof}
The first displayed equation says that the transition maps form a
cochain map.  The second says that the projective cocycle at level
\(R'\) pulls back to the projective cocycle at level \(R\).  If
\(d^1_{R'}b_{R'}=v(c_{o,R'})\), these two equations imply
\[
d^1_R\rho^1_{RR'}(b_{R'})
=\rho^2_{RR'}d^1_{R'}b_{R'}
=\rho^2_{RR'}v(c_{o,R'})
=v(c_{o,R}).
\]
Thus both \(\rho^1_{RR'}(b_{R'})\) and \(b_R\) kill \(c_{o,R}\).
Row \(312\) asks that the chosen lifts agree, not merely that both
exist.  The defect \(\Delta^b_{RR'}\) is therefore a closed
\(1\)-cochain; strict compatibility is exactly the assertion
\(\Delta^b_{RR'}=0\) in the fixed gauge, together with the two-step
composition equations for the \(\rho^i\).  Inspection of the row
\(311\) packet, the transition-orientation packet, and
\(\texttt{certificates/orientation/k3e\_reduced\_orientation/transitions.csv}\)
shows that none of these finite rows are present.
\end{proof}

The packet
\[
\texttt{certificates/orientation/rhomred\_weyl\_orientation\_cocycle\_cochain\_transition}
\]
verifies
\(\texttt{RHOMRED\_WEYL\_ORIENTATION\_COCYCLE\_COCHAIN\_TRANSITION\_OBSTRUCTION\_VERIFIED}\):
it imports the row \(311\) cochain packet and the transition-orientation
ledger, records the equations defining compatibility under
\(R\to R'\), and keeps the transition cochain-map, cocycle-pullback,
cochain-defect, gauge, and two-step coherence rows open.

\begin{proposition}[Computation criterion for the Weyl torsor defects]
\label{prop:weyl-orientation-torsor-defect-computation}
\textnormal{[criterion proved; torsor-defect rows not supplied]}
Fix \(R\), a retained quotient stratum \(S\), and a simple type-II
reflection \(s_{\delta_i}\).  Suppose the row \(306\) lift
\(\tau_{i,R,S}\) and the quotient-orientation packages
\[
\mathfrak q_{R,S},\qquad
\mathfrak q_{R,s_{\delta_i}S}
\]
are supplied by explicit Cech--Borel representatives, null
trivialisations, finite-stabilizer edge representatives, and
linearization characters.  Let
\[
\Theta^1_{i,R,S}\colon
\mathcal Q^1_{R,S}\longrightarrow
\mathcal Q^1_{R,s_{\delta_i}S}
\]
be the degree-one cochain map induced by \(\tau_{i,R,S}\) on the
quotient-orientation complex
\[
\mathcal Q^\bullet_{R,S}:=
C^\bullet(\mathfrak M^{\mathrm{red}}_{R,S};\mathbb F_2)
\oplus
\bigoplus_{H\in\mathcal H(R,S)}C^\bullet(BH;\mathbb F_2),
\]
with the connected \(E\)-free summand already killed by its row \(284\)
null-trivialization.  If
\(\eta_{R,S}\in\mathcal Q^1_{R,S}\) and
\(\eta_{R,s_{\delta_i}S}\in\mathcal Q^1_{R,s_{\delta_i}S}\) denote the
chosen quotient-orientation null trivialisations, the torsor-defect
cochain is
\[
\Delta^\omega_{i,R,S}:=
\Theta^1_{i,R,S}(\eta_{R,S})-\eta_{R,s_{\delta_i}S}.
\]
After the degree-two quotient representatives have been transported
with zero defect, \(\Delta^\omega_{i,R,S}\) is closed.  The row \(313\)
torsor defect is the cohomology class
\[
\omega_{i,R,S}:=
[\Delta^\omega_{i,R,S}]
\in
H^1(\mathfrak M^{\mathrm{red}}_{R,s_{\delta_i}S};\mathbb F_2)
\oplus
\bigoplus_{H\in\mathcal H(R,s_{\delta_i}S)}
H^1(BH;\mathbb F_2),
\]
recorded in a finite basis for the displayed \(H^1\)-space.  Thus row
\(313\) computes the defect vector; row \(314\) is the separate
assertion that this vector is zero for every retained \(i,R,S\).

The present finite tables do not compute row \(313\).  They contain no
row \(306\) Weyl lifts, no quotient Borel representatives, no
finite-stabilizer transport rows, no linearization transport rows, no
\(\mathcal Q^\bullet\)-basis rows, and no torsor-defect vector rows.
\end{proposition}

\begin{proof}
The quotient orientation is not a single line: it is the orientation
line together with null trivialisations of the reduced gerbe, the
free-\(E\) Borel class, the finite-stabilizer Borel classes, and the
zero linearization characters.  Transport along \(\tau_{i,R,S}\)
therefore gives cochain maps on the Cech--Borel complexes carrying
these representatives.  If the degree-two representatives are
transported with zero defect, the difference between the transported
source null-trivialization and the target null-trivialization has zero
coboundary; hence it defines the displayed \(H^1\)-class.  Coordinates
of this class in a chosen finite basis are exactly the torsor-defect
computation.  A target Weyl relation, Maass-character value, or local
Pfaffian sign has no such cochain representative and cannot compute
\(\omega_{i,R,S}\).  Inspection of the reduced-orientation and
row \(306\) packets shows that the required rows are absent.
\end{proof}

The packet
\[
\texttt{certificates/orientation/rhomred\_weyl\_orientation\_torsor\_defects}
\]
verifies
\(\texttt{RHOMRED\_WEYL\_ORIENTATION\_TORSOR\_DEFECTS\_OBSTRUCTION\_VERIFIED}\):
it imports the row \(306\) Weyl-wall transport criterion, the row
\(312\) cochain-transition packet, and the blocked reduced-orientation
ledger; it records the finite cochain formula for
\(\omega_{i,R,S}\) and keeps the quotient-orientation representative,
transport, \(H^1\)-basis, and defect-vector rows open.

\begin{proposition}[Vanishing criterion for the Weyl torsor defects]
\label{prop:weyl-orientation-torsor-defect-vanishing}
\textnormal{[criterion proved; zero-defect rows not supplied]}
Assume the row \(313\) data of
Proposition~\ref{prop:weyl-orientation-torsor-defect-computation} have
been supplied for every retained triple \((i,R,S)\).  Thus each
\(\Delta^\omega_{i,R,S}\) is a closed cochain in the target
quotient-orientation complex and has coordinates
\[
v(\omega_{i,R,S})\in
H^1(\mathfrak M^{\mathrm{red}}_{R,s_{\delta_i}S};\mathbb F_2)
\oplus
\bigoplus_{H\in\mathcal H(R,s_{\delta_i}S)}
H^1(BH;\mathbb F_2)
\]
in a fixed finite basis.  Row \(314\) is proved exactly when the
finite vanishing table contains, for every retained triple
\((i,R,S)\), the following four rows:
\[
d\Delta^\omega_{i,R,S}=0,\qquad
v(\omega_{i,R,S})=0,\qquad
\operatorname{rank}\bigl(v(\omega_{i,R,S})\bigr)=0,
\qquad
\mathrm{cov}\{(i,R,S)\}=\mathrm{cov}_{\mathrm{ret}} .
\]
Equivalently, the row \(313\) defect vectors must be present, closed,
computed in the displayed \(H^1\)-basis, and zero with complete
coverage over the retained Weyl-wall atlas.  Under these hypotheses
the Weyl lift \(\tau_{i,R,S}\) preserves the quotient-orientation
torsor, and \(\omega_{i,R,S}=0\) for all retained \(i,R,S\).

The present finite tables do not prove row \(314\).  The row \(313\)
packet contains no torsor-defect vector rows, no \(H^1\)-basis rows,
no zero-vector certificates, and no coverage rows for retained
triples.  Therefore the vanishing of \(\omega_{i,R,S}\) remains an
open orientation obligation.
\end{proposition}

\begin{proof}
In the fixed \(H^1\)-basis, a closed defect cochain represents the
zero torsor class if and only if its coordinate vector is zero.  Thus
the four displayed finite checks are precisely the statement that
each class \(\omega_{i,R,S}\) vanishes, together with the assertion
that no retained type-II wall has been omitted.  This is the needed
source-side orientation statement: it says that the transported null
trivialisation equals the target null trivialisation in the quotient
orientation torsor.  A target Weyl relation, a Maass character, a
Pfaffian wall sign, or an empty defect table is not a coordinate
calculation in the displayed \(H^1\)-space.  Inspection of the row
\(313\) packet shows that the required defect vectors and zero
certificates are not supplied.
\end{proof}

The packet
\[
\texttt{certificates/orientation/rhomred\_weyl\_orientation\_torsor\_defect\_vanishing}
\]
verifies
\(\texttt{RHOMRED\_WEYL\_ORIENTATION\_TORSOR\_DEFECT\_VANISHING\_OBSTRUCTION\_VERIFIED}\):
it imports the row \(313\) torsor-defect computation criterion,
records the finite zero-vector criterion for row \(314\), and keeps
the computed defect-vector, closedness, zero-certificate, and retained
coverage rows open.

\begin{proposition}[Weyl preservation criterion for the reduced gerbe]
\label{prop:weyl-alpha-red-preservation-criterion}
\textnormal{[criterion proved; preservation rows not supplied]}
Fix \(R\), a retained quotient stratum \(S\), and a simple type-II
reflection \(s_{\delta_i}\).  Suppose the row \(306\) lift
\(\tau_{i,R,S}\) is supplied, and suppose row \(280\) supplies
Cech--Borel cocycles
\[
a^{\mathrm{red}}_{R,S}\in Z^2(\mathcal B_{R,S};\mathbb F_2),
\qquad
a^{\mathrm{red}}_{R,s_{\delta_i}S}
\in Z^2(\mathcal B_{R,s_{\delta_i}S};\mathbb F_2)
\]
representing
\(\alpha^{\mathrm{red}}_{R,S}\) and
\(\alpha^{\mathrm{red}}_{R,s_{\delta_i}S}\).  Let
\[
\Theta^{\bullet,\mathrm{red}}_{i,R,S}\colon
C^\bullet(\mathcal B_{R,S};\mathbb F_2)
\longrightarrow
C^\bullet(\mathcal B_{R,s_{\delta_i}S};\mathbb F_2)
\]
be the cochain map induced by the Weyl transport on the reduced
Cech--Borel groupoids.  The row \(315\) preservation defect is
\[
\Delta^{\mathrm{red}}_{i,R,S}:=
\Theta^{2,\mathrm{red}}_{i,R,S}(a^{\mathrm{red}}_{R,S})
-a^{\mathrm{red}}_{R,s_{\delta_i}S}.
\]
Weyl transport preserves the reduced gerbe on this retained wall
exactly when there is a finite cochain
\[
u^{\mathrm{red}}_{i,R,S}\in
C^1(\mathcal B_{R,s_{\delta_i}S};\mathbb F_2)
\]
with
\[
d u^{\mathrm{red}}_{i,R,S}=\Delta^{\mathrm{red}}_{i,R,S};
\]
equivalently,
\[
[\Delta^{\mathrm{red}}_{i,R,S}]=0
\in H^2(\mathcal B_{R,s_{\delta_i}S};\mathbb F_2).
\]
The finite row proving row \(315\) must record the source and target
\(\alpha^{\mathrm{red}}\)-cocycles, the \(\Theta^{2,\mathrm{red}}\)
matrix, the transported defect vector, the coboundary witness
\(u^{\mathrm{red}}_{i,R,S}\), a zero class-rank certificate, and
coverage of every retained triple \((i,R,S)\).  If strict
representative preservation is imposed as the chosen gauge, the same
row must record the stronger condition
\(\Delta^{\mathrm{red}}_{i,R,S}=0\).

The present finite tables do not prove row \(315\).  The row \(280\)
packet contains no computed \(\alpha^{\mathrm{red}}\)-values, the row
\(281\) packet contains no Cech--Borel null-homotopies, the row \(306\)
packet contains no Weyl-lift cochain maps, and row \(314\) has not
proved vanishing of the quotient-orientation torsor defects.
\end{proposition}

\begin{proof}
Because \(\Theta^{\bullet,\mathrm{red}}_{i,R,S}\) is required to be a
cochain map, the displayed defect is closed whenever the source and
target representatives are closed and the degree-three transport
defect is zero.  The transported reduced gerbe class equals the target
reduced gerbe class precisely when this closed defect is a coboundary.
A chosen cochain \(u^{\mathrm{red}}_{i,R,S}\) with
\(d u^{\mathrm{red}}_{i,R,S}=\Delta^{\mathrm{red}}_{i,R,S}\) is the
finite certificate of equality in \(H^2\).  Complete retained-wall
coverage is part of the statement because \(O1^+\) is a Weyl-equivariant
orientation assertion, not a single-wall computation.  A target Weyl
relation, a scalar multiplier, a Pfaffian sign, or the uncomputed
vanishing of the full torsor defect is not a cochain calculation for
\(\alpha^{\mathrm{red}}\).  Inspection of the row \(280\), row \(281\),
row \(306\), and row \(314\) packets shows that the needed rows are
absent.
\end{proof}

The packet
\[
\texttt{certificates/orientation/rhomred\_weyl\_alpha\_red\_preservation}
\]
verifies
\(\texttt{RHOMRED\_WEYL\_ALPHA\_RED\_PRESERVATION\_OBSTRUCTION\_VERIFIED}\):
it imports the reduced-gerbe, reduced-gerbe null-homotopy, Weyl-wall
transport, and torsor-defect vanishing packets; it records the finite
cochain criterion for row \(315\); and it keeps the source-target
\(\alpha^{\mathrm{red}}\)-cocycle, transport matrix, coboundary
witness, zero class-rank, and retained coverage rows open.

\begin{proposition}[Weyl preservation criterion for the connected
free \(E\)-Borel class]
\label{prop:weyl-alpha-e-free-preservation-criterion}
\textnormal{[criterion proved; preservation rows not supplied]}
Fix \(R\), a retained free \(E\)-quotient stratum \(S\), and a simple
type-II reflection \(s_{\delta_i}\).  Suppose the row \(306\) lift
\(\tau_{i,R,S}\) is supplied and suppose row \(282\) supplies
connected \(BE\)-edge representatives
\[
a^{E,\mathrm{free}}_{R,S},\qquad
a^{E,\mathrm{free}}_{R,s_{\delta_i}S}
\in Z^2(BE;\mathbb F_2)
\]
with coordinates
\[
\alpha^{E,\mathrm{free}}_{R,S}=a_1(R,S)u_1+a_2(R,S)u_2,\qquad
\alpha^{E,\mathrm{free}}_{R,s_{\delta_i}S}
=a_1(R,s_{\delta_i}S)u_1+a_2(R,s_{\delta_i}S)u_2
\]
in the basis of Lemma~\ref{lem:connected-e-descent}.  Let
\[
M^{E}_{i,R,S}\in \operatorname{GL}_2(\mathbb F_2)
\]
be the induced map of the Weyl transport on
\(H^2(BE;\mathbb F_2)=\mathbb F_2u_1\oplus\mathbb F_2u_2\), computed
from the cochain-level transport of the connected \(E\)-translation
rigidification.  The row \(316\) preservation defect is the coordinate
vector
\[
\Delta^{E,\mathrm{free}}_{i,R,S}:=
M^{E}_{i,R,S}
\begin{pmatrix}a_1(R,S)\\ a_2(R,S)\end{pmatrix}
-
\begin{pmatrix}a_1(R,s_{\delta_i}S)\\
a_2(R,s_{\delta_i}S)\end{pmatrix}
\in\mathbb F_2^2 .
\]
Weyl transport preserves the connected free \(E\)-Borel class on this
retained wall exactly when
\[
\Delta^{E,\mathrm{free}}_{i,R,S}=0.
\]
Equivalently, in a cochain model one may record a degree-one witness
for the difference between the transported source representative and
the target representative; after passing to the connected \(BE\)-edge,
this is the same as the displayed zero coordinate vector, since
\(u_1,u_2\) form a basis of \(H^2(BE;\mathbb F_2)\).

The finite row proving row \(316\) must therefore contain the source
and target \(a_1,a_2\) coefficient rows, the induced matrix
\(M^E_{i,R,S}\), the transported coordinate vector, the zero-vector
certificate, the underlying cochain-transport defect if strict
representatives are used, and coverage of every retained triple
\((i,R,S)\).  The present finite tables do not prove row \(316\):
rows \(282\), \(283\), and \(284\) supply no connected \(BE\)-edge
coefficients, zero checks, or null-trivialisations; row \(306\)
supplies no Weyl cochain map; and row \(315\) has not proved
preservation of \(\alpha^{\mathrm{red}}\).
\end{proposition}

\begin{proof}
The connected part of the quotient orientation is controlled by the
degree-two edge of \(H^*_E(-;\mathbb F_2)\).  In degree two the
connected \(BE\)-edge has the fixed basis \(u_1,u_2\).  Therefore a
Weyl transport preserves \(\alpha^{E,\mathrm{free}}\) if and only if
the transported coordinate vector equals the target coordinate vector.
The displayed equation is exactly this equality in
\(\mathbb F_2^2\).  If the comparison is made before passage to the
edge, the difference of cocycle representatives must be recorded as a
coboundary; its edge coordinate is the same defect vector.  Freeness of
the \(E\)-action, \(H^1(BE;\mathbb F_2)=0\), target Weyl relations,
Maass-character values, and Pfaffian signs do not compute the
\(H^2(BE)\) coordinate transport.  Inspection of the row \(282\), row
\(283\), row \(284\), row \(306\), and row \(315\) packets shows that
the required finite rows are absent.
\end{proof}

The packet
\[
\texttt{certificates/orientation/rhomred\_weyl\_alpha\_e\_free\_preservation}
\]
verifies
\(\texttt{RHOMRED\_WEYL\_ALPHA\_E\_FREE\_PRESERVATION\_OBSTRUCTION\_VERIFIED}\):
it imports the connected free \(E\)-Borel computation, vanishing,
null-trivialization, Weyl-wall transport, and reduced-gerbe
preservation packets; it records the finite coordinate criterion for
row \(316\); and it keeps the source-target \(a_1,a_2\) rows, induced
\(H^2(BE)\)-matrix, zero-vector certificate, cochain-transport defect,
and retained coverage rows open.

\begin{proposition}[Weyl preservation criterion for the finite-stabilizer
Borel classes]
\label{prop:weyl-finite-stabilizer-beta-preservation-criterion}
\textnormal{[criterion proved; preservation rows not supplied]}
Fix \(R\), a retained finite-stabilizer stratum \(S\), and a simple
type-II reflection \(s_{\delta_i}\).  Suppose the row \(306\) lift
\(\tau_{i,R,S}\) is supplied and identifies the source and target
finite stabilizers by an isomorphism
\[
\varphi_{i,R,S}\colon H_{R,S}\xrightarrow{\sim}
H_{R,s_{\delta_i}S}.
\]
Suppose row \(285\) has supplied classifying-space edge representatives
\[
\beta^{H_{R,S}}_{R,S}\in H^2(BH_{R,S};\mathbb F_2),\qquad
\beta^{H_{R,s_{\delta_i}S}}_{R,s_{\delta_i}S}
\in H^2(BH_{R,s_{\delta_i}S};\mathbb F_2),
\]
with coordinates in the appropriate odd, Klein-four, or two-primary
basis of Proposition~\ref{prop:finite-stabilizer-beta-class-criterion}.
Let
\[
M^{H}_{i,R,S}\colon
H^2(BH_{R,S};\mathbb F_2)\longrightarrow
H^2(BH_{R,s_{\delta_i}S};\mathbb F_2)
\]
be the matrix induced by \(B\varphi_{i,R,S}\) on the chosen
classifying-space bases.  The row \(317\) preservation defect is
\[
\Delta^{\beta}_{i,R,S}:=
M^{H}_{i,R,S}\,v(\beta^{H_{R,S}}_{R,S})
-v(\beta^{H_{R,s_{\delta_i}S}}_{R,s_{\delta_i}S}).
\]
Weyl transport preserves the finite-stabilizer Borel class on this
retained wall exactly when
\[
\Delta^{\beta}_{i,R,S}=0
\]
for every retained triple \((i,R,S)\).  If the comparison is made at
the Cech--Borel representative level, the row must also record the
transport of \(\widetilde\beta^H\), the edge-reduction compatibility,
and a zero cochain-level transport defect before taking the displayed
classifying-space coordinate vector.

Thus a row proving \(317\) must contain the source and target
finite-stabilizer orientation rows, the stabilizer isomorphism
\(\varphi_{i,R,S}\), the induced \(H^2(BH)\)-matrix, the transported
\(\beta^H\)-coordinate vector, a zero-vector certificate, edge-reduction
compatibility, and complete retained-wall coverage.  The present
finite tables do not prove row \(317\): row \(285\) has no
finite-stabilizer orientation or edge rows, row \(286\) has no
\(\beta^H=0\) rows, row \(287\) has no compatible
null-trivializations, row \(306\) has no Weyl lift on stabilizers, and
row \(316\) has not proved preservation of \(\alpha^{E,\mathrm{free}}\).
\end{proposition}

\begin{proof}
The finite-stabilizer component of the quotient orientation is the
classifying-space edge of the \(H\)-equivariant Borel class
\(\widetilde\beta^H\).  A Weyl lift can preserve this component only
after it identifies the source and target stabilizer groups and carries
the source edge representative to the target edge representative.  In
the fixed group-cohomology bases this condition is precisely the
vanishing of the displayed coordinate defect.  For odd or trivial
stabilizers the target space is zero only after the finite-stabilizer
edge row has been supplied; for rank-two two-primary stabilizers the
mixed coefficient is part of the basis and cannot be read from cyclic
restrictions.  A residual group order, a scalar trace, a Maass
character, a Pfaffian sign, or preservation of the connected free
\(E\)-component is not a finite-stabilizer edge transport calculation.
Inspection of the row \(285\), row \(286\), row \(287\), row \(306\),
and row \(316\) packets shows that the required finite rows are absent.
\end{proof}

The packet
\[
\texttt{certificates/orientation/rhomred\_weyl\_finite\_stabilizer\_beta\_preservation}
\]
verifies
\(\texttt{RHOMRED\_WEYL\_FINITE\_STABILIZER\_BETA\_PRESERVATION\_OBSTRUCTION\_VERIFIED}\):
it imports the finite-stabilizer Borel computation, vanishing,
null-trivialization, Weyl-wall transport, and connected free
\(E\)-preservation packets; it records the finite edge-coordinate
criterion for row \(317\); and it keeps the stabilizer isomorphism,
source-target \(\beta^H\)-coordinates, \(H^2(BH)\)-matrix,
edge-transport defect, zero-vector certificate, and retained coverage
rows open.

\begin{proposition}[Weyl preservation criterion for the zero
finite-stabilizer linearization]
\label{prop:weyl-finite-stabilizer-lambda-zero-preservation-criterion}
\textnormal{[criterion proved; linearization-transport rows not
supplied]}
Fix \(R\), a retained finite-stabilizer stratum \(S\), and a simple
type-II reflection \(s_{\delta_i}\).  Suppose the row \(306\) lift
\(\tau_{i,R,S}\) is supplied and identifies the source and target
finite stabilizers by an isomorphism
\[
\varphi_{i,R,S}\colon H_{R,S}\xrightarrow{\sim}
H_{R,s_{\delta_i}S}.
\]
Suppose row \(288\) has supplied the finite-stabilizer
linearization characters
\[
\lambda^{H_{R,S}}_{R,S}\in H^1(BH_{R,S};\mathbb F_2),
\qquad
\lambda^{H_{R,s_{\delta_i}S}}_{R,s_{\delta_i}S}
\in H^1(BH_{R,s_{\delta_i}S};\mathbb F_2),
\]
with coordinates in the bases of
Proposition~\ref{prop:finite-stabilizer-linearization-character-criterion},
and suppose row \(289\) has supplied the corresponding zero rows.
Let
\[
N^{H}_{i,R,S}\colon
H^1(BH_{R,S};\mathbb F_2)\longrightarrow
H^1(BH_{R,s_{\delta_i}S};\mathbb F_2)
\]
be the matrix induced by \(B\varphi_{i,R,S}\) on the chosen
degree-one classifying-space bases.  The row \(318\)
linearization-transport defect is
\[
\Delta^{\lambda}_{i,R,S}:=
N^H_{i,R,S}\,v(\lambda^{H_{R,S}}_{R,S})
-v(\lambda^{H_{R,s_{\delta_i}S}}_{R,s_{\delta_i}S}).
\]
Weyl transport preserves the zero finite-stabilizer linearization on
this retained wall exactly when
\[
\Delta^{\lambda}_{i,R,S}=0
\]
for every retained triple \((i,R,S)\).  If the comparison is made at
the orientation-line action level, the row must also record the
transport of the sign action of \(H_{R,S}\) to
\(H_{R,s_{\delta_i}S}\) and a zero action-transport defect before
passing to \(H^1(BH;\mathbb F_2)\).

Thus a row proving \(318\) must contain the source and target
\(\lambda^H\)-coordinate rows, the stabilizer isomorphism
\(\varphi_{i,R,S}\), the induced \(H^1(BH)\)-matrix, the transported
coordinate vector, a rank-zero defect row, orientation-line action
transport, and complete retained-wall coverage.  The present finite
tables do not prove row \(318\): row \(288\) has no orientation-line
action rows, generator-basis rows, character values, or computed
\(\lambda^H\)-coordinates; row \(289\) has no zero-coordinate rows;
row \(306\) has no Weyl lift on finite stabilizers; and row \(317\)
has not proved preservation of the degree-two finite-stabilizer Borel
classes.
\end{proposition}

\begin{proof}
After the degree-two finite-stabilizer gerbe has been trivialized, the
remaining descent ambiguity of the orientation line is the sign
character
\[
\lambda^H\in H^1(BH;\mathbb F_2)=\operatorname{Hom}(H,\mathbb F_2).
\]
A Weyl lift can preserve the zero character only after it identifies
the source and target stabilizer groups and transports the source sign
action to the target sign action.  In the chosen \(H^1(BH)\)-bases,
this condition is exactly the vanishing of the displayed coordinate
defect.  If rows \(288\) and \(289\) have supplied zero characters on
both sides, the class-level equality follows from functoriality of
group cohomology, but the transport still has no mathematical meaning
without the stabilizer isomorphism, the induced degree-one matrix, and
the action-level comparison.  The identity
\(H^1(B1;\mathbb F_2)=0\), a vanishing \(\beta^H\)-class, a
null-trivialization of \(\beta^H\), a connected \(E\)-Borel
calculation, a Maass-character value, a Pfaffian wall sign, or a scalar
trace does not compute this degree-one transport.  Inspection of the
row \(288\), row \(289\), row \(306\), and row \(317\) packets shows
that the required finite rows are absent.
\end{proof}

The packet
\[
\texttt{certificates/orientation/rhomred\_weyl\_finite\_stabilizer\_lambda\_zero\_preservation}
\]
verifies
\(\texttt{RHOMRED\_WEYL\_FINITE\_STABILIZER\_LAMBDA\_ZERO\_PRESERVATION\_OBSTRUCTION\_VERIFIED}\):
it imports the finite-stabilizer linearization-character,
linearization-vanishing, Weyl-wall transport, and finite-stabilizer
Borel-preservation packets; it records the degree-one coordinate
criterion for row \(318\); and it keeps the stabilizer isomorphism,
source-target \(\lambda^H\)-coordinates, \(H^1(BH)\)-matrix,
orientation-line action transport, coordinate defect, rank-zero row,
and retained coverage rows open.

\begin{definition}[Orientation character of the strictified Weyl
transport]
\label{def:orientation-character-epsilon-o}
\textnormal{[definition specified; finite character rows not supplied]}
Assume rows \(306\)--\(318\) have supplied a strictified
\(W^{(2)}(\Lambda_{II}^{2,1})\)-transport of the complete
quotient-oriented determinant-line package on a retained finite
height \(R\): the finite partial action groupoid \(\mathcal G_R\),
the lifts \(\tau_{i,R,S}\), the cochain \(b_R\) killing the
projective cocycle \(c_{o,R}\), transition compatibility of \(b_R\),
vanishing torsor defects, and preservation of
\(\alpha^{\mathrm{red}}\), \(\alpha^{E,\mathrm{free}}\),
\(\beta^H\), and \(\lambda^H=0\).  For each simple type-II reflection
\(s_{\delta_i}\) and retained source stratum \(S\), let
\[
\sigma^o_{i,R,S}\in\{\pm1\}
\]
be the residual sign of the strictified lift on the oriented
one-dimensional Pfaffian square-root line after the target line and
the complete quotient-orientation package have been identified with
the transported source package.  A row defining the finite-stage
orientation character must prove that \(\sigma^o_{i,R,S}\) is
independent of \(S\) on every retained connected wall component and is
compatible under \(R'\le R\).  Write the resulting sign as
\(\sigma^o_{i,R}\).

The finite-stage orientation character is then the homomorphism
\[
\epsilon_{o,R}\colon W^{(2)}(\Lambda_{II}^{2,1})\longrightarrow
\{\pm1\}
\]
defined on a reduced word
\[
w=s_{\delta_{i_1}}\cdots s_{\delta_{i_m}}
\]
by
\[
\epsilon_{o,R}(w):=
\prod_{a=1}^{m}\sigma^o_{i_a,R}.
\]
The type-II Coxeter presentation has only the relations
\[
s_{\delta_i}^2=1,\qquad i=1,2,3,
\]
and no braid relations between distinct simple reflections.  Hence the
displayed formula is well-defined once the generator signs are
supplied, since \((\sigma^o_{i,R})^2=1\).  A compatible inverse system
of finite characters defines
\[
\epsilon_o=\varprojlim_R\epsilon_{o,R}.
\]

This is a definition from source-side orientation transport.  It is
not the automorphic Maass character
\(\nu_{\Delta_5}\), not the determinant character of the target Coxeter
group, not the OP scalar branch, and not a consequence of the scalar
trace.  Row \(320\) computes the generator values
\(\epsilon_o(s_{\delta_i})\); row \(321\) compares the resulting
character with the determinant character; row \(322\) compares it with
\(\nu_{\Delta_5}\).  The present finite tables do not define
\(\epsilon_{o,R}\): row \(312\) has no compatible cochain \(b_R\), row
\(314\) has no zero torsor-defect rows, rows \(315\)--\(318\) do not
prove preservation of the four quotient-orientation components, and no
generator-sign or transition-compatible character row is supplied.
\end{definition}

The packet
\[
\texttt{certificates/orientation/rhomred\_weyl\_orientation\_character\_definition}
\]
verifies
\(\texttt{RHOMRED\_WEYL\_ORIENTATION\_CHARACTER\_DEFINITION\_OBSTRUCTION\_VERIFIED}\):
it imports the Coxeter-cochain transition, torsor-defect vanishing,
and four component-preservation packets; it records the word-evaluation
definition of \(\epsilon_{o,R}\); it imports the automorphic Maass
character only as a firewall; and it keeps the strict Weyl action,
generator-sign, stratum-independence, transition-compatibility,
finite-character, and inverse-limit character rows open.

\begin{proposition}[Simple-wall sign criterion for the orientation
character]
\label{prop:orientation-character-simple-wall-sign-criterion}
\textnormal{[local sign formula proved; finite simple-wall sign rows
not supplied]}
Assume the finite-stage orientation character of
Definition~\ref{def:orientation-character-epsilon-o} has been supplied.
Fix a simple type-II reflection \(s_{\delta_i}\), a retained generic
wall stratum \(S\subset \Hpl_{\delta_i}\), and a retained
\((O2_{\mathrm{atlas}})\) wall chart with real normal coordinate
\(u_i\).  Suppose the chart supplies the rank-one normal Pfaffian
crossing
\[
K^{\mathrm{nor}}_{\delta_i,R,S}
\simeq [\mathbb R\xrightarrow{u_i}\mathbb R],
\qquad
\operatorname{Pf}(K^{\mathrm{nor}}_{\delta_i,R,S})
=\upsilon_{i,R,S}u_i,
\]
with
\[
s_{\delta_i}(u_i)=-u_i,\qquad
s_{\delta_i}(\upsilon_{i,R,S})=\upsilon_{i,R,S},\qquad
N_{\delta_i,R,S}^{\mathrm{Pf}}=1.
\]
Then the source-side generator sign of row \(320\) is
\[
\sigma^o_{i,R,S}
=\frac{s_{\delta_i}^{*}(\upsilon_{i,R,S}u_i)}
{\upsilon_{i,R,S}u_i}
=-1
=(-1)^{N_{\delta_i,R,S}^{\mathrm{Pf}}}.
\]
If the row also supplies stratum-independence and transition
compatibility for these signs, then
\[
\epsilon_{o,R}(s_{\delta_i})=-1,\qquad i=1,2,3,
\]
and the compatible inverse-limit character satisfies
\[
\epsilon_o(s_{\delta_i})=-1,\qquad i=1,2,3.
\]

Thus a row proving \(320\) must contain, for all three simple type-II
walls, retained wall-object rows, rank-one wall-chart rows, the normal
coordinate flip \(u_i\mapsto -u_i\), invariant tangent Pfaffian units,
normal rank \(1\), divisor order \(1\), local sign rows, the row \(319\)
orientation-character row, stratum-independence rows, and transition
compatibility rows.  The present finite tables do not prove row
\(320\): the row \(319\) packet has no constructed orientation
character or generator-sign rows, the \((O2_{\mathrm{atlas}})\) packet
has no wall-object, rank-one chart, coordinate-flip, invariant-unit, or
local-sign rows, and the finite Pfaffian packet has no type-II wall
chart or sign rows.
\end{proposition}

\begin{proof}
The local calculation is the rank-one Pfaffian crossing of
Proposition~\ref{prop:appE-local-pfaffian-sign}.  On the retained
chart the tangent Pfaffian unit is invariant and the reflection changes
only the normal coordinate:
\[
s_{\delta_i}^{*}(\upsilon_{i,R,S}u_i)
=\upsilon_{i,R,S}(-u_i)
=-\upsilon_{i,R,S}u_i.
\]
The normal complex contains one odd crossing, so
\(N_{\delta_i,R,S}^{\mathrm{Pf}}=1\) and the sign is
\((-1)^1=-1\).  Squaring removes this sign:
\[
s_{\delta_i}^{*}((\upsilon_{i,R,S}u_i)^2)
=(\upsilon_{i,R,S}u_i)^2.
\]
Hence squared determinants, scalar traces, and OP scalar branches
cannot recover the orientation character.  The Maass character value
\(\nu_{\Delta_5}(s_{\delta_i})=-1\) is an automorphic target value; it
does not supply the source wall chart, invariant unit, normal-rank row,
or generator-sign row.  Inspection of the row \(319\),
\((O2_{\mathrm{atlas}})\), and finite Pfaffian packets shows that those
finite rows are absent.
\end{proof}

The packet
\[
\texttt{certificates/orientation/rhomred\_weyl\_orientation\_simple\_sign\_computation}
\]
verifies
\(\texttt{RHOMRED\_WEYL\_ORIENTATION\_SIMPLE\_SIGN\_COMPUTATION\_OBSTRUCTION\_VERIFIED}\):
it imports the orientation-character definition, the
\((O2_{\mathrm{atlas}})\) wall-atlas ledger, the finite Pfaffian
ledger, and the automorphic Maass character as a firewall; it records
the rank-one formula for row \(320\); and it keeps the wall-object,
rank-one chart, invariant-unit, coordinate-flip, local-sign,
generator-sign, stratum-independence, and transition-compatibility
rows open.

\begin{proposition}[Determinant criterion for the orientation
character]
\label{prop:orientation-character-determinant-criterion}
\textnormal{[group-theoretic criterion proved; source generator signs
not supplied]}
Let
\[
\det_W:W^{(2)}(\La^{2,1}_{II})\longrightarrow\{\pm1\}
\]
be the determinant character of the reflection representation on
\(\La^{2,1}_{II}\otimes_{\mathbb Z}\mathbb R\).  Suppose the
source-side finite-stage character \(\epsilon_{o,R}\) of
Definition~\ref{def:orientation-character-epsilon-o} has been
constructed and the row \(320\) finite sign rows prove
\[
\epsilon_{o,R}(s_{\delta_i})=-1,\qquad i=1,2,3.
\]
Then
\[
\epsilon_{o,R}=\det_W\quad\text{on}\quad
W^{(2)}(\La^{2,1}_{II}).
\]
If the finite characters are transition-compatible, then the inverse
limit character satisfies
\[
\epsilon_o=\det_W\quad\text{on}\quad
W^{(2)}(\La^{2,1}_{II}).
\]

The present finite data do not prove row \(321\).  The row \(319\)
packet does not construct \(\epsilon_{o,R}\), and the row \(320\)
packet contains no generator-sign rows.  The automorphic packet records
\(\nu_{\Delta_5}|_{W^{(2)}}=\det_W\), but that is a target Maass
character computation, not a source orientation-character computation.
\end{proposition}

\begin{proof}
Each \(s_{\delta_i}\) is a real reflection in the retained
rank-three Lorentzian reflection representation, hence
\[
\det_W(s_{\delta_i})=-1,\qquad i=1,2,3.
\]
The type-II Coxeter presentation used in
Definition~\ref{def:orientation-character-epsilon-o} has generators
\(s_{\delta_1},s_{\delta_2},s_{\delta_3}\), relations
\(s_{\delta_i}^2=1\), and no braid relation between distinct
generators.  Therefore a \(\{\pm1\}\)-valued character is determined
by its three generator values.  If row \(320\) supplies the three
source values \(\epsilon_{o,R}(s_{\delta_i})=-1\), then
\(\epsilon_{o,R}\) and \(\det_W\) agree on the generators and hence on
all words in \(W^{(2)}(\La^{2,1}_{II})\).  The same argument passes to
the inverse limit once the transition-compatibility rows for
\(\epsilon_{o,R}\) are supplied.
\end{proof}

The packet
\[
\texttt{certificates/orientation/rhomred\_weyl\_orientation\_determinant\_comparison}
\]
verifies
\(\texttt{RHOMRED\_WEYL\_ORIENTATION\_DETERMINANT\_COMPARISON\_OBSTRUCTION\_VERIFIED}\):
it imports the row \(319\) orientation-character definition, the row
\(320\) simple-sign criterion, and the automorphic determinant
comparison as a firewall; it records the determinant criterion for row
\(321\); and it keeps the source character, source generator signs,
equality-on-generators, transition compatibility, and inverse-limit
comparison rows open.

\begin{proposition}[Maass-character criterion for the orientation
character]
\label{prop:orientation-character-maass-criterion}
\textnormal{[comparison criterion proved; source determinant comparison
not supplied]}
Assume the row \(321\) source comparison has been supplied:
\[
\epsilon_o=\det_W\quad\text{on}\quad
W^{(2)}(\La^{2,1}_{II}).
\]
Assume also the automorphic Maass character computation for the
Borcherds--Gritsenko--Nikulin square root,
\[
\nu_{\Delta_5}\big|_{W^{(2)}(\La^{2,1}_{II})}
=\det_W .
\]
Then
\[
\epsilon_o=\nu_{\Delta_5}\big|_{W^{(2)}(\La^{2,1}_{II})}.
\]

The present finite data do not prove row \(322\).  The Maass character
packet proves the automorphic equality
\(\nu_{\Delta_5}|_{W^{(2)}}=\det_W\), but the row \(321\) packet does
not prove the source equality \(\epsilon_o=\det_W\).  Therefore the
Maass character cannot be used as a substitute for the missing
source-side orientation character or the missing type-II wall sign
rows.
\end{proposition}

\begin{proof}
The assertion is the transitivity of equality of characters on
\(W^{(2)}(\La^{2,1}_{II})\).  If the source orientation character has
already been identified with \(\det_W\), and the automorphic Maass
character restricts to the same determinant character, then both
characters have identical values on every element of
\(W^{(2)}(\La^{2,1}_{II})\).  The second equality is automorphic target
data.  The first equality is the source-side obstruction in row \(321\)
and is not supplied by the current packets.
\end{proof}

The packet
\[
\texttt{certificates/orientation/rhomred\_weyl\_orientation\_maass\_comparison}
\]
verifies
\(\texttt{RHOMRED\_WEYL\_ORIENTATION\_MAASS\_COMPARISON\_OBSTRUCTION\_VERIFIED}\):
it imports the row \(321\) determinant-comparison criterion and the
automorphic Maass-character packet; it records the transitivity
criterion for row \(322\); and it keeps the source determinant
comparison, inverse-limit source character, and orientation-versus-Maass
comparison rows open.

\begin{proposition}[\(S_3\)-chamber automorphisms have trivial Maass
character]
\label{prop:s3-chamber-maass-triviality}
\textnormal{[proved; automorphic character only]}
The restriction of the Maass character of the Igusa square root to the
type-I chamber automorphism group is trivial:
\[
\nu_{\Delta_5}\big|_{\Aut(\Poly_{II})}=1,\qquad
\Aut(\Poly_{II})\simeq S_3.
\]
Equivalently,
\[
\nu_{\Delta_5}(s_{f_{-2}-f_2})
=\nu_{\Delta_5}(s_{f_2-f_3})
=\nu_{\Delta_5}(s_{f_{-2}-f_3})=+1 .
\]
This is an automorphic multiplier statement.  It does not decide
whether \(\epsilon_o\) extends from
\(W^{(2)}(\La^{2,1}_{II})\) to
\(W^{(2)}(\La^{2,1}_{II})\rtimes S_3\); that is a separate
source-side orientation-transport problem.
\end{proposition}

\begin{proof}
Lemma~\ref{lem:bkm-chamber-s3} identifies
\(\Aut(\Poly_{II})\simeq S_3\) and shows that the three type-I
reflections
\[
s_{f_{-2}-f_2},\qquad s_{f_2-f_3},\qquad s_{f_{-2}-f_3}
\]
act as the three transpositions of the type-II walls.  These
transpositions generate \(S_3\).  The Maass multiplier computation of
Lemma~\ref{lem:appE-maass-on-type-ii}, equivalently the automorphic
packet
\(\texttt{certificates/automorphic/delta5\_maass\_character}\), gives
value \(+1\) on each of the three type-I chamber reflections.  Since
\(\nu_{\Delta_5}\) is a character, it is \(+1\) on every product of
these generators, hence trivial on \(\Aut(\Poly_{II})\).
\end{proof}

The packet
\[
\texttt{certificates/automorphic/delta5\_s3\_chamber\_maass\_triviality}
\]
verifies
\(\texttt{DELTA5\_S3\_CHAMBER\_MAASS\_TRIVIALITY\_VERIFIED}\):
it imports the Maass-character packet, records the three type-I
chamber transpositions and their \(+1\) values, checks the
chamber-automorphism relation row, and excludes any conclusion about
the \(S_3\)-extension of the compact orientation character.

\begin{proposition}[Semidirect extension decision for the orientation
character]
\label{prop:orientation-character-s3-extension-decision}
\textnormal{[abstract character criterion proved; source semidirect
Hall lift not supplied]}
Let \(W=W^{(2)}(\La^{2,1}_{II})\) and
\(\Sigma=\Aut(\Poly_{II})\simeq S_3\).  The group \(\Sigma\) acts on
the simple reflections of \(W\) by permutation.  Suppose the
source-side orientation character \(\epsilon_o:W\to\{\pm1\}\) has been
constructed and satisfies
\[
\epsilon_o(s_{\delta_1})=\epsilon_o(s_{\delta_2})
=\epsilon_o(s_{\delta_3})=-1.
\]
Then \(\epsilon_o\) is \(\Sigma\)-invariant and therefore extends as
an abstract character to \(W\rtimes\Sigma\).  The extensions are
\[
\widetilde\epsilon_{\chi}(w,\sigma)=\epsilon_o(w)\chi(\sigma),
\qquad
\chi\in\Hom(\Sigma,\{\pm1\})=\{1,\operatorname{sgn}\}.
\]
The unique extension whose restriction to \(\Sigma\) equals the Maass
character of row \(323\) is the trivial-\(\Sigma\) extension
\[
\widetilde\epsilon_{\mathrm{Maass}}(w,\sigma)=\epsilon_o(w).
\]

This is only the abstract character decision.  A compact-source
extension of the orientation character requires semidirect Hall lifts
of the chamber-permuting type-I reflections, their action on the
orientation line, determinant square-root compatibility, preservation
of the quotient-orientation null-trivialisations, and the semidirect
relations.  The present finite packets do not supply those rows.
\end{proposition}

\begin{proof}
The chamber action of Lemma~\ref{lem:bkm-chamber-s3} sends
\(s_{\delta_i}\) to \(s_{\delta_{\sigma(i)}}\).  Hence, under the
displayed hypothesis,
\[
\epsilon_o(\sigma s_{\delta_i}\sigma^{-1})
=\epsilon_o(s_{\delta_{\sigma(i)}})
=-1
=\epsilon_o(s_{\delta_i})
\]
for every generator \(s_{\delta_i}\) and every
\(\sigma\in\Sigma\).  Since the type-II Coxeter presentation has only
the involutive relations \(s_{\delta_i}^2=1\) and no braid relations,
this equality on the generators gives \(\Sigma\)-invariance of
\(\epsilon_o\).  A character on \(W\) extends to the semidirect product
precisely after this invariance condition holds; any two extensions
differ by a character of \(\Sigma\).  Since
\(\Hom(S_3,\{\pm1\})=\{1,\operatorname{sgn}\}\), there are exactly two
abstract extensions.  Proposition~\ref{prop:s3-chamber-maass-triviality}
shows that \(\nu_{\Delta_5}\) is trivial on \(\Sigma\), so the
Maass-compatible abstract extension is the one with
\(\chi=1\).
\end{proof}

The packet
\[
\texttt{certificates/orientation/rhomred\_weyl\_orientation\_s3\_extension\_decision}
\]
verifies
\(\texttt{RHOMRED\_WEYL\_ORIENTATION\_S3\_EXTENSION\_DECISION\_OBSTRUCTION\_VERIFIED}\):
it imports the determinant-comparison criterion and the \(S_3\) Maass
triviality packet; it records the two abstract semidirect character
extensions and the Maass-compatible choice; and it keeps the
source-side semidirect Hall lifts, orientation-line transport,
square-root compatibility, quotient-orientation transport, semidirect
relations, and compact-source extension rows open.

\begin{proposition}[Semidirect Hall-lift construction criterion]
\label{prop:semidirect-hall-lift-construction-criterion}
\textnormal{[criterion proved; semidirect Hall lifts not supplied]}
Let \(W=W^{(2)}(\La^{2,1}_{II})\) and
\(\Sigma=\Aut(\Poly_{II})\simeq S_3\).  A compact-source lift of the
abstract extension of Proposition~\ref{prop:orientation-character-s3-extension-decision}
is a finite-stage action of \(W\rtimes\Sigma\) on the retained compact
Hall correspondence package.  It consists of:
\begin{enumerate}
\item the populated type-II partial action groupoid for \(W\), with
orientation-line arrows \(\tau_{i,R,S}\) and zero projective Coxeter
defects;
\item for the three chamber transpositions \(\sigma\in\Sigma\), Hall
correspondences
\[
\mathfrak H_{\sigma,R,S}:
\mathcal M_{R,S}\longleftarrow
\mathcal Z_{\sigma,R,S}\longrightarrow
\mathcal M_{R,\sigma S}
\]
over the retained quotient stacks, together with source and target
maps on the Hall product, Hall coproduct, unit, counit, primitive
subspace, pairing, radical, and PBW filtration rows;
\item Picard-groupoid isomorphisms
\[
\theta^o_{\sigma,R,S}:o_{R,S}\xrightarrow{\sim}
\sigma^*o_{R,\sigma S}
\]
whose squares are the determinant-line transports;
\item transport of
\(\alpha^{\mathrm{red}}\), \(\alpha^{E,\mathrm{free}}\),
\(\widetilde\beta^H\), \(\beta^H\), and \(\lambda^H=0\) through the
same chamber correspondences;
\item zero-defect rows for the \(S_3\) relations and for the
semidirect conjugation relations
\[
\theta^o_{\sigma,R,\delta_i}\,\tau_{i,R,S}\,
(\theta^o_{\sigma,R,\delta_i})^{-1}
=\tau_{\sigma(i),R,\sigma S};
\]
\item transition compatibility in \(R\) and \(R^1\varprojlim=0\) for
the semidirect Hall and Picard-groupoid data.
\end{enumerate}
If these rows are supplied, the trivial-\(S_3\) abstract extension of
row \(324\) is represented by compact-source orientation transport on
\(W\rtimes S_3\).  The present finite packets do not supply these
rows: the compact Hall source packet is empty blocked, the type-II
partial action groupoid has no object or arrow rows, the \(\tau_i\)
packet has no orientation-line transport rows, and the transition
orientation packet has no Picard-groupoid transition rows.
\end{proposition}

\begin{proof}
The listed data are exactly the data of an action of the semidirect
product in the oriented Hall correspondence 2-category.  Items \(1\)
and \(2\) give the underlying source correspondences on the retained
Hall stacks.  Item \(3\) lifts the action from line bundles to
Joyce--Upmeier orientations by enforcing the determinant square-root
diagram.  Item \(4\) says that the quotient-orientation structure is
transported with the orientation line.  Item \(5\) is the group law:
the \(S_3\) relations define the chamber part, while the conjugation
relations identify the chamber transport of the type-II Weyl lifts.
Item \(6\) is the inverse-system condition needed to pass from finite
height to the pro-object.  Without these rows there is only an abstract
character extension, not a compact-source Hall action.
\end{proof}

The packet
\[
\texttt{certificates/orientation/rhomred\_weyl\_semidirect\_hall\_lifts}
\]
verifies
\(\texttt{RHOMRED\_WEYL\_SEMIDIRECT\_HALL\_LIFTS\_OBSTRUCTION\_VERIFIED}\):
it imports the row \(324\) extension decision, the compact Hall source
obstruction ledger, the type-II partial-action-groupoid ledger, the
\(\tau_i\)-transport ledger, and the transition-orientation ledger; it
records the semidirect Hall-lift construction criterion; and it keeps
the source Hall correspondences, orientation transports, determinant
square-root rows, quotient-orientation transports, semidirect relation
rows, transition rows, and inverse-limit rows open.

\begin{proposition}[Imaginary divisor monodromy is not Weyl monodromy]
\label{prop:imaginary-divisor-not-weyl-monodromy}
\textnormal{[firewall proved; local divisor monodromy not yet defined]}
The orientation character
\[
\epsilon_o:W^{(2)}(\La^{2,1}_{II})\longrightarrow\{\pm1\}
\]
is a real-root reflection character.  Its domain is generated by the
three real simple reflections
\[
s_{\delta_1},\qquad s_{\delta_2},\qquad s_{\delta_3}.
\]
An imaginary simple root of the Borcherds--Kac--Moody denominator
algebra contributes a multiplicity in the denominator formula; it does
not define an element of \(W^{(2)}(\La^{2,1}_{II})\).  Hence monodromy
around a divisor component indexed by an imaginary root, if such
monodromy is constructed in the compact Hall theory, is not a Weyl
character value and is not a value of \(\epsilon_o\).

Thus no imaginary-divisor monodromy is being called Weyl monodromy in
the orientation-character construction.  Row \(327\) defines local
divisor monodromy as a separate datum; row \(328\) proves that this
separate datum is irrelevant to the Maass character
\(\nu_{\Delta_5}\) on \(W^{(2)}(\La^{2,1}_{II})\).
\end{proposition}

\begin{proof}
Definition~\ref{def:bkm-real-roots} defines the real-root system as
the \(W^{(2)}(\La^{2,1}_{II})\)-orbit of
\(\{\delta_1,\delta_2,\delta_3\}\), and the type-II Weyl group is
generated by the associated real simple reflections.  The imaginary
simple roots are the Borcherds correction recorded in
Table~\ref{tab:bkm-real-root-comparison}: they are absent from the
rank-three Kac--Moody real-root datum and enter \(\fg_{\Delta_5}\) by
the multiplicities \(m(a)\) and \(\tau(a)\) in the denominator
formula.  A Weyl reflection is attached to a real root of norm \(2\);
an imaginary simple root contributes a denominator factor and a root
space, not a reflection generator.  Therefore the notation
\(\epsilon_o(w)\) is meaningful only for
\(w\in W^{(2)}(\La^{2,1}_{II})\), and no loop around an imaginary
divisor component can be read as such a value without adding a new
definition.  That new definition is isolated in row \(327\).
\end{proof}

The packet
\[
\texttt{certificates/orientation/rhomred\_weyl\_imaginary\_divisor\_firewall}
\]
verifies
\(\texttt{RHOMRED\_WEYL\_IMAGINARY\_DIVISOR\_FIREWALL\_VERIFIED}\):
it imports the real-root definition, the BKM real/imaginary comparison,
the Chapter~6 imaginary-root firewall, and the row \(325\) lift
criterion; it records that real-root reflections are the only Weyl
generators; and it keeps local divisor monodromy outside the Weyl
orientation character until row \(327\) defines it.

\begin{definition}[Local divisor monodromy datum]
\label{def:local-divisor-monodromy-datum}
\textnormal{[definition specified; local divisor rows not supplied]}
Fix a finite height \(R\) and a retained compact Hall stratum
\(\mathcal M_{R,S}\).  A local divisor monodromy datum with label
\(\eta\) consists of:
\begin{enumerate}
\item a retained effective Cartier divisor or normal-crossing
component
\[
\mathcal D_{R,S,\eta}\subset \mathcal M_{R,S}
\]
together with the row identifying the source of the label \(\eta\);
if \(\eta\) is an imaginary Borcherds label, the label is only a
divisor or multiplicity label and not a Weyl-reflection generator;
\item an analytic or formal normal tube \(N_{R,S,\eta}\) of
\(\mathcal D_{R,S,\eta}\) in \(\mathcal M_{R,S}\), its punctured tube
\[
N^\times_{R,S,\eta}:=
N_{R,S,\eta}\setminus \mathcal D_{R,S,\eta},
\]
and a base point \(x_{R,S,\eta}\in N^\times_{R,S,\eta}\);
\item a positively oriented meridian class
\[
\gamma_{R,S,\eta}\in
\pi_1(N^\times_{R,S,\eta},x_{R,S,\eta});
\]
\item the oriented reduced coefficient system
\[
K_R:=\Phi^{\mathrm{red}}_R\otimes o_R
\]
restricted to \(N^\times_{R,S,\eta}\), and, when the Pfaffian line is
separately supplied, the corresponding rank-one Pfaffian-line factor;
\item the monodromy representation of this local system and the
operator
\[
T_{R,S,\eta}:=
\rho_{K_R}(\gamma_{R,S,\eta})
\in \Aut\bigl((K_R)_{x_{R,S,\eta}}\bigr).
\]
When a rank-one Pfaffian-line summand is supplied and preserved by
the meridian action, its sign is denoted
\(\mu_{R,S,\eta}\in\{\pm1\}\);
\item for every successor \(R'\le R\), transition rows identifying
the divisor components, normal tubes, meridians, coefficient systems,
and monodromy operators.
\end{enumerate}
The finite-stage local divisor monodromy is the partially defined
assignment
\[
(R,S,\eta)\longmapsto T_{R,S,\eta}
\]
on the supplied rows.  Its domain is a set of local compact Hall
divisor data, not \(W^{(2)}(\Lambda^{2,1}_{II})\).  It is not a value
of \(\epsilon_o\), not a value of the Maass character
\(\nu_{\Delta_5}\), and not determined by the Borcherds
multiplicity attached to \(\eta\).  Without the divisor component,
normal tube, meridian, oriented coefficient system, and transition
rows above, no local divisor monodromy value is defined.
\end{definition}

The packet
\[
\texttt{certificates/orientation/rhomred\_local\_divisor\_monodromy\_definition}
\]
verifies
\(\texttt{RHOMRED\_LOCAL\_DIVISOR\_MONODROMY\_DEFINITION\_OBSTRUCTION\_VERIFIED}\):
it imports the row \(326\) firewall, the compact Hall obstruction
ledger, the reduced-orientation obstruction ledger, the reduced
vanishing-cycle transition ledger, and the finite Pfaffian obstruction
ledger; it records Definition~\ref{def:local-divisor-monodromy-datum};
and it keeps the divisor-component, normal-tube, meridian,
coefficient-system, Pfaffian-line, monodromy-operator, and transition
rows open.

\begin{proposition}[Local divisor monodromy is irrelevant to the
Maass character]
\label{prop:local-divisor-monodromy-maass-irrelevance}
\textnormal{[proved as a domain-separation theorem]}
Let \(\nu_{\Delta_5}\) be the Maass character of the automorphic
Igusa square root, with values on the type-II Weyl generators and on
the \(S_3\) chamber automorphisms as recorded by the automorphic
Maass-character packet.  Let
\((R,S,\eta)\mapsto T_{R,S,\eta}\) be any supplied local divisor
monodromy datum in the sense of
Definition~\ref{def:local-divisor-monodromy-datum}.  Then the values
of \(\nu_{\Delta_5}\) on
\[
W^{(2)}(\Lambda^{2,1}_{II})\quad\text{and}\quad
\Aut(\Poly_{II})\simeq S_3
\]
are independent of the operators \(T_{R,S,\eta}\).

More precisely, there is no expression
\[
\nu_{\Delta_5}(\gamma_{R,S,\eta})
\]
in the present data.  Such an expression would require an additional
comparison map from the local meridian groupoid of the compact Hall
divisor rows to the automorphic or Weyl group on which
\(\nu_{\Delta_5}\) is defined.  No such comparison map is part of
Definition~\ref{def:local-divisor-monodromy-datum}, and no such row is
supplied by the current packets.
\end{proposition}

\begin{proof}
The Maass character is automorphic target data.  Its recorded inputs
are the type-II real-root reflections \(s_{\delta_i}\) and the
chamber automorphisms in \(\Aut(\Poly_{II})\); the packet
\(\texttt{certificates/automorphic/delta5\_maass\_character}\) records
these values and contains no compact Hall source rows.  By
Definition~\ref{def:local-divisor-monodromy-datum}, local divisor
monodromy is instead a meridian action on the restricted oriented
coefficient system \(K_R=\Phi_R^{\mathrm{red}}\otimes o_R\) over a
punctured normal tube of a retained compact Hall divisor component.
Its input is the meridian class \(\gamma_{R,S,\eta}\), not an element
of \(W^{(2)}(\Lambda^{2,1}_{II})\) and not a chamber automorphism.

Therefore the Maass-character computation has no argument on which a
local divisor monodromy value can be evaluated.  Changing or supplying
operators \(T_{R,S,\eta}\) while keeping the automorphic Maass packet
fixed changes only source-side local-system data; it does not change
the automorphic multiplier rows defining \(\nu_{\Delta_5}\).  A
comparison would require a new meridian-to-automorphic map and a proof
that the Maass character pulls back along it.  Those rows are absent.
\end{proof}

The packet
\[
\texttt{certificates/orientation/rhomred\_local\_divisor\_monodromy\_maass\_irrelevance}
\]
verifies
\(\texttt{RHOMRED\_LOCAL\_DIVISOR\_MONODROMY\_MAASS\_IRRELEVANCE\_VERIFIED}\):
it imports Definition~\ref{def:local-divisor-monodromy-datum}, the row
\(322\) Maass-comparison criterion, and the automorphic
Maass-character packet; it records the domain-separation theorem; and
it keeps any meridian-to-automorphic comparison map, local monodromy
operator rows, and row \(329\) simple-wall orientation computations
separate.

\begin{proposition}[Orientation on the three simple wall charts]
\label{prop:three-simple-wall-chart-orientation}
\textnormal{[chart-level computation criterion proved; finite chart
rows not supplied]}
For \(i=1,2,3\), let a retained \((O2_{\mathrm{atlas}})\) chart over
the simple type-II wall \(\Hpl_{\delta_i}\) supply the data
\[
(U_{i,R,S},T_{i,R,S},u_i,\mathscr L^{\mathrm{tan}}_{i,R,S},
\upsilon_{i,R,S},o_{i,R,S})
\]
with
\[
U_{i,R,S}=T_{i,R,S}\times(-\varepsilon,\varepsilon),\qquad
s_{\delta_i}(t,u_i)=(t,-u_i),
\]
rank-one normal Pfaffian block
\[
K^{\mathrm{cross}}_{i,R,S}=[\mathbb R\xrightarrow{u_i}\mathbb R],
\qquad
\operatorname{pf}_{i,R,S}^{\mathrm{loc}}=\upsilon_{i,R,S}u_i,
\]
divisor order \(1\), and invariant tangent Pfaffian unit
\[
s_{\delta_i}(\upsilon_{i,R,S})=\upsilon_{i,R,S}.
\]
Then the source orientation sign on each of the three supplied simple
wall charts is
\[
\omega^{\mathrm{chart}}_{i,R,S}
:=\frac{s_{\delta_i}^*\operatorname{pf}_{i,R,S}^{\mathrm{loc}}}
        {\operatorname{pf}_{i,R,S}^{\mathrm{loc}}}
=-1,\qquad i=1,2,3.
\]
Equivalently, the three-chart orientation vector is
\[
(\omega^{\mathrm{chart}}_{1,R,S},
 \omega^{\mathrm{chart}}_{2,R,S},
 \omega^{\mathrm{chart}}_{3,R,S})
=(-1,-1,-1)
\]
on the supplied rows.

The present finite packets do not supply the three rows to which this
calculation applies.  The \((O2_{\mathrm{atlas}})\) packet has no
retained wall objects and no wall-chart rows; the finite Pfaffian
packet has no type-II wall chart, Pfaffian line, invariant unit, or
local sign row; and the row \(320\) packet contains only the abstract
rank-one sign criterion, not populated chart-orientation rows.
\end{proposition}

\begin{proof}
The computation is componentwise.  On the \(i\)-th supplied rank-one
chart the reflection fixes the tangent factor and sends the normal
coordinate to its negative.  Hence
\[
s_{\delta_i}^*(\upsilon_{i,R,S}u_i)
=\upsilon_{i,R,S}(-u_i)
=-\upsilon_{i,R,S}u_i.
\]
The quotient of the transported local Pfaffian section by the original
section is therefore \(-1\).  The same calculation holds for
\(i=1,2,3\).  The assertion is a source chart-orientation computation:
it uses the normal coordinate, invariant tangent unit, divisor order,
and retained quotient orientation on the chart.  By
Proposition~\ref{prop:local-divisor-monodromy-maass-irrelevance}, no
Maass-character value or local divisor monodromy datum supplies these
rows.
\end{proof}

The packet
\[
\texttt{certificates/orientation/rhomred\_three\_simple\_wall\_chart\_orientation}
\]
verifies
\(\texttt{RHOMRED\_THREE\_SIMPLE\_WALL\_CHART\_ORIENTATION\_OBSTRUCTION\_VERIFIED}\):
it imports the row \(320\) simple-wall sign criterion, the
\((O2_{\mathrm{atlas}})\) wall-atlas ledger, the finite Pfaffian
ledger, and the row \(328\) Maass-irrelevance firewall; it records the
three-chart orientation computation criterion for row \(329\); and it
keeps the three retained chart rows, invariant-unit rows, quotient
orientation rows, local sign rows, and transition rows open.

\begin{proposition}[Orientation on the isotropic \(a_{ij}\)-charts]
\label{prop:isotropic-aij-chart-orientation}
\textnormal{[chart-level parity criterion proved; source charts not
supplied]}
For \(1\le i<j\le3\), let \(a_{ij}:=\delta_i+\delta_j\) be the
primitive isotropic boundary ray of the type-II chamber.  Suppose a
retained compact Hall chart \(\mathcal U_{ij,R,S}\) in degree
\(a_{ij}\) supplies:
\begin{enumerate}
\item a source decomposition of the oriented reduced coefficient block
into \(10\) even and \(0\) odd one-dimensional summands,
compatible with the target \(a_{ij}:10|0\) row;
\item the quotient-orientation line and its determinant-square row on
\(\mathcal U_{ij,R,S}\);
\item identification of the single affine \(A_1^{(1)}\) null-root
direction and the nine Gritsenko--Nikulin isotropic directions
\(u_{ij,1},\ldots,u_{ij,9}\);
\item transition compatibility for \(R'\le R\).
\end{enumerate}
Then the local isotropic chart orientation parity is even and the
chart sign is
\[
\omega^{\mathrm{iso}}_{ij,R,S}
=(-1)^{\dim K^{\mathrm{iso}}_{\overline1,ij,R,S}}
=(-1)^0=+1 .
\]
Thus the three supplied isotropic chart rows would give
\[
(\omega^{\mathrm{iso}}_{12,R,S},
 \omega^{\mathrm{iso}}_{13,R,S},
 \omega^{\mathrm{iso}}_{23,R,S})
=(+1,+1,+1).
\]

The present packets do not supply these source chart rows.  The
\(\texttt{type\_ii\_chamber\_isotropic}\) and
\(\texttt{wle3\_target\_parity}\) packets certify the target arithmetic
\(a_{ij}:10|0\), the decomposition \(10=1+9\), and the three
isotropic rays \(a_{12},a_{13},a_{23}\); they do not construct compact
Hall source charts, determinant-square orientation rows, or transition
rows.  Consequently row \(330\) is presently a chart-orientation
criterion, not an unconditional source orientation computation.
\end{proposition}

\begin{proof}
The target \(W_{\le3}\) parity table records
\[
a_{ij}:10|0,\qquad i<j,
\]
and the chamber-isotropic packet records the decomposition of each
primitive ray into the affine null-root contribution \(1\) and the
Gritsenko--Nikulin isotropic contribution \(9\).  If a source compact
Hall chart realizes these ten directions as even oriented summands and
supplies no odd summand, the parity exponent of the local determinant
orientation is the odd rank.  This rank is \(0\).  Hence the sign is
\((-1)^0=+1\) for each of the three isotropic charts.  The computation
uses source chart and orientation rows; the target parity table alone
does not provide them.
\end{proof}

The packet
\[
\texttt{certificates/orientation/rhomred\_isotropic\_aij\_chart\_orientation}
\]
verifies
\(\texttt{RHOMRED\_ISOTROPIC\_AIJ\_CHART\_ORIENTATION\_OBSTRUCTION\_VERIFIED}\):
it imports the type-II chamber-isotropic packet, the \(W_{\le3}\)
target-parity packet, the row \(329\) simple-wall chart-orientation
packet, and the compact Hall source obstruction ledger; it records the
isotropic chart-orientation criterion for row \(330\); and it keeps
the compact source chart, quotient-orientation, determinant-square,
isotropic-source parity, and transition rows open.

\begin{proposition}[Orientation on the \(\delta_{123}\)-chart]
\label{prop:delta123-chart-orientation}
\textnormal{[chart-level parity criterion proved; source chart not
supplied]}
Let
\[
\delta_{123}:=\delta_1+\delta_2+\delta_3.
\]
Suppose a retained compact Hall chart \(\mathcal U_{123,R,S}\) in
degree \(\delta_{123}\) supplies:
\begin{enumerate}
\item a source decomposition of the oriented reduced coefficient block
into \(29\) even and \(93\) odd one-dimensional summands, compatible
with the target \(\delta_{123}:29|93\) row;
\item source rows realizing the two real--real--real directions, the
\(27\) mixed real--isotropic directions, and the \(93\) odd
imaginary-simple directions of
Proposition~\ref{prop:bkm-delta123-presentation-count};
\item the quotient-orientation line and its determinant-square row on
\(\mathcal U_{123,R,S}\);
\item transition compatibility for \(R'\le R\).
\end{enumerate}
Then the local first-timelike chart orientation parity is odd and the
chart sign is
\[
\omega^{\mathrm{time}}_{123,R,S}
=(-1)^{\dim K^{\mathrm{time}}_{\overline1,123,R,S}}
=(-1)^{93}=-1 .
\]

The present packets do not supply the source chart to which this
calculation applies.  The
\(\texttt{delta123\_presentation\_split}\) and
\(\texttt{wle3\_target\_parity}\) packets certify the target
presentation count \(29|93\), with \(29=2+27\) and \(29-93=-64\);
they do not construct a compact Hall source chart, a source parity
decomposition, a quotient orientation, a determinant-square row, or
transition rows.  Consequently row \(331\) is presently a
chart-orientation criterion, not an unconditional source orientation
computation.
\end{proposition}

\begin{proof}
Proposition~\ref{prop:bkm-delta123-presentation-count} and the finite
target packet
\(\texttt{certificates/targets/delta5\_gn\_kac/delta123\_presentation\_split}\)
compute the first timelike target presentation:
\[
\rootmult_{\overline0}(\delta_{123})=29,\qquad
\rootmult_{\overline1}(\delta_{123})=93,\qquad
29-93=-64.
\]
The even part is the direct presentation count \(2+27\): two
real--real--real directions after the Jacobi relation and twenty-seven
mixed real--isotropic brackets.  The odd part consists of the
ninety-three negative-norm imaginary-simple directions.  If a compact
Hall chart realizes this decomposition as the source reduced
coefficient block and supplies the determinant-square orientation row,
the parity exponent of the local determinant orientation is the odd
rank.  That rank is \(93\).  Hence the chart sign is
\((-1)^{93}=-1\).  The target presentation count alone is not a source
orientation row.
\end{proof}

The packet
\[
\texttt{certificates/orientation/rhomred\_delta123\_chart\_orientation}
\]
verifies
\(\texttt{RHOMRED\_DELTA123\_CHART\_ORIENTATION\_OBSTRUCTION\_VERIFIED}\):
it imports the target \(\delta_{123}\) presentation-split packet, the
\(W_{\le3}\) target-parity packet, the compact Hall source obstruction
ledger, the finite Pfaffian obstruction ledger, and the row \(330\)
isotropic chart-orientation firewall; it records the
\(\delta_{123}\) chart-orientation criterion for row \(331\); and it
keeps the compact source chart, source parity, source presentation
decomposition, quotient-orientation, determinant-square, and transition
rows open.

\begin{proposition}[Orientation-line compatibility with the \(29|93\)
parity split]
\label{prop:orientation-line-delta123-parity-compatibility}
\textnormal{[line-level compatibility criterion proved; source parity
and orientation line not supplied]}
Let \(\mathcal U_{123,R,S}\) be a retained compact Hall chart in degree
\(\delta_{123}\), and let
\[
K^{\mathrm{red}}_{123,R,S}
\]
be the corresponding reduced coefficient block.  Suppose the following
source data are supplied:
\begin{enumerate}
\item a parity involution \(J_{123,R,S}\) on
\(K^{\mathrm{red}}_{123,R,S}\), with eigenspace decomposition
\[
K^{\mathrm{red}}_{123,R,S}
=K_{\overline0,123,R,S}\oplus K_{\overline1,123,R,S},
\qquad
\dim K_{\overline0,123,R,S}=29,\quad
\dim K_{\overline1,123,R,S}=93;
\]
\item source comparison rows identifying the even summand with the two
real--real--real directions and the \(27\) mixed real--isotropic
directions, and identifying the odd summand with the \(93\) odd
imaginary-simple directions of the target
\(\delta_{123}\)-presentation;
\item a Joyce--Upmeier orientation line \(o_{123,R,S}\) and a commuting
determinant-square diagram
\[
o_{123,R,S}^{\otimes2}
\xrightarrow{\ \sim\ }
\det K^{\mathrm{red}}_{123,R,S}
\xrightarrow{\ \sim\ }
\det K_{\overline0,123,R,S}\otimes
\bigl(\det K_{\overline1,123,R,S}\bigr)^{-1};
\]
\item quotient-orientation null-trivialisations on the chart, compatible
with the parity splitting;
\item transition compatibility: for every \(R'\le R\), the transition
map commutes with \(J_{123}\), preserves the determinant-square diagram
above, and has zero Picard-groupoid, parity-commutator, off-diagonal,
composition, and \(R^1\varprojlim\) defects.
\end{enumerate}
Then the orientation line is compatible with the \(29|93\) parity split.
In this case the determinant parity exponent is the source odd rank
\(93\), so the local orientation sign on the first timelike chart is
\[
(-1)^{\dim K_{\overline1,123,R,S}}=(-1)^{93}=-1.
\]

The present packets do not supply these hypotheses.  The target
\(\delta_{123}:29|93\) presentation is verified, and
Proposition~\ref{prop:delta123-chart-orientation} records the local
chart criterion, but the compact Hall source packet has no parity block
rows, the reduced-orientation packet has no orientation-line
square-root rows, the transition-parity packet has no parity involution
or block-diagonal transition rows, and the transition-orientation packet
has no Picard-groupoid transport rows.  Hence row \(332\) is presently
a compatibility criterion and obstruction ledger, not a proof that the
orientation line is compatible with the \(29|93\) split.
\end{proposition}

\begin{proof}
The determinant functor on a \(\mathbb Z/2\)-graded finite block gives
the Berezinian factorisation
\[
\det K^{\mathrm{red}}_{123,R,S}
\simeq
\det K_{\overline0,123,R,S}\otimes
\bigl(\det K_{\overline1,123,R,S}\bigr)^{-1}.
\]
A Joyce--Upmeier orientation on the chart is a square root of this
determinant line in the Picard groupoid, not a scalar dimension.  If the
displayed square diagram commutes, the chosen orientation line is an
orientation of the parity-split determinant line itself.  The target
presentation count \(29|93\) identifies the required ranks; the source
comparison rows in item \(2\) assert that these target ranks are the
actual source parity ranks on \(\mathcal U_{123,R,S}\).  The local sign
is then the parity of the odd determinant factor, namely
\((-1)^{93}=-1\).  The transition rows in item \(5\) make the statement
independent of the finite HN height and prevent a parity split at one
stage from being read as a pro-object orientation.
\end{proof}

The packet
\[
\texttt{certificates/orientation/rhomred\_delta123\_parity\_compatibility}
\]
verifies
\(\texttt{RHOMRED\_DELTA123\_PARITY\_COMPATIBILITY\_OBSTRUCTION\_VERIFIED}\):
it imports the row \(331\) \(\delta_{123}\)-chart criterion, the
reduced determinant-line packet, the square-root obstruction ledger, the
compact Hall source obstruction ledger, the transition-parity
obstruction ledger, and the transition-orientation obstruction ledger;
it records the line-level compatibility criterion for row \(332\); and
it keeps the source parity block, parity involution,
source-to-target comparison, orientation square-root,
determinant-square, quotient-orientation, parity-transition, and
Picard-groupoid transition rows open.

\begin{proposition}[Orientation-line compatibility with negative roots]
\label{prop:orientation-line-negative-root-compatibility}
\textnormal{[negative-root compatibility criterion proved; source
pairing and orientation transport not supplied]}
Let \(\gamma\) be a retained positive degree whose target row carries a
formal negative-dual block, and let \(-\gamma\) denote the corresponding
negative root degree.  Suppose the compact Hall source supplies:
\begin{enumerate}
\item retained source charts \(\mathcal U_{\gamma,R,S}\) and
\(\mathcal U_{-\gamma,R,S}\) and a compact-source duality morphism
\[
\mathbb D_{\gamma,R,S}:\mathcal U_{\gamma,R,S}\longrightarrow
\mathcal U_{-\gamma,R,S}
\]
lifting the retained finite-window dual closure;
\item parity decompositions of
\[
K^{\mathrm{red}}_{\gamma,R,S},\qquad
K^{\mathrm{red}}_{-\gamma,R,S}
\]
with equal even ranks and equal odd ranks, identified by the
positive--negative target pairing block for \(\gamma\);
\item a nondegenerate source Hopf-pairing matrix
\[
G_{\gamma,R,S}:
K^{\mathrm{red}}_{\gamma,R,S}\otimes
K^{\mathrm{red}}_{-\gamma,R,S}\longrightarrow \mathbb C
\]
whose off-parity blocks vanish and whose radical is the recorded
source radical;
\item orientation square roots \(o_{\gamma,R,S}\) and
\(o_{-\gamma,R,S}\), with determinant-square diagrams commuting with
the duality isomorphism
\[
\det K^{\mathrm{red}}_{-\gamma,R,S}
\simeq
\bigl(\det K^{\mathrm{red}}_{\gamma,R,S}\bigr)^{-1}
\otimes \varepsilon^{\mathrm{CY}_3}_{\gamma,R,S},
\]
where \(\varepsilon^{\mathrm{CY}_3}_{\gamma,R,S}\) is the supplied
Calabi--Yau threefold Serre-pairing sign line;
\item quotient-orientation null-trivialisations and transition rows
preserving the duality morphism, the pairing matrix, the parity
decompositions, and the two orientation square roots.
\end{enumerate}
Then the reduced orientation line is compatible with the negative root
\(-\gamma\): the negative-root orientation is the Serre-dual
Picard-groupoid transport of the positive-root orientation, and the
positive--negative pairing is even and nondegenerate modulo the recorded
radical.

The present packets do not supply these hypotheses.  The retained
dual-closure packet proves only that the finite HN window is closed
under duals.  The target presentation packet records formal
positive--negative blocks, but every such row has
\(\texttt{source\_pairing=false}\).  The compact Hall source packet has
no \(G\)-entries, Hopf-pairing identities, radical rows, source
negative-root parity blocks, or Chevalley anti-involution rows; the
orientation packets have no square-root row, no quotient-orientation
transport row, and no Picard-groupoid transition row.  Hence row \(333\)
is presently a negative-root compatibility criterion and obstruction
ledger, not a proof of negative-root orientation compatibility.
\end{proposition}

\begin{proof}
The proof is formal once the listed source data are supplied.  The
pairing matrix \(G_{\gamma,R,S}\) identifies the negative source block
with the dual of the positive source block after quotienting by the
recorded radical.  Since its off-parity blocks vanish, the
identification respects the even and odd determinant factors.  The
displayed determinant-duality isomorphism then identifies the
determinant line of the negative block with the dual determinant line of
the positive block, up to the supplied Calabi--Yau Serre sign line.
Commutativity of the square-root diagrams lifts this determinant
identity to the Joyce--Upmeier orientation lines.  The quotient
null-trivialisations make the statement one on the reduced quotient
chart, and the transition rows make the positive--negative
identification a statement in the pro-object rather than at a single
finite height.  No target formal pairing block supplies the source
pairing matrix, radical, or orientation-line transport.
\end{proof}

The packet
\[
\texttt{certificates/orientation/rhomred\_negative\_root\_orientation\_compatibility}
\]
verifies
\(\texttt{RHOMRED\_NEGATIVE\_ROOT\_ORIENTATION\_COMPATIBILITY\_OBSTRUCTION\_VERIFIED}\):
it imports the retained dual-closure packet, the target presentation
packet, the formal Hall-pairing pushforward packet, the compact Hall
source obstruction ledger, the row \(332\) parity-compatibility packet,
and the transition-orientation obstruction ledger; it records the
negative-root orientation-compatibility criterion for row \(333\); and
it keeps the compact source negative-root charts, source
positive--negative pairing matrices, Hopf-pairing identities, radical
quotient rows, source parity blocks, orientation square roots,
Serre-duality sign line, quotient-orientation transport, and
Picard-groupoid transition rows open.

\begin{proposition}[Chevalley anti-involution on orientation data]
\label{prop:orientation-chevalley-antiinvolution}
\textnormal{[anti-involution criterion proved; source maps not supplied]}
Let \(P^{\mathrm{red}}_{R,S}\) be the reduced compact Hall primitive
quotient at finite stage \((R,S)\), with retained positive and negative
degree pieces and Joyce--Upmeier orientation lines
\(o_{\gamma,R,S}\).  Suppose the compact Hall source supplies:
\begin{enumerate}
\item source simple representatives \(e_i\), \(f_i\), and \(h_i\),
including the imaginary simple generators and the recorded parity
blocks, with \(e_i\) and \(f_i\) lying in opposite retained degrees and
\(h_i\) in the Cartan part;
\item an even involutive linear map
\[
\vartheta_{R,S}:P^{\mathrm{red}}_{R,S}\longrightarrow
P^{\mathrm{red}}_{R,S}
\]
which sends \(P^{\mathrm{red}}_{\gamma,R,S}\) to
\(P^{\mathrm{red}}_{-\gamma,R,S}\), fixes the Cartan part, sends
\(e_i\leftrightarrow f_i\), and preserves the source parity
decomposition;
\item source bracket and relation rows showing
\[
\vartheta_{R,S}([x,y])=[\vartheta_{R,S}(y),\vartheta_{R,S}(x)]
\]
for homogeneous primitive classes, together with zero-defect Cartan,
Chevalley, real-Serre, Borcherds-orthogonality, and super-sign
relation matrices after quotienting by the recorded radical;
\item a source positive--negative Hopf pairing for which
\(\vartheta_{R,S}\) is adjoint, the radical is stable under
\(\vartheta_{R,S}\), and the quotient pairing is nondegenerate and
parity preserving;
\item orientation-line isomorphisms in the Picard groupoid
\[
\vartheta^o_{\gamma,R,S}:
o_{\gamma,R,S}\xrightarrow{\sim}
o_{-\gamma,R,S}^{\vee}\otimes
\varepsilon^{\mathrm{CY}_3}_{\gamma,R,S}
\]
whose determinant-square diagrams commute with the source duality map,
the Hopf pairing, and the Calabi--Yau threefold Serre sign line;
\item quotient-orientation null-trivialisations and transition rows
preserving \(\vartheta_{R,S}\), \(\vartheta^o_{\gamma,R,S}\), the
radical quotient, the parity decomposition, and the determinant-square
diagrams.
\end{enumerate}
Then \((\vartheta_{R,S},\vartheta^o_{R,S})\) is a Chevalley
anti-involution on the oriented reduced compact Hall primitive package.
It exchanges the positive and negative simple root data, fixes the
Cartan orientation, is anti-multiplicative for the Lie superbracket, is
adjoint for the Hopf pairing on the radical quotient, and transports
Joyce--Upmeier orientations by the Serre-dual Picard-groupoid
isomorphism.

The present packets do not supply these hypotheses.  The target
presentation contains Chevalley and Serre language only at the formal
target level.  The compact Hall source packet has no simple
representative rows, no bracket \(B\)-entries, no Chevalley relation
rows, no \(G\)-entries, no radical rows, and no parity blocks.  The
negative-root compatibility packet records the dual-orientation
criterion but supplies no source anti-involution map.  The
orientation-transition packet has no Picard-groupoid transition map.
Hence row \(334\) is presently an anti-involution criterion and
obstruction ledger, not a construction of a Chevalley anti-involution on
orientation data.
\end{proposition}

\begin{proof}
All assertions follow from the listed source identities.  The linear map
\(\vartheta_{R,S}\) is involutive and exchanges the positive and
negative primitive source pieces by item \(2\).  Item \(3\) is exactly
the anti-automorphism identity for the even Lie-superalgebra
anti-involution, with the Cartan, Chevalley, Serre, Borcherds, and
super-sign relations checked after passage to the radical quotient.
Item \(4\) makes this anti-involution adjoint for the source Hopf
pairing and prevents a target formal pairing block from replacing the
compact source pairing.  Item \(5\) lifts the determinant duality to the
Joyce--Upmeier orientation lines and records the Calabi--Yau Serre sign
line.  Item \(6\) makes the construction compatible with the quotient
null-trivialisations and with finite-stage transitions.  These are
precisely the data of a Chevalley anti-involution on orientation data.
\end{proof}

The packet
\[
\texttt{certificates/orientation/rhomred\_orientation\_chevalley\_antiinvolution}
\]
verifies
\(\texttt{RHOMRED\_ORIENTATION\_CHEVALLEY\_ANTIINVOLUTION\_OBSTRUCTION\_VERIFIED}\):
it imports the row \(333\) negative-root compatibility packet, the
compact Hall source obstruction ledger, the target presentation packet,
the retained dual-closure packet, the formal Hall-pairing pushforward
packet, and the transition-orientation obstruction ledger; it records
the source Chevalley anti-involution criterion for row \(334\); and it
keeps the source simple representatives, positive--negative source
charts, bracket matrices, Chevalley relation rows, Cartan orientation
rows, Hopf-pairing matrices, radical quotient, parity preservation,
orientation square roots, determinant-duality square, Serre sign line,
quotient-orientation transport, Picard-groupoid transition, involutivity
identity, and anti-multiplicativity identity open.

\begin{proposition}[Frobenius property of the orientation pairing]
\label{prop:orientation-frobenius-pairing}
\textnormal{[Frobenius criterion proved; source cyclicity not supplied]}
Let \(P^{\mathrm{red}}_{R,S}\) be the oriented reduced compact Hall
primitive quotient at finite stage \((R,S)\).  Suppose the compact Hall
source supplies:
\begin{enumerate}
\item Hall product and coproduct matrices \(M_{R,S}\) and \(D_{R,S}\),
unit and counit rows, and zero-defect associativity,
coassociativity, bialgebra-compatibility, and primitive-closure rows;
\item a source Hopf pairing
\[
\langle-,-\rangle^o_{R,S}:
P^{\mathrm{red}}_{R,S}\otimes P^{\mathrm{red}}_{R,S}
\longrightarrow \ell_{R,S}
\]
with orientation-line factors included in the target line
\(\ell_{R,S}\), together with a source trace functional whose degree is
recorded in the next row;
\item zero-defect Hopf-adjointness rows
\[
\langle xy,z\rangle^o_{R,S}
=
\langle x\otimes y,\Delta z\rangle^o_{R,S}
\]
with the Koszul signs and Thom--Sebastiani orientation signs recorded
at the source;
\item zero-defect primitive Frobenius-cyclicity rows
\[
\langle [x,y],z\rangle^o_{R,S}
=
\langle x,[y,z]\rangle^o_{R,S}
\]
for homogeneous primitive classes, with the parity signs and the
Chevalley anti-involution of
Proposition~\ref{prop:orientation-chevalley-antiinvolution} acting by
adjunction;
\item radical-kernel, quotient-splitting, radical ideal/coideal, and
quotient nondegeneracy rows, so that the pairing descends to a
nondegenerate form on the reduced radical quotient;
\item Picard-groupoid orientation pairing squares, Calabi--Yau Serre
sign-line rows, radical-transition rows, pairing-kernel
Mittag--Leffler rows, and orientation-transition rows making the
cyclicity identities independent of the finite HN height.
\end{enumerate}
Then \(\langle-,-\rangle^o_{R,S}\) is a Frobenius pairing on the
oriented reduced compact Hall primitive quotient: it is nondegenerate
after quotienting by the recorded radical, it is adjoint for the Hall
product and coproduct, and its restriction to primitives is cyclic for
the Lie superbracket with the recorded orientation signs.

The present packets do not supply these hypotheses.  The formal
Hall-pairing pushforward packet proves only degree compatibility for a
supplied pairing; it explicitly does not prove Hopf adjointness,
Frobenius cyclicity, or quotient nondegeneracy.  The compact Hall source
packet has no \(M\)-, \(D\)-, \(B\)-, \(G\)-, \(K\)-, or \(Q\)-entries,
no unit or counit rows, no Hall-bialgebra identities, no
Hopf-pairing identity rows, and no radical ideal/coideal rows.  The row
\(334\) packet records the Chevalley anti-involution only as a
criterion.  Hence row \(335\) is presently a Frobenius-pairing
criterion and obstruction ledger, not a proof that the orientation
pairing is Frobenius.
\end{proposition}

\begin{proof}
The argument is the standard finite-dimensional Frobenius argument once
the listed source data are present.  The product, coproduct, unit,
counit, and bialgebra identities make the finite-stage Hall package a
bialgebra on the retained source quotient.  Hopf adjointness identifies
left multiplication with the adjoint of the coproduct.  Restricting to
primitive classes turns this adjointness into the displayed cyclic
identity for the Lie superbracket; the Koszul signs are exactly the
source parity signs, and the orientation signs are the
Thom--Sebastiani and Serre signs recorded in the Picard groupoid.
Radical ideal/coideal stability and quotient nondegeneracy descend the
pairing to the reduced radical quotient.  The transition and
Mittag--Leffler rows make the construction independent of the finite
height.  No scalar trace, target pairing block, or formal degree
compatibility row supplies these source identities.
\end{proof}

The packet
\[
\texttt{certificates/orientation/rhomred\_orientation\_frobenius\_pairing}
\]
verifies
\(\texttt{RHOMRED\_ORIENTATION\_FROBENIUS\_PAIRING\_OBSTRUCTION\_VERIFIED}\):
it imports the row \(334\) Chevalley anti-involution criterion, the
compact Hall source obstruction ledger, the formal Hall-pairing
pushforward packet, the transition-radical obstruction ledger, the
pairing-kernel \(R^1\varprojlim\) obstruction ledger, and the
transition-orientation obstruction ledger; it records the Frobenius
pairing criterion for row \(335\); and it keeps the Hall product and
coproduct matrices, bialgebra identities, source Hopf-pairing matrix,
Hopf-adjointness row, Frobenius-cyclicity row, quotient
nondegeneracy row, radical kernel, quotient splitting, radical
ideal/coideal, orientation pairing line, orientation cyclic signs, Serre
sign line, radical transition rows, pairing-kernel Mittag--Leffler rows,
orientation-transition rows, and row \(336\) trace-degree computation
open.

\begin{proposition}[Degree of the Frobenius trace]
\label{prop:orientation-frobenius-trace-degree}
\textnormal{[trace-degree criterion proved; source trace not supplied]}
Let \(P^{\mathrm{red}}_{R,S}\) be the oriented reduced compact Hall
primitive quotient, and suppose the Frobenius pairing of
Proposition~\ref{prop:orientation-frobenius-pairing} has been supplied
on the radical quotient.  Suppose in addition that the compact Hall
source supplies:
\begin{enumerate}
\item a source trace functional
\[
\operatorname{tr}^o_{R,S}:
P^{\mathrm{red}}_{R,S,\mathrm{rad}}\longrightarrow
\ell^{\mathrm{tr}}_{R,S}[-3]
\]
on the reduced radical quotient, with the target trace line
\(\ell^{\mathrm{tr}}_{R,S}\) recorded in the orientation Picard
groupoid;
\item a source grading convention row for cohomological degree,
parity, orientation-line degree, and Calabi--Yau threefold Serre shift,
with the trace-degree identity
\[
\deg(\operatorname{tr}^o_{R,S})=-3
\]
having defect rank zero;
\item zero-defect trace--pairing rows
\[
\langle x,y\rangle^o_{R,S}
=
\operatorname{tr}^o_{R,S}(xy)
\]
for homogeneous classes \(x,y\), including the Koszul and
Thom--Sebastiani orientation signs;
\item zero-defect trace cyclicity rows
\[
\operatorname{tr}^o_{R,S}(xy)
=(-1)^{|x||y|}\operatorname{tr}^o_{R,S}(yx)
\]
and compatibility with the Chevalley anti-involution and the
Calabi--Yau Serre sign line;
\item radical-quotient, quotient-orientation, transition, and
Mittag--Leffler rows preserving the trace functional, the trace degree,
and the trace--pairing identity in the pro-object.
\end{enumerate}
Then the Frobenius trace has degree \(-3\).  Equivalently, the
orientation pairing has cohomological degree \(-3\):
\[
\deg\langle x,y\rangle^o_{R,S}
=\deg x+\deg y-3
\]
for homogeneous classes on the reduced radical quotient.

The present packets do not supply these hypotheses.  The row \(335\)
packet records the Frobenius pairing only as a criterion and explicitly
keeps the source trace functional open.  The compact Hall source packet
has no source trace rows, no product or pairing matrices, no
Hopf-pairing identities, and no radical quotient rows.  The formal
Hall-pairing pushforward packet proves only homogeneous degree-zero
compatibility for supplied pairing rows.  The protected-integration
Gram-degree packet records the scalar target monomial \(s\)-degree
\(m_R\) of \(I_R^{\mathrm{prot}}\), not the cohomological degree of a
source Frobenius trace.  The level-\(Z\) protected trace packet is an
empty obstruction ledger.  Hence row \(336\) is presently a trace-degree
criterion and obstruction ledger, not a proof that the Frobenius trace
has degree \(-3\).
\end{proposition}

\begin{proof}
The statement follows immediately from the supplied source trace data.
The trace-degree row fixes the target shift of
\(\operatorname{tr}^o_{R,S}\) to \([-3]\).  The trace--pairing identity
then identifies the Frobenius pairing with the composite
\((x,y)\mapsto \operatorname{tr}^o_{R,S}(xy)\), so a homogeneous product
of degree \(\deg x+\deg y\) is sent to the trace line in degree
\(\deg x+\deg y-3\).  Trace cyclicity and the Chevalley/Serre
compatibility rows show that this degree is the Frobenius trace degree,
not only a degree of one chosen product expression.  Radical and
transition rows descend the identity to the reduced pro-object.  A
scalar protected-integration monomial degree does not provide the
source trace functional or the cohomological shift.
\end{proof}

The packet
\[
\texttt{certificates/orientation/rhomred\_orientation\_frobenius\_trace\_degree}
\]
verifies
\(\texttt{RHOMRED\_ORIENTATION\_FROBENIUS\_TRACE\_DEGREE\_OBSTRUCTION\_VERIFIED}\):
it imports the row \(335\) Frobenius-pairing criterion, the compact Hall
source obstruction ledger, the formal Hall-pairing pushforward packet,
the protected-integration Gram-degree packet, the level-\(Z\) protected
trace obstruction ledger, and the transition-orientation obstruction
ledger; it records the trace-degree criterion for row \(336\); and it
keeps the source trace functional, trace-degree row, trace--pairing
identity, trace cyclicity, orientation trace line, Calabi--Yau Serre
shift, radical quotient, quotient-orientation transport, transition
compatibility, and row \(337\) Serre-sign computation open.

\begin{proposition}[Calabi--Yau threefold Serre signs]
\label{prop:orientation-cy3-serre-signs}
\textnormal{[Serre-sign criterion proved; source sign rows not supplied]}
Let \(K^{\mathrm{red}}_{\gamma,R,S}\) and
\(K^{\mathrm{red}}_{-\gamma,R,S}\) be the reduced oriented compact
Hall complexes in opposite retained degrees.  Suppose the compact Hall
source supplies:
\begin{enumerate}
\item a reduced Serre-duality isomorphism
\[
S_{\gamma,R,S}:
K^{\mathrm{red}}_{\gamma,R,S}
\xrightarrow{\sim}
\bigl(K^{\mathrm{red}}_{-\gamma,R,S}\bigr)^{\vee}[-3]
\]
on the radical quotient, compatible with the source Ext-pairing and
with the quotient null-trivialisations;
\item source rows fixing cohomological degree, parity, determinant-line
ordering, and the Koszul sign convention for the shift \([-3]\), with
a sign-exponent row
\(\sigma^{\mathrm{CY}_3}_{\gamma,R,S}\) of defect rank zero;
\item a Picard-groupoid sign line
\(\varepsilon^{\mathrm{CY}_3}_{\gamma,R,S}\) and a commuting
determinant-duality square
\[
\det K^{\mathrm{red}}_{-\gamma,R,S}
\simeq
\bigl(\det K^{\mathrm{red}}_{\gamma,R,S}\bigr)^{-1}
\otimes \varepsilon^{\mathrm{CY}_3}_{\gamma,R,S};
\]
\item zero-defect comparison rows identifying the
Thom--Sebastiani sign, the Chevalley anti-involution sign, the
Frobenius-cyclicity sign, and the trace-cyclicity sign with
\(\varepsilon^{\mathrm{CY}_3}_{\gamma,R,S}\);
\item radical-quotient, quotient-orientation, finite-stage transition,
and Mittag--Leffler rows preserving \(S_{\gamma,R,S}\), the sign
exponent, and the determinant-duality square.
\end{enumerate}
Then the orientation signs agree with Calabi--Yau threefold Serre
duality: every sign used in negative-root transport, Chevalley
adjunction, Frobenius cyclicity, and trace cyclicity is the sign
obtained from the displayed Serre-duality isomorphism and the recorded
source grading convention.

The present packets do not supply these hypotheses.  The PTVV packet
imports the ambient \((-1)\)-shifted symplectic form on retained finite
derived substacks, but it explicitly does not construct the Joyce
d-critical truncation, reduced \(\operatorname{RHom}\), determinant
line, orientation square root, or quotient orientation.  The row
\(333\), \(334\), \(335\), and \(336\) packets all keep the
Calabi--Yau Serre sign line or shift open.  The compact Hall source
packet has no source pairing matrices, Hopf-pairing identities, parity
blocks, radical rows, determinant-duality squares, or Serre sign rows.
The formal Hall-pairing pushforward and target negative-dual blocks are
not source sign computations.  Hence row \(337\) is presently a
Serre-sign criterion and obstruction ledger, not a proof that the
orientation signs agree with Calabi--Yau threefold Serre duality.
\end{proposition}

\begin{proof}
Once the listed source rows are present, the assertion is a determinant
calculation.  The reduced Serre-duality isomorphism identifies the
negative complex with the shifted dual of the positive complex.  The
source degree and parity rows determine the Koszul exponent attached to
the shift \([-3]\); the determinant-ordering row transports this
exponent to the determinant line.  The determinant-duality square is
therefore the Picard-groupoid form of the Calabi--Yau threefold Serre
sign line.  The comparison rows force the signs used in
Thom--Sebastiani gluing, Chevalley adjunction, Frobenius cyclicity, and
trace cyclicity to be this same line.  Radical and transition rows make
the equality one on the reduced pro-object.  An ambient shifted
symplectic form or a target pairing block gives no source determinant
ordering and therefore no Serre-sign computation.
\end{proof}

The packet
\[
\texttt{certificates/orientation/rhomred\_orientation\_cy3\_serre\_signs}
\]
verifies
\(\texttt{RHOMRED\_ORIENTATION\_CY3\_SERRE\_SIGNS\_OBSTRUCTION\_VERIFIED}\):
it imports the row \(333\) negative-root compatibility criterion, the
row \(334\) Chevalley anti-involution criterion, the row \(335\)
Frobenius-pairing criterion, the row \(336\) trace-degree criterion,
the PTVV finite-substack shifted-symplectic packet, the compact Hall
source obstruction ledger, and the transition-orientation obstruction
ledger; it records the Calabi--Yau Serre-sign criterion for row \(337\);
and it keeps the reduced Serre-duality isomorphism, source
sign-exponent table, orientation sign line, determinant-duality square,
Thom--Sebastiani sign comparison, Chevalley sign comparison,
Frobenius-cyclicity sign comparison, trace-cyclicity sign comparison,
radical quotient, quotient orientation, and transition compatibility
open.

\begin{proposition}[Reduced Calabi--Yau pairing modulo the radical]
\label{prop:orientation-cy3-pairing-nondegeneracy}
\textnormal{[nondegeneracy criterion proved; quotient matrices not supplied]}
Let \(P^{\mathrm{red}}_{\gamma,R,S}\) and
\(P^{\mathrm{red}}_{-\gamma,R,S}\) be the retained opposite-degree
primitive source spaces at finite stage \((R,S)\).  Suppose the compact
Hall source supplies:
\begin{enumerate}
\item a Calabi--Yau threefold Hopf-pairing matrix
\[
G^{\mathrm{CY}_3}_{\gamma,R,S}:
P^{\mathrm{red}}_{\gamma,R,S}\otimes
P^{\mathrm{red}}_{-\gamma,R,S}\longrightarrow
\ell^{\mathrm{CY}_3}_{\gamma,R,S}
\]
with the source Serre sign line of
Proposition~\ref{prop:orientation-cy3-serre-signs} included in the
target line;
\item left- and right-kernel rows
\[
K^L_{\gamma,R,S}=\ker G^{\mathrm{CY}_3}_{\gamma,R,S},\qquad
K^R_{\gamma,R,S}=\ker (G^{\mathrm{CY}_3}_{\gamma,R,S})^{t}
\]
and zero-defect rows identifying these kernels with the recorded
Hopf-pairing radical in the positive and negative source spaces;
\item quotient splitting matrices
\[
Q_{\gamma,R,S}:P^{\mathrm{red}}_{\gamma,R,S}\twoheadrightarrow
\overline P_{\gamma,R,S},\qquad
Q_{-\gamma,R,S}:P^{\mathrm{red}}_{-\gamma,R,S}\twoheadrightarrow
\overline P_{-\gamma,R,S}
\]
whose kernels are exactly the recorded radicals;
\item a quotient-pairing matrix
\[
\overline G_{\gamma,R,S}
   = Q_{\gamma,R,S}\,G^{\mathrm{CY}_3}_{\gamma,R,S}\,
     Q_{-\gamma,R,S}^{t}
\]
with full left rank and full right rank, equivalently determinant or
Smith defect rank zero after choosing finite quotient bases;
\item radical ideal/coideal rows, transition-radical rows, and
pairing-kernel Mittag--Leffler rows proving that the quotient pairing
and its nondegeneracy are independent of the finite HN height.
\end{enumerate}
Then the reduced Calabi--Yau threefold pairing is nondegenerate modulo
the radical:
\[
\overline G_{\gamma,R,S}:
\overline P_{\gamma,R,S}\otimes
\overline P_{-\gamma,R,S}
\longrightarrow \ell^{\mathrm{CY}_3}_{\gamma,R,S}
\]
has zero left and right kernels, with the orientation sign line fixed
by Calabi--Yau threefold Serre duality.

The present packets do not supply these hypotheses.  The compact Hall
source packet has no \(G\)-entries, kernel matrices, quotient
splittings, quotient-nondegeneracy identities, radical ideal/coideal
rows, or quotient transition rows.  The formal Hall-pairing pushforward
packet proves only homogeneous degree-zero compatibility for three
supplied formal pairings and explicitly does not prove quotient
nondegeneracy.  The row \(333\), \(335\), and \(337\) packets record
negative-root compatibility, Frobenius pairing, and Serre signs only as
criteria.  The transition-radical and pairing-kernel Mittag--Leffler
packets are obstruction ledgers with empty transition and kernel
tables.  Hence row \(338\) is presently a nondegeneracy-modulo-radical
criterion and obstruction ledger, not a proof that the reduced
Calabi--Yau pairing is nondegenerate modulo the radical.
\end{proposition}

\begin{proof}
The assertion is the elementary quotient-pairing argument once the
listed finite matrices exist.  By the kernel rows, the left and right
radicals of \(G^{\mathrm{CY}_3}_{\gamma,R,S}\) are exactly the kernels
of \(Q_{\gamma,R,S}\) and \(Q_{-\gamma,R,S}\).  Therefore
\(G^{\mathrm{CY}_3}_{\gamma,R,S}\) factors through
\(\overline P_{\gamma,R,S}\otimes\overline P_{-\gamma,R,S}\).  The
full-rank row for
\(Q_{\gamma,R,S}G^{\mathrm{CY}_3}_{\gamma,R,S}Q_{-\gamma,R,S}^{t}\)
states precisely that the induced form has no left or right kernel.
The radical ideal/coideal rows make the quotient compatible with the
Hall algebra operations; the transition and Mittag--Leffler rows make
the finite-stage statement stable in the pro-object.  The Serre-sign
line fixes the orientation of the target line but does not by itself
imply full rank.
\end{proof}

The packet
\[
\texttt{certificates/orientation/rhomred\_orientation\_cy3\_pairing\_nondegeneracy}
\]
verifies
\(\texttt{RHOMRED\_ORIENTATION\_CY3\_PAIRING\_NONDEGENERACY\_OBSTRUCTION\_VERIFIED}\):
it imports the row \(333\) negative-root compatibility criterion, the
row \(335\) Frobenius-pairing criterion, the row \(337\) Serre-sign
criterion, the compact Hall source obstruction ledger, the formal
Hall-pairing pushforward packet, the transition-radical obstruction
ledger, and the pairing-kernel Mittag--Leffler obstruction ledger; it
records the reduced Calabi--Yau pairing nondegeneracy criterion for
row \(338\); and it keeps the source pairing matrix, left and right
kernel rows, radical identification, quotient splitting matrices,
quotient-pairing full-rank row, radical ideal/coideal rows, radical
transition rows, pairing-kernel Mittag--Leffler rows, and inverse-limit
nondegeneracy open.

\begin{definition}[Orientation-preserving transition datum]
\label{def:transition-orientation-preservation}
Let \(R'\le R\), and suppose the finite-moduli transition geometry of
Definition~\ref{def:finite-moduli-transition-geometry} is supplied.
An orientation-preserving transition datum consists of:
\begin{enumerate}
\item a Picard-groupoid isomorphism over the retained transition stack
\[
\theta^o_{RR'}:(t_{RR'})^*o_{R'}\xrightarrow{\sim}o_R;
\]
\item commutativity of the determinant square-root diagram
\[
\bigl((t_{RR'})^*o_{R'}\bigr)^{\otimes2}
\longrightarrow (t_{RR'})^*\det K^{\mathrm{red}}_{R'}
\longrightarrow \det K^{\mathrm{red}}_R,
\]
with the square \(o_R^{\otimes2}\simeq\det K^{\mathrm{red}}_R\);
\item transport of the null-trivialisations of
\(\alpha^{\mathrm{red}}\), \(\alpha^{E,\mathrm{free}}\),
\(\widetilde\beta^H\), \(\beta^H\), and
\(\lambda^H=0\) for every retained finite stabilizer stratum;
\item compatibility of \(\theta^o_{RR'}\) with the
Thom--Sebastiani orientation isomorphisms on extension and two-step
flag stacks;
\item compatibility with the Weyl lifts \(\tau_i\), the cochain
killing the projective Coxeter cocycle, and the quotient-cocycle
transport rows of \((O1)^+\);
\item zero Picard-groupoid, square-root, null-trivialization,
Thom--Sebastiani, Weyl, and composition defect ranks, together with
\(R^1\varprojlim=0\) for the orientation Picard data.
\end{enumerate}
Proper or closed stack transitions alone do not define this datum.
Scalar traces, squared determinants, OP scalar branches, Maass-character
values, Pfaffian products, and target root windows are not orientation
transport data.
\end{definition}

\begin{proposition}[Transition preservation of Joyce--Upmeier
orientations; \textit{proved}]
\label{prop:transition-orientation-preservation-criterion}
Given an orientation-preserving transition datum of
Definition~\ref{def:transition-orientation-preservation}, the
transition \(R\to R'\) preserves the reduced Joyce--Upmeier
orientation line, the quotient-orientation null-trivialisations, the
Thom--Sebastiani multiplicative structure, and the \((O1)^+\)
Weyl-equivariant determinant-line transport.
\end{proposition}

\begin{proof}
The isomorphism \(\theta^o_{RR'}\) identifies the pulled-back target
orientation line with the source orientation line in the Picard
groupoid.  The square-root diagram in item \(2\) shows that this
identification is an identification of Joyce--Upmeier orientations, not
only of underlying line bundles.  Item \(3\) transports the quotient
Cech--Borel null-trivialisations and the zero finite-stabilizer
linearization characters.  Item \(4\) makes the Hall
Thom--Sebastiani orientation isomorphisms commute with transition, and
item \(5\) makes the type-II Weyl lifts and Coxeter cochain commute
with transition.  The zero-defect and \(R^1\varprojlim\) rows in item
\(6\) give strict composition and inverse-limit exactness for the
orientation Picard data.
\end{proof}

The packet
\(\texttt{certificates/orientation/transition\_orientation\_preservation}\)
records the present state of this datum.  Its verifier status is
\texttt{TRANSITION\_ORIENTATION\_PRESERVATION\_OBSTRUCTION\_VERIFIED}:
the stack-transition input, orientation-line pullback isomorphisms,
determinant square compatibility, quotient Cech--Borel
null-trivialization transport rows, finite-stabilizer and
linearization transport rows, Thom--Sebastiani transition rows,
Weyl-lift transition rows, Coxeter-cochain transition rows, strict
Picard-groupoid transition rows, and orientation Mittag--Leffler rows
are not yet supplied.  Hence transition preservation of orientations
remains open.

\begin{proposition}[Survival of orientation classes in the inverse limit]
\label{prop:orientation-class-inverse-limit-survival}
\textnormal{[inverse-limit survival criterion proved; orientation tower not supplied]}
Let \(\mathcal R\) be the retained cofinal HN index category, and for
each \(R\in\mathcal R\) let \(o_R\) be a Joyce--Upmeier orientation
line on the reduced compact Hall frame at stage \(R\).  Suppose the
compact orientation source supplies:
\begin{enumerate}
\item finite-stage orientation rows
\[
o_R^{\otimes 2}\simeq \det K^{\mathrm{red}}_R,\qquad
\operatorname{cl}(o_R)=0
\]
with zero square-root and orientation-class defect ranks for every
retained stratum;
\item for every successor \(R'\le R\), Picard-groupoid transition
isomorphisms
\[
\theta^o_{RR'}:(t_{RR'})^*o_{R'}\xrightarrow{\sim}o_R
\]
whose determinant-square, quotient-orientation, Thom--Sebastiani, Weyl,
and Coxeter transition defects are zero;
\item strict identity and composition rows for the transitions
\(\theta^o_{RR'}\), so that the \(o_R\) form an inverse system in the
orientation Picard groupoid;
\item an orientation Mittag--Leffler row with
\[
R^1\varprojlim_R\,\pi_0\operatorname{Pic}^{or}_R=0
\]
and compatible image-stabilization witnesses for the orientation
classes;
\item a limit row identifying the pro-orientation class
\[
o_\infty=\varprojlim_R(o_R,\theta^o_{RR'})
\]
with the class used by the Dirac--Igusa pro-object.
\end{enumerate}
Then the finite-stage orientation classes survive the inverse limit:
the compatible system \((o_R,\theta^o_{RR'})\) determines an
orientation class \(o_\infty\) on the pro-object, and every finite
projection of \(o_\infty\) is the prescribed \(o_R\).

The present packets do not supply these hypotheses.  The reduced
orientation ledger has no orientation-line rows, no square-root rows,
no orientation-class rows, and no transition rows.  The square-root
packet records three missing square-root rows and has no orientation
class zero row.  The transition-orientation packet has no
Picard-groupoid pullback maps, no determinant-square transition rows,
no strict composition rows, and no orientation Mittag--Leffler row.
The row \(338\) nondegeneracy packet keeps inverse-limit
nondegeneracy open and does not construct orientation classes.  Hence
row \(339\) is presently an inverse-limit survival criterion and
obstruction ledger, not a proof that orientation classes survive the
inverse limit.
\end{proposition}

\begin{proof}
The claim is the standard exactness statement for an inverse system in
Picard groupoids once the listed rows are present.  The finite-stage
rows give objects \(o_R\) with the prescribed determinant square and
vanishing orientation obstruction.  The transition rows make these
objects compatible under pullback.  Strict identity and composition
rows replace weak pairwise compatibility by an inverse system.  The
Mittag--Leffler row kills the \(R^1\varprojlim\) obstruction to lifting
the compatible finite classes to a pro-class.  The limit row then
identifies the resulting object with the orientation datum used in the
Dirac--Igusa pro-object.  This proves survival of the prescribed
classes.  It does not prove that every orientation class at the limit
comes from a finite-stage class; that is the separate row \(340\)
obligation.
\end{proof}

The packet
\[
\texttt{certificates/orientation/rhomred\_orientation\_class\_limit\_survival}
\]
verifies
\(\texttt{RHOMRED\_ORIENTATION\_CLASS\_LIMIT\_SURVIVAL\_OBSTRUCTION\_VERIFIED}\):
it imports the reduced-orientation obstruction ledger, the square-root
obstruction packet, the transition-orientation obstruction ledger, and
the row \(338\) nondegeneracy criterion; it records the inverse-limit
survival criterion for row \(339\); and it keeps the finite orientation
line rows, square-root rows, orientation-class zero rows, Picard
transition maps, strict composition rows, orientation
Mittag--Leffler rows, limit class row, and row \(340\) no-new-class
statement open.

\begin{proposition}[No new orientation class at the limit]
\label{prop:orientation-no-new-limit-class}
\textnormal{[no-phantom criterion proved; pro-Picard comparison not supplied]}
Let \(\operatorname{Pic}^{or}_\infty\) be the orientation Picard
groupoid of the Dirac--Igusa pro-object, and let
\(\operatorname{Pic}^{or}_R\) be the finite-stage orientation Picard
groupoid at HN height \(R\).  Suppose the compact orientation source
supplies:
\begin{enumerate}
\item a pro-Picard comparison functor
\[
\Xi:\operatorname{Pic}^{or}_\infty
\longrightarrow
\varprojlim_R\operatorname{Pic}^{or}_R
\]
whose object and automorphism components are defined by the finite
stack transitions;
\item finite-stage effectivity rows showing that every orientation line
on the pro-object restricts to the recorded finite-stage orientation
line, square-root, quotient-orientation, Thom--Sebastiani, Weyl, and
Coxeter data;
\item separatedness rows proving that a pro-orientation class whose
finite restrictions are all trivial is itself trivial;
\item automorphism Mittag--Leffler rows with
\[
R^1\varprojlim_R\,\pi_1\operatorname{Pic}^{or}_R=0
\]
and object-class Mittag--Leffler rows with zero \(R^1\varprojlim\)
defect for the relevant orientation-class tower;
\item a zero-defect comparison row
\[
\pi_0\operatorname{Pic}^{or}_\infty
\xrightarrow{\sim}
\varprojlim_R\pi_0\operatorname{Pic}^{or}_R
\]
on the subfunctor cut out by the determinant square, quotient
orientation, Thom--Sebastiani, Weyl, and Coxeter conditions.
\end{enumerate}
Then no new orientation class appears at the limit: every
Dirac--Igusa pro-orientation class is detected by its compatible
finite-stage restrictions, and the class is the image of the finite
tower under Proposition~\ref{prop:orientation-class-inverse-limit-survival}.

The present packets do not supply these hypotheses.  The row \(339\)
packet records only the survival criterion and has no limit-survival
rows.  The reduced-orientation ledger has no finite orientation-line or
transition rows.  The transition-orientation packet has no
Picard-groupoid transition maps, no strict composition rows, no
orientation Mittag--Leffler rows, and no pro-Picard comparison rows.
The square-root packet records missing finite square roots and no
orientation-class zero rows.  Hence row \(340\) is presently a
no-phantom-class criterion and obstruction ledger, not a proof that no
new orientation class appears at the limit.
\end{proposition}

\begin{proof}
Assume the listed comparison and exactness rows.  The functor \(\Xi\)
restricts a pro-orientation class to its finite-stage orientation
tower.  The effectivity rows identify this tower with the same
finite-stage orientation data used in the Dirac--Igusa construction.
Separatedness says that the kernel of \(\pi_0(\Xi)\) is zero: a class
with trivial finite restrictions is trivial at the pro-level.
Vanishing of \(R^1\varprojlim\) on automorphisms and object classes
removes the usual derived-limit ambiguity in passing between Picard
groupoids and their towers.  The displayed zero-defect comparison row
therefore identifies pro-orientation classes with compatible finite
orientation classes.  By
Proposition~\ref{prop:orientation-class-inverse-limit-survival}, each
compatible finite tower gives the prescribed limit class.  No additional
class can occur at the limit.
\end{proof}

The packet
\[
\texttt{certificates/orientation/rhomred\_orientation\_no\_new\_limit\_class}
\]
verifies
\(\texttt{RHOMRED\_ORIENTATION\_NO\_NEW\_LIMIT\_CLASS\_OBSTRUCTION\_VERIFIED}\):
it imports the row \(339\) inverse-limit survival criterion, the
reduced-orientation obstruction ledger, the square-root obstruction
packet, and the transition-orientation obstruction ledger; it records
the no-new-limit-class criterion for row \(340\); and it keeps the
pro-Picard comparison functor, finite-stage effectivity rows,
separatedness rows, automorphism Mittag--Leffler rows, object-class
Mittag--Leffler rows, zero-defect comparison row, and pro-orientation
class table open.

\begin{definition}[Vanishing-cycle-preserving transition datum]
\label{def:transition-vanishing-cycle-preservation}
Let \(R'\le R\).  Suppose the finite-moduli transition geometry of
Definition~\ref{def:finite-moduli-transition-geometry} and the
orientation-preserving transition datum of
Definition~\ref{def:transition-orientation-preservation} are supplied.
A vanishing-cycle-preserving transition datum consists of:
\begin{enumerate}
\item transition maps of oriented d-critical charts, with transition
graph \(\mathfrak G_{RR'}\) and projections
\[
p_R:\mathfrak G_{RR'}\to\mathfrak Z_R,\qquad
p_{R'}:\mathfrak G_{RR'}\to\mathfrak Z_{R'};
\]
\item potential compatibility on charts:
\[
f_R=p_{R'}^*f_{R'}+q_{RR'}
\]
after pullback to the transition graph, where \(q_{RR'}\) is the
recorded non-degenerate quadratic stabilization and its parity shift is
recorded;
\item pullback compatibility of the K3 semiregularity cosections and
a morphism of reduced obstruction complexes with zero cosection and
reduced-complex defect ranks;
\item a constructible-complex isomorphism
\[
\theta^\Phi_{RR'}:
p_R^*\Phi^{\mathrm{red}}_R
\xrightarrow{\sim}
p_{R'}^*\Phi^{\mathrm{red}}_{R'}
\]
on \(\mathfrak G_{RR'}\), with BBDJS perverse normalization and
support defects zero;
\item compatibility of \(\theta^\Phi_{RR'}\) with the
Thom--Sebastiani vanishing-cycle isomorphisms on extension and
two-step flag stacks and with proper base change on retained support;
\item zero d-critical, potential, cosection, vanishing-cycle
transport, Thom--Sebastiani, base-change, and composition defect
ranks, together with \(R^1\varprojlim=0\) for the reduced
vanishing-cycle complexes.
\end{enumerate}
Scalar traces, Euler characteristics, Behrend-weighted numbers,
Pfaffian products, orientation lines alone, target root windows, D0
scalar tests, and signed multiplicities are not vanishing-cycle
transition data.
\end{definition}

\begin{proposition}[Transition preservation of reduced vanishing
cycles; \textit{proved}]
\label{prop:transition-vanishing-cycle-preservation-criterion}
Given a vanishing-cycle-preserving transition datum of
Definition~\ref{def:transition-vanishing-cycle-preservation}, the
transition \(R\to R'\) preserves the reduced BBDJS perverse sheaves of
vanishing cycles, their cosection reduction, their
Thom--Sebastiani products, and their compact-support base-change
domains.
\end{proposition}

\begin{proof}
The d-critical chart transition and potential compatibility identify
the local vanishing-cycle functors after the recorded quadratic
stabilization.  The cosection compatibility identifies the reduced
obstruction theories, so the reduced vanishing-cycle complexes lie on
the same reduced critical locus after pullback to
\(\mathfrak G_{RR'}\).  The isomorphism
\(\theta^\Phi_{RR'}\) identifies the two BBDJS complexes in
\(D^b_c(\mathfrak G_{RR'})\), and the perverse-normalization row
records the shift introduced by stabilization.  The
Thom--Sebastiani and proper-base-change rows make this identification
compatible with extension and two-step flag correspondences.  The
zero-defect and \(R^1\varprojlim\) rows give strict composition and
inverse-limit exactness for the reduced vanishing-cycle system.
\end{proof}

The packet
\(\texttt{certificates/vanishing\_cycles/transition\_vanishing\_cycle\_preservation}\)
records the present state of this datum.  Its verifier status is
\texttt{TRANSITION\_VANISHING\_CYCLE\_PRESERVATION\_OBSTRUCTION\_VERIFIED}:
the stack-transition input, orientation-transition input, d-critical
chart transition rows, potential-compatibility rows, quadratic
stabilization rows, cosection-pullback rows, reduced obstruction
complex maps, graph-pullback vanishing-cycle isomorphisms, perverse
normalization rows, support and base-change rows,
Thom--Sebastiani vanishing-cycle compatibility rows, strict
composition rows, and vanishing-cycle Mittag--Leffler rows are not yet
supplied.  Hence transition preservation of reduced vanishing cycles
remains open.

\subsubsection*{Cosection localization and the cosection-reduced
heart}

After the surjectivity, kernel, and perfectness rows have been
supplied, the K3 semiregularity cosection
\(\operatorname{CS}_{\mathcal C}\) selects the reduced obstruction
theory.  Only then does the corresponding Behrend constructible
function \(\nu^{\mathrm{red}}\) and the reduced vanishing-cycle complex
\(\Phi^{\mathrm{red}}\) enter the Hall product
\cite{BEHREND,JoyceSong}.  The filtered Hall additivity
\(q^*\operatorname{CS}_C=p_1^*\operatorname{CS}_A+p_2^*\operatorname{CS}_B\)
on extension stacks is the later cosection-localization compatibility
required by the Hall correspondence; the Joyce--Upmeier orientation
lines are required to be Thom--Sebastiani-coherent on the resulting
two-step and four-input flag stacks.

\begin{proposition}[Finite d-critical Hall product;
\textit{conditional on the (O1) and atlas data}]
\label{prop:finite-dcritical-hall-product}
Fix $R$ and closed-configuration substacks
$\mathfrak M_{c,R,I}^{\mathrm{loc,cl}}$ in a noetherian abelian Hall
heart, and assume:
\begin{enumerate}
\item finite-type quasi-smooth derived Artin enhancements of the
selected object, extension, and two-step flag stacks, with perfect
obstruction theories;
\item finite residual inertia after scalar, $E$-translation, and
wrapped rigidifications;
\item retained closed flag-Quot/filtration extension and flag stacks,
or an explicitly specified compact-support six-functor model, with
admissible $p^*$ and $q_!$;
\item PTVV $(-1)$-shifted symplectic structures restricting to these
substacks, classical truncations carrying Joyce d-critical
structures \cite{PTVV2013,BBDJS2015};
\item surjective K3 semiregularity cosections with filtered Hall
additivity;
\item Joyce--Upmeier orientation lines and orientation transports
\cite{JoyceUpmeier} with all finite-stabilizer characters
trivialized;
\item Thom--Sebastiani transport for the reduced vanishing-cycle
sheaves satisfying two-step flag pentagons and higher-coloured
configuration coherences.
\end{enumerate}
Then the finite reduced local-local product
\[
m_R^{\mathrm{red}}=q_!\circ\operatorname{TS}^{\mathrm{red}}_R\circ p^*
\]
is an honest morphism on
$H_c^{\bullet}(\mathfrak M_{c,R,I}^{\mathrm{loc,cl}},
\Phi_{c,R,I}^{\mathrm{red}}\otimes o_{c,R,I}^{\mathrm{red}})$
and is associative on the finite stage.
\end{proposition}

The first obstruction to removing these hypotheses is finite inertia:
proper finite-type correspondences do not by themselves eliminate
residual $\mathbb G_m$, direct-sum, or parabolic automorphisms, and a
residual $\mathbb G_m$ already contributes
$H^\bullet(B\mathbb G_m)=\C[u]$ to protected cohomology.  The phrase
\emph{finite stage} here means finite charge support, finite-type
stacks, and finite-dimensional protected cohomology after the stated
quotients and admissibility hypotheses; it does not mean that the
Hall correspondence morphisms are themselves finite morphisms.
The raw exact-triangle stack is not a proper Hall correspondence; the
proper object is the retained compactified filtration stack, or else
the compact-support functor \(q_!\) is an extra datum.

\subsubsection*{Quotient-after-correspondence ordering}

The $E$-quotient is taken \emph{after} the extension correspondences,
not before.  A pseudofunctor
$\mathsf{Corr}^{E,\mathrm{hyb}}_R\to
\mathsf{Corr}^{\mathrm{red},\mathrm{hyb}}_R$ supplies the descent;
applying the vanishing-cycle Borel--Moore functor with orientation
lines included gives $Q_{E,R}$ between the corresponding finite
correspondence module categories.  For each operation
$\mu_R^\bullet=q_!\operatorname{TS}p^*$ a chosen comparison
isomorphism
\[
\theta^Q_{\mu,R}:Q_{E,R}\,q_!\operatorname{TS}p^*
\xrightarrow{\sim}
\overline q_!\overline{\operatorname{TS}}\,\overline p^*
(Q_{E,R}^{\boxtimes k})
\]
is part of the datum.  Its compatibility with the eight-word flag
$2$-morphisms and HN transitions is the (O1) descent calculus.  A
quotient-first construction (object-stack quotient followed by a new
Hall product) is not the same datum.

\subsubsection*{The orientation obstruction class}

The orientation lane carries the residual
\[
\mathfrak O_{\mathrm{or},R}
=\bigl(
\mathfrak o^{\mathrm{red}}_{\mathcal C},
\mathfrak o^{E,\mathrm{free}}_{\mathcal C},
\{\mathfrak o^{H}_{\mathcal C,S}\}_{H,S},
\{\mathfrak o^{\lambda^H}_{\mathcal C,S}\}_{H,S},
[c_o]_R,\,\xi_{c_o,R},\,
\{\mathfrak o^{\tau,2}_{i,R}\}_i,\,
\{\omega_{i,\mathcal C,R}\}_i\bigr)
\]
required to vanish at every finite HN height, with chosen
null-trivializations.  Theorems~\ref{thm:G-O1}
and~\ref{thm:G-O1-plus} of
Appendix~\ref{app:obstruction-discharges} are the quotient-orientation
and Weyl-transport criteria relative to that retained datum.  The
packet
\[
\texttt{certificates/orientation/k3e\_reduced\_orientation/blocked\_obligations.csv}
\]
records the missing finite equations: determinant square-root and
orientation-class rows, Thom--Sebastiani and pentagon rows, the four
quotient Cech--Borel class rows, finite-stabilizer edge reduction,
odd-transfer, Klein-four and two-primary coefficient rows, Weyl lift
involutivity, torsor-defect, quotient-cocycle transport,
projective-cocycle, cochain, Coxeter, transition, Mittag--Leffler, and
scalar-firewall rows.  Its verifier returns
\(\texttt{ORIENTATION\_OBSTRUCTION\_LEDGER\_VERIFIED}\), with
\(\texttt{orientation\_certification=false}\) and
\(\texttt{mathematical\_certification=false}\).  It is the finite list
of equations still required for \((O1)\) and \((O1)^+\), not a proof of
those equations.
The
Klein-four and two-primary finite-stabilizer lemmas there give the
edge coordinates of (O1).  Proposition~\ref{prop:finite-stabilizer-orientation-detection}
is the finite-site test: the Borel spectral sequence must first reduce
the equivariant class to the classifying-space edge, and for
two-primary rank two stabilizers the mixed coefficient must be computed
directly.  The wrapped-stratum Weyl-equivariant transport along
\(W^{(2)}(\La^{2,1}_{II})\) is reduced to the no-braid type-II Coxeter
presentation and the Maass-character match.

\subsubsection*{Comparison with the Calabi--Yau quantum-group stratum}

Vol III's Calabi--Yau quantum-groups stratification observes that
$K3\times E$ is not presented by a global toric/quiver-variety NCCR
atlas of the Schiffmann--Vasserot type, so the orientation lane
cannot be supplied by such an NCCR-induced heart; the
cosection-twisted reduced d-critical structure is forced on the
chosen stability heart instead.  The displayed Borel
calculation at $E[2]$ is the smallest finite stabilizer detected by
the cosection-reduced determinant; even $N\ge4$ stabilizers carry the
mixed coefficient $A_{12}^{(N)}$, which the displayed cyclic
restrictions miss, and which (O1) requires to be null-trivialized as
part of the full degree-two gerbe.

The orientation $o_R$ trivializes the Joyce--Upmeier obstruction to
a compatible Hall product on the cosection-twisted heart.  Once
$o_R$ is in place, the Hall product, the collision coproduct, and
the protected primitive extraction are unambiguous.  Entry
$(\mathrm L 3)$ is the Hall apparatus of
\S\ref{subsec:hall-correspondences}.


\subsection{Hall correspondences: product, coproduct, primitives}
\label{subsec:hall-correspondences}

The third entry of the Dirac--Igusa pro-object is the Hall structure
on the cosection-reduced d-critical heart.  It is the algebraic
content that the orientation lane equips and the hybrid carrier
indexes.  Three operations live here: the Hall product
$m_R^{\mathrm{red}}$, the collision coproduct $\Delta_R^{\mathrm{ch}}$,
and the primitive projection $P_R^{\Prim}$.  After normal-ordered
Gram descent these form at most a graded bialgebra-and-primitive
datum.  The word ``Hopf'' is reserved for the additional antipode data
of Definition~\ref{def:finite-hall-antipode-datum}; recognition with
the Borcherds--Kac--Moody target is the \S5.6 obligation only after
that data and the radical quotient data have been supplied.

The Kontsevich--Soibelman cohomological Hall algebra of a
$3$-Calabi--Yau category \cite{KSCoHA} is the prototype; on
$K3\times E$ the cosection reduction \cite{JoyceSong,JoyceUpmeier}
replaces the unreduced critical CoHA, and the wrapped sector is the
residual the hybrid carrier of \S\ref{subsec:hybrid-ran-carrier}
retains.  The Maulik--Okounkov stable-envelope construction and the
Schiffmann--Vasserot theorems on Drinfeld doubles of CoHA operate
on the same input; their integral form is the Hall side of the
present construction.

\begin{definition}[Finite compact Hall object]
\label{def:finite-geometric-hall-object}
Fix a finite HN height \(R\) and a finite retained charge set
\(\widehat\Gamma_R\).  A finite compact Hall object at height \(R\) is
the \(\widehat\Gamma_R\)-graded object
\[
\mathcal H_R^{\mathrm{geom}}
=
\bigoplus_{\widehat c\in\widehat\Gamma_R}
H_*^{\mathrm{BM}}\!\left(
\mathfrak M^{\mathrm{red}}_{R,\widehat c},
\Phi^{\mathrm{red}}_{R,\widehat c}\otimes
o_{R,\widehat c}\right),
\]
where the summand is defined only after the following data are fixed:
the finite-type retained rigidified stack
\(\mathfrak M^{\mathrm{red}}_{R,\widehat c}\) with finite residual
inertia, the cosection-reduced BBDJS vanishing-cycle complex
\(\Phi^{\mathrm{red}}_{R,\widehat c}\), the Joyce--Upmeier orientation
line \(o_{R,\widehat c}\), and an admissible compact-support
Borel--Moore six-functor model giving finite-rank groups.  The parity
is the cohomological parity of these Borel--Moore groups together with
the orientation parity recorded in the source basis provenance.

The object \(\mathcal H_R^{\mathrm{geom}}\) is only the finite
coefficient object.  It does not include the Hall product, coproduct,
unit, counit, Hopf pairing, radical quotient, PBW filtration, primitive
subspace, normal-ordered Gram descent, or transition maps.  Those are
additional source rows.  A scalar trace, Euler characteristic, signed
root multiplicity, Pfaffian product, target root window, or denominator
product is not a substitute for the displayed Borel--Moore direct sum.
\end{definition}

The packet
\(\texttt{certificates/hall/finite\_geometric\_hall\_object}\)
verifies
\(\texttt{FINITE\_GEOMETRIC\_HALL\_OBJECT\_DEFINITION\_VERIFIED}\):
it records Definition~\ref{def:finite-geometric-hall-object} as the row
\(341\) definition and imports the retained finite-moduli packet, the
reduced-orientation obstruction ledger, the transition
vanishing-cycle obstruction ledger, and the compact Hall source
obstruction ledger.  The present source packets do not yet populate
\(\mathcal H_R^{\mathrm{geom}}\): no orientation-line rows,
vanishing-cycle coefficient rows, object-component rows, source-degree
rows, homogeneous basis provenance rows, or finite-rank rows are
supplied.  Hence row \(341\) is a definition of the finite compact Hall
object, not a proof that the compact \(K3\times E\) source stage has
been populated.

\subsubsection*{The compact Hall product}

\begin{definition}[Compact Hall product on the cosection-reduced
heart]
\label{def:compact-hall-product}
Let $\mathcal C_{\sigma,S,\le R}$ be the finite HN quotient of
\S\ref{subsec:hybrid-ran-carrier}.  Assume the retained Liu--Hilbert
compactified correspondence and the compact-support pushforward of
Proposition~\ref{prop:finite-dcritical-hall-product}, including finite
residual inertia, base-change, projection-formula, Thom--Sebastiani, and
quotient-after-correspondence coherences.  The compact Hall
product at height $R$ is the operation
\[
m_R^{\mathrm{red}}=
q_{c,c'}{}_!\circ\operatorname{TS}^{\mathrm{red}}_{c,c'}\circ p_{c,c'}^*:
H_*^{\mathrm{BM}}\!\bigl(\mathfrak M_c^\sigma\times\mathfrak M_{c'}^\sigma,
\Phi_c^{\mathrm{red}}\boxtimes\Phi_{c'}^{\mathrm{red}}
\otimes o_c\boxtimes o_{c'}\bigr)
\to
H_*^{\mathrm{BM}}\!\bigl(\mathfrak M_{c+c'}^\sigma,
\Phi_{c+c'}^{\mathrm{red}}\otimes o_{c+c'}\bigr)
\]
attached to the retained compactified extension correspondence
\[
\mathfrak M_c^\sigma\times\mathfrak M_{c'}^\sigma
\xleftarrow{p_{c,c'}}\overline{\mathfrak E}_{c,c'}^{\mathrm{ret},\sigma}
\xrightarrow{q_{c,c'}}\mathfrak M_{c+c'}^\sigma,
\]
in the cosection-reduced d-critical heart, with vanishing-cycle
coefficients $\Phi^{\mathrm{red}}$, Joyce--Upmeier orientation lines
$o$, and Thom--Sebastiani transport
$\operatorname{TS}^{\mathrm{red}}_{c,c'}$.  This definition names the
operator once the three functorial pieces \(p^*\),
\(\operatorname{TS}^{\mathrm{red}}\), and \(q_!\) are supplied in the
chosen six-functor model.  It does not prove that \(p^*\) exists, that
\(q_!\) exists for the retained correspondence, or that \(q_!\) is
independent of a compactification; these are separate functorial
obligations.
It also does not supply product matrices, associativity, transition
compatibility, or primitive closure.  Associativity on the finite stage
is the eight-word two-step binary flag pentagon together with the
quotient-after-correspondence and HN-transition coherences of
\S\ref{subsec:hybrid-ran-carrier}.
\end{definition}

The packet
\(\texttt{certificates/hall/finite\_hall\_product\_operator}\)
verifies
\(\texttt{FINITE\_HALL\_PRODUCT\_OPERATOR\_DEFINITION\_VERIFIED}\):
it records the row \(342\) definition
\[
m_R^{\mathrm{red}}=q_!\operatorname{TS}^{\mathrm{red}}p^*
\]
and imports the row \(341\) finite compact Hall object, the
compact-support exceptional pushforward model, the compact Hall source
obstruction ledger, and the transition-Hall-product obstruction ledger.
The present packets give only the formal product expression and the
chosen compact-support notation.  They do not populate retained
extension-correspondence rows, \(p^*\) rows, \(q_!\) rows, reduced
Thom--Sebastiani rows on the Hall correspondence, product matrix rows,
base-change or projection-formula rows, compactification-independence
rows, associativity rows, or product-transition rows.

\begin{proposition}[Pullback leg of the finite Hall product]
\label{prop:finite-hall-product-pullback-existence}
\textnormal{[six-functor criterion proved; retained \(p\)-row not supplied]}
Let
\[
p_{c,c'}:\overline{\mathfrak E}_{R,c,c'}^{\mathrm{ret}}
\longrightarrow
\mathfrak M_{R,c}^{\mathrm{red}}\times
\mathfrak M_{R,c'}^{\mathrm{red}}
\]
be a supplied retained extension-correspondence morphism at finite HN
height \(R\).  Assume that \(p_{c,c'}\) is a finite-type morphism in
the retained Artin-stack category on which the chosen constructible
six-functor formalism is installed, and that the input coefficient
\[
K_{R,c,c'}=
(\Phi^{\mathrm{red}}_{R,c}\otimes o_{R,c})
\boxtimes
(\Phi^{\mathrm{red}}_{R,c'}\otimes o_{R,c'})
\]
is a bounded constructible complex on the input product stack.  Then
the pullback
\[
p_{c,c'}^*K_{R,c,c'}
\in
D^b_c(\overline{\mathfrak E}_{R,c,c'}^{\mathrm{ret}})
\]
exists in the same six-functor formalism.  This is the \(p^*\)-leg of
the finite Hall product \(m_R^{\mathrm{red}}\).

The present packets do not supply the hypotheses for the actual
compact \(K3\times E\) Hall product.  The finite-moduli packet has no
extension-stack rows, finite-type flag rows, source-map rows, or
d-critical coefficient rows on the extension correspondence.  The row
\(341\) finite-Hall-object packet has no object-component or
coefficient rows.  The row \(342\) product-operator packet has no
populated \(p^*\) rows.  The compact-support packet gives notation for
compact-support pushforwards, not a populated pullback leg.  Hence row
\(343\) is presently this pullback-existence criterion and obstruction
ledger, not a proof that the retained Hall product has a populated
\(p^*\)-leg.
\end{proposition}

\begin{proof}
In a constructible six-functor theory on the retained finite-type
Artin-stack category, every morphism \(p\) in the category carries a
pullback functor
\[
p^*:D^b_c(Y)\longrightarrow D^b_c(X).
\]
The finite-type and constructibility hypotheses ensure that bounded
constructible coefficients remain in the retained coefficient category
after pullback.  Applying this to the supplied extension morphism
\(p_{c,c'}\) and to \(K_{R,c,c'}\) gives the displayed object.  The
argument uses only the pullback functor; it does not construct
\(q_!\), prove Thom--Sebastiani compatibility, produce product
matrices, or prove associativity.
\end{proof}

The packet
\(\texttt{certificates/hall/finite\_hall\_product\_pullback\_functor}\)
verifies
\(\texttt{FINITE\_HALL\_PRODUCT\_PULLBACK\_FUNCTOR\_OBSTRUCTION\_VERIFIED}\):
it imports the row \(342\) product-operator packet, the finite-moduli
obstruction ledger, the row \(341\) finite-Hall-object packet, and the
compact-support model; it records the pullback-existence criterion; and
it keeps the retained extension correspondence, finite-type \(p\)-map,
constructible input coefficient, populated \(p^*\)-row, and product
matrix rows open.

\begin{proposition}[Exceptional pushforward leg of the finite Hall
product]
\label{prop:finite-hall-product-pushforward-existence}
\textnormal{[compact-support criterion proved; retained \(q\)-row not
supplied]}
Let
\[
q_{c,c'}:\overline{\mathfrak E}_{R,c,c'}^{\mathrm{ret}}
\longrightarrow
\mathfrak M_{R,c+c'}^{\mathrm{red}}
\]
be a supplied retained extension-correspondence morphism at finite HN
height \(R\), and let
\[
K_{\mathfrak E}\in
D^b_c(\overline{\mathfrak E}_{R,c,c'}^{\mathrm{ret}})
\]
be the bounded constructible coefficient obtained after the preceding
pullback and reduced Thom--Sebastiani transport.  Assume that
\(q_{c,c'}\) lies in the chosen retained finite-type Artin-stack
category with finite residual inertia and that one of the following
two pieces of functorial data is supplied:
\begin{enumerate}
\item \(q_{c,c'}\) is proper in that category; or
\item a compact-support model
\[
j_q:\overline{\mathfrak E}_{R,c,c'}^{\mathrm{ret}}
\hookrightarrow
\overline{\mathfrak E}_{q},\qquad
\overline q:\overline{\mathfrak E}_{q}\to
\mathfrak M_{R,c+c'}^{\mathrm{red}},
\qquad
q_{c,c'}=\overline q\circ j_q
\]
is supplied as in
Definition~\ref{def:compact-support-exceptional-pushforward-model},
with \(j_{q,!}K_{\mathfrak E}\) and
\(\overline q_*j_{q,!}K_{\mathfrak E}\) in the same constructible
six-functor formalism.
\end{enumerate}
Then the exceptional pushforward leg
\[
q^{\mathrm{cs}}_{c,c',!}K_{\mathfrak E}
:=
\overline q_*j_{q,!}K_{\mathfrak E}
\in D^b_c(\mathfrak M_{R,c+c'}^{\mathrm{red}})
\]
exists; in the proper case it is the ordinary proper pushforward
\(q_{c,c',*}K_{\mathfrak E}\).  This is the \(q_!\)-leg of the finite
Hall product \(m_R^{\mathrm{red}}\).

The present packets do not supply the hypotheses for the actual
compact \(K3\times E\) Hall product.  The row \(342\) product operator
records \(q_!\) only as notation.  The row \(343\) packet proves only
the pullback criterion.  The finite-moduli ledger has no extension
flag-stack row, no retained \(q\)-map row, no properness row for
\(q\), and no d-critical coefficient row on the extension
correspondence.  The compact-support model defines
\(\overline q_*j_{q,!}\), but its map-population and
choice-independence tables are empty.  The compact Hall source packet
has no product matrix rows.  Hence row \(344\) is this exceptional
pushforward criterion and obstruction ledger, not a proof that the
retained Hall product has a populated \(q_!\)-leg.
\end{proposition}

\begin{proof}
In the proper case, the six-functor formalism supplies the proper
direct image
\[
q_{c,c',*}:D^b_c(\overline{\mathfrak E}_{R,c,c'}^{\mathrm{ret}})
\longrightarrow
D^b_c(\mathfrak M_{R,c+c'}^{\mathrm{red}}),
\]
and properness preserves constructibility.  In the compact-support
case, the supplied open immersion gives extension by zero
\[
j_{q,!}K_{\mathfrak E}\in D^b_c(\overline{\mathfrak E}_{q}),
\]
and the supplied proper map \(\overline q\) gives the proper direct
image
\[
\overline q_*j_{q,!}K_{\mathfrak E}
\in D^b_c(\mathfrak M_{R,c+c'}^{\mathrm{red}}).
\]
This composite is exactly the compact-support exceptional pushforward
of Definition~\ref{def:compact-support-exceptional-pushforward-model}.
The argument does not compare two compactifications, prove base
change or the projection formula, construct the reduced
Thom--Sebastiani row, produce product matrices, or prove associativity.
\end{proof}

The packet
\(\texttt{certificates/hall/finite\_hall\_product\_pushforward\_functor}\)
verifies
\(\texttt{FINITE\_HALL\_PRODUCT\_PUSHFORWARD\_FUNCTOR\_OBSTRUCTION\_VERIFIED}\):
it imports the row \(342\) product-operator packet, the row \(343\)
pullback packet, the finite-moduli obstruction ledger, the row \(341\)
finite-Hall-object packet, the compact-support model, and the compact
Hall source obstruction ledger; it records the compact-support
pushforward criterion; and it keeps the retained extension
correspondence, retained \(q\)-map, properness or compactification
row, constructible coefficient row, populated \(q_!\)-row,
compactification independence, and product matrix rows open.

\begin{proposition}[Compactification independence criterion for the
Hall \(q_!\)-leg]
\label{prop:finite-hall-product-qbang-compactification-independence}
\textnormal{[common-refinement criterion proved; compactification pair
not supplied]}
Let
\[
q:\mathfrak E\longrightarrow Z
\]
be a supplied retained Hall target morphism at finite HN height \(R\),
and let \(K_{\mathfrak E}\in D^b_c(\mathfrak E)\) be the bounded
constructible coefficient on the correspondence source.  Suppose two
compact-support models for \((q,K_{\mathfrak E})\) are supplied:
\[
j_i:\mathfrak E\hookrightarrow\overline{\mathfrak E}_i,\qquad
\overline q_i:\overline{\mathfrak E}_i\to Z,\qquad
q=\overline q_i\circ j_i,\qquad i=1,2,
\]
with \(\overline q_i\) proper.  Assume further that a common proper
refinement is supplied:
\[
j_{12}:\mathfrak E\hookrightarrow\overline{\mathfrak E}_{12},
\qquad
r_i:\overline{\mathfrak E}_{12}\to\overline{\mathfrak E}_i,\qquad
\overline q_{12}:\overline{\mathfrak E}_{12}\to Z,
\]
where the maps \(r_i\) and \(\overline q_{12}\) are proper and satisfy
\[
r_i\circ j_{12}=j_i,\qquad
\overline q_i\circ r_i=\overline q_{12}.
\]
Finally assume that the extension-by-zero comparison morphisms
\[
r_{i,*}j_{12,!}K_{\mathfrak E}
\xrightarrow{\;\sim\;}
j_{i,!}K_{\mathfrak E},
\qquad i=1,2,
\]
are isomorphisms in the chosen constructible six-functor formalism.
Then the two compact-support pushforwards are canonically identified:
\[
\overline q_{1,*}j_{1,!}K_{\mathfrak E}
\cong
\overline q_{12,*}j_{12,!}K_{\mathfrak E}
\cong
\overline q_{2,*}j_{2,!}K_{\mathfrak E}.
\]
Thus \(q_!^{\mathrm{cs}}K_{\mathfrak E}\) is independent of the two
compactifications whenever the common-refinement datum and the two
comparison isomorphisms are supplied.

The present packets do not supply the hypotheses for the actual
compact \(K3\times E\) Hall product.  The compact-support model has
only the abstract formula
\(\overline q_*j_{q,!}\); its choice-independence table is empty.  The
row \(344\) packet records the existence criterion for a supplied
compactification, not a comparison between two choices.  The
finite-moduli ledger has no retained \(q\)-map row, no compactification
pair row, and no common-refinement row.  Hence row \(345\) is this
common-refinement criterion and obstruction ledger, not a proof that
the actual retained Hall \(q_!\)-leg is already independent of
compactification.
\end{proposition}

\begin{proof}
For \(i=1,2\), proper functoriality gives
\[
\overline q_{i,*}r_{i,*}
=
(\overline q_i\circ r_i)_*
=
\overline q_{12,*}.
\]
Composing this equality with the supplied comparison isomorphism
\[
r_{i,*}j_{12,!}K_{\mathfrak E}\simeq j_{i,!}K_{\mathfrak E}
\]
gives
\[
\overline q_{i,*}j_{i,!}K_{\mathfrak E}
\simeq
\overline q_{i,*}r_{i,*}j_{12,!}K_{\mathfrak E}
=
\overline q_{12,*}j_{12,!}K_{\mathfrak E}.
\]
The two choices \(i=1,2\) are therefore identified through the same
middle object.  The proof uses only proper functoriality and the
supplied extension-by-zero comparison maps; it does not construct the
common refinement, prove base change or the projection formula, produce
product matrices, or prove associativity.
\end{proof}

The packet
\(\texttt{certificates/hall/finite\_hall\_product\_compactification\_independence}\)
verifies
\(\texttt{FINITE\_HALL\_PRODUCT\_COMPACTIFICATION\_INDEPENDENCE\_OBSTRUCTION\_VERIFIED}\):
it imports the row \(342\) product-operator packet, the row \(344\)
pushforward packet, the finite-moduli obstruction ledger, the
compact-support model, and the compact Hall source obstruction ledger;
it records the common-refinement criterion; and it keeps the retained
\(q\)-map, compactification-pair row, common-refinement row,
comparison-isomorphism row, populated choice-independence row, and
product matrix rows open.

\begin{proposition}[Associativity criterion via retained two-step
flags]
\label{prop:finite-hall-product-associativity-two-step-flags}
\textnormal{[two-step-flag criterion proved; populated \(m_R\)-rows
not supplied]}
Let \(c_1,c_2,c_3\in\widehat\Gamma_R\).  Suppose the finite Hall
product legs
\[
m_{a,b}^{\mathrm{red}}
=q_{a,b,!}^{\mathrm{cs}}\operatorname{TS}^{\mathrm{red}}_{a,b}p_{a,b}^*
\]
are supplied, choice-independent, and constructible for the four pairs
\[
(c_1,c_2),\quad (c_1+c_2,c_3),\quad
(c_2,c_3),\quad (c_1,c_2+c_3)
\]
that occur in the two parenthesizations.  Assume also that a retained
two-step extension flag stack
\[
\mathfrak F^{\mathrm{ret}}_{R;c_1,c_2,c_3}
\]
is supplied, with comparison morphisms to the two iterated
correspondence fibre products
\[
\lambda:\mathfrak F^{\mathrm{ret}}_{R;c_1,c_2,c_3}\to
\overline{\mathfrak E}^{\mathrm{ret}}_{R;c_1+c_2,c_3}
\times_{\mathfrak M^{\mathrm{red}}_{R,c_1+c_2}}
\overline{\mathfrak E}^{\mathrm{ret}}_{R;c_1,c_2},
\]
\[
\rho:\mathfrak F^{\mathrm{ret}}_{R;c_1,c_2,c_3}\to
\overline{\mathfrak E}^{\mathrm{ret}}_{R;c_1,c_2+c_3}
\times_{\mathfrak M^{\mathrm{red}}_{R,c_2+c_3}}
\overline{\mathfrak E}^{\mathrm{ret}}_{R;c_2,c_3}.
\]
Assume the squares defining \(\lambda\) and \(\rho\) satisfy the
base-change isomorphisms for the relevant \(p^*\)- and \(q_!^{cs}\)
legs, that the projection formula holds for the coefficient objects,
and that the reduced Thom--Sebastiani identifications on the two
parenthesizations agree on
\(\mathfrak F^{\mathrm{ret}}_{R;c_1,c_2,c_3}\).  Finally assume that
the two compact-support target pushforwards from the flag stack to
\(\mathfrak M^{\mathrm{red}}_{R,c_1+c_2+c_3}\) are identified by the
compactification-independence criterion of
Proposition~\ref{prop:finite-hall-product-qbang-compactification-independence}.

Then for all supplied constructible inputs
\[
K_i\in D^b_c(\mathfrak M^{\mathrm{red}}_{R,c_i}),\qquad i=1,2,3,
\]
there is a canonical associativity isomorphism
\[
m^{\mathrm{red}}_{c_1+c_2,c_3}
\bigl(m^{\mathrm{red}}_{c_1,c_2}(K_1,K_2),K_3\bigr)
\cong
m^{\mathrm{red}}_{c_1,c_2+c_3}
\bigl(K_1,m^{\mathrm{red}}_{c_2,c_3}(K_2,K_3)\bigr).
\]
After Borel--Moore homology and after choosing homogeneous bases, this
is the rowwise associativity identity for the finite product matrix
\(M_R\).

The present packets do not supply the hypotheses for the actual
compact \(K3\times E\) Hall product.  Rows \(342\)--\(345\) define the
operator and record criteria for \(p^*\), \(q_!\), and
compactification independence, but their populated product,
pushforward, and independence tables are empty.  The hybrid two-step
flag packet constructs the eight local/wrapped flag stacks and
comparison maps before the reduced quotient; the wordwise
associativity packets are conditional on functorial rows and do not
populate the compact Hall carrier.  The compact Hall source packet has
no product matrix entries and no Hall-bialgebra associativity rows.
Hence row \(346\) is this associativity criterion and obstruction
ledger, not a proof that the actual retained \(m_R\) is already
associative as a populated finite matrix.
\end{proposition}

\begin{proof}
Expand the left parenthesization by the definition of
\(m^{\mathrm{red}}\).  It is the compact-support pushforward along the
left iterated correspondence after pulling back
\(K_1\boxtimes K_2\boxtimes K_3\) and applying reduced
Thom--Sebastiani twice.  The right parenthesization is the analogous
pushforward along the right iterated correspondence.  Pull both
iterated correspondences back to
\(\mathfrak F^{\mathrm{ret}}_{R;c_1,c_2,c_3}\) by \(\lambda\) and
\(\rho\).  The supplied base-change and projection-formula
isomorphisms identify the two iterated pull--push expressions with
compact-support pushforwards from this common flag stack.  The supplied
Thom--Sebastiani associativity isomorphism identifies the two
coefficient complexes on the flag stack.  The supplied
compactification-independence comparison identifies the two final
compact-support pushforwards to
\(\mathfrak M^{\mathrm{red}}_{R,c_1+c_2+c_3}\).  The resulting
isomorphism is the displayed associativity isomorphism.  Passing to
Borel--Moore homology and to homogeneous bases gives the matrix
identity whenever the product matrix rows are supplied.
\end{proof}

The packet
\(\texttt{certificates/hall/finite\_hall\_product\_associativity\_two\_step\_flags}\)
verifies
\(\texttt{FINITE\_HALL\_PRODUCT\_ASSOCIATIVITY\_TWO\_STEP\_FLAGS\_OBSTRUCTION\_VERIFIED}\):
it imports the row \(342\) product operator, rows \(344\)--\(345\), the
finite-moduli obstruction ledger, the compact-support model, the
eight-word two-step flag-stack packet, the conditional wordwise
associativity packets, and the compact Hall source obstruction ledger;
it records the two-step-flag associativity criterion; and it keeps the
populated compact product rows, compact Hall two-step flag rows,
functorial base-change/projection/Thom--Sebastiani rows, quotient
descent, transition rows, product matrices, and Hall-bialgebra
associativity rows open.

For local-only inputs and ordered mixed and wrapped inputs, the same
construction produces the three operation families
$\mu_R^{\mathrm{loc}}$, $\mu_R^{\mathrm{hyb}}$, $\mu_R^{\mathrm{wr}}$
of \S\ref{subsec:hybrid-ran-carrier}.  Existence of the product as a
genuine morphism on Borel--Moore chains requires the
admissibility/properness residuals of that section to vanish.

\subsubsection*{The collision coproduct on the source coalgebra}

The collision coproduct is defined on the same correspondence
geometry, read in the opposite direction: instead of contracting an
extension to a single middle term, one expands an object into the pairs of
its sub/quotient extensions.  At finite stage, this is the
$2$-step flag stack already supplied as part of the hybrid datum.

\begin{definition}[Collision coproduct]
\label{def:collision-coproduct}
The finite collision coproduct
\[
\Delta_R^{\mathrm{ch}}:C_{X,R}\longrightarrow C_{X,R}\otimes C_{X,R}
\]
is the chain-level pull-push along the $2$-step flag stack
$\mathfrak{Flag}^{(2)}_{c,c',R}$ of subobjects of charge $c$ and
quotients of charge $c'$ in $\mathfrak M_{c+c',R}^\sigma$, computed
in the cosection-reduced d-critical heart, before the reduced
$E$-quotient.  If the higher-coloured residual, unit maps, refinement
coherences, and quotient-after-correspondence residuals of
\S\ref{subsec:hybrid-ran-carrier} vanish, then after $Q_{E,R}$ it
satisfies the coalgebra identities
\[
(\Delta_R^{\mathrm{ch}}\otimes 1)\Delta_R^{\mathrm{ch}}
=(1\otimes\Delta_R^{\mathrm{ch}})\Delta_R^{\mathrm{ch}},
\qquad
(\varepsilon_R^C\otimes 1)\Delta_R^{\mathrm{ch}}
=(1\otimes\varepsilon_R^C)\Delta_R^{\mathrm{ch}}=\mathrm{id},
\]
which are the dual of the Hall pentagon on the same flag stack.  For
$R'\ge R$ the transition map $\rho^C_{R'R}$ commutes with
$\Delta^{\mathrm{ch}}$.
\end{definition}

\begin{definition}[Compact Hall coproduct on the cosection-reduced
heart]
\label{def:compact-hall-coproduct}
Let
\[
\mathfrak M^{\mathrm{red}}_{R,c}\times
\mathfrak M^{\mathrm{red}}_{R,c'}
\xleftarrow{\ p_{c,c'}\ }
\overline{\mathfrak E}^{\mathrm{ret}}_{R,c,c'}
\xrightarrow{\ q_{c,c'}\ }
\mathfrak M^{\mathrm{red}}_{R,c+c'}
\]
be the retained compactified extension correspondence at finite HN
height \(R\), with the same reduced vanishing-cycle and orientation
coefficients as in Definition~\ref{def:compact-hall-product}.  The
compact Hall coproduct operator is the formal pull--inverse
Thom--Sebastiani--push operator
\[
\Delta_R^{\mathrm{red}}
=
p_{c,c',!}^{\mathrm{cs}}\circ
(\operatorname{TS}^{\mathrm{red}}_{c,c'})^{-1}\circ
q_{c,c'}^*.
\]
It sends a coefficient on
\(\mathfrak M^{\mathrm{red}}_{R,c+c'}\) to a coefficient on
\(\mathfrak M^{\mathrm{red}}_{R,c}\times
\mathfrak M^{\mathrm{red}}_{R,c'}\), and after Borel--Moore
Kunneth identification it is read as
\[
\Delta_R^{\mathrm{red}}:
\mathcal H^{\mathrm{geom}}_{R,c+c'}
\longrightarrow
\mathcal H^{\mathrm{geom}}_{R,c}\otimes
\mathcal H^{\mathrm{geom}}_{R,c'}.
\]
This definition names the operator once the three functorial pieces
\(q^*\), \((\operatorname{TS}^{\mathrm{red}})^{-1}\), and
\(p_!^{\mathrm{cs}}\) are supplied in the chosen six-functor model.  It
does not prove that \(q^*\) exists, that \(p_!\) exists for the
reversed correspondence, that \(p_!\) is independent of a
compactification, or that the coproduct is coassociative.  It also does
not supply counit maps, coproduct matrices, transition compatibility,
bialgebra compatibility, or primitive closure.
\end{definition}

The packet
\(\texttt{certificates/hall/finite\_hall\_coproduct\_operator}\)
verifies
\(\texttt{FINITE\_HALL\_COPRODUCT\_OPERATOR\_DEFINITION\_VERIFIED}\):
it records the row \(347\) definition
\[
\Delta_R^{\mathrm{red}}
=p_!(\operatorname{TS}^{\mathrm{red}})^{-1}q^*
\]
and imports the row \(341\) finite compact Hall object, the
compact-support exceptional pushforward model, the compact Hall source
obstruction ledger, and the transition-Hall-coproduct obstruction
ledger.  The present packets give only the formal coproduct expression
and the chosen compact-support notation.  They do not populate retained
splitting-correspondence rows, \(q^*\) rows, \(p_!\) rows, inverse
Thom--Sebastiani rows, compactification-independence rows, coproduct
matrix rows, counit rows, coassociativity rows, bialgebra rows, or
coproduct-transition rows.

\begin{proposition}[Exceptional pushforward leg of the finite Hall
coproduct]
\label{prop:finite-hall-coproduct-pushforward-existence}
\textnormal{[compact-support criterion proved; retained \(p\)-row not
supplied]}
Let
\[
p_{c,c'}:\overline{\mathfrak E}_{R,c,c'}^{\mathrm{ret}}
\longrightarrow
\mathfrak M_{R,c}^{\mathrm{red}}\times
\mathfrak M_{R,c'}^{\mathrm{red}}
\]
be a supplied retained splitting-correspondence morphism at finite HN
height \(R\), and let
\[
K_{\mathfrak E}\in
D^b_c(\overline{\mathfrak E}_{R,c,c'}^{\mathrm{ret}})
\]
be the bounded constructible coefficient obtained after the preceding
pullback \(q_{c,c'}^*\) and inverse reduced Thom--Sebastiani
transport.  Assume that \(p_{c,c'}\) lies in the chosen retained
finite-type Artin-stack category with finite residual inertia and that
one of the following two pieces of functorial data is supplied:
\begin{enumerate}
\item \(p_{c,c'}\) is proper in that category; or
\item a compact-support model
\[
j_p:\overline{\mathfrak E}_{R,c,c'}^{\mathrm{ret}}
\hookrightarrow
\overline{\mathfrak E}_{p},\qquad
\overline p:\overline{\mathfrak E}_{p}\to
\mathfrak M_{R,c}^{\mathrm{red}}\times
\mathfrak M_{R,c'}^{\mathrm{red}},
\qquad
p_{c,c'}=\overline p\circ j_p
\]
is supplied as in
Definition~\ref{def:compact-support-exceptional-pushforward-model},
with \(j_{p,!}K_{\mathfrak E}\) and
\(\overline p_*j_{p,!}K_{\mathfrak E}\) in the same constructible
six-functor formalism.
\end{enumerate}
Then the exceptional pushforward leg
\[
p^{\mathrm{cs}}_{c,c',!}K_{\mathfrak E}
:=
\overline p_*j_{p,!}K_{\mathfrak E}
\in
D^b_c(\mathfrak M_{R,c}^{\mathrm{red}}\times
\mathfrak M_{R,c'}^{\mathrm{red}})
\]
exists; in the proper case it is the ordinary proper pushforward
\(p_{c,c',*}K_{\mathfrak E}\).  This is the \(p_!\)-leg of the finite
Hall coproduct \(\Delta_R^{\mathrm{red}}\).

The present packets do not supply the hypotheses for the actual
compact \(K3\times E\) Hall coproduct.  The row \(347\) coproduct
operator records \(p_!\) only as notation.  The finite-moduli ledger
has no extension flag-stack row, no retained splitting \(p\)-map row,
no properness row for \(p\), and no d-critical coefficient row on the
splitting correspondence.  The compact-support model defines
\(\overline p_*j_{p,!}\), but its map-population and
choice-independence tables are empty.  The compact Hall source packet
has no coproduct matrix rows, and the transition-Hall-coproduct packet
has no coproduct-transport rows.  Hence row \(348\) is this
exceptional pushforward criterion and obstruction ledger, not a proof
that the retained Hall coproduct has a populated \(p_!\)-leg.
\end{proposition}

\begin{proof}
In the proper case, the six-functor formalism supplies the proper
direct image
\[
p_{c,c',*}:
D^b_c(\overline{\mathfrak E}_{R,c,c'}^{\mathrm{ret}})
\longrightarrow
D^b_c(\mathfrak M_{R,c}^{\mathrm{red}}\times
\mathfrak M_{R,c'}^{\mathrm{red}}),
\]
and properness preserves constructibility.  In the compact-support
case, the supplied open immersion gives extension by zero
\[
j_{p,!}K_{\mathfrak E}\in D^b_c(\overline{\mathfrak E}_{p}),
\]
and the supplied proper map \(\overline p\) gives the proper direct
image
\[
\overline p_*j_{p,!}K_{\mathfrak E}
\in
D^b_c(\mathfrak M_{R,c}^{\mathrm{red}}\times
\mathfrak M_{R,c'}^{\mathrm{red}}).
\]
This composite is exactly the compact-support exceptional pushforward
of Definition~\ref{def:compact-support-exceptional-pushforward-model}.
The argument does not compare two compactifications, prove base
change or the projection formula, construct the \(q^*\) row, construct
the inverse reduced Thom--Sebastiani row, produce coproduct matrices,
or prove coassociativity.
\end{proof}

The packet
\(\texttt{certificates/hall/finite\_hall\_coproduct\_pushforward\_functor}\)
verifies
\(\texttt{FINITE\_HALL\_COPRODUCT\_PUSHFORWARD\_FUNCTOR\_OBSTRUCTION\_VERIFIED}\):
it imports the row \(347\) coproduct-operator packet, the
finite-moduli obstruction ledger, the row \(341\) finite-Hall-object
packet, the compact-support model, the compact Hall source obstruction
ledger, and the transition-Hall-coproduct obstruction ledger; it
records the compact-support pushforward criterion; and it keeps the
retained splitting correspondence, retained \(p\)-map, properness or
compactification row, constructible coefficient row, inverse
Thom--Sebastiani row, populated \(p_!\)-row, compactification
independence, coproduct matrix rows, and transition-coproduct rows
open.

\begin{proposition}[Coassociativity criterion via retained two-step
splitting flags]
\label{prop:finite-hall-coproduct-coassociativity-two-step-flags}
\textnormal{[two-step-splitting-flag criterion proved; populated
\(\Delta_R\)-rows not supplied]}
Let \(c_1,c_2,c_3\in\widehat\Gamma_R\).  Suppose the finite Hall
coproduct legs
\[
\Delta_{a,b}^{\mathrm{red}}
=p_{a,b,!}^{\mathrm{cs}}
(\operatorname{TS}_{a,b}^{\mathrm{red}})^{-1}q_{a,b}^*
\]
are supplied, choice-independent, and constructible for the four
pairs
\[
(c_1+c_2,c_3),\quad (c_1,c_2),\quad
(c_1,c_2+c_3),\quad (c_2,c_3)
\]
that occur in the two coproduct parenthesizations.  Assume also that
a retained two-step splitting flag stack
\[
\mathfrak F^{\mathrm{ret}}_{R;c_1,c_2,c_3}
\]
is supplied, with comparison morphisms to the two iterated splitting
correspondence fibre products
\[
\lambda:\mathfrak F^{\mathrm{ret}}_{R;c_1,c_2,c_3}\to
\overline{\mathfrak E}^{\mathrm{ret}}_{R;c_1+c_2,c_3}
\times_{\mathfrak M^{\mathrm{red}}_{R,c_1+c_2}}
\overline{\mathfrak E}^{\mathrm{ret}}_{R;c_1,c_2},
\]
\[
\rho:\mathfrak F^{\mathrm{ret}}_{R;c_1,c_2,c_3}\to
\overline{\mathfrak E}^{\mathrm{ret}}_{R;c_1,c_2+c_3}
\times_{\mathfrak M^{\mathrm{red}}_{R,c_2+c_3}}
\overline{\mathfrak E}^{\mathrm{ret}}_{R;c_2,c_3}.
\]
Assume the squares defining \(\lambda\) and \(\rho\) satisfy the
base-change isomorphisms for the relevant \(q^*\)- and
\(p_!^{\mathrm{cs}}\)-legs, that the projection formula holds for the
coefficient objects, and that the inverse reduced
Thom--Sebastiani identifications on the two parenthesizations agree on
\(\mathfrak F^{\mathrm{ret}}_{R;c_1,c_2,c_3}\).  Finally assume that
the two compact-support pushforwards from the flag stack to
\[
\mathfrak M^{\mathrm{red}}_{R,c_1}\times
\mathfrak M^{\mathrm{red}}_{R,c_2}\times
\mathfrak M^{\mathrm{red}}_{R,c_3}
\]
are identified by a supplied compactification-independence comparison
for the \(p_!^{\mathrm{cs}}\)-legs.

Then for every supplied constructible input
\[
K\in D^b_c(\mathfrak M^{\mathrm{red}}_{R,c_1+c_2+c_3})
\]
there is a canonical coassociativity isomorphism
\[
(\Delta_{c_1,c_2}^{\mathrm{red}}\otimes 1)
\Delta_{c_1+c_2,c_3}^{\mathrm{red}}(K)
\cong
(1\otimes\Delta_{c_2,c_3}^{\mathrm{red}})
\Delta_{c_1,c_2+c_3}^{\mathrm{red}}(K).
\]
After Borel--Moore homology and after choosing homogeneous bases, this
is the rowwise coassociativity identity for the finite coproduct
matrix \(D_R\).

The present packets do not supply the hypotheses for the actual
compact \(K3\times E\) Hall coproduct.  Row \(347\) names the
coproduct operator but has no populated \(\Delta_R\)-rows.  Row
\(348\) proves only the \(p_!\)-existence criterion for a supplied
retained splitting map, and its pushforward table is empty.  The
finite-moduli ledger has no retained two-step splitting flag stack.
The compact-support model has no base-change, projection-formula,
Thom--Sebastiani, or choice-independence rows.  The hybrid two-step
flag packet constructs local/wrapped flag stacks before the reduced
quotient and does not populate the compact Hall carrier.  The compact
Hall source packet has no coproduct matrix entries and no
Hall-bialgebra coassociativity rows.  Hence row \(349\) is this
coassociativity criterion and obstruction ledger, not a proof that the
actual retained \(\Delta_R\) is already coassociative as a populated
finite matrix.
\end{proposition}

\begin{proof}
Expand the left parenthesization by the definition of
\(\Delta^{\mathrm{red}}\).  It is the compact-support pushforward to
\(\mathfrak M^{\mathrm{red}}_{R,c_1}\times
\mathfrak M^{\mathrm{red}}_{R,c_2}\times
\mathfrak M^{\mathrm{red}}_{R,c_3}\) obtained by pulling \(K\) back
along the two relevant \(q\)-maps, applying inverse reduced
Thom--Sebastiani twice, and pushing along the two relevant
\(p\)-maps.  The right parenthesization is the analogous expression
for the other iterated splitting correspondence.  Pull both iterated
correspondences back to
\(\mathfrak F^{\mathrm{ret}}_{R;c_1,c_2,c_3}\) by \(\lambda\) and
\(\rho\).  The supplied base-change and projection-formula
isomorphisms identify the two iterated pull--inverse-TS--push
expressions with compact-support pushforwards from this common flag
stack.  The supplied inverse Thom--Sebastiani coassociativity
isomorphism identifies the two coefficient complexes on the flag
stack.  The supplied compactification-independence comparison
identifies the two compact-support pushforwards to the triple product.
The resulting isomorphism is the displayed coassociativity
isomorphism.  Passing to Borel--Moore homology and to homogeneous
bases gives the matrix identity whenever the coproduct matrix rows are
supplied.
\end{proof}

The packet
\(\texttt{certificates/hall/finite\_hall\_coproduct\_coassociativity\_two\_step\_flags}\)
verifies
\(\texttt{FINITE\_HALL\_COPRODUCT\_COASSOCIATIVITY\_TWO\_STEP\_FLAGS\_OBSTRUCTION\_VERIFIED}\):
it imports the row \(347\) coproduct operator, row \(348\), the
finite-moduli obstruction ledger, the compact-support model, the
eight-word two-step flag-stack packet, the compact Hall source
obstruction ledger, and the transition-Hall-coproduct obstruction
ledger; it records the two-step-splitting-flag coassociativity
criterion; and it keeps the populated compact coproduct rows, compact
Hall two-step splitting flag rows, functorial
base-change/projection/inverse-Thom--Sebastiani rows,
compactification-independence rows, quotient descent, transition rows,
coproduct matrices, and Hall-bialgebra coassociativity rows open.

\begin{proposition}[Bialgebra compatibility criterion for the finite
Hall product and coproduct]
\label{prop:finite-hall-bialgebra-compatibility}
\textnormal{[product--coproduct square criterion proved; populated
bialgebra rows not supplied]}
Let \(a,b,u,v\in\widehat\Gamma_R\) with \(u+v=a+b\), and set
\[
\mathsf S(a,b;u,v)
=\{(a_1,a_2,b_1,b_2)\mid
a=a_1+a_2,\ b=b_1+b_2,\ u=a_1+b_1,\ v=a_2+b_2\}.
\]
The set \(\mathsf S(a,b;u,v)\) is finite at finite HN height.  Suppose
the product legs \(m^{\mathrm{red}}_{a,b}\),
\(m^{\mathrm{red}}_{a_1,b_1}\),
\(m^{\mathrm{red}}_{a_2,b_2}\) and the coproduct legs
\(\Delta^{\mathrm{red}}_{u,v}\),
\(\Delta^{\mathrm{red}}_{a_1,a_2}\),
\(\Delta^{\mathrm{red}}_{b_1,b_2}\) are supplied, constructible, and
choice-independent for every
\((a_1,a_2,b_1,b_2)\in\mathsf S(a,b;u,v)\).  Assume also that a
retained product--splitting square stack
\[
\mathfrak Q^{\mathrm{ret};u,v}_{R;a,b}
\]
is supplied, with comparison morphisms to the iterated correspondence
for \(\Delta_{u,v}m_{a,b}\) and to the finite disjoint union, indexed
by \(\mathsf S(a,b;u,v)\), of the iterated correspondences for
\[
(m_{a_1,b_1}^{\mathrm{red}}\otimes
m_{a_2,b_2}^{\mathrm{red}})
(1\otimes\tau\otimes1)
(\Delta_{a_1,a_2}^{\mathrm{red}}\otimes
\Delta_{b_1,b_2}^{\mathrm{red}}).
\]
Assume the comparison squares satisfy proper base change for the
relevant \(q^*\), \(p^*\), \(q_!^{\mathrm{cs}}\), and
\(p_!^{\mathrm{cs}}\) legs; the projection formula for the coefficient
objects; the product--coproduct Thom--Sebastiani coherence identifying
the product-side and coproduct-side coefficient complexes on
\(\mathfrak Q^{\mathrm{ret};u,v}_{R;a,b}\); the Koszul sign
convention for the middle braiding \(\tau\); and compact-support
choice-independence for all final pushforwards to
\(\mathfrak M^{\mathrm{red}}_{R,u}\times
\mathfrak M^{\mathrm{red}}_{R,v}\).

Then for all supplied constructible inputs
\[
K_a\in D^b_c(\mathfrak M^{\mathrm{red}}_{R,a}),\qquad
K_b\in D^b_c(\mathfrak M^{\mathrm{red}}_{R,b}),
\]
there is a canonical bialgebra-compatibility isomorphism
\[
\Delta^{\mathrm{red}}_{u,v}
m^{\mathrm{red}}_{a,b}(K_a,K_b)
\cong
\bigoplus_{\mathsf S(a,b;u,v)}
(m^{\mathrm{red}}_{a_1,b_1}\otimes
m^{\mathrm{red}}_{a_2,b_2})
(1\otimes\tau\otimes1)
(\Delta^{\mathrm{red}}_{a_1,a_2}K_a\otimes
\Delta^{\mathrm{red}}_{b_1,b_2}K_b).
\]
After Borel--Moore homology and after choosing homogeneous bases, this
is the rowwise matrix identity
\[
D_R^{a+b;u,v}M_R^{a,b}
=
\sum_{\mathsf S(a,b;u,v)}
(M_R^{a_1,b_1}\otimes M_R^{a_2,b_2})
(1\otimes\tau\otimes1)
(D_R^{a;a_1,a_2}\otimes D_R^{b;b_1,b_2}).
\]

The present packets do not supply the hypotheses for the actual
compact \(K3\times E\) Hall bialgebra.  Row \(342\) names the product
operator and row \(347\) names the coproduct operator; neither packet
has populated operation rows.  Rows \(346\) and \(349\) record the
associativity and coassociativity criteria, but their identity tables
are empty.  The finite-moduli ledger has no retained
product--splitting square stack.  The compact-support model has no
base-change, projection-formula, Thom--Sebastiani, or
choice-independence rows for the mixed product--coproduct square.  The
compact Hall source packet has no \(M\)-entries, \(D\)-entries, or
Hall-bialgebra identity rows.  Hence row \(350\) is this bialgebra
compatibility criterion and obstruction ledger, not a proof that the
actual retained finite source already satisfies
\(\Delta m=(m\otimes m)\Delta\).
\end{proposition}

\begin{proof}
Expand the left side by the definitions of \(m^{\mathrm{red}}\) and
\(\Delta^{\mathrm{red}}\).  It is a compact-support pull--TS--push
along the product correspondence followed by a
pull--inverse-TS--push along the splitting correspondence.  Expand the
right side by first splitting the two inputs, applying the Koszul
braiding to exchange the two middle factors, and multiplying the two
output pairs.  Pull the two iterated correspondences to
\(\mathfrak Q^{\mathrm{ret};u,v}_{R;a,b}\).  Proper base change and
the projection formula identify the two iterated pull-push functors on
this common square stack.  The supplied product--coproduct
Thom--Sebastiani coherence identifies the coefficient complexes; the
middle exchange contributes exactly the Koszul sign \(\tau\).  The
compact-support choice-independence comparisons identify the final
pushforwards to
\(\mathfrak M^{\mathrm{red}}_{R,u}\times
\mathfrak M^{\mathrm{red}}_{R,v}\).  The resulting isomorphism is the
displayed bialgebra-compatibility isomorphism.  Applying
Borel--Moore homology, Kunneth, and homogeneous compact-source bases
gives the displayed finite matrix identity whenever \(M_R\), \(D_R\),
and the bialgebra identity rows are supplied.
\end{proof}

The packet
\(\texttt{certificates/hall/finite\_hall\_bialgebra\_compatibility}\)
verifies
\(\texttt{FINITE\_HALL\_BIALGEBRA\_COMPATIBILITY\_OBSTRUCTION\_VERIFIED}\):
it imports the product and coproduct operator packets, the row \(346\)
associativity packet, the row \(349\) coassociativity packet, the
finite-moduli obstruction ledger, the compact-support model, the
compact Hall source obstruction ledger, and the product- and
coproduct-transition obstruction ledgers; it records the
product--coproduct square criterion; and it keeps the populated
product rows, populated coproduct rows, retained product--splitting
square rows, base-change/projection/Thom--Sebastiani rows,
compactification-independence rows, \(M\)- and \(D\)-matrices,
transition rows, and Hall-bialgebra compatibility rows open.

\begin{definition}[Finite Hall antipode datum]
\label{def:finite-hall-antipode-datum}
A finite Hall antipode datum at HN height \(R\) is extra structure on a
supplied finite Hall bialgebra
\[
(\mathcal H_R,m_R,\Delta_R,\eta_R,\epsilon_R)
\]
whose unit, counit, associativity, coassociativity, and
product--coproduct compatibility rows have zero defect.  It consists
of either:
\begin{enumerate}
\item explicit homogeneous antipode matrices
\[
S_{R,\gamma}:\mathcal H_{R,\gamma}\longrightarrow
\mathcal H_{R,\iota(\gamma)}
\]
for a supplied charge involution \(\iota\) preserving the retained
finite window; or
\item a supplied connected conilpotent filtration for which the
standard recursive antipode
\[
S_R(1)=1,\qquad
S_R(x)=-x-\sum S_R(x')\,x''
\quad
\bigl(\widetilde\Delta_Rx=\sum x'\otimes x''\bigr)
\]
terminates in every retained homogeneous block.
\end{enumerate}
In either presentation the convolution identities
\[
m_R(S_R\otimes\mathrm{id})\Delta_R=\eta_R\epsilon_R,\qquad
m_R(\mathrm{id}\otimes S_R)\Delta_R=\eta_R\epsilon_R
\]
must be supplied as zero-defect rows, and \(S_R\) must commute with the
HN transition maps on every block where the pro-object is formed.  A
finite Hall bialgebra becomes a finite Hall Hopf algebra only after
this antipode datum is supplied.

The present compact \(K3\times E\) packets do not supply an antipode
datum.  The source packet has no antipode matrix rows, no charge
involution row, no connected conilpotent recursion row, and no left or
right convolution identity row.  The bialgebra-compatibility packet
records only the criterion for
\(\Delta m=(m\otimes m)(1\otimes\tau\otimes1)(\Delta\otimes\Delta)\).
Hopf-pairing rows, Drinfeld double comparisons, bar-coalgebra
deconcatenation, primitive closure, scalar traces, and signed
multiplicities are not antipode data.
\end{definition}

The packet
\(\texttt{certificates/hall/finite\_hall\_antipode\_datum}\)
verifies
\(\texttt{FINITE\_HALL\_ANTIPODE\_DATUM\_OBSTRUCTION\_VERIFIED}\):
it imports the row \(350\) bialgebra-compatibility packet, the compact
Hall source obstruction ledger, and the product- and
coproduct-transition obstruction ledgers; it records
Definition~\ref{def:finite-hall-antipode-datum}; and it keeps the
antipode matrices, charge involution or connected-conilpotent
filtration, unit and counit rows, bialgebra rows, convolution-identity
rows, and transition-antipode rows open.  Thus row \(351\) defines the
antipode data required by any finite Hall Hopf-algebra claim; it does
not assert that the current compact Hall source is a Hopf algebra.

\begin{proposition}[No Hopf algebra without an antipode; \textit{proved}]
\label{prop:finite-hall-nonhopf-boundary}
Let
\[
(\mathcal H_R,m_R,\Delta_R,\eta_R,\epsilon_R)
\]
be a supplied finite Hall bialgebra stage.  If the finite Hall
antipode datum of Definition~\ref{def:finite-hall-antipode-datum} is
not supplied, then this stage is not a finite Hall Hopf algebra in the
verified source package of the manuscript.  It may be used only as a
bialgebra-level input.  Every downstream occurrence of a Hopf algebra
claim is therefore conditional on the antipode matrices, or on the
terminating connected-conilpotent recursion, together with the two
convolution identities and transition compatibility.

This is not a non-existence theorem for antipodes on the underlying
graded vector space.  It is the source-data boundary used in this
manuscript: absent \(S_R\) and its convolution identities, the Hopf
algebra name is unavailable.
\end{proposition}

\begin{proof}
A Hopf algebra is a bialgebra equipped with an antipode.  In the
finite Hall presentation this means exactly the datum
\(S_R\) and the two identities
\[
m_R(S_R\otimes\mathrm{id})\Delta_R=\eta_R\epsilon_R,\qquad
m_R(\mathrm{id}\otimes S_R)\Delta_R=\eta_R\epsilon_R .
\]
If no such row is supplied, no convolution inverse for the identity
endomorphism of \(\mathcal H_R\) has been constructed.  Bialgebra
compatibility, Hopf-pairing adjointness, a target Drinfeld double, a
bar-coalgebra model, primitive closure, and scalar trace data do not
define that convolution inverse.  The verified source package
therefore remains bialgebra-level until the antipode datum is added.
\end{proof}

The packet
\(\texttt{certificates/hall/finite\_hall\_nonhopf\_boundary}\)
verifies \(\texttt{FINITE\_HALL\_NONHOPF\_BOUNDARY\_VERIFIED}\): it
imports the row \(351\) antipode-datum packet, the row \(350\)
bialgebra-compatibility packet, and the compact Hall source
obstruction ledger; it records
Proposition~\ref{prop:finite-hall-nonhopf-boundary}; it checks that
the current source has no antipode rows, no convolution identity rows,
and no finite Hall Hopf-algebra claim; and it records explicitly that
no non-existence theorem for future antipodes is being asserted.

This is the source coalgebra of \S\ref{subsec:koszul-source}; its
identification with the bar construction
$\bar B_E^{\mathrm{ch}}(\mathcal F_X^{\mathrm{hyb}})$ is the bar
comparison there.

\subsubsection*{Primitive projection and the bialgebra datum}

The primitive subspace is a bialgebra notion.  It uses the
augmentation, unit, counit, and coproduct; its closure under the
supercommutator uses bialgebra compatibility.  It does not make the
source object Hopf.  Hopf-level recognition is available only after
the antipode datum named in
Proposition~\ref{prop:finite-hall-nonhopf-boundary} is supplied.

\begin{definition}[Finite primitive Hall subspace]
\label{def:finite-hall-primitive-subspace}
Let
\((\mathcal H_R,\Delta_R,\eta_R,\epsilon_R)\) be a finite counital
Hall coalgebra stage with supplied homogeneous finite bases.  Put
\[
I_R=\ker\epsilon_R .
\]
The reduced coproduct is the map
\[
\bar\Delta_R:I_R\longrightarrow I_R\otimes I_R,\qquad
\bar\Delta_R(x)=
\Delta_Rx-\eta_R(1)\otimes x-x\otimes\eta_R(1),
\]
equivalently the projection of \(\Delta_R|_{I_R}\) to
\(I_R\otimes I_R\).  The finite primitive Hall subspace is
\[
P_R=\Prim(\mathcal H_R)
=\ker\bigl(\bar\Delta_R:I_R\to I_R\otimes I_R\bigr).
\]
In finite coordinates, row \(353\) consists of the counit matrix,
the augmentation-ideal basis, the reduced-coproduct matrix
\(\bar D_R\), and kernel rows giving a basis, inclusion, and projection
for \(P_R\).  It is a definition of a subspace of the supplied finite
coalgebra.  It does not assert that the compact \(K3\times E\) source
has supplied those rows, and it does not assert primitive closure under
the Hall commutator.
\end{definition}

The packet
\(\texttt{certificates/hall/finite\_hall\_primitive\_subspace\_definition}\)
verifies
\(\texttt{FINITE\_HALL\_PRIMITIVE\_SUBSPACE\_DEFINITION\_OBSTRUCTION\_VERIFIED}\):
it imports the row \(352\) non-Hopf boundary, the row \(350\)
bialgebra-compatibility packet, and the compact Hall source
obstruction ledger; it records
Definition~\ref{def:finite-hall-primitive-subspace}; and it keeps the
source coproduct matrix, unit and counit rows, augmentation-ideal
basis, reduced-coproduct matrix, primitive-kernel rows, primitive
basis, inclusion, projection, and degree--parity split open.  Thus
row \(353\) defines \(P_R=\ker\bar\Delta_R\) for supplied finite data;
it does not populate \(P_R\) in the present compact source.

\begin{definition}[Primitive Hall sub-bracket]
\label{def:primitive-hall-bracket}
Write
\[
\widetilde P_X=H^0\Prim_{\mathrm{prot}}(\mathcal F_X)
=\ker\bigl(\widetilde\Delta:
\mathcal F_X\to\mathcal F_X\otimes\mathcal F_X\bigr)
\]
for the primitive layer of the protected Hall bialgebra object, where
$\widetilde\Delta=\Delta-(\eta\varepsilon\otimes\mathrm{id}+
\mathrm{id}\otimes\eta\varepsilon)$.  The Hall product restricts to
a $\widehat\Gamma_X$-graded bracket
\[
[\,\cdot\,,\,\cdot\,]_X:\widetilde P_{X,\widehat c}\otimes
\widetilde P_{X,\widehat c'}
\longrightarrow\widetilde P_{X,\widehat c\star\widehat c'},
\]
the protected primitive Hall bracket.  After normal-ordered Gram
descent (\S\ref{subsec:hybrid-ran-carrier}) and the Hopf-radical
quotient, the descended bracket is $\Gamma_{\mathrm{gram}}$-graded
and the descended primitive layer is the candidate $P_X^{\Pi,\mathrm{red}}$
to be identified with $\mathfrak g_{\Delta_5}$.
\end{definition}

The finite packet
\(\texttt{certificates/charge/hall\_pairing\_pushforward\_compatibility}\)
supplies only the grading calculation behind this sentence: a supplied
\(\widehat\Gamma_X\)-graded bracket and a supplied homogeneous
degree-zero pairing remain graded after pushforward by the additive
normal-ordered map \(\overline\Pi_X\).  The compact Hall
correspondence, primitive closure, Hopf adjointness, Frobenius cyclic
identity, quotient non-degeneracy, and protected trace remain separate
data.

\begin{proposition}[Product-only Hall data do not define primitives;
\textit{proved}]
\label{prop:hall-product-only-no-primitives}
Let \(\mathcal H_R\) be a finite HN stage carrying the Hall product
\(m_R\) and unit \(\eta_R\).  Without a counital coproduct
\(\Delta_R\) and the bialgebra compatibility of \(m_R\) with
\(\Delta_R\), there is no defined primitive subspace
\[
\Prim(\mathcal H_R)=
\ker\bigl(\Delta_R-\eta_R\varepsilon_R\otimes\mathrm{id}
-\mathrm{id}\otimes\eta_R\varepsilon_R\bigr),
\]
and the supercommutator of \(m_R\) does not define a Lie bracket on
primitives.  Thus a compact Hall product, even if associative and
oriented, is not a primitive-recognition datum.
\end{proposition}

\begin{proof}
The displayed formula is the definition of primitives in a counital
bialgebra.  If \(\Delta_R\) or \(\varepsilon_R\) is absent, the kernel
is not an object.  If \(\Delta_R\) is present but not compatible with
\(m_R\), the bialgebra identity
\[
\Delta_R([x,y])=[\Delta_Rx,\Delta_Ry]
\]
in the graded tensor-product algebra does not hold.  Hence
\(\Delta_R x=x\otimes1+1\otimes x\) and
\(\Delta_R y=y\otimes1+1\otimes y\) no longer imply
\(\Delta_R([x,y])=[x,y]\otimes1+1\otimes[x,y]\).

The product does not determine the coproduct.  Over a characteristic
zero field, the graded algebra \(A=k[x,y]\), \(|y|=2|x|\), has the
primitive coproduct
\(\Delta_0x=x\otimes1+1\otimes x\),
\(\Delta_0y=y\otimes1+1\otimes y\), and also the counital
coassociative algebra coproduct
\[
\Delta_1x=x\otimes1+1\otimes x,\qquad
\Delta_1y=y\otimes1+1\otimes y+x\otimes x.
\]
The multiplication is the same in both bialgebras, but
\(\Prim(A,\Delta_0)\) and \(\Prim(A,\Delta_1)\) are different
subspaces: \(y\) is primitive for \(\Delta_0\), while
\(y-\frac12x^2\) is primitive for \(\Delta_1\).  Hence the Hall
product alone cannot determine \(P_R^{\Prim}\), the Hopf radical, or
the PBW comparison used in primitive recognition.
\end{proof}

\begin{definition}[Finite Hall bialgebra defect tuple]
\label{def:finite-hall-bialgebra-defect-tuple}
Fix homogeneous finite bases on a retained HN stage and write
\((M,D,\eta,\epsilon)\) for the matrices of the Hall product,
coproduct, unit, and counit.  The finite Hall bialgebra defect tuple is
\[
\mathfrak O_{\mathrm{Hbialg},R}
=\bigl(
\mathfrak o^{\mathrm{unit},L}_R,\mathfrak o^{\mathrm{unit},R}_R,
\mathfrak o^{\mathrm{counit},L}_R,\mathfrak o^{\mathrm{counit},R}_R,
\mathfrak o^{\mathrm{assoc}}_R,\mathfrak o^{\mathrm{coassoc}}_R,
\mathfrak o^{\mathrm{bialg}}_R,\mathfrak o^{\mathrm{prim}}_R
\bigr),
\]
where
\[
\mathfrak o^{\mathrm{assoc}}_R
=m_R(m_R\otimes1)-m_R(1\otimes m_R),\qquad
\mathfrak o^{\mathrm{coassoc}}_R
=(\Delta_R\otimes1)\Delta_R-(1\otimes\Delta_R)\Delta_R,
\]
\[
\mathfrak o^{\mathrm{bialg}}_R
=\Delta_Rm_R-(m_R\otimes m_R)(1\otimes\tau\otimes1)
(\Delta_R\otimes\Delta_R),
\]
with \(\tau\) the Koszul sign braiding, and
\[
\mathfrak o^{\mathrm{prim}}_R(x,y)
=\widetilde\Delta_R
\bigl(m_R(x,y)-(-1)^{|x||y|}m_R(y,x)\bigr)
\]
for \(x,y\in\ker\widetilde\Delta_R\cap\ker\epsilon_R\).  The unit and
counit components are the four usual left/right defects
\[
m_R(\eta_R\otimes1)-1,\quad m_R(1\otimes\eta_R)-1,\quad
(\epsilon_R\otimes1)\Delta_R-1,\quad
(1\otimes\epsilon_R)\Delta_R-1.
\]
\end{definition}

\begin{proposition}[Hall primitives are formed only after the bialgebra
defects vanish; \textit{proved}]
\label{prop:hall-bialgebra-defects-primitive-closure}
At a finite HN stage, the matrices
\((M,D,\eta,\epsilon)\) define a counital graded bialgebra whose
supercommutator restricts to a bracket on the primitive subspace if and
only if the first seven components of
\(\mathfrak O_{\mathrm{Hbialg},R}\) vanish; then
\(\mathfrak o^{\mathrm{prim}}_R=0\) follows.  Conversely, a finite
source proof table in which any component of
\(\mathfrak O_{\mathrm{Hbialg},R}\) is absent or non-zero does not
define the primitive Hall Lie superalgebra used in
Theorem~\ref{thm:primitive-recognition}.
\end{proposition}

\begin{proof}
The unit, counit, associativity, coassociativity, and bialgebra
compatibility equations are exactly the first seven displayed defects.
When they vanish, the standard calculation in a graded bialgebra gives,
for primitive \(x,y\),
\[
\Delta_R([x,y])
=[x,y]\otimes1+1\otimes[x,y],
\qquad
\epsilon_R([x,y])=0,
\]
because \(\Delta_R\) is multiplicative and \(\epsilon_R\) is an algebra
homomorphism.  Hence the supercommutator lies in the primitive subspace,
which is the vanishing of \(\mathfrak o^{\mathrm{prim}}_R\).  If one of
the first seven equations is missing or non-zero, the object is not a
counital bialgebra on the chosen finite stage, and the displayed
primitive-bracket calculation is not available.
\end{proof}

The packet
\(\texttt{certificates/hall/finite\_hall\_primitive\_commutator\_closure}\)
verifies
\(\texttt{FINITE\_HALL\_PRIMITIVE\_COMMUTATOR\_CLOSURE\_OBSTRUCTION\_VERIFIED}\):
it imports the row \(353\) primitive-subspace definition, the row
\(350\) bialgebra-compatibility packet, and the compact Hall source
obstruction ledger; it records
Proposition~\ref{prop:hall-bialgebra-defects-primitive-closure} as the
row \(354\) relative proof; and it keeps the product matrix, coproduct
matrix, unit and counit rows, primitive-kernel basis, bialgebra
identity rows, primitive-closure defect rows, and supercommutator
matrix rows open.  Hence the proposition proves the closure theorem for
supplied finite bialgebra data; the present compact \(K3\times E\)
source still has no populated primitive Hall bracket.

\begin{proposition}[The primitive Hall supercommutator satisfies
graded Jacobi; \textit{proved}]
\label{prop:finite-hall-graded-jacobi}
Let \(A\) be an associative \(\mathbb Z/2\)-graded algebra and let
\[
[x,y]=xy-(-1)^{|x||y|}yx
\]
for homogeneous \(x,y\in A\).  Then the bracket is graded
skew-symmetric and satisfies the graded Jacobi identity
\[
(-1)^{|x||z|}[x,[y,z]]
+(-1)^{|y||x|}[y,[z,x]]
+(-1)^{|z||y|}[z,[x,y]]=0 .
\]
Consequently, if the finite Hall bialgebra data of
Proposition~\ref{prop:hall-bialgebra-defects-primitive-closure} are
supplied so that \(P_R\) is closed under the supercommutator, then
\(P_R\) is a finite graded Lie algebra object under the Hall bracket.
\end{proposition}

\begin{proof}
Graded skew-symmetry is immediate from the definition.  For Jacobi,
expand the three displayed summands in the associative algebra \(A\).
The twelve monomials \(xyz,xzy,yxz,yzx,zxy,zyx\), each with its
Koszul sign, occur in cancelling pairs after using only associativity
and the parity rule \(|xy|=|x|+|y|\).  Thus the supercommutator on
any associative graded algebra is a graded Lie bracket.  If
\(x,y,z\in P_R\), Proposition~\ref{prop:hall-bialgebra-defects-primitive-closure}
keeps each inner and outer commutator in \(P_R\), so the same identity
holds in the primitive Hall subspace.
\end{proof}

The packet \(\texttt{certificates/hall/finite\_hall\_graded\_jacobi}\)
verifies
\(\texttt{FINITE\_HALL\_GRADED\_JACOBI\_OBSTRUCTION\_VERIFIED}\): it
imports the row \(354\) primitive-commutator closure packet, the row
\(350\) bialgebra-compatibility packet, and the compact Hall source
obstruction ledger; it records
Proposition~\ref{prop:finite-hall-graded-jacobi}; and it keeps the
source product matrix, primitive bracket matrix, associativity rows,
primitive-closure rows, Jacobi defect rows, parity rows, and
normal-ordered bracket matrices open.  Hence row \(355\) proves
graded Jacobi for supplied associative finite Hall data; it does not
populate a compact \(K3\times E\) primitive bracket.

\begin{proposition}[The Hall bracket preserves normal-ordered Gram
degree; \textit{proved for supplied homogeneous data}]
\label{prop:finite-hall-normal-ordered-gram-degree}
Let \(P_R=\bigoplus_{\widehat c\in\widehat\Gamma_R}P_{R,\widehat c}\)
be a finite primitive Hall stage whose bracket is obtained from a
homogeneous Hall product:
\[
[P_{R,\widehat c},P_{R,\widehat c'}]\subset
P_{R,\widehat c+\widehat c'} .
\]
Let
\[
\overline\Pi_X:\widehat\Gamma_R\longrightarrow\Gamma_R^\Pi
\]
be the additive normal-ordered Gram homomorphism.  Define the
descended fibre-summed pieces
\[
P^\Pi_{R,\gamma}=
\bigoplus_{\overline\Pi_X(\widehat c)=\gamma}P_{R,\widehat c}.
\]
Then
\[
[P^\Pi_{R,\gamma},P^\Pi_{R,\gamma'}]\subset
P^\Pi_{R,\gamma+\gamma'} .
\]
Equivalently, every non-zero finite bracket row has zero
normal-ordered Gram-degree defect
\[
\overline\Pi_X(\widehat e)
-\overline\Pi_X(\widehat c)
-\overline\Pi_X(\widehat c')=0
\]
for its source degrees \(\widehat c,\widehat c'\) and target degree
\(\widehat e=\widehat c+\widehat c'\).
\end{proposition}

\begin{proof}
The Hall product is homogeneous by hypothesis, so both ordered products
\(xy\) and \(yx\) in the supercommutator have source degree
\(\widehat c+\widehat c'\).  Subtracting them with the Koszul sign
cannot change the charge degree.  Applying the additive
normal-ordered homomorphism gives
\[
\overline\Pi_X(\widehat c+\widehat c')
=\overline\Pi_X(\widehat c)+\overline\Pi_X(\widehat c').
\]
If several source charges have the same descended Gram label, the
definition of \(P^\Pi_{R,\gamma}\) is the direct sum over that fibre;
the same additivity sends the whole bracket of the two fibres into the
fibre over \(\gamma+\gamma'\).  This proves the displayed inclusion.
\end{proof}

The packet
\(\texttt{certificates/hall/finite\_hall\_normal\_ordered\_gram\_degree}\)
verifies
\(\texttt{FINITE\_HALL\_NORMAL\_ORDERED\_GRAM\_DEGREE\_OBSTRUCTION\_VERIFIED}\):
it imports the row \(355\) graded-Jacobi packet, the finite
\(\texttt{hall\_pairing\_pushforward\_compatibility}\) packet, the
empty \(K3\times E\) charge-window fixture, and the compact Hall source
obstruction ledger; it records
Proposition~\ref{prop:finite-hall-normal-ordered-gram-degree}; and it
keeps the current source degree rows, normal-ordered lift rows, Hall
product rows, primitive bracket rows, and bracket-degree defect rows
open.  Hence row \(356\) proves the Gram-degree statement for supplied
homogeneous normal-ordered Hall data; it does not populate the compact
\(K3\times E\) bracket or replace the row \(357\) parity check.

\begin{proposition}[The Hall bracket respects parity;
\textit{proved for supplied homogeneous data}]
\label{prop:finite-hall-bracket-parity}
Let \(A=A_{\bar0}\oplus A_{\bar1}\) be a
\(\mathbb Z/2\)-graded associative algebra whose product is even:
\[
A_pA_q\subset A_{p+q}\qquad(p,q\in\mathbb Z/2).
\]
For homogeneous \(x\in A_p\) and \(y\in A_q\), set
\[
[x,y]=xy-(-1)^{pq}yx .
\]
Then \([x,y]\in A_{p+q}\).  Consequently, if the finite primitive Hall
stage \(P_R\) is supplied as a parity-graded subspace closed under this
supercommutator, then
\[
[P_{R,\bar p},P_{R,\bar q}]\subset P_{R,\bar p+\bar q}.
\]
After normal-ordered Gram descent, the same statement holds fibrewise:
\[
[P^\Pi_{R,\gamma,\bar p},P^\Pi_{R,\gamma',\bar q}]
\subset P^\Pi_{R,\gamma+\gamma',\bar p+\bar q}
\]
whenever the row \(356\) degree hypotheses and the parity block rows
are supplied.
\end{proposition}

\begin{proof}
Since the multiplication is even, both \(xy\) and \(yx\) lie in the
same parity summand \(A_{p+q}\).  Their signed difference therefore
also lies in \(A_{p+q}\).  Restricting to \(P_R\) uses only the supplied
closure of \(P_R\) under the supercommutator and the supplied parity
decomposition of \(P_R\).  The normal-ordered statement is obtained by
combining this parity calculation with
Proposition~\ref{prop:finite-hall-normal-ordered-gram-degree}: the
Gram label adds independently of the parity label, while the parity
label adds in \(\mathbb Z/2\).
\end{proof}

The packet \(\texttt{certificates/hall/finite\_hall\_bracket\_parity}\)
verifies
\(\texttt{FINITE\_HALL\_BRACKET\_PARITY\_OBSTRUCTION\_VERIFIED}\): it
imports the row \(355\) graded-Jacobi packet, the row \(356\)
normal-ordered Gram-degree packet, the finite
\(\texttt{hall\_pairing\_pushforward\_compatibility}\) bracket rows,
the finite \(\texttt{parity\_pushforward}\) packet, and the compact
Hall source obstruction ledger; it records
Proposition~\ref{prop:finite-hall-bracket-parity}; and it keeps the
current source parity blocks, homogeneous product rows, primitive
bracket rows, and bracket-parity defect rows open.  Hence row \(357\)
proves parity compatibility for supplied homogeneous finite Hall data;
it does not prove full \(K3\times E\) source parity or the GN--Kac
parity comparison required later in primitive recognition.

\begin{definition}[Compact geometric Hall bialgebra datum]
\label{def:compact-geometric-hall-bialgebra-datum}
At a finite HN height \(R\), a compact geometric Hall bialgebra datum is
the following finite correspondence package on the cosection-reduced
d-critical heart:
\begin{enumerate}
\item finite-type retained substacks
\(\mathfrak M^{\mathrm{red}}_{R,\widehat c}\) for all
\(\widehat c\in\widehat\Gamma_R\), with finite residual inertia,
reduced vanishing-cycle sheaves \(\Phi^{\mathrm{red}}_{R,\widehat c}\),
and Joyce--Upmeier orientation lines \(o_{R,\widehat c}\);
\item compactified extension correspondences
\[
\mathfrak M^{\mathrm{red}}_{R,\widehat c}\times
\mathfrak M^{\mathrm{red}}_{R,\widehat c'}
\xleftarrow{\ p\ }
\overline{\mathfrak E}^{\mathrm{red}}_{R,\widehat c,\widehat c'}
\xrightarrow{\ q\ }
\mathfrak M^{\mathrm{red}}_{R,\widehat c+\widehat c'}
\]
with admissible \(p^*,q_!\) for the product and admissible
\(q^*,p_!\) for the collision coproduct, in either a retained
compactified flag-Quot model or an explicitly supplied compact-support
six-functor model;
\item reduced Thom--Sebastiani isomorphisms
\(\operatorname{TS}^{\mathrm{red}}_{\widehat c,\widehat c'}\) and their
inverse transports, compatible with the multiplicative orientation
cocycle;
\item vacuum correspondences for the zero object, giving the unit
\(\eta_R\) and counit \(\epsilon_R\);
\item three-step and four-step retained flag stacks carrying the
associativity, coassociativity, and bialgebra-compatibility diagrams,
with base-change, projection-formula, Koszul sign, and
Thom--Sebastiani coherence witnesses;
\item reduced \(E\)-quotient and HN-transition compatibilities for all
the preceding correspondences.
\end{enumerate}
The product and coproduct attached to this datum are
\[
m_R=q_!\operatorname{TS}^{\mathrm{red}}p^*,\qquad
\Delta_R=p_!(\operatorname{TS}^{\mathrm{red}})^{-1}q^*.
\]
\end{definition}

\begin{definition}[Hall-product-preserving transition datum]
\label{def:transition-hall-product-preservation}
Let \(R'\le R\).  Suppose the finite-moduli transition geometry of
Definition~\ref{def:finite-moduli-transition-geometry}, the
orientation-preserving transition datum of
Definition~\ref{def:transition-orientation-preservation}, and the
vanishing-cycle-preserving transition datum of
Definition~\ref{def:transition-vanishing-cycle-preservation} are
supplied.  A Hall-product-preserving transition datum consists of:
\begin{enumerate}
\item a transition cube for the retained compactified extension
correspondences
\[
\mathfrak M_R\times\mathfrak M_R
\xleftarrow{\ p_R\ }\overline{\mathfrak E}_R
\xrightarrow{\ q_R\ }\mathfrak M_R,
\qquad
\mathfrak M_{R'}\times\mathfrak M_{R'}
\xleftarrow{\ p_{R'}\ }\overline{\mathfrak E}_{R'}
\xrightarrow{\ q_{R'}\ }\mathfrak M_{R'},
\]
whose input and output squares commute with the stage transitions;
\item transport of the coefficient system
\(\Phi_R^{\mathrm{red}}\otimes o_R\) on the extension graph, including
the Thom--Sebastiani product isomorphism
\[
\operatorname{TS}^{\mathrm{red}}_R:
p_R^*(\Phi_R^{\mathrm{red}}\boxtimes\Phi_R^{\mathrm{red}})
\xrightarrow{\sim}q_R^*\Phi_R^{\mathrm{red}};
\]
\item proper base-change, projection-formula, and compact-support
pushforward isomorphisms comparing the two pull-push functors
\[
q_{R,!}\operatorname{TS}^{\mathrm{red}}_Rp_R^*
\quad\text{and}\quad
q_{R',!}\operatorname{TS}^{\mathrm{red}}_{R'}p_{R'}^*
\]
after applying the transition maps;
\item finite product matrices \(M_R,M_{R'}\) in compact source bases
and transition matrices \(T_R,T_{R'}\) satisfying the product
intertwining identity
\[
T_{\alpha+\beta}\,M_R^{\alpha,\beta}
=M_{R'}^{\alpha,\beta}(T_\alpha\otimes T_\beta)
\]
in every retained degree and parity block;
\item compatibility with the associativity square on the retained
two-step flag stack;
\item zero correspondence, coefficient-system, base-change,
projection-formula, product-intertwining, associativity-transition, and
composition defect ranks, together with \(R^1\varprojlim=0\) for the
Hall product maps.
\end{enumerate}
Scalar traces, denominator products, signed multiplicities, target BKM
products, Pfaffian products, orientation transport alone, and
vanishing-cycle transport alone are not Hall-product transition data.
\end{definition}

\begin{proposition}[Transition preservation of compact Hall products;
\textit{proved}]
\label{prop:transition-hall-product-preservation-criterion}
Given a Hall-product-preserving transition datum of
Definition~\ref{def:transition-hall-product-preservation}, the
finite-stage transition \(R\to R'\) preserves the compact Hall product:
on every retained finite degree block,
\[
T_R\circ m_R=m_{R'}\circ(T_R\otimes T_R).
\]
\end{proposition}

\begin{proof}
The transition cube identifies the input pullbacks \(p_R^*\) and
\(p_{R'}^*\) after passage to the transition graph.  The coefficient
transport identifies \(\Phi^{\mathrm{red}}\otimes o\) and the
Thom--Sebastiani product isomorphisms on the two sides.  Proper base
change, the projection formula, and compact-support pushforward
compatibility then identify the two exceptional pushforwards along
\(q_R\) and \(q_{R'}\).  Hence the two pull-push functors defining
\(m_R\) and \(m_{R'}\) commute with transition.  In finite compact
source bases this is exactly the displayed matrix identity.  The
associativity-transition and \(R^1\varprojlim\) rows give strict
composition and inverse-limit exactness for the product maps.
\end{proof}

The packet
\(\texttt{certificates/hall/transition\_hall\_product\_preservation}\)
records the present state of this datum.  Its verifier status is
\texttt{TRANSITION\_HALL\_PRODUCT\_PRESERVATION\_OBSTRUCTION\_VERIFIED}:
the transition-geometry input, orientation-transition input,
vanishing-cycle transition input, source product matrices, extension
correspondence transition squares, input and output product-stack
squares, coefficient-system transport rows, Thom--Sebastiani product
transport rows, proper base-change rows, projection-formula rows,
compact-support pushforward rows, product-matrix intertwining rows,
associativity-transition rows, and Hall-product Mittag--Leffler rows
are not yet supplied.  Hence transition preservation of compact Hall
products remains open.

\begin{definition}[Hall-coproduct-preserving transition datum]
\label{def:transition-hall-coproduct-preservation}
Let \(R'\le R\).  Suppose the finite-moduli transition geometry of
Definition~\ref{def:finite-moduli-transition-geometry}, the
orientation-preserving transition datum of
Definition~\ref{def:transition-orientation-preservation}, and the
vanishing-cycle-preserving transition datum of
Definition~\ref{def:transition-vanishing-cycle-preservation} are
supplied.  A Hall-coproduct-preserving transition datum consists of:
\begin{enumerate}
\item a transition cube for the same retained compactified extension
correspondences, read in the splitting direction
\[
\mathfrak M_R
\xleftarrow{\ q_R\ }\overline{\mathfrak E}_R
\xrightarrow{\ p_R\ }\mathfrak M_R\times\mathfrak M_R,
\qquad
\mathfrak M_{R'}
\xleftarrow{\ q_{R'}\ }\overline{\mathfrak E}_{R'}
\xrightarrow{\ p_{R'}\ }\mathfrak M_{R'}\times\mathfrak M_{R'},
\]
whose \(q\)-source square and \(p\)-two-output square commute with the
stage transitions;
\item transport of the coefficient system
\(\Phi_R^{\mathrm{red}}\otimes o_R\) on the splitting graph, including
the inverse Thom--Sebastiani coproduct isomorphism
\[
(\operatorname{TS}^{\mathrm{red}}_R)^{-1}:
q_R^*\Phi_R^{\mathrm{red}}
\xrightarrow{\sim}
p_R^*(\Phi_R^{\mathrm{red}}\boxtimes\Phi_R^{\mathrm{red}});
\]
\item proper base-change, projection-formula, and compact-support
pushforward isomorphisms comparing the two pull-push functors
\[
p_{R,!}(\operatorname{TS}^{\mathrm{red}}_R)^{-1}q_R^*
\quad\text{and}\quad
p_{R',!}(\operatorname{TS}^{\mathrm{red}}_{R'})^{-1}q_{R'}^*
\]
after applying the transition maps;
\item finite coproduct matrices \(D_R,D_{R'}\) in compact source bases,
counit matrices \(\epsilon_R,\epsilon_{R'}\), and transition matrices
\(T_\gamma\) satisfying the coproduct intertwining identity
\[
(T_\mu\otimes T_\nu)\,D_R^{\rho;\mu,\nu}
=D_{R'}^{\rho;\mu,\nu}\,T_\rho
\]
in every retained degree and parity block with \(\rho=\mu+\nu\);
\item compatibility with the coassociativity square on the retained
two-step flag stack and with the zero-object counit square;
\item zero correspondence, coefficient-system, base-change,
projection-formula, coproduct-intertwining, coassociativity-transition,
counit-transition, and composition defect ranks, together with
\(R^1\varprojlim=0\) for the Hall coproduct maps.
\end{enumerate}
Scalar traces, denominator products, signed multiplicities, target BKM
coalgebra data, Pfaffian operations, bar-coalgebra deconcatenation,
Hall-product preservation, orientation transport alone, and
vanishing-cycle transport alone are not Hall-coproduct transition data.
\end{definition}

\begin{proposition}[Transition preservation of compact Hall coproducts;
\textit{proved}]
\label{prop:transition-hall-coproduct-preservation-criterion}
Given a Hall-coproduct-preserving transition datum of
Definition~\ref{def:transition-hall-coproduct-preservation}, the
finite-stage transition \(R\to R'\) preserves the compact Hall
coproduct: on every retained finite degree block with \(\rho=\mu+\nu\),
\[
(T_\mu\otimes T_\nu)\circ\Delta_R^{\rho;\mu,\nu}
=\Delta_{R'}^{\rho;\mu,\nu}\circ T_\rho .
\]
\end{proposition}

\begin{proof}
The transition cube identifies the source pullbacks \(q_R^*\) and
\(q_{R'}^*\) after passage to the transition graph.  The coefficient
transport identifies \(\Phi^{\mathrm{red}}\otimes o\) and the inverse
Thom--Sebastiani isomorphisms on the two sides.  Proper base change,
the projection formula, and compact-support pushforward compatibility
then identify the two exceptional pushforwards along \(p_R\) and
\(p_{R'}\).  Hence the two pull-push functors defining \(\Delta_R\) and
\(\Delta_{R'}\) commute with transition.  In finite compact source
bases this is exactly the displayed matrix identity.  The
coassociativity-transition, counit-transition, and \(R^1\varprojlim\)
rows give strict composition, counital compatibility, and inverse-limit
exactness for the coproduct maps.
\end{proof}

The packet
\(\texttt{certificates/hall/transition\_hall\_coproduct\_preservation}\)
records the present state of this datum.  Its verifier status is
\texttt{TRANSITION\_HALL\_COPRODUCT\_PRESERVATION\_OBSTRUCTION\_VERIFIED}:
the transition-geometry input, orientation-transition input,
vanishing-cycle transition input, source coproduct matrices, source
counit maps, splitting correspondence transition squares, source object
and target product stack squares, coefficient-system transport rows,
inverse Thom--Sebastiani coproduct transport rows, proper base-change
rows, projection-formula rows, compact-support pushforward rows,
coproduct-matrix intertwining rows, coassociativity-transition rows,
counit-transition rows, and Hall-coproduct Mittag--Leffler rows are not
yet supplied.  Hence transition preservation of compact Hall coproducts
remains open.

\begin{definition}[Primitive-subspace-preserving transition datum]
\label{def:transition-primitive-subspace-preservation}
Let \(R'\le R\), and suppose the finite Hall stages carry counital
coproducts \((\Delta_R,\epsilon_R)\) and
\((\Delta_{R'},\epsilon_{R'})\).  Put
\[
\widetilde\Delta_R
=\Delta_R-\eta_R\epsilon_R\otimes\mathrm{id}
-\mathrm{id}\otimes\eta_R\epsilon_R,
\qquad
P_R=\ker\widetilde\Delta_R\cap\ker\epsilon_R .
\]
A primitive-subspace-preserving transition datum consists of:
\begin{enumerate}
\item finite compact-source matrices for
\(\Delta_R,\Delta_{R'},\epsilon_R,\epsilon_{R'}\), the associated
reduced coproducts \(\widetilde\Delta_R,\widetilde\Delta_{R'}\), and
the ambient Hall transition \(T_R:\mathcal H_R\to\mathcal H_{R'}\);
\item source and target primitive kernel rows, including inclusions
\(\iota_R:P_R\hookrightarrow\mathcal H_R\),
\(\iota_{R'}:P_{R'}\hookrightarrow\mathcal H_{R'}\) and finite
primitive projection matrices;
\item the reduced-coproduct and counit intertwining identities
\[
\widetilde\Delta_{R'}T_R=(T_R\otimes T_R)\widetilde\Delta_R,\qquad
\epsilon_{R'}T_R=\epsilon_R;
\]
\item a restricted transition matrix
\[
T_R^{\Prim}:P_R\longrightarrow P_{R'}
\]
satisfying \(\iota_{R'}T_R^{\Prim}=T_R\iota_R\), together with zero
primitive-image, projection-commutator, and composition defect ranks.
\end{enumerate}
Hall product preservation, scalar trace data, denominator products,
target root spaces, bar-coalgebra primitives, normal-ordered degree
maps, and primitive ranks without inclusions are not
primitive-subspace transition data.
\end{definition}

\begin{proposition}[Transition preservation of primitive Hall
subspaces; \textit{proved}]
\label{prop:transition-primitive-subspace-preservation-criterion}
Given a primitive-subspace-preserving transition datum of
Definition~\ref{def:transition-primitive-subspace-preservation}, the
ambient Hall transition carries primitives to primitives:
\[
T_R(P_R)\subseteq P_{R'}.
\]
Equivalently, \(T_R|_{P_R}\) factors through the displayed
\(T_R^{\Prim}:P_R\to P_{R'}\).
\end{proposition}

\begin{proof}
Let \(x\in P_R\).  Then
\(\widetilde\Delta_Rx=0\) and \(\epsilon_Rx=0\).  The two
intertwining identities give
\[
\widetilde\Delta_{R'}(T_Rx)
=(T_R\otimes T_R)\widetilde\Delta_Rx=0,\qquad
\epsilon_{R'}(T_Rx)=\epsilon_Rx=0.
\]
Thus \(T_Rx\in\ker\widetilde\Delta_{R'}\cap\ker\epsilon_{R'}=P_{R'}\).
The finite restricted matrix \(T_R^{\Prim}\) is the same statement in
the chosen primitive bases.  The projection and composition defect
rows assert that this factorization is compatible with the finite
primitive summands and with successive stage transitions.
\end{proof}

The packet
\(\texttt{certificates/hall/transition\_primitive\_subspace\_preservation}\)
records the present state of this datum.  Its verifier status is
\texttt{TRANSITION\_PRIMITIVE\_SUBSPACE\_PRESERVATION\_OBSTRUCTION\_VERIFIED}:
the source and target coproduct matrices, source and target counit
maps, reduced-coproduct matrices, primitive kernel rows, ambient Hall
transition matrices, Hall-coproduct transition input, counit-transition
compatibility, kernel-intertwining rows, primitive projection matrices,
restricted primitive transition matrices, and primitive-transition
composition rows are not yet supplied.  Hence transition preservation
of primitive Hall subspaces remains open.

\begin{proposition}[Geometric Hall data give the finite bialgebra
defect rows; \textit{proved relative to the datum}]
\label{prop:geometric-hall-data-give-bialgebra-defects}
Let the compact geometric Hall bialgebra datum of
Definition~\ref{def:compact-geometric-hall-bialgebra-datum} be given at
height \(R\).  Then the matrices of
\((m_R,\Delta_R,\eta_R,\epsilon_R)\) satisfy the unit, counit,
associativity, coassociativity, and bialgebra-compatibility equations
of Definition~\ref{def:finite-hall-bialgebra-defect-tuple}.  Hence the
first seven components of
\(\mathfrak O_{\mathrm{Hbialg},R}\) vanish, and
Proposition~\ref{prop:hall-bialgebra-defects-primitive-closure}
forms the primitive Hall bracket.  Conversely, finite matrices
\((M,D,\eta,\epsilon)\) whose bialgebra defect ranks vanish but which
are not obtained from such a compact geometric Hall bialgebra datum are
not compact \(K3\times E\) source rows; they are only an abstract finite
bialgebra.
\end{proposition}

\begin{proof}
The left and right unit identities are the pull-push identities for the
vacuum extension correspondences.  The counit identities are their
collision analogues after applying the augmentation to the zero-object
factor.  Associativity of \(m_R\) is the equality of the two pull-push
composites on the retained three-step extension flag stack; the
base-change isomorphism, projection formula, and triple
Thom--Sebastiani coherence identify the two composites.  The same
argument applied to the reversed correspondence gives coassociativity
of \(\Delta_R\).

Bialgebra compatibility is the comparison between first multiplying and
then splitting an extension, and first splitting the two inputs and then
multiplying the four resulting pieces.  The retained four-step
flag-of-flags stack is the common refinement of these two operations.
Base change on this common refinement gives the equality of the two
pull-push functors; the middle interchange contributes exactly the
Koszul sign braiding \(\tau\), and the multiplicative
Thom--Sebastiani orientation transport identifies the coefficient
systems.  This is the matrix identity
\[
\Delta_Rm_R
=(m_R\otimes m_R)(1\otimes\tau\otimes1)
(\Delta_R\otimes\Delta_R).
\]
Thus all first seven defect matrices have rank zero.  The primitive
closure is then the algebraic calculation of
Proposition~\ref{prop:hall-bialgebra-defects-primitive-closure}.

The converse sentence is the Beilinson cut: an abstract finite
bialgebra with the same matrices is not a \(K3\times E\) source object
unless its entries are the pull-push matrices of the retained compact
correspondences with the stated orientation, quotient, and transition
data.
\end{proof}

\begin{lemma}[Primitive bracket lifting; \textit{conditional on the
Frobenius defect $\mathfrak o_F=0$}]
\label{lem:primitive-bracket-lifting}
Let $\langle\cdot,\cdot\rangle_X$ be the protected Hopf pairing on
$\widetilde P_X$, non-degenerate on the radical quotient.  The
trilinear Frobenius defect
\[
\mathfrak o_F(x,y,z):=\langle[x,y],z\rangle
+(-1)^{|x|\cdot|y|}\langle y,[x,z]\rangle,
\quad
x,y,z\in\widetilde P_X,\quad c_x+c_y+c_z=0,
\]
viewed as a class in
$C^3_{\mathrm{cyc}}(\widetilde P_X;k)$ after the reduced $E$-quotient
and radical quotient, vanishes precisely when the bracket is
graded-symmetric for the Hopf pairing.  Vanishing is the cyclic-CY-3
Serre-functor identity of
Theorem~\ref{thm:hall-product}.  Non-degeneracy of the
Hopf pairing alone does \emph{not} force $\mathfrak o_F=0$; it only
detects the vanishing once supplied.
\end{lemma}

\begin{definition}[Finite positive-negative Hall pairing]
\label{def:finite-hall-positive-negative-pairing}
Let
\[
P_R^\Pi=\bigoplus_{\gamma\in\Gamma_R^\Pi,\ \bar p\in\mathbb Z/2}
P_{R,\gamma,\bar p}^\Pi
\]
be a finite normal-ordered primitive Hall stage, equipped with the
descended product \(m_R^\Theta\) and a degree-zero trace
\[
\operatorname{tr}_{0,R}:
H_{R,0,\bar0}^{\Pi}\longrightarrow k .
\]
A finite positive-negative Hall pairing is the collection of bilinear
forms
\[
\langle\cdot,\cdot\rangle_{R,\gamma,\bar p}:
P_{R,\gamma,\bar p}^\Pi\otimes
P_{R,-\gamma,\bar p}^\Pi\longrightarrow k
\]
defined, in homogeneous finite bases, by the matrix
\[
G_{\gamma,\bar p;ab}
=\operatorname{tr}_{0,R}
\bigl(m_R^\Theta(e_{\gamma,\bar p,a},e_{-\gamma,\bar p,b})\bigr).
\]
Equivalently, \(G_{\gamma,\bar p;ab}\) is the pull-push coefficient of
the compact positive-negative pairing correspondence
\[
\mathfrak P^{\mathrm{red}}_{R,\gamma,-\gamma;0}
\longrightarrow
\mathfrak M^{\mathrm{red}}_{R,0}
\]
with reduced Thom--Sebastiani and orientation transport, followed by
\(\operatorname{tr}_{0,R}\).  A finite row of type \(G\) consists of
the degree \(\gamma\), the parity \(\bar p\), the positive and negative
primitive basis labels, the scalar coefficient, the pairing
correspondence, and the geometric source reference.

This is a definition of the source matrix \(G\).  It does not assert
that \(G\) is homogeneous, graded-symmetric, invariant under the Hall
bracket, compatible with the coproduct, or non-degenerate after the
radical quotient.  Those are the separate defect rows of
Definition~\ref{def:finite-hopf-pairing-identity-tuple}.
\end{definition}

The packet
\(\texttt{certificates/hall/finite\_hall\_positive\_negative\_pairing}\)
verifies
\(\texttt{FINITE\_HALL\_POSITIVE\_NEGATIVE\_PAIRING\_OBSTRUCTION\_VERIFIED}\):
it imports the formal degree-zero pairing rows of
\(\texttt{hall\_pairing\_pushforward\_compatibility}\), the row \(356\)
normal-ordered Gram-degree packet, the row \(357\) parity packet, and
the compact Hall source obstruction ledger; it records
Definition~\ref{def:finite-hall-positive-negative-pairing}; and it
keeps the current source \(G\)-entries, degree-zero trace row,
positive-negative pairing correspondences, orientation transport rows,
and Hopf-pairing identity rows open.  Hence row \(358\) defines the
finite positive-negative Hall pairing matrix; it does not prove any of
the row \(359\)--\(362\) pairing identities.

\begin{proposition}[The finite Hall pairing is homogeneous;
\textit{proved for supplied source data}]
\label{prop:finite-hall-pairing-homogeneity}
Assume the finite positive-negative Hall pairing of
Definition~\ref{def:finite-hall-positive-negative-pairing} is supplied
from a homogeneous normal-ordered Hall product and a degree-zero trace.
Extend it by zero to all pairs
\[
P_{R,\gamma,\bar p}^{\Pi}\otimes P_{R,\gamma',\bar q}^{\Pi}.
\]
Then
\[
\langle
P_{R,\gamma,\bar p}^{\Pi},
P_{R,\gamma',\bar q}^{\Pi}
\rangle_R=0
\quad\text{unless}\quad
\gamma+\gamma'=0\ \text{and}\ \bar p=\bar q .
\]
Equivalently, every non-zero finite pairing row has total
normal-ordered Gram degree \(0\) and belongs to one parity block.
\end{proposition}

\begin{proof}
For homogeneous inputs \(x\in P_{R,\gamma,\bar p}^{\Pi}\) and
\(y\in P_{R,\gamma',\bar q}^{\Pi}\), the row \(356\) degree
calculation gives
\[
m_R^\Theta(x,y)\in H_{R,\gamma+\gamma',\bar p+\bar q}^{\Pi}.
\]
The trace in Definition~\ref{def:finite-hall-positive-negative-pairing}
is defined on the degree-zero block.  Hence
\(\operatorname{tr}_{0,R}(m_R^\Theta(x,y))\) can contribute only when
\(\gamma+\gamma'=0\).  The positive-negative pairing matrix \(G\) is
defined blockwise on
\(P_{R,\gamma,\bar p}^{\Pi}\otimes P_{R,-\gamma,\bar p}^{\Pi}\); its
extension by zero kills the off-parity blocks.  Thus every non-zero
entry has total normal-ordered Gram degree \(0\) and one fixed parity
label.
\end{proof}

The packet
\(\texttt{certificates/hall/finite\_hall\_pairing\_homogeneity}\)
verifies
\(\texttt{FINITE\_HALL\_PAIRING\_HOMOGENEITY\_OBSTRUCTION\_VERIFIED}\):
it imports the row \(358\) positive-negative pairing definition, the
row \(356\) normal-ordered Gram-degree packet, the formal degree-zero
pairing rows of
\(\texttt{hall\_pairing\_pushforward\_compatibility}\), and the compact
Hall source obstruction ledger; it records
Proposition~\ref{prop:finite-hall-pairing-homogeneity}; and it keeps
the current source \(G\)-entries, trace row, pairing correspondence
rows, and homogeneity defect rows open.  Hence row \(359\) proves the
homogeneity criterion for supplied finite source data; it does not
prove graded supersymmetry, Frobenius invariance, coproduct
adjointness, or quotient non-degeneracy.

\begin{proposition}[The finite Hall pairing is supersymmetric;
\textit{proved for supplied signed-transpose rows}]
\label{prop:finite-hall-pairing-supersymmetry}
Assume the finite positive-negative Hall pairing of
Definition~\ref{def:finite-hall-positive-negative-pairing} is
homogeneous in the sense of
Proposition~\ref{prop:finite-hall-pairing-homogeneity}.  Suppose that
for every retained degree \(\gamma\), parity \(\bar p\), and pair of
finite basis labels \(a,b\), the signed-transpose row
\[
G_{\gamma,\bar p;ab}
-(-1)^{\bar p}G_{-\gamma,\bar p;ba}=0
\]
is supplied.  Then the extended finite pairing is supersymmetric:
\[
\langle x,y\rangle_R
=(-1)^{|x||y|}\langle y,x\rangle_R
\]
for all homogeneous \(x,y\in P_R^\Pi\).

Conversely, a compact-source proof of this identity requires the
source \(G\)-entries on both signs, source parity blocks,
signed-transpose defect rows, and the reduced Calabi--Yau
threefold Serre-sign and orientation-transport data.  Formal
degree-zero pairing rows, target positive-negative blocks, and scalar
traces do not supply these signs.
\end{proposition}

\begin{proof}
Write \(x\in P_{R,\gamma,\bar p}^{\Pi}\) and
\(y\in P_{R,\gamma',\bar q}^{\Pi}\).  By
Proposition~\ref{prop:finite-hall-pairing-homogeneity}, both
\(\langle x,y\rangle_R\) and \(\langle y,x\rangle_R\) vanish unless
\(\gamma'=-\gamma\) and \(\bar q=\bar p\).  On a non-zero block one
therefore has \(|x|=|y|=\bar p\), so
\((-1)^{|x||y|}=(-1)^{\bar p}\).  Expanding in the chosen finite
bases, the asserted supersymmetry equation is exactly
\[
G_{\gamma,\bar p;ab}
=(-1)^{\bar p}G_{-\gamma,\bar p;ba}.
\]
Linearity gives the identity for arbitrary homogeneous vectors.  This
is a finite signed-transpose criterion; the geometric computation of
the signs is the source Serre/orientation problem named in the final
sentence of the statement.
\end{proof}

The packet
\(\texttt{certificates/hall/finite\_hall\_pairing\_supersymmetry}\)
verifies
\(\texttt{FINITE\_HALL\_PAIRING\_SUPERSYMMETRY\_OBSTRUCTION\_VERIFIED}\):
it imports the row \(359\) homogeneity packet, the row \(358\)
positive-negative pairing definition, the reduced Calabi--Yau
threefold Serre-sign criterion packet, and the compact Hall source
obstruction ledger; it records
Proposition~\ref{prop:finite-hall-pairing-supersymmetry}; and it keeps
the current source \(G\)-entries, source parity blocks,
signed-transpose rows, Serre-sign rows, and orientation-transport rows
open.  Hence row \(360\) proves the finite supersymmetry criterion for
supplied source data; it does not prove Frobenius invariance, coproduct
adjointness, or quotient non-degeneracy.

\begin{definition}[Finite Hopf-pairing identity tuple]
\label{def:finite-hopf-pairing-identity-tuple}
Fix homogeneous finite bases on a retained primitive stage and write
\((M,B,D,G,K,Q)\) for the product, bracket, coproduct, Hopf-pairing,
radical, and quotient-splitting matrices.  The finite Hopf-pairing
identity tuple is
\[
\mathfrak O_{\mathrm{Hpair},R}
=\bigl(
\mathfrak o^{\mathrm{adj}}_R,\,
\mathfrak o^{\mathrm{Frob}}_R,\,
\mathfrak o^{\mathrm{qnd}}_R
\bigr),
\]
where
\[
\mathfrak o^{\mathrm{adj}}_R(x,y,z)
=\langle m_R^\Theta(x,y),z\rangle_R
-\langle x\otimes y,\Delta_R^\Theta z\rangle_R,
\]
\[
\mathfrak o^{\mathrm{Frob}}_R(x,y,z)
=\langle[x,y]_\Theta,z\rangle_R
+(-1)^{|x||y|}\langle y,[x,z]_\Theta\rangle_R,
\]
and \(\mathfrak o^{\mathrm{qnd}}_R\) is the finite rank defect of the
pairing induced by \(G\) on
\((P_R^\Pi/\operatorname{Rad}_R^\Pi)_\gamma\otimes
(P_R^\Pi/\operatorname{Rad}_R^\Pi)_{-\gamma}\) in every retained
parity block.  In matrix form these are the adjointness, Frobenius
cyclic, and quotient non-degeneracy rows of the finite source proof
tables.
\end{definition}

\begin{proposition}[The finite Hall pairing is invariant;
\textit{proved for supplied Frobenius rows}]
\label{prop:finite-hall-pairing-invariance}
Assume the finite primitive Hall bracket is the graded-skew
supercommutator and is parity-compatible, the finite Hall pairing is
supersymmetric in the sense of
Proposition~\ref{prop:finite-hall-pairing-supersymmetry}, and the
Frobenius cyclic component of
Definition~\ref{def:finite-hopf-pairing-identity-tuple} vanishes:
\[
\mathfrak o^{\mathrm{Frob}}_R(x,y,z)=0
\]
for every homogeneous retained triple.  Then the finite pairing is
invariant under the Hall bracket:
\[
\langle [x,y]_\Theta,z\rangle_R
=\langle x,[y,z]_\Theta\rangle_R
\]
for all homogeneous \(x,y,z\in P_R^\Pi\).

Conversely, a compact-source proof of this identity requires source
primitive bracket rows, source \(G\)-entries, parity blocks, Frobenius
defect rows, cyclic three-point pairing correspondences, and the
reduced Calabi--Yau threefold Serre-sign/orientation transport data.
The graded Jacobi identity, pairing supersymmetry, or target invariant
form does not supply these source rows.
\end{proposition}

\begin{proof}
Put \(p=|x|\), \(q=|y|\), \(r=|z|\).  Apply the assumed vanishing of
\(\mathfrak o^{\mathrm{Frob}}_R\) to the triple \((y,z,x)\):
\[
0=\langle [y,z]_\Theta,x\rangle_R
  +(-1)^{qr}\langle z,[y,x]_\Theta\rangle_R .
\]
The bracket parity row gives
\(|[y,z]_\Theta|=q+r\) and \(|[x,y]_\Theta|=p+q\).  Supersymmetry of
the pairing gives
\[
\langle [y,z]_\Theta,x\rangle_R
=(-1)^{p(q+r)}\langle x,[y,z]_\Theta\rangle_R .
\]
Using graded skew-symmetry
\([y,x]_\Theta=-(-1)^{pq}[x,y]_\Theta\) and supersymmetry again,
\[
\langle z,[y,x]_\Theta\rangle_R
=-(-1)^{pq+(p+q)r}
\langle [x,y]_\Theta,z\rangle_R .
\]
The two exponents \(p(q+r)\) and \(qr+pq+(p+q)r\) agree in
\(\mathbb Z/2\).  After removing the common sign, the displayed
vanishing becomes
\[
\langle x,[y,z]_\Theta\rangle_R
-\langle [x,y]_\Theta,z\rangle_R=0,
\]
which is the asserted invariance.
\end{proof}

The packet
\(\texttt{certificates/hall/finite\_hall\_pairing\_invariance}\)
verifies
\(\texttt{FINITE\_HALL\_PAIRING\_INVARIANCE\_OBSTRUCTION\_VERIFIED}\):
it imports the row \(360\) supersymmetry packet, the row \(357\)
bracket-parity packet, the row \(355\) graded-Jacobi packet, the
reduced Calabi--Yau threefold Serre-sign criterion packet, and the
compact Hall source obstruction ledger; it records
Proposition~\ref{prop:finite-hall-pairing-invariance}; and it keeps
the current source \(B\)-entries, \(G\)-entries, parity blocks,
Frobenius defect rows, cyclic three-point correspondences,
Serre-sign rows, and orientation-transport rows open.  Hence row
\(361\) proves the finite invariance criterion for supplied source
data; it does not prove coproduct adjointness or quotient
non-degeneracy.

\begin{proposition}[The finite Hall product and coproduct are adjoint;
\textit{proved for supplied Hopf-adjointness rows}]
\label{prop:finite-hall-pairing-coproduct-adjointness}
Assume the finite Hall product \(m_R^\Theta\), coproduct
\(\Delta_R^\Theta\), and positive-negative pairing \(G\) are supplied
in homogeneous finite bases.  Equip \(P_R^\Pi\otimes P_R^\Pi\) with the
Koszul tensor product pairing
\[
\langle x_1\otimes x_2,u_1\otimes u_2\rangle_{R\otimes R}
=(-1)^{|x_2||u_1|}
\langle x_1,u_1\rangle_R\langle x_2,u_2\rangle_R .
\]
If every adjointness component of
Definition~\ref{def:finite-hopf-pairing-identity-tuple} vanishes,
\[
\mathfrak o^{\mathrm{adj}}_R(x,y,z)=0 ,
\]
then
\[
\langle m_R^\Theta(x,y),z\rangle_R
=\langle x\otimes y,\Delta_R^\Theta z\rangle_{R\otimes R}
\]
for all homogeneous retained finite vectors \(x,y,z\) in the domain of
the Hopf-pairing tuple.

In coordinates, for homogeneous basis vectors
\(e_a,e_b,e_d\), this is the rank-zero matrix identity
\[
\sum_c M_{ab}^{c}G_{cd}
-\sum_{u,v}D_{d}^{uv}(-1)^{|e_b||e_u|}G_{au}G_{bv}=0 .
\]
A compact-source proof therefore requires source \(M\)-entries,
\(D\)-entries, \(G\)-entries, tensor-pairing sign rows, Hopf
adjointness defect rows, and the product/coproduct correspondence
orientation data.  Bialgebra compatibility, pairing invariance, target
Hopf forms, and scalar traces do not supply this adjointness matrix.
\end{proposition}

\begin{proof}
Write
\[
m_R^\Theta(e_a,e_b)=\sum_c M_{ab}^{c}e_c,\qquad
\Delta_R^\Theta(e_d)=\sum_{u,v}D_d^{uv}e_u\otimes e_v .
\]
By bilinearity of \(G\),
\[
\langle m_R^\Theta(e_a,e_b),e_d\rangle_R
=\sum_c M_{ab}^{c}G_{cd}.
\]
By the Koszul tensor product rule,
\[
\langle e_a\otimes e_b,\Delta_R^\Theta(e_d)\rangle_{R\otimes R}
=\sum_{u,v}D_d^{uv}(-1)^{|e_b||e_u|}G_{au}G_{bv}.
\]
Thus the displayed coordinate defect is exactly
\(\mathfrak o^{\mathrm{adj}}_R(e_a,e_b,e_d)\).  Its vanishing on the
finite homogeneous basis gives the identity for basis vectors, and
linearity gives it for all homogeneous vectors.
\end{proof}

The packet
\(\texttt{certificates/hall/finite\_hall\_pairing\_coproduct\_adjointness}\)
verifies
\(\texttt{FINITE\_HALL\_PAIRING\_COPRODUCT\_ADJOINTNESS\_OBSTRUCTION\_VERIFIED}\):
it imports the row \(361\) pairing-invariance packet, the row \(342\)
product-operator packet, the row \(347\) coproduct-operator packet, the
row \(350\) bialgebra-compatibility packet, and the compact Hall source
obstruction ledger; it records
Proposition~\ref{prop:finite-hall-pairing-coproduct-adjointness}; and
it keeps the current source \(M\)-entries, \(D\)-entries,
\(G\)-entries, tensor-pairing sign rows, Hopf-adjointness defect rows,
product/coproduct correspondence rows, and orientation-transport rows
open.  Hence row \(362\) proves the coproduct-adjointness criterion for
supplied finite source data; it does not prove quotient
non-degeneracy.

\begin{proposition}[The finite quotient Hopf pairing is non-degenerate;
\textit{proved for supplied quotient-rank rows}]
\label{prop:finite-hall-quotient-pairing-nondegeneracy}
Let \(G_{\gamma,\bar p}\) be a supplied homogeneous finite
positive-negative pairing matrix on
\[
P_{R,\gamma,\bar p}^{\Pi}\otimes P_{R,-\gamma,\bar p}^{\Pi}.
\]
Let \(K_{\gamma,\bar p}\subset P_{R,\gamma,\bar p}^{\Pi}\) and
\(K_{-\gamma,\bar p}\subset P_{R,-\gamma,\bar p}^{\Pi}\) be supplied
left and right radical bases for \(G_{\gamma,\bar p}\), and let
\[
Q_{\gamma,\bar p}:L_{\gamma,\bar p}\longrightarrow
P_{R,\gamma,\bar p}^{\Pi},\qquad
Q_{-\gamma,\bar p}:L_{-\gamma,\bar p}\longrightarrow
P_{R,-\gamma,\bar p}^{\Pi}
\]
be quotient splittings.  Write
\[
\overline G_{\gamma,\bar p}
=Q_{\gamma,\bar p}^{t}G_{\gamma,\bar p}Q_{-\gamma,\bar p}
\]
for the induced quotient pairing matrix.  If the finite rank rows give
\[
\operatorname{rank}\overline G_{\gamma,\bar p}
=\dim L_{\gamma,\bar p}
=\dim L_{-\gamma,\bar p}
\]
for every retained \((\gamma,\bar p)\), then the induced pairing
\[
L_{\gamma,\bar p}\otimes L_{-\gamma,\bar p}\longrightarrow k
\]
is non-degenerate.  This is the quotient non-degeneracy component
\(\mathfrak o^{\mathrm{qnd}}_R=0\) of
Definition~\ref{def:finite-hopf-pairing-identity-tuple}.

Conversely, a compact-source proof requires the source \(G\)-entries,
left and right radical bases \(K_{\pm\gamma,\bar p}\), quotient
splittings \(Q_{\pm\gamma,\bar p}\), quotient-pairing matrices, and
rank or determinant rows proving full rank.  Hopf adjointness,
Frobenius invariance, target radical tables, rank equality alone, and
signed dimensions do not supply quotient non-degeneracy.
\end{proposition}

\begin{proof}
Because \(K_{\gamma,\bar p}\) and \(K_{-\gamma,\bar p}\) are supplied
as the left and right radicals of \(G_{\gamma,\bar p}\), the pairing
descends to
\[
P_{R,\gamma,\bar p}^{\Pi}/K_{\gamma,\bar p}
\quad\text{and}\quad
P_{R,-\gamma,\bar p}^{\Pi}/K_{-\gamma,\bar p}.
\]
The splittings \(Q_{\gamma,\bar p}\) and \(Q_{-\gamma,\bar p}\) give
finite bases of these quotients, and the matrix of the descended
pairing in these bases is precisely
\(\overline G_{\gamma,\bar p}
=Q_{\gamma,\bar p}^{t}G_{\gamma,\bar p}Q_{-\gamma,\bar p}\).  Full
rank equal to both quotient dimensions says that the two adjoint maps
\[
L_{\gamma,\bar p}\longrightarrow L_{-\gamma,\bar p}^{*},
\qquad
L_{-\gamma,\bar p}\longrightarrow L_{\gamma,\bar p}^{*}
\]
are isomorphisms.  This is non-degeneracy of the quotient pairing.
\end{proof}

The packet
\(\texttt{certificates/hall/finite\_hall\_quotient\_pairing\_nondegeneracy}\)
verifies
\(\texttt{FINITE\_HALL\_QUOTIENT\_PAIRING\_NONDEGENERACY\_OBSTRUCTION\_VERIFIED}\):
it imports the row \(362\) coproduct-adjointness packet, the row \(360\)
supersymmetry packet, the row \(358\) positive-negative pairing
definition, the pairing-kernel \(\varprojlim^1\) obstruction packet,
the transition-radical obstruction packet, and the compact Hall source
obstruction ledger; it records
Proposition~\ref{prop:finite-hall-quotient-pairing-nondegeneracy}; and
it keeps the current source \(G\)-entries, \(K\)-entries,
\(Q\)-entries, quotient-pairing matrices, quotient-rank rows,
determinant rows, and radical-identification rows open.  Hence row
\(363\) proves the quotient non-degeneracy criterion for supplied
finite source data; it does not prove transition preservation of
radicals or \(\varprojlim^1\)-vanishing for the radical tower.

\begin{definition}[Compact geometric Hopf-pairing datum]
\label{def:compact-geometric-hopf-pairing-datum}
At finite HN height \(R\), assume the compact geometric Hall
bialgebra datum of
Definition~\ref{def:compact-geometric-hall-bialgebra-datum}.  A
compact geometric Hopf-pairing datum on the descended primitive stage is
the following additional source data:
\begin{enumerate}
\item a degree-zero protected trace
\[
\operatorname{tr}_{0,R}:H_c^\bullet
(\mathfrak M^{\mathrm{red}}_{R,0},\Phi^{\mathrm{red}}_{R,0}\otimes
o_{R,0})\longrightarrow k
\]
compatible with the vacuum counit and with HN transitions;
\item positive-negative compact pairing correspondences
\(\mathfrak P^{\mathrm{red}}_{R,\gamma,-\gamma;0}\) obtained by
restricting the Hall product to total descended Gram degree \(0\), with
orientation and Thom--Sebastiani transport inherited from the Hall
datum;
\item homogeneous pairing matrices
\[
G_{\gamma,ab}
=\operatorname{tr}_{0,R}\!\left(
m_R(e_{\gamma,a},e_{-\gamma,b})\right)
\]
on every retained parity block, together with kernel matrices
\(K_{\gamma,\bar p}\) and quotient splittings \(Q_{\gamma,\bar p}\);
\item triple pairing correspondences for charges
\(\gamma_1+\gamma_2+\gamma_3=0\), with cyclic permutation
isomorphisms induced by the \(3\)-Calabi--Yau Serre pairing and by
Joyce--Upmeier orientation transport;
\item finite rank-zero witnesses for Hopf adjointness and the primitive
Frobenius cyclic identity, and finite rank witnesses that the induced
pairing on
\((P_R^\Pi/\operatorname{Rad}_R^\Pi)_\gamma\otimes
(P_R^\Pi/\operatorname{Rad}_R^\Pi)_{-\gamma}\)
is non-degenerate in every retained parity block.
\end{enumerate}
The radical is
\[
\operatorname{Rad}_{R,\gamma,\bar p}^{\Pi}
=\{x\in P_{R,\gamma,\bar p}^{\Pi}\mid
\langle x,y\rangle_R=0\text{ for all }
y\in P_{R,-\gamma,\bar p}^{\Pi}\},
\]
with both signs included.  The datum is source-side: the target
Gritsenko--Nikulin radical may be used only as a comparison codomain
after the source radical and quotient have been formed.
\end{definition}

The packet
\(\texttt{certificates/hall/finite\_hall\_compact\_geometric\_hopf\_pairing\_datum}\)
verifies
\(\texttt{FINITE\_HALL\_COMPACT\_GEOMETRIC\_HOPF\_PAIRING\_DATUM\_OBSTRUCTION\_VERIFIED}\):
it imports the row \(363\) quotient non-degeneracy packet, the row
\(362\) coproduct-adjointness packet, the row \(361\) Frobenius
invariance packet, the reduced Calabi--Yau threefold Serre-sign
criterion packet, the level-\(\mathsf Z\) protected-trace obstruction
ledger as a firewall, and the compact Hall source obstruction ledger;
it records
Definition~\ref{def:compact-geometric-hopf-pairing-datum}; and it
keeps the source degree-zero trace, positive-negative pairing
correspondences, orientation and Thom--Sebastiani transport,
\(G\)-entries, \(K\)-entries, \(Q\)-entries, cyclic three-point
correspondences, Hopf-adjointness witnesses, Frobenius cyclic
witnesses, and quotient non-degeneracy witnesses open.  Hence row
\(364\) records the compact geometric Hopf-pairing datum as a
source-level datum; it does not construct it, and the
level-\(\mathsf Z\) protected trace does not substitute for the source
degree-zero trace.

\begin{definition}[Radical-preserving transition datum]
\label{def:transition-radical-preservation}
Let \(R'\le R\).  Suppose the primitive-subspace transition datum of
Definition~\ref{def:transition-primitive-subspace-preservation} is
given after normal-ordered descent, so that
\[
T_R^{\Prim}:P_R^\Pi\longrightarrow P_{R'}^\Pi
\]
is defined.  A radical-preserving transition datum consists of:
\begin{enumerate}
\item source and target Hopf-pairing matrices
\(G_{R,\gamma,\bar p}\) and \(G_{R',\gamma,\bar p}\), with the Hopf
adjointness, Frobenius cyclic, and quotient non-degeneracy rows of
Definition~\ref{def:finite-hopf-pairing-identity-tuple};
\item source and target radical kernel rows
\[
K_{R,\gamma,\bar p}=\operatorname{Rad}_{R,\gamma,\bar p}^{\Pi},
\qquad
K_{R',\gamma,\bar p}=\operatorname{Rad}_{R',\gamma,\bar p}^{\Pi},
\]
together with quotient splittings \(Q_{R,\gamma,\bar p}\) and
\(Q_{R',\gamma,\bar p}\);
\item for every target opposite primitive \(y'\in
P_{R',-\gamma,\bar p}^{\Pi}\), a lifted source opposite primitive
\(L_{R'R}^{-\gamma}y'\in P_{R,-\gamma,\bar p}^{\Pi}\) with
\(T_R^{\Prim,-\gamma}L_{R'R}^{-\gamma}y'=y'\);
\item the pairing-transition identity
\[
\bigl\langle T_R^{\Prim,\gamma}x,y'\bigr\rangle_{R'}
=\bigl\langle x,L_{R'R}^{-\gamma}y'\bigr\rangle_R
\]
for every \(x\in P_{R,\gamma,\bar p}^{\Pi}\) and
\(y'\in P_{R',-\gamma,\bar p}^{\Pi}\);
\item a restricted radical transition matrix
\[
T_R^{\operatorname{Rad}}:K_{R,\gamma,\bar p}\longrightarrow
K_{R',\gamma,\bar p}
\]
and a descended quotient transition matrix
\[
\bar T_R:P_{R,\gamma,\bar p}^{\Pi}/K_{R,\gamma,\bar p}
\longrightarrow
P_{R',\gamma,\bar p}^{\Pi}/K_{R',\gamma,\bar p},
\]
with zero radical-image, quotient well-definedness, splitting
compatibility, and composition defect ranks.
\end{enumerate}
Scalar traces, denominator products, signed multiplicities, target
radical tables, pairing matrices without Hopf identities, rank
equalities, primitive transition alone, and normal-ordered degree maps
are not radical-transition data.
\end{definition}

\begin{proposition}[Transition preservation of Hopf radicals;
\textit{proved}]
\label{prop:transition-radical-preservation-criterion}
Given a radical-preserving transition datum of
Definition~\ref{def:transition-radical-preservation}, the primitive
transition carries the Hopf-pairing radical into the Hopf-pairing
radical:
\[
T_R^{\Prim,\gamma}\bigl(\operatorname{Rad}_{R,\gamma,\bar p}^{\Pi}\bigr)
\subseteq
\operatorname{Rad}_{R',\gamma,\bar p}^{\Pi}.
\]
Consequently \(T_R^{\Prim}\) descends to the radical quotients through
the displayed \(\bar T_R\).
\end{proposition}

\begin{proof}
Let \(x\in\operatorname{Rad}_{R,\gamma,\bar p}^{\Pi}\), and let
\(y'\in P_{R',-\gamma,\bar p}^{\Pi}\).  Choose the supplied lift
\(L_{R'R}^{-\gamma}y'\in P_{R,-\gamma,\bar p}^{\Pi}\).  The
pairing-transition identity gives
\[
\bigl\langle T_R^{\Prim,\gamma}x,y'\bigr\rangle_{R'}
=\bigl\langle x,L_{R'R}^{-\gamma}y'\bigr\rangle_R=0,
\]
because \(x\) pairs trivially with every source opposite primitive.
Since this holds for every \(y'\), the vector
\(T_R^{\Prim,\gamma}x\) lies in
\(\operatorname{Rad}_{R',\gamma,\bar p}^{\Pi}\).  The finite restricted
matrix \(T_R^{\operatorname{Rad}}\) records this inclusion in radical bases, and the
quotient well-definedness row is precisely the statement that
\(T_R^{\Prim}\) factors through \(\bar T_R\).  The splitting and
composition rows make the quotient transition independent of the
chosen finite representatives and compatible with successive stages.
\end{proof}

The packet
\(\texttt{certificates/hall/transition\_radical\_preservation}\)
records the present state of this datum.  Its verifier status is
\texttt{TRANSITION\_RADICAL\_PRESERVATION\_OBSTRUCTION\_VERIFIED}: the
primitive-transition input, source and target Hopf-pairing matrices,
source and target Hopf-pairing identity rows, source and target radical
kernels, opposite-primitive lift rows, pairing-transition identities,
radical-image identities, restricted radical transition matrices,
source and target quotient splittings, descended quotient transition
matrices, and radical-transition composition rows are not yet supplied.
Hence transition preservation of Hopf radicals remains open.

\begin{proposition}[Geometric Hopf-pairing data give the finite
Hopf-pairing defect rows; \textit{proved relative to the datum}]
\label{prop:geometric-hopf-pairing-data-give-defects}
Let the compact geometric Hall bialgebra datum and the compact
geometric Hopf-pairing datum be given at height \(R\).  Then
\[
\mathfrak O_{\mathrm{Hpair},R}=0.
\]
Consequently the radical of the source pairing is a coproduct coideal;
if the cyclic Frobenius component of the datum is imposed on primitives,
it is also a Lie ideal, and the quotient carries the descended primitive
Lie bracket and Hopf pairing.  Conversely, a matrix \(G\), a kernel
matrix \(K\), and a quotient matrix \(Q\) with the right ranks are not a
compact source Hopf-pairing datum unless their entries are obtained
from the protected trace, pairing correspondences, and cyclic
three-point coherences above.
\end{proposition}

\begin{proof}
Hopf adjointness is the equality of two pull-push evaluations on the
same three-point extension/collision correspondence.  On one side one
multiplies \(x\) and \(y\), then pairs with \(z\) by
\(\operatorname{tr}_{0,R}\); on the other one splits \(z\) by
\(\Delta_R\), pairs the two resulting factors with \(x\) and \(y\), and
then applies the tensor trace.  The common flag refinement and the
projection formula identify the two Borel--Moore chain maps, while the
Thom--Sebastiani and orientation transports identify the coefficient
systems.  This is
\[
\langle m_R^\Theta(x,y),z\rangle_R
=\langle x\otimes y,\Delta_R^\Theta z\rangle_R.
\]

The Frobenius cyclic identity is the same three-point trace after a
cyclic permutation of the inputs.  The \(3\)-Calabi--Yau Serre pairing
changes the order by the Koszul sign; the supercommutator therefore
gives
\[
\langle[x,y]_\Theta,z\rangle_R
+(-1)^{|x||y|}\langle y,[x,z]_\Theta\rangle_R=0.
\]
The quotient non-degeneracy component is not a consequence of the first
two equations; it is the finite rank condition in the datum after
dividing by \(K_{\gamma,\bar p}\).  Hence the three components of
\(\mathfrak O_{\mathrm{Hpair},R}\) vanish.

The coideal and Lie-ideal conclusions are
Lemma~\ref{lem:hopf-radical-ideal-coideal}.  The converse sentence is a
source/target firewall: a finite bilinear form can have the same kernel
rank as the target radical without being the protected trace pairing of
the compact \(K3\times E\) Hall source.
\end{proof}

The packet
\(\texttt{certificates/hall/finite\_hall\_hopf\_pairing\_defect\_rows}\)
verifies
\(\texttt{FINITE\_HALL\_HOPF\_PAIRING\_DEFECT\_ROWS\_OBSTRUCTION\_VERIFIED}\):
it imports the row \(364\) compact geometric Hopf-pairing datum packet,
the row \(350\) Hall bialgebra compatibility packet, the row \(362\)
Hopf-adjointness criterion, the row \(361\) Frobenius cyclic criterion,
the row \(363\) quotient non-degeneracy criterion, and the compact Hall
source obstruction ledger; it records
Proposition~\ref{prop:geometric-hopf-pairing-data-give-defects}; and it
keeps the present compact source Hopf-adjointness rows, Frobenius
cyclic rows, quotient non-degeneracy rows, \(M\)-, \(D\)-, \(G\)-,
\(K\)-, and \(Q\)-entries, radical ideal/coideal rows, degree-zero
source trace, positive-negative pairing correspondences, and cyclic
three-point coherences open.  Hence row \(365\) proves only the
relative implication from supplied compact geometric Hopf-pairing data
to \(\mathfrak O_{\mathrm{Hpair},R}=0\); it does not construct the
datum and does not prove Hopf-radical descent for the current compact
source.

\begin{proposition}[A pairing matrix is not a Hopf-pairing proof;
\textit{proved}]
\label{prop:pairing-matrix-not-hopf-pairing-proof}
At a finite HN stage, the matrix \(G\), its kernel matrix \(K\), and a
quotient splitting \(Q\) do not by themselves prove that the radical is
the Hopf radical used in primitive recognition.  The finite stage proves
Hopf-radical descent only after
\(\mathfrak O_{\mathrm{Hpair},R}=0\): Hopf adjointness, the primitive
Frobenius cyclic identity, and non-degeneracy of the induced quotient
pairing.  Under these three vanishings, the radical is a coproduct
coideal and a Lie ideal, hence the quotient carries the descended
primitive Lie bracket and Hopf pairing.
\end{proposition}

\begin{proof}
The matrix \(G\) defines a bilinear form and \(K=\ker G\) records its
linear radical.  Coideal stability is not a linear-algebra consequence
of being a kernel: it follows from Hopf adjointness
\(\langle m(x,y),r\rangle=\langle x\otimes y,\Delta r\rangle\) and
non-degeneracy of the quotient tensor pairing.  Lie-ideal stability is
not a consequence of kernel membership either; it follows from the
Frobenius identity
\(\langle[x,r],z\rangle=-(-1)^{|x||r|}\langle r,[x,z]\rangle\).  These
are exactly the first two defects in
\(\mathfrak O_{\mathrm{Hpair},R}\), and quotient non-degeneracy is the
third.  With all three vanishing, Lemma~\ref{lem:hopf-radical-ideal-coideal}
applies.  If any one is absent, \(K\) is only the radical of a finite
matrix, not the Hopf radical required by
Theorem~\ref{thm:primitive-recognition}.
\end{proof}

The packet
\(\texttt{certificates/hall/finite\_hall\_pairing\_matrix\_firewall}\)
verifies
\(\texttt{FINITE\_HALL\_PAIRING\_MATRIX\_FIREWALL\_OBSTRUCTION\_VERIFIED}\):
it imports the row \(365\) Hopf-pairing defect packet, the row \(362\)
Hopf-adjointness criterion, the row \(361\) Frobenius cyclic criterion,
the row \(363\) quotient non-degeneracy criterion, and the compact Hall
source obstruction ledger; it records
Proposition~\ref{prop:pairing-matrix-not-hopf-pairing-proof}; and it
keeps the current source \(G\)-, \(K\)-, and \(Q\)-entries,
Hopf-adjointness rows, Frobenius cyclic rows, quotient
non-degeneracy rows, radical coideal rows, radical Lie-ideal rows, and
quotient Hopf-pairing rows open.  Hence row \(366\) is a firewall
theorem: \(G,K,Q\) and their ranks can identify a finite bilinear
kernel, but they do not prove the Hopf radical used in primitive
recognition.  Hopf-radical descent begins only after
\(\mathfrak O_{\mathrm{Hpair},R}=0\).

\begin{lemma}[Hopf radical is a Lie ideal and a coideal;
\textit{conditional on $\mathfrak o_F=0$ and Hopf adjointness}]
\label{lem:hopf-radical-ideal-coideal}
Let $P_R^\Pi$ be a finite normal-ordered primitive stage with
transported product $m_R^\Theta$, coproduct $\Delta_R^\Theta$, graded
symmetric Hopf pairing $\langle\cdot,\cdot\rangle_R$, and radical
$\operatorname{Rad}_R^\Pi$.  Assume Hopf adjointness
$\langle m_R^\Theta(x,y),z\rangle_R
=\langle x\otimes y,\Delta_R^\Theta z\rangle_R$
and non-degeneracy on
$P_R^\Pi/\operatorname{Rad}_R^\Pi\otimes
P_R^\Pi/\operatorname{Rad}_R^\Pi$.  Then
\[
(\bar q_{\Pi,R}\otimes\bar q_{\Pi,R})\Delta_R^\Theta(r)=0
\quad(r\in\operatorname{Rad}_R^\Pi),
\]
i.e.\ $\operatorname{Rad}_R^\Pi$ is a coideal.  If, in addition,
$\mathfrak o_{F,R}=0$, then $\operatorname{Rad}_R^\Pi$ is a Lie
ideal.
\end{lemma}

\begin{proof}
For $r\in\operatorname{Rad}_R^\Pi$ and $x,y\in P_R^\Pi$, Hopf
adjointness gives
$\langle x\otimes y,\Delta_R^\Theta(r)\rangle_R
=\langle m_R^\Theta(x,y),r\rangle_R=0$;
non-degeneracy of the quotient tensor pairing then kills the class.
The Lie-ideal half is direct: if $\mathfrak o_{F,R}=0$ then
$\langle[x,r]_\Theta,z\rangle_R
=-(-1)^{|x||r|}\langle r,[x,z]_\Theta\rangle_R=0$ for every $z$,
so $[x,r]_\Theta\in\operatorname{Rad}_R^\Pi$.
\end{proof}

The packet
\(\texttt{certificates/hall/finite\_hall\_hopf\_radical\_ideal\_coideal}\)
verifies
\(\texttt{FINITE\_HALL\_HOPF\_RADICAL\_IDEAL\_COIDEAL\_OBSTRUCTION\_VERIFIED}\):
it imports the row \(366\) pairing-matrix firewall, the row \(365\)
Hopf-pairing defect packet, the row \(362\) Hopf-adjointness criterion,
the row \(361\) Frobenius cyclic criterion, the row \(363\) quotient
non-degeneracy criterion, and the compact Hall source obstruction
ledger; it records Lemma~\ref{lem:hopf-radical-ideal-coideal}; and it
keeps the present source Hopf-adjointness rows, Frobenius cyclic rows,
quotient tensor non-degeneracy rows, radical coideal rows, radical
Lie-ideal rows, \(G\)-, \(K\)-, and \(Q\)-entries open.  Hence row
\(367\) proves the conditional algebraic criterion for a supplied
finite Hopf pairing; it does not prove the radical ideal/coideal rows
for the current compact \(K3\times E\) source.

\subsubsection*{Integral form and stable-envelope transport}

Two parallel constructions act on the same compact Hall heart:
\begin{enumerate}
\item the integral-form comparison modeled on the
Schiffmann--Vasserot elliptic Hall algebra; a compact
\(K3\times E\) specialization would require a non-toric
Maulik--Okounkov-type stable-envelope transport;
\item the proposed Drinfeld double
\(\mathcal D_\hbar(\textup{CoHA}_{K3\times E})\) of Vol III, whose
positive half must be the protected primitive Hopf algebra above and
whose Drinfeld pairing must be the Hopf pairing.
\end{enumerate}
The compatibility datum between these two constructions is a comparison
with Calabi--Yau quantum groups.  A different constant or
sign would be a distinct mathematical claim; the Drinfeld double is not
asserted here as an input theorem.

\begin{definition}[Integral-form stable-envelope comparison datum]
\label{def:integral-form-stable-envelope-comparison-datum}
At finite height \(R\), an integral-form stable-envelope comparison
datum for the compact \(K3\times E\) Hall heart consists of the
following finite rows:
\begin{enumerate}
\item an integral lattice
\((Y^+_{R,\mathbb Z}(X),m_R,\Delta_R,\langle-,-\rangle_R)\)
inside the protected Hall bialgebra, stable under product,
coproduct, and the Hopf pairing;
\item a separated and complete topology on \(Y^+_R(X)\), compatible
with the HN transition maps and the Hopf radical quotient;
\item a completed Drinfeld double
\(\mathcal D_\hbar(Y^+_R(X))\) with cross-relations whose structure
constants are the finite Hall pairing rows;
\item a Cartan part in the double, identified with
\(\Lambda_{II}^{2,1}\otimes\mathbb C\), and rows proving that this
Cartan survives the reduced primitive quotient;
\item a compact non-toric stable-envelope correspondence satisfying
support, normalization, degree, and chamber-cocycle identities on the
retained \(K3\times E\) moduli;
\item comparison constants, signs, and degree shifts identifying the
preceding double with the Calabi--Yau quantum-group normalization.
\end{enumerate}
The datum is source-level.  The target BKM denominator, the
Schiffmann--Vasserot affine \(A_0\) theorem, and the
Maulik--Okounkov quiver-variety stable envelope do not supply these
finite compact rows.
\end{definition}

\begin{proposition}[Stable-envelope and Drinfeld-double scope
firewall; \textit{proved}]
\label{prop:integral-form-stable-envelope-transport-firewall}
The Maulik--Okounkov stable-envelope theorem and the
Schiffmann--Vasserot elliptic-Hall comparison may be used here only in
their stated scope.  They do not construct the integral form,
separated completion, compact non-toric stable envelope, Drinfeld
double cross-relations, or Cartan identification required for the
compact \(K3\times E\) Hall heart.  Any assertion
\[
\mathcal D_\hbar(Y^+(K3\times E))^{\mathrm{int}}_{\mathrm{compl}}
 \simeq
Y_\hbar(\widehat{\mathfrak g}_{\Delta_5})
\]
is therefore conditional on the datum of
Definition~\ref{def:integral-form-stable-envelope-comparison-datum}.
\end{proposition}

\begin{proof}
Appendix~\ref{app:stable-envelope} records the precise imported
theorems.  Maulik--Okounkov stable envelopes are constructed for
smooth symplectic varieties in the quiver-variety setting with a torus
action and chamber data; their theorem gives the geometric
\(R\)-matrix in that category.  Schiffmann--Vasserot identify the
\(\mathbb C^3\) cohomological Hall algebra with the positive half in
the affine \(A_0\) elliptic-Hall setting.  Neither theorem contains a
compact non-toric \(K3\times E\) moduli atlas, a retained D0-HN
comparison, an \(O2\) wall atlas, or the Hopf-radical quotient rows
above.

Thus the cited theorems fix the target shape and the normalization
problem, but they do not supply the finite source rows of
Definition~\ref{def:integral-form-stable-envelope-comparison-datum}.
The compact comparison is a separate datum: without those rows the
symbol \(\mathcal D_\hbar(Y^+(K3\times E))^{\mathrm{int}}_{\mathrm{compl}}\)
is a requested construction, not an established input theorem.
\end{proof}

The packet
\(\texttt{certificates/hall/finite\_hall\_integral\_form\_stable\_envelope\_transport}\)
verifies
\(\texttt{FINITE\_HALL\_INTEGRAL\_FORM\_STABLE\_ENVELOPE\_TRANSPORT\_OBSTRUCTION\_VERIFIED}\):
it imports the row \(367\) Hopf-radical ideal/coideal criterion, the
compact Hall source ledger, the D0-HN obstruction ledger, and the
\((O2)\) wall-atlas ledger; it records
Definition~\ref{def:integral-form-stable-envelope-comparison-datum}
and
Proposition~\ref{prop:integral-form-stable-envelope-transport-firewall};
and it keeps the present compact source integral-form rows, product
stability rows, coproduct stability rows, pairing stability rows,
separated-completion rows, complete-completion rows, Drinfeld-double
rows, cross-relation rows, stable-envelope rows, Cartan
identification rows, primitive Cartan survival rows, compact
non-toric replacement theorem, and Calabi--Yau quantum-group
comparison rows open.  Hence row \(368\) proves only the scope
firewall and the required finite datum; it does not construct the
compact \(K3\times E\) integral form, the compact Drinfeld double, or
stable-envelope transport.

\subsubsection*{Conilpotent source coalgebra and Hopf radical
quotient}

\begin{definition}[Finite conilpotent source-coalgebra datum]
\label{def:conilpotent-source-coalgebra}
A finite conilpotent source-coalgebra datum at HN height \(R\) is a
source-side package
\[
(C_{X,R},\eta^C_R,\varepsilon^C_R,
F^{\mathrm{co}}_{R,\bullet},
\Delta_R^{\mathrm{ch}},\rho^C_{R'R})
\]
obtained from the retained collision correspondences on
\(\mathcal F_{X,\sigma,S,\le R}^{\mathrm{hyb}}\), after the reduced
\(E\)-quotient, with the following finite rows:
\begin{enumerate}
\item \(C_{X,R}\) is a coaugmented chiral coalgebra on \(E\), with
coaugmentation \(\eta^C_R\), counit \(\varepsilon^C_R\), and
coalgebra identities for \(\Delta_R^{\mathrm{ch}}\);
\item \(C_{X,R}\) is filtered by a finite increasing source charge
filtration
\[
0=F^{\mathrm{co}}_{R,-1}\subset F^{\mathrm{co}}_{R,0}\subset
\cdots\subset F^{\mathrm{co}}_{R,N_R}=C_{X,R}
\]
whose associated graded pieces are finite dimensional in each
\(\Gamma_R^\Pi\)-degree, where
\(\Gamma_R^\Pi:=\overline\Pi_X(\widehat\Gamma_R)\)
is the image, under the \emph{normal-ordered} Gram homomorphism
\(\overline\Pi_X\) of Chapter~\ref{ch:k3xe-setup}, of the retained
normal-ordered charge window $\widehat\Gamma_R$ at Harder--Narasimhan
height \(R\) (\S\ref{subsec:hybrid-ran-carrier});
\item the additive \(\overline\Pi_X\), not the raw quadratic
\(\Pi_X\), is the grading map;
\item the reduced iterated coproduct is nilpotent:
\[
(\overline\Delta_R^{\mathrm{ch}})^{N_R+1}=0;
\]
\item for every successor \(R'\le R\), the transition map
\(\rho^C_{R'R}\) preserves the filtration, counit, coaugmentation, and
collision coproduct.
\end{enumerate}
Only after these transition rows and the strict Mittag--Leffler rows
are supplied does the inverse limit define the source coalgebra
\[
C_X=\varprojlim_R C_{X,R}.
\]
\end{definition}

\begin{proposition}[Conilpotent source-coalgebra firewall;
\textit{proved}]
\label{prop:conilpotent-source-coalgebra-firewall}
The notation \(C_{X,R}\), the formal bar-coalgebra construction, the
target BKM denominator, the stable-envelope comparison, and the
Drinfeld-double comparison do not construct the finite conilpotent
source-coalgebra datum of
Definition~\ref{def:conilpotent-source-coalgebra}.  A compact
\(K3\times E\) source coalgebra is available only after the
coaugmentation, counit, finite filtration, collision-coproduct,
coalgebra-identity, conilpotence, quotient-after-correspondence,
transition, and Mittag--Leffler rows are supplied.
\end{proposition}

\begin{proof}
A conilpotent coalgebra is not determined by its name.  It requires a
coalgebra object, counit and coaugmentation, a coproduct satisfying
coassociativity and counit identities, and a filtration for which the
reduced iterated coproduct terminates.  In this manuscript those
structures must be produced by retained compact collision
correspondences and their quotient-after-correspondence identities.
The target denominator and the later Drinfeld-double comparison live
on the comparison side; stable-envelope transport is a separate
normalization problem.  None of them supplies the source chain maps or
zero-defect rows listed in
Definition~\ref{def:conilpotent-source-coalgebra}.  The inverse-limit
object \(C_X\) additionally requires transition compatibility and
strict Mittag--Leffler exactness, since otherwise
\(R^1\varprojlim\) remains an obstruction.
\end{proof}

The packet
\(\texttt{certificates/hall/finite\_hall\_conilpotent\_source\_coalgebra}\)
verifies
\(\texttt{FINITE\_HALL\_CONILPOTENT\_SOURCE\_COALGEBRA\_OBSTRUCTION\_VERIFIED}\):
it imports the row \(368\) stable-envelope/Drinfeld-double firewall,
the hybrid carrier ledger, the compact Hall source obstruction
ledger, the D0-HN obstruction ledger, the Hall product and coproduct
transition ledgers, and the primitive-space and pairing-kernel
\(\varprojlim^1\) ledgers; it records
Definition~\ref{def:conilpotent-source-coalgebra} and
Proposition~\ref{prop:conilpotent-source-coalgebra-firewall}; and it
keeps the present compact source coalgebra rows, coaugmentation rows,
counit rows, finite filtration rows, collision-coproduct rows,
coalgebra-identity rows, conilpotence rows, quotient coalgebra
identity rows, transition coalgebra rows, inverse-limit
Mittag--Leffler rows, and bar-cobar comparison rows open.  Hence row
\(369\) proves only the source-data boundary for \(C_{X,R}\) and
\(C_X\); it does not construct the compact \(K3\times E\)
conilpotent source coalgebra.

\begin{proposition}[Sectorial Hall truncation and inverse limit;
\textit{conditional on
$[\textup{K3$\times$E-fin}]$, $[\textup{K3$\times$E-atlas}]$,
and $[\textup{ML}]$}]
\label{prop:sectorial-hall-truncation}
Under the three standing hypotheses
$[\textup{K3$\times$E-fin}]$ (finite-type semistable
substacks at bounded HN height),
$[\textup{K3$\times$E-atlas}]$ (finite d-critical/cosection
Hall atlas with Joyce--Upmeier orientation lines), and
$[\textup{ML}]$ (Mittag--Leffler descent for the HN inverse system),
the finite Hall product, coproduct, primitive projection, and Hopf
pairing are compatible with the transition maps $\rho_{R'R}$, and
their inverse limits define the completed Hall Hopf datum
$\widetilde P_X$, $C_X$.  The Mittag--Leffler hypothesis controls
the $R^1\varprojlim$ obstruction on every finite-window slice; without
it, only a compatible pro-object is obtained.
\end{proposition}

The hypothesis $[\textup{K3$\times$E-fin}]$ is not a
consequence of fibrewise K3 boundedness alone.  A precise sufficient
input is product-stability boundedness in the
Abramovich--Polishchuk/Liu sense
\cite{AbramovichPolishchukTStructures,LiuProductStability}: for a
rational K3 stability condition $\sigma_S=(\mathcal A_S,Z_S)$ and
Liu's product stability $\sigma_{s,t}$ on
$D^b(\operatorname{Coh}(S\times E))$, fixed-class boundedness
\[
\mathfrak M^{ss}_{\sigma_{s,t}}(\gamma)\text{ is finite type over }\C
\]
together with Liu's support property would imply
$[\textup{K3$\times$E-fin}]$ for bounded $h_S$-windows.
Piyaratne--Toda boundedness for threefold Bridgeland/BMT
\cite{PiyaratneTodaModuli3folds} does not automatically supply this
theorem; fibrewise Bayer--Macri boundedness alone does not either.

\subsubsection*{Wall-crossing and chamber dependence}

The lattice cocycle $B(c,c')$ of the dyonic Mukai construction is
chamber-independent.  The compact Hall side $D_0(X;\sigma,S)$ is
chamber-dependent: finite-type semistable substacks change under
wall-crossing, and the Kontsevich--Soibelman wall-crossing formula
\cite{KSCoHA} relates Hall counts in adjacent chambers.  If the
datum $D_0(X;\sigma,S)$ is supplied in chamber $\sigma$, then
$\overline\Pi_X$ remains a homomorphism and the recognition target
is $\mathfrak g_{\Delta_5}$ provided the chamber identifications hold
for the OP/DT scalar import; that the chosen $\sigma$ lies in the
OP/DT chamber, or that the wall-crossing transport is compatible
with $\overline\Pi_X$, is a separate conjectural input.  No
chamber-independence of the recognition target is asserted.

\subsubsection*{The Hall obstruction class}

The Hall lane carries the residuals
\[
\mathfrak O_{\mathrm{Hall},R}
=(\mathfrak o^{\mathrm{stage}}_R,\,
\mathfrak o^{\mathrm{conf}}_R,\,
\mathfrak o^{\mathrm{prop}}_R,\,
\mathfrak o^{\mathrm{TS}}_R,\,
\mathfrak O_{\mathrm{Hbialg},R},\,
\mathfrak O_{\mathrm{Hpair},R},\,
\mathfrak o^{\mathrm{ML}}_{R'R}).
\]
The first four entries assert the existence of the compact geometric
Hall bialgebra datum of
Definition~\ref{def:compact-geometric-hall-bialgebra-datum}: finite
retained substacks, compact-support functoriality, and coherent
Thom--Sebastiani orientation transport.  The tuple
\(\mathfrak O_{\mathrm{Hbialg},R}\) is then the algebraic defect tuple
of Definition~\ref{def:finite-hall-bialgebra-defect-tuple}; it is not
defined before product, coproduct, unit, and counit matrices have been
obtained from the compact correspondences.  The tuple
\(\mathfrak O_{\mathrm{Hpair},R}\) is the Hopf-pairing identity tuple of
Definition~\ref{def:finite-hopf-pairing-identity-tuple}; it is
independent of the bialgebra equations and is not defined before the
compact geometric Hopf-pairing datum of
Definition~\ref{def:compact-geometric-hopf-pairing-datum}: protected
trace, positive-negative pairing correspondences, cyclic three-point
coherences, source radical, quotient splitting, and quotient
non-degeneracy.  On the projection-finite stratum the wrapped product
is the classical d-critical CoHA product after cosection reduction; on
the wrapped stratum the construction is open and governed by the (D0)
obstruction.  PBW compatibility remains a separate recognition datum on
the descended primitive layer.

\subsubsection*{Comparison with Drinfeld doubling}

Vol III's Drinfeld doubling $Y^+(X)\to G(X)$ is the comparison target
for the compact source side; the integral form is
the chiral lift of the Schiffmann--Vasserot elliptic Hall algebra
after the stable-envelope and completion data are fixed.
Maulik--Okounkov stable-envelope transport supplies the integral
form on the projection-finite Ran stratum; the wrapped stratum is
the corresponding stratum-completion.  No black-box equality
$\mathrm{CoHA}=W_\infty$
is asserted: $\mathrm{CoHA}(\C^3)=Y^+(\mathfrak{gl}_1)$ in Vol III,
and $W_{1+\infty}$ appears only after Drinfeld doubling, centre,
vacuum evaluation.  The same discipline applies here: the protected
primitive Hopf algebra of $K3\times E$ becomes the BKM target only
after the Hopf radical quotient, the Frobenius identity, and the
no-extra-relations theorem of \S\ref{subsec:primitive-recognition}.

Under \(\mathfrak O_{\mathrm{Hall},R}=0\) and the Mittag--Leffler
transition datum, the Hall correspondences give a finite-stage
protected primitive Hopf datum on the cosection-twisted heart.  This
datum lives at the operator level and has no automatic comparison to
the scalar denominator
$\mathrm{den}(\autcor)=64^{-1}\Delta_5(2Z)$ imported in
Chapter~\ref{ch:view2-bkm}.  Closing the cross-level identity
\textcircled{2}$\leftrightarrow$\textcircled{3} requires a chain-level
scalar incarnation of the primitive structure: a Pfaffian line
$\mathscr L_{\mathrm{Pf},R}$ on the reduced moduli, a Dirac block
decomposition, and a first-order Pfaffian section whose
$\Sp_4(\Z)$-pullback recovers $\Delta_5$.  This is the fourth entry
$(\mathrm L 4)$, constructed in \S\ref{subsec:dirac-pfaffian}.


\subsection{First-order Dirac block and protected Pfaffian}
\label{subsec:dirac-pfaffian}

The fourth entry of the Dirac--Igusa pro-object is the protected
Pfaffian line and its first-order Dirac block decomposition.  The
Pfaffian, not the determinant, is the object of interest for a precise
reason: the automorphic target $\Delta_5$ is the square root of the
scalar weight-ten form $\Delta_5^2=\chi_{10}$, and the square root
carries the order-two Maass multiplier $\nu_{\Delta_5}$ that the square
sends to the trivial character.  Geometrically the determinant line
$\mathscr L_{\det}$ is canonically trivialisable and forgets
orientation; its square root, the Pfaffian line, retains the
orientation character $\epsilon_o$, and it is precisely the match
$\epsilon_o=\nu_{\Delta_5}$ that distinguishes the first-order object
$\Delta_5$ from its orientation-blind square $\Delta_5^2$.  Two
definitions of this Pfaffian meet.  The geometric definition starts from
the cosection-reduced d-critical heart, carries a Joyce--Upmeier
orientation, descends to the reduced quotient, and associates a
Pfaffian line to the determinant of the reduced obstruction theory in
the Brav--Bussi--Dupont--Joyce--Szendr\H{o}i framework
\cite{BBDJS2015}.  The algebraic construction starts from the imported
Borcherds--Kac--Moody denominator algebra, fixes a positive window in
$\Gamma_{\mathrm{gram}}$, and writes the algebraic super-Pfaffian
\[
\operatorname{sPf}^{\mathrm{alg}}_R(D)
=\prod_{\gamma\in A_R}(1-x^\gamma)^{\sdim V_\gamma}
\]
of the trivial first-order block decomposition $D_\gamma$ on the
target side.  The Pfaffian--Dirac theorem of Chapter~6 is the
identification of these two Pfaffian sections after the chosen
line-comparison $\iota_{\mathrm{aut}}$ and the four-obstruction
datum of Appendix~\ref{app:obstruction-discharges}.

\subsubsection*{Pfaffian line and squaring isomorphism}

\begin{definition}[Pfaffian line on the cosection-reduced quotient]
\label{def:pfaffian-line}
Let $K^{\mathrm{red}}_{\mathcal C}$ denote the cosection-reduced
self-Ext complex on a charge stratum, with reduced determinant
$\det K^{\mathrm{red}}_{\mathcal C}$ and Joyce--Upmeier orientation
$o_{\mathcal C}$.  The Pfaffian line at finite stage is
\[
\mathscr L_{\mathrm{Pf},R}=
\bigotimes_{\mathcal C}\operatorname{Pf}(K^{\mathrm{red}}_{\mathcal C,R})
\]
on the reduced moduli, with squaring isomorphism
\[
\mathscr L_{\mathrm{Pf},R}^{\otimes 2}\simeq\mathscr L_{\det,R}
\]
in the Brav--Bussi--Dupont--Joyce--Szendr\H{o}i sense
\cite{BBDJS2015}.  The successor maps
$\rho^{\mathrm{Pf}}_{R^+R}$ commute with these squaring isomorphisms,
identifying the inverse limit with a Pfaffian line
$\mathscr L_{\mathrm{Pf},X}$ on the inverse-limit object.  A
normalized section $\operatorname{pf}_{X,R}$ at finite stage and
$\operatorname{pf}_X=\varprojlim_R\operatorname{pf}_{X,R}$ are part
of the datum.
\end{definition}

\subsubsection*{Pfaffian-to-automorphic comparison}

The recognition statement is read on the automorphic line.  A
section-level identification $\operatorname{pf}_X=\Delta_5$ requires
a comparison of $\mathscr L_{\mathrm{Pf},X}$ with the pullback of the
automorphic weight-five line to the moduli base.

\begin{definition}[Pfaffian-to-automorphic line comparison]
\label{def:pfaffian-to-automorphic-line-comparison}
A section-level Pfaffian-to-automorphic comparison is a chosen
isomorphism of line bundles
\[
\iota_{\mathrm{aut}}:\mathscr L_{\mathrm{Pf},X}
\xrightarrow{\sim}\mathcal L^5\otimes\nu_{\Delta_5}
\]
to the pullback of the automorphic weight-five line, twisted by the
multiplier $\nu_{\Delta_5}$, compatible with cusp trivialization at
$\mathfrak c_\infty$ and with the squaring map
$\mathscr L_{\mathrm{Pf},X}^{\otimes 2}\simeq\mathscr L_{\det,X}$.
\end{definition}

\begin{definition}[Finite geometric Pfaffian datum
\((P_{\mathrm{fin}})\)]
\label{def:finite-geometric-pfaffian-datum}
At finite HN height \(R\), a finite geometric Pfaffian datum is the
following system on the reduced compact quotient
\(\mathfrak Q_R=(\mathfrak M_R/E)^{\mathrm{red}}\).
\begin{enumerate}
\item[\textup{(P1)}] For every retained stratum \(\mathcal C\), a
cosection-reduced perfect self-Ext complex
\(K^{\mathrm{red}}_{\mathcal C,R}\) with odd skew self-duality
\[
\vartheta_{\mathcal C,R}:K^{\mathrm{red}}_{\mathcal C,R}
\xrightarrow{\sim}
(K^{\mathrm{red}}_{\mathcal C,R})^\vee[1],
\]
compatible with the reduced obstruction theory and the quotient by
\(E\).
\item[\textup{(P2)}] A Joyce--Upmeier orientation
\(o_{\mathcal C,R}\) and the corresponding Pfaffian line
\[
\mathscr L_{\mathrm{Pf},R}
=\bigotimes_{\mathcal C}\operatorname{Pf}
(K^{\mathrm{red}}_{\mathcal C,R}),
\qquad
\mathscr L_{\mathrm{Pf},R}^{\otimes 2}\simeq\mathscr L_{\det,R}.
\]
\item[\textup{(P3)}] A finite-rank skew first-order block
\(D_{R,\gamma}\) on each retained primitive degree
\(\gamma\in\Gamma_R^{\Pi,+}\), with Pfaffian section
\[
\operatorname{pf}_{X,R}\in
H^0(\mathfrak Q_R,\mathscr L_{\mathrm{Pf},R}).
\]
In a cusp frame this section has support exactly on the retained
active degrees, full even/odd parity dimensions
\(p^{\bar0}_{R,\gamma},p^{\bar1}_{R,\gamma}\), and signed exponent
\(p^{\bar0}_{R,\gamma}-p^{\bar1}_{R,\gamma}=a_\Delta(\gamma)\).
\item[\textup{(P4)}] For every retained simple type-II wall chart,
with wall coordinate \(u_{\delta_i,R}\), a local factorization
\[
\operatorname{pf}_{X,R}
=\upsilon_{\delta_i,R}\,u_{\delta_i,R}\cdot
\operatorname{pf}^{\mathrm{tan}}_{i,R},
\qquad
\operatorname{ord}_{u_{\delta_i,R}=0}(\operatorname{pf}_{X,R})=1,
\]
where the unit \(\upsilon_{\delta_i,R}\) is invariant under the chosen
reflection lift and the normal Pfaffian rank is
\(N^{\mathrm{Pf}}_{\delta_i,R}=1\).  No type-I wall is allowed in this
datum.
\item[\textup{(P5)}] A finite-stage comparison
\(\iota_{\mathrm{aut},R}:\mathscr L_{\mathrm{Pf},R}\to
\mathcal L^5\otimes\nu_{\Delta_5}\) on the automorphic chart,
compatible with the cusp trivialization, the leading coefficient
\(\operatorname{lc}_{\mathfrak c_\infty}(\operatorname{pf}_{X,R})=64\),
and squaring to the determinant line.
\item[\textup{(P6)}] Strict transition isomorphisms
\(\rho^{\mathrm{Pf}}_{R^+R}:
\mathscr L_{\mathrm{Pf},R^+}|_{\mathfrak Q_R}
\xrightarrow{\sim}\mathscr L_{\mathrm{Pf},R}\) carrying
\(\operatorname{pf}_{X,R^+}\) to \(\operatorname{pf}_{X,R}\), with
closed images and vanishing \(\lim^1\) for the line and section
systems.
\end{enumerate}
The scalar Borcherds product, the squared determinant
\(\Delta_5^2\), and the OP scalar \(-4096\Delta_5^{-2}\) are not
entries of \((P_{\mathrm{fin}})\).  They are shadows obtained after
choosing the cusp frame and, in the OP case, after squaring and
inverting the Pfaffian section.

After finite homogeneous trivialisations, \((P_{\mathrm{fin}})\) is
recorded by the Pfaffian proof tables
\[
\mathrm{strata},\quad
\mathrm{skew\_complexes},\quad
\mathrm{pfaffian\_lines},\quad
\mathrm{pfaffian\_sections},\quad
\mathrm{wall\_charts},\quad
\mathrm{automorphic\_comparisons},\quad
\mathrm{transitions},\quad
\mathrm{scalar\_firewall}.
\]
The rows are not scalar product rows.  They record retained strata,
cosection-reduced skew perfect complexes, Pfaffian lines and their
squares, finite Pfaffian sections, type-II rank-one wall charts,
Pfaffian-to-automorphic line comparisons, strict transition maps, and
the exclusion of scalar Borcherds products, squared determinants, OP
scalar traces, and exponent tables as entries of \((P_{\mathrm{fin}})\).
Defect ranks, exponent residuals, leading residuals, and
\(R^1\varprojlim\)-ranks must be zero; the type-II normal Pfaffian rank
and divisor order must be one; the simple-reflection sign is \(-1\);
the automorphic comparison has leading coefficient \(64\), weight \(5\),
and character \(\nu_{\Delta_5}\).  Every row must carry geometric
source provenance and a proof reference.
\end{definition}

\begin{proposition}[Scalar Pfaffian data do not determine a Pfaffian section;
\textit{proved}]
\label{prop:scalar-pfaffian-data-not-section}
Let \(U\) be a quotient chart in the type-II wall complement, let
\(\mathcal L^5\) be the automorphic weight-five line on \(U\), and let
\(T_\chi\) be a non-trivial order-two flat line with character
\(\chi:\pi_1(U)\to\{\pm1\}\).  A divisor, a leading coefficient, a
scalar Borcherds product written after choosing a cusp trivialization of
\(\mathcal L^5\), or the square of that product in \(\mathcal L^{10}\)
does not determine a section of \(\mathcal L^5\otimes T_\chi\).  Thus a
section-level Pfaffian identity
\[
\operatorname{pf}_X=\Delta_5
\]
requires, in addition to the scalar product data, the Pfaffian line
\(\mathscr L_{\mathrm{Pf},X}\), the section \(\operatorname{pf}_X\), the
comparison \(\iota_{\mathrm{aut}}\), the order-two character
\(\epsilon_o=\nu_{\Delta_5}\), and the inverse-limit transition data for
these objects.  In particular, the exponent table \(f(nm,l)\), the
squared identity \(\Delta_5^2=\chi_{10}\), and the OP scalar
\(-4096\,\Delta_5^{-2}\) do not construct
\(\operatorname{pf}_{X,R}\) or the orientation character
\(\epsilon_{o,R}\).
\end{proposition}

\begin{proof}
Choose a simply connected open set \(V\subset U\) and a flat frame
\(\tau\) of \(T_\chi|_V\).  For a local section \(s\) of \(\mathcal L^5\),
the section \(s\otimes\tau\) of \(\mathcal L^5\otimes T_\chi\) has the
same local scalar expression as \(s\) after the frame \(\tau\) is set
equal to \(1\).  Under the canonical trivialization
\(T_\chi^{\otimes2}\simeq\mathcal O_U\), the squares of \(s\) and
\(s\otimes\tau\) give the same section \(s^2\) of \(\mathcal L^{10}\).
Their divisors and leading coefficients in the chosen cusp
trivialization are therefore identical.

The two objects are nevertheless different globally.  Along a loop
\(\ell\) with \(\chi(\ell)=-1\), analytic continuation fixes \(s\) and
sends \(s\otimes\tau\) to \(-s\otimes\tau\).  This monodromy is exactly
the information carried by the order-two orientation character.  Hence
the square, the local product, and the scalar trace forget the data
needed to define the Pfaffian section as a section of a twisted
weight-five line.  The line, section, line comparison, character, and
transition compatibilities are therefore separate data, not consequences
of the scalar product.
\end{proof}

\subsubsection*{First-order Dirac block decomposition}

\begin{definition}[First-order Dirac block]
\label{def:first-order-dirac-block}
For $\gamma\in\Gamma_R^{\Pi,+}$, write $P^{\Pi,+}_{R,\gamma}$ for the
cusp-positive primitive piece in degree $\gamma$ and put
$x^\gamma=q^nr^ls^m$ for $\gamma=(n,l,m)$.  The two-term
Pfaffian/Koszul block on
$P^{\Pi,+}_{R,\gamma}\oplus(P^{\Pi,+}_{R,\gamma})^\vee[1]$ is the
skew operator
\[
J_\gamma=
\begin{pmatrix}0&1-x^\gamma\\-(1-x^\gamma)&0\end{pmatrix},
\qquad
\operatorname{Pf}(J_\gamma)=(1-x^\gamma)^{\sdim P^{\Pi,+}_{R,\gamma}}.
\]
For $\gamma\notin\Gamma_R^{\Pi,+}$, set $P^{\Pi,+}_{R,\gamma}=0$.
The first-order block of $\mathfrak D_X^{D_0}$ at height $R$ is the
direct sum
\[
\mathfrak D_{X,R}^{D_0}=\bigoplus_{\gamma\in\Gamma_R^{\Pi,+}\setminus\{0\}}
J_\gamma,
\]
acting on
$\bigoplus_\gamma P^{\Pi,+}_{R,\gamma}\oplus
(P^{\Pi,+}_{R,\gamma})^\vee[1]$.
\end{definition}

The principal symbol of $\mathfrak D_{X,R}^{D_0}$ with respect to the
finite Chevalley--Eilenberg/Koszul filtration is the normal-ordered
Hall bracket of \S\ref{subsec:hall-correspondences}: the Pfaffian
line is the determinant detector, and the symbol carries the Lie
data.  The phrase \emph{first-order} means the Pfaffian/primitive
object before scalar squaring; no analytic Dirac operator on a
constructed state space is asserted.

\subsubsection*{The Pfaffian product at finite stage}

\begin{proposition}[Pfaffian product at finite stage;
\textit{conditional on \((P_{\mathrm{fin}})\)}]
\label{prop:pfaffian-product-finite-stage}
Suppose the orientation lane $(\textup{O1},\textup{O1}^+)$ has
descended to the reduced quotient, the finite geometric Pfaffian datum
\((P_{\mathrm{fin}})\) of
Definition~\ref{def:finite-geometric-pfaffian-datum} is supplied, and
the support and parity residuals
\[
\mathfrak o^{\mathrm{supp}}_{R,\gamma}=
\begin{cases}
P^{\Pi,+}_{R,\gamma},&a_\Delta(\gamma)=0\\
0,&a_\Delta(\gamma)\ne0
\end{cases},
\qquad
\mathfrak o^{\mathrm{par}}_{R,\gamma}
=p^{\bar0}_{R,\gamma}-p^{\bar1}_{R,\gamma}-a_\Delta(\gamma)\in\Z,
\]
together with the leading-coefficient residual
$\mathfrak o^{\mathrm{lead}}_R=\operatorname{lc}_{\mathfrak c_\infty}
(\operatorname{pf}_{X,R})-64$, all vanish on the test window, where
$a_\Delta(\gamma)$ is the Borcherds exponent and
$p^{\bar p}_{R,\gamma}=\dim P^{\Pi,+}_{R,\gamma,\overline p}$.  Then
the Pfaffian section at finite stage has expansion
\[
\operatorname{pf}_{X,R}|_{\mathfrak c_\infty}
=64q^{1/2}r^{1/2}s^{1/2}
\prod_{\substack{\gamma\in\Gamma_R^{\Pi,+}\\\gamma\ne 0}}
(1-x^\gamma)^{\sdim P^{\Pi,+}_{R,\gamma}}.
\]
The cusp-positive Pfaffian half is recorded; negative root
representatives and the Cartan/imaginary completion enter the Lie
recognition statement of \S\ref{subsec:primitive-recognition}.
\end{proposition}

The support residual is a super-vector-space residual, not a $K_0$
class: its vanishing excludes invisible even--odd cancelling
primitive pieces in target-zero degrees.  The parity residual is a
signed superdimension residual.  These finite-stage residuals on the
test window do not by themselves prove compact source coverage; the
inclusion $A_R\subset\Gamma_R^{\Pi,+}$ of active target degrees in the
source image is itself a residual condition, not a consequence of
the Borcherds product.

\subsubsection*{Local Pfaffian wall on type-II walls}

The (O2) obstruction concerns the local Pfaffian normal form on the
three simple type-II walls $\Hpl_{\delta_i}$, $i=1,2,3$, of the
Borcherds--Cartan datum $\mathscr B_{\Delta_5}$.

\begin{definition}[Rank-one type-II Pfaffian crossing]
\label{def:rank-one-type-ii-crossing}
For a simple type-II root $\delta$, let $T_\delta$ be a real analytic
tangent chart with trivial $s_\delta$-action, set
$U_\delta=T_\delta\times(-\varepsilon,\varepsilon)$ with
$s_\delta(t,u)=(t,-u)$, and let
$\mathscr L^{\mathrm{tan}}_\delta$ be an oriented tangent Pfaffian
line with $s_\delta$-invariant unit $\upsilon_\delta$.  The rank-one
odd Pfaffian crossing is the skew family
\[
D_u=\begin{pmatrix}0&u\\-u&0\end{pmatrix},
\qquad K^{\mathrm{cross}}_\delta=
[\R\xrightarrow{u}\R],
\]
in odd Pfaffian parity.  The total local Pfaffian section is
\[
\operatorname{pf}_\delta=\upsilon_\delta\,u,
\qquad
s_\delta^*\operatorname{pf}_\delta=-\operatorname{pf}_\delta,
\qquad
s_\delta^*(\operatorname{pf}_\delta^2)=\operatorname{pf}_\delta^2.
\]
\end{definition}

\begin{proposition}[Rank-one crossing gives the type-II local sign;
\textit{proved}]
\label{prop:rank-one-crossing-o2-sign}
For $i=1,2,3$, the rank-one crossing of
Definition~\ref{def:rank-one-type-ii-crossing} forms the finite
generic type-II local-sign datum with
$N^{\mathrm{Pf}}_{\delta_i}=1$, and at finite HN height $R$, with
successor maps preserving the tangent charts, the wall coordinate
$u_{\delta_i,R}$, and the invariant units, the (O2) residual entries
\[
\mathfrak o^{\mathrm{chart}}_{\delta_i,R},\
\mathfrak o^{\mathrm{wall}}_{\delta_i,R},\
\mathfrak o^{\mathrm{split}}_{\delta_i,R},\
\mathfrak o^{\mathrm{unit}}_{\delta_i,R},\
\mathfrak o^{\mathrm{rank}}_{\delta_i,R},\
\mathfrak o^{\mathrm{tr}}_{\delta_i,R'R}
\]
vanish for these local models, and
$\varepsilon^{\mathrm{loc}}_{\delta_i,R}=-1$.
\end{proposition}

\begin{proof}
The fixed locus of $s_\delta(t,u)=(t,-u)$ is $u=0$, so the wall and
the tangent fixed locus agree.  The determinant factor is split: the
tangent line is $\mathscr L^{\mathrm{tan}}_\delta$, and the normal
factor is $K^{\mathrm{cross}}_\delta$.  The unit $\upsilon_\delta$ is
invariant by construction.  The skew matrix has Pfaffian $u$, so the
normal-mode rank is one.  Reflection sends $u\mapsto -u$ and fixes
$\upsilon_\delta$; hence the local Pfaffian section changes sign and
its square is fixed.  The transition hypothesis is exactly
preservation of these data under successor maps.
\end{proof}

\subsubsection*{Half-Hilbert orbit and the type-II sign sum}

The (O2) obstruction also requires a wall \emph{atlas}: at finite
HN height $R$, each retained component and boundary stratum of
$\Hpl_{\delta_i}$ needs a chart of the rank-one form, and the
normal Pfaffian units, tangent/normal splittings, quotient
orientations, and protected integration maps must agree on chart
overlaps.  A single generic chart computes only a local sign.

The finite O2 wall-atlas proof packet is the table system
\[
\mathrm{wall\_objects},\quad
\mathrm{wall\_charts},\quad
\mathrm{overlaps},\quad
\mathrm{orbit\_transport},\quad
\mathrm{half\_hilbert\_orbit},\quad
\mathrm{compatibilities},\quad
\mathrm{scalar\_separation}.
\]
Its rows must carry the retained wall objects, rank-one chart
coordinates, Cech/cochain overlap equations, Weyl transport, the
twelve-coordinate half-Hilbert orbit, and the component, boundary,
Thom--Sebastiani, HN-transition, quotient-orientation, and protected
integration identities.  The scalar-separation rows record
\(4096=2^{12}=64^2\) without identifying the OP scalar normalization
or the Maass character value with wall-atlas data.

The half-Hilbert orbit of the type-II reflection group on the
relevant retained stratum carries a sign sum of size $4096=2^{12}$,
counting the supplied finite-window normal Pfaffian signs across the
twelve retained components and boundary strata of the type-II wall
configuration; this sum is the signature of the (O2) wall
normal-form datum.  Local rank-one charts contribute the local
factors; the chosen Weyl lifts convert the retained global system
into the orientation character on the
$W^{(2)}(\Lambda_{II}^{2,1})$-orbit.  The full normal-form
theorem with overlap, boundary, and component coherence is relative
to the retained \((O2_{\mathrm{atlas}})\) datum and
Theorem~\ref{thm:G-O2-orbit-transport} of
Appendix~\ref{app:G-O2-discharge}; the local model is the finite
rank-one wall factor.

\begin{proposition}[Generic local sign and the orientation character]
\label{prop:generic-local-sign-orientation-character}
Suppose (O1)$^+$ supplies a normalized lift
$\tau_{i,R,S}:o_{R,S}\to s_{\delta_i}^* o_{R,S}$ of the simple
type-II reflection on every retained orbit, and the (O2)/hybrid
residual for $\delta_i$ vanishes.  In a local chart with
$K^{\mathrm{nor}}_{\delta_i,R}\simeq[\R\xrightarrow{u_i}\R]$, with
normal Pfaffian unit $\upsilon_i$ and $s_{\delta_i}(\upsilon_i)
=\upsilon_i$, let $\ell_{\delta_i,R,S}$ be the path in the reduced
quotient wall-complement carrying $u_i=\epsilon$ to $u_i=-\epsilon$,
closed by $\tau_{i,R,S}$.  Then
\[
M_{\ell_{\delta_i,R,S}}(\upsilon_i u_i)=\upsilon_i(-u_i)
=-\upsilon_i u_i,
\qquad
\epsilon_{o,R}(s_{\delta_i})=(-1)^{N^{\mathrm{Pf}}_{\delta_i,R}}.
\]
For the rank-one type-II model
$\epsilon_{o,R}(s_{\delta_i})=-1=\nu_{\Delta_5}(s_{\delta_i})$.  The
global statement
$\epsilon_o|_{W^{(2)}(\Lambda_{II}^{2,1})}=\nu_{\Delta_5}|
_{W^{(2)}(\Lambda_{II}^{2,1})}$
is the chamber-permuting half of the Pfaffian--Dirac orientation
identity, conditional on the (O2) wall atlas.
\end{proposition}

\subsubsection*{Pfaffian normalization and trace}

The supplied Pfaffian-to-automorphic comparison
$\iota_{\mathrm{aut}}$ then gives, on the imported Borcherds product
side,
\[
\iota_{\mathrm{aut}}\bigl(\operatorname{Pf}_{\mathrm{prot}}
(\mathfrak D_X)\bigr)=\mathcal D_X=\Delta_5,
\]
with cusp-side expansion
\[
\operatorname{pf}_X|_{\mathfrak c_\infty}
=64q^{1/2}r^{1/2}s^{1/2}\prod_{\gamma=(n,l,m)\in
\Gamma_{\mathrm{eff}}}(1-q^nr^ls^m)^{f(nm,l)},
\]
matching the imported Borcherds--Gritsenko--Nikulin product.  The
scalar branch obtained after taking the protected trace of the
inverse squared Pfaffian section is the Oberdieck--Pixton
normalization
\[
Z^X_{\mathrm{OP}}=-4096\,\Delta_5^{-2},
\]
where the negative sign is the OP scalar branch and
\(4096=2^{12}=64^2\) is the scalar equality between the retained
half-Hilbert sign count and the squared theta-leading coefficient.
Neither number defines the orientation character.

\subsubsection*{Comparison with the Pfaffian identity}

The Vol II primitive-to-shadow comparison singles out the protected Pfaffian
$\operatorname{Pf}_{\mathrm{prot}}(\mathfrak D_X)=\Delta_5$ as the
target identity for the retained Pfaffian datum; the \(K3\times E\)
side supplies the Pfaffian line and scalar branch data.
BBDJS \cite{BBDJS2015} provides the orientation-line and Pfaffian
calculus on the cosection-reduced d-critical heart; Joyce--Upmeier
\cite{JoyceUpmeier} provides multiplicativity under reduced
extension correspondences; Joyce--Tanaka--Upmeier
\cite{JoyceTanakaUpmeier} extends the framework to $4$-fold
analogues but is not used directly here.  The constant \(4096=2^{12}\)
is fixed at the orientation level by the retained half-Hilbert wall
atlas and at the scalar level by \(64^2\); the local sign \(-1\) is
fixed by the rank-one Pfaffian wall computation.  The
Oberdieck--Pixton minus sign is a scalar branch convention, not the
source of the orientation character.  No
identification of the Pfaffian $\Delta_5$ with a compact microscopic
state space is asserted: \(\Delta_5\) is the protected Pfaffian section,
the scalar trace is \(\Delta_5^{-2}\), and the operator-level datum is the
pro-object $\mathfrak D_X^{\mathrm{DI}}$, defined at the Pfaffian and
orientation level relative to the retained data of
Appendix~\ref{app:obstruction-discharges}.

The Pfaffian recovers $\Delta_5$ at the target scalar level, but
the cross-level recognition that identifies the primitive Hopf
algebra of \S\ref{subsec:hall-correspondences} with the Borcherds--Kac--Moody
target $\autcor$ requires a separately constructed source coalgebra
and an explicit chiral-Koszul cobar comparison.  The source coalgebra
$C_X$ is the bar of the augmented binary/two-step hybrid correspondence
algebra; it becomes the bar of a hybrid factorization algebra only
after the higher-coloured residual vanishes.  The cobar map
$\Theta_{\mathrm{Kos},R}$ is the quasi-isomorphism to the imported
current envelope $\mathsf{Den}_{\Delta_5,E}$.  This Koszul lane
$(\mathrm L 5)$ is the content of
\S\ref{subsec:koszul-source}.


\subsection{Koszul source coalgebra and chiral MC twisting}
\label{subsec:koszul-source}

The fifth entry of the Dirac--Igusa pro-object is the Koszul lane:
the source coalgebra $C_X$, the bar comparison $b_X$, and the cobar
quasi-isomorphism $\Theta_{\mathrm{Kos}}$ to the imported current
envelope $\mathsf{Den}_{\Delta_5,E}=U_E^{\mathrm{ch}}(\operatorname{Cur}_E(\mathfrak g_{\Delta_5}))$,
the chiral enveloping algebra of the level-zero current algebra of the
BKM denominator algebra over $E$
(Definition~\ref{def:bkm-formal-current-envelope}).  Source--target
separation is structural.  The target counit
\[
\varepsilon_{\Delta,R}:
\Omega_E^{\mathrm{ch}}\bar B_E^{\mathrm{ch}}
(\mathsf{Den}_{\Delta_5,E,\le R})
\longrightarrow
\mathsf{Den}_{\Delta_5,E,\le R}
\]
is internal to the imported denominator current envelope and produces
no source data.  The compact Koszul comparison
$\Omega_E^{\mathrm{ch}}C_X\simeq\mathsf{Den}_{\Delta_5,E}$ is a
source-to-target statement on the supplied hybrid source, conditional
on the chiral Koszul source datum, the finite Koszul comparison datum,
and strict Koszul Mittag--Leffler exactness.

The Beilinson cut applies in its purest form on this lane.  $\bar
B^{\mathrm{ch}}$ is MC twisting / coupling coalgebra; it is not the
level-\(\mathsf Z\) bulk.  The target of the lane is
$\mathsf{Den}_{\Delta_5,E}$, the chiral shadow of the BKM denominator
algebra; the source coalgebra $C_X$ is the supplied compact Hall data
after collision coproduct, conilpotent filtration, and bar-length
truncation, and the cobar map is the comparison.

\subsubsection*{Source--target separation}

\begin{lemma}[Source--target separation for the Koszul lane;
\textit{proved}]
\label{lem:source-target-koszul-separation}
Fix a finite height $R$.  The target counit
$\varepsilon_{\Delta,R}$ above is a morphism internal to the imported
current envelope.  A source Koszul datum is different data: it must
first construct
\[
\mathcal F_{X,\sigma,S,\le R}^{\mathrm{hyb}},\qquad
C_{X,R},\qquad F^{\mathrm{co}}_{R,\bullet},\qquad
\Delta_R^{\mathrm{ch}}
\]
from the finite hybrid correspondence category
$\mathsf{Corr}^{E,\mathrm{hyb}}_R$, the quotient-after-correspondence
pseudofunctor inducing $Q_{E,R}$, orientations, the eight-word two-step
binary flag atlas, and the separate higher-coloured configuration
residual; only then may it
compare
\[
b_{X,R}:C_{X,R}\simeq
\bar B_E^{\mathrm{ch}}\bigl(\mathcal F_{X,\sigma,S,\le R}^{\mathrm{hyb}},
\varepsilon_{X,R}^{\mathrm{hyb}}\bigr),
\quad
\Theta_{\mathrm{Kos},R}:\Omega_E^{\mathrm{ch}}C_{X,R}\simeq
\mathsf{Den}_{\Delta_5,E,\le R}.
\]
Setting $C_{X,R}=\bar B_E^{\mathrm{ch}}(\mathsf{Den}_{\Delta_5,E,\le R})$,
or choosing a homotopy inverse to $\varepsilon_{\Delta,R}$, does not
produce a source datum.
\end{lemma}

\begin{proof}
The domain of $\varepsilon_{\Delta,R}$ contains only target
augmentation, target charge filtration, and target chiral
operations; it does not contain rigidified wrapped prequotients,
anchors, mixed local/wrapped correspondences,
quotient-after-correspondence pseudofunctor, orientation transport,
protected integration, or primitive representatives.  These are the
data from which the source coalgebra and collision coproduct are
defined.
\end{proof}

\subsubsection*{Chiral Koszul source datum}

\begin{definition}[Chiral Koszul source datum]
\label{def:chiral-koszul-source-datum}
Fix $(\sigma,S,R)$ and the Gram image $\Gamma_R^\Pi$.  A finite-stage
chiral Koszul source datum is
\[
\mathfrak K^{\mathrm{Kos}}_{X,R}
=\bigl(
\mathcal F_{X,\sigma,S,\le R}^{\mathrm{hyb}},
C_{X,R},
F^{\mathrm{co}}_{R,\bullet},
\Delta_R^{\mathrm{ch}},
b_{X,R},
\Theta_{\mathrm{Kos},R}\bigr)
\]
with the following clauses; every entry is source-side finite data
except for the terminal comparison with the imported target.
\begin{enumerate}
\item \emph{Source chiral object.}  $\mathcal F_{X,\sigma,S,\le R}^{\mathrm{hyb}}$
is the finite hybrid object of \S\ref{subsec:hybrid-ran-carrier},
with $\mathfrak O_{\mathrm{hyb},R}=0$ and with a supplied wrapped
bar-cobar domain datum.  The wrapped sector is not made Koszul by the
projection-finite \(b=0\) Ran sector.
\item \emph{Conilpotent chiral coalgebra.}  $C_{X,R}$ is a coaugmented
conilpotent chiral coalgebra on $E$, filtered by a finite increasing
source charge filtration $F^{\mathrm{co}}_{R,\bullet}$; reduced iterated
coproduct nilpotent on each finite stage.
\item \emph{Collision coproduct.}  $\Delta_R^{\mathrm{ch}}$ is the
collision coproduct from the finite Hall correspondences (local-local,
ordered local-wrapped, wrapped-wrapped, and eight-word flag stacks),
defined before $Q_{E,R}$; the quotient coalgebra identities are the
vanishing of the higher-coloured, unit, refinement, and quotient
residuals.  Transition maps $\rho^C_{R'R}$ commute with
$\Delta^{\mathrm{ch}}$.
\item \emph{Bar comparison.}  $b_{X,R}:C_{X,R}\simeq
\bar B_E^{\mathrm{ch}}(\mathcal F_{X,\sigma,S,\le R}^{\mathrm{hyb}},
\varepsilon_{X,R}^{\mathrm{hyb}})$
is a charge-preserving quasi-isomorphism of conilpotent chiral
coalgebras compatible with primitive projection, the completed Hopf
pairing, the reduced $E$-quotient, and the normal-ordered descent
$\overline\Pi_{X,*}^\Theta$, separately on the projection-finite and
wrapped summands and on their mixed words.
\item \emph{Cobar comparison with the target envelope.}
$\Theta_{\mathrm{Kos},R}:\Omega_E^{\mathrm{ch}}C_{X,R}\simeq
\mathsf{Den}_{\Delta_5,E,\le R}$
is a charge-preserving quasi-isomorphism of augmented chiral
algebras.  Restricted to constant-current primitives after the
Hopf-radical quotient it gives the primitive recognition morphism
\[
\iota_R^{\Prim}:H^0\Prim_{\mathrm{prot}}(\overline\Pi_{X,*}^\Theta
\mathcal F_{X,\sigma,S,\le R}^{\mathrm{hyb}})
\longrightarrow\mathfrak g_{\Delta_5,\le R}.
\]
\end{enumerate}
The exact finite Koszul obstruction tuple is
\[
\mathfrak O_{\mathrm{Kos},R}
=(\mathfrak o^C_R,\mathfrak o^{\mathrm{fil}}_R,
\mathfrak o^{\mathrm{conil}}_R,
\mathfrak o^{\Delta,\mathrm{ch}}_R,
\{\mathfrak o^{C,\mathrm{tr}}_{R'R}\}_{R'\ge R},
\mathfrak o^b_R,\mathfrak o^\Omega_R,
\mathfrak o^{\mathrm{wrKos}}_R,
\mathfrak o^{\mathrm{hyb}}_R,
\mathfrak o^\Pi_R,
\mathfrak o^{\mathrm{sep}}_R,
\mathfrak o^{\Prim}_R).
\]
\end{definition}

\subsubsection*{Canonical finite source-bar coalgebra}

For a hybrid source with bounded-length augmentation ideal, a
canonical model exists.

\begin{proposition}[Canonical finite source-bar coalgebra;
\textit{conditional on $\mathfrak O_{\mathrm{hyb},R}=0$ and bounded
length}]
\label{prop:canonical-source-bar-coalgebra}
Suppose $\mathfrak O_{\mathrm{hyb},R}=0$ and there is an integer $N_R$
such that every nonzero ordered word of non-vacuum hybrid inputs of
total HN height $\le R$ has length at most $N_R$, with HN height
positive on every non-vacuum colour surviving in the reduced bar
construction.  Suppose an augmentation
$\varepsilon_{X,R}^{\mathrm{hyb}}:\mathcal F_{X,\sigma,S,\le R}^{\mathrm{hyb}}\to k\mathbf 1$
is supplied.  Then the canonical source-bar coalgebra
\[
C^{\mathrm{bar}}_{X,R}
=\bar B_E^{\mathrm{ch}}\bigl(\mathcal F_{X,\sigma,S,\le R}^{\mathrm{hyb}},
\varepsilon_{X,R}^{\mathrm{hyb}}\bigr)
=k\mathbf 1\oplus\bigoplus_{1\le p\le N_R}
\bar B^{\mathrm{ch},p}_E(\overline{\mathcal F}_{X,R}),
\]
filtered by bar-length, with deconcatenation collision coproduct
\[
\Delta_R^{\mathrm{bar}}([x_1|\cdots|x_p])=
\mathbf 1\otimes[x_1|\cdots|x_p]
+[x_1|\cdots|x_p]\otimes\mathbf 1
+\sum_{i=1}^{p-1}[x_1|\cdots|x_i]\otimes[x_{i+1}|\cdots|x_p],
\]
realizes the coalgebra and source-bar part of a chiral Koszul source
datum with
$\mathfrak o^C_R=\mathfrak o^{\mathrm{fil}}_R=
\mathfrak o^{\mathrm{conil}}_R=\mathfrak o^{\Delta,\mathrm{ch}}_R=0$
and $b_{X,R}=\mathrm{id}$.
\end{proposition}

\begin{proof}
Apply the Beilinson--Drinfeld chiral bar construction
\cite[\S3.7]{BD} to the supplied hybrid object.  Since the finite HN
quotient kills charges of height $>R$ and non-vacuum words have
length $\le N_R$, only finitely many bar lengths occur.  The
deconcatenation coproduct on each ordered word is realized by
cutting the corresponding ordered collision flag in
$\mathsf{Corr}^{E,\mathrm{hyb}}_R$; admissibility of pull-push and
Thom--Sebastiani transport are part of $\mathfrak O_{\mathrm{hyb},R}=0$.
The eight-word flag pentagon gives binary coassociativity; the
higher-coloured, unit, refinement, and quotient coherences contained in
$\mathfrak O_{\mathrm{hyb},R}=0$ give the chiral coalgebra identities.
After at most $p\le N_R$ reduced iterations, every branch has length one and
the next reduced coproduct is zero, so $(\overline\Delta_R^{\mathrm{bar}})^{N_R+1}=0$.
\end{proof}

The finite packet
\(\texttt{certificates/charge/bar\_pushforward\_commutation}\) supplies
only the grading fact used in the preceding construction: for the
additive normal-ordered map
\(\overline\Pi_X:\widehat\Gamma_X\to\Gamma_{\mathrm{gram}}\), finite
bar words, deconcatenation, and multiplication terms have the same
Gram degree before and after pushforward.  It does not supply
\(\mathfrak O_{\mathrm{hyb},R}=0\), the wrapped bar-cobar domain, the
source-internal counit, or the source-to-target chiral
quasi-isomorphism.

\subsubsection*{Cobar obstruction}

The canonical source gives a conilpotent chiral bar coalgebra.  A
cobar quasi-isomorphism to the denominator chiral algebra is an
additional datum, consisting of a source-internal bar-cobar counit and
a source-to-target chiral algebra quasi-isomorphism.

\begin{corollary}[Source-to-target cobar obstruction;
\textit{proved}]
\label{cor:source-bar-remaining-cobar-obstruction}
Under the hypotheses of
Proposition~\ref{prop:canonical-source-bar-coalgebra}, a cobar
quasi-isomorphism $\Theta_{\mathrm{Kos},R}$ requires:
\begin{enumerate}
\item \emph{Source-internal bar-cobar counit.}  The supplied hybrid
source, including the wrapped summands and mixed local/wrapped words,
must lie in the Beilinson--Drinfeld/Francis--Gaitsgory completed
bar-cobar domain \cite[Theorem~6.2.2]{FrancisGaitsgoryCKD}, with the
source-internal counit
\[
\varepsilon^{\mathrm{src}}_{\mathrm{bar},R}:
\Omega_E^{\mathrm{ch}}C^{\mathrm{bar}}_{X,R}
\longrightarrow\mathcal F_{X,\sigma,S,\le R}^{\mathrm{hyb}}
\]
a quasi-isomorphism in the chosen completed/coderived conilpotent
category.  Defects:
$\mathfrak o^{\varepsilon,\mathrm{src}}_R=
(\mathfrak o^{\mathrm{aug}}_R,\mathfrak o^{\mathrm{FG}}_R,
\mathfrak o^{\mathrm{wrKos}}_R,
H^\bullet\operatorname{Cone}\varepsilon^{\mathrm{src}}_{\mathrm{bar},R})$.
\item \emph{Source-to-target chiral algebra quasi-isomorphism.}
$\vartheta_R:\mathcal F_{X,\sigma,S,\le R}^{\mathrm{hyb}}
\xrightarrow{\simeq}\mathsf{Den}_{\Delta_5,E,\le R}$
preserving differential, product, unit, transition maps, and
primitive restriction.  Defects:
\[
\mathfrak o^\vartheta_R=\bigl(d_{\Delta}\vartheta_R-\vartheta_R d_{\mathrm{hyb}},
\vartheta_R\mu_R^{\mathrm{hyb}}-\mu_{\Delta,R}(\vartheta_R\otimes\vartheta_R),
\vartheta_R\eta_R^{\mathrm{hyb}}-\eta_{\Delta,R},
\vartheta_R\circ\tau^{E,\mathrm{hyb}}_{R'R}-\pi^\Delta_{R'R}\circ\vartheta_{R'},
H^\bullet\operatorname{Cone}\vartheta_R,
\vartheta_R|_{\Prim}-\iota_R^{\Prim}\bigr).
\]
\end{enumerate}
Then $\Theta_{\mathrm{Kos},R}=\vartheta_R\circ
\varepsilon^{\mathrm{src}}_{\mathrm{bar},R}$, and the surviving cobar
residual is $\mathfrak o^\Omega_R=(\mathfrak o^{\varepsilon,\mathrm{src}}_R,
\mathfrak o^\vartheta_R)$.  No homotopy inverse to the target counit
$\varepsilon_{\Delta,R}$ enters this construction.
\end{corollary}

\begin{definition}[Finite Koszul comparison datum]
\label{def:finite-koszul-comparison-datum}
Fix a finite HN height \(R\), and suppose the finite-stage chiral
Koszul source datum
\(\mathfrak K^{\mathrm{Kos}}_{X,R}\) of
Definition~\ref{def:chiral-koszul-source-datum} is fixed.  A finite
Koszul comparison datum is
\[
\mathfrak K^{\mathrm{cmp}}_{X,R}
=\bigl(
\varepsilon^{\mathrm{src}}_{\mathrm{bar},R},
\vartheta_R,
H^{W}_{R},
H^{\mathrm{Pf}}_{R},
H^{\mathrm{Hall}}_{R},
H^{\mathrm{rad}}_{R},
H^{\mathrm{PBW}}_{R},
H^{\mathrm{tr}}_{R'R}\bigr)
\]
with the following finite chain-level clauses.
\begin{enumerate}
\item[\textup{(K0)}] \emph{Two quasi-isomorphism cones.}
The source-internal bar-cobar counit
\[
\varepsilon^{\mathrm{src}}_{\mathrm{bar},R}:
\Omega_E^{\mathrm{ch}}C_{X,R}\longrightarrow
\mathcal F_{X,\sigma,S,\le R}^{\mathrm{hyb}}
\]
and the source-to-target chiral algebra morphism
\[
\vartheta_R:\mathcal F_{X,\sigma,S,\le R}^{\mathrm{hyb}}
\longrightarrow \mathsf{Den}_{\Delta_5,E,\le R}
\]
are supplied with finite filtrations by HN height, bar length,
local/wrapped word, and parity.  Their cone cohomology groups vanish in
each filtered piece.
\item[\textup{(K1)}] \emph{Weyl compatibility.}
For every retained type-II simple reflection \(s_{\delta_i}\) and
chosen lift \(\tau_{i,R}\), the square comparing the source Weyl
transport on \(C_{X,R}\) with the target Weyl action on
\(\mathsf{Den}_{\Delta_5,E,\le R}\) commutes up to a specified
homotopy \(H^{W}_{i,R}\).  The Coxeter and chamber-overlap defect
complexes of these homotopies are acyclic.
\item[\textup{(K2)}] \emph{Pfaffian orientation compatibility.}
The determinant line of the finite cobar comparison is identified with
the Pfaffian-line transport of \((O1)^+\) and the
\((O2_{\mathrm{atlas}})\) wall datum.  On each type-II wall chart the
induced sign is the rank-one Pfaffian sign
\(\epsilon_{o,R}(s_{\delta_i})\); after the target comparison it is
\(\nu_{\Delta_5}(s_{\delta_i})\).  The difference is the finite defect
of \(H^{\mathrm{Pf}}_{R}\), required to vanish.
\item[\textup{(K3)}] \emph{Hall product, coproduct, and pairing
compatibility.}
On primitive classes, \(\vartheta_R\circ
\varepsilon^{\mathrm{src}}_{\mathrm{bar},R}\) intertwines the Hall
supercommutator, collision coproduct, unit, counit, and Hopf pairing
with the corresponding operations in the finite target presentation.
Equivalently the finite matrices \(M,D,\eta,\epsilon,B,G\) of
Definition~\ref{def:finite-source-proof-tables} are carried to the
target product, coproduct, bracket, and invariant pairing blocks.
\item[\textup{(K4)}] \emph{Radical quotient compatibility.}
The primitive restriction maps the finite source Hopf radical onto the
finite GN/Kac radical block and induces the comparison on radical
quotients:
\[
\operatorname{Rad}^{\Pi}_{X,R}
\xrightarrow{\ \Theta_{\mathrm{Kos},R}|_{\Prim}\ }
\operatorname{Rad}_{\mathrm{GN},R},
\qquad
P^{\Pi,\mathrm{red}}_{X,R}
\xrightarrow{\ \iota_R^{\Prim}\ }
\mathfrak g_{\Delta_5,\le R}.
\]
The kernel and cokernel defect complexes of this quotient square are
zero.
\item[\textup{(K5)}] \emph{PBW and primitive-recognition
compatibility.}
The primitive restriction of
\(\Theta_{\mathrm{Kos},R}\) is the same map
\(\iota_R^{\Prim}\) used in primitive recognition, and it preserves the
PBW filtrations on the source and target enveloping algebras.  The
associated-graded comparison agrees with the PBW rows of
Definition~\ref{def:finite-source-proof-tables}.
\item[\textup{(K6)}] \emph{Strict transitions.}
For \(R'\ge R\), the transition homotopy \(H^{\mathrm{tr}}_{R'R}\)
identifies the two composites obtained by first applying
\(\Theta_{\mathrm{Kos},R'}\) and then truncating, or first truncating
the source data and then applying \(\Theta_{\mathrm{Kos},R}\).  The
same strictness holds for the Weyl, Pfaffian, Hall-pairing, radical,
and PBW compatibility complexes.
\item[\textup{(K7)}] \emph{Koszul Mittag--Leffler exactness with
compatibility.}
The inverse systems of the two quasi-isomorphism cones and of all
defect complexes in \((K1)\)--\((K6)\) are strict
Mittag--Leffler systems; \(R^1\varprojlim=0\) for each of them.
\end{enumerate}
The finite Koszul comparison residual
\[
\mathfrak O_{\mathrm{Kcmp},R}
=\bigl(
H^\bullet\Cone\varepsilon^{\mathrm{src}}_{\mathrm{bar},R},
H^\bullet\Cone\vartheta_R,
\mathfrak o^W_R,\mathfrak o^{\mathrm{Pf}}_R,
\mathfrak o^{\mathrm{Hall}}_R,\mathfrak o^{\mathrm{rad}}_R,
\mathfrak o^{\mathrm{PBW}}_R,\mathfrak o^{\mathrm{tr}}_{R'R},
\mathfrak o^{\mathrm{ML}}_{\mathrm{Kcmp}}\bigr)
\]
vanishes exactly when \((K0)\)--\((K7)\) hold.  A target-internal
bar-cobar counit, a primitive restriction, or a scalar denominator
identity supplies none of these compatibility homotopies.

After choosing finite homogeneous bases, the datum is recorded by three
source tables:
\[
\texttt{koszul\_cones},\qquad
\texttt{koszul\_comparison\_identities},\qquad
\texttt{koszul\_transition\_ml}.
\]
The first table records the two acyclic cones
\(\texttt{source\_bar\_cobar\_counit}\) and
\(\texttt{source\_to\_target\_quasi\_isomorphism}\), with zero
cohomology rank in every HN, bar-length, word, and parity block.  The
second table records zero finite defects for
\[
\texttt{weyl\_action},\quad
\texttt{pfaffian\_orientation},\quad
\texttt{hall\_product},\quad
\texttt{hall\_coproduct},\quad
\texttt{hopf\_pairing},\quad
\texttt{radical\_quotient},\quad
\texttt{pbw}.
\]
The third table records strict transition and Mittag--Leffler
vanishing for
\[
\texttt{source\_cone},\quad \texttt{target\_cone},\quad
\texttt{weyl\_action},\quad \texttt{pfaffian\_orientation},\quad
\texttt{hall\_pairing},\quad \texttt{radical\_quotient},\quad
\texttt{pbw},
\]
with \(R^1\varprojlim\)-rank zero.  These rows require geometric source
provenance and proof references, just as the Hall and
primitive-recognition tables do.
\end{definition}

\begin{proposition}[Finite Koszul comparison compatibility;
\textit{conditional on \(\mathfrak K^{\mathrm{cmp}}_{X,R}\)}]
\label{prop:finite-koszul-comparison-compatibility}
If the finite Koszul comparison datum
\(\mathfrak K^{\mathrm{cmp}}_{X,R}\) is supplied and
\(\mathfrak O_{\mathrm{Kcmp},R}=0\), then
\[
\Theta_{\mathrm{Kos},R}
=\vartheta_R\circ\varepsilon^{\mathrm{src}}_{\mathrm{bar},R}
:\Omega_E^{\mathrm{ch}}C_{X,R}\xrightarrow{\simeq}
\mathsf{Den}_{\Delta_5,E,\le R}
\]
is a quasi-isomorphism of augmented chiral algebras whose primitive
restriction respects the Weyl action, the Pfaffian orientation
character, the Hall product/coproduct, the Hopf pairing, the radical
quotient, and the PBW associated graded on the finite window.
\end{proposition}

\begin{proof}
Clause (K0) gives the quasi-isomorphism by acyclicity of the two cones.
Clauses (K1)--(K2) identify the Weyl and Pfaffian-orientation squares,
with zero Coxeter and wall-sign defects.  Clause (K3) says that the
source Hall operations and Hopf pairing are carried to the finite
target blocks, so the primitive bracket and pairing used by
recognition are the restrictions of the same cobar comparison.  Clause
(K4) identifies the radical quotient, and (K5) identifies the PBW
associated gradeds.  The conclusion is therefore a finite diagram
chase in the common homogeneous bases.  No scalar Borcherds product or
target-internal counit participates in these diagrams.
\end{proof}

\begin{proposition}[Primitive restriction is not a Koszul
quasi-isomorphism; \textit{proved}]
\label{prop:primitive-restriction-not-koszul}
Let
\[
\vartheta_R:\mathcal F_{X,\sigma,S,\le R}^{\mathrm{hyb}}
\longrightarrow \mathsf{Den}_{\Delta_5,E,\le R}
\]
be a morphism of augmented chiral algebras whose restriction to the
constant-current primitive summand agrees with
\(\iota_R^{\Prim}\).  This primitive restriction does not imply that
\(\vartheta_R\) is a quasi-isomorphism.  The required finite
obstruction is the full cohomology
\[
H^\bullet\operatorname{Cone}(\vartheta_R)
\]
on every HN weight, bar length, local/wrapped word, and parity block.
Consequently the chiral Koszul comparison
\(\Theta_{\mathrm{Kos},R}\) is proved only after the source-internal
counit cone and the source-to-target cone are acyclic, compatibly in
\(R\).
\end{proposition}

\begin{proof}
The assertion already fails for ordinary augmented dg algebras, hence
also for constant chiral algebras on \(E\).  Let \(B\) be any augmented
dg algebra and let \(M\) be a non-zero dg \(B\)-bimodule placed in a
nonprimitive HN weight, with zero differential.  The square-zero
extension \(A=B\oplus M\) has multiplication
\[
(b,m)(b',m')=(bb',bm'+mb'),\qquad MM=0,
\]
and the projection \(p:A\to B\) is a morphism of augmented dg
algebras.  On the chosen primitive summand of \(B\subset A\), \(p\) is
the identity.  But
\[
H^\bullet\operatorname{Cone}(p)\cong H^\bullet(M)[1]\ne0.
\]
Thus agreement on primitives detects only the selected weight-one
summand.  It does not see square-zero, wrapped, mixed, higher
bar-length, or parity-cancelling classes in the source.  In the
Dirac--Igusa notation these missing classes are exactly the cone
terms in
\(\mathfrak o^{\varepsilon,\mathrm{src}}_R\) and
\(\mathfrak o^\vartheta_R\), together with the Koszul
Mittag--Leffler obstruction in the inverse system.
\end{proof}

\subsubsection*{Class M completion and chiral Koszulness}

The bar/cobar comparison takes place in the completed ambient
category.  Class~M holds in the completed ambient: analytic
HS-sewing, coderived BV $=$ bar, weight-completed, pro-, $J$-adic.
On chain-level ordinary complexes, the chiral bar/cobar pair is the
comparison input, not the closed Koszulness theorem; chiral
Koszulness in the homotopy setting is a separate theorem of
Francis--Gaitsgory \cite{FrancisGaitsgoryCKD}.

\subsubsection*{Koszul Mittag--Leffler exactness}

The completed bar/cobar inverse system is controlled by a
finite-slice Mittag--Leffler condition on every relevant tower.

\begin{proposition}[Consequence of the chiral Koszul source datum;
\textit{conditional on Koszul Mittag--Leffler}]
\label{prop:chiral-koszul-source-datum-consequence}
Assume $\mathfrak K^{\mathrm{Kos}}_{X,R}$ is supplied for all $R$ with
compatible transition maps and
$\mathfrak o^{C,\mathrm{tr}}_{R'R}=0$, and that the Koszul
Mittag--Leffler condition holds: $R^1\varprojlim=0$ for the inverse
systems of source coalgebras, completed chiral bar/cobar
constructions, target truncations, and the cones of $b_{X,R}$ and
$\Theta_{\mathrm{Kos},R}$ \cite[\S3.5]{WeibelHA}.  Then
$C_X=\varprojlim_R C_{X,R}$ is a completed pro-conilpotent chiral
coalgebra attached to the hybrid compact source, conilpotent on
every finite HN quotient, and the maps $b_{X,R}$ and
$\Theta_{\mathrm{Kos},R}$ induce
\[
b_X:C_X\xrightarrow{\simeq}
\bar B_E^{\mathrm{ch}}(\mathcal F_X^{\mathrm{hyb}}),
\qquad
\Theta_{\mathrm{Kos}}:\Omega_E^{\mathrm{ch}}C_X
\xrightarrow{\simeq}\mathsf{Den}_{\Delta_5,E}.
\]
If, in addition, the finite Koszul comparison data
\(\mathfrak K^{\mathrm{cmp}}_{X,R}\) of
Definition~\ref{def:finite-koszul-comparison-datum} are supplied for a
cofinal HN subsystem and satisfy the strict compatibility
Mittag--Leffler condition \((K7)\), then the inverse-limit map
\(\Theta_{\mathrm{Kos}}\) respects the Weyl action, Pfaffian
orientation character, Hall product/coproduct, Hopf pairing, radical
quotient, and PBW associated graded.
This identifies the supplied source cobar algebra with the formal
target current envelope; it does not construct the $E_3$-source, the
hybrid wrapped Hall correspondences, the orientation line, the
primitive recognition data, or $C_X$ itself from the target-internal
counit.
\end{proposition}

\begin{proof}
Inverse-limit existence follows from transition compatibility, the
bounded exhaustive finite charge filtration, finite coassociative and
counital collision coproduct, and continuity hypotheses.  The
coradical filtration is finite on every $\Gamma_R^\Pi$-truncation;
the Koszul Mittag--Leffler hypothesis kills $R^1\varprojlim$ on the
cones of $b_{X,R}$ and $\Theta_{\mathrm{Kos},R}$, giving completed
quasi-isomorphisms.  If the finite comparison data are present, the
same argument applied to the Weyl, Pfaffian, Hall-pairing, radical, and
PBW defect complexes in \((K1)\)--\((K7)\) passes the zero finite
defects to the inverse limit.  Without that exactness one obtains only
a compatible pro-quasi-isomorphism, with no theorem-level control of
these recognition compatibilities.
\end{proof}

\subsubsection*{Chiral MC twisting / coupling coalgebra}

Bar is twisting; bar is not the level-\(\mathsf Z\) bulk.  The reduced bar coalgebra
$\bar B_E^{\mathrm{ch}}(\mathcal F_X^{\mathrm{hyb}})$ is the
Maurer--Cartan coupling coalgebra of the supplied hybrid chiral
algebra: its primitive elements are the chain-level inputs of MC
twisting in the Beilinson--Drinfeld/Francis--Gaitsgory framework,
and its grouplike completion records the homotopy-coherent twisting
data.  The level-\(\mathsf Z\) bulk in Vol~II is
$Z^{\mathrm{der}}_{\mathrm{ch}}(A)$ for the relevant boundary algebra
$A$; the visible target is its BKM chiral shadow
$\mathsf{Den}_{\Delta_5,E}$.  The Koszul comparison maps $C_X$, an
instance of MC coupling, into that target chiral algebra by the cobar
construction.

\subsubsection*{The Koszul obstruction class}

The Koszul lane carries the residuals
\[
\mathfrak O_{\mathrm{Kos},R},\qquad
\mathfrak O_{\mathrm{Kcmp},R},\qquad
\mathfrak o^{\mathrm{ML},\mathrm{Kos}}_{R'R}.
\]
With $\mathfrak O_{\mathrm{hyb},R}=0$ and the canonical source-bar
construction, the coalgebra defects vanish; only the cobar residual
$\mathfrak o^\Omega_R=(\mathfrak o^{\varepsilon,\mathrm{src}}_R,
\mathfrak o^\vartheta_R)$ remains, with the wrapped Koszul defect
\(\mathfrak o^{\mathrm{wrKos}}_R\) included in
\(\mathfrak o^{\varepsilon,\mathrm{src}}_R\).  Vanishing of the source-internal
counit defect and construction of the source-to-target chiral algebra quasi-isomorphism
operates in the completed Beilinson--Drinfeld/Francis--Gaitsgory
bar-cobar domain.  The comparison residual
\(\mathfrak O_{\mathrm{Kcmp},R}\) is separate: it records whether this
quasi-isomorphism respects the Weyl action, Pfaffian orientation,
Hall pairing, radical quotient, and PBW filtration used by primitive
recognition.  The Koszul Mittag--Leffler hypothesis controls the
inverse-system passage for both the quasi-isomorphism cones and these
compatibility complexes.

\subsubsection*{Source-target separation}

The source--target separation is the $K3\times E$-side instance of
the Vol II Beilinson cut: bar is twisting, bar is not bulk; the
boundary algebra is $A=\mathcal F_X^{\mathrm{hyb}}$ at finite stage,
and the bulk is $Z^{\mathrm{der}}_{\mathrm{ch}}(A)$ which receives
the cobar of the bar coalgebra after MC coupling.  Vol I master
patterns prohibit the identification of the chiral bar with the bulk
chiral algebra and require the Koszul comparison to be a separate
theorem in the homotopy setting; the same separation is enforced
above.  The compact Koszul comparison and the chiral
Koszulness theorem are not identical; the present construction lives
on the comparison side, with chiral Koszulness in the
Francis--Gaitsgory sense as the abstract input.

Under the Koszul source datum and the finite Koszul comparison datum,
\(\Theta_{\mathrm{Kos},R}\) identifies
the cobar of the source coalgebra with the imported current envelope of
\(\autcor\), compatibly with the Weyl, Pfaffian, Hall-pairing,
radical, and PBW structures.  Under the recognition datum
\((S1)\)--\((S10)\), the
source-side protected primitive \(P_X^{\Pi,\mathrm{red}}\), extracted
from \(C_X\) via the Hopf radical quotient, the Frobenius identity, and
the no-extra-relations criterion, is identified with the
Borcherds--Kac--Moody algebra \(\mathfrak g_{\Delta_5}\).  Entry
\((\mathrm L 6)\) is this recognition apparatus, in
\S\ref{subsec:primitive-recognition}.


\subsection{Primitive recognition apparatus}
\label{subsec:primitive-recognition}

The sixth entry of the Dirac--Igusa pro-object is the primitive
recognition criterion.  Its conclusion, after the compact source data
are supplied, is an isomorphism of Gram-graded Lie superalgebras
\[
P_X^{\Pi,\mathrm{red}}\cong\mathfrak g_{\Delta_5}
\]
between the cosection-twisted reduced Hall primitive layer of
$K3\times E$ and the imported Borcherds--Kac--Moody denominator
algebra.  Signed superdimensions alone do not suffice for this
theorem.  The Borcherds product fixes
$\sdim P^{\Pi,\mathrm{red}}_{X,\gamma}=f(nm,l)$ for $\gamma=(n,l,m)$;
this is one integer in each degree.  A source--target identification
fixes two integers in every degree (the parity dimensions), the
relation ideal, the radical equality, the no-extra-relations theorem,
the completed PBW filtration, and the Hopf-radical coideal stability.

The determinant admits the zero-bracket graded super vector space with
the same signed dimensions: the denominator is identical, while the
Chevalley diagonal, the non-orthogonal brackets, and the Hopf pairing
datum have disappeared.  Recognition therefore requires
representatives, relations, pairings, kernels, and PBW data, not
numerics.

\subsubsection*{Borcherds--Cartan target datum}

The target side imports the Borcherds--Kac superdatum
$\mathscr B_{\Delta_5}=(H,I,p,\alpha_i,(\,,\,))$ of
Gritsenko--Nikulin \cite{GN,GNPartI,GNII} after the
\cite{KAC:1,BorcherdsGKM88} presentation conventions.  The Cartan is
$H=\Lambda_{II}^{2,1}\otimes\C$, the three real simple roots
$\delta_1,\delta_2,\delta_3$ form a hyperbolic triangle with Gram
matrix
\[
((\delta_i,\delta_j))_{i,j}=
\begin{pmatrix}2&-2&-2\\-2&2&-2\\-2&-2&2\end{pmatrix},
\]
and the imaginary fibres of $I\to H^*$ carry the
Gritsenko--Nikulin parity split
$\mu_{\overline 0}(a),\mu_{\overline 1}(a)$ for imaginary simple
roots.  The degree map is
\[
\alpha(n,l,m)=2nf_2-lf_3+2mf_{-2}
=n\delta_1+m\delta_2+(n+m-l)\delta_3,
\]
with $\widehat{\mathfrak F}(\mathscr B_{\Delta_5})$ the completed
free Lie superalgebra on the parity-homogeneous generators
$e_i,f_i$, $i\in I$, and $J_{\mathrm{BK}}$ the closed
Borcherds--Kac ideal generated by Cartan, Chevalley, real Serre,
orthogonality, and super sign relations.  The target denominator
algebra is
\[
G(\mathscr B_{\Delta_5})
=\widehat{\mathfrak F}(\mathscr B_{\Delta_5})/
\overline{J_{\mathrm{BK}}+\operatorname{Rad}_{\mathrm{GN}}}
=\mathcal A_{\mathrm{den}}
=\mathfrak g_{\Delta_5}.
\]

\subsubsection*{Source-side data}

\begin{theorem}[Primitive recognition from finite source
representatives; \textit{conditional on (S1)--(S10)}]
\label{thm:primitive-recognition}
Let a compact Hall/factorization construction give a completed
$\widehat\Gamma_X$-graded primitive Hall Lie superalgebra
$\widetilde P_X$ and a completed Hopf pairing.  Its normal-ordered
Gram descent
\[
P_X^{\Pi,\mathrm{desc}}
=H^0\Prim_{\mathrm{prot}}(\overline\Pi_{X,*}^\Theta\mathcal F_X)
=\overline\Pi_{X,*}^\Theta\widetilde P_X,
\quad
P_X^{\Pi,\mathrm{red}}
=P_X^{\Pi,\mathrm{desc}}/\overline{\operatorname{Rad}\langle\,,\,\rangle_X}.
\]
Suppose the following ten source-side data are constructed.

\begin{enumerate}
\item[\textup{(S1)}] \emph{Charge and normal-ordered descent.}  The
chain-level descent $\overline\Pi_{X,*}^\Theta$ categorifies the
homomorphism $\overline\Pi_X(c,T)=\Pi_X(c)-T$, so that the descended
bracket is $\Gamma_{\mathrm{gram}}$-graded.  Each descended homogeneous
parity piece is finite dimensional.

\item[\textup{(S2)}] \emph{Cartan and root-degree matching.}  An even
Cartan identification $\iota_H:P^{\Pi,\mathrm{red}}_{X,0}
\xrightarrow{\sim}\Lambda_{II}^{2,1}\otimes\C$ aligning the descended
Gram grading with the GN root grading via $\alpha(n,l,m)$.

\item[\textup{(S3)}] \emph{Simple primitive representatives and
parities.}  For every simple-root index $i\in I$, choose
$\gamma_i\in\Gamma_{\mathrm{gram}}$ with
$\alpha(\gamma_i)=\alpha_i$.  The source supplies parity-homogeneous
representatives $e_i^X,f_i^X,h_i^X$ in
$P^{\Pi,\mathrm{red}}_{X,\gamma_i}$,
$P^{\Pi,\mathrm{red}}_{X,-\gamma_i}$, and
$P^{\Pi,\mathrm{red}}_{X,0}$, with $|e_i^X|=|f_i^X|=p(i)$,
$h_i^X=\iota_H^{-1}(\alpha_i)$, and the GN multiplicity fibres
$\mu_{\overline 0}(a),\mu_{\overline 1}(a)$ for imaginary simple
roots.  Real representatives for $\delta_1,\delta_2,\delta_3$ are even.

\item[\textup{(S4)}] \emph{Borcherds--Kac relations in the Hall
bracket.}  In the protected Hall/factorization bracket the
representatives satisfy:
\[
[h,h']=0,\quad [h,e_i^X]=(h,\alpha_i)e_i^X,\quad
[h,f_i^X]=-(h,\alpha_i)f_i^X,\quad
[e_i^X,f_j^X]=\delta_{ij}h_i^X,
\]
the real Serre relations
\[
(\operatorname{ad}e_i^X)^{1-2(\alpha_i,\alpha_j)/(\alpha_i,\alpha_i)}
e_j^X=0,
\quad
(\operatorname{ad}f_i^X)^{1-2(\alpha_i,\alpha_j)/(\alpha_i,\alpha_i)}
f_j^X=0\qquad(\text{real }i,\ i\ne j),
\]
the Borcherds orthogonality relations
$(\alpha_i,\alpha_j)=0\Rightarrow[e_i^X,e_j^X]=[f_i^X,f_j^X]=0$, and
the super sign rule.  In the rank-three real subsystem, the Serre
exponent is $3$ ($(\delta_i,\delta_j)=-2$); first brackets
$[e_i^X,e_j^X]$ are not forced to vanish.  For the primitive isotropic
degrees $a_{ij}=\delta_i+\delta_j$ ($i<j$), choose
$\gamma_{ij}\in\Gamma_{\mathrm{gram}}$ with
$\alpha(\gamma_{ij})=a_{ij}$.  The source supplies
the nine Gritsenko--Nikulin isotropic simple representatives
$u^{\pm,X}_{ij,r}\in P^{\Pi,\mathrm{red}}_{X,\pm\gamma_{ij}}$,
$1\le r\le 9$, satisfying orthogonality and Chevalley zero with
adjacent real roots, with real-string relations
$(\operatorname{ad}e_k^X)^5 u^{+,X}_{ij,r}=0$ for the complementary
$\{i,j,k\}=\{1,2,3\}$.  The tenth even vector in
$\mathfrak g_{a_{ij}}$ is the real bracket $[e_i^X,e_j^X]$.

\item[\textup{(S5)}] \emph{Counital Hall bialgebra, Hopf pairing, and
radical quotient.}  A continuous completed Hall product, unit,
coproduct, and counit satisfy associativity, coassociativity, the unit
and counit identities, and the bialgebra compatibility
\[
\Delta_X m_X
=(m_X\otimes m_X)(1\otimes\tau\otimes1)(\Delta_X\otimes\Delta_X)
\]
with the Koszul sign braiding \(\tau\).  The bracket is the
supercommutator of this product restricted to primitives.  The Hopf
pairing is continuous, homogeneous, invariant for bracket and coproduct,
has vanishing Frobenius cyclic defect, is non-degenerate after the
radical quotient, and has closed homogeneous radical
$\operatorname{Rad}_X=\{x\mid\langle x,y\rangle_X=0\ \forall y\}$ a Lie
ideal and coideal.  After quotienting by
$\overline{\operatorname{Rad}_X}$, the source radical quotient is
identified with the GN/Kac quotient
$\overline{\operatorname{Rad}_{\mathrm{GN}}}$.  At every finite stage
the Lie-ideal and coideal stability hold and the inverse-limit kernel
is exactly $\operatorname{Rad}_X$.

\item[\textup{(S6)}] \emph{No extra relations.}  If
$\pi:\widehat{\mathfrak F}(\mathscr B_{\Delta_5})\to
P_X^{\Pi,\mathrm{red}}$ is the continuous map determined by the
representatives,
$\ker\pi=\overline{J_{\mathrm{BK}}+\operatorname{Rad}_{\mathrm{GN}}}$.
Equivalently, $G(\mathscr B_{\Delta_5})\to P_X^{\Pi,\mathrm{red}}$ is
injective.  This is checked finite stage by finite stage on saturated
windows, with kernel identifications commuting with HN transitions.

\item[\textup{(S7)}] \emph{Generation.}  $P_X^{\Pi,\mathrm{red}}$ is
topologically generated by $P^{\Pi,\mathrm{red}}_{X,0}$ and the simple
primitive representatives.

\item[\textup{(S8)}] \emph{Completed PBW.}  PBW filtration
compatibility:
\[
\operatorname{gr}_{\mathrm{PBW}}U(P_X^{\Pi,\mathrm{red}})_W
\xrightarrow{\sim}
\operatorname{gr}_{\mathrm{PBW}}U(\mathfrak g_{\Delta_5})_W
\]
on every saturated finite window $W$, with strict transition maps
and inverse limit in the HN topology of
\S\ref{subsec:hall-correspondences}.

\item[\textup{(S9)}] \emph{Exact completion.}  Inverse limit exact on
the saturated finite root windows: transition maps strict, images
closed, and passage to the completed radical quotient commuting with
finite-window kernels, cokernels, PBW associated gradeds, and parity
pieces.  Equivalently, no new relation, generator, radical element,
or parity-cancelling summand appears only after completion.

\item[\textup{(S10)}] \emph{Full parity dimensions.}  For every
nonzero $\gamma$ with $\alpha(\gamma)\in\mathcal R$,
\[
\dim P^{\Pi,\mathrm{red}}_{X,\gamma,\overline p}
=\rootmult_{\overline p}(\alpha(\gamma)),\qquad p=0,1.
\]
The same equality is required in degree \(-\gamma\), with the target
parity dimensions transported by the Borcherds--Kac Chevalley
anti-involution; it is not inferred from the positive half or from the
signed denominator.
On the first timelike channel $\gamma=(1,1,1)$,
\[
d^{\mathrm{GN}}_{(1,1,1),\overline 0}=29,\qquad
d^{\mathrm{GN}}_{(1,1,1),\overline 1}=93.
\]
These are presentation-level target dimensions, not an inference
from $29-93=-64$.
\end{enumerate}
Then $P_X^{\Pi,\mathrm{red}}\cong\mathfrak g_{\Delta_5}$ as completed
Lie superalgebras graded by \(\Gamma_{\mathrm{gram}}\) through
\(\alpha\), with denominator
$64^{-1}\Delta_5(2Z)$.  If in addition the chiral Koszul source datum
and the finite Koszul comparison datum of
\S\ref{subsec:koszul-source} are constructed, including the strict
Koszul Mittag--Leffler exactness for the compatibility complexes, then
\(\Omega_E^{\mathrm{ch}}C_X\simeq\mathsf{Den}_{\Delta_5,E}\) through a
map respecting the Weyl action, Pfaffian orientation character, Hall
pairing, radical quotient, and PBW associated graded.
\end{theorem}

\begin{proof}
The simple primitive representatives define a continuous homomorphism
$\pi:\widehat{\mathfrak F}(\mathscr B_{\Delta_5})\to
P_X^{\Pi,\mathrm{red}}$.  (S1) and (S2) make $\pi$ a homomorphism of
the descended Gram-graded completions.  (S3) fixes the real and
imaginary simple generators with their parities and multiplicity
fibres.  (S4) puts $J_{\mathrm{BK}}$ in the kernel.  (S5) identifies
the compact Hopf radical with the GN/Kac radical, so $\pi$ descends
to $G(\mathscr B_{\Delta_5})\to P_X^{\Pi,\mathrm{red}}$.  (S7) gives
surjectivity; (S6) gives injectivity after the completed radical
quotient.  (S8) gives PBW-compatible finite-height identifications;
(S9), in the form of
Proposition~\ref{prop:strict-ml-completion-primitive-recognition},
gives exactness in the inverse limit; (S10) rules out hidden
zero-superdimension root spaces and cancelling ideals invisible to the
determinant.  The chiral statement uses $C_X$, $b_X$,
$\Theta_{\mathrm{Kos}}$, the finite Koszul comparison datum
\(\mathfrak K^{\mathrm{cmp}}_{X,R}\), and the Koszul
Mittag--Leffler exactness of \S\ref{subsec:koszul-source}.
\end{proof}

\subsubsection*{Cofinal finite-window primitive recognition}

The Borcherds product fixes signed target superdimensions; the
Gritsenko--Nikulin/Kac presentation fixes full target parity
dimensions only in the degrees where the presentation has been
computed.  A compact $K3\times E$ source is identified with the
denominator algebra only by source representatives, pairings,
radicals, relation ideals, PBW, and exact transition maps.

The imported target on the first window $W_{\le 3}$ has parity table
\[
\delta_i:1|0,\quad
a_{ij}:10|0,\quad
2\delta_i+\delta_j:1|0,\quad
\delta_{123}:29|93,
\]
so $29-93=-64=f(1,1)$.  The row $a_{ij}:10|0$ is the real bracket
$[e_i,e_j]$ together with the nine Gritsenko--Nikulin isotropic simple
directions.  These are target data; they do not construct
compact source primitives, source pairings, radical kernels, or
source bracket constants.
The finite target packet \(\texttt{wle3\_target\_parity}\) verifies
this table from the theta-series coefficients and the
Gritsenko--Nikulin/Kac presentation: \(a_{ij}\) is isotropic,
the real-string exponent for \(2\delta_i+\delta_j\) is \(3\), and
\(\delta_{123}\) splits as \(29=2+27\) even directions and \(93\) odd
imaginary-simple directions.  It is a target-reference table only.

On the first base terminal and saturation-extension target rows the
coefficient table
\(\texttt{k3e\_first\_relation\_window\_coefficients}\) computes
\[
f(4,4)=10,\quad f(2,2)=108,\quad
f(2,1)=f(4,3)=-513,\quad
f(3,3)=-64,\quad f(4,2)=4016,
\]
together with their discriminant-orbit translates and the terminal
zero rows \(f(5,5)=f(1,-3)=0\).  These equalities match the signed
dimensions in the A071 target presentation.  They remain one-integer
target arithmetic; they do not supply the two parity dimensions required
by (R3), or any compact-source row in (R0)--(R8).
The A071 target-presentation check verifies the target rows
\[
2a_{ij}:10|0,\qquad C_{k,3}:29|93,\qquad
C_{k,4}:10|0,\qquad C_{k,5}:0|0,
\]
and verifies that \(C_{k,2}\), \(D_k=\delta_k+2a_{ij}\), and
\(2\delta_{123}\) are signed-only: they have no target basis,
positive-negative pairing block, radical block, or PBW row.  This is a
target-reference computation, not a compact-source theorem.

The first relation-closed compact-source window is therefore a separate
finite computation.  Its row packet is
\[
\mathrm{window\_closure},\quad
\mathrm{target\_source\_representatives},\quad
\mathrm{chevalley\_serre\_matrices},\quad
\mathrm{hall\_boundary\_complex},\quad
\mathrm{spectral\_sequence},\quad
\mathrm{kernel\_matrices},\quad
\mathrm{pbw\_graded},\quad
\mathrm{first\_window\_theorems},\quad
\mathrm{transitions},\quad
\mathrm{scalar\_firewall}.
\]
These rows must prove downward saturation and relation closure,
geometric source representatives with full parity ranks, Cartan,
Chevalley, real-Serre, Borcherds-orthogonality, and super-sign
relations, the Hall--Chevalley boundary \(d_1\), strong convergence of
the finite-window spectral sequence, computed kernel bases, no-extra
kernel equality, PBW associated-graded equality, and strict
Mittag--Leffler transition.  Target labels, signed dimensions, scalar
traces, Pfaffian products, denominator products, target PBW rows,
arbitrary matrices, and status-only rows are excluded by the scalar
firewall.

\begin{proposition}[Target-presentation rows are a prerequisite;
\textit{proved}]
\label{prop:target-presentation-prerequisite}
Let \(W\subset\mathcal R_+\) be a finite downward-saturated target-root
window.  If \(W\) is used in a primitive-recognition datum satisfying
(R0)--(R8), then the target restriction
\(\mathfrak g_{\Delta_5}|_{S_W}\) must be represented by a finite
target-presentation datum in the sense of
Definition~\ref{def:bkm-finite-target-presentation-datum}.  Signed
rows \(\smult(\beta)\) and imaginary-simple integers \(m(\beta)\)
alone are not admissible codomains for the comparison maps
\(A_{\beta,\overline p}\).
\end{proposition}

\begin{proof}
Clause (R2) compares source relation rows with the Borcherds--Kac
target relations, so the Cartan, Chevalley, real-Serre,
orthogonality, and super-sign rows must be present in the target
window.  Clause (R3) compares two parity dimensions in each sign, not
the signed difference.  Clause (R4) compares the source
positive-negative Hopf pairing and radical quotient with the target
invariant-pairing kernel.  Clause (R7) compares PBW associated gradeds.
Thus the target side must carry the parity, relation, pairing,
radical, and PBW rows of
Definition~\ref{def:bkm-finite-target-presentation-datum}.  A scalar
row \(\smult(\beta)\) or \(m(\beta)\) supplies none of the pairing,
radical, relation, or PBW codomains, and supplies only one integer
where (R3) requires two.
\end{proof}

\begin{definition}[Cofinal finite-window primitive-recognition datum]
\label{def:cofinal-primitive-recognition-datum}
A finite set \(W\subset\mathcal R_+\) is \emph{downward saturated} if
\(\beta\in W\), \(\beta'\in\mathcal R_+\), \(0<\beta'\le\beta\), and
\(\beta-\beta'\in Q_+\) imply \(\beta'\in W\).  The canonical
height-window convention is
\[
W_\nu=\{\beta=b_1\delta_1+b_2\delta_2+b_3\delta_3\in\mathcal R_+
\mid b_1+b_2+b_3\le \nu\}.
\]
A cofinal finite-window primitive-recognition datum for
$P_X^{\Pi,\mathrm{red}}$ is an increasing sequence of finite
downward-saturated target-root windows
$W_1\subset W_2\subset\cdots$ with
$\bigcup_\nu W_\nu=\mathcal R_+$, together with the following finite
data on every $W_\nu$, compatible under $W_\nu\subset W_{\nu+1}$.  Put
\[
\Gamma(W_\nu):=\{\gamma\in\Gamma_{\mathrm{gram}}\mid
\alpha(\gamma)\in W_\nu\}.
\]

\begin{enumerate}
\item[(R0)] \emph{Compact source labels.}  Every source basis vector
of degree in $\Gamma(W_\nu)\cup(-\Gamma(W_\nu))\cup\{0\}$ satisfies
the compact source admissibility condition.  Target labels from
$G(\mathscr B_{\Delta_5})$ appear only in target rows or as
codomains of comparison maps.
\item[(R1)] \emph{Representatives and parities.}  Simple-root
representatives $e_i^X,f_i^X,h_i^X$ whose target roots lie in $W_\nu$
are supplied in source degrees $\pm\gamma_i$ with
$\alpha(\gamma_i)=\alpha_i$, with the GN parity fibres for imaginary
simple roots and even real representatives for $\delta_1,\delta_2,\delta_3$.
\item[(R2)] \emph{Relations.}  Cartan, Chevalley, real Serre,
orthogonality, and super sign relations checked for every relation
whose displayed target roots lie in
$W_\nu\cup(-W_\nu)\cup\{0\}$ and whose source lifts lie in
$\Gamma(W_\nu)\cup(-\Gamma(W_\nu))\cup\{0\}$.
\item[(R3)] \emph{Full finite parity dimensions.}  For every
$\gamma\in\Gamma(W_\nu)$ and for its negative \(-\gamma\),
$\dim P^{\Pi,\mathrm{red}}_{X,\gamma,\overline p}
=\rootmult_{\overline p}^{\mathrm{GN}}(\alpha(\gamma))$ ($p=0,1$),
with the negative-root parity dimensions carried by the
Borcherds--Kac Chevalley anti-involution, and with zero on source
degrees whose image is not a target root.
Two-integer equality, not signed-superdimension.
\item[(R4)] \emph{Hall bialgebra, Hopf pairing, and radical.}  The
finite product \(M\), coproduct \(D\), unit \(\eta\), and counit
\(\epsilon\) satisfy associativity, coassociativity, unit, counit, and
bialgebra compatibility identities on the window; the supercommutator
bracket sends primitives to primitives.  Positive-negative pairing
matrices are written in parity blocks; kernels agree with GN/Kac
invariant-pairing kernels; Hopf adjointness, the primitive Frobenius
cyclic identity, and non-degeneracy on the radical quotient are verified
by finite rank-zero defect rows; stable under finite brackets in the
window; coideal for finite coproduct; carried to the next window by
transitions; inverse-limit radical the closure of finite kernels.
\item[(R5)] \emph{No extra relations.}  Finite kernel equality
$\ker\pi_\nu=
(\overline{J_{\mathrm{BK}}+\operatorname{Rad}_{\mathrm{GN}}})_{\le W_\nu}$.
\item[(R6)] \emph{Generation.}  Source primitive spaces over
$\Gamma(W_\nu)\cup(-\Gamma(W_\nu))$ are generated by Cartan and
supplied simple representatives through brackets whose intermediate
target roots stay in $W_\nu$.
\item[(R7)] \emph{Completed PBW compatibility.}  Strict PBW
filtration preservation under transitions, inverse-limit PBW
compatibility.
\item[(R8)] \emph{Exact triangular completion.}  For any source stage
\(R\) assigned to \(W_\nu\), words whose total Gram degree lies in the
finite target test window \(\Gamma_R^{\mathrm{test}}\) of
Definition~\ref{def:finite-target-test-window} only through
positive-negative cancellation are not added as extra finite generators;
transition maps strict, images closed, \(\lim^1\) vanishing on the
triangular finite-window systems.
\end{enumerate}
The datum is finite at each $\nu$, but not numerical: it contains
representatives, relation verifications, pairing matrices, kernel
equalities, and PBW transition data.  The Borcherds product and the
OP scalar square provide none of these source-side objects.  The further
assertion that every relation-closed target degree occurs at some finite
HN stage is the finite-\(R\) exhaustion datum of
Definition~\ref{def:relation-closed-target-finite-r-exhaustion}; it is
not a consequence of target relation closure alone.
\end{definition}

\begin{proposition}[Height root windows and their Gram preimages;
\textit{proved}]
\label{prop:height-root-windows-gamma}
The sets \(W_\nu\) of
Definition~\ref{def:cofinal-primitive-recognition-datum} are finite,
downward saturated, and cofinal in \(\mathcal R_+\).  If
\(\beta=b_1\delta_1+b_2\delta_2+b_3\delta_3\in W_\nu\), then
\[
\gamma_\beta=(b_1,\ b_1+b_2-b_3,\ b_2)\in\Gamma_{\mathrm{gram}}
\]
is the unique element with
\[
\alpha(\gamma_\beta)
=b_1\delta_1+b_2\delta_2+b_3\delta_3.
\]
Hence
\[
\Gamma(W_\nu)
=\{(b_1,b_1+b_2-b_3,b_2)\mid
b_1\delta_1+b_2\delta_2+b_3\delta_3\in W_\nu\}.
\]
Moreover \(\gamma_\beta\) lies in the type-II product chamber:
\[
b_2>0\Rightarrow m>0,\qquad
b_2=0,\ b_1>0\Rightarrow m=0,\ n>0,\qquad
b_1=b_2=0,\ b_3>0\Rightarrow m=n=0,\ l<0.
\]
\end{proposition}

\begin{proof}
Finiteness follows because \(b_1+b_2+b_3\le\nu\) gives only finitely
many triples in \(Q_+\).  If
\(\beta'\in\mathcal R_+\), \(0<\beta'\le\beta\), and
\(\beta\in W_\nu\), then
\(\operatorname{ht}(\beta')\le\operatorname{ht}(\beta)\le\nu\); hence
\(\beta'\in W_\nu\).  Cofinality follows because every positive root
has finite height.

For \(\gamma=(n,l,m)\), the product-chamber root map is
\[
\alpha(n,l,m)=n\delta_1+m\delta_2+(n+m-l)\delta_3.
\]
Solving \(\alpha(n,l,m)=b_1\delta_1+b_2\delta_2+b_3\delta_3\) gives
\(n=b_1\), \(m=b_2\), and \(l=b_1+b_2-b_3\), proving uniqueness and the
displayed formula for \(\Gamma(W_\nu)\).  The three chamber cases are
then immediate from \(b_i\ge0\) and \(\beta\ne0\).
\end{proof}

The packet
\(\texttt{certificates/lattice/downward\_root\_windows}\) records
Proposition~\ref{prop:height-root-windows-gamma}.  It verifies
\texttt{DOWNWARD\_ROOT\_WINDOWS\_VERIFIED}: the height-window formula,
the inverse formula for \(\Gamma(W_\nu)\), the product-chamber sector
rule, and the ambient positive-cone audit for \(1\le\nu\le7\).  It is a
target-side root-window convention, not a compact-source recognition
packet, not a Pfaffian orientation, and not a protected trace.

\begin{proposition}[Active height windows exhaust the positive roots;
\textit{proved}]
\label{prop:active-root-window-exhaustion}
Let \(W_\nu^{\mathrm{act}}:=W_\nu\).  Then
\[
\bigcup_{\nu\ge1}W_\nu^{\mathrm{act}}=\mathcal R_+ .
\]
Equivalently, every positive root
\(\beta=b_1\delta_1+b_2\delta_2+b_3\delta_3\) is witnessed in the
finite active window \(W_{\operatorname{ht}(\beta)}\), where
\(\operatorname{ht}(\beta)=b_1+b_2+b_3\).
\end{proposition}

\begin{proof}
By definition \(W_\nu\subset\mathcal R_+\) for every \(\nu\).  Conversely,
if \(\beta\in\mathcal R_+\), then its coefficients in the positive
root monoid are non-negative and have finite sum
\(\operatorname{ht}(\beta)\).  Hence
\(\beta\in W_{\operatorname{ht}(\beta)}^{\mathrm{act}}\).  The two
inclusions give the displayed equality.
\end{proof}

The packet
\(\texttt{certificates/lattice/active\_root\_window\_exhaustion}\)
records Proposition~\ref{prop:active-root-window-exhaustion}.  It
verifies \texttt{ACTIVE\_ROOT\_WINDOW\_EXHAUSTION\_VERIFIED}: the
height-witness formula
\(\beta\in\mathcal R_+\Rightarrow\beta\in W_{\operatorname{ht}(\beta)}\),
the low active membership samples, the inverse \(\gamma_\beta\) formula,
and the type-II chamber sector rule.  It is not the compact-source
support exhaustion required before moduli transitions, orientation
transport, vanishing-cycle transport, Hall operations, radical
transport, or PBW transport.

\begin{definition}[Source-window compact-support exhaustion datum]
\label{def:source-window-compact-support-exhaustion}
A source-window compact-support exhaustion datum consists of:
\begin{enumerate}
\item a set \(\mathscr S_X^{\mathrm{cpt}}\) of compact-source support
rows, each carrying a source degree, a compact geometric provenance row,
proper support, finite stabilizer data, and the orientation row used by
the source Hall construction;
\item a height function
\[
h_X:\mathscr S_X^{\mathrm{cpt}}\longrightarrow \mathbb Z_{\ge0}
\]
with finite fibres;
\item finite source windows
\[
\mathscr S_{X,\le R}^{\mathrm{cpt}}
=\{s\in\mathscr S_X^{\mathrm{cpt}}\mid h_X(s)\le R\};
\]
\item a cofinality witness proving that every
\(s\in\mathscr S_X^{\mathrm{cpt}}\) lies in
\(\mathscr S_{X,\le h_X(s)}^{\mathrm{cpt}}\);
\item transition rows
\(\mathscr S_{X,\le R}^{\mathrm{cpt}}\hookrightarrow
\mathscr S_{X,\le R'}^{\mathrm{cpt}}\) for \(R\le R'\), later refined
to closed embeddings or proper maps before orientation, vanishing-cycle,
Hall, radical, and PBW transport is asserted.
\end{enumerate}
Target root labels, signed multiplicities, Pfaffian products, and
protected traces are not entries of this datum.
\end{definition}

\begin{proposition}[Source-window compact-support exhaustion criterion;
\textit{proved}]
\label{prop:source-window-compact-support-criterion}
If a datum of
Definition~\ref{def:source-window-compact-support-exhaustion} is
supplied, then the finite source windows exhaust compact support:
\[
\bigcup_{R\ge0}\mathscr S_{X,\le R}^{\mathrm{cpt}}
=\mathscr S_X^{\mathrm{cpt}}.
\]
\end{proposition}

\begin{proof}
The inclusion from left to right is part of the definition of the
finite source windows.  Conversely, for
\(s\in\mathscr S_X^{\mathrm{cpt}}\), the cofinality witness gives
\(s\in\mathscr S_{X,\le h_X(s)}^{\mathrm{cpt}}\).  Thus every compact
source support row occurs in one finite source window.
\end{proof}

The packet
\(\texttt{certificates/sources/source\_window\_compact\_support}\)
records the present state of this datum.  Its verifier status is
\texttt{SOURCE\_WINDOW\_COMPACT\_SUPPORT\_OBSTRUCTION\_VERIFIED}: the
obstruction ledger lists the missing compact support index set, source
height function, finite source windows, compact provenance rows,
cofinality witness, source-window transitions, and target-source
compatibility rows.  Its core source tables are empty, so
the compact-source support-exhaustion theorem remains open.  This is
the firewall separating the target-root exhaustion of
Proposition~\ref{prop:active-root-window-exhaustion} from compact-source
support exhaustion.

\begin{definition}[Target/source window compatibility datum]
\label{def:target-source-window-compatibility}
Let \(W_\nu\) be the target height window and
\(\mathscr S_{X,\le R}^{\mathrm{cpt}}\) a finite compact-source
window.  A target/source window compatibility datum consists of:
\begin{enumerate}
\item a degree map
\[
d_R:\mathscr S_{X,\le R}^{\mathrm{cpt}}\longrightarrow
\Gamma_{\mathrm{gram}}
\]
recording the normal-ordered Gram degree of every compact support row;
\item a target index \(\nu(R)\) such that
\[
d_R(s)\in\Gamma(W_{\nu(R)})\qquad
\text{for every }s\in\mathscr S_{X,\le R}^{\mathrm{cpt}};
\]
\item the row-wise equality
\[
\alpha(d_R(s))\in W_{\nu(R)}
\]
with zero degree-map defect;
\item compatibility with compact provenance: the support row, its
source degree, its compact geometric provenance, and its Gram degree
all carry the same support identifier;
\item compatibility under transitions
\[
\mathscr S_{X,\le R}^{\mathrm{cpt}}\hookrightarrow
\mathscr S_{X,\le R'}^{\mathrm{cpt}},\qquad
W_{\nu(R)}\subset W_{\nu(R')},
\]
so that the source degree map commutes with passage to larger windows.
\end{enumerate}
\end{definition}

\begin{proposition}[Compatibility criterion for target and source windows;
\textit{proved}]
\label{prop:target-source-window-compatibility-criterion}
Given a datum of
Definition~\ref{def:target-source-window-compatibility}, every compact
source support row in \(\mathscr S_{X,\le R}^{\mathrm{cpt}}\) has a
well-defined target codomain in
\(\Gamma(W_{\nu(R)})\cup(-\Gamma(W_{\nu(R)}))\cup\{0\}\), and the
finite comparison maps on that source window are degree-compatible with
the target window \(W_{\nu(R)}\).
\end{proposition}

\begin{proof}
For a positive source degree the codomain is \(\alpha(d_R(s))\in
W_{\nu(R)}\) by the compatibility datum.  The negative degree is carried
to \(-\Gamma(W_{\nu(R)})\), and the Cartan degree is \(0\).  The
zero-defect equality makes the target degree independent of notation,
and the transition condition shows that enlarging \(R\) does not move a
previously defined source support row outside the enlarged target
window.  Thus every finite comparison row has a target codomain in the
chosen window.
\end{proof}

The packet
\(\texttt{certificates/sources/target\_source\_window\_compatibility}\)
records the current state of this datum.  Its verifier status is
\texttt{TARGET\_SOURCE\_WINDOW\_COMPATIBILITY\_OBSTRUCTION\_VERIFIED}:
the source degree maps, \(\Gamma(W_\nu)\)-membership rows, target-window
matching rows, compact-provenance compatibility rows, cofinal-index
compatibility rows, transition compatibility rows, and zero
compatibility-defect rows are not yet supplied.  Hence target/source
window compatibility remains open.

\begin{proposition}[Finite-window presentation recognition;
\textit{proved}]
\label{prop:finite-window-presentation-recognition}
Fix a downward-saturated finite target-root window \(W\) and put
\(S_W=W\cup(-W)\cup\{0\}\).  For a graded Lie superalgebra \(L\), write
\[
L|_{S_W}:=\bigoplus_{\beta\in S_W}L_\beta
\]
with the bracket retained only on pairs whose total degree remains in
\(S_W\).  Suppose a compact \(K3\times E\) source supplies
\((\mathrm{R0})\)--\((\mathrm{R8})\) on \(W\).  Then the finite
comparison matrices \(A_{\beta,\bar p}\) determine an isomorphism of
truncated Gram-graded Lie superalgebras
\[
A_W:\ P^{\Pi,\mathrm{red}}_X|_{\alpha^{-1}(S_W)}
\xrightarrow{\ \sim\ }
\mathfrak g_{\Delta_5}|_{S_W},
\]
compatible with the Cartan identification, the simple-root
representatives, the Hopf-pairing radical quotient, and the
associated-graded PBW filtration on \(W\).  Conversely, such an
isomorphism, together with compact-source basis provenance, parity
blocks, simple representatives, product, coproduct, unit, counit,
Hopf-pairing, radical, quotient, and transition maps from which its
matrices are obtained, gives \((\mathrm{R0})\)--\((\mathrm{R8})\).
\end{proposition}

\begin{proof}
The representatives of (R1) define a homomorphism from the free
Borcherds--Kac presentation to the source truncation.  The relation
rows of (R2) put every Cartan, Chevalley, real-Serre, orthogonality,
and super-sign relation whose degree lies in \(S_W\) in the kernel.
The Hopf bialgebra and pairing data of (R4) make the source quotient a
truncated Lie superalgebra: product and coproduct are compatible, the
supercommutator closes on primitives, and the radical is a Lie ideal
and a coproduct coideal.  The no-extra relation row (R5) is exactly the
finite equality
\[
\ker\pi_W=
\bigl(\overline{J_{\mathrm{BK}}+\operatorname{Rad}_{\mathrm{GN}}}
\bigr)|_{S_W}.
\]
Thus the presentation map descends injectively from
\(\mathfrak g_{\Delta_5}|_{S_W}\) to the source radical quotient.
Generation (R6) gives surjectivity on every source degree in
\(\alpha^{-1}(S_W)\).  The parity dimensions of (R3) identify the
even and odd dimensions degree by degree in both signs, so no
zero-superdimension summand or parity-cancelling kernel can remain.
PBW compatibility (R7) identifies the associated gradeds of the
universal enveloping algebras on \(W\).  Exact triangular completion
(R8) says that the finite truncation is stable under passage to the
next window and that no relation is created only by completion.  The
matrix \(A_W\) is therefore a degree-preserving bijection preserving
all brackets, pairing quotients, and PBW associated gradeds present in
the finite window.

Conversely, a truncated isomorphism with specified compact-source basis
provenance, parity blocks, simple representatives, product, coproduct,
unit, counit, Hopf pairing, radical, quotient, and transition maps
yields the relation rows, radical identities, no-extra equality,
generation spans, PBW rows, and transition rows by writing those finite
maps in homogeneous bases.  Without the named compact-source maps one
has only a linear isomorphism of graded vector spaces, not (R0)--(R8).
\end{proof}

\begin{definition}[PBW-transition datum]
\label{def:transition-pbw-filtration-datum}
Let \(W_\nu\subset W_{\nu+1}\) be two downward-saturated finite
windows.  Put
\[
L_\nu=P^{\Pi,\mathrm{red}}_{X,\nu}/\operatorname{Rad}^{\Pi}_{X,\nu},
\qquad
L_{\nu+1}=P^{\Pi,\mathrm{red}}_{X,\nu+1}/
\operatorname{Rad}^{\Pi}_{X,\nu+1}.
\]
A PBW-transition datum for \(L_{\nu+1}\to L_\nu\) consists of the
following finite rows:
\begin{enumerate}
\item radical-quotient transition matrices \(T^Q_{\nu+1,\nu}\) in
homogeneous quotient bases;
\item ordered PBW word bases for \(F^{\mathrm{PBW}}_{\le d}U(L_\nu)\)
and \(F^{\mathrm{PBW}}_{\le d}U(L_{\nu+1})\), with their
associated-graded ranks;
\item source and target relation-rewrite rows which identify the
normal forms in the enveloping algebras;
\item the multiplicative word transition induced by
\(T^Q_{\nu+1,\nu}\), together with the zero-defect relation-rewrite
identity saying that it descends to
\[
U(T^Q_{\nu+1,\nu}):U(L_{\nu+1})\longrightarrow U(L_\nu);
\]
\item for every retained \(d\), a zero-defect filtered-image row
\[
U(T^Q_{\nu+1,\nu})
\bigl(F^{\mathrm{PBW}}_{\le d}U(L_{\nu+1})\bigr)
\subset
F^{\mathrm{PBW}}_{\le d}U(L_\nu);
\]
\item the induced matrices on
\(\operatorname{gr}^{\mathrm{PBW}}_dU(L_{\nu+1})\to
\operatorname{gr}^{\mathrm{PBW}}_dU(L_\nu)\), the PBW-symbol
intertwining identity, the \(A_\beta\)-PBW transition square, and the
composition law for three consecutive windows.
\end{enumerate}
The datum is absent if only primitive transitions, radical transitions,
PBW rank equalities, target PBW rows, Hilbert series, scalar traces,
or the formal normal-ordered PBW-pushforward packet are supplied.
\end{definition}

\begin{proposition}[Transition preservation of PBW filtrations;
\textit{proved}]
\label{prop:transition-pbw-filtration-preservation-criterion}
A PBW-transition datum in the sense of
Definition~\ref{def:transition-pbw-filtration-datum} determines a
filtered algebra homomorphism
\[
U(T^Q_{\nu+1,\nu}):
\bigl(U(L_{\nu+1}),F^{\mathrm{PBW}}_{\le\bullet}\bigr)
\longrightarrow
\bigl(U(L_\nu),F^{\mathrm{PBW}}_{\le\bullet}\bigr)
\]
and hence maps
\[
\operatorname{gr}^{\mathrm{PBW}}_d U(L_{\nu+1})
\longrightarrow
\operatorname{gr}^{\mathrm{PBW}}_d U(L_\nu)
\]
for every retained PBW degree \(d\).  These associated-graded maps are
compatible with the finite comparison matrices \(A_\beta\), and they
compose functorially on triples of windows.  Without the rows of
Definition~\ref{def:transition-pbw-filtration-datum}, transition
preservation of PBW filtrations is not established.
\end{proposition}

\begin{proof}
The quotient transition matrix sends degree-one primitive generators of
\(L_{\nu+1}\) to degree-one elements of \(L_\nu\).  Multiplication gives
a word-level map from the tensor algebra on \(L_{\nu+1}\) to the tensor
algebra on \(L_\nu\).  The relation-rewrite zero-defect row is exactly
the statement that the image of every source enveloping relation is a
target enveloping relation; hence the word-level map descends to the
universal enveloping quotients.

The filtered-image row says, in the chosen finite PBW bases, that the
matrix of the descended map has no component from words of length
\(\le d\) to target normal forms of length \(>d\).  This is precisely
\(U(T^Q_{\nu+1,\nu})(F^{\mathrm{PBW}}_{\le d})\subset
F^{\mathrm{PBW}}_{\le d}\).  Therefore the standard quotient
\(F_{\le d}/F_{\le d-1}\) produces the displayed associated-graded
matrix.  The PBW-symbol row identifies this quotient map with the
symbol of the degree-one transition, and the \(A_\beta\)-square is the
same matrix identity after applying the finite source-to-target
comparison.  The composition row gives equality of the two matrix
products for three windows, so the associated graded maps form an
inverse system.  No scalar or rank-only datum contains these matrix
identities.
\end{proof}

The packet
\[
\texttt{certificates/hall/transition\_pbw\_filtration\_preservation}
\]
records the current state of this datum.  Its verifier returns
\(\texttt{TRANSITION\_PBW\_FILTRATION\_PRESERVATION\_OBSTRUCTION\_VERIFIED}\):
the radical-quotient transition input, source and target PBW
filtration rows, associated-graded rank rows, ordered word bases,
relation-rewrite systems, quotient transition matrices, multiplicative
word transitions, relation-rewrite compatibility rows, filtered-image
identities, restricted filtered transition matrices, associated-graded
transition matrices, PBW-symbol identities, \(A_\beta\)-PBW transition
squares, and PBW-transition composition rows are not yet supplied.

\begin{definition}[Parity-transition datum]
\label{def:transition-parity-decomposition-datum}
Let \(V_{\nu,\gamma}\) and \(V_{\nu+1,\gamma}\) be the retained
finite primitive source spaces, after the finite radical quotient if
recognition is being tested after \textup{(R4)}.  A parity-transition
datum for
\[
T_{\nu+1,\nu,\gamma}:V_{\nu+1,\gamma}\longrightarrow V_{\nu,\gamma}
\]
consists of parity-homogeneous bases
\[
V_{\nu,\gamma}=V_{\nu,\gamma,\overline 0}\oplus
V_{\nu,\gamma,\overline 1},
\qquad
V_{\nu+1,\gamma}=V_{\nu+1,\gamma,\overline 0}\oplus
V_{\nu+1,\gamma,\overline 1},
\]
the finite fermion-parity involutions
\[
J_{\nu,\gamma}=+1\ \text{on }V_{\nu,\gamma,\overline 0},\qquad
J_{\nu,\gamma}=-1\ \text{on }V_{\nu,\gamma,\overline 1},
\]
and the transition matrix of \(T_{\nu+1,\nu,\gamma}\), with the
following zero-defect rows:
\[
J_{\nu,\gamma}T_{\nu+1,\nu,\gamma}
=T_{\nu+1,\nu,\gamma}J_{\nu+1,\gamma},
\]
equivalently
\[
T_{\overline 0\overline 1}=0,\qquad
T_{\overline 1\overline 0}=0.
\]
The datum also includes the restricted matrices
\[
T_{\overline p}:V_{\nu+1,\gamma,\overline p}
\longrightarrow V_{\nu,\gamma,\overline p},
\qquad p=0,1,
\]
the same rows for \(-\gamma\) transported by the Borcherds--Kac
Chevalley anti-involution, the comparison square with the finite
\(A_\beta\)-matrices on each parity piece, and the composition law for
three consecutive windows.  Signed superdimension, target parity
tables, parity ranks without transition matrices, primitive/radical/PBW
transition rows, and the formal parity-pushforward packet are not such
a datum.
\end{definition}

\begin{proposition}[Transition preservation of parity decompositions;
\textit{proved}]
\label{prop:transition-parity-decomposition-preservation-criterion}
A parity-transition datum in the sense of
Definition~\ref{def:transition-parity-decomposition-datum} makes each
finite transition \(T_{\nu+1,\nu,\gamma}\) an even linear map.  Hence
\[
T_{\nu+1,\nu,\gamma}
\bigl(V_{\nu+1,\gamma,\overline p}\bigr)
\subset
V_{\nu,\gamma,\overline p},\qquad p=0,1,
\]
and the even and odd retained root-space towers are defined
separately.  The same conclusion holds for the negative-root spaces
and is compatible with the finite comparison maps \(A_\beta\).  Without
the rows of
Definition~\ref{def:transition-parity-decomposition-datum}, transition
preservation of parity decompositions is not established.
\end{proposition}

\begin{proof}
The commutator identity
\(J_{\nu,\gamma}T_{\nu+1,\nu,\gamma}
=T_{\nu+1,\nu,\gamma}J_{\nu+1,\gamma}\)
is the coordinate-free statement that the transition is even.  If
\(x\in V_{\nu+1,\gamma,\overline p}\), then
\[
J_{\nu,\gamma}T_{\nu+1,\nu,\gamma}x
=T_{\nu+1,\nu,\gamma}J_{\nu+1,\gamma}x
=(-1)^pT_{\nu+1,\nu,\gamma}x.
\]
Thus \(T_{\nu+1,\nu,\gamma}x\) lies in the
\((-1)^p\)-eigenspace of \(J_{\nu,\gamma}\), namely
\(V_{\nu,\gamma,\overline p}\).  In parity-homogeneous bases this is
equivalent to the two off-diagonal block identities
\(T_{\overline 0\overline 1}=T_{\overline 1\overline 0}=0\).  The
restricted matrices are the diagonal blocks.  The Chevalley
anti-involution and the \(A_\beta\)-comparison square are additional
finite matrix identities in the datum, so the same parity preservation
holds on negative roots and after comparison with the target
presentation.  The composition row gives functoriality of the even and
odd restricted maps.
\end{proof}

The packet
\[
\texttt{certificates/hall/transition\_parity\_decomposition\_preservation}
\]
records the current state of this datum.  Its verifier returns
\(\texttt{TRANSITION\_PARITY\_DECOMPOSITION\_PRESERVATION\_OBSTRUCTION\_VERIFIED}\):
the primitive-transition input, source and target-stage parity blocks,
parity-homogeneous bases, source and target fermion-parity involutions,
ambient transition matrices, parity-commutator identities,
off-diagonal zero identities, even and odd restricted transition
matrices, negative-root parity transport rows, comparison parity
squares, and parity-transition composition rows are not yet supplied.

\begin{definition}[Primitive-space Mittag--Leffler datum]
\label{def:primitive-space-ml-datum}
Fix a cofinal sequence \(W_1\subset W_2\subset\cdots\) as in
Definition~\ref{def:cofinal-primitive-recognition-datum}.  For each
retained Gram degree \(\gamma\) with
\(\alpha(\gamma)\in W_\nu\cup(-W_\nu)\), and for each parity
\(\overline p\), let
\[
P_{\nu,\gamma,\overline p}
:=P^{\Pi,\mathrm{red}}_{X,\gamma,\overline p}
\bigm|_{\alpha^{-1}(W_\nu\cup(-W_\nu))}
\]
denote the finite primitive Hall space after normal-ordered descent and
before quotienting by the Hopf radical.  When
\(W_\nu\subset W_{\nu'}\), transition preservation of primitive
subspaces gives a restricted map
\[
t_{\nu'\nu,\gamma,\overline p}:
P_{\nu',\gamma,\overline p}\longrightarrow
P_{\nu,\gamma,\overline p}.
\]
A primitive-space Mittag--Leffler datum consists of:
\begin{enumerate}
\item primitive-space rows for every
\((\nu,\gamma,\overline p)\), including a finite basis, parity, sign,
target root, and finite rank;
\item restricted primitive transition matrices \(t_{\nu'\nu}\), together
with identity and composition rows;
\item coverage rows for every retained positive and negative root
degree, for both parities, on the cofinal window sequence;
\item for every base window \(W_{\nu_0}\), every retained
\((\gamma,\overline p)\), and every later tower containing that degree,
an image-stabilization witness: there is \(\nu_1\ge \nu_0\) such that
for all \(\nu\ge\nu_1\),
\[
\operatorname{im}\bigl(
P_{\nu,\gamma,\overline p}\to
P_{\nu_0,\gamma,\overline p}\bigr)
=
\operatorname{im}\bigl(
P_{\nu_1,\gamma,\overline p}\to
P_{\nu_0,\gamma,\overline p}\bigr),
\]
with the stabilized image written in the chosen basis of
\(P_{\nu_0,\gamma,\overline p}\);
\item the defect rows
\[
\mathrm{strict}=1,\qquad \mathrm{ML}=1,\qquad
\operatorname{rank}R^1\varprojlim_\nu
P_{\nu,\gamma,\overline p}=0
\]
for every retained primitive-space tower, and compatibility rows for
the \(A_\beta\)-comparison maps.
\end{enumerate}
Transition preservation of primitive subspaces, finite rank, equality of
signed multiplicities, target-window containment, radical-transition
data, PBW-transition data, and denominator-product identities are not a
primitive-space Mittag--Leffler datum.
\end{definition}

\begin{proposition}[Vanishing of \(R^1\varprojlim\) on primitive spaces;
\textit{proved}]
\label{prop:primitive-space-lim1-vanishing-criterion}
For every tower \(P_{\nu,\gamma,\overline p}\) carrying the datum of
Definition~\ref{def:primitive-space-ml-datum}, the inverse system of
primitive spaces is Mittag--Leffler.  Consequently
\[
R^1\varprojlim_\nu P_{\nu,\gamma,\overline p}=0
\]
for every retained Gram degree, target root, sign, and parity.  Without
the rows of Definition~\ref{def:primitive-space-ml-datum}, the passage
from finite primitive Hall spaces to the completed primitive source is
not exact.
\end{proposition}

\begin{proof}
The primitive-space rows and restricted transition matrices define an
inverse system of finite-rank vector spaces once the identity and
composition rows have zero defect.  The image-stabilization row is
precisely the Mittag--Leffler condition: for each fixed base window,
the images of all sufficiently large windows inside the base primitive
space are constant.  The standard derived inverse-limit criterion for
inverse systems of vector spaces therefore gives
\(R^1\varprojlim=0\) on each primitive-space tower.  The conclusion is
degreewise and paritywise.  It cannot be inferred from pointwise
primitive kernels, dimension counts, signed Borcherds multiplicities, or
transition preservation without the stabilized-image rows.
\end{proof}

The packet
\[
\texttt{certificates/hall/primitive\_space\_lim1\_vanishing}
\]
records the current state of this datum.  Its verifier returns
\(\texttt{PRIMITIVE\_SPACE\_LIM1\_VANISHING\_OBSTRUCTION\_VERIFIED}\):
the cofinal root-window sequence, primitive-space rows, primitive bases,
finite-rank rows, parity and sign coverage rows, restricted transition
matrices, transition-functoriality rows, image subspace rows,
image-stabilization witnesses, strict Mittag--Leffler rows, zero
\(R^1\varprojlim\) rows, degreewise coverage rows, and comparison-tower
compatibility rows are not yet supplied.

\begin{definition}[Pairing-kernel Mittag--Leffler datum]
\label{def:pairing-kernel-ml-datum}
Assume the primitive-space rows of
Definition~\ref{def:primitive-space-ml-datum} have been supplied.  For
each retained \((\nu,\gamma,\overline p)\), a finite Hopf-pairing row is
a matrix for
\[
\lambda_{\nu,\gamma,\overline p}:
P_{\nu,\gamma,\overline p}\longrightarrow
P_{\nu,-\gamma,\overline p}^{\vee},
\qquad
x\longmapsto \langle x,-\rangle_\nu ,
\]
together with its opposite transpose
\[
\rho_{\nu,\gamma,\overline p}:
P_{\nu,-\gamma,\overline p}\longrightarrow
P_{\nu,\gamma,\overline p}^{\vee},
\qquad
y\longmapsto \langle -,y\rangle_\nu .
\]
Put
\[
K^{\ell}_{\nu,\gamma,\overline p}:=\ker\lambda_{\nu,\gamma,\overline p},
\qquad
K^{r}_{\nu,\gamma,\overline p}:=\ker\rho_{\nu,\gamma,\overline p}.
\]
The Hopf radical is formed later from these pairing kernels and the
radical-intersection/quotient data; a pairing-kernel datum by itself is
not a radical quotient.

A pairing-kernel Mittag--Leffler datum consists of:
\begin{enumerate}
\item finite Hopf-pairing matrices \(\lambda_{\nu,\gamma,\overline p}\)
and \(\rho_{\nu,\gamma,\overline p}\), with Hopf adjointness,
Frobenius-cyclicity, and quotient-nondegeneracy identity rows marking
that these are the pairings used in recognition;
\item finite bases and ranks for every left kernel
\(K^\ell_{\nu,\gamma,\overline p}\) and right kernel
\(K^r_{\nu,\gamma,\overline p}\);
\item transition maps between the left kernels and between the right
kernels, induced by the primitive transition matrices and the
pairing-transition identity, with identity and composition rows;
\item coverage rows for every retained positive and negative root
degree, both parities, and both kernel sides;
\item for every base window \(W_{\nu_0}\), retained
\((\gamma,\overline p)\), and side \(s\in\{\ell,r\}\), an
image-stabilization witness: there is \(\nu_1\ge\nu_0\) such that for
all \(\nu\ge\nu_1\),
\[
\operatorname{im}\bigl(K^s_{\nu,\gamma,\overline p}\to
K^s_{\nu_0,\gamma,\overline p}\bigr)
=
\operatorname{im}\bigl(K^s_{\nu_1,\gamma,\overline p}\to
K^s_{\nu_0,\gamma,\overline p}\bigr),
\]
with the stabilized image written in a finite kernel basis;
\item the defect rows
\[
\mathrm{strict}=1,\qquad \mathrm{ML}=1,\qquad
\operatorname{rank}R^1\varprojlim_\nu
K^s_{\nu,\gamma,\overline p}=0
\]
for every retained pairing-kernel tower, and compatibility rows with
the \(A_\beta\)-comparison to the target invariant-pairing kernels.
\end{enumerate}
Hopf-pairing matrices, rank equalities, quotient non-degeneracy,
primitive-space Mittag--Leffler exactness, radical-transition
preservation, target invariant kernels, and denominator identities are
not a pairing-kernel Mittag--Leffler datum.
\end{definition}

\begin{proposition}[Vanishing of \(R^1\varprojlim\) on pairing kernels;
\textit{proved}]
\label{prop:pairing-kernel-lim1-vanishing-criterion}
For every tower \(K^s_{\nu,\gamma,\overline p}\), \(s\in\{\ell,r\}\),
carrying the datum of
Definition~\ref{def:pairing-kernel-ml-datum}, the inverse system of
pairing kernels is Mittag--Leffler.  Consequently
\[
R^1\varprojlim_\nu K^s_{\nu,\gamma,\overline p}=0
\]
for every retained Gram degree, target root, sign, parity, and kernel
side.  Without the rows of
Definition~\ref{def:pairing-kernel-ml-datum}, kernel formation for the
Hopf pairing need not commute with passage from finite windows to the
completed primitive source.
\end{proposition}

\begin{proof}
The pairing-kernel rows and restricted transition matrices define an
inverse system of finite-rank vector spaces once identity and
composition defects vanish.  The image-stabilization row is exactly the
Mittag--Leffler condition for this kernel system.  The standard derived
inverse-limit criterion for inverse systems of vector spaces gives
\(R^1\varprojlim=0\) for each left and right pairing-kernel tower.  This
is a statement about the kernels as subspaces with their restricted
transition maps.  It does not follow from the existence of a Hopf
pairing matrix, from rank equality, from non-degeneracy after quotient,
or from Mittag--Leffler exactness of the ambient primitive spaces without
strict kernel-transition and stabilized-image rows.
\end{proof}

The packet
\[
\texttt{certificates/hall/pairing\_kernel\_lim1\_vanishing}
\]
records the current state of this datum.  Its verifier returns
\(\texttt{PAIRING\_KERNEL\_LIM1\_VANISHING\_OBSTRUCTION\_VERIFIED}\):
the cofinal root-window sequence, primitive-space input, Hopf-pairing
matrix rows, Hopf-pairing identity rows, left and right kernel rows,
kernel bases, finite-rank rows, pairing-transition input, restricted
kernel transition matrices, transition-functoriality rows, image
subspace rows, image-stabilization witnesses, strict Mittag--Leffler
rows, zero \(R^1\varprojlim\) rows, degreewise coverage rows, and
comparison-kernel compatibility rows are not yet supplied.

\begin{definition}[PBW associated-graded Mittag--Leffler datum]
\label{def:pbw-associated-graded-ml-datum}
Assume the radical quotient transition data and the PBW-transition
datum of Definition~\ref{def:transition-pbw-filtration-datum} have been
supplied.  Let
\[
L_\nu:=P^{\Pi,\mathrm{red}}_{X}|_{\alpha^{-1}(S_{W_\nu})}
/\operatorname{Rad}^{\Pi}_{X,\nu}
\]
be the finite primitive radical quotient on the window \(W_\nu\).  For
each PBW degree \(d\), retained Gram degree \(\chi\), and parity
\(\overline p\), put
\[
G^{\mathrm{PBW}}_{\nu,d,\chi,\overline p}
:=\Bigl(\operatorname{gr}^{\mathrm{PBW}}_dU(L_\nu)\Bigr)_{\chi,\overline p}.
\]
The associated-graded transition maps of
Proposition~\ref{prop:transition-pbw-filtration-preservation-criterion}
give maps
\[
g_{\nu'\nu,d,\chi,\overline p}:
G^{\mathrm{PBW}}_{\nu',d,\chi,\overline p}
\longrightarrow
G^{\mathrm{PBW}}_{\nu,d,\chi,\overline p}
\qquad (\nu\le\nu').
\]
A PBW associated-graded Mittag--Leffler datum consists of:
\begin{enumerate}
\item finite PBW filtration rows \(F^{\mathrm{PBW}}_{\le d}U(L_\nu)\),
lower filtration rows \(F^{\mathrm{PBW}}_{\le d-1}U(L_\nu)\), and
quotient rows for every
\(G^{\mathrm{PBW}}_{\nu,d,\chi,\overline p}\);
\item ordered word bases, relation-rewrite rows, associated-graded
bases, and associated-graded rank rows;
\item associated-graded transition matrices \(g_{\nu'\nu}\), together
with PBW-symbol intertwining, \(A_\beta\)-PBW comparison squares,
identity rows, and composition rows;
\item coverage rows for every retained PBW degree, Gram degree, parity,
side, and cofinal window;
\item for every base window \(W_{\nu_0}\), every retained
\((d,\chi,\overline p)\), and every side, an image-stabilization
witness: there is \(\nu_1\ge\nu_0\) such that for all \(\nu\ge\nu_1\),
\[
\operatorname{im}\bigl(
G^{\mathrm{PBW}}_{\nu,d,\chi,\overline p}\to
G^{\mathrm{PBW}}_{\nu_0,d,\chi,\overline p}\bigr)
=
\operatorname{im}\bigl(
G^{\mathrm{PBW}}_{\nu_1,d,\chi,\overline p}\to
G^{\mathrm{PBW}}_{\nu_0,d,\chi,\overline p}\bigr),
\]
with the stabilized image written in a finite associated-graded basis;
\item the defect rows
\[
\mathrm{strict}=1,\qquad \mathrm{ML}=1,\qquad
\operatorname{rank}R^1\varprojlim_\nu
G^{\mathrm{PBW}}_{\nu,d,\chi,\overline p}=0
\]
for every retained PBW associated-graded tower.
\end{enumerate}
PBW filteredness, associated-graded rank equality, target PBW rows,
Hilbert series, formal PBW pushforward, denominator identities, and
primitive or pairing-kernel Mittag--Leffler exactness are not a PBW
associated-graded Mittag--Leffler datum.
\end{definition}

\begin{proposition}[Vanishing of \(R^1\varprojlim\) on PBW associated
gradeds; \textit{proved}]
\label{prop:pbw-associated-graded-lim1-vanishing-criterion}
For every tower
\(G^{\mathrm{PBW}}_{\nu,d,\chi,\overline p}\) carrying the datum of
Definition~\ref{def:pbw-associated-graded-ml-datum}, the inverse system
of PBW associated-graded pieces is Mittag--Leffler.  Consequently
\[
R^1\varprojlim_\nu
G^{\mathrm{PBW}}_{\nu,d,\chi,\overline p}=0
\]
for every retained PBW degree, Gram degree, parity, and side.  Without
the rows of Definition~\ref{def:pbw-associated-graded-ml-datum}, PBW
associated gradeds need not commute with passage from finite windows to
the completed enveloping algebra.
\end{proposition}

\begin{proof}
The associated-graded piece rows and associated-graded transition
matrices define an inverse system of finite-rank vector spaces once the
identity and composition rows have zero defect.  The
image-stabilization row is precisely the Mittag--Leffler condition on
that inverse system.  The standard derived inverse-limit criterion for
inverse systems of vector spaces gives \(R^1\varprojlim=0\) on every
retained PBW associated-graded tower.  Filteredness alone gives maps of
filtered pieces; it does not imply that the images on the quotients
\(F_{\le d}/F_{\le d-1}\) stabilize.  Rank equality and Hilbert-series
agreement likewise do not compute the derived inverse limit.
\end{proof}

The packet
\[
\texttt{certificates/hall/pbw\_associated\_graded\_lim1\_vanishing}
\]
records the current state of this datum.  Its verifier returns
\(\texttt{PBW\_ASSOCIATED\_GRADED\_LIM1\_VANISHING\_OBSTRUCTION\_VERIFIED}\):
the cofinal root-window sequence, radical-quotient input,
PBW-transition input, PBW filtration rows, ordered word bases,
relation-rewrite rows, associated-graded piece rows, graded bases,
associated-graded rank rows, associated-graded transition maps,
PBW-symbol intertwining rows, \(A_\beta\)-PBW squares,
transition-functoriality rows, image subspace rows,
image-stabilization witnesses, strict Mittag--Leffler rows, zero
\(R^1\varprojlim\) rows, degreewise coverage rows, and
comparison-graded compatibility rows are not yet supplied.

\begin{proposition}[Strict Mittag--Leffler completion for primitive
recognition; \textit{proved}]
\label{prop:strict-ml-completion-primitive-recognition}
Let
\(\{W_\nu\}_{\nu\ge 1}\) be a cofinal increasing sequence of
downward-saturated finite root windows.  For each \(\nu\), let
\[
A_\nu:\ P^{\Pi,\mathrm{red}}_{X}|_{\alpha^{-1}(S_{W_\nu})}
\xrightarrow{\sim}
\mathfrak g_{\Delta_5}|_{S_{W_\nu}}
\]
be the finite-window isomorphism of
Proposition~\ref{prop:finite-window-presentation-recognition}.  Suppose
the transition maps for \(W_{\nu+1}\to W_\nu\) satisfy the following
strict Mittag--Leffler conditions:
\begin{enumerate}
\item the source primitive-space towers of
Definition~\ref{def:primitive-space-ml-datum} and the target root-space
parity towers are Mittag--Leffler, and \(A_\nu\) commutes with the
transition maps on each parity piece;
\item the pairing-kernel towers of
Definition~\ref{def:pairing-kernel-ml-datum}, the kernel towers of the
finite presentation maps, the Hopf radical intersection towers, and the
relation-radical quotient towers are strict subsystems: transition maps
carry each subspace onto the corresponding image and induce the quotient
transition maps;
\item the PBW filtrations are strict under transition, and the induced
towers on the associated graded pieces carry the
Mittag--Leffler datum of
Definition~\ref{def:pbw-associated-graded-ml-datum};
\item the triangular finite-window towers for positive-negative
cancellation have \(R^1\varprojlim=0\).
\end{enumerate}
Then inverse limit is exact on the finite recognition diagram.  In
particular,
\[
\varprojlim_\nu
\bigl(P^{\Pi,\mathrm{red}}_{X}|_{\alpha^{-1}(S_{W_\nu})}
/\operatorname{Rad}^{\Pi}_{X,\nu}\bigr)
\cong
\bigl(\varprojlim_\nu
P^{\Pi,\mathrm{red}}_{X}|_{\alpha^{-1}(S_{W_\nu})}\bigr)
\big/
\overline{\varprojlim_\nu\operatorname{Rad}^{\Pi}_{X,\nu}},
\]
the completed kernel is the closure of the finite kernels, the
completed PBW associated graded is the inverse limit of the finite PBW
associated gradeds, and the maps \(A_\nu\) assemble to an isomorphism
of completed Gram-graded Lie superalgebras
\[
\varprojlim_\nu
P^{\Pi,\mathrm{red}}_{X}|_{\alpha^{-1}(S_{W_\nu})}
\xrightarrow{\sim}
\varprojlim_\nu \mathfrak g_{\Delta_5}|_{S_{W_\nu}}.
\]
Without these strict Mittag--Leffler conditions, finite-window
isomorphisms do not by themselves identify the completed primitive
source with the completed BKM target.
\end{proposition}

\begin{proof}
For inverse systems of vector spaces, the derived exact sequence for
\(\varprojlim\) shows that exactness of
\[
0\to K_\nu\to V_\nu\to Q_\nu\to0
\]
passes to inverse limits exactly when \(R^1\varprojlim K_\nu=0\), or
more generally when the tower is Mittag--Leffler.  Apply this to the
presentation-kernel tower, the Hopf-radical tower, the quotient tower,
and each even/odd root-space tower.  Strictness says that taking
subobjects, quotients, and finite PBW filtered pieces commutes with the
transition maps; hence the inverse-limit quotient is the quotient by
the closure of the finite radicals, and the inverse-limit kernel is the
closure of the finite kernels.  The strict filtered statement gives
\[
\operatorname{gr}_{\mathrm{PBW}}\varprojlim_\nu U(P_\nu)
\cong
\varprojlim_\nu\operatorname{gr}_{\mathrm{PBW}}U(P_\nu)
\]
on every retained degree.

The maps \(A_\nu\) commute with the transition maps, so they form an
isomorphism of inverse systems.  Since inverse limit is a functor and
the kernel and cokernel towers of \(A_\nu\) are zero, the inverse-limit
map has zero kernel and cokernel.  If the Mittag--Leffler hypotheses are
removed, the \(R^1\varprojlim\) term can contribute a hidden kernel,
cokernel, radical element, PBW class, or parity-cancelling summand
which is invisible on each fixed finite window.  This is exactly the
completion failure prohibited by (R8).
\end{proof}

\begin{definition}[Finite source proof tables]
\label{def:finite-source-proof-tables}
After choosing homogeneous bases on a finite window \(W_\nu\), the
finite-window primitive-recognition datum of
Definition~\ref{def:cofinal-primitive-recognition-datum} is equivalent
to the following source proof tables:
\[
\begin{gathered}
\mathrm{degrees},\ \mathrm{parity\_blocks},\
\mathrm{basis\_provenance},\
\mathrm{simple\_representatives},\
M,\ D,\ \mathrm{unit\_counit},\ B,\ G,\
\mathrm{hopf\_pairing\_identities},\ K,\ Q,\ A,\\
\mathrm{hall\_bialgebra\_identities},\
\mathrm{radical\_ideal\_coideal},\
\mathrm{relation\_rows},\
\mathrm{no\_extra},\
\mathrm{generation},\
\mathrm{pbw},\
\mathrm{transitions},\
\mathrm{A}_{\beta}\mathrm{\_comparison\_maps}.
\end{gathered}
\]
When the same packet is used to assert the current-envelope statement
\(\Omega_E^{\mathrm{ch}}C_X\simeq\mathsf{Den}_{\Delta_5,E}\), it must
also contain the finite Koszul comparison tables
\[
\mathrm{koszul\_cones},\qquad
\mathrm{koszul\_comparison\_identities},\qquad
\mathrm{koszul\_transition\_ml}
\]
of Definition~\ref{def:finite-koszul-comparison-datum}.  They are not
implied by the \(A_\beta\)-comparison maps: \(A_\beta\) records the
primitive radical-quotient comparison, whereas the Koszul tables record
the acyclic cones, compatibility homotopies, and strict
Mittag--Leffler systems of the chiral cobar comparison.

Their mathematical meanings are fixed as follows.  The degree table
records the Gram labels \(\gamma\), their signs, window membership, and
compact source-admissibility.  The provenance table records the retained
charge lift, source stratum, vanishing-cycle or intersection-complex
summand, reduced orientation, quotient orientation,
Thom--Sebastiani transport, finite-stabilizer linearization, protected
integration, and transition compatibility of every basis vector.  The
parity table records the two ranks in each parity, not only their
difference.  The simple-representative table names the source basis
vectors realizing \(e_i^X,f_i^X,h_i^X\), together with their parity and
target reference.  The matrices \(M,D,B\) record, respectively, the
source Hall product, source coproduct, and supercommutator bracket on
the chosen source bases.  The \(\mathrm{unit\_counit}\) table records
the unit and counit maps.  The Hall bialgebra identity table records the
vanishing of the finite defects for associativity, coassociativity,
left and right unit, left and right counit, product-coproduct
compatibility, and primitive closure of the supercommutator.  The
matrix \(G\) records the Hopf pairing.  The Hopf-pairing identity table
records the finite defects for Hopf adjointness
\(\langle m(x,y),z\rangle=\langle x\otimes y,\Delta z\rangle\), the
Frobenius cyclic identity for the primitive bracket, and
non-degeneracy of the induced pairing on the radical quotient.  The
matrices \(K,Q\) record the radical kernel and a quotient splitting.
The matrix \(A\) records the comparison from the source radical quotient
to the imported target basis; target labels are allowed only in this
codomain.

Each row is \emph{geometrically generated}.  It carries a source
identifier for the compact correspondence, reduced stratum, vanishing-
cycle or intersection-complex summand, orientation transport,
Hopf-pairing construction, radical computation, quotient splitting,
comparison map, or transition morphism from which the row is obtained,
together with a proof reference for the relevant theorem, equation,
finite computation, or source certificate.  A row whose matrices satisfy
the displayed rank identities but are not obtained from these
compact-source data is not a row of the finite source proof tables.  In
particular, arbitrary matrices with the correct dimensions, target-only
rows, signed-only rows, and status-only rows do not define the compact
source primitive-recognition datum.

The remaining tables record identities.  The
\(\mathrm{radical\_ideal\_coideal}\) rows state that \(K\) is stable as
a Lie ideal and as a coproduct coideal.  The
\(\mathrm{relation\_rows}\) table records the Cartan, Chevalley, real
Serre, Borcherds orthogonality, and super sign relations.  The
\(\mathrm{no\_extra}\) table records the finite rank identity
\(\ker\pi_\nu=(\overline{J_{\mathrm{BK}}+\operatorname{Rad}_{\mathrm{GN}}})_{\le W_\nu}\).
The \(\mathrm{generation}\) table records, for each source basis vector
in the radical quotient, a bracket-word span by the Cartan and simple
primitive representatives with all intermediate degrees remaining in
\(W_\nu\).  The \(\mathrm{pbw}\) table records the associated-graded PBW
rank comparison.  The \(\mathrm{transitions}\) table records strictness
and Mittag--Leffler exactness under \(W_{\nu+1}\to W_\nu\).  The
\(\mathrm{A}_{\beta}\)-comparison table records that the comparison
matrix \(A\) intertwines the source bracket, coproduct, pairing, radical
quotient, and PBW filtration with the target blocks.
Admissibility of these finite tables is numerical: every defect-rank
row above must have rank \(0\), and every rank-comparison row must have
equal ranks on the two sides and on their combined span.  A non-zero
defect rank or a failed rank equality is a failed finite-window
recognition datum.  The identity rows are also exhaustive.  In the
finite source proof tables the canonical labels are
\[
\begin{gathered}
\texttt{unit\_left},\texttt{unit\_right},\texttt{counit\_left},
\texttt{counit\_right},\texttt{associativity},\texttt{coassociativity},
\texttt{bialgebra\_compatibility},\texttt{primitive\_closure},\\
\texttt{hopf\_adjointness},\texttt{frobenius\_cyclic},
\texttt{quotient\_nondegenerate},\quad
\texttt{lie\_ideal},\texttt{coproduct\_coideal},\\
\texttt{cartan},\texttt{chevalley},\texttt{real\_serre},
\texttt{borcherds\_orthogonality},\texttt{super\_sign},\quad
\texttt{bracket},\texttt{coproduct},\texttt{pairing},
\texttt{radical\_quotient},\texttt{pbw}.
\end{gathered}
\]
Omitting one family is a failed finite-window recognition datum.
If the current-envelope comparison is also asserted, the Koszul table
families listed above are part of the same finite packet; omitting one
of them is a failed finite Koszul comparison datum, even if
\((\mathrm{R0})\)--\((\mathrm{R8})\) hold algebraically.
\end{definition}

For the compact \(K3\times E\) Hall source packet used in this
manuscript, the companion obstruction ledger
\[
\texttt{certificates/sources/k3e\_compact\_hall/blocked\_obligations.csv}
\]
records precisely the missing source artifacts.  Its verifier returns
\(\texttt{COMPACT\_HALL\_OBSTRUCTION\_LEDGER\_VERIFIED}\), while still
recording \(\texttt{compact\_source\_recognition=false}\) and
\(\texttt{mathematical\_certification=false}\).  The ledger is therefore
not evidence for primitive recognition.  It is the finite list of
objects which must be supplied before \(P_X^{\Pi,\mathrm{red}}\) exists
as a constructed compact-source primitive Hall Lie superalgebra:
the coproduct \(D\), unit and counit, bialgebra compatibility, primitive
closure, Hopf pairing, radical ideal/coideal, no-extra kernel equality,
PBW comparison, strict Mittag--Leffler transition, the
\(A_\beta\)-intertwining rows, and the finite Koszul comparison rows.

\begin{proposition}[The finite source proof tables are exactly the
finite recognition datum; \textit{proved}]
\label{prop:finite-source-proof-tables-equivalence}
For a fixed downward-saturated finite window \(W_\nu\), a populated
system of finite source proof tables in the sense of
Definition~\ref{def:finite-source-proof-tables}, with all stated rank
and intertwining identities verified over the chosen coefficient ring
and with every row geometrically generated from the compact-source
correspondences and orientation/Hopf data named there,
is equivalent to clauses \textup{(R0)}--\textup{(R8)} of
Definition~\ref{def:cofinal-primitive-recognition-datum} after a choice
of homogeneous bases.  In particular, a header-only table system, or a
table system whose rows contain status text but no mathematical payload,
or a table system of arbitrary matrices without compact-source
provenance, does not define a finite-window primitive-recognition datum.
\end{proposition}

\begin{proof}
Choosing homogeneous bases turns each finite-dimensional source
primitive space, pairing block, radical, quotient, and comparison map
into matrices, but only after the corresponding finite Hall
correspondence, reduced orientation, Hopf pairing, radical quotient, and
transition maps have been constructed.  The row-level geometric source
and proof reference record this construction before coordinates are
chosen.  The degree and provenance tables are precisely (R0).
The simple-representative table is (R1).  The parity table is (R3).
The representative rows together with the relation table are
(R1)--(R2).
The \(M,D,\eta,\epsilon,B\) matrices and the Hall bialgebra identity
table are the product-coproduct and primitive-closure part of (R4).
The \(G\) matrix and the Hopf-pairing identity table are the pairing
adjointness, Frobenius, and quotient non-degeneracy part of (R4).  The
\(K,Q\) matrices and the
\(\mathrm{radical\_ideal\_coideal}\) identities are the radical part of
(R4).  The
\(\mathrm{no\_extra}\) rank identity is (R5).  The generation span table
is (R6).  The PBW table is (R7), and the transition table is (R8).  The
\(\mathrm{A}_{\beta}\)-comparison rows assert that the source quotient
is compared to the same target blocks in every operation, so the matrix
statements are independent of the chosen coordinates.  Conversely, any
finite-window datum satisfying (R0)--(R8) yields these tables after
choosing bases and writing the finite maps and identities in matrix
form, with the geometric source and proof reference inherited from the
finite maps themselves.  Empty rows have no source vectors, maps, ranks,
or identities; arbitrary matrices with no compact-source origin have no
Hall correspondence, orientation datum, pairing construction, or
transition map.  They satisfy none of the clauses.
\end{proof}

\begin{proposition}[Global recognition is exactly the cofinal datum;
\textit{proved}]
\label{prop:global-recognition-cofinal-datum}
Assume (S1)--(S4) of Theorem~\ref{thm:primitive-recognition}.  Then
the global source-side assertions (S5)--(S10) are equivalent to
a cofinal finite-window primitive-recognition datum
$\{W_\nu,(\textup{R0})\text{--}(\textup{R8})\}$ in the sense of
Definition~\ref{def:cofinal-primitive-recognition-datum}, and
conversely.
\end{proposition}

\begin{proof}
The two sides record the same data: each $W_\nu$ is a saturated
finite window in (S5)--(S9), and (R3) is exactly the
two-integer parity equality of (S10) on $W_\nu$.  The Hopf pairing
matrices, kernel equalities, PBW filtrations, no-extra-relations
relations, and transition strictness are exactly the cofinal
contents of (S5)--(S9), and the cusp-positive arithmetic side of
$f(nm,l)$ is recovered from the parity equalities (R3) on the active
target degrees.  Conversely, given (R0)--(R8) on a cofinal system,
the strict Mittag--Leffler completion theorem
Proposition~\ref{prop:strict-ml-completion-primitive-recognition}
gives the global statements (S5)--(S9), and (R3) on a cofinal
exhaustion gives (S10).
\end{proof}

\subsubsection*{Recognition criterion at the window $W_{\le 3}$}

For a finite window $W$ with $S_W=W\cup(-W)\cup\{0\}$, a compact
source datum consists of compact-source bases
$b^X_{R,\gamma_\beta,\overline p,a}$, parity pieces, product,
coproduct, bracket, pairing, radical, and quotient matrices
$M,D,B,G,K,Q$, comparison maps $A_{\beta,\overline p}$, relation
identities, the no-extra kernel equality, PBW associated gradeds,
and strict Mittag--Leffler transition maps.  The target labels
$e_i,E_{ij},u_{ij,r},T_i,M_{ij,r},w_s$ occur only in target blocks
or as codomain coordinates of $A_{\beta,\overline p}$.

The window $W_{\le 3}$ is the first arithmetic window, not a
relation-closed degree window.  A relation-closed degree test must include
real-Serre terminal rows, doubled isotropic rows, $\delta_{123}$-real
strings, odd--odd $\delta_{123}$ rows, and complementary
height-seven strings, or pass to the downward-saturated closure.  In
the target-reference packet \(\texttt{wle3\_target\_parity}\),
\(W_{\le3}\) has the verified rows \(\delta_i:1|0\),
\(a_{ij}:10|0\), \(2\delta_i+\delta_j:1|0\), and
\(\delta_{123}:29|93\); its scalar firewall excludes source parity,
Hall brackets, pairing radicals, PBW comparison, and primitive
recognition.
The finite packet \(\texttt{parity\_pushforward}\) verifies a different
formal statement: once finite source parity blocks are supplied, the
additive normal-ordered Gram pushforward preserves the even and odd
blocks separately.  It supplies neither the source parity dimensions
required in (R3) nor the Gritsenko--Nikulin/Kac target parity equality;
those remain primitive-recognition data.
The finite packet \(\texttt{pbw\_pushforward}\) is likewise formal: once
a finite PBW filtration is supplied, additive normal-ordered Gram
pushforward preserves its filtered pieces and associated graded.  It
does not supply the compact-source PBW comparison required in (R7), the
target PBW rows, or the no-extra-relations theorem.
The finite packet \(\texttt{triangular\_pushforward}\) is the analogous
formal statement for a supplied positive/Cartan/negative splitting: the
normal-ordered pushforward preserves the three direct-sum parts and the
displayed bracket signs.  It does not supply the exact triangular
completion required in (R8), target triangular completion, or the
Mittag--Leffler exactness of triangular finite-window systems.
The finite target degree-closure packet
\(\texttt{wrel\_degree\_closure}\) verifies that the enlarged target
degree set has exactly the \(35\) rows
\[
W_{\le3}^{act}\cup R_4\cup I_4\cup C_{\le7}\cup\{2\delta_{123}\},
\]
with \(R_4\) and \(C_{k,5}\) terminal zero rows.  Its
downward-saturation audit finds exactly three nonzero proper
subdegree defects, namely
\[
D_k=\delta_k+2a_{ij},\qquad \{i,j,k\}=\{1,2,3\},
\]
below \(2\delta_{123}\), each with signed multiplicity \(-513\).
This is still only the target degree closure.
In the larger target degree window,
Remark~\ref{rem:bkm-relation-window-target-parity}
supplies full parity only for \(2a_{ij}\), \(C_{k,3}\), \(C_{k,4}\),
and \(C_{k,5}\).  The rows \(C_{k,2}\), \(D_k\), and
\(2\delta_{123}\) remain signed target rows; they cannot serve as
source comparison codomains until the finite target presentation is
computed in those degrees.  The base-string terminal degree set is
therefore not yet a primitive-recognition window in the sense of
(R0)--(R8).
The A071 target-presentation verifier records this boundary: it checks
the \(2a_{ij}\), \(C_{k,3}\), \(C_{k,4}\), and \(C_{k,5}\) parity rows,
and rejects \(C_{k,2}\), \(D_k\), and \(2\delta_{123}\) as target
basis, pairing, radical, or PBW rows.  The failure is structural, not
a missing scalar coefficient.  The minimal nonzero downward-saturation
extension \(W_{\mathrm{rel}}^+\cup\{D_1,D_2,D_3\}\) is now
target-recorded; a primitive-recognition window must still supply the
missing target-presentation and compact-source data.  The
post-saturation real-string audit records \(21\) terminal codomains
outside the saturated extension, \(10\) with nonzero signed
multiplicity.  The terminal-layer audit groups these codomains into
\(21\) distinct terminal degrees; adjoining them creates \(114\)
nonzero subdegree-defect incidences across \(26\) distinct subdegrees.
The next target-degree computation adjoins exactly these \(26\) rows and
gives an \(85\)-row degree set with zero remaining nonzero
proper-subdegree defect.  This closes downward saturation for the
displayed terminal layer, but full finite relation closure remains a
further target closure problem: if these \(26\) rows are promoted to
target-generator rows, the real-root strings force \(55\) terminal
checks, \(54\) terminal codomains outside the \(85\)-row set, and
\(12\) missing codomains with nonzero signed multiplicity.  Adjoining
those \(55\) terminal codomains would create \(546\) nonzero
proper-subdegree defect incidences across \(98\) distinct subdegrees.
Adjoining those \(98\) rows gives a \(237\)-row target degree set with
zero remaining nonzero proper-subdegree defect for that layer.  If this
\(98\)-row layer is promoted to target-generator rows, the real-root
strings force \(206\) terminal checks, \(182\) new terminal degrees, and
\(24\) missing nonzero terminal degrees.  The exact terminal-layer table
records all \(206\) codomains, and the exact saturation-extension table
records the \(50\) distinct nonzero proper subdegrees created by
adjoining them, with \(624\) incidences.  Adjoining both layers gives a
\(469\)-row target degree set with zero remaining nonzero
proper-subdegree defects for that layer.  Promoting the newly adjoined
\(50\)-row layer forces \(96\) real-string terminal checks, \(90\) new
terminal degrees, no missing nonzero terminal degree, and a \(12\)-row
nonzero saturation extension with \(44\) incidences.  Adjoining these
rows gives a \(571\)-row target degree set with zero remaining nonzero
proper-subdegree defects for that layer.  Promoting the newly adjoined
\(12\)-row layer forces \(20\) real-string terminal checks, all with
zero signed terminal multiplicity, and a \(6\)-row nonzero saturation
extension with \(8\) incidences.  Adjoining these rows gives a
\(597\)-row target degree set with zero remaining nonzero
proper-subdegree defects for that layer.  Promoting the newly adjoined
\(6\)-row layer forces \(10\) real-string terminal checks, all with
zero signed terminal multiplicity, and a \(6\)-row nonzero saturation
extension with \(18\) incidences.  Adjoining these rows gives a
\(613\)-row target degree set with zero remaining nonzero
proper-subdegree defects for that layer.  Promoting the next newly
adjoined \(6\)-row layer forces \(10\) real-string terminal checks, all
with zero signed terminal multiplicity, and a \(6\)-row nonzero
saturation extension with \(8\) incidences.  Adjoining these rows gives
a \(629\)-row target degree set with zero remaining nonzero
proper-subdegree defects for that layer.  Promoting the next newly
adjoined \(6\)-row layer forces \(10\) real-string terminal checks, all
with zero signed terminal multiplicity, and a \(6\)-row nonzero
saturation extension with \(18\) incidences.  Adjoining these rows gives
a \(645\)-row target degree set with zero remaining nonzero
proper-subdegree defects for that layer.  The tail is not a finite
accident.  For every even \(N\ge 14\), put
\[
\begin{aligned}
A_N&=\{(0,N,N-1),(0,N,N),(0,N+1,N),
      (N,0,N-1),(N,0,N),(N+1,0,N)\},\\
B_N&=\{(0,N,N+1),(0,N+1,N+1),(0,N+1,N+2),
      (N,0,N+1),(N+1,0,N+1),(N+1,0,N+2)\}.
\end{aligned}
\]
The verifier
\(\texttt{wrel\_tail\_staircase}\) proves symbolically, using the
Lorentzian pairing and the \(q^0\)-term \(y^{-1}+10+y\) of
\(\phi_{0,1}\), that the real-root terminal audit and downward
saturation force
\[
A_N\Longrightarrow B_N,\qquad B_N\Longrightarrow A_{N+2}.
\]
Its finite-anchor table checks the six rows of \(A_{14}\) inside the
\(\texttt{wrel\_degree\_closure}\) packet.  Hence the target-degree
audit contains an infinite six-row staircase.  Its
\(\texttt{unbounded\_tail.csv}\) table records the subsequences
\(A_{14+2k}\) and \(B_{14+2k}\) with a strictly increasing coordinate.
The source comparison can therefore be formulated only through finite
presentation data, not by closing a finite target-degree enumeration.
To promote this target audit to the first relation-closed compact-source
theorem, one must populate the first-window packet
\[
\mathrm{window\_closure},\quad
\mathrm{target\_source\_representatives},\quad
\mathrm{chevalley\_serre\_matrices},\quad
\mathrm{hall\_boundary\_complex},\quad
\mathrm{spectral\_sequence},\quad
\mathrm{kernel\_matrices},\quad
\mathrm{pbw\_graded},\quad
\mathrm{first\_window\_theorems},\quad
\mathrm{transitions},\quad
\mathrm{scalar\_firewall}.
\]
The theorem rows include the first-window instances of \(O1\),
\(O1^+\), \(O2\), \(P_{\mathrm{fin}}\), Hall product, Hall coproduct,
Hopf pairing, Koszul comparison, PBW comparison, primitive
recognition, and Pfaffian equality.  A target relation-closed degree
set without these compact-source rows is not a source theorem.
At present the packet
\(\texttt{certificates/first\_window/k3e\_relation\_closed\_window}\)
verifies only the obstruction ledger
\(\texttt{FIRST\_WINDOW\_OBSTRUCTION\_LEDGER\_VERIFIED}\), with
\(\texttt{first\_window\_certification=false}\).  Thus the first
relation-closed source theorem remains open until these finite rows are
supplied.

\begin{proposition}[Target-admissibility obstruction for the
height-seven closure; \textit{proved}]
\label{prop:target-admissibility-obstruction-wrel}
Let \(W\) be a downward-saturated finite target-root window containing
the height-\(7\) closure \(W_{\mathrm{rel}}^+\) of
Chapter~\ref{ch:pfaffian-dirac}.  If \(W\) is a
primitive-recognition window satisfying (R0)--(R8), then the target
restriction must contain a finite target-presentation datum.  In
particular, for every row \(C_{k,2}\), \(D_k=\delta_k+2a_{ij}\), and
\(2\delta_{123}\), it must supply the full even/odd parity dimensions,
the positive-negative pairing block, the target radical block, and the
PBW associated-graded row.  The signed target identities
\[
 \sdim(C_{k,2})=108,\quad m(C_{k,2})=90,\qquad
 \sdim(D_k)=-513,\quad m(D_k)=54,\qquad
 \sdim(2\delta_{123})=4016,\quad m(2\delta_{123})=-540
\]
do not supply these data.  Hence the presently available target data
do not make the saturated extension of \(W_{\mathrm{rel}}^+\) a
primitive-recognition window.
\end{proposition}

\begin{proof}
The rows \(C_{k,2}\) and \(2\delta_{123}\) lie in
\(W_{\mathrm{rel}}^+\), and \(D_k\) lies in every downward-saturated
window containing \(W_{\mathrm{rel}}^+\), since
\[
0<D_k=\delta_k+2a_{ij}<2\delta_{123}.
\]
By
Proposition~\ref{prop:target-presentation-prerequisite}, they must be
target-presentation rows.  Clause (R3) requires the two integers
\(\rootmult_{\overline 0}^{\mathrm{GN}}(\beta)\) and
\(\rootmult_{\overline 1}^{\mathrm{GN}}(\beta)\) for every
\(\beta\in W\), and also in the negative degree by the
Borcherds--Kac Chevalley anti-involution.  A signed multiplicity gives
only their difference.  The integer \(m(\beta)\) records the
Gritsenko--Nikulin imaginary-simple contribution in the denominator
presentation; it is not the full quotient root-space parity split and
does not determine an invariant positive-negative pairing.

Clauses (R4), (R5), and (R7) then require the target pairing kernel,
radical block, and PBW associated-graded comparison in the same
degrees.  Without those target blocks the source-to-target matrices
\(A_{\beta,\overline p}\) have no defined codomain in these rows, and
the equations expressing Hopf radical equality, no-extra-relations, and
PBW compatibility are not well-formed.  Thus signed data in
\(C_{k,2}\), \(D_k\), and \(2\delta_{123}\) are not admissible target
data for (R0)--(R8).
\end{proof}

\subsubsection*{What signed dimensions miss}

\begin{lemma}[Signed dimensions are not recognition;
\textit{proved}]
\label{lem:signed-dimensions-miss}
Let $V_\bullet$ be a $\widehat\Gamma_X$-graded super vector space
with zero bracket and zero coproduct, with the same signed
superdimensions $\sdim V_\gamma=f(nm,l)$ as
$P_X^{\Pi,\mathrm{red}}$.  Then $V_\bullet$ has the same Borcherds
denominator $64^{-1}\Delta_5(2Z)$ as $\mathfrak g_{\Delta_5}$ but is
not isomorphic to it as a Lie superalgebra.
\end{lemma}

\begin{proof}
The denominator depends only on signed superdimensions; the bracket
data, no-extra-relations theorem, Hopf-radical coideal stability, and
PBW filtration are independent.  If the Cartan is mapped nontrivially,
the Chevalley relation \([e_i,f_i]=h_i\) fails in \(V_\bullet\).  If
the Cartan is killed, the root-degree Cartan identification fails.  In
either case the kernel of the free presentation map contains all
non-Cartan brackets, hence is strictly larger than
\(\overline{J_{\mathrm{BK}}+\operatorname{Rad}_{\mathrm{GN}}}\) in
degrees where the BKM bracket is nonzero.  With the zero Hopf pairing
the radical is all of \(V_\bullet\), so the radical quotient is zero.
\end{proof}

This lemma is the counterexample that prohibits the determinant route
to recognition.  It is also the structural reason the ten clauses
(S1)--(S10) do not collapse to their signed-dimension shadow:
the signed dimensions are necessary and not sufficient; (S1)--(S9)
supply the source presentation, and (S10) supplies the missing parity
split.

\begin{lemma}[Finite-window scalar exclusions;
\textit{proved}]
\label{lem:first-window-scalar-exclusions}
Let \(W\) be a finite target-root window.  None of the following data
implies a finite-window primitive-recognition datum on \(W\) in the
sense of Definition~\ref{def:cofinal-primitive-recognition-datum}:
target labels without compact representatives, signed dimensions, the
scalar trace \(Z_{BPS}^{K3\times E}\), the Pfaffian product, the
Borcherds denominator product, target PBW rows without source PBW
strictness, arbitrary matrices with the right ranks but no compact
provenance, or rows whose only content is a verification status.
\end{lemma}

\begin{proof}
Clauses (R0)--(R8) require source objects before comparison: compact
basis vectors, source representatives, Hall product and coproduct,
unit and counit, Hopf pairing, radical and quotient matrices, relation
rows, no-extra kernel equality, generation spans, PBW strictness, and
Mittag--Leffler transition maps.  Target labels give only codomain
coordinates for \(A_{\beta,\bar p}\), not the source vectors mapped to
them.  Signed dimensions and the denominator product give one integer
per degree; (R3) requires the two parity dimensions, while (R4)--(R8)
require operations and kernels.  Lemma~\ref{lem:signed-dimensions-miss}
exhibits a zero-bracket object with the same denominator and no
recognition isomorphism.  The scalar trace and the Pfaffian product
live after taking a protected scalar shadow; they forget the bracket,
coproduct, Hopf radical, PBW filtration, and finite transition system.
Target PBW rows test only the imported target presentation and do not
give source PBW strictness or a source-to-target intertwining matrix.
Arbitrary matrices with correct ranks do not supply a Hall
correspondence, reduced orientation, compact cycle, pairing
construction, or transition morphism, so they fail the geometric
provenance part of (R0) and the construction clauses of (R4)--(R8).
Status-only rows contain no vector space, map, matrix, kernel, rank, or
identity.  Hence each listed substitute omits at least one structural
component of (R0)--(R8), and none defines a finite-window
primitive-recognition datum.
\end{proof}

\subsubsection*{Obstruction classes}

The recognition lane carries the residuals
\[
\mathfrak O_{\mathrm{Rec},\nu}=
\bigl(
\mathfrak o^{(\textup{R1})}_\nu,
\ldots,
\mathfrak o^{(\textup{R8})}_\nu,
\mathfrak o^{\mathrm{tr}}_{\nu+1,\nu}
\bigr),
\]
required to vanish on every cofinal window.  The first arithmetic
window $W_{\le 3}$ is achievable from the GN target side; the first
relation-closed degree window includes the real-Serre terminal rows and the
$\delta_{123}$ odd--odd rows, but some added target rows are still only
signed target rows by
Remark~\ref{rem:bkm-relation-window-target-parity}.  The
finite-stage components of the
(R0)--(R8) criterion, including source-radical/Hopf-coideal stability,
remain source-recognition data.  Appendix~\ref{app:obstruction-discharges}
supplies the orientation lane; it does not supply the source matrices
\(M,D,B,G,K,Q\), the comparison maps \(A_{\beta,\bar p}\), the kernel
equality, the PBW associated-graded comparison, or the finite Koszul
comparison residual
\(\mathfrak O_{\mathrm{Kcmp},R}\).  Construction of the
source-internal recognition criterion at any one relation-closed degree
window beyond \(W_{\le 3}\), together with the Koszul compatibility
checks of Definition~\ref{def:finite-koszul-comparison-datum}, is the
unresolved compact-source datum of Chapter~\ref{ch:open-obstructions}.

\subsubsection*{Comparison with the denominator algebra}

The recognition theorem is the source side of the Vol III Drinfeld
double construction.  The target presentation uses the generalized
Kac--Moody convention of Borcherds \cite{BorcherdsGKM88} and Kac
\cite{KAC:1}; the automorphic correction algebra is the
Gritsenko--Nikulin algebra of \cite[Sections~3--4,
Proposition~3.1]{GN}, and the product-lift data are those of
\cite[Theorem~2.1, \S5.1.1]{GNII} together with
Borcherds' automorphic product \cite{Borcherds95,B}.  The identity
$\mathrm{CoHA}(\C^3)=W_{1+\infty}$ belongs to the local toric case;
the present recognition criterion is at the level of the protected primitive Hall
algebra, before Drinfeld doubling; after the compact source datum is
supplied, it identifies that algebra with $\mathfrak g_{\Delta_5}$, not
with the centred or
vacuum-evaluated double.

Entries $(\mathrm L 1)$--$(\mathrm L 6)$ name the constructed
objects.  The finite HN system in
\S\ref{subsec:eight-finite-constructions} supplies the charge windows,
Hall correspondences, orientation transports, Pfaffian lines, Koszul
source coalgebras, primitive-recognition data, and their local
obstruction classes.


\subsection{Eight finite constructions}
\label{subsec:eight-finite-constructions}

The seventh entry of the Dirac--Igusa pro-object is the finite HN
system of charge windows, Hall correspondences, orientation transports,
Pfaffian lines, Koszul source coalgebras, and primitive-recognition
data.  At HN height \(R\) these data are finite-type or
finite-dimensional, and their obstruction classes vanish exactly when
the transition system defines the inverse limit
\(\mathfrak D_X^{\mathrm{DI}}\).

At fixed HN height all objects lie on finite-type stacks and
finite-dimensional Borel--Moore chain spaces; the inverse limit is taken
after the finite stages have compatible transition maps.  For
$R\in\mathcal R_{\mathrm{HN}}$, the constructions assume:
$[\textup{K3$\times$E-fin}]$ (finite-type semistable substacks
at bounded HN height),
$[\textup{K3$\times$E-atlas}]$ (finite d-critical/cosection
Hall atlas with Joyce--Upmeier orientation lines), and
$[\textup{ML}]$ (Mittag--Leffler descent).

At finite height the first two hypotheses are the table system
\[
\mathrm{hn\_type\_bounds},\quad
\mathrm{semistable\_substacks},\quad
\mathrm{derived\_enhancements},\quad
\mathrm{universal\_complexes},\quad
\mathrm{scalar\_rigidifications},\quad
\mathrm{stratifications},\quad
\mathrm{extension\_flag\_stacks},\quad
\mathrm{cosection\_atlas},\quad
\mathrm{transitions},\quad
\mathrm{scalar\_firewall}.
\]
The scalar firewall excludes Liu stability alone, the formal charge
window, a target window, a Hilbert-scheme scalar test, or a Pfaffian
product as finite-moduli data.
The finite-moduli packet
\(\texttt{certificates/moduli/k3e\_finite\_moduli}\) currently records
the HN-bound, semistable-substack, derived-enhancement, and
rigidification-inertia rows, together with the finite class-bound
and universal-complex rows and the \(E\)-translation-rigidification
rows, retained closed-substack rows, and retained extension-closure
rows, retained HN-factor-closure rows, and retained dual-closure rows.
The remaining rows stay in the obstruction ledger
\(\texttt{MODULI\_OBSTRUCTION\_LEDGER\_VERIFIED}\), with
\(\texttt{moduli\_certification=false}\).  It is still not a
construction of the full \([K3\times E\text{-fin}]\) package or
\([K3\times E\text{-atlas}]\).

\begin{definition}[Finite-moduli transition geometry datum]
\label{def:finite-moduli-transition-geometry}
Let \(R'\le R\).  A finite-moduli transition geometry datum from height
\(R\) to height \(R'\) consists of the following rows.
\begin{enumerate}
\item Transition morphisms
\[
t^{\mathrm{obj}}_{RR'}:\mathfrak M^{\mathrm{obj}}_R\to
\mathfrak M^{\mathrm{obj}}_{R'},\qquad
t^{\mathrm{ext}}_{RR'}:\mathfrak M^{\mathrm{ext}}_R\to
\mathfrak M^{\mathrm{ext}}_{R'},\qquad
t^{\mathrm{fl}}_{RR'}:\mathfrak M^{\mathrm{fl}}_R\to
\mathfrak M^{\mathrm{fl}}_{R'}
\]
on the retained object, extension, and two-step flag stacks.
\item Retained closed substacks
\[
\mathfrak Z^{a}_R\subset \mathfrak M^{a}_R,\qquad
\mathfrak Z^{a}_{R'}\subset \mathfrak M^{a}_{R'},
\qquad a\in\{\mathrm{obj},\mathrm{ext},\mathrm{fl}\},
\]
such that \(t^a_{RR'}(\mathfrak Z^a_R)\subset \mathfrak Z^a_{R'}\).
\item For each \(a\), the restricted morphism
\[
t^a_{RR'}|\mathfrak Z^a_R:\mathfrak Z^a_R\to\mathfrak Z^a_{R'}
\]
is either proper or a closed immersion of finite-type stacks.  In the
closed-immersion case the image, equivalently the graph, is recorded as
a closed substack.
\item The universal complexes, scalar rigidifications, stabilizer
stratifications, cosections, vanishing-cycle domains, and orientation
domains pull back to the corresponding rows on the source stage.
\item The base-change, retained-locus, stabilizer, orientation-domain,
vanishing-cycle-domain, and composition defect ranks are zero.
\end{enumerate}
Root windows, source-window indices, signed multiplicities, Pfaffian
products, and scalar traces are not transition morphisms of moduli
stacks.
\end{definition}

\begin{proposition}[Proper or closed finite-moduli transitions;
\textit{proved}]
\label{prop:finite-moduli-transition-geometry-criterion}
Given a finite-moduli transition geometry datum of
Definition~\ref{def:finite-moduli-transition-geometry}, the
\(\textup{(M2)}\) component of
Definition~\ref{def:finite-stage-dirac-igusa-morphism} is a system of
proper or closed transition morphisms on the retained object,
extension, and two-step flag stacks, compatible with the cosection,
vanishing-cycle, and Joyce--Upmeier orientation domains.
\end{proposition}

\begin{proof}
The datum supplies the three stack morphisms and their restrictions to
the retained closed substacks.  Item \(3\) says that each restricted
component is proper or a closed immersion of finite-type stacks.  The
closed-locus condition in item \(2\) makes these restrictions the
transition morphisms of the retained finite stages.  Item \(4\) gives
the required pullback compatibility for universal complexes,
cosections, vanishing-cycle domains, and orientation domains.  The
zero-defect rows in item \(5\) give strict base-change and composition.
Thus the moduli part of the finite-stage transition is exactly the
\(\textup{(M2)}\) datum.
\end{proof}

The packet
\(\texttt{certificates/moduli/finite\_moduli\_transition\_geometry}\)
records the present state of this datum.  Its verifier status is
\texttt{FINITE\_MODULI\_TRANSITION\_GEOMETRY\_OBSTRUCTION\_VERIFIED}:
the object-stack, extension-stack, and two-step flag-stack transition
morphisms, the proper-or-closed component rows, the finite-type
source-target rows, the closed image or graph rows, and the
base-change, retained-locus, composition, orientation-domain, and
vanishing-cycle-domain zero-defect rows are not yet supplied.  Hence
the finite-moduli transition-geometry theorem remains open.

\begin{definition}[Finite-stage Dirac--Igusa category]
\label{def:finite-stage-dirac-igusa-morphism}
For \(R\in\mathcal R_{\mathrm{HN}}\), a finite-stage
Dirac--Igusa datum is the tuple
\[
\mathfrak D^{\mathrm{DI}}_{X,R}=
\bigl(
\mathcal A_{X,R}^{E_3},F_{X,R}^{\mathrm{hyb}},
\Gamma_{X,R},\Pi_X,\Phi_R,o_R,H_R,P_R^\Pi,
C_{X,R},\Theta_{\mathrm{Kos},R},
\mathscr L_{\mathrm{Pf},R},\operatorname{pf}_{X,R},
\varepsilon_{o,R},\mathrm{Rec}_R
\bigr)
\]
with the finite data and residual equations specified in
\S\ref{subsec:eight-finite-constructions}.  The category
\(\mathsf{DI}^{\mathrm{fin}}_{X,R}\) has these tuples as objects and
the morphisms below as arrows.  The finite-stage category
\(\mathsf{DI}^{\mathrm{fin}}_X\) is the disjoint union over retained
heights, together with transition arrows for comparable heights.
For \(R'\le R\), a transition morphism
\[
\rho_{RR'}:\mathfrak D^{\mathrm{DI}}_{X,R}\longrightarrow
\mathfrak D^{\mathrm{DI}}_{X,R'}
\]
is the following compatible system.
\begin{enumerate}
\item[\textup{(M1)}] A restriction of charge windows and central
extensions \(\widehat\Gamma_R\to\widehat\Gamma_{R'}\) preserving
\(\Pi_X\), the normal-ordering cocycle \(B\), and the target
Gram labels.
\item[\textup{(M2)}] Proper or closed transition morphisms of the
object, extension, and two-step flag stacks, with the corresponding
cosection, vanishing-cycle, and orientation domains, as in
Proposition~\ref{prop:finite-moduli-transition-geometry-criterion},
and with reduced vanishing-cycle transport as in
Proposition~\ref{prop:transition-vanishing-cycle-preservation-criterion}.
\item[\textup{(M3)}] A morphism of hybrid local/wrapped
factorization carriers over \(\operatorname{Ran}^{\mathrm{hyb}}(E)\),
compatible with the eight-word flag atlas, quotient-after-correspondence
pseudofunctor, and protected integration.
\item[\textup{(M4)}] Isomorphisms of quotient-orientation data,
Weyl lifts, and type-II wall charts preserving the orientation
cocycle, normal Pfaffian units, wall coordinates, and the local rank
\(N_{\delta_i}^{\mathrm{Pf}}=1\), with orientation transport as in
Proposition~\ref{prop:transition-orientation-preservation-criterion}.
\item[\textup{(M5)}] A morphism of finite Hall bialgebra data
preserving product, coproduct, unit, counit, Hopf pairing, primitive
projection, radical quotient, and normal-ordered Gram descent, with
product preservation as in
Proposition~\ref{prop:transition-hall-product-preservation-criterion}
and coproduct preservation as in
Proposition~\ref{prop:transition-hall-coproduct-preservation-criterion},
and primitive-subspace preservation as in
Proposition~\ref{prop:transition-primitive-subspace-preservation-criterion},
and radical preservation as in
Proposition~\ref{prop:transition-radical-preservation-criterion}.
\item[\textup{(M6)}] A Pfaffian-line transition
\(\rho^{\mathrm{Pf}}_{RR'}:\mathscr L_{\mathrm{Pf},R}|_{\mathfrak Q_{R'}}
\xrightarrow{\sim}\mathscr L_{\mathrm{Pf},R'}\) carrying
\(\operatorname{pf}_{X,R}\) to \(\operatorname{pf}_{X,R'}\), compatible
with squaring and with the automorphic comparison.
\item[\textup{(M7)}] A filtered coalgebra morphism
\(C_{X,R}\to C_{X,R'}\) preserving the bar differential, augmentation,
Koszul source datum, and the comparison
\(\Theta_{\mathrm{Kos}}\).
\item[\textup{(M8)}] A morphism of primitive-recognition tables
preserving representatives, relation rows, parity blocks, product and
coproduct matrices, Hopf-pairing matrices, radical and quotient
splittings, comparison maps, no-extra rows, PBW filtrations, and
Mittag--Leffler transition rows.  PBW-filtration preservation is
Proposition~\ref{prop:transition-pbw-filtration-preservation-criterion};
parity-decomposition preservation is
Proposition~\ref{prop:transition-parity-decomposition-preservation-criterion};
primitive-space \(R^1\varprojlim\)-vanishing is
Proposition~\ref{prop:primitive-space-lim1-vanishing-criterion};
pairing-kernel \(R^1\varprojlim\)-vanishing is
Proposition~\ref{prop:pairing-kernel-lim1-vanishing-criterion}; and PBW
associated-graded \(R^1\varprojlim\)-vanishing is
Proposition~\ref{prop:pbw-associated-graded-lim1-vanishing-criterion}.
\end{enumerate}
All diagrams expressing these compatibilities commute in the relevant
derived, Borel--Moore, filtered, or completed category.  Composition is
componentwise composition.  Thus the identity and associativity
conditions for \(\mathsf{DI}^{\mathrm{fin}}_X\) are precisely the
zero-defect rows in \(\mathrm{cofinal\_subsystems}\),
\(\mathrm{stage\_morphisms}\), and
\(\mathrm{composition\_laws}\).  An isomorphism of finite stages is a
transition morphism whose components are equivalences of stacks where applicable,
quasi-isomorphisms on chain and coalgebra objects, isomorphisms on
orientation and Pfaffian lines, isomorphisms of Hall bialgebra and
recognition matrices on the retained windows, and equality of the
orientation characters.
\end{definition}

\begin{definition}[Dirac--Igusa pro-object and pro-isomorphism]
\label{def:dirac-igusa-pro-isomorphism}
A Dirac--Igusa pro-object is an object of
\[
\operatorname{Pro}_{\mathcal R_{\mathrm{HN}}}(\mathsf{DI}^{\mathrm{fin}}_X),
\]
represented by a functor
\[
\mathfrak D_X^{\mathrm{DI}}:
\mathcal R_{\mathrm{HN}}^{\mathrm{op}}\longrightarrow
\mathsf{DI}^{\mathrm{fin}}_X
\]
from a cofinal HN subsystem to the finite-stage category of
Definition~\ref{def:finite-stage-dirac-igusa-morphism}.  The symbol
\[
\varprojlim_R\mathfrak D^{\mathrm{DI}}_{X,R}
\]
means this pro-object.  It is not a componentwise inverse limit unless
the corresponding \(\mathrm{ml\_exactness}\) rows give strict
Mittag--Leffler towers and \(R^1\varprojlim=0\).  When a component is
realized as a completed module or complex, its topology is the
projection topology: a fundamental system of neighbourhoods of zero is
given by the kernels of the projections to finite retained heights.
For Picard-groupoid components, this means compatible line bundles and
trivialisations on a cofinal subsystem.  For completed Lie, Hall, and
Koszul components, it means the completed filtered object determined by
the finite quotients.

A pro-isomorphism between two Dirac--Igusa pro-objects is an
isomorphism in this pro-category: after passing to a common cofinal
subsystem, it is represented by finite-stage isomorphisms commuting
with all transition morphisms and with two-sided inverse defects zero.
Equality of scalar Pfaffian products, equality of squared determinants,
or equality of signed denominator exponents is not a pro-isomorphism
unless the morphism data \textup{(M1)}--\textup{(M8)} are supplied.
\end{definition}

\begin{proposition}[Universal property and cofinal presentation invariance]
\label{prop:dirac-igusa-pro-object-universal-property}
Let \(\mathsf C_X=\mathsf{DI}^{\mathrm{fin}}_X\), and suppose the
finite morphism rows of Definition~\ref{def:finite-stage-dirac-igusa-morphism}
give an inverse system
\[
\mathfrak D=(\mathfrak D_R,\rho_{RR'}):
\mathcal R_{\mathrm{HN}}^{\mathrm{op}}\longrightarrow \mathsf C_X .
\]
For a finite-stage object \(T\in\mathsf C_X\), regarded as a constant
pro-object, the Dirac--Igusa pro-object satisfies
\[
\Hom_{\operatorname{Pro}(\mathsf C_X)}
 (T,\mathfrak D)
\cong
\varprojlim_{R\in\mathcal R_{\mathrm{HN}}}
\Hom_{\mathsf C_X}(T,\mathfrak D_R).
\]
For another Dirac--Igusa pro-object
\(\mathfrak E=(\mathfrak E_S)_{S\in\mathcal S}\),
\[
\Hom_{\operatorname{Pro}(\mathsf C_X)}
 (\mathfrak E,\mathfrak D)
\cong
\varprojlim_{R\in\mathcal R_{\mathrm{HN}}}
\varinjlim_{S\in\mathcal S}
\Hom_{\mathsf C_X}(\mathfrak E_S,\mathfrak D_R).
\]
These two formulas are the universal property of
\(\mathfrak D_X^{\mathrm{DI}}\).  They are statements in the
pro-category; they do not assert componentwise inverse limits of
complexes, line bundles, Hall objects, or Lie superalgebras.

If
\(i:\mathcal R'_{\mathrm{HN}}\hookrightarrow\mathcal R_{\mathrm{HN}}\)
is a cofinal subsystem, the restriction \(i^\ast\mathfrak D\) defines a
canonically isomorphic pro-object.  More generally, two finite
presentations of the Dirac--Igusa datum over cofinal HN subsystems
define the same pro-object exactly when, after passage to a common
cofinal refinement, the finite-stage isomorphisms commute with
\textup{(M1)}--\textup{(M8)} and have zero two-sided inverse defects.
This is the independence of finite presentation choices used in the
notation
\(\mathfrak D_X^{\mathrm{DI}}=\varprojlim_R\mathfrak D_{X,R}^{\mathrm{DI}}\).
\end{proposition}

\begin{proof}
The first displayed identity is the special case of the standard
pro-category formula with the source represented by a constant system.
The second is the defining Hom formula in
\(\operatorname{Pro}(\mathsf C_X)\):
morphisms \((\mathfrak E_S)_S\to(\mathfrak D_R)_R\) are compatible
families of finite-stage morphisms after allowing a cofinal refinement
of the source index.  Cofinal restriction leaves these Hom functors
unchanged, hence gives the canonical pro-isomorphism by Yoneda.  The
finite morphism tables supply the only non-formal input: they ensure
that the stage maps exist, compose strictly, and have zero inverse
defects on the common refinement.  The Mittag--Leffler rows enter only
when one replaces a pro-object by an actual completed component; they
are not part of the pro-object universal property itself.
\end{proof}

\begin{definition}[Moduli-cohomology Mittag--Leffler datum]
\label{def:moduli-cohomology-ml-datum}
Fix a cofinal sequence in the retained HN system.  For each retained
object, extension, mixed, wrapped, and two-step flag stack
\(\mathfrak M_{R,a}\), for each coefficient system
\(\mathcal K_{R,a}\) used in the finite Hall or reduced
vanishing-cycle construction, and for each cohomological degree \(q\),
put
\[
C^q_{R,a,\mathcal K}:=
H^q_c(\mathfrak M_{R,a},\mathcal K_{R,a})
\]
or the corresponding Borel--Moore cohomology group, according to the
six-functor convention of the finite Hall operation under discussion.
A moduli-cohomology Mittag--Leffler datum consists of:
\begin{enumerate}
\item finite-rank rows for every \(C^q_{R,a,\mathcal K}\), with the
stack kind, coefficient system, support block, and topology recorded;
\item transition maps
\[
u^q_{RR',a,\mathcal K}:C^q_{R,a,\mathcal K}\longrightarrow
C^q_{R',a,\mathcal K}
\]
induced by the proper pushforward, closed restriction, or specified
six-functor operation along the retained moduli transition;
\item identity and composition rows for the maps \(u^q_{RR'}\);
\item for every base stage \(R_0\), stack kind \(a\), coefficient
\(\mathcal K\), and degree \(q\), an image-stabilization witness:
there is \(R_1\ge R_0\) such that for all \(R\ge R_1\),
\[
\operatorname{im}\bigl(C^q_{R,a,\mathcal K}\to
C^q_{R_0,a,\mathcal K}\bigr)
=
\operatorname{im}\bigl(C^q_{R_1,a,\mathcal K}\to
C^q_{R_0,a,\mathcal K}\bigr),
\]
with the stabilized image written in a finite basis;
\item the defect rows
\[
\mathrm{strict}=1,\qquad \mathrm{ML}=1,\qquad
\operatorname{rank}R^1\varprojlim_R C^q_{R,a,\mathcal K}=0
\]
for every retained tower.
\end{enumerate}
Finite type of the stacks, properness or closedness of transition
morphisms, bounded HN types, Euler characteristics, protected traces,
and target-root windows are not a moduli-cohomology
Mittag--Leffler datum.
\end{definition}

\begin{proposition}[Vanishing of \(R^1\varprojlim\) on moduli
cohomology; \textit{proved}]
\label{prop:moduli-cohomology-lim1-vanishing-criterion}
For every tower \(C^\bullet_{R,a,\mathcal K}\) carrying the datum of
Definition~\ref{def:moduli-cohomology-ml-datum}, the inverse system is
Mittag--Leffler in each cohomological degree.  Consequently
\[
R^1\varprojlim_R C^q_{R,a,\mathcal K}=0
\]
for every retained stack kind \(a\), coefficient system
\(\mathcal K\), support block, and degree \(q\).  Without the rows of
Definition~\ref{def:moduli-cohomology-ml-datum}, the passage from
finite retained moduli cohomology to componentwise inverse-limit
cohomology is not established.
\end{proposition}

\begin{proof}
The identity and composition rows make
\(\{C^q_{R,a,\mathcal K},u^q_{RR'}\}\) an inverse system of finite-rank
modules.  The image-stabilization row is exactly the
Mittag--Leffler condition: for every fixed base stage \(R_0\), the
images inside \(C^q_{R_0,a,\mathcal K}\) are eventually constant.
The standard derived-inverse-limit criterion for inverse systems of
modules then gives \(R^1\varprojlim=0\) on each such tower.  The
vanishing is degreewise and coefficientwise; it cannot be inferred
from finite type, from a scalar Euler characteristic, or from the
existence of proper or closed transition morphisms without the
stabilized-image rows.
\end{proof}

The packet
\[
\texttt{certificates/moduli/moduli\_cohomology\_lim1\_vanishing}
\]
records the current state of this datum.  Its verifier returns
\(\texttt{MODULI\_COHOMOLOGY\_LIM1\_VANISHING\_OBSTRUCTION\_VERIFIED}\):
the cofinal HN sequence, retained moduli stack inputs, transition
geometry input, coefficient systems, finite cohomology groups,
cohomology transition maps, functoriality rows, image subspace rows,
image-stabilization witnesses, strict Mittag--Leffler status rows,
zero \(R^1\varprojlim\) rows, and degreewise coverage rows are not yet
supplied.

\paragraph{Finite proof tables for morphisms and pro-isomorphisms.}
The finite-stage functor is supplied by the tables
\[
\mathrm{cofinal\_subsystems},\quad
\mathrm{stage\_objects},\quad
\mathrm{stage\_morphisms},\quad
\mathrm{component\_maps},\quad
\mathrm{composition\_laws},\quad
\mathrm{ml\_exactness},\quad
\mathrm{pro\_isomorphisms},\quad
\mathrm{scalar\_firewall}.
\]
The table \(\mathrm{stage\_objects}\) contains one row for each of
\[
\mathcal A^{E_3},\ F^{\mathrm{hyb}},\ \Gamma,\ \Pi,\ \Phi,\ o,\ H,\
P^\Pi,\ C,\ \Theta_{\mathrm{Kos}},\ \mathscr L_{\mathrm{Pf}},\
\operatorname{pf},\ \varepsilon_o,\ \mathrm{Rec}
\]
at each retained height.  The table \(\mathrm{component\_maps}\)
contains rows for \(\textup{(M1)}\)--\(\textup{(M8)}\), with zero
compatibility, kernel, and cokernel ranks in the corresponding
derived, Borel--Moore, filtered, or completed category.  The tables
\(\mathrm{cofinal\_subsystems}\), \(\mathrm{stage\_morphisms}\), and
\(\mathrm{composition\_laws}\) record directed cofinality, strict
transition morphisms, identity laws, and
\(\rho_{R'R''}\rho_{RR'}=\rho_{RR''}\).  The table
\(\mathrm{ml\_exactness}\) records \(R^1\varprojlim=0\) for the
component towers.  The table \(\mathrm{pro\_isomorphisms}\) records
finite-stage isomorphism families over a common cofinal subsystem,
with commuting squares and two-sided inverse defects zero.  The table
\(\mathrm{scalar\_firewall}\) excludes scalar Pfaffian products,
denominator products, squared determinants, signed exponent tables,
and protected traces as morphism or pro-isomorphism data.
The current packet
\(\texttt{certificates/morphisms/k3e\_dirac\_igusa\_pro\_object}\)
verifies only the obstruction ledger
\(\texttt{MORPHISM\_OBSTRUCTION\_LEDGER\_VERIFIED}\), with
\(\texttt{pro\_object\_certification=false}\).  It records the missing
cofinal subsystem, stage-object, transition, component-map,
composition-law, Mittag--Leffler, and pro-isomorphism rows; it does not
construct the functor
\(\mathcal R_{\mathrm{HN}}^{\mathrm{op}}\to
\mathsf{DI}^{\mathrm{fin}}_X\).

\paragraph{Finite proof tables for the compact \(E_3\)-source.}
The inputs \([\textup{$E_3$-HolFA}]\) and
\([\textup{K3$\to$E-sp}]\) are supplied at finite height by
\[
\mathrm{formal\_target\_charts},\quad
\mathrm{field\_complexes},\quad
\mathrm{e3\_operations},\quad
\mathrm{bv\_qme},\quad
\mathrm{anomaly\_framing},\quad
\mathrm{compact\_support},\quad
\mathrm{factorization\_descent},\quad
\mathrm{k3\_to\_e\_specialization},\quad
\mathrm{transitions},\quad
\mathrm{scalar\_firewall}.
\]
These rows record the finite formal target charts, field complexes,
holomorphic \(E_3\)-operations, BV quantum master equation, anomaly
cancellation, framing and formality, compact retained support,
prefactorization descent, \(K3\)-to-\(E\) chain specialization,
strict transitions, and \(R^1\varprojlim=0\).  The scalar firewall
excludes the automorphic section, BKM denominator, target current
envelope, protected trace, and hybrid-carrier-only data as entries of
\(\mathcal A^{E_3}_{X,R}\).
The current finite \(E_3\)-source packet records these rows only by the
obstruction ledger
\(\texttt{E3\_SOURCE\_OBSTRUCTION\_LEDGER\_VERIFIED}\), with
\(\texttt{e3\_source\_certification=false}\).  It is an absence ledger,
not a construction of \(\mathcal A^{E_3}_{X,R}\).

\subsubsection*{Charge window}

\begin{definition}[Finite target test window]
\label{def:finite-target-test-window}
Fix an HN height \(R\).  Let
\(\Gamma_R^{HN}\subset F_{\sigma,S}(R)\subset\Gamma_X^{\mathrm{form}}\)
be the retained finite HN charge set.  For \(c\in\Gamma_R^{HN}\), set
\[
\mathcal T_R(c)
=\Bigl\{\sum_{i<j}B(c_i,c_j)\Bigm|c=\sum_i c_i,\
c_i\in F_{\sigma,S}(R),\ \sum_i h_S(c_i)\le R\Bigr\}
\subset\Gamma_{\mathrm{gram}},
\]
and
\[
\widehat\Gamma_R
:=\{(c,T)\mid c\in\Gamma_R^{HN},\ T\in\mathcal T_R(c)\}
\subset\widehat\Gamma_X .
\]
The finite target test window is
\[
\Gamma_R^{\mathrm{test}}
:=\overline\Pi_X(\widehat\Gamma_R)
=\{\Pi_X(c)+T\mid c\in\Gamma_R^{HN},\ T\in\mathcal T_R(c)\}
\subset\Gamma_{\mathrm{gram}} .
\]
It is finite once \(\Gamma_R^{HN}\) and the translate sets
\(\mathcal T_R(c)\) are finite.  For \(R\le R'\), compatible inclusions
of HN charge sets and translate sets induce
\(\Gamma_R^{\mathrm{test}}\subset\Gamma_{R'}^{\mathrm{test}}\).
This is a target window in \(\Gamma_{\mathrm{gram}}\), not a compact
source window and not the (O1) orientation datum.
\end{definition}

\textit{Hypotheses.} Stability datum $(\sigma,S)$, finite HN height
\(R\), and a retained HN charge set \(\Gamma_R^{HN}\).

\textit{Finite datum.} Finite saturated downward window
$F_{\sigma,S}(R)\subset\Gamma_X^{\mathrm{HN}}(R)$, the dyonic
formal Mukai/effective-charge support, the Gram map
$\Pi_X:\Gamma_X^{\mathrm{form}}\to\Gamma_{\mathrm{gram}}$, the cocycle
$B$, the finite normal-ordered lift \(\widehat\Gamma_R\), and the target
test window \(\Gamma_R^{\mathrm{test}}\subset\Gamma_{\mathrm{gram}}\) of
Definition~\ref{def:finite-target-test-window}.

\textit{Construction.}
\begin{enumerate}
\item Form the Liu charge $L_F(n)=Z_S(Rp_{S,*}(F\otimes p_E^*H^n))$
\cite{LiuProductStability} for $\sigma_S$ a rational K3 stability
condition and a degree-one line bundle $H$ on $E$.  Restrict to
charges with bounded $h_S$-height $\le R$; the resulting set is
$F_{\sigma,S}(R)$.
\item Compute the Gram map and cocycle $B(c,c')$ on
$F_{\sigma,S}(R)\times F_{\sigma,S}(R)$ via the Mukai dictionary of
\S\ref{subsec:d6-d2-d0-dictionary-appendixC}; check the symmetric
bilinear $2$-cocycle condition.
\item Form the finite normal-ordered lift set
\(\widehat\Gamma_R\) and its image
\(\Gamma_R^{\mathrm{test}}=\overline\Pi_X(\widehat\Gamma_R)\).
\end{enumerate}
At table level this input would be supplied by
\[
\mathrm{finite\_test\_window\_definition},\quad
\mathrm{active\_support},\quad
\mathrm{window\_saturation},\quad
\mathrm{liu\_hn\_window},\quad
\mathrm{gram\_map},\quad
\mathrm{cocycle\_identities},\quad
\mathrm{normal\_ordered\_lifts},\quad
\mathrm{target\_windows},\quad
\mathrm{primitive\_lifts},\quad
\mathrm{transitions},\quad
\mathrm{scalar\_firewall}.
\]
When supplied, these rows verify the active \(\phi_{0,1}\)-support,
downward saturation, Liu height cutoff, Mukai--Gram map, cocycle
identities, normal-ordered lift orbit, target test-window image,
primitive formal lift, strict transition, and \(R^1\varprojlim=0\).  The
finite-test-window definition row is separated in
\[
\texttt{certificates/charge/finite\_test\_window}.
\]
Its verifier returns
\(\texttt{FINITE\_TEST\_WINDOW\_DEFINITION\_VERIFIED}\), with populated
test-window certification and relation-closed target exhaustion both
false.  It supplies the formula for \(\Gamma_R^{\mathrm{test}}\), not the
HN charge rows, normal-ordered translate rows, image-membership rows,
finiteness witnesses, transition inclusions, or the assertion that every
relation-closed target degree appears at finite \(R\).  The scalar
firewall rejects target-presentation rows, signed exponent lists,
Pfaffian products, Hilbert-scheme scalar shadows, orientation data, and
source-rank-only data as charge-window evidence.

\begin{definition}[Relation-closed target finite-\(R\) exhaustion datum]
\label{def:relation-closed-target-finite-r-exhaustion}
Let \(\mathcal D_{\mathrm{rel}}\subset\Gamma_{\mathrm{gram}}\) be the
target-degree set on which relation closure is asserted.  In the present
target audit this includes the finite \(W_{\mathrm{rel}}\)-cascade
recorded by
\(\texttt{certificates/targets/delta5\_gn\_kac/wrel\_degree\_closure}\)
and the symbolic tail \(A_N\Rightarrow B_N\Rightarrow A_{N+2}\) for even
\(N\ge14\) recorded by
\(\texttt{certificates/targets/delta5\_gn\_kac/wrel\_tail\_staircase}\).
A finite-\(R\) exhaustion datum for \(\mathcal D_{\mathrm{rel}}\)
consists of:
\begin{enumerate}
\item a target-scope table listing the finite prefix and the symbolic
tail families;
\item for every finite-prefix degree \(\gamma\), a finite height
\(R(\gamma)\), a charge \(c_\gamma\in\Gamma_{R(\gamma)}^{HN}\), and a
translate \(T_\gamma\in\mathcal T_{R(\gamma)}(c_\gamma)\);
\item for each symbolic tail family, formulae \(R_N\), \(c_N\), and
\(T_N\) valid for every allowed parameter \(N\);
\item the image identities
\[
\Pi_X(c_\gamma)+T_\gamma=\gamma,\qquad
\Pi_X(c_N)+T_N=\gamma_N,
\]
with zero defect in \(\Gamma_{\mathrm{gram}}\);
\item compatibility rows for \(R\le R'\), and a zero coverage-defect row
for the finite prefix and every symbolic family.
\end{enumerate}
Target relation closure, the infinite tail staircase, active-root
cofinality, signed multiplicities, and the formula for
\(\Gamma_R^{\mathrm{test}}\) are not finite-\(R\) exhaustion data.
\end{definition}

\begin{proposition}[Finite-\(R\) realization of relation-closed target
degrees; \textit{proved}]
\label{prop:relation-closed-target-finite-r-exhaustion-criterion}
If \(\mathcal D_{\mathrm{rel}}\) carries the datum of
Definition~\ref{def:relation-closed-target-finite-r-exhaustion}, then
for every \(\gamma\in\mathcal D_{\mathrm{rel}}\) there exists a finite
HN height \(R\) such that
\[
\gamma\in\Gamma_R^{\mathrm{test}}.
\]
\end{proposition}

\begin{proof}
For a finite-prefix degree \(\gamma\), the datum gives
\(R(\gamma)\), \(c_\gamma\in\Gamma_{R(\gamma)}^{HN}\), and
\(T_\gamma\in\mathcal T_{R(\gamma)}(c_\gamma)\) with
\(\Pi_X(c_\gamma)+T_\gamma=\gamma\).  By
Definition~\ref{def:finite-target-test-window},
\(\gamma\in\Gamma_{R(\gamma)}^{\mathrm{test}}\).  The same argument
applies to a symbolic tail degree \(\gamma_N\) using \(R_N,c_N,T_N\).
The zero coverage-defect row says that these two cases exhaust
\(\mathcal D_{\mathrm{rel}}\).  No target-only table supplies the
charges \(c\) or translates \(T\), so target relation closure alone does
not prove the proposition.
\end{proof}

The packet
\[
\texttt{certificates/charge/relation\_closed\_target\_finite\_R\_exhaustion}
\]
verifies
\(\texttt{RELATION\_CLOSED\_TARGET\_FINITE\_R\_EXHAUSTION\_OBSTRUCTION\_VERIFIED}\):
the finite prefix degree-to-\(R\) rows, symbolic tail \(R_N\)-formulae,
HN charge witnesses, normal-ordered translate witnesses, image-equality
rows, finite-\(R\) bounds, transition compatibility rows, compact source
compatibility rows, and zero coverage-defect row are not supplied.

\textit{Obstruction classes.}
$[\textup{K3$\times$E-fin}]$ on the Liu charge bound; $B$
cocycle property.

\subsubsection*{Finite Hall stage and atlas}

\textit{Hypotheses.} The charge window, plus the cosection-twisted reduced
d-critical heart of \S\ref{subsec:cosection-twisted-orientation}.

\textit{Finite datum.} The finite Hall stage
$\mathcal C_{\sigma,S,\le R}$, finite-type derived Artin enhancements
$\mathfrak M_{\widehat c,R},
\mathfrak E_{\widehat c,\widehat c',R},
\mathfrak F^{(2)}_{\widehat c,\widehat c',\widehat c'',R}$ of object,
extension, and two-step flag stacks; reduced cosections; reduced
vanishing cycles $\Phi^{\mathrm{red}}$; orientation lines
$o^{\mathrm{red}}$.
At table level this datum is supplied by
\(\mathrm{hn\_type\_bounds}\), \(\mathrm{semistable\_substacks}\),
\(\mathrm{derived\_enhancements}\), \(\mathrm{universal\_complexes}\),
\(\mathrm{scalar\_rigidifications}\), \(\mathrm{stratifications}\),
\(\mathrm{extension\_flag\_stacks}\), \(\mathrm{cosection\_atlas}\),
\(\mathrm{transitions}\), and \(\mathrm{scalar\_firewall}\).
The current packet verifies only
\(\texttt{MODULI\_OBSTRUCTION\_LEDGER\_VERIFIED}\), with
\(\texttt{moduli\_certification=false}\); no finite Hall stage is thereby
constructed.

\textit{Construction.}
\begin{enumerate}
\item Restrict the chosen stability heart to charges in $F_{\sigma,S}(R)$;
form the finite-type derived enhancements with PTVV $(-1)$-shifted
symplectic structures \cite{PTVV2013} on the truncations; descend to
Joyce d-critical structures \cite{BBDJS2015}.
\item Compute the K3 semiregularity cosection
$\operatorname{CS}_{\mathcal C}$ on each retained charge stratum.
Surjectivity on reduced strata and filtered Hall additivity on
extension stacks are separate rows.
\item Choose Joyce--Upmeier orientation lines \cite{JoyceUpmeier} with
chosen square roots, multiplicative under reduced extension
correspondences.
\end{enumerate}

\textit{Obstruction classes.}
$[\textup{K3$\times$E-atlas}]$; finite closed inertia stratifications
after the wrapped quotient steps, together with closure under
extensions, HN factors, and duals where needed.

\subsubsection*{Hybrid wrapped factorization datum}

\textit{Hypotheses.} The finite Hall stage.

\textit{Finite datum.} The geometric elliptic-degree map
\(b_R^{\mathrm{geom}}\) of
Definition~\ref{def:geometric-elliptic-degree-map}; the hybrid
factorization carrier
$\mathcal F_{X,\sigma,S,\le R}^{\mathrm{hyb}}$ of
\S\ref{subsec:hybrid-ran-carrier}: the finite-stage prestack
\(\operatorname{Ran}^{\mathrm{hyb}}_{R,\mathrm{pre}}(E)\);
rigidified wrapped prequotients
$\mathcal M_{\eta,R}^{\mathrm{wr,rig}}$ with anchors
$\lambda_{\eta,R}$; closed-configuration local sectors
$\mathcal F_R^{\mathrm{loc}}$; wrapped sectors
$\mathcal G_R^{\mathrm{wr}}$; the eight-word two-step binary flag
atlas; the higher-coloured residual
$\mathfrak o^{\mathrm{col}}_R$; the quotient-after-correspondence
pseudofunctor inducing $Q_{E,R}$ and comparison isomorphisms
$\theta^Q_{\mu,R}$; protected integration
$I_R^{\mathrm{prot}}$.
At table level this datum would be supplied by
\[
\mathrm{hybrid\_ran\_prestack\_definition},\quad
\mathrm{local\_stratum\_b0\_prestack\_definition},\quad
\mathrm{geometric\_elliptic\_degree\_map},\quad
\mathrm{geometric\_elliptic\_degree\_additivity},\quad
\mathrm{geometric\_borcherds\_degree\_separation},\quad
\mathrm{positive\_elliptic\_degree\_projection},\quad
\mathrm{hybrid\_base},\quad
\mathrm{local\_configurations},\quad
\mathrm{wrapped\_stratum\_bpositive\_prestack\_definition},\quad
\mathrm{local\_local\_correspondence\_definition},\quad
\mathrm{mixed\_local\_wrapped\_correspondence\_definition},\quad
\mathrm{wrapped\_wrapped\_correspondence\_definition},\quad
\mathrm{two\_sided\_mixed\_correspondence\_definition},\quad
\mathrm{source\_target\_anchor\_memory\_definition},\quad
\mathrm{determinant\_anchor\_construction},\quad
\mathrm{determinant\_anchor\_translation\_weight},\quad
\mathrm{determinant\_anchor\_chi\_zero\_degeneracy},\quad
\mathrm{determinant\_anchor\_extra\_anchor\_data},\quad
\mathrm{wrapped\_prequotients},\quad
\mathrm{correspondences},\quad
\mathrm{flag\_atlas},\quad
\mathrm{higher\_coloured},\quad
\mathrm{quotient\_pseudofunctor},\quad
\mathrm{protected\_integration},\quad
\mathrm{transitions},\quad
\mathrm{scalar\_firewall}.
\]
When supplied, these rows record the hybrid prestack definition, the
\(b=0\) local subprestack, the \(b>0\) wrapped subprestack,
elliptic-degree map, additivity criterion,
Borcherds-coordinate
separation, positive-degree projection criterion, and split, local closed
configuration stacks, wrapped anchors, local/mixed/wrapped
correspondences, source/target anchor memory before quotient, all eight
words in \(\{\mathrm L,\mathrm W\}^3\), determinant anchors,
determinant-anchor weights,
hypersurface \(\chi=0\) determinant-anchor degeneracy,
hypersurface \(\chi=0\) Jordan--Hoelder extra-anchor data,
higher-coloured coherences, quotient-after-correspondence, protected
integration, and strict HN transitions.
The degree-map definition is isolated in
\[
\texttt{certificates/hybrid/geometric\_elliptic\_degree\_map}.
\]
Its verifier returns
\(\texttt{GEOMETRIC\_ELLIPTIC\_DEGREE\_MAP\_OBSTRUCTION\_VERIFIED}\):
\(b_R^{\mathrm{geom}}\) is the coefficient of \([E]\) in
\((p_E)_*\operatorname{cyc}_1(\gamma)\), not
\(\operatorname{pr}_3\overline\Pi_X\).  The retained HN class rows,
one-cycle rows, pushforward rows, local/wrapped split rows, and
transition rows are not supplied.  The additivity criterion is isolated
in
\[
\texttt{certificates/hybrid/geometric\_elliptic\_degree\_additivity}.
\]
Its verifier returns
\(\texttt{GEOMETRIC\_ELLIPTIC\_DEGREE\_ADDITIVITY\_OBSTRUCTION\_VERIFIED}\):
Definition~\ref{def:geometric-elliptic-degree-additivity-datum} and
Proposition~\ref{prop:geometric-elliptic-degree-additivity-criterion}
give the finite-stage criterion, while extension-pair, one-cycle
additivity, pushforward additivity, integer degree additivity, and
transition rows are not supplied.  The retained-extension-closure packet
does not by itself prove additivity of \(b_R^{\mathrm{geom}}\).
The geometric/Borcherds separation packet is
\[
\texttt{certificates/hybrid/geometric\_borcherds\_degree\_separation}.
\]
Its verifier returns
\(\texttt{GEOMETRIC\_BORCHERDS\_DEGREE\_SEPARATION\_OBSTRUCTION\_VERIFIED}\):
Proposition~\ref{prop:geometric-borcherds-degree-separation} separates
the support degree \(b_R^{\mathrm{geom}}\) from the trace coordinate
\(m_R^{\mathrm{Bch}}=\operatorname{pr}_3\overline\Pi_X\); branch
comparison rows are not supplied.
The positive elliptic-degree projection packet is
\[
\texttt{certificates/hybrid/positive\_elliptic\_degree\_projection}.
\]
Its verifier returns
\(\texttt{POSITIVE\_ELLIPTIC\_DEGREE\_PROJECTION\_VERIFIED}\):
Proposition~\ref{prop:positive-elliptic-degree-projects} proves that an
effective one-cycle with positive \(E\)-degree projects onto all of
\(E\), while retained object support-realization rows are not supplied.
The hybrid Ran prestack definition packet is
\[
\texttt{certificates/hybrid/hybrid\_ran\_prestack\_definition}.
\]
Its verifier returns
\(\texttt{HYBRID\_RAN\_PRESTACK\_DEFINITION\_VERIFIED}\):
Definition~\ref{def:hybrid-ran-prestack} defines the finite-stage
carrier as a prestack functor; local component, wrapped component,
descent, stackification, and transition rows are not supplied.
The local \(b=0\) stratum packet is
\[
\texttt{certificates/hybrid/local\_stratum\_b0\_prestack\_definition}.
\]
Its verifier returns
\(\texttt{LOCAL\_STRATUM\_B0\_PRESTACK\_DEFINITION\_VERIFIED}\):
Definition~\ref{def:hybrid-local-stratum-prestack} defines the full
\(J=\varnothing\) subprestack with colours in
\(\Gamma_R^{\mathrm{loc}}\); closed-configuration, local sheaf,
descent, and transition rows are not supplied.
The wrapped \(b>0\) stratum packet is
\[
\texttt{certificates/hybrid/wrapped\_stratum\_bpositive\_prestack\_definition}.
\]
Its verifier returns
\(\texttt{WRAPPED\_STRATUM\_BPOSITIVE\_PRESTACK\_DEFINITION\_VERIFIED}\):
Definition~\ref{def:hybrid-wrapped-stratum-prestack} defines the full
\(I=\varnothing\) subprestack with colours in
\(\Gamma_R^{\mathrm{wr}}\); the \(E\)-equivariant prequotient
definition is separated in the next packet, while anchor,
support-realization, descent, and transition rows are not supplied.
The \(E\)-equivariant wrapped prequotient packet is
\[
\texttt{certificates/hybrid/equivariant\_wrapped\_prequotient\_definition}.
\]
Its verifier returns
\(\texttt{E\_EQUIVARIANT\_WRAPPED\_PREQUOTIENT\_DEFINITION\_VERIFIED}\):
Definition~\ref{def:e-equivariant-wrapped-prequotient-stack} defines
\(\mathcal M_{\eta,R}^{\mathrm{wr,rig}}\) before the reduced
\(E\)-quotient and the translation action with origin \(0_E\);
finite type is separated in the next packet, while properness, anchor
population, anchor losslessness, quotient descent, aggregate
hybrid-carrier population, and transition rows are not supplied.
The wrapped-prequotient finite-type packet is
\[
\texttt{certificates/hybrid/wrapped\_prequotient\_finite\_type}.
\]
Its verifier returns
\(\texttt{WRAPPED\_PREQUOTIENT\_FINITE\_TYPE\_VERIFIED}\):
Proposition~\ref{prop:wrapped-prequotient-finite-type} proves finite
type for \(\mathcal M_{\eta,R}^{\mathrm{wr,rig}}\) at fixed \(R\) from
the retained finite class set, finite-type semistable substacks, and
scalar rigidification; properness, anchor population, quotient
descent, aggregate hybrid-carrier population, and transition rows are
not supplied.
The local/local correspondence packet is
\[
\texttt{certificates/hybrid/local\_local\_correspondence\_definition}.
\]
Its verifier returns
\(\texttt{LOCAL\_LOCAL\_CORRESPONDENCE\_DEFINITION\_VERIFIED}\):
Definition~\ref{def:local-local-correspondence-datum} defines the
ordered \(LL\) projection-finite diagram; extension stack,
source/target map, descent, properness, Thom--Sebastiani, collision,
and transition rows are not supplied by that definition packet.
The local/local target-properness packet is
\[
\texttt{certificates/hybrid/local\_local\_extension\_properness}.
\]
Its verifier returns
\(\texttt{LOCAL\_LOCAL\_TARGET\_PROPERNESS\_VERIFIED}\):
Proposition~\ref{prop:local-local-target-properness} realizes the
\(LL\) extension stack as a closed relative Quot locus and proves
properness of the target map \(q^{\mathrm{LL}}\).  It does not assert
properness of the source map \(p^{\mathrm{LL}}\), and descent,
collision compatibility, Thom--Sebastiani transport, transition rows,
and aggregate hybrid-carrier population are not supplied.
The mixed local/wrapped correspondence packet is
\[
\texttt{certificates/hybrid/mixed\_local\_wrapped\_correspondence\_definition}.
\]
Its verifier returns
\(\texttt{MIXED\_LOCAL\_WRAPPED\_CORRESPONDENCE\_DEFINITION\_VERIFIED}\):
Definition~\ref{def:mixed-local-wrapped-correspondence-datum} defines
the one-sided \(LW\) and \(WL\) unreduced mixed diagrams; extension
stack, source/target map, populated anchor-memory, admissibility,
base-change, projection-formula, Thom--Sebastiani, and transition rows
are not supplied by that definition packet.
The mixed local/wrapped target-admissibility packet is
\[
\texttt{certificates/hybrid/mixed\_local\_wrapped\_extension\_admissibility}.
\]
Its verifier returns
\(\texttt{MIXED\_LOCAL\_WRAPPED\_ADMISSIBILITY\_VERIFIED}\):
Proposition~\ref{prop:mixed-local-wrapped-target-admissibility}
realizes the \(LW\) and \(WL\) extension stacks as closed relative
Quot loci over the wrapped target prequotient and proves properness of
the target maps \(q^{\mathrm{LW}}\) and \(q^{\mathrm{WL}}\).  It does
not assert properness of \(p^{\mathrm{LW}}\) or \(p^{\mathrm{WL}}\),
and compact-support pushforward, populated anchor-memory rows, base
change, projection formula, Thom--Sebastiani transport, quotient
descent, transition rows, and aggregate hybrid-carrier population are
not supplied.
The wrapped/wrapped correspondence packet is
\[
\texttt{certificates/hybrid/wrapped\_wrapped\_correspondence\_definition}.
\]
Its verifier returns
\(\texttt{WRAPPED\_WRAPPED\_CORRESPONDENCE\_DEFINITION\_VERIFIED}\):
Definition~\ref{def:wrapped-wrapped-correspondence-datum} defines the
ordered \(WW\) unreduced diagram; extension stack, source/target map,
populated anchor-memory, admissibility, Thom--Sebastiani,
symmetric-descent, and transition rows are not supplied by that
definition packet.
The wrapped/wrapped target-admissibility packet is
\[
\texttt{certificates/hybrid/wrapped\_wrapped\_extension\_admissibility}.
\]
Its verifier returns
\(\texttt{WRAPPED\_WRAPPED\_ADMISSIBILITY\_VERIFIED}\):
Proposition~\ref{prop:wrapped-wrapped-target-admissibility} realizes
the \(WW\) extension stack as a closed relative Quot locus over the
wrapped target prequotient and proves properness of the target map
\(q^{\mathrm{WW}}\).  It does not assert properness of
\(p^{\mathrm{WW}}\), and compact-support pushforward, populated
anchor-memory rows, symmetric descent, Thom--Sebastiani transport,
quotient descent, transition rows, and aggregate hybrid-carrier
population are not supplied.
The compact-support exceptional pushforward model packet is
\[
\texttt{certificates/hybrid/compact\_support\_exceptional\_pushforward\_model}.
\]
Its verifier returns
\(\texttt{COMPACT\_SUPPORT\_EXCEPTIONAL\_PUSHFORWARD\_MODEL\_DEFINED}\):
Definition~\ref{def:compact-support-exceptional-pushforward-model}
defines a chosen compactification
\(\mathfrak E\hookrightarrow\overline{\mathfrak E}_f\to Z\) and the
formula \(f^{\mathrm{cs}}_!K=\overline f_*j_{f,!}K\), with the
proper-map specialization \(f^{\mathrm{cs}}_!=f_*\).  It does not prove
choice independence, base change, projection formula,
Thom--Sebastiani transport, symmetric descent, quotient descent,
transition compatibility, associativity, or aggregate hybrid-carrier
population.
The mixed-correspondence base-change packet is
\[
\texttt{certificates/hybrid/mixed\_correspondence\_base\_change}.
\]
Its verifier returns
\(\texttt{MIXED\_CORRESPONDENCE\_BASE\_CHANGE\_VERIFIED}\):
Proposition~\ref{prop:mixed-correspondence-base-change} proves the
Cartesian compact-support base-change isomorphisms for the \(LW\) and
\(WL\) target maps.  It does not prove compactification-choice
independence, projection formula, Thom--Sebastiani transport, quotient
descent, transition compatibility, associativity, or aggregate
hybrid-carrier population.
The mixed-correspondence projection-formula packet is
\[
\texttt{certificates/hybrid/mixed\_correspondence\_projection\_formula}.
\]
Its verifier returns
\(\texttt{MIXED\_CORRESPONDENCE\_PROJECTION\_FORMULA\_VERIFIED}\):
Proposition~\ref{prop:mixed-correspondence-projection-formula} proves
the compact-support projection formula for the \(LW\) and \(WL\)
target maps with retained finite-Tor target coefficients.  It does not
prove compactification-choice independence, Thom--Sebastiani
transport, quotient descent, transition compatibility, associativity,
or aggregate hybrid-carrier population.
The mixed-correspondence Thom--Sebastiani packet is
\[
\texttt{certificates/hybrid/mixed\_correspondence\_thom\_sebastiani}.
\]
Its verifier returns
\(\texttt{MIXED\_THOM\_SEBASTIANI\_TRANSPORT\_CONDITIONAL\_VERIFIED}\):
Proposition~\ref{prop:mixed-correspondence-thom-sebastiani} proves the
binary \(LW/WL\) reduced Thom--Sebastiani transport from supplied
additive d-critical charts and zero-defect orientation transport.  It
does not prove \((O1)\), quotient descent, transition compatibility,
two-step flag associativity, or aggregate hybrid-carrier population.
The eight-word binary vocabulary packet is
\[
\texttt{certificates/hybrid/eight\_word\_binary\_vocabulary}.
\]
Its verifier returns
\(\texttt{EIGHT\_WORD\_BINARY\_VOCABULARY\_DEFINED}\):
Definition~\ref{def:eight-local-wrapped-binary-words} fixes the
alphabet \(\{\mathrm L,\mathrm W\}\), the binary type product, the
eight length-three words, and their two parenthesization symbols.  It
does not construct the two-step flag stacks or prove associativity.
The eight-word two-step flag-stack packet is
\[
\texttt{certificates/hybrid/eight\_word\_two\_step\_flag\_stacks}.
\]
Its verifier returns
\(\texttt{EIGHT\_WORD\_TWO\_STEP\_FLAG\_STACKS\_CONSTRUCTED}\):
Proposition~\ref{prop:eight-two-step-flag-stacks} constructs the
retained two-step flag stacks and the two comparison maps for all
eight words.  It does not prove associativity, pentagon coherence,
quotient descent, transition compatibility, or aggregate population.
The \(LLL\) associativity packet is
\[
\texttt{certificates/hybrid/lll\_associativity}.
\]
Its verifier returns
\(\texttt{LLL\_ASSOCIATIVITY\_CONDITIONAL\_VERIFIED}\):
Proposition~\ref{prop:lll-associativity} proves the pure local
associativity defect vanishes on the \(LLL\) two-step flag stack,
under the stated local base-change, projection-formula, and reduced
Thom--Sebastiani coefficient rows.  It does not prove the seven mixed
or wrapped word associativity rows or the four-input pentagon.
The \(LLW\) associativity packet is
\[
\texttt{certificates/hybrid/llw\_associativity}.
\]
Its verifier returns
\(\texttt{LLW\_ASSOCIATIVITY\_CONDITIONAL\_VERIFIED}\):
Proposition~\ref{prop:llw-associativity} proves the \(LLW\)
associativity defect vanishes on the local/local--mixed two-step flag
stack, under the stated local and mixed functorial coefficient rows.
It does not prove the six remaining word associativity rows or the
four-input pentagon.
The \(LWL\) associativity packet is
\[
\texttt{certificates/hybrid/lwl\_associativity}.
\]
Its verifier returns
\(\texttt{LWL\_ASSOCIATIVITY\_CONDITIONAL\_VERIFIED}\):
Proposition~\ref{prop:lwl-associativity} proves the \(LWL\)
associativity defect vanishes on the ordered mixed--mixed two-step flag
stack, under the stated mixed functorial coefficient rows.  It does not
prove the five remaining word associativity rows or the four-input
pentagon.
The \(WLL\) associativity packet is
\[
\texttt{certificates/hybrid/wll\_associativity}.
\]
Its verifier returns
\(\texttt{WLL\_ASSOCIATIVITY\_CONDITIONAL\_VERIFIED}\):
Proposition~\ref{prop:wll-associativity} proves the \(WLL\)
associativity defect vanishes on the mixed--local/local two-step flag
stack, under the stated local and mixed functorial coefficient rows.
It does not prove the four \(WW\)-dependent word associativity rows or
the four-input pentagon.
The \(LWW\) associativity packet is
\[
\texttt{certificates/hybrid/lww\_associativity}.
\]
Its verifier returns
\(\texttt{LWW\_ASSOCIATIVITY\_CONDITIONAL\_VERIFIED}\):
Proposition~\ref{prop:lww-associativity} proves the \(LWW\)
associativity defect vanishes on the mixed--wrapped/wrapped two-step
flag stack, under the stated mixed functorial rows and
wrapped/wrapped functorial hypotheses.  It does not prove the three
remaining \(WW\)-dependent word associativity rows or the four-input
pentagon.
The \(WLW\) associativity packet is
\[
\texttt{certificates/hybrid/wlw\_associativity}.
\]
Its verifier returns
\(\texttt{WLW\_ASSOCIATIVITY\_CONDITIONAL\_VERIFIED}\):
Proposition~\ref{prop:wlw-associativity} proves the \(WLW\)
associativity defect vanishes on the mixed--wrapped/wrapped two-step
flag stack, under the stated mixed functorial rows and
wrapped/wrapped functorial hypotheses.  It does not prove the two
remaining \(WW\)-dependent word associativity rows or the four-input
pentagon.
The \(WWL\) associativity packet is
\[
\texttt{certificates/hybrid/wwl\_associativity}.
\]
Its verifier returns
\(\texttt{WWL\_ASSOCIATIVITY\_CONDITIONAL\_VERIFIED}\):
Proposition~\ref{prop:wwl-associativity} proves the \(WWL\)
associativity defect vanishes on the wrapped/wrapped--mixed two-step
flag stack, under the stated mixed functorial rows and
wrapped/wrapped functorial hypotheses.  It does not prove the
remaining \(WWW\) word associativity row or the four-input pentagon.
The \(WWW\) associativity packet is
\[
\texttt{certificates/hybrid/www\_associativity}.
\]
Its verifier returns
\(\texttt{WWW\_ASSOCIATIVITY\_CONDITIONAL\_VERIFIED}\):
Proposition~\ref{prop:www-associativity} proves the \(WWW\)
associativity defect vanishes on the wrapped/wrapped two-step flag
stack, under the stated wrapped/wrapped functorial hypotheses.  The
eight length-three wordwise associativity defects are then zero under
the stated functorial rows.  The four-input pentagon remains separate.
The four-input pentagon packet is
\[
\texttt{certificates/hybrid/four\_input\_pentagon\_coherence}.
\]
Its verifier returns
\(\texttt{FOUR\_INPUT\_PENTAGON\_CONDITIONAL\_VERIFIED}\):
Proposition~\ref{prop:four-input-pentagon-coherence} constructs the
retained four-step flag stack for every word in
\(\{\mathrm L,\mathrm W\}^4\) and proves the Mac Lane pentagon for
the binary hybrid product on that common refinement, under the stated
functorial hypotheses.  It does not define the higher-coloured tree
category or prove quotient descent, transition compatibility, or
aggregate population.
The higher-coloured tree-category packet is
\[
\texttt{certificates/hybrid/higher\_coloured\_tree\_category\_definition}.
\]
Its verifier returns
\(\texttt{HIGHER\_COLOURED\_TREE\_CATEGORY\_DEFINED}\):
Definition~\ref{def:higher-coloured-hybrid-tree-category} defines the
nonunital higher-coloured hybrid tree category and the six components
of \(\mathfrak o^{\mathrm{col}}_R\).  It does not prove
\(\mathfrak o^{\mathrm{col}}_R=0\), unit compatibility, symmetric
descent, wrapped order conventions, quotient descent, transition
compatibility, or aggregate population.
The higher-coloured residual criterion packet is
\[
\texttt{certificates/hybrid/higher\_coloured\_residual\_vanishing\_criterion}.
\]
Its verifier returns
\(\texttt{HIGHER\_COLOURED\_RESIDUAL\_CONDITIONAL\_CRITERION\_VERIFIED}\):
Proposition~\ref{prop:higher-coloured-residual-vanishing-criterion}
proves the finite-stage implication from the six named component
rows to \(\mathfrak o^{\mathrm{col}}_R=0\).  The tree component is
discharged by the eight associativity rows and the four-input
pentagon.  The unit and symmetric-relabelling rows are supplied by
separate packets.  The full non-tree tuple remains incomplete:
refinement, ordered-chart descent, overlap, quotient, transition, and
aggregate rows are not supplied.
The hybrid unit packet is
\[
\texttt{certificates/hybrid/hybrid\_unit\_object\_and\_vacuum\_correspondences}.
\]
Its verifier returns
\(\texttt{HYBRID\_UNIT\_OBJECT\_AND\_VACUUM\_CORRESPONDENCES\_DEFINED}\):
Definition~\ref{def:hybrid-unit-object-vacuum-correspondences}
defines the local zero unit object and the four split vacuum
correspondences for local and wrapped inputs.  It does not prove the
local or wrapped unit identity laws, quotient descent, transition
compatibility, or aggregate population.
The local unit compatibility packet is
\[
\texttt{certificates/hybrid/local\_unit\_compatibility}.
\]
Its verifier returns
\(\texttt{LOCAL\_UNIT\_COMPATIBILITY\_VERIFIED}\):
Proposition~\ref{prop:local-unit-compatibility} proves that the left
and right \(LL\) split vacuum correspondences act as identity
pull-push maps on local coefficients.  It does not prove wrapped unit
compatibility, quotient descent, transition compatibility, or
aggregate population.
The wrapped unit compatibility packet is
\[
\texttt{certificates/hybrid/wrapped\_unit\_compatibility}.
\]
Its verifier returns
\(\texttt{WRAPPED\_UNIT\_COMPATIBILITY\_VERIFIED}\):
Proposition~\ref{prop:wrapped-unit-compatibility} proves that the
left \(LW\) and right \(WL\) split vacuum correspondences act as
identity pull-push maps on wrapped coefficients, with wrapped anchor
memory unchanged before the reduced \(E\)-quotient.  It does not prove
quotient descent, transition compatibility, or aggregate population.
The local configuration symmetric descent packet is
\[
\texttt{certificates/hybrid/local\_configuration\_symmetric\_descent}.
\]
Its verifier returns
\(\texttt{LOCAL\_CONFIGURATION\_SYMMETRIC\_DESCENT\_VERIFIED}\):
Proposition~\ref{prop:local-configuration-symmetric-descent} proves
finite symmetric-group descent for projection-finite closed local
configuration charts and their reduced coefficients before the reduced
\(E\)-quotient.  It does not prove collision compatibility,
closed-chart refinement, ordered-chart descent to the symmetric
finite-colour colimit, quotient descent, transition compatibility, or
aggregate population.
The wrapped insertion order convention packet is
\[
\texttt{certificates/hybrid/wrapped\_insertion\_order\_conventions}.
\]
Its verifier returns
\(\texttt{WRAPPED\_INSERTION\_ORDER\_CONVENTIONS\_VERIFIED}\):
Proposition~\ref{prop:wrapped-insertion-order-conventions} fixes the
planar wrapped source order, the \(WW\) quotient/subobject convention,
and the order preservation of the eight two-step flag comparison maps
before the reduced \(E\)-quotient.  It does not prove unordered
wrapped-pair symmetric descent, ordered-chart descent, overlap
compatibility, quotient descent, transition compatibility, or aggregate
population.
The ordered two-sided mixed correspondence packet is
\[
\texttt{certificates/hybrid/two\_sided\_mixed\_correspondence\_definition}.
\]
Its verifier returns
\(\texttt{TWO\_SIDED\_MIXED\_CORRESPONDENCE\_DEFINITION\_VERIFIED}\):
Definition~\ref{def:two-sided-mixed-correspondence-datum} defines the
ordered \(LWL\) unreduced flag diagram; flag-stack, source/target map,
populated anchor-memory, admissibility, base-change, projection-formula,
Thom--Sebastiani, and transition rows are not supplied.
The source/target anchor-memory packet is
\[
\texttt{certificates/hybrid/source\_target\_anchor\_memory\_definition}.
\]
Its verifier returns
\(\texttt{SOURCE\_TARGET\_ANCHOR\_MEMORY\_DEFINITION\_VERIFIED}\):
Definition~\ref{def:source-target-anchor-memory-datum} defines the
\(LW\), \(WL\), \(WW\), and \(LWL\) anchor-memory maps on unreduced
correspondence stacks before the \(E\)-quotient.  The determinant
anchor is constructed by the next packet and its translation weight by
the following packet.  The \(\chi=0\) degeneracy is isolated by the
chi-zero packet.  The degree-zero Jordan--Hoelder extra-anchor datum is
recorded by the extra-anchor-data packet; anchor-residual, full
losslessness, quotient-descent, stabilizer linearization, populated
finite-stack, and transition rows are not supplied by the
source/target-memory packet.
The determinant-anchor construction packet is
\[
\texttt{certificates/hybrid/determinant\_anchor\_construction}.
\]
Its verifier returns
\(\texttt{DETERMINANT\_ANCHOR\_CONSTRUCTION\_VERIFIED}\):
Definition~\ref{def:determinant-anchor} constructs
\(\lambda^{\det}_{\eta,R}\) as the normalized determinant of
cohomology, checks that it lands in \(\operatorname{Pic}^0(E)\simeq E\),
and records determinant-functor base-change; unit-weight, \(\chi=0\)
degeneracy, Jordan--Hoelder extra-anchor data, anchor-residual,
losslessness, quotient-descent, populated wrapped-prequotient, and
transition rows are not supplied.  The determinant-anchor
translation-weight packet
\[
\texttt{certificates/hybrid/determinant\_anchor\_translation\_weight}
\]
returns
\(\texttt{DETERMINANT\_ANCHOR\_TRANSLATION\_WEIGHT\_VERIFIED}\): under
the support-translation convention it computes
\(\lambda^{\det}_{\eta,R}(a\cdot\mathcal F)
=\lambda^{\det}_{\eta,R}(\mathcal F)+\chi_\eta a\).
The determinant-anchor chi-zero degeneracy packet
\[
\texttt{certificates/hybrid/determinant\_anchor\_chi\_zero\_degeneracy}
\]
returns
\(\texttt{DETERMINANT\_ANCHOR\_CHI\_ZERO\_DEGENERACY\_VERIFIED}\):
on the branch \(\chi_\eta=0\), the determinant anchor is constant on
the full retained \(E\)-translation orbit and has rank-one separation
defect.  The determinant-anchor extra-anchor-data packet
\[
\texttt{certificates/hybrid/determinant\_anchor\_extra\_anchor\_data}
\]
returns
\(\texttt{DETERMINANT\_ANCHOR\_EXTRA\_ANCHOR\_DATA\_VERIFIED}\):
the degree-zero Jordan--Hoelder divisor refines the determinant anchor
and separates the rank-two pair \(\mathcal O_E^{\oplus2}\) and
\(L\oplus L^{-1}\), \(L\ne\mathcal O_E\).  Losslessness,
anchor-residual, quotient-descent, stabilizer linearization, populated
wrapped-prequotient, and transition rows are not supplied.
The quotient-after-correspondence pseudofunctor definition packet is
\[
\texttt{certificates/hybrid/quotient\_after\_correspondence\_pseudofunctor\_definition}.
\]
Its verifier returns
\(\texttt{QUOTIENT\_AFTER\_CORRESPONDENCE\_PSEUDOFUNCTOR\_DEFINED}\):
Definition~\ref{def:quotient-after-correspondence-pseudofunctor}
defines the object assignment \(Y\mapsto [Y/E]\), the generating-arrow
assignment
\((Y_0\leftarrow\mathfrak E_e\to Y_1)\mapsto
([Y_0/E]\leftarrow[\mathfrak E_e/E]\to[Y_1/E])\), the
coefficient-descent interface, and the quotient residual components.
The composition component is supplied by
\[
\texttt{certificates/hybrid/quotient\_after\_correspondence\_composition}.
\]
Proposition~\ref{prop:quotient-after-correspondence-preserves-composition}
proves preservation of finite-height composition of retained spans.
The base-change square component is supplied by
\[
\texttt{certificates/hybrid/quotient\_after\_correspondence\_base\_change}.
\]
Proposition~\ref{prop:quotient-after-correspondence-preserves-base-change}
proves preservation of retained finite-height cartesian squares after
quotient.  The Thom--Sebastiani descent component is supplied by
\[
\texttt{certificates/hybrid/quotient\_after\_correspondence\_thom\_sebastiani}.
\]
Proposition~\ref{prop:quotient-after-correspondence-preserves-thom-sebastiani}
proves descent of equivariant reduced Thom--Sebastiani transport
through the quotient correspondence.  The quotient-first exclusion
component is supplied by
\[
\texttt{certificates/hybrid/quotient\_after\_correspondence\_quotient\_first\_exclusion}.
\]
Proposition~\ref{prop:quotient-after-correspondence-excludes-quotient-first}
proves that every reduced generating arrow used by
\(\mathcal Q^{\mathrm{corr}}_{E,R}\) is the quotient of a preformed
unreduced \(E\)-equivariant correspondence.  The Borel--Moore chain
functor component is supplied by
\[
\texttt{certificates/hybrid/quotient\_after\_correspondence\_bm\_chain\_functor}.
\]
Proposition~\ref{prop:quotient-after-correspondence-borel-moore-chain-functor}
constructs \(Q_{E,R}\) on quotient-presented Borel--Moore chain
objects.  The operation-comparison maps are supplied by
\[
\texttt{certificates/hybrid/quotient\_after\_correspondence\_theta\_mu\_comparison}.
\]
Proposition~\ref{prop:quotient-after-correspondence-theta-mu-comparison}
constructs \(\theta^Q_{\mu,R}\) for quotient-presented operation
rows.  The quotient two-step flag associativity component is supplied
by
Proposition~\ref{prop:quotient-after-correspondence-theta-mu-associativity}.
The quotient product-coproduct compatibility component is supplied by
Proposition~\ref{prop:quotient-after-correspondence-theta-mu-coproduct}.
The quotient primitive-projection compatibility component is supplied
by
Proposition~\ref{prop:quotient-after-correspondence-theta-mu-primitives}.
The HN-transition compatibility of \(Q_{E,R}\) is supplied by
Proposition~\ref{prop:quotient-after-correspondence-hn-transition-compatibility}.
Compact-support pushforward base change, protected-integration
transition compatibility, and aggregate population are not supplied.  Hall
product, Hall coproduct, and primitive-subspace transition preservation
are not supplied by row 256.
The definition of the protected integration map is supplied separately
by
Proposition~\ref{prop:protected-integration-definition}.
Its product compatibility is supplied separately by
Proposition~\ref{prop:protected-integration-product-compatibility}.
Its coproduct compatibility is supplied separately by
Proposition~\ref{prop:protected-integration-coproduct-compatibility}.
Its primitive-projection compatibility is supplied separately by
Proposition~\ref{prop:protected-integration-primitive-compatibility}.
Its Gram degree is supplied separately by
Proposition~\ref{prop:protected-integration-gram-degree}.
Its normal-ordered label comparison is supplied separately by
Proposition~\ref{prop:protected-integration-pi-x-separation}.

\textit{Construction.}
\begin{enumerate}
\item Split $\Gamma_R$ into $\Gamma_R^{\mathrm{loc}}$
($b_R^{\mathrm{geom}}=0$) and $\Gamma_R^{\mathrm{wr}}$
($b_R^{\mathrm{geom}}>0$).
\item Form the local carrier
\(\operatorname{Ran}^{\mathrm{loc}}_{R,\mathrm{pre}}(E)\) as the
\(J=\varnothing\) subprestack of
\(\operatorname{Ran}^{\mathrm{hyb}}_{R,\mathrm{pre}}(E)\), with local
colour labels in \(\Gamma_R^{\mathrm{loc}}\).
\item Form the wrapped carrier
\(\operatorname{Ran}^{\mathrm{wr}}_{R,\mathrm{pre}}(E)\) by imposing
\(I=\varnothing\) in the hybrid prestack and allowing only colour maps
\(\eta:J\to\Gamma_R^{\mathrm{wr}}\), with wrapped prequotient families
and their anchors retained.
\item For $\alpha\in\Gamma_R^{\mathrm{loc}}$: form the
closed-configuration incidence stacks
$\mathfrak M_{\alpha,R,I}^{\mathrm{loc,cl}}$ and the configuration
maps $\pi_{\alpha,R,I}:\mathfrak M_{\alpha,R,I}^{\mathrm{loc,cl}}
\to E^I$, with vanishing-cycle/orientation pushforwards
$\mathscr H_{\alpha,R,I}^{\mathrm{loc}}$.  Define
$\mathcal F_{\alpha,R}^{\mathrm{loc}}(U)$ by symmetric-group descent
on $E^I$.
\item Retain the ordered \(LL\) local/local correspondence diagram of
Definition~\ref{def:local-local-correspondence-datum}; its
target-proper extension stack is
Proposition~\ref{prop:local-local-target-properness}.  Source
properness is not claimed.  Descent, collision compatibility,
Thom--Sebastiani, compact-support pushforward, aggregate population,
and transition rows remain separate obligations.
\item For $\eta\in\Gamma_R^{\mathrm{wr}}$: form the rigidified
prequotient $\mathcal M_{\eta,R}^{\mathrm{wr,rig}}\to
\mathcal M_{\eta,R}^{\mathrm{wr}}\) of
Definition~\ref{def:e-equivariant-wrapped-prequotient-stack}, with its
\(E\)-translation action before the reduced quotient is taken.  The
finite-type assertion is
Proposition~\ref{prop:wrapped-prequotient-finite-type}.  Then
attach the determinant anchor \(\lambda^{\det}_{\eta,R}\) of
Definition~\ref{def:determinant-anchor}.  Its translation weight,
computed in
Proposition~\ref{prop:determinant-anchor-translation-weight}, is
\(\chi_\eta\).  On the \(\chi_\eta=0\) branch,
Proposition~\ref{prop:determinant-anchor-chi-zero-degeneracy} gives a
rank-one determinant-anchor separation defect.  Attach the
degree-zero Jordan--Hoelder divisor of
Proposition~\ref{prop:determinant-anchor-extra-data-separates-rank-two}
as extra anchor data; it separates the determinant-zero rank-two test
pair but is not full anchor losslessness.  Unit-weight normalization,
losslessness, quotient descent, stabilizer linearization, transition
compatibility, and finite-stage population remain separate rows.  The
finite anchor residual \(\mathfrak o^\lambda_R\) is
Definition~\ref{def:finite-anchor-residual}; its vanishing is a later
row.
\item Retain the one-sided \(LW\) and \(WL\) mixed local/wrapped
correspondence diagrams of
Definition~\ref{def:mixed-local-wrapped-correspondence-datum}; their
target-admissible extension-stack and target-map rows are supplied by
Proposition~\ref{prop:mixed-local-wrapped-target-admissibility}.  Their
populated anchor-memory, compact-support pushforward, base-change,
projection-formula, Thom--Sebastiani, quotient-descent, and transition
rows remain separate obligations.
\item Retain the ordered \(WW\) wrapped/wrapped correspondence diagram
of Definition~\ref{def:wrapped-wrapped-correspondence-datum}; its
target-admissible extension-stack and target-map rows are supplied by
Proposition~\ref{prop:wrapped-wrapped-target-admissibility}.  Its
populated anchor-memory, compact-support pushforward, symmetric
descent, Thom--Sebastiani, quotient-descent, and transition rows
remain separate obligations.
\item Retain the compact-support exceptional pushforward model of
Definition~\ref{def:compact-support-exceptional-pushforward-model}.
It supplies the chosen formula \(f^{\mathrm{cs}}_!=\overline f_*j_!\)
and the proper-map specialization.  It does not populate every
correspondence map with a compactification and does not prove base
change, projection formula, Thom--Sebastiani, symmetric descent,
quotient descent, transition compatibility, or associativity.
\item Retain the \(LW/WL\) mixed base-change isomorphisms of
Proposition~\ref{prop:mixed-correspondence-base-change}.  They apply to
Cartesian compact-support squares with pulled coefficient complexes.
Projection formula, Thom--Sebastiani, quotient descent, transitions,
associativity, and aggregate population remain separate rows.
\item Retain the \(LW/WL\) mixed projection-formula isomorphisms of
Proposition~\ref{prop:mixed-correspondence-projection-formula}.  They
apply to retained finite-Tor target coefficients in the chosen
compact-support model.  Thom--Sebastiani, quotient descent,
transitions, associativity, and aggregate population remain separate
rows.
\item Retain the \(LW/WL\) mixed Thom--Sebastiani transports of
Proposition~\ref{prop:mixed-correspondence-thom-sebastiani}.  They
apply on the supplied reduced additive d-critical charts with
zero-defect Joyce--Upmeier orientation transport.  Quotient descent,
transitions, two-step flag associativity, and aggregate population
remain separate rows.
\item Retain the eight-word binary vocabulary of
Definition~\ref{def:eight-local-wrapped-binary-words}.  This fixes
\(\{\mathrm{LLL},\mathrm{LLW},\mathrm{LWL},\mathrm{WLL},
\mathrm{LWW},\mathrm{WLW},\mathrm{WWL},\mathrm{WWW}\}\), the
\(\mathrm L/\mathrm W\) type product, and the two parenthesization
symbols.  The actual two-step flag stacks, associativity defects,
pentagon, quotient descent, transitions, and aggregate population
remain separate rows.
\item Retain the eight two-step flag stacks and their left/right
comparison maps from Proposition~\ref{prop:eight-two-step-flag-stacks}.
Associativity defects, the four-input pentagon, quotient descent,
transitions, and aggregate population remain separate rows.
\item Retain the \(LLL\) associativity identity of
Proposition~\ref{prop:lll-associativity}.  The seven remaining word
associativity rows, the four-input pentagon, quotient descent,
transitions, and aggregate population remain separate rows.
\item Retain the \(LLW\) associativity identity of
Proposition~\ref{prop:llw-associativity}.  The six remaining word
associativity rows, the four-input pentagon, quotient descent,
transitions, and aggregate population remain separate rows.
\item Retain the \(LWL\) associativity identity of
Proposition~\ref{prop:lwl-associativity}.  The five remaining word
associativity rows, the four-input pentagon, quotient descent,
transitions, and aggregate population remain separate rows.
\item Retain the \(WLL\) associativity identity of
Proposition~\ref{prop:wll-associativity}.  The four remaining
\(WW\)-dependent associativity rows, the four-input pentagon, quotient
descent, transitions, and aggregate population remain separate rows.
\item Retain the \(LWW\) associativity identity of
Proposition~\ref{prop:lww-associativity}.  The three remaining
\(WW\)-dependent associativity rows, the four-input pentagon, quotient
descent, transitions, and aggregate population remain separate rows.
\item Retain the \(WLW\) associativity identity of
Proposition~\ref{prop:wlw-associativity}.  The two remaining
\(WW\)-dependent associativity rows, the four-input pentagon, quotient
descent, transitions, and aggregate population remain separate rows.
\item Retain the \(WWL\) associativity identity of
Proposition~\ref{prop:wwl-associativity}.  The remaining \(WWW\)
associativity row, the four-input pentagon, quotient descent,
transitions, and aggregate population remain separate rows.
\item Retain the \(WWW\) associativity identity of
Proposition~\ref{prop:www-associativity}.  This completes the
length-three wordwise associativity list.  The four-input pentagon,
quotient descent, transitions, and aggregate population remain
separate rows.
\item Retain the four-input pentagon coherence of
Proposition~\ref{prop:four-input-pentagon-coherence}.  The
higher-coloured tree category, remaining unit compatibilities,
quotient descent, transitions, and aggregate population remain
separate rows.
\item Retain the nonunital higher-coloured tree category of
Definition~\ref{def:higher-coloured-hybrid-tree-category}.  The
vanishing of \(\mathfrak o^{\mathrm{col}}_R\), remaining unit
compatibilities, quotient descent, transitions, and aggregate
population remain separate rows.
\item Retain the conditional higher-coloured residual criterion of
Proposition~\ref{prop:higher-coloured-residual-vanishing-criterion}.
The tree component is discharged by associativity and pentagon
coherence.  The unit and symmetric-relabelling rows are supplied by
separate proof packets.  Refinement, ordered-chart descent, overlap,
quotient descent, transitions, and aggregate population rows remain
separate rows.
\item Retain the hybrid unit object and split vacuum correspondences
of Definition~\ref{def:hybrid-unit-object-vacuum-correspondences}.
The local and wrapped unit identity laws are supplied by the next two
proof rows; quotient descent, transitions, and aggregate population
remain separate rows.
\item Retain the local unit compatibility theorem of
Proposition~\ref{prop:local-unit-compatibility}.  The wrapped unit
identity law, quotient descent, transitions, and aggregate population
remain separate rows.
\item Retain the wrapped unit compatibility theorem of
Proposition~\ref{prop:wrapped-unit-compatibility}.  The unit component
of the higher-coloured residual is now discharged at fixed height
before quotient.  Quotient descent, transitions, and aggregate
population remain separate rows.
\item Retain the local configuration symmetric descent theorem of
Proposition~\ref{prop:local-configuration-symmetric-descent}.  The
symmetry component of the higher-coloured residual is now discharged
for projection-finite local configuration charts at fixed height before
quotient.  Collision compatibility, refinement, ordered-chart descent,
overlap, quotient descent, transitions, and aggregate population remain
separate rows.
\item Retain the wrapped insertion order convention theorem of
Proposition~\ref{prop:wrapped-insertion-order-conventions}.  The
wrapped order input row of the higher-coloured residual criterion is
now supplied at fixed height before quotient.  Unordered wrapped-pair
symmetric descent, ordered-chart descent, overlap, quotient descent,
transitions, and aggregate population remain separate rows.
\item Retain the ordered \(LWL\) two-sided mixed correspondence diagram
of Definition~\ref{def:two-sided-mixed-correspondence-datum}; its
flag-stack, source/target-map, populated anchor-memory, admissibility,
base-change, projection-formula, Thom--Sebastiani, and transition rows
remain separate obligations.
\item Retain the source/target anchor-memory maps
\(a^{\mathrm{LW}},a^{\mathrm{WL}},a^{\mathrm{WW}}\), and
\(a^{\mathrm{LWL}}\) of
Definition~\ref{def:source-target-anchor-memory-datum} on the
unreduced mixed and wrapped correspondence types.  This step assumes
the abstract wrapped anchors \(\lambda_{\eta,R}\).  The
determinant-anchor construction and translation-weight computation are
separate packets, and the \(\chi=0\) determinant-anchor degeneracy is
handled by its own packet.  The degree-zero Jordan--Hoelder
extra-anchor datum and the residual definition \(o^\lambda_R\) are
separate packets; residual vanishing, losslessness, quotient descent,
and transition rows remain separate obligations.
\item Retain the quotient-after-correspondence pseudofunctor definition
of Definition~\ref{def:quotient-after-correspondence-pseudofunctor}.
This supplies the object, arrow, and coefficient-descent interface
rows before Borel--Moore chains.
\item Retain the composition comparison of
Proposition~\ref{prop:quotient-after-correspondence-preserves-composition}.
This supplies \(\mathfrak o^{Q,\mathrm{comp}}_R=0\) for finite-height
retained spans.
\item Retain the base-change square comparison of
Proposition~\ref{prop:quotient-after-correspondence-preserves-base-change}.
This supplies \(\mathfrak o^{Q,\mathrm{bc}}_R=0\) for finite-height
retained cartesian squares.  Compact-support pushforward base change
and aggregate population remain separate rows; the
quotient flag component is treated below.
\item Retain the Thom--Sebastiani descent comparison of
Proposition~\ref{prop:quotient-after-correspondence-preserves-thom-sebastiani}.
This supplies \(\mathfrak o^{Q,\mathrm{TS}}_R=0\) for equivariant
finite-height reduced Thom--Sebastiani rows.
\item Retain the quotient-first exclusion of
Proposition~\ref{prop:quotient-after-correspondence-excludes-quotient-first}.
This supplies the row-250 firewall: an independent object-quotient-
first Hall product is not used unless it is identified with
\([\mathfrak E_e/E]\) for a preformed unreduced correspondence.
\item Retain the Borel--Moore chain functor of
Proposition~\ref{prop:quotient-after-correspondence-borel-moore-chain-functor}.
This supplies \(\mathfrak o^{Q,\mathrm{BM}}_R=0\) on
quotient-presented coefficient objects.
\item Retain the operation-comparison maps of
Proposition~\ref{prop:quotient-after-correspondence-theta-mu-comparison}.
This supplies \(\mathfrak o^{Q,\theta}_R=0\) for quotient-presented
finite-height operation rows.
\item Retain the associativity compatibility of
Proposition~\ref{prop:quotient-after-correspondence-theta-mu-associativity}.
This supplies \(\mathfrak o^{Q,\mathrm{flag}}_R=0\) for the eight
wordwise two-step flag associativity rows.
\item Retain the coproduct compatibility of
Proposition~\ref{prop:quotient-after-correspondence-theta-mu-coproduct}.
This supplies \(\mathfrak o^{Q,\Delta}_R=0\) for supplied finite
product-coproduct rows.
\item Retain the primitive-projection compatibility of
Proposition~\ref{prop:quotient-after-correspondence-theta-mu-primitives}.
This supplies \(\mathfrak o^{Q,\Prim}_R=0\) for supplied finite
primitive kernel and projection rows.
\item Retain the HN-transition compatibility of \(Q_{E,R}\) from
Proposition~\ref{prop:quotient-after-correspondence-hn-transition-compatibility}.
This supplies \(\mathfrak o^{Q,\mathrm{tr}}_{R'R}=0\) for supplied
quotient-presented finite HN object and coefficient transition rows.
Protected integration and aggregate population remain separate rows.
Hall product, Hall coproduct, and primitive-subspace transition
preservation remain outside row 256.
\item Retain the eight-word vocabulary
\(\mathcal W^{(3)}_{\mathrm{hyb}}\) of
Definition~\ref{def:eight-local-wrapped-binary-words} and the
two-step binary flag stacks
\(\mathfrak F^{(2),w}_{\xi_1,\xi_2,\xi_3,R}\) of
Proposition~\ref{prop:eight-two-step-flag-stacks}.  The vanishing of
\(\mathfrak o^{\mathrm{assoc}}_{\mathrm{LLL},R}\) is supplied by
Proposition~\ref{prop:lll-associativity}, and
\(\mathfrak o^{\mathrm{assoc}}_{\mathrm{LLW},R}\) is supplied by
Proposition~\ref{prop:llw-associativity}; the \(LWL\) associativity
defect is zero by Proposition~\ref{prop:lwl-associativity}; and the
\(WLL\) associativity defect is zero by
Proposition~\ref{prop:wll-associativity}.  The \(LWW\) associativity
defect is zero by Proposition~\ref{prop:lww-associativity}.  The
\(WLW\) associativity defect is zero by
Proposition~\ref{prop:wlw-associativity}.  The \(WWL\) associativity
defect is zero by Proposition~\ref{prop:wwl-associativity}; and the
\(WWW\) associativity defect is zero by
Proposition~\ref{prop:www-associativity}.  The four-input pentagon is
supplied by Proposition~\ref{prop:four-input-pentagon-coherence}.
The unit, refinement, symmetry, descent, and higher-coloured tree data
remain in \(\mathfrak o^{\mathrm{col}}_R\).
\item Use the supplied quotient-after-correspondence pseudofunctor
definition
\(\mathsf{Corr}^{E,\mathrm{hyb}}_R\to
\mathsf{Corr}^{\mathrm{red},\mathrm{hyb}}_R\), then supply the
separate compact-support pushforward base-change and transition rows
before treating the quotient residual as zero.
\item Define $I_R^{\mathrm{prot}}$ on the reduced functor by
Proposition~\ref{prop:protected-integration-definition}.  Its product
compatibility is supplied by
Proposition~\ref{prop:protected-integration-product-compatibility}.
Its coproduct compatibility is supplied by
Proposition~\ref{prop:protected-integration-coproduct-compatibility}.
Its primitive-projection compatibility is supplied by
Proposition~\ref{prop:protected-integration-primitive-compatibility}.
Its Gram degree is supplied by
Proposition~\ref{prop:protected-integration-gram-degree}.  The
normal-ordered label comparison is supplied by
Proposition~\ref{prop:protected-integration-pi-x-separation}.  The
transition compatibility row remains separate.
\end{enumerate}

\textit{Obstruction classes.}  $\mathfrak O_{\mathrm{hyb},R}$
(\S\ref{subsec:hybrid-ran-carrier}).

\subsubsection*{Orientation descent and Weyl transport}

\textit{Hypotheses.} The hybrid wrapped factorization datum.

\textit{Finite datum.} The (O1) and (O1)$^+$ data: null-trivializations of
the reduced and equivariant orientation gerbes,
$E[N]$-linearizations with vanishing degree-one characters, and
Weyl-equivariant lifts $\tau_i$ ($i=1,2,3$) compatible with the
quotient-orientation cocycle $\mathfrak q_{\mathcal C}$.

\textit{Construction.}
\begin{enumerate}
\item For each charge stratum $\mathcal C$: compute the reduced
gerbe $\alpha^{\mathrm{red}}_{\mathcal C}$, free $E$-class
$\alpha^{E,\mathrm{free}}_{\mathcal C}$, and equivariant
finite-stabilizer classes $\widetilde\beta^{H}_{\mathcal C,S}$ for
the actual stabilizers $H\subset E_{\mathrm{tors}}$ on
finite-stabilizer strata.
\item Verify the Klein-four condition at $H=E[2]$:
$(b_{20},b_{11},b_{02})=(r_1,r_1+r_2+r_3,r_2)$ with
$r_1=r_2=r_3=0$; for full-level even $N\ge 4$, verify the
two-primary residual on $E[N]_{(2)}$ including the mixed coefficient
$A_{12}^{(N)}$.
\item Choose zero $H^1(BH;\mathbb F_2)$-linearization characters
$\lambda^H_{\mathcal C,S}=0$.
\item For each simple type-II reflection $s_{\delta_i}$, supply the
lift $\tau_i:o_{\mathcal C}\to s_{\delta_i}^* o_{\mathcal C}$ on the
finite partial action groupoid $\mathcal G_R$; compute the
$2$-cocycle $c_{o,R}\in Z^2(\mathcal G_R;\underline{\mathbb F}_2)$
and verify $[c_{o,R}]=0$ with chosen $1$-cochain.
\item Verify $\tau_i^2=1$ and the Coxeter relations on the
$W^{(2)}(\Lambda_{II}^{2,1})$-orbit; verify
$\omega_{i,\mathcal C}=0$ for each torsor defect of the
quotient-orientation transport.
\end{enumerate}

\textit{Obstruction classes.}  $\mathfrak O_{\mathrm{or},R}$
(\S\ref{subsec:cosection-twisted-orientation}); the (O2) wall-atlas
enters the Pfaffian wall normal form.

\subsubsection*{Hall product, coproduct, and primitive projection}

\textit{Hypotheses.} The finite Hall atlas, hybrid carrier, and orientation transport.

\textit{Finite datum.} The compact Hall product $m_R^{\mathrm{red}}$,
collision coproduct $\Delta_R^{\mathrm{ch}}$, primitive projection
$P_R^{\Prim}$, and Hopf pairing $\langle\,,\,\rangle_R$ on the
cosection-reduced d-critical heart, with normal-ordered Gram descent
$\overline\Pi_{X,*}^\Theta$.  The product alone is not a primitive
datum by Proposition~\ref{prop:hall-product-only-no-primitives}.

\textit{Construction.}
\begin{enumerate}
\item Form $m_R^{\mathrm{red}}=q_!\circ\operatorname{TS}^{\mathrm{red}}_R
\circ p^*$ on each binary correspondence (LL, LW, WL, WW); verify
two-step pentagon coherence on the eight-word atlas.  This operation is
defined only in the retained compactified correspondence model, or in
an explicitly supplied exceptional compact-support model as in
Definition~\ref{def:compact-support-exceptional-pushforward-model};
for \(LW/WL\) Cartesian target refinements, base change is
Proposition~\ref{prop:mixed-correspondence-base-change}; for
\(LW/WL\) retained finite-Tor target coefficients, the projection
formula is Proposition~\ref{prop:mixed-correspondence-projection-formula};
for \(LW/WL\) binary coefficient transport, Thom--Sebastiani is
Proposition~\ref{prop:mixed-correspondence-thom-sebastiani}.
\item Form $\Delta_R^{\mathrm{ch}}$ by pull-push on the $2$-step flag
stack; verify coassociativity and counitality after $Q_{E,R}$ using
the higher-coloured, unit, refinement, quotient, and transition
coherences.
\item Form $\langle\,,\,\rangle_R$ on the primitive layer; verify
Hopf adjointness with the coproduct on the radical quotient.
\item Compute the Frobenius defect $\mathfrak o_{F,R}$ via the
CY-3 Serre-functor identity \cite{KSCoHA}; verify
$\mathfrak o_{F,R}=0$.
\item Compute the radical $\operatorname{Rad}_R^\Pi$; verify
Hopf-coideal and Lie-ideal stability.
\item Form the chain-level normal-ordered Gram descent
$\overline\Pi_{X,*}^\Theta$ via the negative-cyclic
Loday--Quillen--Tsygan route or the PTVV--Joyce orientation route
\cite{LodayQuillen1984,Tsygan1983,LodayCyclicHomology,PTVV2013};
verify the triple homotopy compatibility
$d_{\mathrm{Hoch}}\Theta_\Pi=B_{\mathrm{ch}}$ and the
$\widehat\Gamma_R$-grading absorption.
\end{enumerate}

\textit{Obstruction classes.}  $\mathfrak O_{\mathrm{Hall},R}$,
including $\mathfrak o_{F,R}$, the Hopf-coideal residual, and the
$\Theta$-pentagon residual.

\subsubsection*{Dirac block, Pfaffian line, and type-II wall atlas}

\textit{Hypotheses.} The orientation transport and Hall product.

\textit{Finite datum.} The first-order Dirac block decomposition
$\mathfrak D_{X,R}^{D_0}=\bigoplus_\gamma J_\gamma$, the Pfaffian
line $\mathscr L_{\mathrm{Pf},R}$ with squaring isomorphism, the
Pfaffian-to-automorphic comparison $\iota_{\mathrm{aut},R}$, the
normalized section $\operatorname{pf}_{X,R}$, and the type-II wall
atlas (rank-one normal model on every retained component, boundary
stratum, and overlap).

\textit{Construction.}
\begin{enumerate}
\item Form the Pfaffian line
$\mathscr L_{\mathrm{Pf},R}=\bigotimes_{\mathcal C}
\operatorname{Pf}(K^{\mathrm{red}}_{\mathcal C,R})$ on the
cosection-reduced quotient \cite{BBDJS2015}; verify the squaring
isomorphism.
\item Form $J_\gamma$ on
$P^{\Pi,+}_{R,\gamma}\oplus(P^{\Pi,+}_{R,\gamma})^\vee[1]$ for each
$\gamma\in\Gamma_R^{\Pi,+}\setminus\{0\}$; verify Pfaffian
contribution $(1-x^\gamma)^{\sdim P^{\Pi,+}_{R,\gamma}}$.
\item Verify the support, parity, and leading-coefficient residuals
$\mathfrak o^{\mathrm{supp}}_{R,\gamma}=
\mathfrak o^{\mathrm{par}}_{R,\gamma}=
\mathfrak o^{\mathrm{lead}}_R=0$ on the test window.
\item For each simple type-II wall $\Hpl_{\delta_i}$ ($i=1,2,3$):
take the rank-one Pfaffian crossing chart on every retained
component and boundary stratum; verify
$N_{\delta_i}^{\mathrm{Pf}}=1$ and the local sign
$\varepsilon^{\mathrm{loc}}_{\delta_i,R}=-1$; verify chart-overlap
agreement of the normal Pfaffian units, tangent/normal splittings,
quotient orientations, and protected integration maps; verify the
half-Hilbert orbit sign sum (size $4096=2^{12}$).
\item Supply the Pfaffian-to-automorphic comparison
$\iota_{\mathrm{aut},R}:\mathscr L_{\mathrm{Pf},R}\to\mathcal L^5
\otimes\nu_{\Delta_5}$ compatible with cusp trivialization and
squaring.
\end{enumerate}

\textit{Obstruction classes.}  (D0-Pf), (O2)/hybrid compatibility
residuals, and the Pfaffian-to-automorphic comparison residuals.

\subsubsection*{Koszul source coalgebra and cobar comparison}

\textit{Hypotheses.} The hybrid carrier and Hall product, an augmentation of the hybrid
source, and a wrapped bar-cobar domain datum for the wrapped summands
and mixed local/wrapped words.

\textit{Finite datum.} The conilpotent chiral source coalgebra $C_{X,R}$,
finite filtration $F^{\mathrm{co}}_{R,\bullet}$, collision coproduct
$\Delta_R^{\mathrm{ch}}$, bar comparison $b_{X,R}$, and cobar
comparison $\Theta_{\mathrm{Kos},R}$.

\textit{Construction.}
\begin{enumerate}
\item Choose an augmentation
$\varepsilon_{X,R}^{\mathrm{hyb}}:\mathcal F_{X,\sigma,S,\le R}^{\mathrm{hyb}}
\to k\mathbf 1$ and verify bounded length $N_R$.
\item Form the canonical source-bar coalgebra
$C^{\mathrm{bar}}_{X,R}=\bar B_E^{\mathrm{ch}}(\mathcal F_{X,\sigma,S,\le R}^{\mathrm{hyb}})$
\cite[\S3.7]{BD} with bar-length filtration and deconcatenation
collision coproduct.
\item Verify binary coassociativity on the eight-word atlas, the
higher-coloured and unit coherences, bar-length conilpotence, and
transition compatibility; set $b_{X,R}=\mathrm{id}$.
\item Verify that the projection-finite summands, wrapped summands,
and mixed local/wrapped words lie in the chosen completed
Beilinson--Drinfeld/Francis--Gaitsgory bar-cobar domain.  Under that
wrapped Koszul datum, verify the source-internal counit
$\varepsilon^{\mathrm{src}}_{\mathrm{bar},R}:\Omega_E^{\mathrm{ch}}
C^{\mathrm{bar}}_{X,R}\to\mathcal F_{X,\sigma,S,\le R}^{\mathrm{hyb}}$
is a quasi-isomorphism in the chosen completed/coderived category
\cite[Theorem~6.2.2]{FrancisGaitsgoryCKD}.
\item Construct the source-to-target chiral algebra
quasi-isomorphism $\vartheta_R:
\mathcal F_{X,\sigma,S,\le R}^{\mathrm{hyb}}\xrightarrow{\simeq}
\mathsf{Den}_{\Delta_5,E,\le R}$ preserving differential, product,
unit, transition maps, and primitive restriction.
\item Supply the finite Koszul comparison datum
\(\mathfrak K^{\mathrm{cmp}}_{X,R}\) of
Definition~\ref{def:finite-koszul-comparison-datum}: verify the Weyl
transport square, Pfaffian-orientation sign square, Hall
product/coproduct and Hopf-pairing square, radical quotient square, PBW
associated-graded square, and their transition homotopies.
\item Set $\Theta_{\mathrm{Kos},R}=
\vartheta_R\circ\varepsilon^{\mathrm{src}}_{\mathrm{bar},R}$;
verify the Koszul Mittag--Leffler condition on the cones of $b$ and
$\Theta_{\mathrm{Kos}}$ and on the comparison-defect complexes of
\(\mathfrak K^{\mathrm{cmp}}_{X,R}\).
\end{enumerate}

\textit{Obstruction classes.}  $\mathfrak O_{\mathrm{Kos},R}$,
including $\mathfrak o^\Omega_R=
(\mathfrak o^{\varepsilon,\mathrm{src}}_R,
\mathfrak o^\vartheta_R)$ and the wrapped Koszul defect
\(\mathfrak o^{\mathrm{wrKos}}_R\), together with
\(\mathfrak O_{\mathrm{Kcmp},R}\) for the Weyl, Pfaffian, Hall-pairing,
radical, PBW, transition, and Mittag--Leffler comparison defects.

\subsubsection*{Primitive recognition in cofinal finite windows}

\textit{Hypotheses.} The Hall product, Pfaffian datum, and Koszul source coalgebra, plus the imported Borcherds--Kac
target datum $\mathscr B_{\Delta_5}$.

\textit{Conditional datum.} For every cofinal saturated finite window
$W_\nu$, the recognition data (R0)--(R8); after these data are supplied
on a cofinal relation-closed system, the inverse-limit isomorphism
$P_X^{\Pi,\mathrm{red}}\cong\mathfrak g_{\Delta_5}$, and after the
Koszul source datum is also supplied, the chiral identification
$\Omega_E^{\mathrm{ch}}C_X\simeq\mathsf{Den}_{\Delta_5,E}$.

\textit{Construction.}
\begin{enumerate}
\item For each $\nu$, identify the saturated downward target-root window
$W_\nu=\{\beta\in\mathcal R_+\mid\operatorname{ht}\beta\le\nu\}$ and
the source window
$\Gamma(W_\nu)=\{\gamma\in\Gamma_{\mathrm{gram}}\mid
\alpha(\gamma)\in W_\nu\}$.  In coordinates
\(\beta=b_1\delta_1+b_2\delta_2+b_3\delta_3\), this means
\(\gamma_\beta=(b_1,b_1+b_2-b_3,b_2)\).  The packet
\(\texttt{certificates/lattice/downward\_root\_windows}\) verifies
\texttt{DOWNWARD\_ROOT\_WINDOWS\_VERIFIED}; it does not supply compact
source representatives or the recognition rows below.  The packet
\(\texttt{certificates/lattice/active\_root\_window\_exhaustion}\)
verifies \texttt{ACTIVE\_ROOT\_WINDOW\_EXHAUSTION\_VERIFIED}:
\(\bigcup_\nu W_\nu=\mathcal R_+\) by the height witness
\(\beta\in W_{\operatorname{ht}(\beta)}\).  This is target-root
exhaustion, not compact-source support exhaustion.  The packet
\(\texttt{certificates/sources/source\_window\_compact\_support}\)
verifies \texttt{SOURCE\_WINDOW\_COMPACT\_SUPPORT\_OBSTRUCTION\_VERIFIED}:
the compact support index set, source height function, finite source
windows, compact provenance rows, cofinality witness, transitions, and
target-source compatibility rows are still absent.  The packet
\(\texttt{certificates/sources/target\_source\_window\_compatibility}\)
verifies
\texttt{TARGET\_SOURCE\_WINDOW\_COMPATIBILITY\_OBSTRUCTION\_VERIFIED}:
source degree maps into \(\Gamma(W_\nu)\), target-window matching,
compact-provenance compatibility, cofinal-index compatibility,
transition compatibility, and zero compatibility-defect rows are still
absent.
\item Choose the simple-root representatives $e_i^X,f_i^X,h_i^X$ in
source degrees $\pm\gamma_i$, with
$\alpha(\gamma_i)=\alpha_i\in W_\nu$; verify parity assignments via
the GN multiplicity fibres.
\item Verify the Borcherds--Kac relations (Cartan, Chevalley, real
Serre, orthogonality, super sign) on every relation whose source and
target degrees lie in
$\Gamma(W_\nu)\cup(-\Gamma(W_\nu))\cup\{0\}$ and
$W_\nu\cup(-W_\nu)\cup\{0\}$, respectively.
\item Compute the full finite parity dimensions
$\dim P^{\Pi,\mathrm{red}}_{X,\gamma,\overline p}$ ($p=0,1$) for
$\gamma\in\Gamma(W_\nu)$; verify two-integer equality with
$\rootmult_{\overline p}^{\mathrm{GN}}(\alpha(\gamma))$.
\item Compute the Hopf pairing matrices in parity blocks; verify
kernel agreement with GN/Kac, coideal stability, and transition
compatibility.
\item Verify the no-extra-relations equality
$\ker\pi_\nu=
(\overline{J_{\mathrm{BK}}+\operatorname{Rad}_{\mathrm{GN}}})_{\le W_\nu}$.
\item Verify generation by Cartan and simple representatives;
compute PBW associated gradeds and verify strict transition
preservation.
\item Verify exact triangular completion: $\lim^1=0$ on the
triangular finite-window systems.
\end{enumerate}

\textit{Obstruction classes.}  $\mathfrak O_{\mathrm{Rec},\nu}$
(\S\ref{subsec:primitive-recognition}); for the chiral
identification, the finite Koszul comparison residual
\(\mathfrak O_{\mathrm{Kcmp},R}\) and the inverse-limit Koszul
Mittag--Leffler condition of \S\ref{subsec:koszul-source}.

\subsubsection*{Finite HN order}

At HN height \(R\) the finite datum has the following order:
\[
\begin{array}{rcl}
(\sigma,S,R) & \longmapsto &
F_{\sigma,S}(R),\widehat\Gamma_R\\
& \longmapsto &
\mathcal C_{\sigma,S,\le R},\Phi^{\mathrm{red}},o^{\mathrm{red}}\\
& \longmapsto &
\mathcal F_{X,\sigma,S,\le R}^{\mathrm{hyb}}\\
& \longmapsto &
\textup{(O1)},\textup{(O1)}^+\\
& \longmapsto &
m_R^{\mathrm{red}},\Delta_R^{\mathrm{ch}},
\langle\,,\,\rangle_R,\overline\Pi_{X,*}^\Theta\\
& \longmapsto &
\mathfrak D_{X,R}^{D_0},\mathscr L_{\mathrm{Pf},R},
\operatorname{pf}_{X,R},\textup{(O2)}\\
& \longmapsto &
C_{X,R},b_{X,R},\Theta_{\mathrm{Kos},R}\\
& \longmapsto &
\textup{(R0)\text{--}(R8)\ on\ }W_\nu
\end{array}
\]
The pro-object
\[
\mathfrak D_X^{\mathrm{DI}}
=\varprojlim_R\bigl(
\mathcal A_{X,R}^{E_3},
F_{X,R}^{\mathrm{hyb}},
\Gamma_{X,R},
\Pi_X,
\Phi_R,
o_R,
H_R,
P_R^\Pi,
C_{X,R},
\Theta_{\mathrm{Kos},R},
\mathscr L_{\mathrm{Pf},R},
\operatorname{pf}_{X,R},
\varepsilon_{o,R},
\mathrm{Rec}_R\bigr)
\]
is the inverse limit when every finite obstruction class vanishes
compatibly.

\subsubsection*{Finite-stage theorems}

\begin{theorem}[Finite-stage Dirac--Igusa pro-object;
\textit{conditional on the finite obstruction classes}]
\label{thm:finite-stage-dirac-igusa-pro-object}
Suppose the finite charge window, Hall atlas, hybrid carrier,
orientation transport, Hall product, Pfaffian datum, Koszul source
coalgebra, and primitive-recognition datum are defined at heights
\(R\in\mathcal R_{\mathrm{HN}}\), all finite obstruction classes
vanish, the HN transitions are strict, and the finite morphism proof
tables above have zero defects.  Then these data determine
a functor
\[
R\longmapsto \mathfrak D_{X,R}^{\mathrm{DI}}
\]
from \(\mathcal R_{\mathrm{HN}}^{\mathrm{op}}\) to the finite-stage
category of Definition~\ref{def:finite-stage-dirac-igusa-morphism}.
Equivalently, for every \(R'\le R\) there is a morphism
\(\rho_{RR'}\) satisfying \textup{(M1)}--\textup{(M8)}, and the
composition law
\(\rho_{R'R''}\rho_{RR'}=\rho_{RR''}\) holds for \(R''\le R'\le R\).
The inverse limit
\(\mathfrak D_X^{\mathrm{DI}}=\varprojlim_R
\mathfrak D_{X,R}^{\mathrm{DI}}\) is the resulting
Dirac--Igusa pro-object of
Definition~\ref{def:dirac-igusa-pro-isomorphism}.
\end{theorem}

\begin{theorem}[Pfaffian--Dirac identification at finite stage;
\textit{with local orientation and wall data at $R$}]
\label{thm:pfaffian-dirac-finite-stage}
Suppose the finite geometric Pfaffian datum, orientation transport,
type-II wall atlas, and primitive-recognition datum have vanishing
obstruction classes at height \(R\).  Then on the cusp expansion
\[
\operatorname{pf}_{X,R}|_{\mathfrak c_\infty}
=64q^{1/2}r^{1/2}s^{1/2}\prod_{\gamma\in\Gamma_R^{\Pi,+}\setminus\{0\}}
(1-x^\gamma)^{\sdim P^{\Pi,+}_{R,\gamma}},
\]
the Pfaffian section matches the Borcherds product on the test
window, with the orientation character
$\varepsilon_{o,R}|_{W^{(2),\le R}}=\nu_{\Delta_5}|_{W^{(2),\le R}}$
on the type-II reflection truncation.  The recognition map
$\iota_R^{\Prim}:H^0\Prim_{\mathrm{prot}}(\overline\Pi_{X,*}^\Theta
\mathcal F_{X,\sigma,S,\le R}^{\mathrm{hyb}})\to\mathfrak g_{\Delta_5,\le R}$
is an isomorphism on the saturated finite window.
\end{theorem}

\subsubsection*{Relative theorem on \(K3\times E\)}

The finite Hall atlas, hybrid carrier, and orientation transport at the test window use the retained D0-HN datum
through Theorem~\ref{thm:G-D0}; the orientation descent at (O1) by the retained
quotient-orientation datum and Theorem~\ref{thm:G-O1}; these data at
(O1)$^+$ by the chosen Weyl lifts and Theorem~\ref{thm:G-O1-plus};
the Pfaffian wall atlas at (O2) by the retained
\((O2_{\mathrm{atlas}})\) datum and
Theorem~\ref{thm:G-O2-orbit-transport}; the cofinal recognition
criterion beyond the base terminal and saturated target audits is the
unresolved compact-source recognition datum of
Chapter~\ref{ch:open-obstructions}.

The eight finite constructions give the Dirac--Igusa pro-object only when the
obstruction classes vanish.  Entry $(\mathrm L 8)$ names $(D0)$,
$(O1)$, $(O1)^+$, and \((O2_{\mathrm{atlas}})\), locates each in the
construction, and carries the retained data and the wall-atlas transport through
Appendix~\ref{app:obstruction-discharges} in
\S\ref{subsec:four-named-obstructions-ch5}.


\subsection{Four named obstructions}
\label{subsec:four-named-obstructions-ch5}

The eighth entry of the Dirac--Igusa pro-object is the obstruction
theory.  The cross-level identity ②$\leftrightarrow$③ between the
imported Borcherds--Kac--Moody denominator algebra and the compact
protected Pfaffian passes through four obstruction classes and the
finite geometric Pfaffian datum.  The first obstruction is reduced on
$K3\times E$ to the retained
D0-HN datum in
Appendix~\ref{app:obstruction-discharges}
(Theorem~\ref{thm:G-D0}); \((O1)\) and \((O1)^+\) are reduced there
to the retained quotient-orientation and Weyl-transport data
(Theorems~\ref{thm:G-O1}, \ref{thm:G-O1-plus}); the fourth is reduced
to the retained type-II
wall-atlas datum \((O2_{\mathrm{atlas}})\) by
Theorem~\ref{thm:G-O2-orbit-transport}.  The finite geometric Pfaffian
datum \((P_{\mathrm{fin}})\) of
Definition~\ref{def:finite-geometric-pfaffian-datum} supplies the
cosection-reduced skew perfect complexes, Pfaffian line, finite
Pfaffian section, type-II wall divisor order, automorphic
line-comparison, support and parity rows, leading coefficient, and
Pfaffian-line Mittag--Leffler compatibility for the cofinal HN product.
The principal
Pfaffian and orientation theorem
\[
\operatorname{Pf}_{\mathrm{prot}}(\mathfrak D_X^{\mathrm{DI}})
=\Delta_5,
\qquad
\varepsilon_o=\nu_{\Delta_5}
\]
holds relative to the first five components of the retained
obstruction datum
\(\mathfrak R_X^{\mathrm{DI}}\), namely
\((D0_{\mathrm{HN}},O1_{\mathrm{quot}},\mathcal T_W,
O2_{\mathrm{atlas}},P_{\mathrm{fin}})\).
The Lie-superalgebra recognition \(P_X\cong\mathfrak g_{\Delta_5}\)
requires the finite compact-source recognition datum \((S1)\)--\((S10)\)
of \S\ref{subsec:primitive-recognition}.
On \(K3\times E\) the D0 criterion, the quotient-orientation criteria,
and the conditional wall-atlas transport are stated in
Appendix~\ref{app:obstruction-discharges}; the compact-source
recognition datum is a separate hypothesis.
The retained obstruction data are the components of the tuple
\(\mathfrak R_X^{\mathrm{DI}}\) of
Definition~\ref{def:G-retained-obstruction-datum}.  Thus
``retained datum'' in this section means a projection of that tuple,
not a scalar normalization or a theorem label.

The fourth obstruction is the type-II wall-atlas datum; the
Appendix~G theorem transports that datum along the Weyl orbit and
compares its character with \(\nu_{\Delta_5}\).

\subsubsection*{(D0) D0-degeneration limit of the cosection-reduced
d-critical orientation theory}

\textit{Obstruction.}  At every finite HN height $R$ and every retained
charge stratum $\mathcal C\in\mathcal C_{\sigma,S,\le R}$, the
cosection-reduced d-critical structure on the K3-fibre direction must
extend through a retained flat \(D0\)-degeneration family to the
zero-dimensional sector on \(K3\times E\).  Equivalently, the datum
\((D0_{\mathrm{HN}})\) of
Definition~\ref{def:G-D0-HN-system} supplies the finite derived
degeneration family, the HN transition morphisms, the reduced
obstruction theory, the additive K3 semiregularity cosection
\(\operatorname{CS}_{\mathcal C,R}\), the reduced vanishing-cycle
specialization, the Joyce--Upmeier orientation specialization, the
Pfaffian-line and protected-integration specialization, and strict
Mittag--Leffler exactness.  The filtered Hall additivity equation
\[
q^*\operatorname{CS}_C
=p_1^*\operatorname{CS}_A+p_2^*\operatorname{CS}_B
\]
is part of this datum on extension stacks throughout the limit.
The local/wrapped boundary \(b_R^{\mathrm{geom}}\to0\) belongs to the
hybrid Ran carrier; it enters \((D0)\) only through compatibility of
the retained Hall correspondences with the zero-dimensional limit.

\textit{Finite site.}  The finite Hall atlas and hybrid carrier of
\S\ref{subsec:eight-finite-constructions}.  Failure of (D0) prevents
construction of the reduced Hall product $m_R^{\mathrm{red}}$, the
orientation gerbe descent of (O1), and the Pfaffian line of (D0-Pf)
on the wrapped sector.

\textit{Criterion.}  Criterion on $K3\times E$ in
Theorem~\ref{thm:G-D0}, relative to the retained D0-HN datum of
Definition~\ref{def:G-D0-HN-system}.  The
fibrewise K3 framework of \cite{JoyceSong,JoyceUpmeier} supplies the
orientation technology on the projection-finite stratum; the wrapped
stratum limit is part of the retained datum of
Appendix~\ref{app:G-D0-discharge}.
The finite D0 proof tables of
Definition~\ref{def:G-D0-HN-system} are the checkable form of this
datum.  A Hilbert-scheme scalar specialization, a K3 scalar test, or a
scalar trace row is not a D0-HN transition row.

\subsubsection*{(O1) Strong reduced orientation on the
$K3\times E$-quotient self-Ext frame bundle}

\textit{Obstruction.}  On the cosection-reduced moduli
$\mathfrak M^{\mathrm{red}}_{\mathcal C,R}$ at every finite HN height
$R$, the orientation line $o_{\mathcal C}$ has chosen square root
$o_{\mathcal C}^{\otimes 2}\simeq\det K^{\mathrm{red}}_{\mathcal C}$
\cite{JoyceUpmeier}, multiplicative under reduced extension
correspondences, with the quotient-descent equations: vanishing of the
reduced gerbe class
\[
\alpha^{\mathrm{red}}_{\mathcal C}
=w_2(\det K^{\mathrm{red}}_{\mathcal C})
+\mathrm{cs}^{*}w_2(\det\operatorname{coker}\mathrm{cs})
\in H^2(\mathfrak M^{\mathrm{red}}_{\mathcal C};\mathbb F_2),
\]
the free $E$-class
$\alpha^{E,\mathrm{free}}_{\mathcal C}=
a_1(\mathcal C)u_1+a_2(\mathcal C)u_2\in H^2(BE;\mathbb F_2)$,
and, for each finite-stabilizer stratum $S$ with stabilizer
$H\subset E_{\mathrm{tors}}$, the equivariant gerbe class
$\widetilde\beta^{H}_{\mathcal C,S}\in H^2_H(S;\mathbb F_2)$ together
with the choice of zero $H^1(BH;\mathbb F_2)$-linearization character
$\lambda^H_{\mathcal C,S}=0$.  The statement applies on every charge
stratum used by the Hall correspondences and is compatible with HN
transitions.

\textit{Finite site.}  The orientation descent datum of
\S\ref{subsec:eight-finite-constructions}; (O1) datum of
\S\ref{subsec:cosection-twisted-orientation}.  Failure of (O1)
prevents quotient-after-correspondence descent and obstructs the
reduced Hall product, the Pfaffian line on the reduced quotient,
and the orientation character.

\textit{Criterion.}  Theorem~\ref{thm:G-O1} is a
quotient-orientation criterion relative to the retained
\((O1_{\mathrm{quot}})\) datum.  The Klein-four condition at
$H=E[2]$ reduces the edge class to \(r_1=r_2=r_3=0\), but quotient
orientation also requires the zero degree-one linearization character.
For even $N\ge 4$ the full two-primary residual includes the mixed
coefficient \(A_{12}^{(N)}\), which cyclic restrictions do not detect.

\subsubsection*{(O1)$^+$ Weyl-equivariant transport along
$W^{(2)}(\Lambda_{II}^{2,1})$}

\textit{Obstruction.}  The orientation line of (O1) carries
Weyl-equivariant lifts of the three type-II wall transports
$\tau_i:o_{\mathcal C}\to s_{\delta_i}^* o_{\mathcal C}$
($i=1,2,3$) such that:
\begin{enumerate}
\item the $2$-cocycle $c_{o,R}\in Z^2(\mathcal G_R;
\underline{\mathbb F}_2)$ on the finite partial action groupoid
$\mathcal G_R$ has vanishing cohomology class
$[c_{o,R}]=0$ in
$H^2(W^{(2)}(\Lambda_{II}^{2,1});\mathbb F_2)$ on each complete
orbit, with chosen $1$-cochain killing $c_{o,R}$;
\item $\tau_i^2=1$ for $i=1,2,3$ and the Coxeter relations of
$W^{(2)}(\Lambda_{II}^{2,1})$ hold on the lifts;
\item the quotient-orientation cocycle $\mathfrak q_{\mathcal C}$
transports along $\tau_i$ between charge strata over Weyl-related
Gram degrees, with vanishing torsor defect
$\omega_{i,\mathcal C}=0$.
\end{enumerate}

\textit{Finite site.}  The Weyl-transport datum of
\S\ref{subsec:eight-finite-constructions}; (O1)$^+$ datum of
\S\ref{subsec:cosection-twisted-orientation}.  Failure of (O1)$^+$
obstructs the orientation character $\varepsilon_o$ and prevents
recognition of the Weyl orientation pattern as
$\nu_{\Delta_5}|_{W^{(2)}(\Lambda_{II}^{2,1})}$.

\textit{Conditional theorem.}  The equivariance part of
Theorem~\ref{thm:G-O1-plus} is conditional on the retained
quotient-orientation datum and chosen Weyl lifts.
The local rank-one Pfaffian crossing
supplies the local sign $\varepsilon^{\mathrm{loc}}_{\delta_i,R}=-1$
on a generic chart; the character comparison also uses the retained
\((O2_{\mathrm{atlas}})\) datum.  The global Coxeter coherence on
$W^{(2)}(\Lambda_{II}^{2,1})$ orbits is established by the
no-braid type-II Coxeter presentation and Maass-character match of
Appendix~\ref{app:G-O1-plus-discharge}.

\subsubsection*{(O2) Local Pfaffian wall normal form on type-II
walls}

\textit{Obstruction.}  At every finite HN height $R$ and every retained
component, boundary stratum, and chart overlap of the three simple
type-II walls $\Hpl_{\delta_i}$, $i=1,2,3$, the cosection-reduced
self-Ext complex splits as
\[
K^{\mathrm{red}}_{\delta_i,R}\simeq
K^{\mathrm{tan}}_{\delta_i,R}\oplus K^{\mathrm{nor}}_{\delta_i,R},
\]
with $s_{\delta_i}$-equivariant tangent factor and the rank-one
normal form
\[
K^{\mathrm{nor}}_{\delta_i,R}\simeq
[\R\xrightarrow{u_{\delta_i,R}}\R],
\qquad
N_{\delta_i,R}^{\mathrm{Pf}}=1,
\]
with normal Pfaffian unit
$\operatorname{Pf}(K^{\mathrm{nor}}_{\delta_i,R})
=\upsilon_{\delta_i,R}u_{\delta_i,R}$,
$s_{\delta_i}(\upsilon_{\delta_i,R})=\upsilon_{\delta_i,R}$,
$s_{\delta_i}(u_{\delta_i,R})=-u_{\delta_i,R}$.  The normal Pfaffian
units, tangent/normal splittings, quotient orientations, and
protected integration maps agree on chart overlaps.  The half-Hilbert
orbit sign sum of size $4096=2^{12}$ on the retained components and
boundary strata of the $W^{(2)}(\Lambda_{II}^{2,1})$-orbit reproduces
the reflection character $\nu_{\Delta_5}|_{W^{(2)}(\Lambda_{II}^{2,1})}$.

\textit{Finite site.}  The Pfaffian wall-atlas datum of
\S\ref{subsec:eight-finite-constructions}; rank-one Pfaffian crossing
of \S\ref{subsec:dirac-pfaffian}.  The finite proof surface is
\(\mathrm{wall\_objects}\), \(\mathrm{wall\_charts}\),
\(\mathrm{overlaps}\), \(\mathrm{orbit\_transport}\),
\(\mathrm{half\_hilbert\_orbit}\), \(\mathrm{compatibilities}\), and
\(\mathrm{scalar\_separation}\); these rows distinguish the rank-one
local sign, the global half-Hilbert atlas, and the scalar equality
\(2^{12}=64^2\).  Failure of (O2) prevents the
local Pfaffian section
$\operatorname{pf}_{\delta_i,R}=\upsilon_{\delta_i,R}u_{\delta_i,R}$
from descending to the orientation character and obstructs the
section-level identification
$\operatorname{pf}_X=\Delta_5$ on the wall complement.

\textit{Conditional theorem.}  Conditional on
\((O2_{\mathrm{atlas}})\) in Definition~\ref{def:G-O2-atlas},
Theorem~\ref{thm:G-O2-orbit-transport}.  The local rank-one crossing
model is established (Proposition~\ref{prop:rank-one-crossing-o2-sign});
component, boundary, and overlap compatibility on the full
half-Hilbert orbit is part of the retained wall-atlas datum and is
transported by Theorem~\ref{thm:G-O2-orbit-transport}.  The
$4096=2^{12}$ sign count is the scalar shadow of that retained datum,
as separated in Appendix~\ref{app:G-O2-discharge}.  The full
normal-form theorem with this sign count is relative to
Appendix~\ref{app:pfaffian-wall-normal-form}.

\subsubsection*{Joint theorem with recognition datum}

\begin{theorem}[Pfaffian--Dirac identification with primitive recognition datum]
\label{thm:pfaffian-dirac-conditional}
On \(K3\times E\), suppose (D0), the retained
\((O1_{\mathrm{quot}})\) datum, the \((O1)^+\) Weyl lifts, and
\((O2_{\mathrm{atlas}})\) hold at every finite HN height $R$ with
strict HN transitions as in
Appendix~\ref{app:obstruction-discharges}.  Suppose the finite
geometric Pfaffian datum \((P_{\mathrm{fin}})\) of
Definition~\ref{def:finite-geometric-pfaffian-datum} is supplied on a
cofinal HN system.  Suppose the recognition data (R0)--(R8) of
\S\ref{subsec:primitive-recognition} are supplied on a cofinal sequence
of saturated downward finite windows.  Suppose the chiral Koszul source
datum of \S\ref{subsec:koszul-source} is supplied, together with the
finite Koszul comparison datum of
Definition~\ref{def:finite-koszul-comparison-datum} and the strict
Koszul Mittag--Leffler exactness hypothesis.
Then
\[
\operatorname{Pf}_{\mathrm{prot}}(\mathfrak D_X^{\mathrm{DI}})
=\Delta_5,
\qquad
\varepsilon_o=\nu_{\Delta_5}|_{W^{(2)}(\Lambda_{II}^{2,1})},
\qquad
P_X\cong\mathfrak g_{\Delta_5},
\]
and $\Omega_E^{\mathrm{ch}}C_X\simeq\mathsf{Den}_{\Delta_5,E}$.
\end{theorem}

The eight finite constructions of
\S\ref{subsec:eight-finite-constructions}, the recognition theorem
of \S\ref{subsec:primitive-recognition}, and the Koszul comparison
of \S\ref{subsec:koszul-source} give the displayed implication.  The Pfaffian and orientation
conclusion holds by the retained D0-HN datum, the retained
\((O1_{\mathrm{quot}})\) datum, the chosen Weyl lifts, and the
retained \((O2_{\mathrm{atlas}})\) datum; the primitive
Lie-superalgebra and current-envelope conclusions require the cofinal
recognition datum, the Koszul source datum, and the finite Koszul
comparison datum.
The current-envelope conclusion uses the finite comparison
compatibility of
Proposition~\ref{prop:finite-koszul-comparison-compatibility}; without
that datum, the source cobar quasi-isomorphism need not preserve the
Weyl action, Pfaffian orientation character, Hall pairing, radical
quotient, or PBW filtration used in recognition.

\subsubsection*{The 4096 sign sum and the (O2) wall atlas}

The constant $4096=2^{12}$ is the size of the half-Hilbert
orbit of the type-II reflection group on the retained finite-window
components and boundary strata of the type-II wall configuration.
At the local rank-one chart, each retained component contributes a
sign $\varepsilon^{\mathrm{loc}}_{\delta_i,R}=-1$; the global sign
system is the retained (O2) wall atlas of Appendix~E transported along
the chosen Weyl lifts, and its character is
\(\nu_{\Delta_5}|_{W^{(2)}(\Lambda_{II}^{2,1})}\).
The Oberdieck--Pixton branch
$Z^X_{\mathrm{OP}}=-4096\,\Delta_5^{-2}$ is a scalar convention and is
not the source of the orientation character.  The equality
$\varepsilon_o=\nu_{\Delta_5}$ on
$W^{(2)}(\Lambda_{II}^{2,1})$ is relative to the retained quotient
orientation, the chosen Weyl lifts, and the (O2) wall atlas; the
local rank-one Pfaffian computation gives only the seed sign.

\subsubsection*{Compatibility with the companion volumes}

The four obstruction classes are source-side geometric data.  Their
automorphic shadow is smaller: the Gritsenko--Cl\'ery catalogue fixes
the Borcherds products, weights, characters, and denominator rows; it
does not record the \(K3\times E\) quotient orientation, the
Weyl-equivariant determinant-line lifts, or the Pfaffian wall atlas.
Thus (D0) is the cosection-reduced \(K3\)-fibre degeneration datum,
(O1) is the reduced orientation gerbe with finite-stabilizer descent,
(O1)\(^+\) is the Weyl-equivariant determinant-line transport, and
(O2) is the retained type-II Pfaffian wall atlas.  The
relation-closed window tests compatibility of these source data with
the target BKM relations; it is not an equality theorem between
geometric obstruction classes and automorphic catalogue rows.

The inverse limit on the HN system \(\mathcal R_{\mathrm{HN}}\)
contains the hybrid carrier, the cosection-twisted orientation, the
Hall correspondences, the protected Pfaffian, the Koszul source
coalgebra, the primitive recognition apparatus, the finite constructions,
and the retained orientation and wall-atlas data into
\[
\mathfrak D_X^{\mathrm{DI}}
=\varprojlim_R(\mathcal A_{X,R}^{E_3},
F_{X,R}^{\mathrm{hyb}},\dots,\mathrm{Rec}_R).
\]
The Pfaffian and orientation part of
\(\textcircled{2}\leftrightarrow\textcircled{3}\) is then the
retained-data theorem of Chapter~\ref{ch:pfaffian-dirac}; the
Lie-superalgebra recognition remains conditional on \((S1)\)--\((S10)\).


% Chapter 5 internally \input's chapters/05_subsections/05_1_hybrid.tex
% through chapters/05_subsections/05_8_obstructions.tex.

% ------------------------------------------------------------------
% Part III: Certified — the cross-level identity ②↔③
% ------------------------------------------------------------------
\part{Certified}

\chapter[The cross-level identity (2) <-> (3): Pfaffian-Dirac, orientation character, primitive recognition]{The cross-level identity \protect\textcircled{\scriptsize 2}\,\(\leftrightarrow\)\,\protect\textcircled{\scriptsize 3}: Pfaffian--Dirac, orientation character, primitive recognition}
\label{ch:pfaffian-dirac}

The Igusa cusp form \(\Delta_5\) of weight five on
\(\Sp(2,\Z)\backslash\mathcal H_2\) carries three names, each rooted in a
distinct categorical-dimensional level:

\begin{itemize}
\item \(\mathcal D_X=\Delta_5\) — the Borcherds--Gritsenko--Nikulin section of the
  weight-\(5\) automorphic line over the genus-two Siegel space (View
  \textcircled{\scriptsize 1});
\item \(\operatorname{den}(\mathfrak g_{\Delta_5})=64^{-1}\Delta_5(2Z)\) —
  the denominator of the generalized Borcherds--Kac--Moody (BKM)
  superalgebra \(\mathfrak g_{\Delta_5}\) on
  \(\La^{2,1}_{II}\) (View \textcircled{\scriptsize 2});
\item \(\operatorname{Pf}_{\mathrm{prot}}(\mathfrak D_X^{\mathrm{DI}})\) —
  the protected Pfaffian of the Dirac--Igusa pro-object on
  \(X=K3\times E\) (View \textcircled{\scriptsize 3}).
\end{itemize}

The cross-level identity
\textcircled{\scriptsize 1}\(\leftrightarrow\)\textcircled{\scriptsize 2}
is Borcherds 1995 and Gritsenko--Nikulin 1998.  The edge
\textcircled{\scriptsize 2}\(\leftrightarrow\)\textcircled{\scriptsize 3}
has two pieces.  The four obstruction classes and the finite geometric
Pfaffian datum \((P_{\mathrm{fin}})\) give the Pfaffian; the
orientation classes give the orientation character on \(K3\times E\).
The Lie-superalgebra recognition has the additional finite
compact-source datum \((S1)\)--\((S10)\) and the finite Koszul
comparison datum of
Definition~\ref{def:finite-koszul-comparison-datum}.  The three
statements have separate hypotheses:

\begin{enumerate}
\item the \emph{Pfaffian--Dirac theorem}: under (D0), (O1), (O1)\(^+\),
  (O2), and the finite geometric Pfaffian datum
  \((P_{\mathrm{fin}})\) on a cofinal HN system,
  \[
    \operatorname{Pf}_{\mathrm{prot}}(\mathfrak D_X^{\mathrm{DI}})
    \;=\;\Delta_5
  \]
  as sections of \(\mathcal L^5\otimes\nu_{\Delta_5}\);
\item the \emph{orientation character identity}: under (O1)+(O1)\(^+\)+(O2),
  \[
    \epsilon_o\;=\;\nu_{\Delta_5}
    \quad\text{on}\quad W^{(2)}(\La^{2,1}_{II});
  \]
\item the \emph{primitive recognition theorem}: under the finite
  compact-source recognition datum (Cartan, simple representatives,
  Chevalley, real Serre, imaginary orthogonality, generation,
  completed PBW, Hopf radical equality, no-extra-relations, full parity
  dimensions),
  \[
    H^0\Prim_{\mathrm{prot}}\!\bigl(
    \overline\Pi_{X,*}^{\Theta}\mathcal F_X\bigr)
    \big/\overline{\operatorname{Rad}\langle\,\cdot\,,\,\cdot\,\rangle_X}
    \;\cong\;\mathfrak g_{\Delta_5}
  \]
  after passage to the inverse limit over the cofinal root windows, with
  the chiral cobar
  \(\Omega_E^{\mathrm{ch}}C_X\) identified as the formal current envelope
  \(\mathsf{Den}_{\Delta_5,E}\) only after the chiral Koszul source
  datum is also fixed.
\end{enumerate}

The theorem-dependency ledger for these three assertions is
\[
\mathrm{beilinson\_levels},\quad
\mathrm{objects},\quad
\mathrm{morphisms},\quad
\mathrm{theorem\_dependencies},\quad
\mathrm{dependency\_edges},\quad
\mathrm{forbidden\_edges},\quad
\mathrm{status\_separation},\quad
\mathrm{level\_equalities},\quad
\mathrm{equality\_sites},\quad
\mathrm{definition\_separation},\quad
\mathrm{obstruction\_independence},\quad
\mathrm{transitions},\quad
\mathrm{scalar\_firewall}.
\]
It is the finite graph recording which retained packets enter each
proof and which edges are forbidden.  In particular, the Pfaffian and
orientation proofs do not use the mirror conjecture, the gravity-line
residual, the scalar trace, or the OP scalar branch; primitive
recognition does not use the scalar trace, the denominator product, a
target counit, or target labels as source representatives.  Claim
statuses do not collapse: imported BKM and automorphic identities,
retained-data Pfaffian identities, conditional recognition, mirror
conjectures, scalar trace identities, and open gravity-line residuals
occupy distinct rows.
The definition-separation rows assert non-circularity at the level of
finite data: the Pfaffian line is not defined by the automorphic
section or by its desired value, the orientation character is not
defined by scalar signs, the chiral Koszul source is not the target
counit, the primitive source is not defined by target multiplicities,
and the Dirac--Igusa pro-object is not defined by a scalar shadow.
The obstruction-independence rows assert that \((D0)\), \((O1)\),
\((O1)^+\), \((O2_{\mathrm{atlas}})\), \(P_{\mathrm{fin}}\), the
Koszul source comparison, and primitive recognition are separate finite
data; typing dependencies are not logical implications between
vanishing statements.
The certificate
\(\texttt{certificates/dependencies/k3e\_theorem\_dependencies}\)
is schema-complete only: its verifier returns
\(\texttt{SCHEMA\_COMPLETE\_SCHEMA\_ONLY}\) and records
\(\texttt{dependency\_certification=false}\) and
\(\texttt{mathematical\_certification=false}\).  It verifies that the
finite ledger is typed, populated, separated by status, and guarded by
forbidden-edge, definition-separation, and obstruction-independence
rows; it does not prove the retained Pfaffian,
orientation, Koszul, recognition, mirror, scalar-trace, or gravity-line
claims.

\begin{proposition}[Hypothesis-deletion residues]
\label{prop:ch6-hypothesis-deletion-residues}
The theorem lane has the following exact residues when one named datum
is omitted.
\begin{enumerate}
\item Without the finite compact-source recognition datum
\((S1)\)--\((S10)\), the Pfaffian--Dirac section identity and the
orientation character identity remain under their own hypotheses, but
there is no theorem identifying
\[
P_X^{\Pi,\mathrm{red}}\quad\text{with}\quad\mathfrak g_{\Delta_5}.
\]
The surviving primitive object is only the retained Hall primitive
object with its supplied product, coproduct, pairing, radical, and PBW
rows; denominator exponents or signed multiplicities do not determine
its Lie bracket.

\item Without the chiral Koszul source datum, the finite Koszul
comparison datum, and strict Koszul Mittag--Leffler exactness, the
source coalgebra \(C_X\) is not identified with the target current
envelope.  The theorem no longer supplies a quasi-isomorphism
\[
\Omega_E^{\mathrm{ch}} C_X
\longrightarrow
\mathsf{Den}_{\Delta_5,E}
\]
with acyclic cone.  The exact residue is the compact
Hall primitive source, if supplied, before the source-to-target chiral
comparison.

\item Without the finite geometric Pfaffian datum
\((P_{\mathrm{fin}})\), the obstruction data may still define
orientation and local wall signs, but there is no global object
\[
\operatorname{pf}_{X,R}\in H^0(\mathfrak Q_R,\mathscr L_{\mathrm{Pf},R})
\]
whose inverse limit can be compared with \(\mathcal D_X=\Delta_5\).
The residue is a reduced orientation/wall datum, not a Pfaffian section
identity.

\item Without \((O2_{\mathrm{atlas}})\), the retained D0-HN datum,
quotient orientation, and Weyl-lift torsor do not compute the local
rank-one Pfaffian signs.  The residue is a Weyl-transported
orientation line on the retained quotient, but no theorem gives
\(\epsilon_o(s_{\delta_i})=-1\), no divisor order along the type-II
walls is computed, and the Pfaffian section is not identified with
\(\Delta_5\).

\item Without \((O1)^+\), the quotient orientation may exist, but it is
not transported by a coherent action of
\(W^{(2)}(\Lambda_{II}^{2,1})\).  The residue is an oriented reduced
quotient on the retained finite stages; there is no global character
\[
\epsilon_o:W^{(2)}(\Lambda_{II}^{2,1})\longrightarrow\{\pm1\}.
\]

\item Without \((O1)\), the quotient square root and reduced
orientation line on the \(K3\times E\)-quotient self-Ext frame bundle
are absent.  The residue is the finite d-critical/cosection and
D0-specialization data before orientation; Hall products and Pfaffian
sections are unsigned and do not define the protected oriented
Pfaffian.

\item Without \((D0)\), the cosection-reduced orientation theory has no
D0-degeneration endpoint on the cofinal HN pro-system.  The residue is
the formal Mukai--Gram lattice, the imported automorphic/BKM/theta
triangle, and any finite stage objects supplied independently; it is
not a compact \(K3\times E\) Pfaffian--Dirac theorem.
\end{enumerate}
Each residue is a positive statement about the surviving object.  None
of the rows permits a scalar trace, an automorphic target, or a
denominator product to replace the deleted datum.
\end{proposition}

\begin{proof}
The conclusion is a direct reading of the source and target of each
comparison.  Primitive recognition needs the source representatives,
relations, radical, pairing, parity, and PBW rows; without them there
is no Lie-superalgebra isomorphism.  The Koszul comparison is the only
map from the source coalgebra to the target current envelope.  The
finite geometric Pfaffian datum supplies the Pfaffian line and section,
so deleting it removes the domain of the section identity.  The
\((O2_{\mathrm{atlas}})\) datum supplies the local rank-one signs at
the type-II walls; deleting it removes the generator values of the
orientation character.  The \((O1)^+\) datum supplies coherent Weyl
transport; deleting it removes the homomorphism from the Weyl group.
The \((O1)\) datum supplies the quotient orientation line; deleting it
removes the oriented Pfaffian input.  The \((D0)\) datum supplies the
D0-degeneration endpoint and transition-compatible reduced orientation
theory; deleting it removes the compact pro-limit in which the
Pfaffian--Dirac comparison lives.  These are object-level failures in
the Beilinson chain, not failures of notation.
\end{proof}

\begin{proposition}[Implication and non-redundancy of obstruction rows]
\label{prop:ch6-obstruction-independence}
Let
\[
\mathsf{Obs}^{\mathrm{DI}}_{\mathrm{fin}}
=
\mathsf{D0}\times\mathsf{O1}\times\mathsf{O1}^+
\times\mathsf{O2}\times\mathsf{Pfin}\times
\mathsf{Kos}\times\mathsf{Rec}
\]
be the finite theorem-dependency product whose seven factors are the
row categories for \(D0_{\mathrm{HN}}\), \(O1_{\mathrm{quot}}\),
\(\mathcal T_W\), \(O2_{\mathrm{atlas}}\), \(P_{\mathrm{fin}}\), the
chiral Koszul comparison datum, and the compact-source primitive
recognition datum.  Each vanishing predicate is a predicate on its own
factor.  The only implication used by the Pfaffian--Dirac lane is the
conjunction of the relevant vanishing predicates into the corresponding
theorem:
\[
D0\wedge O1\wedge O1^+\wedge O2\wedge P_{\mathrm{fin}}
\Longrightarrow
\text{Pfaffian section and orientation character},
\]
and, separately,
\[
\text{compact source}\wedge\text{Koszul comparison}\wedge
\text{primitive recognition}
\Longrightarrow
P_X^{\Pi,\mathrm{red}}\cong\mathfrak g_{\Delta_5}.
\]
There is no theorem-level implication from any one named vanishing
predicate to any other.  The type dependencies recorded in
\(\texttt{obstruction\_independence.csv}\) say only that a later datum
is not meaningful until its earlier input object exists; they do not
prove the earlier datum.

For each of the seven rows, deleting that row and leaving the other row
slots formally zero gives a finite theorem-dependency witness to
non-redundancy.  The residual theorem in the deleted row is exactly the
corresponding item of
Proposition~\ref{prop:ch6-hypothesis-deletion-residues}.  These are
counterexamples to illegitimate implications inside the finite
dependency calculus.  They are not asserted to be geometric
\(K3\times E\) families with all other analytic hypotheses realized.
\end{proposition}

\begin{proof}
In the product category
\(\mathsf{Obs}^{\mathrm{DI}}_{\mathrm{fin}}\), a vanishing predicate is
pulled back from one projection.  A predicate pulled back from one
projection cannot force a predicate pulled back from another projection
unless such an implication is added as a row of the dependency ledger.
The table
\(\texttt{obstruction\_independence.csv}\) requires
\(\texttt{logical\_implications=none}\) for the seven named data and
records the type dependencies separately.  Thus the row \(O1^+\) being
typed over \(O1\) means that Weyl transport needs an orientation line;
it is not a proof that the orientation line exists.  The same applies
to \(O2_{\mathrm{atlas}}\), \(P_{\mathrm{fin}}\), the Koszul
comparison, and primitive recognition.

Deleting one row removes the source, target, or comparison named in
Proposition~\ref{prop:ch6-hypothesis-deletion-residues}, while leaving
the remaining row slots available as formal data.  This supplies the
finite-row countermodel to every attempted proof that the deleted row
is redundant.  Since each deleted row removes a different object or
comparison map, no two rows have the same residue; the obstruction
classes are non-redundant in the theorem-dependency calculus.
\end{proof}

The eight finite constructions of \S5.7 define, at every finite stage \(R\)
in the Harder--Narasimhan inverse system, the data \(\mathcal A_{X,R}^{E_3}\),
\(F_{X,R}^{\mathrm{hyb}}\), \(\Gamma_{X,R}\), \(\Pi_X\), \(\Phi_R\),
\(o_R\), \(H_R\), \(P_R^\Pi\), \(C_{X,R}\), \(\Theta_{\mathrm{Kos},R}\),
\(\mathscr L_{\mathrm{Pf},R}\), \(\operatorname{pf}_{X,R}\),
\(\varepsilon_{o,R}\), and \(\mathrm{Rec}_R\) of the pro-object
\(\mathfrak D_X^{\mathrm{DI}}\); the four obstruction classes of \S5.8
are exactly (D0), (O1), (O1)\(^+\), (O2).  These four classes act on
the Pfaffian line and the orientation character; the assertion that
\(\mathrm{Rec}_R\) is eventually an isomorphism is the separate
compact-source recognition datum.  The relevant geometric data are the
open factorization chart of the vacuum \(b\), the protected Pfaffian
functor, the Hall--Drinfeld correspondence, the chosen
Sp\(^{\mathrm{ch}}\) specialization, the pro-object ambient, the
D0-degeneration endpoint, the Bussi--Brav--Dupont--Joyce--Szendr\H{o}i
orientation, the non-degeneracy of the Pfaffian line, and the type-II
wall normal form.

The Pfaffian and
orientation theorems are retained-data theorems on \(K3\times E\), as in
Appendix~\ref{app:obstruction-discharges}; the primitive recognition
theorem is conditional on the compact-source datum
\((S1)\)--\((S10)\).  The BKM denominator identity (View
\textcircled{\scriptsize 2}) is
\emph{proved}, imported from
\cite{Borcherds95,GN,GNII,BorcherdsGKM88,KAC:1}; the automorphic
identification (View \textcircled{\scriptsize 1}) is \emph{proved},
imported from \cite{Igusa1962,Igusa1964II}; the operator-level
Pfaffian datum \(\mathfrak D_X^{\mathrm{DI}}\) is \emph{proved on
\(K3\times E\) relative to the retained D0-HN datum, the retained
quotient-orientation datum, the chosen Weyl lifts,
\((O2_{\mathrm{atlas}})\), and the retained finite geometric Pfaffian datum
\((P_{\mathrm{fin}})\)}
of
Appendix~\ref{app:obstruction-discharges}; the chiral cobar
identification is \emph{conditional} on the Koszul source datum of
Definition~\ref{def:chiral-koszul-source-datum}.  The
\(\kappa\)-tuple of the Calabi--Yau quantum-groups companion,
\(
(\kappa_{\mathrm{cat}},\kappa_{\mathrm{ch}}^{\mathrm{Hodge}},
\kappa_{\mathrm{ch}}^{\mathrm{Heis}},\kappa_{\mathrm{BKM}},
\kappa_{\mathrm{fiber}})=(0,0,3,5,24),
\)
fixes the level of \(\Delta_5\): the cusp-form weight is
\(\kappa_{\mathrm{BKM}}=5\) and the elliptic genus weight is
\(\chi_{\mathrm{top}}(K3)=\kappa_{\mathrm{fiber}}=24\); the additive
formula \(\kappa_{\mathrm{BKM}}=\kappa_{\mathrm{ch}}+\chi(\mathcal O_{\mathrm{fiber}})\)
has no invariant meaning on the \(\kappa\)-tuple.


\section{Data and Obstructions}
\label{sec:ch6-data}

\subsection{The pro-object \(\mathfrak D_X^{\mathrm{DI}}\) and its eight finite-stage components}
\label{subsec:ch6-pro-object}

Three names of \(\Delta_5\): the section \(\mathcal D_X=\Delta_5\) (View
\textcircled{\scriptsize 1}); the BKM denominator
\(\operatorname{den}(\mathfrak g_{\Delta_5})=64^{-1}\Delta_5(2Z)\) (View
\textcircled{\scriptsize 2}); the protected Pfaffian
\(\operatorname{Pf}_{\mathrm{prot}}(\mathfrak D_X^{\mathrm{DI}})\) (View
\textcircled{\scriptsize 3}, relative to retained data).  The handle
``\(\Delta_5\)'' is reserved for the working symbol; theorem-level
sentences specify the named view.

The Dirac--Igusa pro-object is
\[
\mathfrak D_X^{\mathrm{DI}}
\;=\;\lim_R\bigl(
\mathcal A_{X,R}^{E_3},\,
F_{X,R}^{\mathrm{hyb}},\,
\Gamma_{X,R},\,
\Pi_X,\,
\Phi_R,\,
o_R,\,
H_R,\,
P_R^\Pi,\,
C_{X,R},\,
\Theta_{\mathrm{Kos},R},\,
\mathscr L_{\mathrm{Pf},R},\,
\operatorname{pf}_{X,R},\,
\varepsilon_{o,R},\,
\mathrm{Rec}_R\bigr).
\]
The components are the values of the eight finite constructions of \S5.7,
with the ambient category the cofinal HN inverse system of
Definition~\ref{def:normal-ordered-charge-lattice-appendixC}; \(R\)
ranges over a chosen cofinal subsystem
\(\mathcal R_{\mathrm{HN}}\subset\Lambda_{\mathrm{ample}}\).  The
transition maps are the finite-stage morphisms of
Definition~\ref{def:finite-stage-dirac-igusa-morphism}; isomorphism of
two such pro-objects means pro-isomorphism in the sense of
Definition~\ref{def:dirac-igusa-pro-isomorphism}.  At finite stage,
these transition and pro-isomorphism claims are the
\(\mathrm{cofinal\_subsystems}\), \(\mathrm{stage\_objects}\),
\(\mathrm{stage\_morphisms}\), \(\mathrm{component\_maps}\),
\(\mathrm{composition\_laws}\), \(\mathrm{ml\_exactness}\),
\(\mathrm{pro\_isomorphisms}\), and \(\mathrm{scalar\_firewall}\)
tables of \S\ref{subsec:eight-finite-constructions}.  The label
\(\gamma\) is the pro-object ambient.
The universal property of this pro-object is the Hom formula of
Proposition~\ref{prop:dirac-igusa-pro-object-universal-property}; in
particular, changing the finite HN presentation by a cofinal subsystem
does not change \(\mathfrak D_X^{\mathrm{DI}}\) in the pro-category.

\begin{remark}[The three \(\Delta_5\)-names]
\label{rem:ch6-three-delta-names}
The equality
\(\operatorname{Pf}_{\mathrm{prot}}(\mathfrak D_X^{\mathrm{DI}})=\mathcal D_X\) is a
cross-level identity, not a tautology.  The right-hand side is an
automorphic section; the left-hand side is the inverse limit of finite
Pfaffian sections \(\operatorname{pf}_{X,R}\) of the line bundles
\(\mathscr L_{\mathrm{Pf},R}\) on the reduced quotient stacks
\(\mathfrak Q_R\); the two are related by the fixed comparison
\(\iota_{\mathrm{aut}}:\mathscr L_{\mathrm{Pf},X}\xrightarrow{\sim}
\mathcal L^5\otimes\nu_{\Delta_5}\).  The denominator name
\(\operatorname{den}(\mathfrak g_{\Delta_5})=64^{-1}\Delta_5(2Z)\) is
the same global section \(\mathcal D_X\) written in the type-II chamber
polarization that exposes the Weyl--Kac--Borcherds product expansion;
the prefactor \(64^{-1}\) and the doubling \(2Z\) of the modular
variable enter from the theta normalization of the Borcherds product
on the index-$4$ sublattice $\Lambda^{2,1}_{II}\subset\Lambda^{2,1}$.
The source side \(P_X^{\Pi,\mathrm{red}}\) identifies with
\(\mathfrak g_{\Delta_5}\) only after the recognition criterion of
\S\ref{sec:ch6-primitive-recognition} is established.
\end{remark}

\subsection{The four obstruction classes: (D0), (O1), (O1)\(^+\), (O2)}
\label{subsec:ch6-obstructions}

The four classes match the eight finite constructions one-to-one only on the
orientation lane.  They are the obstruction data of
Theorem~\ref{thm:pfaffian-dirac-finite-stage}.  On \(K3\times E\), Appendix~\ref{app:obstruction-discharges}
reduces \((D0)\) to the retained D0-HN datum by
Theorem~\ref{thm:G-D0}; reduces \((O1)\) to the
retained quotient-orientation datum by Theorem~\ref{thm:G-O1}; reduces
\((O1)^+\) to the chosen Weyl lifts by Theorem~\ref{thm:G-O1-plus};
and reduces \((O2)\) to the retained type-II wall-atlas datum
\((O2_{\mathrm{atlas}})\) by Theorem~\ref{thm:G-O2-orbit-transport}.  The
\((R1)\)--\((R2)\) scalar separation for the type-II middle-wall
constant \(4096=2^{12}=64^2\) inside \((O2)\) is
Theorem~\ref{thm:G-R1-R2}.
The finite D0 packet
\(\texttt{certificates/d0/k3e\_d0\_hn}\) is presently verified only
through its obstruction ledger:
\(\texttt{D0\_OBSTRUCTION\_LEDGER\_VERIFIED}\), with
\(\texttt{d0\_certification=false}\).  Hence the D0-dependent
assertions in this chapter are relative theorems until the flat
degeneration, retained-substack, HN-transition, derived-enhancement,
semiregularity-cosection, vanishing-cycle, orientation, Pfaffian,
protected-integration, and Mittag--Leffler rows are supplied.
The finite geometric Pfaffian packet
\(\texttt{certificates/pfaffian/k3e\_finite\_pfaffian}\) is likewise
verified only through its obstruction ledger:
\(\texttt{PFAFFIAN\_OBSTRUCTION\_LEDGER\_VERIFIED}\), with
\(\texttt{pfaffian\_certification=false}\).  Hence the protected
Pfaffian identity below is a relative section identity until the skew
perfect complexes, Pfaffian lines, sections, wall charts,
automorphic-comparison rows, and Pfaffian-line transition rows are
constructed.
The finite O2 wall-atlas packet
\(\texttt{certificates/wall\_atlas/k3e\_o2\_atlas}\) is also verified
only through its obstruction ledger:
\(\texttt{O2\_OBSTRUCTION\_LEDGER\_VERIFIED}\), with
\(\texttt{o2\_certification=false}\).  Hence every assertion using
\((O2_{\mathrm{atlas}})\) is relative until the simple-wall coverage,
wall-object, wall-chart, overlap, Weyl-transport, half-Hilbert-orbit,
compatibility, type-I-exclusion, D0-limit, and scalar-separation rows
are supplied.
The finite orientation packet
\(\texttt{certificates/orientation/k3e\_reduced\_orientation}\)
is presently verified only through its obstruction ledger:
\(\texttt{ORIENTATION\_OBSTRUCTION\_LEDGER\_VERIFIED}\), with
\(\texttt{orientation\_certification=false}\).  Thus the
orientation-dependent assertions in this chapter are relative
theorems until the square-root, quotient Cech--Borel,
finite-stabilizer, Weyl-lift, Coxeter, and transition rows are supplied.

\paragraph{(D0) — D0-degeneration endpoint of the cosection-reduced
d-critical orientation theory.}  At every
\(R\), the retained \(D0_{\mathrm{HN}}\) datum of
Definition~\ref{def:G-D0-HN-system} is fixed: a flat
D0-degeneration family, retained finite-type object/extension/flag
stacks, proper or closed HN transition morphisms, quasi-smooth
d-critical derived enhancements, surjective additive K3 semiregularity
cosections, reduced vanishing-cycle specialization, Joyce--Upmeier
orientation specialization, Pfaffian-line and protected-integration
specialization, and strict Mittag--Leffler exactness.  Compatibly in
\(R\), these clauses give the HN-completed state object and the finite
first-order block \(\mathfrak D_X\).
The finite first-order block also carries the exponent-matching clause
\[
\sdim P^\Pi_{X,(n,l,m)}=f(nm,l)
\qquad ((n,l,m)\in\Gamma_{\mathrm{eff}}),
\]
where \(f(n,l)\) is the Fourier coefficient of the weak Jacobi form
\(\phi_{0,1}\) of weight zero and index one.  This clause is not a
consequence of the reduced orientation alone; it is part of the
retained finite-stage Dirac block.  The cosection contraction
of the Pantev--Toen--Vaquie--Vezzosi (PTVV)
\((-1)\)-shifted symplectic form by the K3 semiregularity cosection
\(\sigma_{K3}\) is the Maulik--Toda / Oberdieck reduced obstruction theory
\cite{MaulikTodaGV,OberdieckReducedPT}; this absorbs the parity flip
\(\mathrm{vd}^{\mathrm{red}}=\mathrm{vd}^{\mathrm{std}}+1\) into the
reduced orientation line as the \(\Z/2\)-grading shift
\(o^{\mathrm{red}}_{\mathcal C}=o^{\mathrm{std}}_{\mathcal C}\otimes\Pi\).

\paragraph{(O1) — strong multiplicative reduced orientation datum on the
reduced \(K3\times E\)-quotient self-Ext frame bundle.}  A Joyce--Upmeier-type orientation line
\(o_{\mathcal C}\in\operatorname{Pic}(\mathfrak M^{or}_X)\) with
\(o_{\mathcal C}^{\otimes 2}\simeq
\det\operatorname{RHom}_{\mathrm{red}}(\mathcal C,\mathcal C)\),
extension-multiplicative under
\(o_{\mathcal C\oplus\mathcal C'}\simeq
o_{\mathcal C}\otimes o_{\mathcal C'}\), and equipped with the full
quotient-orientation datum
\[
\theta_{\mathcal C}=
\bigl(\theta^{\mathrm{red}}_{\mathcal C},\theta^{E,\mathrm{free}}_{\mathcal C},
\{(\theta^{H}_{\mathcal C},\lambda^{H}_{\mathcal C}=0)\}_{H\in\mathcal H(\mathcal C)}\bigr)
\]
on every retained charge stratum.  The vanishings
\(\alpha^{\mathrm{red}}_{\mathcal C}=0\),
\(\alpha^{E,\mathrm{free}}_{\mathcal C}=0\),
\(\widetilde\beta^{H}_{\mathcal C}=0\),
\(\lambda^{H}_{\mathcal C}=0\) are the orientation-line, the Borel-equivariant
free-action, the residual finite-stabilizer Borel, and the residual
finite-stabilizer linearization conditions; for \(N=2\), the Klein-four
calculation
\(\beta^{E,E[2]}_{\mathcal C}=b_{20}x_1^2+b_{11}x_1x_2+b_{02}x_2^2\)
specializes the residual.  For \(2^a\parallel N\), \(a\ge2\), the
rank-two coefficient \(A_{12}^{(N)}x_1x_2\) is not detected by cyclic
restrictions and is part of the finite-stabilizer orientation test of
Proposition~\ref{prop:finite-stabilizer-orientation-detection}.  The
orientation lift technology is
BBDJS \cite{BBDJS2015} on the unreduced side and Joyce--Upmeier
\cite{JoyceUpmeier,JoyceTanakaUpmeier} on the multiplicative side; the
reduced \(K3\times E\)-quotient theorem is separate.

\paragraph{(O1)\(^+\) — Weyl-equivariant determinant-line transport along
\(W^{(2)}(\La^{2,1}_{II})\) reflections.}
The chosen lifts
\(\tau_i:o_{\mathcal C}\to s_{\delta_i}^{*}o_{\mathcal C}\), \(i=1,2,3\),
of the three type-II wall transports satisfy
\([c_o]=0\) in
\(H^2(W^{(2)}(\La^{2,1}_{II});\mathbb F_2)\) (equivalently, in the
finite action-groupoid cohomology
\(H^2(\mathcal G_R;\underline{\mathbb F}_2)\) at every \(R\)),
\(\widetilde\tau_i^{\,2}=1\), and the quotient-orientation transport
defects \(\omega_{i,\mathcal C}=0\) for every simple reflection.  Then
\(\widetilde\tau_i\) define an honest Weyl action on the reduced quotient
orientation; the passage from the three generator signs to a character on
\(W^{(2)}(\La^{2,1}_{II})\) is the group-theoretic part of
Definition~\ref{def:O1-plus-weyl-transport}.

\paragraph{(O2) — Pfaffian local model and wall atlas on type-II walls.}
At a generic point of each
\(\Hpl_{\delta_i}\), the chosen real local equation
\(u_i=u_{\delta_i}\) with \(s_{\delta_i}(u_i)=-u_i\), the normal
determinant-complex splitting
\[
K^{\mathrm{nor}}_{\delta_i,R}\simeq
\bigl[\R\xrightarrow{u_i}\R\bigr],
\qquad
\operatorname{Pf}(K^{\mathrm{nor}}_{\delta_i,R})=\upsilon_{\delta_i,R}u_i,
\qquad
s_{\delta_i}(\upsilon_{\delta_i,R})=\upsilon_{\delta_i,R},
\]
and the rank-one normal-form rank \(N_{\delta_i}^{\mathrm{Pf}}=1\) hold.
The wall charts cover every retained component and boundary stratum of
\(\Hpl_{\delta_i}\); on overlaps the tangent/normal splittings, normal
Pfaffian units, quotient orientations, and protected integration maps
agree after a zero \v Cech sign class.  The retained atlas datum must
contain, for the middle wall \(\delta_2\leftrightarrow(0,1,1)\), a
wrapped wall object
\(x_{\delta_2,R}\in\mathcal M_{\eta,R}^{\mathrm{wr,rig}}\) with
\(b_R^{\mathrm{geom}}(\eta)>0\) and \(\overline\Pi_X(\eta)=(0,1,1)\).
Proposition~\ref{prop:appE-local-pfaffian-sign} supplies only the
rank-one local Pfaffian sign on such a chart; the component, boundary,
overlap, quotient, and protected-integration compatibilities are part of
\((O2_{\mathrm{atlas}})\).

\subsection{The OP minus sign is a scalar branch convention}
\label{subsec:ch6-op-scalar-branch}

The Oberdieck--Pixton (OP) signed scalar is
\(C_{\mathrm{OP},X}=-4096\).  Its absolute square-normalization is
\(C_{\square,X}=4096=64^2\), and its sign is the OP scalar branch
recorded by the fermionic shift \((-p_{\mathrm{DT}})^n\) of the
OP/DT chamber.  It is not the source of the orientation character
\(\epsilon_o\); orientation lines are killed by squaring.  The
orientation character is computed by monodromy
of the Pfaffian section around the type-II walls (Lemma
\ref{prop:appE-local-pfaffian-sign}); equivalently, by the
\((-1)^{N_{\delta_i}^{\mathrm{Pf}}}\) values from (O2).

\section{The Pfaffian--Dirac theorem}
\label{sec:ch6-pfaffian-dirac}

The protected Pfaffian of the Dirac--Igusa pro-object is compared with
the automorphic section.

\begin{theorem}[Pfaffian--Dirac section identity]
\label{thm:ch6-pfaffian-dirac}
Granted the four obstruction classes of
Section~\ref{subsec:ch6-obstructions}, with \((O2)\) understood as the
retained wall-atlas datum \((O2_{\mathrm{atlas}})\), and granted the
finite geometric Pfaffian datum \((P_{\mathrm{fin}})\) of
Definition~\ref{def:finite-geometric-pfaffian-datum} on a cofinal HN
system, including skew perfect complexes, Pfaffian line and section,
support, parity, active-support coverage, type-II wall divisor order,
leading coefficient, Pfaffian-line comparison, and Pfaffian-line
Mittag--Leffler compatibilities,
the protected Pfaffian of the Dirac--Igusa pro-object
\(\mathfrak D_X^{\mathrm{DI}}\) equals the automorphic section
\(\mathcal D_X=\Delta_5\):
\[
\boxed{\;
\iota_{\mathrm{aut}}\!\left(
\operatorname{Pf}_{\mathrm{prot}}(\mathfrak D_X^{\mathrm{DI}})\right)
\;=\;\mathcal D_X\;=\;\Delta_5
\;}
\]
as sections of \(\mathcal L^5\otimes\nu_{\Delta_5}\) over the genus-two
Siegel space \(\Sp(2,\Z)\backslash\mathcal H_2\), where
\[
\iota_{\mathrm{aut}}:\mathscr L_{\mathrm{Pf},X}
\xrightarrow{\,\sim\,}\mathcal L^5\otimes\nu_{\Delta_5}
\]
is the fixed Pfaffian-to-automorphic line comparison, compatible with
the cusp trivialization at \(\mathfrak c_\infty\) and with squaring to
\(\mathcal L^{10}\).  Equivalently, in the type-II Borcherds chamber,
\[
\iota_{\mathrm{aut}}\!\left(
\operatorname{Pf}_{\mathrm{prot}}(\mathfrak D_X^{\mathrm{DI}})\right)
\big|_{\mathfrak c_\infty}
=64\,q^{1/2}r^{1/2}s^{1/2}
\!\!\!\!\prod_{(n,l,m)\in\Gamma_{\mathrm{eff}}}\!\!\!
(1-q^nr^ls^m)^{f(nm,l)}.
\]
After squaring, the orientation forgets and the OP scalar branch
identifies the level-\(n\) protected scalar shadow with the inverse
squared Pfaffian:
\[
 Z^X_{\mathrm{OP}}\;=\;
 -\,64^2\,(\Pf_{\mathrm{prot}}(\mathfrak D_X^{\mathrm{DI}}))^{-2}
 \;=\;-4096\,\Delta_5^{-2},
\]
in agreement with \cite{OB}, where \(4096 = 64^2 =
([q^{1/2}r^{1/2}s^{1/2}]\Delta_5)^2\) is the squared theta-normalization
and the leading minus sign is the OP scalar branch convention.  The
BPS normalization rescales by \(-4096^{-1}\):
\(Z^{K3\times E}_{\mathrm{BPS}} = -4096^{-1}Z^X_{\mathrm{OP}}
= \Delta_5^{-2}\).
\end{theorem}

\begin{proof}
The inverse-limit assembly consists of a Pfaffian product, a leading
coefficient identification with a known weight-five Siegel form, and a
\(q\)-expansion principle on a punctured neighbourhood of the type-II
cusp \(\mathfrak c_\infty\).

\smallskip\emph{Finite-stage Pfaffian product
\((\delta+\varepsilon)\).}  The retained D0-HN datum supplies the
cosection-reduced d-critical pro-limit on the retained HN system.  The
finite geometric Pfaffian datum \((P_{\mathrm{fin}})\) supplies, at
every cofinal \(R\), the cosection-reduced skew perfect complexes
\(K^{\mathrm{red}}_{\mathcal C,R}\), the first-order block
\(\mathfrak D_{X,R}\) on the reduced compact quotient
\(\mathfrak Q_R=(\mathfrak M_R/E)^{\mathrm{red}}\), the Pfaffian line
\(\mathscr L_{\mathrm{Pf},R}\), the section \(\operatorname{pf}_{X,R}\),
active-support coverage, the support and parity residual vanishings,
the type-II wall divisor order, and the normalized leading coefficient
\(\mathfrak o^{\mathrm{lead}}_R=0\).  Hypothesis (O1) supplies
the orientation datum whose square-root line \(o_{\mathcal C}\) satisfies
\(o_{\mathcal C}^{\otimes 2}\simeq
\det\operatorname{RHom}_{\mathrm{red}}(\mathcal C,\mathcal C)\) and
extension-multiplicativity
\(o_{\mathcal C\oplus\mathcal C'}\simeq
o_{\mathcal C}\otimes o_{\mathcal C'}\), descended under the reduced
\(E\)-quotient.  The protected Pfaffian functor of (\(\beta\)) is then the
graded Pfaffian product of the determinant of
\(\operatorname{RHom}_{\mathrm{red}}\) over the retained
charge strata.  The exponent-matching clause of the retained
finite-stage Dirac block gives
\[
\sdim P^\Pi_{X,(n,l,m)}=f(nm,l)
\quad((n,l,m)\in\Gamma_{\mathrm{eff}}),
\]
so the graded Pfaffian product, expanded in the type-II cusp polarization
\((q,r,s)=(e^{2\pi iz_1},e^{2\pi iz_2},e^{2\pi iz_3})\), is
\[
\operatorname{Pf}_{\mathrm{prot}}(\mathfrak D_{X,R})
=v_0\!\!\prod_{(n,l,m)\in\Gamma_{\mathrm{eff,R}}}\!\!
(1-q^nr^ls^m)^{f(nm,l)},
\qquad v_0=64q^{1/2}r^{1/2}s^{1/2},
\]
where \(\Gamma_{\mathrm{eff,R}}\) is the active support cut by HN height
\(R\).  The Weyl monomial \(v_0\) is the theta normalization
\([q^{1/2}r^{1/2}s^{1/2}]\Delta_5=64\) of \cite[Theorem~4.1]{GN}.
The finite normalization table \(\texttt{delta5\_theta\_leading}\)
records the scalar equalities \(64^2=4096\), \(D_5=64^{-1}\Delta_5\),
and \(\chi_{10}^{\mathrm{OP}}=4096^{-1}\Delta_{10}\), together with
separate convention rows for
\(\Delta_{10}\), \(\chi_{10}\), \(\Phi_{10}^{\mathrm{un}}\), the OP
monic form, and the inverse scalar branches; it has no
Pfaffian-line, orientation-character, or O2 wall-atlas content.
This product is the cusp expansion of the supplied Pfaffian section, not
a replacement for that section: by
Proposition~\ref{prop:scalar-pfaffian-data-not-section}, the scalar
product, its square, and the OP scalar do not determine
\(\mathscr L_{\mathrm{Pf},R}\), \(\operatorname{pf}_{X,R}\), or
\(\epsilon_{o,R}\) without the Pfaffian-line and line-comparison data.

\smallskip\emph{Inverse limit \((\gamma)\).}  The
pro-object ambient is the cofinal HN system; the retained D0-HN datum
and finite geometric Pfaffian datum supply strict transition maps with
vanishing successor residuals
\(\mathfrak o^{\mathrm{atlas}}_{R^+R}=
\mathfrak o^{\mathrm{tr}}_{R^+R}=0\); the finite Pfaffian sections
\(\operatorname{pf}_{X,R}\) form a Mittag-Leffler system on the active
support, and the inverse limit is the formal Borcherds product
\[
\operatorname{Pf}_{\mathrm{prot}}(\mathfrak D_X^{\mathrm{DI}})
\big|_{\mathfrak c_\infty}
=\lim_R\operatorname{Pf}_{\mathrm{prot}}(\mathfrak D_{X,R})
=64\,q^{1/2}r^{1/2}s^{1/2}
\!\!\!\prod_{(n,l,m)\in\Gamma_{\mathrm{eff}}}\!\!
(1-q^nr^ls^m)^{f(nm,l)}.
\]
Class \(M\) completion (chain-level after weight-completion, in the
sense of Vol I's class \(M\) global triangle) is exactly the ambient in
which \(\lim_R\) commutes with the graded Pfaffian product on the
active support; this is the \(\gamma\) clause.

\smallskip\emph{Comparison with the imported Borcherds product
\((\beta)\).}  By
Proposition~\ref{prop:bgn-identification-recap},
the formal product
\(64\,q^{1/2}r^{1/2}s^{1/2}\prod(1-q^nr^ls^m)^{f(nm,l)}\) equals the
Fourier expansion of the automorphic section
\(\mathcal D_X=\Delta_5\) at \(\mathfrak c_\infty\); this is the
Borcherds--Gritsenko--Nikulin theorem
\cite{Borcherds95,GN,GNII}.  The fixed
\(\iota_{\mathrm{aut}}:\mathscr L_{\mathrm{Pf},X}\xrightarrow{\sim}
\mathcal L^5\otimes\nu_{\Delta_5}\) of (\(\beta\)), compatible with cusp
trivialization and squaring, transports the inverse-limit Pfaffian to a
section of \(\mathcal L^5\otimes\nu_{\Delta_5}\).

\smallskip\emph{\(q\)-expansion principle \((\alpha)\).}
The chart \(\alpha\) is the open factorization category
\(\mathcal C\) on \((X,D,\tau)\) with boundary vacuum \(b\): on the
punctured neighbourhood of \(\mathfrak c_\infty\) inside Sp\((2,\Z)
\backslash\mathcal H_2\) the chart algebra
\(A_b=\mathrm{End}_{\mathcal C}(b)\) supplies the local
factorization-algebra structure within which the Pfaffian section is
holomorphic.  Two holomorphic genus-two Siegel modular forms with the
same Maass character \(\nu_{\Delta_5}\) and the same Fourier expansion at
\(\mathfrak c_\infty\) are equal: this is the \(q\)-expansion principle
for holomorphic Siegel modular forms
\cite{Igusa1962,Igusa1964II}.  Both sides
have weight \(5\) and Maass character \(\nu_{\Delta_5}\): the right-hand
side by the Borcherds--Gritsenko--Nikulin theorem; the left-hand side
because the orientation character of the retained
\(\mathfrak D_X^{\mathrm{DI}}\), under (O1)+(O1)\(^+\)+(O2), is exactly
\(\nu_{\Delta_5}\) on \(W^{(2)}(\La^{2,1}_{II})\) (this is
Theorem~\ref{thm:ch6-orientation-character}).  Hence
\(\iota_{\mathrm{aut}}(\operatorname{Pf}_{\mathrm{prot}}(\mathfrak D_X^{\mathrm{DI}}))=\Delta_5\)
as sections of \(\mathcal L^5\otimes\nu_{\Delta_5}\).

\smallskip\emph{Squared scalar \((\beta)\).}
Squaring forgets orientation: the squared line
\(\bigl(\mathcal L^5\otimes\nu_{\Delta_5}\bigr)^{\otimes 2}=\mathcal L^{10}\)
is trivially charactered, and \(\Delta_5^2=\Delta_{10}\) is scalar.
The OP scalar normalization of Convention~\ref{conv:op-chamber} fixes the partition function
\(Z^X_{\mathrm{OP}}=-(\chi_{10}^{\mathrm{OP}})^{-1}=-4096\,\Delta_5^{-2}\),
with \(\chi_{10}^{\mathrm{OP}}=D_5^2=4096^{-1}\Delta_5^2\) and
\(D_5=64^{-1}\Delta_5\).  Composition of the squared Pfaffian section
with the OP scalar branch gives the OP chamber scalar representative
\[
 Z^X_{\mathrm{OP}}\;=\;-\,(\Pf_{\mathrm{prot}}(\mathfrak D_X^{\mathrm{DI}}))^{-2}\cdot
 C_{\square,X},\qquad C_{\square,X}=4096,
\]
identifying \(Z^X_{\mathrm{OP}}=-4096\,\Delta_5^{-2}\); the minus sign
is the OP scalar branch convention, recorded by
\((-p_{\mathrm{DT}})^n\), and is not a feature of the orientation
character.  The protected scalar trace is
\(Z_{\mathrm{BPS}}^{K3\times E}=\Delta_5^{-2}\).
\end{proof}

\begin{remark}[Why (D0)+(O1)+(O1)\(^+\)+(O2)+\((P_{\mathrm{fin}})\)]
\label{rem:ch6-why-four}
Each named datum is essential.  (D0) gives the existence and
the finite first-order block; without it there is no
\(\mathfrak D_{X,R}\) and no Pfaffian product to take.  (O1) gives
the orientation line and the reduced \(E\)-quotient descent; without
the line the Pfaffian section has no domain.  (O1)\(^+\) gives the
honest Weyl action on the reduced quotient orientation; without it the
type-II reflection signs are projective and the orientation character is
projective, not a homomorphism.  (O2) gives the local Pfaffian normal
form; without it the rank \(N_{\delta_i}^{\mathrm{Pf}}\) is undetermined
and the local sign cannot be computed.  The finite geometric Pfaffian
datum \((P_{\mathrm{fin}})\) gives the skew perfect complexes,
Pfaffian line, section, support, parity, active-support coverage, wall
divisor order, leading coefficient, line comparison, and
Pfaffian-line Mittag--Leffler compatibility; without it the inverse-limit
product has no section-level domain and no determinant-line object.  If
one of these data is omitted, one of the proof steps loses its object.
\end{remark}

\section{The orientation character identity}
\label{sec:ch6-orientation-character}

The second theorem identifies the geometric orientation monodromy character
\(\epsilon_o\) of the Dirac--Igusa pro-object with the automorphic Maass
character \(\nu_{\Delta_5}\) on the type-II Weyl group.  The proof
splits into a local Pfaffian wall computation, a group-theoretic
extension, and an independent automorphic check.

\begin{theorem}[Orientation character]
\label{thm:ch6-orientation-character}
Recall the type-II simple roots
\(\delta_1=2f_2-f_3,\delta_2=2f_{-2}-f_3,\delta_3=f_3\) on
\(\La^{2,1}_{II}\subset\La^{2,1}=\Z f_2\oplus\Z f_3\oplus\Z f_{-2}\)
(Equation~\eqref{eq:bkm-type-ii-simple-roots}).  Granted the orientation
classes (O1), (O1)\(^+\), (O2), the geometric orientation monodromy
character of the Dirac--Igusa pro-object equals the automorphic Maass
character on the type-II Weyl group:
\[
\boxed{\;
\epsilon_o\;=\;\nu_{\Delta_5}
\quad\text{on}\quad W^{(2)}(\La^{2,1}_{II})\;}.
\]
The character takes the value \(-1\) on each of the three type-II simple
reflections \(s_{\delta_1},s_{\delta_2},s_{\delta_3}\), and the
identity is reduced to these three values by the abelianization
\(W^{(2)}(\La^{2,1}_{II})_{\mathrm{ab}}\simeq(\Z/2)^3\)
(Definition~\ref{def:O1-plus-weyl-transport}).
\end{theorem}

The identity is a comparison of two characters defined independently:
\(\epsilon_o\) is computed by transport of the Pfaffian section under
the chosen Weyl action of (O1)\(^+\); \(\nu_{\Delta_5}\) is the Maass
multiplier of the holomorphic Siegel form \(\mathcal D_X=\Delta_5\) computed from
the multiplier system of the corresponding theta lift
\cite{MAASS:1,GN,Gritsenko99}.  No scalar branch convention enters this
sign comparison: replacing \(\Delta_5\) by \(c\Delta_5\),
\(c\in\C^\times\), changes neither the Maass character nor the divisor
orders.  The factor \(64\) names the monic Borcherds product
\(D_5=64^{-1}\Delta_5\); the OP/DT scalar \(-4096\) is unused.

\begin{proof}
Three inputs combine:
Proposition~\ref{prop:appE-local-pfaffian-sign} computes the local Pfaffian
signs on the three simple type-II walls;
Definition~\ref{def:O1-plus-weyl-transport} extends a generator
sign-tuple \((\epsilon_1,\epsilon_2,\epsilon_3)\in\{\pm1\}^3\) to a
unique character on \(W^{(2)}(\La^{2,1}_{II})\); and
Lemma~\ref{lem:bkm-maass-restriction} computes the automorphic
target values \(\nu_{\Delta_5}(s_{\delta_i})=-1\) directly from the
Maass multiplier and the Borcherds divisor order
\(d_{\delta_i}^{\mathrm{aut}}=1\).

\smallskip\emph{Local Pfaffian signs from \((O2)\).}  By (O2), at a
generic point of each \(\Hpl_{\delta_i}\), \(i=1,2,3\), the chosen
local equation \(u_i\) satisfies \(s_{\delta_i}(u_i)=-u_i\) and the
normal determinant-complex splits as
\(K^{\mathrm{nor}}_{\delta_i}\simeq[\R\xrightarrow{u_i}\R]\) with
\(\operatorname{Pf}=\upsilon_{\delta_i}u_i\) and
\(s_{\delta_i}(\upsilon_{\delta_i})=\upsilon_{\delta_i}\).  Transport
along the closed path
\(\ell_{\delta_i}\) (the path in the reduced quotient wall-complement
carrying \(u_i=\epsilon\) to \(u_i=-\epsilon\), closed by the (O1)\(^+\)
lift \(\widetilde\tau_i\)) sends
\(\upsilon_{\delta_i}u_i\) to \(-\upsilon_{\delta_i}u_i\), giving
\[
\epsilon_o(s_{\delta_i})=(-1)^{N_{\delta_i}^{\mathrm{Pf}}}=-1,
\qquad i=1,2,3.
\]
This is Proposition~\ref{prop:appE-local-pfaffian-sign}.  Squaring eliminates
the sign and hence eliminates the OP scalar branch ambiguity: the squared
Pfaffian \(\operatorname{Pf}(K^{\mathrm{nor}}_{\delta_i})^2=u_i^2\) is
fixed by \(s_{\delta_i}\).

\smallskip\emph{Group-theoretic extension via \((O1)^+\).}  The
type-II Weyl group has Coxeter presentation
\[
W^{(2)}(\La^{2,1}_{II})
\simeq
\langle s_{\delta_1},s_{\delta_2},s_{\delta_3}\mid
s_{\delta_1}^2=s_{\delta_2}^2=s_{\delta_3}^2=1\rangle,
\]
with no braid relations — the off-diagonal entries
\((\delta_i,\delta_j)=-2\) for \(i\ne j\) give Coxeter exponents
\(m_{ij}=\infty\).  Hence
\(W^{(2)}(\La^{2,1}_{II})_{\mathrm{ab}}\simeq(\Z/2)^3\) with one
generator per \(s_{\delta_i}\), and any \(\{\pm1\}\)-valued assignment
on the generators extends uniquely to a homomorphism.  By (O1)\(^+\),
\([c_o]=0\), \(\widetilde\tau_i^{\,2}=1\), and
\(\omega_{i,\mathcal C}=0\); the lifts \(\widetilde\tau_i\) define an
honest action on the reduced quotient orientation, and the induced
character has \(\epsilon_o(s_{\delta_i})=-1\) for \(i=1,2,3\).  This
is Definition~\ref{def:O1-plus-weyl-transport}.

\smallskip\emph{Independent automorphic computation of
\(\nu_{\Delta_5}\).}  The Maass multiplier
\(\nu_{\Delta_5}\) of the singular theta lift of
\(\phi^{\mathrm{Weil}}_{0,1}\) is
computed from the multiplier system of the Borcherds product (Borcherds
\cite{Borcherds95}, Gritsenko--Nikulin
\cite[Section~3]{GN},\cite{Gritsenko99}, Maass \cite{MAASS:1}).
On the type-II simple roots
\(\delta_1=2f_2-f_3\), \(\delta_2=2f_{-2}-f_3\), \(\delta_3=f_3\), the
Borcherds divisor orders are
\[
d_{\delta_i}^{\mathrm{aut}}
:=\ord_{\Hpl_{\delta_i}}(\Delta_5)=f(0,\pm1)=1,
\qquad i=1,2,3,
\]
so \(\nu_{\Delta_5}(s_{\delta_i})=(-1)^{d_{\delta_i}^{\mathrm{aut}}}=-1\)
for \(i=1,2,3\).  This is
Lemma~\ref{lem:bkm-maass-restriction}.  The chamber-permuting
roots \(f_{-2}-f_2,f_2-f_3,f_{-2}-f_3\) lie in
\(\La^{2,1}\setminus\La^{2,1}_{II}\) and are not type-II walls; their
Maass values are \(+1\) but no Hall orientation sign is asserted on them
unless a semidirect-product Hall lift is constructed
(Remark~\ref{rmk:ch6-S3-extension}).

\smallskip\emph{The comparison \(\epsilon_o=\nu_{\Delta_5}\).}
Two \(\{\pm1\}\)-valued characters of \(W^{(2)}(\La^{2,1}_{II})\) that
agree on \(s_{\delta_1},s_{\delta_2},s_{\delta_3}\) agree everywhere
by the Coxeter presentation above.  The local Pfaffian computation and
the automorphic computation give that both \(\epsilon_o\) and
\(\nu_{\Delta_5}\) take the value \(-1\) on each of the three simple
reflections.  Therefore
\(\epsilon_o=\nu_{\Delta_5}\) on
\(W^{(2)}(\La^{2,1}_{II})\).
\end{proof}

\begin{remark}[The OP minus sign is not the source of \(\epsilon_o\)]
\label{rmk:ch6-op-not-orientation-source}
The local Pfaffian computation and the automorphic multiplier
computation are independent.  The former uses no automorphic data; the
latter uses no Pfaffian local model.  In particular,
the OP scalar
\(C_{\mathrm{OP},X}=-4096=(-1)\cdot C_{\square,X}\) does
not source the orientation character: squaring forgets orientation, so
\(C_{\mathrm{OP},X}\) carries no information about \(\epsilon_o\).  The
minus sign \(\operatorname{sign}(C_{\mathrm{OP},X})=-1\) is the OP scalar
branch convention of \cite{OB}, recorded by the fermionic
shift \((-p_{\mathrm{DT}})^n\) of the OP/DT chamber.  The orientation
character is a \emph{geometric} monodromy datum, computed by transport
of the Pfaffian line around type-II walls; it is fixed by the data of
(O1)+(O1)\(^+\)+(O2), and is independently checked against the
automorphic target by the Maass multiplier computation.
\end{remark}

\begin{remark}[Imaginary roots are not Weyl reflections]
\label{rmk:ch6-imaginary-roots}
The Weyl group \(W^{(2)}(\La^{2,1}_{II})\) of the type-II BKM datum is
generated by real simple reflections only.  The imaginary simple roots
contribute to the Weyl--Kac--Borcherds denominator through the
multiplicity coefficients \(m(a)\) and \(\tau(a)\) of
Gritsenko--Nikulin \cite[Theorem~2.1]{GNII},\cite[Section~3]{GN}; they
do not define reflections in
\(W^{(2)}(\La^{2,1}_{II})\).  Hence
\(\epsilon_o:W^{(2)}(\La^{2,1}_{II})\to\{\pm 1\}\) is a real-root
reflection datum.  Local divisor monodromy of a compact Hall theory
around imaginary components of the Borcherds divisor is not a Weyl
character and is not asserted here.
\end{remark}

\begin{remark}[Extension to \(W^{(2)}(\La^{2,1}_{II})\rtimes\Aut(\Poly_{II})\)]
\label{rmk:ch6-S3-extension}
The chamber-permuting symmetry group
\(\Aut(\Poly_{II})\simeq S_3\) acts on the Coxeter generators
\(s_{\delta_1},s_{\delta_2},s_{\delta_3}\) by permutation; its three
reflections \(s_{f_{-2}-f_2},s_{f_2-f_3},s_{f_{-2}-f_3}\) have
\(\nu_{\Delta_5}=+1\).  An extension of \(\epsilon_o\) to
\(W^{(2)}(\La^{2,1}_{II})\rtimes\Aut(\Poly_{II})\) would require a
separate Hall semidirect-product lift of the chamber-permuting
reflections.
\end{remark}

\section{The primitive recognition theorem}
\label{sec:ch6-primitive-recognition}

The third theorem is a criterion for identifying the source-side
primitive Lie superalgebra of the Dirac--Igusa pro-object with the BKM
denominator algebra \(\mathfrak g_{\Delta_5}\).  The conditions form the
\emph{recognition criterion}: ten finite-stage data carrying compact
geometric provenance, none of which is inferable from the Borcherds
product or signed determinant alone.

\begin{theorem}[Primitive recognition]
\label{thm:ch6-primitive-recognition}
Granted the recognition criterion (S1)--(S10) of
Theorem~\ref{thm:primitive-recognition}, the cofinal finite-window
datum of Definition~\ref{def:cofinal-primitive-recognition-datum},
the chiral Koszul source datum of
Definition~\ref{def:chiral-koszul-source-datum}, the finite Koszul
comparison datum of
Definition~\ref{def:finite-koszul-comparison-datum}, and the retained
orientation data \((O1_{\mathrm{quot}})\), the chosen Weyl lifts, and
\((O2_{\mathrm{atlas}})\),
suppose the compact \(K3\times E\) Hall/factorization construction
supplies, on a cofinal sequence of downward-saturated finite root windows
\(W_1\subset W_2\subset\cdots\) with
\(\bigcup_\nu W_\nu=\mathcal R_+\):
\begin{itemize}
\item \emph{(S1) Chain-level descent endpoint.}  The (D0)-degenerated
  source coalgebra \((C_{X,W_\nu},Q_{X,W_\nu})\) of
  Definition~\ref{def:chiral-koszul-source-datum}, descended chain-level
  to its primitive part on every finite window
  \(W_\nu\), carrying both the orientation lane and the
  Frobenius--Serre primitive structure feeding the recognition window.
\item \emph{(S2) Cartan and root-degree matching.}  An even Cartan
  identification \(\iota_H:P^{\Pi,\mathrm{red}}_{X,0}\xrightarrow{\sim}
  \La^{2,1}_{II}\otimes\C\), with Gram-grading
  \(\alpha(n,l,m)=2nf_2-lf_3+2mf_{-2}\).
\item \emph{(S3) Simple primitive representatives.}  Parity-homogeneous
  \(e_i^X,f_i^X,h_i^X\) for every simple-root index \(i\in I\) of
  \(\mathscr B_{\Delta_5}\), with Gritsenko--Nikulin parity dimensions
  \(\mu_{\overline 0}(a),\mu_{\overline 1}(a)\) on imaginary fibres.
\item \emph{(S4) Borcherds--Kac relations.}  Cartan, Chevalley
  (\([e_i^X,f_j^X]=\delta_{ij}h_i^X\)), real Serre, Borcherds
  orthogonality, super sign, and isotropic-string relations satisfied in
  the protected Hall bracket.
\item \emph{(S5) Completed Hopf radical equality.}  The completed Hopf
  pairing has closed homogeneous radical \(\operatorname{Rad}_X\) which
  is a Lie ideal and a coideal, and identifies under the radical
  quotient with \(\operatorname{Rad}_{\mathrm{GN}}\).
\item \emph{(S6) No-extra-relations.}  The presentation map
  \(\pi:\widehat{\mathfrak F}(\mathscr B_{\Delta_5})\to
  P_X^{\Pi,\mathrm{red}}\) has
  \(\ker\pi=\overline{J_{\mathrm{BK}}+\operatorname{Rad}_{\mathrm{GN}}}\),
  on every finite window separately.
\item \emph{(S7) Generation.}  \(P_X^{\Pi,\mathrm{red}}\) is
  topologically generated by \(P^{\Pi,\mathrm{red}}_{X,0}\) and
  \(\{e_i^X,f_i^X\}_{i\in I}\).
\item \emph{(S8) Completed PBW.}  The associated graded
  \(\operatorname{gr}_{\mathrm{PBW}}U(P_X^{\Pi,\mathrm{red}})_W
  \xrightarrow{\sim}
  \operatorname{gr}_{\mathrm{PBW}}U(\mathfrak g_{\Delta_5})_W\) is an
  isomorphism on every finite window, with strict transitions.
\item \emph{(S9) Exact triangular completion.}  The HN inverse limit is
  exact on the saturated finite root windows: \(\lim^1\) vanishes for
  the kernels, radicals, PBW associated gradeds, and root-space parity
  pieces.
\item \emph{(S10) Full parity dimensions.}  Two integers per nonzero
  root degree, in both signs:
  \(\dim P^{\Pi,\mathrm{red}}_{X,\gamma,\overline p}
  =\rootmult_{\overline p}^{\mathrm{GN}}(\alpha(\gamma))\),
  \(p=0,1\), and the same equality in degree \(-\gamma\) through the
  Borcherds--Kac Chevalley anti-involution; not only their signed
  difference.
\end{itemize}
Then the recognition datum determines an isomorphism of completed
\(\Gamma_{\mathrm{gram}}\)-graded Lie superalgebras
\[
\boxed{\;
H^0\Prim_{\mathrm{prot}}\!\bigl(
\overline\Pi_{X,*}^{\Theta}\mathcal F_X\bigr)
\big/\overline{\operatorname{Rad}\langle\,\cdot\,,\,\cdot\,\rangle_X}
\;\cong\;\mathfrak g_{\Delta_5}
\;}
\]
and the Weyl--Kac--Borcherds denominator of
\(P_X^{\Pi,\mathrm{red}}\) is \(\operatorname{den}(\mathfrak g_{\Delta_5})
=64^{-1}\Delta_5(2Z)\).  If in addition the chiral Koszul source datum of
Definition~\ref{def:chiral-koszul-source-datum} is fixed, including
the finite Koszul comparison datum of
Definition~\ref{def:finite-koszul-comparison-datum} and the strict
Koszul Mittag--Leffler exactness hypothesis of
Proposition~\ref{prop:chiral-koszul-source-datum-consequence}, then the
chiral cobar of the source coalgebra \(C_X\) is identified with the
formal current envelope compatibly with the Weyl action, Pfaffian
orientation character, Hall pairing, radical quotient, and PBW
associated graded:
\[
\Omega_E^{\mathrm{ch}}C_X\;\simeq\;\mathsf{Den}_{\Delta_5,E}.
\]
\end{theorem}

\begin{proof}
The argument is the inverse-limit assembly of (S1)--(S10) on each
finite window, together with the chiral cobar identification.  The
window-wise theorem is Theorem~\ref{thm:primitive-recognition}, with
the algebraic finite-window step isolated in
Proposition~\ref{prop:finite-window-presentation-recognition}; the
additional input is the chiral-cobar comparison
\(\Omega_E^{\mathrm{ch}}C_X\simeq\mathsf{Den}_{\Delta_5,E}\), with
the compatibility statement supplied by
Proposition~\ref{prop:finite-koszul-comparison-compatibility} and its
Mittag--Leffler passage in
Proposition~\ref{prop:chiral-koszul-source-datum-consequence}.

\smallskip\emph{Presentation map and kernel
\((\alpha+\beta)\).}  (S2) and (S3) define the continuous Lie-superalgebra
homomorphism
\[
\pi:\widehat{\mathfrak F}(\mathscr B_{\Delta_5})
\longrightarrow P_X^{\Pi,\mathrm{red}}
\]
on the descended Gram-graded completion.  (S4) puts the Borcherds--Kac
ideal \(J_{\mathrm{BK}}\) in the kernel.  (S6) gives the kernel equality
\(\ker\pi=\overline{J_{\mathrm{BK}}+\operatorname{Rad}_{\mathrm{GN}}}\)
on every finite window, hence injectivity of the descended map
\(G(\mathscr B_{\Delta_5})\to P_X^{\Pi,\mathrm{red}}\).  (S7) gives
surjectivity.  Thus the descended map is an isomorphism on every finite
window by
Proposition~\ref{prop:finite-window-presentation-recognition}.

\smallskip\emph{Radical descent \((\varepsilon)\).}
(S5) identifies the source Hopf radical
\(\overline{\operatorname{Rad}_X}\) with the target radical
\(\overline{\operatorname{Rad}_{\mathrm{GN}}}\).  The required Hopf
adjointness, ideal/coideal property, quotient-tensor non-degeneracy, and
closedness on the inverse system are part of (S5), equivalently the
finite matrix conditions (R4)--(R5).  The orientation lane supplies the
trace pairing and Frobenius formalism on which those equations are
written; it does not prove the source matrices \(M,D,B,G,K,Q\), the
kernel equality, or the PBW comparison.

\smallskip\emph{Full parity dimensions \((\delta)\).}
(S10) is the two-integer equality
\(\dim P^{\Pi,\mathrm{red}}_{X,\gamma,\overline p}=
\rootmult_{\overline p}^{\mathrm{GN}}(\alpha(\gamma))\) for \(p=0,1\), not
just the signed difference.  This rules out hidden zero-superdimension
source summands, wrong parity lifts invisible to the determinant, and
parity-cancelling ideals.  A putative compact source matching every
Hall bracket template of \(\mathfrak g_{\Delta_5}\) on a finite window
can still fail (S5): a kernel element \(k\in K_{\gamma_\rho}\) with
\((Q_{\gamma_\mu}\otimes Q_{\gamma_\nu})D^{\mu,\nu}_\rho(k)\ne 0\) is
invisible to the quotient target table yet breaks the radical
coideal property, so the Hopf quotient is not defined
(Theorem~\ref{thm:primitive-recognition}).  This is the finite
tensor-Hopf obstruction hidden by signed determinant data; (S5)+(S10)
together block it.

\smallskip\emph{PBW comparison and exact completion
\((\gamma)\).}  (S8) is the associated-graded isomorphism on every
finite window, with strict transitions.  (S9) is the exact-completion
clause: \(\lim^1\) vanishes for kernels, radicals, PBW associated
gradeds, and root-space parity pieces.  Equivalently, no new generator,
relation, radical element, or parity-cancelling summand appears only
after completion.  Proposition~\ref{prop:strict-ml-completion-primitive-recognition}
is the homological algebra that passes from the finite-window
isomorphisms to the completed source and target.  Class \(M\)
completion is the ambient: chain-level identifications hold after
weight-completion; this is the \(\gamma\) clause.

\smallskip\emph{Passage to the inverse limit.}  Since the
windows are cofinal (\(\bigcup_\nu W_\nu=\mathcal R_+\)) and the radical
in (S5) is the closure of the finite kernels, the finite isomorphisms
of the finite-window clauses determine the inverse-limit isomorphism
\[
H^0\Prim_{\mathrm{prot}}\!\bigl(
\overline\Pi_{X,*}^{\Theta}\mathcal F_X\bigr)
\big/\overline{\operatorname{Rad}\langle\,\cdot\,,\,\cdot\,\rangle_X}
=P_X^{\Pi,\mathrm{red}}
\xrightarrow{\,\sim\,}\mathfrak g_{\Delta_5}.
\]
The Weyl--Kac--Borcherds denominator of
\(P_X^{\Pi,\mathrm{red}}\) is therefore
\(\operatorname{den}(\mathfrak g_{\Delta_5})=64^{-1}\Delta_5(2Z)\), by
Theorem~\ref{thm:bkm-denominator-algebra-identity}.

\smallskip\emph{Chiral cobar identification.}  Suppose the
chiral Koszul source datum of
Definition~\ref{def:chiral-koszul-source-datum} is fixed:
\((\mathcal F_{X,\sigma,S,\le R}^{\mathrm{hyb}},C_{X,R},
F^{\mathrm{co}}_{R,\bullet},\Delta_R^{\mathrm{ch}},
b_{X,R},\Theta_{\mathrm{Kos},R})\) with vanishing residual
\(\mathfrak O_{\mathrm{Kos},R}=0\), together with the finite Koszul
comparison datum
\(\mathfrak K^{\mathrm{cmp}}_{X,R}\) and
\(\mathfrak O_{\mathrm{Kcmp},R}=0\), compatibly in \(R\) and with the
strict Mittag--Leffler exactness hypothesis of
Proposition~\ref{prop:chiral-koszul-source-datum-consequence}.  The
target side is the formal current envelope
\(\mathsf{Den}_{\Delta_5,E}=U_E^{\mathrm{ch}}\!
\bigl(\operatorname{Cur}_E(\mathcal A_{\mathrm{den}})\bigr)\)
of Proposition~\ref{prop:bkm-formal-current-envelope}.  The cobar comparison
of Propositions~\ref{prop:finite-koszul-comparison-compatibility}
and~\ref{prop:chiral-koszul-source-datum-consequence}
identifies
\[
\Omega_E^{\mathrm{ch}}C_X\;\simeq\;\mathsf{Den}_{\Delta_5,E},
\]
through the source-internal bar/cobar admissibility (the residual
\(\mathfrak o^{\varepsilon,\mathrm{src}}_R=0\)) and the source-to-target
quasi-isomorphism (\(\mathfrak o^{\vartheta}_R=0\)), and it preserves
the Weyl, Pfaffian-orientation, Hall-pairing, radical-quotient, and
PBW structures used above.  Agreement on the
primitive summand is not enough by
Proposition~\ref{prop:primitive-restriction-not-koszul}; the full cone
of \(\vartheta_R\) and all compatibility defect complexes must be
acyclic.  The construction avoids the
target-internal bar/cobar counit; this is the \(\beta\)-separation
between observables and states.
\end{proof}

\begin{remark}[The zero-bracket graded super-vector-space counterexample]
\label{rmk:ch6-zero-bracket-counterexample}
The criterion (S1)--(S10) blocks the trivial counterexample of a graded
super vector space \(V_\bullet\) with all brackets zero and parity
dimensions matching \(\rootmult_{\overline p}^{\mathrm{GN}}\) on every
window.  Such a \(V_\bullet\) may satisfy the target parity table and
the signed-superdimension residual on the active support.  It fails
(S4): the Chevalley relation
\([e_i^X,f_j^X]=\delta_{ij}h_i^X\) fails on the diagonal.  It fails
(S6): the kernel of the presentation map
\(\pi:\widehat{\mathfrak F}(\mathscr B_{\Delta_5})\to V_\bullet\) is not
\(\overline{J_{\mathrm{BK}}+\operatorname{Rad}_{\mathrm{GN}}}\) but
contains every commutator \([e_i,e_j]\) with \((\alpha_i,\alpha_j)\ne 0\),
so the induced map \(G(\mathscr B_{\Delta_5})\to V_\bullet\) is not
injective.  It fails (S5) for the zero pairing because the radical is
all of \(V_\bullet\), so the radical quotient is zero and not the BKM
target.  If one installs a nonzero pairing by hand, the denominator
still does not supply its invariant Hopf adjointness, coideal radical,
or equality with \(\operatorname{Rad}_{\mathrm{GN}}\).  The
counterexample shows that signed-superdimension matching alone is not
recognition; the bracket-level data of (S4)+(S5)+(S6) is the
mathematical content.
\end{remark}

\begin{remark}[Recognition without (D0)+(O1)+(O1)\(^+\)+(O2)]
\label{rmk:ch6-recognition-without-orientation}
The recognition criterion of Theorem~\ref{thm:ch6-primitive-recognition}
nominally lists no orientation hypothesis among (S1)--(S10).  The
orientation classes enter through (S1) — the chain-level descent
\(\overline\Pi_{X,*}^{\Theta}\mathcal F_X\) is conditional on the
orientation lane via Definition~\ref{def:O1-strong-reduced-orientation} and the
Frobenius identity from the CY-3 Serre functor — and through (S5),
where the Hopf radical's coideal property is fixed by the same Frobenius
identity (Theorem~\ref{thm:hall-product}).  In the
chapter's combined statement
(Theorem~\ref{thm:ch6-pfaffian-dirac}+Theorem~\ref{thm:ch6-orientation-character}+Theorem~\ref{thm:ch6-primitive-recognition})
the orientation classes enter all three theorems.
\end{remark}

\section{Compatibility with chiral Koszul duality and Calabi--Yau quantum groups}
\label{sec:ch6-chiral-compatibility}

The three theorems sit inside chiral Koszul duality and
Calabi--Yau quantum groups as the View~\textcircled{2} /
View~\textcircled{3} face of the Calabi--Yau-to-chiral comparison.  Three
compatibilities are required.

\subsection{Primitive-to-shadow comparison}
\label{subsec:ch6-primitive-to-shadow}

The primitive-to-shadow comparison is the cross-level identity
\textcircled{\scriptsize 2}\(\leftrightarrow\)\textcircled{\scriptsize 3}
applied to the primitive-to-shadow part of the chain.  The Vol II primitive
datum on \((X,D,\tau)\) is the bicoloured nine-tuple
\[
\bigl(\mathcal C\big|_{(X,D,\tau)},\,
\mathcal D^{\mathrm{cl}}\big|_X;\,
b,\,A_b,\,F^{\mathrm{cl}};\,
\sigma_{\mathrm{half}},\,\operatorname{tr}^{\mathrm{cl}}_X;\,
\operatorname{Tr}_{\mathcal C}\bigr).
\]
The chart of vacuum \(b\) is the \(\alpha\)-datum; the chart
algebra \(A_b=\mathrm{End}_{\mathcal C}(b)\) is the open primitive
algebra; the closed-colour vacuum factorization algebra
\(F^{\mathrm{cl}}\) gives the closed-side trace
\(\operatorname{tr}^{\mathrm{cl}}_X\) which evaluates at the type-II
cusp \(\mathfrak c_\infty\) to the weight-\(5\) Borcherds--Gritsenko--Nikulin
section \(\mathcal D_X=\Delta_5\) (View \textcircled{\scriptsize 1}).  The
comparison
\(A_b=\Phi^{(\Sigma_{d-1},C)}_d(\mathcal C_X)\) is the chain-fusion datum of
chiral Koszul duality and Calabi--Yau quantum groups at level
\(2\): under that datum the Calabi--Yau-side chiral algebra
\(A_X\) is compared with the open-side chart algebra \(A_b\) on the same
curve \(C\) with boundary vacuum induced by
\((\mathcal C_X,\Sigma_{d-1},C)\).  In the present setting, the
comparison concerns the chiral algebra on \(E\) attached to the Hall
source of \(K3\times E\); conditional on the chiral Koszul source
datum, the finite Koszul comparison datum, and strict Koszul
Mittag--Leffler exactness, the chiral cobar
\(\Omega_E^{\mathrm{ch}}C_X\simeq\mathsf{Den}_{\Delta_5,E}\) of
Theorem~\ref{thm:ch6-primitive-recognition} is the source-side chiral
expression at level \(2\), and its constant-current Lie superalgebra is
\(\mathfrak g_{\Delta_5}\) (View \textcircled{\scriptsize 2}).  The
protected Pfaffian \(\operatorname{Pf}_{\mathrm{prot}}(\mathfrak D_X^{\mathrm{DI}})\)
is the pre-trace Pfaffian section.  After squaring and applying the
retained level-\(\mathsf Z\) datum, Vol~I Theorem~H
(chiral-Hochschild concentration on the Koszul locus) and Vol~II's
cyclic-Hochschild Beilinson stratification \textnormal{(i)+(ii)}
identify the derived-centre trace pairing of degree
\(-\dim_\C(K3\times E)\) whose scalar is \(\Delta_5^{-2}\)
(Appendix~\ref{app:obstruction-discharges},
Theorem~\ref{thm:G-vol2-discharge}).  The level-\(\mathsf A\)
gravity-line operator algebra acting on the boundary is the open
Hall--Borcherds residual of Vol~II and is not constructed by the
level-\(\mathsf Z\) trace criterion.

\subsection{Vol III chain fusion at level \(2\)}
\label{subsec:ch6-vol3-chain-fusion}

The Vol III \(\kappa\)-stratification supplies the
\(\kappa\)-tuple
\((\kappa_{\mathrm{cat}},\kappa_{\mathrm{ch}}^{\mathrm{Hodge}},
\kappa_{\mathrm{ch}}^{\mathrm{Heis}},\kappa_{\mathrm{BKM}},
\kappa_{\mathrm{fiber}})=(0,0,3,5,24)\) at \(K3\times E\); the row
\(\kappa_{\mathrm{BKM}}=5\) is the cusp-form weight realized by
\(\Delta_5\), and the row \(\kappa_{\mathrm{fiber}}=24\) is
\(\chi_{\mathrm{top}}(K3)=24\).  The additive expression
\(\kappa_{\mathrm{BKM}}=\kappa_{\mathrm{ch}}+\chi(\mathcal O_{\mathrm{fiber}})\)
has no invariant meaning: the Hodge reading gives \(0+2=2\) at \(N=1\),
and the Heisenberg reading gives the accidental equality \(3+2=5\) only
on that row.  The CHL frame has Borcherds weights
\((5,3,2,\tfrac32,1)\), fixed by \(c_N(0)/2\) with
\(c_N(0)=(10,6,4,3,2)\) per Jatkar--Sen and Govindarajan--Krishna
\cite{JS06Dyon,GK09CHL,GK10Composite},
not by a fixed additive
formula.  Thus \(\kappa_{\mathrm{fiber}}\) is the Euler characteristic
of the topological fibre and not of \(\mathcal O_{\mathrm{fiber}}\); for
K3, \(\chi_{\mathrm{top}}(K3)=24\) but
\(\chi(\mathcal O_{K3})=2\).  Under
the Vol III chain-fusion datum at level \(2\), the chiral algebra
\(A_X\) on the CY face is compared with the chart algebra
\(A_b\) on the Vol II open face.  In the \(K3\times E\) specialization
this comparison becomes
\(A_X=A_b=\Omega_E^{\mathrm{ch}}C_X\simeq\mathsf{Den}_{\Delta_5,E}\) only
under the hypotheses of Theorem~\ref{thm:ch6-primitive-recognition}.

\begin{remark}[Coefficient projection is not a Hall comparison]
\label{rmk:ch6-coefficient-projection}
A Hall--Borcherds comparison map cannot be defined by coefficient
projection alone, and a coefficient projection respecting charge
addition because both products use the same cone is not a recognition
theorem.  A coefficient projection is not a Hall product/bracket matrix
comparison, and does not prove Borcherds--Serre rows, radical descent,
no-extra-relations, or PBW.  The recognition criterion (S4)--(S10) of
Theorem~\ref{thm:ch6-primitive-recognition} is the mathematical
statement:
the source-side bracket matrices, Hopf radical, kernel equality, and
PBW associated-graded comparison are independent data, none of which is
inferable from a coefficient projection.
\end{remark}

\subsection{BCOV / mixed holomorphic-topological strings}
\label{subsec:ch6-bcov}

The mixed holomorphic-topological strings companion provides the local
model \(\R^2_{\mathrm{top}}\times\C^2_{\mathrm{hol}}\) with brane
completion as one realization of the factorization-algebra datum
\(\Phi_d^{\mathrm{FA}}\); the global obstruction is the holomorphic
de Rham obstruction class.  The level-\(n\) scalar shadow
\(Z_{\mathrm{BPS}}^{K3\times E}=\Delta_5^{-2}\) on the closed
\(K3\times E\mapsto\mathrm{point}\) cobordism is the BPS partition
function (Chapter~\ref{ch:zbps-realization}).  A three-dimensional
gravity-line interpretation with internal compact factor \(K3\times E\)
requires the Hall--Borcherds residual of Vol II and is not claimed by
Theorem~\ref{thm:ch6-pfaffian-dirac}.  The identity
\(\operatorname{Pf}_{\mathrm{prot}}(\mathfrak D_X^{\mathrm{DI}})=\mathcal D_X=\Delta_5\)
is a protected Pfaffian identity before trace; its inverse square is
the level-\(n\) protected scalar trace.  The Pfaffian and orientation
parts of the operator-level datum \(\mathfrak D_X^{\mathrm{DI}}\) are
defined on \(K3\times E\)
relative to the retained data of
Appendix~\ref{app:obstruction-discharges}; the primitive Lie-superalgebra
part is the compact-source recognition problem.

\section{Residual obstructions and open subroutes}
\label{sec:ch6-residuals}

The Pfaffian and orientation theorems are retained-data theorems on
\(K3\times E\), relative to Appendix~\ref{app:obstruction-discharges}.
The residuals in this
chapter are the compact-source recognition datum \((S1)\)--\((S10)\)
and the gravity-line Hall--Borcherds operator algebra.

\paragraph{(D0).}  Theorem~\ref{thm:G-D0} is the criterion saying that
the retained D0-HN datum supplies the cosection-twisted reduced
d-critical orientation through the D0-degeneration pro-limit and the
cofinal HN inverse system.  The \(D_0\)-state, \(D_0\)-quotient,
\(D_0\)-observable, \(D_0\)-descent, and \(D_0\)-Pfaffian rows are
therefore part of that datum.

\paragraph{(O1).}  Theorem~\ref{thm:G-O1} supplies the strong reduced
orientation on the \(K3\times E\)-quotient self-Ext frame bundle
relative to the retained quotient-orientation datum.  The
finite-stabilizer linearization characters and the connected
\(BE\)-edge class are part of that datum, not consequences of the
primitive recognition theorem.

\paragraph{(O1)\(^+\).}  The equivariance part of
Theorem~\ref{thm:G-O1-plus} is relative to
the chosen Weyl lifts, the vanishing projective cocycle, and the zero
torsor defects on the retained groupoid.  It gives the honest
\(W^{(2)}(\La^{2,1}_{II})\)-transport of the
orientation character.  The comparison with \(\nu_{\Delta_5}\) also
uses the rank-one wall signs from \((O2_{\mathrm{atlas}})\).  Neither
statement supplies the relation-closed source matrices of the
recognition datum.

\paragraph{(O2).}  Relative to the retained type-II wall-atlas datum
\((O2_{\mathrm{atlas}})\), the rank-one normal-form rank
\(N_{\delta_i}^{\mathrm{Pf}}=1\) holds for \(i=1,2,3\)
(Proposition~\ref{prop:appE-local-pfaffian-sign}).  The type-II middle-wall
constant \(4096=2^{12}=64^2\) admits two readings: as the squared
theta-leading on the type-II Borcherds product (where the \(64=2^6\)
prefactor of \(D_5=64^{-1}\Delta_5\) doubles by \(\Delta_5^2=\Delta_{10}\))
and as the \(2^{12}\) sign assignments over the twelve binary wall
coordinates of the half-Hilbert orbit.  The equality of these two
readings is scalar only, by Theorem~\ref{thm:G-R1-R2}.  The
OP scalar branch sign
\(-1\) is recorded by \((-p_{\mathrm{DT}})^n\) of the OP/DT chamber,
not by the orientation line.  The middle wall
\(\delta_2\leftrightarrow(0,1,1)\) is represented by a wrapped wall
object with positive geometric \(b\)-degree only after the retained
\((O2_{\mathrm{atlas}})\) datum supplies that object.
Proposition~\ref{prop:appE-local-pfaffian-sign}
is the rank-one local sign calculation; the component, boundary,
overlap, quotient, and protected-integration compatibilities are the
global content of \((O2_{\mathrm{atlas}})\).  Higher-rank wall atlases
are not used in the proof of Theorem~\ref{thm:ch6-pfaffian-dirac}.

\paragraph{Recognition criterion (S1)--(S10).}  The window
\(W_{\le 3}\) is the first arithmetic window but is not relation-closed.
Put \(a_{ij}=\delta_i+\delta_j\),
\(\delta_{123}=\delta_1+\delta_2+\delta_3\), and
\(C_{k,s}=a_{ij}+s\delta_k\) for \(\{i,j,k\}=\{1,2,3\}\).
The first base-string terminal degree window is the height-\(7\) closure
\[
W_{\mathrm{rel}}^+
=W_{\le 3}^{\mathrm{act}}\cup R_4\cup I_4\cup C_{\le 7}
\cup\{2\delta_{123}\},
\]
where \(R_4=\{3\delta_i+\delta_j:i\ne j\}\),
\(I_4=\{2a_{ij}:i<j\}\), and
\(C_{\le7}=\{C_{k,s}:2\le s\le5\}\).
It adds the real-Serre terminal rows, doubled isotropic rows,
\(\delta_{123}\)-real strings, odd--odd \(\delta_{123}\) rows, and
complementary height-seven strings.  In this enlarged window,
the finite target degree-closure packet
\(\texttt{wrel\_degree\_closure}\) checks the \(35\) rows, their
\(\delta\)-basis coordinates, heights, norms, Jacobi signed
coefficients, six real-Serre terminal zero rows, and the three
complementary-string terminal zero rows.  Its downward-saturation
audit proves the precise defect: \(W_{\mathrm{rel}}^+\) is not
downward saturated, and the only nonzero proper subdegrees outside
the \(35\)-row set are
\[
D_k=\delta_k+2a_{ij},\qquad \{i,j,k\}=\{1,2,3\},
\]
all below \(2\delta_{123}\), with signed multiplicity
\(\smult(D_k)=-513\).
The same packet verifies that adjoining these three rows gives the
minimal nonzero downward-saturation extension
\[
W_{\mathrm{sat}}^+
=W_{\mathrm{rel}}^+\cup\{D_1,D_2,D_3\},
\]
with zero remaining nonzero proper-subdegree defect.
Remark~\ref{rem:bkm-relation-window-target-parity} supplies full target
parity for \(I_4=\{2a_{ij}\}\), \(C_{k,3}\), \(C_{k,4}\), and
\(C_{k,5}\), while \(C_{k,2}\), \(D_k\), and \(2\delta_{123}\) remain
signed target rows until the finite target presentation is computed.
The finite coefficient table
\(\texttt{k3e\_first\_relation\_window\_coefficients}\) checks the
Jacobi arithmetic for these rows:
\[
f(4,4)=10,\qquad f(2,2)=108,\qquad
f(2,1)=f(4,3)=-513,\qquad
f(3,3)=-64,\qquad f(4,2)=4016,
\]
with the terminal rows \(f(5,5)=f(1,-3)=0\).  This is a target
signed-dimension check, not a parity split or a source comparison.
The A071 target-presentation verifier separately checks the target
parity rows \(2a_{ij}:10|0\), \(C_{k,3}:29|93\),
\(C_{k,4}:10|0\), and \(C_{k,5}:0|0\), and records
\(C_{k,2}\), \(D_k\), and \(2\delta_{123}\) as signed-only rows with
no target basis, pairing, radical, or PBW blocks.
Proposition~\ref{prop:target-admissibility-obstruction-wrel} proves
that these signed rows have no admissible parity, pairing, radical, or
PBW target blocks for (R0)--(R8).  Thus \(W_{\mathrm{rel}}^+\) is a
base-string terminal degree set with a three-dimensional saturation
defect, and \(W_{\mathrm{sat}}^+\) is a target-recorded saturated
extension with signed-only rows.  The post-saturation real-string audit
then finds \(21\) terminal codomains outside \(W_{\mathrm{sat}}^+\),
ten with nonzero signed multiplicity.  Full finite relation closure
therefore starts beyond \(W_{\mathrm{sat}}^+\).  The terminal-layer
audit groups these codomains into \(21\) distinct terminal degrees;
adjoining them creates \(114\) nonzero subdegree-defect incidences
across \(26\) distinct subdegrees.  The next target-degree computation
adjoins exactly these \(26\) rows and gives an \(85\)-row degree set with
zero remaining nonzero proper-subdegree defect.  This closes downward
saturation for the displayed terminal layer, not the real-string
relation audit and not the compact-source comparison.  If these \(26\)
rows are promoted to target-generator rows, the real-root strings force
\(55\) terminal checks; \(54\) terminal codomains lie outside the
\(85\)-row set, and \(12\) of the missing codomains have nonzero signed
multiplicity.  If those \(55\) terminal codomains are adjoined, the next
downward-saturation audit already has \(546\) nonzero proper-subdegree
defect incidences across \(98\) distinct subdegrees.  Adjoining those
\(98\) rows gives a \(237\)-row target degree set with zero remaining
nonzero proper-subdegree defect for that layer.  If the \(98\)-row layer
is promoted to target-generator rows, the real-root strings force
\(206\) terminal checks, \(182\) new terminal degrees, and \(24\)
missing nonzero terminal degrees.  The exact terminal-layer table
records those \(206\) codomains, and the exact saturation-extension
table records \(50\) nonzero proper subdegrees with \(624\) incidences;
adjoining both gives a \(469\)-row target degree set with zero remaining
nonzero proper-subdegree defects for that layer.  Promoting the
\(50\)-row layer forces \(96\) real-string checks, \(90\) new terminal
degrees, no missing nonzero terminal degree, and a \(12\)-row saturation
extension with \(44\) incidences; adjoining both gives a \(571\)-row
target degree set with zero remaining nonzero proper-subdegree defects
for that layer.  Promoting the \(12\)-row layer forces \(20\) zero
terminal checks and a \(6\)-row saturation extension with \(8\)
incidences; adjoining both gives a \(597\)-row target degree set with
zero remaining nonzero proper-subdegree defects for that layer.
Promoting the newly adjoined \(6\)-row layer forces \(10\) zero
terminal checks and a \(6\)-row saturation extension with \(18\)
incidences; adjoining both gives a \(613\)-row target degree set with
zero remaining nonzero proper-subdegree defects for that layer.
Promoting the next newly adjoined \(6\)-row layer forces \(10\) zero
terminal checks and a \(6\)-row saturation extension with \(8\)
incidences; adjoining both gives a \(629\)-row target degree set with
zero remaining nonzero proper-subdegree defects for that layer.
Promoting the next newly adjoined \(6\)-row layer forces \(10\) zero
terminal checks and a \(6\)-row saturation extension with \(18\)
incidences; adjoining both gives a \(645\)-row target degree set with
zero remaining nonzero proper-subdegree defects for that layer.  The
\(\texttt{wrel\_tail\_staircase}\) verifier then proves the symbolic
tail \(A_N\Rightarrow B_N\Rightarrow A_{N+2}\) for every even
\(N\ge14\), and its finite-anchor table checks \(A_{14}\) inside the
\(\texttt{wrel\_degree\_closure}\) packet.  Hence the target-degree
audit has an infinite six-row staircase; its unbounded-tail table
records a strictly increasing coordinate subsequence.  No finite prefix
of this target-only audit is a primitive-recognition theorem.  Neither
displayed target window is yet a primitive-recognition
window satisfying (R0)--(R8).  At
\(W_{\le 3}\)
the target parity
table
\(\delta_i:1|0\), \(a_{ij}:10|0\), \(2\delta_i+\delta_j:1|0\),
\(\delta_{123}:29|93\) is computed; the row \(a_{ij}:10|0\) is the
real bracket \([e_i,e_j]\) together with the nine Gritsenko--Nikulin
isotropic simple directions.  The
\(\texttt{wle3\_target\_parity}\) verifier checks this target table and
its scalar firewall; a compact source implementing
(S3)--(S10) on \(W_{\le 3}\) is open, and the source identification at
\(W_{\mathrm{rel}}^+\) requires a finite source-and-target
presentation computation.
The first relation-closed source theorem is exactly the finite packet
\[
\mathrm{window\_closure},\quad
\mathrm{target\_source\_representatives},\quad
\mathrm{chevalley\_serre\_matrices},\quad
\mathrm{hall\_boundary\_complex},\quad
\mathrm{spectral\_sequence},\quad
\mathrm{kernel\_matrices},\quad
\mathrm{pbw\_graded},\quad
\mathrm{first\_window\_theorems},\quad
\mathrm{transitions},\quad
\mathrm{scalar\_firewall}.
\]
It must compute the compact representatives, Chevalley and Serre
matrices, Hall--Chevalley boundary, kernel bases, no-extra equality,
PBW associated gradeds, and first-window theorem rows.  It is not
provided by the target parity table, the signed Borcherds exponents,
the Pfaffian product, or the scalar trace.
The current packet
\(\texttt{certificates/first\_window/k3e\_relation\_closed\_window}\)
records these missing rows by the obstruction ledger
\(\texttt{FIRST\_WINDOW\_OBSTRUCTION\_LEDGER\_VERIFIED}\), with
\(\texttt{first\_window\_certification=false}\).  Hence Chapter~\ref{ch:pfaffian-dirac}
uses no first-window primitive-recognition theorem beyond this
checked absence status.

\paragraph{Vol III Hall--Drinfeld--Borcherds spectral sequence on compact
non-toric CY3.}  The Hall--Drinfeld--Borcherds spectral sequence
abutting from Vol III's CoHA-type algebra to the BKM denominator on
\(K3\times E\) is open on the compact non-toric trivial-canonical threefold,
even after (D0), (O1), (O1)\(^+\), and (O2).  The
finite-stage Hall source kernel of
Theorem~\ref{thm:finite-stage-dirac-igusa-pro-object} is the
finite-stage substitute.

\paragraph{Protected determinant and character level.}
The Pfaffian and orientation theorems identify the
\textcircled{\scriptsize 2}\(\leftrightarrow\)\textcircled{\scriptsize 3}
edge at the protected determinant and character level on \(K3\times E\).
The primitive Lie-superalgebra recognition remains the finite
compact-source datum \((S1)\)--\((S10)\).  The
cross-level identities
\textcircled{\scriptsize 1}\(\leftrightarrow\)\textcircled{\scriptsize 2}
and \textcircled{\scriptsize 2}\(\leftrightarrow\)\textcircled{\scriptsize 5}
are imported from Borcherds--Gritsenko--Nikulin and the regularized
theta lift.  The
operator-level datum
\(\mathfrak D_X^{\mathrm{DI}}\) is defined at the Pfaffian and
orientation level relative to the retained data of
Appendix~\ref{app:obstruction-discharges}; the
scalar shadow \(Z_{\mathrm{BPS}}^{K3\times E}=\Delta_5^{-2}\) is the
level-\(n\) protected trace.  The compact microscopic state-space
interpretation requires the compact-source recognition datum; the
\(3d\) gravitational path integral requires the Hall--Borcherds
residual; the Calabi--Yau realization requires a
construction external to the Pfaffian theorem.


% ------------------------------------------------------------------
% Part IV: Shadowed — scalar shadows and frontiers
% ------------------------------------------------------------------
\part{Shadowed}


\providecommand{\ClaimStatusConjectured}{}
\providecommand{\ClaimStatusHeuristic}{}
\providecommand{\ClaimStatusConditional}{}
\providecommand{\ClaimStatusProvedElsewhere}{}
\providecommand{\ClaimStatusOpen}{}
\providecommand{\ClaimStatusFolklore}{}

\chapter{The mirror discriminant}
\label{ch:view4-mirror}

\section{The mirror discriminant}
\label{sec:view4-mirror-discriminant}

The automorphic, BKM, Pfaffian, and theta-lift constructions place
\(\Delta_5\) on chains that begin with explicit modular or
representation-theoretic data: the
Borcherds--Gritsenko--Nikulin section of the weight-\(5\) automorphic
line over \(\Sp_4(\Z) \backslash \mathcal H_2\) (View~\textcircled{1}),
the denominator
\(\operatorname{den}(\mathfrak g_{\Delta_5}) = 64^{-1} \Delta_5(2Z)\) of
the Borcherds--Kac--Moody Lie superalgebra on \(\Lambda_{II}^{2,1}\)
(View~\textcircled{2}), the protected Pfaffian
\(\operatorname{Pf}_{\mathrm{prot}}(\mathfrak D_X^{\mathrm{DI}})\) of the
Dirac--Igusa pro-object (View~\textcircled{3}), and the singular theta
lift of the weak Jacobi form \(\phi_{0,1}\), namely the Borcherds
regularized theta integral on the
\(\mathrm O(3,2)\) type-IV domain, with \(\Mp_2\)-Weil input and the
\(\PSp_4\simeq\SO(3,2)_+\) exterior-square bridge
(View~\textcircled{5}).
The period-side assertion is geometric. Let \(\mathcal P_{K3\times E}^{(3)}\) be
the rank-\(3\) Mukai-invariant period quotient transported to
\(\mathcal H_2/\Sp_4(\Z)\).  Let \(\mathfrak H_{\mathrm{II}}\) denote
the type-II Heegner divisor on this quotient.  The automorphic
calculation gives the divisor equality
\begin{equation}
\label{eq:view4-conjecture}
[\,\Delta_5^{-2}\,]
\;=\;
2[\mathfrak H_{\mathrm{II}}]
\qquad
\text{(equality of divisor classes on the }
(K3 \times E)\text{-component)},
\end{equation}
where \([\Delta_5^{-2}]\) is the polar divisor of the meromorphic
section \(\Delta_5^{-2}\) of the automorphic line.  The conjectural
content is the existence of a rank-\(3\) Mukai--Hodge quotient of a
mirror period family whose discriminant support is
\(\mathfrak H_{\mathrm{II}}\), together with the comparison of the
automorphic line with the period Hodge line.

This identification is conjectural. The cross-level identities
\textcircled{4}\(\leftrightarrow\)\textcircled{1},
\textcircled{4}\(\leftrightarrow\)\textcircled{2},
\textcircled{4}\(\leftrightarrow\)\textcircled{5} are open.  The
automorphic identity, the BKM denominator identity, the conditional
Pfaffian identity, and the singular theta-lift identity do not use the
period-discriminant conjecture.

The handle \(\Delta_5\) splits into three names. The automorphic section over
\(\Sp_4(\Z) \backslash \mathcal H_2\) is denoted
\(\mathcal D_X = \Delta_5\) (View~\textcircled{1}); the BKM denominator
is
\(\operatorname{den}(\mathfrak g_{\Delta_5}) = 64^{-1} \Delta_5(2Z)\)
(View~\textcircled{2}); the protected Pfaffian is
\(\operatorname{Pf}_{\mathrm{prot}}(\mathfrak D_X^{\mathrm{DI}})\)
(View~\textcircled{3}). The mirror conjecture concerns the first: the
divisor identification~\eqref{eq:view4-conjecture} is a statement about
\(\mathcal D_X\) as an automorphic section, transported to the period
side.

\subsection{The K3 mirror-period family}
\label{subsec:k3-mirror-period}

The lattice-polarised K3 mirror is fixed by Dolgachev~\cite{DolgachevMirrorLP}
and tracked, throughout, through the Mukai lattice. Let
\[
\Lambda_{K3}
\;=\;
3U \oplus 2 E_8(-1),
\qquad
\widetilde \Lambda_{K3}
\;=\;
U \oplus \Lambda_{K3}
\;\simeq\;
4U \oplus 2 E_8(-1),
\]
the K3 lattice of signature \((3,19)\) and the Mukai lattice of
signature \((4,20)\) (Mukai~\cite{Mukai1984,Mukai1987}). A
\emph{lattice polarisation} of a complex K3 surface \(S\) is a
primitive embedding \(\iota \colon M \hookrightarrow \Pic(S)\) of an
even hyperbolic lattice \(M\), of signature \((1, t)\), such that
\(M^\perp \subset \Lambda_{K3}\) admits a primitive embedding of the
hyperbolic plane \(U\); the pair \((S, \iota)\) is then
\emph{Nikulin-admissible}, in the convention of
Gritsenko--Nikulin~\cite{GNII}\,\S\,3 and the rank-3 sublattice
constructions of Borcherds~\cite{Borcherds95}\,\S\,15. Here
\[
t=\rk M-1
\]
is the negative index of \(M\), so that
\(\rk M^\perp=21-t\) and
\(\mathrm{sign}(M^\perp)=(2,19-t)\).  The convention is
\(\mathrm{sign}(L)=(\#\text{positive},\#\text{negative})\).

The associated period domain
\[
\Omega_M
\;=\;
\bigl\{
[\omega] \in \P\!\bigl(M^\perp \otimes \C\bigr)
\;:\;
\langle \omega, \omega \rangle = 0,
\;
\langle \omega, \bar\omega\rangle > 0
\bigr\}^+
\]
is the type-IV symmetric domain of signature \((2, 19 - t)\) and
dimension \(19-t\), with arithmetic group
\(\Gamma_M = \mathrm O^+(M^\perp)\) acting properly discontinuously.
The \emph{Dolgachev mirror lattice} is the
sub-lattice
\[
\check M
\;=\;
M^\perp \cap (U)^\perp
\;\subset\;
\Lambda_{K3},
\qquad
M^\perp
\;=\;
U \oplus \check M,
\]
so that \(\check M\) carries signature \((1,18-t)\) and is hyperbolic.
The mirror K3 surface \(S^\vee\) is lattice-polarised by \(\check M\).
Dolgachev mirror symmetry gives a correspondence between the
\(M\)-polarised and \(\check M\)-polarised moduli quotients; it is an
involution on a single quotient only in the self-mirror cases.

The Aspinwall--Morrison form of the same construction~\cite{AspinwallMorrisonStringy}
uses the Mukai lattice to compare complex and complexified K\"ahler
directions at large-volume and large-complex-structure cusps.  The full
K3 sigma-model moduli space is larger than \(\Omega_M/\Gamma_M\);
the lattice correspondence is the only input needed for the
\(\Delta_5\) discriminant statement.  The type-IV lattice of the mirror
quotient is
\[
\widetilde M
\;=\;
U_{\mathrm{mir}} \oplus M
\;\simeq\;
\check M^\perp
\;\subset\;
\Lambda_{K3},
\qquad
\mathrm{sign}(\widetilde M)=(2,t+1).
\]
Here \(U_{\mathrm{mir}}\) is the hyperbolic plane split from
\(M^\perp=U_{\mathrm{mir}}\oplus\check M\).  In the full Mukai lattice
\(\widetilde\Lambda_{K3}=U_{\mathrm{Muk}}\oplus\Lambda_{K3}\), the
orthogonal complement of \(U_{\mathrm{Muk}}\oplus M\) is
\(M^\perp=U_{\mathrm{mir}}\oplus\check M\); the extra Mukai plane is
contained in the augmented sublattice.  No canonical action on
\(\widetilde \Lambda_{K3}\) is needed here; only the primitive
decomposition and the induced period domains enter the argument.

\paragraph{The rank-three case.} The polarisation relevant to View~\textcircled{4}
is
\[
M
\;=\;
U \oplus \langle -2 \rangle,
\qquad
\mathrm{rank}(M) = 3,
\qquad
\mathrm{sign}(M) = (1, 2),
\]
with Dolgachev mirror lattice
\[
\check M
\;=\;
\bigl(M^\perp \cap U^\perp\bigr)_{\Lambda_{K3}}
\;\simeq\;
U \oplus E_7(-1) \oplus E_8(-1),
\]
of rank \(17\) and signature \((1, 16)\) inside \(\Lambda_{K3}\); the
convention is
\(\mathrm{sign} = (\#\text{positive}, \#\text{negative})\) on the
bilinear form of the lattice. The decomposition follows from the
embedding \(\langle -2 \rangle \hookrightarrow E_8(-1)\) on a single
root, after which the orthogonal complement of the root inside
\(E_8(-1)\) is \(E_7(-1)\). The rank-five mirror-domain lattice is
\[
\widetilde M
\;=\;
U \oplus M
\;=\;
2 U \oplus \langle -2 \rangle,
\qquad
\mathrm{rank}(\widetilde M) = 5,
\qquad
\mathrm{sign}(\widetilde M) = (2, 3).
\]
Under the exterior-square convention (the model of
\(\PSp_4(\R)\simeq\SO(3,2)_+\) at \S\,\ref{ssec:exceptional-iso})
the same lattice is written
\(\Lambda^{3,2} = U \oplus U \oplus \langle 2 \rangle\), with the
hyperbolic-plane summand
\(\HypPlan = \begin{pmatrix} 0 & -1 \\ -1 & 0 \end{pmatrix}\) carrying
signature \((1, 1)\) and the rank-1 summand carrying form value \(+2\)
on a generator. The two writings are the same lattice expressed in
opposite sign conventions on the bilinear form: the convention
\((p, q) = (3, 2)\) is signature ``three positive, two negative''
for the fixed quadratic form, while the
Dolgachev convention writes the same signature ``two positive, three
negative'' on the form with reversed sign. The
\((+2)\)-vectors in the fixed convention are the \((-2)\)-vectors of the
Dolgachev convention; the Heegner divisors and the period domain
are independent of the sign choice.

In the Dolgachev sign on \(\widetilde M\), the period domain attached
to \(\widetilde M\) is the type-IV symmetric domain
\[
\Omega_{\widetilde M}
\;=\;
\bigl\{
[Z] \in \P\!\bigl(\widetilde M \otimes \C\bigr)
\;:\;
(Z, Z) = 0,\;
(Z, \bar Z) > 0
\bigr\}^+,
\qquad
\dim_{\C} \Omega_{\widetilde M} = 3,
\]
whereas the exterior-square sign of \(\Lambda^{3,2}\) writes the same
domain with \((Z,\bar Z)<0\).  The two inequalities define the same
oriented two-plane after multiplying the bilinear form by \(-1\).
The arithmetic group is
\[
\Gamma_{\widetilde M}
:=\wedge^2\Sp_4(\Z)/\{\pm\Id_4\}
\subset \mathrm P\mathrm O^+(\widetilde M)
\]
acting properly discontinuously, where the component is fixed by the
positive cone \(\mathcal C(\Lambda^{2,1})_+\)
(\S\,\ref{ssec:exceptional-iso}). The exceptional isogeny
\[
\PSp_4(\R)
\;\simeq\;
\SO(3, 2)_+,
\qquad
\Sp_4(\Z) / \{\pm 1\}
\;\simeq\;
\Gamma_{\widetilde M},
\]
imported as View~\textcircled{5} (the \(\PSp_4\simeq\SO(3,2)_+\)
bridge), gives the identification
\begin{equation}
\label{eq:view4-period-bridge}
\Omega_{\widetilde M} \big/ \Gamma_{\widetilde M}
\;\;\simeq\;\;
\mathcal H_2 \big/ \Sp_4(\Z).
\end{equation}
The form
\(\Delta_5 \in M_5(\Sp_4(\Z), \nu_{\Delta_5})\) of
View~\textcircled{1}, viewed as a section on the right-hand side of
\eqref{eq:view4-period-bridge}, is the Borcherds lift of the weak Jacobi
form \(\phi_{0,1}\) under the data
\((\widetilde M, \phi_{0,1})\) on the left-hand side
(Borcherds~\cite[Thm.~13.3]{B} applied to signature
\((2, 3)\); Gritsenko--Nikulin~\cite[Thm.~1.1]{GNII} for the
identification with the additive Gritsenko lift; the rank-3 row of
the Vol~III lattice-polarised \(\mathfrak g_L\) catalogue at
\(L = U \oplus \langle -2 \rangle\)).

The K3-side period domain on the rank-3 polarisation is a different,
larger object:
\[
\Omega_M
\;=\;
\bigl\{
[\omega] \in \P\!\bigl(M^\perp \otimes \C\bigr)
\;:\;
\langle \omega, \omega \rangle = 0,\;
\langle \omega, \bar\omega \rangle > 0
\bigr\}^+,
\qquad
\dim_{\C} \Omega_M = 17,
\]
where \(M^\perp \subset \Lambda_{K3}\) has signature \((2, 17)\) and
parametrises the \emph{transcendental} Hodge filtration of a
lattice-polarised K3 surface. The Borcherds lift used in
View~\textcircled{4} is on \(\Omega_{\widetilde M}\), not on
\(\Omega_M\): the difference is the adjoining of one Mukai hyperbolic
plane to the rank-\(3\) K3 lattice, not an identification of the
even elliptic cohomology with the elliptic period.  The two period spaces are not
interchangeable: \(\Omega_{\widetilde M}\) of dimension \(3\) is the
genus-two Siegel space and supports \(\Delta_5\) directly via the
exceptional isogeny, while \(\Omega_M\) of dimension \(17\) is the
K3-only period space, on which the rank-\(3\) lattice-polarised lift
of \(\phi_{0,1}\) lands only after the additional Mukai-\(U\)
augmentation.

\subsection{The K3 \texorpdfstring{$\times$}{x} E Mukai--Hodge quotient}
\label{subsec:k3xe-period-family}

The full Calabi--Yau-threefold variation of Hodge structure on
\(X=S\times E_\tau\) is
\[
H^3(X,\Z)
\;=\;
H^2(S,\Z)\otimes H^1(E_\tau,\Z),
\]
of rank \(44\).  The elliptic modulus \(\tau\) belongs to
\(H^1(E_\tau)\).  It does not belong to the even elliptic lattice
\(H^0(E,\Z)\oplus H^2(E,\Z)\), whose Hodge type is constant.

The Mukai lattice
\(\widetilde H(S,\Z)=H^*(S, \Z)\) of signature \((4, 20)\) carries the Mukai
pairing
\(\langle (r, c, s), (r', c', s') \rangle_{\mathrm{Muk}} = c \cdot c' - r s' - r' s\)
(Definition~\ref{def:mukai-lattice}).  The even cohomology of the elliptic curve
is the hyperbolic plane
\(U_E = H^{\mathrm{even}}(E, \Z) = H^0(E,\Z)\oplus H^2(E,\Z)\)
with the standard hyperbolic form. Since \(H^{\mathrm{odd}}(S)=0\), the even
integral cohomology is, as an integral group with its Mukai tensor pairing,
\[
H^{\mathrm{even}}(K3 \times E, \Z)
\;=\;
\widetilde H(S,\Z) \otimes U_E
\;=\;
\widetilde H(S,\Z) \otimes U,
\]
a lattice of rank \(48\).  This even Mukai lattice is not the
Calabi--Yau-threefold period lattice; it is the lattice from which the
rank-\(3\) Mukai-invariant quotient is cut.

For \((K3 \times E)\)-pairs \((S, E_\tau)\) with \(S\)
lattice-polarised by \(M = U \oplus \langle -2 \rangle\), the retained
rank-\(3\) Mukai--Hodge quotient is
\begin{equation}
\label{eq:view4-k3e-moduli}
\mathcal P_{K3 \times E}^{(3)}
\;\;:=\;\;
\Omega_{\widetilde M} \big/ \Gamma_{\widetilde M}
\;\;\overset{\eqref{eq:view4-period-bridge}}{\simeq}\;\;
\mathcal H_2 \big/ \Sp_4(\Z),
\end{equation}
where the augmented lattice \(\widetilde M=U_{\mathrm{mir}}\oplus M\)
is a rank-\(5\) type-IV lattice.  The summand \(U_{\mathrm{mir}}\) is
an abstract Mukai hyperbolic plane in the period quotient.  It is not
the constant even elliptic lattice \(U_E\).  The elliptic period enters
through \(H^1(E_\tau)\) and then through the exterior-square
\(\PSp_4\simeq\SO(3,2)_+\) coordinate system on
\(\Omega_{\widetilde M}\).
Equation~\eqref{eq:view4-k3e-moduli} is the period quotient on which
View~\textcircled{4} is stated.  The quotient is the
\(\PSp_4\simeq\SO(3,2)_+\) model; the assertion that it is the base of
the compact mirror discriminant is Conjecture~\ref{conj:mirror-discriminant}.

The K3-only period domain \(\Omega_M\) of dimension \(17\) does not
enter \eqref{eq:view4-k3e-moduli}: its role is to parametrise the
transcendental Hodge filtration of \(S\) alone. The rank-\(3\)
\((K3 \times E)\) Mukai--Hodge quotient is on
\(\Omega_{\widetilde M}\), of dimension \(3\). The relation
\(\widetilde M = U_{\mathrm{mir}} \oplus M\) is the rank-five
type-IV bridge, not an equality with the full \(H^3(K3\times E)\)
period domain.

\subsection{The discriminant locus}
\label{subsec:view4-discriminant}

The candidate period quotient carries a natural Heegner divisor.  It
becomes the support of the discriminant divisor of a mirror period
family only after the comparison asserted in
Conjecture~\ref{conj:mirror-discriminant}.
On the type-IV symmetric domain \(\Omega_{\widetilde M}\), the Heegner
divisors are
\[
\mathfrak H_{(\delta)}
\;=\;
\bigl\{
[Z] \in \Omega_{\widetilde M}
\;:\;
(Z, \delta) = 0
\bigr\},
\qquad
\delta \in \widetilde M,
\;\;
(\delta, \delta) = 2,
\]
indexed by primitive square-\(2\) vectors of the fixed form on
\(\widetilde M = \Lambda^{3,2}\)
(equivalently, square-\((-2)\) vectors of the Dolgachev form after the
sign flip recorded above).  Such a hyperplane records an extra
algebraic Mukai vector; on a K3 slice its \(H^2\)-projection is the
ordinary nodal \((-2)\)-class when the projection is nonzero.  The
type-II Heegner divisor is
\begin{equation}
\label{eq:view4-discriminant-divisor}
\mathfrak H_{\mathrm{II}}
\;\;=\;\;
\bigcup_{\substack{\delta \in \widetilde M \\ (\delta, \delta) = 2\\
(\delta,\widetilde M)=2\Z}}
\mathfrak H_{(\delta)} \big/ \Gamma_{\widetilde M}
\;\;\subset\;\;
\Omega_{\widetilde M} / \Gamma_{\widetilde M}.
\end{equation}
The \(\Gamma_{\widetilde M}\)-orbit structure on primitive square-\(2\)
vectors of \(\widetilde M\) breaks into type-I and type-II classes
(\(\delta \cdot \widetilde M = \Z\) versus \(\delta \cdot \widetilde M = 2\Z\),
the dichotomy already imported as the type-II versus type-I
distinction at the \(W^{(2)}(\Lambda_{II}^{2,1})\)-orbit level); the union in
\eqref{eq:view4-discriminant-divisor} runs over the type-II orbit,
which is the orbit relevant to the \(\Sp_4(\Z)\)-image of the
exceptional isogeny. The image in the quotient is one irreducible
Heegner-type divisor
(Bruinier--Funke~\cite[\S\,2]{BruinierFunkeHeegner} for the general
lattice-Heegner formulation in signature \((2, 3)\); the index-1 case
is implicit in Borcherds~\cite[Thm.~10.1]{B}). The same divisor,
transported across the bridge~\eqref{eq:view4-period-bridge}, is the
\emph{Humbert surface}
\[
\mathcal H_2^{(\mathrm{Hum})}
\;\;=\;\;
\bigl\{
[Z] \in \mathcal H_2 \big/ \Sp_4(\Z)
\;:\;
Z = \begin{pmatrix} \tau_1 & 0 \\ 0 & \tau_2 \end{pmatrix} \;\text{up to } \Sp_4(\Z)
\bigr\},
\]
the diagonal locus on which the principally polarised abelian surface
\(\C^2 / (\Z^2 + Z \Z^2)\) decomposes as a product of two elliptic
curves.  Under the lattice identification
\eqref{eq:view4-period-bridge}, this Humbert surface is the Heegner
support of a norm-\((-2)\) class in the rank-\(3\) Mukai-invariant
period domain.  The assertion that it is the full discriminant support
of a compact mirror family is part of Conjecture~\ref{conj:mirror-discriminant}.

The form \(\Delta_5\) of View~\textcircled{1} vanishes precisely along
this divisor: \(\Delta_5 \in M_5(\Sp_4(\Z), \nu_{\Delta_5})\) has a
simple zero on \(\mathcal H_2^{(\mathrm{Hum})}\)
\cite[Thm.~3.1]{GNPartI}\cite[\S\,3]{GNII}, with multiplier
\(\nu_{\Delta_5}\) accounting for the half-integral shift. Concretely,
the Borcherds product expansion of View~\textcircled{1} gives
\[
\Delta_5(Z)
\;=\;
64\, q^{1/2} r^{1/2} s^{1/2}
\prod_{(n, l, m) \in \Gamma_{\mathrm{eff}}}
(1 - q^n r^l s^m)^{f(nm, l)},
\]
and the automorphic zero divisor is exactly
\(\mathcal H_2^{(\mathrm{Hum})}\), counted with multiplicity \(1\).
Thus the theorem-level divisor identity is
\begin{equation}
\label{eq:view4-divisor-of-Delta5}
[\,\operatorname{div}(\Delta_5)\,]
\;\;=\;\;
[\,\mathcal H_2^{(\mathrm{Hum})}\,]
\quad \text{(divisor classes on } \mathcal H_2 / \Sp_4(\Z)).
\end{equation}
The additional equality with the mirror discriminant
\(\mathfrak D_{\mathrm{mir}}\) is not contained in the automorphic
theorem; it is the support part of the mirror-discriminant conjecture.

\subsection{The conjecture}
\label{subsec:view4-conjecture}

The discriminant of a period quotient carries multiplicities, not just
a support.  Let
\(\mathfrak D_{\mathrm{mir}}\) denote the conjectural reduced
discriminant divisor of that quotient, and let
\(\lambda_{\mathrm{per}}=\mathcal F^2\) denote the descended K3
period Hodge line on the same quotient.  In the rank-\(3\)
Mukai-invariant component the expected section of
\(\lambda_{\mathrm{per}}^{-10}\) has polar order \(2\) along the type-II Heegner divisor:
locally this is the Picard--Lefschetz contribution of a nodal K3
vanishing cycle to the Hodge filtration, and automorphically it is the
square of the simple zero of \(\Delta_5\).  This multiplicity comparison
is part of the conjecture until the rank-\(3\) quotient and the
comparison line have been constructed.

\begin{conjecture}[The mirror discriminant identity]
\label{conj:mirror-discriminant}
\ClaimStatusConjectured
On the rank-\(3\) \((K3 \times E)\) Mukai--Hodge quotient
parametrised by \eqref{eq:view4-k3e-moduli}, the meromorphic section
\[
\Delta_5^{\,-2}
\;\in\;
\Gamma\!\left(
\mathcal H_2 / \Sp_4(\Z),
\;\;
\mathcal L^{-10}_{\,\mathrm{aut}}
\right)
\]
of the inverse \(10\)th tensor power of the automorphic line bundle
\(\mathcal L_{\mathrm{aut}}\) coincides, as a meromorphic section of
\(\lambda_{\mathrm{per}}^{-10}\) on
\(\mathcal P_{K3 \times E}^{(3)}\), with the
discriminant section of the rank-\(3\) quotient. Equivalently, the polar
divisor of \(\Delta_5^{-2}\) is twice the quotient discriminant
\(\mathfrak D_{\mathrm{mir}}\), the support of
\(\mathfrak D_{\mathrm{mir}}\) is \(\mathfrak H_{\mathrm{II}}\), and the comparison map
\[
\mathcal L^{-10}_{\,\mathrm{aut}}
\;\;\xrightarrow{\;\;\sim\;\;}\;\;
\lambda_{\mathrm{per}}^{-10}
\]
on the open complement \(\mathcal P_{K3 \times E}^{\circ}
= \mathcal P_{K3 \times E}^{(3)} \setminus \mathfrak D_{\mathrm{mir}}\)
extends to a global identification with \(\Delta_5^{-2}\) as a
section of \(\lambda_{\mathrm{per}}^{-10}\).
\end{conjecture}

The conjecture is, at the level of divisor classes,
\begin{equation}
\label{eq:view4-divisor-equation}
[\,\Delta_5^{\,-2}\,]
\;\;=\;\;
2\,[\,\mathfrak D_{\mathrm{mir}}\,]
\;\;=\;\;
2\,[\,\mathfrak H_{\mathrm{II}}\,]
\;\;=\;\;
2\,[\,\mathcal H_2^{(\mathrm{Hum})}\,];
\end{equation}
the identity~\eqref{eq:view4-divisor-of-Delta5} fixes the right-hand
side, and the conjectural content is the matching of multiplicities
and the trivialisation of the comparison line bundle.

\paragraph{Evidence.}
The genus-two Bryan--Pandharipande generating function for
disconnected genus-two BPS states on \(K3 \times E\) predicts the cusp
branch governed by the reciprocal Igusa form~\cite[Conj.~1]{BRYAN};
Oberdieck--Pixton prove the normalized OP/DT/PT scalar representative
\[
Z^X_{\mathrm{OP}}=-4096\,\Delta_5^{-2}
\]
\cite[Thm.~1]{OB}.  In the trace convention of this manuscript
\[
Z_{\mathrm{BPS}}^{K3\times E}:=-4096^{-1}Z^X_{\mathrm{OP}}
=\Delta_5^{-2}.
\]
This is the input to View~\textcircled{1} on the curve-counting side.  The local
Picard--Lefschetz model and Saito's Hodge-filtration formalism predict
the same order-two degeneration as the automorphic divisor
\([\Delta_5^{-2}]\).  The unproved point is the global comparison: the
line
\(\mathcal L^{-10}_{\,\mathrm{aut}} \otimes \lambda_{\mathrm{per}}^{10}\)
must be trivialised and the identification of the two sides must extend
across the discriminant divisor.

A separate piece of evidence is the rank-\(4\) precedent set by
Gritsenko--Nikulin~\cite[Thm.~4.1]{GNII} in the lattice-polarised
setting at the polarisation \(L\) producing a Borcherds product
identified with the quasi-pull-back of the discriminant section of
the universal lattice-polarised Hodge bundle on
\(\Omega_L / \Gamma_L\). Conjecture~\ref{conj:mirror-discriminant} is
the rank-\(3\) analogue of the same identification, evaluated at the
type-II rank-\(3\) Mukai sublattice
\(M = U \oplus \langle -2 \rangle\) of \(\Lambda^{2,1}\)
(see Appendix~\ref{appendix:mukai-gram-dictionary}); the Fake-Monster Borcherds
product \(\Phi_{12}\) on the Leech-extension lattice
\(\mathrm{II}_{25,1}\) of \cite{BorcherdsGKM88, Borcherds95} is a
distinct, signature-\((26,2)\) object and is not invoked here.

\paragraph{Epistemic strength.}
\ClaimStatusConjectured\ The identification claim of
Conjecture~\ref{conj:mirror-discriminant} is conjectured. The
comparison line bundle is a standard Hodge-theoretic object; the
automorphic Heegner identity~\eqref{eq:view4-divisor-of-Delta5} is a
theorem; the equality
\(\operatorname{Supp}(\mathfrak D_{\mathrm{mir}})=\mathfrak H_{\mathrm{II}}\),
the multiplicity \(2\), and the global extension across the
discriminant divisor are conjectural.

\paragraph{Consequences.}
If proved, View~\textcircled{4} embeds the \(\Delta_5\) identity inside
mirror symmetry on \(K3 \times E\) at the level of period-side
discriminants.  The identification would give a period-side derivation
of \(\Delta_5\) independent of the Borcherds product
(View~\textcircled{1}), the BKM denominator (View~\textcircled{2}),
the Pfaffian of the Dirac--Igusa pro-object (View~\textcircled{3}), and
the singular theta lift (View~\textcircled{5}).  It is a conjectural
fifth chain.

A proof of Conjecture~\ref{conj:mirror-discriminant} would identify
\(\Delta_5^{-2}\) with the section of \(\lambda_{\mathrm{per}}^{-10}\)
on the rank-\(3\)
\((K3 \times E)\) Mukai--Hodge quotient; together with the Bryan--Oberdieck--Pixton
theorem~\cite[Thm.~1]{OB} that
\(Z^X_{\mathrm{OP}}=-4096\,\Delta_5^{-2}\), this would give a
period-side interpretation of the unscaled trace
\(Z_{\mathrm{BPS}}^{K3\times E}:=-4096^{-1}Z^X_{\mathrm{OP}}\).

\subsection{Range of the period conjecture}
\label{subsec:view4-not-asserted}

The Beilinson cut is in force: a statement is not primitive if it
holds only after passing to a trace, choosing a chamber, completing a
category, or installing descent data.
The conjectural identification is on the period side: it concerns
the section \(\Delta_5^{-2}\) of an automorphic line bundle and the
discriminant divisor of the rank-\(3\) \(K3\times E\) mirror period
component. It is not a
claim about the protected operator-level datum
\(\mathfrak D_X^{\mathrm{DI}}\) of View~\textcircled{3}; that datum
is established at the Pfaffian and orientation level only relative to
the retained data in the
View~\textcircled{2}\(\leftrightarrow\)\textcircled{3} chapters and
Appendix~\ref{app:obstruction-discharges}. View~\textcircled{4} does
not enter the operator chain.

In the trace convention of Chapter~\ref{ch:zbps-realization}, the
level-\(n\) scalar partition function
\(Z^{K3 \times E}_{\mathrm{BPS}}:=\Delta_5^{-2}\) is the
orientation-forgetful trace of the Pfaffian datum;
its three-dimensional gravity-line realization requires the
Hall--Borcherds residual (Vol~II, \cite{ChiralBarCobarVolII}). The
mirror conjecture above says nothing about the gravitational lift; it
addresses only the geometric origin of the section
\(\Delta_5^{-2}\) on the rank-\(3\) mirror period component.

The Borcherds--Kac--Moody denominator
\(\operatorname{den}(\mathfrak g_{\Delta_5}) = 64^{-1} \Delta_5(2 Z)\)
is on the algebraic side; it is a theorem of
Gritsenko--Nikulin~\cite[\S\,3, Thm.~3.1]{GNII}. The mirror
conjecture is on the period side and does not enter the algebraic
chain.

\subsection{Lattice-convention check}
\label{subsec:view4-lattice-check}

The lattice convention adopted above is the Dolgachev convention
\(\Lambda_{K3} = 3 U \oplus 2 E_8(-1)\), \(M = U \oplus \langle -2 \rangle\),
mirror lattice
\(\check M \simeq U \oplus E_7(-1) \oplus E_8(-1)\) (rank \(17\),
signature \((1,16)\)). Aspinwall--Morrison's original mirror-pair
construction~\cite{AspinwallMorrisonStringy} adopts the same lattice up
to the choice of which hyperbolic plane in \(\Lambda_{K3}\) is split
off; the two conventions agree on the K3 mirror lattice. The
Mukai-lattice form
\(\widetilde M = U \oplus M = 2U \oplus \langle -2\rangle\)
of rank 5 and signature \((2, 3)\) is the convention adopted in
Vol~III~\cite[Prop.~6.4]{CYQGVolIII}; it is the convention
relevant to the period-domain bridge~\eqref{eq:view4-period-bridge}.
The lattice \(\widetilde M\) is identified with
\(\Lambda^{3,2}\) of \S\,\ref{ssec:exceptional-iso} by the sign-flip
on the rank-\(1\) summand recorded in \S\,\ref{subsec:k3-mirror-period}.
Disagreement on the lattice convention appears
only as a sign convention on the Mukai pairing; it does not affect the
discriminant divisor.

The K3 vs.\ \((K3 \times E)\) period-space matching
\eqref{eq:view4-k3e-moduli} adjoins the distinguished hyperbolic
summand \(U_{\mathrm{mir}}\) to the rank-\(3\) K3 lattice and lets the
elliptic variation enter through \(H^1(E_\tau)\).  It does not
identify \(U_{\mathrm{mir}}\) with the constant even elliptic lattice
\(U_E\), and it is not an abstract equality between the
K3-mirror period family and the \((K3 \times E)\)-mirror period family.
The possible discrepancy is exactly the comparison between the
\(\Sp_4\)-side and the \(\mathrm O^+(2,3)\)-side; the
isomorphism~\eqref{eq:view4-period-bridge} is
exactly the \(\PSp_4\simeq\SO(3,2)_+\) bridge of
View~\textcircled{5}.

The word discriminant names a section, not only a support, once
multiplicities are present. There are two natural divisors on the
period side: the Heegner support \(\mathfrak H_{\mathrm{II}}\), and
the polar divisor of the section of \(\lambda_{\mathrm{per}}^{-10}\).
The conjecture concerns the second:
\(\Delta_5^{-2}\) is conjectured to be that section, with polar
divisor twice the support divisor \(\mathfrak H_{\mathrm{II}}\).  The
factor of \(2\) is predicted by the
local Picard--Lefschetz model and is matched, on the automorphic side,
by the square of \(\Delta_5\); a mismatch at the level of multiplicity
would falsify the conjecture by a factor that the Borcherds product
does not allow.

The relation to the Borcherds--Gritsenko--Nikulin quasi-pull-back
theorems is the cleanest comparison.  There the Borcherds form
\(\Phi_{12}\) on the even unimodular lattice of signature \((2,26)\)
is pulled back to lattice-polarised type-IV domains, and its divisor is
read as a discriminant section of the corresponding Hodge bundle
\cite[Thm.~4.1]{GNII}.  The rank-\(3\) lattice
\(L = U \oplus \langle -2 \rangle\) requires a separate comparison line
and a separate Heegner-divisor analysis; the \(\Phi_{12}\) theorem is
evidence, not a proof.

\subsection{Components of the period conjecture}
\label{subsec:view4-period-components}

\begin{enumerate}
\item[\textnormal{(M1)}]
\ClaimStatusOpen\ \emph{Trivialisation of the comparison line.} The
comparison line bundle
\(\mathcal L^{-10}_{\,\mathrm{aut}} \otimes \lambda_{\mathrm{per}}^{10}\)
on \(\mathcal P_{K3 \times E}^{\circ}\) must be trivialised.  Since
\(\nu_{\Delta_5}^2=1\), the inverse-square section has no residual
character; the obstruction is an ordinary Hodge line class.  A
trivialisation of this line, together with the divisor and boundary
conditions below, strengthens the divisor comparison to an equality of
meromorphic sections.

\item[\textnormal{(M2)}]
\ClaimStatusOpen\ \emph{Extension and multiplicity.}
The section comparison is first defined on
\(\mathcal P_{K3 \times E}^{\circ}\), where the Heegner divisor has
been removed.  The period condition is the extension across
\(\mathfrak H_{\mathrm{II}}\) with polar order \(2\), together with
the extension to the rank-2 cusp with boundary exponent governed by
the Borcherds singular weight \(c_1(0)/2=5\)
(\cite[\S\,15]{Borcherds95}).  Thus the divisor of the extended
section, not the open part itself, is compared with
\([\Delta_5^{-2}]\).

\item[\textnormal{(M3)}]
\ClaimStatusOpen\ \emph{The mirror correspondence on the section.}
Dolgachev mirror symmetry gives a correspondence between the
\(M\)-polarised and \(\check M\)-polarised period quotients, not a
canonical involution of \(\Omega_{\widetilde M}/\Gamma_{\widetilde M}\).
The period condition is that a chosen correspondence identifies the
period Hodge line with the automorphic line and carries the
type-II Heegner divisor to the type-II Heegner divisor.  After this
comparison, \(\Delta_5^{-2}\) is required to be the same meromorphic
section of the compared inverse Hodge line.

\item[\textnormal{(M4)}]
\ClaimStatusOpen\ \emph{Rank \(\geq 5\) lattice compatibility.}
The Vol~III rank-\(\geq 5\) lattice-polarised \(\mathfrak g_L\)
family~\cite{CYQGVolIIILatticePolarized} is the conjectural extension of
Conjecture~\ref{conj:mirror-discriminant} to higher Picard rank.
At rank \(\geq 5\) the discriminant identity is conditional on the
Gritsenko--Cl\'ery 2018 admissibility conjecture
(Conj.~5.1 of \cite{GC}); the rank-\(3\) statement is the corresponding
conjectural statement for the K3-Mukai-invariant lattice and is the
rank relevant to the \(\Delta_5\) row.
\end{enumerate}

The four period conditions \((\mathrm M 1)\)--\((\mathrm M 4)\) are
distinct from the four obstruction classes \((\mathrm{D0})\),
\((\mathrm{O1})\), \((\mathrm{O1}^+)\), \((\mathrm{O2})\) of
View~\textcircled{3}; the latter govern the operator-level chain
and are the conditioning for \textcircled{2}\(\leftrightarrow\)\textcircled{3},
while the former govern the period-side chain and are the
conditioning for \textcircled{4}\(\leftrightarrow\)\textcircled{1}.
The two systems lie on different chains: \((\mathrm M 1)\)--\((\mathrm M 4)\)
are period-line and Heegner-divisor conditions, whereas
\((\mathrm{D0})\), \((\mathrm{O1})\), \((\mathrm{O1}^+)\), and
\((\mathrm{O2})\) are Pfaffian-line and orientation conditions.

\subsection{Compatibility with the Vol III lattice convention}
\label{subsec:view4-volIII-lattice-convention}

The lattice convention
\(M = U \oplus \langle -2 \rangle\),
\(\check M \simeq U \oplus E_7(-1) \oplus E_8(-1)\),
\(\widetilde M = U \oplus M\) of
\S\,\ref{subsec:k3-mirror-period} is the convention of
Vol~III~\cite[\S\,6.4 + Prop.~6.4 + Thm.~6.7]{CYQGVolIII}; on the
\((K3 \times E)\)-component the Vol~III \(\kappa\)-tuple
\(\bigl(\kappa_{\mathrm{cat}}, \kappa_{\mathrm{ch}}^{\mathrm{Hodge}},
\kappa_{\mathrm{ch}}^{\mathrm{Heis}}, \kappa_{\mathrm{BKM}}, \kappa_{\mathrm{fiber}}\bigr)
= (0, 0, 3, 5, 24)\) is recorded with each entry sourced from a
distinct construction. The mirror discriminant of
Conjecture~\ref{conj:mirror-discriminant} is on the
\(\kappa_{\mathrm{BKM}} = 5\) row of the \(\kappa\)-tuple; the
period-side datum is the discriminant divisor on the
lattice-polarised period stratum of the Vol~III chart
classification (Vol~III~\cite[\S\,V]{CYQGVolIII}: the
``lattice-polarised period domain'' equivariance stratum). The
\((K3 \times E)\)-fiber rank \(\kappa_{\mathrm{fiber}} = 24\) is the
Mukai-lattice rank of \(\widetilde \Lambda_{K3}\); it is the rank of
the Heisenberg algebra of the Mukai lattice, not a discriminant
invariant.

The mixed-holomorphic-topological-strings companion volume
\cite{MixedHTSCompanion} records that the modular line bundle on
\(\overline{\mathcal A}_g\) at level \(g = 2\) is
\(\Omega_{\mathrm{central}}\), and that
\(\Phi_{10}^{\mathrm{un}} = \chi_{10} = \Delta_5^{\,2}\) is the genus-\(2\)
section.  The mirror-discriminant conjecture is consistent with that statement:
the polar divisor of \(\Delta_5^{-2} = (\Phi_{10}^{\mathrm{un}})^{-1}\) on
\(\overline{\mathcal A}_2\) is twice the divisor of \(\Delta_5\),
and the mirror discriminant identification is precisely the
geometric origin of that section.

The Vol~II Hall--Borcherds residual is silent on
View~\textcircled{4}: the period-side discriminant identification
is independent of the gravitational-lift bridge, and the
conjecture above does not require the residual.
Solving the Hall--Borcherds residual on the
\((K3 \times E)\)-component would not by itself prove
Conjecture~\ref{conj:mirror-discriminant}; conversely, a proof
of Conjecture~\ref{conj:mirror-discriminant} would not
automatically construct the Hall--Borcherds residual.

\medskip

View~\textcircled{4} remains conjectural.  The automorphic, BKM,
Pfaffian, and theta-lift views form the theorem-level quadrilateral at
the stated categorical dimensions.  The mirror-discriminant identity is
the remaining period-side equality.


\chapter{The level-\(n\) protected scalar shadow and the \(K3\times E\) realization question}
\label{ch:zbps-realization}

The denominator algebra and the protected Pfaffian sit at the
primitive and chart levels of the Beilinson chain
\(\mathsf{P}\to\mathsf{C}\to\mathsf{S}\to\mathsf{Z}\to\mathsf{A}\).
After taking the protected trace, the same \(\Delta_5\) survives,
squared, as the partition function on the closed
\(K3\times E\to\mathrm{pt}\) cobordism, equivalently the level-\(n\)
modular-functor scalar trace of the realization formalism of Vol IV
\cite{ChiralBarCobarVolIV}.  The scalar shadow is exactly the object
allowed by the Beilinson cut: an automorphic identity,
not an algebra; a trace, not an operator; an orientation-forgetful
invariant, not an orientation lift.  The orientation-level object is
\(\mathfrak D_X^{\mathrm{DI}}\) of Chapter~5; the scalar identity is
the level-\(n\) shadow it casts.

Three distinct meanings of \(\Delta_5\) coexist and are tagged on
first occurrence: the automorphic section
\(\mathcal D_X=\Delta_5\) (View I --- Borcherds--Gritsenko--Nikulin); the
Borcherds--Kac--Moody denominator
\(\operatorname{den}(\mathfrak g_{\Delta_5})=64^{-1}\Delta_5(2Z)\)
(View II --- Gritsenko--Nikulin); the protected Pfaffian
\(\mathrm{Pf}_{\mathrm{prot}}(\mathfrak D_X^{\mathrm{DI}})\) (View
III --- Chapter~5).  The handle ``\(\Delta_5\)'' alone never
substitutes for any of the three.

\section{The level-\(n\) protected scalar shadow on \(K3\times E\to\mathrm{pt}\)}
\label{sec:zbps-scalar-shadow}

\subsection{The Oberdieck--Pixton scalar identity}

Let \(X=S\times E\), with \(S\) a smooth complex K3 surface and \(E\) a
smooth elliptic curve.  The reduced Donaldson--Thomas / stable-pair /
Gromov--Witten correspondence on \(X\), in the primitive K3-class
chamber, gives the normalized scalar representative
\[
Z^X_{\mathrm{OP/DT/PT}}=-4096\,\Delta_5^{-2}
\]
on the genus-\(2\) Siegel upper half-space \(\mathfrak H_2\).  The
unscaled level-\(n\) trace retained below is
\(Z_{\mathrm{BPS}}^{K3\times E}=\Delta_5^{-2}\).

\begin{convention}[OP chamber, sign branch, primitive class]
\label{conv:op-chamber}
The variables
\((T_{\mathrm{OP}},Q_{\mathrm{OP}},P_{\mathrm{OP}})\) are the
Oberdieck--Pixton variables of \cite[\S0]{OB}:
\(T_{\mathrm{OP}}\) is the K3-surface degree,
\(Q_{\mathrm{OP}}\) the K3-curve class, and
\(P_{\mathrm{OP}}\) the elliptic-genus variable
(equivalently \(P_{\mathrm{OP}}=-p_{\mathrm{DT}}\), giving the
\((-p_{\mathrm{DT}})^n\) DT-chamber).  The K3 curve class
\(\beta_h\) is taken \emph{primitive} of self-intersection
\(\beta_h\cdot\beta_h=2h-2\) (\cite{MaulikPandharipande}); the
identity \eqref{eq:zbps-op-dt} is the primitive-class branch.
The leading \(-1\) of \(Z^X_{\mathrm{OP}}=-(\chi_{10}^{\mathrm{OP}})^{-1}\)
is the OP scalar convention; the squared theta-leading
\(4096=64^2\) is the normalization conversion
\(\chi_{10}^{\mathrm{OP}}=4096^{-1}\Delta_5^2\).  None of these is the
orientation character.
\end{convention}

The scalar identity \eqref{eq:zbps-op-dt} is an identity in the OP chamber.  If the
retained compact Hall source lies in a different stability chamber,
comparison with \eqref{eq:zbps-op-dt} requires the chamber-compatible
Hall descent datum of Definition~\ref{def:chamber-compatible-hall-descent};
the normal-ordered lattice map alone does not transport Hall brackets.

Under Convention~\ref{conv:op-chamber} and \cite[Theorem~1]{OB},
\cite[Proposition~5]{OPand}, \cite{MaulikPandharipande},
\begin{equation}
\label{eq:zbps-op-dt}
Z^X_{\mathrm{DT}}
\;=\;
Z^X_{\mathrm{OP}}(P_{\mathrm{OP}},Q_{\mathrm{OP}},T_{\mathrm{OP}})
\;=\;
-\,4096\,\Delta_5(Z)^{-2},
\qquad
Z=\Omega\in\mathfrak H_2.
\end{equation}
The constant \(4096=64^2\) is the squared theta-leading coefficient
\([q^{1/2}r^{1/2}s^{1/2}]\Delta_5\) of the automorphic section, and
the leading sign is the Oberdieck--Pixton scalar convention
\(Z^X_{\mathrm{OP}}=-(\chi_{10}^{\mathrm{OP}})^{-1}\).  Both are
normalization data of the imported reduced theory.

\begin{theorem}[Oberdieck--Pixton scalar branch and unscaled Igusa shadow]
\label{thm:zbps-level-n-shadow}
\ClaimStatusProvedElsewhere
With variables
\((T_{\mathrm{OP}},Q_{\mathrm{OP}},P_{\mathrm{OP}})\) normalized in the
Oberdieck--Pixton chamber of the Oberdieck--Pixton normalization
\cite[Theorem~1]{OB}, the proved reduced DT/PT scalar branch is
\[
Z^X_{\mathrm{DT}}=Z^X_{\mathrm{OP}}=Z^X_{\mathrm{PT}}
=-4096\,\Delta_5(Z)^{-2}.
\]
\ClaimStatusDefinition\ In the unscaled trace convention of this
manuscript, the level-\(n\) protected scalar shadow on the closed
\(K3\times E\to\mathrm{pt}\) cobordism is defined by
\begin{equation}
\label{eq:zbps-equality}
Z_{\mathrm{BPS}}^{K3\times E}(T_{\mathrm{OP}},Q_{\mathrm{OP}},P_{\mathrm{OP}})
\;=\;
\bigl(\Phi_{10}^{\mathrm{un}}(Z)\bigr)^{-1}
\;=\;
\Delta_5(Z)^{-2}.
\end{equation}
Equivalently, in Oberdieck--Pixton normalization
\(\chi_{10}^{\mathrm{OP}}=2^{-12}\Delta_5^2\),
\(\Phi_{10}^{\mathrm{un}}=\chi_{10}=\Delta_5^2\), and the unscaled
form \((\Phi_{10}^{\mathrm{un}})^{-1}\) is the retained level-\(n\)
trace convention on each primitive closed K3-class fibre in the OP
chamber.
The normalization packet
\(\texttt{certificates/normalizations/delta5\_theta\_leading}\)
checks the named convention rows
\(\Delta_{10}\), \(\chi_{10}\), \(\Phi_{10}^{\mathrm{un}}\),
\(\chi_{10}^{\mathrm{OP}}\), \((\Phi_{10}^{\mathrm{un}})^{-1}\),
\((\chi_{10}^{\mathrm{OP}})^{-1}\), and \(Z^X_{\mathrm{OP}}\)
separately; the scalar equalities below are relations among those rows,
not a collapse of the conventions.
The divisor packet
\(\texttt{certificates/automorphic/delta5\_humbert\_divisor}\)
checks that the polar divisor of the inverse square
\(\Delta_5^{-2}\) has order \(2\) on the same Humbert support on which
\(\Delta_5\) has order \(1\); this is automorphic divisor arithmetic,
not the mirror-discriminant assertion of Chapter~\ref{ch:view4-mirror}.
The Behrend-weighted reduced quotient
integration of Oberdieck's reduced PT/DT scalar-quotient integration
\cite{OberdieckReducedPT} gives the OP chamber branch
\[
Z^X_{\mathrm{DT}}=Z^X_{\mathrm{OP}}=Z^X_{\mathrm{PT}}
=-4096\,Z_{\mathrm{BPS}}^{K3\times E}.
\]
Thus the OP branch is a normalized scalar representative of the same
unscaled trace, not the definition of \(Z_{\mathrm{BPS}}^{K3\times E}\).
\end{theorem}

\begin{proof}
Equation~\eqref{eq:zbps-equality} is the dictionary between the
Oberdieck--Pixton primitive series and the unscaled Igusa form
\(\Phi_{10}^{\mathrm{un}}\): the Oberdieck--Pixton normalization
\(\chi_{10}^{\mathrm{OP}}=D_5^2=64^{-2}\Delta_5^2=2^{-12}\Delta_5^2\)
of \cite[\S0]{OB} together with the Igusa identification
\(\Phi_{10}^{\mathrm{un}}=\chi_{10}=\Delta_5^2\) of
\cite{Igusa1962,Igusa1964II} gives
\((\Phi_{10}^{\mathrm{un}})^{-1}=\Delta_5^{-2}\).  Multiplying by
\(2^{12}=4096\) and the OP scalar sign restores
\eqref{eq:zbps-op-dt}.  The reduced DT/PT comparison is
\cite[Theorem~3(i)]{OberdieckReducedPT}, the reduced PT/GW comparison
on primitive K3 classes is
\cite[Proposition~5]{OPand}, and the OP product identity is
\cite[Theorem~1]{OB}; their composition gives the displayed equality
in the OP chamber.
\end{proof}

\subsection{Three readings of the same exponent}

The integers \(f(n,l)\) of the weight-zero index-one weak Jacobi form
\(\phi_{0,1}\) drive both sides of the OP/Borcherds dictionary.  At
the primitive datum they are virtual \(K_0\)-superdimensions of
\(\mathbb U_{n,l}\), and after taking scalar shadows they appear twice:
once as exponents of the Borcherds product on the automorphic side,
once as multiplicities in the Mukai dimension count of the
Hilbert-scheme stratification on the geometric side.  In
\(\mathfrak g_{\Delta_5}\), the same \(f(nm,l)\) is the signed
root supermultiplicity at \(\alpha(n,l,m)=2nf_2-lf_3+2mf_{-2}\).

The variable map of the Oberdieck--Pixton normalization \cite[Theorem~1]{OB} encodes
this triple reading.  A factor \((1-q^nr^ls^m)^{f(nm,l)}\) of the
Borcherds product on the automorphic side is, on the Oberdieck--Pixton
side, a factor
\((1-p_{\mathrm{OP}}^kq_{\mathrm{OP}}^h\widetilde q_{\mathrm{OP}}^d)^{c(4hd-k^2)}\)
of \(\chi_{10}^{\mathrm{OP}}\), the K3-elliptic-genus exponents
\(c(4hd-k^2)=2f(hd,k)\) recording the squared automorphic line.  The
multiplicities \(2f\) are exactly the K3 elliptic-genus coefficients
of \(Z_{\mathrm{K3}}=2\phi_{0,1}\); the squaring is part of the
identification \(\chi_{10}^{\mathrm{OP}}=D_5^2\) of
the Oberdieck--Pixton/Pandharipande scalar identity \cite[Theorem~1]{OB}.

\begin{remark}[Constants are normalization, not orientation]
\label{rem:zbps-constants-not-orientation}
Three constants intervene in \eqref{eq:zbps-op-dt}.  The leading
\(64=[q^{1/2}r^{1/2}s^{1/2}]\Delta_5\) is the theta-constant
normalization of the Igusa section; the monic Borcherds product is
\(D_5=64^{-1}\Delta_5\).  The factor \(4096=64^2\) is the
squared theta-leading coefficient produced by the OP convention
\(\chi_{10}^{\mathrm{OP}}=D_5^2\).  The leading minus sign is the OP
scalar branch \(Z^X_{\mathrm{OP}}=-(\chi_{10}^{\mathrm{OP}})^{-1}\).
None of the three is the orientation character.  The protected
orientation character \(\epsilon_o\) and the automorphic reflection
character \(\nu_{\Delta_5}\) are determined elsewhere: \(\nu_{\Delta_5}\)
by the Maass transformation law of \(\Delta_5\), \(\epsilon_o\) by the
rank-one type-II Pfaffian signs together with the retained
half-Hilbert wall atlas.  The number \(4096=2^{12}\) occurs in
two different structures, not as one orientation datum: the
Pfaffian-side \(2^{12}\) is the cardinality of the retained sign
orbit, while the OP-side \(2^{12}\) is
the squared theta-leading coefficient in the imported scalar identity.
\end{remark}

\subsection{Worked Fourier-coefficient check}

The leading Fourier coefficients of \(\Delta_5\) and \(\Delta_5^2\)
provide the elementary verification that the constants in
\eqref{eq:zbps-op-dt} are normalization, not enumerative content.
Recall \(\phi_{0,1}\) has leading expansion
\begin{equation}
\label{eq:zbps-phi01-leading}
\phi_{0,1}(\tau,z)
\;=\;
(r^{-1}+10+r)+q\bigl(10r^{-2}-64r^{-1}+108-64r+10r^2\bigr)+O(q^2),
\end{equation}
with coefficients
\(f(0,-1)=f(0,1)=1\), \(f(0,0)=10\), \(f(1,1)=f(1,-1)=-64\),
\(f(1,0)=108\), \(f(1,2)=f(1,-2)=10\), \(f(1,3)=f(1,-3)=0\).
The Borcherds--Gritsenko--Nikulin product
\begin{equation}
\label{eq:zbps-bgn-product}
\Delta_5(Z)
\;=\;
64\,q^{1/2}r^{1/2}s^{1/2}
\prod_{(n,l,m)\in\Gamma_{\mathrm{eff}}}
(1-q^nr^ls^m)^{f(nm,l)}
\end{equation}
opens with the theta-leading term
\([q^{1/2}r^{1/2}s^{1/2}]\Delta_5=64\); the next non-trivial
timelike coefficient belongs to the monic product after the leading
monomial is removed:
\[
[qrs]\,D_5/(qrs)^{1/2}=93,\qquad D_5=64^{-1}\Delta_5 .
\]
The isolated factor \((1-qrs)^{f(1,1)}=(1-qrs)^{-64}\) contributes
\(64qrs\), and the remaining root factors contribute the other
\(29qrs\), as in Appendix~\ref{app:borcherds-product-expansion}.
Squaring \eqref{eq:zbps-bgn-product},
\begin{equation}
\label{eq:zbps-delta5-squared-leading}
\Delta_5^2(Z)
\;=\;
4096\,q\,r\,s
\prod(1-q^nr^ls^m)^{2f(nm,l)},
\end{equation}
with theta-leading \(64^2=4096\).  The OP product
\(\chi_{10}^{\mathrm{OP}}=p\,q\,\widetilde q\prod(1-p^kq^h\widetilde
q^d)^{c(4hd-k^2)}\), with K3-elliptic exponents \(c(4hd-k^2)=2f(hd,k)\),
is exactly the right-hand side of \eqref{eq:zbps-delta5-squared-leading}
divided by \(4096\).  Thus
\begin{equation}
\label{eq:zbps-leading-coeff-check}
\chi_{10}^{\mathrm{OP}}=4096^{-1}\,\Delta_5^2,
\qquad
[pq\widetilde q]\,\chi_{10}^{\mathrm{OP}}=1,
\qquad
[q^{1/2}r^{1/2}s^{1/2}]\,(64^{-1}\Delta_5)=1.
\end{equation}
The squared theta-leading coefficient \(4096\) is therefore the
single explicit normalization constant in the OP/DT scalar
identity \(Z^X_{\mathrm{OP}}=-(\chi_{10}^{\mathrm{OP}})^{-1}\); the
sign \(-1\) is OP convention; together they produce the constant
\(-4096\) of \eqref{eq:zbps-op-dt}.  The Fourier coefficient at
\(q^{1/2}r^{1/2}s^{1/2}\) identifies the scalar \(4096\) as
\(\bigl([q^{1/2}r^{1/2}s^{1/2}]\Delta_5\bigr)^2\), a leading-coefficient
square, not an enumerative invariant.

\subsection{Cobordism positioning and the closed \(K3\times E\to\mathrm{pt}\)
cobordism}

The cobordism \(K3\times E\to\mathrm{pt}\) is the closed
Calabi--Yau-threefold target viewed in the retained trace convention.
The trace is the Laurent expansion of the
meromorphic automorphic section \(\Delta_5^{-2}\); evaluation at a
point of \(\mathfrak H_2\) outside the polar divisor gives a complex
number.  The ``level-\(n\)'' label is the curve-counting integer
\[
n=\chi(\mathcal O_Y)
\]
for a one-dimensional subscheme \(Y\subset K3\times E\) of class
\((\beta_h,d)\).  The target has
\[
\chi(\mathcal O_{K3\times E})
=\chi(\mathcal O_{K3})\chi(\mathcal O_E)=2\cdot 0=0,
\]
and this target Euler characteristic is not the exponent \(n\).  The
Laurent variables record \(n\), the elliptic degree \(d\), and the K3
norm \(h-1\).  The dual notation
\(Z_{\mathrm{BPS}}^{K3\times E}(T_{\mathrm{OP}},Q_{\mathrm{OP}},P_{\mathrm{OP}})\)
stores the elliptic and K3
gradings as a Laurent expansion of the inverse Igusa cusp form on
\(\mathfrak H_2\).

The \(K3\times E\) target is a threefold; the perfect obstruction
theories and reduced virtual fundamental classes live on the
stable-pair and DT moduli spaces
\(P_n(X,(\beta_h,d))\) and \(I_n(X,(\beta_h,d))\), not on the target
itself and not on any auxiliary fourfold.  The
closed K3-class fibre at primitive curve class \(\beta_h\) of norm
\(2h-2\) supplies the K3 component of the coefficient indexed by
\((h,n,d)\), not a new target dimension.

\section{Spin \(L\)-factor normalization for \(\Delta_5^2\)}
\label{sec:spin-l-factor}

\subsection{The Saito--Kurokawa eigenline}

The unscaled Igusa form \(\Phi_{10}^{\mathrm{un}}=\Delta_5^2\) is the
weight-\(10\) Saito--Kurokawa lift of the unique normalized cusp form
\(g=\Delta E_6\in S_{18}(\SL_2(\Z))\); the Hecke eigenline is the
one-dimensional \(\C\Delta_{10}^{\theta}=\C\chi_{10}^{\mathrm{OP}}\),
independent of the scalar normalization of its generator.

\begin{proposition}[Hecke eigenvalues on the Saito--Kurokawa
eigenline]
\label{prop:zbps-sk-hecke}
For each prime \(p\), let \(T_1(p)\) and \(T_2(p^2)\) denote the
degree-one and degree-two genus-two Hecke operators on
\(S_{10}(\Sp_4(\Z))\) of \cite[\S 4]{Andr74}.  The Saito--Kurokawa
form \(\Delta_{10}^{\theta}=\Delta_5^2\) is a simultaneous
eigenform with eigenvalues
\[
 T_1(p)\,\Delta_{10}^{\theta}
 \;=\;
 \lambda_1^{\mathrm{SK}}(p)\,\Delta_{10}^{\theta},
 \qquad
 T_2(p^2)\,\Delta_{10}^{\theta}
 \;=\;
 \lambda_2^{\mathrm{SK}}(p^2)\,\Delta_{10}^{\theta},
\]
where the eigenvalues are determined by the Saito--Kurokawa
correspondence from the elliptic Hecke eigenvalues of
\(g=\Delta E_6\):
\[
 \lambda_1^{\mathrm{SK}}(p)\;=\;\lambda_g(p)+p^8+p^9,
 \qquad
 \lambda_2^{\mathrm{SK}}(p^2)\;=\;\lambda_g(p)^2+\lambda_g(p)\bigl(p^8+p^9\bigr)+p^{18},
\]
read off the \(X\)- and \(X^2\)-coefficients of the local Euler factor
associated with \eqref{eq:zbps-spinor-factor}, with \(\lambda_g(p)\) the normalized \(p\)-th Hecke eigenvalue of
\(g=\Delta E_6\), the unique normalized cusp form in
\(S_{18}(\SL_2(\Z))\): from \(\Delta E_6 = q-528q^2-4284q^3+\cdots\),
\(\lambda_g(2)=-528\), \(\lambda_g(3)=-4284\) \cite[Table~1]{ZagierSK}.
Concretely \(\lambda_1^{\mathrm{SK}}(2)=-528+2^8+2^9=240\) and
\(\lambda_1^{\mathrm{SK}}(3)=-4284+3^8+3^9=21960\), while
\(\lambda_2^{\mathrm{SK}}(4)=135424\) and
\(\lambda_2^{\mathrm{SK}}(9)=293343849\); the
Ramanujan bound is violated (\(\lvert\lambda_1^{\mathrm{SK}}(p)\rvert
\sim p^9\), whereas tempered classical eigenvalues at weight \(10\)
have size \(p^{17/2}\)), the signature of the
Saito--Kurokawa CAP form.
\end{proposition}

\begin{proof}
The Saito--Kurokawa correspondence \cite[\S 1]{ZagierSK} produces, for
each weight \(2k\) elliptic newform \(g\in S_{2k}(\SL_2(\Z))\), a
weight-\(k+1\) Siegel cusp form \(F_g\in S_{k+1}(\Sp_4(\Z))\) with
spinor \(L\)-function
\(L_{\mathrm{spin}}(s,F_g)=\zeta(s-k+1)\zeta(s-k)L(s,g)\) (which is
exactly \eqref{eq:zbps-spinor-factor} at \(k=9\)).  The Euler product
factorisation determines the Hecke eigenvalues at each prime.  In
Andrianov's polynomial the Siegel weight is \(K=k+1\), so
\(p^{2K-4}=p^{2k-2}\), \(p^{2K-3}=p^{2k-1}\), and
\(p^{4K-6}=p^{4k-2}\).  Writing
\(L_{\mathrm{spin}}(s,F_g)=\prod_p L_p(p^{-s},F_g)^{-1}\) with
\[
 L_p(X,F_g)
 \;=\;(1-p^{k-1}X)(1-p^kX)(1-\lambda_g(p)X+p^{2k-1}X^2)
\]
and comparing with the Andrianov factorisation
\cite[\S 4.5]{Andr74}
\[
L_p(X,F_g)=
1-\lambda_1^{\mathrm{SK}}(p)X+
(\lambda_1^{\mathrm{SK}}(p)^2-\lambda_2^{\mathrm{SK}}(p^2)-p^{2k-2})X^2-
\lambda_1^{\mathrm{SK}}(p)p^{2k-1}X^3+p^{4k-2}X^4,
\]
the \(X\)-coefficient gives
\(\lambda_1^{\mathrm{SK}}(p)=\lambda_g(p)+p^{k-1}+p^k\).  The
\(X^2\)-coefficient of the Euler product is
\[
2p^{2k-1}+(p^{k-1}+p^k)\lambda_g(p).
\]
Since the Andrianov coefficient of \(X^2\) is
\(\lambda_1^{\mathrm{SK}}(p)^2-\lambda_2^{\mathrm{SK}}(p^2)-p^{2k-2}\),
one obtains
\[
\begin{aligned}
\lambda_2^{\mathrm{SK}}(p^2)
&=(\lambda_g(p)+p^{k-1}+p^k)^2-p^{2k-2}
  -2p^{2k-1}-(p^{k-1}+p^k)\lambda_g(p) \\
&=\lambda_g(p)^2+\lambda_g(p)(p^{k-1}+p^k)+p^{2k}.
\end{aligned}
\]
At \(k=9\) this is the displayed formula.  The dimension formula
\(\dim S_{18}(\SL_2(\Z))=1\) forces \(g=\Delta E_6\) at \(k=9\), and
the eigenvalues at \(p=2,3\) are tabulated in
\cite[Table~1]{ZagierSK}.
\end{proof}

\begin{remark}[The SK-correspondence onto the full-level Maa{\ss} line]
\label{rem:zbps-sk-hecke-level0}
The SK-correspondence \(g\mapsto F_g\) is a Hecke-equivariant
injection \(S_{2k}(\SL_2(\Z))\hookrightarrow S_{k+1}(\Sp_4(\Z))\)
whose image is the Maa{\ss} subspace
\(\mathrm{Maass}_{k+1}\subset S_{k+1}(\Sp_4(\Z))\).  At \(k=9\),
\(\dim\mathrm{Maass}_{10}=\dim S_{18}(\SL_2(\Z))=1\) and
\(\mathrm{Maass}_{10}=\C\Delta_5^2\); the Hecke eigenvalues
\(\lambda_i^{\mathrm{SK}}\) of
Proposition~\ref{prop:zbps-sk-hecke} record the explicit
arithmetic translation of the full-level elliptic Hecke action on
\(g=\Delta E_6\) under the SK-correspondence.  No new modular
forms or arithmetic data enter; the Saito--Kurokawa eigenline of
\(\Delta_5^2\) is fully determined by the full-level arithmetic of
\(\Delta E_6\) and the SK-correspondence.
\end{remark}
The OP
monic normalization is
\begin{equation}
\label{eq:zbps-op-monic}
\chi_{10}^{\mathrm{OP}}
\;=\;
2^{-12}\,\Delta_{10}^{\theta}
\;=\;
4096^{-1}\,\Delta_5^2,
\end{equation}
and the level-\(n\) shadow~\eqref{eq:zbps-equality} is the inverse of
the \emph{theta-normalized} Hecke eigenline element
\(\Delta_{10}^{\theta}=\Delta_5^2\).

\begin{theorem}[Andrianov factorization of the spinor \(L\)-function]
\label{thm:zbps-andrianov-factor}
\ClaimStatusProvedElsewhere
For the Saito--Kurokawa representation generated by
\(\Delta_{10}^{\theta}=\Delta_5^2\),
\begin{equation}
\label{eq:zbps-spinor-factor}
L_{\mathrm{spin}}(s,\Delta_{10}^{\theta})
\;=\;
\zeta(s-8)\,\zeta(s-9)\,L(s,g),
\qquad g=\Delta E_6.
\end{equation}
Equivalently, for Andrianov's local polynomial
\[
Q_{p,F}(t)=
1-\lambda_F(p)t+
\bigl(\lambda_F(p)^2-\lambda_F(p^2)-p^{16}\bigr)t^2
-\lambda_F(p)p^{17}t^3+p^{34}t^4
\]
at weight \(10\), the Euler product
\(\prod_p Q_{p,F}(p^{-s})^{-1}\) is the left-hand side of
\eqref{eq:zbps-spinor-factor}.  Andrianov's completed function is
\[
\Psi_F(s)
=(2\pi)^{-s}\Gamma(s)\Gamma(s-8)
L_{\mathrm{spin}}(s,F).
\]
\end{theorem}

This is Andrianov's spin \(L\)-factor factorization
\cite[Theorem~3.1.1]{Andr74}; the polynomial \(Q_{p,F}\) is
Andrianov's genus-two Hecke polynomial in weight \(10\), with
coefficient
\(\lambda_F(p)^2-\lambda_F(p^2)-p^{2\cdot 10-4}\) at \(t^2\).
The proof uses Andrianov's spinor \(L\)-factor formula for the
degree-\(2\) Saito--Kurokawa lift attached to the weight-\(18\)
elliptic cusp form \(g\)
\cite{Andr74}, together with the dimension count
\(\dim S_{18}(\SL_2(\Z))=1\) which forces the source form to be
\(g=\Delta E_6\) \cite{ZagierSK}.

\begin{proposition}[Polar coefficient at \(s=10\)]
\label{prop:zbps-spin-pole-residue}
\ClaimStatusProvedElsewhere
With \(\Lambda(g,s)=(2\pi)^{-s}\Gamma(s)L(s,g)\), the functional
equation \(\Lambda(g,s)=-\Lambda(g,18-s)\) gives \(L(9,g)=0\), so the
elliptic central zero cancels the possible \(\zeta(s-8)\)-pole in the
spin product at \(s=9\).  At \(s=10\) the polar coefficient is
\begin{equation}
\label{eq:zbps-residue}
\operatorname*{Res}_{s=10}L_{\mathrm{spin}}(s,\Delta_{10}^{\theta})
\;=\;
\zeta(2)\,L(10,g)
\;=\;
\frac{\pi^2}{6}\,L(10,g),
\end{equation}
with
\[
L(10,g)/(2\pi i)^{10}\Omega^+(g)\in\Q
\]
by Manin--Deligne periods.  The period theorem fixes the
normalisation of the critical value; it is not a non-vanishing theorem.
Thus \(L_{\mathrm{spin}}(s,\Delta_{10}^{\theta})\) has a simple pole at
\(s=10\) exactly when \(L(10,g)\ne0\).
\end{proposition}

This is the Saito--Kurokawa critical-value identity \cite{ZagierSK} together with
the Manin--Deligne period decomposition \cite{Manin72,Del79}; the proof uses the
Eichler--Shimura--Manin period theorem for the rational newform \(g\)
\cite{Manin72,KZ}, in Deligne's period normalization \cite{Del79}.

\begin{remark}[Hecke eigenline is normalization-independent]
The spinor \(L\)-function depends on the Hecke eigenline
\(\C\Delta_{10}^{\theta}=\C\chi_{10}^{\mathrm{OP}}\), not on the scalar
used to label its generator.  Theorem~\ref{thm:zbps-andrianov-factor}
and Proposition~\ref{prop:zbps-spin-pole-residue} are statements about
the eigenline.  The OP scalar constant \(2^{-12}\) of
\eqref{eq:zbps-op-monic} relabels the generator from
\(\Delta_5^2\) to \(\chi_{10}^{\mathrm{OP}}\); it does not change
\(L_{\mathrm{spin}}\), the residue, or the cancellation of the
elliptic central zero against the zeta pole at \(s=9\).
\end{remark}

\subsection{Independence from the Dirac--Igusa construction}

The Saito--Kurokawa factorization of \(L_{\mathrm{spin}}(s,\Delta_5^2)\)
and the OP/DT scalar identity~\eqref{eq:zbps-op-dt} are imported
results: the first from Andrianov \cite{Andr74}, the second from
Oberdieck--Pixton \cite{OB} and Oberdieck \cite{OberdieckReducedPT}.
Neither uses the protected Pfaffian \(\mathrm{Pf}_{\mathrm{prot}}\) of
Chapter~5, the orientation character \(\epsilon_o\) of (O1)/(O1\(^+\)),
the type-II wall normal form of (O2), or the primitive recognition
\(P_X\cong\mathfrak g_{\Delta_5}\) problem of Chapter~6.  Conversely, the
construction of \(\mathfrak D_X^{\mathrm{DI}}\) and the
Pfaffian--Dirac theorem \(\mathrm{Pf}_{\mathrm{prot}}(\mathfrak
D_X^{\mathrm{DI}})=\Delta_5\) do not use the spinor \(L\)-factor or
the scalar inverse Igusa form.  The two routes meet at the squared
section \(\Delta_5^2=\Phi_{10}^{\mathrm{un}}\) and at no earlier
point.

\section{What the scalar shadow does not supply}
\label{sec:scalar-trace-not-gravity}

\subsection{\texorpdfstring{$Z_{\mathrm{BPS}}$ is the level-$n$ protected trace, not a path integral}{Z\_BPS is the level-n protected trace, not a path integral}}

The squared section \(\Delta_5^2=\Phi_{10}^{\mathrm{un}}\) is the
determinant section, and its inverse
\(Z_{\mathrm{BPS}}^{K3\times E}=\Delta_5^{-2}\) is the protected
scalar trace in the Beilinson chain, after the
primitive object at level \(\mathsf{P}\) (the open factorization
dg-category, the boundary vacuum \(b\), and \(A_b=\mathrm{End}_{\mathcal
C}(b)\)) and after the chart object at level \(\mathsf{C}\)
(the disk algebra and its bar / cobar / chiral Hochschild data).  The
scalar partition function is the \emph{trace} of the protected
data; it does not by itself construct the level-\(\mathsf A\) acting
algebra.

At finite height the protected trace is a separate row-level datum:
\[
\mathrm{trace\_categories},\quad
\mathrm{protected\_trace\_functors},\quad
\mathrm{trace\_operators},\quad
\mathrm{scalar\_normalizations},\quad
\mathrm{trace\_identities},\quad
\mathrm{forgetful\_maps},\quad
\mathrm{gravity\_residuals},\quad
\mathrm{transitions},\quad
\mathrm{scalar\_firewall}.
\]
These rows record the level-\(\mathsf Z\) trace category, the closed
\(K3\times E\to\mathrm{pt}\) trace functor, the operator whose
determinant square is \(\Phi_{10}^{\mathrm{un}}\), the constants
\(64\), \(4096\), and the OP sign \(-1\), the identity
\(\Delta_5^{-2}=(\Phi_{10}^{\mathrm{un}})^{-1}\), and the forgetful
maps killing orientation, primitive bracket, Hopf pairing, PBW, and
parity data.  The gravity-residual rows must mark the Hall--Borcherds
level-\(\mathsf A\) operator algebra as conjectural, not constructed,
and unused in the trace proof.  The scalar firewall excludes
orientation characters, Pfaffian lines, primitive brackets, Hopf
pairings, PBW bases, parity splittings, source Koszul maps, and
gravitational path integrals as data recovered from the scalar trace.

\begin{scope}[No scalar trace is a 3D gravitational path integral]
\label{scope:zbps-no-3d-gravity}
\ClaimStatusOpen
The identity~\eqref{eq:zbps-op-dt}, proved through OP/DT/PT primary
literature, is a closed automorphic equation between scalar power
series.  Neither this scalar identity, nor the automorphic section of
Chapter~3, nor the protected Pfaffian of Chapter~5 identifies
\(Z_{\mathrm{BPS}}^{K3\times E}\) to a \(3\)-dimensional
gravitational path integral on \(\mathrm{AdS}_3\times S^3\times K3\),
to a metric Euclidean path-integral over the full conformal class of
the boundary, or to a holographic dual partition function in any sense
beyond the boundary-CFT reading recorded by Vol~II.  The equality
\(\Phi_{10}^{\mathrm{un}}=\Delta_5^2\)
gives the determinant section of a candidate gravity-line theory; the
partition character would be the reciprocal section
\(\Delta_5^{-2}\).  The construction of the gravity-line operator
algebra is open.
\end{scope}

\begin{remark}[Protected trace versus path integral]
The assertion that \(Z_{\mathrm{BPS}}\) is a \(3d\) gravitational
path integral identifies a protected scalar trace with a
level-\(\mathsf A\) acting theory.  The precise statement is that
\(Z_{\mathrm{BPS}}^{K3\times E}\) is the level-\(n\) protected trace of
the retained \(K3\times E\) data; the
gravitational path integral, if such a thing exists at this boundary,
requires the Hall--Borcherds residual of Vol~II
\textup{(\S\ref{sec:hall-borcherds-residual})}.  The residual is open;
the path-integral realization is not claimed.  In particular, no
microscopic state space \(\mathcal H_{K3\times E}\) is asserted to
support \(\dim\mathcal H_{K3\times E,\gamma}=f(\gamma)\) outside the
charge-graded virtual reading recorded in Chapter~2.
\end{remark}

\subsection{The Hall--Borcherds residual}
\label{sec:hall-borcherds-residual}

The passage from a scalar partition function to a chain-level
gravity-line operator algebra requires a comparison datum that lives
neither in the automorphic / OP reduced theory nor in the compact
Pfaffian datum.  Vol~II names this comparison.

\begin{conjecture}[Hall--Borcherds residual; gravity-line comparison]
\label{conj:zbps-hall-borcherds-residual}
\ClaimStatusConjectured
Let \(X=K3\times E\) and assume the oriented critical Hall datum used
for the Vol~III positive-half comparison
\begin{equation}
\label{eq:zbps-hall-to-borch}
\Theta_{\mathrm{Hall}\to\mathrm{Borch}}\colon
\widehat{\mathrm{CoHA}}^{\mathrm{red}}_{\Lambda^{2,1}_{II}}(X)
\longrightarrow
\widehat{U^{\mathrm{ch}}\bigl(\mathfrak n_+(\mathfrak g_{\Delta_5})\bigr)}.
\end{equation}
The Hall--Borcherds residual is the conjectural extension of
\eqref{eq:zbps-hall-to-borch}, after Drinfeld doubling, current
enveloping along \(E\), and weight-completion, to a filtered
\(\mathsf{SC}^{\mathrm{ch,top}}\) morphism from the \(K3\times E\) Hall
boundary object to the gravitational boundary line algebra, whose
derived-centre trace has scalar character
\((\Phi_{10}^{\mathrm{un}})^{-1}=\Delta_5^{-2}\).
\end{conjecture}

The residual is the Vol~II gravity-line Hall--Borcherds residual
\cite{ChiralBarCobarVolII}: it asks for the gravity-line operator
algebra acting on the boundary.  It is not claimed here; its existence
and the construction of its witnessing morphism are open.

\begin{remark}[Compatibility with the \(\kappa\)-stratum]
\label{rem:zbps-kappa-stratum}
The \(\kappa\)-tuple of the K3\(\times\)E witness consists of the structural data
behind the scalar shadow:
\begin{equation}
\label{eq:zbps-kappa-tuple}
\bigl(\kappa_{\mathrm{cat}},\,\kappa_{\mathrm{ch}}^{\mathrm{Hodge}},\,
\kappa_{\mathrm{ch}}^{\mathrm{Heis}},\,
\kappa_{\mathrm{BKM}},\,\kappa_{\mathrm{fiber}}\bigr)_{K3\times E}
\;=\;
(0,\,0,\,3,\,5,\,24).
\end{equation}
The fourth entry \(\kappa_{\mathrm{BKM}}=5\) is the cusp-form weight
\(\mathrm{wt}(\Delta_5)\), the \(\Delta_5\) row;
the fifth entry \(\kappa_{\mathrm{fiber}}=\chi_{\mathrm{top}}(K3)=24\)
is the topological Euler characteristic of the K3 fibre, distinct from
\(\chi(\mathcal O_{K3})=2\).  The additive formula
\(\kappa_{\mathrm{BKM}}=\kappa_{\mathrm{ch}}^{\mathrm{Hodge}}+\chi(\mathcal
O_{\mathrm{fiber}})\) confuses these two: with
\(\kappa_{\mathrm{ch}}^{\mathrm{Hodge}}=0\) and
\(\chi(\mathcal O_{K3})=2\), it predicts \(\kappa_{\mathrm{BKM}}=2\),
which fails because the universal Borcherds weight at \(N=1\) is
\(\kappa_{\mathrm{BKM}}(\Phi_1)=c_1(0)/2=10/2=5\) (Borcherds 1995,
Gritsenko 1999).  The five entries of the tuple are five \emph{distinct}
constructions; they cannot be combined additively.  The Vol~III
counterexample at \(N=1\) is the controlling instance.
\end{remark}

\subsection{Local-to-global in mixed holomorphic--topological language}

A second route from local protected data to a global compact
factorization algebra
runs through the mixed holomorphic--topological model of Costello and
the holomorphic de Rham obstruction class.  Formal local Hamiltonian BF
data imply such a global object only after the global holomorphic de
Rham obstruction class
\([\mathrm{obs}_{HT,\mathrm{glob}}]\in H^1(X,\mathcal O_X)\),
plus descent, quantum master equation, anomaly, and locality data,
with the Hamiltonian condition supplied by the local model.  For
\(X=K3\times E\) the global obstruction group is
\(H^1(K3\times E,\mathcal O_{K3\times E})\cong\C\) (non-zero), so the
formal-local-Hamiltonian-BF route does not by itself identify
\(Z_{\mathrm{BPS}}^{K3\times E}=\Delta_5^{-2}\) to a global compact
factorization algebra.  Both bridges --- Hall--Borcherds for the Vol~II gravity
line, mixed-HT for the Costello local-to-global ascent --- are open;
the only established object here is the scalar shadow.

\section{The \(K3\times E\) threefold realization question}
\label{sec:k3-e-threefold-realization}

The square root in the title is meant in Dirac's sense.  The OP/DT
scalar branch \(Z^X_{\mathrm{OP/DT}}=-4096\,\Delta_5^{-2}\) is
orientation-forgetful: the sign character is lost under squaring, and the
inverse-square scalar cannot recover the sign by which a Pfaffian-line
section distinguishes \(\Delta_5\) from \(-\Delta_5\).  The missing
first-order datum is the protected Pfaffian, whose square is the
inverse of the protected scalar shadow.

\subsection{From the scalar trace to the retained operator datum}

The Dirac square-root pattern compares the chiral-shadow Pfaffian with
the scalar trace.  Before taking the trace one has the protected
Pfaffian of View~\textcircled{3}; after the trace one has the inverse
square as scalar shadow.  The Pfaffian carries the orientation data;
inverting its square forgets it and gives the orientation-forgetful
scalar invariant.

\begin{theorem}[Orientation-forgetting scalar consequence of the Pfaffian]
\label{thm:zbps-square-root-consequence}
\ClaimStatusConditional
On \(K3\times E\), after the retained D0-HN datum, retained
quotient-orientation datum, chosen Weyl lifts, and retained
\((O2_{\mathrm{atlas}})\) wall datum, together with the retained finite
geometric Pfaffian datum \((P_{\mathrm{fin}})\), of
Appendix~\ref{app:obstruction-discharges}, the
protected Pfaffian \(\mathrm{Pf}_{\mathrm{prot}}(\mathfrak
D_X^{\mathrm{DI}})\) of Chapter~5 is the automorphic section
\(\mathcal D_X=\Delta_5\) of Chapter~3, and its inverse square is the
level-\(n\) protected scalar shadow:
\begin{equation}
\label{eq:zbps-square-root}
\bigl(\mathrm{Pf}_{\mathrm{prot}}(\mathfrak D_X^{\mathrm{DI}})\bigr)^{-2}
\;=\;
\Delta_5(Z)^{-2}
\;=\;
\bigl(\Phi_{10}^{\mathrm{un}}(Z)\bigr)^{-1}
\;=\;
Z_{\mathrm{BPS}}^{K3\times E}(T_{\mathrm{OP}},Q_{\mathrm{OP}},P_{\mathrm{OP}}),
\end{equation}
and the OP/DT scalar branch is
\(-4096\,Z_{\mathrm{BPS}}^{K3\times E}\).
\end{theorem}

\begin{proof}
The Pfaffian--Dirac theorem of Chapter~6, with
(D0), (O1), (O1\(^+\)), and \((O2_{\mathrm{atlas}})\) represented by
the retained data of Theorem~\ref{thm:G-four-obstruction-criterion}, and
with the retained finite geometric Pfaffian datum \((P_{\mathrm{fin}})\),
asserts
\(\mathrm{Pf}_{\mathrm{prot}}(\mathfrak D_X^{\mathrm{DI}})=\mathcal D_X=\Delta_5\).
Squaring,
\(\mathrm{Pf}_{\mathrm{prot}}(\mathfrak D_X^{\mathrm{DI}})^2
=\Delta_5^2=\Phi_{10}^{\mathrm{un}}\); inverting,
\(\mathrm{Pf}_{\mathrm{prot}}(\mathfrak D_X^{\mathrm{DI}})^{-2}
=(\Phi_{10}^{\mathrm{un}})^{-1}\).  By Theorem~\ref{thm:zbps-level-n-shadow}
the right-hand side equals the level-\(n\) protected scalar shadow.
The Oberdieck--Pixton/Pandharipande scalar identity
\cite[Theorem~1]{OB} says that the OP chamber representative is
\(-4096\) times this trace; the factor is the OP convention, the
squared theta-leading coefficient, and the sign branch of
\eqref{eq:zbps-op-dt}.
\end{proof}

\subsection{The orientation datum the inverse-square scalar cannot recover}

The square-root sign is invisible to the scalar.  An orientation-line
sign \(\varepsilon\in\{\pm 1\}\) acting by \(\Delta_5\mapsto\varepsilon\Delta_5\)
satisfies \(\varepsilon^2=1\); the scalar \(\Delta_5^{-2}\) is therefore
\(\varepsilon\)-invariant.  Selecting the sign requires a pre-trace
Pfaffian-orientation datum of View~\textcircled{3}: the protected orientation character
\(\epsilon_o\) of the type-II Pfaffian wall normal form (O2) and its
restriction to the type-II Weyl group \(W^{(2)}(\Lambda^{2,1}_{II})\),
matched against the automorphic reflection character \(\nu_{\Delta_5}\).
Equality
\begin{equation}
\label{eq:zbps-orientation-equality}
\epsilon_o\bigm|_{W^{(2)}(\Lambda^{2,1}_{II})}
\;=\;
\nu_{\Delta_5}\bigm|_{W^{(2)}(\Lambda^{2,1}_{II})}
\end{equation}
is the orientation-monodromy-character identity of
Chapter~\ref{ch:pfaffian-dirac}; Appendix~\ref{app:obstruction-discharges}
states the retained-data criteria for (D0), (O1), (O1\(^+\)), and (O2)
on \(K3\times E\).  No constant
appearing in \eqref{eq:zbps-op-dt} --- not the leading \(64\), not
the squared \(4096\), not the OP scalar minus sign --- supplies
\(\epsilon_o\).  The local sign is checked on a rank-one Pfaffian wall;
the global character requires the retained quotient orientation and
the chosen Weyl transport, not a normalization constant.

\subsection{The four obstructions and the realization question}

Let \(\mathfrak D_X^{\mathrm{DI}}\) denote the Dirac--Igusa pro-object
\begin{equation}
\label{eq:zbps-dirac-igusa-pro-object}
\mathfrak D_X^{\mathrm{DI}}
\;=\;
\lim_R\bigl(
\mathcal A^{E_3}_{X,R},\,F^{\mathrm{hyb}}_{X,R},\,\Gamma_{X,R},\,\Pi_X,\,
\Phi_R,\,o_R,\,H_R,\,P_R^\Pi,\,C_{X,R},\,\Theta_{\mathrm{Kos},R},\,
\mathcal L_{\mathrm{Pf},R},\,\mathrm{pf}_{X,R},\,\varepsilon_{o,R},\,
\mathrm{Rec}_R\bigr)
\end{equation}
of Chapter~5: the inverse limit, over finite root windows \(R\),
of fourteen pieces of data including the chosen finite
\(E_3\)-prefactorization algebra \(\mathcal A^{E_3}_{X,R}\), the hybrid
\(\operatorname{Ran}^{\mathrm{hyb}}(E)\) carrier
\(F^{\mathrm{hyb}}_{X,R}\), the cosection-twisted reduced d-critical
orientation \(o_R\), the Hall correspondence datum \(H_R\), the
Pfaffian line \(\mathcal L_{\mathrm{Pf},R}\), the Pfaffian section
\(\mathrm{pf}_{X,R}\), and the recognition map \(\mathrm{Rec}_R\).
The retained D0-HN datum, retained quotient-orientation datum, chosen
Weyl lifts, retained \((O2_{\mathrm{atlas}})\) datum, and retained finite
geometric Pfaffian datum \((P_{\mathrm{fin}})\) are the
Pfaffian and orientation data of Appendix~\ref{app:obstruction-discharges}
for
Theorem~\ref{thm:zbps-square-root-consequence}.

\begin{figure}[!ht]
\centering
\fbox{\parbox{0.94\textwidth}{
\textbf{Realization question (the K3\(\times\)E threefold) — retained scalar-trace level.}\;
The square-root identity~\eqref{eq:zbps-square-root} is a
theorem on \(K3\times E\) at the Pfaffian and scalar-trace level relative to the
retained D0-HN datum, retained quotient-orientation datum, chosen Weyl lifts,
retained \((O2_{\mathrm{atlas}})\) wall datum, and retained finite
geometric Pfaffian datum \((P_{\mathrm{fin}})\) in
Appendix~\ref{app:obstruction-discharges}:

\smallskip
\noindent
\((D0)\) \emph{D0-degeneration limit of the cosection-reduced
d-critical orientation theory.}
\hfill
\textbf{Criterion}: Theorem~\ref{thm:G-D0}, relative to the retained
D0-HN datum of Definition~\ref{def:G-D0-HN-system}: flat
D0-degeneration family, additive K3 semiregularity cosections,
vanishing-cycle and orientation specialization, Pfaffian-line
extension, protected-integration compatibility, and strict HN
transitions.

\smallskip
\noindent
\((O1)\) \emph{Strong reduced orientation on the
\(K3\times E\)-quotient self-Ext frame bundle.}
\hfill
\textbf{Relative to \((O1_{\mathrm{quot}})\)}:
Theorem~\ref{thm:G-O1}.

\smallskip
\noindent
\((O1^+)\) \emph{Weyl-equivariant transport of the orientation along
\(W^{(2)}(\Lambda^{2,1}_{II})\) reflections.}
\hfill
\textbf{Relative to the chosen Weyl lifts}:
Theorem~\ref{thm:G-O1-plus}.

\smallskip
\noindent
\((O2)\) \emph{Local Pfaffian wall normal form on type-II walls; the
\(2^{12}=4096\) sign sum on the half-Hilbert orbit, with
the \((R1)\)--\((R2)\) scalar separation.}
\hfill
\textbf{Relative to \((O2_{\mathrm{atlas}})\)}:
Theorems~\ref{thm:G-O2-orbit-transport},
\ref{thm:G-R1-R2}.

\smallskip
\noindent
\((P_{\mathrm{fin}})\) \emph{Finite geometric Pfaffian datum on a
cofinal HN system: skew perfect complexes, Pfaffian lines and sections,
type-II wall divisor orders, automorphic line comparisons, support and
parity rows, leading coefficient, and Pfaffian-line
Mittag--Leffler compatibility.}
\hfill
\textbf{Relative to \((P_{\mathrm{fin}})\)}:
Definition~\ref{def:G-finite-pfaffian-section}.

\smallskip
\noindent
The level-\(\mathsf Z\) Hall--Borcherds trace comparison is relative to
the retained \((Z_{\mathrm{HB}})\) datum and
Theorem~\ref{thm:G-vol2-discharge}: the retained trace datum places the
compact source in the level-\(\mathsf Z\) chiral-Hochschild domain, and
Vol~I Theorem~H \cite{ChiralBarCobarVolI} together with Vol~II
\cite[chiral-Hochschild Beilinson stratification]{ChiralBarCobarVolII}
clauses~(i) and~(ii) give the chain-level chiral-Hochschild
trace
\(\mathrm{Tr}^{\mathrm{bulk}}_n=\Delta_5^{-2}\) on
\(Z^{\mathrm{der}}_{\mathrm{ch}}(C_X)\).  The Z\(_{\mathrm{BPS}}\)
square-root identity is therefore a retained-datum theorem on
\(K3\times E\) at level \(\mathsf Z\), granted \((Z_{\mathrm{HB}})\)
and those two cited theorems; the
level-\(\mathsf A\) gravity-line operator algebra remains open in
Vol~II.  The finite packet
\(\texttt{certificates/trace/k3e\_protected\_trace}\) currently records
\((Z_{\mathrm{HB}})\) only by the obstruction ledger
\(\texttt{TRACE\_OBSTRUCTION\_LEDGER\_VERIFIED}\), with
\(\texttt{protected\_trace\_certification=false}\).}}
\end{figure}

\subsection{The threefold target, not a fourfold}

The K3\(\times\)E target is a complex threefold:
\(\dim_\C K3\times E=3\).  The integer \(n\) in the
``level-\(n\) shadow'' is
\[
n=\chi(\mathcal O_Y)
\]
for the one-dimensional subscheme \(Y\subset K3\times E\), whereas
\(\chi(\mathcal O_{K3\times E})=2\cdot 0=0\) is the target
holomorphic Euler characteristic.  The perfect
obstruction theories in Chapter~5 are threefold theories; their
virtual classes are reduced virtual classes on the moduli of objects
on the threefold; the Pfaffian \(\mathrm{Pf}_{\mathrm{prot}}\) of
Chapter~5 is the Pfaffian of the threefold's protected operator data.
The Borcherds product formula and the Gritsenko--Nikulin denominator
identity give the Mukai-vector dictionary on the threefold by
projection along the elliptic factor; the OP/DT/PT comparison gives
the scalar shadow on the threefold; the protected Pfaffian gives the
square root.  No fourfold target intervenes anywhere in the chain.

\subsection{Comparison with the determinant-only construction}

The virtual \(K_0\)-determinant identity
\(\mathcal D_X(Z)=\Delta_5(Z)\) is a Grothendieck-class
identity: it sees only the signed superdimensions of the underlying
virtual super vector spaces.  It does not see brackets, not
parities of representatives, not relation constants, not completed
PBW data, not Hopf-pairing radicals, not orientation gerbes, not the
Pfaffian line \(\mathcal L_{\mathrm{Pf}}\).  Squaring the determinant
gives \(\Delta_5^2\); inverting gives the scalar shadow.  All the
information about the retained operator datum \(\mathfrak D_X^{\mathrm{DI}}\)
that distinguishes the retained-data identity
\(\mathrm{Pf}_{\mathrm{prot}}(\mathfrak
D_X^{\mathrm{DI}})=\Delta_5\) from a non-Pfaffian square-root choice
\(\Delta_5\mapsto -\Delta_5\) lives in the Pfaffian datum and the
orientation character.  The inverse-square scalar shadow
\eqref{eq:zbps-equality}, whose OP representative is
\eqref{eq:zbps-op-dt}, is the maximal information visible after the
trace; the Pfaffian and the orientation are the stronger pre-trace data.

\section{Open structures}
\label{sec:zbps-open-problems}

\subsection{The Hall--Borcherds residual revisited}

Conjecture~\ref{conj:zbps-hall-borcherds-residual} asks for the
filtered \(\mathsf{SC}^{\mathrm{ch,top}}\)
morphism from the \(K3\times E\) Hall boundary object to the
gravitational boundary line algebra of Vol~II, whose derived-centre
trace has scalar character \(\Delta_5^{-2}\).  Its trace part is the
bridge from the level-\(\mathsf{S}\) scalar identity
\(Z_{\mathrm{BPS}}^{K3\times E}=\Delta_5^{-2}\) to a
level-\(\mathsf{Z}\) chiral-Hochschild trace datum.  Its acting part is
the further level-\(\mathsf A\) gravity-line operator algebra.

The residual consists of three completions.  The Drinfeld doubling
\(\widehat{\mathrm{CoHA}}^{\mathrm{red}}\to D(\widehat{\mathrm{CoHA}}^{\mathrm{red}})\)
of the Vol~III positive-half pairing: the pairing, completion,
integral form, and stable-envelope transport that convert the
Hall positive half \(\mathsf{Y}^+(K3\times E)\) into the full quantum
group \(G(K3\times E)\).  The current enveloping along \(E\)
and weight-completion: the chiral specialization
\(\mathrm{Sp}^{\mathrm{ch}}_{\Sigma_{d-1},C}\) along the reference
curve \(C=E\), together with the weight-completed ambient required to
host the formal generators of \(\mathfrak g_{\Delta_5}\).  The
filtered \(\mathsf{SC}^{\mathrm{ch,top}}\) morphism property: the
resulting boundary object must be a filtered morphism in the
two-coloured Swiss-cheese chiral-topological operad, not merely a
graded-vector-space comparison.  Their existence is the Vol~II
Hall--Borcherds residual conjecture.

\subsection{The gravity-line operator algebra and the Pentagon-face
determinant}

The Vol~II gravity-line problem asks for a chain-level operator algebra
\(\mathfrak{Op}_{\mathrm{grav}}^{K3\times E}\) whose Pentagon-face
determinant section is \(\Phi_{10}^{\mathrm{un}}=\Delta_5^2\).  If such
an algebra is constructed, its partition character is the reciprocal
section \((\Phi_{10}^{\mathrm{un}})^{-1}=\Delta_5^{-2}\).  The
Brown--Henneaux dictionary \(c=3\ell/(2G_N)\) fixes the Virasoro
boundary datum at \(\mathrm{Vir}_c\); it does not construct the
\(K3\times E\) gravity-line algebra and does not evaluate the
Pentagon-face determinant.  A K3-Heisenberg witness would have to
specialize the Virasoro \(\lambda\)-bracket
\(\{T_\lambda T\}=\partial T+2T\lambda+(c/12)\lambda^3\) to the closed
K3 boundary at \(\kappa_{\mathrm{BKM}}(\Delta_5)=5\) and identify the
resulting determinant with \(\Phi_{10}^{\mathrm{un}}=\Delta_5^2\).
Existence of the gravity-line operator algebra, agreement of its
reciprocal partition character with the unscaled trace
\eqref{eq:zbps-equality}, equivalently with the OP representative
\eqref{eq:zbps-op-dt} after OP normalization on each primitive closed
K3-class fibre, and factorisation through the protected Pfaffian
\(\mathrm{Pf}_{\mathrm{prot}}(\mathfrak D_X^{\mathrm{DI}})=\mathcal D_X=\Delta_5\)
of Chapter~5 are open level-\(\mathsf A\) conditions.

\subsection{Mathieu and umbral moonshine adjacency}

The K3 elliptic genus
\(Z_{\mathrm{K3}}=2\phi_{0,1}\) admits the
\(N=4\) superconformal-character decomposition of
Eguchi--Ooguri--Tachikawa \cite{EOT}: after the massless characters are
removed, the holomorphic mock-modular series
\[
\Sigma(\tau)=-2q^{-1/8}+\sum_{n\ge1}A_nq^{n-1/8}
\]
has coefficients \(A_n\) which are dimensions of, and after twining
characters of, \(M_{24}\)-modules \cite{GHV,GannonUmbral}.  The
decomposition is the infinite massive-sector decomposition of the
\(N=4\) algebra.  The coefficients \(f(n,l)\) of
\(\phi_{0,1}\) which drive the Borcherds product are the untwined
Jacobi coefficients; an \(M_{24}\)-twining would replace
\(\phi_{0,1}\) by the corresponding twined K3 elliptic genus before
the singular theta lift.  The umbral moonshine extension of
Cheng--Duncan--Harvey \cite{CDH} is indexed by the \(23\) Niemeier root
systems with non-empty root system, with \(M_{24}\) the
\(A_1^{24}\) case.  The interaction between this sporadic-group
structure on the K3 elliptic genus and the automorphic / BKM
denominator structure on the
\(K3\times E\) Igusa cusp form is the Mathieu / umbral moonshine
adjacency.

\subsection{The closed-cobordism partition reads as one}

The scalar shadow is the trace, the Pfaffian is the pre-trace
operator-level section, and the primitive input is the retained
\(E_3\)-prefactorization input whose Pfaffian line is presented by
\(\mathfrak D_X^{\mathrm{DI}}\).  The closed
\(K3\times E\to\mathrm{pt}\) cobordism sees only the orientation-forgetful
scalar trace.  At that level the inverse of the squared Pfaffian section is the Laurent expansion
\(\Delta_5^{-2}=(\Phi_{10}^{\mathrm{un}})^{-1}\) of the inverse Igusa
cusp form, and the orientation information that distinguishes
\(\Delta_5\) from \(-\Delta_5\) lives in the Pfaffian datum at level
\(\mathsf{Z}\).  The closed-cobordism
partition function reads as one because \(\Delta_5\) is the same
section throughout: the Borcherds--Gritsenko--Nikulin section of the
weight-\(5\) automorphic line, the Borcherds--Kac--Moody denominator
on the root lattice \(\Lambda^{2,1}_{II}\), and the protected Pfaffian
of the Dirac--Igusa pro-object on \(K3\times E\), relative to the
retained orientation and wall-atlas data.
The squared shadow of this same section is the protected partition
function on the closed cobordism \(K3\times E\to\mathrm{pt}\); the
Hall--Borcherds residual of Vol~II, when constructed, will lift it to
the gravity-line operator algebra at level \(\mathsf{A}\); the
threefold realization question of
\S\ref{sec:k3-e-threefold-realization} is the retained-data
construction of the Dirac--Igusa pro-object and the remaining
compact-source recognition problem.

\section*{Bibliographic notes}

The OP/DT scalar identity is Oberdieck--Pixton \cite{OB},
Oberdieck--Pandharipande \cite{OPand}, and Oberdieck
\cite{OberdieckReducedPT}; the variable conventions follow Bryan
\cite{BRYAN}.  The Oberdieck--Pixton/Pandharipande scalar identity
\cite[Theorem~1]{OB} is the consolidated form used here.  The
Andrianov factorization of the spinor \(L\)-function
is \cite{Andr74}; the Saito--Kurokawa correspondence at weight
\(10\) is \cite{ZagierSK}; the Manin--Deligne period theorem is
\cite{Manin72}, \cite{KZ}, \cite{Del79}.  The dyon-counting
interpretation in CHL backgrounds is Dijkgraaf--Verlinde--Verlinde
\cite{DVV} and Borcherds \cite{B}; these inputs are the
boundary-row physical hosts and are kept separate from the
\(K3\times E\) protected scalar shadow throughout.

The Hall--Borcherds residual conjecture in
\S\ref{sec:hall-borcherds-residual} is Vol~II's Construction
Problem~2 \cite{ChiralBarCobarVolII}, the gravity-line Hall--Borcherds
comparison.
The protected D-brane index conjecture
\(\Omega_X(h,d,n)\stackrel{\mathrm{phys}}{=}\mathbf{DT}^X_{n,(\beta_h,d)}\)
is recorded as the protected D-brane-index conjecture of
Chapter~5; it is not used in the OP/DT scalar proof.  The
Mukai dictionary, the K3-class chamber convention, and the
\(\Lambda^{2,1}_{II}\) lattice data are imported from Chapter~2.  The
protected Pfaffian \(\mathrm{Pf}_{\mathrm{prot}}\) and the Dirac--Igusa
pro-object \(\mathfrak D_X^{\mathrm{DI}}\) are specified in
Chapter~5; their Pfaffian and orientation parts are established
relative to the retained data of Appendix~\ref{app:obstruction-discharges};
the Pfaffian--Dirac theorem is Chapter~6; the four
obstruction classes (D0), (O1), (O1\(^+\)), (O2) are defined and
tracked in \S5.8 of Chapter~5.

The \(\kappa\)-tuple \((0,0,3,5,24)\) and its non-additivity at \(N=1\) (the
controlling instance for \(\Delta_5\)) are recorded in Vol~III; the
Hall positive half \(\mathsf{Y}^+(K3\times E)\) and its Drinfeld
doubling datum toward \(G(K3\times E)\), as well as the 6d holomorphic
Chern--Simons quartic obstruction
\(\int_X\mathrm{Tr}_{\mathrm{ad}}(A(F_A)^3)\), are also Vol~III.  The
local model \(\R^2_{\mathrm{top}}\times\C^2_{\mathrm{hol}}\) and the
holomorphic de Rham obstruction class are mixed-HT-strings.  None of
these supply a 3D gravitational path integral.

\chapter{Open obstruction classes}
\label{ch:open-obstructions}

\section{Open obstruction classes}
\label{sec:open-obstructions-and-subsequent-constructions}

The cross-level identity \(\textcircled{2}\!\leftrightarrow\!\textcircled{3}\)
has two logically distinct parts.  The Pfaffian and
orientation part,
\(\mathrm{Pf}_{\mathrm{prot}}(\mathfrak D_X^{\mathrm{DI}})=\mathcal D_X=\Delta_5\)
and \(\epsilon_o=\nu_{\Delta_5}\) on \(W^{(2)}(\La^{2,1}_{II})\),
holds on \(K3\times E\) relative to the retained D0-HN datum, retained
quotient-orientation datum, chosen Weyl lifts, and retained
\((O2_{\mathrm{atlas}})\) wall datum, together with the retained finite
geometric Pfaffian datum \((P_{\mathrm{fin}})\), of
Appendix~\ref{app:obstruction-discharges}.  The primitive recognition
\(P_X\cong\mathfrak g_{\Delta_5}\) is conditional on the finite
compact-source recognition datum \((S1)\)--\((S10)\).
The exact theorem that remains after deleting any one of these data is
Proposition~\ref{prop:ch6-hypothesis-deletion-residues}; no deletion is
filled by a scalar trace, automorphic section, or denominator product.
The implication and non-redundancy statement is
Proposition~\ref{prop:ch6-obstruction-independence}.
The finite Pfaffian packet
\(\texttt{certificates/pfaffian/k3e\_finite\_pfaffian}\) currently
records \((P_{\mathrm{fin}})\) only by its obstruction ledger
\(\texttt{blocked\_obligations.csv}\): skew perfect complexes,
Pfaffian lines, finite sections, wall charts, automorphic line
comparisons, transitions, and scalar firewalls remain missing, and the
verifier records \(\texttt{pfaffian\_certification=false}\).
The finite O2 wall-atlas packet
\(\texttt{certificates/wall\_atlas/k3e\_o2\_atlas}\) likewise records
\((O2_{\mathrm{atlas}})\) only by its obstruction ledger: simple-wall
coverage, wall objects, rank-one charts, overlaps, Weyl transports,
half-Hilbert-orbit rows, compatibility rows, type-I exclusion,
D0-limit extension, and scalar firewalls remain missing, and the
verifier records \(\texttt{o2\_certification=false}\).
The automorphic identity
\(\textcircled{1}\!\leftrightarrow\!\textcircled{2}\) is theorem
\cite{Borcherds95,GN}; the singular-theta-lift identity
\(\textcircled{2}\!\leftrightarrow\!\textcircled{5}\) is theorem
\cite{Borcherds95}; the BKM denominator computation
\(\operatorname{den}(\mathfrak g_{\Delta_5})=64^{-1}\Delta_5(2Z)\) is
theorem; the Lie-superalgebra recognition remains conditional
from the finite-stage compact \(K3\times E\) source to the closed
denominator algebra.

Each of the four obstructions sits at a definite categorical-dimensional
level on the universal chain
\(\mathsf{P}\to\mathsf{C}\to\mathsf{S}\to\mathsf{Z}\to\mathsf{A}\).
The remaining structures beyond the theorem are the Hall--Borcherds
residual comparing the level-\(n\) protected scalar shadow with a
chain-level trace datum, the chain-fusion conjecture for compact non-toric \(K3\times E\),
and the gravity-line operator algebra required to have Pentagon-face
determinant section \(\Delta_5^2\).  The Mathieu / umbral moonshine adjacency
stands between the compact source and these structures:
its connection to \(\mathfrak g_{\Delta_5}\) is conjectural.  The
conjecture asks for the twined denominator algebra and its reciprocal
partition character. The
finite protected-trace packet
\(\texttt{certificates/trace/k3e\_protected\_trace}\) presently records
the level-\(\mathsf Z\) trace datum only by its obstruction ledger:
trace categories, protected trace functors, operators, scalar
normalizations, identities, forgetful maps, gravity-residual separation,
transitions, and scalar firewalls remain missing, and the verifier
records \(\texttt{protected\_trace\_certification=false}\).
Beilinson cut separates the retained-data Pfaffian theorem from the
primitive source recognition and the gravity-line operator algebra.

The three names of \(\Delta_5\) are kept distinct on first occurrence:
the automorphic section
\(\mathcal D_X=\Delta_5\in H^0(\mathcal H_2,\mathcal L^5\otimes\nu_{\Delta_5})\)
of view \textcircled{1}; the BKM denominator
\(\operatorname{den}(\mathfrak g_{\Delta_5})=64^{-1}\Delta_5(2Z)\) of
view \textcircled{2}; and the compact protected Pfaffian
\(\mathrm{Pf}_{\mathrm{prot}}(\mathfrak D_X^{\mathrm{DI}})\) of view
\textcircled{3}, relative to the retained data of
Appendix~\ref{app:obstruction-discharges}.

\subsection{The four named obstructions}
\label{subsec:four-named-obstructions}

The four obstructions act on distinct mathematical objects before any
scalar trace is taken.  The compact target is the Dirac--Igusa pro-object
\(\mathfrak D_X^{\mathrm{DI}}=\varprojlim_R(\mathcal A_{X,R}^{E_3},
F_{X,R}^{\mathrm{hyb}},\Gamma_{X,R},\Pi_X,\Phi_R,o_R,H_R,P_R^\Pi,
C_{X,R},\Theta_{\mathrm{Kos},R},\mathcal L_{\mathrm{Pf},R},
\mathrm{pf}_{X,R},\epsilon_{o,R},\mathrm{Rec}_R)\). The level
discipline is \emph{primitive object first, cross-level identity second,
scalar shadow last}. The four obstructions and \((P_{\mathrm{fin}})\)
are pre-trace Pfaffian-orientation data of View~\textcircled{3};
the level-\(\mathsf Z\) object is the derived-centre trace pairing.
They do not identify the compact primitive Lie superalgebra; that is
the separate datum \((S1)\)--\((S10)\).

\subsubsection{(D0): D0-degeneration limit of the cosection-reduced d-critical orientation theory}
\label{subsubsec:obstruction-D0}

\paragraph{D0-degeneration.} The cosection-reduced d-critical orientation
theory on the retained Hall stack
\(\mathcal H_R^{\mathrm{or}}\) of finite HN height \(R\) extends through
the typed D0-HN datum of Definition~\ref{def:G-D0-HN-system}: a flat
D0-degeneration family, retained object/extension/flag stacks, proper
or closed HN transitions, additive K3 semiregularity cosections,
reduced vanishing-cycle specialization, Joyce--Upmeier orientation
specialization, Pfaffian-line and protected-integration compatibility,
and strict Mittag--Leffler exactness.  In particular, the orientation
line of the reduced determinant complex
\[
\mathscr D^{\mathrm{red}}_{R,\mathcal C}
=\det\operatorname{RHom}_{\mathrm{red}}(\mathcal C,\mathcal C)
\]
admits a Brav--Bussi--Dupont--Joyce--Szendr\H{o}i \cite{BBDJS2015} square-root
gerbe section over the closure
\(\overline{\mathfrak M^{\mathrm{st}}_{R,\mathcal C}}\subset
\overline{\mathcal H_R^{\mathrm{or}}}\) including the D0-degeneration
stratum, compatible with the cosection contraction inherited from
\(K3\times E\). The retained Hall datum supplies an oriented critical
finite-stage Pandharipande--Thomas/Gholampour--Sheshmani
inverse system \(\widehat{\mathcal H}^{\mathrm{or}}_{\sigma,S}
=\varprojlim_R\mathsf F_R\mathcal H_{\sigma,S}^{\mathrm{or}}\) on the
finite HN quotients of
Proposition~\ref{prop:sectorial-hall-truncation};
\textup{(D0)} asks that this system carry a coherent orientation through
the D0-degeneration limit.
The finite packet
\(\texttt{certificates/d0/k3e\_d0\_hn}\) records the missing rows in
\(\texttt{blocked\_obligations.csv}\): flat degeneration families,
retained substacks, HN transitions, derived and d-critical charts,
semiregularity cosections, vanishing-cycle and orientation
specializations, Pfaffian/protected-integration specializations,
Mittag--Leffler rows, and scalar firewalls.  Its verifier returns
\(\texttt{D0\_OBSTRUCTION\_LEDGER\_VERIFIED}\), with
\(\texttt{d0\_certification=false}\).  This is the finite D0-HN
obstruction list, not a D0-HN construction.

\paragraph{Criterion.} The retained D0-HN datum of
Definition~\ref{def:G-D0-HN-system} supplies the D0-degeneration
orientation through the retained inverse system via
Theorem~\ref{thm:G-D0}, and
Theorems~\ref{thm:G-O1}, \ref{thm:G-O1-plus}, and
\ref{thm:G-O2-orbit-transport} transport it to the type-II Weyl
character.  The Maass character \(\nu_{\Delta_5}\) and the geometric
orientation \(\epsilon_o\) then agree on
\(W^{(2)}(\La^{2,1}_{II})\).  BBDJS \cite{BBDJS2015}
vanishing-cycle data require an oriented
d-critical locus, not boundedness alone, and Joyce--Upmeier
\cite{JoyceUpmeier} square-root data on CY\(_3\) moduli do not
automatically extend to the reduced \(K3\times E\) quotient: the
cosection-induced shift in the determinant complex must be shown
compatible with descent.  The connected \(BE\) Borel/edge class
\(\alpha^{E,\mathrm{free}}=a_1u_1+a_2u_2\in H^2(BE;\mathbb F_2)\)
is separated from finite \(BE[N]\) Borel terms in the Borel
spectral sequence; the quotient-orientation descent across
object, extension, mixed, wrapped, two-step flag, overlap,
Thom--Sebastiani, and \(R'\to R\) transition strata is conditional,
with overlap/correspondence/flag/transition compatibilities named
separately.

\paragraph{Argument.} The retained Hall boundedness datum
\textup{[K3\(\times\)E-fin]}, motivated by
\cite{LiuProductStability,PiyaratneTodaModuli3folds}, supplies a
finite-type retained slice; the cited boundedness results do not by
themselves prove this slice for the compact \(K3\times E\) Hall sector.
The datum also does not give quasi-smoothness, a d-critical chart,
BBDJS coefficients, an orientation line, or compact-support
six-functor admissibility through the D0 limit.
The finite rows required for this assertion are
\(\mathrm{hn\_type\_bounds}\), \(\mathrm{semistable\_substacks}\),
\(\mathrm{derived\_enhancements}\), \(\mathrm{universal\_complexes}\),
\(\mathrm{scalar\_rigidifications}\), \(\mathrm{stratifications}\),
\(\mathrm{extension\_flag\_stacks}\), \(\mathrm{cosection\_atlas}\),
\(\mathrm{transitions}\), and \(\mathrm{scalar\_firewall}\).
The finite-moduli packet
\(\texttt{certificates/moduli/k3e\_finite\_moduli}\) presently records
the HN-bound, semistable-substack, derived-enhancement, and
rigidification-inertia rows, together with the finite class-bound
and universal-complex rows and the \(E\)-translation-rigidification
rows, retained closed-substack rows, and retained extension-closure
rows, retained HN-factor-closure rows, and retained dual-closure rows,
while its obstruction ledger keeps the flag stacks, cosection charts,
transitions, and scalar firewalls open; the verifier records
\(\texttt{moduli\_certification=false}\).
Pantev--To{\"e}n--Vaqui{\'e}--Vezzosi \cite{PTVV2013} shifted symplectic
structures cover ambient derived
moduli under Calabi--Yau hypotheses; the retained atlas is not produced
by PTVV alone. Joyce--Upmeier \cite{JoyceUpmeier} square-root orientation
on Calabi--Yau\(_3\) moduli requires exact-sequence/direct-sum
compatibility; the reduced \(K3\times E\)-quotient inherits this only
after the cosection-induced shift is shown compatible with descent.

\paragraph{Quotient-orientation classes.} The retained
quotient-orientation datum consists of
\(\alpha^{\mathrm{red}}_{R,\mathcal C}=w_2(\mathscr D^{\mathrm{red}}_{R,\mathcal C})
\in H^2(S;\mathbb F_2)\), \(\alpha^{E,\mathrm{free}}_{R,\mathcal C}\),
\(\widetilde\beta^{H}_{R,\mathcal C}\in H^2_H(S;\mathbb F_2)\), and
\(\lambda^{H}_{R,\mathcal C}\in H^1(BH;\mathbb F_2)\) on every retained
object, extension, mixed, wrapped, and flag stratum \(\mathcal C\), with
killing of the Borel mixed terms and edge-spectral-sequence
differentials at every node of the descent diagram.  On the rank-one
D6--D2--D0 stratum \(\mathcal C=(1,\beta,n)\), with \(\beta\) the
K3-fibre primitive class, the BBDJS line carries the cosection
contraction explicitly, and the four obstruction classes
\((\alpha^{\mathrm{red}},\alpha^{E,\mathrm{free}},\widetilde\beta^{H},
\lambda^{H})\) are the retained quotient-orientation datum of
Appendix~\ref{app:G-O1-discharge}.  The remaining source-side problems
are this finite-stabilizer quotient orientation and the finite
compact-source recognition datum for the Hall primitive algebra.
The finite orientation packet
\(\texttt{certificates/orientation/k3e\_reduced\_orientation}\) records
this as an obstruction ledger, not as an orientation proof.  Its
\(\texttt{blocked\_obligations.csv}\) table lists the missing
square-root, Thom--Sebastiani, quotient Borel, finite-stabilizer,
Weyl-lift, Coxeter, transition, and scalar-firewall rows, and the
verifier reports \(\texttt{ORIENTATION\_OBSTRUCTION\_LEDGER\_VERIFIED}\)
with \(\texttt{orientation\_certification=false}\).

\subsubsection{(O1): strong reduced orientation on the K3{\(\times\)}E-quotient self-Ext frame bundle}
\label{subsubsec:obstruction-O1}

\paragraph{Reduced quotient orientation.} The BBDJS \cite{BBDJS2015} orientation gerbe on the moduli
of stable sheaves on \(K3\times E\) descends to the
\(K3\times E\)-quotient self-Ext frame bundle, with reduced orientation
gauge: on the stable stratum
\(\mathfrak M^{\mathrm{st}}_{R,\mathcal C}\subset\mathcal H_R^{\mathrm{or}}\)
the obstruction tuple
\[
\mathfrak o^{\mathrm{or}}_{R,\mathcal C}
=\bigl(\alpha^{\mathrm{red}}_{R,\mathcal C},\,
\alpha^{E,\mathrm{free}}_{R,\mathcal C},\,
\{(\widetilde\beta^{H}_{R,\mathcal C},\lambda^{H}_{R,\mathcal C})\}_{H},\,
\mu^{\mathrm{or}}_{R,e},\,\dots\bigr)
\]
vanishes simultaneously, with all overlap, correspondence, flag, and
\(R'\to R\) transition data trivialized coherently. \textup{(O1)}
demands the strong reduced descent: the connected-\(BE\) class
\(\alpha^{E,\mathrm{free}}\), the finite-\(BE[N]\) gerbe classes
\(\widetilde\beta^H\), and the residual finite linearization characters
\(\lambda^H\) all vanish, separately, on every retained stratum, and the
spectral-sequence differentials and Borel mixed terms are killed.

\paragraph{Criterion.} The Hall well-formedness criterion is prior to
recognition: by
Proposition~\ref{prop:hall-product-only-no-primitives}, an associative
Hall product without a counital coproduct, bialgebra compatibility, Hopf
pairing, radical, quotient, and strict transition data has no defined
primitive source object.  The primitive recognition criterion then reads:
target arithmetic on \(W_{\le3}\) is target
Borcherds--Gritsenko--Nikulin--Kac data only, and recognition
requires an actual compact source datum supplying
\(M,D,B,G,K,Q,A_\beta\), relation rows, Hopf radical/coideal kernel
equality, PBW associated-graded comparison, and strict
transition/Mittag--Leffler rows.  The Koszul comparison adds the finite
comparison residual
\(\mathfrak O_{\mathrm{Kcmp},R}\): by
Proposition~\ref{prop:primitive-restriction-not-koszul}, agreement of
the source-to-target map on constant-current primitives does not imply
the quasi-isomorphism
\(\Omega_E^{\mathrm{ch}}C_X\simeq\mathsf{Den}_{\Delta_5,E}\), and by
Definition~\ref{def:finite-koszul-comparison-datum} even that
quasi-isomorphism must still respect Weyl transport, Pfaffian
orientation, Hall pairing, radical quotient, PBW filtration, and HN
transitions.  The pro-object comparison also requires the finite
morphism rows
\(\mathrm{cofinal\_subsystems}\), \(\mathrm{stage\_objects}\),
\(\mathrm{stage\_morphisms}\), \(\mathrm{component\_maps}\),
\(\mathrm{composition\_laws}\), \(\mathrm{ml\_exactness}\),
\(\mathrm{pro\_isomorphisms}\), and \(\mathrm{scalar\_firewall}\);
without them a scalar Pfaffian product or denominator identity does not
define a Dirac--Igusa pro-object.
Theorem~\ref{thm:G-O1} of
Appendix~\ref{app:obstruction-discharges} states what the retained
\textup{(O1)} datum carries
to the recognition theorem. The connected \(E\) and finite \(E[N]\)
cohomologies are \emph{distinct} obstruction spaces:
\(H^*(BE;\mathbb F_2)=\mathbb F_2[u_1,u_2]\) with
\(|u_i|=2\), so \(H^1(BE;\mathbb F_2)=0\); for
\(E[2]\cong(\mathbb Z/2)^2\),
\(H^*(BE[2];\mathbb F_2)=\mathbb F_2[x_1,x_2]\) with
\(|x_i|=1\), and the residual linearization
\(\lambda=\lambda_1x_1+\lambda_2x_2\) is independent data after the
degree-two gerbe class is killed.

\paragraph{Argument.} Translation invariance under \(E\)-action does
not produce orientation descent: a fixed underlying line may carry a
nontrivial finite character, and ordinary translation invariance is
\emph{not} a substitute for vanishing of
\(\lambda^H\) on full two-primary generators of
\(H^*(B(\mathbb Z/2^a)^2;\mathbb F_2)
=\Lambda(x_1,x_2)\otimes\mathbb F_2[y_1,y_2]\) at \(N=2^a\). The
quotient-orientation criterion is finite \v{C}ech--Borel descent on the
retained groupoid, not a vanishing theorem.

\paragraph{Quotient-orientation theorem.} Theorem~\ref{thm:G-O1} is the
quotient-orientation criterion relative to the retained
\((O1_{\mathrm{quot}})\) datum.  That datum contains the BBDJS gerbe
sections on the self-Ext frame bundle, their null classes on the
rank-one D6--D2--D0 stratum and its extension, mixed, wrapped, and
two-step flag completions, and the overlap and correspondence
transitions on the cofinal HN system.  The data
\(M,D,B,G,K,Q,A_\beta\), the relation rows, the Hopf radical equality,
and the PBW comparison required by primitive recognition are still
separate compact-source data.

\subsubsection{(O1)\(^+\): Weyl-equivariant transport of the orientation along W\((^2)\)(\(\La_{II}^{2,1}\)) reflections}
\label{subsubsec:obstruction-O1plus}

\paragraph{Weyl transport.} The orientation character on the type-II Weyl
group \(W^{(2)}(\La^{2,1}_{II})\) is Weyl-equivariant: the orientation
line of \textup{(O1)} is equipped with coherent lifts
\(\tau_i:o_{\mathcal C}\to s_{\delta_i}^{*}o_{\mathcal C}\) of the
three type-II simple reflections, each acting with sign \(-1\), and
the projective determinant-line cocycle
\(c_o\in Z^2(W^{(2)}(\La^{2,1}_{II});\mathbb F_2)\) of the lift is a
coboundary, with quotient torsor defects
\(\omega_{i,\mathcal C}=0\) on every retained groupoid component. At
finite height \(R\), the partial action groupoid \(\mathcal G_R\)
carries
\(c_{o,R}\in Z^2(\mathcal G_R;\underline{\mathbb F}_2)\), and
\textup{(O1)\(^+\)} asks that
\([c_{o,R}]=0\), \(\tau_i^2=1\), and \(\omega_{i,\mathcal C}=0\) hold
coherently on every stratum and at every transition height.

\paragraph{Criterion.} The local sign-extraction criterion reads:
only after the unit character is killed and
\(N_\delta^{\mathrm{Pf}}=1\) is computed from the reduced normal
self-Ext chart does the local sign equal \(-1\), and Maass values
\(\nu_{\Delta_5}(s_{\delta_i})=-1\) plus divisor orders
\(d^{\mathrm{aut}}_{\delta_i}=f(0,\pm1)=1\) supply target data
only --- they do not compute \(N_\delta^{\mathrm{Pf}}\),
\(\chi_\upsilon\), or \(\epsilon_o\). The Coxeter presentation
\[
W^{(2)}(\La^{2,1}_{II})
=\langle s_{\delta_1},s_{\delta_2},s_{\delta_3}\mid s_{\delta_i}^2=1\rangle
\]
implies that any orientation character on \(W^{(2)}(\La^{2,1}_{II})\)
is determined by three generator signs once \textup{(O1)\(^+\)} gives
the honest Weyl lift; the three rank-one signs established by
\((O2_{\mathrm{atlas}})\) are \(-1\), and then propagate to the full
character.  The
Weyl-equivariant transport is a separate identity, distinct from the
no-extra-relations and PBW identities, established relative to the Weyl lifts in
Theorem~\ref{thm:G-O1-plus}.

\paragraph{Argument.} Three local signs on the simple reflections
imply a character on the full type-II Weyl group only after
\([c_{o,R}]=0\) and \(\omega_{i,\mathcal C}=0\): the Weyl
determinant-line cocycle measures whether the lift is honest, and the
quotient torsor defects measure whether the descent commutes with the
reflection action on the retained groupoid. The cocycle reduces to the
classical group-cohomology class only on a complete
\(W^{(2)}(\La^{2,1}_{II})\)-orbit with constant coefficients; finite
height \(R\) is a partial action groupoid whose component data must be
identified on intersections.

\paragraph{Conditional theorem.} Theorem~\ref{thm:G-O1-plus} is the
criterion for finite-height lifts \(\tau_{i,R}\) whose squares are
one, whose projective cocycle class \([c_{o,R}]\) is zero, and whose
torsor defects \(\omega_{i,\mathcal C}\) vanish on the retained
groupoid components.  These data give Weyl-equivariant
orientation transport once supplied.  In the current finite
orientation packet they are exactly the missing Weyl-lift and
Coxeter rows in the obstruction ledger.  They do not supply the
height-seven source matrices required for the relation-closed
primitive recognition window.

\subsubsection{(O2): local Pfaffian wall normal form on type-II walls}
\label{subsubsec:obstruction-O2}

\paragraph{Type-II wall normal form.} On the three type-II walls
\(\Hpl_{\delta_i}\) for
\(\delta_1=2f_2-f_3,\,\delta_2=2f_{-2}-f_3,\,\delta_3=f_3\), the
retained wall-atlas datum has half-Hilbert sign support of size
\(4096=2^{12}\).  The local computation in Appendix~E is only the
rank-one crossing. Concretely, in an
\(s_\delta\)-equivariant real-analytic d-critical/Kuranishi chart
\((U_\delta,x_\delta,u_\delta,K^{\mathrm{red}}_\delta)\), with
\(K^{\mathrm{red}}_\delta\) the real form of the cosection-reduced
self-Ext complex
\(\operatorname{RHom}_{\mathrm{red}}(\mathcal C_\delta,\mathcal C_\delta)\),
the splitting
\(K^{\mathrm{red}}_\delta\simeq K^{\mathrm{tan}}_\delta\oplus
K^{\mathrm{nor}}_\delta\) is equipped with the extension class
\(\eta_\delta\in
\operatorname{Ext}^1_{D^{s_\delta}(U_\delta)}
(K^{\mathrm{nor}}_\delta,K^{\mathrm{tan}}_\delta)\) trivial after
the reduced \(E\)-quotient and quotient-orientation trivializations,
and the normal Pfaffian admits the rank-one form
\(\operatorname{Pf}(K^{\mathrm{nor}}_\delta)
=\upsilon_\delta\,u_\delta^{N_\delta^{\mathrm{Pf}}}\) with
\(\upsilon_\delta(x_\delta)\ne0\),
\(s_\delta(\upsilon_\delta)=\upsilon_\delta\), and
\(N_\delta^{\mathrm{Pf}}=1\).

\paragraph{The 4096 sign sum.} The integer \(4096=2^{12}\) is the
half-Hilbert sign count retained in \((O2_{\mathrm{atlas}})\): twelve
sign degrees of freedom, each \(\mathbb Z/2\), give \(2^{12}=4096\)
signed states.  The type-II wall locus realizes this count only after
the component, boundary, overlap, quotient, and protected-integration
compatibilities of the retained atlas hold. The arithmetic
match \(-4096=-(64)^2\) with the OP/DT scalar
\(Z^X_{\mathrm{OP}}=-4096\,\Delta_5^{-2}\) is the squared-theta
normalization, fixed by Borcherds normalization and the OP scalar
sign convention; it is \emph{not} a derivation of \(\epsilon_o\).

\paragraph{Criterion.} The local model and the exact separation
from OP and Maass scalars are the following: scalar
constants \(D_5=64^{-1}\Delta_5\),
\(\chi_{10}^{\mathrm{OP}}=4096^{-1}\Delta_5^2\),
\(Z^X_{\mathrm{OP}}=-4096\,\Delta_5^{-2}\) are scalar normalizations
and do not define \(\epsilon_o\); the Pfaffian normal form is the
content.  Proposition~\ref{prop:scalar-pfaffian-data-not-section}
strengthens this separation: the scalar product, its square, and the OP
trace also do not determine \(\operatorname{pf}_{X,R}\) as a section of
the Pfaffian line, because the order-two monodromy is invisible after
squaring.  Theorem~\ref{thm:G-O2-orbit-transport} is relative to
\((O2_{\mathrm{atlas}})\), whose matrices over \(\mathbb Z\) or
\(\mathbb Q\) record the Pfaffian-rank identity
\(N_\delta^{\mathrm{Pf}}=1\), the splitting \(\eta_\delta=0\), the
unit \(\upsilon_\delta\) with \(s_\delta\)-character trivial, and
overlap/transition compatibility across the cofinal HN system.
The packet
\(\texttt{certificates/wall\_atlas/k3e\_o2\_atlas}\) presently supplies
only the obstruction ledger
\(\texttt{O2\_OBSTRUCTION\_LEDGER\_VERIFIED}\), with
\(\texttt{o2\_certification=false}\).  It records the missing atlas
data; it does not supply the wall atlas.

\paragraph{Argument.} The condition \((\delta,\delta)=2\), the
automorphic divisor order
\(d^{\mathrm{aut}}_{\delta_i}=1\), and the Maass character
\(\nu_{\Delta_5}(s_{\delta_i})=-1\) are target data that do not imply
the Pfaffian rank assertion
\(N_\delta^{\mathrm{Pf}}=1\). Graph-isogeny atom candidates
\(B_2=i_{\Gamma_\varphi,*}L\) are Koszul-Ext-local; they are
\emph{candidates}, not constructions, until semistability, full
\(\widehat\Gamma_X\) charge match, reduced normal Ext splitting,
invariant unit, and atlas overlap compatibility are established. The
``higher terms do not alter'' phrase in the rank-one Ext criterion
requires the explicit equivariant real/parametric Morse lemma and
unit check before it operates as a theorem.

\paragraph{Wall-atlas transport.}
Definition~\ref{def:G-O2-atlas} gives the retained type-II wall-atlas
datum.  Theorem~\ref{thm:G-O2-orbit-transport} transports that datum
over the half-Hilbert orbit; it does not construct the graph-isogeny
or reducible-curve atoms.  Theorem~\ref{thm:G-R1-R2} identifies the
\(2^{12}\) sign count with the squared theta-leading normalization
only after scalarization.
The remaining relation-closed source comparison is the recognition
datum \((S1)\)--\((S10)\), not a type-II wall normal-form problem.

\subsection{Mathieu and umbral moonshine adjacency}
\label{subsec:moonshine-adjacency}

The K3 elliptic genus has a precise representation-theoretic
decomposition. Eguchi, Ooguri, and Tachikawa \cite{EOT} observed in 2010 that
the elliptic genus
\(Z_{\mathrm{K3}}(\tau,z)=2\phi_{0,1}(\tau,z)\) admits an
\(M_{24}\)-graded character expansion: the mock-modular completion
\(\Sigma(\tau)\) of the elliptic genus is a graded \(M_{24}\)-module,
and the multiplicities of irreducible \(M_{24}\)-representations in
the first several Fourier coefficients of the holomorphic part match
the dimensions of irreducible \(M_{24}\)-representations on the nose.
Gannon \cite{GannonUmbral} proves the arithmetic identity; the
geometric origin via a sigma-model action of \(M_{24}\) on the K3
elliptic-genus state space remains open outside the classes accounted
for by known K3 sigma-model symmetries.  The \(M_{24}\) action
visible at the level of the elliptic genus does not lift to a global
geometric symmetry of \(D^b(\mathrm{Coh}(K3))\); the work of
Gaberdiel--Hohenegger--Volpato \cite{GHV} on Mathieu twining characters
identifies the precise refinement at the level of supersymmetric
sigma-model conformal field theory data.

The umbral moonshine of Cheng--Duncan--Harvey \cite{CDH} extends Mathieu
moonshine to a family of \(23\) cases indexed by Niemeier root
systems, one case per Niemeier lattice having a non-empty root system.
For a Niemeier lattice \(N^X\) with root system \(X\), the finite
group is the umbral group
\[
G^X=\operatorname{Aut}(N^X)/W^X,
\]
not the discriminant group.  Mathieu moonshine is the \(X=A_1^{24}\)
case.  The Leech lattice is the \(24\)th Niemeier lattice and has empty
root system; it belongs to Conway moonshine, not to the \(23\)-case
umbral list.  The umbral moonshine conjecture of
Cheng--Duncan--Harvey \cite{CDH}, in the module-existence form proved by
Gannon \cite{GannonUmbral}, asserts that each pair \((X,G^X)\) carries a
graded module whose graded character is the prescribed vector-valued
mock modular form.

\paragraph{The conjectural connection to \(\mathfrak g_{\Delta_5}\).}
The Mathieu moonshine module lies on the \(N=4\) superconformal
character side of the K3 elliptic genus.  The Eguchi--Ooguri--Tachikawa
module is graded by the massive level \(n\); the \(z\)-dependence and
the \(l\)-index are carried by the \(N=4\) character, not by a primitive
\((n,l)\)-grading on the \(M_{24}\)-module.  The Beilinson-cut
statement is therefore a conjectural extension of the twined K3
elliptic genus to the Borcherds--Kac--Moody datum on
\(\La^{2,1}_{II}\), producing a \(g\)-twined refinement of
\(\mathfrak g_{\Delta_5}\) with denominator
\(\operatorname{den}(\mathfrak g_{\Delta_5}^{(g)})\) for each
\(M_{24}\)-conjugacy class \(g\). The umbral cases of
\cite{CDH,ChengHarrison} extend the conjecture to a family of \(23\)
Borcherds--Kac--Moody-style algebras associated to the \(23\) Niemeier
root systems with non-empty root sets.

\paragraph{Precise conjecture.}
For each conjugacy class \(g\in M_{24}\), let
\(\phi^{(g)}_{0,1}(\tau,z)\) denote the \(g\)-twined K3 elliptic
genus of \cite{GHV}.  The Cheng--Duncan--Harvey--Gannon umbral
moonshine conjecture, refined to the BKM stratum of
\(\La^{2,1}_{II}\), asserts the existence of a \(g\)-twined
Borcherds--Kac--Moody algebra
\(\mathfrak g_{\Delta_5}^{(g)}\) with denominator
\[
 \operatorname{den}\bigl(\mathfrak g_{\Delta_5}^{(g)}\bigr)
 \;=\;\Theta\bigl(\phi^{(g),\mathrm{Weil}}_{0,1}\bigr)\!\bigm|_{\Hpl_2},
\]
the singular theta lift of the \(g\)-twined Weil-representation
extension of \(\phi^{(g)}_{0,1}\), with the chamber and Weyl-vector
data determined by the relevant umbral pair \((X,G^X)\), when such a
pair is part of the twining datum.  No map from all \(M_{24}\)-classes
to Niemeier root systems is asserted here.

\paragraph{The order-seven and order-fifteen classes.}
The four classes \(\{7A,7B,15A,15B\}\) lie outside the diagonal
Gritsenko--Cl\'ery/Borcherds frame used here, namely the rows selected
by \(\varphi(N)\mid2\); this does not exclude their occurrence in
other CHL dyon towers.  Their cyclotomic fields have degrees \(6\) and
\(8\).  Their cycle shapes on the standard \(24\)-dimensional permutation
representation are
\[
 7A:\ 1^3\!\cdot 7^3,\qquad
 7B:\ 1^3\!\cdot 7^3,\qquad
 15A:\ 1\!\cdot 3\!\cdot 5\!\cdot 15,\qquad
 15B:\ 1\!\cdot 3\!\cdot 5\!\cdot 15,
\]
the \(7A\)/\(7B\) pair distinguished by Galois conjugation under
\(\mathrm{Gal}(\Q(\zeta_7)/\Q)\) and the \(15A\)/\(15B\) pair by
\(\mathrm{Gal}(\Q(\zeta_{15})/\Q)\); the rational character on each
pair is integer-valued, but individual class characters take values
in the cyclotomic field of order equal to the class order.  These
cycle shapes have different geometric status: \(7A\) and \(7B\) have
six orbits and pass the Mukai--Kondo orbit test, while \(15A\) and
\(15B\) have four orbits and are non-geometric in that test
\cite[\S2]{GaberdielVolpatoOrbifoldK3}.  The orders
\(|g|\in\{7,15\}\) satisfy \(\varphi(|g|)\nmid 2\)
(\(\varphi(7)=6\), \(\varphi(15)=8\)), placing them outside the
Gritsenko--Cl\'ery frame \(N\in\{1,2,3,4,6\}\) of \(\kappa_{\mathrm{BKM}}\)-universality
selected by \(\varphi(N)\mid 2\) of
Theorem~\ref{thm:bkm-kappa-universal}.  A Borcherds--Kac--Moody
construction for these twined genera therefore requires data beyond
that denominator frame.

\begin{conjecture}[Order-seven and order-fifteen twined denominators]
\label{conj:order-seven-fifteen-twined-denominators}
For each \(g\in\{7A,7B,15A,15B\}\):
\begin{enumerate}
\item[\textnormal{(M${}^\star$\!7A)}] The Gaberdiel--Hohenegger--Volpato
\(7A\)-twined K3 elliptic genus
\(\phi^{(7A)}_{0,1}\) is a weight-\(0\), index-\(1\) weak Jacobi form
with multiplier for the appropriate congruence subgroup attached to
the \(7A\) class.  Its singular theta lift, if the Weil extension and
Weyl chamber data are fixed, has Borcherds weight
\(c_{7A}(0)/2\) and multiplier, and equals
\(\operatorname{den}(\mathfrak g_{\Delta_5}^{(7A)})\) for a
conjectural \(7A\)-twined Borcherds--Kac--Moody algebra.
\item[\textnormal{(M${}^\star$\!7B)}] Identical statement for the
Galois conjugate class \(7B\), with
\(\phi^{(7B)}_{0,1}=\phi^{(7A)}_{0,1}\circ
\sigma_{\zeta_7\to\zeta_7^3}\) the Galois-conjugate character.
\item[\textnormal{(M${}^\star$\!15A)}] The
\(15A\)-twined character \(\phi^{(15A)}_{0,1}\) is a weight-\(0\),
index-\(1\) weak Jacobi form with multiplier for the appropriate
congruence subgroup attached to the \(15A\) class.  Its singular theta
lift, if the Weil extension and Weyl chamber data are fixed, has
Borcherds weight \(c_{15A}(0)/2\) and multiplier, and equals
\(\operatorname{den}(\mathfrak g_{\Delta_5}^{(15A)})\).
\item[\textnormal{(M${}^\star$\!15B)}] Identical statement for the
Galois conjugate class \(15B\).
\end{enumerate}
The four clauses have finite necessary conditions in their first
Fourier coefficients, but no finite truncation proves the denominator
identity.
\end{conjecture}

\begin{remark}[Fourier conditions and the parity obstruction]
\label{rem:order-seven-fifteen-cross-checks}
Three Fourier conditions constrain each clause
\((\mathrm M^\star\!g)\) at finite truncation depth.

\textup{(i)} The leading Fourier coefficient
\(c_g(0)=[q^0 r^0]\phi^{(g)}_{0,1}\) is the
Gaberdiel--Hohenegger--Volpato McKay--Thompson value of \(g\) at the
\(q^0 r^0\) position of the K3 elliptic genus, tabulated in
\cite[Table~1]{GHV} for all 26 conjugacy classes.  Outside the CHL
frame \(\{1A,2A,3A,4A,6A\}\), this leading coefficient need not be an
even integer in the Borcherds lattice normalisation; for the
classes \(\{7A,7B,15A,15B\}\) it is not one of the values
\(c_N(0)\in\{10,6,4,3,2\}\) of
Theorem~\ref{thm:bkm-kappa-universal} \textup{(within the frame the
\(N=4\) entry \(c_4(0)=3\) is already odd, with half-integral weight
\(\tfrac32\))}.

\textup{(ii)} The Borcherds weight formula gives
\(\kappa_{\mathrm{BKM}}(\Phi^{(g)})=c_g(0)/2\) only after the twined
Borcherds datum and its lattice have been fixed.  For the
order-seven and order-fifteen classes this weight need not be an integer; the lift must
therefore be treated as a multiplier-bearing automorphic product, not
as a weight-\(5\) paramodular form.  This is the arithmetic obstruction
to extending the level-\(\mathsf Z\) denominator construction beyond
the CHL frame.

\textup{(iii)} Mathieu-moonshine consistency with the EOT
\(M_{24}\)-multiplicity predictions through the first eight Fourier
coefficients of \(\Sigma(\tau)\) imposes a finite necessary condition
on each clause; proving
\((\mathrm M^\star\!g)\) requires identifying the geometric origin
of the \(M_{24}\)-action on the K3 elliptic-genus state space and its
umbral refinement in the sense of Cheng--Harrison \cite{ChengHarrison}.
\end{remark}

\paragraph{Independence from the four Pfaffian--Dirac obstructions.}
The Pfaffian theorem~\ref{thm:ch6-pfaffian-dirac} and the
orientation-character theorem~\ref{thm:ch6-orientation-character} are
retained-data theorems on \(K3\times E\), after the retained D0-HN datum,
retained quotient-orientation datum, chosen Weyl lifts,
\((O2_{\mathrm{atlas}})\), and the retained finite geometric Pfaffian
datum \((P_{\mathrm{fin}})\) of
Theorem~\ref{thm:G-four-obstruction-criterion}.  The primitive
recognition theorem remains conditional on \((S1)\)--\((S10)\), and
the level-\(\mathsf A\) gravity-line operator algebra remains conditional on
the Vol~II Hall--Borcherds residual.  The moonshine
adjacency is independent of both: it lives at level \(\mathsf{Z}\)
(centre/denominator) and at level \(\mathsf{A}\) (acting object on
the K3 elliptic-genus state space).  A proven identification of the
moonshine module as a \(g\)-twined character of
\(\mathfrak g_{\Delta_5}\) does not by itself solve the Vol~II
Hall--Borcherds residual; conversely the Hall--Borcherds residual
does not prove the moonshine adjacency.

\paragraph{Twined denominator problem.} The \(g\)-twining operation would
produce \(\operatorname{den}(\mathfrak g_{\Delta_5}^{(g)})\) from
\(\operatorname{den}(\mathfrak g_{\Delta_5})\) and an \(M_{24}\)-action
on a denominator-compatible lift of the Mathieu module.  The ordinary
massive-level grading of the \(M_{24}\)-module supplies the mock modular
coefficients; the Borcherds product requires the additional Jacobi
variable and \(\La^{2,1}_{II}\)-root grading.  That lift, and the
denominator-compatible action on it, are not part of the EOT theorem.

\paragraph{Compatibility with Vol III.} The Vol~III dimension-stratified
BKM catalogue \cite{CYQGVolIII} places \(\mathfrak g_{\Delta_5}\) as
the \(d=3\) sibling, with companions at \(d=3\) (Borcherds Monster
\(V^\natural\) on \(\mathrm{II}_{1,1}\)) and \(d=5\) (Fake Monster on
\(\mathrm{II}_{25,1}\) via \(K3\times K3\times E\)).
The Mathieu/umbral moonshine adjacency to \(\mathfrak g_{\Delta_5}\)
is a level-\(\mathsf{Z}\) cross-stratum conjecture: the denominator-compatible
\(M_{24}\) action on the K3 elliptic-genus state space for
\(\{7A,7B,15A,15B\}\) is the missing piece on the K3-fibre side; the
\(g\)-twined denominator
\(\operatorname{den}(\mathfrak g_{\Delta_5}^{(g)})\) is the missing
piece on the BKM side.

\paragraph{Twined denominator identity.} For each of the \(26\)
conjugacy classes of \(M_{24}\), the expected object is a
\(g\)-twined Borcherds product on \(\La^{2,1}_{II}\).  Its
specialization to the \(N=4\) character expansion must recover the
twined McKay--Thompson coefficients of EOT--Gannon.  A coefficientwise
refinement of the EOT multiplicities by BKM denominator coefficients is
therefore a consequence only after the denominator-compatible Jacobi
lift and the \(\La^{2,1}_{II}\)-root grading have been fixed.  The
required identity is the consistency of the \(g\)-twined
denominator with the OP scalar
\(Z^X_{\mathrm{OP}}=-4096\,\Delta_5^{-2}\) under the \(g\)-twining
construction, including the squared sign convention.

\subsection{Residual operator structures}
\label{subsec:subsequent-constructions}

The Beilinson cut separates the remaining operator data by the object
on which the missing structure must live: the Hall--Borcherds operator
lift, the compact non-toric chain-fusion equivalence, and the
gravity-line algebra.

\subsubsection{The Hall--Borcherds residual}
\label{subsubsec:open-hall-borcherds}

\paragraph{Hall--Borcherds residual.} The level-\(n\) protected scalar shadow
\(Z^{K3\times E}_{\mathrm{BPS}}=(\Phi_{10}^{\mathrm{un}})^{-1}=
\Delta_5^{-2}\) on the closed
\(K3\times E\mapsto\,\mathrm{point}\) cobordism is the BPS partition
function.  A 3d gravitational path-integral interpretation with a
genuine operator-level lift requires the Hall--Borcherds residual: an
operator-valued completion of the scalar trace whose trace recovers the
level-\(n\) shadow and whose algebraic content carries the metric
degrees of freedom.  This residual remains open in Vol~II.

\paragraph{Comparison with Vol II.} Vol~II \cite{ChiralBarCobarVolII}
names the gravity-line operator algebra required to have Pentagon-face
determinant section \(\Phi_{10}^{\mathrm{un}}=\Delta_5^2\), whose
partition character is \((\Phi_{10}^{\mathrm{un}})^{-1}=\Delta_5^{-2}\), and identifies the
Hall--Borcherds residual as the comparison arrow from the bulk
\(Z^{\mathrm{der}}_{\mathrm{ch}}(A_b)\) to the line-operator algebra.
The Vol~II residual asks for the construction of the
gravity-line operator algebra with Pentagon-face determinant
\(\det_{\mathrm{Pentagon}}=\Phi_{10}^{\mathrm{un}}=\Delta_5^2\)
on the closed cobordism, whose primitive square-root Borcherds weight is
\(\omega_{\mathrm{Borcherds}}(\Delta_5)=c_1(0)/2=5\) at the
K3-Heisenberg witness.
The Igusa-side datum is the same automorphic determinant identity
\(\Delta_5^2=\Phi_{10}^{\mathrm{un}}\), with the Appendix~G Pfaffian datum
\(\mathfrak D_X^{\mathrm{DI}}\) carrying the retained orientation and
Weyl-transport data on the chiral side.

\paragraph{Argument.} The protected graded trace
\[
\mathrm{Tr}_{\mathrm{prot}}\bigl((-1)^F
q^{L_0}r^{J_0}s^{\widetilde{L}_0}\bigr)=\Delta_5^{-2}
\]
forgets orientation; the OP chamber representative is
\(-4096\) times this trace.  The
sign-character \(\epsilon_o\) is invisible in the squared form and is
part of the Pfaffian--Dirac orientation theorem, not of the
gravity-line residual.  The Hall--Borcherds residual asks instead for
an acting operator algebra whose trace recovers the scalar shadow.
The realization as a three-dimensional gravitational path integral is conditional on
Brown--Henneaux, modular invariance, vacuum dominance, and saddle
dominance \cite{HM}; no such realization is asserted.

\paragraph{Scalar non-faithfulness.} The structural reason
\(\Delta_5^{-2}\) does not by itself recover the operator-level
gravity-line algebra is the scalar non-faithfulness of the graded
Euler character on filtered \(SC^{\mathrm{ch},\mathrm{top}}\)-objects,
proved as a lemma in Vol~III~\cite{CYQGVolIII}: an equality
\(\chi(C_1) = \chi(C_2)\) between filtered
\(SC^{\mathrm{ch},\mathrm{top}}\)-objects does not determine a
quasi-isomorphism, an \(SC^{\mathrm{ch},\mathrm{top}}\)-morphism, a
Hall product, a Drinfeld pairing, a current-envelope locality
structure, or a derived-centre trace --- adding a filtered acyclic
summand \(K\) with \(\chi(K) = 0\) leaves the scalar unchanged while
admitting non-trivial extensions of the operadic action, and a
filtered differential may be perturbed within the associated graded
without altering \(\chi\). The missing chain-level content underneath
\(\Delta_5^{-2}\) sits under Vol~III's Hall--Borcherds datum: a
positive-half Hall--Borcherds bialgebra morphism on \(E\), a completed
Drinfeld-double extension, a current-envelope morphism along \(E\), and
trace compatibility with \(\Delta_5^{-2}\) \cite{CYQGVolIII}.  The
Igusa scalar is only the trace shadow of this datum.  The
Hall--Borcherds residual is the missing chain-level lift
together with its finite-window identities.

\paragraph{Pentagon-face cyclic trace.} The required cyclic trace is
\(\mathrm{Tr}_{\mathrm{Pentagon}}^{\mathrm{cl}}\) on the closed-colour
factorization algebra \(F^{\mathrm{cl}}\) at the
\(K3\times E\)-Heisenberg witness.  Its evaluation must recover
\(c_1(0)/2=5\) on the universal Borcherds-weight identity, and its
operator-level lift to the open-colour boundary algebra
\(A_b=\mathrm{End}_{\mathcal C}(b)\) must factor through the
chain-fusion datum below.  The compatibility condition is the
consistency of the Pentagon-face scalar with the unscaled trace
\(\Delta_5^{-2}\) under the squared-trace operation.  The OP chamber
representative is \(-4096\) times this trace.

\subsubsection{Chain fusion for compact non-toric K3{\(\times\)}E}
\label{subsubsec:open-chain-fusion}

\paragraph{Chain-fusion datum.} For a Calabi--Yau\(_d\) category
\(\mathcal C_X\) and a chart datum \((\Sigma_{d-1},C)\), a
chain-fusion datum consists of a boundary vacuum
\(b(X,\Sigma,C)\) in an open factorization dg-category on
\((C,D_C,\tau_C)\), a comparison morphism
\[
 \Phi^{(\Sigma_{d-1},C)}_d(\mathcal C_X)
 \longrightarrow
 \mathrm{End}_{\mathcal C}(b(X,\Sigma,C)),
\]
and the compatibility of this morphism with clutching, trace, and the
Calabi--Yau orientation.  The chain-fusion conjecture asserts that this
morphism is a quasi-isomorphism for the admissible compact chart.  The
affine toric examples \(\mathbb C^3\), local \(\mathbb P^2\), and the
conifold furnish local comparison models; they do not construct the
compact non-toric boundary vacuum for \(K3\times E\).  At \(d=3\), the
retained Dirac--Igusa data supply the protected scalar shadow on
\(K3\times E\); the boundary-vacuum construction required for chain
fusion remains open.

\paragraph{Comparison with Vol III.} Vol~III's catalogue
\cite{CYQGVolIII} names this conjecture in the two-stage
factorization frame and identifies, after chain fusion, the
level-\(\mathsf Z\) Drinfeld double obligation
\(G(X)=D(Y^+(X))\) for compact non-toric \(\mathrm{CY}_3\), together
with the conditional TCFT/framing/anomaly/completion datum needed
for compact non-formal \(\mathrm{CY}_3\) in general. Chain fusion
remains a conjecture, not a theorem; the assertion of \(\Delta_5\) as
a protected Pfaffian section on
\(K3\times E\) is the retained-data identity
\(\mathrm{Pf}_{\mathrm{prot}}(\mathfrak D_X^{\mathrm{DI}})=\mathcal D_X=\Delta_5\)
of view \textcircled{3}, not an underlying microscopic state space.

\paragraph{Argument.} The compact \(K3\times E\) case at \(d=3\) is a
candidate instance of chain fusion via the rectified finite
\(K3\)-to-\(E\) specialization
(Definition~\ref{def:O1-strong-reduced-orientation}).
The proposed chiral specialization, if constructed from the retained
finite-stage datum, has target form
\(\mathrm{Sp}^{\mathrm{ch}}_{K3,E}(\mathfrak A_{K3\times E})
=U^{\mathrm{ch}}(\mathfrak{heis}_{\mathrm{Muk}})\otimes
U^{\mathrm{ch}}(\mathfrak g^{\mathrm{BPS}}_{K3})\)
on \(\mathrm{Ran}(E)\).  The map
\(\mathrm{Sp}^{\mathrm{ch}}_{K3,E}\) is the chosen
\([K3{\to}E\text{-}\mathrm{sp}]\) datum, not a canonical functor.  The
finite proof rows for that datum are
\(\mathrm{formal\_target\_charts}\), \(\mathrm{field\_complexes}\),
\(\mathrm{e3\_operations}\), \(\mathrm{bv\_qme}\),
\(\mathrm{anomaly\_framing}\), \(\mathrm{compact\_support}\),
\(\mathrm{factorization\_descent}\),
\(\mathrm{k3\_to\_e\_specialization}\), \(\mathrm{transitions}\), and
\(\mathrm{scalar\_firewall}\); without them the compact source is only
a named input.  The finite packet
\(\texttt{certificates/e3\_source/k3e\_e3\_holfa}\) currently records
these rows only by the obstruction ledger
\(\texttt{E3\_SOURCE\_OBSTRUCTION\_LEDGER\_VERIFIED}\), with
\(\texttt{e3\_source\_certification=false}\).  The
boundary algebra \(A_b=\mathrm{End}_{\mathcal C}(b)\) on the open
factorization side can coincide with this target only after the
chain-fusion boundary vacuum \(b(K3\times E,\Sigma_2,E)\) exists.
Compact non-toric
\(K3\times E\) generalizations to other compact non-toric
\(\mathrm{CY}_3\)-data (e.g.\ Borcea--Voisin threefolds, Schoen's
small resolutions) require independent boundary vacua.

\paragraph{Boundary vacuum condition.} The missing object is
\(b(K3\times E,\Sigma_2,E)\), a boundary vacuum in an open
factorization dg-category on
\((E,D_E,\tau_E)\) with \(D_E\) at the Igusa cusp
\(\mathfrak c_\infty\) of the Siegel space.  The chain-fusion
isomorphism on \(W_{\le3}\) at finite HN height \(R=3\) must identify
the chiral algebra with
\(A_{b(K3\times E,\Sigma_2,E)}\). The necessary
compatibility condition is the consistency of the chain-fusion isomorphism with
the unscaled trace \(\Delta_5^{-2}\), whose OP representative is
\(-4096\,\Delta_5^{-2}\), and with the universal Borcherds-weight
identity \(\kappa_{\mathrm{BKM}}(\Phi_1)=c_1(0)/2=5\); see the
\(\kappa\)-tuple discussion below.

\subsubsection{The gravity-line operator algebra}
\label{subsubsec:open-gravity-line}

\paragraph{Gravity-line algebra.} A gravity-line operator algebra, if constructed,
is an operator algebra on the closed-colour factorization algebra
\(F^{\mathrm{cl}}\) on \(K3\times E\) whose Pentagon-face determinant
section is required to be \(\Delta_5^2=\Phi_{10}^{\mathrm{un}}\), a
weight-\(10\) automorphic section.  Its partition character is the
reciprocal section \(\Delta_5^{-2}\).  The primitive square-root row has
Borcherds weight
\[
\omega_{\mathrm{Borcherds}}(\Delta_5)
=\kappa_{\mathrm{BKM}}(\Delta_5)=c_1(0)/2=5.
\]
This is the same Hall--Borcherds residual and remains open in Vol~II.

\paragraph{Comparison with Vol II.} Vol~II \cite{ChiralBarCobarVolII}
names the gravity-line Hall--Borcherds comparison conjecture and the
Pentagon-face determinant;
Vol~III \cite{CYQGVolIII} names
\(\kappa_{\mathrm{BKM}}(\Phi_N)=c_N(0)/2\) as the universal
Borcherds-weight identity \cite{Borcherds95}, with
\(\kappa_{\mathrm{BKM}}(\Phi_1)=c_1(0)/2=5\) at the paramodular
\(\Delta_5\) datum.

\paragraph{The \(\kappa\)-tuple.} The five-fold \(\kappa\)-tuple
\((\kappa_{\mathrm{cat}},\kappa_{\mathrm{ch}}^{\mathrm{Hodge}},
\kappa_{\mathrm{ch}}^{\mathrm{Heis}},\kappa_{\mathrm{BKM}},\kappa_{\mathrm{fiber}})
=(0,0,3,5,24)\) at \(K3\times E\) records five separate invariants,
distinct in construction and meaning, never additively combined. Here
\(\kappa_{\mathrm{cat}}=\chi(\mathcal O_{K3\times E})=
\chi(\mathcal O_{K3})\cdot\chi(\mathcal O_E)=2\cdot 0=0\) is the
K{\"u}nneth-multiplicative categorical Euler;
\(\kappa_{\mathrm{ch}}^{\mathrm{Hodge}}=\sum_q(-1)^qh^{0,q}(K3\times E)=0\)
by Serre duality on the compact total space;
\(\kappa_{\mathrm{ch}}^{\mathrm{Heis}}=3\) is the rank of the
Heisenberg--Cartan in the chiral Hall specialization on
\(\mathrm{Ran}(E)\) and equals the Cartan rank of
\(\mathfrak g_{\Delta_5}\);
\(\kappa_{\mathrm{BKM}}=c_1(0)/2=5\) is the universal Borcherds-weight
identity at the Gritsenko--Cl{\'e}ry row whose Jacobi seed is
\(\phi_{0,1}\) and whose Borcherds product is \(\Delta_5\), with
\(c_1(0)=10\);
\(\kappa_{\mathrm{fiber}}=24\) is the Mukai-lattice rank
\(\mathrm{rk}\,\widetilde H^*(K3,\mathbb Z)=24\) of the K3 fibre, equal
to \(\chi_{\mathrm{top}}(K3)\). The additive identity
\(\kappa_{\mathrm{BKM}}=\kappa_{\mathrm{ch}}+\chi(\mathcal O_{\mathrm{fiber}})\)
fails because it confuses
\(\chi_{\mathrm{top}}(K3)=24\) with
\(\chi(\mathcal O_{K3})=2\); the universal identity
\(\kappa_{\mathrm{BKM}}(\Phi_N)=c_N(0)/2\) is the correct frame.

\paragraph{Argument.} The Pentagon-face determinant is the
modularity-witness of the open-side cyclic trace plus clutching;
modularity is a property of the open trace, not of the closed
algebra.  On the \(\Delta_5\) row, \(c_1(0)/2=5\) is the weight of the
primitive Borcherds factor.  The determinant required by chain fusion
is its square
\(\Delta_5^2=\Phi_{10}^{\mathrm{un}}\), hence a form of total weight \(10\)
on the closed cobordism; the level-\(n\) protected scalar shadow is
the reciprocal Laurent expansion \(\Delta_5^{-2}\).

\paragraph{Gravity-line trace condition.} The required object is the
gravity-line operator algebra
\(\mathrm{GravLine}_{K3\times E}\) on the open factorization
dg-category on \((E,D_E,\tau_E)\) at the Igusa cusp, with
Pentagon-face determinant identity
\(\det_{\mathrm{Pentagon}}^{\mathrm{cl}}(\mathrm{GravLine}_{K3\times E})
=\Delta_5^2\) and compatibility with the universal
Borcherds-weight identity
\(\kappa_{\mathrm{BKM}}(\Phi_1)=5\) and with the \(\kappa\)-tuple
\((0,0,3,5,24)\) at the K3-Heisenberg witness. The compatibility
condition is the existence of chain fusion at
\(b(K3\times E,\Sigma_2,E)\) and the Hall--Borcherds residual lifting
the protected trace
\(Z_{\mathrm{BPS}}^{K3\times E}=\Delta_5^{-2}\) toward the acting
operator-algebra level.

The finite compact-source recognition datum \((S1)\)--\((S10)\) on the
first relation-closed window is the packet
\[
\mathrm{window\_closure},\quad
\mathrm{target\_source\_representatives},\quad
\mathrm{chevalley\_serre\_matrices},\quad
\mathrm{hall\_boundary\_complex},\quad
\mathrm{spectral\_sequence},\quad
\mathrm{kernel\_matrices},\quad
\mathrm{pbw\_graded},\quad
\mathrm{first\_window\_theorems},\quad
\mathrm{transitions},\quad
\mathrm{scalar\_firewall}.
\]
The compact Hall source packet presently has a verified obstruction
ledger,
\[
\texttt{certificates/sources/k3e\_compact\_hall/blocked\_obligations.csv},
\]
which lists the missing \(M,D,\eta,\epsilon,B,G,K,Q,A\) matrices, the
Hall bialgebra identity rows, Hopf-pairing identity rows,
radical ideal/coideal rows, relation rows, no-extra equality,
generation, PBW, strict transition, \(A_\beta\)-intertwining, and
finite Koszul comparison rows.  Its verifier records
\(\texttt{compact\_source\_recognition=false}\).  This ledger is the
finite obstruction list for the source theorem, not a substitute for
the source theorem.
The first relation-closed window has its own obstruction ledger,
\[
\texttt{certificates/first\_window/k3e\_relation\_closed\_window/blocked\_obligations.csv},
\]
with status
\(\texttt{FIRST\_WINDOW\_OBSTRUCTION\_LEDGER\_VERIFIED}\) and
\(\texttt{first\_window\_certification=false}\).  It records the
missing window, representative, relation, Hall-boundary, spectral,
kernel, PBW, theorem, and transition rows for the first finite
primitive-recognition theorem.
The Dirac--Igusa pro-object morphism packet has its own obstruction
ledger,
\[
\texttt{certificates/morphisms/k3e\_dirac\_igusa\_pro\_object/blocked\_obligations.csv},
\]
with status
\(\texttt{MORPHISM\_OBSTRUCTION\_LEDGER\_VERIFIED}\) and
\(\texttt{pro\_object\_certification=false}\).  It records the missing
cofinal HN subsystem, fourteen finite-stage component rows, strict
transition morphisms, component maps for \((M1)\)--\((M8)\),
composition laws, Mittag--Leffler exactness, and pro-isomorphism rows.
The theorem-dependency ledger
\[
\mathrm{beilinson\_levels},\quad
\mathrm{objects},\quad
\mathrm{morphisms},\quad
\mathrm{theorem\_dependencies},\quad
\mathrm{dependency\_edges},\quad
\mathrm{forbidden\_edges},\quad
\mathrm{status\_separation},\quad
\mathrm{level\_equalities},\quad
\mathrm{equality\_sites},\quad
\mathrm{definition\_separation},\quad
\mathrm{obstruction\_independence},\quad
\mathrm{transitions},\quad
\mathrm{scalar\_firewall}
\]
is now a populated schema-complete ledger.  Its verifier checks
Beilinson-level sites, typed objects, morphisms, theorem rows, allowed
proof edges, forbidden proof edges, claim-status separation,
Beilinson-level equality scopes, equality-site rows, transition rows,
definition-separation rows, obstruction-independence rows, and
scalar-firewall coverage; it still records
\(\texttt{dependency\_certification=false}\) and
\(\texttt{mathematical\_certification=false}\).  Thus the proof-level
exclusions are finite rows, but the retained construction packets,
recognition matrices, mirror comparison, and gravity-line operator
algebra remain separate mathematical obligations.  Together with these
packets, the chain-fusion boundary vacuum
\(b(K3\times E,\Sigma_2,E)\), the Hall--Borcherds residual lifting
\(Z_{\mathrm{BPS}}^{K3\times E}=\Delta_5^{-2}\), the gravity-line
operator algebra with Pentagon-face determinant \(\Delta_5^2\), and the
Mathieu-twined denominators
\(\operatorname{den}(\mathfrak g_{\Delta_5}^{(g)})\) on the \(26\)
\(M_{24}\)-conjugacy classes are the unresolved data.  The four
obstructions concern, respectively, the D0-degeneration of the reduced
orientation theory, quotient orientation, Weyl transport of the
orientation line, and the type-II Pfaffian wall atlas.  The remaining
operator obstructions concern the Hall--Borcherds operator lift, compact
non-toric chain fusion, and the gravity-line algebra.  The Beilinson cut
separates these objects before any scalar trace is taken.


% ------------------------------------------------------------------
% Appendices
% ------------------------------------------------------------------
\appendix

\section{Lattice computations}
\label{app:lattices}

The seven integral lattices used in the manuscript are recorded here in
one place: the hyperbolic plane $U=\mathrm{II}_{1,1}$, the Lorentzian
root lattice $\La_{II}^{2,1}$ of the Borcherds--Kac--Moody algebra
$\autcor$, the K3 cohomology lattice $\La_{K3}=\mathrm{II}_{3,19}$, the
lattice-polarized K3 sublattice $\La_S$, the Mukai lattice
$\widetilde H(S,\Z)\simeq\mathrm{II}_{4,20}$, the formal dyonic charge
lattice $\Gamma_X^{\mathrm{form}}=\La_S\oplus\La_S$, and the
signature-$(3,2)$ lattice $\La^{3,2}$ that bridges $\Sp_4(\Z)$ to
$\SO_+(3,2)$ via $\PSp_4(\Z)\simeq\SO_+(\La^{3,2})$.  Each is given by
its Gram matrix or root data, signature, discriminant, and
automorphism group, with primary citations (Conway--Sloane
1999~\S2.5,~\S4,~\S16; Nikulin 1980 Theorem~1.1.2 and
Theorem~1.14.4; Mukai 1984; Gritsenko--Nikulin 1998
Part~II~\S2,~\S5).  Ranks, signatures, and the Cartan matrix of the
type-II simple roots are verified against the rational computation
\texttt{compute/verify\_lattice.py}.

This appendix is upstream of every chapter: load-bearing identities
about $\Delta_5$, $\autcor$, the orientation character $\nu_{\Delta_5}$,
and the Pfaffian wall normal form on
$W^{(2)}(\La_{II}^{2,1})$ refer back to the lattice data fixed
here, not to a sign or scaling chosen in the body.

\begin{convention}[Sign of the bilinear form; orientation of the
time-like cone]
\label{conv:lattice-signs}
The standard mathematical convention for the hyperbolic plane $U$ is
$U=(\Z^2,((0,1),(1,0)))$, even unimodular of signature $(1,1)$ and
discriminant $-1$ (Conway--Sloane 1999~\S2.5).  This appendix takes
this convention as primary.

The Borcherds--Gritsenko--Nikulin packaging used in
\S\ref{ssec:exceptional-iso} of the body and the exterior-square
model of $\Sp_4(\Z)$ employs the negated form
\[
\HypPlan=\Big(\Z^2,\begin{pmatrix}0&-1\\-1&0\end{pmatrix}\Big),
\]
so that the time-like cone $\Cone(\La^{2,1})_+\subset
\La^{2,1}\otimes\R$ has \emph{negative} square $(x,x)<0$.  The
relation is $\HypPlan=U(-1)$.  Every Gram matrix below is given in the
standard convention; a sign flip transports it to the body
convention without changing rank, discriminant absolute value, or
automorphism group.
\end{convention}

\subsection{The hyperbolic plane $U=\mathrm{II}_{1,1}$}
\label{app:lattice-U}

The hyperbolic plane is the unique even unimodular indefinite lattice
of signature $(1,1)$ (Conway--Sloane 1999~\S2.5; Nikulin 1980
Theorem~1.1.2).

\begin{definition}[$U=\mathrm{II}_{1,1}$]
\label{def:hyperbolic-plane}
\textnormal{[definitional, $\alpha$]}
Let $U:=\Z e\oplus\Z f$ with Gram matrix
\[
G_U=\begin{pmatrix}0&1\\1&0\end{pmatrix},
\qquad e^2=f^2=0,\quad e\cdot f=1.
\]
\end{definition}

\begin{verification}[Rank, signature, discriminant of $U$]
\label{ver:U-invariants}
\textnormal{[computed, $\gamma$]}
$\rk U=2$.  The eigenvalues of $G_U$ are $\pm1$, so the signature is
$(1,1)$.  The discriminant is
\[
\operatorname{disc} U=\det G_U=-1.
\]
$U$ is even (every $G_U$-diagonal entry is $0\equiv 0\pmod2$) and
unimodular ($|\det G_U|=1$).  These are the defining invariants of
$\mathrm{II}_{1,1}$ in the Conway--Sloane classification (Conway--Sloane
1999~\S2.5, Theorem~2.5.1).
\end{verification}

\begin{proposition}[$\mathrm{Aut}(U)$]
\label{prop:U-aut}
\textnormal{[proved, $\alpha$]}
The orthogonal group of $U$ is
\[
\mathrm{O}(U)=\{\pm\Id\}\times\langle\sigma\rangle
\simeq(\Z/2)^{2},
\qquad \sigma:e\leftrightarrow f.
\]
The proper subgroup $\mathrm{O}^+(U)\simeq\Z/2$ swaps $e$ and $f$ but
fixes the spinor norm.
\end{proposition}

\begin{proof}
An isometry $\phi\in\mathrm{O}(U)$ preserves the isotropic vectors of
$U$.  In rank $2$ over $\Z$, the isotropic primitive vectors are
exactly $\pm e,\pm f$, so $\phi$ permutes these four vectors.  The
permutations preserving the bilinear form give the four-element group
above (Conway--Sloane 1999~\S4.1).
\end{proof}

\begin{remark}[Standard sign convention versus body convention]
\label{rem:U-vs-HypPlan}
The body of the manuscript uses
$\HypPlan=U(-1)$
in the cusp polarization of \S\ref{ssec:exceptional-iso} and in the
exterior-square model of the automorphic group, where the time-like
cone of $\La^{2,1}$ has negative square.
The map
$U\to U(-1)$, $(e,f)\mapsto(e,f)$
is a bijection of underlying $\Z$-modules; only the bilinear form
flips sign.  Lattices that occur in the body
($\La^{2,1}$, $\La^{3,2}$) are presented in body coordinates so that
the Gram matrices match line-for-line with the manuscript text;
even-unimodular ambient lattices ($\La_{K3}$, $\widetilde H(S,\Z)$)
are presented in the standard Conway--Sloane convention.  Each
section flags its choice explicitly.
\end{remark}

\subsection{The Lorentzian root lattice $\La_{II}^{2,1}$}
\label{app:lattice-II21}

The Borcherds--Kac--Moody algebra $\autcor$ of
\cite[Theorem~2.1, Section~5.1.1, case $(t=1,II,\overline 0)$]{GNII}
is graded by an integral lattice of signature $(2,1)$.  This is the
lattice on which the type-II simple-root Cartan matrix is
$2(\Id_3-\J_3)$ where $\J_3$ is the all-ones matrix.

\begin{definition}[Ambient $\La^{2,1}=U(-1)\oplus\langle 2\rangle$]
\label{def:lambda-21-ambient}
\textnormal{[definitional, $\alpha$]}
Set
\[
\La^{2,1}:=U(-1)\oplus\langle 2\rangle
=\Z e\oplus\Z f\oplus\Z h,
\qquad e\cdot f=-1,\ e^2=f^2=0,\ h^2=2.
\]
The Gram matrix in the basis $(e,f,h)$ is
\[
G_{2,1}=\begin{pmatrix}0&-1&0\\ -1&0&0\\ 0&0&2\end{pmatrix},
\qquad \det G_{2,1}=-2,\qquad
\operatorname{disc}\La^{2,1}=-2.
\]
The body's lattice $\La^{2,1}=\HypPlan\oplus[2]$ in
$(f_2,f_3,f_{-2})$-coordinates uses exactly this Gram, with
$\HypPlan=U(-1)$ on the $(f_2,f_{-2})$-block and $[2]=\langle 2\rangle$
on the $f_3$-axis.
\end{definition}

\begin{definition}[$\La_{II}^{2,1}$ as the index-four type-II sublattice of $\La^{2,1}$]
\label{def:lambda-II-21}
\textnormal{[definitional, $\alpha$]}
The type-II sublattice
\[
\La_{II}^{2,1}\;:=\;
\{m f_2+l f_3+n f_{-2}\in\La^{2,1}\mid m\equiv 0\,(2),\ n\equiv 0\,(2)\}
\;\subset\;\La^{2,1}
\]
is the index-$4$ sublattice on which the type-II Cartan matrix
$2(\Id_3-\J_3)$ is the Gram of the simple roots
$(\delta_1,\delta_2,\delta_3)$ of
Proposition~\ref{prop:type-II-cartan}.  Abstractly it is the lattice
\[
\La_{II}^{2,1}\;\simeq\;U(4)\oplus\langle 2\rangle,\qquad
\operatorname{disc}\La_{II}^{2,1}=-32,
\]
matching the body terminology of Chapter~\ref{ch:pentadic} and the
type-II sublattice convention of Chapter~\ref{ch:view2-bkm}.
\end{definition}

\begin{verification}[Rank, signature, discriminant of $\La_{II}^{2,1}$]
\label{ver:II21-invariants}
\textnormal{[computed, $\gamma$]}
$\rk\La_{II}^{2,1}=3$.  The Gram of $\La_{II}^{2,1}$ in
$(\delta_1,\delta_2,\delta_3)$ of
Proposition~\ref{prop:type-II-cartan} is the Cartan matrix
$2(\Id_3-\J_3)$ of equation~\eqref{eq:bkm-cartan-matrix}; its
determinant is
\[
\det\bigl(2(\Id_3-\J_3)\bigr)=-32.
\]
The hyperbolic-plus-rank-one block of the ambient $\La^{2,1}$ has
signature $(2,1)$; the index-$4$ sublattice $\La_{II}^{2,1}$ inherits
this signature, so
\[
\sigma(\La_{II}^{2,1})=(2,1).
\]
The discriminant group $(\La_{II}^{2,1})^*/\La_{II}^{2,1}$ has order
$32$ as required.  The lattice $\La_{II}^{2,1}$ is even because every
$\delta_i^2=2$ and the off-diagonal couplings $(\delta_i,\delta_j)=-2$
are integral.  Switching to the standard hyperbolic-plane convention
$U=((0,1),(1,0))$ on the ambient $\La^{2,1}$ presents
$\La^{2,1}\simeq U\oplus\langle -2\rangle$ and
$\La_{II}^{2,1}\simeq U(4)\oplus\langle 2\rangle$; rank, signature,
discriminant group order, automorphism group, and Cartan datum are
unchanged.
\end{verification}

\begin{remark}[Identification with the body lattice]
\label{rem:II21-vs-La21}
The body uses
$\La^{2,1}=\HypPlan\oplus[2]
=\Z f_2\oplus\Z f_3\oplus\Z f_{-2}$
with Gram in $(f_2,f_3,f_{-2})$
\[
\begin{pmatrix}0&0&-1\\ 0&2&0\\ -1&0&0\end{pmatrix},
\]
matching $G_{2,1}$ of Definition~\ref{def:lambda-21-ambient} up to
the basis swap $e\mapsto f_2$, $f\mapsto f_{-2}$, $h\mapsto f_3$.
The type-II sublattice $\La_{II}^{2,1}\subset\La^{2,1}$ of
Definition~\ref{def:lambda-II-21} is the index-$4$ sublattice on
which the simple roots $\delta_1=2f_2-f_3$, $\delta_2=2f_{-2}-f_3$,
$\delta_3=f_3$ of Proposition~\ref{prop:type-II-cartan} live; it is
the even Lorentzian root lattice of $\autcor$.
\end{remark}

\begin{proposition}[Type-II simple-root Cartan]
\label{prop:type-II-cartan}
\textnormal{[proved, $\gamma$]}
In the body coordinates fixed at the start of the Weyl-chamber
section, set
\[
\delta_1=2f_2-f_3,\qquad
\delta_2=2f_{-2}-f_3,\qquad
\delta_3=f_3.
\]
Then $\delta_i^2=2$ for $i=1,2,3$ and
\[
\bigl((\delta_i,\delta_j)\bigr)_{i,j=1,2,3}
=\begin{pmatrix} 2&-2&-2\\ -2&2&-2\\ -2&-2&2\end{pmatrix}
=2(\Id_3-\J_3),
\]
where $\J_3$ is the $3\times 3$ all-ones matrix.  Each $\delta_i$ is
type-II in the sense that $(\delta_i,\La^{2,1})=2\Z$.
\end{proposition}

\begin{proof}
Direct computation in the body Gram.  The companion script
\texttt{compute/verify\_lattice.py} at lines 92--97 verifies the matrix
$2(\Id_3-\J_3)$ from the same data; the print output records
\texttt{Type-II simple-root Gram matrix} equal to the displayed Cartan.
The type-II divisibility statement follows because each $\delta_i$
pairs against $\La^{2,1}$ with even values: $(\delta_i,f_2)$,
$(\delta_i,f_3)$, $(\delta_i,f_{-2})\in 2\Z$ for each $i$.
This is the GN root datum of \cite[Section 5.1.1, case
$(1,II,\overline0)$]{GNII}.
\end{proof}

\begin{proposition}[Weyl vector $\rho$]
\label{prop:weyl-vector-II21}
\textnormal{[proved, $\gamma$]}
The Weyl vector
\[
\rho:=\tfrac12(\delta_1+\delta_2+\delta_3)
=f_2-\tfrac12 f_3+f_{-2}\in\La^{2,1}\otimes\Q
\]
satisfies $(\rho,\delta_i)=-1$ for $i=1,2,3$ and
\[
(\rho,\rho)=-\tfrac32,\qquad \rho\in\Cone(\La^{2,1})_+.
\]
In particular $\rho$ lies in the time-like cone, so the chamber
$\Poly_{II}=\{x\mid(x,\delta_i)\le 0,\,i=1,2,3\}$ has finite hyperbolic
volume.
\end{proposition}

\begin{proof}
The simple-root condition $(\rho,\delta_i)=-(\delta_i^2)/2=-1$
follows from Kac~\cite[\S2.5, Lemma~2.5]{KAC:1} for any choice of
simple roots of square $2$; for the GN type-II datum see
\cite[(2.7)]{GNII}.  The script
\texttt{compute/verify\_lattice.py} asserts $\rho=(1,-1/2,1)$ in
$(f_2,f_3,f_{-2})$ coordinates (line~99) and
$(\rho,\delta_i)=-1$ at line~100.  Direct computation in the body
Gram of Definition~\ref{def:lambda-II-21}:
\[
(\rho,\rho)
=2(f_2,f_{-2})+\tfrac14(f_3,f_3)
=2(-1)+\tfrac14(2)=-\tfrac32.
\]
Since $(\rho,\rho)<0$, $\rho\in\Cone(\La^{2,1})_+$.
\end{proof}

\begin{proposition}[Reflection group; Nikulin's theorem]
\label{prop:nikulin-II21-reflection}
\textnormal{[proved, $\alpha$; Nikulin 1980 Theorem~1.4.1]}
Let $W^{(2)}(\La^{2,1})$ be the subgroup of $\Orth(\La^{2,1})$
generated by reflections in primitive roots of square $2$.  Then
\[
\Orth(\La^{2,1})_+=W^{(2)}(\La^{2,1}),
\]
where the subscript $+$ denotes the index-two subgroup preserving the
time-like cone $\Cone(\La^{2,1})_+$.  Equivalently, every isometry of
$\La^{2,1}$ preserving the chosen positive cone is a product of
square-$2$ reflections.
\end{proposition}

\begin{proof}
Nikulin~\cite[Theorem~1.4.1]{NIKULIN} establishes for the rank-three
hyperbolic lattice $\HypPlan\oplus[2]$ the equality
$\Orth(\La^{2,1})_+=W^{(2)}(\La^{2,1})$; see also
Gritsenko--Nikulin~\cite[Proposition 1.5]{GN}.  The body of the
manuscript invokes this at line~5408--5410.
\end{proof}

\begin{verification}[Type-II reflection orbit]
\label{ver:type-II-reflection-orbit}
\textnormal{[computed, $\gamma$]}
For each $\delta\in\{\delta_1,\delta_2,\delta_3\}$, the reflection
$s_\delta:x\mapsto x-(x,\delta)\delta$ satisfies
$s_\delta(\delta)=-\delta$ and preserves the Cartan matrix on the
$\delta$-orbit (\texttt{compute/verify\_lattice.py} lines~103--108).
Concretely, $s_{\delta_1}(\delta_1)=-\delta_1$,
$s_{\delta_1}(\delta_2)=\delta_2+2\delta_1$, etc.; the bilinear form
on $\{\delta_1,\delta_2,\delta_3\}$ is reflection-invariant by direct
matrix computation.
\end{verification}

\begin{remark}[The two type-$2$ divisibility classes]
\label{rem:type-I-type-II}
Each $\delta\in\La^{2,1}$ with $\delta^2=2$ has divisibility
$(\delta,\La^{2,1})\in\{\Z,2\Z\}$.  The two cases are recorded as
\emph{type I} ($(\delta,\La^{2,1})=\Z$) and \emph{type II}
($(\delta,\La^{2,1})=2\Z$); both classes generate index-finite normal
subgroups of $W^{(2)}(\La^{2,1})$:
\[
[W^{(2)}(\La^{2,1}):W^{(2)}(\La^{2,1}_I)]=2,\qquad
[W^{(2)}(\La^{2,1}):W^{(2)}(\La^{2,1}_{II})]=6,
\]
following \cite[Sections 4--5]{GN}.  The denominator algebra
$\autcor$ uses the type-II reflection chamber.
\end{remark}

\subsection{The K3 lattice $\La_{K3}=\mathrm{II}_{3,19}$}
\label{app:lattice-K3}

For $S$ a complex K3 surface, the integral cohomology
$H^2(S,\Z)$ with cup-product pairing is the unique even unimodular
indefinite lattice of signature $(3,19)$
(Conway--Sloane 1999~\S16.4; Nikulin 1980 Theorem~1.1.4).

\begin{definition}[$\La_{K3}=\mathrm{II}_{3,19}$]
\label{def:K3-lattice}
\textnormal{[definitional, $\alpha$]}
The K3 cohomology lattice is
\[
\La_{K3}:=H^2(S,\Z)\simeq -E_8\oplus -E_8\oplus U\oplus U\oplus U
\equiv -2E_8\oplus 3U,
\]
where $E_8$ is the rank-$8$ root lattice with the standard positive-
definite Cartan pairing (Conway--Sloane 1999~\S4.8) and $-E_8$
denotes $E_8$ with bilinear form negated.
\end{definition}

\begin{verification}[Rank, signature, discriminant of $\La_{K3}$]
\label{ver:K3-invariants}
\textnormal{[computed, $\alpha$]}
\[
\rk\La_{K3}=8+8+2+2+2=22,
\qquad
\sigma(\La_{K3})=(3,19).
\]
The three $U$-summands contribute signature $(3,3)$ and the two
$-E_8$-summands contribute signature $(0,16)$; total $(3,19)$.  The
discriminant is
\[
\operatorname{disc}\La_{K3}=
(\det -E_8)^2\cdot(\det U)^3=1\cdot(-1)^3=-1,
\]
so $|\operatorname{disc}|=1$ and the lattice is unimodular
(Conway--Sloane 1999~\S16.4, Theorem~16.4.1).  Even unimodularity
follows because each summand is even unimodular.
\end{verification}

\begin{proposition}[Aut and Hodge data on $\La_{K3}$]
\label{prop:K3-aut}
\textnormal{[proved, $\alpha$]}
The orthogonal group $\Orth(\La_{K3})$ acts transitively on the set
of primitive vectors of any fixed square (Eichler 1952;
Nikulin~\cite[Theorem 1.14.4]{Nikulin1980}).  For a polarized K3 with
period $\omega\in H^{2,0}(S)\subset\La_{K3}\otimes\C$, the
transcendental and algebraic sublattices fit into the orthogonal
splitting
\[
\La_{K3}=\NS(S)\oplus T(S),
\]
where $\NS(S)\subset H^{1,1}(S,\Z)$ is the N\'eron--Severi
(algebraic) sublattice of signature $(1,\rho-1)$ and $T(S)$ is the
transcendental lattice of signature $(2,20-\rho)$, with
$\rho=\operatorname{rk}\NS(S)$ the Picard rank
(Beauville--Bourguignon--Demazure 1985~\S6;
\cite[\S2]{Nikulin1980}).
\end{proposition}

\begin{proof}
The transitivity statement is Eichler's theorem on indefinite
unimodular lattices, sharpened by
Nikulin~\cite[Theorem~1.14.4]{Nikulin1980}: for an even unimodular
indefinite lattice $L$ of rank $\ge 3$, $\Orth(L)$ acts transitively
on primitive vectors of fixed square in each squared-class.
Decomposition by Hodge type partitions $\La_{K3}\otimes\C$ into
$H^{2,0}\oplus H^{1,1}\oplus H^{0,2}$; intersecting with $\La_{K3}$
gives the transcendental--algebraic splitting at the integral level.
\end{proof}

\subsection{The lattice-polarized K3 sublattice $\La_S$}
\label{app:lattice-Lambda-S}

A lattice polarization is a primitive embedding of an integral
lattice $\La_S$ into $\La_{K3}$ whose orthogonal complement carries
period structure.

\begin{definition}[Lattice polarization $\La_S\hookrightarrow\La_{K3}$]
\label{def:lattice-polarization}
\textnormal{[definitional, $\alpha$]}
A \emph{lattice polarization} of signature $(1,\rho-1)$ on a K3
surface $S$ is a primitive embedding
\[
\La_S\hookrightarrow\La_{K3}=\mathrm{II}_{3,19},\qquad
\sigma(\La_S)=(1,\rho-1),\quad 1\le\rho\le 19,
\]
where $\La_S\subset H^{1,1}(S,\Z)$ contains an ample class.  The
orthogonal complement
\[
\La_S^\perp\subset\La_{K3},\qquad \sigma(\La_S^\perp)=(2,20-\rho),
\]
is the lattice on which the period map of the family
$\mathcal F_{\La_S}\to\mathcal D_{\La_S^\perp}$ is defined
(Dolgachev~1996; Nikulin 1980~Proposition~1.6.1).
\end{definition}

\begin{verification}[Numerics of $\La_S$ and $\La_S^\perp$]
\label{ver:lambda-S-invariants}
\textnormal{[computed, $\alpha$]}
\[
\rk\La_S=\rho,\quad \sigma(\La_S)=(1,\rho-1),\quad
\rk\La_S^\perp=22-\rho,\quad \sigma(\La_S^\perp)=(2,20-\rho).
\]
The discriminant pairs:
\[
\operatorname{disc}\La_S\cdot\operatorname{disc}\La_S^\perp
=\operatorname{disc}\La_{K3}=-1,
\]
so $\operatorname{disc}\La_S^\perp=-\operatorname{disc}\La_S^{-1}\cdot
\det A^2$ for $A$ the change-of-basis matrix realizing the primitive
embedding (Nikulin 1980~Corollary~1.6.2).  In particular
$\La_S^\perp/\La_S^*$ and $\La_{K3}/\La_S\oplus\La_S^\perp$ are
isomorphic as discriminant groups.
\end{verification}

\begin{proposition}[Nikulin embedding criterion]
\label{prop:nikulin-embedding}
\textnormal{[proved; Nikulin 1980 Theorem~1.14.4]}
A primitive embedding $\La_S\hookrightarrow\La_{K3}$ exists if and
only if there is an isomorphism
\[
(A_{\La_S},q_{\La_S})\simeq(A_{\La_S^\perp},-q_{\La_S^\perp})
\]
of finite quadratic forms, where $A_L:=L^*/L$ is the discriminant
group and $q_L$ its quadratic form (Nikulin 1980 Definition~1.3.1
and Theorem~1.14.4; Conway--Sloane 1999~\S15.5).  The embedding is
unique up to $\Orth(\La_{K3})$ when
\[
\rk\La_S+\ell(A_{\La_S})\le 22-\rk\La_S
\]
and the discriminant form satisfies the Nikulin uniqueness criterion
(Nikulin 1980 Theorem~1.14.4).
\end{proposition}

\begin{proof}
This is the gluing-via-discriminant-form theorem
(\cite[\S1.14]{Nikulin1980}; \cite[\S15.5]{ConwaySloane}).  An
embedding into $\La_{K3}$ exists iff the discriminant form on
$\La_S^*/\La_S$ is realized by the orthogonal complement; uniqueness
is the Nikulin orbit-uniqueness statement under the corank
inequality.
\end{proof}

\begin{remark}[Mirror lattice convention]
\label{rem:mirror-lattice-convention}
The Aspinwall--Morrison mirror to a $\La_S$-polarized K3 family is the
$\check\La_S$-polarized family with
\[
\check\La_S=U\oplus(\La_S^\perp\ominus U),
\]
the splitting depending on a choice of isotropic primitive
$U\hookrightarrow\La_S^\perp$ (Dolgachev 1996~\S5; Aspinwall 1997).
The two factors $\La_S$ and $\La_S^\perp\ominus U$ play the roles of
algebraic--symplectic exchange under mirror symmetry.  No mirror
identification is asserted at this lattice level; the assertion that
$\Delta_5^{-2}$ is the discriminant of the mirror period family
(View~\textcircled{4}) is conjectural and not used in
Chapter~\ref{chap:dirac-igusa-pro-object}.
\end{remark}

\subsection{The Mukai lattice $\widetilde H(S,\Z)=\mathrm{II}_{4,20}$}
\label{app:lattice-Mukai}

For $S$ a complex K3 surface, the full integral cohomology
$H^*(S,\Z)$ with the Mukai pairing is an even unimodular lattice of
signature $(4,20)$ (Mukai 1984~Theorem~1.5; Mukai 1987).

\begin{definition}[Mukai lattice $\widetilde H(S,\Z)$]
\label{def:mukai-lattice-appA}
\textnormal{[definitional, $\alpha$]}
The Mukai lattice of $S$ is
\[
\widetilde H(S,\Z)
:=H^0(S,\Z)\oplus H^2(S,\Z)\oplus H^4(S,\Z),
\qquad
\widetilde H(S,\Z)\simeq U\oplus U\oplus U\oplus U\oplus -2E_8
\equiv 4U\oplus 2(-E_8),
\]
graded by even cohomological degree.  The Mukai pairing is
\[
\langle (r,D,s),(r',D',s')\rangle_{\mathrm{Muk}}
:=D\cdot D'-rs'-r's
\]
for $(r,D,s),(r',D',s')\in
\Z\oplus H^2(S,\Z)\oplus\Z$
(Mukai 1984~Theorem~1.5).  The pairing of an $H^0$-class with an
$H^4$-class carries a minus sign relative to the cup product, so that
the rank-$4$ component of $H^*(S,\Z)\otimes H^*(S,\Z)$ contributes
$-rs'-r's$.
\end{definition}

\begin{verification}[Rank, signature, discriminant of
$\widetilde H(S,\Z)$]
\label{ver:mukai-invariants}
\textnormal{[computed, $\alpha$]}
\[
\rk\widetilde H(S,\Z)
=1+22+1=24,
\qquad
\sigma(\widetilde H(S,\Z))=(4,20).
\]
The four $U$-summands give signature $(4,4)$ and the $-2E_8$ pieces
give signature $(0,16)$, total $(4,20)$.  Discriminant
\[
\operatorname{disc}\widetilde H(S,\Z)
=(\det U)^4\cdot(\det -E_8)^2
=(-1)^4\cdot 1=1,
\]
so $\widetilde H(S,\Z)$ is even unimodular, denoted
$\mathrm{II}_{4,20}$ (Conway--Sloane 1999~\S16.4).  The body uses the
notation $\La_S=\widetilde H(S,\Z)$ at line~4589 and~4714 of
\texttt{main.tex}; this is consistent with the convention here.
\end{verification}

\begin{proposition}[Mukai vector]
\label{prop:mukai-vector}
\textnormal{[proved, $\alpha$; Mukai 1987]}
For $F$ a coherent sheaf on $S$, the Mukai vector is
\[
v(F):=\ch(F)\sqrt{\td(S)}
=(\rk F,\,c_1(F),\,\ch_2(F)+\rk F\cdot[\mathrm{pt}])
\in\widetilde H(S,\Z),
\]
where $\sqrt{\td(S)}=(1,0,1)$ in Mukai notation on a K3 surface
(\cite[Theorem~3]{Mukai1987}).  The Mukai pairing on
$\widetilde H(S,\Z)$ recovers the Euler characteristic of $\Hom$:
\[
\langle v(F),v(G)\rangle_{\mathrm{Muk}}
=-\chi(F,G)
=-\sum_i(-1)^i\dim\Ext^i(F,G).
\]
\end{proposition}

\begin{proof}
Mukai~\cite[Theorem~3]{Mukai1987} computes $\sqrt{\td(S)}$ for K3
surfaces.  The Hirzebruch--Riemann--Roch formula gives
$\chi(F,G)=\int_S\ch(F^\vee)\ch(G)\td(S)
=-\langle v(F),v(G)\rangle_{\mathrm{Muk}}$ once the sign of the
$H^0$--$H^4$ pairing is taken into account; see Mukai 1984~\S2.
\end{proof}

\begin{remark}[Sign of the Mukai pairing on rank--$0$ components]
\label{rem:mukai-rank-0-sign}
The Mukai pairing
$\langle(r,D,s),(r',D',s')\rangle_{\mathrm{Muk}}
=D\cdot D'-rs'-r's$
restricts to the cup-product pairing on the middle $H^2$-component
\[
\langle(0,D,0),(0,D',0)\rangle_{\mathrm{Muk}}=D\cdot D',
\]
and to the negative cup-product on the rank--Euler component
\[
\langle(r,0,s),(r',0,s')\rangle_{\mathrm{Muk}}=-rs'-r's.
\]
The relative minus sign is forced by Hirzebruch--Riemann--Roch
\cite[Theorem~5.1.1]{HuybrechtsLehn}: $-\chi(F,G)$ contains a
$-\rk F\cdot\chi(\O,G)-\rk G\cdot\chi(\O,F)$ term coming from
$\sqrt{\td(S)}=(1,0,1)$.  An apparent attack route is to negate the
$H^0$--$H^4$ pairing or to drop the minus sign; this would convert
the Mukai lattice to an isotropic-rank-$2$ unimodular form not of
signature $(4,20)$, in disagreement with the Mukai 1984 classification
(\cite[Theorem~1.5]{Mukai1984}) and the unique even unimodular
type II$_{4,20}$.  The sign as stated is forced by these
requirements.
\end{remark}

\begin{proposition}[Aut of the Mukai lattice]
\label{prop:mukai-aut}
\textnormal{[proved, $\alpha$]}
$\Orth(\widetilde H(S,\Z))$ acts transitively on the set of primitive
vectors of any fixed square (Eichler 1952; Nikulin 1980
Theorem~1.14.4).  In particular, the Mukai vector $v=(1,0,1)$
(corresponding to the structure sheaf $\mathcal O_S$) and the
generator $(0,0,1)$ of $H^4(S,\Z)\cap\widetilde H(S,\Z)$ are
$\Orth(\widetilde H(S,\Z))$-conjugate to any other primitive vector
of the same square.  The derived autoequivalence group of $D^b(S)$
maps to $\Orth(\widetilde H(S,\Z))$ via the action on Mukai vectors
(Mukai 1987~\S2; Orlov 1997).
\end{proposition}

\begin{proof}
Eichler's theorem (\cite[Theorem~1.14.4]{Nikulin1980}) gives the
transitivity statement.  Mukai~\cite[\S2]{Mukai1987} constructs the
representation of derived autoequivalences on
$\widetilde H(S,\Z)$ via Fourier--Mukai kernels.
\end{proof}

\subsection{The formal dyonic charge lattice
$\Gamma_X^{\mathrm{form}}=\La_S\oplus\La_S$}
\label{app:lattice-formal-dyonic}

The formal additive dyonic lattice on $X=S\times E$ packages the
electric--magnetic charge data on which the Gram-shadow projection
$\Pi_X$ is defined.  Throughout $\La_S=\widetilde H(S,\Z)$ in the
sense of \S\ref{app:lattice-Mukai}.

\begin{definition}[$\Gamma_X^{\mathrm{form}}$ and its Gram pairing]
\label{def:gamma-X-form}
\textnormal{[definitional, $\alpha$; from
Definition~\ref{def:mukai-lattice}]}
For $X=S\times E$ with $S$ a K3 surface and $E$ an elliptic curve,
\[
\Gamma_X^{\mathrm{form}}:=\La_S\oplus\La_S,\qquad
c=(Q,P)\in\Gamma_X^{\mathrm{form}},\quad Q,P\in\La_S.
\]
The Gram pairing
$\langle\,,\,\rangle_{\Gamma}:\Gamma_X^{\mathrm{form}}
\times\Gamma_X^{\mathrm{form}}\to\Z$ is inherited from the Mukai
pairing on each factor:
\[
\langle (Q,P),(Q',P')\rangle_\Gamma
:=\langle Q,Q'\rangle_{\mathrm{Muk}}
+\langle P,P'\rangle_{\mathrm{Muk}}.
\]
The Gram-index projection
\[
\Pi_X:\Gamma_X^{\mathrm{form}}\to\Gamma_{\mathrm{gram}}=\Z^3,\qquad
(Q,P)\mapsto\Big(\tfrac{Q^2}{2},\,Q\cdot P,\,\tfrac{P^2}{2}\Big),
\]
records the genus-two Gram triple of the dyonic class.  This map is
quadratic, not additive: the obstruction is the symmetric bilinear
$2$-cocycle $B$ of Lemma~\ref{lem:gram-additivity-defect}.
\end{definition}

\begin{verification}[Rank, signature of $\Gamma_X^{\mathrm{form}}$]
\label{ver:gamma-X-form-invariants}
\textnormal{[computed, $\alpha$]}
\[
\rk\Gamma_X^{\mathrm{form}}=2\cdot\rk\La_S=48,
\qquad
\sigma(\Gamma_X^{\mathrm{form}})=2\cdot(4,20)=(8,40).
\]
The discriminant is
\[
\operatorname{disc}\Gamma_X^{\mathrm{form}}
=(\operatorname{disc}\La_S)^2=1,
\]
so $\Gamma_X^{\mathrm{form}}$ is even unimodular.  In Conway--Sloane
notation it is $\mathrm{II}_{8,40}$ (Conway--Sloane 1999~\S16.4).
\end{verification}

\begin{remark}[Relation to the K3$\times$E Mukai geometry]
\label{rem:gamma-X-vs-K3E}
Under the K\"unneth decomposition for $X=S\times E$, the integral
even cohomology decomposes as
\[
H^{\mathrm{ev}}(X,\Z)
=H^*(S,\Z)\otimes H^*(E,\Z)
=H^*(S,\Z)\otimes(\Z[1]\oplus\Z[\omega_E]),
\]
where $\omega_E$ is the fundamental class of $E$.  Every Mukai vector
on $X$ in the formal even sector decomposes as
$v_X=P\otimes 1_E+Q\otimes\omega_E$ with $Q,P\in\La_S$, and
$\Gamma_X^{\mathrm{form}}=\La_S\oplus\La_S$ is the additive lattice
on which the $(P,Q)$-grading lives
(Definition~\ref{def:gram-index-lattice}).
The Mukai pairing on $\widetilde H(X,\Z)$ pairs against the elliptic
$U$-summand $H^0(E,\Z)\oplus H^2(E,\Z)\simeq U$, giving
\[
\langle v_X,v_X'\rangle_{\widetilde H(X)}
=Q\cdot P'+Q'\cdot P,
\]
where each summand is the Mukai pairing on $\La_S$.  The resulting
genus-two Gram triple $(Q^2/2,Q\cdot P,P^2/2)$ is the Gram-index
data $\Pi_X(Q,P)$ used in the Borcherds product
(Theorem~\ref{thm:d6-d2-d0-dictionary-appendixC}).
\end{remark}

\begin{lemma}[Two orthogonal hyperbolic planes inside $\La_S$]
\label{lem:two-hyperbolic-planes-appA}
\textnormal{[proved, $\alpha$]}
$\La_S=\widetilde H(S,\Z)\simeq 4U\oplus 2(-E_8)$ contains four
mutually orthogonal copies of $U$.  In particular, two of them give
\[
U_1\oplus U_2\subset\La_S,\qquad U_i\simeq U,\ U_1\perp U_2,
\]
which is the hypothesis of the formal primitive lift
Lemma~\ref{lem:primitive-lift-every-gram-triple}: every
Gram triple $(n,l,m)\in\Z^3$ has a primitive saturated formal lift
$c=(Q,P)\in\Gamma_X^{\mathrm{form}}$ with $\Pi_X(Q,P)=(n,l,m)$.
\end{lemma}

\begin{proof}
By Definition~\ref{def:mukai-lattice}, $\La_S\simeq 4U\oplus 2(-E_8)$
contains four mutually orthogonal $U$-summands.  Selecting any two of
them gives the displayed sublattice, which is itself isomorphic to
$U\oplus U\simeq\mathrm{II}_{2,2}$ as a sublattice of $\La_S$.  The
Gram-triple lift then follows from
Lemma~\ref{lem:primitive-lift-every-gram-triple}: take
$Q=e_1+nf_1$, $P=lf_1+e_2+mf_2$ in $U_1\oplus U_2$ to obtain
$\Pi_X(Q,P)=(n,l,m)$.
\end{proof}

\subsection{The signature-$(3,2)$ lattice $\La^{3,2}$ and the
$\Sp_4\simeq\Spin(3,2)$ bridge}
\label{app:lattice-32}

The exceptional isomorphism $\PSp_4(\R)\simeq\SO_+(3,2)$ identifies
the Siegel modular variety with the fourth-orthogonal modular variety
of an integral lattice $\La^{3,2}$ of signature $(3,2)$
(\cite[Section~3]{GN}; Borcherds 1995~\S5).  The Igusa cusp form
$\Delta_5$ lives on the Siegel side; the Borcherds product expansion
lives on the orthogonal side.

\begin{definition}[$\La^{3,2}$ in the exterior-square model]
\label{def:lambda-32}
\textnormal{[definitional, $\alpha$]}
Let $\La^4=\Z e_1\oplus\Z e_2\oplus\Z e_3\oplus\Z e_4$ and consider
the exterior square $\wedge^2\La^4$ with the symmetric pairing
\[
u\wedge v\cdot u'\wedge v'=
[\det(u,v,u',v')]_{e_1\wedge e_2\wedge e_3\wedge e_4}\in\Z.
\]
Then $(\wedge^2\La^4,(\,,\,))\simeq\mathrm{II}_{3,3}$ as an even
unimodular lattice of signature $(3,3)$.  Fix
$q_0:=e_1\wedge e_3+e_2\wedge e_4\in\wedge^2\La^4$ and define
\[
\La^{3,2}:=q_0^\perp\subset\wedge^2\La^4.
\]
\end{definition}

\begin{verification}[Rank, signature, discriminant of $\La^{3,2}$]
\label{ver:lambda-32-invariants}
\textnormal{[computed, $\alpha$; verified by
\texttt{compute/verify\_lattice.py}]}
$\rk\wedge^2\La^4=\binom{4}{2}=6$.  Removing one isotropic ray from
$q_0$ gives $\rk\La^{3,2}=5$.  The \texttt{expected\_gram} array of
\texttt{compute/verify\_lattice.py} at lines~82--88 is the explicit Gram
matrix of $\La^{3,2}$ in the basis $(f_1,f_2,f_3,f_{-2},f_{-1})$:
\[
G_{3,2}=\begin{pmatrix}
0&0&0&0&-1\\
0&0&0&-1&0\\
0&0&2&0&0\\
0&-1&0&0&0\\
-1&0&0&0&0
\end{pmatrix},
\]
where $f_1=e_1\wedge e_2$, $f_2=e_2\wedge e_3$,
$f_3=e_1\wedge e_3-e_2\wedge e_4$, $f_{-2}=e_4\wedge e_1$,
$f_{-1}=e_4\wedge e_3$.  Direct computation: the matrix is
antidiagonal except for the central $h^2=2$ entry, with off-diagonal
entries $-1,-1,-1,-1$.  Expansion gives
\[
\det G_{3,2}=2\cdot(-1)^4\cdot\operatorname{sgn}(1,2,3,4\mapsto4,3,2,1)
=2\cdot 1\cdot(+1)=+2,
\]
since the permutation $(1,2,3,4)\mapsto(4,3,2,1)$ has six inversions
and even parity.  Therefore
\[
\sigma(\La^{3,2})=(3,2),\qquad
\operatorname{disc}\La^{3,2}=2.
\]
The script asserts
\texttt{gram == expected\_gram} at line~89.
\end{verification}

\begin{proposition}[$\La^{3,2}\simeq U\oplus U\oplus\langle 2\rangle$
in the manuscript convention; $\La^{3,2}\simeq U(-1)\oplus
U(-1)\oplus\langle 2\rangle$ in the standard convention]
\label{prop:lambda-32-decomposition}
\textnormal{[proved, $\alpha$]}
In the body convention of the exterior-square model,
\[
\La^{3,2}\simeq\HypPlan\oplus\HypPlan\oplus[2],
\]
with $\HypPlan=U(-1)$ and $[2]$ the rank-$1$ lattice with Gram
$(2)$.  The Gram presentation displayed in
Verification~\ref{ver:lambda-32-invariants} arises from the
$f$-basis: $(f_1,f_{-1})$ pairs as a copy of $\HypPlan$
($(f_1,f_{-1})=-1$), $(f_2,f_{-2})$ pairs as a second copy of
$\HypPlan$ ($(f_2,f_{-2})=-1$), and $f_3$ contributes
$(f_3,f_3)=2$.
\end{proposition}

\begin{proof}
The pairings $(f_1,f_{-1})=-1$, $(f_2,f_{-2})=-1$, $(f_3,f_3)=2$ are
read off from $G_{3,2}$.  The orthogonal decomposition is direct:
$\Z f_1+\Z f_{-1}\perp\Z f_2+\Z f_{-2}\perp\Z f_3$ inside
$\La^{3,2}$ because every cross-pair vanishes in $G_{3,2}$.  The
identification with $\HypPlan\oplus\HypPlan\oplus[2]$ follows
because each rank-$2$ block has Gram $((0,-1),(-1,0))=\HypPlan$
and the rank-$1$ block has Gram $(2)=[2]$.
\end{proof}

\begin{proposition}[$\PSp_4(\Z)\simeq\SO_+(\La^{3,2})$]
\label{prop:psp4-iso-so32}
\textnormal{[proved, $\alpha$; exterior-square model of $\Sp_4(\Z)$]}
The exterior-square representation
$\wedge^2:\Sp_4(\Z)/\{\pm\Id_4\}\to\SO_+(\La^{3,2})\simeq
\Orth(\La^{3,2})_+/\{\pm\Id_5\}$
is an isomorphism.  Concretely, the integral generators
\[
g_0=\begin{pmatrix}0&\Id_2\\-\Id_2&0\end{pmatrix},\quad
g_M=\begin{pmatrix}\Id_2&M\\ 0&\Id_2\end{pmatrix},\quad
g_A=\begin{pmatrix}{}^tA^{-1}&0\\ 0&A\end{pmatrix}
\quad(A\in\GL_2(\Z))
\]
of $\Sp_4(\Z)$ map under $\wedge^2$ to integral isometries of
$\La^{3,2}$, displayed as the three explicit $5\times 5$ matrices
$\wedge^2(g_0)$, $\wedge^2(g_M)$, $\wedge^2(g_A)$ in the body's
exterior-square model.  The full image is $\SO_+(\La^{3,2})$.
The Siegel space $\Hpl_2$ embeds into the type-IV upper half-space
$\Hpl^{\IV}_+\subset\Proj(\La^{3,2}\otimes\C)$ via
\[
Z=\begin{pmatrix}z_1&z_2\\ z_2&z_3\end{pmatrix}
\mapsto
{}^t(z_2^2-z_1z_3,\,z_3,\,z_2,\,z_1,\,1)z_0,
\]
and the action of $\Sp_4(\Z)$ on $\Hpl_2$ corresponds to the action
of $\SO_+(\La^{3,2})$ on $\Hpl^{\IV}_+$.
\end{proposition}

\begin{proof}
The Pfaffian pairing on $\wedge^2\La^4$ is preserved by
$\wedge^2\SL_4$.  The subgroup fixing $q_0$ is $\Sp_4(\Z)$ by
definition of the symplectic form
$B_{q_0}(x,y)e_1\wedge e_2\wedge e_3\wedge e_4=-x\wedge y\wedge q_0$.
The kernel on $q_0^\perp=\La^{3,2}$ is $\{\pm\Id_4\}$.  The
exceptional Lie-algebra isomorphism $\mathfrak{sp}_4(\R)\simeq\mathfrak{so}(3,2)$
extends to the integral form on the displayed generators; this is
the content of \cite[Section~3]{GN} and the manuscript's
exterior-square model of the automorphic group.
\end{proof}

\begin{remark}[Embedding $\La_{II}^{2,1}\subset\La^{3,2}$]
\label{rem:II21-in-32}
The primitive sublattice
\[
\La^{2,1}=\HypPlan\oplus[2]
=\Z f_2\oplus\Z f_3\oplus\Z f_{-2}\subset\La^{3,2}
\]
of the body's Weyl-chamber section is the BKM root lattice: discarding the
$(f_1,f_{-1})$-block of $\La^{3,2}\simeq\HypPlan\oplus\HypPlan\oplus
[2]$ leaves $\HypPlan\oplus[2]\simeq\La^{2,1}$.  The type-II
sublattice $\La_{II}^{2,1}=\{mf_2+lf_3+nf_{-2}\mid
m,n\equiv 0\!\pmod 2\}$ is the lattice of even integral type-II
classes, on which the simple-root Cartan
\[
\bigl((\delta_i,\delta_j)\bigr)
=\begin{pmatrix}2&-2&-2\\-2&2&-2\\-2&-2&2\end{pmatrix}
\]
of Proposition~\ref{prop:type-II-cartan} is realized via
$\delta_1=2f_2-f_3$, $\delta_2=2f_{-2}-f_3$, $\delta_3=f_3$.  Every
isometry of $\La^{2,1}$ extends to an isometry of $\La^{3,2}$ by
the identity on $(\La^{2,1})^\perp\subset\La^{3,2}$.
\end{remark}

\subsection{Cross-references and consistency check}
\label{app:lattice-cross-refs}

The lattices recorded above sit in the manuscript as follows.

\begin{itemize}
\item $U=\mathrm{II}_{1,1}$: every $U$-summand of
$\La^{3,2}$, $\La_{K3}$, and $\widetilde H(S,\Z)$.  The relation
$\HypPlan=U(-1)$ packages the body's negated convention.
\item $\La_{II}^{2,1}\simeq U(-1)\oplus\langle 2\rangle\simeq
U\oplus\langle-2\rangle$ (sign flip on the rank-$2$ hyperbolic block):
the BKM root lattice on which the type-II simple-root Cartan
$2(\Id_3-\J_3)$ and the Weyl vector
$\rho=\frac12(\delta_1+\delta_2+\delta_3)$ are realized.  The
denominator algebra
$\autcor=\mathfrak g_{\Delta_5}$ is graded by $\La_{II}^{2,1}$
(\cite[Theorem~2.1]{GNII}).
\item $\La_{K3}=\mathrm{II}_{3,19}$: K3 cohomology with
intersection pairing.  Polarized K3 sublattices $\La_S$ embed
primitively into $\La_{K3}$.
\item $\La_S$: signature $(1,\rho-1)$ for $\rho\le 19$ (the
algebraic Picard rank).  Used in
Definition~\ref{def:mukai-lattice} as the K3 factor of
$\Gamma_X^{\mathrm{form}}$ and in the lattice-polarization setup of
the body.
\item $\widetilde H(S,\Z)=\mathrm{II}_{4,20}$: the Mukai lattice on
which the D6--D2--D0 Mukai vectors $v_X(\mathcal I_Y)$ live; the
Gram-index map $\Pi_X$ is computed via this pairing
(Theorem~\ref{thm:d6-d2-d0-dictionary-appendixC}).
\item $\Gamma_X^{\mathrm{form}}=\La_S\oplus\La_S$: the formal
additive dyonic charge lattice of $X=S\times E$, on which the
quadratic Gram-shadow projection $\Pi_X$ and the symmetric bilinear
$2$-cocycle $B$ are defined
(Definition~\ref{def:mukai-lattice}).
\item $\La^{3,2}\simeq U(-1)\oplus U(-1)\oplus\langle 2\rangle$:
the signature-$(3,2)$ orthogonal lattice that bridges
$\PSp_4(\Z)$ to $\SO_+(3,2)$ via the exterior-square
representation; on this lattice the Borcherds product expansion of
$\Delta_5$ is written.
\end{itemize}

\begin{verification}[End-to-end computation check]
\label{ver:lattice-script-check}
\textnormal{[computed, $\gamma$]}
The script \texttt{compute/verify\_lattice.py} verifies, in one
self-contained $\Q$-arithmetic computation:
(i) the rank-$5$ Gram matrix $G_{3,2}$ in the
$(f_1,f_2,f_3,f_{-2},f_{-1})$-basis (lines~82--89);
(ii) the orthogonality of the symplectic invariant
$q_0=e_1\wedge e_3+e_2\wedge e_4$ to the basis (line~90);
(iii) the type-II simple-root Cartan
$2(\Id_3-\J_3)$ at lines~92--97;
(iv) the Weyl vector $\rho=(1,-1/2,1)$ in $(f_2,f_3,f_{-2})$
coordinates and the $(\rho,\delta_i)=-1$ relation
(lines~99--101);
(v) the type-II reflection action and Cartan invariance under each
$s_{\delta_i}$ (lines~103--108).
The \texttt{All lattice checks passed.} line printed at the end of the
script's output confirms agreement of the Gram, Cartan, and
Weyl-vector data of \S\ref{app:lattice-II21} and \S\ref{app:lattice-32}
with the $\Q$-arithmetic computation.
\end{verification}

\begin{remark}[Discriminant convention; sign of $\La^{3,2}$]
\label{rem:discriminant-conventions}
In the body convention $\HypPlan\oplus\HypPlan\oplus[2]$, both
hyperbolic-plane summands have Gram determinant $-1$ and the $[2]$
summand contributes $+2$, so
$\operatorname{disc}\La^{3,2}=(-1)^2\cdot 2=+2$.  In the standard
convention $U\oplus U\oplus\langle-2\rangle$, the two $U$-summands
contribute determinant $-1$ each and $\langle-2\rangle$ contributes
$-2$, so $\operatorname{disc}\La^{3,2}=(-1)^2\cdot(-2)=-2$.  The two
presentations differ by an overall sign on the rank-$2$ hyperbolic
blocks; both carry an order-$2$ discriminant group with single
nontrivial generator $f_3/2\in(\La^{3,2})^*/\La^{3,2}$.  The same
analysis applies to $\La_{II}^{2,1}$: discriminant $-2$ (body) or
$+2$ (standard), discriminant group $\Z/2$ generated by $h/2$.  In
every case the order of the discriminant group is $2$ and the
type-II divisibility statement
$(\delta,\La^{2,1})\in\{\Z,2\Z\}$ depends on the order, not the sign
of the determinant.
\end{remark}

\begin{remark}[Cited primary literature]
\label{rem:lattice-primary-lit}
The primary references behind every load-bearing lattice claim above
are: Conway--Sloane 1999~\S2.5,~\S4.8,~\S15.5,~\S16.4 (classification
of even unimodular indefinite lattices, $E_8$ root system,
discriminant-form gluing, K3 lattice as $\mathrm{II}_{3,19}$);
Nikulin 1980~Theorem~1.1.2,~Theorem~1.4.1 (= \cite{NIKULIN}
Theorem~1.4.1 in the K3-paper version),
Theorem~1.6.1,~Theorem~1.14.4 (signature classification, reflection
group of square-$2$ vectors, primitive embeddings, discriminant-form
uniqueness); Mukai 1984~Theorem~1.5 (Mukai pairing, signature
$(4,20)$), Mukai 1987~Theorem~3 ($\sqrt{\td(S)}=(1,0,1)$ on K3);
Gritsenko--Nikulin 1998 Part~II~Theorem~2.1, Section~5.1.1, case
$(t=1,II,\overline 0)$ (BKM root datum for $\autcor$, Weyl vector,
type-II simple-root Cartan).
\end{remark}

\section{Borcherds-product expansion: \texorpdfstring{$\phi_{0,1}$}{phi01}, the
\texorpdfstring{$\Delta_5$}{Delta5} infinite product, and the
Gritsenko--Cl\'ery eight-row catalogue}
\label{app:borcherds-product-expansion}

The Fourier expansion of the
weight-zero index-one weak Jacobi form $\phi_{0,1}$, the
Borcherds--Gritsenko--Nikulin product expansion of the Igusa cusp form
$\Delta_5$ in the type-II cusp chamber, the theta-constant normalization
constant $64=2^{6}$, the doubling identity that turns the monic Borcherds
product into the denominator of the Lorentzian Borcherds--Kac--Moody
algebra $\autcor$, and the full eight-row Gritsenko--Cl\'ery catalogue of
diagonal-divisor Siegel modular forms are the automorphic data at the
type-II cusp.  Every displayed numerical Fourier
coefficient is produced by exact rational expansion of
the theta quotient; every identity is
attributed to author, year, and theorem of the primary literature.
Orientation data, Pfaffian wall normal forms, and compact Hall input do
not enter these automorphic identities; they live entirely on the
automorphic / target side, prior to the geometric content of (D0),
(O1), (O1)\(^+\), and (O2).

Three names for $\Delta_5$ appear and are kept distinct: the
theta-constant automorphic section $\mathcal D_X=\Delta_5$ (View \textcircled{\scriptsize 1}); the
Borcherds--Kac--Moody denominator
$\operatorname{den}(\autcor)=64^{-1}\Delta_5(2Z)$
(View \textcircled{\scriptsize 2}); and the protected Pfaffian
$\operatorname{Pf}_{\mathrm{prot}}(\mathfrak D_X^{\mathrm{DI}})$ of the
pro-object $\mathfrak D_X^{\mathrm{DI}}$ (View \textcircled{\scriptsize 3},
after the retained D0-HN datum, retained quotient-orientation datum,
chosen Weyl lifts, retained \((O2_{\mathrm{atlas}})\) wall datum, and
finite geometric Pfaffian datum \((P_{\mathrm{fin}})\) of
Appendix~\ref{app:obstruction-discharges}).  The first and third
sections coincide as weight-five automorphic sections; the second
coincides with their Borcherds denominator after the normalization
\(D_5=64^{-1}\Delta_5\) and the pullback \(Z\mapsto2Z\).  The
orientation character that distinguishes
$\Delta_5$ from $-\Delta_5$ enters only at View \textcircled{\scriptsize 3}.

\subsection*{B.1 The weak Jacobi form \texorpdfstring{$\phi_{0,1}$}{phi01}}

The K3 elliptic genus is the modular input.  Eichler--Zagier
\cite[\S2.2 and Theorem~9.3]{EZ:1} fix the weight-zero index-one
weak-Jacobi normalization
\[
\phi_{0,1}(\tau,z)\in J_{0,1}^{\mathrm w},
\qquad
\phi_{0,1}(\tau,z)=\sum_{n\ge0,\;l\in\Z}f(n,l)\,q^n r^l,
\qquad
q=e^{2\pi i\tau},\ r=e^{2\pi iz}.
\]
The relation to the K3 right-moving elliptic genus is
$Z_{\mathrm{K3}}(\tau,z)=2\phi_{0,1}(\tau,z)$ \cite[\S2.2]{EZ:1};
the half is the Borcherds-lift representative for the Borcherds product.
The cusp normalization is
\[
\phi_{0,1}(\tau,z)=r^{-1}+10+r+O(q).
\]
The Jacobi weak-form support is
$f(n,l)=0$ whenever $4n-l^{2}<-1$, with isolated extreme entries
$f(0,\pm1)=1$ on the discriminant hyperbola $4n-l^{2}=-1$.
\emph{(proved; \cite[\S2.2, Theorem~9.3]{EZ:1};
\cite[Example~2.4]{GNII}.)}

\paragraph{Verified Fourier table.}
The exact coefficients $f(n,l)$ for $n\le2$ are the rational integers
computed by the theta-quotient identity
\(\phi_{0,1}=4\sum_{i=2,3,4}\theta_i(\tau,z)^2/\theta_i(\tau,0)^2\) of
\cite[\S2.2]{EZ:1} by rational arithmetic:
\[
\begin{array}{c|ccccccc}
& l=-3 & l=-2 & l=-1 & l=0 & l=1 & l=2 & l=3 \\ \hline
n=0 & 0 & 0 & 1 & 10 & 1 & 0 & 0 \\
n=1 & 0 & 10 & -64 & 108 & -64 & 10 & 0 \\
n=2 & 1 & 108 & -513 & 808 & -513 & 108 & 1
\end{array}
\]
\emph{(computed; cross-checked with \cite[Example~2.4]{GNII} and
\cite[Theorem~4.1]{GN}.)}
The finite table
\(\texttt{certificates/jacobi/k3e\_first\_relation\_window\_coefficients}\)
records this leading normalization together with the checked evenness
and discriminant-dependence rows used in the first relation-closed
window.
The negative entries reflect the index-one Jacobi-form theta-decomposition
$\phi_{0,1}=\theta_{1,0}h_0+\theta_{1,1}h_1$ with theta-coefficient
expansions $h_0=2(45 E_4 E_{4,1}-50 E_6 E_{6,1})/\eta^{24}+\cdots$ in
\cite[\S2.2]{EZ:1}; the integer values are not target signs but signed
multiplicities of the Borcherds source.

\paragraph{Fourier identities.}
\[
\begin{aligned}
f(0,0)&=10,&\qquad f(0,\pm1)&=1,&\qquad f(0,l)&=0\ \ (|l|\ge2),\\
f(1,0)&=108,&\qquad f(1,\pm1)&=-64,&\qquad f(1,\pm2)&=10,\\
f(2,0)&=808,&\qquad f(2,\pm1)&=-513,&\qquad f(2,\pm2)&=108,\\
&&\qquad f(2,\pm3)&=1,&\qquad f(1,\pm3)&=0.
\end{aligned}
\]
The diagonal entry $f(0,0)=10$ controls the Borcherds-lift weight
\(\operatorname{wt}(\Delta_5)=f(0,0)/2=5\); the entries $f(0,\pm1)=1$
control the simple type-II divisor orders
$d_{\delta_i}^{\mathrm{aut}}=1$; the entry $f(1,1)=-64$ controls the
multiplicity of the first timelike imaginary root, computed and used
below.

\subsection*{B.2 The cusp chamber and the effective semigroup}

Set $Z=\bigl(\begin{smallmatrix}z_1&z_2\\z_2&z_3\end{smallmatrix}\bigr)
\in\Hpl_2$, $q=e^{2\pi iz_1}$, $r=e^{2\pi iz_2}$, $s=e^{2\pi iz_3}$.
The standard type-II cusp polarization is the ordered limit
\[
0<|s|\ll|q|\ll|r|^{-1}\ll1,
\]
which determines the effective semigroup
\[
\Gamma_{\mathrm{eff}}
=\{(n,l,m)\in\Z^3\mid m>0,\ n\ge0\}
\cup\{(n,l,0)\mid n>0\}
\cup\{(0,l,0)\mid l<0\}.
\]
Inside the Gram-index lattice $\Gamma_{\mathrm{ind}}=\Z^3$, the
\emph{active support}
\[
\Gamma_{\mathrm{act}}
:=\{(n,l,m)\in\Gamma_{\mathrm{eff}}\mid f(nm,l)\ne0\}
\]
is the set of triples that contribute a visible product factor; the
remaining triples in $\Gamma_{\mathrm{eff}}$ have zero exponent and fix
ordering only.  The signed Gram pairing
\((\alpha(n,l,m),\alpha(n,l,m))=-2(4nm-l^{2})\)
under
\[
\alpha:(n,l,m)\longmapsto 2nf_2-lf_3+2mf_{-2}\in\La^{2,1}_{II}
\]
restricts to the Lorentzian root lattice of $\autcor$.
The finite packet
\(\texttt{certificates/lattice/product\_chamber\_root\_map}\)
gives a computed check of the three sector definitions of
\(\Gamma_{\mathrm{eff}}\), their pairwise closure under addition, the
active rows selected by \(f(nm,l)\ne0\), the inactive zero-exponent rows,
and the identity
\[
\alpha(n,l,m)=n\delta_1+m\delta_2+(n+m-l)\delta_3
\]
used in the type-II chamber.  This is target-side lattice arithmetic; it
does not construct a compact \(K3\times E\) source, a Pfaffian
orientation, an \(O2\) wall atlas, a mirror discriminant, or a protected
trace.

\begin{proposition}[Images of the three simple type-II walls]
\label{prop:type-ii-wall-images}
\textnormal{[computed]}
Let
\[
\gamma_{\mathrm{wall}}:
\{\delta_1,\delta_2,\delta_3\}\longrightarrow \Gamma_{\mathrm{eff}}
\subset\Z^3
\]
be the inverse of the product-chamber root map
\[
\alpha(n,l,m)
=n\delta_1+m\delta_2+(n+m-l)\delta_3
\]
on the three simple type-II roots.  Then
\[
\gamma_{\mathrm{wall}}(\delta_1)=(1,1,0),\qquad
\gamma_{\mathrm{wall}}(\delta_2)=(0,1,1),\qquad
\gamma_{\mathrm{wall}}(\delta_3)=(0,-1,0).
\]
Equivalently, under
\[
\alpha(n,l,m)=2nf_2-lf_3+2mf_{-2},
\]
these triples map respectively to
\[
2f_2-f_3,\qquad 2f_{-2}-f_3,\qquad f_3 .
\]
This is a target-side wall-image computation.  It does not construct
compact \(K3\times E\) representatives, retained type-II wall objects,
wall charts, Pfaffian signs, protected traces, or a Hall carrier.
\end{proposition}

\begin{proof}
For a triple \((n,l,m)\), the coefficient of
\((\delta_1,\delta_2,\delta_3)\) in the type-II chamber is
\[
(n,\ m,\ n+m-l).
\]
For \(\delta_1\) this coefficient vector must be \((1,0,0)\), hence
\[
n=1,\qquad m=0,\qquad 1-l=0,
\]
so \(l=1\) and \(\gamma_{\mathrm{wall}}(\delta_1)=(1,1,0)\).
For \(\delta_2\) the coefficient vector is \((0,1,0)\), hence
\[
n=0,\qquad m=1,\qquad 1-l=0,
\]
so \(l=1\) and \(\gamma_{\mathrm{wall}}(\delta_2)=(0,1,1)\).
For \(\delta_3\) the coefficient vector is \((0,0,1)\), hence
\[
n=0,\qquad m=0,\qquad -l=1,
\]
so \(l=-1\) and \(\gamma_{\mathrm{wall}}(\delta_3)=(0,-1,0)\).
Substituting the three triples into
\(\alpha(n,l,m)=2nf_2-lf_3+2mf_{-2}\) gives
\[
2f_2-f_3,\qquad -f_3+2f_{-2},\qquad f_3,
\]
which are precisely \(\delta_1,\delta_2,\delta_3\).
\end{proof}

The finite packet
\[
\texttt{certificates/lattice/type\_ii\_wall\_images}
\]
verifies
\(\texttt{TYPE\_II\_WALL\_IMAGES\_VERIFIED}\): the three inverse
images above have zero root-map defect and lie in the product chamber
sectors.  The packet records compact source representatives, retained
wall objects, wall charts, Pfaffian signs, protected traces, and Hall
carrier data as excluded scalar substitutes.

\subsection*{B.3 The Borcherds product formula for
\texorpdfstring{$\Delta_5$}{Delta5}}

\begin{theorem}[Borcherds--Gritsenko--Nikulin product expansion]
\label{thm:borcherds-product-delta5}
On the type-II cusp polarization $0<|s|\ll|q|\ll|r|^{-1}\ll1$, the
weight-five Igusa cusp form admits the absolutely convergent infinite
product
\[
\Delta_5(Z)
=64\,q^{1/2}r^{1/2}s^{1/2}
\!\!\prod_{(n,l,m)\in\Gamma_{\mathrm{eff}}}\!\!
(1-q^n r^l s^m)^{f(nm,l)},
\]
with theta-leading coefficient
\[
[q^{1/2}r^{1/2}s^{1/2}]\Delta_5\;=\;64\;=\;2^{6}.
\]
The monic product is $D_5:=64^{-1}\Delta_5$; the constant $64$ is the
theta-constant Igusa normalization, not a separately chosen scalar.
\emph{(proved; \cite[Theorem~10.1]{Borcherds95},
\cite[Theorem~2.1, Example~2.4, equations (2.15)--(2.16)]{GNII},
\cite[Theorem~4.1]{GN}.)}
\end{theorem}

\begin{proof}
The general Borcherds singular theta-lift sends a weakly holomorphic
vector-valued modular form of negative weight for the Weil
representation of a lattice of signature \((2,n)\) to a meromorphic
orthogonal automorphic form realized as an infinite product
\cite[Theorem~10.1]{Borcherds95}; the regularisation and convergence on
divisors with singularities in higher rank is in
\cite[\S\S6--13]{B}.  The genus-two specialization at the
Siegel cusp uses the exterior-square isogeny
$\PSp_4(\R)\simeq\SO(3,2)_+$ and the
Eichler--Zagier dictionary $J^{\mathrm w}_{0,1}\leftrightarrow$
vector-valued Mp$_{2}$-modular forms \cite[\S5]{EZ:1};
the ordered limit $0<|s|\ll|q|\ll|r|^{-1}\ll1$ selects the type-II cusp
\cite[\S2, Example~2.4]{GNII}.  Gritsenko--Nikulin \cite[(2.15)--(2.16)]{GNII}
write the resulting product in the unnormalized form
\[
D_5(Z)
=q^{1/2}r^{1/2}s^{1/2}
\prod_{(n,l,m)\in\Gamma_{\mathrm{eff}}}(1-q^n r^l s^m)^{f(nm,l)},
\]
and identify $D_5$ with $64^{-1}$ times the theta-constant cusp form
$\Delta_5$ of Igusa \cite[Theorem~4.1]{GN}, \cite{Igusa1962},
\cite{Igusa1964II}; the leading constant $2^{6}=64$ is the explicit
theta-product coefficient.
\end{proof}

\paragraph{The constant $2^{6}$ as a theta-product leading coefficient.}
Igusa's classical product formula displays $\Delta_5$ as the product of
the ten even theta constants of genus two,
\[
\Delta_5(Z)
=\prod_{\substack{a,b\in(\Z/2\Z)^2\\{}^{t}\!ab\equiv0\pmod2}}
\nu_{a,b}(Z),
\qquad
\nu_{a,b}(Z)
=\sum_{l\in\Z^{2}}\!\!\exp\!\bigl(\pi i(Z[l+\tfrac12 a]+{}^{t}\!bl)\bigr),
\]
\cite{Igusa1962}, \cite[Theorem~4.1]{GN}, \cite[\S6]{FREITAG:1}.
In the type-II cusp polarization, with $z_2$ the off-diagonal modulus,
the four characteristics $a=(0,0)$ start with $1$, the two
characteristics $a=(1,0)$ start with $2q^{1/8}$, the two
characteristics $a=(0,1)$ start with $2s^{1/8}$, and the two
characteristics $a=(1,1)$ start with $2q^{1/8}r^{1/4}s^{1/8}$; the
total leading monomial is $q^{1/2}r^{1/2}s^{1/2}$ with constant
$1\cdot 2^{2}\cdot 2^{2}\cdot 2^{2}=2^{6}=64$
\emph{(proved; \cite{Igusa1962}, \cite[Theorem~4.1]{GN}.)}
This is not a normalization choice: the constant $64$ is read off the
theta product, with no scalar freedom.  Replacing $\Delta_5$ by
$c\Delta_5$, $c\in\C^{\times}$, would change neither the Maass character
\(\nu_{\Delta_5}\) nor the divisor orders
\(d_{\delta}^{\mathrm{aut}}\) of \cite[Theorem~4.1]{GN}.

\subsection*{B.4 Sample exponent verifications}

The exponent $f(nm,l)$ in the product formula is read directly from the
Jacobi-coefficient table of B.1.  The following sample is verified
exactly by the displayed coefficient table:

\begin{center}
\small
\begin{tabular}{c|cccccc|c}
\hline
factor $(n,l,m)$ & \multicolumn{6}{c|}{Lorentzian Gram label
$\alpha=2nf_{2}-lf_{3}+2mf_{-2}$} & exponent $f(nm,l)$ \\
\hline
$(1,1,0)$ & \multicolumn{6}{c|}{$\delta_{1}=2f_{2}-f_{3}$} & $f(0,1)=1$ \\
$(0,1,1)$ & \multicolumn{6}{c|}{$\delta_{2}=2f_{-2}-f_{3}$} & $f(0,1)=1$ \\
$(0,-1,0)$ & \multicolumn{6}{c|}{$\delta_{3}=f_{3}$} & $f(0,-1)=1$ \\
$(0,0,m)\ (m\ge1)$ & \multicolumn{6}{c|}{cusp ray
$2mf_{-2}$} & $f(0,0)=10$ \\
$(1,0,1)$ & \multicolumn{6}{c|}{$2f_{2}+2f_{-2}$ (timelike,
$(\alpha,\alpha)=-8$)} & $f(1,0)=108$ \\
$(1,1,1)$ & \multicolumn{6}{c|}{$2f_{2}-f_{3}+2f_{-2}$ (timelike,
$(\alpha,\alpha)=-6$)} & $f(1,1)=-64$ \\
$(1,2,1)$ & \multicolumn{6}{c|}{$2f_{2}-2f_{3}+2f_{-2}$ (isotropic,
$(\alpha,\alpha)=0$)} & $f(1,2)=10$ \\
$(1,3,1)$ & \multicolumn{6}{c|}{$2f_{2}-3f_{3}+2f_{-2}$
(spacelike, $(\alpha,\alpha)=10$)} & $f(1,3)=0$ \\
$(2,1,1)$ & \multicolumn{6}{c|}{$4f_{2}-f_{3}+2f_{-2}$
($(\alpha,\alpha)=-14$)} & $f(2,1)=-513$ \\
$(1,2,2)$ & \multicolumn{6}{c|}{$2f_{2}-2f_{3}+4f_{-2}$
($(\alpha,\alpha)=-12$)} & $f(2,2)=108$ \\
\hline
\end{tabular}
\end{center}

The triple $(1,1,1)$ gives the first timelike imaginary root: its
exponent $f(1,1)=-64$ is the signed superdimension of the root space
$\mathfrak g_{\Delta_5,\alpha(1,1,1)}$ on the active support, and is the
content of the lattice cross-check
$f(1,1)=29-93=-64$ from the
free-Lie/imaginary-simple split computed by
the Gritsenko--Nikulin target presentation.  The triple $(1,3,1)$ has
exponent zero because $f(1,3)=0$: this triple lies in
$\Gamma_{\mathrm{eff}}\setminus\Gamma_{\mathrm{act}}$ and contributes
ordering, not a visible factor.  These are signed superdimensions, not
positive dimensions; negative entries record an odd excess in the
Borcherds source, and only after the (O1), (O1)\(^+\), (O2) orientation
data become available do they translate to operator parities.

\paragraph{Theta-leading coefficient.}
The coefficient \(64\) is the coefficient of the leading monomial of
\(\Delta_5\):
\[
\Delta_5(Z)=64\,q^{1/2}r^{1/2}s^{1/2}\bigl(1+\cdots\bigr).
\]
After the leading monomial is divided out, the first \(qrs\)-coefficient
of \(D_5=64^{-1}\Delta_5\) is \(93\), not \(64\).  The factor
\((1-qrs)^{-64}\) contributes \(64\,qrs\), and the remaining
root factors contribute the other \(29\,qrs\); this is the identity
\[
[qrs]\,D_5/(qrs)^{1/2}=93,\qquad 29-93=-64=f(1,1).
\]
\emph{(computed by exact rational expansion of the Borcherds product.)}

\subsection*{B.5 The doubling identity
\texorpdfstring{$\operatorname{den}(\autcor)=64^{-1}\Delta_5(2Z)$}
{den(g\_Delta5)=64\^{}-1 Delta5(2Z)}}

\begin{lemma}[Two exponential normalizations]
\label{lem:appB-doubling}
The product chamber convention writes monomials as $q^{n}r^{l}s^{m}$
with $q^{n}r^{l}s^{m}=\exp(-\pi i(\alpha(n,l,m),z))$, while the standard
denominator-side convention of a Borcherds--Kac--Moody algebra writes
formal characters as $\exp(-2\pi i(\alpha,z))$.  Replacing $Z$ by $2Z$
converts the first character to the second, so
\[
\operatorname{den}(\autcor)
\;=\;D_{5}(2Z)
\;=\;64^{-1}\Delta_{5}(2Z).
\]
\emph{(proved; \cite[Proposition~3.1]{GN}, \cite[Section~5.1.1]{GNII}
together with Theorem~\ref{thm:borcherds-product-delta5}.)}
\end{lemma}

\begin{proof}
The product side of Theorem~\ref{thm:borcherds-product-delta5} gives
\(D_{5}(Z)=q^{1/2}r^{1/2}s^{1/2}\prod(1-q^{n}r^{l}s^{m})^{f(nm,l)}\)
with leading monomial $q^{1/2}r^{1/2}s^{1/2}$.  Doubling
$Z\mapsto 2Z$ replaces $q^{n}r^{l}s^{m}$ by $q^{2n}r^{2l}s^{2m}$ and
hence $\exp(-\pi i(\alpha,z))$ by $\exp(-2\pi i(\alpha,z))$.  The
Lorentzian Borcherds--Kac--Moody algebra $\autcor$ of
Gritsenko--Nikulin \cite[Sections~3--4, Proposition~3.1]{GN} carries the
denominator identity
\[
\operatorname{den}(\autcor)
=\sum_{w\in W^{(2)}(\La^{2,1}_{II})}\det(w)\,
\Bigl(e^{-2\pi i(w\rho,z)}
-\sum_{a}m(a)\,e^{-2\pi i(w(\rho+a),z)}\Bigr)
\]
on the imaginary-positive cone, where \(m(a)\) is determined by
the negative-norm additive coefficients and by
\(1-\sum_{t\ge1}m(ta)q^t=\prod_{t\ge1}(1-q^t)^{\tau(ta)}\) on isotropic
rays.  The right-hand side is equal to
$D_{5}(2Z)$ in the doubled exponential.  Equivalently
$\operatorname{den}(\autcor)=64^{-1}\Delta_{5}(2Z)$.
\end{proof}

The factor $64$ in $\operatorname{den}(\autcor)=64^{-1}\Delta_{5}(2Z)$
is the same theta-constant $2^{6}$ from the leading-coefficient
calculation, transported through
$Z\mapsto 2Z$; it is not a separate normalization, an integration
constant, or a free scalar.  The three identifications
\(\mathcal D_{X}=\Delta_{5}\),
\(\operatorname{den}(\autcor)=64^{-1}\Delta_{5}(2Z)\),
\(\operatorname{Pf}_{\mathrm{prot}}(\mathfrak D_{X}^{DI})\) are kept
distinct; the doubling lemma supplies the denominator-character
normalization.

\subsection*{B.6 The eight-row Gritsenko--Cl\'ery catalogue}

Gritsenko and Cl\'ery \cite[Theorem~1.4]{GC} classify all Siegel modular
forms of genus two, with character or multiplier system on a
Hecke-type congruence subgroup $\Gamma_{t}(N)\subset\operatorname{Sp}(4,\Q)$,
whose divisor is supported on the $\Gamma_{t}(N)$-translates of the
diagonal Humbert surface
$\mathcal H_{\mathrm{diag}}=\{Z=\operatorname{diag}(a,b):a,b\in\mathbb H_{1}\}$
with multiplicity one.  The classification produces exactly eight rows.
\emph{(proved; \cite[Theorem~1.4 and Theorem~3.1]{GC}.)}

The rows are:

\begin{center}
\renewcommand{\arraystretch}{1.16}
\small
\begin{tabular}{c|c|c|c|c|c|c}
\hline
$j$ & $F_{j}$ & $(t,N)$ & $\Gamma_{j}$ &
$\operatorname{wt}(F_{j})$ & multiplier $\nu_{j}$ & $c_{j}(0)$ \\ \hline
$1$ & $\Delta_{5}$ & $(1,1)$ & $\Gamma_{1}=\operatorname{Sp}(4,\Z)$ & $5$ &
$v_{\eta}^{12}\times v_{H}$ & $10$ \\
$2$ & $\Delta_{2}$ & $(2,1)$ & $\Gamma_{2}^{+}$ & $2$ &
$v_{\eta}^{6}\times v_{H}$ & $4$ \\
$3$ & $\Delta_{1}$ & $(3,1)$ & $\Gamma_{3}^{+}$ & $1$ &
$v_{\eta}^{4}\times v_{H}$ & $2$ \\
$4$ & $\Delta_{1/2}$ & $(4,1)$ & $\Gamma_{4}^{+}$ & $\tfrac12$ &
$v_{\eta}^{3}\times v_{H}$ (degree 8) & $1$ \\
$5$ & $\nabla_{3}$ & $(1,2)$ & $\Gamma_{0}^{(2)}(2)$ & $3$ &
$\chi_{2}^{(2)}\times v_{H}$ & $6$ \\
$6$ & $\nabla_{2}$ & $(1,3)$ & $\Gamma_{0}^{(2)}(3)$ & $2$ &
$\chi_{2}^{(3)}\times v_{H}$ & $4$ \\
$7$ & $\nabla_{3/2}$ & $(1,4)$ & $\Gamma_{0}^{(2)}(4)$ & $\tfrac32$ &
$\chi_{4}^{(4)}\times v_{H}$ (order 4) & $3$ \\
$8$ & $Q_{1}$ & $(2,2)$ & $\Gamma_{2}(2)$ & $1$ &
$\chi_{4}^{(2)}\times v_{H}$ & $2$ \\
\hline
\end{tabular}
\end{center}

In the displayed \(F_j\)-order the weights and constant Fourier
coefficients are
\[
(5,2,1,\tfrac12,3,2,\tfrac32,1),\qquad
(10,4,2,1,6,4,3,2).
\]
In the input-pair order
\[
(1,1),(2,1),(1,2),(3,1),(1,3),(4,1),(1,4),(2,2)
\]
the same data read
\[
(5,2,3,1,2,\tfrac12,\tfrac32,1),\qquad
(10,4,6,2,4,1,3,2).
\]
Each row satisfies the Borcherds--Gritsenko universal weight identity
\[
\operatorname{wt}(F_{j})\;=\;\frac{c_{j}(0)}{2}
\;=\;\kappa_{\mathrm{BKM}}(F_{j}),
\qquad j=1,\ldots,8 .
\]
\emph{(proved; \cite{Borcherds95},
\cite[Theorem~3.1]{GC}; the $\kappa_{\mathrm{BKM}}$ relabelling is the
universal Borcherds-weight identity used in the Calabi--Yau
quantum-groups $\kappa$-stratification.)}

The half-integer weights $\tfrac12$ (row $4$) and $\tfrac32$ (row $7$)
are realized by multiplier systems $v_{\eta}^{k}\times v_{H}$ on the
ambient paramodular cover, not by passing to a metaplectic double
cover.  The multiplier $v_{\eta}$ is the Dedekind eta multiplier on
$\operatorname{SL}_{2}(\Z)$; the multiplier $v_{H}$ is the Heisenberg
multiplier on the Heisenberg lift of $\operatorname{Sp}(4,\Z)$; and
$v_{\eta}^{3}\times v_{H}$ has order $8$ on the row-$4$ cover, while
$\chi_{4}^{(4)}\times v_{H}$ has order $4$ on the row-$7$ cover.
\emph{(proved; \cite[Theorem~3.1]{GC}.)}

\paragraph{Denominator content of the rows.}
Rows $1$--$4$ carry Gritsenko--Nikulin denominator theorems
\cite[Theorem~2.1, equations (2.11), (2.20), (2.21)]{GNII}; in the
$e_{t}^{(n,l,m)}=q^{n}r^{l}s^{tm}$ exponential normalization of
\cite[\S5.1.1]{GNII}, the products are
\[
\begin{aligned}
\Delta_{2}(Z)&=q^{1/4}r^{1/2}s^{1/2}\!\!
\prod_{(n,l,m)\in\Gamma_{2}^{\mathrm{prod}}}\!\!
(1-q^{n}r^{l}s^{2m})^{f_{2}(nm,l)},\\
\Delta_{1}(Z)&=q^{1/6}r^{1/2}s^{1/2}\!\!
\prod_{(n,l,m)\in\Gamma_{3}^{\mathrm{prod}}}\!\!
(1-q^{n}r^{l}s^{3m})^{f_{3}(nm,l)},\\
\Delta_{1/2}(Z)&=q^{1/8}r^{1/2}s^{1/2}\!\!
\prod_{(n,l,m)\in\Gamma_{4}^{\mathrm{prod}}}\!\!
(1-q^{n}r^{l}s^{4m})^{f_{4}(nm,l)},
\end{aligned}
\]
with cusp Jacobi inputs $\phi_{0,2},\phi_{0,3},\phi_{0,4}$ supplied
explicitly by \cite[\S5]{GNII}.  Rows $5$--$8$ carry Cl\'ery--Gritsenko
products on the corresponding congruence covers, with cusp Jacobi inputs
$\phi_{j}$ as in \cite[Theorem~3.1]{GC}, and the additional
denominator data (Weyl chamber, simple real roots, reflection character,
Weyl alternating sum) supplied by Govindarajan's twisted-CHL theorem
\cite{Gov11BKM}, \cite{GK09CHL}, \cite{CD09BKM} for $F_{5}=\nabla_{3}$,
$F_{6}=\nabla_{2}$, $F_{7}=\nabla_{3/2}$, $F_{8}=Q_{1}$.  For $F_{7}$ the
denominator theorem lives on the order-four multiplier cover
$\Gamma_{0}^{(2)}(4)\langle\chi_{4}^{(4)}\times v_{H}\rangle$.
\emph{(proved; \cite{Gov11BKM},
\cite[Theorem~2.1]{GNII}, \cite[Theorem~3.1]{GC}.)}

\paragraph{Level separation.}
The eight-row catalogue is a list of \emph{automorphic products plus
denominator theorems on the rows where the data have been supplied}.
It is not a list of compact Hall sources, reduced curve-counting
backgrounds, or chiral hosts.  For a row \(j\), the automorphic product
\(F_j\) acquires a compact physical host only after separate choices of
an automorphic line, a denominator algebra, a scalar trace, a chiral
source, and a Hall realization; these choices are row-dependent.  The
row $j=1$
($F_{1}=\Delta_{5}$) is the only row with a proved reduced CY/DT scalar
host (Oberdieck--Pixton, Oberdieck--Pandharipande:
$Z^{X}_{\mathrm{OP}}=-4096\Delta_{5}^{-2}$); even row $1$ has open
compact Hall/chiral host at the primitive-recognition level, governed
by \((S1)\)--\((S10)\).

\subsection*{B.7 The Borcherds-weight coordinate \texorpdfstring{$\kappa_{\mathrm{BKM}}=5$}{kappa\_BKM=5}}

The universal Borcherds-weight identity
\(\kappa_{\mathrm{BKM}}(\Phi_{N})=c_{N}(0)/2\) is the row-by-row
specialization of the Borcherds singular-theta-lift weight formula
\cite[Theorem~10.1]{Borcherds95} together with the Gritsenko--Nikulin
unification \cite[\S2]{GN}, \cite{Gritsenko99}.  The relevant row is
$j=1$, with
\[
\kappa_{\mathrm{BKM}}(\Delta_{5})=\operatorname{wt}(\Delta_{5})=5
=\frac{c_{1}(0)}{2}=\frac{f(0,0)}{2}=\frac{10}{2}.
\]
\emph{(proved; \cite[Theorem~10.1]{Borcherds95},
\cite[\S2.2]{EZ:1}, \cite[Theorem~3.1]{GC}.)}

This is the Borcherds-weight invariant: the same integer
$\kappa_{\mathrm{BKM}}=5$ identifies the \(\Delta_5\) row inside the
companion Calabi--Yau quantum-groups $\kappa$-stratification, and it
is the entry that distinguishes
$\kappa_{\mathrm{BKM}}=\operatorname{wt}(\Delta_{5})$ (the
Igusa cusp-form weight) from the unrelated quantities
$\kappa_{\mathrm{ch}}^{\mathrm{Hodge}}=0$,
$\kappa_{\mathrm{ch}}^{\mathrm{Heis}}=3$, $\kappa_{\mathrm{cat}}=0$,
$\kappa_{\mathrm{fiber}}=24=\chi_{\mathrm{top}}(\mathrm{K3})$.  The
additive formula $\kappa_{\mathrm{BKM}}=\kappa_{\mathrm{ch}}+
\chi(\mathcal O_{\mathrm{fiber}})$ would identify
$24=\chi_{\mathrm{top}}(\mathrm{K3})$ with
$2=\chi(\mathcal O_{\mathrm{K3}})$.  The Heisenberg reading gives an
accidental equality at \(N=1\), but it does not extend to the CHL
weights \(c_N(0)/2=(5,3,2,\tfrac32,1)\); the appendix
asserts $\kappa_{\mathrm{BKM}}=5$ as the row entry of the
Borcherds-weight identity, not as the value of any additive formula.

\subsection*{B.8 Normalizations and residual obstructions}

The following normalizations fix the Borcherds product without residual
scalar, chamber, multiplier, or coefficient ambiguity.

\begin{itemize}
\item \emph{Theta-constant normalization.}  The constant \(64\) is not
a free scalar on \(\Delta_{5}\).  The identity \(64=2^{6}\) is the
explicit theta-product leading coefficient
\cite{Igusa1962}, \cite[Theorem~4.1]{GN}; with the theta normalization
fixed, no residual scalar freedom remains.  Cross-checked against
the consistency
$f(1,1)=-64$ and the \(qrs\)-coefficient identity \(c_{\Delta}(1,1,1)=64\)
\cite[Theorem~4.1]{GN}.
\item \emph{Cusp chamber definition \(\Gamma_{\mathrm{eff}}\).}
The semigroup $\Gamma_{\mathrm{eff}}$ is the ordered-limit semigroup of
the type-II cusp polarization $0<|s|\ll|q|\ll|r|^{-1}\ll1$
\cite[Example~2.4]{GNII}; not a free choice.  Triples in
$\Gamma_{\mathrm{eff}}\setminus\Gamma_{\mathrm{act}}$ have zero exponent
and therefore contribute no visible product factor.
\item \emph{Multiplier system for half-integer weights.}
Half-integer weights $\tfrac12$ and $\tfrac32$ are realized by
$v_{\eta}^{k}\times v_{H}$, not by the Weil-representation metaplectic
double cover; the integer weight $5$ of $\Delta_{5}$ uses the abelian
multiplier $v_{\eta}^{12}\times v_{H}$ (which is automorphic-trivial
modulo the Maass character $\nu_{\Delta_{5}}$ on the type-II Weyl
generators).  No metaplectic factor enters the Borcherds weight
identity.  See \cite[Theorem~3.1]{GC} for the full classification.
\item \emph{Doubling factor $64^{-1}\Delta_{5}(2Z)$.}
The factor $64^{-1}$ in $\operatorname{den}(\autcor)=64^{-1}
\Delta_{5}(2Z)$ is the same theta-constant $2^{6}$ from
Theorem~\ref{thm:borcherds-product-delta5}, transported through
$Z\mapsto 2Z$; no new constant.  See
Lemma~\ref{lem:appB-doubling}.
\item \emph{The $\kappa_{\mathrm{BKM}}=5$ row.}
The integer $\kappa_{\mathrm{BKM}}=5$ is the row entry of the
Borcherds-weight identity $\operatorname{wt}(F_{j})=c_{N}(0)/2$ in the
catalogue of \cite[Theorem~3.1]{GC}; it agrees with
$\operatorname{wt}(\Delta_{5})=f(0,0)/2=5$ from
\cite[\S2.2]{EZ:1}; it agrees with the entry
$(c_{1}(0),\operatorname{wt})=(10,5)$ in the companion Calabi--Yau
quantum-groups $\kappa$-stratification.
\item \emph{Eichler--Zagier vs Gritsenko--Nikulin convention.}  The
present appendix uses the Eichler--Zagier convention
$\phi_{0,1}=Z_{\mathrm{K3}}/2$ \cite[\S2.2]{EZ:1}; this matches the
Gritsenko--Nikulin Jacobi convention in
\cite[Example~2.4]{GNII} and \cite[\S5.1.1]{GNII} on the
$(n,l)$-coefficient table given in B.1, and matches the singular
theta-lift normalization \cite[Theorem~10.1]{Borcherds95} after the
Mp$_{2}$-vector-valued bridge.  No residual factor of $2$ enters the
exponent $f(nm,l)$ of the Borcherds product.
\end{itemize}

\paragraph{Residual obstructions.}
The four obstruction classes (D0), (O1), (O1)\(^+\), (O2) are not
part of the Borcherds product calculation; they are geometric
content on the source side of $\Delta_{5}$.  Within the automorphic calculation, no
residual Borcherds-product obstruction remains.  The
universal identity $\operatorname{wt}(F_{j})=c_{N}(0)/2$ fixes the row
$j=1$, $\kappa_{\mathrm{BKM}}=5$, and any
disagreement with the Vol~III catalogue, the Borcherds singular
theta-lift weight formula, or the Gritsenko--Cl\'ery classification
would be a contradiction among the stated comparison theorems.

\paragraph{Computations.}
Fourier coefficients are obtained from the rational theta-quotient
expansion; lattice and Gram-pairing identities are the computations of
Appendix~\ref{app:lattice-II21};
the type-II cusp chamber, active support, and Borcherds product are in
Proposition~\ref{prop:bgn-identification-recap}; the doubling identity is
Remark~\ref{rem:bkm-doubling-Z}; the catalogue tables
are in Appendix~\ref{app:eight-finite-constructions}.

\section{Mukai--Gram dictionary on \texorpdfstring{$K3\times E$}{K3 x E}}
\label{appendix:mukai-gram-dictionary}

Let \(X=S\times E\) be the smooth trivial-canonical threefold, with
\(S\) a smooth complex K3 surface and \(E\) a smooth complex elliptic curve. The
Mukai vector and Mukai pairing on \(\widetilde H(S,\Z)\), together with
the K\"unneth decomposition on \(H^*(X,\Z)\), produce the Gram-degree
projection \(\Pi_X:\Gamma_X^{\mathrm{form}}\to\Gamma_{\mathrm{gram}}\)
and its normal-ordered companion \(\overline\Pi_X\) on the central
extension \(\widehat\Gamma_X\). The Borcherds expansion variables
\(q^nr^ls^m\) of an Igusa cusp expansion record Gram-degree shadows in
\(\Gamma_{\mathrm{gram}}\).  When the shadow is carried by an actual
one-dimensional subscheme \(Y\subset X\), its dyonic charge
\(c=(Q_Y,P_Y)\) satisfies the formula below.  The existence of such a
subscheme is the algebraic-effective Hall-support condition, not a
formal consequence of the Fourier monomial. For \(Y\) of class
\([Y]=(\beta_h,d)\) and \(\chi(\mathcal O_Y)=n\),
\[
\Pi_X(Q_Y,P_Y)=(h-1,\,n,\,d-1).
\]
The threefold arithmetic Euler characteristic
\(n=\chi(\mathcal O_Y)\) is the direct structure-sheaf invariant; no
``Todd correction'' enters, and no ambient fourfold appears.

The Mukai--Gram construction sits at Beilinson chart layer
\(\alpha\) (intrinsic charge data) and
\(\gamma\) (ambient Mukai/Gram lattice). The
\(\kappa\)-tuple, orientation character \(\epsilon_o\), Pfaffian
section, and recognition theorem live at higher Beilinson levels and
depend on this chart datum.

\subsection{Mukai vector and Mukai pairing on a K3 surface}
\label{subsec:mukai-vector-pairing}

\begin{definition}[Mukai lattice and Mukai vector]
\label{def:mukai-lattice-vector-appendixC}
\ClaimStatusDefinition\
Let \(S\) be a smooth complex K3 surface. The Mukai lattice is the
total cohomology
\[
\widetilde H(S,\Z)
:=H^0(S,\Z)\oplus H^2(S,\Z)\oplus H^4(S,\Z),
\]
graded by even cohomological degree, an even unimodular lattice of
signature \((4,20)\) and rank \(24\) when equipped with the Mukai
pairing below. For a coherent sheaf \(F\in\mathrm{Coh}(S)\),
respectively for an object \(F\in D^b(S)\), the Mukai vector is
\[
v_S(F)
:=\ch(F)\cdot\sqrt{\td(S)}
=\bigl(\rk F,\;c_1(F),\;\ch_2(F)+\rk F\cdot[\mathrm{pt}]\bigr)
\;\in\;\widetilde H(S,\Z).
\]
The square root \(\sqrt{\td(S)}=(1,0,1)\) on the K3 surface uses the
Todd class \(\td(S)=(1,0,2[\mathrm{pt}])\), so
\(\sqrt{\td(S)}=(1,0,[\mathrm{pt}])\); identifying
\([\mathrm{pt}]\) with the integer \(1\) in the displayed coordinates
gives the displayed vector. The triple
\(v_S(F)=(r,D,s)\) records the rank \(r\in H^0\), the first Chern class
\(D\in H^2\), and the Mukai-corrected second component
\(s=\ch_2(F)+r\in H^4\). Mukai introduced this lattice and pairing in
\cite[Theorem~1.5]{Mukai1987}, with the symplectic structure of the
moduli space of sheaves established in \cite{Mukai1984}.
\end{definition}

\begin{proposition}[Mukai pairing on the K3 surface]
\label{prop:mukai-pairing-appendixC}
\ClaimStatusProvedElsewhere\
The Mukai pairing on \(\widetilde H(S,\Z)\) is the symmetric \(\Z\)-bilinear
form
\[
\bigl\langle(r,D,s),(r',D',s')\bigr\rangle_{\mathrm{Muk}}
=D\cdot D'-rs'-r's,
\]
equivalently
\[
\bigl\langle v,w\bigr\rangle_{\mathrm{Muk}}
=-\int_S v^{\vee}\cup w,
\qquad
v^\vee:=(r,-D,s),
\]
where \(v_2\cup w_2\) is the K3 cup product on \(H^2(S,\Z)\) and
\(v_0\cup w_4\), \(v_4\cup w_0\) record the cross-degree pairings.
The form is symmetric, even, unimodular, of signature \((4,20)\); its
hyperbolic-plane summands are
\[
\widetilde H(S,\Z)
\simeq U^{\oplus4}\oplus E_8(-1)^{\oplus2}
\]
as integral lattices.  For a polarized K3 of Picard rank \(\rho\), the
N\'eron--Severi lattice \(\mathrm{NS}(S)\) has signature
\((1,\rho-1)\), while the algebraic Mukai lattice
\[
\widetilde H^{1,1}_{\mathrm{alg}}(S,\Z)
=H^0(S,\Z)\oplus\mathrm{NS}(S)\oplus H^4(S,\Z)
\]
has signature \((2,\rho)\).  This is Mukai
\cite[Theorem~1.5(ii)]{Mukai1987}; see also Huybrechts--Lehn
\cite[Chapter~6]{HuybrechtsLehn}.
\end{proposition}

\begin{proof}
The bilinear identity is the definition; its symmetry is immediate from
the symmetry of the K3 cup product and the identification of degree-zero
and degree-four components. Even unimodularity follows from the K3
intersection lattice on \(H^2\) (even unimodular of signature
\((3,19)\)) and the unimodular pairing
\(H^0\times H^4\to\Z\) under \(rs'\mapsto rs'\). The signature
\((4,20)\) follows by Hodge index plus the hyperbolic factor on
\(H^0\oplus H^4\).  The same hyperbolic factor changes
\(\mathrm{NS}(S)\) of signature \((1,\rho-1)\) into
\(\widetilde H^{1,1}_{\mathrm{alg}}(S,\Z)\) of signature
\((2,\rho)\).  The lattice isomorphism is
\cite[Theorem~1.5(ii)]{Mukai1987}.
\end{proof}

\begin{remark}[Mukai pairing is symmetric on both parts]
\label{rem:mukai-symmetric-both-parts}
\ClaimStatusProved\
The Mukai pairing on \(\widetilde H(S,\Z)\) is symmetric in both the
algebraic and transcendental parts. There is no orthosymplectic
\(\mathfrak{osp}\) refinement at the level of the lattice. The
Hodge \(\Z/2\)-super-extension \((\mathbb C,\omega_S)\mapsto
\mathbb C\oplus\omega_S\) used in Mukai self-mirror
language imposes a separate \(\Z/2\)-grading on a Yangian-type quantum
group constructed from this lattice; that grading does not change the
underlying lattice pairing, which remains symmetric. The classical Lie
limit of the Mukai-form Yangian is the orthogonal Lie algebra
\(\mathfrak{so}(4,20)\), not \(\mathfrak{osp}(4\mid20)\); see
\cite[\S2]{SchiffmannVasserot2013} for the Mukai-side Yangian
classical limit and \cite{KSCoHA} for the cohomological Hall algebra
realization.
\end{remark}

\begin{example}[Mukai vector of a point and of an ideal sheaf]
\label{ex:mukai-vector-point-and-ideal}
\ClaimStatusDerivation\
\(v_S(\mathcal O_p)=(0,0,1)\) for the structure sheaf of a point
\(p\in S\); \(v_S(\mathcal O_S)=(1,0,1)\) for the trivial line bundle;
\(v_S(L)=(1,c_1(L),c_1(L)^2/2+1)\) for a line bundle \(L\). For an ideal
sheaf \(\mathcal I_C\) of a smooth curve \(C\subset S\) of genus \(g\),
the adjunction formula on the K3 surface gives
\(C^2=c_1(\mathcal O_S(C))^2=2g-2\), and
\[
v_S(\mathcal I_C)=(1,-c_1(\mathcal O_S(C)),\ch_2(\mathcal I_C)+1).
\]
The Mukai pairing
\(\langle v_S(\mathcal O_p),v_S(\mathcal O_p)\rangle_{\mathrm{Muk}}
=0-0\cdot1-0\cdot1=0\), recording that points have zero self-pairing in
the Mukai form. The pairing
\(\langle v_S(\mathcal O_S),v_S(\mathcal O_p)\rangle_{\mathrm{Muk}}
=0-1\cdot1-0\cdot0=-1\) records the Euler characteristic
\(\chi(\mathcal O_S,\mathcal O_p)=1\) up to the Mukai sign convention
\(\chi(F,G)=-\langle v_S(F),v_S(G)\rangle_{\mathrm{Muk}}\) of
\cite[Proposition~2.2]{Mukai1987}.
\end{example}

\subsection{K\"unneth on \texorpdfstring{$K3\times E$}{K3 x E} and the rank-3 Gram lattice}
\label{subsec:kunneth-k3-times-e-appendixC}

\begin{proposition}[K\"unneth for \(X=K3\times E\)]
\label{prop:kunneth-k3-times-e-appendixC}
\ClaimStatusProvedElsewhere\
Let \(S\) be a smooth complex K3 surface and \(E\) a smooth complex
elliptic curve. The even integral cohomology of \(X=S\times E\) splits as
\[
H^{\mathrm{even}}(X,\Z)
\simeq H^*(S,\Z)\otimes_{\Z}H^{\mathrm{even}}(E,\Z),
\]
the odd--odd K\"unneth term \(H^{\mathrm{odd}}(S)\otimes H^{\mathrm{odd}}(E)\)
absent because \(S\) carries no odd cohomology, both factors free
\(\Z\)-modules. Fixing a generator \(\omega_E\in H^2(E,\Z)\) with
\(\int_E\omega_E=1\),
\(H^{\mathrm{even}}(E,\Z)=\Z\cdot1_E\oplus\Z\cdot\omega_E\), and every
even class on \(X\) decomposes uniquely as
\[
v_X=v^{(0)}\otimes 1_E+v^{(\omega)}\otimes\omega_E,
\qquad v^{(0)},v^{(\omega)}\in H^*(S,\Z).
\]
The formal K\"unneth lattice
\[
H^{\mathrm{even}}(X,\Z)
\simeq \widetilde H(S,\Z)\otimes U_E,\qquad
U_E:=H^0(E,\Z)\oplus H^2(E,\Z),
\]
carries the tensor K\"unneth pairing of rank \(48\) and signature
\((24,24)\), obtained from the \((4,20)\)-Mukai pairing on
\(\widetilde H(S,\Z)\) and the \((1,1)\)-hyperbolic pairing on \(U_E\).
This cross-pairing is distinct from the direct-sum componentwise Gram
pairing on \(\Gamma_X^{\mathrm{form}}\) used by \(\Pi_X\).  The odd
K\"unneth term \(H^{\mathrm{odd}}(S)\otimes H^{\mathrm{odd}}(E)\) is
absent because the K3 surface \(S\) carries no odd cohomology. See
Huybrechts \cite[Chapter~3]{HuybrechtsLehn}.
\end{proposition}

\begin{proof}
\(E\) and \(S\) have torsion-free integral cohomology, so K\"unneth holds
integrally and produces the displayed tensor decomposition. The
tensor-pairing factorization is bilinearity of the cup product over the
K\"unneth isomorphism. The lattice signature comes from
\((4,20)\) on \(\widetilde H(S,\Z)\) (Proposition~\ref{prop:mukai-pairing-appendixC})
tensored with \((1,1)\) on \(U_E\).
\end{proof}

\begin{definition}[Gram-degree projection on \(X=K3\times E\)]
\label{def:gram-degree-projection-appendixC}
\ClaimStatusDefinition\
The formal additive even-Mukai sector is
\[
\Gamma_X^{\mathrm{form}}
:=\widetilde H(S,\Z)\oplus\widetilde H(S,\Z),\qquad
c=(Q,P)\in\Gamma_X^{\mathrm{form}},
\]
with \(Q,P\in\widetilde H(S,\Z)\). The Gram-index lattice is
\[
\Gamma_{\mathrm{gram}}:=\Z^3,
\]
and the Gram-degree projection is the quadratic map
\[
\Pi_X:\Gamma_X^{\mathrm{form}}\longrightarrow\Gamma_{\mathrm{gram}},
\qquad
(Q,P)\longmapsto
\Bigl(\tfrac{(Q,Q)_{\mathrm{Muk}}}2,\,
(Q,P)_{\mathrm{Muk}},\,
\tfrac{(P,P)_{\mathrm{Muk}}}2\Bigr).
\]
The triple \((n,l,m)=\Pi_X(Q,P)\) is the Gram-degree shadow; it is the
target of the Igusa Fourier expansion at the standard cusp. It is
\emph{not} an additive Hall grading: the additivity defect is the
symmetric bilinear cocycle \(B\) of \S\ref{subsec:gram-cocycle-appendixC}
below.
\end{definition}

\subsection{D6--D2--D0 dictionary}
\label{subsec:d6-d2-d0-dictionary-appendixC}

The K\"unneth decomposition of \(v_X(\mathcal I_Y)\) on the threefold
\(X=K3\times E\) produces the D6--D2--D0 dictionary
\(\Pi_X(Q_Y,P_Y)=(h-1,n,d-1)\) of
Theorem~\ref{thm:d6-d2-d0-dictionary-appendixC} as a chart-data
identity. The Mukai-vector calculation isolates the K3-fibre rank-one
component \(P_Y=(1,0,1-d)\) on the \(1_E\)-summand and the curve-class
component \(Q_Y=(0,-\beta^\vee,-n)\) on the \(\omega_E\)-summand, with
\(\beta^\vee\in H^2(S,\Z)\) the Poincar\'e dual of \(\beta\); the Mukai
pairing on the K3 components delivers the Gram-degree shadow.

\begin{theorem}[Mukai vector of an ideal sheaf on the threefold $K3\times E$]
\label{thm:mukai-vector-ideal-sheaf-threefold}
\ClaimStatusProvedHere\
Let \(X=S\times E\) with \(S\) a smooth complex K3 surface and \(E\) a
smooth complex elliptic curve, so \(K_X\simeq\mathcal O_X\) and \(X\)
is a smooth trivial-canonical threefold. Let \(Y\subset X\) be a
one-dimensional closed subscheme of
class \([Y]=(\beta,d)\in H_2(S,\Z)\oplus H_2(E,\Z)\), write
\(\beta^\vee\in H^2(S,\Z)\) for the Poincar\'e dual of \(\beta\), and let
\(n=\chi(\mathcal O_Y)\in\Z\). The Mukai vector of the
ideal sheaf \(\mathcal I_Y\in D^b(X)\) decomposes against the K\"unneth
basis as
\[
v_X(\mathcal I_Y)
=(1,0,1-d)\otimes1_E
+(0,-\beta^\vee,-n)\otimes\omega_E.
\]
\end{theorem}

\begin{proof}
The trivial-canonical threefold condition gives
\(\sqrt{\td(X)}=\pi_S^*\sqrt{\td(S)}\,\pi_E^*\sqrt{\td(E)}=(1,0,1)\otimes1\),
since the elliptic curve is one-dimensional with trivial Todd class
\(\td(E)=1\). The ideal-sheaf identity is
\(\ch(\mathcal I_Y)=1-\ch(\mathcal O_Y)\). For the structure sheaf of a
one-dimensional subscheme of a trivial-canonical threefold,
\(\ch_0(\mathcal O_Y)=0\), \(\ch_1(\mathcal O_Y)=0\),
\(\ch_2(\mathcal O_Y)=\operatorname{PD}[Y]\), and
\(\int_X\ch_3(\mathcal O_Y)=\chi(\mathcal O_Y)=n\). Decomposing
\([Y]=(\beta,d)\) on the K\"unneth basis,
\(\operatorname{PD}[Y]=\beta^\vee\otimes\omega_E+d[\mathrm{pt}_S]\otimes 1_E\),
where \([\mathrm{pt}_S]\) is the K3 point class; the fibrewise
elliptic-degree component \(d[\mathrm{pt}_S]\otimes 1_E\) records the
elliptic degree. Hence
\[
\ch(\mathcal O_Y)
=(0,0,d)\otimes1_E+(0,\beta^\vee,n)\otimes\omega_E,
\]
and
\(1-\ch(\mathcal O_Y)
=(1,0,-d)\otimes1_E+(0,-\beta^\vee,-n)\otimes\omega_E\).
Multiplying by
\(\sqrt{\td(X)}=(1,0,1)\otimes1_E\) using the K3 cup product on each
K\"unneth summand,
\((1,0,1)\cup(1,0,-d)=(1,0,1-d)\) since \([\mathrm{pt}_S]\cup1=[\mathrm{pt}_S]\),
and \((1,0,1)\cup(0,-\beta^\vee,-n)=(0,-\beta^\vee,-n)\) since
\(\beta^\vee\cup[\mathrm{pt}_S]=0\) in \(H^*(S,\Z)\), yields the displayed
Mukai vector
\(v_X(\mathcal I_Y)=(1,0,1-d)\otimes1_E+(0,-\beta^\vee,-n)\otimes\omega_E\).
\end{proof}

\begin{remark}[Threefold Euler characteristic, no Todd correction]
\label{rem:threefold-euler-characteristic-no-todd-correction}
\ClaimStatusDerivation\
The integer \(n=\chi(\mathcal O_Y)\in\Z\) is the holomorphic Euler
characteristic of the structure sheaf \(\mathcal O_Y\) on the threefold
\(X=S\times E\); the trivial-canonical threefold condition makes
\(\sqrt{\td(X)}\) act as identity on the third Chern character so that
\(\int_X\ch_3(\mathcal O_Y)=n\) directly. There is no auxiliary
``Todd correction'' to install on top of \(n\); the Mukai vector
\((0,-\beta^\vee,-n)\otimes\omega_E\) records \(n\) in the K3
\(H^4\)-coordinate on the elliptic \(\omega_E\)-summand. The sign is fixed by
\(\ch(\mathcal I_Y)=1-\ch(\mathcal O_Y)\) and the K\"unneth orientation
\(\int_E\omega_E=+1\).
\end{remark}

\begin{theorem}[D6--D2--D0 Mukai--Gram dictionary, restated as chart data]
\label{thm:d6-d2-d0-dictionary-appendixC}
\ClaimStatusProvedHere\
With the standing data of
Theorem~\ref{thm:mukai-vector-ideal-sheaf-threefold}, set
\[
P_Y:=(1,0,1-d),\qquad
Q_Y:=(0,-\beta^\vee,-n)\;\in\;\widetilde H(S,\Z),
\]
the K\"unneth components of \(v_X(\mathcal I_Y)\). Then the Gram-degree
projection of Definition~\ref{def:gram-degree-projection-appendixC} reads
\[
\Pi_X(Q_Y,P_Y)
=\Bigl(\tfrac{(Q_Y,Q_Y)_{\mathrm{Muk}}}2,\,
(Q_Y,P_Y)_{\mathrm{Muk}},\,
\tfrac{(P_Y,P_Y)_{\mathrm{Muk}}}2\Bigr)
=\Bigl(\tfrac{(\beta^\vee)^2}2,\,n,\,d-1\Bigr).
\]
For the elliptically fibered K3 surface \(S\to\mathbb P^1\) with section
\(s_1\) and curve class
\(\beta=\beta_h=s_1+hF\) (\(h\in\Z_{\ge0}\), \(F\) the fibre class), the
self-intersection
\(\beta_h^2=2h-2\) gives
\[
\Pi_X(Q_Y,P_Y)=(h-1,\,n,\,d-1).
\]
\end{theorem}

\begin{proof}
Apply the Mukai pairing of
Proposition~\ref{prop:mukai-pairing-appendixC} component by component:
\begin{align*}
(Q_Y,Q_Y)_{\mathrm{Muk}}
&=(-\beta^\vee)\cdot(-\beta^\vee)-0\cdot(-n)-0\cdot(-n)=(\beta^\vee)^2,\\
(Q_Y,P_Y)_{\mathrm{Muk}}
&=(-\beta^\vee)\cdot 0-0\cdot(1-d)-1\cdot(-n)=n,\\
(P_Y,P_Y)_{\mathrm{Muk}}
&=0\cdot0-1\cdot(1-d)-1\cdot(1-d)=2(d-1),
\end{align*}
giving the half-pairings \(((\beta^\vee)^2/2,n,d-1)\). For the elliptic class
\(\beta_h=s_1+hF\), identified with its Poincar\'e dual, the K3 intersection \(s_1^2=-2\) on a section of an
elliptic K3 surface (genus formula \(2g(s_1)-2=s_1^2+s_1\cdot K_S=-2+0\),
giving \(g(s_1)=0\)), \(s_1\cdot F=1\), \(F^2=0\), so
\(\beta_h^2=s_1^2+2h(s_1\cdot F)+h^2 F^2=-2+2h=2h-2\). The Gram-degree
shadow follows. The cross-level identity with the OP/DT scalar
normalization is the
\(Q_Y^2/2=h-1\) line of Theorem~\ref{thm:d6-d2-d0-dictionary-appendixC}.
\end{proof}

\begin{remark}[Algebraic vs primitive D2 charges]
\label{rem:algebraic-versus-primitive-d2-charges}
\ClaimStatusDerivation\
The K3 curve class \(\beta\in H_2(S,\Z)\) appearing through its
Poincar\'e dual in \(Q_Y=(0,-\beta^\vee,-n)\) is algebraic:
\(\beta^\vee\in\mathrm{NS}(S)\), since \(Y\subset X\) is a
closed subscheme. Primitivity of \(\beta\) — that \(\beta\) is not a
non-trivial integer multiple of another effective class — is a separate
condition; the dictionary holds for any algebraic \(\beta\), not only
primitive ones. The Mukai-pairing formula
\((Q_Y,Q_Y)_{\mathrm{Muk}}=(\beta^\vee)^2\) records the
\(\mathrm{NS}(S)\)-self-intersection regardless of primitivity. The
Bryan--Oberdieck--Pandharipande--Yin/Maulik--Pandharipande conventions
\cite{OPand,BRYAN,MaulikTodaGV} take \(\beta\) algebraic; the elliptic
degree \(d\in\Z_{\ge0}\) is unconstrained beyond non-negativity in the
effective semigroup.
\end{remark}

\begin{remark}[Chart-data layer; not the additive Hall grading]
\label{rem:chart-data-not-additive-hall-grading}
\ClaimStatusProved\
The triple \((n,l,m)=\Pi_X(Q_Y,P_Y)\) records the Gram-degree
invariants of the dyonic charge \(c=(Q,P)\). It is the target of the
Igusa Fourier expansion and the Borcherds product variables
\(q^nr^ls^m\). It is \emph{not} the additive Hall grading on a
hypothetical compact \(K3\times E\) Hall sector; the additive grading is
the upstairs charge \(c\in\Gamma_X^{\mathrm{form}}\) (or its
normal-ordered lift \(\widehat\Gamma_X\) below), and \(\Pi_X\) is a
quadratic map onto Gram invariants. The polarization defect that obstructs
the upstairs additive grading from descending directly is the symmetric
bilinear cocycle \(B\) of \S\ref{subsec:gram-cocycle-appendixC}.
\end{remark}

\subsection{Symmetric bilinear Gram cocycle}
\label{subsec:gram-cocycle-appendixC}

\begin{lemma}[Gram additivity defect, restated as chart structure]
\label{lem:gram-additivity-defect-appendixC}
\ClaimStatusProvedElsewhere\
For \(c_i=(Q_i,P_i)\in\Gamma_X^{\mathrm{form}}\), \(i=1,2\),
\[
\Pi_X(c_1+c_2)
=\Pi_X(c_1)+\Pi_X(c_2)+B(c_1,c_2),
\]
where the symmetric bilinear cocycle \(B\) is
\[
B(c,c')
:=\bigl((Q,Q')_{\mathrm{Muk}},\;
(Q,P')_{\mathrm{Muk}}+(Q',P)_{\mathrm{Muk}},\;
(P,P')_{\mathrm{Muk}}\bigr).
\]
The map \(B\) is symmetric and \(\Z\)-bilinear with values in
\(\Gamma_{\mathrm{gram}}=\Z^3\), satisfies the polarization identity
\[
B(c,c)=2\Pi_X(c),
\]
and equals \(-\delta\Pi_X\) as an additive coboundary,
\((\delta\Pi_X)(c,c'):=\Pi_X(c)+\Pi_X(c')-\Pi_X(c+c')\). Thus the
ordinary additive cohomology class \([B]=0\); the obstruction recorded
here is the impossibility of an \emph{additive linear} trivialization,
not a non-zero additive cohomology class.
\end{lemma}

\begin{proof}
This is Lemma~\ref{lem:gram-additivity-defect}.  Symmetry,
bilinearity, integrality, polarization, and \(\delta\)-cocycle property
are Lemma~\ref{lem:gram-additivity-defect-appendixC}(i)--(iii). The
splitting \(B=-\delta\Pi_X\) by the quadratic refinement \(\Pi_X\) is
clause~(vi) of the same lemma; clause~(vii) shows no additive linear
\(L:\Gamma_X^{\mathrm{form}}\to\Gamma_{\mathrm{gram}}\) can satisfy
\(B=\delta L\), since \(B(c,c)=2\Pi_X(c)\) is non-zero on charges with
non-trivial Mukai self-pairings.
\end{proof}

\begin{remark}[Cleanly stated polarization identity]
\label{rem:cleanly-stated-polarization-identity}
\ClaimStatusDerivation\
The polarization identity
\(B(c,c)=2\Pi_X(c)\) is the symmetric companion to
\(\Pi_X(c+c')=\Pi_X(c)+\Pi_X(c')+B(c,c')\); evaluated at
\(c=c'\) the latter gives
\(\Pi_X(2c)=2\Pi_X(c)+B(c,c)\), and \(\Pi_X(2c)=4\Pi_X(c)\) by
quadraticity of \(\Pi_X\), hence \(B(c,c)=2\Pi_X(c)\). This is purely
chart-data structure; no Heisenberg or central extension hypothesis on
the source is invoked.
\end{remark}

\subsection{Normal-ordered Gram map}
\label{subsec:normal-ordered-gram-map-appendixC}

\begin{definition}[Normal-ordered charge lattice]
\label{def:normal-ordered-charge-lattice-appendixC}
\ClaimStatusDefinition\
Set
\[
\widehat\Gamma_X
:=\Gamma_X^{\mathrm{form}}\oplus_B\Gamma_{\mathrm{gram}},
\]
the underlying set
\(\Gamma_X^{\mathrm{form}}\times\Gamma_{\mathrm{gram}}\) equipped with
the abelian-group law
\[
(c,T)\star(c',T'):=\bigl(c+c',\,T+T'+B(c,c')\bigr).
\]
The identity is \((0,0)\); the inverse of \((c,T)\) is
\((-c,\,-T+B(c,c))=(-c,\,-T+2\Pi_X(c))\). The split section
\[
s:\Gamma_X^{\mathrm{form}}\longrightarrow\widehat\Gamma_X,
\qquad s(c):=(c,\Pi_X(c)),
\]
is a homomorphism of abelian groups; the raw placement
\[
i_0:\Gamma_X^{\mathrm{form}}\longrightarrow\widehat\Gamma_X,
\qquad i_0(c):=(c,0),
\]
is not a homomorphism unless \(B\equiv0\).
\end{definition}

\begin{proposition}[Normal-ordered Gram homomorphism]
\label{prop:normal-ordered-gram-homomorphism-appendixC}
\ClaimStatusProvedElsewhere\
The map
\[
\overline\Pi_X:\widehat\Gamma_X\longrightarrow\Gamma_{\mathrm{gram}},
\qquad
\overline\Pi_X(c,T):=\Pi_X(c)-T,
\]
is a homomorphism of abelian groups, \(\overline\Pi_X\colon
(\widehat\Gamma_X,\star)\to(\Gamma_{\mathrm{gram}},+)\); equivalently,
\[
\overline\Pi_X\bigl((c,T)\star(c',T')\bigr)
=\overline\Pi_X(c,T)+\overline\Pi_X(c',T').
\]
On the split section,
\(\overline\Pi_X(s(c))=0\); on the raw placement,
\(\overline\Pi_X(i_0(c))=\Pi_X(c)\).
\end{proposition}

\begin{proof}
Direct calculation,
\begin{align*}
\overline\Pi_X((c,T)\star(c',T'))
&=\overline\Pi_X(c+c',T+T'+B(c,c'))\\
&=\Pi_X(c+c')-T-T'-B(c,c')\\
&=\bigl(\Pi_X(c)+\Pi_X(c')+B(c,c')\bigr)-T-T'-B(c,c')\\
&=\bigl(\Pi_X(c)-T\bigr)+\bigl(\Pi_X(c')-T'\bigr)\\
&=\overline\Pi_X(c,T)+\overline\Pi_X(c',T'),
\end{align*}
using Lemma~\ref{lem:gram-additivity-defect-appendixC}. The split-section
and raw-placement identities are the definitions.
\end{proof}

The finite formal packet
\(\texttt{certificates/charge/mukai\_gram\_cocycle}\) records explicit
rank-one Mukai-slice rows for the three simple Gram shadows
\(\delta_1,\delta_2,\delta_3\) and one test charge.  Its verifier checks
the fixed Mukai sign convention, integrality of \(\Pi_X\), symmetry and
bilinearity of \(B\), the identity
\[
B(c,c')=\Pi_X(c+c')-\Pi_X(c)-\Pi_X(c'),
\]
the polarization identity \(B(c,c)=2\Pi_X(c)\), additivity of the split
section \(s(c)=(c,\Pi_X(c))\), non-additivity of the raw placement
\((c,0)\), and the homomorphism property of \(\overline\Pi_X\).  This
packet is formal chart arithmetic; it is not a finite HN charge window,
a compact source construction, a Hall bracket computation, a Pfaffian
orientation, an \(O2\) wall atlas, a mirror discriminant, or a protected
trace.

\begin{proposition}[Compact Mukai representatives for the three simple wall charges]
\label{prop:type-ii-compact-mukai-representatives}
\textnormal{[constructed in the compact \(K3\) Mukai chart]}
Let \(S\) be a smooth projective \(K3\) surface.  Choose a closed point
\(p\in S\) and a zero-dimensional closed subscheme \(Z\subset S\) of
length \(2\).  With the convention
\[
v(E)=\ch(E)\sqrt{\td(S)},\qquad \sqrt{\td(S)}=(1,0,1),
\]
one has
\[
v(\mathcal I_Z)=(1,0,-1),\qquad
v(\mathcal I_p)=(1,0,0),\qquad
v(\mathcal O_p)=(0,0,1).
\]
Define compact Mukai-pair representatives in
\(\widetilde H(S,\Z)\oplus\widetilde H(S,\Z)\) by
\[
x_{\delta_1}:=(\mathcal I_Z,\mathcal I_p),\qquad
x_{\delta_2}:=(\mathcal I_p,\mathcal I_Z),\qquad
x_{\delta_3}:=(\mathcal I_p,\mathcal O_p).
\]
Then
\[
\Pi_X(v(x_{\delta_1}))=(1,1,0),\qquad
\Pi_X(v(x_{\delta_2}))=(0,1,1),\qquad
\Pi_X(v(x_{\delta_3}))=(0,-1,0).
\]
Thus the three triples of
Proposition~\ref{prop:type-ii-wall-images} have compact representatives
at the Mukai-chart level.  This construction does not supply retained
HN semistability, a compact Hall source basis, reduced orientations,
vanishing-cycle summands, \(E\)-quotient data, O2 wall charts, Pfaffian
signs, or protected integration.
\end{proposition}

\begin{proof}
For a zero-dimensional closed subscheme \(W\subset S\) of length \(r\),
\(\ch(\mathcal I_W)=1-r[\mathrm{pt}]\), hence
\[
v(\mathcal I_W)=(1,0,1-r).
\]
For \(r=2\) this gives \(v(\mathcal I_Z)=(1,0,-1)\), and for \(r=1\)
it gives \(v(\mathcal I_p)=(1,0,0)\).  The skyscraper sheaf has
\(\ch(\mathcal O_p)=(0,0,1)\), hence
\(v(\mathcal O_p)=(0,0,1)\).

Use the rank-one Mukai pairing
\[
\bigl((r,d,s),(r',d',s')\bigr)=2dd'-rs'-r's.
\]
The three pair representatives have Mukai vectors
\[
\begin{array}{c|cc}
&Q&P\\
\hline
x_{\delta_1}&(1,0,-1)&(1,0,0)\\
x_{\delta_2}&(1,0,0)&(1,0,-1)\\
x_{\delta_3}&(1,0,0)&(0,0,1).
\end{array}
\]
Therefore
\[
\Pi_X(Q,P)=\Bigl(\frac{(Q,Q)}2,(Q,P),\frac{(P,P)}2\Bigr)
\]
gives respectively
\[
\Bigl(\frac2 2,1,\frac0 2\Bigr)=(1,1,0),\qquad
\Bigl(\frac0 2,1,\frac2 2\Bigr)=(0,1,1),\qquad
\Bigl(\frac0 2,-1,\frac0 2\Bigr)=(0,-1,0).
\]
All three sheaves are coherent sheaves on the projective surface \(S\),
so their supports are proper over \(\C\).  The final sentence records
the extra structure not produced by the Mukai-vector calculation.
\end{proof}

The finite packet
\[
\texttt{certificates/charge/type\_ii\_compact\_mukai\_representatives}
\]
verifies
\(\texttt{TYPE\_II\_COMPACT\_MUKAI\_REPRESENTATIVES\_VERIFIED}\): the
three compact Mukai-pair representatives above have the required
Gram-degree shadows.  Its scalar firewall excludes reading these rows
as compact Hall source basis vectors or retained O2 wall objects.

\begin{remark}[Polarization defect]
\label{rem:resolution-polarization-defect-appendixC}
\ClaimStatusProved\
The composite
\(\overline\Pi_X\circ i_0=\Pi_X\) records the quadratic Gram-degree
shadow of the bare upstairs charge; the composite
\(\overline\Pi_X\circ s\equiv0\) records the split-section convention
that absorbs the entire Gram-degree into the central coordinate
\(T\). The choice of section is the choice of normal ordering: bare
placements collect cross-pairing \(B(c_i,c_j)\) in the central
coordinate (Lemma~\ref{lem:gram-additivity-defect-appendixC}), while
\(s\)-placements record \(\overline\Pi_X=0\) on every component and the
Gram-degree information lives entirely in the upstairs \(\Pi_X(c)\). The
two presentations agree on the homomorphism \(\overline\Pi_X\). This is
the additive group law on which a hypothetical compact Hall product can
descend without a chain-level associator at the lattice level.
\end{remark}

\begin{remark}[Chart data, not Hall support]
\label{rem:chart-data-not-hall-support-appendixC}
\ClaimStatusProved\
The lattice \(\widehat\Gamma_X\) is the formal additive group on which
charges and their Gram-degree shadows are tracked. It is not the
algebraic effective semigroup \(\Gamma_X^{\mathrm{eff}}\) of curve
classes representable by closed subschemes \(Y\subset X\); membership
in the effective sub-semigroup is a separate algebraic-effective
condition (Definition~\ref{def:algebraic-effective-hall-support}). The
Beilinson chart layer \(\alpha\) (charge data) and \(\gamma\) (Mukai
lattice) license \(\widehat\Gamma_X\) as ambient; the Hall-support
restriction is part of the categorical-input layer at higher Beilinson
level, not the Mukai--Gram chart datum.
\end{remark}

\subsection{Three Igusa names and the Mukai--Gram chart}
\label{subsec:three-igusa-names-chart-data}

Three Igusa names are distinct mathematical objects and remain distinct under
the Mukai--Gram chart datum:
\begin{itemize}
\item the automorphic section \(\mathcal D_X=\Delta_5\), a section of
  the weight-\(5\) automorphic line over the genus-\(2\) Siegel space
  with Maass character \(\nu_{\Delta_5}\);
\item the BKM denominator
  \(\mathrm{den}(\mathfrak g_{\Delta_5})=64^{-1}\Delta_5(2Z)\)
  on the root lattice \(\Lambda_{II}^{2,1}\); and
\item the compact protected Pfaffian
  \(\operatorname{Pf}_{\mathrm{prot}}(\mathfrak D_X)\), a section of
  \(\mathcal L^5\otimes\nu_{\Delta_5}\) on the conditional pro-object
  \(\mathfrak D_X^{\mathrm{DI}}\) of \S5.
\end{itemize}
The Mukai--Gram dictionary supplies the Gram-degree shadow
\(\Pi_X:\Gamma_X^{\mathrm{form}}\to\Gamma_{\mathrm{gram}}\) on which
the cusp-Fourier expansion of \(\Delta_5\), the root-degree map of
\(\mathrm{den}(\mathfrak g_{\Delta_5})\), and the wall data of
\(\operatorname{Pf}_{\mathrm{prot}}\) all read.
The Mukai pairing on \(\widetilde H(S,\Z)\) is symmetric in both
algebraic and transcendental parts; the ungraded Mukai-form classical
limit of an associated Yangian-type construction is the orthogonal Lie
algebra \(\mathfrak{so}(4,20)\), not the orthosymplectic
\(\mathfrak{osp}(4\mid20)\). The Hodge \(\Z/2\)-super-extension that
appears in Yangian conventions imposes a separate
\(\Z/2\)-grading on the algebraic side; it does not enter the lattice
pairing above. The Bryan/Oberdieck--Pandharipande/
Maulik--Toda conventions \cite{OPand,BRYAN,MaulikTodaGV} take \(\beta\)
algebraic and \(d\) elliptic-degree, agreeing with
Theorem~\ref{thm:d6-d2-d0-dictionary-appendixC}.

The Mukai--Gram chart datum does not touch the orientation character
\(\epsilon_o\) of
the Weyl-transport condition \ref{def:O1-plus-weyl-transport}, the Pfaffian
section \(\operatorname{Pf}_{\mathrm{prot}}(\mathfrak D_X)\), or the
recognition theorem of
Theorem~\ref{thm:pfaffian-dirac-finite-stage}; those structures live
at higher Beilinson levels and depend on the chart datum.

\section{Hall correspondences on \texorpdfstring{$\mathfrak M^+_{\mathrm{eff}}(K3\times E)$}{the moduli of effective sheaves on K3 times E}}
\label{app:hall-correspondences}

The Hall side of the Dirac--Igusa construction rests on the following
imported theorems and on their \(K3\times E\) specialization.  The Hall
correspondences are the convolution correspondences through which the
Beilinson--Drinfeld factorization structure on the elliptic carrier
acts on the retained \(K3\times E\) Hall stack \cite{BD}.  Equivalently,
they comprise the chart-augmented \(A_b\) data for the open factorization dg-category
\((X, D, \tau)\) at a chosen boundary vacuum \(b\), with chiral shadow
on the elliptic carrier supplied separately by
\S~\ref{ch:zbps-realization}.

The cited theorems fix the exact hypotheses needed for the
\(K3\times E\) specialization: finite retained substacks, proper
extension correspondences, compact-support functoriality, vanishing
cycles, Thom--Sebastiani maps, and orientation coefficients.

The standing data are: \(S\) is a smooth complex projective \(K3\)
surface, \(E\) is a smooth elliptic curve, \(X = S\times E\) is the
trivial-canonical threefold, \(\sigma_S\in H^0(S, K_S)\) is the holomorphic
two-form trivializing \(K_S\), and \(\sigma = \mathrm{pr}_S^*\sigma_S
\in H^0(X, p_S^*K_S)\) is the inflated cosection of
\S\ref{ch:zbps-realization}.  The algebraic Mukai lattice used by the
Hall support is
\[
\Lambda_{S,\mathrm{alg}} = H^0(S,\Z)\oplus
\mathrm{NS}(S)\oplus H^4(S,\Z),
\]
the algebraic sublattice of the full Mukai lattice
\(\widetilde H(S,\Z)=H^0(S,\Z)\oplus H^2(S,\Z)\oplus H^4(S,\Z)\),
endowed with the restricted Mukai pairing of \cite{Mukai1984}.
\(\Gamma_{\mathrm{eff}}\) is the algebraic/effective Hall support of
Definition~\ref{def:algebraic-effective-hall-support}; and
\(\mathfrak M^+_{\mathrm{eff}}(X)\) is the moduli stack of perfect
complexes of effective dyonic class semistable for the supplied
stability datum on \(X\).

\subsection*{D.1\quad The Hall product on the moduli of stable sheaves}
\label{app:hall-product}

The first imported theorem is the existence and associativity of a
Hall product on the equivariant cohomology of the moduli stack of
semistable sheaves, with vanishing-cycle coefficients, after Joyce and
Song.

\begin{theorem}[Joyce--Song generalised \(\mathrm{DT}\) Hall product;
\(\gamma\), proved on \(\mathrm{CY}_3\) projective; conditional on
(O1), (O1)\(^+\) on \(K3\times E\)]
\label{thm:hall-product}
Let \(Y\) be a smooth complex projective Calabi--Yau threefold.  The
moduli stack \(\mathfrak M(Y)\) of compactly supported coherent
sheaves on \(Y\) carries a critical \emph{Hall product}
\[
m\;:\;
H^\bullet_c\bigl(\mathfrak M(Y),\Phi_W\otimes o\bigr)
\otimes
H^\bullet_c\bigl(\mathfrak M(Y),\Phi_W\otimes o\bigr)
\longrightarrow
H^\bullet_c\bigl(\mathfrak M(Y),\Phi_W\otimes o\bigr),
\]
with \(\Phi_W\) the perverse sheaf of vanishing cycles of the
holomorphic Chern--Simons potential and \(o\) the orientation line
supplied by Joyce--Upmeier orientation data.  The product is defined
by pull-push along the universal extension stack
\[
\mathfrak M(Y)\xleftarrow{\;p\;}
\mathfrak{Ext}^1(\mathfrak M(Y),\mathfrak M(Y))
\xrightarrow{\;q\;}
\mathfrak M(Y),
\]
where \(\mathfrak{Ext}^1\) is the moduli stack of short exact sequences
\(0\to A\to C\to B\to 0\), \(p(0\to A\to C\to B\to 0) = (A,B)\), and
\(q(0\to A\to C\to B\to 0) = C\).  The product
\(m = q_!\circ\mathrm{TS}\circ p^*\) is associative, where
\(\mathrm{TS}\) is the Massey--Saito--Joyce Thom--Sebastiani
isomorphism on the universal extension stack
\cite[Theorem~5.4]{JoyceSong}, \cite{MasseyST}, \cite{BBDJS2015}.  Its
unit is the class of the empty object.
\end{theorem}

\begin{proof}
The product is the pull-push form of three standard structures.  The
moduli stack \(\mathfrak M(Y)\) is locally the critical locus of the Joyce--Song
holomorphic Chern--Simons potential
\(W:\mathrm{Crit}(W)\hookrightarrow Y_W\) where \(Y_W\) is a Darboux
chart of the \((-1)\)-shifted symplectic structure on \(\mathfrak M(Y)\)
supplied by Pantev--To\"en--Vaqui\'e--Vezzosi
\cite[Theorem~0.3]{PTVV2013}; the perverse sheaf \(\Phi_W\) of
vanishing cycles is well defined on the oriented d-critical locus by
Brav--Bussi--Dupont--Joyce--Szendr\H{o}i
\cite[Theorem~6.9]{BBDJS2015}.  The universal extension stack
is locally a critical locus of \(W_A\boxplus W_B\) on
\(Y_A\times Y_B\), and the Massey--Saito Thom--Sebastiani isomorphism
\cite[Theorem~1.1]{MasseyST}, \cite[Theorem~5.10]{BBDJS2015} identifies
\(\Phi_{W_A\boxplus W_B}\) with \(\Phi_{W_A}\boxtimes\Phi_{W_B}\) on
\(\mathrm{Crit}(W_A)\times\mathrm{Crit}(W_B)\).  The
multiplicative orientation cocycle of Joyce--Upmeier
\cite[Theorem~4.1]{JoyceUpmeier},
\cite[Theorem~1.11]{JoyceTanakaUpmeier} trivializes the orientation
gerbe on extensions, supplying the line bundle factor.  When the
exceptional pushforward is admissible, pulling \(\Phi_W\otimes o\)
along \(p\), tensoring by Thom--Sebastiani, and pushing along \(q_!\)
defines \(m\).  Associativity reduces, via the
two-step flag stack
\(\mathfrak{Flag}(\mathfrak M(Y),\mathfrak M(Y),\mathfrak M(Y))\), to
the pentagon coherence of Thom--Sebastiani triple stratification
\cite[Theorem~5.10(iv)]{BBDJS2015} together with the Joyce--Upmeier
multiplicative cocycle on triple extensions
\cite[\S4]{JoyceUpmeier}.  See Joyce--Song
\cite[\S 5.2,~Theorem~5.5]{JoyceSong} for the original construction
and Davison \cite[Theorem~A]{Davison2017} for the cohomological
amplification.
\end{proof}

\noindent\emph{\(K3\times E\) specialization.} On \(Y = K3\times E\),
the cosection \(\sigma = \mathrm{pr}_S^*\sigma_S\) of the
holomorphic two-form on the K3 factor reduces the obstruction
theory through Kiem--Li (Appendix~\ref{app:cosection}); the supplied
moduli is \(\mathfrak M^{\mathrm{red}}_{\mathrm{eff}}(X)\subset
\mathfrak M(X)\), and the product is the reduced critical product
\[
m^{\mathrm{red}}\;=\;q_!^{\mathrm{red}}\circ\mathrm{TS}^{\mathrm{red}}
\circ p^*
\]
of Proposition~\ref{prop:finite-dcritical-hall-product}.  Vanishing of
the seven obstruction classes
\(\mathfrak o^{\mathrm{conf}}_R\),
\(\mathfrak o^{\mathrm{prop},\mathrm{LL}}_R\),
\(\mathfrak o^{\mathrm{prop},\mathrm{mix}}_R\),
\(\mathfrak o^{\mathrm{prop},\mathrm{WW}}_R\),
\(\mathfrak o^{\mathrm{conf\text{-}prop}}_R\),
\(\mathfrak o^{\mathrm{assoc}}_{w,R}\), and
\(\mathfrak o^{I,\mathrm{prod}}_{w,R}\), together with finite residual
inertia, retained compactified flag-Quot correspondences or an
explicit compact-support six-functor model, base-change,
projection-formula, and quotient-after-correspondence coherences, is
the data \((D0)\)+\((O1)\), absorbed into \(\mathrm{Bnd}\),
\(\mathrm{dCrit}\), \(\mathrm{RedOr}\), and \(\mathrm{Six}\) of
Definition~\ref{def:compact-hall-product}.  \emph{The
existence of a chain-level Hall product on \(K3\times E\) without
those four data is open.} The resulting invariant
\(Y^+(X) := H^\bullet_c(\mathfrak M^+_{\mathrm{eff}}(X), \Phi_W\otimes o)\)
is the positive-half cohomological Hall algebra (CoHA) of
Kontsevich--Soibelman \cite{KSCoHA} restricted to effective dyonic
class; its Drinfeld double, regulated tensor product completion, and
integral form constitute the level-3 object \(G(X)\), conditional on
the Hall pairing and stable-envelope transport.  Thus \(G(X)\) arises
from \(Y^+(X)\) by Hall pairing, integral-form completion, and
stable-envelope transport, as in the Vol~III \(\Delta_5\) row
\cite{CYQGVolIII}.

\subsection*{D.2\quad The Hall coproduct via collisions}
\label{app:hall-coproduct}

The Hall product does not by itself supply a coproduct.  A collision
coproduct is additional six-functor data on the reversed extension
correspondence: properness or compact-support replacement for the
left leg, inverse Thom--Sebastiani transport, quotient descent, vacuum
unit/counit data, and the four-step flag-of-flags compatibility that
compares \(\Delta m\) with
\((m\otimes m)(1\otimes\tau\otimes1)(\Delta\otimes\Delta)\).  The
Hopf pairing and its radical quotient are separate data.  When the
coproduct and bialgebra compatibility are supplied, \(Y^+(X)\) becomes
a graded bialgebra object; after the Hopf pairing, Drinfeld doubling,
completion, integral form, and stable-envelope transport, it is a
candidate for the level-3 quantum group \(G(X)=D_\hbar(Y^+(X))\) of
Schiffmann--Vasserot type \cite{SchiffmannVasserot2013}.

\begin{theorem}[Collision coproduct datum, formal conditional]
\label{thm:hall-coproduct}
Let \(Y\) be a smooth projective Calabi--Yau threefold and suppose
the Joyce--Song Hall product of Theorem~\ref{thm:hall-product} is
given.  Assume, in addition, that the reversed extension
correspondence carries admissible compact-support six-functor data,
inverse Thom--Sebastiani transport, vacuum unit/counit data, and
three-step and four-step flag coherences for coassociativity and
bialgebra compatibility.  Then the critical Hall positive half
\(Y^+(Y) = H^\bullet_c(\mathfrak M^+_{\mathrm{eff}}(Y),\Phi_W\otimes o)\)
carries a collision coproduct
\[
\Delta\;:\;Y^+(Y)\longrightarrow Y^+(Y)\otimes Y^+(Y),
\qquad
\Delta\;=\;p_!^{\mathrm{coll}}\circ
\mathrm{TS}^{-1}\circ q^{\mathrm{coll},*},
\]
defined by the collision/extension correspondence
\[
\mathfrak M^+_{\mathrm{eff}}(Y)
\xleftarrow{\;q^{\mathrm{coll}}\;}
\mathfrak{Ext}^{1,\mathrm{coll}}(\mathfrak M^+_{\mathrm{eff}}(Y))
\xrightarrow{\;p^{\mathrm{coll}}\;}
\mathfrak M^+_{\mathrm{eff}}(Y)\times\mathfrak M^+_{\mathrm{eff}}(Y),
\]
with \(\mathfrak{Ext}^{1,\mathrm{coll}}\) the moduli of short exact
sequences \(0\to A\to C\to B\to 0\) read with \(C\) on the left and
\((A,B)\) on the right, \(q^{\mathrm{coll}}(0\to A\to C\to B\to 0) = C\),
\(p^{\mathrm{coll}}(0\to A\to C\to B\to 0) = (A,B)\).  The
coproduct \(\Delta\) is graded coassociative.  Together with the Hall
product \(m\) of Theorem~\ref{thm:hall-product},
\((Y^+(Y), m, \Delta)\) is a graded bialgebra.
\end{theorem}

\begin{proof}
The formula is the formal reverse of the Hall product.  Its existence
requires the extra six-functor hypotheses because the reversed leg is
not automatically proper in the compact threefold problem.
Given those hypotheses, coassociativity is the inverse
Thom--Sebastiani pentagon on the three-step flag stack.  The
compatibility of \(m\) and \(\Delta\) is not a consequence of the
product theorem alone; it is the equality of the two pull-push functors
on the four-step flag-of-flags common refinement.  On quivers with
potential this is the Davison--Meinhardt bialgebra theorem
\cite[Theorem~A]{DavisonMeinhardtCoHA}; on the compact
\(K3\times E\) threefold it is the finite correspondence datum of
Definition~\ref{def:compact-geometric-hall-bialgebra-datum} and
Proposition~\ref{prop:geometric-hall-data-give-bialgebra-defects}.
\end{proof}

\noindent\emph{Drinfeld double.} The Hall pairing
\(\langle\,\cdot\,,\,\cdot\,\rangle_{\mathrm{Hall}}\), when supplied
and non-degenerate after the completed radical quotient, pairs
\(Y^+\) with itself; the Drinfeld double \(D_\hbar(Y^+(X))\) is the
quantum double in the sense of Drinfeld \cite[\S 13]{KAC:1}, supplying
the negative half \(Y^-(X)\) and the Cartan generators \(K_\gamma\),
\(\gamma\in\Gamma_{\mathrm{eff}}\), with relations encoding the Hall
pairing.  The \emph{level-3 \(K3\times E\) quantum group} is then
\[
G(X)\;:=\;D_\hbar(Y^+(X))_{\mathrm{compl}}^{\mathrm{int}}
\]
where \((-)_{\mathrm{compl}}^{\mathrm{int}}\) is the integral-form
completion under the Hall pairing, conditional on the
stable-envelope transport of \S~D.6.  In Schiffmann--Vasserot
\cite{SchiffmannVasserot2013} the analogous identification on
toric \(\mathrm{CY}_3\) ambient is established;
\(K3\times E\) is not equipped here with a global toric or
quiver-variety NCCR model of that kind: there is no finite
fixed-point McKay presentation, no compatible product-polarized HPD
self-dual model, and no Schiffmann--Vasserot stable-envelope atlas on
the compact non-toric ambient.  The cosection-twisted reduced
obstruction theory of \S~D.4 is the replacement; the global
\(G(X)\) on the compact non-toric \(K3\times E\) is not constructed
from these data without the Hall pairing, integral form, completion,
and stable-envelope transport.

\subsection*{D.3\quad The Hall primitives and the formal recognition criterion}
\label{app:hall-primitives}

For a connected graded bialgebra \(B = \bigoplus_{c}B_c\), the
primitive elements at charge \(c\) are
\[
\mathrm{Prim}_c(B)\;=\;
\{x\in B_c\mid \Delta(x) = x\otimes 1 + 1\otimes x\},
\]
equivalently the kernel of the reduced coproduct \(\bar\Delta\).  On a
graded connected cocommutative Hopf algebra in characteristic zero,
the Milnor--Moore theorem \cite[Theorem~5.18]{LurieHA} identifies
\(B\simeq U(\mathrm{Prim}_\bullet(B))\) as Hopf algebras, where
\(\mathrm{Prim}_\bullet(B)\) is the Lie algebra of primitives under
the bracket \([x,y] = xy - (-1)^{|x|\,|y|}yx\).

\begin{theorem}[Hall primitives, formal level; \(\beta\), formal]
\label{thm:hall-primitives-formal}
Let \(\mathfrak M\) carry a critical Hall positive half
\(Y^+(Y) = (B, m, \Delta)\) of
Theorem~\ref{thm:hall-coproduct}.  The graded subspace
\[
\mathrm{Prim}_\bullet(Y^+(Y))\;=\;
\bigoplus_{c\in\Gamma_{\mathrm{eff}}}
\mathrm{Prim}_c(Y^+(Y))
\]
is a graded Lie subalgebra of \(Y^+(Y)\) under the antisymmetrised
product \([\,\cdot\,,\,\cdot\,]\).  At a charge \(c\) such that the
moduli \(\mathfrak M^+_c(Y)\) is connected and the \(C\)-component of
\(\Delta\) on \(\mathfrak M^+_c(Y)\) factors through extensions whose
graded pieces lie in strictly lower charges, the inclusion
\(\mathrm{Prim}_c\hookrightarrow Y^+_c\) is the kernel of
\(\bar\Delta_c\).  When \(Y^+(Y)\) is graded connected and
cocommutative, Milnor--Moore identifies
\(Y^+(Y)\simeq U(\mathrm{Prim}_\bullet(Y^+(Y)))\) as graded Hopf
algebras.
\end{theorem}

\begin{proof}
Primitivity of \(\mathrm{Prim}_\bullet\) under the antisymmetrised
product is verified by direct computation of
\(\Delta([x,y])\) using the cocommutative-up-to-twist axiom and the
multiplicativity of \(\Delta\); the resulting expression collapses to
\([x,y]\otimes 1 + 1\otimes [x,y]\).  On a connected graded
cocommutative Hopf algebra in characteristic zero, the Milnor--Moore
identification is \cite[\S 21.1]{LurieHA}; on graded connected
non-cocommutative Hopf algebras, the Cartier--Quillen--Milnor--Moore
theorem \cite[Theorem~21.2.1]{LurieHA} requires either
cocommutativity or a chain-level lift through bar / cobar, which is
exactly the chiral Koszul separation theorem of Vol I
\cite{FrancisGaitsgoryCKD}.
\end{proof}

\noindent\emph{Recognition.}
The recognition theorem
of Theorem~\ref{thm:primitive-recognition} identifies the completed
formal primitive Lie superalgebra with the imported
Borcherds--Kac--Moody Lie superalgebra
\(\mathfrak g_{\Delta_5}\) of
Theorem~\ref{thm:bkm-denominator-algebra-identity}, conditional on the
finite compact-source datum \((S1)\)--\((S10)\): the completed PBW
filtration, the Hopf pairing radical equality, and the parity
dimension recovery.  In
particular, on \(K3\times E\) the formal primitive subspace
\(\mathrm{Prim}_\bullet(Y^+(X))\) is the \emph{candidate} primitive
Lie superalgebra; its identification with the imported BKM
denominator algebra \(\mathfrak g_{\Delta_5}\) is the primitive
recognition part of View \(\textcircled{3}\).  The formal Hopf algebra
structure alone gives no recognition theorem; recognition acts on it.

\subsection*{D.4\quad The cosection-twist on \texorpdfstring{$K3\times E$}{K3 x E}: Kiem--Li localisation}
\label{app:cosection}

On \(K3\times E\), the holomorphic two-form on the K3 factor supplies
the semiregularity cosection on the retained nonzero K3-class
branches of the obstruction theory of \(\mathfrak M(X)\).  The
Kiem--Li \cite{KiemLi2013} localisation theorem then localises the
virtual fundamental class to the zero-locus of the cosection when the
surjectivity hypothesis holds.  The localisation is the source of the
\emph{reduced} virtual class on which the Hall product of
\S~D.1 carries chain-level coherence.

\begin{definition}[Cosection]
\label{def:cosection-app-D}
Let \(\mathfrak M\) be a Deligne--Mumford stack equipped with a
perfect obstruction theory \(\phi:\mathbb E\to\mathbb L_{\mathfrak M}\)
of amplitude \([-1,0]\), and let \(\mathrm{Ob}(\mathfrak M) =
h^1(\mathbb E^\vee)\) be the coherent obstruction sheaf.  A
\emph{cosection} of the obstruction theory is an
\(\mathcal O_{\mathfrak M}\)-linear morphism
\[
\sigma\;:\;\mathrm{Ob}(\mathfrak M)\longrightarrow
\mathcal O_{\mathfrak M};
\]
its \emph{vanishing locus} is the closed substack
\(\mathfrak M^{\mathrm{red}} = (\sigma = 0)\) on which the obstruction
sheaf restricts to the kernel \(\ker\sigma\subset\mathrm{Ob}(\mathfrak M)\).
\end{definition}

\begin{theorem}[Kiem--Li cosection localisation;
\cite{KiemLi2013}~Theorem~5.0.1; \(\gamma\), proved]
\label{thm:cosection-localisation}
Let \(\mathfrak M\) be a Deligne--Mumford stack with perfect
obstruction theory \(\phi\) and surjective cosection
\(\sigma:\mathrm{Ob}(\mathfrak M)\twoheadrightarrow\mathcal O_{\mathfrak M}\)
on \(\mathfrak M\setminus\mathfrak M^{\mathrm{red}}\), where
\(\mathfrak M^{\mathrm{red}} = (\sigma = 0)\).  The virtual
fundamental class \([\mathfrak M]^{\mathrm{vir}}\in
A_{\mathrm{vd}}(\mathfrak M)\) lifts to a cosection-localised class
\[
[\mathfrak M]^{\mathrm{vir}}_{\mathrm{loc}}
\;\in\;A_{\mathrm{vd}}(\mathfrak M^{\mathrm{red}})
\]
satisfying \(\iota_*[\mathfrak M]^{\mathrm{vir}}_{\mathrm{loc}} =
[\mathfrak M]^{\mathrm{vir}}\), where
\(\iota:\mathfrak M^{\mathrm{red}}\hookrightarrow\mathfrak M\) is the
inclusion.  If the cosection is used separately to quotient a trivial
obstruction direction and form a reduced obstruction theory, that
reduced theory has virtual dimension
\(\mathrm{vd}^{\mathrm{red}}=\mathrm{vd}+1\); this dimension shift is
not the Kiem--Li localized cycle itself.
\end{theorem}

\begin{proof}
The cosection-kernel cone \(C(\sigma) = \ker(\sigma:\mathrm{Ob}\to
\mathcal O)\) is a perfect obstruction cone of amplitude \([-1,0]\)
on \(\mathfrak M^{\mathrm{red}}\) with
\(\mathrm{rk}\,C(\sigma) = \mathrm{rk}\,\mathrm{Ob} - 1\).  The
intrinsic normal cone \(\mathfrak C\) of \(\mathfrak M\) sits in the
total obstruction bundle on \(\mathfrak M\setminus\mathfrak M^{\mathrm{red}}\),
where \(\sigma\) is surjective and the cone has codimension one in
\(\mathrm{Ob}\); on \(\mathfrak M^{\mathrm{red}}\), the cone embeds
into \(C(\sigma)\) by \cite[Proposition~4.3]{KiemLi2013}.
Intersection with the zero section of \(C(\sigma)\) on
\(\mathfrak M^{\mathrm{red}}\), pushed forward to \(\mathfrak M\)
through \(\iota\), yields \([\mathfrak M]^{\mathrm{vir}}\) by
\cite[Theorem~5.0.1]{KiemLi2013}.  The localized Gysin map preserves
the original virtual dimension.  The shift
\(\mathrm{vd}^{\mathrm{red}} = \mathrm{vd} + 1\) enters only after a
separate reduced obstruction theory is formed by removing the trivial
obstruction summand.
\end{proof}

\noindent\emph{\(K3\times E\) cosection.} On \(X = K3\times E\), the
cosection of the obstruction theory of \(\mathfrak M(X)\) is induced
by the pullback \(\sigma=\mathrm{pr}_S^*\sigma_S\) of the K3
holomorphic two-form.
Concretely, for a perfect complex \(F\) on \(X\), the
chain-level cosection acts on \(\mathrm{Ext}^2(F,F)\) by
the semiregularity map of Buchweitz--Flenner /
Bloch \cite{MaulikTodaGV},
\[
\mathrm{cs}_F\;:\;\mathrm{Ext}^2(F,F)\longrightarrow
\C,
\qquad
\mathrm{cs}_F(\xi)\;=\;\int_X
\sigma\wedge\mathrm{Tr}\bigl((\xi\otimes 1_{\Omega_X^1})\circ
\mathrm{At}(F)\bigr),
\]
where \(\mathrm{At}(F)\in\mathrm{Ext}^1(F, F\otimes\Omega^1_X)\) is
the Atiyah class.  The cosection vanishes precisely on the
\emph{reduced} moduli
\(\mathfrak M^{\mathrm{red}}(X)\subset\mathfrak M(X)\) consisting of
sheaves whose Atiyah class is annihilated by \(\sigma\); equivalently,
sheaves whose deformation space along the K3 direction is
constrained.  See Maulik--Toda
\cite[Lemma~2.6, Theorem~2.7]{MaulikTodaGV} and Oberdieck
\cite[Theorem~1, Theorem~3]{OberdieckReducedPT} for the explicit
identification on the DT/PT side.  Surjectivity is an additional
row-level condition: the retained stratum must be identified with the
nonzero K3-class branch and must record a cokernel-rank-zero
semiregularity row.  On the \(s=0\) Hall cusp the retained charge can
lose the K3 component that makes the semiregularity map surjective; the
holomorphic two-form \(\sigma_S\) does not vanish.  The localisation then passes to the
\(D0\)-degeneration limit named at obstruction \((D0)\) of
\S~\ref{ch:zbps-realization}.

\begin{theorem}[Reduced virtual class on \(K3\times E\); primitive
stable-pair branch]
\label{thm:reduced-virtual-class-K3xE}
On the stable-pair or ideal-sheaf Deligne--Mumford branch carrying a
perfect obstruction theory and a surjective K3 semiregularity
cosection, Kiem--Li localization gives
\[
[\mathfrak M_c(X)]^{\mathrm{vir}}_{\mathrm{loc}}\in
A_{\mathrm{vd}}(\mathfrak M_c(X)^{\mathrm{red}}).
\]
The reduced obstruction theory obtained by quotienting the trivial
obstruction direction gives a distinct reduced class
\[
[\mathfrak M_c(X)]^{\mathrm{red}}\in
A_{\mathrm{vd}+1}(\mathfrak M_c(X)^{\mathrm{red}}).
\]
The reduced DT/PT theory of the primitive nonzero K3-class branch uses
this reduced class and gives the Oberdieck--Pixton scalar
\(Z^X_{\mathrm{OP}}=-4096\,\Delta_5^{-2}\)
\cite[Theorem~3]{OberdieckReducedPT}.  The semistable Hall stack
requires the retained compact-support six-functor and d-critical Hall
datum of \S\ref{subsec:hall-correspondences}.
On the \(s = 0\) cusp the cosection degenerates; \emph{the
cosection-reduced theory there is the \((D0)\) degeneration limit
named at \S~\ref{ch:zbps-realization} and controlled by the retained
D0-HN criterion of
Theorem~\ref{thm:G-D0}}.\quad[\(\delta\), \((D0)\)]
\end{theorem}

\subsection*{D.5\quad BBDJS reduced d-critical orientation}
\label{app:bbdjs}

The Hall product of \S~D.1 requires not just a virtual fundamental
class but a perverse-sheaf-of-vanishing-cycles refinement and an
orientation gerbe; the source is the Brav--Bussi--Dupont--Joyce--Szendr\H{o}i
\cite{BBDJS2015} extension of the Joyce
\cite{JoyceSong} d-critical-locus framework.

\begin{theorem}[BBDJS perverse-sheaf-of-vanishing-cycles; oriented
\(\mathrm{CY}_3\) version of \cite{BBDJS2015}~Theorem~6.9; \(\gamma\), proved]
\label{thm:bbdjs}
Let \(\mathfrak M\) be a derived Artin stack with a \((-1)\)-shifted
symplectic structure (PTVV \cite{PTVV2013}) whose classical truncation
is a d-critical locus in the sense of Joyce
\cite[Definition~2.5]{JoyceSong}.  Suppose \(\mathfrak M\) carries an
\emph{orientation}, i.e.~a square root \(o\) of the Joyce determinant
line \(K_{\mathfrak M}^{\mathrm{vir}} = (\det\mathbb L_{\mathfrak M})^{1/2}\),
in the sense of \cite[Theorem~6.9]{BBDJS2015}.  There is a perverse
sheaf
\[
\Phi_{\mathfrak M}\;\in\;\mathrm{Perv}(\mathfrak M),
\]
the \emph{perverse sheaf of vanishing cycles} of the d-critical
locus; locally on a Joyce--Song Darboux chart \(\mathrm{Crit}(W)\)
inside a smooth ambient \(Y_W\), \(\Phi_{\mathfrak M}|_{\mathrm{Crit}(W)}
\simeq\Phi_W[d_{Y_W}]\) is the standard perverse sheaf of vanishing
cycles of the holomorphic potential \(W:Y_W\to\C\), shifted by the
ambient dimension.  The chart-comparison cocycle is canonical up to
the orientation choice; the orientation determines a multiplicative
cocycle of Joyce--Upmeier \cite[Theorem~4.1]{JoyceUpmeier} on the
moduli of objects.
\end{theorem}

\begin{proof}
Three layers stack.  PTVV \cite[Theorem~0.3]{PTVV2013} produces the
\((-1)\)-shifted symplectic form on the derived moduli of
\(\mathrm{CY}_3\) coherent sheaves; Bussi--Brav--Joyce
\cite[\S 5]{BBDJS2015} produce the d-critical-locus chart structure
on the classical truncation; Bussi--Brav--Dupont--Joyce--Szendr\H{o}i
\cite[Theorem~6.9]{BBDJS2015} produce the perverse sheaf of
vanishing cycles with its chart cocycle.  The orientation is the datum needed
to unify chart-local \(\Phi_W\) into a global perverse sheaf
\(\Phi_{\mathfrak M}\); without it, the chart-cocycle has a sign
ambiguity that obstructs the descent to a perverse sheaf
\cite[\S 6.4]{BBDJS2015}.  Joyce--Upmeier
\cite[Theorem~1.11]{JoyceUpmeier} construct the multiplicative orientation
cocycle on the universal extension stack; this is the
\(\mathrm{CY}_3\) projective theorem for the retained Hall orientation.
\end{proof}

\noindent\emph{Reduced refinement on \(K3\times E\).}
On \(X = K3\times E\), the \((-1)\)-shifted symplectic structure
\(\omega_{\mathfrak M}\) on \(\mathfrak M(X)\) does not restrict to a
\((-1)\)-shifted symplectic structure on
\(\mathfrak M^{\mathrm{red}}(X)\); the cosection-twist of \S~D.4
contracts \(\omega_{\mathfrak M}\) by \(\sigma\), and the
\emph{reduced} \((-1)\)-shifted form
\(\omega_{\mathfrak M}^{\mathrm{red}}\) is non-degenerate only after
restriction to the cosection vanishing locus \cite[\S 2.6]{MaulikTodaGV}.
The corresponding reduced d-critical-locus structure is supplied by
the Maulik--Toda \cite[Theorem~2.7]{MaulikTodaGV} extension; the
reduced perverse sheaf of vanishing cycles
\(\Phi_{\mathfrak M}^{\mathrm{red}}\) and the reduced orientation
line \(o^{\mathrm{red}}_{\mathcal C}\) are exactly the data
\((O1)\) of \S~\ref{ch:zbps-realization}.  In
the notation of (P4), the reduced
orientation is twisted by a parity-flip line \(\Pi\) generated by the
cosection cokernel,
\[
o^{\mathrm{red}}_{\mathcal C}\;=\;
o^{\mathrm{std}}_{\mathcal C}\otimes\Pi,
\]
absorbing the dimension shift \(\mathrm{vd}^{\mathrm{red}} =
\mathrm{vd}^{\mathrm{std}} + 1\) into a \(\Z/2\)-grading shift on the
oriented BBDJS line.

\begin{proposition}[BBDJS reduced d-critical orientation on
\(K3\times E\); \(\delta\), \((O1)+(O1)^+\)]
\label{prop:bbdjs-reduced-K3xE}
Let \(\mathfrak M^{\mathrm{red},+}_c(X)\) be the cosection-reduced
moduli of semistable sheaves of effective dyonic class \(c\) on
\(X = K3\times E\).  The hypotheses \emph{strong reduced orientation
on the \(K3\times E\)-quotient self-Ext frame bundle (O1) and
Weyl-equivariant transport along \(W^{(2)}(\Lambda_{II}^{2,1})\)
reflections (O1)\(^+\)} of
\S~\ref{ch:zbps-realization} are exactly the
following.  There exist a multiplicative orientation cocycle of
Joyce--Upmeier type extending equivariantly across the reduced
\(E\)-quotient, the \(\sigma_S\)-cosection-induced parity flip line
\(\Pi\), and a connected \(BE\)-equivariant trivialisation of the
small-orbit gerbe class \(\overline\beta^E_{\mathcal C}\) on the
finite-stabilizer strata, such that
\[
\Phi^{\mathrm{red}}_{\mathfrak M^{\mathrm{red},+}_c(X)}
\otimes
o^{\mathrm{red}}_c
\]
descends through the free \(E\)-translation quotient and trivializes
the Pin\(^c\)-lift gerbe class of the reduced derived self-Ext frame
bundle.\quad The strong reduced orientation \((O1)\) is established
from the retained quotient-orientation datum.  Joyce--Upmeier
\cite[Theorem~1.11]{JoyceUpmeier} supplies the projective
Calabi--Yau-threefold orientation technology used here; the
quasi-projective and quotient descent belongs to the retained
\((O1_{\mathrm{quot}})\) datum.
The Cao--Kool--Monavari CY\(_4\) theorem
\cite[Theorem~1.6]{CaoKoolMonavari} is a fourfold analogue, not an
ingredient of the \(K3\times E\) CY\(_3\) reduction.  The free
\(E\)-quotient, finite-stabilizer, and Weyl-transport theorem on
\(K3\times E\) follows only relative to the retained
quotient-orientation datum and the chosen Weyl lifts of
Theorems~\ref{thm:G-O1} and~\ref{thm:G-O1-plus}.
\end{proposition}

\begin{proof}
The construction of the reduced perverse sheaf
\(\Phi^{\mathrm{red}}_{\mathfrak M^{\mathrm{red},+}_c(X)}\) on the
cosection-reduced moduli is by Maulik--Toda
\cite[Theorem~2.7]{MaulikTodaGV}; the orientation gerbe lifts to the
reduced Joyce d-critical locus by the BBDJS
\cite[\S 6.4]{BBDJS2015} chart-cocycle calculation.  The
multiplicative \emph{equivariant} orientation cocycle across the
\(E\)-quotient and across the finite-stabilizer locus
\(BE[N]\)-strata is the joint vanishing of \((O1)\) and \((O1)^+\):
on the connected \(BE\)-equivariant trivialisation of
\(\alpha^E_{\mathcal C}\in H^2_E(\mathfrak M^{\mathrm{red}}_{\mathcal C};\mathbb F_2)\)
of Definition~\ref{def:O1-plus-weyl-transport} it suffices to trivialize the
small-orbit gerbe class \(\overline\beta^E_{\mathcal C}\) on
\(BE[N]\)-strata, governed by the Klein-four
detection criterion of Lemma~\ref{lem:G-Klein-four}.  The Weyl-equivariant transport across
\(W^{(2)}(\Lambda_{II}^{2,1})\) reflections is the orientation
character calculation of
Theorem~\ref{thm:pfaffian-dirac-finite-stage}
on the three type-II walls.  Theorems~\ref{thm:G-O1}
and~\ref{thm:G-O1-plus} supply these transports on \(K3\times E\);
the Pin\(^c\)-lift theory of
\cite{JoyceUpmeier,JoyceTanakaUpmeier} supplies the analytic
technology in the Calabi--Yau-threefold setting.
\end{proof}

\subsection*{D.6\quad Stable-envelope transport: from Hall to Yangian}
\label{app:stable-envelope}

The level-3 quantum-group structure on \(G(X) = D_\hbar(Y^+(X))\) of
\S~D.2 is, on toric ambient or quiver varieties, identified by
Maulik--Okounkov \cite{MaulikOkounkov} with the geometric \(R\)-matrix
of the \emph{stable envelope}, providing the link from the Hall side
to the Yangian / quantum-group side.

\begin{definition}[Stable envelope]
\label{def:stable-envelope}
Let \(\mathfrak M\) be a smooth symplectic variety with a torus
\(T\)-action and a finite \(T\)-fixed locus \(\mathfrak M^T\); let
\(\mathcal C\subset\mathrm{Lie}(T)\) be a chamber of the
\(T\)-equivariant cohomology of \(\mathfrak M\).  The
\emph{stable envelope} \cite[\S 3]{MaulikOkounkov} is the
\(T\)-equivariant cohomological correspondence
\[
\mathrm{Stab}_{\mathcal C}\;:\;
H^\bullet_T(\mathfrak M^T)\longrightarrow
H^\bullet_T(\mathfrak M),
\]
characterised by support, normalization, and degree axioms
\cite[Theorem~3.3.4]{MaulikOkounkov}.  The geometric \(R\)-matrix is
the comparison
\[
R^{\mathrm{MO}}_{\mathcal C, \mathcal C'}\;=\;
\mathrm{Stab}_{\mathcal C'}^{-1}\circ\mathrm{Stab}_{\mathcal C}
\;\in\;\mathrm{End}(H^\bullet_T(\mathfrak M^T))
\]
across two chambers \(\mathcal C, \mathcal C'\); it satisfies the
Yang--Baxter equation by \cite[Theorem~4.6.1]{MaulikOkounkov}.
\end{definition}

\begin{theorem}[Maulik--Okounkov; \cite{MaulikOkounkov}~Theorem~10.2.1;
\(\beta\), proved on quiver varieties; conditional on \(K3\times E\)
compact non-toric]
\label{thm:maulik-okounkov}
Let \(\mathfrak M = \mathfrak M_{Q,\mathbf v,\mathbf w}\) be a Nakajima
quiver variety \cite{Nakajima1994,Nakajima1998} attached to a finite
quiver \(Q\) with dimension vectors \(\mathbf v, \mathbf w\).  The
stable envelope \(\mathrm{Stab}_{\mathcal C}\) makes
\(H^\bullet_T(\mathfrak M_{Q,\mathbf v,\mathbf w})\) a module over a
Yangian \(Y_\hbar(\mathfrak g_Q)\) attached to \(Q\), with
\(\mathfrak g_Q\) the Kac--Moody Lie algebra of the quiver, and the
geometric \(R\)-matrix \(R^{\mathrm{MO}}\) is the rational
\(R\)-matrix of \(Y_\hbar(\mathfrak g_Q)\).
\end{theorem}

\begin{proof}
The Maulik--Okounkov construction proceeds through:
(i) the existence and uniqueness of the stable envelope (axioms of
support, polarisation, normalization), \cite[\S 3.3]{MaulikOkounkov};
(ii) the Yang--Baxter equation as a chamber-wall cocycle,
\cite[\S 4]{MaulikOkounkov}; (iii) the identification of the Bethe
algebra generated by the matrix elements of \(\mathrm{Stab}\) with a
Yangian, \cite[\S 5--10]{MaulikOkounkov}.  The Schiffmann--Vasserot
identification \(\mathrm{CoHA}(\C^3) = Y^+(\widehat{\mathfrak{gl}}_1)\)
of \cite[Theorem~A]{SchiffmannVasserot2013} is the affine \(A_0\)
specialization; \(\mathcal W_{1+\infty}\) arises \emph{only after}
Drinfeld doubling and Fock-evaluation.  The three-arrow chain
\(\mathrm{CoHA} = Y^+\hookrightarrow Y\xrightarrow{\mathrm{ev}_\lambda}
\mathrm{End}(\mathcal W_{1+\infty}[\lambda]\text{-vac})\) of
\cite{KSCoHA}, \cite{SchiffmannVasserot2013} traverses the level-2
\(\to\) level-3 \(\to\) level-3-evaluated passage; the named identity
on each arrow measures the structural distance from
\(\mathrm{CoHA}(\C^3)\) to \(\mathcal W_{1+\infty}\).
\end{proof}

\noindent\emph{\(K3\times E\) specialization: open.}
On \(X = K3\times E\), the moduli \(\mathfrak M^+_{\mathrm{eff}}(X)\)
is not a Nakajima quiver variety, the ambient is compact non-toric,
and the global stable envelope is unconstructed.  The five
obstructions to a global toric/quiver-variety NCCR model of the
Schiffmann--Vasserot type obstruct the literal
\cite{MaulikOkounkov} construction.  The cosection-twist of \S~D.4 supplies the
\emph{Serre-equivariant reduced substitute}: the
cosection-reduced obstruction theory of
Theorem~\ref{thm:reduced-virtual-class-K3xE} restores a
\(\mathrm{CY}_3\)-symmetric reduced obstruction theory on the
\(s\ne 0\) sector, and the Hall-to-Yangian comparison is then
expected at the level of the reduced theory and conditional on the
stable-envelope transport, after the retained D0-HN datum, retained
quotient-orientation datum, chosen Weyl lifts, and retained
\((O2_{\mathrm{atlas}})\) wall datum of
Appendix~\ref{app:obstruction-discharges}.

The Hall--Yangian comparison on \(K3\times E\) is therefore
conjectural:
\[
G(X)\;\overset{?}{=}\;
D_\hbar(\mathrm{CoHA}^{\mathrm{red},+}_{K3\times E})_{\mathrm{compl}}^{\mathrm{int}}
\;\overset{?}{=}\;
Y_\hbar(\widehat{\mathfrak g}_{\Delta_5})
\]
where \(\widehat{\mathfrak g}_{\Delta_5}\) is the central extension of
the imported BKM Lie superalgebra of
Theorem~\ref{thm:bkm-denominator-algebra-identity}, and the equality is
conditional on a \(K3\times E\) stable-envelope transport.  The
Hall--Drinfeld double \(\mathcal D_\hbar(\mathrm{CoHA}_{K3\times E})\)
of Kontsevich--Soibelman \cite{KSCoHA} places this identification at
the same level-3 quantum group seen from the geometric side; it is
the open compact non-toric instance of the chain-fusion conjecture.

\subsection*{D.7\quad Level separations}
\label{app:hall-level-separations}

The Beilinson cut separates the following objects.

\begin{enumerate}
\item \(Y^+(X)\) is the cohomological-Hall positive half
\(H^\bullet_c(\mathfrak M^+_{\mathrm{eff}}(X), \Phi_W\otimes o)\) with
the Joyce--Song product/coproduct of
Theorems~\ref{thm:hall-product},~\ref{thm:hall-coproduct}; \(G(X) =
D_\hbar(Y^+(X))_{\mathrm{compl}}^{\mathrm{int}}\) is the Drinfeld
double after Hall pairing, completion, integral form, stable-envelope
transport, and descent.

\item \(\mathrm{CoHA}(\C^3)=Y^+(\widehat{\mathfrak{gl}}_1)\)
\cite[Theorem~A]{SchiffmannVasserot2013}; \(\mathcal W_{1+\infty}\)
appears only after Drinfeld doubling and Fock evaluation.

\item No global Schiffmann--Vasserot/Nakajima atlas is known for compact non-toric
\(K3\times E\).  The cosection-twist of \S~D.4 supplies the
Serre-equivariant reduced substitute used here.

\item \(\mathrm{Prim}_\bullet(Y^+(X))\) is the formal primitive Lie
superalgebra of \(Y^+(X)\); its identification with
\(\mathfrak g_{\Delta_5}\) is the primitive-recognition theorem of
Theorem~\ref{thm:primitive-recognition}, conditional on the finite
compact-source datum \((S1)\)--\((S10)\).

\item BBDJS \cite[Theorem~6.9]{BBDJS2015} requires an
\emph{oriented} d-critical locus; the orientation is the square root
of the Joyce determinant line, and the multiplicative cocycle is
Joyce--Upmeier \cite[Theorem~4.1]{JoyceUpmeier}.

\item Kiem--Li \cite[Theorem~5.0.1]{KiemLi2013} requires the cosection to be
surjective off the localised locus; on \(K3\times E\), the cosection
degenerates on the \(s = 0\) cusp, which is the \((D0)\)
degeneration limit controlled by the retained D0-HN criterion of
Theorem~\ref{thm:G-D0}.

\item The cosection contracts the \((-1)\)-shifted symplectic form;
\(\omega^{\mathrm{red}}\) is non-degenerate only after
restriction to the cosection vanishing locus
\cite[\S 2.6]{MaulikTodaGV}, and the orientation line carries the
\(\Z/2\)-parity flip \(\Pi\) absorbing the dimension shift
\(\mathrm{vd}^{\mathrm{red}} = \mathrm{vd}^{\mathrm{std}} + 1\).

\item \(K3\times E\) is a trivial-canonical threefold.  The scalar-trace
index is \(n=\chi(\mathcal O_Y)\) on the
closed target, whereas
\(\chi(\mathcal O_{K3\times E})=\chi(\mathcal O_{K3})\chi(\mathcal O_E)
=2\cdot0=0\), and the BBDJS / Joyce--Upmeier
machinery applies in the \(\mathrm{CY}_3\) form
\cite[\S 6]{BBDJS2015}, not in the \(\mathrm{CY}_4\) Borisov--Joyce
\cite{BorisovJoyce} or Oh--Thomas \cite{OhThomasI,OhThomasII} form.
\end{enumerate}

\subsection*{D.8\quad Four local obstruction calculations}
\label{app:hall-local-obstruction-calculations}

Four local calculations enter the Hall correspondence.

\textbf{(F1)\quad Hall associativity.} The Joyce--Song associativity
of Theorem~\ref{thm:hall-product} requires the pentagon coherence of
Thom--Sebastiani triple stratification
\cite[Theorem~5.10(iv)]{BBDJS2015}.  The
chart-cocycle of the d-critical locus carries the orientation sign
whose pentagon coherence is measured by the BBDJS orientation gerbe.
That gerbe \cite[\S 6.4]{BBDJS2015} trivializes
exactly this ambiguity; the multiplicative cocycle of Joyce--Upmeier
\cite[\S 4]{JoyceUpmeier} extends the trivialisation to the universal
extension stack.  On \(K3\times E\), the gerbe descends through the
free \(E\)-quotient and trivializes the small-orbit class by
Theorems~\ref{thm:G-O1} and~\ref{thm:G-O1-plus}.

\textbf{(F2)\quad Cosection surjectivity.} Kiem--Li
\cite[Theorem~5.0.1]{KiemLi2013} requires
\(\sigma:\mathrm{Ob}\twoheadrightarrow\mathcal O\) on
\(\mathfrak M\setminus\mathfrak M^{\mathrm{red}}\).  On the \(s = 0\)
cusp of the dyonic charge lattice of \(K3\times E\),
the K3-class component needed for semiregularity can disappear and
surjectivity fails, while the K3 holomorphic two-form itself remains
nonzero.  The smooth primitive nonzero K3-class branch satisfies
surjectivity by Maulik--Toda
\cite[\S 2]{MaulikTodaGV}; the cusp branch is the \((D0)\)
degeneration limit named at obstruction \((D0)\) of
\S~\ref{ch:zbps-realization}, controlled by the retained D0-HN
criterion of Theorem~\ref{thm:G-D0}.
The Oberdieck--Pixton scalar
\(Z^X_{\mathrm{OP}} = -4096\,\Delta_5^{-2}\) is proved on the
\(s\ne 0\) branch \cite[Theorem~3]{OberdieckReducedPT}; on the cusp
the chain-level theory degenerates to the Borcherds
imaginary-simple-root contribution.

\textbf{(F3)\quad BBDJS orientation existence on compact non-toric.}
The Joyce--Upmeier multiplicative cocycle \cite[Theorem~4.1]{JoyceUpmeier}
is established on the \(\mathrm{CY}_3\) projective ambient; its
extension across the cosection-reduced quotient ambient is part of the
retained \((O1_{\mathrm{quot}})\) datum.  On the \(K3\times E\)
quotient by free \(E\)-translation, descent of the orientation line
requires the small-orbit gerbe class on \(BE[N]\)-strata to
trivialize.  The joint vanishing of the connected
\(BE\)-equivariant degree-two obstruction
\(\alpha^E_{\mathcal C}\in H^2_E(\mathfrak M^{\mathrm{red}}_{\mathcal C};\mathbb F_2)\)
the finite-stabilizer Cech--Borel gerbe classes, and the residual
degree-one stabilizer characters is exactly \((O1)\);
Theorem~\ref{thm:G-O1} of Appendix~\ref{app:obstruction-discharges}
is the quotient-orientation criterion relative to this retained datum.
The remaining open data here are the retained quotient orientation and
the compact-source recognition rows.

\textbf{(F4)\quad \((-1)\)-shifted vs \((-2)\)-shifted symplectic.}
\(K3\times E\) is a trivial-canonical threefold; the moduli
\(\mathfrak M(K3\times E)\) carries a \((-1)\)-shifted symplectic
structure by PTVV \cite[Theorem~0.3]{PTVV2013}.  The Borisov--Joyce
\cite{BorisovJoyce} and Oh--Thomas \cite{OhThomasI,OhThomasII}
theories are \(\mathrm{CY}_4\) theories.  The Calabi--Yau dimension
of \(X = K3\times E\) is \(\dim_{\C}(K3) + \dim_{\C}(E) = 2+1 = 3\),
the holomorphic three-form is \(\sigma_S\wedge dz_E\), and the
shifted symplectic is \([-1]\)-shifted by \cite[\S 2.1]{PTVV2013}.
The cosection of \S~D.4 removes one obstruction direction and
increases the reduced virtual dimension by one, but does not reduce
the shifted-symplectic degree; \(\omega^{\mathrm{red}}\) is still
\((-1)\)-shifted, on the cosection vanishing locus, by the
Maulik--Toda \cite[Theorem~2.7]{MaulikTodaGV} extension.  The
cosection-twist of \S~D.4 is the Serre-equivariant reduced substitute
for the missing global toric/quiver-variety atlas.

The four calculations above are the local Hall forms of
\((D0)\), \((O1)\), \((O1)^+\), and \((O2)\).  Their vanishing is the
hypothesis of the conditional cross-level identity
\(\textcircled{2}\)\(\leftrightarrow\)\(\textcircled{3}\).  The Hall
correspondences above supply the compact-support functorial maps; the
recognition
\(\mathrm{Prim}_\bullet(Y^+(X))\simeq\mathfrak g_{\Delta_5}\) of
Theorem~\ref{thm:primitive-recognition} remains conditional on
\((S1)\)--\((S10)\).  The protected Pfaffian theorem compares the
View~\textcircled{3} Pfaffian section with \(\mathcal D_X=\Delta_5\); the
level-\(\mathsf Z\) scalar trace is a later trace operation.



\section{Pfaffian wall normal form on type-II walls}
\label{app:pfaffian-wall-normal-form}

The three local Pfaffian wall computations supply the
sign data used in Chapter \ref{ch:pfaffian-dirac} for the
orientation-character match
\[
\epsilon_o = \nu_{\Delta_5}\quad\text{on}\quad
W^{(2)}(\La^{2,1}_{II}),
\]
after the global orientation transport of
Theorems~\ref{thm:G-O1}, \ref{thm:G-O1-plus}, and
\ref{thm:G-O2-orbit-transport}. Every numerical
identity of \S\ref{app:pfaffian-wall-three-walls} and
\S\ref{app:pfaffian-wall-recognition} is independently verified by
exact rational lattice arithmetic and exact rational expansion of the
Borcherds product.

\subsection{The three simple type-II walls}
\label{app:pfaffian-wall-three-walls}

The hyperbolic root lattice of signature \((2,1)\) is
\(\La^{2,1} = \HypPlan \oplus \langle 2 \rangle\), with basis
\((f_2, f_3, f_{-2})\) and bilinear form
\((f_2, f_{-2}) = -1\), \((f_3, f_3) = 2\). Its index-four type-II
sublattice is
\[
\La^{2,1}_{II}
= \{m f_2 + l f_3 + n f_{-2} \in \La^{2,1} \mid
m \equiv n \equiv 0 \pmod 2\},
\]
abstractly \(\La^{2,1}_{II}\simeq U(4)\oplus\langle 2\rangle\) with
discriminant \(-32\); the reflection-group index
\([W^{(2)}(\La^{2,1}):W^{(2)}(\La^{2,1}_{II})]=6\) tracks the orbit
data on the chamber-walls and is independent of the sublattice index.
The three simple type-II roots are
\[
\delta_1 = 2 f_2 - f_3, \qquad
\delta_2 = 2 f_{-2} - f_3, \qquad
\delta_3 = f_3,
\]
with pairwise pairings forming the type-II Cartan matrix
\[
A = \bigl((\delta_i, \delta_j)\bigr)_{1 \le i, j \le 3}
=
\begin{pmatrix}
2 & -2 & -2 \\
-2 & 2 & -2 \\
-2 & -2 & 2
\end{pmatrix},
\qquad
\det A = -32.
\]
The matrix is symmetric generalized Cartan, hyperbolic (one negative
eigenvalue, two positive), and the type-II Weyl group
\(W^{(2)}(\La^{2,1}_{II})\) generated by
\(s_{\delta_1}, s_{\delta_2}, s_{\delta_3}\) is the Coxeter group with
relations \(s_{\delta_i}^2=1\) and no braid relation between distinct
generators; the off-diagonal value \(-2\) is the infinite Coxeter
exponent.  It acts on the hyperbolic plane \(\R\Poly_{II}\). The Weyl
vector \(\rho\) determined
by \((\rho, \delta_i) = -1\) for \(i = 1, 2, 3\) is
\[
\rho = \tfrac{1}{2}(\delta_1 + \delta_2 + \delta_3)
     = f_2 - \tfrac{1}{2} f_3 + f_{-2}.
\]
The Cartan matrix, the Weyl vector, the reflection identity
\(s_{\delta_i}(\delta_i) = -\delta_i\), and the preservation of the
form under \(s_{\delta_i}\) are all direct computations in the BKM
Gram.

\begin{remark}[Type-II versus type-I]
The lattice \(\La^{2,1}\) admits both type-I roots
(divisibility \((\delta, \La^{2,1}) = \Z\), simple system
\(\{f_{-2} - f_2, f_2 - f_3, f_{-2} - f_3\}\), trivial multiplier) and
type-II roots (divisibility \(2\Z\), simple system
\(\{\delta_1, \delta_2, \delta_3\}\) above, multiplier \(-1\) on each
generator). The Borcherds--Gritsenko--Nikulin denominator algebra
\(\fg_{\Delta_5}\) is the type-II algebra, with denominator
\(64^{-1} \Delta_5(2Z)\) on \(\La^{2,1}_{II}\); the doubling \(2Z\) is
the passage from the type-I to the type-II lattice, and the constant
\(64 = 2^6\) is the theta-normalization
\([q^{1/2} r^{1/2} s^{1/2}] \Delta_5\). The Pfaffian wall calculation
is type-II throughout.
\end{remark}

\subsection{The rank-one type-II Pfaffian crossing}
\label{app:pfaffian-wall-rank-one}

For each simple type-II wall
\(\delta \in \{\delta_1, \delta_2, \delta_3\}\), the retained
rank-one chart in the datum \((O2_{\mathrm{atlas}})\) has local model
\[
U_\delta = T_\delta \times (-\varepsilon, \varepsilon),
\qquad
s_\delta(t, u) = (t, -u),
\]
with wall locus \(T_\delta \times \{0\}\) coinciding with the fixed
locus of \(s_\delta\).  The reduced normal self-Ext factor on this
retained chart is the rank-one odd Pfaffian crossing
\[
K^{\mathrm{cross}}_\delta = [\R \xrightarrow{\,u\,} \R],
\qquad
D_u = \begin{pmatrix} 0 & u \\ -u & 0 \end{pmatrix},
\qquad
\mathrm{Pf}(D_u) = u.
\]
Choose an oriented tangent Pfaffian line
\(\mathscr L^{\mathrm{tan}}_\delta\) on \(T_\delta\) with
\(s_\delta\)-invariant unit \(\upsilon_\delta\):
\[
s_\delta(\upsilon_\delta) = \upsilon_\delta.
\]
The total local Pfaffian section is
\[
\mathrm{pf}_\delta = \upsilon_\delta\, u.
\]

\begin{proposition}[Local Pfaffian sign; conditional on the retained rank-one wall chart]
\label{prop:appE-local-pfaffian-sign}
For the local rank-one type-II Pfaffian crossing
\((U_\delta, \mathrm{pf}_\delta)\) of
\S\ref{app:pfaffian-wall-rank-one},
\[
s_\delta^*\,\mathrm{pf}_\delta = -\mathrm{pf}_\delta,
\qquad
s_\delta^*\,(\mathrm{pf}_\delta^2) = \mathrm{pf}_\delta^2.
\]
Equivalently, the local Hall/Pfaffian sign is
\[
\epsilon_o^{\mathrm{loc}}(s_\delta) = (-1)^{N_\delta^{\mathrm{Pf}}} = -1,
\qquad
N_\delta^{\mathrm{Pf}} = 1.
\]
\end{proposition}

\begin{proof}
The fixed locus of \(s_\delta(t, u) = (t, -u)\) is \(\{u = 0\}\), and
the wall coincides with the tangent fixed locus, so the determinant
factor is split: \(\mathscr L^{\mathrm{tan}}_\delta\) on \(T_\delta\)
is the tangent factor and \(K^{\mathrm{cross}}_\delta\) is the normal
factor. The unit \(\upsilon_\delta\) is invariant by construction; the
skew matrix \(D_u\) has Pfaffian \(u\); and \(s_\delta\) sends
\(u \mapsto -u\) while fixing \(\upsilon_\delta\). Therefore
\[
s_\delta^*(\upsilon_\delta\,u) = \upsilon_\delta\,(-u) = -\upsilon_\delta\,u
= -\mathrm{pf}_\delta,
\]
and squaring gives the second equality. The Pfaffian rank exponent
\(N_\delta^{\mathrm{Pf}} = 1\) follows because the normal factor is the
single odd crossing \([\R \to \R]\). The displayed sign is
\((-1)^1 = -1\).
\end{proof}

\subsection{Match with the Maass multiplier}
\label{app:pfaffian-wall-maass-match}

The Maass multiplier \(\nu_{\Delta_5}\) of the weight-\(5\) Igusa cusp
form is the multiplier system of the section
\(\Delta_5 \in M_5(\Sp_4(\Z), \nu_{\Delta_5})\) of the automorphic
line bundle \(\mathcal L^5 \otimes \nu_{\Delta_5}\) on the
genus-two Siegel space \(\Sp(2,\Z) \backslash \Hpl_2\). Its values on
the type-II reflections in
\(\Orth(\La^{2,1})_+ = W^{(2)}(\La^{2,1})\) are computed in
\cite[Proposition 2.1]{GN} (citing Maass \cite[(1.3)--(1.5)]{MAASS:1})
via the exterior-square identification
\(\PSp_4(\Z) \simeq \Gamma_{3,2}\subset \mathrm P\Orth(\La^{3,2})_+\):

\begin{lemma}[Maass multiplier on simple type-II reflections]
\label{lem:appE-maass-on-type-ii}
The Maass character of \cite[Proposition 2.1]{GN}, in the
normalization of \cite[(1.3)--(1.5)]{MAASS:1}, evaluates on the
simple type-II reflections as
\[
\nu_{\Delta_5}(s_{\delta_i}) = -1, \qquad i = 1, 2, 3,
\]
and trivially on the three \(\mathrm{Aut}(\Poly_{II}) \simeq S_3\)
reflections \(s_{f_{-2} - f_2}, s_{f_2 - f_3}, s_{f_{-2} - f_3}\):
\[
\nu_{\Delta_5}(s_{f_{-2} - f_2})
= \nu_{\Delta_5}(s_{f_2 - f_3})
= \nu_{\Delta_5}(s_{f_{-2} - f_3}) = +1.
\]
The character is invariant under rescaling \(\Delta_5\) by a nonzero
constant: the theta-normalization
\([q^{1/2} r^{1/2} s^{1/2}] \Delta_5 = 64\) and the monic
\(D_5 = 64^{-1} \Delta_5\) give the same multiplier values on
generators.

\medskip
\emph{Derivation.}  Each \(s_{\delta_i}\) lifts under the exterior-square
isomorphism \(\PSp_4(\Z)\xrightarrow{\sim}\Gamma_{3,2}\) of
Theorem~\ref{thm:exceptional-iso} to a paramodular involution acting
on the genus-two Siegel space; \(\Delta_5\) transforms by the Maass
character on these involutions, valued \(-1\) on each type-II
generator (Maass \cite[(1.3)--(1.5)]{MAASS:1}, Gritsenko--Nikulin
\cite[Proposition 2.1]{GN}).  The trivial \(S_3\)-reflections
\(s_{f_{-2}-f_2}, s_{f_2-f_3}, s_{f_{-2}-f_3}\) lift to inner Weyl
permutations on the level-one paramodular and act with multiplier
\(+1\).
\end{lemma}

Combining Proposition \ref{prop:appE-local-pfaffian-sign} and Lemma
\ref{lem:appE-maass-on-type-ii}:

\begin{theorem}[Local orientation match on type-II generators]
\label{thm:appE-local-orientation-match}
Composing Proposition~\ref{prop:appE-local-pfaffian-sign} and
Lemma~\ref{lem:appE-maass-on-type-ii}, on the three simple type-II
reflections,
\[
\epsilon_o^{\mathrm{loc}}(s_{\delta_i}) = -1 = \nu_{\Delta_5}(s_{\delta_i}),
\qquad i = 1, 2, 3.
\]
The local Pfaffian sign and the Maass multiplier coincide on
generators of the type-II reflection group, relative to the retained
rank-one wall charts.
\end{theorem}

The match \(\epsilon_o^{\mathrm{loc}} = \nu_{\Delta_5}\) on the three
type-II simple reflections is the local seed of the global
orientation-character identity. Its extension to the whole
type-II Weyl group requires global Hall-stack transport of the
local crossing, the \((O2)\) obstruction.

\subsection{Extension to the type-II Weyl group}
\label{app:pfaffian-wall-W2-extension}

\subsubsection{Generation and a presentation}

The type-II Weyl group \(W^{(2)}(\La^{2,1}_{II})\) is generated by
\(s_{\delta_1}, s_{\delta_2}, s_{\delta_3}\) (Nikulin's reflection
theorem for \(\HypPlan \oplus \langle 2 \rangle\),
\cite[Theorem 4.1]{GN}). The generator pairings
\((\delta_i, \delta_j) = -2\) for \(i \ne j\) are obtuse and equal to
\(-2\); the corresponding Coxeter exponent
\(m_{ij}\) is the order of \(s_{\delta_i} s_{\delta_j}\), which is
infinite because the rotation \(s_{\delta_i} s_{\delta_j}\) acts on
the hyperbolic plane by a hyperbolic translation. Thus
\[
W^{(2)}(\La^{2,1}_{II})
= \langle s_{\delta_1}, s_{\delta_2}, s_{\delta_3} \mid
s_{\delta_i}^2 = 1 \rangle
\]
is the free product of three involutions, equivalently the
\((\infty, \infty, \infty)\)-triangle reflection group on the
hyperbolic plane.

\subsubsection{The determinant character}

The determinant character
\(\det : W^{(2)}(\La^{2,1}_{II}) \to \{\pm 1\}\) is the unique
homomorphism with \(\det(s_{\delta_i}) = -1\) on each generator. The relations \(s_{\delta_i}^2 = 1\) are respected since
\((-1)^2 = 1\), and there are no further relations, so \(\det\)
extends unambiguously. The Maass multiplier
\(\nu_{\Delta_5}|_{W^{(2)}(\La^{2,1}_{II})}\) is also a
\(\{\pm 1\}\)-character that equals \(-1\) on each generator (Lemma
\ref{lem:appE-maass-on-type-ii}); it therefore equals \(\det\):
\[
\nu_{\Delta_5}|_{W^{(2)}(\La^{2,1}_{II})}
= \det = (\text{sign character}).
\]
This is the global Maass character on the type-II Weyl group.

\subsubsection{The orientation character is a $\{\pm 1\}$-character}

The orientation character \(\epsilon_o\) on the Hall-stack moduli is,
by construction, a homomorphism
\(W^{(2)}(\La^{2,1}_{II}) \to \{\pm 1\}\) provided the orientation
cocycle on the half-Hilbert orbit
trivializes consistently with the local rank-one crossings of
\S\ref{app:pfaffian-wall-rank-one} across all generators. Under that
trivialisation, the global orientation character is determined by its
values on \(s_{\delta_1}, s_{\delta_2}, s_{\delta_3}\), each equal to
\(-1\) by Theorem \ref{thm:appE-local-orientation-match}. Hence
\[
\epsilon_o = \det = \nu_{\Delta_5}
\quad\text{on}\quad W^{(2)}(\La^{2,1}_{II}).
\]
This is the global orientation-character match relative to the retained
quotient orientation and the chosen Weyl transport.  Its proof in this
appendix establishes the \((O2)\) wall-atlas part: the trivialisation
of the orientation cocycle on the half-Hilbert orbit.

\subsection{The (O2) obstruction: half-Hilbert orbit transport}
\label{app:pfaffian-wall-O2}

\subsubsection{Statement of (O2)}

\begin{definition}[(O2) local Pfaffian wall normal form, global form]
\label{def:appE-O2}
The obstruction \((O2)\) is the assertion that, for every retained HN
height \(R\) and every retained type-II wall transport
\(s_\delta \in W^{(2)}(\La^{2,1}_{II})\), the local rank-one crossing
of \S\ref{app:pfaffian-wall-rank-one}, defined on the generic
wall chart \(U_{\delta, R}\), extends to a strictly compatible system
of charts on the half-Hilbert orbit
\(\mathfrak H^{\mathrm{half}}_R\), with the orientation cocycle on
this orbit trivialized by an explicit cochain compatible with the local
Pfaffian sign on each component, the boundary of each component, and
the overlap of every adjacent pair of components.
\end{definition}

\subsubsection{Local computation and global transport}

Theorem \ref{thm:appE-local-orientation-match} proves the local sign
identity on each of the three simple type-II generators by a single
chart computation: \emph{three local sign computations, one per
generator}. The global statement of
\S\ref{app:pfaffian-wall-W2-extension},
\(\epsilon_o = \nu_{\Delta_5}\) on the entire Weyl group, requires
that these three local seeds extend to a strictly coherent system on
the half-Hilbert orbit, namely the entire orbit of the half-Hilbert
moduli stack under the \(W^{(2)}\)-action. The component, boundary,
and overlap compatibility data of Definition \ref{def:appE-O2} are
\emph{not} produced by the local rank-one crossings alone.  They are
the content of the retained \((O2_{\mathrm{atlas}})\) datum transported
by Theorem~\ref{thm:G-O2-orbit-transport}.

\subsubsection{Two appearances of the constant $4096$}
\label{app:pfaffian-wall-O2-4096}

The integer \(4096\) appears in two structures, and the equality is
scalar only. On the orientation side, the
half-Hilbert orbit of the type-II reflection group on the
cosection-reduced self-Ext stack carries \(12\) binary signs, so the
orientation-cocycle support carries \(2^{12} = 4096\) sign assignments;
this is the reading on which the
\((O2_{\mathrm{atlas}})\) transport of Appendix~\ref{app:obstruction-discharges},
Theorem~\ref{thm:G-O2-orbit-transport}, operates. On the scalar side,
the Oberdieck--Pixton normalization
\[
Z^X_{\mathrm{OP}}(P_{\mathrm{OP}}, Q_{\mathrm{OP}}, T_{\mathrm{OP}})
= -\bigl(\chi_{10}^{\mathrm{OP}}\bigr)^{-1}
= -4096\, \Delta_5(Z)^{-2},
\qquad
\chi_{10}^{\mathrm{OP}} = D_5^2 = (64^{-1} \Delta_5)^2
= 4096^{-1} \Delta_5^2,
\]
exhibits \(4096 = 64^2 = ([q^{1/2} r^{1/2} s^{1/2}] \Delta_5)^2\) as the
squared theta-leading coefficient, with \(-4096 = (-1)\cdot 64^2\) the
OP scalar branch convention times the squared theta-leading
(\cite{OB,OPand}; arithmetic check by the product expansion).

The two values coincide because \(2^{12}=64^2\): the half-Hilbert orbit
sign count and the squared theta-leading are the same integer after
orientation has been forgotten.
Theorem~\ref{thm:G-R1-R2} of Appendix~\ref{app:obstruction-discharges}
makes the scalar separation precise: the orientation character
\(\epsilon_o = \nu_{\Delta_5}\) is fixed, relative to the retained
orientation and Weyl-transport data, by wall transport of the
sign-count, while the OP scalar magnitude \(4096\) enters
\(Z^X_{\mathrm{OP}} = -4096\,\Delta_5^{-2}\) through the squared
theta-leading.

\subsubsection{The retained (O2) wall-atlas route}

The retained \((O2_{\mathrm{atlas}})\) datum consists of:
\begin{enumerate}
\item[(O2.a)] A strictly compatible system of charts \(U_{\delta, R}\),
\(\delta \in W^{(2)}(\La^{2,1}_{II})\), on the half-Hilbert orbit at
height \(R\), each presenting the rank-one type-II Pfaffian crossing of
\S\ref{app:pfaffian-wall-rank-one}.
\item[(O2.b)] An orientation cocycle
\(c_o \in Z^2(W^{(2)}(\La^{2,1}_{II}); \mathbb F_2)\) on the orbit, with
explicit cochain trivialisation
\(\eta_o \in C^1(W^{(2)}(\La^{2,1}_{II}); \mathbb F_2)\) satisfying
\(\delta \eta_o = c_o\), compatible with the chosen charts.
\item[(O2.c)] Strict compatibility of the trivialisation with the
component decomposition, the boundary of each component, and the
overlap of every adjacent pair of components — equivalently, the
component, boundary, and overlap residual entries of the retained
wall-atlas datum vanish with explicit null-trivialisations.
\end{enumerate}
Once (O2.a)--(O2.c), the retained quotient orientation, and the chosen
Weyl transport are fixed, Theorem
\ref{thm:appE-local-orientation-match} extends to the global
Weyl-group character match
\(\epsilon_o = \nu_{\Delta_5}\) on \(W^{(2)}(\La^{2,1}_{II})\).

\subsection{First-timelike parity split as a recognition obligation}
\label{app:pfaffian-wall-recognition}

The orientation-character match of \S\ref{app:pfaffian-wall-W2-extension}
is the conditional \emph{multiplier} comparison. The recognition
theorem \(\mathcal P_X \cong \fg_{\Delta_5}\) of Chapter
\ref{ch:pfaffian-dirac} requires further data, namely full
parity dimensions per root degree, not just signed dimensions. The
first nontrivial test occurs at the imaginary simple root
\(\delta_{123} = \delta_1 + \delta_2 + \delta_3\):

\begin{lemma}[First-timelike parity split]
\label{lem:appE-first-timelike-parity}
The signed root multiplicity at \(\delta_{123}\) is
\[
\smult\,\fg_{\Delta_5, \alpha(\delta_{123})} = f(1, 1) = -64,
\]
and the full target parity computed from the Gritsenko--Nikulin
presentation is
\[
\bigl(\rootmult_{\overline 0}(\delta_{123}),\,
      \rootmult_{\overline 1}(\delta_{123})\bigr)
= (29, 93),
\qquad
29 - 93 = -64.
\]
The decomposition of the even part is
\[
29 = 2 + 3 \cdot 9,
\]
where \(2 = \tfrac{1}{3}\,\mu_{\mathrm{Mob}}(1)\,3!\) is the Witt dimension of the
free Lie algebra on three generators in multidegree \((1, 1, 1)\) and
\(3 \cdot 9\) is the three real--isotropic mixings each contributing the
eta-9 multiplicity \(\tau = 9\) of the doubled isotropic ray. The odd
part \(93\) is the absolute value of the timelike Borcherds correction
\(m(\delta_{123})\), with
\[
m(\delta_{123}) =
-\bigl[\,q^1 r^1 s^1\,\bigr]\bigl((qrs)^{-1/2}64^{-1}\Delta_5\bigr)
= -93.
\]
\end{lemma}

\begin{proof}
The signed multiplicity is \(f(1, 1) = -64\) by direct expansion of
\(\phi_{0,1}\) in the Eichler--Zagier normalization
\cite[\S2.2]{EZ:1}. The full parity decomposition
is the Gritsenko--Nikulin presentation reading: the even part counts
free Lie words in three real-root letters, the odd part counts simple
imaginary timelike roots. Witt's dimension formula
\(W_{\mathrm{multidegree}}(1,1,1)
= \tfrac{1}{n} \sum_{d \mid \gcd} \mu_{\mathrm{Mob}}(d)\,(n/d)! / \prod (n_i/d)!\) with
\(n = 3\), \(\gcd(1,1,1) = 1\) gives \(\tfrac{1}{3} \cdot 1 \cdot 6 = 2\).
The three real--isotropic mixings each contribute
\(\tau = 9\) on the primitive isotropic rays (the eta-9 expansion
\(\prod_{k \ge 1}(1 - q^k)^9\) of \cite[Theorem 2.1, (5.1)]{GNII}); summing over
the three independent isotropic directions gives \(3 \cdot 9 = 27\).
Hence \(29 = 2 + 27\). The odd part is \(-m(\delta_{123}) = 93\) by the
Gritsenko--Nikulin imaginary simple-root formula.
\end{proof}

\begin{corollary}[Signed dimension is not recognition]
\label{cor:appE-signed-not-recognition}
A graded super vector space with parity dimensions
\((\dim_{\overline 0}, \dim_{\overline 1}) = (29, 93)\) and zero
bracket has the same signed dimension \(-64\) as the
\(\fg_{\Delta_5, \alpha(\delta_{123})}\) component but the wrong Lie
superalgebra structure. Recognition of \(\mathcal P_X\) with
\(\fg_{\Delta_5}\) at degree \(\delta_{123}\) requires both parity
dimensions \(29\) and \(93\) and the Borcherds--Kac--Moody bracket; the
determinant identity alone does not imply this recognition statement.
\end{corollary}

\subsubsection{Higher-degree consistency}

The same Borcherds--Kac arithmetic gives higher-degree signed checks
and the finite rows where full target parity is known:

\begin{itemize}
\item \emph{Doubled isotropic.} For \(\delta = 2 a_{ij}\), \(i \ne j\),
\[
(\smult,\, \tau,\, \text{gap}) = (10,\, 9,\, 1),
\]
the signed multiplicity \(10\) splits as eta-9 multiplicity \(9\) plus
the Kac null-root residual \(1\) of the rank-two affine subroot
system, for each of the three pairs \((i,j)\).

\item \emph{Height-four timelike.} For the three rows
\[
C_{k,2}=a_{ij}+2\delta_k,\qquad \{i,j,k\}=\{1,2,3\},
\]
namely coefficient vectors \((2,1,1),(1,2,1),(1,1,2)\),
\[
(\smult,\, m,\, \text{gap}) = (108,\, 90,\, 18),
\]
each with the same triple of integers.  These rows are not Weyl
translates of \(\delta_{123}\); rather
\[
s_{\delta_k}(\delta_{123})=a_{ij}+3\delta_k=C_{k,3}.
\]
Thus \(C_{k,2}\) is a signed target row.  The integer \(18=108-90\)
is the signed lower-bracket residual at height four, not a full parity
block.

\item \emph{Real-string exponents.} The Borcherds--Kac real-root strings
through \(\delta_{ij}\)-style mixed degrees give
\((\text{real-real exponent}) = 3\) and
\((\text{complement-isotropic exponent}) = 5\); the timelike string
through \(\delta_{123}\) has exponent \(3\). All three are pairwise
consistent under Weyl reflection.
\end{itemize}
These integers are the degree-wise Borcherds--Kac arithmetic used in
the recognition criterion.

\subsection{Local Pfaffian sign computation}
\label{app:pfaffian-wall-sign-computation}

Let \(\delta\in\{\delta_1,\delta_2,\delta_3\}\subset\La^{2,1}\), with
\((f_2,f_{-2})=-1\) and \((f_3,f_3)=2\).  On a normal chart
\[
U_\delta=T_\delta\times(-\epsilon,\epsilon),\qquad
s_\delta(t,u)=(t,-u),\qquad
\Hpl_\delta\cap U_\delta=T_\delta\times\{0\},
\]
choose an oriented tangent Pfaffian unit \(\upsilon_\delta\) satisfying
\(s_\delta^*\upsilon_\delta=\upsilon_\delta\).  The rank-one normal
complex is
\[
K^{\mathrm{nor}}_\delta=[\mathbb R\xrightarrow{u}\mathbb R],
\qquad
D_u=
\begin{pmatrix}
0&u\\
-u&0
\end{pmatrix},
\qquad
\operatorname{Pf}(D_u)=u .
\]
Thus
\[
\operatorname{pf}_\delta=\upsilon_\delta u,\qquad
s_\delta^*\operatorname{pf}_\delta=-\operatorname{pf}_\delta,\qquad
s_\delta^*\operatorname{pf}_\delta^2=\operatorname{pf}_\delta^2 .
\]
The normal Pfaffian exponent is \(N_\delta^{\mathrm{Pf}}=1\), hence
\[
\epsilon_o^{\mathrm{loc}}(s_\delta)
=(-1)^{N_\delta^{\mathrm{Pf}}}
=-1
=\nu_{\Delta_5}(s_\delta)
\]
for the three simple type-II reflections, by the Maass-character
calculation of Gritsenko--Nikulin \cite[Proposition~2.1]{GN}.

The \((O2_{\mathrm{atlas}})\)-transport formulas of
Appendix~\ref{app:eight-finite-constructions}
extends this local computation to the global half-Hilbert orbit by
component, boundary, and overlap compatibility.

\subsection{O2 relative summary}
\label{app:pfaffian-wall-open-obligations}

\begin{itemize}
\item \textbf{(O2.a) Strictly compatible chart system on the half-Hilbert
orbit.} The local rank-one crossing is fixed at each simple
type-II generator (\S\ref{app:pfaffian-wall-rank-one}); strict
compatibility across all elements of \(W^{(2)}(\La^{2,1}_{II})\) is
the half-Hilbert orbit chart system retained in
\((O2_{\mathrm{atlas}})\) and transported by
Theorem~\ref{thm:G-O2-orbit-transport}.
\item \textbf{(O2.b) Orientation cocycle trivialisation.}
\(c_o \in Z^2(W^{(2)}(\La^{2,1}_{II}); \mathbb F_2)\) is the
orientation cocycle; its trivializing cochain
\(\eta_o \in C^1(W^{(2)}(\La^{2,1}_{II}); \mathbb F_2)\) is part of
the chosen Weyl-transport datum of Theorem~\ref{thm:G-O1-plus} and
is transported in
Theorem~\ref{thm:G-O2-orbit-transport}.
\item \textbf{(O2.c) Component, boundary, overlap compatibility.}
The residual entries
\(\mathfrak o^{\mathrm{chart}}_{\delta, R}\),
\(\mathfrak o^{\mathrm{wall}}_{\delta, R}\),
\(\mathfrak o^{\mathrm{split}}_{\delta, R}\),
\(\mathfrak o^{\mathrm{unit}}_{\delta, R}\),
\(\mathfrak o^{\mathrm{rank}}_{\delta, R}\),
\(\mathfrak o^{\mathrm{tr}}_{\delta, R^+ R}\) of
Proposition \ref{prop:rank-one-crossing-o2-sign} are fixed at finite
height \(R\) for the local rank-one models; the global compatibility
across the half-Hilbert orbit is retained in
\((O2_{\mathrm{atlas}})\) and transported by
Theorem~\ref{thm:G-O2-orbit-transport}.
\item \textbf{\((R1)\)--\((R2)\) scalar separation.} Separated in
Appendix~\ref{app:obstruction-discharges},
Theorem~\ref{thm:G-R1-R2}: the orientation-cocycle support on the
half-Hilbert orbit has size \(2^{12}\) (R1), and this integer is equal
to \(64^2\), the theta-leading square on the scalar side (R2).  This is
not an identification of the two data.
\end{itemize}

Relative to the chosen Weyl lifts and \((O2_{\mathrm{atlas}})\),
Theorems~\ref{thm:G-O1-plus}, \ref{thm:G-O2-orbit-transport}, and
\ref{thm:G-R1-R2} extend Theorem
\ref{thm:appE-local-orientation-match} from the three local generators
to the global character match
\(\epsilon_o = \nu_{\Delta_5}\) on \(W^{(2)}(\La^{2,1}_{II})\).

% !TEX root = ../main.tex
%
% Appendix F — Eight build algorithms in pseudocode for the Dirac--Igusa
% pro-object \mathfrak D_X^{\mathrm{DI}}.  Authored from main.tex line
% ranges 5887--7460 (compact realization opening; finite algebraic
% Dirac--Igusa target system; chiral Koszul source datum) and
% 12147--14346 (Dirac--Igusa realization datum; primitive recognition
% theorem; cofinal finite-window primitive-recognition datum; finite
% retained K3 \times E Hall kernel; geometric hybrid local/wrapped/mixed
% atlas; geometric source chiral Koszul coalgebra).  Cross-repo anchors:
% \texttt{chiral-bar-cobar-vol2} Construction Problem 1 (the source of
% \operatorname{Pf}_{\mathrm{prot}}(\mathfrak D_X) = \Delta_5); \texttt{calabi-yau-quantum-groups}
% Hall--Drinfeld correspondence-programme remark (the K3-side BKM
% Hall--Drinfeld double).
%
% Loaded from main.tex via % !TEX root = ../main.tex
%
% Appendix F — Eight build algorithms in pseudocode for the Dirac--Igusa
% pro-object \mathfrak D_X^{\mathrm{DI}}.  Authored from main.tex line
% ranges 5887--7460 (compact realization opening; finite algebraic
% Dirac--Igusa target system; chiral Koszul source datum) and
% 12147--14346 (Dirac--Igusa realization datum; primitive recognition
% theorem; cofinal finite-window primitive-recognition datum; finite
% retained K3 \times E Hall kernel; geometric hybrid local/wrapped/mixed
% atlas; geometric source chiral Koszul coalgebra).  Cross-repo anchors:
% \texttt{chiral-bar-cobar-vol2} Construction Problem 1 (the source of
% \operatorname{Pf}_{\mathrm{prot}}(\mathfrak D_X) = \Delta_5); \texttt{calabi-yau-quantum-groups}
% Hall--Drinfeld correspondence-programme remark (the K3-side BKM
% Hall--Drinfeld double).
%
% Loaded from main.tex via % !TEX root = ../main.tex
%
% Appendix F — Eight build algorithms in pseudocode for the Dirac--Igusa
% pro-object \mathfrak D_X^{\mathrm{DI}}.  Authored from main.tex line
% ranges 5887--7460 (compact realization opening; finite algebraic
% Dirac--Igusa target system; chiral Koszul source datum) and
% 12147--14346 (Dirac--Igusa realization datum; primitive recognition
% theorem; cofinal finite-window primitive-recognition datum; finite
% retained K3 \times E Hall kernel; geometric hybrid local/wrapped/mixed
% atlas; geometric source chiral Koszul coalgebra).  Cross-repo anchors:
% \texttt{chiral-bar-cobar-vol2} Construction Problem 1 (the source of
% \operatorname{Pf}_{\mathrm{prot}}(\mathfrak D_X) = \Delta_5); \texttt{calabi-yau-quantum-groups}
% Hall--Drinfeld correspondence-programme remark (the K3-side BKM
% Hall--Drinfeld double).
%
% Loaded from main.tex via \input{appendices/F_build_algorithms} after
% \appendix has been opened at main.tex line 17081, so that
% \thesection = \Alph{section}.  No \chapter, no second \appendix.

\section{Eight build algorithms for the Dirac--Igusa pro-object}
\label{app:eight-build-algorithms}

The Dirac--Igusa pro-object
\[
\mathfrak D_X^{\mathrm{DI}}
=\varprojlim_R
\bigl(
\mathcal A^{E_3}_{X,R},
\mathcal F_{X,R}^{\mathrm{hyb}},
\Gamma_{X,R},
\Pi_X,
\Phi_R,
o_R,
\mathcal H_R,
P_R^\Pi,
C_{X,R},
\Theta_{\mathrm{Kos},R},
\mathscr L_{\mathrm{Pf},R},
\operatorname{pf}_{X,R},
\epsilon_{o,R},
\mathrm{Rec}_R
\bigr)
\]
of Definition~\ref{thm:finite-stage-dirac-igusa-pro-object} and
Definition~\ref{def:O1-strong-reduced-orientation} is the conditional source
operator-level package whose protected Pfaffian, after the four
obstructions \((D_0)\), \((O1)\), \((O1)^+\), \((O2)\) are discharged,
realizes the Igusa cusp form,
\[
\operatorname{Pf}_{\mathrm{prot}}(\mathfrak D_X^{\mathrm{DI}})=\Delta_5,\qquad
\epsilon_o=\nu_{\Delta_5}\big|_{W^{(2)}(\La^{2,1}_{II})},\qquad
P_X^{\Pi,\mathrm{red}}\cong\fg_{\Delta_5}.
\]
Construction at finite Harder--Narasimhan height \(R\) proceeds through
eight construction steps.  Each step is a procedure on input data
already supplied by earlier steps and feeds named output data into the
next.  The eight steps stratify the realization datum
\(\mathfrak K_X=(L_X,H_X,O_X,D_X,R_X)\) of
Definition~\ref{thm:finite-stage-dirac-igusa-pro-object}: steps 1--3
realize \((L_X)\) and \((H_X)\); step 4 realizes the local content of
\((O_X)\); step 5 produces the Hall correspondence object; step 6 fixes
the Pfaffian--Dirac line and feeds the orientation character; steps 7--8
realize \((D_X)\) and \((R_X)\).

Throughout, \(X=S\times E\) with \(S\) a smooth complex K3 surface and
\(E\) a smooth elliptic curve.  \(\sigma=\sigma_{s,t}\) is Liu's product
stability condition on the Abramovich--Polishchuk heart in
\(\mathsf{D}^b(\mathrm{Coh}(X))\)
\cite{AbramovichPolishchukTStructures,LiuProductStability}.  The HN
sector \((\sigma,S)\) is fixed and the cofinal subsystem
\(\mathcal R_{\mathrm{HN}}=\{R_0<R_1<\cdots\}\) is the one of
Definition~\ref{def:O1-strong-reduced-orientation}, governed by
\textup{[K3$\times$E-fin]}, \textup{[K3$\times$E-atlas]}, and
\textup{[ML]}.

Each construction step carries a licensing tag pulled from the
five-type universal stage chain
\(\mathsf P\to\mathsf C\to\mathsf S\to\mathsf Z\to\mathsf A\) of
\texttt{chiral-bar-cobar-vol2:CLAUDE.md~§9}: \(\alpha\) primitive,
\(\beta\) functorial passage, \(\gamma\) shadow specialization,
\(\delta\) endpoint convergence, \(\epsilon\) terminal scalar.  The
ordering charge\(\to\)moduli\(\to\)hybrid\(\to\)orientation\(\to\)Hall
\(\to\)Dirac+Pfaffian\(\to\)Koszul\(\to\)recognition is the unique
ordering compatible with input/output type discipline; reordering any
two adjacent steps fails an input check.

The four obstructions are attributed to specific steps:
\((D_0)\) is the cofinal-trivialization datum of the entire eight-step
construction at every finite \(R\) (steps 1--6 supply its components);
\((O1)\) lives at step 4; \((O1)^+\) lives at step 4 with a Coxeter
coherence check feeding step 6; \((O2)\) lives at step 6.  Step 8
inherits the obstruction tower discharged by steps 1--7 and adds the
exact-completion residual \textup{(R8)} of
Definition~\ref{def:cofinal-primitive-recognition-datum}.

\subsection{Construction~1 — Charge window}
\label{app:build-charge-window}

\paragraph{Tag.} \(\alpha\)-primitive (lattice datum).

\paragraph{Input.}
HN height \(R\in\mathcal R_{\mathrm{HN}}\); the formal dyonic Mukai
lattice \(\Gamma_X^{\mathrm{form}}\) of
Definition~\ref{def:gram-index-lattice}; the Gram map
\(\Pi_X:\Gamma_X^{\mathrm{form}}\to\Gamma_{\mathrm{gram}}\) and the
cocycle \(B\) of Lemma~\ref{lem:gram-additivity-defect}; the active
support \(\Gamma_{\mathrm{act}}=\{\gamma\in\Gamma_{\mathrm{eff}}\mid
f(nm,l)\ne0\}\) of the Igusa Jacobi form \(\phi_{0,1}\); the cofinal
subsystem \(\mathcal R_{\mathrm{HN}}\).

\paragraph{Output.}
A finite cusp-positive active window
\[
A_R\subset\Gamma_{\mathrm{act}}\subset\Gamma_{\mathrm{gram}},
\qquad
A_R\hbox{ downward saturated for the cofinal target-window order},
\]
together with the normal-ordered extension
\[
\widehat\Gamma_R=
\{(c,T)\mid c\in\Gamma_R^{HN},\ T\in\mathcal T_R(c)\}
\subset\widehat\Gamma_X,
\qquad
\Gamma_R^\Pi=\overline\Pi_X(\widehat\Gamma_R)\subset\Gamma_{\mathrm{gram}},
\]
the lift \(L:\Gamma_{\mathrm{gram}}\to\Gamma_X^{\mathrm{form}}\) of
Lemma~\ref{lem:primitive-lift-every-gram-triple}, and the finite test
window \(\Gamma_R^{\mathrm{test}}\) of
Definition~\ref{def:O1-strong-reduced-orientation}.

\paragraph{Procedure.}
\begin{enumerate}
\item Compute the active support cone
\[
\Gamma_{\mathrm{act}}=\{(n,l,m)\in\Gamma_{\mathrm{eff}}\mid f(nm,l)\ne 0\}
\]
from the Fourier expansion of the index-one weak Jacobi form
\(\phi_{0,1}(\tau,z)=\sum_{n,l}c(n,l)q^nr^l\) of
\cite{EZ:1}, with the index substitution
\(f(nm,l):=c(nm,l)\) tabulated in
Appendix~\ref{sec:n1-boundary-compatibility-conditions} for the
\(N=1\) Igusa boundary.

\item Set
\[
A_R^{\circ}
=\{\gamma=(n,l,m)\in\Gamma_{\mathrm{act}}\mid n+|l|+m\le R\}.
\]
Form its downward saturation
\[
A_R
=\{\gamma'\in\Gamma_{\mathrm{act}}\mid
\exists\gamma\in A_R^{\circ},\ 0<\gamma'\le\gamma,\
\gamma-\gamma'\in Q_+\},
\]
in the sense of Definition~\ref{def:cofinal-primitive-recognition-datum}.

\item Normal-ordered lift.  For each \(c\in\Gamma_R^{HN}\) compute the
\(B\)-translate orbit
\[
\mathcal T_R(c)
=\Bigl\{\sum_{i<j}B(c_i,c_j)\Bigm|c=\sum_i c_i,\
c_i\in F_{\sigma,S}(R),\ \sum_i h_S(c_i)\le R\Bigr\}.
\]
Set \(\widehat\Gamma_R=\{(c,T):c\in\Gamma_R^{HN},\ T\in\mathcal T_R(c)\}\).

\item Finite target test window:
\(\Gamma_R^{\mathrm{test}}=\overline\Pi_X(\widehat\Gamma_R)\), compatibly
with \(R\to R'\) by inclusion.

\item Output \((A_R,\widehat\Gamma_R,\Gamma_R^\Pi,\Gamma_R^{\mathrm{test}},L)\).
\end{enumerate}

\paragraph{Type checks.}
\(A_R\) is finite and downward saturated;
\(\overline\Pi_X(\widehat\Gamma_R)=\Gamma_R^\Pi\) lies in
\(\Gamma_{\mathrm{gram}}\); \(\Gamma_R^{HN}\subset
F_{\sigma,S}(R)\) is finite by \textup{[K3$\times$E-fin]}; the cocycle
identity
\(\Pi_X(c+c')=\Pi_X(c)+\Pi_X(c')+B(c,c')\) is an axiom of the cocycle.

\paragraph{Obstructions discharged or invoked.}
None; this step is purely lattice combinatorics from the imported
Gritsenko--Nikulin Borcherds product
\cite{GN,GNPartI,GNII,BorcherdsGKM88,Borcherds95}.
The downward saturation is the input requirement of \textup{(R0)} in
Definition~\ref{def:cofinal-primitive-recognition-datum}.

\subsection{Construction~2 — Finite moduli}
\label{app:build-finite-moduli}

\paragraph{Tag.} \(\alpha\)-primitive (geometric source).

\paragraph{Input.}
The active window \(A_R\) and lift \(L\) of
Construction~\ref{app:build-charge-window}; Liu's product stability
condition \(\sigma_{s,t}\); the standing hypotheses
\textup{[K3$\times$E-fin]} and \textup{[K3$\times$E-atlas]}.

\paragraph{Output.}
A finite class set \(C_R\) and lift
\(\ell_R:\Gamma_R^\Pi\to C_R\) with
\(\Pi_X(\ell_R(\gamma))=\gamma\); the family of finite-type quasi-smooth
derived Artin substacks
\[
\mathfrak M^{ss}_{R,c}
=\mathfrak M^{ss}_{\sigma_{s,t}}(X,c)_{\le R}
\subset\operatorname{Perf}(X),
\qquad c\in C_R,
\]
of \(\sigma_{s,t}\)-semistable objects of full Liu numerical class
\(c\) and HN height \(\le R\); the scalar-rigidified truncation
\(\mathfrak M^{\mathrm{sc}}_{R,c}\); the universal perfect complex
\(\mathcal A_{R,c}\) on \(X\times\mathfrak M^{\mathrm{sc}}_{R,c}\); the
finite stratification by retained closed substacks of finite residual
inertia.

\paragraph{Procedure.}
\begin{enumerate}
\item Form \(C_R=\{c=L\gamma:\gamma\in A_R\}\) closed under retained
extension sums (taken in \(\Gamma_X^{\mathrm{form}}\)).

\item For each \(c\in C_R\) cut the substack of
\(\sigma_{s,t}\)-semistable objects of class \(c\) by HN height \(\le
R\):
\[
\mathfrak M^{ss}_{R,c}=
\bigl\{[A]\in\operatorname{Perf}(X)\mid
\sigma_{s,t}\hbox{-semistable},\
\hn(A)\le R,\
\nu(A)=c\bigr\}.
\]

\item Verify finite-typeness: \textup{[K3$\times$E-fin]} (Liu product
stability + product-boundedness theorem at fixed full Liu charge with
finite fibres over the active charge) gives the finite-type assertion.
PTVV \cite{PTVV2013} on the Calabi--Yau threefold \(X\) gives the
\((-1)\)-shifted symplectic structure on
\(\operatorname{Perf}(X)\), restricted to the retained semistable
sector.

\item Quasi-smoothness: the universal perfect complex
\(\mathcal A_{R,c}\) supplies a relative cotangent complex of bounded
Tor-amplitude on the retained charts.  Tangent at \(A\) is
\(T_A\mathfrak M^{ss}_{R,c}\simeq\operatorname{RHom}_X(A,A)[1]\).

\item Scalar rigidification: form \(\mathfrak M^{\mathrm{sc}}_{R,c}\) by
killing the scalar \(\mathbb G_m\)-automorphism of
\(c\)-stable factors.

\item Two-step flag stack
\(\mathfrak F^{(2)}_{R,c,c',c''}\) and extension stack
\(\mathfrak E_{R,c,c'}\) constructed as iterated derived vector stacks
of \(R\pi_*\mathcal RHom_X(\mathcal B,\mathcal A[1])\); both finite type
and quasi-smooth in the retained sector by
Proposition~\ref{prop:finite-dcritical-hall-product}.

\item Output \((C_R,\ell_R,\{\mathfrak M^{ss}_{R,c}\},
\{\mathfrak E_{R,c,c'}\},\{\mathfrak F^{(2)}_{R,c,c',c''}\})\) with
universal perfect complexes and the finite stratification.
\end{enumerate}

\paragraph{Type checks.}
\(\mathfrak M^{ss}_{R,c}\) is a finite-type quasi-smooth derived Artin
substack of \(\operatorname{Perf}(X)\); the scalar rigidification leaves
only finite residual inertia
\(H_S\subset E_{\mathrm{tors}}\) on each connected stabilizer
component, by
Proposition~\ref{prop:finite-dcritical-hall-product};
PTVV \((-1)\)-shifted symplectic structure restricts.

\paragraph{Obstructions discharged or invoked.}
\textup{[K3$\times$E-fin]} is invoked here in its product-boundedness
form and supplies the finite-type assertion not in the cited literature
for the K3\(\times\)E threefold; it is part of the standing datum.
\textup{[K3$\times$E-atlas]} supplies the quasi-smooth d-critical
charts.  These enter \((D_0)\) at the level of the finite Hall atlas
of Theorem~\ref{thm:finite-stage-dirac-igusa-pro-object}.

\subsection{Construction~3 — Hybrid wrapped factorization carrier}
\label{app:build-hybrid-wrapped}

\paragraph{Tag.} \(\beta\)-functorial passage (Stage-1 carrier).

\paragraph{Input.}
The finite moduli stacks
\(\{\mathfrak M^{\mathrm{sc}}_{R,c}\}\) of
Construction~\ref{app:build-finite-moduli}; the
projection-finite/wrapped split by elliptic degree
\(b_R^{\mathrm{geom}}:\Gamma_{\sigma,S}^{HN}(R)\to\Z_{\ge0}\) of
Definition~\ref{def:hybrid-wrapped-factorization-datum}; a fixed
origin \(0_E\in E\); the Abramovich--Polishchuk/Liu wrapped anchor
candidate.

\paragraph{Output.}
The finite hybrid local/wrapped factorization datum
\[
\mathcal F_{X,\sigma,S,\le R}^{\mathrm{hyb}}
=
(\mathcal C_{\sigma,S,\le R},
b_R^{\mathrm{geom}},
\mathcal F_R^{\mathrm{loc}},
\mathcal G_R^{\mathrm{wr}},
\mathfrak E_R^{\mathrm{loc}},
\mathfrak E_R^{\mathrm{hyb}},
\mathfrak E_R^{\mathrm{wr}},
\mathfrak{Flag}_R^{\mathrm{hyb}},
\mathsf{Corr}_R^{E,\mathrm{hyb}},
\lambda_R,
Q_{E,R},
\theta^Q_R,
o_R,
I_R^{\mathrm{prot}})
\]
of Definition~\ref{def:hybrid-wrapped-factorization-datum},
together with the hybrid Ran base
\(\operatorname{Ran}^{\mathrm{hyb}}(E)\), the eight-word two-step
binary flag atlas \(\{w\in\{\mathrm L,\mathrm W\}^3\}\), and the
quotient-after-correspondence functor \(Q_{E,R}\).

\paragraph{Procedure.}
\begin{enumerate}
\item \emph{Elliptic-degree split.}  Decompose
\(\Gamma_{\sigma,S}^{HN}(R)
=\Gamma_R^{\mathrm{loc}}\sqcup\Gamma_R^{\mathrm{wr}}\) where
\(\Gamma_R^{\mathrm{loc}}=\{b_R^{\mathrm{geom}}=0\}\) and
\(\Gamma_R^{\mathrm{wr}}=\{b_R^{\mathrm{geom}}>0\}\).  Verify additivity
\(b_R^{\mathrm{geom}}(\gamma+\gamma')
=b_R^{\mathrm{geom}}(\gamma)+b_R^{\mathrm{geom}}(\gamma')\) on retained
extensions remaining in the quotient.

\item \emph{Projection-finite local sector.}  For
\(\alpha\in\Gamma_R^{\mathrm{loc}}\) and finite index set \(I\), build
the closed-support incidence stack over \(E^I\):
\[
\mathfrak M^{\mathrm{loc,cl}}_{\alpha,R,I}
=
\Bigl\{(A,\mathbf e)\in\mathfrak M^{\mathrm{sc}}_{R,L\alpha}\times E^I
\Bigm|p_E(\mathrm{Supp}\,A)\subset\bigcup_{i\in I}S\times\{e_i\}\Bigr\}.
\]
Form
\(\mathscr H^{\mathrm{loc}}_{\alpha,R,I}=
(\pi_{\alpha,R,I})_!(\Phi_{\alpha,R,I}^{\mathrm{loc}}\otimes
o_{\alpha,R,I}^{\mathrm{loc}})\) and the open-support sheaf
\(\mathcal F_{\alpha,R}^{\mathrm{loc}}(U)\) over \(U\subset E\) by
restricted Borel--Moore configuration sums.

\item \emph{Wrapped prequotient sector.}  For
\(\eta\in\Gamma_R^{\mathrm{wr}}\), form the rigidified prequotient
\(\mathcal M_{\eta,R}^{\mathrm{wr,rig}}\to\mathcal M_{\eta,R}^{\mathrm{wr}}\)
by the Abramovich--Polishchuk / determinant anchor
\[
\lambda_\eta:\mathcal M_{\eta,R}^{\mathrm{wr,rig}}\longrightarrow E,
\qquad
\lambda_\eta(A)=\det Rp_{E,*}A
\otimes\mathcal O_E(-\chi(A)\cdot 0_E),
\]
divided by a finite cover when \(\chi(A)\ne 1\) so that
\(\lambda_\eta(t\cdot A)=\lambda_\eta(A)+t\) holds.  Verify the
losslessness residual \(\mathfrak o^{\lambda,\mathrm{loss}}_R=0\) by
the Ext-distinguishing test of
Definition~\ref{def:hybrid-wrapped-factorization-datum}, separating
\(W_0=i_{s,*}(\mathcal O_E\oplus\mathcal O_E)\) from
\(W_L=i_{s,*}(L\oplus L^{-1})\).

\item \emph{Mixed local/wrapped correspondences.}  For an ordered
hybrid input
\(\alpha_-\in\Gamma_R^{\mathrm{loc}}\),
\(\eta\in\Gamma_R^{\mathrm{wr}}\),
\(\beta_+\in\Gamma_R^{\mathrm{loc}}\), form
\(\mathfrak E^{\mathrm{mix}}_{\alpha_-;\eta;\beta_+,R}\) before the
reduced \(E\)-quotient.

\item \emph{Eight-word two-step flag atlas.}  Build flag stacks
\(\mathfrak F^{(2),w}_R\) for each
\(w\in\{\mathrm L,\mathrm W\}^3=
\{\mathrm{LLL},\mathrm{LLW},\mathrm{LWL},\mathrm{WLL},\mathrm{LWW},
\mathrm{WLW},\mathrm{WWL},\mathrm{WWW}\}\), classifying filtrations
with three associated graded factors of those types.  Verify the
four-input pentagon coherence and the quotient-after-correspondence
descent.

\item \emph{Hybrid correspondence-module functor.}  Assemble the
correspondence category
\(\mathsf{Corr}^{E,\mathrm{hyb}}_R\) and the functor
\(\mathcal B_R^{E,\mathrm{hyb}}:\mathsf{Corr}^{E,\mathrm{hyb}}_R
\to\mathsf{Ch}_{\C}\) by
\(\mathcal B_R^{E,\mathrm{hyb}}(e)=q_!\,\mathrm{TS}^{\mathrm{red}}_e\,p^*\)
for every generating arrow
\(e=(Y\xleftarrow{p}\mathfrak E_e\xrightarrow{q}Z)\).

\item \emph{Quotient functor \(Q_{E,R}\).}  Apply
\(E\)-equivariant compactly supported vanishing-cycle chains and the
orientation descent of step~4 to produce the reduced object
\[
\mathcal F^{\mathrm{geom}}_{X,\sigma,S,\le R}
=Q_{E,R}\circ\mathcal B_R^{E,\mathrm{hyb}}.
\]

\item Output the hybrid datum with all coherences and the protected
integration map \(I_R^{\mathrm{prot}}\), with homogeneous
\(s\)-degree the Borcherds trace exponent
\(m_R=\operatorname{pr}_3\overline\Pi_X\).
\end{enumerate}

\paragraph{Type checks.}
The carrier
\(\mathcal F_{X,\sigma,S,\le R}^{\mathrm{hyb}}\) is a hybrid
correspondence-module functor on
\(\mathsf{Corr}^{E,\mathrm{hyb}}_R\) with values in
\(\mathsf{Ch}_{\C}\); the four-input pentagon is supplied by
the eight-word two-step flag atlas; the anchor residual tuple
\(\mathfrak o^\lambda_R=
(\mathfrak o^{\lambda,\mathrm{ex}}_R,
\mathfrak o^{\lambda,\mathrm{unit}}_R,
\mathfrak o^{\lambda,\mathrm{loss}}_R,
\mathfrak o^{\lambda,\mathrm{multi}}_R,
\mathfrak o^{\lambda,\mathrm{tr}}_{R'R})\) vanishes; the hybrid
residual \(\mathfrak O_{\mathrm{hyb},R}=0\) at every height \(R\),
compatibly with transition maps.

\paragraph{Obstructions discharged or invoked.}
The middle type-II wall lies in
\(\delta_2\leftrightarrow(0,1,1)\), elliptic degree
\(d=m+1=2\) by Lemma~\ref{app:pfaffian-wall-three-walls}; it
cannot be represented by a projection-finite local stack and demands
the wrapped prequotient.  This step's correctness — the existence of
the wrapped representative
\(x_{\delta_2,R}\in\mathcal M_{\eta,R}^{\mathrm{wr,rig}}\) with
\(\overline\Pi_X(x_{\delta_2,R})=(0,1,1)\) and its compatible anchor
\(\lambda_{\eta,R}\) — is precisely the hybrid input to \((O2)\) at
the middle wall.

\subsection{Construction~4 — Reduced d-critical orientation
and Joyce--Upmeier transport}
\label{app:build-orientation}

\paragraph{Tag.} \(\beta\)-functorial passage (oriented carrier).

\paragraph{Input.}
The hybrid carrier
\(\mathcal F_{X,\sigma,S,\le R}^{\mathrm{hyb}}\) of
Construction~\ref{app:build-hybrid-wrapped}; the K3 holomorphic
\(2\)-form \(\sigma_S\) on \(S\); BBDJS d-critical chart machinery
\cite{BBDJS2015}; the Joyce--Upmeier multiplicative orientation
formalism \cite{JoyceUpmeier,JoyceTanakaUpmeier}; the cosection
localization machinery of Kiem--Li and the cosection-reduced
obstruction theory; the type-II Weyl group
\(W^{(2)}(\La^{2,1}_{II})\) generated by
\(s_{\delta_1},s_{\delta_2},s_{\delta_3}\).

\paragraph{Output.}
The cosection
\(\sigma=p_S^*\sigma_S\); the reduced tangent complex
\[
\operatorname{RHom}_{\mathrm{red}}(A,A)
=
\operatorname{Cone}\!\Bigl(
\operatorname{RHom}_X(A,A)
\xrightarrow{(\mathrm{tr},\mathrm{cs})}
R\Gamma(X,\mathcal O_X)\oplus H^2(S,\mathcal O_S)[-1]
\Bigr)[-1];
\]
the BBDJS vanishing-cycle complex \(\Phi_{R,c}\); the reduced
orientation line
\[
o_{R,c}
=\det\operatorname{Ext}^1_{\mathrm{red}}(A,A)^{-1}
\quad\text{on Darboux charts},\qquad
o_{R,c}^{\otimes 2}\simeq\det\operatorname{RHom}_{\mathrm{red}}(A,A);
\]
the Joyce--Upmeier multiplicative transport
\(o_{c\oplus c'}\simeq o_c\otimes o_{c'}\) compatible with the reduced
extension correspondences; the quotient orientation
\(\mathfrak Q^{or}_R\) of
Theorem~\ref{prop:generic-local-sign-orientation-character}; the
Weyl-equivariant lifts
\(\tau_{i,R,S}:o_{R,S}\to s_{\delta_i}^*o_{R,S}\) for \(i=1,2,3\) of
\((O1)^+\); the local generic wall-sign data \(\epsilon_{o,R}^{\mathrm{loc}}\).

\paragraph{Procedure.}
\begin{enumerate}
\item \emph{Cosection.}  Pull back the K3 semiregularity cosection
\(\mathrm{cs}\) along \(p_S:X\to S\) to obtain the 2-step cosection
\(\sigma=p_S^*\sigma_S\) on the K3 factor; verify additivity in exact
triangles by \textup{(RedOr)} of
Theorem~\ref{thm:finite-stage-dirac-igusa-pro-object}.

\item \emph{Reduced obstruction.}  Form the reduced tangent complex by
the displayed cone above (Kiem--Li cosection localization).  PTVV
\cite{PTVV2013} \((-1)\)-shifted symplectic on
\(\operatorname{RHom}_X(A,A)\) descends, after the cosection
contraction, to the reduced complex.

\item \emph{BBDJS d-critical chart.}  On each Darboux chart of
\(\mathfrak M^{\mathrm{sc}}_{R,c}\), build the BBDJS vanishing-cycle
complex \(\Phi_{R,c}\) \cite{BBDJS2015}; verify reduced
Thom--Sebastiani transport coherent on two-step flag stacks of
Construction~\ref{app:build-hybrid-wrapped}.

\item \emph{Reduced orientation line.}  Define
\(o_{R,c}=\det\operatorname{Ext}^1_{\mathrm{red}}(A,A)^{-1}\); verify
the squared isomorphism
\(o_{R,c}^{\otimes 2}\simeq\det\operatorname{RHom}_{\mathrm{red}}(A,A)\)
by \(3\)-Calabi--Yau Serre duality
\(\operatorname{Ext}^2_{\mathrm{red}}(A,A)
\simeq\operatorname{Ext}^1_{\mathrm{red}}(A,A)^\vee\).

\item \emph{Quotient-orientation descent on the K3\(\times\)E-quotient
self-Ext frame bundle.}  Compute the four Cech--Borel obstruction
classes for every retained stratum \(S\):
\[
\alpha^{\mathrm{red}}_{R,S}=w_2(\mathscr D_{R,S}),
\quad
\alpha^{E,\mathrm{free}}_{R,S}=
\mathrm{edge}_{2,0}(w_2^E(\mathscr D_{R,S})),
\quad
\widetilde\beta^H_{R,S}=w_2^G(\mathscr D_{R,S}),
\quad
\lambda^H_{R,S}\in H^1_G(S;\mathbb F_2),
\]
and verify simultaneous vanishing with compatible
null-trivializations on overlaps, on Hall correspondence pullbacks,
on two-step flag stacks, and under successor maps.  The full
calculation proceeds primary stabilizer by stabilizer; for cyclic
\(\Z/2^a\) the absent mixed coefficient requires the rank-two
two-primary check.  Output the descended quotient orientation system
\(\mathfrak Q^{or}_R\).

\item \emph{Multiplicative Joyce--Upmeier transport.}  Transport
\(o_{c\oplus c'}\simeq o_c\otimes o_{c'}\) coherently across two-step
flag stacks
\cite{JoyceUpmeier,JoyceTanakaUpmeier}; verify the reduced
Thom--Sebastiani identification on three-input flags.

\item \emph{Weyl-equivariant lifts \((O1)^+\).}  For each
\(i\in\{1,2,3\}\) construct the lift
\(\tau_{i,R,S}:o_{R,S}\to s_{\delta_i}^*o_{R,S}\); verify the
projective Weyl cocycle
\(c_{o,R}\in Z^2(\mathcal G_R;\underline{\mathbb F}_2)\) on the
retained action groupoid; null-trivialize \([c_{o,R}]\) and choose a
\(1\)-cochain killing it; verify \(\tau_{i,R,S}^2=1\); verify the
Coxeter relations
\((s_{\delta_i}s_{\delta_j})^3=1\) for \(i\ne j\) (the
\((3,3,3)\)-hyperbolic triangle relations of
Appendix~\ref{app:pfaffian-wall-three-walls}); verify that the torsor
defects \(\omega_{i,\mathcal C}\) of
Definition~\ref{thm:finite-stage-dirac-igusa-pro-object}~$(O_X)$ vanish.

\item \emph{Local wall-sign datum.}  At each generic type-II wall point
\(\Hpl_{\delta_i}\), with the rank-one local model
\(K_{\delta_i}^{\mathrm{cross}}=[\R\xrightarrow{u_i}\R]\), record the
local generic sign \(\epsilon_{o,R}^{\mathrm{loc}}(s_{\delta_i})=-1\)
of Proposition~\ref{prop:rank-one-crossing-o2-sign}.

\item Output
\((o_R,\mathfrak Q^{or}_R,\{\tau_{i,R,S}\},\epsilon_{o,R}^{\mathrm{loc}})\).
\end{enumerate}

\paragraph{Type checks.}
\(o_R\) is a Picard-line section squaring to the reduced determinant;
the four Cech--Borel classes vanish with compatible
null-trivializations; the lifts \(\tau_{i,R}\) satisfy
\(\tau_{i,R}^2=1\) and the \((3,3,3)\)-Coxeter relations; multiplicative
transport agrees on flag stacks.

\paragraph{Obstructions discharged or invoked.}
This step is the entire site of \((O1)\) (strong reduced orientation
on the K3\(\times\)E-quotient self-Ext frame bundle) and of
\((O1)^+\) (Weyl-equivariant transport along
\(W^{(2)}(\La^{2,1}_{II})\) reflections).  Discharging \((O1)\)
requires the simultaneous vanishing and compatible
null-trivialization of the four Cech--Borel classes on every retained
stratum; discharging \((O1)^+\) requires \([c_{o,R}]=0\), the
Coxeter coherence, and the vanishing of all torsor defects
\(\omega_{i,\mathcal C}\).  The local wall-sign datum is the input to
\((O2)\) at step~6.

\subsection{Construction~5 — Hall correspondences and primitive part}
\label{app:build-hall}

\paragraph{Tag.} \(\gamma\)-shadow specialization
(Stage-2 chiral Hall shadow).

\paragraph{Input.}
The oriented hybrid carrier from
Construction~\ref{app:build-orientation}; the extension stack
\(\mathfrak E^{\mathrm{geom}}_{R,\widehat c,\widehat c'}\); the
two-step flag stack
\(\mathfrak F^{(2),\mathrm{geom}}_{R,\widehat c,\widehat c',\widehat c''}\);
the BBDJS reduced Thom--Sebastiani transport; six-functor admissibility
of \(p^*,q_!,p_!\) (Construction~\ref{app:build-finite-moduli}~step
\textup{(Six)}).

\paragraph{Output.}
The finite Hall object
\[
\mathcal H_R^{\mathrm{geom}}
=\bigoplus_{\widehat c\in\widehat\Gamma_R}
H_c^\bullet\!\bigl(
\mathfrak M^{\mathrm{geom}}_{R,\widehat c},
\Phi_{R,c}\otimes o_{R,c}\bigr);
\]
the Hall product \(m_R=q_!\,\mathrm{TS}^{\mathrm{red}}\,p^*\) and Hall
coproduct \(\Delta_R=p_!\,(\mathrm{TS}^{\mathrm{red}})^{-1}\,q^*\); the
finite Hall bialgebra
\((\mathcal H_R^{\mathrm{geom}},m_R,\Delta_R,\eta_R,\epsilon_R)\); the
primitive part
\[
P_R^{\mathrm{geom}}
=\ker(\Delta_R-1\otimes\mathrm{id}-\mathrm{id}\otimes 1)\cap\ker\epsilon_R;
\]
the normal-ordered Gram descent
\(P_R^{\Pi,\mathrm{geom}}
=\overline\Pi_{X,*}^{\Theta_R}P_R^{\mathrm{geom}}\); the
\(\Gamma_R^\Pi\)-graded super-Lie bracket
\([\,\bar x,\bar y\,]_{\mathrm{red}}\) of
Proposition~\ref{thm:hall-primitives-formal};
the Hopf pairing matrix \(G_{\gamma,ab}\) and radical \(K_\gamma\).

\paragraph{Procedure.}
\begin{enumerate}
\item \emph{Pullback.}  For
\(\widehat c,\widehat c'\in\widehat\Gamma_R\), pull back along
\(p:\mathfrak E^{\mathrm{geom}}_{R,\widehat c,\widehat c'}
\to\mathfrak M^{\mathrm{geom}}_{R,\widehat c}
\times\mathfrak M^{\mathrm{geom}}_{R,\widehat c'}\); apply reduced
Thom--Sebastiani \(\mathrm{TS}^{\mathrm{red}}\) to combine vanishing
cycle data \cite{BBDJS2015}; pushforward exceptionally along
\(q:\mathfrak E^{\mathrm{geom}}_{R,\widehat c,\widehat c'}
\to\mathfrak M^{\mathrm{geom}}_{R,\widehat c\star\widehat c'}\) to
obtain
\(m_R(x\otimes y)=q_!\,\mathrm{TS}^{\mathrm{red}}\,p^*(x\boxtimes y)\).

\item \emph{Coproduct.}  Dually pull back along \(q\), apply
inverse reduced Thom--Sebastiani, pushforward along \(p\) to obtain
\(\Delta_R(z)=p_!\,(\mathrm{TS}^{\mathrm{red}})^{-1}\,q^*(z)\).
The compact geometric coproduct correspondence
\(\mathfrak C^{\mathrm{geom}}_{\rho;\mu,\nu}\) of
Proposition~\ref{thm:hall-primitives-formal} is the
support data for splitting.

\item \emph{Bialgebra check.}  Verify associativity of \(m_R\) and
coassociativity of \(\Delta_R\) on the two-step flag stack
\(\mathfrak F^{(2),\mathrm{geom}}_{R,\widehat c,\widehat c',\widehat c''}\)
by base change and Thom--Sebastiani coherence.  Verify the bialgebra
identity (Hopf compatibility) via the Massey--Saito Joyce--Upmeier
multiplicative orientation transport.

\item \emph{Primitive extraction.}  Compute
\(P_R^{\mathrm{geom}}=\ker(\Delta_R-1\otimes\mathrm{id}
-\mathrm{id}\otimes 1)\cap\ker\epsilon_R\); since
\(\Delta_R-1\otimes\mathrm{id}-\mathrm{id}\otimes 1\) is degree-graded
on \(\widehat\Gamma_R\), this reduces to a finite linear computation
in each charge.

\item \emph{Normal-ordered descent.}  Apply
\(\overline\Pi_{X,*}^{\Theta_R}\) of
Theorem~\ref{thm:hall-product} to obtain the
\(\Gamma_R^\Pi\)-graded
\(P_R^{\Pi,\mathrm{geom}}\); verify additivity
\(\overline\Pi_X(\widehat c\star\widehat c')
=\overline\Pi_X(\widehat c)+\overline\Pi_X(\widehat c')\).

\item \emph{Bracket and pairing matrices.}  Compute the matrices
\[
M^{\alpha+\beta,c}_{\alpha,a;\beta,b}
=\int_{\mathfrak E^{\mathrm{geom}}_{\alpha,\beta;\alpha+\beta}}
q^!(e_{\alpha+\beta,c}^\vee)\cap
\mathrm{TS}^{\mathrm{red}}(p^*(e_{\alpha,a}\boxtimes e_{\beta,b})),
\quad
B_{\alpha,\beta}(x\otimes y)
=M_{\alpha,\beta}(x\otimes y)
-(-1)^{|x||y|}M_{\beta,\alpha}\mathsf{sw}(x\otimes y),
\]
\[
G_{\gamma,ab}=
\mathrm{tr}_0(q_!\,\mathrm{TS}^{\mathrm{red}}\,p^*
(e_{\gamma,a}\boxtimes e_{-\gamma,b}))
\quad\hbox{on}\quad
\mathfrak P^{\mathrm{geom}}_{\gamma,-\gamma;0}.
\]
Form the radical
\(K_{\gamma,\bar p}
=\ker G_{\gamma,\bar p}\cap\ker G_{-\gamma,\bar p}^t\) and the
quotient \(L_{\gamma,\bar p}\); fix splittings.

\item Output
\((\mathcal H_R^{\mathrm{geom}},m_R,\Delta_R,P_R^{\mathrm{geom}},
P_R^{\Pi,\mathrm{geom}},
\{B_{\alpha,\beta}\},\{G_{\gamma,ab}\},\{K_\gamma\},\{L_\gamma\})\).
\end{enumerate}

\paragraph{Type checks.}
\(\mathcal H_R^{\mathrm{geom}}\) is a finite \(\widehat\Gamma_R\)-graded
Hall bialgebra; \(m_R\) is associative on the two-step flag stack;
\(\Delta_R\) is coassociative; the supercommutator of primitives is
primitive (the bialgebra identity); the descended bracket is
\(\Gamma_R^\Pi\)-graded after \(\overline\Pi_{X,*}^{\Theta_R}\);
finite-dimensional in each homogeneous parity piece by
finite-typeness.

\paragraph{Obstructions discharged or invoked.}
\((D_0)\) is partially discharged: the finite Hall package
\((\mathcal H_R^{\mathrm{geom}},m_R,\Delta_R,P_R)\) is constructed.
The seven obstruction classes
\(\mathfrak o_\Theta^{\mathrm{top}},
\mathfrak o_\Theta^{\mathrm{Hoch}},
\mathfrak o_\Theta^{\mathrm{tri}},
\mathfrak o_\Theta^{\mathrm{Jac}},
\mathfrak o_F,
\mathfrak o_\Theta^{\mathrm{cyc}},
\mathfrak o_\Theta^{\mathrm{rad}}\) of
Theorem~\ref{thm:hall-product} are required to vanish on
this finite Hall stage and compatibly in the inverse limit.

\paragraph{Cross-volume anchor.}  The output
\((\mathcal H_R^{\mathrm{geom}},m_R)\) is the finite-stage realization
of the K3-side BKM Hall--Drinfeld double
\(\mathcal D_\hbar(\mathrm{CoHA}_{K3\times E})\) of
\texttt{calabi-yau-quantum-groups:CLAUDE.md~§6.7}.  In the universal
stage chain \(\mathsf P\to\mathsf C\to\mathsf S\) the Hall positive half
\(Y^+(X)\) is precisely \(P_R^{\mathrm{geom}}\) at finite \(R\); the
Drinfeld double \(G(X)=D(Y^+(X))\) requires Hopf pairing, completion,
integral form, and stable-envelope transport which are supplied by
\(G_{\gamma,ab}\) and \((\mathrm{ML})\).  The κ-tuple coordinate
\(\kappa_{\mathrm{BKM}}=5\) is the weight of \(\Delta_5\) and the
output of the level-4 scalar trace; it is not a level-2 invariant of
the carrier itself.

\subsection{Construction~6 — Dirac operator and Pfaffian line}
\label{app:build-dirac-pfaffian}

\paragraph{Tag.} \(\gamma\)-shadow specialization
(orientation gerbe + Pfaffian).

\paragraph{Input.}
The oriented Hall package
\((\mathcal H_R^{\mathrm{geom}},o_R,\{\tau_{i,R,S}\})\) of
Constructions~\ref{app:build-orientation}--\ref{app:build-hall}; the
local wall-sign datum
\(\epsilon_{o,R}^{\mathrm{loc}}\) of
Construction~\ref{app:build-orientation}; the rank-one type-II
Pfaffian crossing datum of
Definition~\ref{def:rank-one-type-ii-crossing} and
Appendix~\ref{app:pfaffian-wall-rank-one}; the half-Hilbert orbit
under \(W^{(2)}(\La^{2,1}_{II})\) of size \(2^{12}=4096\).

\paragraph{Output.}
The Dirac operator on the orientation gerbe
\(D_R\); the Pfaffian line
\(\mathscr L_{\mathrm{Pf},R}=
\det\operatorname{Ext}^1_{\mathrm{red}}(A,A)^{-1}\) (squaring to the
reduced determinant); the protected Pfaffian section
\(\operatorname{pf}_{X,R}\); the global wall-sign character
\(\epsilon_{o,R}:W^{(2)}(\La^{2,1}_{II})\to\{\pm1\}\); the Pfaffian
section asymptotic
\[
\operatorname{pf}_{X,R}
=
64q^{1/2}r^{1/2}s^{1/2}
\prod_{\gamma=(n,l,m)\in A_R}
(1-q^nr^ls^m)^{f(nm,l)}+\hbox{higher-window correction},
\]
which exhausts to \(\Delta_5\) as \(R\to\infty\).

\paragraph{Procedure.}
\begin{enumerate}
\item \emph{Dirac operator.}  On each Darboux chart of
\(\mathfrak M^{\mathrm{sc}}_{R,c}\), with reduced normal self-Ext
splitting \(K^{\mathrm{nor}}_{R,c}\), build the odd Dirac operator
\(D_{R,c}\) acting on \(\Omega^{\mathrm{sp}}\otimes o_{R,c}\)
(spinors on the orientation gerbe); verify that, at a generic
type-II wall point, the local rank-one normal block is
\(D_{\delta_i}=\begin{psmallmatrix}0&u\\-u&0\end{psmallmatrix}\) with
\(\operatorname{Pf}(D_{\delta_i})=u\)
(Definition~\ref{def:rank-one-type-ii-crossing}).

\item \emph{Pfaffian line.}  Set
\(\mathscr L_{\mathrm{Pf},R}=
\det\operatorname{Ext}^1_{\mathrm{red}}(A,A)^{-1}\); verify
\(\mathscr L_{\mathrm{Pf},R}^{\otimes 2}\simeq
\det\operatorname{RHom}_{\mathrm{red}}(A,A)\) by
\(3\)-Calabi--Yau Serre duality.

\item \emph{Local wall sign.}  At each generic point of
\(\Hpl_{\delta_i}\) with normal model
\(K^{\mathrm{nor}}_{\delta_i,R}=[\R\xrightarrow{u_i}\R]\), compute
the monodromy of the Pfaffian section
\(\operatorname{pf}_{\delta_i,R}=\upsilon_iu_i\) along the
wall-reflection loop \(\ell_{\delta_i,R,S}\) closed by the
\((O1)^+\) lift \(\tau_{i,R,S}\):
\[
M_{\ell_{\delta_i,R,S}}(\upsilon_iu_i)
=\upsilon_i(-u_i)=-\upsilon_iu_i,
\qquad
\epsilon_{o,R}^{\mathrm{loc}}(s_{\delta_i})
=(-1)^{N^{\mathrm{Pf}}_{\delta_i,R}}=-1.
\]

\item \emph{Half-Hilbert orbit transport.}  The reflection
\(s_{\delta_i}\) acts on the half-Hilbert orbit (compact source labels
of size \(2^{12}=4096\) parametrized by 12 binary parities tied to the
\(2^{12}\) odd-spin structures on the lattice
\(\La^{2,1}_{II}/2\La^{2,1}_{II}\)).  Sum the local rank-one
contributions over the orbit:
\[
\epsilon_{o,R}(s_{\delta_i})
=\prod_{x\in\hbox{orbit}/2^{12}}
\epsilon_{o,R}^{\mathrm{loc}}(s_{\delta_i};x).
\]
The output is the global character
\(\epsilon_{o,R}:W^{(2)}(\La^{2,1}_{II})\to\{\pm 1\}\).

\item \emph{Compatibility package.}  Verify the component, boundary,
overlap, quotient, and protected integration compatibility package of
Proposition~\ref{prop:appE-local-pfaffian-sign} on every
retained type-II wall stratum.

\item \emph{Pfaffian section.}  Define
\(\operatorname{pf}_{X,R}\) as the protected integration of the
oriented Hall product over
\(\mathcal F_{X,\sigma,S,\le R}^{\mathrm{hyb}}\):
\[
\operatorname{pf}_{X,R}
=
64q^{1/2}r^{1/2}s^{1/2}
I_R^{\mathrm{prot}}\!\left(
\bigotimes_{\gamma\in A_R}\operatorname{pf}_{\gamma,R}\right),
\qquad
\operatorname{pf}_{\gamma,R}=(1-q^nr^ls^m)^{f(nm,l)}.
\]

\item Output
\((D_R,\mathscr L_{\mathrm{Pf},R},\operatorname{pf}_{X,R},
\epsilon_{o,R})\), with the protected integration compatibility data.
\end{enumerate}

\paragraph{Type checks.}
\(\mathscr L_{\mathrm{Pf},R}\) is a Picard-line whose square
isomorphism is supplied by reduced Serre duality;
\(\operatorname{pf}_{X,R}\) is a section of
\(\mathscr L_{\mathrm{Pf},R}\) with the displayed cusp expansion; the
character \(\epsilon_{o,R}\) takes values in \(\{\pm 1\}\) and
satisfies the Coxeter relations of
\(W^{(2)}(\La^{2,1}_{II})\); on each simple type-II reflection
\(\epsilon_{o,R}(s_{\delta_i})=-1\) by the rank-one local model.

\paragraph{Obstructions discharged or invoked.}
This step is the entire site of \((O2)\) (local Pfaffian wall normal
form on type-II walls; the \(2^{12}=4096\) sign sum on the
half-Hilbert orbit transport).  The local rank-one Pfaffian crossing
of Appendix~\ref{app:pfaffian-wall-rank-one} discharges \((O2)\)
generically; the global half-Hilbert orbit transport remains the open
\((O2)\) component of
Open~Problem~\ref{def:O1-plus-weyl-transport}.  The character
identity
\(\epsilon_o=\nu_{\Delta_5}\big|_{W^{(2)}(\La^{2,1}_{II})}\) becomes a
theorem only after \((D_0)\), \((O1)\), \((O1)^+\), and \((O2)\) are
discharged compatibly in \(R\).

\subsection{Construction~7 — Chiral Koszul source coalgebra}
\label{app:build-koszul}

\paragraph{Tag.} \(\gamma\)-shadow specialization
(\emph{Bar = twisting; not bulk}; the bulk is
\(Z^{\mathrm{der}}_{\mathrm{ch}}(\mathcal F_X^{\mathrm{hyb}})\),
constructed elsewhere).

\paragraph{Input.}
The oriented Hall primitives
\(P_R^{\Pi,\mathrm{geom}}\) of
Construction~\ref{app:build-hall}; the hybrid carrier
\(\mathcal F_{X,\sigma,S,\le R}^{\mathrm{hyb}}\) of
Construction~\ref{app:build-hybrid-wrapped}; the augmentation
\(\varepsilon_{X,R}^{\mathrm{hyb}}:
\mathcal F_{X,\sigma,S,\le R}^{\mathrm{hyb}}\to k\mathbf 1\); the
hybrid residual \(\mathfrak O_{\mathrm{hyb},R}=0\); the bounded
length \(N_R\) (positivity of HN height on every non-vacuum colour);
the Beilinson--Drinfeld/Francis--Gaitsgory bar-cobar formalism
\cite{BD,FrancisGaitsgoryCKD}; the conilpotent completed domain on
the elliptic curve \(E\).

\paragraph{Output.}
The conilpotent chiral coalgebra on \(E\):
\[
C_{X,R}^{\mathrm{geom}}
=\bar B_E^{\mathrm{ch}}
(\mathcal F_{X,\sigma,S,\le R}^{\mathrm{hyb}},
\varepsilon_{X,R}^{\mathrm{hyb}})
=k\mathbf 1\oplus\bigoplus_{1\le p\le N_R}
\bar B_E^{\mathrm{ch},p}(\overline{\mathcal F}_{X,R}^{\mathrm{geom}});
\]
the bar-length filtration
\(F^{\mathrm{co}}_{R,p}\); the chiral collision coproduct
\(\Delta_R^{\mathrm{ch}}\); the bar comparison
\(b_{X,R}=\mathrm{id}_{\bar B_E^{\mathrm{ch}}(\cdot)}\); the source
cobar counit
\(\varepsilon^{\mathrm{src}}_{\mathrm{bar},R}:
\Omega_E^{\mathrm{ch}}C_{X,R}^{\mathrm{geom}}
\xrightarrow{\simeq}
\mathcal F_{X,\sigma,S,\le R}^{\mathrm{hyb}}\); after a separate
source-to-target quasi-isomorphism
\(\vartheta_R:\mathcal F_{X,\sigma,S,\le R}^{\mathrm{hyb}}
\xrightarrow{\simeq}\mathsf{Den}_{\Delta_5,E,\le R}\), the cobar
comparison
\(\Theta_{\mathrm{Kos},R}=\vartheta_R\circ\varepsilon^{\mathrm{src}}_{\mathrm{bar},R}\).

\paragraph{Procedure.}
\begin{enumerate}
\item \emph{Augmentation kernel.}  Set
\(\overline{\mathcal F}_{X,R}^{\mathrm{geom}}
=\ker(\varepsilon_{X,R}^{\mathrm{hyb}})\); verify the bounded-length
hypothesis (positivity of HN height on every non-vacuum surviving
colour, hence at most \(N_R\) factors in any non-zero word).

\item \emph{Bar coalgebra.}  Form the reduced chiral bar
\[
C_{X,R}^{\mathrm{geom}}
=k\mathbf 1\oplus\bigoplus_{1\le p\le N_R}
\bar B_E^{\mathrm{ch},p}
(\overline{\mathcal F}_{X,R}^{\mathrm{geom}}),
\]
length \(p\) summand generated by ordered \(p\)-fold inputs from the
augmentation ideal of total HN height \(\le R\) (height \(>R\) outputs
zero); coaugmentation by the inclusion of the vacuum summand; counit by
projection.

\item \emph{Bar-length filtration.}
\[
F^{\mathrm{co}}_{R,p}C_{X,R}^{\mathrm{geom}}
=k\mathbf 1\oplus\bigoplus_{1\le j\le p}
\bar B_E^{\mathrm{ch},j}
(\overline{\mathcal F}_{X,R}^{\mathrm{geom}}),\qquad
0\le p\le N_R.
\]

\item \emph{Collision coproduct.}  For
\(x=[x_1|\cdots|x_p]\), define the cut-map coproduct
\[
\Delta_R^{\mathrm{ch}}(x)
=\mathbf1\otimes x+x\otimes\mathbf1
+\sum_{i=1}^{p-1}[x_1|\cdots|x_i]\otimes[x_{i+1}|\cdots|x_p],
\]
with each cut realized by the corresponding finite collision/flag
pull-push map in
\(\mathsf{Corr}^{E,\mathrm{hyb}}_{R,\mathrm{geom}}\) of
Construction~\ref{app:build-hybrid-wrapped}.  Verify
coassociativity by the four-input pentagon and the eight-word flag
coherences; verify counit identities via vacuum collision flags.

\item \emph{Conilpotence.}  Verify
\((\overline\Delta_R^{\mathrm{bar}})^{N_R+1}=0\): each reduced cut
strictly lowers bar length on each branch.

\item \emph{Bar comparison.}  Take \(b_{X,R}\) to be the identity of
\(\bar B_E^{\mathrm{ch}}
(\mathcal F_{X,\sigma,S,\le R}^{\mathrm{hyb}},
\varepsilon_{X,R}^{\mathrm{hyb}})\), eliminating the coalgebra-side
finite Koszul residuals
\(\mathfrak o^C_R, \mathfrak o^{\mathrm{fil}}_R,
\mathfrak o^{\mathrm{conil}}_R, \mathfrak o^{\Delta,\mathrm{ch}}_R\)
of Definition~\ref{def:chiral-koszul-source-datum}.

\item \emph{Source cobar counit.}  Apply Francis--Gaitsgory chiral
Koszul duality
\cite[Theorem~6.2.2]{FrancisGaitsgoryCKD} in the conilpotent completed
domain to obtain
\(\varepsilon^{\mathrm{src}}_{\mathrm{bar},R}:
\Omega_E^{\mathrm{ch}}C_{X,R}^{\mathrm{geom}}
\xrightarrow{\simeq}
\mathcal F_{X,\sigma,S,\le R}^{\mathrm{hyb}}\); verify it is a
source-internal quasi-isomorphism of chiral algebras on \(E\), not a
homotopy inverse of the target counit.

\item \emph{Source-to-target identification.}  Construct
\(\vartheta_R\) from the geometric primitive representatives, Hall
products, coproducts, pairings, and HN transition compatibility;
verify the residual
\(\mathfrak o^\vartheta_R\) of
Corollary~\ref{cor:source-bar-remaining-cobar-obstruction} vanishes;
set \(\Theta_{\mathrm{Kos},R}=\vartheta_R\circ
\varepsilon^{\mathrm{src}}_{\mathrm{bar},R}\).

\item \emph{Inverse limit.}  Verify the Koszul Mittag--Leffler
condition \(R^1\varprojlim=0\) for the inverse systems of source
coalgebras, completed chiral bar/cobar constructions, target
truncations, and the cones of \(b_{X,R}\) and
\(\Theta_{\mathrm{Kos},R}\) \cite[\S3.5]{WeibelHA}.  Form
\[
C_X^{\mathrm{geom}}
=\varprojlim_R C_{X,R}^{\mathrm{geom}},\qquad
\Theta_{\mathrm{Kos}}:
\Omega_E^{\mathrm{ch}}C_X^{\mathrm{geom}}
\xrightarrow{\simeq}\mathsf{Den}_{\Delta_5,E}.
\]

\item Output
\((C_{X,R}^{\mathrm{geom}},F^{\mathrm{co}}_{R,\bullet},
\Delta_R^{\mathrm{ch}},b_{X,R},\Theta_{\mathrm{Kos},R})\) at every
\(R\), with successor compatibility.
\end{enumerate}

\paragraph{Type checks.}
\(C_{X,R}^{\mathrm{geom}}\) is a coaugmented conilpotent chiral
coalgebra on \(E\) with bounded bar-length filtration of length at
most \(N_R\); \(\Delta_R^{\mathrm{ch}}\) is coassociative and counital
on the chosen flag stacks; the bar-cobar inversion is supplied by
Francis--Gaitsgory chiral Koszul; the cobar comparison
\(\Theta_{\mathrm{Kos},R}\) is source-internal, not the target
bar-cobar counit (the source/target separation of
Lemma~\ref{lem:source-target-koszul-separation}).

\paragraph{Obstructions discharged or invoked.}
The Koszul Mittag--Leffler exactness hypothesis is the
\textup{(D0)}-component for the chiral Koszul lane.  The total finite
Koszul residual
\(\mathfrak O_{\mathrm{Kos},R}=
(\mathfrak o^C_R,\mathfrak o^{\mathrm{fil}}_R,
\mathfrak o^{\mathrm{conil}}_R,\mathfrak o^{\Delta,\mathrm{ch}}_R,
\{\mathfrak o^{C,\mathrm{tr}}_{R'R}\},\mathfrak o^b_R,
\mathfrak o^\Omega_R,\mathfrak o^{\mathrm{hyb}}_R,
\mathfrak o^\Pi_R,\mathfrak o^{\mathrm{sep}}_R,
\mathfrak o^{\Prim}_R)\)
vanishes after this step.

\paragraph{Cross-volume firewall.}  The bar coalgebra
\(\bar B_E^{\mathrm{ch}}
(\mathcal F_{X,R}^{\mathrm{hyb}})\) is a \emph{twisting / coupling}
coalgebra (Maurer--Cartan classification), \emph{not} the bulk
\(Z^{\mathrm{der}}_{\mathrm{ch}}(\mathcal F_X^{\mathrm{hyb}})
=\mathrm{ChirHoch}^\bullet(\mathcal F_X^{\mathrm{hyb}},
\mathcal F_X^{\mathrm{hyb}})\); cf.\
\texttt{chiral-bar-cobar-vol2:CLAUDE.md~§6} (universal stage chain
\(\mathsf P\to\mathsf C\to\mathsf S\to\mathsf Z\)).
The chiral Koszul comparison
\(\Theta_{\mathrm{Kos}}:\Omega_E^{\mathrm{ch}}C_X^{\mathrm{geom}}
\xrightarrow{\simeq}\mathsf{Den}_{\Delta_5,E}\) is a comparison
between two algebras of the same level; the homotopy inverse to the
target-internal counit
\(\Omega_E^{\mathrm{ch}}\bar B_E^{\mathrm{ch}}
(\mathsf{Den}_{\Delta_5,E})\to\mathsf{Den}_{\Delta_5,E}\) is not used
in either direction.

\subsection{Construction~8 — Primitive recognition}
\label{app:build-recognition}

\paragraph{Tag.} \(\delta\)-endpoint convergence
(source-to-target identification).

\paragraph{Input.}
The descended source primitive Lie superalgebra
\(P_X^{\Pi,\mathrm{red}}\) of
Construction~\ref{app:build-hall} (after the radical quotient
\(\overline{\mathrm{Rad}\langle\,\cdot,\cdot\rangle_X}\)); the imported
target denominator algebra
\(\fg_{\Delta_5}=G(\mathscr B_{\Delta_5})\) on
\(\La^{2,1}_{II}\) of Definition~\ref{def:bkm-imaginary-simple-roots} and
Definition~\ref{def:bkm-relations}; the increasing cofinal
sequence of finite downward-saturated windows
\(W_1\subset W_2\subset\cdots\) with
\(\bigcup_\nu W_\nu=\mathcal R_+\); the Borcherds--Kac--Moody data
\((H,I,p,\alpha_i,(\,,\,))\) of
Theorem~\ref{thm:primitive-recognition}.

\paragraph{Output.}
The cofinal finite-window primitive-recognition datum
\(\{\textup{(R0),(R1),(R2),(R3),(R4),(R5),(R6),(R7),(R8)}\}\) of
Definition~\ref{def:cofinal-primitive-recognition-datum} for every
\(W_\nu\), compatible under \(W_\nu\subset W_{\nu+1}\); the source
representatives
\(e_i^X,f_i^X,h_i^X\); the source presentation map
\(\pi_W:\widehat{\mathfrak F}(\mathscr B_{\Delta_5})_W
\to P_X^{\Pi,\mathrm{red}}\big|_W\); the kernel equality
\(\ker\pi_W=
(\overline{J_{\mathrm{BK}}+\mathrm{Rad}_{\mathrm{GN}}})_W\); the
isomorphism
\[
\boxed{\quad
P_X^{\Pi,\mathrm{red}}\cong\fg_{\Delta_5},\quad
\Omega_E^{\mathrm{ch}}C_X\simeq\mathsf{Den}_{\Delta_5,E}.
\quad}
\]

\paragraph{Procedure.}
For every cofinal \(W_\nu\):
\begin{enumerate}
\item[\textup{(R0)}] \emph{Compact source labels.}  Verify every
finite source basis vector in \(W_\nu\cup(-W_\nu)\cup\{0\}\) satisfies
the compact source-admissibility condition of
Theorem~\ref{thm:finite-stage-dirac-igusa-pro-object}.

\item[\textup{(R1)}] \emph{Representatives and parities.}  Supply
parity-homogeneous primitive representatives
\(e_i^X\in P^{\Pi,\mathrm{red}}_{X,\alpha_i}\),
\(f_i^X\in P^{\Pi,\mathrm{red}}_{X,-\alpha_i}\),
\(h_i^X=\iota_H^{-1}(\alpha_i)\) for every simple-root index \(i\) with
\(\alpha_i\in W_\nu\), with the Gritsenko--Nikulin parity fibres
\(\mu_{\overline 0}(a),\mu_{\overline 1}(a)\) for imaginary simple
roots and even real representatives for \(\delta_1,\delta_2,\delta_3\).

\item[\textup{(R2)}] \emph{Relations.}  Verify in the protected Hall
super-bracket the Cartan, Chevalley, real Serre, Borcherds
orthogonality, and super sign relations for every relation whose
displayed source and target degrees lie in
\(W_\nu\cup(-W_\nu)\cup\{0\}\).  In the real rank-three subsystem the
Cartan matrix is
\(((\delta_i,\delta_j))=\mathrm{diag}(2)-2(J-\mathrm{diag}(1))\) with
real Serre exponent~3.

\item[\textup{(R3)}] \emph{Full finite parity dimensions.}  For every
\(\beta\in W_\nu\) verify
\(\dim P^{\Pi,\mathrm{red}}_{X,\beta,\overline p}
=\rootmult_{\overline p}^{\mathrm{GN}}(\beta)\) for both parities
\(p=0,1\).  In the first timelike channel
\(\beta=\delta_{123}\),
Proposition~\ref{prop:bkm-delta123-presentation-count} requires
\((d^{\mathrm{GN}}_{(1,1,1),\overline 0},
d^{\mathrm{GN}}_{(1,1,1),\overline 1})=(29,93)\).  Determinant
information \(29-93=-64=f(1,1)\) is consistent but not the
recognition.

\item[\textup{(R4)}] \emph{Hopf pairing and radical.}  Compute the
positive-negative pairing matrices in parity blocks; verify the
kernel agrees with the
Gritsenko--Nikulin/Kac invariant-pairing kernel; verify Lie-ideal and
coideal stability under all finite brackets; verify carrying under
HN transitions to the next window.

\item[\textup{(R5)}] \emph{No extra relations.}  Verify the kernel
equality
\(\ker\pi_\nu=(\overline{J_{\mathrm{BK}}+\mathrm{Rad}_{\mathrm{GN}}})_{\le W_\nu}\) on the finite triangular span generated by
\(W_\nu\cup(-W_\nu)\cup\{0\}\).

\item[\textup{(R6)}] \emph{Generation.}  Verify that every source
primitive root space in \(W_\nu\cup(-W_\nu)\), after the finite
radical quotient, is generated by the Cartan and the supplied
simple-root representatives through brackets whose intermediate
positive degrees stay in \(W_\nu\).

\item[\textup{(R7)}] \emph{Completed PBW compatibility.}  Verify the
PBW-graded isomorphism
\(\mathrm{gr}_{\mathrm{PBW}}U(P_X^{\Pi,\mathrm{red}})_{\le W_\nu}
\xrightarrow{\sim}
\mathrm{gr}_{\mathrm{PBW}}U(\fg_{\Delta_5})_{\le W_\nu}\); verify
strict preservation under \(W_{\nu+1}\to W_\nu\) transition maps.

\item[\textup{(R8)}] \emph{Exact triangular completion.}  Verify
strictness of the transition maps and closedness of images on every
finite window, so that passage to the inverse limit commutes with the
finite kernels, cokernels, radicals, PBW associated gradeds, and
parity pieces.  Equivalently the relevant \(\lim^1\) terms vanish for
these triangular finite-window systems.

\item Output the global isomorphism
\(P_X^{\Pi,\mathrm{red}}\cong\fg_{\Delta_5}\) by passage to the inverse
limit (compatible cofinal datum); together with the chiral Koszul
comparison of Construction~\ref{app:build-koszul} this gives
\(\Omega_E^{\mathrm{ch}}C_X\simeq\mathsf{Den}_{\Delta_5,E}\) and hence
the protected Pfaffian identification
\(\operatorname{Pf}_{\mathrm{prot}}(\mathfrak D_X^{\mathrm{DI}})=\Delta_5\).
\end{enumerate}

\paragraph{Type checks.}
The source presentation map \(\pi:\widehat{\mathfrak F}(\mathscr B_{\Delta_5})\to P_X^{\Pi,\mathrm{red}}\) is
continuous and degree-preserving; \textup{(R5)} gives injectivity
after the completed radical quotient; \textup{(R6)} gives
surjectivity; \textup{(R7)} gives PBW compatibility; \textup{(R8)}
gives exactness in the inverse-limit step.  No claim on signed
superdimensions \(\sdim=f(nm,l)\) substitutes for the two-integer
parity dimensions of \textup{(R3)}.

\paragraph{Obstructions discharged or invoked.}
This step inherits the obstruction tower discharged by
Constructions~1--7 and adds the cofinal exact-completion residual
\textup{(R8)}.  The Beilinson--Drinfeld bar-cobar firewall is
respected throughout: \textup{(R5)} is a source kernel-equality
theorem, not an inversion of the target bar-cobar counit;
\textup{(R7)} is a PBW graded comparison of source and target
enveloping algebras, not a graph-equality argument from the scalar
square.

\paragraph{Cross-volume anchor.}  The output is
\texttt{chiral-bar-cobar-vol2:CLAUDE.md~§9} Construction
Problem~1: the operator-level construction of
\(\mathfrak D_X^{\mathrm{DI}}\) for \(K3\times E\) with
\(\operatorname{Pf}_{\mathrm{prot}}(\mathfrak D_X^{\mathrm{DI}})=\Delta_5\), under the
four obstructions \((D_0),(O1),(O1)^+,(O2)\).  The Drinfeld doubling
\(Y^+(X)\to G(X)=D(Y^+(X))\) of
\texttt{calabi-yau-quantum-groups:CLAUDE.md} requires the Hopf
pairing of \textup{(R4)}, completion by \textup{(R8)}, integral form,
and stable-envelope transport which are supplied by the same
data.  The κ-tuple coordinate \(\kappa_{\mathrm{BKM}}=5\) is the
weight of \(\Delta_5\) and lives at level~4 of the universal stage
chain; it is a consequence of the construction, not an input to it.

\subsection{The eight-step audit}
\label{app:build-audit}

The eight construction steps form a strict ordered sequence of
input/output type discipline.  An audit of any candidate construction of
\(\mathfrak D_X^{\mathrm{DI}}\) reduces to verifying, for every
\(R\in\mathcal R_{\mathrm{HN}}\):
\begin{enumerate}
\item the lattice combinatorics of
Construction~\ref{app:build-charge-window};
\item the finite-type quasi-smooth derived stacks of
Construction~\ref{app:build-finite-moduli};
\item the hybrid local/wrapped factorization datum and the eight-word
two-step flag atlas of Construction~\ref{app:build-hybrid-wrapped};
\item the four Cech--Borel orientation classes and the
\((O1)^+\)~Coxeter coherence of
Construction~\ref{app:build-orientation};
\item the finite Hall bialgebra and primitive matrices of
Construction~\ref{app:build-hall};
\item the rank-one Pfaffian crossing and the half-Hilbert orbit
transport of Construction~\ref{app:build-dirac-pfaffian};
\item the bar coalgebra, conilpotent collision coproduct, and the
Francis--Gaitsgory cobar quasi-isomorphism of
Construction~\ref{app:build-koszul};
\item the cofinal datum
\(\textup{(R0)}\)--\(\textup{(R8)}\) of
Construction~\ref{app:build-recognition};
\end{enumerate}
together with the compatible vanishing of the obstruction tuples
\((\mathfrak O^{E_3-\mathrm{src}}_R,
\mathfrak O_{\mathrm{hyb},R},
\mathfrak O_{\mathrm{Kos},R},
\mathfrak O^{\mathrm{O1}}_R,
\mathfrak O^{\mathrm{O1}^+}_R,
\mathfrak O^{\mathrm{O2}}_R,
\mathfrak O_R^{D_0})\) at every height \(R\), and the
Mittag--Leffler exactness condition \textup{[ML]} of
Definition~\ref{def:O1-strong-reduced-orientation} on the cofinal HN inverse
system.  Each obstruction is attributable to the unique step
in which its components are constructed.

\paragraph{Order rigidity.}  Each of the seven adjacent input/output
dependencies fails type-checking under exchange:
\begin{enumerate}
\item charge window $\to$ finite moduli: $\Lambda_{\mathrm{ample}}$
  parametrises the HN strata; without the window, the moduli stack
  has no finite cofinal subsystem.
\item finite moduli $\to$ hybrid carrier: $\mathrm{Ran}^{\mathrm{hyb}}(E)$
  needs the finite Hall stack as factor on which to graft the
  wrapped/local atlas.
\item hybrid carrier $\to$ reduced orientation: the cosection-twisted
  Joyce--Upmeier line bundle is defined on the carrier, not on the
  raw moduli (the cosection lives on $\mathrm{Ran}^{\mathrm{hyb}}$).
\item reduced orientation $\to$ Hall correspondences: the Pfaffian
  line bundle on the Hall stack is the orientation-twisted sign data;
  without orientation, the Hall product is unsigned.
\item Hall correspondences $\to$ Pfaffian line + Dirac operator: the
  primitive Hall bracket $[-,-]_X$ defines the Hall--Drinfeld
  correspondence whose Pfaffian section is the Dirac--Igusa Pfaffian.
\item Pfaffian line $\to$ chiral Koszul source coalgebra: the source
  $C_X$ is the bar of the augmented hybrid factorization algebra
  pulled back along the Pfaffian comparison
  $\iota_{\mathrm{aut}}$.
\item chiral Koszul source $\to$ primitive recognition: the primitive
  part $P_X^{\Pi,\mathrm{red}}$ is the Frobenius--Serre primitive
  filtration on the source, which the recognition theorem matches to
  $\mathfrak g_{\Delta_5}$.
\end{enumerate}
Two diagnostic exchanges: orientation before hybrid leaves the middle
type-II wall \(\delta_2\leftrightarrow(0,1,1)\) without a
representative, by
Lemma~\ref{app:pfaffian-wall-three-walls}; recognition before the
Koszul cobar leaves the identification
\(\Omega_E^{\mathrm{ch}}C_X\simeq\mathsf{Den}_{\Delta_5,E}\)
unconstructed.  All other adjacent exchanges fail by the same
input/output type criterion.

\paragraph{Conditional theorem.}  Granted the eight constructions and
the four obstructions \((D_0),(O1),(O1)^+,(O2)\),
Theorem~\ref{thm:pfaffian-dirac-finite-stage} reads
\[
\operatorname{Pf}_{\mathrm{prot}}(\mathfrak D_X^{\mathrm{DI}})=\Delta_5,\qquad
\epsilon_o=\nu_{\Delta_5}\big|_{W^{(2)}(\La^{2,1}_{II})},\qquad
P_X^{\Pi,\mathrm{red}}\cong\fg_{\Delta_5}.
\]
The scalar square
\(Z_{\mathrm{BPS}}^{K3\times E}=
(\Phi_{10}^{\mathrm{un}})^{-1}=\Delta_5^{-2}\) is the level-4 trace
of the constructed protected algebra; promotion to the 3d
gravitational path integral is not claimed and requires the
Hall--Borcherds residual of
\texttt{chiral-bar-cobar-vol2}.

%%% End of Appendix F.
 after
% \appendix has been opened at main.tex line 17081, so that
% \thesection = \Alph{section}.  No \chapter, no second \appendix.

\section{Eight build algorithms for the Dirac--Igusa pro-object}
\label{app:eight-build-algorithms}

The Dirac--Igusa pro-object
\[
\mathfrak D_X^{\mathrm{DI}}
=\varprojlim_R
\bigl(
\mathcal A^{E_3}_{X,R},
\mathcal F_{X,R}^{\mathrm{hyb}},
\Gamma_{X,R},
\Pi_X,
\Phi_R,
o_R,
\mathcal H_R,
P_R^\Pi,
C_{X,R},
\Theta_{\mathrm{Kos},R},
\mathscr L_{\mathrm{Pf},R},
\operatorname{pf}_{X,R},
\epsilon_{o,R},
\mathrm{Rec}_R
\bigr)
\]
of Definition~\ref{thm:finite-stage-dirac-igusa-pro-object} and
Definition~\ref{def:O1-strong-reduced-orientation} is the conditional source
operator-level package whose protected Pfaffian, after the four
obstructions \((D_0)\), \((O1)\), \((O1)^+\), \((O2)\) are discharged,
realizes the Igusa cusp form,
\[
\operatorname{Pf}_{\mathrm{prot}}(\mathfrak D_X^{\mathrm{DI}})=\Delta_5,\qquad
\epsilon_o=\nu_{\Delta_5}\big|_{W^{(2)}(\La^{2,1}_{II})},\qquad
P_X^{\Pi,\mathrm{red}}\cong\fg_{\Delta_5}.
\]
Construction at finite Harder--Narasimhan height \(R\) proceeds through
eight construction steps.  Each step is a procedure on input data
already supplied by earlier steps and feeds named output data into the
next.  The eight steps stratify the realization datum
\(\mathfrak K_X=(L_X,H_X,O_X,D_X,R_X)\) of
Definition~\ref{thm:finite-stage-dirac-igusa-pro-object}: steps 1--3
realize \((L_X)\) and \((H_X)\); step 4 realizes the local content of
\((O_X)\); step 5 produces the Hall correspondence object; step 6 fixes
the Pfaffian--Dirac line and feeds the orientation character; steps 7--8
realize \((D_X)\) and \((R_X)\).

Throughout, \(X=S\times E\) with \(S\) a smooth complex K3 surface and
\(E\) a smooth elliptic curve.  \(\sigma=\sigma_{s,t}\) is Liu's product
stability condition on the Abramovich--Polishchuk heart in
\(\mathsf{D}^b(\mathrm{Coh}(X))\)
\cite{AbramovichPolishchukTStructures,LiuProductStability}.  The HN
sector \((\sigma,S)\) is fixed and the cofinal subsystem
\(\mathcal R_{\mathrm{HN}}=\{R_0<R_1<\cdots\}\) is the one of
Definition~\ref{def:O1-strong-reduced-orientation}, governed by
\textup{[K3$\times$E-fin]}, \textup{[K3$\times$E-atlas]}, and
\textup{[ML]}.

Each construction step carries a licensing tag pulled from the
five-type universal stage chain
\(\mathsf P\to\mathsf C\to\mathsf S\to\mathsf Z\to\mathsf A\) of
\texttt{chiral-bar-cobar-vol2:CLAUDE.md~§9}: \(\alpha\) primitive,
\(\beta\) functorial passage, \(\gamma\) shadow specialization,
\(\delta\) endpoint convergence, \(\epsilon\) terminal scalar.  The
ordering charge\(\to\)moduli\(\to\)hybrid\(\to\)orientation\(\to\)Hall
\(\to\)Dirac+Pfaffian\(\to\)Koszul\(\to\)recognition is the unique
ordering compatible with input/output type discipline; reordering any
two adjacent steps fails an input check.

The four obstructions are attributed to specific steps:
\((D_0)\) is the cofinal-trivialization datum of the entire eight-step
construction at every finite \(R\) (steps 1--6 supply its components);
\((O1)\) lives at step 4; \((O1)^+\) lives at step 4 with a Coxeter
coherence check feeding step 6; \((O2)\) lives at step 6.  Step 8
inherits the obstruction tower discharged by steps 1--7 and adds the
exact-completion residual \textup{(R8)} of
Definition~\ref{def:cofinal-primitive-recognition-datum}.

\subsection{Construction~1 — Charge window}
\label{app:build-charge-window}

\paragraph{Tag.} \(\alpha\)-primitive (lattice datum).

\paragraph{Input.}
HN height \(R\in\mathcal R_{\mathrm{HN}}\); the formal dyonic Mukai
lattice \(\Gamma_X^{\mathrm{form}}\) of
Definition~\ref{def:gram-index-lattice}; the Gram map
\(\Pi_X:\Gamma_X^{\mathrm{form}}\to\Gamma_{\mathrm{gram}}\) and the
cocycle \(B\) of Lemma~\ref{lem:gram-additivity-defect}; the active
support \(\Gamma_{\mathrm{act}}=\{\gamma\in\Gamma_{\mathrm{eff}}\mid
f(nm,l)\ne0\}\) of the Igusa Jacobi form \(\phi_{0,1}\); the cofinal
subsystem \(\mathcal R_{\mathrm{HN}}\).

\paragraph{Output.}
A finite cusp-positive active window
\[
A_R\subset\Gamma_{\mathrm{act}}\subset\Gamma_{\mathrm{gram}},
\qquad
A_R\hbox{ downward saturated for the cofinal target-window order},
\]
together with the normal-ordered extension
\[
\widehat\Gamma_R=
\{(c,T)\mid c\in\Gamma_R^{HN},\ T\in\mathcal T_R(c)\}
\subset\widehat\Gamma_X,
\qquad
\Gamma_R^\Pi=\overline\Pi_X(\widehat\Gamma_R)\subset\Gamma_{\mathrm{gram}},
\]
the lift \(L:\Gamma_{\mathrm{gram}}\to\Gamma_X^{\mathrm{form}}\) of
Lemma~\ref{lem:primitive-lift-every-gram-triple}, and the finite test
window \(\Gamma_R^{\mathrm{test}}\) of
Definition~\ref{def:O1-strong-reduced-orientation}.

\paragraph{Procedure.}
\begin{enumerate}
\item Compute the active support cone
\[
\Gamma_{\mathrm{act}}=\{(n,l,m)\in\Gamma_{\mathrm{eff}}\mid f(nm,l)\ne 0\}
\]
from the Fourier expansion of the index-one weak Jacobi form
\(\phi_{0,1}(\tau,z)=\sum_{n,l}c(n,l)q^nr^l\) of
\cite{EZ:1}, with the index substitution
\(f(nm,l):=c(nm,l)\) tabulated in
Appendix~\ref{sec:n1-boundary-compatibility-conditions} for the
\(N=1\) Igusa boundary.

\item Set
\[
A_R^{\circ}
=\{\gamma=(n,l,m)\in\Gamma_{\mathrm{act}}\mid n+|l|+m\le R\}.
\]
Form its downward saturation
\[
A_R
=\{\gamma'\in\Gamma_{\mathrm{act}}\mid
\exists\gamma\in A_R^{\circ},\ 0<\gamma'\le\gamma,\
\gamma-\gamma'\in Q_+\},
\]
in the sense of Definition~\ref{def:cofinal-primitive-recognition-datum}.

\item Normal-ordered lift.  For each \(c\in\Gamma_R^{HN}\) compute the
\(B\)-translate orbit
\[
\mathcal T_R(c)
=\Bigl\{\sum_{i<j}B(c_i,c_j)\Bigm|c=\sum_i c_i,\
c_i\in F_{\sigma,S}(R),\ \sum_i h_S(c_i)\le R\Bigr\}.
\]
Set \(\widehat\Gamma_R=\{(c,T):c\in\Gamma_R^{HN},\ T\in\mathcal T_R(c)\}\).

\item Finite target test window:
\(\Gamma_R^{\mathrm{test}}=\overline\Pi_X(\widehat\Gamma_R)\), compatibly
with \(R\to R'\) by inclusion.

\item Output \((A_R,\widehat\Gamma_R,\Gamma_R^\Pi,\Gamma_R^{\mathrm{test}},L)\).
\end{enumerate}

\paragraph{Type checks.}
\(A_R\) is finite and downward saturated;
\(\overline\Pi_X(\widehat\Gamma_R)=\Gamma_R^\Pi\) lies in
\(\Gamma_{\mathrm{gram}}\); \(\Gamma_R^{HN}\subset
F_{\sigma,S}(R)\) is finite by \textup{[K3$\times$E-fin]}; the cocycle
identity
\(\Pi_X(c+c')=\Pi_X(c)+\Pi_X(c')+B(c,c')\) is an axiom of the cocycle.

\paragraph{Obstructions discharged or invoked.}
None; this step is purely lattice combinatorics from the imported
Gritsenko--Nikulin Borcherds product
\cite{GN,GNPartI,GNII,BorcherdsGKM88,Borcherds95}.
The downward saturation is the input requirement of \textup{(R0)} in
Definition~\ref{def:cofinal-primitive-recognition-datum}.

\subsection{Construction~2 — Finite moduli}
\label{app:build-finite-moduli}

\paragraph{Tag.} \(\alpha\)-primitive (geometric source).

\paragraph{Input.}
The active window \(A_R\) and lift \(L\) of
Construction~\ref{app:build-charge-window}; Liu's product stability
condition \(\sigma_{s,t}\); the standing hypotheses
\textup{[K3$\times$E-fin]} and \textup{[K3$\times$E-atlas]}.

\paragraph{Output.}
A finite class set \(C_R\) and lift
\(\ell_R:\Gamma_R^\Pi\to C_R\) with
\(\Pi_X(\ell_R(\gamma))=\gamma\); the family of finite-type quasi-smooth
derived Artin substacks
\[
\mathfrak M^{ss}_{R,c}
=\mathfrak M^{ss}_{\sigma_{s,t}}(X,c)_{\le R}
\subset\operatorname{Perf}(X),
\qquad c\in C_R,
\]
of \(\sigma_{s,t}\)-semistable objects of full Liu numerical class
\(c\) and HN height \(\le R\); the scalar-rigidified truncation
\(\mathfrak M^{\mathrm{sc}}_{R,c}\); the universal perfect complex
\(\mathcal A_{R,c}\) on \(X\times\mathfrak M^{\mathrm{sc}}_{R,c}\); the
finite stratification by retained closed substacks of finite residual
inertia.

\paragraph{Procedure.}
\begin{enumerate}
\item Form \(C_R=\{c=L\gamma:\gamma\in A_R\}\) closed under retained
extension sums (taken in \(\Gamma_X^{\mathrm{form}}\)).

\item For each \(c\in C_R\) cut the substack of
\(\sigma_{s,t}\)-semistable objects of class \(c\) by HN height \(\le
R\):
\[
\mathfrak M^{ss}_{R,c}=
\bigl\{[A]\in\operatorname{Perf}(X)\mid
\sigma_{s,t}\hbox{-semistable},\
\hn(A)\le R,\
\nu(A)=c\bigr\}.
\]

\item Verify finite-typeness: \textup{[K3$\times$E-fin]} (Liu product
stability + product-boundedness theorem at fixed full Liu charge with
finite fibres over the active charge) gives the finite-type assertion.
PTVV \cite{PTVV2013} on the Calabi--Yau threefold \(X\) gives the
\((-1)\)-shifted symplectic structure on
\(\operatorname{Perf}(X)\), restricted to the retained semistable
sector.

\item Quasi-smoothness: the universal perfect complex
\(\mathcal A_{R,c}\) supplies a relative cotangent complex of bounded
Tor-amplitude on the retained charts.  Tangent at \(A\) is
\(T_A\mathfrak M^{ss}_{R,c}\simeq\operatorname{RHom}_X(A,A)[1]\).

\item Scalar rigidification: form \(\mathfrak M^{\mathrm{sc}}_{R,c}\) by
killing the scalar \(\mathbb G_m\)-automorphism of
\(c\)-stable factors.

\item Two-step flag stack
\(\mathfrak F^{(2)}_{R,c,c',c''}\) and extension stack
\(\mathfrak E_{R,c,c'}\) constructed as iterated derived vector stacks
of \(R\pi_*\mathcal RHom_X(\mathcal B,\mathcal A[1])\); both finite type
and quasi-smooth in the retained sector by
Proposition~\ref{prop:finite-dcritical-hall-product}.

\item Output \((C_R,\ell_R,\{\mathfrak M^{ss}_{R,c}\},
\{\mathfrak E_{R,c,c'}\},\{\mathfrak F^{(2)}_{R,c,c',c''}\})\) with
universal perfect complexes and the finite stratification.
\end{enumerate}

\paragraph{Type checks.}
\(\mathfrak M^{ss}_{R,c}\) is a finite-type quasi-smooth derived Artin
substack of \(\operatorname{Perf}(X)\); the scalar rigidification leaves
only finite residual inertia
\(H_S\subset E_{\mathrm{tors}}\) on each connected stabilizer
component, by
Proposition~\ref{prop:finite-dcritical-hall-product};
PTVV \((-1)\)-shifted symplectic structure restricts.

\paragraph{Obstructions discharged or invoked.}
\textup{[K3$\times$E-fin]} is invoked here in its product-boundedness
form and supplies the finite-type assertion not in the cited literature
for the K3\(\times\)E threefold; it is part of the standing datum.
\textup{[K3$\times$E-atlas]} supplies the quasi-smooth d-critical
charts.  These enter \((D_0)\) at the level of the finite Hall atlas
of Theorem~\ref{thm:finite-stage-dirac-igusa-pro-object}.

\subsection{Construction~3 — Hybrid wrapped factorization carrier}
\label{app:build-hybrid-wrapped}

\paragraph{Tag.} \(\beta\)-functorial passage (Stage-1 carrier).

\paragraph{Input.}
The finite moduli stacks
\(\{\mathfrak M^{\mathrm{sc}}_{R,c}\}\) of
Construction~\ref{app:build-finite-moduli}; the
projection-finite/wrapped split by elliptic degree
\(b_R^{\mathrm{geom}}:\Gamma_{\sigma,S}^{HN}(R)\to\Z_{\ge0}\) of
Definition~\ref{def:hybrid-wrapped-factorization-datum}; a fixed
origin \(0_E\in E\); the Abramovich--Polishchuk/Liu wrapped anchor
candidate.

\paragraph{Output.}
The finite hybrid local/wrapped factorization datum
\[
\mathcal F_{X,\sigma,S,\le R}^{\mathrm{hyb}}
=
(\mathcal C_{\sigma,S,\le R},
b_R^{\mathrm{geom}},
\mathcal F_R^{\mathrm{loc}},
\mathcal G_R^{\mathrm{wr}},
\mathfrak E_R^{\mathrm{loc}},
\mathfrak E_R^{\mathrm{hyb}},
\mathfrak E_R^{\mathrm{wr}},
\mathfrak{Flag}_R^{\mathrm{hyb}},
\mathsf{Corr}_R^{E,\mathrm{hyb}},
\lambda_R,
Q_{E,R},
\theta^Q_R,
o_R,
I_R^{\mathrm{prot}})
\]
of Definition~\ref{def:hybrid-wrapped-factorization-datum},
together with the hybrid Ran base
\(\operatorname{Ran}^{\mathrm{hyb}}(E)\), the eight-word two-step
binary flag atlas \(\{w\in\{\mathrm L,\mathrm W\}^3\}\), and the
quotient-after-correspondence functor \(Q_{E,R}\).

\paragraph{Procedure.}
\begin{enumerate}
\item \emph{Elliptic-degree split.}  Decompose
\(\Gamma_{\sigma,S}^{HN}(R)
=\Gamma_R^{\mathrm{loc}}\sqcup\Gamma_R^{\mathrm{wr}}\) where
\(\Gamma_R^{\mathrm{loc}}=\{b_R^{\mathrm{geom}}=0\}\) and
\(\Gamma_R^{\mathrm{wr}}=\{b_R^{\mathrm{geom}}>0\}\).  Verify additivity
\(b_R^{\mathrm{geom}}(\gamma+\gamma')
=b_R^{\mathrm{geom}}(\gamma)+b_R^{\mathrm{geom}}(\gamma')\) on retained
extensions remaining in the quotient.

\item \emph{Projection-finite local sector.}  For
\(\alpha\in\Gamma_R^{\mathrm{loc}}\) and finite index set \(I\), build
the closed-support incidence stack over \(E^I\):
\[
\mathfrak M^{\mathrm{loc,cl}}_{\alpha,R,I}
=
\Bigl\{(A,\mathbf e)\in\mathfrak M^{\mathrm{sc}}_{R,L\alpha}\times E^I
\Bigm|p_E(\mathrm{Supp}\,A)\subset\bigcup_{i\in I}S\times\{e_i\}\Bigr\}.
\]
Form
\(\mathscr H^{\mathrm{loc}}_{\alpha,R,I}=
(\pi_{\alpha,R,I})_!(\Phi_{\alpha,R,I}^{\mathrm{loc}}\otimes
o_{\alpha,R,I}^{\mathrm{loc}})\) and the open-support sheaf
\(\mathcal F_{\alpha,R}^{\mathrm{loc}}(U)\) over \(U\subset E\) by
restricted Borel--Moore configuration sums.

\item \emph{Wrapped prequotient sector.}  For
\(\eta\in\Gamma_R^{\mathrm{wr}}\), form the rigidified prequotient
\(\mathcal M_{\eta,R}^{\mathrm{wr,rig}}\to\mathcal M_{\eta,R}^{\mathrm{wr}}\)
by the Abramovich--Polishchuk / determinant anchor
\[
\lambda_\eta:\mathcal M_{\eta,R}^{\mathrm{wr,rig}}\longrightarrow E,
\qquad
\lambda_\eta(A)=\det Rp_{E,*}A
\otimes\mathcal O_E(-\chi(A)\cdot 0_E),
\]
divided by a finite cover when \(\chi(A)\ne 1\) so that
\(\lambda_\eta(t\cdot A)=\lambda_\eta(A)+t\) holds.  Verify the
losslessness residual \(\mathfrak o^{\lambda,\mathrm{loss}}_R=0\) by
the Ext-distinguishing test of
Definition~\ref{def:hybrid-wrapped-factorization-datum}, separating
\(W_0=i_{s,*}(\mathcal O_E\oplus\mathcal O_E)\) from
\(W_L=i_{s,*}(L\oplus L^{-1})\).

\item \emph{Mixed local/wrapped correspondences.}  For an ordered
hybrid input
\(\alpha_-\in\Gamma_R^{\mathrm{loc}}\),
\(\eta\in\Gamma_R^{\mathrm{wr}}\),
\(\beta_+\in\Gamma_R^{\mathrm{loc}}\), form
\(\mathfrak E^{\mathrm{mix}}_{\alpha_-;\eta;\beta_+,R}\) before the
reduced \(E\)-quotient.

\item \emph{Eight-word two-step flag atlas.}  Build flag stacks
\(\mathfrak F^{(2),w}_R\) for each
\(w\in\{\mathrm L,\mathrm W\}^3=
\{\mathrm{LLL},\mathrm{LLW},\mathrm{LWL},\mathrm{WLL},\mathrm{LWW},
\mathrm{WLW},\mathrm{WWL},\mathrm{WWW}\}\), classifying filtrations
with three associated graded factors of those types.  Verify the
four-input pentagon coherence and the quotient-after-correspondence
descent.

\item \emph{Hybrid correspondence-module functor.}  Assemble the
correspondence category
\(\mathsf{Corr}^{E,\mathrm{hyb}}_R\) and the functor
\(\mathcal B_R^{E,\mathrm{hyb}}:\mathsf{Corr}^{E,\mathrm{hyb}}_R
\to\mathsf{Ch}_{\C}\) by
\(\mathcal B_R^{E,\mathrm{hyb}}(e)=q_!\,\mathrm{TS}^{\mathrm{red}}_e\,p^*\)
for every generating arrow
\(e=(Y\xleftarrow{p}\mathfrak E_e\xrightarrow{q}Z)\).

\item \emph{Quotient functor \(Q_{E,R}\).}  Apply
\(E\)-equivariant compactly supported vanishing-cycle chains and the
orientation descent of step~4 to produce the reduced object
\[
\mathcal F^{\mathrm{geom}}_{X,\sigma,S,\le R}
=Q_{E,R}\circ\mathcal B_R^{E,\mathrm{hyb}}.
\]

\item Output the hybrid datum with all coherences and the protected
integration map \(I_R^{\mathrm{prot}}\), with homogeneous
\(s\)-degree the Borcherds trace exponent
\(m_R=\operatorname{pr}_3\overline\Pi_X\).
\end{enumerate}

\paragraph{Type checks.}
The carrier
\(\mathcal F_{X,\sigma,S,\le R}^{\mathrm{hyb}}\) is a hybrid
correspondence-module functor on
\(\mathsf{Corr}^{E,\mathrm{hyb}}_R\) with values in
\(\mathsf{Ch}_{\C}\); the four-input pentagon is supplied by
the eight-word two-step flag atlas; the anchor residual tuple
\(\mathfrak o^\lambda_R=
(\mathfrak o^{\lambda,\mathrm{ex}}_R,
\mathfrak o^{\lambda,\mathrm{unit}}_R,
\mathfrak o^{\lambda,\mathrm{loss}}_R,
\mathfrak o^{\lambda,\mathrm{multi}}_R,
\mathfrak o^{\lambda,\mathrm{tr}}_{R'R})\) vanishes; the hybrid
residual \(\mathfrak O_{\mathrm{hyb},R}=0\) at every height \(R\),
compatibly with transition maps.

\paragraph{Obstructions discharged or invoked.}
The middle type-II wall lies in
\(\delta_2\leftrightarrow(0,1,1)\), elliptic degree
\(d=m+1=2\) by Lemma~\ref{app:pfaffian-wall-three-walls}; it
cannot be represented by a projection-finite local stack and demands
the wrapped prequotient.  This step's correctness — the existence of
the wrapped representative
\(x_{\delta_2,R}\in\mathcal M_{\eta,R}^{\mathrm{wr,rig}}\) with
\(\overline\Pi_X(x_{\delta_2,R})=(0,1,1)\) and its compatible anchor
\(\lambda_{\eta,R}\) — is precisely the hybrid input to \((O2)\) at
the middle wall.

\subsection{Construction~4 — Reduced d-critical orientation
and Joyce--Upmeier transport}
\label{app:build-orientation}

\paragraph{Tag.} \(\beta\)-functorial passage (oriented carrier).

\paragraph{Input.}
The hybrid carrier
\(\mathcal F_{X,\sigma,S,\le R}^{\mathrm{hyb}}\) of
Construction~\ref{app:build-hybrid-wrapped}; the K3 holomorphic
\(2\)-form \(\sigma_S\) on \(S\); BBDJS d-critical chart machinery
\cite{BBDJS2015}; the Joyce--Upmeier multiplicative orientation
formalism \cite{JoyceUpmeier,JoyceTanakaUpmeier}; the cosection
localization machinery of Kiem--Li and the cosection-reduced
obstruction theory; the type-II Weyl group
\(W^{(2)}(\La^{2,1}_{II})\) generated by
\(s_{\delta_1},s_{\delta_2},s_{\delta_3}\).

\paragraph{Output.}
The cosection
\(\sigma=p_S^*\sigma_S\); the reduced tangent complex
\[
\operatorname{RHom}_{\mathrm{red}}(A,A)
=
\operatorname{Cone}\!\Bigl(
\operatorname{RHom}_X(A,A)
\xrightarrow{(\mathrm{tr},\mathrm{cs})}
R\Gamma(X,\mathcal O_X)\oplus H^2(S,\mathcal O_S)[-1]
\Bigr)[-1];
\]
the BBDJS vanishing-cycle complex \(\Phi_{R,c}\); the reduced
orientation line
\[
o_{R,c}
=\det\operatorname{Ext}^1_{\mathrm{red}}(A,A)^{-1}
\quad\text{on Darboux charts},\qquad
o_{R,c}^{\otimes 2}\simeq\det\operatorname{RHom}_{\mathrm{red}}(A,A);
\]
the Joyce--Upmeier multiplicative transport
\(o_{c\oplus c'}\simeq o_c\otimes o_{c'}\) compatible with the reduced
extension correspondences; the quotient orientation
\(\mathfrak Q^{or}_R\) of
Theorem~\ref{prop:generic-local-sign-orientation-character}; the
Weyl-equivariant lifts
\(\tau_{i,R,S}:o_{R,S}\to s_{\delta_i}^*o_{R,S}\) for \(i=1,2,3\) of
\((O1)^+\); the local generic wall-sign data \(\epsilon_{o,R}^{\mathrm{loc}}\).

\paragraph{Procedure.}
\begin{enumerate}
\item \emph{Cosection.}  Pull back the K3 semiregularity cosection
\(\mathrm{cs}\) along \(p_S:X\to S\) to obtain the 2-step cosection
\(\sigma=p_S^*\sigma_S\) on the K3 factor; verify additivity in exact
triangles by \textup{(RedOr)} of
Theorem~\ref{thm:finite-stage-dirac-igusa-pro-object}.

\item \emph{Reduced obstruction.}  Form the reduced tangent complex by
the displayed cone above (Kiem--Li cosection localization).  PTVV
\cite{PTVV2013} \((-1)\)-shifted symplectic on
\(\operatorname{RHom}_X(A,A)\) descends, after the cosection
contraction, to the reduced complex.

\item \emph{BBDJS d-critical chart.}  On each Darboux chart of
\(\mathfrak M^{\mathrm{sc}}_{R,c}\), build the BBDJS vanishing-cycle
complex \(\Phi_{R,c}\) \cite{BBDJS2015}; verify reduced
Thom--Sebastiani transport coherent on two-step flag stacks of
Construction~\ref{app:build-hybrid-wrapped}.

\item \emph{Reduced orientation line.}  Define
\(o_{R,c}=\det\operatorname{Ext}^1_{\mathrm{red}}(A,A)^{-1}\); verify
the squared isomorphism
\(o_{R,c}^{\otimes 2}\simeq\det\operatorname{RHom}_{\mathrm{red}}(A,A)\)
by \(3\)-Calabi--Yau Serre duality
\(\operatorname{Ext}^2_{\mathrm{red}}(A,A)
\simeq\operatorname{Ext}^1_{\mathrm{red}}(A,A)^\vee\).

\item \emph{Quotient-orientation descent on the K3\(\times\)E-quotient
self-Ext frame bundle.}  Compute the four Cech--Borel obstruction
classes for every retained stratum \(S\):
\[
\alpha^{\mathrm{red}}_{R,S}=w_2(\mathscr D_{R,S}),
\quad
\alpha^{E,\mathrm{free}}_{R,S}=
\mathrm{edge}_{2,0}(w_2^E(\mathscr D_{R,S})),
\quad
\widetilde\beta^H_{R,S}=w_2^G(\mathscr D_{R,S}),
\quad
\lambda^H_{R,S}\in H^1_G(S;\mathbb F_2),
\]
and verify simultaneous vanishing with compatible
null-trivializations on overlaps, on Hall correspondence pullbacks,
on two-step flag stacks, and under successor maps.  The full
calculation proceeds primary stabilizer by stabilizer; for cyclic
\(\Z/2^a\) the absent mixed coefficient requires the rank-two
two-primary check.  Output the descended quotient orientation system
\(\mathfrak Q^{or}_R\).

\item \emph{Multiplicative Joyce--Upmeier transport.}  Transport
\(o_{c\oplus c'}\simeq o_c\otimes o_{c'}\) coherently across two-step
flag stacks
\cite{JoyceUpmeier,JoyceTanakaUpmeier}; verify the reduced
Thom--Sebastiani identification on three-input flags.

\item \emph{Weyl-equivariant lifts \((O1)^+\).}  For each
\(i\in\{1,2,3\}\) construct the lift
\(\tau_{i,R,S}:o_{R,S}\to s_{\delta_i}^*o_{R,S}\); verify the
projective Weyl cocycle
\(c_{o,R}\in Z^2(\mathcal G_R;\underline{\mathbb F}_2)\) on the
retained action groupoid; null-trivialize \([c_{o,R}]\) and choose a
\(1\)-cochain killing it; verify \(\tau_{i,R,S}^2=1\); verify the
Coxeter relations
\((s_{\delta_i}s_{\delta_j})^3=1\) for \(i\ne j\) (the
\((3,3,3)\)-hyperbolic triangle relations of
Appendix~\ref{app:pfaffian-wall-three-walls}); verify that the torsor
defects \(\omega_{i,\mathcal C}\) of
Definition~\ref{thm:finite-stage-dirac-igusa-pro-object}~$(O_X)$ vanish.

\item \emph{Local wall-sign datum.}  At each generic type-II wall point
\(\Hpl_{\delta_i}\), with the rank-one local model
\(K_{\delta_i}^{\mathrm{cross}}=[\R\xrightarrow{u_i}\R]\), record the
local generic sign \(\epsilon_{o,R}^{\mathrm{loc}}(s_{\delta_i})=-1\)
of Proposition~\ref{prop:rank-one-crossing-o2-sign}.

\item Output
\((o_R,\mathfrak Q^{or}_R,\{\tau_{i,R,S}\},\epsilon_{o,R}^{\mathrm{loc}})\).
\end{enumerate}

\paragraph{Type checks.}
\(o_R\) is a Picard-line section squaring to the reduced determinant;
the four Cech--Borel classes vanish with compatible
null-trivializations; the lifts \(\tau_{i,R}\) satisfy
\(\tau_{i,R}^2=1\) and the \((3,3,3)\)-Coxeter relations; multiplicative
transport agrees on flag stacks.

\paragraph{Obstructions discharged or invoked.}
This step is the entire site of \((O1)\) (strong reduced orientation
on the K3\(\times\)E-quotient self-Ext frame bundle) and of
\((O1)^+\) (Weyl-equivariant transport along
\(W^{(2)}(\La^{2,1}_{II})\) reflections).  Discharging \((O1)\)
requires the simultaneous vanishing and compatible
null-trivialization of the four Cech--Borel classes on every retained
stratum; discharging \((O1)^+\) requires \([c_{o,R}]=0\), the
Coxeter coherence, and the vanishing of all torsor defects
\(\omega_{i,\mathcal C}\).  The local wall-sign datum is the input to
\((O2)\) at step~6.

\subsection{Construction~5 — Hall correspondences and primitive part}
\label{app:build-hall}

\paragraph{Tag.} \(\gamma\)-shadow specialization
(Stage-2 chiral Hall shadow).

\paragraph{Input.}
The oriented hybrid carrier from
Construction~\ref{app:build-orientation}; the extension stack
\(\mathfrak E^{\mathrm{geom}}_{R,\widehat c,\widehat c'}\); the
two-step flag stack
\(\mathfrak F^{(2),\mathrm{geom}}_{R,\widehat c,\widehat c',\widehat c''}\);
the BBDJS reduced Thom--Sebastiani transport; six-functor admissibility
of \(p^*,q_!,p_!\) (Construction~\ref{app:build-finite-moduli}~step
\textup{(Six)}).

\paragraph{Output.}
The finite Hall object
\[
\mathcal H_R^{\mathrm{geom}}
=\bigoplus_{\widehat c\in\widehat\Gamma_R}
H_c^\bullet\!\bigl(
\mathfrak M^{\mathrm{geom}}_{R,\widehat c},
\Phi_{R,c}\otimes o_{R,c}\bigr);
\]
the Hall product \(m_R=q_!\,\mathrm{TS}^{\mathrm{red}}\,p^*\) and Hall
coproduct \(\Delta_R=p_!\,(\mathrm{TS}^{\mathrm{red}})^{-1}\,q^*\); the
finite Hall bialgebra
\((\mathcal H_R^{\mathrm{geom}},m_R,\Delta_R,\eta_R,\epsilon_R)\); the
primitive part
\[
P_R^{\mathrm{geom}}
=\ker(\Delta_R-1\otimes\mathrm{id}-\mathrm{id}\otimes 1)\cap\ker\epsilon_R;
\]
the normal-ordered Gram descent
\(P_R^{\Pi,\mathrm{geom}}
=\overline\Pi_{X,*}^{\Theta_R}P_R^{\mathrm{geom}}\); the
\(\Gamma_R^\Pi\)-graded super-Lie bracket
\([\,\bar x,\bar y\,]_{\mathrm{red}}\) of
Proposition~\ref{thm:hall-primitives-formal};
the Hopf pairing matrix \(G_{\gamma,ab}\) and radical \(K_\gamma\).

\paragraph{Procedure.}
\begin{enumerate}
\item \emph{Pullback.}  For
\(\widehat c,\widehat c'\in\widehat\Gamma_R\), pull back along
\(p:\mathfrak E^{\mathrm{geom}}_{R,\widehat c,\widehat c'}
\to\mathfrak M^{\mathrm{geom}}_{R,\widehat c}
\times\mathfrak M^{\mathrm{geom}}_{R,\widehat c'}\); apply reduced
Thom--Sebastiani \(\mathrm{TS}^{\mathrm{red}}\) to combine vanishing
cycle data \cite{BBDJS2015}; pushforward exceptionally along
\(q:\mathfrak E^{\mathrm{geom}}_{R,\widehat c,\widehat c'}
\to\mathfrak M^{\mathrm{geom}}_{R,\widehat c\star\widehat c'}\) to
obtain
\(m_R(x\otimes y)=q_!\,\mathrm{TS}^{\mathrm{red}}\,p^*(x\boxtimes y)\).

\item \emph{Coproduct.}  Dually pull back along \(q\), apply
inverse reduced Thom--Sebastiani, pushforward along \(p\) to obtain
\(\Delta_R(z)=p_!\,(\mathrm{TS}^{\mathrm{red}})^{-1}\,q^*(z)\).
The compact geometric coproduct correspondence
\(\mathfrak C^{\mathrm{geom}}_{\rho;\mu,\nu}\) of
Proposition~\ref{thm:hall-primitives-formal} is the
support data for splitting.

\item \emph{Bialgebra check.}  Verify associativity of \(m_R\) and
coassociativity of \(\Delta_R\) on the two-step flag stack
\(\mathfrak F^{(2),\mathrm{geom}}_{R,\widehat c,\widehat c',\widehat c''}\)
by base change and Thom--Sebastiani coherence.  Verify the bialgebra
identity (Hopf compatibility) via the Massey--Saito Joyce--Upmeier
multiplicative orientation transport.

\item \emph{Primitive extraction.}  Compute
\(P_R^{\mathrm{geom}}=\ker(\Delta_R-1\otimes\mathrm{id}
-\mathrm{id}\otimes 1)\cap\ker\epsilon_R\); since
\(\Delta_R-1\otimes\mathrm{id}-\mathrm{id}\otimes 1\) is degree-graded
on \(\widehat\Gamma_R\), this reduces to a finite linear computation
in each charge.

\item \emph{Normal-ordered descent.}  Apply
\(\overline\Pi_{X,*}^{\Theta_R}\) of
Theorem~\ref{thm:hall-product} to obtain the
\(\Gamma_R^\Pi\)-graded
\(P_R^{\Pi,\mathrm{geom}}\); verify additivity
\(\overline\Pi_X(\widehat c\star\widehat c')
=\overline\Pi_X(\widehat c)+\overline\Pi_X(\widehat c')\).

\item \emph{Bracket and pairing matrices.}  Compute the matrices
\[
M^{\alpha+\beta,c}_{\alpha,a;\beta,b}
=\int_{\mathfrak E^{\mathrm{geom}}_{\alpha,\beta;\alpha+\beta}}
q^!(e_{\alpha+\beta,c}^\vee)\cap
\mathrm{TS}^{\mathrm{red}}(p^*(e_{\alpha,a}\boxtimes e_{\beta,b})),
\quad
B_{\alpha,\beta}(x\otimes y)
=M_{\alpha,\beta}(x\otimes y)
-(-1)^{|x||y|}M_{\beta,\alpha}\mathsf{sw}(x\otimes y),
\]
\[
G_{\gamma,ab}=
\mathrm{tr}_0(q_!\,\mathrm{TS}^{\mathrm{red}}\,p^*
(e_{\gamma,a}\boxtimes e_{-\gamma,b}))
\quad\hbox{on}\quad
\mathfrak P^{\mathrm{geom}}_{\gamma,-\gamma;0}.
\]
Form the radical
\(K_{\gamma,\bar p}
=\ker G_{\gamma,\bar p}\cap\ker G_{-\gamma,\bar p}^t\) and the
quotient \(L_{\gamma,\bar p}\); fix splittings.

\item Output
\((\mathcal H_R^{\mathrm{geom}},m_R,\Delta_R,P_R^{\mathrm{geom}},
P_R^{\Pi,\mathrm{geom}},
\{B_{\alpha,\beta}\},\{G_{\gamma,ab}\},\{K_\gamma\},\{L_\gamma\})\).
\end{enumerate}

\paragraph{Type checks.}
\(\mathcal H_R^{\mathrm{geom}}\) is a finite \(\widehat\Gamma_R\)-graded
Hall bialgebra; \(m_R\) is associative on the two-step flag stack;
\(\Delta_R\) is coassociative; the supercommutator of primitives is
primitive (the bialgebra identity); the descended bracket is
\(\Gamma_R^\Pi\)-graded after \(\overline\Pi_{X,*}^{\Theta_R}\);
finite-dimensional in each homogeneous parity piece by
finite-typeness.

\paragraph{Obstructions discharged or invoked.}
\((D_0)\) is partially discharged: the finite Hall package
\((\mathcal H_R^{\mathrm{geom}},m_R,\Delta_R,P_R)\) is constructed.
The seven obstruction classes
\(\mathfrak o_\Theta^{\mathrm{top}},
\mathfrak o_\Theta^{\mathrm{Hoch}},
\mathfrak o_\Theta^{\mathrm{tri}},
\mathfrak o_\Theta^{\mathrm{Jac}},
\mathfrak o_F,
\mathfrak o_\Theta^{\mathrm{cyc}},
\mathfrak o_\Theta^{\mathrm{rad}}\) of
Theorem~\ref{thm:hall-product} are required to vanish on
this finite Hall stage and compatibly in the inverse limit.

\paragraph{Cross-volume anchor.}  The output
\((\mathcal H_R^{\mathrm{geom}},m_R)\) is the finite-stage realization
of the K3-side BKM Hall--Drinfeld double
\(\mathcal D_\hbar(\mathrm{CoHA}_{K3\times E})\) of
\texttt{calabi-yau-quantum-groups:CLAUDE.md~§6.7}.  In the universal
stage chain \(\mathsf P\to\mathsf C\to\mathsf S\) the Hall positive half
\(Y^+(X)\) is precisely \(P_R^{\mathrm{geom}}\) at finite \(R\); the
Drinfeld double \(G(X)=D(Y^+(X))\) requires Hopf pairing, completion,
integral form, and stable-envelope transport which are supplied by
\(G_{\gamma,ab}\) and \((\mathrm{ML})\).  The κ-tuple coordinate
\(\kappa_{\mathrm{BKM}}=5\) is the weight of \(\Delta_5\) and the
output of the level-4 scalar trace; it is not a level-2 invariant of
the carrier itself.

\subsection{Construction~6 — Dirac operator and Pfaffian line}
\label{app:build-dirac-pfaffian}

\paragraph{Tag.} \(\gamma\)-shadow specialization
(orientation gerbe + Pfaffian).

\paragraph{Input.}
The oriented Hall package
\((\mathcal H_R^{\mathrm{geom}},o_R,\{\tau_{i,R,S}\})\) of
Constructions~\ref{app:build-orientation}--\ref{app:build-hall}; the
local wall-sign datum
\(\epsilon_{o,R}^{\mathrm{loc}}\) of
Construction~\ref{app:build-orientation}; the rank-one type-II
Pfaffian crossing datum of
Definition~\ref{def:rank-one-type-ii-crossing} and
Appendix~\ref{app:pfaffian-wall-rank-one}; the half-Hilbert orbit
under \(W^{(2)}(\La^{2,1}_{II})\) of size \(2^{12}=4096\).

\paragraph{Output.}
The Dirac operator on the orientation gerbe
\(D_R\); the Pfaffian line
\(\mathscr L_{\mathrm{Pf},R}=
\det\operatorname{Ext}^1_{\mathrm{red}}(A,A)^{-1}\) (squaring to the
reduced determinant); the protected Pfaffian section
\(\operatorname{pf}_{X,R}\); the global wall-sign character
\(\epsilon_{o,R}:W^{(2)}(\La^{2,1}_{II})\to\{\pm1\}\); the Pfaffian
section asymptotic
\[
\operatorname{pf}_{X,R}
=
64q^{1/2}r^{1/2}s^{1/2}
\prod_{\gamma=(n,l,m)\in A_R}
(1-q^nr^ls^m)^{f(nm,l)}+\hbox{higher-window correction},
\]
which exhausts to \(\Delta_5\) as \(R\to\infty\).

\paragraph{Procedure.}
\begin{enumerate}
\item \emph{Dirac operator.}  On each Darboux chart of
\(\mathfrak M^{\mathrm{sc}}_{R,c}\), with reduced normal self-Ext
splitting \(K^{\mathrm{nor}}_{R,c}\), build the odd Dirac operator
\(D_{R,c}\) acting on \(\Omega^{\mathrm{sp}}\otimes o_{R,c}\)
(spinors on the orientation gerbe); verify that, at a generic
type-II wall point, the local rank-one normal block is
\(D_{\delta_i}=\begin{psmallmatrix}0&u\\-u&0\end{psmallmatrix}\) with
\(\operatorname{Pf}(D_{\delta_i})=u\)
(Definition~\ref{def:rank-one-type-ii-crossing}).

\item \emph{Pfaffian line.}  Set
\(\mathscr L_{\mathrm{Pf},R}=
\det\operatorname{Ext}^1_{\mathrm{red}}(A,A)^{-1}\); verify
\(\mathscr L_{\mathrm{Pf},R}^{\otimes 2}\simeq
\det\operatorname{RHom}_{\mathrm{red}}(A,A)\) by
\(3\)-Calabi--Yau Serre duality.

\item \emph{Local wall sign.}  At each generic point of
\(\Hpl_{\delta_i}\) with normal model
\(K^{\mathrm{nor}}_{\delta_i,R}=[\R\xrightarrow{u_i}\R]\), compute
the monodromy of the Pfaffian section
\(\operatorname{pf}_{\delta_i,R}=\upsilon_iu_i\) along the
wall-reflection loop \(\ell_{\delta_i,R,S}\) closed by the
\((O1)^+\) lift \(\tau_{i,R,S}\):
\[
M_{\ell_{\delta_i,R,S}}(\upsilon_iu_i)
=\upsilon_i(-u_i)=-\upsilon_iu_i,
\qquad
\epsilon_{o,R}^{\mathrm{loc}}(s_{\delta_i})
=(-1)^{N^{\mathrm{Pf}}_{\delta_i,R}}=-1.
\]

\item \emph{Half-Hilbert orbit transport.}  The reflection
\(s_{\delta_i}\) acts on the half-Hilbert orbit (compact source labels
of size \(2^{12}=4096\) parametrized by 12 binary parities tied to the
\(2^{12}\) odd-spin structures on the lattice
\(\La^{2,1}_{II}/2\La^{2,1}_{II}\)).  Sum the local rank-one
contributions over the orbit:
\[
\epsilon_{o,R}(s_{\delta_i})
=\prod_{x\in\hbox{orbit}/2^{12}}
\epsilon_{o,R}^{\mathrm{loc}}(s_{\delta_i};x).
\]
The output is the global character
\(\epsilon_{o,R}:W^{(2)}(\La^{2,1}_{II})\to\{\pm 1\}\).

\item \emph{Compatibility package.}  Verify the component, boundary,
overlap, quotient, and protected integration compatibility package of
Proposition~\ref{prop:appE-local-pfaffian-sign} on every
retained type-II wall stratum.

\item \emph{Pfaffian section.}  Define
\(\operatorname{pf}_{X,R}\) as the protected integration of the
oriented Hall product over
\(\mathcal F_{X,\sigma,S,\le R}^{\mathrm{hyb}}\):
\[
\operatorname{pf}_{X,R}
=
64q^{1/2}r^{1/2}s^{1/2}
I_R^{\mathrm{prot}}\!\left(
\bigotimes_{\gamma\in A_R}\operatorname{pf}_{\gamma,R}\right),
\qquad
\operatorname{pf}_{\gamma,R}=(1-q^nr^ls^m)^{f(nm,l)}.
\]

\item Output
\((D_R,\mathscr L_{\mathrm{Pf},R},\operatorname{pf}_{X,R},
\epsilon_{o,R})\), with the protected integration compatibility data.
\end{enumerate}

\paragraph{Type checks.}
\(\mathscr L_{\mathrm{Pf},R}\) is a Picard-line whose square
isomorphism is supplied by reduced Serre duality;
\(\operatorname{pf}_{X,R}\) is a section of
\(\mathscr L_{\mathrm{Pf},R}\) with the displayed cusp expansion; the
character \(\epsilon_{o,R}\) takes values in \(\{\pm 1\}\) and
satisfies the Coxeter relations of
\(W^{(2)}(\La^{2,1}_{II})\); on each simple type-II reflection
\(\epsilon_{o,R}(s_{\delta_i})=-1\) by the rank-one local model.

\paragraph{Obstructions discharged or invoked.}
This step is the entire site of \((O2)\) (local Pfaffian wall normal
form on type-II walls; the \(2^{12}=4096\) sign sum on the
half-Hilbert orbit transport).  The local rank-one Pfaffian crossing
of Appendix~\ref{app:pfaffian-wall-rank-one} discharges \((O2)\)
generically; the global half-Hilbert orbit transport remains the open
\((O2)\) component of
Open~Problem~\ref{def:O1-plus-weyl-transport}.  The character
identity
\(\epsilon_o=\nu_{\Delta_5}\big|_{W^{(2)}(\La^{2,1}_{II})}\) becomes a
theorem only after \((D_0)\), \((O1)\), \((O1)^+\), and \((O2)\) are
discharged compatibly in \(R\).

\subsection{Construction~7 — Chiral Koszul source coalgebra}
\label{app:build-koszul}

\paragraph{Tag.} \(\gamma\)-shadow specialization
(\emph{Bar = twisting; not bulk}; the bulk is
\(Z^{\mathrm{der}}_{\mathrm{ch}}(\mathcal F_X^{\mathrm{hyb}})\),
constructed elsewhere).

\paragraph{Input.}
The oriented Hall primitives
\(P_R^{\Pi,\mathrm{geom}}\) of
Construction~\ref{app:build-hall}; the hybrid carrier
\(\mathcal F_{X,\sigma,S,\le R}^{\mathrm{hyb}}\) of
Construction~\ref{app:build-hybrid-wrapped}; the augmentation
\(\varepsilon_{X,R}^{\mathrm{hyb}}:
\mathcal F_{X,\sigma,S,\le R}^{\mathrm{hyb}}\to k\mathbf 1\); the
hybrid residual \(\mathfrak O_{\mathrm{hyb},R}=0\); the bounded
length \(N_R\) (positivity of HN height on every non-vacuum colour);
the Beilinson--Drinfeld/Francis--Gaitsgory bar-cobar formalism
\cite{BD,FrancisGaitsgoryCKD}; the conilpotent completed domain on
the elliptic curve \(E\).

\paragraph{Output.}
The conilpotent chiral coalgebra on \(E\):
\[
C_{X,R}^{\mathrm{geom}}
=\bar B_E^{\mathrm{ch}}
(\mathcal F_{X,\sigma,S,\le R}^{\mathrm{hyb}},
\varepsilon_{X,R}^{\mathrm{hyb}})
=k\mathbf 1\oplus\bigoplus_{1\le p\le N_R}
\bar B_E^{\mathrm{ch},p}(\overline{\mathcal F}_{X,R}^{\mathrm{geom}});
\]
the bar-length filtration
\(F^{\mathrm{co}}_{R,p}\); the chiral collision coproduct
\(\Delta_R^{\mathrm{ch}}\); the bar comparison
\(b_{X,R}=\mathrm{id}_{\bar B_E^{\mathrm{ch}}(\cdot)}\); the source
cobar counit
\(\varepsilon^{\mathrm{src}}_{\mathrm{bar},R}:
\Omega_E^{\mathrm{ch}}C_{X,R}^{\mathrm{geom}}
\xrightarrow{\simeq}
\mathcal F_{X,\sigma,S,\le R}^{\mathrm{hyb}}\); after a separate
source-to-target quasi-isomorphism
\(\vartheta_R:\mathcal F_{X,\sigma,S,\le R}^{\mathrm{hyb}}
\xrightarrow{\simeq}\mathsf{Den}_{\Delta_5,E,\le R}\), the cobar
comparison
\(\Theta_{\mathrm{Kos},R}=\vartheta_R\circ\varepsilon^{\mathrm{src}}_{\mathrm{bar},R}\).

\paragraph{Procedure.}
\begin{enumerate}
\item \emph{Augmentation kernel.}  Set
\(\overline{\mathcal F}_{X,R}^{\mathrm{geom}}
=\ker(\varepsilon_{X,R}^{\mathrm{hyb}})\); verify the bounded-length
hypothesis (positivity of HN height on every non-vacuum surviving
colour, hence at most \(N_R\) factors in any non-zero word).

\item \emph{Bar coalgebra.}  Form the reduced chiral bar
\[
C_{X,R}^{\mathrm{geom}}
=k\mathbf 1\oplus\bigoplus_{1\le p\le N_R}
\bar B_E^{\mathrm{ch},p}
(\overline{\mathcal F}_{X,R}^{\mathrm{geom}}),
\]
length \(p\) summand generated by ordered \(p\)-fold inputs from the
augmentation ideal of total HN height \(\le R\) (height \(>R\) outputs
zero); coaugmentation by the inclusion of the vacuum summand; counit by
projection.

\item \emph{Bar-length filtration.}
\[
F^{\mathrm{co}}_{R,p}C_{X,R}^{\mathrm{geom}}
=k\mathbf 1\oplus\bigoplus_{1\le j\le p}
\bar B_E^{\mathrm{ch},j}
(\overline{\mathcal F}_{X,R}^{\mathrm{geom}}),\qquad
0\le p\le N_R.
\]

\item \emph{Collision coproduct.}  For
\(x=[x_1|\cdots|x_p]\), define the cut-map coproduct
\[
\Delta_R^{\mathrm{ch}}(x)
=\mathbf1\otimes x+x\otimes\mathbf1
+\sum_{i=1}^{p-1}[x_1|\cdots|x_i]\otimes[x_{i+1}|\cdots|x_p],
\]
with each cut realized by the corresponding finite collision/flag
pull-push map in
\(\mathsf{Corr}^{E,\mathrm{hyb}}_{R,\mathrm{geom}}\) of
Construction~\ref{app:build-hybrid-wrapped}.  Verify
coassociativity by the four-input pentagon and the eight-word flag
coherences; verify counit identities via vacuum collision flags.

\item \emph{Conilpotence.}  Verify
\((\overline\Delta_R^{\mathrm{bar}})^{N_R+1}=0\): each reduced cut
strictly lowers bar length on each branch.

\item \emph{Bar comparison.}  Take \(b_{X,R}\) to be the identity of
\(\bar B_E^{\mathrm{ch}}
(\mathcal F_{X,\sigma,S,\le R}^{\mathrm{hyb}},
\varepsilon_{X,R}^{\mathrm{hyb}})\), eliminating the coalgebra-side
finite Koszul residuals
\(\mathfrak o^C_R, \mathfrak o^{\mathrm{fil}}_R,
\mathfrak o^{\mathrm{conil}}_R, \mathfrak o^{\Delta,\mathrm{ch}}_R\)
of Definition~\ref{def:chiral-koszul-source-datum}.

\item \emph{Source cobar counit.}  Apply Francis--Gaitsgory chiral
Koszul duality
\cite[Theorem~6.2.2]{FrancisGaitsgoryCKD} in the conilpotent completed
domain to obtain
\(\varepsilon^{\mathrm{src}}_{\mathrm{bar},R}:
\Omega_E^{\mathrm{ch}}C_{X,R}^{\mathrm{geom}}
\xrightarrow{\simeq}
\mathcal F_{X,\sigma,S,\le R}^{\mathrm{hyb}}\); verify it is a
source-internal quasi-isomorphism of chiral algebras on \(E\), not a
homotopy inverse of the target counit.

\item \emph{Source-to-target identification.}  Construct
\(\vartheta_R\) from the geometric primitive representatives, Hall
products, coproducts, pairings, and HN transition compatibility;
verify the residual
\(\mathfrak o^\vartheta_R\) of
Corollary~\ref{cor:source-bar-remaining-cobar-obstruction} vanishes;
set \(\Theta_{\mathrm{Kos},R}=\vartheta_R\circ
\varepsilon^{\mathrm{src}}_{\mathrm{bar},R}\).

\item \emph{Inverse limit.}  Verify the Koszul Mittag--Leffler
condition \(R^1\varprojlim=0\) for the inverse systems of source
coalgebras, completed chiral bar/cobar constructions, target
truncations, and the cones of \(b_{X,R}\) and
\(\Theta_{\mathrm{Kos},R}\) \cite[\S3.5]{WeibelHA}.  Form
\[
C_X^{\mathrm{geom}}
=\varprojlim_R C_{X,R}^{\mathrm{geom}},\qquad
\Theta_{\mathrm{Kos}}:
\Omega_E^{\mathrm{ch}}C_X^{\mathrm{geom}}
\xrightarrow{\simeq}\mathsf{Den}_{\Delta_5,E}.
\]

\item Output
\((C_{X,R}^{\mathrm{geom}},F^{\mathrm{co}}_{R,\bullet},
\Delta_R^{\mathrm{ch}},b_{X,R},\Theta_{\mathrm{Kos},R})\) at every
\(R\), with successor compatibility.
\end{enumerate}

\paragraph{Type checks.}
\(C_{X,R}^{\mathrm{geom}}\) is a coaugmented conilpotent chiral
coalgebra on \(E\) with bounded bar-length filtration of length at
most \(N_R\); \(\Delta_R^{\mathrm{ch}}\) is coassociative and counital
on the chosen flag stacks; the bar-cobar inversion is supplied by
Francis--Gaitsgory chiral Koszul; the cobar comparison
\(\Theta_{\mathrm{Kos},R}\) is source-internal, not the target
bar-cobar counit (the source/target separation of
Lemma~\ref{lem:source-target-koszul-separation}).

\paragraph{Obstructions discharged or invoked.}
The Koszul Mittag--Leffler exactness hypothesis is the
\textup{(D0)}-component for the chiral Koszul lane.  The total finite
Koszul residual
\(\mathfrak O_{\mathrm{Kos},R}=
(\mathfrak o^C_R,\mathfrak o^{\mathrm{fil}}_R,
\mathfrak o^{\mathrm{conil}}_R,\mathfrak o^{\Delta,\mathrm{ch}}_R,
\{\mathfrak o^{C,\mathrm{tr}}_{R'R}\},\mathfrak o^b_R,
\mathfrak o^\Omega_R,\mathfrak o^{\mathrm{hyb}}_R,
\mathfrak o^\Pi_R,\mathfrak o^{\mathrm{sep}}_R,
\mathfrak o^{\Prim}_R)\)
vanishes after this step.

\paragraph{Cross-volume firewall.}  The bar coalgebra
\(\bar B_E^{\mathrm{ch}}
(\mathcal F_{X,R}^{\mathrm{hyb}})\) is a \emph{twisting / coupling}
coalgebra (Maurer--Cartan classification), \emph{not} the bulk
\(Z^{\mathrm{der}}_{\mathrm{ch}}(\mathcal F_X^{\mathrm{hyb}})
=\mathrm{ChirHoch}^\bullet(\mathcal F_X^{\mathrm{hyb}},
\mathcal F_X^{\mathrm{hyb}})\); cf.\
\texttt{chiral-bar-cobar-vol2:CLAUDE.md~§6} (universal stage chain
\(\mathsf P\to\mathsf C\to\mathsf S\to\mathsf Z\)).
The chiral Koszul comparison
\(\Theta_{\mathrm{Kos}}:\Omega_E^{\mathrm{ch}}C_X^{\mathrm{geom}}
\xrightarrow{\simeq}\mathsf{Den}_{\Delta_5,E}\) is a comparison
between two algebras of the same level; the homotopy inverse to the
target-internal counit
\(\Omega_E^{\mathrm{ch}}\bar B_E^{\mathrm{ch}}
(\mathsf{Den}_{\Delta_5,E})\to\mathsf{Den}_{\Delta_5,E}\) is not used
in either direction.

\subsection{Construction~8 — Primitive recognition}
\label{app:build-recognition}

\paragraph{Tag.} \(\delta\)-endpoint convergence
(source-to-target identification).

\paragraph{Input.}
The descended source primitive Lie superalgebra
\(P_X^{\Pi,\mathrm{red}}\) of
Construction~\ref{app:build-hall} (after the radical quotient
\(\overline{\mathrm{Rad}\langle\,\cdot,\cdot\rangle_X}\)); the imported
target denominator algebra
\(\fg_{\Delta_5}=G(\mathscr B_{\Delta_5})\) on
\(\La^{2,1}_{II}\) of Definition~\ref{def:bkm-imaginary-simple-roots} and
Definition~\ref{def:bkm-relations}; the increasing cofinal
sequence of finite downward-saturated windows
\(W_1\subset W_2\subset\cdots\) with
\(\bigcup_\nu W_\nu=\mathcal R_+\); the Borcherds--Kac--Moody data
\((H,I,p,\alpha_i,(\,,\,))\) of
Theorem~\ref{thm:primitive-recognition}.

\paragraph{Output.}
The cofinal finite-window primitive-recognition datum
\(\{\textup{(R0),(R1),(R2),(R3),(R4),(R5),(R6),(R7),(R8)}\}\) of
Definition~\ref{def:cofinal-primitive-recognition-datum} for every
\(W_\nu\), compatible under \(W_\nu\subset W_{\nu+1}\); the source
representatives
\(e_i^X,f_i^X,h_i^X\); the source presentation map
\(\pi_W:\widehat{\mathfrak F}(\mathscr B_{\Delta_5})_W
\to P_X^{\Pi,\mathrm{red}}\big|_W\); the kernel equality
\(\ker\pi_W=
(\overline{J_{\mathrm{BK}}+\mathrm{Rad}_{\mathrm{GN}}})_W\); the
isomorphism
\[
\boxed{\quad
P_X^{\Pi,\mathrm{red}}\cong\fg_{\Delta_5},\quad
\Omega_E^{\mathrm{ch}}C_X\simeq\mathsf{Den}_{\Delta_5,E}.
\quad}
\]

\paragraph{Procedure.}
For every cofinal \(W_\nu\):
\begin{enumerate}
\item[\textup{(R0)}] \emph{Compact source labels.}  Verify every
finite source basis vector in \(W_\nu\cup(-W_\nu)\cup\{0\}\) satisfies
the compact source-admissibility condition of
Theorem~\ref{thm:finite-stage-dirac-igusa-pro-object}.

\item[\textup{(R1)}] \emph{Representatives and parities.}  Supply
parity-homogeneous primitive representatives
\(e_i^X\in P^{\Pi,\mathrm{red}}_{X,\alpha_i}\),
\(f_i^X\in P^{\Pi,\mathrm{red}}_{X,-\alpha_i}\),
\(h_i^X=\iota_H^{-1}(\alpha_i)\) for every simple-root index \(i\) with
\(\alpha_i\in W_\nu\), with the Gritsenko--Nikulin parity fibres
\(\mu_{\overline 0}(a),\mu_{\overline 1}(a)\) for imaginary simple
roots and even real representatives for \(\delta_1,\delta_2,\delta_3\).

\item[\textup{(R2)}] \emph{Relations.}  Verify in the protected Hall
super-bracket the Cartan, Chevalley, real Serre, Borcherds
orthogonality, and super sign relations for every relation whose
displayed source and target degrees lie in
\(W_\nu\cup(-W_\nu)\cup\{0\}\).  In the real rank-three subsystem the
Cartan matrix is
\(((\delta_i,\delta_j))=\mathrm{diag}(2)-2(J-\mathrm{diag}(1))\) with
real Serre exponent~3.

\item[\textup{(R3)}] \emph{Full finite parity dimensions.}  For every
\(\beta\in W_\nu\) verify
\(\dim P^{\Pi,\mathrm{red}}_{X,\beta,\overline p}
=\rootmult_{\overline p}^{\mathrm{GN}}(\beta)\) for both parities
\(p=0,1\).  In the first timelike channel
\(\beta=\delta_{123}\),
Proposition~\ref{prop:bkm-delta123-presentation-count} requires
\((d^{\mathrm{GN}}_{(1,1,1),\overline 0},
d^{\mathrm{GN}}_{(1,1,1),\overline 1})=(29,93)\).  Determinant
information \(29-93=-64=f(1,1)\) is consistent but not the
recognition.

\item[\textup{(R4)}] \emph{Hopf pairing and radical.}  Compute the
positive-negative pairing matrices in parity blocks; verify the
kernel agrees with the
Gritsenko--Nikulin/Kac invariant-pairing kernel; verify Lie-ideal and
coideal stability under all finite brackets; verify carrying under
HN transitions to the next window.

\item[\textup{(R5)}] \emph{No extra relations.}  Verify the kernel
equality
\(\ker\pi_\nu=(\overline{J_{\mathrm{BK}}+\mathrm{Rad}_{\mathrm{GN}}})_{\le W_\nu}\) on the finite triangular span generated by
\(W_\nu\cup(-W_\nu)\cup\{0\}\).

\item[\textup{(R6)}] \emph{Generation.}  Verify that every source
primitive root space in \(W_\nu\cup(-W_\nu)\), after the finite
radical quotient, is generated by the Cartan and the supplied
simple-root representatives through brackets whose intermediate
positive degrees stay in \(W_\nu\).

\item[\textup{(R7)}] \emph{Completed PBW compatibility.}  Verify the
PBW-graded isomorphism
\(\mathrm{gr}_{\mathrm{PBW}}U(P_X^{\Pi,\mathrm{red}})_{\le W_\nu}
\xrightarrow{\sim}
\mathrm{gr}_{\mathrm{PBW}}U(\fg_{\Delta_5})_{\le W_\nu}\); verify
strict preservation under \(W_{\nu+1}\to W_\nu\) transition maps.

\item[\textup{(R8)}] \emph{Exact triangular completion.}  Verify
strictness of the transition maps and closedness of images on every
finite window, so that passage to the inverse limit commutes with the
finite kernels, cokernels, radicals, PBW associated gradeds, and
parity pieces.  Equivalently the relevant \(\lim^1\) terms vanish for
these triangular finite-window systems.

\item Output the global isomorphism
\(P_X^{\Pi,\mathrm{red}}\cong\fg_{\Delta_5}\) by passage to the inverse
limit (compatible cofinal datum); together with the chiral Koszul
comparison of Construction~\ref{app:build-koszul} this gives
\(\Omega_E^{\mathrm{ch}}C_X\simeq\mathsf{Den}_{\Delta_5,E}\) and hence
the protected Pfaffian identification
\(\operatorname{Pf}_{\mathrm{prot}}(\mathfrak D_X^{\mathrm{DI}})=\Delta_5\).
\end{enumerate}

\paragraph{Type checks.}
The source presentation map \(\pi:\widehat{\mathfrak F}(\mathscr B_{\Delta_5})\to P_X^{\Pi,\mathrm{red}}\) is
continuous and degree-preserving; \textup{(R5)} gives injectivity
after the completed radical quotient; \textup{(R6)} gives
surjectivity; \textup{(R7)} gives PBW compatibility; \textup{(R8)}
gives exactness in the inverse-limit step.  No claim on signed
superdimensions \(\sdim=f(nm,l)\) substitutes for the two-integer
parity dimensions of \textup{(R3)}.

\paragraph{Obstructions discharged or invoked.}
This step inherits the obstruction tower discharged by
Constructions~1--7 and adds the cofinal exact-completion residual
\textup{(R8)}.  The Beilinson--Drinfeld bar-cobar firewall is
respected throughout: \textup{(R5)} is a source kernel-equality
theorem, not an inversion of the target bar-cobar counit;
\textup{(R7)} is a PBW graded comparison of source and target
enveloping algebras, not a graph-equality argument from the scalar
square.

\paragraph{Cross-volume anchor.}  The output is
\texttt{chiral-bar-cobar-vol2:CLAUDE.md~§9} Construction
Problem~1: the operator-level construction of
\(\mathfrak D_X^{\mathrm{DI}}\) for \(K3\times E\) with
\(\operatorname{Pf}_{\mathrm{prot}}(\mathfrak D_X^{\mathrm{DI}})=\Delta_5\), under the
four obstructions \((D_0),(O1),(O1)^+,(O2)\).  The Drinfeld doubling
\(Y^+(X)\to G(X)=D(Y^+(X))\) of
\texttt{calabi-yau-quantum-groups:CLAUDE.md} requires the Hopf
pairing of \textup{(R4)}, completion by \textup{(R8)}, integral form,
and stable-envelope transport which are supplied by the same
data.  The κ-tuple coordinate \(\kappa_{\mathrm{BKM}}=5\) is the
weight of \(\Delta_5\) and lives at level~4 of the universal stage
chain; it is a consequence of the construction, not an input to it.

\subsection{The eight-step audit}
\label{app:build-audit}

The eight construction steps form a strict ordered sequence of
input/output type discipline.  An audit of any candidate construction of
\(\mathfrak D_X^{\mathrm{DI}}\) reduces to verifying, for every
\(R\in\mathcal R_{\mathrm{HN}}\):
\begin{enumerate}
\item the lattice combinatorics of
Construction~\ref{app:build-charge-window};
\item the finite-type quasi-smooth derived stacks of
Construction~\ref{app:build-finite-moduli};
\item the hybrid local/wrapped factorization datum and the eight-word
two-step flag atlas of Construction~\ref{app:build-hybrid-wrapped};
\item the four Cech--Borel orientation classes and the
\((O1)^+\)~Coxeter coherence of
Construction~\ref{app:build-orientation};
\item the finite Hall bialgebra and primitive matrices of
Construction~\ref{app:build-hall};
\item the rank-one Pfaffian crossing and the half-Hilbert orbit
transport of Construction~\ref{app:build-dirac-pfaffian};
\item the bar coalgebra, conilpotent collision coproduct, and the
Francis--Gaitsgory cobar quasi-isomorphism of
Construction~\ref{app:build-koszul};
\item the cofinal datum
\(\textup{(R0)}\)--\(\textup{(R8)}\) of
Construction~\ref{app:build-recognition};
\end{enumerate}
together with the compatible vanishing of the obstruction tuples
\((\mathfrak O^{E_3-\mathrm{src}}_R,
\mathfrak O_{\mathrm{hyb},R},
\mathfrak O_{\mathrm{Kos},R},
\mathfrak O^{\mathrm{O1}}_R,
\mathfrak O^{\mathrm{O1}^+}_R,
\mathfrak O^{\mathrm{O2}}_R,
\mathfrak O_R^{D_0})\) at every height \(R\), and the
Mittag--Leffler exactness condition \textup{[ML]} of
Definition~\ref{def:O1-strong-reduced-orientation} on the cofinal HN inverse
system.  Each obstruction is attributable to the unique step
in which its components are constructed.

\paragraph{Order rigidity.}  Each of the seven adjacent input/output
dependencies fails type-checking under exchange:
\begin{enumerate}
\item charge window $\to$ finite moduli: $\Lambda_{\mathrm{ample}}$
  parametrises the HN strata; without the window, the moduli stack
  has no finite cofinal subsystem.
\item finite moduli $\to$ hybrid carrier: $\mathrm{Ran}^{\mathrm{hyb}}(E)$
  needs the finite Hall stack as factor on which to graft the
  wrapped/local atlas.
\item hybrid carrier $\to$ reduced orientation: the cosection-twisted
  Joyce--Upmeier line bundle is defined on the carrier, not on the
  raw moduli (the cosection lives on $\mathrm{Ran}^{\mathrm{hyb}}$).
\item reduced orientation $\to$ Hall correspondences: the Pfaffian
  line bundle on the Hall stack is the orientation-twisted sign data;
  without orientation, the Hall product is unsigned.
\item Hall correspondences $\to$ Pfaffian line + Dirac operator: the
  primitive Hall bracket $[-,-]_X$ defines the Hall--Drinfeld
  correspondence whose Pfaffian section is the Dirac--Igusa Pfaffian.
\item Pfaffian line $\to$ chiral Koszul source coalgebra: the source
  $C_X$ is the bar of the augmented hybrid factorization algebra
  pulled back along the Pfaffian comparison
  $\iota_{\mathrm{aut}}$.
\item chiral Koszul source $\to$ primitive recognition: the primitive
  part $P_X^{\Pi,\mathrm{red}}$ is the Frobenius--Serre primitive
  filtration on the source, which the recognition theorem matches to
  $\mathfrak g_{\Delta_5}$.
\end{enumerate}
Two diagnostic exchanges: orientation before hybrid leaves the middle
type-II wall \(\delta_2\leftrightarrow(0,1,1)\) without a
representative, by
Lemma~\ref{app:pfaffian-wall-three-walls}; recognition before the
Koszul cobar leaves the identification
\(\Omega_E^{\mathrm{ch}}C_X\simeq\mathsf{Den}_{\Delta_5,E}\)
unconstructed.  All other adjacent exchanges fail by the same
input/output type criterion.

\paragraph{Conditional theorem.}  Granted the eight constructions and
the four obstructions \((D_0),(O1),(O1)^+,(O2)\),
Theorem~\ref{thm:pfaffian-dirac-finite-stage} reads
\[
\operatorname{Pf}_{\mathrm{prot}}(\mathfrak D_X^{\mathrm{DI}})=\Delta_5,\qquad
\epsilon_o=\nu_{\Delta_5}\big|_{W^{(2)}(\La^{2,1}_{II})},\qquad
P_X^{\Pi,\mathrm{red}}\cong\fg_{\Delta_5}.
\]
The scalar square
\(Z_{\mathrm{BPS}}^{K3\times E}=
(\Phi_{10}^{\mathrm{un}})^{-1}=\Delta_5^{-2}\) is the level-4 trace
of the constructed protected algebra; promotion to the 3d
gravitational path integral is not claimed and requires the
Hall--Borcherds residual of
\texttt{chiral-bar-cobar-vol2}.

%%% End of Appendix F.
 after
% \appendix has been opened at main.tex line 17081, so that
% \thesection = \Alph{section}.  No \chapter, no second \appendix.

\section{Eight build algorithms for the Dirac--Igusa pro-object}
\label{app:eight-build-algorithms}

The Dirac--Igusa pro-object
\[
\mathfrak D_X^{\mathrm{DI}}
=\varprojlim_R
\bigl(
\mathcal A^{E_3}_{X,R},
\mathcal F_{X,R}^{\mathrm{hyb}},
\Gamma_{X,R},
\Pi_X,
\Phi_R,
o_R,
\mathcal H_R,
P_R^\Pi,
C_{X,R},
\Theta_{\mathrm{Kos},R},
\mathscr L_{\mathrm{Pf},R},
\operatorname{pf}_{X,R},
\epsilon_{o,R},
\mathrm{Rec}_R
\bigr)
\]
of Definition~\ref{thm:finite-stage-dirac-igusa-pro-object} and
Definition~\ref{def:O1-strong-reduced-orientation} is the conditional source
operator-level package whose protected Pfaffian, after the four
obstructions \((D_0)\), \((O1)\), \((O1)^+\), \((O2)\) are discharged,
realizes the Igusa cusp form,
\[
\operatorname{Pf}_{\mathrm{prot}}(\mathfrak D_X^{\mathrm{DI}})=\Delta_5,\qquad
\epsilon_o=\nu_{\Delta_5}\big|_{W^{(2)}(\La^{2,1}_{II})},\qquad
P_X^{\Pi,\mathrm{red}}\cong\fg_{\Delta_5}.
\]
Construction at finite Harder--Narasimhan height \(R\) proceeds through
eight construction steps.  Each step is a procedure on input data
already supplied by earlier steps and feeds named output data into the
next.  The eight steps stratify the realization datum
\(\mathfrak K_X=(L_X,H_X,O_X,D_X,R_X)\) of
Definition~\ref{thm:finite-stage-dirac-igusa-pro-object}: steps 1--3
realize \((L_X)\) and \((H_X)\); step 4 realizes the local content of
\((O_X)\); step 5 produces the Hall correspondence object; step 6 fixes
the Pfaffian--Dirac line and feeds the orientation character; steps 7--8
realize \((D_X)\) and \((R_X)\).

Throughout, \(X=S\times E\) with \(S\) a smooth complex K3 surface and
\(E\) a smooth elliptic curve.  \(\sigma=\sigma_{s,t}\) is Liu's product
stability condition on the Abramovich--Polishchuk heart in
\(\mathsf{D}^b(\mathrm{Coh}(X))\)
\cite{AbramovichPolishchukTStructures,LiuProductStability}.  The HN
sector \((\sigma,S)\) is fixed and the cofinal subsystem
\(\mathcal R_{\mathrm{HN}}=\{R_0<R_1<\cdots\}\) is the one of
Definition~\ref{def:O1-strong-reduced-orientation}, governed by
\textup{[K3$\times$E-fin]}, \textup{[K3$\times$E-atlas]}, and
\textup{[ML]}.

Each construction step carries a licensing tag pulled from the
five-type universal stage chain
\(\mathsf P\to\mathsf C\to\mathsf S\to\mathsf Z\to\mathsf A\) of
\texttt{chiral-bar-cobar-vol2:CLAUDE.md~§9}: \(\alpha\) primitive,
\(\beta\) functorial passage, \(\gamma\) shadow specialization,
\(\delta\) endpoint convergence, \(\epsilon\) terminal scalar.  The
ordering charge\(\to\)moduli\(\to\)hybrid\(\to\)orientation\(\to\)Hall
\(\to\)Dirac+Pfaffian\(\to\)Koszul\(\to\)recognition is the unique
ordering compatible with input/output type discipline; reordering any
two adjacent steps fails an input check.

The four obstructions are attributed to specific steps:
\((D_0)\) is the cofinal-trivialization datum of the entire eight-step
construction at every finite \(R\) (steps 1--6 supply its components);
\((O1)\) lives at step 4; \((O1)^+\) lives at step 4 with a Coxeter
coherence check feeding step 6; \((O2)\) lives at step 6.  Step 8
inherits the obstruction tower discharged by steps 1--7 and adds the
exact-completion residual \textup{(R8)} of
Definition~\ref{def:cofinal-primitive-recognition-datum}.

\subsection{Construction~1 — Charge window}
\label{app:build-charge-window}

\paragraph{Tag.} \(\alpha\)-primitive (lattice datum).

\paragraph{Input.}
HN height \(R\in\mathcal R_{\mathrm{HN}}\); the formal dyonic Mukai
lattice \(\Gamma_X^{\mathrm{form}}\) of
Definition~\ref{def:gram-index-lattice}; the Gram map
\(\Pi_X:\Gamma_X^{\mathrm{form}}\to\Gamma_{\mathrm{gram}}\) and the
cocycle \(B\) of Lemma~\ref{lem:gram-additivity-defect}; the active
support \(\Gamma_{\mathrm{act}}=\{\gamma\in\Gamma_{\mathrm{eff}}\mid
f(nm,l)\ne0\}\) of the Igusa Jacobi form \(\phi_{0,1}\); the cofinal
subsystem \(\mathcal R_{\mathrm{HN}}\).

\paragraph{Output.}
A finite cusp-positive active window
\[
A_R\subset\Gamma_{\mathrm{act}}\subset\Gamma_{\mathrm{gram}},
\qquad
A_R\hbox{ downward saturated for the cofinal target-window order},
\]
together with the normal-ordered extension
\[
\widehat\Gamma_R=
\{(c,T)\mid c\in\Gamma_R^{HN},\ T\in\mathcal T_R(c)\}
\subset\widehat\Gamma_X,
\qquad
\Gamma_R^\Pi=\overline\Pi_X(\widehat\Gamma_R)\subset\Gamma_{\mathrm{gram}},
\]
the lift \(L:\Gamma_{\mathrm{gram}}\to\Gamma_X^{\mathrm{form}}\) of
Lemma~\ref{lem:primitive-lift-every-gram-triple}, and the finite test
window \(\Gamma_R^{\mathrm{test}}\) of
Definition~\ref{def:O1-strong-reduced-orientation}.

\paragraph{Procedure.}
\begin{enumerate}
\item Compute the active support cone
\[
\Gamma_{\mathrm{act}}=\{(n,l,m)\in\Gamma_{\mathrm{eff}}\mid f(nm,l)\ne 0\}
\]
from the Fourier expansion of the index-one weak Jacobi form
\(\phi_{0,1}(\tau,z)=\sum_{n,l}c(n,l)q^nr^l\) of
\cite{EZ:1}, with the index substitution
\(f(nm,l):=c(nm,l)\) tabulated in
Appendix~\ref{sec:n1-boundary-compatibility-conditions} for the
\(N=1\) Igusa boundary.

\item Set
\[
A_R^{\circ}
=\{\gamma=(n,l,m)\in\Gamma_{\mathrm{act}}\mid n+|l|+m\le R\}.
\]
Form its downward saturation
\[
A_R
=\{\gamma'\in\Gamma_{\mathrm{act}}\mid
\exists\gamma\in A_R^{\circ},\ 0<\gamma'\le\gamma,\
\gamma-\gamma'\in Q_+\},
\]
in the sense of Definition~\ref{def:cofinal-primitive-recognition-datum}.

\item Normal-ordered lift.  For each \(c\in\Gamma_R^{HN}\) compute the
\(B\)-translate orbit
\[
\mathcal T_R(c)
=\Bigl\{\sum_{i<j}B(c_i,c_j)\Bigm|c=\sum_i c_i,\
c_i\in F_{\sigma,S}(R),\ \sum_i h_S(c_i)\le R\Bigr\}.
\]
Set \(\widehat\Gamma_R=\{(c,T):c\in\Gamma_R^{HN},\ T\in\mathcal T_R(c)\}\).

\item Finite target test window:
\(\Gamma_R^{\mathrm{test}}=\overline\Pi_X(\widehat\Gamma_R)\), compatibly
with \(R\to R'\) by inclusion.

\item Output \((A_R,\widehat\Gamma_R,\Gamma_R^\Pi,\Gamma_R^{\mathrm{test}},L)\).
\end{enumerate}

\paragraph{Type checks.}
\(A_R\) is finite and downward saturated;
\(\overline\Pi_X(\widehat\Gamma_R)=\Gamma_R^\Pi\) lies in
\(\Gamma_{\mathrm{gram}}\); \(\Gamma_R^{HN}\subset
F_{\sigma,S}(R)\) is finite by \textup{[K3$\times$E-fin]}; the cocycle
identity
\(\Pi_X(c+c')=\Pi_X(c)+\Pi_X(c')+B(c,c')\) is an axiom of the cocycle.

\paragraph{Obstructions discharged or invoked.}
None; this step is purely lattice combinatorics from the imported
Gritsenko--Nikulin Borcherds product
\cite{GN,GNPartI,GNII,BorcherdsGKM88,Borcherds95}.
The downward saturation is the input requirement of \textup{(R0)} in
Definition~\ref{def:cofinal-primitive-recognition-datum}.

\subsection{Construction~2 — Finite moduli}
\label{app:build-finite-moduli}

\paragraph{Tag.} \(\alpha\)-primitive (geometric source).

\paragraph{Input.}
The active window \(A_R\) and lift \(L\) of
Construction~\ref{app:build-charge-window}; Liu's product stability
condition \(\sigma_{s,t}\); the standing hypotheses
\textup{[K3$\times$E-fin]} and \textup{[K3$\times$E-atlas]}.

\paragraph{Output.}
A finite class set \(C_R\) and lift
\(\ell_R:\Gamma_R^\Pi\to C_R\) with
\(\Pi_X(\ell_R(\gamma))=\gamma\); the family of finite-type quasi-smooth
derived Artin substacks
\[
\mathfrak M^{ss}_{R,c}
=\mathfrak M^{ss}_{\sigma_{s,t}}(X,c)_{\le R}
\subset\operatorname{Perf}(X),
\qquad c\in C_R,
\]
of \(\sigma_{s,t}\)-semistable objects of full Liu numerical class
\(c\) and HN height \(\le R\); the scalar-rigidified truncation
\(\mathfrak M^{\mathrm{sc}}_{R,c}\); the universal perfect complex
\(\mathcal A_{R,c}\) on \(X\times\mathfrak M^{\mathrm{sc}}_{R,c}\); the
finite stratification by retained closed substacks of finite residual
inertia.

\paragraph{Procedure.}
\begin{enumerate}
\item Form \(C_R=\{c=L\gamma:\gamma\in A_R\}\) closed under retained
extension sums (taken in \(\Gamma_X^{\mathrm{form}}\)).

\item For each \(c\in C_R\) cut the substack of
\(\sigma_{s,t}\)-semistable objects of class \(c\) by HN height \(\le
R\):
\[
\mathfrak M^{ss}_{R,c}=
\bigl\{[A]\in\operatorname{Perf}(X)\mid
\sigma_{s,t}\hbox{-semistable},\
\hn(A)\le R,\
\nu(A)=c\bigr\}.
\]

\item Verify finite-typeness: \textup{[K3$\times$E-fin]} (Liu product
stability + product-boundedness theorem at fixed full Liu charge with
finite fibres over the active charge) gives the finite-type assertion.
PTVV \cite{PTVV2013} on the Calabi--Yau threefold \(X\) gives the
\((-1)\)-shifted symplectic structure on
\(\operatorname{Perf}(X)\), restricted to the retained semistable
sector.

\item Quasi-smoothness: the universal perfect complex
\(\mathcal A_{R,c}\) supplies a relative cotangent complex of bounded
Tor-amplitude on the retained charts.  Tangent at \(A\) is
\(T_A\mathfrak M^{ss}_{R,c}\simeq\operatorname{RHom}_X(A,A)[1]\).

\item Scalar rigidification: form \(\mathfrak M^{\mathrm{sc}}_{R,c}\) by
killing the scalar \(\mathbb G_m\)-automorphism of
\(c\)-stable factors.

\item Two-step flag stack
\(\mathfrak F^{(2)}_{R,c,c',c''}\) and extension stack
\(\mathfrak E_{R,c,c'}\) constructed as iterated derived vector stacks
of \(R\pi_*\mathcal RHom_X(\mathcal B,\mathcal A[1])\); both finite type
and quasi-smooth in the retained sector by
Proposition~\ref{prop:finite-dcritical-hall-product}.

\item Output \((C_R,\ell_R,\{\mathfrak M^{ss}_{R,c}\},
\{\mathfrak E_{R,c,c'}\},\{\mathfrak F^{(2)}_{R,c,c',c''}\})\) with
universal perfect complexes and the finite stratification.
\end{enumerate}

\paragraph{Type checks.}
\(\mathfrak M^{ss}_{R,c}\) is a finite-type quasi-smooth derived Artin
substack of \(\operatorname{Perf}(X)\); the scalar rigidification leaves
only finite residual inertia
\(H_S\subset E_{\mathrm{tors}}\) on each connected stabilizer
component, by
Proposition~\ref{prop:finite-dcritical-hall-product};
PTVV \((-1)\)-shifted symplectic structure restricts.

\paragraph{Obstructions discharged or invoked.}
\textup{[K3$\times$E-fin]} is invoked here in its product-boundedness
form and supplies the finite-type assertion not in the cited literature
for the K3\(\times\)E threefold; it is part of the standing datum.
\textup{[K3$\times$E-atlas]} supplies the quasi-smooth d-critical
charts.  These enter \((D_0)\) at the level of the finite Hall atlas
of Theorem~\ref{thm:finite-stage-dirac-igusa-pro-object}.

\subsection{Construction~3 — Hybrid wrapped factorization carrier}
\label{app:build-hybrid-wrapped}

\paragraph{Tag.} \(\beta\)-functorial passage (Stage-1 carrier).

\paragraph{Input.}
The finite moduli stacks
\(\{\mathfrak M^{\mathrm{sc}}_{R,c}\}\) of
Construction~\ref{app:build-finite-moduli}; the
projection-finite/wrapped split by elliptic degree
\(b_R^{\mathrm{geom}}:\Gamma_{\sigma,S}^{HN}(R)\to\Z_{\ge0}\) of
Definition~\ref{def:hybrid-wrapped-factorization-datum}; a fixed
origin \(0_E\in E\); the Abramovich--Polishchuk/Liu wrapped anchor
candidate.

\paragraph{Output.}
The finite hybrid local/wrapped factorization datum
\[
\mathcal F_{X,\sigma,S,\le R}^{\mathrm{hyb}}
=
(\mathcal C_{\sigma,S,\le R},
b_R^{\mathrm{geom}},
\mathcal F_R^{\mathrm{loc}},
\mathcal G_R^{\mathrm{wr}},
\mathfrak E_R^{\mathrm{loc}},
\mathfrak E_R^{\mathrm{hyb}},
\mathfrak E_R^{\mathrm{wr}},
\mathfrak{Flag}_R^{\mathrm{hyb}},
\mathsf{Corr}_R^{E,\mathrm{hyb}},
\lambda_R,
Q_{E,R},
\theta^Q_R,
o_R,
I_R^{\mathrm{prot}})
\]
of Definition~\ref{def:hybrid-wrapped-factorization-datum},
together with the hybrid Ran base
\(\operatorname{Ran}^{\mathrm{hyb}}(E)\), the eight-word two-step
binary flag atlas \(\{w\in\{\mathrm L,\mathrm W\}^3\}\), and the
quotient-after-correspondence functor \(Q_{E,R}\).

\paragraph{Procedure.}
\begin{enumerate}
\item \emph{Elliptic-degree split.}  Decompose
\(\Gamma_{\sigma,S}^{HN}(R)
=\Gamma_R^{\mathrm{loc}}\sqcup\Gamma_R^{\mathrm{wr}}\) where
\(\Gamma_R^{\mathrm{loc}}=\{b_R^{\mathrm{geom}}=0\}\) and
\(\Gamma_R^{\mathrm{wr}}=\{b_R^{\mathrm{geom}}>0\}\).  Verify additivity
\(b_R^{\mathrm{geom}}(\gamma+\gamma')
=b_R^{\mathrm{geom}}(\gamma)+b_R^{\mathrm{geom}}(\gamma')\) on retained
extensions remaining in the quotient.

\item \emph{Projection-finite local sector.}  For
\(\alpha\in\Gamma_R^{\mathrm{loc}}\) and finite index set \(I\), build
the closed-support incidence stack over \(E^I\):
\[
\mathfrak M^{\mathrm{loc,cl}}_{\alpha,R,I}
=
\Bigl\{(A,\mathbf e)\in\mathfrak M^{\mathrm{sc}}_{R,L\alpha}\times E^I
\Bigm|p_E(\mathrm{Supp}\,A)\subset\bigcup_{i\in I}S\times\{e_i\}\Bigr\}.
\]
Form
\(\mathscr H^{\mathrm{loc}}_{\alpha,R,I}=
(\pi_{\alpha,R,I})_!(\Phi_{\alpha,R,I}^{\mathrm{loc}}\otimes
o_{\alpha,R,I}^{\mathrm{loc}})\) and the open-support sheaf
\(\mathcal F_{\alpha,R}^{\mathrm{loc}}(U)\) over \(U\subset E\) by
restricted Borel--Moore configuration sums.

\item \emph{Wrapped prequotient sector.}  For
\(\eta\in\Gamma_R^{\mathrm{wr}}\), form the rigidified prequotient
\(\mathcal M_{\eta,R}^{\mathrm{wr,rig}}\to\mathcal M_{\eta,R}^{\mathrm{wr}}\)
by the Abramovich--Polishchuk / determinant anchor
\[
\lambda_\eta:\mathcal M_{\eta,R}^{\mathrm{wr,rig}}\longrightarrow E,
\qquad
\lambda_\eta(A)=\det Rp_{E,*}A
\otimes\mathcal O_E(-\chi(A)\cdot 0_E),
\]
divided by a finite cover when \(\chi(A)\ne 1\) so that
\(\lambda_\eta(t\cdot A)=\lambda_\eta(A)+t\) holds.  Verify the
losslessness residual \(\mathfrak o^{\lambda,\mathrm{loss}}_R=0\) by
the Ext-distinguishing test of
Definition~\ref{def:hybrid-wrapped-factorization-datum}, separating
\(W_0=i_{s,*}(\mathcal O_E\oplus\mathcal O_E)\) from
\(W_L=i_{s,*}(L\oplus L^{-1})\).

\item \emph{Mixed local/wrapped correspondences.}  For an ordered
hybrid input
\(\alpha_-\in\Gamma_R^{\mathrm{loc}}\),
\(\eta\in\Gamma_R^{\mathrm{wr}}\),
\(\beta_+\in\Gamma_R^{\mathrm{loc}}\), form
\(\mathfrak E^{\mathrm{mix}}_{\alpha_-;\eta;\beta_+,R}\) before the
reduced \(E\)-quotient.

\item \emph{Eight-word two-step flag atlas.}  Build flag stacks
\(\mathfrak F^{(2),w}_R\) for each
\(w\in\{\mathrm L,\mathrm W\}^3=
\{\mathrm{LLL},\mathrm{LLW},\mathrm{LWL},\mathrm{WLL},\mathrm{LWW},
\mathrm{WLW},\mathrm{WWL},\mathrm{WWW}\}\), classifying filtrations
with three associated graded factors of those types.  Verify the
four-input pentagon coherence and the quotient-after-correspondence
descent.

\item \emph{Hybrid correspondence-module functor.}  Assemble the
correspondence category
\(\mathsf{Corr}^{E,\mathrm{hyb}}_R\) and the functor
\(\mathcal B_R^{E,\mathrm{hyb}}:\mathsf{Corr}^{E,\mathrm{hyb}}_R
\to\mathsf{Ch}_{\C}\) by
\(\mathcal B_R^{E,\mathrm{hyb}}(e)=q_!\,\mathrm{TS}^{\mathrm{red}}_e\,p^*\)
for every generating arrow
\(e=(Y\xleftarrow{p}\mathfrak E_e\xrightarrow{q}Z)\).

\item \emph{Quotient functor \(Q_{E,R}\).}  Apply
\(E\)-equivariant compactly supported vanishing-cycle chains and the
orientation descent of step~4 to produce the reduced object
\[
\mathcal F^{\mathrm{geom}}_{X,\sigma,S,\le R}
=Q_{E,R}\circ\mathcal B_R^{E,\mathrm{hyb}}.
\]

\item Output the hybrid datum with all coherences and the protected
integration map \(I_R^{\mathrm{prot}}\), with homogeneous
\(s\)-degree the Borcherds trace exponent
\(m_R=\operatorname{pr}_3\overline\Pi_X\).
\end{enumerate}

\paragraph{Type checks.}
The carrier
\(\mathcal F_{X,\sigma,S,\le R}^{\mathrm{hyb}}\) is a hybrid
correspondence-module functor on
\(\mathsf{Corr}^{E,\mathrm{hyb}}_R\) with values in
\(\mathsf{Ch}_{\C}\); the four-input pentagon is supplied by
the eight-word two-step flag atlas; the anchor residual tuple
\(\mathfrak o^\lambda_R=
(\mathfrak o^{\lambda,\mathrm{ex}}_R,
\mathfrak o^{\lambda,\mathrm{unit}}_R,
\mathfrak o^{\lambda,\mathrm{loss}}_R,
\mathfrak o^{\lambda,\mathrm{multi}}_R,
\mathfrak o^{\lambda,\mathrm{tr}}_{R'R})\) vanishes; the hybrid
residual \(\mathfrak O_{\mathrm{hyb},R}=0\) at every height \(R\),
compatibly with transition maps.

\paragraph{Obstructions discharged or invoked.}
The middle type-II wall lies in
\(\delta_2\leftrightarrow(0,1,1)\), elliptic degree
\(d=m+1=2\) by Lemma~\ref{app:pfaffian-wall-three-walls}; it
cannot be represented by a projection-finite local stack and demands
the wrapped prequotient.  This step's correctness — the existence of
the wrapped representative
\(x_{\delta_2,R}\in\mathcal M_{\eta,R}^{\mathrm{wr,rig}}\) with
\(\overline\Pi_X(x_{\delta_2,R})=(0,1,1)\) and its compatible anchor
\(\lambda_{\eta,R}\) — is precisely the hybrid input to \((O2)\) at
the middle wall.

\subsection{Construction~4 — Reduced d-critical orientation
and Joyce--Upmeier transport}
\label{app:build-orientation}

\paragraph{Tag.} \(\beta\)-functorial passage (oriented carrier).

\paragraph{Input.}
The hybrid carrier
\(\mathcal F_{X,\sigma,S,\le R}^{\mathrm{hyb}}\) of
Construction~\ref{app:build-hybrid-wrapped}; the K3 holomorphic
\(2\)-form \(\sigma_S\) on \(S\); BBDJS d-critical chart machinery
\cite{BBDJS2015}; the Joyce--Upmeier multiplicative orientation
formalism \cite{JoyceUpmeier,JoyceTanakaUpmeier}; the cosection
localization machinery of Kiem--Li and the cosection-reduced
obstruction theory; the type-II Weyl group
\(W^{(2)}(\La^{2,1}_{II})\) generated by
\(s_{\delta_1},s_{\delta_2},s_{\delta_3}\).

\paragraph{Output.}
The cosection
\(\sigma=p_S^*\sigma_S\); the reduced tangent complex
\[
\operatorname{RHom}_{\mathrm{red}}(A,A)
=
\operatorname{Cone}\!\Bigl(
\operatorname{RHom}_X(A,A)
\xrightarrow{(\mathrm{tr},\mathrm{cs})}
R\Gamma(X,\mathcal O_X)\oplus H^2(S,\mathcal O_S)[-1]
\Bigr)[-1];
\]
the BBDJS vanishing-cycle complex \(\Phi_{R,c}\); the reduced
orientation line
\[
o_{R,c}
=\det\operatorname{Ext}^1_{\mathrm{red}}(A,A)^{-1}
\quad\text{on Darboux charts},\qquad
o_{R,c}^{\otimes 2}\simeq\det\operatorname{RHom}_{\mathrm{red}}(A,A);
\]
the Joyce--Upmeier multiplicative transport
\(o_{c\oplus c'}\simeq o_c\otimes o_{c'}\) compatible with the reduced
extension correspondences; the quotient orientation
\(\mathfrak Q^{or}_R\) of
Theorem~\ref{prop:generic-local-sign-orientation-character}; the
Weyl-equivariant lifts
\(\tau_{i,R,S}:o_{R,S}\to s_{\delta_i}^*o_{R,S}\) for \(i=1,2,3\) of
\((O1)^+\); the local generic wall-sign data \(\epsilon_{o,R}^{\mathrm{loc}}\).

\paragraph{Procedure.}
\begin{enumerate}
\item \emph{Cosection.}  Pull back the K3 semiregularity cosection
\(\mathrm{cs}\) along \(p_S:X\to S\) to obtain the 2-step cosection
\(\sigma=p_S^*\sigma_S\) on the K3 factor; verify additivity in exact
triangles by \textup{(RedOr)} of
Theorem~\ref{thm:finite-stage-dirac-igusa-pro-object}.

\item \emph{Reduced obstruction.}  Form the reduced tangent complex by
the displayed cone above (Kiem--Li cosection localization).  PTVV
\cite{PTVV2013} \((-1)\)-shifted symplectic on
\(\operatorname{RHom}_X(A,A)\) descends, after the cosection
contraction, to the reduced complex.

\item \emph{BBDJS d-critical chart.}  On each Darboux chart of
\(\mathfrak M^{\mathrm{sc}}_{R,c}\), build the BBDJS vanishing-cycle
complex \(\Phi_{R,c}\) \cite{BBDJS2015}; verify reduced
Thom--Sebastiani transport coherent on two-step flag stacks of
Construction~\ref{app:build-hybrid-wrapped}.

\item \emph{Reduced orientation line.}  Define
\(o_{R,c}=\det\operatorname{Ext}^1_{\mathrm{red}}(A,A)^{-1}\); verify
the squared isomorphism
\(o_{R,c}^{\otimes 2}\simeq\det\operatorname{RHom}_{\mathrm{red}}(A,A)\)
by \(3\)-Calabi--Yau Serre duality
\(\operatorname{Ext}^2_{\mathrm{red}}(A,A)
\simeq\operatorname{Ext}^1_{\mathrm{red}}(A,A)^\vee\).

\item \emph{Quotient-orientation descent on the K3\(\times\)E-quotient
self-Ext frame bundle.}  Compute the four Cech--Borel obstruction
classes for every retained stratum \(S\):
\[
\alpha^{\mathrm{red}}_{R,S}=w_2(\mathscr D_{R,S}),
\quad
\alpha^{E,\mathrm{free}}_{R,S}=
\mathrm{edge}_{2,0}(w_2^E(\mathscr D_{R,S})),
\quad
\widetilde\beta^H_{R,S}=w_2^G(\mathscr D_{R,S}),
\quad
\lambda^H_{R,S}\in H^1_G(S;\mathbb F_2),
\]
and verify simultaneous vanishing with compatible
null-trivializations on overlaps, on Hall correspondence pullbacks,
on two-step flag stacks, and under successor maps.  The full
calculation proceeds primary stabilizer by stabilizer; for cyclic
\(\Z/2^a\) the absent mixed coefficient requires the rank-two
two-primary check.  Output the descended quotient orientation system
\(\mathfrak Q^{or}_R\).

\item \emph{Multiplicative Joyce--Upmeier transport.}  Transport
\(o_{c\oplus c'}\simeq o_c\otimes o_{c'}\) coherently across two-step
flag stacks
\cite{JoyceUpmeier,JoyceTanakaUpmeier}; verify the reduced
Thom--Sebastiani identification on three-input flags.

\item \emph{Weyl-equivariant lifts \((O1)^+\).}  For each
\(i\in\{1,2,3\}\) construct the lift
\(\tau_{i,R,S}:o_{R,S}\to s_{\delta_i}^*o_{R,S}\); verify the
projective Weyl cocycle
\(c_{o,R}\in Z^2(\mathcal G_R;\underline{\mathbb F}_2)\) on the
retained action groupoid; null-trivialize \([c_{o,R}]\) and choose a
\(1\)-cochain killing it; verify \(\tau_{i,R,S}^2=1\); verify the
Coxeter relations
\((s_{\delta_i}s_{\delta_j})^3=1\) for \(i\ne j\) (the
\((3,3,3)\)-hyperbolic triangle relations of
Appendix~\ref{app:pfaffian-wall-three-walls}); verify that the torsor
defects \(\omega_{i,\mathcal C}\) of
Definition~\ref{thm:finite-stage-dirac-igusa-pro-object}~$(O_X)$ vanish.

\item \emph{Local wall-sign datum.}  At each generic type-II wall point
\(\Hpl_{\delta_i}\), with the rank-one local model
\(K_{\delta_i}^{\mathrm{cross}}=[\R\xrightarrow{u_i}\R]\), record the
local generic sign \(\epsilon_{o,R}^{\mathrm{loc}}(s_{\delta_i})=-1\)
of Proposition~\ref{prop:rank-one-crossing-o2-sign}.

\item Output
\((o_R,\mathfrak Q^{or}_R,\{\tau_{i,R,S}\},\epsilon_{o,R}^{\mathrm{loc}})\).
\end{enumerate}

\paragraph{Type checks.}
\(o_R\) is a Picard-line section squaring to the reduced determinant;
the four Cech--Borel classes vanish with compatible
null-trivializations; the lifts \(\tau_{i,R}\) satisfy
\(\tau_{i,R}^2=1\) and the \((3,3,3)\)-Coxeter relations; multiplicative
transport agrees on flag stacks.

\paragraph{Obstructions discharged or invoked.}
This step is the entire site of \((O1)\) (strong reduced orientation
on the K3\(\times\)E-quotient self-Ext frame bundle) and of
\((O1)^+\) (Weyl-equivariant transport along
\(W^{(2)}(\La^{2,1}_{II})\) reflections).  Discharging \((O1)\)
requires the simultaneous vanishing and compatible
null-trivialization of the four Cech--Borel classes on every retained
stratum; discharging \((O1)^+\) requires \([c_{o,R}]=0\), the
Coxeter coherence, and the vanishing of all torsor defects
\(\omega_{i,\mathcal C}\).  The local wall-sign datum is the input to
\((O2)\) at step~6.

\subsection{Construction~5 — Hall correspondences and primitive part}
\label{app:build-hall}

\paragraph{Tag.} \(\gamma\)-shadow specialization
(Stage-2 chiral Hall shadow).

\paragraph{Input.}
The oriented hybrid carrier from
Construction~\ref{app:build-orientation}; the extension stack
\(\mathfrak E^{\mathrm{geom}}_{R,\widehat c,\widehat c'}\); the
two-step flag stack
\(\mathfrak F^{(2),\mathrm{geom}}_{R,\widehat c,\widehat c',\widehat c''}\);
the BBDJS reduced Thom--Sebastiani transport; six-functor admissibility
of \(p^*,q_!,p_!\) (Construction~\ref{app:build-finite-moduli}~step
\textup{(Six)}).

\paragraph{Output.}
The finite Hall object
\[
\mathcal H_R^{\mathrm{geom}}
=\bigoplus_{\widehat c\in\widehat\Gamma_R}
H_c^\bullet\!\bigl(
\mathfrak M^{\mathrm{geom}}_{R,\widehat c},
\Phi_{R,c}\otimes o_{R,c}\bigr);
\]
the Hall product \(m_R=q_!\,\mathrm{TS}^{\mathrm{red}}\,p^*\) and Hall
coproduct \(\Delta_R=p_!\,(\mathrm{TS}^{\mathrm{red}})^{-1}\,q^*\); the
finite Hall bialgebra
\((\mathcal H_R^{\mathrm{geom}},m_R,\Delta_R,\eta_R,\epsilon_R)\); the
primitive part
\[
P_R^{\mathrm{geom}}
=\ker(\Delta_R-1\otimes\mathrm{id}-\mathrm{id}\otimes 1)\cap\ker\epsilon_R;
\]
the normal-ordered Gram descent
\(P_R^{\Pi,\mathrm{geom}}
=\overline\Pi_{X,*}^{\Theta_R}P_R^{\mathrm{geom}}\); the
\(\Gamma_R^\Pi\)-graded super-Lie bracket
\([\,\bar x,\bar y\,]_{\mathrm{red}}\) of
Proposition~\ref{thm:hall-primitives-formal};
the Hopf pairing matrix \(G_{\gamma,ab}\) and radical \(K_\gamma\).

\paragraph{Procedure.}
\begin{enumerate}
\item \emph{Pullback.}  For
\(\widehat c,\widehat c'\in\widehat\Gamma_R\), pull back along
\(p:\mathfrak E^{\mathrm{geom}}_{R,\widehat c,\widehat c'}
\to\mathfrak M^{\mathrm{geom}}_{R,\widehat c}
\times\mathfrak M^{\mathrm{geom}}_{R,\widehat c'}\); apply reduced
Thom--Sebastiani \(\mathrm{TS}^{\mathrm{red}}\) to combine vanishing
cycle data \cite{BBDJS2015}; pushforward exceptionally along
\(q:\mathfrak E^{\mathrm{geom}}_{R,\widehat c,\widehat c'}
\to\mathfrak M^{\mathrm{geom}}_{R,\widehat c\star\widehat c'}\) to
obtain
\(m_R(x\otimes y)=q_!\,\mathrm{TS}^{\mathrm{red}}\,p^*(x\boxtimes y)\).

\item \emph{Coproduct.}  Dually pull back along \(q\), apply
inverse reduced Thom--Sebastiani, pushforward along \(p\) to obtain
\(\Delta_R(z)=p_!\,(\mathrm{TS}^{\mathrm{red}})^{-1}\,q^*(z)\).
The compact geometric coproduct correspondence
\(\mathfrak C^{\mathrm{geom}}_{\rho;\mu,\nu}\) of
Proposition~\ref{thm:hall-primitives-formal} is the
support data for splitting.

\item \emph{Bialgebra check.}  Verify associativity of \(m_R\) and
coassociativity of \(\Delta_R\) on the two-step flag stack
\(\mathfrak F^{(2),\mathrm{geom}}_{R,\widehat c,\widehat c',\widehat c''}\)
by base change and Thom--Sebastiani coherence.  Verify the bialgebra
identity (Hopf compatibility) via the Massey--Saito Joyce--Upmeier
multiplicative orientation transport.

\item \emph{Primitive extraction.}  Compute
\(P_R^{\mathrm{geom}}=\ker(\Delta_R-1\otimes\mathrm{id}
-\mathrm{id}\otimes 1)\cap\ker\epsilon_R\); since
\(\Delta_R-1\otimes\mathrm{id}-\mathrm{id}\otimes 1\) is degree-graded
on \(\widehat\Gamma_R\), this reduces to a finite linear computation
in each charge.

\item \emph{Normal-ordered descent.}  Apply
\(\overline\Pi_{X,*}^{\Theta_R}\) of
Theorem~\ref{thm:hall-product} to obtain the
\(\Gamma_R^\Pi\)-graded
\(P_R^{\Pi,\mathrm{geom}}\); verify additivity
\(\overline\Pi_X(\widehat c\star\widehat c')
=\overline\Pi_X(\widehat c)+\overline\Pi_X(\widehat c')\).

\item \emph{Bracket and pairing matrices.}  Compute the matrices
\[
M^{\alpha+\beta,c}_{\alpha,a;\beta,b}
=\int_{\mathfrak E^{\mathrm{geom}}_{\alpha,\beta;\alpha+\beta}}
q^!(e_{\alpha+\beta,c}^\vee)\cap
\mathrm{TS}^{\mathrm{red}}(p^*(e_{\alpha,a}\boxtimes e_{\beta,b})),
\quad
B_{\alpha,\beta}(x\otimes y)
=M_{\alpha,\beta}(x\otimes y)
-(-1)^{|x||y|}M_{\beta,\alpha}\mathsf{sw}(x\otimes y),
\]
\[
G_{\gamma,ab}=
\mathrm{tr}_0(q_!\,\mathrm{TS}^{\mathrm{red}}\,p^*
(e_{\gamma,a}\boxtimes e_{-\gamma,b}))
\quad\hbox{on}\quad
\mathfrak P^{\mathrm{geom}}_{\gamma,-\gamma;0}.
\]
Form the radical
\(K_{\gamma,\bar p}
=\ker G_{\gamma,\bar p}\cap\ker G_{-\gamma,\bar p}^t\) and the
quotient \(L_{\gamma,\bar p}\); fix splittings.

\item Output
\((\mathcal H_R^{\mathrm{geom}},m_R,\Delta_R,P_R^{\mathrm{geom}},
P_R^{\Pi,\mathrm{geom}},
\{B_{\alpha,\beta}\},\{G_{\gamma,ab}\},\{K_\gamma\},\{L_\gamma\})\).
\end{enumerate}

\paragraph{Type checks.}
\(\mathcal H_R^{\mathrm{geom}}\) is a finite \(\widehat\Gamma_R\)-graded
Hall bialgebra; \(m_R\) is associative on the two-step flag stack;
\(\Delta_R\) is coassociative; the supercommutator of primitives is
primitive (the bialgebra identity); the descended bracket is
\(\Gamma_R^\Pi\)-graded after \(\overline\Pi_{X,*}^{\Theta_R}\);
finite-dimensional in each homogeneous parity piece by
finite-typeness.

\paragraph{Obstructions discharged or invoked.}
\((D_0)\) is partially discharged: the finite Hall package
\((\mathcal H_R^{\mathrm{geom}},m_R,\Delta_R,P_R)\) is constructed.
The seven obstruction classes
\(\mathfrak o_\Theta^{\mathrm{top}},
\mathfrak o_\Theta^{\mathrm{Hoch}},
\mathfrak o_\Theta^{\mathrm{tri}},
\mathfrak o_\Theta^{\mathrm{Jac}},
\mathfrak o_F,
\mathfrak o_\Theta^{\mathrm{cyc}},
\mathfrak o_\Theta^{\mathrm{rad}}\) of
Theorem~\ref{thm:hall-product} are required to vanish on
this finite Hall stage and compatibly in the inverse limit.

\paragraph{Cross-volume anchor.}  The output
\((\mathcal H_R^{\mathrm{geom}},m_R)\) is the finite-stage realization
of the K3-side BKM Hall--Drinfeld double
\(\mathcal D_\hbar(\mathrm{CoHA}_{K3\times E})\) of
\texttt{calabi-yau-quantum-groups:CLAUDE.md~§6.7}.  In the universal
stage chain \(\mathsf P\to\mathsf C\to\mathsf S\) the Hall positive half
\(Y^+(X)\) is precisely \(P_R^{\mathrm{geom}}\) at finite \(R\); the
Drinfeld double \(G(X)=D(Y^+(X))\) requires Hopf pairing, completion,
integral form, and stable-envelope transport which are supplied by
\(G_{\gamma,ab}\) and \((\mathrm{ML})\).  The κ-tuple coordinate
\(\kappa_{\mathrm{BKM}}=5\) is the weight of \(\Delta_5\) and the
output of the level-4 scalar trace; it is not a level-2 invariant of
the carrier itself.

\subsection{Construction~6 — Dirac operator and Pfaffian line}
\label{app:build-dirac-pfaffian}

\paragraph{Tag.} \(\gamma\)-shadow specialization
(orientation gerbe + Pfaffian).

\paragraph{Input.}
The oriented Hall package
\((\mathcal H_R^{\mathrm{geom}},o_R,\{\tau_{i,R,S}\})\) of
Constructions~\ref{app:build-orientation}--\ref{app:build-hall}; the
local wall-sign datum
\(\epsilon_{o,R}^{\mathrm{loc}}\) of
Construction~\ref{app:build-orientation}; the rank-one type-II
Pfaffian crossing datum of
Definition~\ref{def:rank-one-type-ii-crossing} and
Appendix~\ref{app:pfaffian-wall-rank-one}; the half-Hilbert orbit
under \(W^{(2)}(\La^{2,1}_{II})\) of size \(2^{12}=4096\).

\paragraph{Output.}
The Dirac operator on the orientation gerbe
\(D_R\); the Pfaffian line
\(\mathscr L_{\mathrm{Pf},R}=
\det\operatorname{Ext}^1_{\mathrm{red}}(A,A)^{-1}\) (squaring to the
reduced determinant); the protected Pfaffian section
\(\operatorname{pf}_{X,R}\); the global wall-sign character
\(\epsilon_{o,R}:W^{(2)}(\La^{2,1}_{II})\to\{\pm1\}\); the Pfaffian
section asymptotic
\[
\operatorname{pf}_{X,R}
=
64q^{1/2}r^{1/2}s^{1/2}
\prod_{\gamma=(n,l,m)\in A_R}
(1-q^nr^ls^m)^{f(nm,l)}+\hbox{higher-window correction},
\]
which exhausts to \(\Delta_5\) as \(R\to\infty\).

\paragraph{Procedure.}
\begin{enumerate}
\item \emph{Dirac operator.}  On each Darboux chart of
\(\mathfrak M^{\mathrm{sc}}_{R,c}\), with reduced normal self-Ext
splitting \(K^{\mathrm{nor}}_{R,c}\), build the odd Dirac operator
\(D_{R,c}\) acting on \(\Omega^{\mathrm{sp}}\otimes o_{R,c}\)
(spinors on the orientation gerbe); verify that, at a generic
type-II wall point, the local rank-one normal block is
\(D_{\delta_i}=\begin{psmallmatrix}0&u\\-u&0\end{psmallmatrix}\) with
\(\operatorname{Pf}(D_{\delta_i})=u\)
(Definition~\ref{def:rank-one-type-ii-crossing}).

\item \emph{Pfaffian line.}  Set
\(\mathscr L_{\mathrm{Pf},R}=
\det\operatorname{Ext}^1_{\mathrm{red}}(A,A)^{-1}\); verify
\(\mathscr L_{\mathrm{Pf},R}^{\otimes 2}\simeq
\det\operatorname{RHom}_{\mathrm{red}}(A,A)\) by
\(3\)-Calabi--Yau Serre duality.

\item \emph{Local wall sign.}  At each generic point of
\(\Hpl_{\delta_i}\) with normal model
\(K^{\mathrm{nor}}_{\delta_i,R}=[\R\xrightarrow{u_i}\R]\), compute
the monodromy of the Pfaffian section
\(\operatorname{pf}_{\delta_i,R}=\upsilon_iu_i\) along the
wall-reflection loop \(\ell_{\delta_i,R,S}\) closed by the
\((O1)^+\) lift \(\tau_{i,R,S}\):
\[
M_{\ell_{\delta_i,R,S}}(\upsilon_iu_i)
=\upsilon_i(-u_i)=-\upsilon_iu_i,
\qquad
\epsilon_{o,R}^{\mathrm{loc}}(s_{\delta_i})
=(-1)^{N^{\mathrm{Pf}}_{\delta_i,R}}=-1.
\]

\item \emph{Half-Hilbert orbit transport.}  The reflection
\(s_{\delta_i}\) acts on the half-Hilbert orbit (compact source labels
of size \(2^{12}=4096\) parametrized by 12 binary parities tied to the
\(2^{12}\) odd-spin structures on the lattice
\(\La^{2,1}_{II}/2\La^{2,1}_{II}\)).  Sum the local rank-one
contributions over the orbit:
\[
\epsilon_{o,R}(s_{\delta_i})
=\prod_{x\in\hbox{orbit}/2^{12}}
\epsilon_{o,R}^{\mathrm{loc}}(s_{\delta_i};x).
\]
The output is the global character
\(\epsilon_{o,R}:W^{(2)}(\La^{2,1}_{II})\to\{\pm 1\}\).

\item \emph{Compatibility package.}  Verify the component, boundary,
overlap, quotient, and protected integration compatibility package of
Proposition~\ref{prop:appE-local-pfaffian-sign} on every
retained type-II wall stratum.

\item \emph{Pfaffian section.}  Define
\(\operatorname{pf}_{X,R}\) as the protected integration of the
oriented Hall product over
\(\mathcal F_{X,\sigma,S,\le R}^{\mathrm{hyb}}\):
\[
\operatorname{pf}_{X,R}
=
64q^{1/2}r^{1/2}s^{1/2}
I_R^{\mathrm{prot}}\!\left(
\bigotimes_{\gamma\in A_R}\operatorname{pf}_{\gamma,R}\right),
\qquad
\operatorname{pf}_{\gamma,R}=(1-q^nr^ls^m)^{f(nm,l)}.
\]

\item Output
\((D_R,\mathscr L_{\mathrm{Pf},R},\operatorname{pf}_{X,R},
\epsilon_{o,R})\), with the protected integration compatibility data.
\end{enumerate}

\paragraph{Type checks.}
\(\mathscr L_{\mathrm{Pf},R}\) is a Picard-line whose square
isomorphism is supplied by reduced Serre duality;
\(\operatorname{pf}_{X,R}\) is a section of
\(\mathscr L_{\mathrm{Pf},R}\) with the displayed cusp expansion; the
character \(\epsilon_{o,R}\) takes values in \(\{\pm 1\}\) and
satisfies the Coxeter relations of
\(W^{(2)}(\La^{2,1}_{II})\); on each simple type-II reflection
\(\epsilon_{o,R}(s_{\delta_i})=-1\) by the rank-one local model.

\paragraph{Obstructions discharged or invoked.}
This step is the entire site of \((O2)\) (local Pfaffian wall normal
form on type-II walls; the \(2^{12}=4096\) sign sum on the
half-Hilbert orbit transport).  The local rank-one Pfaffian crossing
of Appendix~\ref{app:pfaffian-wall-rank-one} discharges \((O2)\)
generically; the global half-Hilbert orbit transport remains the open
\((O2)\) component of
Open~Problem~\ref{def:O1-plus-weyl-transport}.  The character
identity
\(\epsilon_o=\nu_{\Delta_5}\big|_{W^{(2)}(\La^{2,1}_{II})}\) becomes a
theorem only after \((D_0)\), \((O1)\), \((O1)^+\), and \((O2)\) are
discharged compatibly in \(R\).

\subsection{Construction~7 — Chiral Koszul source coalgebra}
\label{app:build-koszul}

\paragraph{Tag.} \(\gamma\)-shadow specialization
(\emph{Bar = twisting; not bulk}; the bulk is
\(Z^{\mathrm{der}}_{\mathrm{ch}}(\mathcal F_X^{\mathrm{hyb}})\),
constructed elsewhere).

\paragraph{Input.}
The oriented Hall primitives
\(P_R^{\Pi,\mathrm{geom}}\) of
Construction~\ref{app:build-hall}; the hybrid carrier
\(\mathcal F_{X,\sigma,S,\le R}^{\mathrm{hyb}}\) of
Construction~\ref{app:build-hybrid-wrapped}; the augmentation
\(\varepsilon_{X,R}^{\mathrm{hyb}}:
\mathcal F_{X,\sigma,S,\le R}^{\mathrm{hyb}}\to k\mathbf 1\); the
hybrid residual \(\mathfrak O_{\mathrm{hyb},R}=0\); the bounded
length \(N_R\) (positivity of HN height on every non-vacuum colour);
the Beilinson--Drinfeld/Francis--Gaitsgory bar-cobar formalism
\cite{BD,FrancisGaitsgoryCKD}; the conilpotent completed domain on
the elliptic curve \(E\).

\paragraph{Output.}
The conilpotent chiral coalgebra on \(E\):
\[
C_{X,R}^{\mathrm{geom}}
=\bar B_E^{\mathrm{ch}}
(\mathcal F_{X,\sigma,S,\le R}^{\mathrm{hyb}},
\varepsilon_{X,R}^{\mathrm{hyb}})
=k\mathbf 1\oplus\bigoplus_{1\le p\le N_R}
\bar B_E^{\mathrm{ch},p}(\overline{\mathcal F}_{X,R}^{\mathrm{geom}});
\]
the bar-length filtration
\(F^{\mathrm{co}}_{R,p}\); the chiral collision coproduct
\(\Delta_R^{\mathrm{ch}}\); the bar comparison
\(b_{X,R}=\mathrm{id}_{\bar B_E^{\mathrm{ch}}(\cdot)}\); the source
cobar counit
\(\varepsilon^{\mathrm{src}}_{\mathrm{bar},R}:
\Omega_E^{\mathrm{ch}}C_{X,R}^{\mathrm{geom}}
\xrightarrow{\simeq}
\mathcal F_{X,\sigma,S,\le R}^{\mathrm{hyb}}\); after a separate
source-to-target quasi-isomorphism
\(\vartheta_R:\mathcal F_{X,\sigma,S,\le R}^{\mathrm{hyb}}
\xrightarrow{\simeq}\mathsf{Den}_{\Delta_5,E,\le R}\), the cobar
comparison
\(\Theta_{\mathrm{Kos},R}=\vartheta_R\circ\varepsilon^{\mathrm{src}}_{\mathrm{bar},R}\).

\paragraph{Procedure.}
\begin{enumerate}
\item \emph{Augmentation kernel.}  Set
\(\overline{\mathcal F}_{X,R}^{\mathrm{geom}}
=\ker(\varepsilon_{X,R}^{\mathrm{hyb}})\); verify the bounded-length
hypothesis (positivity of HN height on every non-vacuum surviving
colour, hence at most \(N_R\) factors in any non-zero word).

\item \emph{Bar coalgebra.}  Form the reduced chiral bar
\[
C_{X,R}^{\mathrm{geom}}
=k\mathbf 1\oplus\bigoplus_{1\le p\le N_R}
\bar B_E^{\mathrm{ch},p}
(\overline{\mathcal F}_{X,R}^{\mathrm{geom}}),
\]
length \(p\) summand generated by ordered \(p\)-fold inputs from the
augmentation ideal of total HN height \(\le R\) (height \(>R\) outputs
zero); coaugmentation by the inclusion of the vacuum summand; counit by
projection.

\item \emph{Bar-length filtration.}
\[
F^{\mathrm{co}}_{R,p}C_{X,R}^{\mathrm{geom}}
=k\mathbf 1\oplus\bigoplus_{1\le j\le p}
\bar B_E^{\mathrm{ch},j}
(\overline{\mathcal F}_{X,R}^{\mathrm{geom}}),\qquad
0\le p\le N_R.
\]

\item \emph{Collision coproduct.}  For
\(x=[x_1|\cdots|x_p]\), define the cut-map coproduct
\[
\Delta_R^{\mathrm{ch}}(x)
=\mathbf1\otimes x+x\otimes\mathbf1
+\sum_{i=1}^{p-1}[x_1|\cdots|x_i]\otimes[x_{i+1}|\cdots|x_p],
\]
with each cut realized by the corresponding finite collision/flag
pull-push map in
\(\mathsf{Corr}^{E,\mathrm{hyb}}_{R,\mathrm{geom}}\) of
Construction~\ref{app:build-hybrid-wrapped}.  Verify
coassociativity by the four-input pentagon and the eight-word flag
coherences; verify counit identities via vacuum collision flags.

\item \emph{Conilpotence.}  Verify
\((\overline\Delta_R^{\mathrm{bar}})^{N_R+1}=0\): each reduced cut
strictly lowers bar length on each branch.

\item \emph{Bar comparison.}  Take \(b_{X,R}\) to be the identity of
\(\bar B_E^{\mathrm{ch}}
(\mathcal F_{X,\sigma,S,\le R}^{\mathrm{hyb}},
\varepsilon_{X,R}^{\mathrm{hyb}})\), eliminating the coalgebra-side
finite Koszul residuals
\(\mathfrak o^C_R, \mathfrak o^{\mathrm{fil}}_R,
\mathfrak o^{\mathrm{conil}}_R, \mathfrak o^{\Delta,\mathrm{ch}}_R\)
of Definition~\ref{def:chiral-koszul-source-datum}.

\item \emph{Source cobar counit.}  Apply Francis--Gaitsgory chiral
Koszul duality
\cite[Theorem~6.2.2]{FrancisGaitsgoryCKD} in the conilpotent completed
domain to obtain
\(\varepsilon^{\mathrm{src}}_{\mathrm{bar},R}:
\Omega_E^{\mathrm{ch}}C_{X,R}^{\mathrm{geom}}
\xrightarrow{\simeq}
\mathcal F_{X,\sigma,S,\le R}^{\mathrm{hyb}}\); verify it is a
source-internal quasi-isomorphism of chiral algebras on \(E\), not a
homotopy inverse of the target counit.

\item \emph{Source-to-target identification.}  Construct
\(\vartheta_R\) from the geometric primitive representatives, Hall
products, coproducts, pairings, and HN transition compatibility;
verify the residual
\(\mathfrak o^\vartheta_R\) of
Corollary~\ref{cor:source-bar-remaining-cobar-obstruction} vanishes;
set \(\Theta_{\mathrm{Kos},R}=\vartheta_R\circ
\varepsilon^{\mathrm{src}}_{\mathrm{bar},R}\).

\item \emph{Inverse limit.}  Verify the Koszul Mittag--Leffler
condition \(R^1\varprojlim=0\) for the inverse systems of source
coalgebras, completed chiral bar/cobar constructions, target
truncations, and the cones of \(b_{X,R}\) and
\(\Theta_{\mathrm{Kos},R}\) \cite[\S3.5]{WeibelHA}.  Form
\[
C_X^{\mathrm{geom}}
=\varprojlim_R C_{X,R}^{\mathrm{geom}},\qquad
\Theta_{\mathrm{Kos}}:
\Omega_E^{\mathrm{ch}}C_X^{\mathrm{geom}}
\xrightarrow{\simeq}\mathsf{Den}_{\Delta_5,E}.
\]

\item Output
\((C_{X,R}^{\mathrm{geom}},F^{\mathrm{co}}_{R,\bullet},
\Delta_R^{\mathrm{ch}},b_{X,R},\Theta_{\mathrm{Kos},R})\) at every
\(R\), with successor compatibility.
\end{enumerate}

\paragraph{Type checks.}
\(C_{X,R}^{\mathrm{geom}}\) is a coaugmented conilpotent chiral
coalgebra on \(E\) with bounded bar-length filtration of length at
most \(N_R\); \(\Delta_R^{\mathrm{ch}}\) is coassociative and counital
on the chosen flag stacks; the bar-cobar inversion is supplied by
Francis--Gaitsgory chiral Koszul; the cobar comparison
\(\Theta_{\mathrm{Kos},R}\) is source-internal, not the target
bar-cobar counit (the source/target separation of
Lemma~\ref{lem:source-target-koszul-separation}).

\paragraph{Obstructions discharged or invoked.}
The Koszul Mittag--Leffler exactness hypothesis is the
\textup{(D0)}-component for the chiral Koszul lane.  The total finite
Koszul residual
\(\mathfrak O_{\mathrm{Kos},R}=
(\mathfrak o^C_R,\mathfrak o^{\mathrm{fil}}_R,
\mathfrak o^{\mathrm{conil}}_R,\mathfrak o^{\Delta,\mathrm{ch}}_R,
\{\mathfrak o^{C,\mathrm{tr}}_{R'R}\},\mathfrak o^b_R,
\mathfrak o^\Omega_R,\mathfrak o^{\mathrm{hyb}}_R,
\mathfrak o^\Pi_R,\mathfrak o^{\mathrm{sep}}_R,
\mathfrak o^{\Prim}_R)\)
vanishes after this step.

\paragraph{Cross-volume firewall.}  The bar coalgebra
\(\bar B_E^{\mathrm{ch}}
(\mathcal F_{X,R}^{\mathrm{hyb}})\) is a \emph{twisting / coupling}
coalgebra (Maurer--Cartan classification), \emph{not} the bulk
\(Z^{\mathrm{der}}_{\mathrm{ch}}(\mathcal F_X^{\mathrm{hyb}})
=\mathrm{ChirHoch}^\bullet(\mathcal F_X^{\mathrm{hyb}},
\mathcal F_X^{\mathrm{hyb}})\); cf.\
\texttt{chiral-bar-cobar-vol2:CLAUDE.md~§6} (universal stage chain
\(\mathsf P\to\mathsf C\to\mathsf S\to\mathsf Z\)).
The chiral Koszul comparison
\(\Theta_{\mathrm{Kos}}:\Omega_E^{\mathrm{ch}}C_X^{\mathrm{geom}}
\xrightarrow{\simeq}\mathsf{Den}_{\Delta_5,E}\) is a comparison
between two algebras of the same level; the homotopy inverse to the
target-internal counit
\(\Omega_E^{\mathrm{ch}}\bar B_E^{\mathrm{ch}}
(\mathsf{Den}_{\Delta_5,E})\to\mathsf{Den}_{\Delta_5,E}\) is not used
in either direction.

\subsection{Construction~8 — Primitive recognition}
\label{app:build-recognition}

\paragraph{Tag.} \(\delta\)-endpoint convergence
(source-to-target identification).

\paragraph{Input.}
The descended source primitive Lie superalgebra
\(P_X^{\Pi,\mathrm{red}}\) of
Construction~\ref{app:build-hall} (after the radical quotient
\(\overline{\mathrm{Rad}\langle\,\cdot,\cdot\rangle_X}\)); the imported
target denominator algebra
\(\fg_{\Delta_5}=G(\mathscr B_{\Delta_5})\) on
\(\La^{2,1}_{II}\) of Definition~\ref{def:bkm-imaginary-simple-roots} and
Definition~\ref{def:bkm-relations}; the increasing cofinal
sequence of finite downward-saturated windows
\(W_1\subset W_2\subset\cdots\) with
\(\bigcup_\nu W_\nu=\mathcal R_+\); the Borcherds--Kac--Moody data
\((H,I,p,\alpha_i,(\,,\,))\) of
Theorem~\ref{thm:primitive-recognition}.

\paragraph{Output.}
The cofinal finite-window primitive-recognition datum
\(\{\textup{(R0),(R1),(R2),(R3),(R4),(R5),(R6),(R7),(R8)}\}\) of
Definition~\ref{def:cofinal-primitive-recognition-datum} for every
\(W_\nu\), compatible under \(W_\nu\subset W_{\nu+1}\); the source
representatives
\(e_i^X,f_i^X,h_i^X\); the source presentation map
\(\pi_W:\widehat{\mathfrak F}(\mathscr B_{\Delta_5})_W
\to P_X^{\Pi,\mathrm{red}}\big|_W\); the kernel equality
\(\ker\pi_W=
(\overline{J_{\mathrm{BK}}+\mathrm{Rad}_{\mathrm{GN}}})_W\); the
isomorphism
\[
\boxed{\quad
P_X^{\Pi,\mathrm{red}}\cong\fg_{\Delta_5},\quad
\Omega_E^{\mathrm{ch}}C_X\simeq\mathsf{Den}_{\Delta_5,E}.
\quad}
\]

\paragraph{Procedure.}
For every cofinal \(W_\nu\):
\begin{enumerate}
\item[\textup{(R0)}] \emph{Compact source labels.}  Verify every
finite source basis vector in \(W_\nu\cup(-W_\nu)\cup\{0\}\) satisfies
the compact source-admissibility condition of
Theorem~\ref{thm:finite-stage-dirac-igusa-pro-object}.

\item[\textup{(R1)}] \emph{Representatives and parities.}  Supply
parity-homogeneous primitive representatives
\(e_i^X\in P^{\Pi,\mathrm{red}}_{X,\alpha_i}\),
\(f_i^X\in P^{\Pi,\mathrm{red}}_{X,-\alpha_i}\),
\(h_i^X=\iota_H^{-1}(\alpha_i)\) for every simple-root index \(i\) with
\(\alpha_i\in W_\nu\), with the Gritsenko--Nikulin parity fibres
\(\mu_{\overline 0}(a),\mu_{\overline 1}(a)\) for imaginary simple
roots and even real representatives for \(\delta_1,\delta_2,\delta_3\).

\item[\textup{(R2)}] \emph{Relations.}  Verify in the protected Hall
super-bracket the Cartan, Chevalley, real Serre, Borcherds
orthogonality, and super sign relations for every relation whose
displayed source and target degrees lie in
\(W_\nu\cup(-W_\nu)\cup\{0\}\).  In the real rank-three subsystem the
Cartan matrix is
\(((\delta_i,\delta_j))=\mathrm{diag}(2)-2(J-\mathrm{diag}(1))\) with
real Serre exponent~3.

\item[\textup{(R3)}] \emph{Full finite parity dimensions.}  For every
\(\beta\in W_\nu\) verify
\(\dim P^{\Pi,\mathrm{red}}_{X,\beta,\overline p}
=\rootmult_{\overline p}^{\mathrm{GN}}(\beta)\) for both parities
\(p=0,1\).  In the first timelike channel
\(\beta=\delta_{123}\),
Proposition~\ref{prop:bkm-delta123-presentation-count} requires
\((d^{\mathrm{GN}}_{(1,1,1),\overline 0},
d^{\mathrm{GN}}_{(1,1,1),\overline 1})=(29,93)\).  Determinant
information \(29-93=-64=f(1,1)\) is consistent but not the
recognition.

\item[\textup{(R4)}] \emph{Hopf pairing and radical.}  Compute the
positive-negative pairing matrices in parity blocks; verify the
kernel agrees with the
Gritsenko--Nikulin/Kac invariant-pairing kernel; verify Lie-ideal and
coideal stability under all finite brackets; verify carrying under
HN transitions to the next window.

\item[\textup{(R5)}] \emph{No extra relations.}  Verify the kernel
equality
\(\ker\pi_\nu=(\overline{J_{\mathrm{BK}}+\mathrm{Rad}_{\mathrm{GN}}})_{\le W_\nu}\) on the finite triangular span generated by
\(W_\nu\cup(-W_\nu)\cup\{0\}\).

\item[\textup{(R6)}] \emph{Generation.}  Verify that every source
primitive root space in \(W_\nu\cup(-W_\nu)\), after the finite
radical quotient, is generated by the Cartan and the supplied
simple-root representatives through brackets whose intermediate
positive degrees stay in \(W_\nu\).

\item[\textup{(R7)}] \emph{Completed PBW compatibility.}  Verify the
PBW-graded isomorphism
\(\mathrm{gr}_{\mathrm{PBW}}U(P_X^{\Pi,\mathrm{red}})_{\le W_\nu}
\xrightarrow{\sim}
\mathrm{gr}_{\mathrm{PBW}}U(\fg_{\Delta_5})_{\le W_\nu}\); verify
strict preservation under \(W_{\nu+1}\to W_\nu\) transition maps.

\item[\textup{(R8)}] \emph{Exact triangular completion.}  Verify
strictness of the transition maps and closedness of images on every
finite window, so that passage to the inverse limit commutes with the
finite kernels, cokernels, radicals, PBW associated gradeds, and
parity pieces.  Equivalently the relevant \(\lim^1\) terms vanish for
these triangular finite-window systems.

\item Output the global isomorphism
\(P_X^{\Pi,\mathrm{red}}\cong\fg_{\Delta_5}\) by passage to the inverse
limit (compatible cofinal datum); together with the chiral Koszul
comparison of Construction~\ref{app:build-koszul} this gives
\(\Omega_E^{\mathrm{ch}}C_X\simeq\mathsf{Den}_{\Delta_5,E}\) and hence
the protected Pfaffian identification
\(\operatorname{Pf}_{\mathrm{prot}}(\mathfrak D_X^{\mathrm{DI}})=\Delta_5\).
\end{enumerate}

\paragraph{Type checks.}
The source presentation map \(\pi:\widehat{\mathfrak F}(\mathscr B_{\Delta_5})\to P_X^{\Pi,\mathrm{red}}\) is
continuous and degree-preserving; \textup{(R5)} gives injectivity
after the completed radical quotient; \textup{(R6)} gives
surjectivity; \textup{(R7)} gives PBW compatibility; \textup{(R8)}
gives exactness in the inverse-limit step.  No claim on signed
superdimensions \(\sdim=f(nm,l)\) substitutes for the two-integer
parity dimensions of \textup{(R3)}.

\paragraph{Obstructions discharged or invoked.}
This step inherits the obstruction tower discharged by
Constructions~1--7 and adds the cofinal exact-completion residual
\textup{(R8)}.  The Beilinson--Drinfeld bar-cobar firewall is
respected throughout: \textup{(R5)} is a source kernel-equality
theorem, not an inversion of the target bar-cobar counit;
\textup{(R7)} is a PBW graded comparison of source and target
enveloping algebras, not a graph-equality argument from the scalar
square.

\paragraph{Cross-volume anchor.}  The output is
\texttt{chiral-bar-cobar-vol2:CLAUDE.md~§9} Construction
Problem~1: the operator-level construction of
\(\mathfrak D_X^{\mathrm{DI}}\) for \(K3\times E\) with
\(\operatorname{Pf}_{\mathrm{prot}}(\mathfrak D_X^{\mathrm{DI}})=\Delta_5\), under the
four obstructions \((D_0),(O1),(O1)^+,(O2)\).  The Drinfeld doubling
\(Y^+(X)\to G(X)=D(Y^+(X))\) of
\texttt{calabi-yau-quantum-groups:CLAUDE.md} requires the Hopf
pairing of \textup{(R4)}, completion by \textup{(R8)}, integral form,
and stable-envelope transport which are supplied by the same
data.  The κ-tuple coordinate \(\kappa_{\mathrm{BKM}}=5\) is the
weight of \(\Delta_5\) and lives at level~4 of the universal stage
chain; it is a consequence of the construction, not an input to it.

\subsection{The eight-step audit}
\label{app:build-audit}

The eight construction steps form a strict ordered sequence of
input/output type discipline.  An audit of any candidate construction of
\(\mathfrak D_X^{\mathrm{DI}}\) reduces to verifying, for every
\(R\in\mathcal R_{\mathrm{HN}}\):
\begin{enumerate}
\item the lattice combinatorics of
Construction~\ref{app:build-charge-window};
\item the finite-type quasi-smooth derived stacks of
Construction~\ref{app:build-finite-moduli};
\item the hybrid local/wrapped factorization datum and the eight-word
two-step flag atlas of Construction~\ref{app:build-hybrid-wrapped};
\item the four Cech--Borel orientation classes and the
\((O1)^+\)~Coxeter coherence of
Construction~\ref{app:build-orientation};
\item the finite Hall bialgebra and primitive matrices of
Construction~\ref{app:build-hall};
\item the rank-one Pfaffian crossing and the half-Hilbert orbit
transport of Construction~\ref{app:build-dirac-pfaffian};
\item the bar coalgebra, conilpotent collision coproduct, and the
Francis--Gaitsgory cobar quasi-isomorphism of
Construction~\ref{app:build-koszul};
\item the cofinal datum
\(\textup{(R0)}\)--\(\textup{(R8)}\) of
Construction~\ref{app:build-recognition};
\end{enumerate}
together with the compatible vanishing of the obstruction tuples
\((\mathfrak O^{E_3-\mathrm{src}}_R,
\mathfrak O_{\mathrm{hyb},R},
\mathfrak O_{\mathrm{Kos},R},
\mathfrak O^{\mathrm{O1}}_R,
\mathfrak O^{\mathrm{O1}^+}_R,
\mathfrak O^{\mathrm{O2}}_R,
\mathfrak O_R^{D_0})\) at every height \(R\), and the
Mittag--Leffler exactness condition \textup{[ML]} of
Definition~\ref{def:O1-strong-reduced-orientation} on the cofinal HN inverse
system.  Each obstruction is attributable to the unique step
in which its components are constructed.

\paragraph{Order rigidity.}  Each of the seven adjacent input/output
dependencies fails type-checking under exchange:
\begin{enumerate}
\item charge window $\to$ finite moduli: $\Lambda_{\mathrm{ample}}$
  parametrises the HN strata; without the window, the moduli stack
  has no finite cofinal subsystem.
\item finite moduli $\to$ hybrid carrier: $\mathrm{Ran}^{\mathrm{hyb}}(E)$
  needs the finite Hall stack as factor on which to graft the
  wrapped/local atlas.
\item hybrid carrier $\to$ reduced orientation: the cosection-twisted
  Joyce--Upmeier line bundle is defined on the carrier, not on the
  raw moduli (the cosection lives on $\mathrm{Ran}^{\mathrm{hyb}}$).
\item reduced orientation $\to$ Hall correspondences: the Pfaffian
  line bundle on the Hall stack is the orientation-twisted sign data;
  without orientation, the Hall product is unsigned.
\item Hall correspondences $\to$ Pfaffian line + Dirac operator: the
  primitive Hall bracket $[-,-]_X$ defines the Hall--Drinfeld
  correspondence whose Pfaffian section is the Dirac--Igusa Pfaffian.
\item Pfaffian line $\to$ chiral Koszul source coalgebra: the source
  $C_X$ is the bar of the augmented hybrid factorization algebra
  pulled back along the Pfaffian comparison
  $\iota_{\mathrm{aut}}$.
\item chiral Koszul source $\to$ primitive recognition: the primitive
  part $P_X^{\Pi,\mathrm{red}}$ is the Frobenius--Serre primitive
  filtration on the source, which the recognition theorem matches to
  $\mathfrak g_{\Delta_5}$.
\end{enumerate}
Two diagnostic exchanges: orientation before hybrid leaves the middle
type-II wall \(\delta_2\leftrightarrow(0,1,1)\) without a
representative, by
Lemma~\ref{app:pfaffian-wall-three-walls}; recognition before the
Koszul cobar leaves the identification
\(\Omega_E^{\mathrm{ch}}C_X\simeq\mathsf{Den}_{\Delta_5,E}\)
unconstructed.  All other adjacent exchanges fail by the same
input/output type criterion.

\paragraph{Conditional theorem.}  Granted the eight constructions and
the four obstructions \((D_0),(O1),(O1)^+,(O2)\),
Theorem~\ref{thm:pfaffian-dirac-finite-stage} reads
\[
\operatorname{Pf}_{\mathrm{prot}}(\mathfrak D_X^{\mathrm{DI}})=\Delta_5,\qquad
\epsilon_o=\nu_{\Delta_5}\big|_{W^{(2)}(\La^{2,1}_{II})},\qquad
P_X^{\Pi,\mathrm{red}}\cong\fg_{\Delta_5}.
\]
The scalar square
\(Z_{\mathrm{BPS}}^{K3\times E}=
(\Phi_{10}^{\mathrm{un}})^{-1}=\Delta_5^{-2}\) is the level-4 trace
of the constructed protected algebra; promotion to the 3d
gravitational path integral is not claimed and requires the
Hall--Borcherds residual of
\texttt{chiral-bar-cobar-vol2}.

%%% End of Appendix F.


% ------------------------------------------------------------------
% Bibliography
% ------------------------------------------------------------------
\bibliographystyle{alpha}
\bibliography{proj}

\end{document}
